\documentclass[11pt,a4paper]{article}
\pdfextension mapfile{+cm.map}
\pdfextension mapfile{+cmextra.map}
\pdfextension mapfile{+symbols.map}
\pdfextension mapfile{+euler.map}
\pdfextension mapfile{+latxfont.map}
\pdfextension mapfile{+rsfs.map}
\usepackage[margin=25mm,headheight=15pt]{geometry}
\usepackage{fontspec}
\setmainfont{Latin Modern Roman}
\usepackage{amsmath,amssymb,amsthm,amscd,mathtools,mathrsfs}
\usepackage{microtype,booktabs,longtable,array,xcolor,graphicx,enumitem,needspace,fancyhdr}
\usepackage[round,authoryear]{natbib}
\usepackage[unicode,colorlinks=true,linkcolor=blue!40!black,citecolor=blue!40!black,urlcolor=blue!40!black]{hyperref}
\usepackage{xurl}
\usepackage{listings}
\usepackage{tabularx}
\lstset{basicstyle=\ttfamily\small,breaklines=true,columns=fullflexible,keepspaces=true}
\usepackage{tikz}
\usepackage{tikz-cd}
\usetikzlibrary{arrows.meta,calc,patterns,positioning}
\usepackage{aliascnt}
\usepackage[nameinlink,noabbrev]{cleveref}
\allowdisplaybreaks[2]
\setcounter{MaxMatrixCols}{32}
\setlength{\emergencystretch}{3em}
\setlength{\parskip}{2pt}
\newtheorem{theorem}{Theorem}[section]
\newaliascnt{proposition}{theorem}
\newtheorem{proposition}[proposition]{Proposition}
\aliascntresetthe{proposition}
\crefname{proposition}{proposition}{propositions}
\newaliascnt{lemma}{theorem}
\newtheorem{lemma}[lemma]{Lemma}
\aliascntresetthe{lemma}
\crefname{lemma}{lemma}{lemmas}
\newaliascnt{corollary}{theorem}
\newtheorem{corollary}[corollary]{Corollary}
\aliascntresetthe{corollary}
\crefname{corollary}{corollary}{corollarys}
\newaliascnt{claim}{theorem}
\newtheorem{claim}[claim]{Claim}
\aliascntresetthe{claim}
\crefname{claim}{claim}{claims}
\newaliascnt{statement}{theorem}
\newtheorem{statement}[statement]{Statement}
\aliascntresetthe{statement}
\crefname{statement}{statement}{statements}
\theoremstyle{definition}
\newaliascnt{definition}{theorem}
\newtheorem{definition}[definition]{Definition}
\aliascntresetthe{definition}
\crefname{definition}{definition}{definitions}
\newaliascnt{remark}{theorem}
\newtheorem{remark}[remark]{Remark}
\aliascntresetthe{remark}
\crefname{remark}{remark}{remarks}
\newaliascnt{example}{theorem}
\newtheorem{example}[example]{Example}
\aliascntresetthe{example}
\crefname{example}{example}{examples}
\newaliascnt{problem}{theorem}
\newtheorem{problem}[problem]{Open problem}
\aliascntresetthe{problem}
\crefname{problem}{open problem}{open problems}
\newaliascnt{question}{theorem}
\newtheorem{question}[question]{Question}
\aliascntresetthe{question}
\crefname{question}{question}{questions}
\newaliascnt{auditrecord}{theorem}
\newtheorem{auditrecord}[auditrecord]{Verification record}
\aliascntresetthe{auditrecord}
\crefname{auditrecord}{verification record}{verification records}
\newaliascnt{test}{theorem}
\newtheorem{test}[test]{Verification}
\aliascntresetthe{test}
\crefname{test}{verification}{verifications}
\newaliascnt{conjecture}{theorem}
\newtheorem{conjecture}[conjecture]{Conjecture}
\aliascntresetthe{conjecture}
\crefname{conjecture}{conjecture}{conjectures}
\newcommand{\C}{\mathbb C}\newcommand{\R}{\mathbb R}\newcommand{\Z}{\mathbb Z}\newcommand{\Q}{\mathbb Q}\newcommand{\N}{\mathbb N}\newcommand{\F}{\mathbb F}
\newcommand{\CC}{\mathbb C}\newcommand{\RR}{\mathbb R}\newcommand{\ZZ}{\mathbb Z}\newcommand{\QQ}{\mathbb Q}\newcommand{\NN}{\mathbb N}
\newcommand{\PP}{\mathbb P}\newcommand{\Pp}{\mathbb P}\newcommand{\Pone}{\mathbb P^1}\newcommand{\HH}{\mathbb H}
\newcommand{\RP}{\mathbb{RP}}\newcommand{\CP}{\mathbb{CP}}
\newcommand{\OO}{\mathcal O}\newcommand{\cO}{\mathcal O}\newcommand{\cM}{\mathcal M}\newcommand{\cF}{\mathcal F}
\newcommand{\h}{\mathfrak h}\newcommand{\hz}{\mathfrak h_z}\newcommand{\Acon}{\mathfrak A}\newcommand{\wt}[1]{\widehat{#1}}
\newcommand{\red}{\mathrm{red}}\newcommand{\tor}{\mathrm{tor}}\newcommand{\dP}{\mathrm{dP}}
\newcommand{\gh}{\widehat\gamma}\newcommand{\uh}{\widehat u}\newcommand{\wh}{\widehat w}\newcommand{\dhat}{\widehat\delta}\renewcommand{\dh}{\widehat\delta}
\newcommand{\dd}{\mathrm d}\newcommand{\ii}{\mathrm i}\newcommand{\eps}{\varepsilon}\newcommand{\partt}{\mathrm{part}}\newcommand{\Bcirc}{B^\circ}\newcommand{\ddbar}{\partial\bar\partial}
\DeclareMathOperator{\Pic}{Pic}\DeclareMathOperator{\Alb}{Alb}\DeclareMathOperator{\Aut}{Aut}\DeclareMathOperator{\NS}{NS}
\DeclareMathOperator{\Hom}{Hom}\DeclareMathOperator{\End}{End}\DeclareMathOperator{\Ext}{Ext}\DeclareMathOperator{\Tor}{Tor}
\DeclareMathOperator{\Spec}{Spec}\DeclareMathOperator{\Proj}{Proj}\DeclareMathOperator{\Sym}{Sym}\DeclareMathOperator{\Diff}{Diff}
\DeclareMathOperator{\GL}{GL}\DeclareMathOperator{\SL}{SL}\DeclareMathOperator{\PSL}{PSL}\DeclareMathOperator{\Sp}{Sp}
\DeclareMathOperator{\im}{im}\DeclareMathOperator{\coker}{coker}\DeclareMathOperator{\rank}{rank}\DeclareMathOperator{\diag}{diag}
\DeclareMathOperator{\ord}{ord}\DeclareMathOperator{\vol}{vol}\DeclareMathOperator{\pr}{pr}\DeclareMathOperator{\id}{id}
\DeclareMathOperator{\divisor}{div}\DeclareMathOperator{\Tr}{Tr}\DeclareMathOperator{\Nm}{Nm}\DeclareMathOperator{\Sha}{Sha}
\DeclareMathOperator{\Imn}{Im}\DeclareMathOperator{\Ren}{Re}\DeclareMathOperator{\trdeg}{tr.deg}
\newcommand{\Uses}[1]{\par\smallskip\noindent\textit{Uses:} #1.\par\smallskip}
\newcommand{\Verified}{}\newcommand{\AuthorSelfCheck}{}\newcommand{\NeedsCheck}{}\newcommand{\Open}{\textbf{Open.}\ }

\newcommand{\StatementRecord}[4]{\par\smallskip\begingroup\upshape\raggedright\noindent\textbf{Label:} \nolinkurl{#1}.\par\noindent\textbf{Status:} #2.\par\noindent\textbf{Check record:} #3.\par\noindent\textbf{Uses:} #4.\par\endgroup\smallskip}
\newcommand{\RungZero}[1]{\par\noindent\textbf{Result and mechanism.} #1\par\smallskip}
\newcommand{\RungOne}{\par\Needspace{6\baselineskip}\noindent\textbf{Proof at the first rung.}\par\nobreak\smallskip}
\newcommand{\CurrentChecker}{GPT-6 Astra (OpenAI, ChatGPT Pro mode, effort Pro), 18 November 2026}
\newaliascnt{cited}{theorem}
\newtheorem{cited}[cited]{Cited result}
\aliascntresetthe{cited}
\crefname{cited}{cited result}{cited results}
\newaliascnt{refuted}{theorem}
\newtheorem{refuted}[refuted]{Correction}
\aliascntresetthe{refuted}
\crefname{refuted}{correction}{corrections}
\newaliascnt{computed}{theorem}
\newtheorem{computed}[computed]{Computation}
\aliascntresetthe{computed}
\crefname{computed}{computation}{computations}
\newaliascnt{openq}{theorem}
\newtheorem{openq}[openq]{Open question}
\aliascntresetthe{openq}
\crefname{openq}{open question}{open questions}

\providecommand{\rk}{\operatorname{rk}}
\providecommand{\Ann}{\operatorname{Ann}}
\providecommand{\Log}{\operatorname{Log}}
\providecommand{\vir}{\mathrm{vir}}


\newcommand{\ExternalPaperRef}[3]{\ifmmode\text{\upshape #1#2 (Paper~#3)}\else #1#2 (Paper~#3)\fi}
\newcommand{\CompendiumRef}[2]{\ifmmode\text{\upshape S6\_09, Paper~#1, label~\nolinkurl{#2}}\else S6\_09, Paper~#1, label~\nolinkurl{#2}\fi}

\begin{document}

\title{The $(3,4,\infty)$ family: periods, moduli and the two ends}\author{AI-produced mathematical text}\date{19 November 2026}\maketitle
\paragraph{AI production and check provenance.}
This paper is AI-produced. No human has checked its mathematics.
The recorded stages are GPT-5.6 Sol (OpenAI), in Codex at reasoning
effort ultra, 24 August--9 September 2026, with some contributions by
GPT-5.6 Sol (OpenAI) in ChatGPT Pro mode; GPT-6 Astra (OpenAI), in
ChatGPT Pro mode at effort Pro, for the research runs of 12--16
November, S6\_09 of 16--17 November, the continuing S6\_06 re-derivation
of 16--18 November, and this writing of 17--19 November 2026;
and Claude Fable 5.2 (Anthropic)
at effort high or max and Claude Opus 5.5 (Anthropic) at effort max
for the readings, checks and plans of 12--18 November 2026.
S6\_06\_BLIND is credited as a re-derivation with access to preceding
material. Each statement records the actual check scope and preserves
the names, effort settings and dates of its historical checks.

\paragraph{Effort names.}
The archive's literal settings are retained with their model, provider,
interface and date. OpenAI documents reasoning allocation controls
\cite[sections ``Reasoning effort'' and ``Reasoning mode'']{OpenAIReasoning2026}
and maps ChatGPT's Pro selection to GPT-6 Pro, powered by GPT-6 Astra
\cite[section ``Choose a thinking level'']{OpenAIChatGPTPro2026}.
Anthropic documents effort as control of response-token use, including
thinking \cite[sections ``How effort works'' and ``Effort levels'']{AnthropicEffort2026}.
The check event states what was read, re-derived or executed.

\paragraph{Public-domain dedication.}
This AI-produced text is dedicated to the public domain under
CC0 1.0 Universal \cite{CC01Universal}.


\section{Results and their mechanisms}
\label{sec:I:opening}

\RungZero{We start from the compact complex six-spheres constructed by Claude
and L.~Alp\"oge and Engel's elliptic-surface description: in Paper~I,
theta transformations and integral root norms determine the periods
and moduli, while corner meridians and the Albanese nome obstruct
extension at the two ends; in Paper~II, the torsion term in the
differential sequence identifies the missing vanishing step and gives
the exact replacement criterion; in Paper~III, actual fibre framings
and integral attaching maps determine the homotopy classes and
resulting manifolds. In Paper~IV, the full cusp charts, tensor-frame
return and relative Pontryagin calculation determine the smooth
monodromy on every retained fan, with an explicit family giving a
generator of $\Theta_7$; in Paper~V, rank-six lattice data and their
period extension give the complex fourfold on
$3\mathbb H\mathbb P^2\#(-\mathbb H\mathbb P^2)$; in Paper~VI,
normalization, local Milnor fibres and compact-stratum cohomology
give the Euler sign for all degenerations in its stated
torus-fibration class.}

The fixed data are the period solution with
$(\beta^{\rm part}+\tau)(p_0)=0$, the ordered lattice
$\Lambda=\Z\langle\widehat\gamma,\widehat u,\widehat w,\widehat\delta\rangle$,
clockwise finite root generators, the specified periodic cusp fan, and
the stated affine translations. We write $d=\widehat\delta$,
$E=d\otimes\gamma$, $w=\exp(-\pi i c)$ and $u=w^2$.
The map $f:X\to\Pp^1_t$ has special values $0,1,\infty$.
The construction and its smooth six-sphere are those of Claude and
L.~Alp\"oge, \cite[Theorems~6.2 and 8.1, pp.~36 and 63--64]{AC2026},
and Engel's elliptic-surface description. The contradiction with
\cite[Corollary~2.3, p.~680]{CDP2020} is credited to Claude and
L.~Alp\"oge, \cite[Section~1.2, p.~5; Theorem~10.5 and
Section~10.5, Lemma~10.7, pp.~90--91]{AC2026}; Paper~II gives the exact
replacement for the vanishing step.

\subsection{The period and its exact domain}
\RungZero{For the zero $[\mu]$ of $\wp_\tau$ selected by the cusp
branch, the missing period is
\[
 \beta^{\rm part}=\frac{6i}{\pi}\Log
  \frac{\vartheta_1(\mu,\tau)}{2\sqrt3\,\eta_D(\tau)}.
\]
The two theta transformation identities give the prescribed affine
period laws, and the cusp value fixes the additive constant
(\ref{fam:src:C.1.1}).}
\RungOne
\begin{enumerate}[label=\arabic*.]
\item The pointed elliptic-curve map sends $[\mu]$ to the signed
section with $x=0$. The section has height $1/6$, so the
Mordell--Weil comparison identifies the required generator
(\ref{fam:src:B.1.1}, \ref{fam:src:B.1.3}).
\item The quotient $\vartheta_1(\mu,\tau)/\eta_D(\tau)$ never vanishes:
$\wp$ has its pole at the lattice and its zeros away from that lattice.
It therefore has the specified logarithm on the simply connected
period cover.
\item The transformed quotients acquire the factors
$e^{-\pi i/3+\pi i(1-\mu)^2/\tau}$ and
$i e^{\pi i\mu^2/\tau}$. Multiplication of their logarithms by
$6i/\pi$ gives the two affine laws, with the integer defects fixed at
the finite fixed points.
\item The Fourier equation at the cusp gives
$e^{2\pi i\mu_0}=-(5+2\sqrt6)$ and
$2\sin(\pi\mu_0)=2\sqrt3$. The product then gives
$\beta^{\rm part}+\tau\to0$.
\item The difference from the prescribed solution descends to a
holomorphic function on $\Pp^1$ and has zero cusp value, hence is zero.
\end{enumerate}

\RungZero{The exact domain is $0<|w|<216$. The normalized theta norm
at a zero of $\wp$ has maximum $\sqrt2$, attained precisely over
$t=1$; the real period determinant therefore stays nonzero precisely
on that disc (\ref{fam:src:C.2.3}).}
\RungOne
\begin{enumerate}[label=\arabic*.]
\item The full real determinant has the nonzero factor
$D_c=\Im\beta_c-6(\Im\mu)^2/\Im\tau$, and
$D_{\rm part}=(6/\pi)\log(N(P)/(2\sqrt3))$
(\ref{fam:src:B.1.2}).
\item The isogeny distribution identity yields the division formula
\[
N(P)^{n^2}=\frac{N(nP)|\Delta_E|^{(n^2-1)/12}}{|\psi_n(P)|}.
\]
(\ref{ari:C.1.1}, \ref{ari:C.4.2}).
\item The $n=2,4,8$ formulas and their exact polynomial denominators
bound $N(P)$ strictly below $\sqrt2$ on the three annuli
$|t-1|\ge9/10$, $1/28\le|t-1|\le9/10$, and
$1/500\le|t-1|\le1/28$
(\ref{ari:C.8.1}, \ref{ari:C.8.4}).
\item The local theta expansion gives
$\log N-\log\sqrt2\le-\pi|h|^2/2+150|h|^4<0$
away from $h=0$ on the remaining disc. The two applications of
Rouch\'e's theorem cover that entire $t$-disc
(\ref{ari:C.8.2}, \ref{ari:C.8.3}).
\item At $t=1$, $N=\sqrt2$ and
$M=\max D_{\rm part}=-(3/\pi)\log6$.
At the cusp $D_{\rm part}\to-\infty$. Thus $D_c$ has no zero exactly
when $\Im c<-M$, giving $|w|<e^{-\pi M}=216$.
\end{enumerate}

\subsection{All discrete data and the parameter}
\RungZero{For
$w_1=r_1\varepsilon+k_1d$, $w_2=r_2\varepsilon'+k_2d$,
$3\nmid r_1$, $2\nmid r_2$, and cusp winding $l_0$, the complete
invariants are
\[
 \left(|12l_0-4r_1-3r_2|,\ e^{-\pi i(c+k_2)}\right).
\]
The finite norm equations determine the parity obstruction and the
remaining cusp winding determines the integer $p$
(\ref{fam:src:B.2.1}).}
\RungOne
\begin{enumerate}[label=\arabic*.]
\item The anticanonical pencil and the three labelled special fibres
recover the base coordinate. A biholomorphism consequently intertwines
the two monodromy matrices over that fixed base.
\item Their rational centralizer is $\Q I+\Q E$. The period equation
$L\Pi_c=\Pi_{c'}H$ forces $L=aI_2$, $a=\pm1$, and
$c'-c=n\in\Z$ (\ref{fam:src:A.2.1}).
\item On the root covers the exact equation is
$N_j\ell_j=-(I+nE)w'_j+aw_j$.
The norm images are $3\Z\varepsilon+\Z d$ and
$4\Z\varepsilon'+2\Z d$; hence the last finite obstruction is
$n\equiv k_2-k'_2\pmod2$.
\item A cusp extension has the actual character map
$(x,q)\mapsto(q^v a(q)x^a,q)$, with $v\in\Z^2$ and a holomorphic
unit $a(q)$. Its winding agrees with the two finite defects exactly
when $p'=ap$.
\item The differences of these local maps form a torsor under
$0\to\OO\to\mathcal V\to\OO(-1)\to0$.
Its $H^1$ vanishes, so the displayed local maps and their inverses
glue. This proves sufficiency while retaining the complete collar
section (\ref{fam:src:A.3.1}, \ref{fam:src:B.2.1}).
\end{enumerate}

\RungZero{At each fixed nonzero $u=w^2$ there are two sphere
isomorphism classes, exchanged by a positive $u$ circuit. The
order-four norm quotient detects their difference, and the family
with fixed logarithmic data closes over the punctured $w$-disc
(\ref{fam:src:B.2.2}, \ref{fam:src:B.2.3}).}
\RungOne
\begin{enumerate}[label=\arabic*.]
\item The datum $|p|=1$ is carried to the canonical one by the
explicit finite and cusp maps of \ref{fam:src:B.2.1}.
\item Its remaining identification subgroup is $2\Z$, because
$d$ has order two in $\Lambda^{A_2}/N_2\Lambda$.
\item The logarithms at fixed $u$ differ by integers, producing
exactly the classes $c+2\Z$ and $c+1+2\Z$.
\item A positive circuit sends $c$ to $c-1$, and exchanges these
classes. A descent datum would have to identify the corresponding
fibres; the computed order-two obstruction excludes that datum.
\item The even-shift maps satisfy their exact cocycle law and give
the proper family over $w$. Local inverse logarithms give the
corresponding families on simply connected $u$-open sets.
\end{enumerate}

\RungZero{For coprime multiplicities $(3k,4l)$, the fixed-data period
subgroup is $2kl\Z$, giving $2kl$ classes at fixed $u$.
The two norm images are multiplied by $k$ and $l$, so their central
congruences become $k\mid n$ and $2l\mid n$
(\ref{fam:src:B.3.1}, \ref{push:src:E.1.2}).}
\RungOne
\begin{enumerate}[label=\arabic*.]
\item The vanishing of $H^2(S^6;\R)$ forces every map to
$\Pp^1$ to be constant on a regular K\"ahler torus fibre by the
integral of its pulled-back Fubini--Study form. A biholomorphism
therefore preserves the base and its multiplicities.
\item The root actions have the same linear parts and orders
$3k,4l$, giving $N_{3k}=kN_1$ and $N_{4l}=lN_2$.
\item In the primitive sphere data, the finite equations impose
$k\mid n$ and $2l\mid n$. Their least common multiple is $2kl$
because $\gcd(k,2l)=1$.
\item The cusp equation is $p'=ap$ with
$p=12kl\,l_0-4l\,r_1-3k\,r_2$. Its integral adjustment and the
same vanishing of $H^1(\mathcal V)$ give every allowed comparison.
\item The Chinese remainder map to the two central congruence
groups has $2kl$ elements, and $c\mapsto c-1$ acts by a generator.
\end{enumerate}

\subsection{Division values and the Weierstrass comparison}
\RungZero{At every torsion specialization of $P$, the normalized
norm is algebraic with an exact division-polynomial value,
including the five values below and the complete orders four
through eight: the distribution identity fixes the multiplicative
constant, and cancellation at $mP=\pm P$ gives the algebraic
expression (\ref{ari:C.4.2}, \ref{ari:C.4.3}, \ref{ari:C.4.4}).
\[
\begin{array}{c|c|c}
t&\operatorname{ord}(P)&N(P)^{12}\\ \hline
0&3&27\\
1&2&64\\
32/5&4&10\\
64(47-27\sqrt3)/11&6&17+10\sqrt3\\
64(47+27\sqrt3)/11&6&10\sqrt3-17
\end{array}
\]
The orders at $t=0,1$ are those on the smooth models after
the specified base changes.}
\RungOne
\begin{enumerate}[label=\arabic*.]
\item The theta Green function has zero Haar mean and expansion
$U_\Lambda(z)=\log|z|+\log|\Delta_\Lambda|/12+o(1)$.
These two normalizations determine the isogeny distribution identity.
\item The divisor and leading coefficient of $\psi_m$ cancel the
singularities of $m^2U(z)-U(mz)+\log|\psi_m(z)|$.
The remaining harmonic function on the compact torus is constant,
and the expansion computes that constant.
\item The exact model map
$(x,y)\mapsto(2x,-2\sqrt2y)$ carries the arithmetic cubic and its
section to the geometric cubic and its selected section.
Its discriminant scales by $64$, while the division polynomial
scales by $(-\sqrt2)^{m^2-1}$; the displayed norm identity is
therefore preserved (\ref{I:weierstrass-comparison}).
\item At exact order $n$ choose $m\equiv\pm1\pmod n$.
Then $N(mP)=N(P)$, and
$N(P)^{12}=|\Delta_E|/|\psi_m(P)|^{12/(m^2-1)}$.
\item The nonzero polynomial $\psi_n(P_t)$ makes $t$ algebraic.
The positive radical in the preceding formula proves algebraicity
of $N(P)$ and the stated field inclusions.
\item Evaluate the exact division recurrence at the selected section
and remove the proper-order factors. The resulting squarefree
polynomials on the smooth base are
\[
\begin{aligned}
F_4&=5t-32,\\
F_5&=89t^2+4736t-4096,\\
F_6&=11t^2-6016t+8192,\\
F_7&=881t^4+4745344t^3-12079104t^2\\
   &\qquad+24641536t-16777216,\\
F_8&=49t^4-2614336t^3-7864320t^2\\
   &\qquad+19398656t-8388608.
\end{aligned}
\]
Their exact gcd computations with their derivatives and the
proper-order factors show that their roots give every point of
exact order $4,5,6,7,8$, including the nonreal order-seven roots
(\ref{ari:C.4.4}).
\item For a root of $F_n$, put $r=N(P)^{12}$. Substitution of
$m=n-1$ into the norm identity gives the positive value determined by
\[
\begin{array}{c|l}
n&\text{equation for }r\\ \hline
4&r=64/|t|\\
5&r^5=432^5/((1-t)(5t-32)^4)\\
6&r^2=1728^2/F_5(t)\\
7&r^{35}=2^{192}3^{81}/(|t|^8|t-1|^3|F_6(t)|^{12})\\
8&r^4=1728^4/|F_7(t)|.
\end{array}
\]
At order six the exact identity
$F_5(64(47\mp27\sqrt3)/11)(17\pm10\sqrt3)^2=1728^2$
selects the two positive radicals in the opening table;
$64/(32/5)=10$ gives order four, and the two central theta
values give the remaining entries
(\ref{ari:C.1.2}, \ref{ari:C.4.4}).
\end{enumerate}

\subsection{The multiple fibres and real structures}
\RungZero{The Albanese nomes of the two reduced multiple fibres are
$q_2=iw/216$ and $q_1=-w^2/48^3=(27/64)q_2^2$.
The invariant-coordinate lattice, including its affine half-period,
and the two exact theta values determine these constants
(\ref{ari:C.2.1}).}
\RungOne
\begin{enumerate}[label=\arabic*.]
\item On the order-four product cover the fixed-coordinate lattice
is $\frac12\Z+\frac{\kappa_2}{4}\Z$. Multiplication by two gives
the Albanese lattice $\Z+\frac{\kappa_2}{2}\Z$.
\item The corresponding order-three calculation gives the lattice
with positive-imaginary modulus $-\kappa_1$; both calculations retain
the finite affine deck translations (\ref{fam:src:A.4.1}).
\item Evaluation of the theta expression gives
$\kappa_1=c+\frac12-\frac{3i}{2\pi}\log48$ and
$\kappa_2=c+\frac32-\frac{3i}{\pi}\log6$.
\item Exponentiation in the two displayed Albanese moduli gives
the claimed nomes and their exact relation. Thus $q_2$ identifies
the parameter domain with the punctured unit disc.
\end{enumerate}

\RungZero{The conjugate member has parameter $1-\bar c$, so real
structures occur precisely on the imaginary $w$-diameter, in two
holomorphic-conjugacy classes. The explicit reflection and cap
meridians show that every nonempty fixed manifold is $S^3$
(\ref{fam:src:B.4.2}, \ref{mapsI:src:A.3.1},
\ref{mapsI:src:A.3.2}, \ref{mapsI:src:A.3.3}).}
\RungOne
\begin{enumerate}[label=\arabic*.]
\item The integral matrix
$H=\left(\begin{smallmatrix}1&0&0&0\\0&1&0&0\\-6&0&-1&0\\0&-1&0&-1\end{smallmatrix}\right)$
intertwines the reflected period matrices. The finite affine defects
and the complete cusp unit shift its target from $-\bar c$ to
$1-\bar c$.
\item The moduli invariant sends that parameter to $-\bar w$.
Equality with $w$ is therefore exactly $w\in i\R$ in the punctured
domain. The displayed translations in \ref{mapsI:real-eq-2} make
the reflection square an integral period.
\item Its action on $\operatorname{Aut}(X)=\C^*$ is
$z\mapsto1/\bar z$. The involution equation gives $z\in\R^*$,
and conjugation multiplies $z$ by $|y|^2$; precisely the two signs
remain.
\item The equations $(I-R_\bullet)x=b+\epsilon d/2\pmod\Lambda$
and the cusp equations $z\bar z=K_v$ give every entry of the fixed
fibre table. One negative-imaginary real form maps onto all
$\RP^1$; the other nonempty images are the two specified closed arcs.
\item On each negative-imaginary real form, the semifree circle
action has a fixed circle. Smith homology makes its orbit surface
a disc, and the explicit attachment is
$(D^2\times S^1)\cup(S^1\times D^2)=S^3$.
\item On the positive-imaginary nonempty real form the two cap
meridians, transported by the actual odd-root matrix, have
intersection one. Their genus-one Heegaard gluing gives $S^3$;
the free circle action has Euler number of absolute value one.
\end{enumerate}

\subsection{Integral monodromy and the period domain}
\RungZero{The monodromy group has index twelve in the integral
stabilizer of $d$. It is the kernel of the character
$\chi(B,x,t)=t+6(x_1+x_2+x_1x_2)-\alpha(B)\pmod{12}$;
the commutator and its two conjugates generate the required
Heisenberg lattice (\ref{fam:src:F.1.1}).}
\RungOne
\begin{enumerate}[label=\arabic*.]
\item In the displayed symplectic integral basis every stabilizer
matrix has the unique form $g(B,x,t)$ of
\ref{fam:src:F.1.1}, with $B\in\operatorname{SL}_2(\Z)$,
$x\in\Z^2$, $t\in\Z$.
\item Multiplication gives
$(B,x,t)(C,y,s)=(BC,x+By,t+s-6x^tJ_2By)$.
Substitution proves that $\chi$ is a homomorphism, using the
abelianization character $\alpha(S)=9$, $\alpha(T)=1$.
\item Both geometric generators have $\chi=0$ and their Levi
projections generate $\operatorname{SL}_2(\Z)$.
\item Their commutator word and conjugates supply
$U=(1,0,-6)$, $V=(0,1,-6)$ and $Z=(0,0,12)$ in the unipotent
kernel, with $[U,V]=Z^{-1}$. These generate exactly the
integral solutions of $t+6(x_1+x_2+x_1x_2)=0\pmod{12}$.
\item Reducing any element of $\ker\chi$ first in its Levi
projection and then by $U,V,Z$ puts it in the geometric group.
The central coordinate maps onto $\Z/12$, proving the index.
\end{enumerate}

\RungZero{On the null locus of $\psi(d,Jd)$ there is a dense set of points
with infinite monodromy stabilizer, because the rational parabolic
conjugates of the explicit unipotent stabilizer are dense there.
On either nonzero locus the real parabolic stabilizer is compact,
and the family period map embeds the $t$-line properly in the
quotient: the modular quotient $j/1728$ is its left inverse and
controls the inverse images of compact sets
(\ref{fam:src:F.2.2}, \ref{fam:src:F.2.3}, \ref{fam:src:F.3.1}).}
\RungOne
\begin{enumerate}[label=\arabic*.]
\item The explicit null complex structure is fixed by an infinite
unipotent subgroup. Rational parabolic conjugates remain rational;
the index-twelve lattice calculation supplies an infinite subgroup
of the geometric monodromy in each such stabilizer.
\item Rational parabolic points are dense in the real parabolic
orbit, which is the null locus. These stabilizers give the stated
failure of properness there.
\item At nonzero sign the parabolic is transitive with stabilizer
isomorphic to a compact product. A discrete subgroup therefore acts
properly on that homogeneous space.
\item In the period coordinates the sign invariant is exactly
$\psi(d,Jd)=-6/D_c$. The modular function of the associated
two-dimensional quotient satisfies $j=1728t$.
\item The function $j/1728$ descends to the monodromy quotient
and is a left inverse to the period map. It proves injectivity
and immersion, including the finite root charts. The inverse
image of a compact subset has bounded $t$ and is closed, proving
properness.
\end{enumerate}

\RungZero{The N\'eron--Severi lattice of every regular fibre is
the integral kernel of one displayed period equation, and its
intersection square is $-2\operatorname{Pf}(E)$.
The same Hermitian form gives the projective and elliptic-quotient
criteria with the actual fibre orientation
(\ref{fam:src:D.1.4}, \ref{fam:src:D.1.6}).}
\RungOne
\begin{enumerate}[label=\arabic*.]
\item An integral alternating form $E$ has type $(1,1)$ exactly
when its restriction to the period kernel vanishes. Substitution of
the two kernel vectors gives
$E_{12}-E_{13}\tau+(6E_{23}-E_{14})\mu+E_{24}\beta
+E_{34}(6\mu^2-\tau\beta)=0$.
\item The determinant of the real period map has sign $D_c<0$.
Transporting the integral wedge square through this actual map
therefore gives $E^2=-2\operatorname{Pf}(E)$.
\item The associated Hermitian form is definite exactly when
$\operatorname{Pf}(E)<0$; reversing its sign then gives a positive
integral form and the projective embedding criterion.
\item A nonzero rank-one semipositive integral form has rational
complex kernel. Quotienting by that kernel gives the elliptic
quotient; pulling back an elliptic polarization proves the reverse
implication. Rank one is the equation $\operatorname{Pf}(E)=0$.
\item The monodromy differences of the period functions show that
the generic integral kernel is $\Z\eta$, with $\eta^2=-12$.
The explicit special integral forms in \ref{fam:src:D.1.5}
produce projective fibres for every fixed ordinary $t$.
\end{enumerate}

\subsection{Cohomology and deformations}
\RungZero{The Hodge numbers are the following table, with rows
indexed by $q$ from $3$ to $0$ and columns by $p$ from $0$ to $3$:
\[
\begin{pmatrix}0&0&0&1\\0&1&2&1\\1&2&1&0\\1&0&0&0\end{pmatrix}.
\]
The cotangent direct images reduce them to line bundles on
$\Pp^1$; the last connecting rank is fixed by the nonzero scalar
$\mu_s(z_2)=4i\sqrt\pi/(3^{3/4}\Gamma(1/4)^2)$
(\ref{inv:A2.2}, \ref{inv:A3.1}).}
\RungOne
\begin{enumerate}[label=\arabic*.]
\item The logarithmic relative cotangent sequence is
$0\to N\to E\to\OO\to0$, with
$N=\omega_f(-2S_1-3S_2)$. The invariant logarithmic vertical
field has no zero in that locally free sequence.
\item Relative duality and the root characters compute all
direct images of $N$ and $E$; the extension groups of the
resulting line bundles on $\Pp^1$ vanish
(\ref{inv:A1.1}, \ref{inv:A1.2}).
\item The determinant of the Kodaira--Spencer matrix descends
with simple zero at each finite special value. The theta
half-period derivative evaluates the remaining order-four
entry as the displayed nonzero scalar.
\item The residue sequence, retaining its local torsion
quotients, gives the direct images of $\Omega_X^1$;
duality gives those of $\Omega_X^2$
(\ref{inv:A1.3}, \ref{inv:A2.4}).
\item Leray over the curve has only its two base columns.
The cohomology of the computed line bundles gives precisely
the displayed table and its Serre-dual entries.
\end{enumerate}

\RungZero{The Fr\"olicher sequence satisfies
$E_1\ne E_2=E_\infty$, and the positive-degree ideal in the
Bott--Chern algebra has square zero. The Hodge table and the
Betti numbers force the differential ranks; the fibre class
and its explicit primitives determine the products
(\ref{inv:C1.1}, \ref{inv:C2.1}, \ref{inv:C3.3}).}
\RungOne
\begin{enumerate}[label=\arabic*.]
\item The vanishing of intermediate Betti numbers, together
with the displayed Hodge table, forces the four nonzero
$d_1$ ranks to be one. Every possible later arrow has zero
source or zero target.
\item The resulting double-complex summands determine the
Bott--Chern and Aeppli dimensions through the cited
decomposition theorem, with conjugation and Serre duality
fixing both diagonals.
\item The global Dolbeault representative $Q$ restricts to
the invariant fibre class $\eta$, with $\eta^2=-12$.
Its nonzero restrictions identify the required products
before passage to Bott--Chern cohomology.
\item The displayed $\partial\bar\partial$ primitives in
\ref{inv:C3.3} kill the positive-degree products. The
remaining top-degree candidate has zero integral by Stokes,
which determines its class. Thus the algebra is $\C\oplus N$
with $N^2=0$.
\end{enumerate}

\RungZero{The relative one-forms have torsion exactly along
$S_1\cup S_2$; the cusp module is torsion-free and fails to be
locally free. The explicit modules for $t=s^m$, $t=xy$ and
$t=xyz$ give this correction to the printed torsion description
in Claude--Alp\"oge, Remarks~1.3 and 9.21
(\ref{inv:A4.1}, \ref{inv:A4.2}).}
\RungOne
\begin{enumerate}[label=\arabic*.]
\item At $t=s^m$, the relation is $ms^{m-1}ds=0$,
giving the torsion summand $\OO/(s^{m-1})\,ds$.
\item At $t=xy$, the relation vector $(y,x,0)$ is
primitive in codimension one. Its quotient embeds in the
explicit ideal presentation and has no torsion.
\item At $t=xyz$, the same calculation with
$(yz,xz,xy)$ gives a torsion-free module of rank two
whose fibre dimension jumps on the double locus.
\item Taking the exterior square retains the separate
triple-point torsion module
$(Ru\oplus Rv)/(xu,y(u-v),zv)$. This records exactly
which differential degree each local conclusion concerns.
\end{enumerate}

\RungZero{Every tangent twist has an explicit cohomology
formula. Integral twists are root-line bundles on $\Pp(3,4)$;
nonintegral twists reduce to finite Koszul modules at the
resonant fibres and a length-one cusp term
(\ref{inv:D4.1}, \ref{inv:D4.4}, \ref{inv:D4.5}).}
\RungOne
\begin{enumerate}[label=\arabic*.]
\item For the root generator $\mathcal R$, write
$r_1=n\bmod3$, $r_2=3n\bmod4$ and
$d(n)=(n-4r_1-3r_2)/12$. The coarse invariant sections
of $\mathcal R^n$ form $\OO(d(n))$.
\item The logarithmic tangent Koszul maps have matrices
$(0,y,y,x)^t$ and
$\left(\begin{smallmatrix}-y&0\\-x&y\end{smallmatrix}\right)$.
Their cokernels, followed by the cusp normal quotient,
give the three integral direct images printed in
\ref{inv:D4.1}.
\item A nonintegral twist has nontrivial cusp character;
the normalization complex is acyclic there except for
the stated normal length-one contribution.
\item At a regular resonance the finite relative Dolbeault
complex has $H^1=(A/s)^2$, $H^2=A/s^2$. The root-three
and root-four matrices in \ref{inv:D4.4}, followed by
their exact character invariants, give the other lengths.
\item Coherence makes the resonance support finite for each
twist. Summing its lengths gives $J_z$ and hence
$h^q(T_X\otimes M_z)=(0,J_z,1+J_z,0)$ for nonintegral $z$;
Leray gives the integral formula from the line bundles.
\end{enumerate}

\RungZero{For the standard data, $h^q(T_X)=(1,1,1,0)$, the Kuranishi germ
is smooth of dimension one, and $\operatorname{Aut}(X)=\C^*$:
these credited results of DGKX version~1 are re-derived here through
the root-curve Koszul complex and the integral translation sheaf.
The same root calculation gives $(1,1,1,0)$ for every coprime
$(3k,4l)$ sphere family (\ref{fam:st:A.5.2},
\ref{fam:src:B.4.1}, \ref{push:src:E.2.1}).}
\RungOne
\begin{enumerate}[label=\arabic*.]
\item The weighted root line satisfies
$K_{\Pp(3,4)}=\mathcal R^{-7}$, and its cohomology is
the count of monomials of weights $3$ and $4$.
\item The explicit Kodaira--Spencer matrix and the cusp
normal exact sequence yield the tangent direct images;
their coarse cohomology is $(1,1,1,0)$.
\item Differentiation of the actual period family in $c$
gives the nonzero class $[\xi\otimes\gamma]$.
The induced map to the one-dimensional Kuranishi germ
has nonzero derivative, forcing its defining ideal to
vanish.
\item Every automorphism preserves the recovered base.
The centralizer and finite affine equations force its
linear part to be the identity. The global translation
space modulo its integral periods is $\C e/\Z e=\C^*$.
\item For $(3k,4l)$ the root characters and the same
local Koszul matrices have the invariant degrees
computed in \ref{push:src:E.2.1}; Leray gives the same
four tangent dimensions.
\end{enumerate}
The credited appendices distinguish the two histories. The
working edition contained our proof of the standard-data
$2\Z$ statement by 8 September 2026. DGKX version~1,
dated 23 September 2026, states it as Theorem~1.4;
version~2 does not contain it. Our tangent-cohomology
argument was re-derived after reading DGKX version~1.

\subsection{The puncture and the outer circle}
\RungZero{The untwisted pinching has explicit semistable
and nodal fillings, while its proposed smoothing to the
six-sphere family is excluded by the virtual Euler number.
The plumbing has virtual Euler number zero and a smooth
six-sphere has Euler number two
(\ref{push:src:A.1.2}, \ref{push:src:A.2.2},
\ref{push:src:A.3.2}, \ref{night:P2.10}, \ref{night:P2.11}).}
\RungOne
\begin{enumerate}[label=\arabic*.]
\item The two line-bundle charts are glued by
$z_0w_1=r\varphi$ and $z_1w_0=r$, with the actual unit
$\varphi$ supplied by the translation of the elliptic
section. The annular identifications give a proper family.
\item The blowup of $(z,r)$ and the two small resolutions
of the resulting node give the two triangular fans and
their semistable fillings; contraction gives the nodal one.
\item For the one-component pinching, normalization and the
conductor formula give
$e(Y)-2e(S)+\int_S c_1(S)c_1(L)-c_1(L)^2=24-24=0$.
\item The relative cotangent complex of an arbitrary flat
local-complete-intersection smoothing fixes this virtual
Chern number. On a smooth fibre it is the topological
Euler number, so it excludes the proposed Euler-two
smoothing and gives the correction to Engel's Conjecture~8.2.
\end{enumerate}

\RungZero{After the cover $w=r^6$ the sphere family fills
over $(\Pp^1\times\Delta_r)\setminus\{(\infty,0)\}$,
with a twelve-cycle over the ordinary central base.
The equations $X_iY_i=r$ compactify the long period,
and the actual finite translations permute the components
freely (\ref{push:src:B.4.19}, \ref{night:P2.3}).}
\RungOne
\begin{enumerate}[label=\arabic*.]
\item Put $X_i=W/r^i$, $Y_i=r^{i+1}/W$.
The relations $X_iY_i=r$ and the annular transition
maps give the full infinite chain, including both axes.
\item The long-period deck transformation changes the
index by twelve and translates the elliptic coordinate
by the specified $6\mu$. Its quotient has twelve
central components and compact fibres.
\item On the two finite root charts the component shifts
are $4l_1$ and $3l_2$. Admissibility makes the relevant
cyclic actions free, so the smooth quotients extend there.
\item The cover $y^{12}=t^{4l_1}(t-1)^{3l_2}$ and the
full toric units identify these charts with the original
family for $r\ne0$. The omitted corner is precisely
where both long-period directions meet.
\end{enumerate}

\RungZero{No flat proper filling exists at $w=0$, after any finite cover,
for any admissible twist or any coprime $(3k,4l)$ family.
The actual cusp deck map has $\gamma$-coefficient $-pn$;
commuting meridians propagate the corresponding ratio along the
connected corner fibre, whose two ends give $0$ and $-p\ne0$
(\ref{push:src:B.4.33}, \ref{night:P2.7}, \ref{night:P2.9}).}
\RungOne
\begin{enumerate}[label=\arabic*.]
\item The intrinsic pluricanonical pencil, its labelled
multiple fibres and the complete cusp marking recover
the base map in a putative filling. After modification
and finite cover the two boundary axes are normal crossings.
\item Modulo $\Z d$, a central meridian has finite orbit
under a finite-index subgroup of the base monodromy.
The explicit two unipotent equations have common kernel
zero, so its radial class is zero
(\ref{night:P2.5}, \ref{night:P2.6}).
\item The actual cusp-origin comparison gives
$\Delta_RF=an\,\gamma-6an(n-1)d$, with
$a=-p=4l_1+3l_2-12l_0$. In particular its
$\gamma$-coefficient is $-pn$, with its sign fixed by
the displayed deck maps (\ref{night:P2.2}).
\item For adjacent divisors with orders
$(\alpha_i,\beta_i)$, their commuting coordinate circles
give
$(m\alpha_i+n\beta_i)\kappa_j
=(m\alpha_j+n\beta_j)\kappa_i$.
Thus the ratio is constant across every meeting pair.
\item Properness and connectedness of the corner fibre
join the horizontal divisor to the cusp divisor by
such a chain. The respective ratios are $0$ and $a$,
forcing $p=0$, contrary to admissibility.
\item For $(3k,4l)$ replace twelve by $12kl$ and use
$p=12kl\,l_0-4l\,l_1-3k\,l_2$.
The primitive root equations still make $p\ne0$;
the same connected-chain calculation applies.
\end{enumerate}

\RungZero{No flat proper extension crosses any point of
$|w|=216$. The intrinsic anticanonical divisor carries
an Albanese curve with invariant $J(iw/216)$, whose
rational-cusp growth makes the whole circle a natural
boundary (\ref{push:src:D.1.3}).}
\RungOne
\begin{enumerate}[label=\arabic*.]
\item The unique anticanonical divisor is $2S_2$.
Its reduced support is intrinsic under every fibrewise
biholomorphism, and its Albanese nome is $iw/216$.
\item At a reduced rational $b/d$, modular inversion
gives $\log|j(b/d+iy)|=2\pi/(d^2y)+o(1)$.
This growth excludes a meromorphic continuation there;
density of rational cusps excludes one across any arc.
\item In a putative flat proper extension, remove the
discrete nonsmooth image and the torsion support of
$R^1\pi_*K^{-1}$. The canonical base-change map then
gives a relative anticanonical section of rank one.
\item Its reduced divisor is a smooth proper surface
family away from a countable set. Its Albanese invariant
is holomorphic there, by the explicit weight-one period
construction (\ref{push:src:D.1.2}).
\item An arc has a point outside that countable set.
At that point the invariant continues $J(iw/216)$,
contradicting the growth calculation.
\end{enumerate}

\RungZero{At a boundary parameter the total root quotient
has a Hausdorff central fibre equal to a circle for rational
$a=\Re c+3/2$ and a point for irrational $a$.
All nearby additive-group limits are determined by an
explicit surviving dual lattice; the quadratic term in
the period matrix produces the additional nonseparation
(\ref{push:src:D.2.3}, \ref{push:src:D.3.1}).}
\RungOne
\begin{enumerate}[label=\arabic*.]
\item The integral basis and complex-linear change in
\ref{push:src:D.2.1} give the four periods
$(1,\frac12)$, $(\tau,\frac12+h)$, $(0,1)$,
$(6h,a+B(h))$, where $B(h)=-3ih^2+O(h^4)$.
\item Inverting the transpose of this real period
matrix gives the exact covectors in
\ref{eq:d-dual-exact}. Every bounded sequence has
eventually fixed integral labels.
\item The labels that survive form a primitive subgroup
$K$ of $L_a$; their last coefficients tend to a
homomorphism $\ell:K\to\R$. Pontryagin--Chabauty
duality gives the annihilator formula in
\ref{push:src:D.3.1}.
\item The explicit paths with
$V_n=nv+v_0$ and $d_n(1+\|V_n\|^2)\to0$
realize every such pair $(K,\ell)$ inside the admissible
side. This specifies the maps between the marked limits.
\item On a generic fixed-boundary ray only
$d_0\Re\xi$ survives when $a$ is rational, and no
nonzero label survives when it is irrational. The
resulting period limits prove exactly the asserted
nonseparation relation. Its quotient is Hausdorff and
has the universal factorization property.
\end{enumerate}

\subsection{Arithmetic values and all level curves}
\RungZero{At the eleven rational nonspecial CM parameters,
the theta norm is an explicit positive algebraic factor
times $e^{-\widehat h_O(P)}$. The local division identity
computes every finite-place factor on an integral minimal
model, and the product formula supplies the archimedean
term (\ref{ari:C.6.1}, \ref{ari:C.6.3}).}
\RungOne
\begin{enumerate}[label=\arabic*.]
\item The exact Weierstrass changes in the CM table
send the displayed rational point to the printed
point on each integral minimal model.
\item At each bad prime, the point or its double
lies in the nonsingular identity component. The
valuation of its coordinate and of $\psi_m(P)$ gives
the local height by \ref{ari:C.6.2}.
\item The eleven finite factors are respectively
$\sqrt6$, $2^{2/3}3^{1/4}$, $2^{3/4}$, $2$, $2$,
$\sqrt3$, $1$, $3^{5/12}$, $1$, $1$, $1$ in the
order of the printed CM table.
\item At infinity the same normalization is
$\lambda_\infty(P)=-\log N(P)$.
The product formula for the division identities
therefore gives $N(P)=A_t e^{-\widehat h_O(P)}$.
\item The Mordell--Weil and local calculations in
\ref{ari:C.6.4}--\ref{ari:C.6.5} specify the
non-torsion specializations to which the height
description applies; they do not decide their norm
transcendence.
\end{enumerate}

\RungZero{The norm has one coarse critical point, its
maximum at $t=1$, and each lower level is a Jordan
curve strictly star-shaped about $1$. The exact
derivative is negative along every ray because the
normalized critical factor satisfies $|1-X|<1$
(\ref{ari:C.8.15}, \ref{push:src:D.4.8}).}
\RungOne
\begin{enumerate}[label=\arabic*.]
\item The heat equation gives
$(t-1)L_t=X(X-2)/24$, with
$X=3W\wp'(\mu)/g_2$. The real radial derivative is
$\Re(X(X-2))/12$ (\ref{ari:C.8.12}).
\item The convergent cusp Fourier series gives
$|1-X|<17/20$ when $\Im\tau\ge7/5$
(\ref{ari:C.8.13}).
\item On the remaining Cayley disc, divide out the
known equianharmonic zero. Exact modular recurrences,
Cauchy tails and the displayed positive Bernstein
coefficients give $|1-X|<1$
(\ref{ari:C.8.14}).
\item The two regions cover a closed modular
fundamental domain. The derivative is therefore
strictly negative along each ray from $1$.
The root expansion at $t=0$ gives
$L_t(0)=7/384$, so that point is regular on the
coarse base.
\item The value tends to $\log\sqrt2$ at $1$ and
to $-\infty$ uniformly at infinity. Each ray meets
each lower level exactly once; the nonzero derivative
gives a continuous, locally regular radial graph.
That graph is the asserted Jordan curve.
\end{enumerate}

\subsection{The exact remaining questions}
\RungZero{Two questions remain open in the retained
scope: a literal coefficient atlas for the global
Bott--Chern representatives, and transcendence of the
normalized theta norm at the specified non-torsion
arithmetic points. The known global classes and the
exact height identities reduce these questions to
the concrete gaps recorded in the final part
(\ref{inv:C3.8}, \ref{ari:C.7.1}--\ref{ari:C.7.12}).}
\RungOne
\begin{enumerate}[label=\arabic*.]
\item The sheaf and double-complex calculations
already construct the global classes and determine
their dimensions and products. The remaining atlas
problem asks for the specified coefficient functions
on every finite and cusp chart with the stated
overlap identities.
\item At torsion points the division identity
already determines an algebraic value of the norm.
At the non-torsion points the proved relation is
$N(P)=A_t e^{-\widehat h_O(P)}$, with $A_t$
explicit and algebraic.
\item The period comparisons and the conductor-361
calculation identify the exact further
algebraic-independence input used by the attempted
transcendence route. That input is recorded as an
open step, together with every established identity;
no transcendence conclusion is asserted from it.
\end{enumerate}


\section{Common conventions and the marked construction}
\label{def:common-conventions}

\section*{Conventions and the notation dictionary}
We use $B=\Pp^1_t$, $p_1=0$, $p_2=1$, $p_0=\infty$, and
$B^\circ=B\setminus\{0,1,\infty\}$. Base meridians are counterclockwise
in a local holomorphic coordinate. The analytic generators are clockwise
root deck transformations, $s_j(g_jz)=e^{-2\pi i/m_j}s_j(z)$, with
$g_1g_2g_0=1$. A positive base loop ends at the inverse deck translate;
identifying its endpoint with the initial fibre applies $g_j$, giving
homology monodromy $A_j$. The retained multiplicities are $m_1=3,m_2=4$;
$f^*p_j=m_jS_j$ includes the nonreduced structure, whereas $S_j$ is reduced.
The cusp fibre $W=f^{-1}(\infty)$ is reduced.
\[
 \Lambda=\Z\langle\gh,\uh,\wh,\dh\rangle,\qquad
 \Pi_c=\begin{pmatrix}6\mu&\tau&1&0\\\beta^{\partt}+c&\mu&0&1\end{pmatrix},
 \qquad w=e^{-\pi ic}.
\]
We fix $(\beta^{\partt}+\tau)(p_0)=0$. Write
$D_{\partt}=\Im\beta^{\partt}-6(\Im\mu)^2/\Im\tau$ and
$M=\sup D_{\partt}$. Engel's line bundle is
$M_S=\OO_S(P-O)$; his indices $0,1,\infty$ correspond to our $1,2,0$.
In his common integral basis, the columns of the following matrix are
our ordered basis vectors:
\[
 P_{\rm mark}=\begin{pmatrix}
 0&-2&-1&-2\\2&-1&0&-1\\0&2&1&3\\-1&0&0&0
 \end{pmatrix},\quad
 T_0P_{\rm mark}=P_{\rm mark}A_1,\quad
 T_1P_{\rm mark}=P_{\rm mark}A_2,\quad
 \delta_{\rm E}=P_{\rm mark}\dh.
\]
For our finite twist vectors $v_1,v_2$, his positive-root affine vectors
are $a_0=-P_{\rm mark}v_1/3$, $a_1=-P_{\rm mark}v_2/4$.
In additive fibre coordinates $(\xi,\eta)$, put $v=e^{2\pi i\eta}$.
The first period acts by
$(\xi,v)\mapsto(\xi+6\mu,e^{2\pi i\beta_c}v)$; Engel's contracting
generator is its inverse, over $-6P$ with multiplier $e^{-2\pi i\beta_c}$.
We choose his reference linearization to be this transported inverse
period lift at $c=0$. Then $\lambda_0=1$ and $u=w^2$.
The comparison maps and the effect of changing this reference are
proved in \ref{fam:section:B.1} and B.2. Homology has integral coefficients unless stated
otherwise; tori and the threefold have their complex orientations.

Each statement below has its own proof or an explicit unresolved target
with its complete partial proof and the precise remaining step.


\section{Invariant and arithmetic notation in the family marking}
\label{inv:notation-dictionary}

\paragraph{Base, fibres, and the two peripheral orientations.}
We use $B=\mathbb P^1_t$, $p_1=0$, $p_2=1$, $p_0=\infty$,
$\Sigma=\{p_0,p_1,p_2\}$, and $D=W+S_1+S_2$ with reduced
components. The full fibres are $f^*p_1=3S_1$, $f^*p_2=4S_2$,
and $f^*p_0=W$. A positive base loop is counterclockwise in a
holomorphic local coordinate. The analytic root deck generator is
clockwise, $s_j(g_jz)=e^{-2\pi i/m_j}s_j(z)$, with
$g_1g_2g_0=1$. The lifted positive base loop ends at the inverse deck
translate; its endpoint identification applies $g_j$. Thus its
homology action is $A_j$, not $A_j^{-1}$.
The clockwise boundary orientation used in the collar integration
is stated separately in \cref{inv:C3.2}. References there to a
peripheral generator refer to this fixed analytic generator.

\paragraph{Integral marking and coordinate order.}
The homology basis is
$(\hat\gamma,\hat u,\hat w,\hat\delta)$, and
$(\gamma,u,w,\delta)$ denotes its dual cohomology basis.
The letter $w$ in that cohomology basis is a covector, whereas the
scalar family parameter below is $e^{-\pi ic}$.
In the ordered additive coordinates $(\zeta_1,\zeta_2)$,
\[
 \Pi_c=\begin{pmatrix}6\mu&\tau&1&0\\
             \beta_c&\mu&0&1\end{pmatrix},
 \qquad \beta_c=\beta^{\rm part}+c,
 \qquad (\beta^{\rm part}+\tau)(p_0)=0.
\]
Its rows are $\sigma_1=6\mu\gamma+\tau u+w$ and
$\sigma_2=\beta_c\gamma+\mu u+\delta$.
The actual monodromy matrices in this column basis are
\[
 A_1=\begin{pmatrix}1&0&0&0\\6&0&1&0\\-6&-1&-1&0\\-2&1&0&1\end{pmatrix},
 \quad
 A_2=\begin{pmatrix}1&0&0&0\\0&0&-1&0\\-6&1&0&0\\3&0&1&1\end{pmatrix},
 \quad A_0=(A_1A_2)^{-1}.
\]
No row interchange is implicit: if a local calculation orders the
coordinates as $(\zeta_2,\zeta_1)$, its period matrix is
$J\Pi_c$, $J=\left(\begin{smallmatrix}0&1\\1&0\end{smallmatrix}\right)$,
and its Hermitian coefficient matrix is $JH J^{\mathsf t}$.
Both the main invariant proofs and the retained working-edition
fibre formulas use the displayed $(\zeta_1,\zeta_2)$ order.

\paragraph{The exact Engel and reversed-meridian transitions.}
Engel's base indices $(0,1,\infty)$ correspond to $(1,2,0)$ here.
Writing his common integral basis as column coordinates, the family
basis is given by the columns of
\[
 P_{\rm mark}=\begin{pmatrix}
 0&-2&-1&-2\\2&-1&0&-1\\0&2&1&3\\-1&0&0&0
 \end{pmatrix},\qquad
 T_0P_{\rm mark}=P_{\rm mark}A_1,\quad
 T_1P_{\rm mark}=P_{\rm mark}A_2.
\]
The invariant source's reversed-meridian matrix is
\[
 Q=P_{\rm mark}S,\qquad
 S=\begin{pmatrix}-1&0&0&0\\-6&0&-1&0\\6&-1&0&0\\2&1&1&1\end{pmatrix},
 \qquad
 Q=\begin{pmatrix}2&-1&0&-2\\2&-1&0&-1\\0&2&1&3\\1&0&0&0\end{pmatrix}.
\]
It satisfies $T_0Q=QA_1^{-1}$ and $T_1Q=QA_2^{-1}$.
This is the explicit convention change of \cref{inv:B1.0}.
For invariant finite twist vectors $v_1,v_2$, Engel's positive-root
affine vectors are $a_0=-P_{\rm mark}v_1/3$ and
$a_1=-P_{\rm mark}v_2/4$; his distinguished cycle is
$P_{\rm mark}\hat\delta$.

\paragraph{Multiplicative parameters and special-fibre coordinates.}
Set $v=e^{2\pi i\zeta_2}$. The first period acts by
$(\zeta_1,v)\mapsto(\zeta_1+6\mu,e^{2\pi i\beta_c}v)$.
Engel's contracting generator is its inverse over $-6P$, on the
complement of the zero section of $\mathcal O(P-O)$.
With the reference inverse lift transported from $c=0$,
\[
 w=e^{-\pi ic},\qquad u_{\rm Engel}=w^2,
 \qquad 0<|w|<216,\quad 0<|u_{\rm Engel}|<46656.
\]
A different reference lift multiplies $u_{\rm Engel}$ by a fixed
$\lambda_0\in\mathbb C^*$. Thus $u_{\rm Engel}$ is distinct from
the integral covector $u$ and from an elliptic logarithm denoted $u$
in the period-transcendence section.
The special values use $\rho=e^{\pi i/3}$ rather than
$e^{2\pi i/3}=\rho-1$; the two conventions differ by the modular
translation $\tau\mapsto\tau+1$. The Albanese parameters and their
upper-half-plane representatives are
\[
 \begin{array}{c|c|c|c}
 &\kappa&\text{upper-half-plane period}&\text{nome}\\\hline
 S_1&c+\frac12-\frac{3i}{2\pi}\log48&-\kappa_1&
             q_1=e^{-2\pi i\kappa_1}=-w^2/48^3\\
 S_2&c+\frac32-\frac{3i}{\pi}\log6&-\kappa_2/2&
             q_2=e^{-\pi i\kappa_2}=iw/216.
 \end{array}
\]
The factor $1/2$ is the full order-four Albanese lattice projection.
In particular $q_1=(27/64)q_2^2$.

\paragraph{Determinants, orientations, and cohomology generators.}
Write $y=\Im\tau>0$, $m=\Im\mu$, $r=\Im\beta_c$, and
$D_c=r-6m^2/y=D_{\rm part}+\Im c$. The symbol $\Delta$ used
for this real quantity in the invariant fibre formulas means $D_c$,
not a Weierstrass discriminant. In admissible fibres $D_c<0$.
The real period determinant in the row order
$(\Re\zeta_1,\Im\zeta_1,\Re\zeta_2,\Im\zeta_2)$ is $yD_c$.
Consequently the marked real orientation is opposite to the complex
orientation, and for
$\eta_F=u\wedge w+6\gamma\wedge\delta$,
\[
 \int_{F,\operatorname{or}_{\mathbb C}}\eta_F^2=-12,
 \qquad
 \Gamma=\frac{i}{2D_c}
       \left(d\bar\zeta_2-\frac myd\bar\zeta_1\right)
       \quad\text{represents }\gamma^{0,1}.
\]
The subscript on $\eta_F$ separates this integral fibre class from
Dedekind's eta function and from Weierstrass quasi-periods.
The primitive global class is $\mathfrak g=g_\gamma$.
For standard data $(l_0,l_1,l_2)=(0,1,-1)$ and
$p=12l_0-4l_1-3l_2=-1$,
\[
 \exp(12\mathfrak g)=f^*\mathcal O_B(1),\quad
 \mathcal O_X(S_1-S_2)=\exp(\mathfrak g),\quad
 [G_0]=12\mathfrak g.
\]
In general the first equality is
$\exp(12g_\gamma)=f^*\mathcal O_B(-p)$.
Here $\exp$ means the sheaf map $z\mapsto e^{2\pi iz}$,
$\int_B\omega_{FS}=1$, $c_{1,BC}(L)=[iF_L/(2\pi)]$, and the
unnormalized Chern--Ricci form represents $2\pi c_{1,BC}(T_X)$.
Thus the prescribed Chern--Ricci form is $\pi f^*\omega_{FS}$.

\paragraph{Root and coarse line bundles.}
For the root curve $\mathcal B$ of orders $(3,4)$ put
$R=R_1R_2^{-1}$, so
$R_1=R^4$, $R_2=R^3$, $\rho^*\mathcal O_B(1)=R^{12}$,
and $K_{\mathcal B}=R^{-7}$.
The notation $\mathcal O_{\mathcal B}(n)$ means $R^n$;
$\mathcal O_B(n)$ is the ordinary coarse degree-$n$ line bundle.
They have different degrees and stabilizer characters.
The logarithmic sheaves are
$\mathcal L=\Omega_B^1(\log\Sigma)$,
$\mathcal A=\Omega_X^1(\log D)$, and
$\mathcal E=\mathcal A/f^*\mathcal L$.
Every coherent base-change assertion uses the full scheme fibre.
The finite cyclic covers and subsequent invariant functor retain its
nilpotent normal directions.

\paragraph{Theta norm, modular coordinate, and heights.}
The theta variable has periods $(1,\tau)$; software using a first
argument of period $\pi$ receives $\pi z$. With $q=e^{2\pi i\tau}$,
\[
 N(z,\tau)=\left|\frac{\vartheta_1(z,\tau)}{\eta(\tau)}\right|
       e^{-\pi(\Im z)^2/\Im\tau},\qquad
 D_{\rm part}=\frac6\pi\log\frac{N(\mu,\tau)}{2\sqrt3},
 \qquad t=\frac{j(\tau)}{1728}.
\]
The modular discriminant is
$\Delta_0=q\prod_{n\ge1}(1-q^n)^{24}$, so
$t=E_4^3/(1728\Delta_0)$ and $t-1=E_6^2/(1728\Delta_0)$.
For the rational family
\[
 E_t:Y^2=X^3-\tfrac34t^3(t-1)X+\tfrac14t^4(t-1)^2,
 \qquad P_t=(0,t^2(t-1)/2),
\]
the invariant differential is $dX/(2Y)$ and
$\Delta_E=27t^8(t-1)^3$.
In the lattice model $X=\wp(u)$ and $Y=\wp'(u)/2$;
the lattice discriminant is
$g_2^3-27g_3^2=(2\pi)^{12}\eta(\tau)^{24}$ when the
period lattice is $\mathbb Z+\tau\mathbb Z$.
We use the height of $(O)$,
$\widehat h_O(P)=\frac12\lim n^{-2}h(x(nP))$.
PARI's \texttt{ellheight}, the surface self-pairing, and the
rank-one BSD regulator use $2\widehat h_O$.
The full real period is $\Omega=\int_{E(\mathbb R)}|\omega|$,
including both components if disconnected.
The archimedean local height is exactly $-\log N$.
At finite places $|p|_v=p^{-1}$ and $n_v=[K_v:\mathbb Q_p]$;
complex embeddings are counted separately in the global sum.

For period arguments, $\eta_j$ is defined by
$\zeta(z+\omega_j)=\zeta(z)+\eta_j$, with
$\omega_1\eta_2-\omega_2\eta_1=-2\pi i$.
The real-linear extension of these quasi-periods is denoted $\eta(u)$.
In the critical-point proof, the Cayley coordinate is
$\chi=(\tau-i)/(\tau+i)$, $z=k\chi^2$, and
$\beta_{\rm cusp}=\log(5+2\sqrt6)/(2\pi)$;
neither symbol is the family parameter $w$ or the period $\beta_c$.

\paragraph{Shared statements.}
The maximum has the single label \cref{fam:src:C.2.3}, with aliases
\texttt{ari:C.8.5} and \texttt{push:src:E.4.3}; the independent
interval, duplication, and radial arguments are all retained.
Tangent cohomology, the parameter Kodaira--Spencer map, and the
one-dimensional Kuranishi germ are collected in
\cref{fam:src:A.5.2}, with aliases
\texttt{inv:D1.3}, \texttt{inv:D2.1}, and \texttt{inv:D3.1}.
The exact parameter classification is \cref{fam:src:A.5.1};
the canonical-pencil statement is \cref{fam:src:A.1.1}.
The working Hodge-line proof is \CompendiumRef{I}{mapsI:hodge:line},
and the finite normal-order calculation is
\CompendiumRef{I}{wb:ffcert:thm:normal-trivialization}.
Working structure-sheaf and cusp statements share \cref{inv:A1.0};
working Picard restriction shares \cref{inv:C3.1};
the global fibre-class extension shares \cref{inv:A3.2};
the Bott--Chern multiplication statement shares \cref{inv:C3.3}.



\begin{lemma}[The exact comparison of the two Weierstrass equations]
\label{I:weierstrass-comparison}
\StatementRecord{I:weierstrass-comparison}{proved}{No earlier separately identified check is recorded. Checked only in this writing. Pre-continuation writing events: GPT-6 Astra (OpenAI), ChatGPT Pro mode, effort Pro, S6\_10\_PAPERS, 18 November 2026: targeted proof check. Continuation metadata events: GPT-6 Astra (OpenAI), ChatGPT Pro mode, effort Pro, S6\_10\_PAPERS continuation, 19 November 2026: used and rechecked in this run; status-vocabulary reconciliation; metadata reconciliation only; no new mathematical check. \emph{Pre-continuation recorded proof scope.} Checked only in this writing by GPT-6 Astra (OpenAI), ChatGPT Pro mode, effort Pro, 18 November 2026; direct substitution and the lattice homothety below}{\ref{ari:C.1.1}; the displayed equations}
Put $C=t^3(t-1)$ and $D=t^4(t-1)^2$, with $t\ne0,1$.
The pointed map
\[
 \Phi:E_t^{\rm ar}\longrightarrow E_t^{\rm geom},\qquad
 (x,y)\longmapsto(2x,-2\sqrt2y)
\]
is an isomorphism from
$y^2=x^3-3Cx/4+D/4$ to $Y^2=X^3-3CX+2D$.
It sends $(0,t^2(t-1)/2)$ to
$(0,-\sqrt2t^2(t-1))$ and preserves the normalized theta norm
under the induced pointed elliptic-curve identification.
The discriminants satisfy $\Delta_{\rm geom}=64\Delta_{\rm ar}$.
\end{lemma}
\begin{proof}
\begin{enumerate}[label=\arabic*.]
\item Substitution gives
$Y^2=8y^2=8x^3-6Cx+2D=X^3-3CX+2D$.
The inverse is $(X,Y)\mapsto(X/2,-Y/(2\sqrt2))$ and both maps
extend over the unique point at infinity. The displayed section
identity follows by the same substitution.
\item Set $a=-\sqrt2$. In Weierstrass uniformization the equations
are related by $X=a^2x$, $Y=a^3y$ and
$\Lambda_{\rm geom}=a^{-1}\Lambda_{\rm ar}$.
The induced elliptic logarithm map is $z\mapsto a^{-1}z$.
Writing $\Lambda_{\rm ar}=\omega(\Z+\tau\Z)$ shows directly
that the argument $z/\omega$ in the defining theta quotient is
unchanged. Thus the norms agree through this map.
\item The discriminant scales by $a^{12}=64$.
The divisor and leading term of the division polynomial give
$\psi_m(\Phi(Q))=a^{m^2-1}\psi_m(Q)$.
The two absolute scaling factors in
$|\Delta_E|^{(m^2-1)/12}/|\psi_m(Q)|$ cancel exactly.
\end{enumerate}
\end{proof}

\subsection{A typed fibre dictionary}\label{fam:section:B.1}

\begin{proposition}\label{fam:src:B.1.1}\label{fam:st:B.1.1}
\StatementRecord{fam:src:B.1.1}{proved}{Historical events: Claude Fable 5.2 (Anthropic), effort high, 12 November 2026: targeted proof check. Claude Fable 5.2 (Anthropic), effort high, 16 November 2026: used and rechecked in this run. GPT-6 Astra (OpenAI), ChatGPT Pro mode, effort Pro, S6\_09\_PAPERS, 17 November 2026: targeted proof check. Also checked in this writing. Pre-continuation writing events: GPT-6 Astra (OpenAI), ChatGPT Pro mode, effort Pro, S6\_10\_PAPERS, 18 November 2026: targeted proof check. Continuation metadata events: GPT-6 Astra (OpenAI), ChatGPT Pro mode, effort Pro, S6\_10\_PAPERS continuation, 19 November 2026: used and rechecked in this run; status-vocabulary reconciliation; metadata reconciliation only; no new mathematical check. \emph{Pre-continuation recorded proof scope.} Read complete source proof and re-derived the displayed centralizer, finite norm, cusp winding, cohomological gluing, period or reflection argument as applicable; external input locators are separately being verified.}{the hypotheses, notation and definitions stated here; \cref{fam:src:A.3.1}}

Let $E_\tau=\mathbb C/(\mathbb Z+\tau\mathbb Z)$, $P=[\mu]$, and $L_\mu=\mathcal O_{E_\tau}(P-O)$. The holomorphic map
\[
 [\xi,\eta]\longmapsto[\xi,e^{2\pi i\eta}]
\]
identifies $\mathbb C^2/\Pi_c\Lambda$ with the quotient of $L_\mu^\times$ by the lift
$A(\xi,W)=(\xi+6\mu,e^{2\pi i\beta_c}W)$. Engel's generator, in Definition 2.3 of the attached paper, is its inverse, with translation $-6P$ and multiplier $e^{-2\pi i\beta_c}$. The section $P=[\mu]$ is a generator of the Mordell--Weil group of the stated pointed Weierstrass surface.
\end{proposition}


\begin{proof}
\begin{enumerate}[label=\arabic*.,leftmargin=*]
\item First quotient by columns two, three and four. Exponentiation kills column four; columns three and two then give
\[
 (\xi,W)\sim(\xi+1,W)\sim(\xi+\tau,e^{2\pi i\mu}W).
\]
For the odd theta function with periods $1,\tau$, the quotient $\theta(\xi-\mu)/\theta(\xi)$ is a meromorphic section of this line bundle: its factors are $1,e^{2\pi i\mu}$ and its divisor is $P-O$. Sending its divisor section to this quotient is the asserted line-bundle isomorphism. Quotienting the remaining first column gives exactly $A$. On sufficiently small simply connected charts the logarithm of $W$ is a holomorphic inverse; changes of logarithm are column-four periods. Thus the displayed map is a biholomorphism of the two quotient complex spaces, with its inverse explicitly defined by local logarithms. Choosing $A^{-1}$ changes the generator of the same infinite cyclic group and gives precisely Engel's declared inverse-translation convention.

\item For clarity the elliptic surface in this dictionary is the displayed cubic, rather than only a curve with the same $j$-value. Let $f_{\rm sc}$ be the branch satisfying
\[
 f_{\rm sc}^4=\frac{12t^3(t-1)}{g_2(\tau)},\qquad
 f_{\rm sc}^6=-\frac{8t^4(t-1)^2}{g_3(\tau)},\qquad
 f_{\rm sc}/t\longrightarrow i\sqrt3/\pi.
\]
The identities $g_2^3/(g_2^3-27g_3^2)=t$ make these equations compatible. Their zero and pole orders give a holomorphic branch on the period cover; the indicated cusp sign fixes it. The map $\xi\mapsto(f_{\rm sc}^2\wp_\tau(\xi),\frac12f_{\rm sc}^3\wp'_\tau(\xi))$ is a pointed isomorphism to
$y^2=x^3-3t^3(t-1)x+2t^4(t-1)^2$. Its factors are $f_{\rm sc}(g_1z)=-f_{\rm sc}/\tau$, $f_{\rm sc}(g_2z)=f_{\rm sc}/\tau$; therefore this isomorphism descends over $B^\circ$.

\item We first verify that $P=[\mu]$ gives a rational section of the resolved elliptic surface. On $B^\circ$ its cubic coordinates descend by the displayed scale and affine laws. At the two finite centres, $\mu_j$ is not a lattice point, while $f_{\rm sc}$ has orders two and one in the root parameter. Its cubic coordinates are therefore holomorphic and invariant, hence holomorphic in the base coordinate. At the cusp $\mu$ is holomorphic and bounded, and the convergent multiplicative expansion of $\wp$ gives meromorphic cubic coordinates with finite pole orders. Even when its exponential meets one, the pole is controlled by its finite vanishing order; it cannot be identically one because the two finite values of $\mu$ differ. The coordinates are thus meromorphic on $\mathbb P^1$ and rational. The closure of its graph on the proper resolved surface is an irreducible curve; its normalization maps finitely with degree one to the smooth base, so is the base itself. This gives the unique section of that surface extending $P$.

\item We also verify that the specified section is a Mordell--Weil generator using its full logarithm. A rational section of this cubic pulls back to a holomorphic section of the good elliptic family at both finite covering points: multiply its projective coordinates by a common power to remove poles, and use the nonsingular limiting Weierstrass equation. At the cusp put $v=1/t$, $X=v^2x,Y=v^3y$. Its equation is $Y^2=X^3-3(1-v)X+2(1-v)^2$. A section cannot pass through the node $(1,0)$: substitution of $X=1+O(v)$, $Y=O(v)$ leaves the coefficient $+1$ of $v$ on the left-hand side of the defining equation. At every other central point the multiplicative uniformization is a local isomorphism. Explicitly, at $v=0$ it is
\[
 X=1+\frac{12U}{(U-1)^2},\qquad
 Y=\frac{12\sqrt3 U(U+1)}{(U-1)^3},\qquad U=e^{2\pi i\xi},
\]
with inverse $U=(Y/(X-1)+\sqrt3)/(Y/(X-1)-\sqrt3)$ away from $O$; a local parameter at $O$ is $X/Y$. The convergent lattice expansion of $\wp$ gives the same holomorphic map for small $v$. Hence a section has a holomorphic unit $U(v)$ and a bounded logarithm at the cusp. Lifting the section on the simply connected cover gives a global logarithm $\nu$, with this cusp branch. Its only remaining ambiguity is an integer, since a nonzero multiple of $\tau$ is unbounded.

\item Its two affine laws have integral defects:
\[
 \nu(g_1z)=-\nu/\tau+\alpha_1+\beta_1\tau(g_1z),\quad
 \nu(g_2z)=\nu/\tau+\alpha_2+\beta_2\tau(g_2z).
\]
Composing them and using cusp invariance gives $\alpha_2=-\beta_1$, $\beta_2=\alpha_1+\beta_1$. Put $n=-\alpha_1-2\beta_1$ and $a=\alpha_1+\beta_1$. Substitution shows that $n\mu+a$ has exactly these defects. Their difference is a global section of the first-coordinate sheaf $\mathcal O_B(-p_0)$ from \ref{fam:src:A.3.1}, so is zero. Every rational section is therefore $nP$. The same cusp argument and the distinct fixed values of $\mu$ show that $P$ has infinite order. This proves $\operatorname{MW}(S)=\mathbb ZP$.
\end{enumerate}
\end{proof}

\begin{proposition}\label{fam:src:B.1.2}\label{fam:st:B.1.2}
\StatementRecord{fam:src:B.1.2}{proved}{Historical events: Claude Fable 5.2 (Anthropic), effort high, 12 November 2026: targeted proof check. Claude Fable 5.2 (Anthropic), effort high, 16 November 2026: used and rechecked in this run. GPT-6 Astra (OpenAI), ChatGPT Pro mode, effort Pro, S6\_09\_PAPERS, 17 November 2026: targeted proof check. Also checked in this writing. Pre-continuation writing events: GPT-6 Astra (OpenAI), ChatGPT Pro mode, effort Pro, S6\_10\_PAPERS, 18 November 2026: targeted proof check. Continuation metadata events: GPT-6 Astra (OpenAI), ChatGPT Pro mode, effort Pro, S6\_10\_PAPERS continuation, 19 November 2026: used and rechecked in this run; status-vocabulary reconciliation; metadata reconciliation only; no new mathematical check. \emph{Pre-continuation recorded proof scope.} Read complete source proof and re-derived the displayed centralizer, finite norm, cusp winding, cohomological gluing, period or reflection argument as applicable; external input locators are separately being verified.}{the hypotheses, notation and definitions stated here; the explicit constructions and calculations in the proof}

The exact set of nonzero parameters for which every ordinary fibre and both normalized finite root fibres of Engel's period-reference quotient are compact complex two-tori is
\[
 0<|u|<e^{-2\pi M}.
\]
This is the parameter domain for the compact family with both smooth finite root models. The condition on ordinary fibres alone at equality is determined by whether a maximizing point of $D_{\rm part}$ is ordinary; this boundary issue is treated in Part C. For an independent reference, the radius is $|\lambda_0|e^{-2\pi M}$.
\end{proposition}


\begin{proof}
\begin{enumerate}[label=\arabic*.,leftmargin=*]
\item The function $N(\xi,W)=|W|\exp(2\pi\operatorname{Im}\mu\,\operatorname{Im}\xi/\operatorname{Im}\tau)$ is invariant under both line-bundle period identifications. The inverse lift multiplies it by $e^{2\pi D_c}$. More directly the real determinant of the four period columns is a nonzero sign times $(\operatorname{Im}\tau)D_c$. Thus the four columns generate a full lattice exactly when $D_c\ne0$, and then the quotient is a compact torus by a closed fundamental parallelepiped. At $D_c=0$ their real span has dimension three; a rank-four lattice in $\mathbb R^4$ cannot result, so the quotient cannot be a compact complex two-torus. The function $D_{\rm part}$ descends continuously to the connected base with the cusp removed and tends to $-\infty$ at the cusp. Its image is therefore $(-\infty,M]$: the maximum is attained, and the intermediate value theorem fills all smaller values. Consequently $D_c$ is nonzero at every base point exactly when $\operatorname{Im}c<-M$. Since $|u|=e^{2\pi\operatorname{Im}c}$, the assertion follows, with the indicated reference scaling.
\end{enumerate}
\end{proof}

\begin{proposition}\label{fam:src:B.1.3}\label{fam:st:B.1.3}
\StatementRecord{fam:src:B.1.3}{proved}{Historical events: Claude Fable 5.2 (Anthropic), effort high, 12 November 2026: targeted proof check. Claude Fable 5.2 (Anthropic), effort high, 16 November 2026: used and rechecked in this run. GPT-6 Astra (OpenAI), ChatGPT Pro mode, effort Pro, S6\_09\_PAPERS, 17 November 2026: targeted proof check. Also checked in this writing. Pre-continuation writing events: GPT-6 Astra (OpenAI), ChatGPT Pro mode, effort Pro, S6\_10\_PAPERS, 18 November 2026: targeted proof check. Continuation metadata events: GPT-6 Astra (OpenAI), ChatGPT Pro mode, effort Pro, S6\_10\_PAPERS continuation, 19 November 2026: used and rechecked in this run; status-vocabulary reconciliation; metadata reconciliation only; no new mathematical check. \emph{Pre-continuation recorded proof scope.} Read complete source proof and re-derived the displayed centralizer, finite norm, cusp winding, cohomological gluing, period or reflection argument as applicable; external input locators are separately being verified.}{the hypotheses, notation and definitions stated here; \cref{fam:src:B.1.1}}

For the original affine logarithm $\mu$, we have $\wp_\tau(\mu)=0$ identically. Equivalently the generator $P=[\mu]$ of \ref{fam:src:B.1.1} is one of the two sections $P_*=(0,\sqrt2t^2(t-1))$ and $-P_*$. This assertion does not use a numerical evaluation of $\mu$ or of $\beta$.
\end{proposition}


\begin{proof}
\begin{enumerate}[label=\arabic*.,leftmargin=*]
\item We give the intersection calculation that identifies the generator. At $t=1$, put $u=t-1$ and blow up the origin. In the $u$-chart $x=uX,y=uY$, the strict equation is
$Y^2=uX^3-3(1+u)^3X+2(1+u)^4$. Its exceptional curve is smooth, and $P_*$ meets it at $(X,Y)=(0,\sqrt2)$. In the other chart $u=xU,y=xY$, the equation is $Y^2=x-3(1+xU)^3U+2(1+xU)^4U^2$. The implicit function theorem solves for $x$ at the common point; the two fibre branches $U=0,x=0$ meet with intersection multiplicity two. These charts cover the tangent conic and resolve the singularity. The two components have multiplicity one; the one met by $P_*$ is the nonidentity component $D$, with $D^2=-2$ from intersection with the full fibre.

\item At $t=0$, two blow-ups in the successive $t$-charts give
\[
 Y^2=tX^3+3t^2(1-t)X+2t^2(1-t)^2,
\quad
 Y_2^2=t^2X_2^3+3t(1-t)X_2+2(1-t)^2.
\]
The second equation has two smooth multiplicity-one components $E_{2,\pm}$ at $Y_2=\pm\sqrt2$, and $P_*$ meets one of them. The remaining singular point, in the second $X$-chart, has
\[
 V^2=H(X,T),\quad H=X^2T+(2+3X)T^2-(4X+3X^2)T^3+2X^2T^4,
 \qquad t=XT.
\]
Let $T_0(X)$ solve $H_T(X,T_0)=0$. It has $T_0=-X^2/4+O(X^3)$. Taylor's formula gives
$H-H(X,T_0)=(T-T_0)^2h_0$, $h_0(0,0)=2$, and $H(X,T_0)=X^4e(X)$, $e(0)=-1/8$. These identities follow by substitution in the displayed polynomial, not by discarding its higher terms. With $Z=(T-T_0)\sqrt{h_0}$, $a=V-Z$, $b=(V+Z)/e(X)$, the equation becomes exactly $ab=X^4$ through a local analytic isomorphism. Its resolution has charts
\[
 a=U_i^{i+1}V_i^i,\quad b=U_i^{3-i}V_i^{4-i},\quad X=U_iV_i
 \quad(0\le i\le3),
\]
with changes $U_{i+1}=1/V_i,V_{i+1}=U_iV_i^2$. They give a chain of three $(-2)$-curves. These are the charts of the two successive point blow-ups resolving $ab=X^4$, so the resolution is proper. Retaining $t=XT$ in the comparison gives multiplicities $2,3,2$ on that chain. The two new multiplicity-one components meet its ends; the first exceptional curve meets its middle component, and the original strict fibre meets the other end of that first exceptional curve. All these statements follow by substituting the displayed monomials and their orders; the only fibre components of multiplicity one are the strict component and $E_{2,\pm}$.

\item Order the six nonidentity components by the graph
$\Theta_1-\Theta_2-\Theta_3-\Theta_4$ and $\Theta_2-\Theta_5-\Theta_6$, with $E_{2,\pm}=\Theta_4,\Theta_6$. Their intersection matrix is $-C$, where
\[
 C=\begin{pmatrix}2&-1&0&0&0&0\\-1&2&-1&0&-1&0\\0&-1&2&-1&0&0\\0&0&-1&2&0&0\\0&-1&0&0&2&-1\\0&0&0&0&-1&2\end{pmatrix},\quad
 C^{-1}=\frac13\begin{pmatrix}6&9&6&3&6&3\\9&18&12&6&12&6\\6&12&10&5&8&4\\3&6&5&4&4&2\\6&12&8&4&10&5\\3&6&4&2&5&4\end{pmatrix}.
\]
Multiplication verifies this inverse. Self-intersections follow as well from the complete multiplicity vector and $F\cdot\Theta_i=0$.

\item The form $dx\wedge dt/y$ is holomorphic and nonvanishing on the finite resolved charts. In the last charts its expression is a unit times $dX\wedge da/a=-dU_i\wedge dV_i$; the earlier blow-up equations cancel the apparent denominators at $y=0$. At infinity it is $-v^{-1}dX\wedge dv/Y$ and has one fibre pole. Thus $K_S=-F$. Adjunction gives $Q^2=-1$ for every section $Q$, since $F\cdot Q=1$. A section meets a unique multiplicity-one component transversely at a smooth fibre point: the pullback of the full base equation to the section has order one, excluding all other possibilities.

\item Let $v_Q$ be the intersection vector with $\Theta_i$ and let $d_Q=Q\cdot D$. The possibilities are $v_Q=0,e_4,e_6$ and $d_Q=0,1$. Define the rational divisor class
\[
 \varphi(Q)=Q-O-(Q\cdot O+1)F+\sum\Theta_i(C^{-1}v_Q)_i+\tfrac12d_QD.
\]
It is orthogonal to $O,F,\Theta_i,D$ by direct multiplication. The chord-and-tangent rational function on the generic cubic has divisor $(Q)+(R)-(Q+R)-(O)$; extending it to the surface adds only vertical components. Applying the displayed orthogonal projection therefore proves $\varphi(Q+R)=\varphi(Q)+\varphi(R)$. Expansion gives
\[
 h(Q)=-\varphi(Q)^2=2+2Q\cdot O-v_Q^tC^{-1}v_Q-\tfrac12d_Q.
\]
For $Q\ne O$, the nonnegative intersection $Q\cdot O$ gives $h(Q)\ge2-4/3-1/2=1/6$. The visible $P_*$ is disjoint from $O$ everywhere, and the explicit blow-up charts place it on both nonidentity components. Hence $h(P_*)=1/6$. The logarithm argument in \ref{fam:src:B.1.1} already proved $P_*=nP$ for an integer $n\ne0$, without identifying the visible section. Additivity gives $1/6=n^2h(P)\ge n^2/6$, so $n=\pm1$. The pointed cubic map in \ref{fam:src:B.1.1} has $x(P)=f_{\rm sc}^2\wp_\tau(\mu)$ with $f_{\rm sc}\ne0$ on the ordinary locus. This proves the asserted zero identity.
\end{enumerate}
\end{proof}

\begin{proposition}\label{fam:src:B.1.4}\label{fam:st:B.1.4}
\StatementRecord{fam:src:B.1.4}{proved}{Historical events: Claude Fable 5.2 (Anthropic), effort high, 12 November 2026: targeted proof check. Claude Fable 5.2 (Anthropic), effort high, 16 November 2026: used and rechecked in this run. GPT-6 Astra (OpenAI), ChatGPT Pro mode, effort Pro, S6\_09\_PAPERS, 17 November 2026: targeted proof check. Also checked in this writing. Pre-continuation writing events: GPT-6 Astra (OpenAI), ChatGPT Pro mode, effort Pro, S6\_10\_PAPERS, 18 November 2026: targeted proof check. Continuation metadata events: GPT-6 Astra (OpenAI), ChatGPT Pro mode, effort Pro, S6\_10\_PAPERS continuation, 19 November 2026: used and rechecked in this run; status-vocabulary reconciliation; metadata reconciliation only; no new mathematical check. \emph{Pre-continuation recorded proof scope.} Read complete source proof and re-derived the displayed centralizer, finite norm, cusp winding, cohomological gluing, period or reflection argument as applicable; external input locators are separately being verified.}{the hypotheses, notation and definitions stated here; \cref{fam:src:B.1.3}; \cref{fam:src:B.1.1}; \cite[Lemma 2.1]{Engel2026}; \cite[Corollary 2.2]{Engel2026}}

The line-bundle comparison of the period family with Engel's family over $B^\circ$ transports its inverse period lift to a global section of $\operatorname{Hom}(t_{6P}^*M_S,M_S)$ with divisor $F_\infty$. Take this transported section for $c=0$ as the reference linearization. Then $\lambda_0=1$ exactly and $u=e^{-2\pi ic}=w^2$. For another independently fixed reference the precise constant is its scalar ratio with this section. With the cusp branch fixed in Part C and $f_{\rm sc}/t\to i\sqrt3/\pi$, the section is $P=(0,-\sqrt2t^2(t-1))$.
\end{proposition}


\begin{proof}
\begin{enumerate}[label=\arabic*.,leftmargin=*]
\item The zero identity $\wp_\tau(\mu)=0$, proved without numerical input in \ref{fam:src:B.1.3}, identifies $P$ with the signed visible section. With the cusp branch in Part C and the specified scale, substitution of $U=-(5+2\sqrt6)$ in the nodal uniformization of \ref{fam:src:B.1.1} gives $Y=-\sqrt2$, so $P=(0,-\sqrt2t^2(t-1))$.

\item To compare the line bundles as families, use their common divisor $P-O$ and rigidify at the zero section. Put $F_\theta(\xi)=\theta(\xi-\mu)/\theta(\xi)$, $r_\theta=\theta(-\mu)/\theta'(0)$, and let $s_{P-O}$ be the canonical meromorphic divisor section. The map
\[
 [\xi,W]\longmapsto\frac{2r_\theta}{f_{\rm sc}}
        W\frac{s_{P-O}(\xi)}{F_\theta(\xi)}
\]
is a holomorphic line-bundle isomorphism, including at $P,O$ after cancellation of the equal divisors. At $O$ it sends the chosen frame to $dx/y$, since $F_\theta\sim r_\theta/\xi$ and $x/y\sim-\xi/f_{\rm sc}$ with the corresponding canonical divisor frame. Two such maps that agree at $O$ differ by a holomorphic scalar on each compact elliptic fibre and hence agree everywhere. This proves descent under the actual period identifications.

\item Engel's \cite[Lemma 2.1]{Engel2026}, which we have read together with its proof, identifies $\operatorname{Hom}(t_{6P}^*M_S,M_S)$ with $\pi_S^*\mathcal O_B(1)$. The lemma applies because the preceding pointed cubic map and Mordell--Weil calculation identify his $(S,O,P)$ with ours, including the chosen sign. To determine which section is transported, it suffices to compute its order at one generic smooth point of each special fibre, since a section of this pulled-back line is a base section. At a finite centre write $6\mu=p+r\tau+\delta$ with $(p,r)=(4,-2),(3,-3)$ and $\delta/f_{\rm sc}$ holomorphic; this divisibility follows by linearizing the affine laws, whose homogeneous orders are two and one. Evaluate at $\xi=f_{\rm sc}z$ for generic fixed $z$. The factor $(2r_\theta/f_{\rm sc})/F_\theta(f_{\rm sc}z)$ tends to $2z$, a unit. Restoring the lattice shift $p+r\tau$ at the target changes the fibre multiplier to $e^{-2\pi i(\beta_c-r\mu)}$, whose finite limits are $e^{-2\pi i\kappa_1}$ and $e^{-2\pi i\kappa_2}$. The source and target theta ratios remain units. Therefore the base section has order zero at both finite points. At the cusp $f_{\rm sc}$ has a simple pole and $r_\theta$ has a nonzero limit; the two factors $2r_\theta/f_{\rm sc}$ cancel between source and target. The remaining multiplier is
$e^{-2\pi i\beta_c}=q e^{-2\pi i(\beta_c+\tau)}$, of exact order one. Thus the transported section has divisor $F_\infty$. Engel's \cite[Corollary 2.2]{Engel2026} now identifies it with a nonzero scalar multiple of his reference, since two base sections with the same divisor have constant ratio. Changing $\beta$ by $c$ multiplies it exactly by $e^{-2\pi ic}$. With our stated period-reference normalization, this scalar is one. For any independently chosen reference $\phi_{\rm ref}$, the scalar is instead the precisely defined ratio $\lambda_0=\phi_{\rm part}/\phi_{\rm ref}$, evaluated on any ordinary fibre. Scaling the reference by $a\in\mathbb C^*$ changes $\lambda_0$ to $\lambda_0/a$; no scalar value is invariant under that declared choice.
\end{enumerate}
\end{proof}


\section{The normalized periods, quantitative bounds, and the boundary}\label{fam:section:C}

\subsection{A theta expression for the missing period}\label{fam:section:C.1}
We use the odd Jacobi function
\[
 \vartheta_1(v,\tau)=2\sum_{n\geq0}(-1)^n
 e^{\pi i\tau(n+1/2)^2}\sin((2n+1)\pi v)
\]
and the Dedekind function
\(\eta_D(\tau)=e^{\pi i\tau/12}\prod_{n\geq1}(1-e^{2\pi in\tau})\).
The subscript distinguishes this function from the invariant integral class.
We choose the cusp lift reached from \(\tau=iT\), \(T>1\), on which
\(\mu=1/2-iy(T)\), \(0<y(T)<T/2\). Put
\(a=5+2\sqrt6\) and \(\rho=e^{\pi i/3}\).

\begin{proposition}\label{fam:src:C.1.1}\label{fam:st:C.1.1}
\StatementRecord{fam:src:C.1.1}{proved}{Historical events: Claude Fable 5.2 (Anthropic), effort high, 12 November 2026: targeted proof check. Claude Fable 5.2 (Anthropic), effort high, 16 November 2026: used and rechecked in this run. Claude Fable 5.2 (Anthropic), effort high, 17 November 2026: carried as background without fault found. GPT-6 Astra (OpenAI), ChatGPT Pro mode, effort Pro, S6\_09\_PAPERS, 17 November 2026: targeted proof check. Also checked in this writing. Pre-continuation writing events: GPT-6 Astra (OpenAI), ChatGPT Pro mode, effort Pro, S6\_10\_PAPERS, 18 November 2026: targeted proof check. Continuation metadata events: GPT-6 Astra (OpenAI), ChatGPT Pro mode, effort Pro, S6\_10\_PAPERS continuation, 19 November 2026: used and rechecked in this run; status-vocabulary reconciliation; metadata reconciliation only; no new mathematical check. \emph{Pre-continuation recorded proof scope.} Read source proof; re-derived theta transformation, determinant/norm comparison, distribution constant or stated analytic inequality as applicable. Numerical replay scopes are recorded with their exact statement labels; the analytic argument is checked here.}{the hypotheses, notation and definitions stated here; \cref{fam:src:B.1.3}}

For the period section identified in Statement \ref{fam:src:B.1.3}, the particular solution
with \((\beta^{\rm part}+\tau)(p_0)=0\) is
\begin{equation}\label{fam:eq:cd-beta}
 \beta^{\rm part}=\frac{6i}{\pi}\operatorname{Log}
 \frac{\vartheta_1(\mu,\tau)}{2\sqrt3\,\eta_D(\tau)}.
\end{equation}
Here the logarithm is continued on the simply connected uniformizing
half-plane from its cusp expansion. Its special values are
\[
 \mu_0=\frac12-\frac{i\log a}{2\pi},\qquad
 \mu(z_1)=\frac{2-\rho}{3},\qquad \mu(z_2)=\frac{1-i}{2},
\]
\[
 \beta^{\rm part}(z_2)=i\left(\frac32-\frac3\pi\log6\right),
 \qquad \operatorname{Re}\beta^{\rm part}(z_1)=-\frac12.
\]
In particular \(e^{2\pi i\mu_0}=-a\).
\end{proposition}


\begin{proof}
\begin{enumerate}[label=\arabic*.,leftmargin=*]
\item Statement \ref{fam:src:B.1.3} identifies the elliptic point represented by \(\mu\) with
one of the sections \((0,\pm\sqrt2t^2(t-1))\). The uniformization map
\((v,\tau)\mapsto(f^2\wp_\tau(v),f^3\wp'_\tau(v)/2)\) therefore gives
\(\wp_\tau(\mu)=0\). No zero of \(\wp\) is a lattice point. Thus the numerator
of \eqref{fam:eq:cd-beta} is nonzero; the denominator is nonzero as well.
A holomorphic logarithm exists on the simply connected cover.

\item For completeness the cusp value follows directly from the normally
convergent Fourier expansion, rather than from a quoted Eichler--Zagier
constant:
\begin{equation}\label{fam:eq:cd-wp}
 \frac{\wp_\tau(v)}{\pi^2}
 =\csc^2(\pi v)-\frac13+
 8\sum_{n\geq1}\frac{nq^n}{1-q^n}(1-\cos(2\pi nv)),
 \qquad q=e^{2\pi i\tau}.
\end{equation}
This identity follows by differentiating the logarithm of the Jacobi
product twice and using
\(E_2=1-24\sum nq^n/(1-q^n)\). Boundedness of the chosen elliptic logarithm
at the cusp permits passage to the limit. We obtain
\(\sin^2(\pi\mu_0)=3\), equivalently
\(u^2+10u+1=0\), \(u=e^{2\pi i\mu_0}\).
On the chosen imaginary-axis branch \(\operatorname{Im}\mu<0\), so \(u=-a\).
One can see the branch without a numerical convention: on the rectangular
elliptic curve with \(\tau=iT\), \(T>1\), \(\wp\) is real on the segment
\(1/2-iy\), \(0\leq y\leq T/2\). Its endpoint values are the largest and
middle real roots of its cubic, respectively positive and negative.
Its derivative has no zero in the open segment, since the zeros of
\(\wp'\) are the three nonzero half-periods. There is exactly one zero
there; this gives the indicated continuation. Consequently
\(2\sin(\pi\mu_0)=2\sqrt3\).

\item Put \(R=\vartheta_1/\eta_D\). Poisson summation for the defining Gaussian
series and its index shift give
\[
 \begin{aligned}
 R(v,\tau+1)&=e^{\pi i/6}R(v,\tau),\\
 R(v/\tau,-1/\tau)&=-i e^{\pi iv^2/\tau}R(v,\tau),\\
 R(1-v,\tau)&=R(v,\tau),\qquad R(v+1,\tau)=-R(v,\tau).
 \end{aligned}
\]
Inserting the actual transformed arguments gives
\[
 \begin{aligned}
 R(\mu\circ g_1,\tau\circ g_1)
 &=e^{-\pi i/3}e^{\pi i(1-\mu)^2/\tau}R(\mu,\tau),\\
 R(\mu\circ g_2,\tau\circ g_2)
 &=i e^{\pi i\mu^2/\tau}R(\mu,\tau).
 \end{aligned}
\]
Multiplication of the logarithms by \(6i/\pi\) gives exactly
\(2-6(1-\mu)^2/\tau\) and \(-3-6\mu^2/\tau\), modulo constants in
\(12\mathbb Z\). These constants vanish: each defect is constant on the
connected cover, and evaluation at the corresponding fixed point makes
both its two sides zero. The fixed-point equations themselves give
\(\mu(z_1)=(2-\rho)/3\) and \(\mu(z_2)=(1-i)/2\).
The Jacobi product has cusp expansion
\(R(\mu,\tau)=2q^{1/12}\sin(\pi\mu)(1+O(q))\).
Equation \eqref{fam:eq:cd-beta} consequently gives
\(\beta^{\rm part}+\tau\to0\). The difference from the specified particular
solution is an invariant holomorphic function extending across all three
points of \(B\), hence a constant, and the cusp normalization makes it zero.

\item At \(\tau=i\), the half-period translation formula gives
\(\vartheta_1((1-i)/2,i)=e^{\pi/4}\vartheta_3(0,i)\).
The identities \(2\eta_D^3=\vartheta_2\vartheta_3\vartheta_4\),
\(\vartheta_2(0,i)=\vartheta_4(0,i)=2^{-1/4}\vartheta_3(0,i)\)
give \(R=e^{\pi/4}\sqrt2\), proving the displayed value of \(\beta\).
At \(\rho\), \(q\) is negative real, \(\mu=1/2-i/(2\sqrt3)\), and the
product for \(R\) has phase \(\pi/12\), with all its remaining factors
positive real. This gives \(\operatorname{Re}\beta=-1/2\).
\end{enumerate}
\end{proof}

\begin{proposition}\label{fam:src:C.1.2}\label{fam:st:C.1.2}
\StatementRecord{fam:src:C.1.2}{proved}{Historical events: Claude Fable 5.2 (Anthropic), effort high, 12 November 2026: targeted proof check. Claude Fable 5.2 (Anthropic), effort high, 16 November 2026: used and rechecked in this run. GPT-6 Astra (OpenAI), ChatGPT Pro mode, effort Pro, S6\_09\_PAPERS, 17 November 2026: targeted proof check. Also checked in this writing. Pre-continuation writing events: GPT-6 Astra (OpenAI), ChatGPT Pro mode, effort Pro, S6\_10\_PAPERS, 18 November 2026: targeted proof check. Continuation metadata events: GPT-6 Astra (OpenAI), ChatGPT Pro mode, effort Pro, S6\_10\_PAPERS continuation, 19 November 2026: used and rechecked in this run; status-vocabulary reconciliation; metadata reconciliation only; no new mathematical check. \emph{Pre-continuation recorded proof scope.} Read source proof; re-derived theta transformation, determinant/norm comparison, distribution constant or stated analytic inequality as applicable. Numerical replay scopes are recorded with their exact statement labels; the analytic argument is checked here.}{the hypotheses, notation and definitions stated here; \cref{fam:src:C.1.1}; \cref{fam:eq:cd-beta}}

The following decimal evaluations of the exact expressions in \ref{fam:src:C.1.1} agree
in two independent evaluations through the displayed twenty decimal places:
\[
\begin{array}{c|l}
\mu_0&0.5-0.36485183191106784792\,i\\
\mu(z_1)&0.5-0.28867513459481288225\,i\\
\mu(z_2)&0.5-0.5\,i\\
\beta^{\rm part}(z_1)&-0.5-1.27101206057537346287\,i\\
\beta^{\rm part}(z_2)&-0.21100425815613411418\,i
\end{array}
\]
These are evaluations of explicit convergent expressions, with working
precision ninety decimal digits. The formula \eqref{fam:eq:cd-beta}, together
with \(j(\tau)=1728t\) and \(\wp_\tau(\mu)=0\), determines the periods
throughout a fundamental domain.
\end{proposition}


\begin{proof}
\begin{enumerate}[label=\arabic*.,leftmargin=*]
\item The executable \texttt{cd\_numerics.py} evaluates \(\mu_0\) first by its
algebraic logarithm and then by solving
\(\pi^2\csc^2(\pi v)-\pi^2/3=0\). For \(\beta\), its first calculation uses
the Jacobi theta series and the eta product; its second uses
\[
 R(v,\tau)=2e^{\pi i\tau/6}\sin(\pi v)
 \prod_{n\geq1}(1-2q^n\cos(2\pi v)+q^{2n}).
\]
The product is truncated after 120 factors. At the two specified points
\(|q|\leq e^{-\pi\sqrt3}\) and \(|\operatorname{Im}v|\leq1/2\), so the
omitted logarithm is bounded by
\(\sum_{n>120}(2|q|^n\cosh\pi+|q|^{2n})/(1-2|q|^n\cosh\pi-|q|^{2n})\),
which is less than \(10^{-280}\). The measured discrepancies of the
ninety-digit calculations are below \(10^{-88}\). We do not promote floating-point rounding to an interval certificate.
A further sixty-digit interval computation, \texttt{cd\_special\_intervals.py},
evaluates the real positive product at $\rho$, including its forty-factor
tail, and the exact logarithmic expression at $i$. It encloses both
imaginary parts within intervals narrower than $3\cdot10^{-59}$ and
therefore certifies the displayed twenty decimal places. It also proves
$-0.2<\operatorname{Im}(\kappa_1^{\rm part}-\kappa_2^{\rm part})<-0.1$.
The exact formulas, product-tail bounds, independent algorithms, and
full outputs accompany the decimals. The special values of \(\mu\) are also checked directly in
the fixed-point equations, an exact algebraic calculation.
\end{enumerate}
\end{proof}

\subsection{Bounds for the sharp parameter disc}\label{fam:section:C.2}
\begin{proposition}\label{fam:src:C.2.1}\label{fam:st:C.2.1}
\StatementRecord{fam:src:C.2.1}{proved}{Historical events: Claude Fable 5.2 (Anthropic), effort high, 12 November 2026: targeted proof check. GPT-6 Astra (OpenAI), ChatGPT Pro mode, effort Pro, S6\_09\_PAPERS, 17 November 2026: targeted proof check. Also checked in this writing. Pre-continuation writing events: GPT-6 Astra (OpenAI), ChatGPT Pro mode, effort Pro, S6\_10\_PAPERS, 18 November 2026: targeted proof check. Continuation metadata events: GPT-6 Astra (OpenAI), ChatGPT Pro mode, effort Pro, S6\_10\_PAPERS continuation, 19 November 2026: used and rechecked in this run; status-vocabulary reconciliation; metadata reconciliation only; no new mathematical check. \emph{Pre-continuation recorded proof scope.} Read source proof; re-derived theta transformation, determinant/norm comparison, distribution constant or stated analytic inequality as applicable. Numerical replay scopes are recorded with their exact statement labels; the analytic argument is checked here.}{the hypotheses, notation and definitions stated here; \cref{fam:src:C.1.1}; \cref{fam:eq:cd-beta}; \cref{fam:eq:cd-norm}; \cref{fam:eq:cd-wp}; \cref{fam:src:C.2.3}}

For the normalization of \ref{fam:src:C.1.1}, the exact admissibility maximum is finite
and satisfies
\[
 -\frac3\pi\log6\ \leq M<2.
\]
More precisely the proof below gives \(M<1.635\). Consequently the sharp
radius \(R_w=e^{-\pi M}\) satisfies
\(e^{-2\pi}<R_w\leq216\), and for the reference linearization with
\(u=w^2\), \(e^{-4\pi}<R_u\leq46656\).
\end{proposition}


\begin{proof}
\begin{enumerate}[label=\arabic*.,leftmargin=*]
\item Equation \eqref{fam:eq:cd-beta} gives
\begin{equation}\label{fam:eq:cd-norm}
 D_{\rm part}=\frac6\pi\log\frac{\mathcal N(\mu,\tau)}{2\sqrt3},
 \qquad
 \mathcal N(v,\tau)=e^{-\pi(\operatorname{Im}v)^2/\operatorname{Im}\tau}
 |R(v,\tau)|.
\end{equation}
The theta transformation formulas prove that \(\mathcal N\) is unchanged
by integral period translations of \(v\), sign change of \(v\), and the
simultaneous modular change of \(v,\tau\). We may therefore use the
comparison isomorphism \(\mathbb C/(\mathbb Z+\tau\mathbb Z)\to
\mathbb C/(\mathbb Z+\tau'\mathbb Z)\), \([v]\mapsto[v/(c\tau+d)]\), to
place \(\tau\) in the standard modular domain and then choose
\(|\operatorname{Im}v|\leq T/2\), \(T=\operatorname{Im}\tau\geq\sqrt3/2\).
Both operations preserve the equation \(\wp(v)=0\) and \eqref{fam:eq:cd-norm}.

\item For \(\sqrt3/2\leq T\leq2\), the theta series and
\(\prod(1-r^n)\geq\exp(-r/(1-r)^2)\) imply
\[
 \mathcal N\leq
 2e^{\pi/6}\left(1+\frac{e^{-\pi\sqrt3/2}}
 {1-e^{-3\pi\sqrt3/2}}\right)
 \exp\frac{e^{-\pi\sqrt3}}{(1-e^{-\pi\sqrt3})^2}
 <3.615.
\]
Indeed the absolute value of the \(n\)-th sine term, including the
Gaussian factor, is at most
\(2\exp(-\pi T(n+1/2-|\operatorname{Im}v|/T)^2)\), which is at most
\(2e^{-\pi Tn^2}\); and \(n^2\geq1+3(n-1)\) for \(n\geq1\).

\item For \(T\geq2\), a zero of \(\wp\) in this period strip has
\(|\operatorname{Im}v|<1\). To prove this, assume its absolute value is at
least one and apply \eqref{fam:eq:cd-wp}. With \(r=e^{-4\pi}\), the absolute
value of the terms other than \(-1/3\) is at most
\[
 \frac1{\sinh^2\pi}+8\left\{
 \frac r{(1-r)^3}+\frac1{2(1-r)}
 \left(\frac{e^{-2\pi}}{(1-e^{-2\pi})^2}
 +\frac{e^{-6\pi}}{(1-e^{-6\pi})^2}\right)\right\}<0.016<\frac13.
\]
Here \(|\sin\pi v|\geq\sinh\pi\), and
\(|q|^n\cosh(2\pi n\operatorname{Im}v)\leq
(e^{-\pi Tn}+e^{-3\pi Tn})/2\). This contradicts \(\wp(v)=0\).
The product expression for \(R\) now gives the uniform bound
\[
 \mathcal N\leq2\cosh\pi\,e^{-\pi/3}
 \exp\left(\frac{2\cosh(2\pi)r}{1-r}+\frac{r^2}{1-r^2}\right)
 <8.151.
\]
Insertion in \eqref{fam:eq:cd-norm} proves the stated upper bound. The script
\texttt{cd\_bounds.py} independently evaluates these monotone bounds by
fifty-digit interval arithmetic: the larger resulting bound for \(D\) is
contained, after outward rounding for display, in
\[
 [1.63421874605891012363,\;1.63421874605891012365].
\]
The output file retains the full fifty-digit endpoints.
This encloses the coarse upper bound; the exact maximum is proved in \ref{fam:src:C.2.3}.

\item At the cusp \(D_{\rm part}=-T+o(1)\), so it tends to minus infinity;
continuity on the rest of the compact orbifold base gives attainment of
its maximum. At \(z_2\), \ref{fam:src:C.1.1} gives
\(D_{\rm part}=-(3/\pi)\log6\), proving the lower bound. Addition of \(c\)
adds \(\operatorname{Im}c\) to \(D\); a zero is forced by continuity
between a maximum and the cusp whenever \(\operatorname{Im}c\geq-M\).
For \(\operatorname{Im}c<-M\), the real determinant \(TD\) is everywhere
nonzero, proving the exact asserted form of the admissible disc.
\end{enumerate}
\end{proof}

\begin{proposition}\label{fam:src:C.2.2}\label{fam:st:C.2.2}
\StatementRecord{fam:src:C.2.2}{proved}{Historical events: Claude Fable 5.2 (Anthropic), effort high, 12 November 2026: targeted proof check. GPT-6 Astra (OpenAI), ChatGPT Pro mode, effort Pro, S6\_09\_PAPERS, 17 November 2026: targeted proof check. Also checked in this writing. Pre-continuation writing events: GPT-6 Astra (OpenAI), ChatGPT Pro mode, effort Pro, S6\_10\_PAPERS, 18 November 2026: targeted proof check. Continuation metadata events: GPT-6 Astra (OpenAI), ChatGPT Pro mode, effort Pro, S6\_10\_PAPERS continuation, 19 November 2026: used and rechecked in this run; status-vocabulary reconciliation; metadata reconciliation only; no new mathematical check. \emph{Pre-continuation recorded proof scope.} Read source proof; re-derived theta transformation, determinant/norm comparison, distribution constant or stated analytic inequality as applicable. Numerical replay scopes are recorded with their exact statement labels; the analytic argument is checked here.}{the hypotheses, notation and definitions stated here; the explicit constructions and calculations in the proof}

The point \(t=1\) is a strict local maximum of \(D_{\rm part}\).
A double-precision calculation on 5151 points of the half modular domain
\(0\leq\operatorname{Re}\tau\leq1/2\), \(|\tau|\geq1\),
\(\operatorname{Im}\tau\leq3\), has largest sampled value at \(\tau=i\).
This finite computation does not establish the global maximum.
\end{proposition}


\begin{proof}
\begin{enumerate}[label=\arabic*.,leftmargin=*]
\item We give a quantitative neighbourhood, which will also specify the local
region needed for a global interval cover. Put
\[
 h=v-(1-\tau)/2,\quad x=\tau-i,\quad
 G(h,\tau)=\wp_\tau((1-\tau)/2+h),\quad
 H(h,\tau)=\log\frac{\vartheta_3(h,\tau)}{\eta_D(\tau)}.
\]
The half-period theta identity gives exactly
$\log\mathcal N=\operatorname{Re}H-\pi(\operatorname{Im}h)^2/T$.
On the complex polydisc $|h|\leq1/5$, $|x|\leq1/10$, the theta series
has $|\vartheta_3-1|<1$, so the logarithm is holomorphic. There are no
poles of $G$ there: for every integer lattice point its imaginary
difference from $(1-\tau)/2+h$ has absolute value at least
$0.9/2-0.2>0$. Both functions are even in $h$.

\item The following uniform bounds use only positive Gaussian sums:
\[
 |G|<20,\qquad |H|<0.6.
\]
For explicit verification put
$A_k=\sum_{n\geq1}n^k e^{-0.9\pi n^2}\cosh(0.4\pi n)$,
$a_0=2A_0$, $a_1=4\pi A_1$, $a_2=8\pi^2A_2$, and $r=e^{-1.8\pi}$.
The bounds used are
\[
 \begin{aligned}
 |G|&\leq\frac{a_2}{1-a_0}+\frac{a_1^2}{(1-a_0)^2}
 +\frac{\pi^2}3\left(1+\frac{24r}{(1-r)^3}\right)<18.373,\\
 |H|&\leq-\log(1-a_0)+\frac{1.1\pi}{12}+
 \frac r{(1-r)^2}<0.547.
 \end{aligned}
\]
The derivative identity behind the first bound is
$G=-\partial_h^2\log\vartheta_3-\pi^2E_2/3$.
The script \texttt{cd\_local\_bound.py} evaluates these bounds by sixty-digit
interval arithmetic. It truncates the positive sums at $n=10$ and bounds
the remaining terms by geometric sums using
$n^2\geq121+23(n-11)$; the polynomial factors $n^k$ are summed explicitly.

\item At $(0,i)$ the theta heat equation gives
\[
 G=0,\quad G_{h^2}=-g_2(i)/4,\quad G_x=id,
 \qquad H=\log\sqrt2,\quad H_x=0,\quad H_{h^2}=-\pi/2.
\]
Subscripts $h^2$ denote power-series coefficients. The first two follow
also from $\wp''=6\wp^2-g_2/2$. Here $d$ is the real positive derivative coefficient of
$e_2(\tau)=(\pi^2/3)(\vartheta_2(0,\tau)^4-
\vartheta_4(0,\tau)^4)$ at $i$. Differentiating the convergent theta
series, with explicit Gaussian tails, gives $10.03<d<10.04$.
The value $H_x=0$ follows by differentiating the modular $S$-invariance
of $\vartheta_3(0,\tau)/\eta_D(\tau)$ at $i$. Differentiating
$\eta_D(-1/\tau)=\sqrt{-i\tau}\eta_D(\tau)$ gives $E_2(i)=3/\pi$,
which with $\wp((1-i)/2,i)=0$ gives $H_{h^2}=-\pi/2$.
The positive Eisenstein series for $g_2$, with an explicit tail in the
same script, gives $47.26<g_2(i)/4<47.30$.

\item Set $\epsilon=3/20000$, $r_h=3/100$, $s=\epsilon/(1/10)=0.0015$,
and $r_2=(r_h/(1/5))^2=0.0225$.
Cauchy's estimates for the even power series of $G$ give, on
$|h|=r_h$, $x=0$,
\[
 |G(h,i)|\geq
 \left(47.26-\frac{500r_2}{1-r_2}\right)r_h^2>0.03217.
\]
For $|x|\leq\epsilon$, the change of $G$ on this circle is at most
$20s/((1-s)(1-r_2))<0.03074$. Rouche's theorem therefore places both
zeros of $G(\cdot,\tau)$ in $|h|<r_h$. They represent the two zeros of
$\wp$ modulo its lattice, since an elliptic function with one double
pole has two zeros with multiplicity.

\item Write $G=e_2(\tau)+h^2Q(h,\tau)$. The same Cauchy bounds give
\[
 |Q+g_2(i)/4|\leq
 \frac{500(r_2+s/(1-s))}{1-r_2},
\quad
 |e_2(\tau)-idx|\leq
 \frac{2000|x|^2}{1-s}.
\]
At any zero they imply $|x|\leq7|h|^2$ and $|h|^2<1/20000$;
the interval output records bounds $6.124$ and $0.00004434$,
respectively. For $H$, its stated coefficients and Cauchy's estimates
now bound the remainder, together with the change of $1/T$ in the
Gaussian, by $5000|h|^4$. For example a direct upper bound for its
coefficient is
\[
 \frac{0.6(70^2+25\cdot70+25^2)}
 {(1-s)(1-25/20000)}+\frac{7\pi}{1-\epsilon}<5000.
\]
Consequently, for $0<|x|\leq3/20000$,
\[
 \log\mathcal N-\log\sqrt2
 \leq-\frac\pi2|h|^2+5000|h|^4<0.
\]
At $h=0$ the equation $G=0$ forces $x=0$ in this disc by the lower
bound for $e_2$. This proves strict local maximality, with an explicit
radius, and covers the square
$|\operatorname{Re}\tau|\leq10^{-4}$,
$|\operatorname{Im}\tau-1|\leq10^{-4}$.

\item The grid calculation uses thirty Fourier terms for $\wp$, complex
Newton iteration, and the theta product. Five sample values were
independently recomputed with eighty Fourier terms and the theta series
at ninety-digit precision; discrepancies are below $7\cdot10^{-16}$.
The source, precise grid rule, sampled maximum, and comparisons are in
\texttt{cd\_numerics.py} and its output.
\end{enumerate}
\end{proof}

\begin{theorem}[Certified bound]\label{fam:src:C.2.3}\label{fam:st:C.2.3}\label{ari:C.8.5}\label{push:src:E.4.3}
\StatementRecord{fam:src:C.2.3}{proved}{Historical events: Claude Fable 5.2 (Anthropic), effort high, 12 November 2026: targeted proof check. Claude Fable 5.2 (Anthropic), effort high, 15 November 2026: targeted proof check. Claude Fable 5.2 (Anthropic), effort high, 16 November 2026: carried as background without fault found; used and rechecked in this run. Claude Fable 5.2 (Anthropic), effort high, 17 November 2026: carried as background without fault found. Claude Opus 5.5 (Anthropic), effort max, sweep 1, scope s6, 18 November 2026: computation or formal certificate. GPT-6 Astra (OpenAI), ChatGPT Pro mode, effort Pro, S6\_09\_PAPERS, 17 November 2026: targeted proof check. Also checked in this writing. Pre-continuation writing events: GPT-6 Astra (OpenAI), ChatGPT Pro mode, effort Pro, S6\_10\_PAPERS, 18 November 2026: computation or formal certificate; targeted proof check. Continuation metadata events: GPT-6 Astra (OpenAI), ChatGPT Pro mode, effort Pro, S6\_10\_PAPERS continuation, 19 November 2026: used and rechecked in this run; status-vocabulary reconciliation; metadata reconciliation only; no new mathematical check. \emph{Pre-continuation recorded proof scope.} Read source proof; re-derived theta transformation, determinant/norm comparison, distribution constant or stated analytic inequality as applicable. Numerical replay scopes are recorded with their exact statement labels; the analytic argument is checked here. PASS: recomputed all 19701 norm-box contractions and strict norm bounds from interval formulas; independently reconstructed the exact 19911-leaf rational partition (36 local squares, 174 excluded below-circle boxes), area 23/40, no pending leaf. code/I/replay\_interval\_leaves.output.json and replay\_family/f\_interval\_verify.stdout. Analytic local/cusp proofs recorded separately.}{the hypotheses, notation and definitions stated here; \cref{fam:src:C.2.2}; \cref{fam:eq:cd-norm}; \cref{fam:src:C.2.1}; \cref{fam:eq:cd-wp}; \cref{fam:src:C.1.1}; \cref{fam:src:C.1.2}}

For the cusp normalization $\beta^{\rm part}+\tau\to0$,
\[
 M=-\frac3\pi\log6=-1.71100425815613411418\ldots,
 \qquad \operatorname*{arg\,max}_{B\setminus\{p_0\}}D_{\rm part}=\{1\}.
\]
The exact parameter discs are
\[
 0<|w|<216,
 \qquad 0<|u|<46656
\]
for the reference linearization $u=w^2$; an arbitrary reference constant
$\lambda_0$ multiplies the latter radius by $|\lambda_0|$.
The proof combines an analytic cusp bound, the quantitative local proof
\ref{fam:src:C.2.2}, and a completed outward-rounded interval cover of the remaining
compact domain.
\end{theorem}


\begin{proof}
\begin{enumerate}[label=\arabic*.,leftmargin=*]
\item By \eqref{fam:eq:cd-norm} it is enough to prove
$\mathcal N(v,\tau)\leq\sqrt2$ at every zero of $\wp_\tau$, with equality
only in the modular orbit of $i$. Modular transformations, period
translations and sign change preserve this norm. Complex conjugation
also preserves it, through
$(v,\tau)\mapsto(\bar v,-\bar\tau)$, as follows by conjugating the
normally convergent theta and eta series. It therefore suffices to work
on $0\leq\operatorname{Re}\tau\leq1/2$, $|\tau|\geq1$.

\item For $T\geq2$, the zero-strip argument in \ref{fam:src:C.2.1} improves to
$|\operatorname{Im}v|<b=43/100$. Replace $1/\sinh^2\pi$ in its displayed
error bound by $1/\sinh^2(\pi b)$; the resulting upper bound is less
than $0.316<1/3$, as certified independently in
\texttt{cd\_special\_intervals.py}. Retain the Gaussian in the product
estimate. The function
$\cosh(\pi y)\exp(-\pi y^2/T)$ increases on $0\leq y\leq b$, $T\geq2$:
its logarithmic derivative is $\pi\tanh(\pi y)-2\pi y/T\geq0$, since
$\tanh(\pi y)\geq y$ there. The latter follows from concavity and its
endpoint values. With $r=e^{-4\pi}$ we get
\[
 \mathcal N\leq
 2\cosh(\pi b)e^{-\pi/3-\pi b^2/2}
 \exp\left(\frac{2\cosh(2\pi b)r}{1-r}+
                  \frac{r^2}{1-r^2}\right)
 <1.082<\sqrt2.
\]
The bound decreases for $T\geq2$: the logarithmic derivative of the
exponential $e^{-\pi T/6-\pi b^2/T}$ is negative and the product
majorant decreases with $|q|$. This excludes the whole cusp region,
without a numerical truncation of the base.

\item We now specify the compact computation as a proof by enclosure.
Its rational starting rectangle is
$\mathcal R=[0,1/2]+i[17/20,2]$, which contains the compact half modular
domain. For every rational rectangle $Q$ retained in this domain, the
program selects a decimal complex centre $v_0$, a decimal nonzero
preconditioner $A$, and a closed complex rectangle $Z$ around $v_0$.
Floating-point Newton iteration only selects these numbers. The proof
uses the subsequent interval tests
\[
 v_0-A\wp(v_0,Q)+(1-A\partial_v\wp(Z,Q))(Z-v_0)
       \subset\operatorname{int}Z,
 \qquad
 \|1-A\partial_v\wp(Z,Q)\|_\infty<1.
\]
Here the norm of multiplication by a complex number is bounded by the
sum of the absolute real and imaginary parts. All displayed sets are
outward-rounded intervals. The mean-value integral along a segment in
$Z$ proves that, for every $\tau\in Q$, the map
$v\mapsto v-A\wp_\tau(v)$ maps $Z$ into itself and is a contraction.
Banach's theorem supplies a genuine zero in $Z$. Both zeros of $\wp$
are related by sign and lattice translations, so their norms agree;
this one enclosure therefore controls the actual period section.

\item For reproducibility the Fourier remainders are explicit. With
$q=e^{2\pi i\tau}$, $p=q e^{2\pi iv}$, $m=q e^{-2\pi iv}$, let
$r\geq|q|$ and $s\geq\max(|p|,|m|)$ be the interval upper endpoints.
The code requires $s<1$. For $N=24$ put
\[
 S_1(r)=r^{N+1}\left(\frac{N+1}{1-r}+\frac r{(1-r)^2}\right),
\]
\[
 S_2(r)=r^{N+1}\left(\frac{(N+1)^2}{1-r}+
 \frac{2(N+1)r}{(1-r)^2}+\frac{r(1+r)}{(1-r)^3}\right).
\]
The errors in $\wp$ and $\partial_v\wp$ are bounded respectively by
$8\pi^2(S_1(r)+S_1(s))/(1-r)$ and $16\pi^3S_2(s)/(1-r)$.
These follow by absolute summation of \eqref{fam:eq:cd-wp} and its derivative.
They are added as rectangles enclosing the corresponding complex discs.
For the norm, use $h=v-(1-\tau)/2$ and the even theta series
$\vartheta_3=1+2\sum_{n\geq1}e^{\pi i n^2\tau}\cos(2\pi nh)$.
With $T_*\leq\operatorname{Im}\tau$ and $H_*\geq|\operatorname{Im}h|$,
its omitted terms after $n=10$ have absolute value at most
\[
 \frac{2e^{-121\pi T_*+22\pi H_*}}
 {1-e^{-23\pi T_*+2\pi H_*}},
\]
and the denominator is tested positive. The logarithm of the eta
product tail after $n=17$ is bounded by
$r^{18}/((1-r)(1-r^{18}))$. The resulting enclosure of
$e^{-2\pi(\operatorname{Im}h)^2/T}|\vartheta_3/\eta_D|^2$
is required to have upper endpoint strictly below two.

\item The completed computation \texttt{f\_interval\_parallel.py}, using
\texttt{f\_interval\_cert.py}, runs these tests at thirty decimal digits
and records every rational box. It accepted 19701 norm boxes. A further
36 boxes lie inside the local square of \ref{fam:src:C.2.2} and are covered by that
analytic proof. Another 174 boxes satisfy
$x_{\rm upper}^2+y_{\rm upper}^2<1$ exactly and lie outside the modular
domain. There are zero pending or undecided boxes. The independent
program \texttt{f\_interval\_verify.py} reconstructs the entire binary
partition from its leaves, checks all reported strict interval bounds,
and verifies the exclusions by rational arithmetic. It consumes all
19911 leaves, with exact total area $23/40$ and maximum depth 38.
The largest reported squared-norm upper bound is the exact decimal
rational
\[
 \frac{9999999982782673130994674616307871}
 {5000000000000000000000000000000000}<2.
\]
The JSON certificate, complete outputs, and independently runnable
source programs are included. Thus the compact domain is fully covered;
the exploratory grid is not used as a substitute for any enclosure.
The local proof gives equality only at $\tau=i$, and the cusp and other
compact boxes give strict inequality. This proves the global bound and
its equality set.

\item Finally \ref{fam:src:C.1.1} evaluates $D$ at $i$ exactly. Since $j(\tau)=1728t$, its
modular orbit maps precisely to $t=1$. Exponentiation gives
$e^{-\pi M}=e^{3\log6}=216$ and $e^{-2\pi M}=6^6=46656$.
The two independent theta calculations and the sixty-digit logarithmic
interval computation in \ref{fam:src:C.1.2} verify the numerical value of $M$ as well
as the exact evaluation; the integer radii are checked against interval
exponentiation in \texttt{cd\_special\_intervals.py}.
\end{enumerate}
\end{proof}


\section{Special values, local expansions and torsion}
\subsection{The normalized Green function and special arguments}
We write $\rho=(1+i\sqrt3)/2$, $T=\operatorname{Im}\tau$, and
\[
 U_\tau(z)=\log N(z,\tau)
 =\log\left|\frac{\vartheta _1(z,\tau)}{\eta(\tau)}\right|
       -\frac{\pi(\operatorname{Im}z)^2}{T}.
\]
Our theta variable has periods $1,\tau$; thus the standard theta series
is evaluated at $\pi z$ when its first argument has period $\pi$.
For a general lattice $\Lambda=\omega(\mathbb Z+\tau\mathbb Z)$ we set
$U_\Lambda(z)=U_\tau(z/\omega)$. The transformation laws of the theta
function show that this definition is independent of the oriented basis.

\begin{proposition}[distribution under isogenies]\label{ari:C.1.1}
\StatementRecord{ari:C.1.1}{proved}{Historical events: Claude Fable 5.2 (Anthropic), effort high, 16 November 2026: used and rechecked in this run. GPT-6 Astra (OpenAI), ChatGPT Pro mode, effort Pro, S6\_09\_PAPERS, 17 November 2026: targeted proof check. Also checked in this writing. Pre-continuation writing events: GPT-6 Astra (OpenAI), ChatGPT Pro mode, effort Pro, S6\_10\_PAPERS, 18 November 2026: targeted proof check. Continuation metadata events: GPT-6 Astra (OpenAI), ChatGPT Pro mode, effort Pro, S6\_10\_PAPERS continuation, 19 November 2026: used and rechecked in this run; status-vocabulary reconciliation; metadata reconciliation only; no new mathematical check. \emph{Pre-continuation recorded proof scope.} Read source proof; re-derived theta transformation, determinant/norm comparison, distribution constant or stated analytic inequality as applicable. Numerical replay scopes are recorded with their exact statement labels; the analytic argument is checked here.}{the hypotheses, notation and definitions stated here; the explicit constructions and calculations in the proof}

Let $\Lambda,\Lambda'\subset\mathbb C$ be lattices and let
$\alpha\ne0$ satisfy $\alpha\Lambda\subset\Lambda'$. Write
$\phi:\mathbb C/\Lambda\to\mathbb C/\Lambda'$ for $z\mapsto\alpha z$,
$K=\alpha^{-1}\Lambda'/\Lambda$, and
$d=[\Lambda':\alpha\Lambda]$. Then
\begin{equation}\label{ari:eq:cs-isogeny}
 U_{\Lambda'}(\alpha z)=\sum_{v\in K}U_\Lambda(z+v).
\end{equation}
The equality holds away from the zeros and as an equality of locally
integrable functions at them. If $\Delta_\Lambda=g_2(\Lambda)^3-
27g_3(\Lambda)^2$, then
\[
 \prod_{0\ne v\in K}N_\Lambda(v)
 =|\alpha|\left|\frac{\Delta_{\Lambda'}}{\Delta_\Lambda}\right|^{1/12}.
\]
In particular, for every nonzero endomorphism of one elliptic curve,
$\prod_{v\in\ker\alpha,\ v\ne0}N(v)=|\alpha|$.
\end{proposition}


\begin{proof}
\begin{enumerate}[label=\arabic*.,leftmargin=*]
\item The Jacobi product, with $q=e^{2\pi i\tau}$ and
$X=e^{2\pi iz}$, gives
\[
 \vartheta _1(z,\tau)/\eta(\tau)
 =i q^{1/12}X^{-1/2}(1-X)
     \prod_{m\ge1}(1-q^mX)(1-q^mX^{-1}),
\]
where a harmless fixed unit changes with the theta convention. For
$0<y<T$, integration of the logarithm over $z=x+iy$, $0\le x\le1$,
gives $-\pi T/6+\pi y-\pi y^2/T$. Each product factor has mean logarithm
zero by the absolutely convergent logarithmic series; the endpoint
singularities are integrable. Integrating in $y$ proves
$\int U_\Lambda\,d\mu_\Lambda=0$ for Haar probability measure.
Moreover
\[
 \Delta_{\rm Eucl}U_\Lambda
   =2\pi\delta_0-\frac{2\pi}{\operatorname{area}\Lambda},\qquad
 U_\Lambda(z)=\log|z|+\frac1{12}\log|\Delta_\Lambda|+o(1).
\]
The second formula follows from
$\vartheta _1'(0,\tau)=2\pi\eta(\tau)^3$ and
$\Delta_{\mathbb Z+\tau\mathbb Z}=(2\pi)^{12}\eta(\tau)^{24}$,
then scaling by $\omega$. The two sides of
\eqref{ari:eq:cs-isogeny} have the same logarithmic masses at $K$ and the
same constant Laplacian, since
$d=|\alpha|^2\operatorname{area}\Lambda/
\operatorname{area}\Lambda'$. Their difference is a smooth harmonic
function on the compact torus, hence constant. Both means vanish:
translations preserve Haar measure and an isogeny sends Haar
probability measure to Haar probability measure. The constant is zero.
Taking $z\to0$ and cancelling the logarithmic singularity gives the
product. This also checks the precise normalization of the special-value calculation in Proposition~\ref{push:src:D.4.5};
its kernel, Laplacian and mean-zero argument are all valid.
\end{enumerate}
\end{proof}
\begin{proposition}[the phases and the two periods]\label{ari:C.1.2}
\StatementRecord{ari:C.1.2}{proved}{Historical events: Claude Fable 5.2 (Anthropic), effort high, 16 November 2026: used and rechecked in this run. GPT-6 Astra (OpenAI), ChatGPT Pro mode, effort Pro, S6\_09\_PAPERS, 17 November 2026: targeted proof check. Also checked in this writing. Pre-continuation writing events: GPT-6 Astra (OpenAI), ChatGPT Pro mode, effort Pro, S6\_10\_PAPERS, 18 November 2026: targeted proof check. Continuation metadata events: GPT-6 Astra (OpenAI), ChatGPT Pro mode, effort Pro, S6\_10\_PAPERS continuation, 19 November 2026: used and rechecked in this run; status-vocabulary reconciliation; metadata reconciliation only; no new mathematical check. \emph{Pre-continuation recorded proof scope.} Re-derived the E6 tail bound and continuation of the zero through the stated component, the endpoint phases and Gaussian factors; verified the explicit isogeny and kernel index N cubed.}{the hypotheses, notation and definitions stated here; \cref{ari:C.1.1}}

With the branch continued from the specified cusp, the exact values are
\begin{align*}
 \frac{\vartheta _1((2-\rho)/3,\rho)}{\eta(\rho)}
   &=e^{\pi i/12}\,3^{1/4}e^{\pi/(6\sqrt3)},\\
 \frac{\vartheta _1((1-i)/2,i)}{\eta(i)}
   &=\sqrt2\,e^{\pi/4},\\
 \beta^{\rm part}(z_1)
   &=-\frac12+i\left(\frac1{\sqrt3}-\frac3{2\pi}\log48\right),\\
 \beta^{\rm part}(z_2)
   &=i\left(\frac32-\frac3\pi\log6\right).
\end{align*}
Consequently $N(z_1)=3^{1/4}$, $N(z_2)=\sqrt2$,
$D_{\rm part}(z_1)=-3\log48/(2\pi)$, and
$D_{\rm part}(z_2)=-3\log6/\pi$.
\end{proposition}


\begin{proof}
\begin{enumerate}[label=\arabic*.,leftmargin=*]
\item The special-value calculation in Proposition~\ref{push:src:D.4.5}
uses multiplication by $1+\rho$ has matrix
$\left(\begin{smallmatrix}1&-1\\1&2\end{smallmatrix}\right)$,
its kernel is $\{0,\mu,-\mu\}$, and $(1+\rho)\mu=1$ for
$\mu=(2-\rho)/3$. Statement \ref{ari:C.1.1} gives $N(\mu)^2=\sqrt3$.
For the requested parallel calculation at $i$, multiplication by $1+i$
has kernel $\{0,(1-i)/2\}$ and absolute value $\sqrt2$; the same
identity gives $N((1-i)/2,i)=\sqrt2$ without a separate theta-constant
evaluation.

\item At $\rho$, $\mu=1/2-i/(2\sqrt3)$ and $q=-e^{-\pi\sqrt3}$. In
\[
 R(z,\tau)=2e^{\pi i\tau/6}\sin(\pi z)
 \prod_{m\ge1}\bigl(1-2q^m\cos(2\pi z)+q^{2m}\bigr),
\]
$\sin(\pi\mu)$ is positive real. The factors are
$(1+q^m e^{\pi/\sqrt3})(1+q^m e^{-\pi/\sqrt3})$, both positive:
when $m$ is odd the larger subtracted number is at most
$e^{-2\pi/\sqrt3}<1$. Thus the phase is exactly $\pi/12$ modulo
$2\pi$. There is no winding ambiguity in the continued logarithm.
Indeed continue from the cusp along $\tau=1/2+iT$,
$T\ge\sqrt3/2$. Real conjugation of the rhombic lattice keeps the
chosen zero on $\mu=1/2-iy$, $0<y<T$. It cannot leave that component:
a collision of the two zeros would require $g_3=0$, which does not
occur on this ray. Indeed, writing $q=-r$, $0<r<1/200$,
the tail after the first term of the $E_6$ series has absolute value
at most $504\cdot32r^2/((1-r^2)(1-8r))<1/2$, while its
first term is positive. Thus $E_6>1/2$. Hitting $y=0$ would again give a half-period zero, and
hitting $y=T$ would give another half-period zero. Either event
would force $g_3=0$. Each displayed product factor remains
positive because $e^{-2\pi T}e^{2\pi y}<1$. The cusp phase tends to
$\pi/12$, so the continued argument has that value. On the imaginary
ray $\tau=iT$, $T\ge1$, the same branch is
$\mu=1/2-iy$ with $0<y\le T/2$; every product factor is positive,
so its argument at $i$ is zero. These checks supply the branch step
used in R3. The Gaussian factors are respectively
$e^{\pi/(6\sqrt3)}$ and $e^{\pi/4}$. Substitution in the definition of
$\beta^{\rm part}$ proves all four identities.
\end{enumerate}
\end{proof}

\begin{proposition}[numerical cross-check]\label{ari:C.1.3}
\StatementRecord{ari:C.1.3}{proved}{Historical events: GPT-6 Astra (OpenAI), ChatGPT Pro mode, effort Pro, S6\_09\_PAPERS, 17 November 2026: targeted proof check. Also checked in this writing. Pre-continuation writing events: GPT-6 Astra (OpenAI), ChatGPT Pro mode, effort Pro, S6\_10\_PAPERS, 18 November 2026: computation or formal certificate; targeted proof check. Continuation metadata events: GPT-6 Astra (OpenAI), ChatGPT Pro mode, effort Pro, S6\_10\_PAPERS continuation, 19 November 2026: used and rechecked in this run; status-vocabulary reconciliation; metadata reconciliation only; no new mathematical check. \emph{Pre-continuation recorded proof scope.} Read the complete numerical crosscheck and its normalizations. Retained diagnostic numerical status; no directed-rounding certificate is claimed. Fully read and freshly replayed unchanged Cspecial\_torsion\_numerics.py; exit zero. Output and SHA256 are retained in code/I/replay\_special/replay\_manifest.json. Numerical outputs retain numerical scope.}{the hypotheses, notation and definitions stated here; \cref{ari:C.1.2}}

The independent analytic computation in
\path{code/Cspecial_torsion_numerics.py} gives, to fifty decimals,
\begin{align*}
 N(z_1)&=1.31607401295249246081921890179699905516006859020582\ldots,\\
 N(z_2)&=1.41421356237309504880168872420969807856967187537695\ldots,\\
 \operatorname{Im}\beta^{\rm part}(z_1)
  &=-1.27101206057537346286579673538412682778809507094949\ldots,\\
 \operatorname{Im}\beta^{\rm part}(z_2)
  &=-0.21100425815613411418178314271018934802718590587233\ldots.
\end{align*}
The theta/eta evaluations and exact expressions differ by less than
$10^{-100}$ at working precision $110$ decimal digits.
\end{proposition}


\begin{proof}
\begin{enumerate}[label=\arabic*.,leftmargin=*]
\item The program directly evaluates the theta series and the eta product,
then compares the displayed closed forms. The eta product has $160$
factors. Since $|q|\le e^{-\pi\sqrt3}<1/200$, the omitted logarithm is
at most $\sum_{m>160}|q|^m/(1-|q|^m)<10^{-368}$.
The output records the measured floating-point discrepancies. The
exact values are established by \ref{ari:C.1.2}; this computation is their
independent numerical check, not an additional rounding certificate.
\end{enumerate}
\end{proof}


\subsection{Division identities and the entire torsion locus}
We use throughout this subsection the rational model
\[
 E_t:\quad y^2=x^3+A(t)x+B(t),\qquad
 A=-\tfrac34t^3(t-1),\quad B=\tfrac14t^4(t-1)^2,
 \quad P=(0,t^2(t-1)/2).
\]
Its discriminant is $\Delta(t)=27t^8(t-1)^3$.
The change of complex Weierstrass model does not change $N$.

\begin{cited}[Siegel functions and local heights]\label{ari:C.4.1}
\StatementRecord{ari:C.4.1}{proved}{Historical events: GPT-6 Astra (OpenAI), ChatGPT Pro mode, effort Pro, S6\_09\_PAPERS, 17 November 2026: targeted proof check. This writing records metadata reconciliation only; no new mathematical check. Continuation metadata events: GPT-6 Astra (OpenAI), ChatGPT Pro mode, effort Pro, S6\_10\_PAPERS continuation, 19 November 2026: used and rechecked in this run; status-vocabulary reconciliation; metadata reconciliation only; no new mathematical check. \emph{Pre-continuation recorded proof scope.} Current-writing verification remains in progress. The earlier checks named above retain their recorded scope.}{the hypotheses, notation and definitions stated here; \citet[Remark~2.2 and Proposition~2.3]{Shin2015}; \CompendiumRef{I}{app:rederive:ari-C.4.1}}

For $a=(a_1,a_2)\in\mathbb R^2\setminus\mathbb Z^2$ and
$v=a_1\tau+a_2$, define
\[
 g_a(\tau)=-e^{\pi ia_2(a_1-1)}q^{B_2(a_1)/2}
 (1-e^{2\pi iv})\prod_{r\ge1}
 (1-q^re^{2\pi iv})(1-q^re^{-2\pi iv}),
 \quad B_2(x)=x^2-x+\tfrac16.
\]
Then
\[
 g_a(\tau)=i e^{\pi ia_1v}
       \frac{\vartheta_1(v,\tau)}{\eta(\tau)},\qquad
 N(v,\tau)=|g_a(\tau)|.
\]
When $a\in n^{-1}\mathbb Z^2\setminus\mathbb Z^2$, $n\ge2$,
$g_a^{12n}$ is a modular unit for $\Gamma(n)$ with Fourier coefficients
in $\mathbb Q(\zeta_n)$ in the variable $q^{1/n}$. On every complex
elliptic curve the model-independent Néron local function for $(O)$ is
\[
 \lambda_\infty(Q)=-\log N(Q).
\]
Thus $D_{\rm part}=-(6/\pi)(\lambda_\infty(P)+\log(2\sqrt3))$.
\end{cited}
The expanded source derivation is retained in the proof compendium, Paper~I, label \nolinkurl{ari:C.4.1}.




The Siegel-function identities are those of \citet[Remark~2.2 and Proposition~2.3]{Shin2015}, with the local-height normalization proved explicitly. A complete re-derivation is given in \CompendiumRef{I}{app:rederive:ari-C.4.1}.


\begin{proposition}[the full division identity]\label{ari:C.4.2}
\StatementRecord{ari:C.4.2}{proved}{Historical events: Claude Fable 5.2 (Anthropic), effort high, 15 November 2026: carried as background without fault found. Claude Fable 5.2 (Anthropic), effort high, 16 November 2026: carried as background without fault found; used and rechecked in this run. GPT-6 Astra (OpenAI), ChatGPT Pro mode, effort Pro, S6\_09\_PAPERS, 17 November 2026: targeted proof check. Also checked in this writing. Pre-continuation writing events: GPT-6 Astra (OpenAI), ChatGPT Pro mode, effort Pro, S6\_10\_PAPERS, 18 November 2026: targeted proof check. Continuation metadata events: GPT-6 Astra (OpenAI), ChatGPT Pro mode, effort Pro, S6\_10\_PAPERS continuation, 19 November 2026: used and rechecked in this run; status-vocabulary reconciliation; metadata reconciliation only; no new mathematical check. \emph{Pre-continuation recorded proof scope.} Read source proof; re-derived theta transformation, determinant/norm comparison, distribution constant or stated analytic inequality as applicable. Numerical replay scopes are recorded with their exact statement labels; the analytic argument is checked here.}{the hypotheses, notation and definitions stated here; \cref{ari:C.1.1}}

Let $\psi_m$ be the standard division polynomial of a smooth
Weierstrass model, normalized by $\psi_1=1$, $\psi_2=2y$.
For $m\ge2$ and $Q,mQ\ne O$,
\begin{equation}\label{ari:eq:cs-division}
 N(Q)^{m^2}
       =N(mQ)\frac{|\Delta_E|^{(m^2-1)/12}}{|\psi_m(Q)|}.
\end{equation}
For the displayed family this specializes to
\[
 N(P)^4=3^{3/4}|t-1|^{-1/4}N(2P),\qquad
 N(P)^9=16|t|^{-2/3}N(3P).
\]
If $P_t$ has exact order $n$ and $m\ge2$ satisfies
$m\equiv1$ or $-1\pmod n$, then
\begin{equation}\label{ari:eq:cs-torsionvalue}
 N(P_t)^{12}
       =\frac{|\Delta(t)|}{|\psi_m(P_t)|^{12/(m^2-1)}}.
\end{equation}
The positive value on the right is independent of $m$.
\end{proposition}


\begin{proof}
\begin{enumerate}[label=\arabic*.,leftmargin=*]
\item Use the complex uniformization $x=\wp_\Lambda(z)$,
$y=\wp'_\Lambda(z)/2$. The rational function $\psi_m$ has divisor
$\sum_{0\ne Q\in E[m]}(Q)-(m^2-1)(O)$ and leading Laurent term
$(-1)^{m-1}m z^{1-m^2}$. One way to verify both normalizations without
a model convention is to form
$(-1)^{m-1}\sigma_\Lambda(mz)/\sigma_\Lambda(z)^{m^2}$.
Its lattice factors cancel, its zeros are exactly the nonzero
$m$-division points, and its expansion at zero is the displayed
one. The addition formula for $\sigma$ gives the standard
division-polynomial recurrence with initial values $1,2y$, so it is
the stated rational function.

\item The function
\[
 m^2U_\Lambda(z)-U_\Lambda(mz)
                    +\log|\psi_m(Q(z))|
\]
is harmonic after its logarithmic singularities cancel. Its
Laplacian has zero constant term as well, so it is constant.
The expansion in \ref{ari:C.1.1} and the leading term of $\psi_m$ make that
constant $(m^2-1)\log|\Delta_E|/12$. This proves
\eqref{ari:eq:cs-division}. Direct substitution of
\[
 \psi_2(P)=t^2(t-1),\qquad
 \psi_3(P)=-9t^6(t-1)^2/16
\]
gives the two special identities. For $m\equiv\pm1\pmod n$,
$mP=\pm P$ and $N$ is even. Cancelling $N(P)$ in
\eqref{ari:eq:cs-division} gives \eqref{ari:eq:cs-torsionvalue}.
Every right side is the same positive number $N(P)^{12}$, which
proves independence of $m$, not merely numerical agreement.
\end{enumerate}
\end{proof}

\begin{proposition}[algebraicity and explicit fields for every order]\label{ari:C.4.3}
\StatementRecord{ari:C.4.3}{proved}{Historical events: Claude Fable 5.2 (Anthropic), effort high, 16 November 2026: used and rechecked in this run. GPT-6 Astra (OpenAI), ChatGPT Pro mode, effort Pro, S6\_09\_PAPERS, 17 November 2026: targeted proof check. Also checked in this writing. Pre-continuation writing events: GPT-6 Astra (OpenAI), ChatGPT Pro mode, effort Pro, S6\_10\_PAPERS, 18 November 2026: computation or formal certificate; targeted proof check. Continuation metadata events: GPT-6 Astra (OpenAI), ChatGPT Pro mode, effort Pro, S6\_10\_PAPERS continuation, 19 November 2026: used and rechecked in this run; status-vocabulary reconciliation; metadata reconciliation only; no new mathematical check. \emph{Pre-continuation recorded proof scope.} Re-derived the nonvanishing division-polynomial specialization from the nonconstant torsion argument, the radical/resultant relations and formal-group multiplicities. Fully read and freshly replayed unchanged Cspecial\_division.py; exit zero. Output and SHA256 are retained in code/I/replay\_special/replay\_manifest.json. Numerical outputs retain numerical scope.}{the hypotheses, notation and definitions stated here; \cref{ari:eq:cs-torsionvalue}}

Let $t_0\in\mathbb C\setminus\{0,1\}$ and suppose $P_{t_0}$ has
exact order $n\ge2$. Then $t_0$ and $N(P_{t_0})$ are algebraic.
Put $K_0=\mathbb Q(t_0,\bar t_0)$,
\[
 W=\Delta(t_0)\Delta(\bar t_0)>0,\qquad
 V=\psi_{n+1}(P_{t_0})\psi_{n+1}(P_{\bar t_0})>0.
\]
With positive real radicals,
\[
 N(P_{t_0})^{12n}=W^{n/2}V^{-6/(n+2)}
 \ \in\ K_0\bigl(W^{1/2},V^{1/(n+2)}\bigr).
\]
This supplies an explicit field, without a claim that it is minimal.
More generally, if $m\equiv\pm1\pmod n$, then
\[
 \lambda_\infty(P_{t_0})
 =\frac{\log|\psi_m(P_{t_0})|}{m^2-1}
       -\frac1{12}\log|\Delta(t_0)|.
\]
These formulas determine the height and $N^{12}$ on the entire
torsion locus.
\end{proposition}


\begin{proof}
\begin{enumerate}[label=\arabic*.,leftmargin=*]
\item The division polynomial $\psi_n(P_t)$ is a nonzero rational
polynomial in $t$. To check nonvanishing for every $n$, a hypothetical
torsion identity would express the elliptic logarithm as
$a\tau+b$ with constant rational $a,b$. Its bounded cusp limit
forces $a=0$, but that limit has nonzero imaginary part, a
contradiction. Thus $t_0$ is a zero of a nonzero polynomial over
$\mathbb Q$. Formula \eqref{ari:eq:cs-torsionvalue} shows algebraicity
of $N$ using only positive rational powers of nonzero algebraic
absolute values. Taking $m=n+1$, raising to the $n$th power, and
using $|\Delta|=W^{1/2}$, $|\psi_{n+1}|=V^{1/2}$ gives the exact
field statement. The local-height formula is its logarithm.

\item Here is a polynomial relation involving only $t$ and $r=N^{12}$,
if one prefers that form to positive radicals. Remove the factors
at $0,1$ from $\psi_n(P_t)$ and divide successively by the
corresponding polynomials for proper divisors of $n$; let $F_n$ be
the resulting primitive polynomial. At a point of exact order
$d\mid n$, the formal group gives
$z([n/d]Q)=(n/d)z(Q)+O(z(Q)^2)$, so the vanishing multiplicities for
$\psi_d$ and $\psi_n$ agree. Hence this cancellation leaves exactly
the exact-order-$n$ divisor, with its multiplicities. Set
$m=n+1$ and $e=m^2-1$. Clearing rational denominators in
\[
 \operatorname{Res}_S\left(
 F_n(S),\
 r^{2e}\{\psi_m(P_T)\psi_m(P_S)\}^{12}
              -\{\Delta(T)\Delta(S)\}^{e}\right)
\]
gives a polynomial in $\mathbb Q[T,r]$ that vanishes at
$(t_0,N(P_{t_0})^{12})$, since $S=\bar t_0$ is a common root.
It is nonzero: at an exact-order-$n$ root no $\psi_m$ vanishes,
so the leading coefficient in $r$ of the resultant is nonzero.
This is an explicit recurrence-and-resultant prescription for
every $n$, including nonreal torsion parameters.
\end{enumerate}
\end{proof}

\begin{proposition}[the five values and all orders through eight]\label{ari:C.4.4}
\StatementRecord{ari:C.4.4}{proved}{Historical events: Claude Fable 5.2 (Anthropic), effort high, 16 November 2026: used and rechecked in this run. GPT-6 Astra (OpenAI), ChatGPT Pro mode, effort Pro, S6\_09\_PAPERS, 17 November 2026: targeted proof check. Also checked in this writing. Pre-continuation writing events: GPT-6 Astra (OpenAI), ChatGPT Pro mode, effort Pro, S6\_10\_PAPERS, 18 November 2026: computation or formal certificate; targeted proof check. Continuation metadata events: GPT-6 Astra (OpenAI), ChatGPT Pro mode, effort Pro, S6\_10\_PAPERS continuation, 19 November 2026: used and rechecked in this run; status-vocabulary reconciliation; metadata reconciliation only; no new mathematical check. \emph{Pre-continuation recorded proof scope.} Re-derived the division recurrences through order eight, squarefree elimination and the five closed torsion values; exact companion replay follows in a separate event. Fully read and freshly replayed unchanged Cspecial\_division.py; exit zero. Output and SHA256 are retained in code/I/replay\_special/replay\_manifest.json. Numerical outputs retain numerical scope.}{the hypotheses, notation and definitions stated here; \cref{ari:eq:cs-torsionvalue}; \cref{ari:C.1.2}; \cref{ari:C.4.3}}

The exact special values are
\[
 \begin{array}{c|c|c}
 t&\operatorname{ord}(P)&N(P)^{12}\\ \hline
 0&3&27\\
 1&2&64\\
 32/5&4&10\\
 64(47-27\sqrt3)/11&6&17+10\sqrt3\\
 64(47+27\sqrt3)/11&6&10\sqrt3-17 .
 \end{array}
\]
At $t=0,1$ the orders refer to the smooth models after the specified
base changes. On the smooth base, $F_2=F_3=1$ and
\[
 \begin{aligned}
 F_4&=5t-32,\\
 F_5&=89t^2+4736t-4096,\\
 F_6&=11t^2-6016t+8192,\\
 F_7&=881t^4+4745344t^3-12079104t^2\\
    &\hspace{13mm}+24641536t-16777216,\\
 F_8&=49t^4-2614336t^3-7864320t^2\\
    &\hspace{13mm}+19398656t-8388608.
 \end{aligned}
\]
All these polynomials are squarefree. At each root of exact order
$n=4,5,6,7,8$, the positive number $r=N^{12}$ is respectively
determined by
\[
 \begin{array}{c|l}
 n&\text{equation for }r\\ \hline
 4&r=64/|t|\\
 5&r^5=432^5/\bigl((1-t)(5t-32)^4\bigr)\\
 6&r^2=1728^2/F_5(t)\\
 7&r^{35}=2^{192}3^{81}/
       \bigl(|t|^8|t-1|^3|F_6(t)|^{12}\bigr)\\
 8&r^4=1728^4/|F_7(t)|.
 \end{array}
\]
The order-five roots are real and below one; the order-six roots
are real and above one. The two complex order-seven roots are
included in the formula.
\end{proposition}


\begin{proof}
\begin{enumerate}[label=\arabic*.,leftmargin=*]
\item The division-polynomial recurrence, evaluated at $x=0$,
$y=t^2(t-1)/2$, gives
\[
 \begin{aligned}
 \psi_4(P)&=-t^{10}(t-1)^4F_4/32,\\
 \psi_5(P)&=t^{16}(t-1)^6F_5/4096,\\
 \psi_6(P)&=9t^{24}(t-1)^9F_6/65536,\\
 \psi_7(P)&=-t^{32}(t-1)^{12}F_7/16777216,\\
 \psi_8(P)&=-t^{42}(t-1)^{16}F_4F_8/268435456.
 \end{aligned}
\]
The recursion and all coefficients are checked over $\mathbb Q[t]$
by \path{Cspecial_division.py}. The same program computes
$\gcd(F_n,F_n')=1$ and excludes common factors for proper orders.
Thus no point or multiplicity is omitted. Applying
\eqref{ari:eq:cs-torsionvalue} with $m=n-1$ gives every displayed
formula. The signs for $F_5,F_6$ follow directly from their quadratic
formula; in particular $F_5$ is positive at the two roots of $F_6$.
At order six, substituting
$t=64(47\mp27\sqrt3)/11$ gives the exact identities
\[
 F_5(t)(17\pm10\sqrt3)^2=1728^2.
\]
The positive roots are $17+10\sqrt3$ and $10\sqrt3-17$; their
product is $11$. The substitution is checked in
\path{Cspecial_relations.py}. The order-four value follows from
$64/(32/5)=10$, and the two central values are \ref{ari:C.1.2}.

\item For clarity, the recurrence used by the exact program is
\[
 \begin{aligned}
 \psi_{2m+1}&=\psi_{m+2}\psi_m^3-\psi_{m-1}\psi_{m+1}^3,\\
 \psi_{2m}&=\frac{\psi_m}{\psi_2}
   (\psi_{m+2}\psi_{m-1}^2-\psi_{m-2}\psi_{m+1}^2).
 \end{aligned}
\]
Its initial values at $P$ are $\psi_0=0,\psi_1=1,\psi_2=2y$,
$\psi_3=-A^2$, $\psi_4=4y(-8B^2-A^3)$. No approximate algebra
enters this calculation. The same code records $\psi_9$, and the
separate file \texttt{Cspecial\_psi16.out} records the factored
$\psi_{16}$ for the duplication ladder. A rational polynomial
relation for the nonreal order-seven values is explicitly computed
by \path{Cspecial_relations.py}, by the resultant prescription
of \ref{ari:C.4.3}.
\end{enumerate}
\end{proof}

\begin{proposition}[complete numerical verification through order eight]\label{ari:C.4.5}
\StatementRecord{ari:C.4.5}{proved}{Historical events: Claude Fable 5.2 (Anthropic), effort high, 15 November 2026: carried as background without fault found. GPT-6 Astra (OpenAI), ChatGPT Pro mode, effort Pro, S6\_09\_PAPERS, 17 November 2026: targeted proof check. Also checked in this writing. Pre-continuation writing events: GPT-6 Astra (OpenAI), ChatGPT Pro mode, effort Pro, S6\_10\_PAPERS, 18 November 2026: computation or formal certificate; targeted proof check. Continuation metadata events: GPT-6 Astra (OpenAI), ChatGPT Pro mode, effort Pro, S6\_10\_PAPERS continuation, 19 November 2026: used and rechecked in this run; status-vocabulary reconciliation; metadata reconciliation only; no new mathematical check. \emph{Pre-continuation recorded proof scope.} Read the full independent high-precision period/theta crosscheck for thirteen roots and the special points; this remains a numerical consistency check. Fully read and freshly replayed unchanged Cspecial\_torsion\_numerics.py; exit zero. Output and SHA256 are retained in code/I/replay\_special/replay\_manifest.json. Numerical outputs retain numerical scope.}{the hypotheses, notation and definitions stated here; the explicit constructions and calculations in the proof}

All $13$ smooth exact-order roots for $2\le n\le8$ were evaluated
at $110$ decimal digits independently of the algebraic norm formula.
The discrepancy in $N^{12}$ is below $10^{-90}$ at every root;
the period coordinates also verify $nP=O$ to that tolerance.
There are $0,0,1,2,2,4,4$ roots for $n=2,\ldots,8$.
\end{proposition}


\begin{proof}
\begin{enumerate}[label=\arabic*.,leftmargin=*]
\item The script \path{Cspecial_torsion_numerics.py} computes the
roots of $F_n$ to $115$ digits and constructs periods of
$y^2=x^3+Ax+B$ using its three roots $e_1,e_2,e_3$. If
$k=(e_2-e_3)/(e_1-e_3)$ and $K$ denotes the complete elliptic
integral with parameter $k$, it uses
\[
 \tau=iK(1-k)/K(k),\qquad
 v=\frac{F\bigl(\arcsin\sqrt{(e_1-e_3)/(-e_3)},k\bigr)}{2K(k)}.
\]
The identity
$\wp(u)=e_3+(e_1-e_3)/\operatorname{sn}^2
(\sqrt{e_1-e_3}\,u,k)$ verifies that this is a zero of $\wp$.
Permutations of the roots are tried until $\operatorname{Im}\tau>0$;
modular reduction and lattice translations improve numerical
conditioning and leave $N$ unchanged. The program then evaluates
the theta quotient directly. It checks both the norm formula and
that the coefficients of $nv$ in the basis $(\tau,1)$ are integral.
The full output gives at least $70$ digits for every $t,\tau,v,N$
and $N^{12}$, including both nonreal conjugates. The largest measured
norm discrepancy is below $2\cdot10^{-108}$. These are numerical
checks of the exact theorem, not numerical proofs of it.
\end{enumerate}
\end{proof}


\section{An analytic proof of the elliptic maximum}

Throughout this section $q=e^{2\pi i\tau}$, $T=\Im\tau$,
$\Delta_0=q\prod_{n\ge1}(1-q^n)^{24}$, and
$t=E_4^3/(1728\Delta_0)$. Thus $t-1=E_6^2/(1728\Delta_0)$.
The discriminant $\Delta_0$ is the modular discriminant and must not
be confused with the discriminant $27t^8(t-1)^3$ of the rational
Weierstrass equation. We use the division identity
\eqref{ari:eq:cs-division}, already proved above, and the modular
invariance of $N$. No interval cover or numerical root computation
enters the following proof.

\begin{proposition}[three uniform bounds on the elliptic norm]\label{ari:C.8.1}
\StatementRecord{ari:C.8.1}{proved}{Historical events: Claude Fable 5.2 (Anthropic), effort high, 15 November 2026: carried as background without fault found. Claude Fable 5.2 (Anthropic), effort high, 16 November 2026: used and rechecked in this run. GPT-6 Astra (OpenAI), ChatGPT Pro mode, effort Pro, S6\_09\_PAPERS, 17 November 2026: targeted proof check. Also checked in this writing. Pre-continuation writing events: GPT-6 Astra (OpenAI), ChatGPT Pro mode, effort Pro, S6\_10\_PAPERS, 18 November 2026: targeted proof check. Continuation metadata events: GPT-6 Astra (OpenAI), ChatGPT Pro mode, effort Pro, S6\_10\_PAPERS continuation, 19 November 2026: used and rechecked in this run; status-vocabulary reconciliation; metadata reconciliation only; no new mathematical check. \emph{Pre-continuation recorded proof scope.} Read source proof; re-derived theta transformation, determinant/norm comparison, distribution constant or stated analytic inequality as applicable. Numerical replay scopes are recorded with their exact statement labels; the analytic argument is checked here.}{the hypotheses, notation and definitions stated here; \cref{ari:eq:ca-sup}; \cref{ari:eq:ca-exp}}

For every $z\in\C$ and every $\tau$ in the standard modular
fundamental region, we have
\[
 \begin{array}{c|ccc}
 \text{restriction}&T\le2&T\le7/5&T\le11/10\\ \hline
 N(z,\tau)&<17/10&<3/2&<143/100.
 \end{array}
\]
If $\wp_\tau(z)=0$ and $T\ge2$, then $N(z,\tau)<541/500$.
\end{proposition}


\begin{proof}
\begin{enumerate}[label=\arabic*.,leftmargin=*]
\item Taking absolute values in the theta series and completing the square
gives
\[
 e^{-\pi(\Im z)^2/T}|\vartheta_1(z,\tau)|
 \le\sum_{n\in\Z}
 e^{-\pi T(n+1/2+\Im z/T)^2}
 \le\sum_{n\in\Z}e^{-\pi Tn^2}.
\]
The last inequality follows from Poisson summation: the Fourier
coefficients of the periodized Gaussian are positive, so its maximum
is at an integer. The eta product and $n^2\ge1+3(n-1)$ for $n\ge1$
therefore imply, on $a\le T\le b$,
\begin{equation}\label{ari:eq:ca-sup}
 N(z,\tau)\le
 \exp\!\left(\frac{r_a}{(1-r_a)^2}\right)
 \max\{F(a),F(b)\},\qquad
 F(T)=e^{\pi T/12}
 \left(1+\frac{2e^{-\pi T}}{1-e^{-3\pi T}}\right),
\end{equation}
where $r_a=e^{-2\pi a}$. Indeed,
$-\sum\log(1-r^n)\le r/(1-r)^2$. The function $F$ is convex,
as follows by expanding its denominator as a geometric series;
each summand is an exponential of a real linear function of $T$.
We use $6/7<\sqrt3/2$ and the following rational bounds:
\[
\begin{array}{c|@{\,}c@{\,}c@{\,}c@{\,}c@{\,}c}
T&6/7&1&11/10&7/5&2\\ \hline
F(T)&<14211/10000&<14116/10000&<14180/10000&<14782/10000&<16944/10000.
\end{array}
\]
The eta factors in \eqref{ari:eq:ca-sup} are less than $10047/10000$
for $a=6/7$ and less than $10019/10000$ for $a=1$.
Multiplying these fractions proves the asserted bounds, using
$[6/7,1]$ and $[1,2]$ separately for the first one.

\item We specify elementary arithmetic for all the scalar estimates in
this proof. One may use $314159/100000<\pi<355/113$ and, for
$0\le x\le20$, the Taylor bounds
\begin{equation}\label{ari:eq:ca-exp}
 S_{80}(x)\le e^x\le
 S_{80}(x)+\frac{x^{81}}{81!}\frac1{1-x/82},
 \qquad S_{80}(x)=\sum_{j=0}^{80}\frac{x^j}{j!}.
\end{equation}
Taking reciprocals bounds negative exponentials. The remainder
bound follows by comparison of consecutive terms, not by numerical
sampling. These inequalities verify each displayed rational bound;
much shorter Taylor sums also suffice. Thus the decimal diagnostics
in the accompanying programs are dispensable.

\item For the cusp we check directly the analytic step of supplied K6.
Choose $|\Im z|\le T/2$ by translation, put $b_0=43/100$ and
$r_0=e^{-4\pi}$. The Fourier series
\[
 \frac{\wp_\tau(z)}{\pi^2}
 =\csc^2(\pi z)-\frac13
   +8\sum_{n\ge1}\frac{nq^n}{1-q^n}
          \bigl(1-\cos(2\pi nz)\bigr)
\]
shows, if $|\Im z|\ge b_0$, that
\begin{align*}
 \left|\frac{\wp_\tau(z)}{\pi^2}+\frac13\right|
 &\le \frac1{\sinh^2(\pi b_0)}
 +8\left\{\frac{r_0}{(1-r_0)^3}
 +\frac1{2(1-r_0)}
 \left(\frac{e^{-2\pi}}{(1-e^{-2\pi})^2}
       +\frac{e^{-6\pi}}{(1-e^{-6\pi})^2}\right)\right\}
 <\frac13.
\end{align*}
The two terms in the last parentheses bound respectively
$r^n e^{2\pi n|\Im z|}$ and $r^n e^{-2\pi n|\Im z|}$.
This inequality excludes a zero. Hence a zero has
$|\Im z|<b_0$. The Jacobi product, $\log(1+x)\le x$, and
the monotonicity of $\cosh(\pi y)e^{-\pi y^2/T}$ on
$0\le y\le b_0$ give
\[
 N(z,\tau)\le
 2\cosh(\pi b_0)e^{-\pi/3-\pi b_0^2/2}
 \exp\!\left(
 \frac{2\cosh(2\pi b_0)r_0}{1-r_0}
 +\frac{r_0^2}{1-r_0^2}\right)<\frac{541}{500}.
\]
For the stated monotonicity, the concave function
$\tanh(\pi y)-y$ is zero at zero and positive at $b_0$.
For $T\ge2$ the right side before substituting $T=2$ decreases:
the derivative of $-\pi T/6-\pi b_0^2/T$ is negative and
both product errors decrease. The two strict scalar inequalities
follow from \eqref{ari:eq:ca-exp}. This proves the cusp assertion
without using any certificate from K6.
\end{enumerate}
\end{proof}

\begin{proposition}[an explicit local maximum on a larger disc]\label{ari:C.8.2}
\StatementRecord{ari:C.8.2}{proved}{Historical events: Claude Fable 5.2 (Anthropic), effort high, 16 November 2026: used and rechecked in this run. GPT-6 Astra (OpenAI), ChatGPT Pro mode, effort Pro, S6\_09\_PAPERS, 17 November 2026: targeted proof check. Also checked in this writing. Pre-continuation writing events: GPT-6 Astra (OpenAI), ChatGPT Pro mode, effort Pro, S6\_10\_PAPERS, 18 November 2026: targeted proof check. Continuation metadata events: GPT-6 Astra (OpenAI), ChatGPT Pro mode, effort Pro, S6\_10\_PAPERS continuation, 19 November 2026: used and rechecked in this run; status-vocabulary reconciliation; metadata reconciliation only; no new mathematical check. \emph{Pre-continuation recorded proof scope.} Read source proof; re-derived theta transformation, determinant/norm comparison, distribution constant or stated analytic inequality as applicable. Numerical replay scopes are recorded with their exact statement labels; the analytic argument is checked here.}{the hypotheses, notation and definitions stated here; \cref{ari:eq:ca-exp}; \cref{ari:eq:ca-derivatives}; \cref{ari:eq:ca-GHbounds}; \cref{ari:eq:ca-H}; \cref{ari:eq:ca-delta-h}}

If $|\tau-i|\le1/50$ and $\wp_\tau(z)=0$, then
$N(z,\tau)\le\sqrt2$, with equality exactly when $\tau=i$.
\end{proposition}


\begin{proof}
\begin{enumerate}[label=\arabic*.,leftmargin=*]
\item Write $\delta=\tau-i$, $h=z-(1-\tau)/2$ and
\[
 G(h,\tau)=\wp_\tau((1-\tau)/2+h),\qquad
 H(h,\tau)=\log\frac{\vartheta_3(h,\tau)}{\eta(\tau)}.
\]
We choose the holomorphic logarithm whose value at $(0,i)$ is
$\log\sqrt2$. The half-period theta identity gives
\begin{equation}\label{ari:eq:ca-H}
 \log N(z,\tau)=\Re H(h,\tau)-\frac{\pi(\Im h)^2}{T}.
\end{equation}
Both $G$ and $H$ are even in $h$. On the polydisc
$|h|\le1/4$, $|\delta|\le1/10$, they satisfy
\begin{equation}\label{ari:eq:ca-GHbounds}
 |G|<30,\qquad |H|<2/3.
\end{equation}
Here are sufficient series estimates, including their margins.
Put
\[
 A_k=\sum_{n\ge1}n^k e^{-9\pi n^2/10}\cosh(\pi n/2),
 \quad a_0=2A_0,\quad a_1=4\pi A_1,\quad a_2=8\pi^2A_2.
\]
Then $a_0<3/10$, $a_1<19/10$, $a_2<12$.
The first two summands suffice to verify these inequalities: for
$n\ge3$, each of the two exponentials in the hyperbolic cosine
has successive-term ratio at most $e^{-29\pi/5}$, and the
factor $n^k$ increases this ratio by at most $(4/3)^2$.
We use \eqref{ari:eq:ca-exp} for the first terms and this geometric
bound for the tail. With $r=e^{-9\pi/5}$, the identities
\[
 G=-\partial_h^2\log\vartheta_3-\frac{\pi^2}{3}E_2,
 \qquad
 |E_2|\le1+\frac{24r}{(1-r)^3}
\]
give
\[
 |G|\le\frac{a_2}{1-a_0}+\frac{a_1^2}{(1-a_0)^2}
       +\frac{\pi^2}{3}\left(1+\frac{24r}{(1-r)^3}\right)<30.
\]
Also
\[
 |H|\le-\log(1-a_0)+\frac{11\pi}{120}
                    +\frac{r}{(1-r)^2}<\frac23.
\]
The logarithm is well defined since $|\vartheta_3-1|<3/10$;
the displayed eta product fixes its branch. A geometric bound
for $-\log(1-x)=\sum_{n\ge1}x^n/n$ verifies the last estimate.

\item Let
\[
 e(\tau)=\frac{\pi^2}{3}(\vartheta_2(0,\tau)^4-
                              \vartheta_4(0,\tau)^4),\qquad
 E(\tau)=\frac{\pi^2}{3}E_2(\tau),\qquad
 g=\frac{g_2(i)}4.
\]
On $|\delta|\le1/50$ the following loose bounds hold:
\begin{equation}\label{ari:eq:ca-derivatives}
 |e'|<14,\quad |e''|<60,\quad |g_2'|<500,
 \quad |E'|<2,\quad |E_6''|<100,
 \qquad 47<g<95/2.
\end{equation}
For completeness these are scalar Fourier-series bounds, not
interval computations. At $T\ge49/50$, upper bounds for
$|\vartheta_j|,|\vartheta_j'|,|\vartheta_j''|$ are
\[
\begin{array}{c|ccc}
j&0\text{ derivatives}&1\text{ derivative}&2\text{ derivatives}\\ \hline
2&929/1000&742/1000&670/1000\\
4&1093/1000&290/1000&911/1000.
\end{array}
\]
They follow by differentiating the absolutely convergent theta
series, taking absolute values, retaining two terms, and using
the same successive-term bound. Differentiating the fourth
powers proves the first two bounds in \eqref{ari:eq:ca-derivatives}.
For $r\le1/470$, the remaining derivative majorants are
\begin{align*}
 |g_2'|&\le640\pi^5\frac{r(1+11r+11r^2+r^3)}{(1-r)^7}<500,\\
 |E'|&\le16\pi^3\frac{r(1+r)}{(1-r)^5}<2,\\
 |E_6''|&\le2016\pi^2
 \frac{r(1+r)(1+120r+1191r^2+2416r^3+1191r^4+120r^5+r^6)}
 {(1-r)^{11}}<100.
\end{align*}
These follow by summing $\sum n^k r^n$ and bounding
$(1-r^n)^{-1}$ by $(1-r)^{-1}$. Substitution of $r=1/470$
and $\pi<355/113$ verifies the rational inequalities.
Finally $g=\pi^4E_4(i)/3$ and
$E_4=1+240\sum n^3q^n/(1-q^n)$ give $47<g<95/2$.
The half-period identity used in Lemma~\ref{push:src:D.2.1},
obtained by differentiating $e$ and using the theta heat equation, reads
\[
 e(i)=0,\qquad e'(i)=id,\qquad d=\frac{2g}{3\pi}>99/10.
\]

\item The differential equation $\wp''=6\wp^2-g_2/2$ determines the
first coefficients in the even expansion:
\[
 G(h,\tau)=e+a h^2+ea h^4+\mathcal R_G,
 \qquad a=3e^2-g_2/4.
\]
Consequently, on $|\delta|\le1/50$,
\[
 |e|\le7/25,\quad |a+g|\le\frac52+3(7/25)^2=\frac{3419}{1250},
 \quad |a|\le\frac{62794}{1250}.
\]
Cauchy's estimate from \eqref{ari:eq:ca-GHbounds} yields
\[
 |\mathcal R_G|\le
 \frac{30(4|h|)^6}{1-16|h|^2}.
\]
On $|h|=1/10$ the term $ah^2$ has modulus at least
$55331/125000>0.4426$, whereas the sum of the other terms is at most
\[
 \frac7{25}+\frac7{25}\frac{62794}{1250}\frac1{10000}
       +\frac{30(4/25)^3}{1-4/25}<0.4277.
\]
Rouch\'e's theorem puts two zeros, with multiplicity, in
$|h|<1/10$. The disc injects into the torus and $\wp$ has exactly
two zeros on the torus, so it contains both of them.

\item Write $G=e+h^2Q$. The preceding estimates give, for $|h|\le1/10$,
\[
 |Q+g|\le
 \frac{3419}{1250}+\frac7{25}\frac{62794}{1250}\frac1{100}
       +\frac{30(4/25)^3}{(1-4/25)/100}<\frac{88}{5}.
\]
Taylor's formula and $|e''|<60$ give
$|e|\ge(99/10-3/5)|\delta|=(93/10)|\delta|$.
At either zero, therefore,
\begin{equation}\label{ari:eq:ca-delta-h}
 |\delta|<7|h|^2,
\end{equation}
because $(95/2+88/5)/(93/10)=7$.

\item Expand
$H=H_0(\tau)+H_2(\tau)h^2+H_4(\tau)h^4+\mathcal R_H$.
The heat equation and the identity for $G$ give
\[
 H_0'=\frac{ie}{4\pi},\quad
 H_2=-\frac{e+E}{2},\quad H_4=-\frac a{12}.
\]
Thus $H_0(i)=\log\sqrt2$, $H_0'(i)=0$,
$|H_0''|<6/5$, $H_2(i)=-\pi/2$, $|H_2'|<8$, and
$|H_4|<21/5$. Cauchy's estimate gives
\[
 |\mathcal R_H|\le
 \frac{(2/3)(4|h|)^6}{1-16|h|^2}
 \le\frac{2048}{63}|h|^4
 \quad (|h|\le1/10).
\]
Combining \eqref{ari:eq:ca-H} and \eqref{ari:eq:ca-delta-h}, we obtain
\begin{align*}
 \log N-\log\sqrt2
 &\le-\frac\pi2\Re(h^2)-\pi(\Im h)^2+C|h|^4\\
 &=-\frac\pi2|h|^2+C|h|^4,\\
 C&=\frac{147}{5}+56+\frac{21}{5}
                 +\frac{2048}{63}+\frac{1100}{49}<150.
\end{align*}
The final term of $C$ bounds the change of $1/T$ from $1$:
$7\pi/(1-1/50)<1100/49$. Since $150|h|^2<3/2<\pi/2$,
the result is strict unless $h=0$. Equation
\eqref{ari:eq:ca-delta-h} then gives $\delta=0$, where equality
holds by the checked isogeny computation at $i$.
\end{enumerate}
\end{proof}

\begin{proposition}[the comparison of the two local discs]\label{ari:C.8.3}
\StatementRecord{ari:C.8.3}{proved}{Historical events: Claude Fable 5.2 (Anthropic), effort high, 15 November 2026: carried as background without fault found. Claude Fable 5.2 (Anthropic), effort high, 16 November 2026: used and rechecked in this run. GPT-6 Astra (OpenAI), ChatGPT Pro mode, effort Pro, S6\_09\_PAPERS, 17 November 2026: targeted proof check. Also checked in this writing. Pre-continuation writing events: GPT-6 Astra (OpenAI), ChatGPT Pro mode, effort Pro, S6\_10\_PAPERS, 18 November 2026: targeted proof check. Continuation metadata events: GPT-6 Astra (OpenAI), ChatGPT Pro mode, effort Pro, S6\_10\_PAPERS continuation, 19 November 2026: used and rechecked in this run; status-vocabulary reconciliation; metadata reconciliation only; no new mathematical check. \emph{Pre-continuation recorded proof scope.} Read source proof; re-derived theta transformation, determinant/norm comparison, distribution constant or stated analytic inequality as applicable. Numerical replay scopes are recorded with their exact statement labels; the analytic argument is checked here.}{the hypotheses, notation and definitions stated here; \cref{ari:eq:ca-derivatives}; \cref{ari:C.8.2}}

Every complex number $t_0$ with $|t_0-1|\le1/500$ has a
representative $\tau$ with $|\tau-i|<1/50$ and $j(\tau)=1728t_0$.
Consequently every zero of $\wp$ for that modulus satisfies
$N\le\sqrt2$, with equality only for $t_0=1$.
\end{proposition}


\begin{proof}
\begin{enumerate}[label=\arabic*.,leftmargin=*]
\item At $i$, $E_6=0$ and the Ramanujan differential identity gives
$|E_6'(i)|=\pi E_4(i)^2>13/2$, since $E_4(i)>29/20$.
On $|\delta|=1/50$, \eqref{ari:eq:ca-derivatives} gives
\[
 |E_6(\tau)|>\frac{13}{100}-\frac1{50}=\frac{11}{100}.
\]
On the whole disc, $|q|<1/470$ and
$|\Delta_0|\le |q|\exp(24|q|/(1-|q|))<1/400$.
It follows on the boundary that
\[
 |t-1|=\frac{|E_6|^2}{1728|\Delta_0|}
 >\frac{121}{43200}>\frac1{500}.
\]
Rouch\'e applied to $E_6'(i)\delta$ shows that $E_6$ has
exactly one zero in the disc. Hence $t-1$ has two zeros there,
with multiplicity, and another application of Rouch\'e shows
that $t-t_0$ has two for each $|t_0-1|\le1/500$.
This argument proves coverage of the $t$-disc, rather than
inferring coverage merely from a boundary bound. Elliptic
curves over $\C$ with the same $j$ are isomorphic. In terms of
lattices the isomorphism is a homothety; it carries the two
zeros of $\wp$ to the two zeros of the scaled $\wp$ and
preserves $N$. Statement \ref{ari:C.8.2} applies to either zero.
\end{enumerate}
\end{proof}

\begin{proposition}[the first three normalized duplication steps]\label{ari:C.8.4}
\StatementRecord{ari:C.8.4}{proved}{Historical events: Claude Fable 5.2 (Anthropic), effort high, 15 November 2026: carried as background without fault found. Claude Fable 5.2 (Anthropic), effort high, 16 November 2026: used and rechecked in this run. GPT-6 Astra (OpenAI), ChatGPT Pro mode, effort Pro, S6\_09\_PAPERS, 17 November 2026: targeted proof check. Also checked in this writing. Pre-continuation writing events: GPT-6 Astra (OpenAI), ChatGPT Pro mode, effort Pro, S6\_10\_PAPERS, 18 November 2026: targeted proof check. Continuation metadata events: GPT-6 Astra (OpenAI), ChatGPT Pro mode, effort Pro, S6\_10\_PAPERS continuation, 19 November 2026: used and rechecked in this run; status-vocabulary reconciliation; metadata reconciliation only; no new mathematical check. \emph{Pre-continuation recorded proof scope.} Read source proof; re-derived theta transformation, determinant/norm comparison, distribution constant or stated analytic inequality as applicable. Numerical replay scopes are recorded with their exact statement labels; the analytic argument is checked here.}{the hypotheses, notation and definitions stated here; \cref{ari:C.4.2}; \cref{ari:eq:cs-division}; \cref{ari:C.3.1}}

Put $u=t-1$. For $n=2^k$, $k=1,2,3$, and $nP\ne O$,
\begin{equation}\label{ari:eq:ca-ladder}
 \left(\frac{N(P)}{\sqrt2}\right)^{n^2}
 =\frac{3^{3/4}}{2n}\,
       \frac{N(nP)|u|^{-1/4}}{|A_k(u)|},
\end{equation}
where
\begin{align*}
 A_1(u)&=1,\qquad A_2(u)=1-5u/27,\\
 A_3(u)&=1-\frac{217}{27}u-\frac{20486}{729}u^2
       +\frac{10910}{19683}u^3
       +\frac{161383}{177147}u^4
       -\frac{245}{14348907}u^5.
\end{align*}
On $|u|\le1/28$ we have $|A_3(u)|\ge1-10|u|$.
For every $k\ge1$, not just these first three steps, define
\[
 E_k(t)=\frac{|\Delta_E|^{(4^k-1)/12}|t-1|^{1/4}}
 {2^{4^k/2}|\psi_{2^k}(P)|}.
\]
Then $(N(P)/\sqrt2)^{4^k}=E_k(t)N(2^kP)|t-1|^{-1/4}$,
and the removable limiting value is
\[
 \lim_{t\to1}E_k(t)=\frac{3^{3/4}}{2^{k+1}}
                 =2^{1-k}\frac{3^{3/4}}4.
\]
\end{proposition}


\begin{proof}
\begin{enumerate}[label=\arabic*.,leftmargin=*]
\item Substitute the exact division polynomials from Statement \ref{ari:C.4.2}
in \eqref{ari:eq:cs-division}. Equivalently, remove
$t^{2(n^2-1)/3}(t-1)^{n^2/4}$ from $\psi_n(P)$ and divide
the resulting polynomial by its value at $t=1$. For $n=2,4,8$
this gives the displayed polynomials and the constant
$3^{3/4}/(2n)$. All operations are rational polynomial operations;
\path{code/C_ladder_algebra.py} reproduces them and records
the next two steps as well. The general identity is \ref{ari:C.4.2}.
For its limit write $\mu=(1-\tau)/2+h$ as in \ref{ari:C.3.1}. Since
$n=2^k$ is even, $n\mu$ is congruent to $nh$ modulo the
period lattice. Hence
\[
 N(nP)=2\pi|\eta(i)|^2 n|h|(1+o(1)),\qquad
 t-1=\tfrac94\pi^4E_4(i)h^4(1+o(1)).
\]
The identity $E_4(i)^3=1728\eta(i)^{24}$ then gives
$N(nP)|t-1|^{-1/4}\to2n/3^{3/4}$.
Since the left side of the general ladder tends to one,
this proves the stated limit for every $k$.
The triangle inequality gives
$|A_3(u)|\ge1-|u|M$, where
\[
 M\le\frac{81}{10}+\frac{29}{28}
                +\frac1{28^2}+\frac1{28^3}+\frac1{28^4}<10.
\]
Here we bounded respectively the five absolute coefficients by
$81/10,29,1,1,1$. This is a rational inequality on one disc,
not a subdivision or a root search.
\end{enumerate}
\end{proof}

\begin{proof}[Division-ladder proof of Theorem~\ref{fam:st:C.2.3}]
\begin{enumerate}[label=\arabic*.,leftmargin=*]
\item Reduce to the standard modular region. The portion $T\ge2$
is covered by Statement \ref{ari:C.8.1}. Suppose $T\le2$ and put
$r=|t-1|$. For $r\ge9/10$, the step $n=2$ in
\eqref{ari:eq:ca-ladder} and $N(2P)<17/10$ give strict inequality,
because
\[
 \frac{27(17/10)^4}{256}<\frac9{10}.
\]
There is no $2$-torsion point of the section on the smooth base.

\item If $r\le9/10$, then $T<7/5$. Indeed, $T\ge7/5$ would give
$|q|<1/6500$ and
\[
 |E_4|>24/25,\qquad |\Delta_0|<1/6400,
 \qquad |t|>\frac{6400(24/25)^3}{1728}>19/10,
\]
contrary to $|t|\le19/10$. The estimates follow from
\[
 \sum_{n\ge1}\sigma_3(n)r^n
 \le\frac{r(1+4r+r^2)}{(1-r)^5},\qquad
 |\Delta_0|\le r\exp\frac{24r}{1-r}.
\]
For $1/28\le r\le9/10$, the $n=4$ step and
$N(4P)<3/2$ now apply. The function
$(1-5r/27)r^{1/4}$ is increasing up to $r=27/25$.
At its left endpoint its fourth power is
\[
 \frac1{28}\left(\frac{751}{756}\right)^4
 >\frac{2187}{65536}
 =\left(\frac{3^{3/4}}8\frac32\right)^4.
\]
This proves strict inequality on the whole annulus. The lower
bound for $A_2$ is positive there, so no division by a zero has
occurred.

\item If $r\le1/28$, then $T<11/10$. For otherwise $|q|<1/1000$,
and the Lambert series with
\[
 \sum_{n\ge1}\sigma_5(n)r^n
 \le\frac{r(1+26r+66r^2+26r^3+r^4)}{(1-r)^7}
\]
gives $|E_6|>9/20$, $|\Delta_0|<1/900$, and hence
$|t-1|>900(9/20)^2/1728>1/28$.
On $1/500\le r\le1/28$, the $n=8$ step,
$N(8P)<143/100$, and $|A_3|\ge1-10r$ prove the claim.
Indeed $(1-10r)r^{1/4}$ has a single critical point in this
interval, a maximum at $r=1/50$, so its minimum is at an endpoint.
At both endpoints its fourth power exceeds $9/5000$, whereas
\[
 \left(\frac{3^{3/4}}{16}\frac{143}{100}\right)^4
 <\frac7{4000}<\frac9{5000}.
\]
Again the displayed positive bound excludes a division-polynomial
zero. The remaining disc $r\le1/500$ is covered by Statement
\ref{ari:C.8.3}. The smooth limiting fibre at $t=0$ has
$N=3^{1/4}<\sqrt2$; the formula also extends there by
continuity, so the singular rational model causes no omitted
modulus. The only equality occurs at $t=1$.
\end{enumerate}
\end{proof}

Theorems~\ref{fam:src:C.2.3} and \cref{push:src:E.4.1}, \cref{push:src:E.4.2} establish this inequality by
computer-assisted proofs. Statements \ref{ari:C.8.1}--\ref{ari:C.8.5} give another
proof with a three-step division ladder and the explicit radius
$1/50$ in the $\tau$-plane. The programs
\path{code/C_ladder_estimates.py} and
\path{code/C_local_analytic_bounds.py} record numerical diagnostics
of the constants; the proof uses the displayed rational estimates
and their series remainder bounds alone.




\section{The retained family and its intrinsic classification}\label{fam:section:A}

Throughout this section $\mathcal P=\{c:\operatorname{Im}c<-M\}$, and $X_c$ means the actual glued space with $v_1=\varepsilon=(1,2,-4,0)^t$, $v_2=-\varepsilon'=(-1,-3,3,0)^t$ and the untwisted cusp collar. We write $d=\widehat\delta$, $e=(0,1)^t$, $E=d\gamma$, and use column vectors. Our peripheral generators are clockwise: $s_j(g_jz)=e^{-2\pi i/m_j}s_j(z)$ and $g_1g_2g_0=1$. Thus $s(g_0z)=s(z)-1$ at the cusp. All period comparisons use
\[
 A_1=\begin{pmatrix}1&0&0&0\\6&0&1&0\\-6&-1&-1&0\\-2&1&0&1\end{pmatrix},\quad
 A_2=\begin{pmatrix}1&0&0&0\\0&0&-1&0\\-6&1&0&0\\3&0&1&1\end{pmatrix},\qquad
 R_g(z)\Pi_c(z)=\Pi_c(gz)A_g.
\]
Here $R_1=\left(\begin{smallmatrix}-1/\tau&0\\(1-\mu)/\tau&1\end{smallmatrix}\right)$, $R_2=\left(\begin{smallmatrix}1/\tau&0\\-\mu/\tau&1\end{smallmatrix}\right)$, and $R_0=I$. The distinction between source and target markings in each equation below is retained by its displayed period equation.

\subsection{Canonical sections and the finite covers}\label{fam:section:A.1}

\begin{proposition}\label{fam:src:A.1.1}\label{fam:st:A.1.1}
\StatementRecord{fam:src:A.1.1}{proved}{Historical events: Claude Fable 5.2 (Anthropic), effort high, 12 November 2026: targeted proof check. Claude Fable 5.2 (Anthropic), effort high, 16 November 2026: used and rechecked in this run. GPT-6 Astra (OpenAI), ChatGPT Pro mode, effort Pro, S6\_09\_PAPERS, 17 November 2026: targeted proof check. Also checked in this writing. Pre-continuation writing events: GPT-6 Astra (OpenAI), ChatGPT Pro mode, effort Pro, S6\_10\_PAPERS, 18 November 2026: computation or formal certificate; targeted proof check. Continuation metadata events: GPT-6 Astra (OpenAI), ChatGPT Pro mode, effort Pro, S6\_10\_PAPERS continuation, 19 November 2026: used and rechecked in this run; status-vocabulary reconciliation; metadata reconciliation only; no new mathematical check. \emph{Pre-continuation recorded proof scope.} Read complete source proof and re-derived the displayed centralizer, finite norm, cusp winding, cohomological gluing, period or reflection argument as applicable; external input locators are separately being verified. PASS exact rational and polynomial verification of matrix orders, centralizer, norm images, even-shift solutions, cusp equation and conjugate-family matrix: code/I/verify\_integral\_classification.py and output JSON. The computation checks the recorded algebra, not analytic extension by itself.}{the hypotheses, notation and definitions stated here; \citet[Theorem~9.1]{AC2026}; \CompendiumRef{I}{app:rederive:fam-src-A.1.1}}

For the constructed holomorphic map $f_c:X_c\to\mathbb P^1$, and also for every admissible choice of its affine translations and cusp collar, we have
\[
 \omega_{X_c}\simeq\mathcal O_{X_c}(-2S_2),\qquad
 f_{c*}\mathcal O_{X_c}(kS_2)=\mathcal O_B(\lfloor k/4\rfloor p_2)\quad(k\in\mathbb Z).
\]
The latter equality is an equality of subsheaves of meromorphic functions under pullback. The anticanonical section ring is $\mathbb C[U,V]$, with $\deg U=1$, $\deg V=2$. The complete degree-two map is $f_c$ followed by a projective coordinate isomorphism of its base. Every biholomorphism between two such threefolds is over the identity of $B$.
\end{proposition}



The canonical pencil is due to \citet[Theorem~9.1]{AC2026}. A complete re-derivation is given in \CompendiumRef{I}{app:rederive:fam-src-A.1.1}.


\begin{proposition}\label{fam:src:A.1.2}\label{fam:st:A.1.2}
\StatementRecord{fam:src:A.1.2}{proved}{Historical events: Claude Fable 5.2 (Anthropic), effort high, 12 November 2026: targeted proof check. Claude Fable 5.2 (Anthropic), effort high, 16 November 2026: used and rechecked in this run. GPT-6 Astra (OpenAI), ChatGPT Pro mode, effort Pro, S6\_09\_PAPERS, 17 November 2026: targeted proof check. Also checked in this writing. Pre-continuation writing events: GPT-6 Astra (OpenAI), ChatGPT Pro mode, effort Pro, S6\_10\_PAPERS, 18 November 2026: targeted proof check. Continuation metadata events: GPT-6 Astra (OpenAI), ChatGPT Pro mode, effort Pro, S6\_10\_PAPERS continuation, 19 November 2026: used and rechecked in this run; status-vocabulary reconciliation; metadata reconciliation only; no new mathematical check. \emph{Pre-continuation recorded proof scope.} Read complete source proof and re-derived the displayed centralizer, finite norm, cusp winding, cohomological gluing, period or reflection argument as applicable; external input locators are separately being verified.}{the hypotheses, notation and definitions stated here; the explicit constructions and calculations in the proof}

The normalized base change $N_j\times_{D_j}\Delta_j\to\Delta_j$, $t_j=s_j^{m_j}$, is the original smooth torus family $J'_j$. The maps $E_1\times F_1\to S_1$ and $E_2\times F_2\to S_2$ described below are finite etale covers of degrees $9$ and $8$.
\end{proposition}


\begin{proof}
\begin{enumerate}[label=\arabic*.,leftmargin=*]
\item On $J'_j$, the deck generator is
$(s,\zeta)\mapsto(\zeta_js,R_j(s)\zeta+\Pi_c(g_js)v_j/m_j)$. It has exact order $m_j$. Its $r$th power changes the invariant circle coordinate $\gamma$ by $r\gamma(v_j)/m_j$; this is nonintegral for $0<r<m_j$, so the action is free. The quotient map is therefore a local biholomorphism. Its local base function is $s^{m_j}$ and its fibre multiplicity is exactly $m_j$. The map
\[
 J'_j\longrightarrow N_j\times_{D_j}\Delta_j,\qquad y\longmapsto(q_j(y),s_j(y))
\]
is finite, is an isomorphism over $s\ne0$, and has normal source. Locally at the central fibre the base-change ring has equation $z^{m_j}-s^{m_j}=0$. Its normalization is the disjoint union of branches $z=\zeta_j^rs$. The displayed map gives those branches in the covering charts. It is thus the normalization everywhere, with its specified root coordinate.

\item Put $\rho=e^{\pi i/3}$, $\mu_1=(2-\rho)/3$, $\mu_2=(1-i)/2$, $\kappa_1=\beta_c(z_1)+2\mu_1$, and $\kappa_2=\beta_c(z_2)+3\mu_2$. The fixed-point transformation laws for $\mu$ give the two displayed values directly. We have $\operatorname{Im}\kappa_j=D_c(z_j)<0$. Define
\[
 E_j=\mathbb C/(\mathbb Z+\kappa_j\mathbb Z),\quad
 F_1=\mathbb C/(\mathbb Z(\rho-2)-3\mathbb Z),\quad
 F_2=\mathbb C/(\mathbb Z(i-1)-2\mathbb Z).
\]
Let $e_1=(0,1,-2,0)^t,f_1=(0,0,-3,1)^t,e_2=(0,1,-1,0)^t,f_2=(0,0,-2,1)^t$. Direct solution of the fixed equations and inspection of the columns give
\[
 \Lambda^{A_1}=\mathbb Z\varepsilon+\mathbb Zd,\quad
 \Lambda^{A_2}=\mathbb Z\varepsilon'+\mathbb Zd,\quad
 (A_j-I)\Lambda=\mathbb Ze_j+\mathbb Zf_j.
\]
The image lattices are saturated: they are the integer solutions of $\gamma=0,2u+w+3\delta=0$ and $\gamma=0,u+w+2\delta=0$, respectively. The determinants of the fixed-plus-moving basis matrices are $-3,-2$.

\item The complex-linear maps $L_1(\xi,\eta)=(\eta,\xi-\eta/3)$ and $L_2(\xi,\eta)=(\eta,\xi-\eta/2)$ carry the product periods to those fixed and moving sublattices. They induce isogenies to the original central tori, with kernel generated respectively by $(2/3,\rho)$ and $(1/2,i)$, of exact orders three and two. For example $\widehat u=2d/3+e_1-2f_1/3=d/2+e_2-f_2/2$ gives both kernel representatives.

\item The affine deck maps lift to
\[
 (\xi,\eta)\longmapsto(\xi+\kappa_1/3,(\rho-1)\eta),\qquad
 (\xi,\eta)\longmapsto(\xi-\kappa_2/4,-i\eta).
\]
Substitution in $L_j$ gives exactly $R_jL_j+\Pi_cv_j/m_j$. They commute with their kernel translations: the moving differences are periods, since $(\rho-2)\rho=-1-\rho=-(\rho-2)-3$ can equally be checked in the lattice basis, and $(-i-1)i=1-i=-(i-1)$. More directly these assertions follow by the two-dimensional period matrices of the indicated rotations. A nonidentity element of either resulting product group translates the $E_j$ coordinate by $2a/3+b\kappa_1/3$ or $a/2-b\kappa_2/4$, respectively. Real independence of $1,\kappa_j$ proves that this translation vanishes only for the identity element. Hence the groups $C_3\times C_3$ and $C_2\times C_4$ act freely and have quotient precisely $S_j$.
\end{enumerate}
\end{proof}

\subsection{The integral obstruction}\label{fam:section:A.2}

\begin{proposition}\label{fam:src:A.2.1}\label{fam:st:A.2.1}
\StatementRecord{fam:src:A.2.1}{proved}{Historical events: Claude Fable 5.2 (Anthropic), effort high, 12 November 2026: targeted proof check. Claude Fable 5.2 (Anthropic), effort high, 16 November 2026: used and rechecked in this run. GPT-6 Astra (OpenAI), ChatGPT Pro mode, effort Pro, S6\_09\_PAPERS, 17 November 2026: targeted proof check. Also checked in this writing. Pre-continuation writing events: GPT-6 Astra (OpenAI), ChatGPT Pro mode, effort Pro, S6\_10\_PAPERS, 18 November 2026: computation or formal certificate; targeted proof check. Continuation metadata events: GPT-6 Astra (OpenAI), ChatGPT Pro mode, effort Pro, S6\_10\_PAPERS continuation, 19 November 2026: used and rechecked in this run; status-vocabulary reconciliation; metadata reconciliation only; no new mathematical check. \emph{Pre-continuation recorded proof scope.} Read complete source proof and re-derived the displayed centralizer, finite norm, cusp winding, cohomological gluing, period or reflection argument as applicable; external input locators are separately being verified. PASS exact rational and polynomial verification of matrix orders, centralizer, norm images, even-shift solutions, cusp equation and conjugate-family matrix: code/I/verify\_integral\_classification.py and output JSON. The computation checks the recorded algebra, not analytic extension by itself.}{the hypotheses, notation and definitions stated here; \cref{fam:src:A.1.2}}

The rational centralizer of $(A_1,A_2)$ is $\mathbb QI+\mathbb QE$. Every biholomorphism $X_c\to X_{c'}$ has $c'-c=n\in2\mathbb Z$, complex-linear fibre part $I_2$, and source-to-target lattice map $I-nE$. In the full fixed lattices,
\[
 \Lambda^{A_1}/N_1\Lambda\simeq\mathbb Z/3,\qquad
 \Lambda^{A_2}/N_2\Lambda\simeq\mathbb Z/4\oplus\mathbb Z/2,
\]
with $d$ generating the second factor of the latter group and vanishing in the former.
\end{proposition}


\begin{proof}
\begin{enumerate}[label=\arabic*.,leftmargin=*]
\item The common invariant vector space is $\mathbb Qd$, and the common invariant covector space is $\mathbb Q\gamma$; this follows by intersecting the fixed bases in \ref{fam:src:A.1.2} and the rows $\langle\gamma,2u+w+3\delta\rangle$, $\langle\gamma,u+w+2\delta\rangle$. A commuting endomorphism preserves $\mathbb Qd\subset\ker\gamma$. Its induced matrix on $\ker\gamma/\mathbb Qd$ commutes with
$Q_1=\left(\begin{smallmatrix}0&1\\-1&-1\end{smallmatrix}\right)$ and $Q_2=\left(\begin{smallmatrix}0&-1\\1&0\end{smallmatrix}\right)$. Commutation with $Q_2$ makes it $\left(\begin{smallmatrix}b&r\\-r&b\end{smallmatrix}\right)$, and the $Q_1$ equation sets $r=0$. Write the remaining columns as
$H\widehat\gamma=(a,x,y,k)^t$, $H\widehat u=b\widehat u+u_0d$, $H\widehat w=b\widehat w+w_0d$, $Hd=d_0d$. The equations on $\widehat u,\widehat w$ give $d_0-w_0=b+u_0$, $u_0=2w_0$ and $u_0=w_0$, hence $u_0=w_0=0,d_0=b$. The first-column equations give $x=2(a-b),y=-4(a-b),x=-y$. Thus $a=b,x=y=0$, proving the assertion. Conversely $A_jE=EA_j=E$.

\item A fibre biholomorphism lifts to an affine map on $\mathbb C^2$: the derivative of a lifted map is lattice periodic, descends to a holomorphic function on a compact torus, and is constant. Its lattice map is locally constant and commutes with both monodromies. Write it as $H=aI+bE$; integrality and $\det H=a^4=1$ give $a=\pm1,b\in\mathbb Z$. The equation $L\Pi_c=\Pi_{c'}H$, in columns three and four, gives $L=aI_2$; column one gives $a(c'-c)+b=0$. Thus $n=c'-c\in\mathbb Z$.

\item The normalized base-change construction in \ref{fam:src:A.1.2} is functorial for biholomorphisms over the base: the base-changed isomorphism transports integral elements of the total quotient ring to integral elements, and its inverse does the same. Its lift is over the identity of the root disc. It commutes with the prescribed deck generator because both compositions lift the same map $(x,s)\mapsto(F(x),\zeta_js)$ on the unnormalized base change. This proves exact commutation, rather than commutation after an unchosen power of the generator.

\item Lift the translation of the resulting affine torus map holomorphically to a vector $B_j(s)$ on the simply connected root disc. Its deck equation, expressed in the source periods, is
\[
 B_j(g_jz)-R_j(z)B_j(z)=\Pi_c(g_jz)
 \left(\frac{(1-a)v_j+n\gamma(v_j)d}{m_j}+\ell_j\right),\qquad\ell_j\in\Lambda.
\]
The integral vector is constant by continuity and real invertibility of the period matrix. Iterating the equation $m_j$ times gives
\[
 N_j\ell_j=-(1-a)v_j-n\gamma(v_j)d.
\]
Exact multiplication yields
\[
 N_1=\begin{pmatrix}3&0&0&0\\6&0&0&0\\-12&0&0&0\\0&2&1&3\end{pmatrix},\quad
 N_2=\begin{pmatrix}4&0&0&0\\12&0&0&0\\-12&0&0&0\\0&2&2&4\end{pmatrix}.
\]
In fixed bases their images are $3\mathbb Z\varepsilon+\mathbb Zd$ and $4\mathbb Z\varepsilon'+2\mathbb Zd$. These explicit generators prove the asserted quotient groups. Applying $\gamma$ to the norm equation excludes $a=-1$ already at $m_1=3$, since $3\gamma(\ell_1)=-2$. For $a=1$, the second norm equation says exactly $2\mid n$; the first imposes no condition. Program \texttt{code/ab\_exact.py} verifies the matrices, Smith forms, centralizer and membership independently by Hermite-form comparisons. All its computations are over $\mathbb Z$ or $\mathbb Q$.
\end{enumerate}
\end{proof}

\subsection{Extension and coherent descent}\label{fam:section:A.3}

\begin{proposition}\label{fam:src:A.3.1}\label{fam:st:A.3.1}
\StatementRecord{fam:src:A.3.1}{proved}{Historical events: Claude Fable 5.2 (Anthropic), effort high, 12 November 2026: targeted proof check. Claude Fable 5.2 (Anthropic), effort high, 16 November 2026: used and rechecked in this run. GPT-6 Astra (OpenAI), ChatGPT Pro mode, effort Pro, S6\_09\_PAPERS, 17 November 2026: targeted proof check. Also checked in this writing. Pre-continuation writing events: GPT-6 Astra (OpenAI), ChatGPT Pro mode, effort Pro, S6\_10\_PAPERS, 18 November 2026: computation or formal certificate; targeted proof check. Continuation metadata events: GPT-6 Astra (OpenAI), ChatGPT Pro mode, effort Pro, S6\_10\_PAPERS continuation, 19 November 2026: used and rechecked in this run; status-vocabulary reconciliation; metadata reconciliation only; no new mathematical check. \emph{Pre-continuation recorded proof scope.} Read complete source proof and re-derived the displayed centralizer, finite norm, cusp winding, cohomological gluing, period or reflection argument as applicable; external input locators are separately being verified. PASS exact rational and polynomial verification of matrix orders, centralizer, norm images, even-shift solutions, cusp equation and conjugate-family matrix: code/I/verify\_integral\_classification.py and output JSON. The computation checks the recorded algebra, not analytic extension by itself.}{the hypotheses, notation and definitions stated here; \cref{fam:src:A.2.1}}

For each $n\in2\mathbb Z$ there is a biholomorphism $F_{n,c}:X_c\to X_{c+n}$, holomorphic jointly in $c$, with $F_{2h,c+2k}F_{2k,c}=F_{2(h+k),c}$. The quotient family over $0<|w|<e^{-\pi M}$, $w=e^{-\pi ic}$, is proper and holomorphic, with pairwise non-biholomorphic fibres.
\end{proposition}


\begin{proof}
\begin{enumerate}[label=\arabic*.,leftmargin=*]
\item Integer shifts satisfy $\Pi_{c+n}=\Pi_c(I+nE)$, so the middle lattices agree. Set
\[
 r_1=-\tfrac n3(\widehat u+\widehat w),\quad
 r_2=\tfrac n4(\widehat u+\widehat w),\quad
 \ell_1=-n\widehat w,\quad \ell_2=\tfrac n2\widehat u,\quad b_j=\Pi_cr_j.
\]
The four-entry equations are $(I-A_1)r_1=(0,0,-n,n/3)^t$ and $(I-A_2)r_2=(0,n/2,0,-n/4)^t$. Consequently $B_j=b_j$ solves the preceding deck equation for $a=1$. The map $(z,\zeta)\mapsto(z,\zeta+b_j(z))$ extends on the entire root family and descends, together with its inverse, to each finite filling.

\item Its middle translation on the punctured overlap is
\[
 a_j=b_j+n\gamma(v_j)\frac{\log s_j}{2\pi i}e.
\]
This follows from the actual overlap $h_{j,c}(\zeta)=\zeta+\sigma_{j,c}$, where $\sigma_{j,c}=(\log s_j)/(2\pi i)\,\Pi_cv_j$: precisely $\sigma_{j,c+n}-\sigma_{j,c}=n\gamma(v_j)(\log s_j)/(2\pi i)e$. Changes of logarithm add periods. Since the chosen generator changes $(\log s_j)/(2\pi i)$ by $-1/m_j$, we have $a_j(g_jz)=R_ja_j(z)+\Pi_c(g_jz)\ell_j$.

\item On the universal cover of $B^\circ$, extend $\ell_1,\ell_2$ by the integral cocycle law $\ell_{gh}=\ell_g+A_g\ell_h$. Freeness of $\pi_1(B^\circ)$ makes this well defined. At $g_0=(g_1g_2)^{-1}$ it gives
\[
 \ell_0=-A_0(\ell_1+A_1\ell_2)=\tfrac{3n}2\widehat w-\tfrac n2d.
\]
Put $q_n=(-3n/2,n/2)^t$ and $a_0=sq_n$. Then $a_0(g_0z)-a_0(z)=\Pi_c\ell_0$. In cusp torus coordinates its map is $(x,t_c)\mapsto(t_c^{q_n}x,t_c)$. The integral linear map $(y,h)\mapsto(y+hq_n,h)$ of the fan lattice translates each height-one triangle by an integer vector. It therefore permutes every cone of the prescribed periodic triangulation. The map and its inverse extend across all strata, including both triple points and all three double curves after quotient. Moreover the cusp deck multipliers are identical at $c$ and $c+n$, since $C_{c+n}\lambda-C_c\lambda=(0,n\lambda_1)^t$ has integral exponential. Thus the extended map descends with all its nonconstant units retained.

\item We supply the global analytic correction. Quotient the trivial affine vector bundle on the universal cover by
$(z,y)\mapsto(gz,R_g(z)y+\Pi_c(gz)\ell_g)$. Its action law is exactly the cocycle law by period covariance. Specify at the finite points sections of the form $a_j+\eta_j$, with $\eta_j$ holomorphic equivariant on the full root disc; at the cusp specify $a_0+\eta_0(t_c)$ with $\eta_0$ holomorphic on the full disc. These prescriptions define a locally nonempty torsor under a rank-two locally free sheaf $\mathcal V$ on $B$.

\item Projection to the first coordinate gives
\[
 0\longrightarrow\mathcal O_Be\longrightarrow\mathcal V
 \longrightarrow\mathcal O_B(-p_0)\longrightarrow0.
\]
Indeed its factors are $-1/\tau,1/\tau,1$. Multiplication by $\Theta$ identifies this first-coordinate sheaf with holomorphic functions at ordinary and finite points and functions vanishing at $p_0$: the finite vanishing orders are $2,1$, and the cusp pole is one. At a finite stalk, any equivariant first component has an equivariant vector lift by averaging a preliminary lift $(u,0)$ over the finite group. This proves stalkwise surjectivity and local freeness, including the off-diagonal entries of $R_j$.

\item Both $H^1(\mathbb P^1,\mathcal O)$ and $H^1(\mathbb P^1,\mathcal O(-1))$ vanish. To make the analytic gluing explicit, split a smooth cocycle with a partition of unity, solve $\bar\partial$ locally by the Cauchy integral, and compare the two resulting solutions on an annulus. A Laurent series $\sum a_nt^n$ is the difference of its nonnegative and negative parts for $\mathcal O$; for $\mathcal O(-1)$ use the second frame $t^{-1}$, so the negative terms become $t^{-1}(-\sum_{n<0}a_nt^{n+1})$. Both series converge on their respective discs and leave no exponent unsolved. This kills the cocycle for any refinement. The exact sequence gives $H^1(B,\mathcal V)=0$. Project the torsor cocycle to $\mathcal O(-1)$, solve it, lift the solutions, and solve the remaining scalar cocycle. We obtain one global section $a$, with $a=a_j+\eta_j$ at finite points and $a=sq_n+\eta_0$ at the cusp.

\item The three actual maps are now
\[
 \zeta\mapsto\zeta+a,\qquad
 \zeta\mapsto\zeta+b_j+\eta_j,\qquad
 (x,t_c)\mapsto(e^{2\pi i\eta_0(t_c)}t_c^{q_n}x,t_c).
\]
The finite overlap identity is the equality of sums
$\sigma_{j,c}+a=\sigma_{j,c+n}+b_j+\eta_j$. The cusp overlap identity is $e^{2\pi ia}=t_c^{q_n}e^{2\pi i\eta_0}$. The holomorphic-unit factors act on every toric chart. The local inverses satisfy the same identities; hence the maps glue to a biholomorphism.

\item For coherence, perform the additive solution only for $n=2$ and multiply that solution by $k$ for $n=2k$. The cocycles have $\gamma(\ell_g)=0$, so their affine terms $\Pi_c\ell_g$ are independent of $c$. The reference translations, $\mathcal V$, and the correction can all be chosen independent of $c$. Adding the translations and multiplying the toric factors proves the stated cocycle identity exactly.

\item Finally the original gluing constructs a holomorphic family over any parameter disc compactly contained in $\mathcal P$. In ordinary and finite charts, real invertibility of $\Pi_c$ is uniform on compact parameter sets, so a fixed compact real parallelepiped proves properness and the lattice quotient is holomorphic. At the cusp write $Z=sB_0+C(t_c,c)$ with $B_0=\left(\begin{smallmatrix}0&1\\-1&0\end{smallmatrix}\right)$. On a compact parameter disc $C$ is uniformly bounded; shrink the cusp radius until
$\|(-2\pi\operatorname{Im}C)/\log|t_c|\|<1/2$. The effective logarithmic lattice is then within $1/2$ of the unimodular $B_0$, uniformly. Reduce its logarithmic valuation modulo the effective logarithmic lattice to a bounded parallelogram. Choose a height-one triangle containing that reduced valuation. Its barycentric coordinates are nonnegative and sum to one; the corresponding three unimodular cone characters therefore have modulus at most one. This puts the point in the closed unit polydisc of that cone. Only finitely many triangles meet the fixed bounded parallelogram, uniformly on the parameter compactum. Boundary orbits have the same representatives by translating their height-one components to these triangles. Thus every orbit has a representative in this finite union of cone polydiscs. Those closed polydiscs times a compact parameter set form a compact surjective representative set. This proves properness at the cusp and proper discontinuity by the same uniform lower bound on nonzero lattice translations. The quotient charts are open subsets of a fixed smooth toric variety times the parameter disc, so projection to the parameter is a submersion. The explicit overlap maps are jointly holomorphic. These local families agree after shrinking cusp radii, and glue over $\mathcal P$. The coherent $2\mathbb Z$ action is properly discontinuous because its action on $\mathcal P$ is. Local inverse logarithms for $w=e^{-\pi ic}$ identify the quotient family with these proper local families. \ref{fam:src:A.2.1} gives pairwise non-biholomorphic fibres.
\end{enumerate}
\end{proof}

\subsection{Albanese invariants and the derivative of the deformation}\label{fam:section:A.4}

\begin{proposition}\label{fam:src:A.4.1}\label{fam:st:A.4.1}
\StatementRecord{fam:src:A.4.1}{proved}{Historical events: Claude Fable 5.2 (Anthropic), effort high, 12 November 2026: targeted proof check. Claude Fable 5.2 (Anthropic), effort high, 16 November 2026: used and rechecked in this run. GPT-6 Astra (OpenAI), ChatGPT Pro mode, effort Pro, S6\_09\_PAPERS, 17 November 2026: targeted proof check. Also checked in this writing. Pre-continuation writing events: GPT-6 Astra (OpenAI), ChatGPT Pro mode, effort Pro, S6\_10\_PAPERS, 18 November 2026: targeted proof check. Continuation metadata events: GPT-6 Astra (OpenAI), ChatGPT Pro mode, effort Pro, S6\_10\_PAPERS continuation, 19 November 2026: used and rechecked in this run; status-vocabulary reconciliation; metadata reconciliation only; no new mathematical check. \emph{Pre-continuation recorded proof scope.} Read complete source proof and re-derived the displayed centralizer, finite norm, cusp winding, cohomological gluing, period or reflection argument as applicable; external input locators are separately being verified.}{the hypotheses, notation and definitions stated here; \cref{fam:src:A.1.2}; \cref{fam:src:A.2.1}}

With $\kappa_j(c)=\kappa_j^{\rm part}+c$, the Albanese elliptic curves of the reduced finite fibres are
\[
 \operatorname{Alb}(S_1)=\mathbb C/(\mathbb Z+\kappa_1(c)\mathbb Z),\qquad
 \operatorname{Alb}(S_2)=\mathbb C/(\mathbb Z+\tfrac12\kappa_2(c)\mathbb Z).
\]
Writing $J(q)=q^{-1}+744+196884q+\cdots$, their invariants are
\[
 j_1(w)=J(e^{-2\pi i\kappa_1^{\rm part}}w^2),\qquad
 j_2(w)=J(e^{-\pi i\kappa_2^{\rm part}}w).
\]
These are convergent formulas because the displayed nome moduli are below one. Outside the explicitly defined countable exceptional set in the proof, equality of both Albanese invariants forces $c'-c\in2\mathbb Z$. Thus the generic Albanese invariants recover the parameter modulo its exact period subgroup, independently of the centralizer calculation.
\end{proposition}


\begin{proof}
\begin{enumerate}[label=\arabic*.,leftmargin=*]
\item In the product covers of \ref{fam:src:A.1.2} the invariant holomorphic one-forms are exactly the multiples of $d\xi$: the moving factor is multiplied by a primitive third or fourth root. Integration thus identifies the Albanese with the quotient of the fixed line by the fixed-coordinate projections of every deck translation. These are, respectively,
\[
 \mathbb Z+\kappa_1\mathbb Z+\tfrac23\mathbb Z+\tfrac{\kappa_1}3\mathbb Z
   =\tfrac13\mathbb Z+\tfrac{\kappa_1}3\mathbb Z,
\quad
 \mathbb Z+\kappa_2\mathbb Z+\tfrac12\mathbb Z+\tfrac{\kappa_2}4\mathbb Z
   =\tfrac12\mathbb Z+\tfrac{\kappa_2}4\mathbb Z.
\]
Multiplication by three and two gives the stated maps of complex tori. In particular the Albanese computation includes the kernel of the product isogeny and the affine quotient, rather than just its fixed eigenspace. Using the positive-imaginary representatives $-\kappa_1$ and $-\kappa_2/2$ in the standard modular $j$-series gives the formulas. At sufficiently small nome $J$ is injective: $1/J(q)=q+O(q^2)$ has nonzero derivative at zero. Thus $j_2$ alone distinguishes parameters within a sufficiently small punctured subdisc when both parameters lie there. For the full generic assertion put $x=-\kappa_2(c)/2$ and $d_*=\kappa_1^{\rm part}-\kappa_2^{\rm part}$, so $-\kappa_1(c)=2x-d_*$. The independently certified interval in \texttt{cd\_special\_intervals.py} gives $-0.2<\operatorname{Im}d_*<-0.1$; in particular it is nonzero. Equality of the two elliptic invariants at $c,c'$ gives matrices $M,N\in\operatorname{PSL}_2(\mathbb Z)$ with $x'=M(x)$ and
\[
 N(2x-d_*)=2M(x)-d_*.
\]
Define the exceptional set to be the $c\in\mathcal P$ obtained from roots of these equations for all pairs for which the two rational functions are not identical. Each equation has at most two roots after clearing denominators, so this is a countable, explicitly specified set. For a pair giving an identity, let $C(x)=2x-d_*$. Then $N=CMC^{-1}$. The real matrix $M$ preserves both the real circle $\mathbb R\cup\{\infty\}$ and $C^{-1}(\mathbb R\cup\{\infty\})$, the parallel line with imaginary part $\operatorname{Im}d_*/2$ together with infinity. Their unique intersection is infinity, so $M$ fixes infinity. An element of $\operatorname{PSL}_2(\mathbb Z)$ fixing infinity is $x\mapsto x+b$, $b\in\mathbb Z$, and then $N(y)=y+2b$. It follows that $c'-c=-2b$. This proves the generic necessity with its actual exceptional set; \ref{fam:src:A.2.1} proves the necessity also on that set.
\end{enumerate}
\end{proof}

\begin{proposition}\label{fam:src:A.4.2}\label{fam:st:A.4.2}
\StatementRecord{fam:src:A.4.2}{proved}{Historical events: Claude Fable 5.2 (Anthropic), effort high, 12 November 2026: targeted proof check. Claude Fable 5.2 (Anthropic), effort high, 16 November 2026: used and rechecked in this run. GPT-6 Astra (OpenAI), ChatGPT Pro mode, effort Pro, S6\_09\_PAPERS, 17 November 2026: targeted proof check. Also checked in this writing. Pre-continuation writing events: GPT-6 Astra (OpenAI), ChatGPT Pro mode, effort Pro, S6\_10\_PAPERS, 18 November 2026: targeted proof check. Continuation metadata events: GPT-6 Astra (OpenAI), ChatGPT Pro mode, effort Pro, S6\_10\_PAPERS continuation, 19 November 2026: used and rechecked in this run; status-vocabulary reconciliation; metadata reconciliation only; no new mathematical check. \emph{Pre-continuation recorded proof scope.} Read complete source proof and re-derived the displayed centralizer, finite norm, cusp winding, cohomological gluing, period or reflection argument as applicable; external input locators are separately being verified.}{the hypotheses, notation and definitions stated here; \cref{fam:src:A.1.1}}

For every $c\in\mathcal P$, the Kodaira--Spencer class of $\partial/\partial c$ in $H^1(X_c,T_{X_c})$ is nonzero. On any fixed ordinary fibre, put $T=\operatorname{Im}\tau$, $m=\operatorname{Im}\mu$, $a=m/T$, $D=\operatorname{Im}\beta_c-6m^2/T$. In the convention $\mathrm{KS}(\partial_c)=[\bar\partial h]$ for a smooth $(1,0)$ lift $h$, its restricted torus-deformation representative is
\[
 -\frac{1}{2iD}\,\partial_{\zeta_2}\otimes
       (d\bar\zeta_2-a\,d\bar\zeta_1).
\]
\end{proposition}


\begin{proof}
\begin{enumerate}[label=\arabic*.,leftmargin=*]
\item Retain the same real marking $r\in\mathbb R^4$ and set $\zeta=\Pi_cr$. The two imaginary-part equations give
$r_1=(\operatorname{Im}\zeta_2-a\operatorname{Im}\zeta_1)/D$. Thus at $c+\delta$ the marked real-linear identification is
\[
 w=\zeta+\frac{\delta}{2iD}P(\zeta-\bar\zeta),\qquad
 P=\begin{pmatrix}0&0\\-a&1\end{pmatrix},\qquad P^2=P.
\]
It is exactly $\Pi_{c+\delta}\Pi_c^{-1}$ as a real map, so it sends every marked lattice vector to the corresponding vector. Its determinant is $(D+\operatorname{Im}\delta)/D>0$. The transported $(0,1)$ graph coefficient is $\nu_\delta=\delta P/(2iD+\delta)$. The derivative at zero is $P/(2iD)$. Equivalently the smooth lift of $\partial_c$ in moving coordinates has vertical part $r_1\partial_{w_2}$ and $\bar\partial r_1=-(d\bar\zeta_2-a d\bar\zeta_1)/(2iD)$, giving the stated sign.

\item A nonzero constant tangent-valued $(0,1)$ form on a compact torus cannot be $\bar\partial$ of a global vector field. Write that field in the invariant holomorphic frame and integrate each derivative of its periodic coefficients against invariant probability volume. Every such derivative has integral zero, whereas the coefficient of $\partial_{\zeta_2}\otimes d\bar\zeta_2$ above has nonzero integral. This proves nonvanishing for the fibre deformation.

\item We explain its transfer to the threefold. Over the dual numbers $A=\mathbb C[\epsilon]/(\epsilon^2)$ at $c$, the chart proof of \ref{fam:src:A.1.1} gives the relative line isomorphism $\omega^{-2}_{X_A/A}=f_A^*\mathcal O_B(p_2)$. Extra parameter differentials disappear in relative forms and all determinant characters are independent of $c$. The two pulled-back sections form an $A$-basis of $H^0(X_A,\omega^{-2})$: reduce a section modulo $\epsilon$, subtract its linear combination of the two lifts, and use the exact sequence $0\to L\xrightarrow{\epsilon}L_A\to L\to0$ to write the remainder uniquely as $\epsilon$ times another such combination. The same argument proves independence.

\item A first-order trivialization would preserve this complete system and give a base automorphism reducing to the identity. It fixes all three special values even over $A$. On the finite etale charts the relative critical ideal is $(s_j^{m_j-1})$, and $t_j=s_j^{m_j}$ is zero on that scheme; its constant parameter ring $A$ injects. At a cusp triple chart, the critical ideal of $t_c=z_0z_1z_2$ is $(z_0z_1,z_0z_2,z_1z_2)$ and has the same property. There are no other critical values. Therefore an infinitesimal displacement of a special value would remain nonzero in its critical chart, a contradiction. Writing a base infinitesimal automorphism as $t+\epsilon(b_0+b_1t+b_2t^2)$, fixing $0,1,\infty$ gives $b_0=b_1=b_2=0$. The trivialization would consequently trivialize the displayed nonzero torus deformation. Finally a zero Kodaira--Spencer class produces such a first-order trivialization by subtracting from local product charts the vector fields whose coboundary is the differentiated transition cocycle. This proves the claim.
\end{enumerate}
\end{proof}

\begin{proposition}\label{fam:src:A.4.3}\label{fam:st:A.4.3}
\StatementRecord{fam:src:A.4.3}{proved}{Historical events: Claude Fable 5.2 (Anthropic), effort high, 12 November 2026: targeted proof check. GPT-6 Astra (OpenAI), ChatGPT Pro mode, effort Pro, S6\_09\_PAPERS, 17 November 2026: targeted proof check. Also checked in this writing. Pre-continuation writing events: GPT-6 Astra (OpenAI), ChatGPT Pro mode, effort Pro, S6\_10\_PAPERS, 18 November 2026: targeted proof check. Continuation metadata events: GPT-6 Astra (OpenAI), ChatGPT Pro mode, effort Pro, S6\_10\_PAPERS continuation, 19 November 2026: used and rechecked in this run; status-vocabulary reconciliation; metadata reconciliation only; no new mathematical check. \emph{Pre-continuation recorded proof scope.} Read complete source proof and re-derived the displayed centralizer, finite norm, cusp winding, cohomological gluing, period or reflection argument as applicable; external input locators are separately being verified.}{the hypotheses, notation and definitions stated here; \cite[Proposition 4.7]{ViaclovskyLog2026}; \cref{fam:src:A.4.2}; \cref{fam:src:A.1.1}; \cite[Lemmas 6.6--6.9 and Proposition 6.15]{Yang2026}; \cref{fam:src:A.2.1}; \cref{fam:src:A.3.1}}

The non-isotriviality argument of Viaclovsky, \cite[Proposition 4.7]{ViaclovskyLog2026}, follows for the present family from \ref{fam:src:A.4.2} and \ref{fam:src:A.1.1}. The cited countability argument proves local non-isotriviality and positive Kuranishi dimension; the computation in \ref{fam:src:A.4.2} additionally proves nonzero derivative at every admissible parameter. The arithmetic method of Yang, \cite[Lemmas 6.6--6.9 and Proposition 6.15]{Yang2026}, specializes to the same centralizer, finite-norm and collar-gluing operations after replacing his displayed matrices by $A_1,A_2$; those replacements are evaluated in \ref{fam:src:A.2.1} and \ref{fam:src:A.3.1}.
\end{proposition}


\begin{proof}
\begin{enumerate}[label=\arabic*.,leftmargin=*]
\item On a fixed ordinary fibre, the three last periods remain fixed and span $\mathbb C^2$, whereas the first changes by $\delta e$. A marked complex-linear identification fixing those three periods must be the identity, so the marked period map is nonconstant. The set of integral markings of one torus is countable. A holomorphic nonconstant map from a disc has uncountable image; hence the unmarked fibre family is locally non-isotrivial. \ref{fam:src:A.1.1} transfers this to total spaces. This is exactly the implication used in the cited Proposition 4.7, whose statement and proof we have read. A nonconstant classifying map gives positive dimension of its local target. It does not by itself determine the derivative at the base point, as the explicit holomorphic map $z\mapsto z^2$ has nonconstant image and derivative zero at zero. \ref{fam:src:A.4.2} computes the derivative here.

\item Yang's cited lemmas use an integral commuting lattice map, an affine fibre lift, normalization of the root base change, and a norm-lattice defect. We have exhibited the respective maps and norm equations for the present matrices. His Proposition 6.15 obtains $6/r$ from the first-column scaling and its order-six norm image. Here the first-column scaling is one and the two computed norm images impose parity, yielding two. The common proof operations transfer through these explicitly stated matrix substitutions; no conclusion about one family is inferred merely from the label of the other. His nonzero absolute Kodaira--Spencer conclusion is obtained in the present setting by the explicit computation and intrinsic-pencil transfer above.
\end{enumerate}
\end{proof}

\subsection{The assembled theorem}\label{fam:section:A.5}

\begin{theorem}\label{fam:src:A.5.1}\label{fam:st:A.5.1}\label{wb:pshift:thm:classification}
\StatementRecord{fam:src:A.5.1}{proved}{Historical events: Claude Fable 5.2 (Anthropic), effort high, 12 November 2026: targeted proof check. Claude Fable 5.2 (Anthropic), effort high, 16 November 2026: used and rechecked in this run. Claude Fable 5.2 (Anthropic), effort high, 17 November 2026: carried as background without fault found. GPT-6 Astra (OpenAI), ChatGPT Pro mode, effort Pro, S6\_09\_PAPERS, 17 November 2026: targeted proof check. Also checked in this writing. Pre-continuation writing events: GPT-6 Astra (OpenAI), ChatGPT Pro mode, effort Pro, S6\_10\_PAPERS, 18 November 2026: targeted proof check; used and rechecked in this run. Continuation metadata events: GPT-6 Astra (OpenAI), ChatGPT Pro mode, effort Pro, S6\_10\_PAPERS continuation, 19 November 2026: used and rechecked in this run; status-vocabulary reconciliation; metadata reconciliation only; no new mathematical check. \emph{Pre-continuation recorded proof scope.} Read complete source proof and re-derived the displayed centralizer, finite norm, cusp winding, cohomological gluing, period or reflection argument as applicable; external input locators are separately being verified. Opened DGKXv1 Theorem1.4 p3 and the pair/bare equivalence proof passages pp92,95 in the original107-page PDF, matching the exact supplied-source labels. The receiving complete proof and its8September2026 working-edition record remain separately credited. This closes the named statement/page-source gate only. Checked this credited theorem statement against its already reviewed fam:src:A.5.1 proof and the current complete centralizer/parity/filling arguments; this is an alias scope record, not an additional independent derivation.}{the hypotheses, notation and definitions stated here; \citet[Theorem~1.4, p.~3]{DuEtAl2026v1}; \cref{app:rederive:fam-src-A.5.1}}

Fix the original cusp-regular period solution, the clockwise peripheral marking, the prescribed periodic triangular cusp fan, the affine vectors $(\varepsilon,-\varepsilon')$, and the untwisted cusp collar. For every $c,c'\in\mathcal P$, the constructed compact complex threefolds satisfy
\[
 X_c\simeq X_{c'}\quad\Longleftrightarrow\quad c'-c\in2\mathbb Z.
\]
Every biholomorphism is over the identity of the original $\mathbb P^1$. The proper holomorphic family of their isomorphism classes has parameter $0<|w|<e^{-\pi M}$ and nonzero Kodaira--Spencer class at every point.
\end{theorem}



The standard classification is \citet[Theorem~1.4, p.~3]{DuEtAl2026v1}. A complete re-derivation is given in Appendix~\ref{app:rederive:fam-src-A.5.1}.


\begin{theorem}\label{fam:src:A.5.2}\label{fam:st:A.5.2}\label{inv:D1.3}\label{inv:D2.1}\label{inv:D3.1}\label{cdp:st:E9}
\StatementRecord{fam:st:A.5.2}{proved}{Historical events: Claude Fable 5.2 (Anthropic), effort high, 16 November 2026: carried as background without fault found. Claude Fable 5.2 (Anthropic), effort high, 17 November 2026: carried as background without fault found. GPT-6 Astra (OpenAI), ChatGPT Pro mode, effort Pro, S6\_09\_PAPERS, 17 November 2026: targeted proof check. Also checked in this writing. Pre-continuation writing events: GPT-6 Astra (OpenAI), ChatGPT Pro mode, effort Pro, S6\_10\_PAPERS, 18 November 2026: computation or formal certificate; used and rechecked in this run. Continuation metadata events: GPT-6 Astra (OpenAI), ChatGPT Pro mode, effort Pro, S6\_10\_PAPERS continuation, 19 November 2026: used and rechecked in this run; status-vocabulary reconciliation; metadata reconciliation only; no new mathematical check. \emph{Pre-continuation recorded proof scope.} Read complete source proof and re-derived the displayed centralizer, finite norm, cusp winding, cohomological gluing, period or reflection argument as applicable; external input locators are separately being verified. Re-read and re-derived the precise original-cusp dependency used by II inv:E2.2/E15. In the complete receiving Koszul/direct-image proof, the kernel of the restricted ambient cotangent map is nu\_*O\_dP6 times the common conormal dq; its global section space is one-dimensional. Gorenstein Serre duality gives h2(W,T\_X|W)=1. Top-degree base change and generic rank1 force the cusp R2T stalk free over C\{q\}, hence the upper modification with quotient C\_p0; the branch-normal quotient is nu\_*K\_dP6 by the explicit coefficient map and invariant volume form. This is a renewed local application check, not another whole-proof credit. Opened DGKXv1 Lemma6.11 and full proof p81, Theorem6.12 and full proof pp81–82, and Proposition6.16 and full proof p85; matched branch-normal adjunction, cohomology and the actual injective local-ring substitution. The current receiving rederivation and renewed cusp module check are retained as separate events. Exact co-located alias labels share this same source reading, without additional independent credit. PASS: unchanged d\_exact.py and d\_koszul.py replays. Two independent weighted cohomology counts for 481 integers, tangent Euler characteristics, root-character invariants, polynomial Smith forms and joint order-three coefficient 123i/(64pi). These are finite algebra checks; analytic comparison is in the targeted proof event. Verified exact label identity with fam:src:A.5.2: all six co-located labels occupy the same statement at source characters44182–45155 in S6\_09\_PAPERS\_I.tex, SHA2565089ddf726828eea0c2d6a1542e7777ae4cc9694a0e06338db13875c56773534. The complete receiving Koszul, direct-image, Kodaira–Spencer and Kuranishi proof already reviewed under fam:src:A.5.2 therefore supplies this identical statement. This is an identity propagation, with no additional independent proof credit.}{the hypotheses, notation and definitions stated here; \cref{fam:src:A.5.1}; \citet[Theorem~6.12, pp.~81--82, and Proposition~6.16, p.~85]{DuEtAl2026v1}; \cref{app:rederive:fam-src-A.5.2}; \cref{fam:section:B.2}; \cref{fam:section:B.3}; \cite[Sections 2--3 and Proposition 6.1]{ViaclovskyTwo2026}; \cref{fam:eq:ab-normal-quotient}}

For every admissible $c$, the tangent cohomology of the actual compact
member is
\[
 \bigl(h^0(X_c,T_{X_c}),h^1(X_c,T_{X_c}),h^2(X_c,T_{X_c}),
 h^3(X_c,T_{X_c})\bigr)=(1,1,1,0).
\]
Its Kuranishi germ is a smooth reduced disc. The family of \ref{fam:src:A.5.1} is
miniversal at every parameter: its Kodaira--Spencer map is an isomorphism,
and every sufficiently small deformation is induced from this family.
More precisely, for the intrinsic fibration $f:X_c\to\mathbb P^1$,
\[
 f_*T_{X_c}=\mathcal O\oplus\mathcal O(-1),\qquad
 R^1f_*T_{X_c}=\mathcal O\oplus\mathcal O(-1)^2,\qquad
 R^2f_*T_{X_c}=\mathcal O.\tag{D.4}
\]
The Leray class $\xi\otimes[\gamma]$ generates $H^1(X_c,T_{X_c})$.
The primary Kuranishi bracket is zero; the one-dimensional obstruction
space $H^2(X_c,T_{X_c})$ therefore does not imply a nonzero obstruction
map.

\end{theorem}



The tangent dimensions and Kuranishi germ are re-derived from \citet[Theorem~6.12, pp.~81--82, and Proposition~6.16, p.~85]{DuEtAl2026v1}. A complete re-derivation is given in Appendix~\ref{app:rederive:fam-src-A.5.2}.


\paragraph{Scope and a precise changed-input test.}
The hypotheses of \ref{fam:src:A.5.1}--\ref{fam:src:A.5.2} fix the original period system and
multiplicities $(3,4)$; \ref{fam:section:B.2} treats its arbitrary affine and cusp data,
and \ref{fam:section:B.3} classifies its stipulated higher-multiplicity variants.
A concrete extension beyond these hypotheses begins with
$y^2=x^3+tx+1$ in \cite[Sections 2--3 and Proposition 6.1]{ViaclovskyTwo2026}.
Already one rank-one filling over a specified root $p_1$ of
$4t^3+27=0$ replaces the hexagonal normalization calculation by
\[
 \begin{gathered}
 E'=\mathbb C^*/\langle c_0\rangle,\qquad
 R=\mathbb P_{E'}(\mathcal O_{E'}\oplus L),\\
 R=(\mathbb P^1\times\mathbb C^*)/
       \langle(z,v)\mapsto(\gamma_0z,c_0v)\rangle.
 \end{gathered}
\]
with its boundary sections identified by $v\mapsto\beta_0v$.
Here $c_0,\gamma_0,\beta_0$ are exactly the nonzero central multipliers
in that paper's equation (3.1). The invariant nodal volume form there
trivializes the ambient canonical line, so the branch-normal map has
quotient $\nu_*K_R$ by the same local coefficient map as
\eqref{fam:eq:ab-normal-quotient}. The ruling has
$\pi_*\mathcal O_R=\mathcal O_{E'}$ and $R^1\pi_*\mathcal O_R=0$;
Leray and Serre duality therefore give $h^\bullet(K_R)=(0,1,1)$.
The first unevaluated global datum for transferring the deformation
argument is the connecting map from this single
$H^1(R,K_R)=\mathbb C$ summand into $H^2(T_X(-\log D))$, where
$D$ is the union of the three reduced singular fibres of that model.
Its direct attack is to compute the boundary-translation Cech cocycle
using the full holomorphic unit functions in that equation (3.1),
their central values $(\gamma_0,c_0,\beta_0)$, and the corresponding
vertical frames.
The global identification argument also requires the integral
centralizer and collar defects of that model's four monodromy matrices.
These are specified substitutions to perform for the stated one-fibre
test; no classification or deformation conclusion for that changed
period system is included in \ref{fam:src:A.5.1}--\ref{fam:src:A.5.2}.

\section{Discrete data and the geometric parameter}\label{fam:section:B}

\subsection{All affine vectors and cusp sections}\label{fam:section:B.2}

\begin{proposition}\label{fam:src:B.2.1}\label{fam:st:B.2.1}
\StatementRecord{fam:src:B.2.1}{proved}{Historical events: Claude Fable 5.2 (Anthropic), effort high, 12 November 2026: targeted proof check. Claude Fable 5.2 (Anthropic), effort high, 16 November 2026: used and rechecked in this run. Claude Fable 5.2 (Anthropic), effort high, 17 November 2026: carried as background without fault found. GPT-6 Astra (OpenAI), ChatGPT Pro mode, effort Pro, S6\_09\_PAPERS, 17 November 2026: targeted proof check. Also checked in this writing. Pre-continuation writing events: GPT-6 Astra (OpenAI), ChatGPT Pro mode, effort Pro, S6\_10\_PAPERS, 18 November 2026: computation or formal certificate; targeted proof check. Continuation metadata events: GPT-6 Astra (OpenAI), ChatGPT Pro mode, effort Pro, S6\_10\_PAPERS continuation, 19 November 2026: used and rechecked in this run; status-vocabulary reconciliation; metadata reconciliation only; no new mathematical check. \emph{Pre-continuation recorded proof scope.} Read complete source proof and re-derived the displayed centralizer, finite norm, cusp winding, cohomological gluing, period or reflection argument as applicable; external input locators are separately being verified. PASS exact rational and polynomial verification of matrix orders, centralizer, norm images, even-shift solutions, cusp equation and conjugate-family matrix: code/I/verify\_integral\_classification.py and output JSON. The computation checks the recorded algebra, not analytic extension by itself.}{the hypotheses, notation and definitions stated here; \cite[Section 6.4]{AC2026}; \cref{fam:src:A.1.1}; \cref{fam:src:A.3.1}}

Let
\[
 w_1=r_1\varepsilon+k_1d,\qquad w_2=r_2\varepsilon'+k_2d,
 \qquad 3\nmid r_1,\quad 2\nmid r_2,
\]
and use an arbitrary holomorphic cusp collar section $\sigma$ in the construction of \cite[Section 6.4]{AC2026}. Write $l_0=\gamma(c(\sigma))$ in its positive cusp-winding convention, and set $p=12l_0-4r_1-3r_2$. Let primes denote a second such datum with the same period solution and cusp fan. Then
\[
 X(c;w_1,w_2,\sigma)\simeq X(c';w'_1,w'_2,\sigma')
\]
holds exactly for the following arithmetic conditions:
\[
 n=c'-c\in\mathbb Z,\qquad
 p'=a p\text{ for some }a\in\{1,-1\},\qquad
 n\equiv k_2-k'_2\pmod2.
\]
All maps are over $B$. A map with the indicated $a$ has fibre-linear part $aI_2$. In particular the complete invariants are
\[
 \bigl(|p|,\ e^{-\pi i(c+k_2)}\bigr),\qquad
 |p|\ge1,\quad\gcd(|p|,6)=1.
\]
Every value of this positive integer occurs. The full collar section enters the construction, but its additional analytic part and its second integral winding contribute no further invariant.
\end{proposition}


\begin{proof}
\begin{enumerate}[label=\arabic*.,leftmargin=*]
\item The determinant calculation in \ref{fam:src:A.1.1} was independent of translations, so the base is still intrinsic. The centralizer proof gives $c'=c+n$ and $aI_2$ exactly as before. Put $S_n=I+nE$ and compare all target periods in source coordinates. On the finite root cover the exact deck equation is
\[
 B_j(g_jz)-R_jB_j(z)=\Pi_c(g_jz)\left(\frac{d_j}{m_j}+\ell_j\right),
 \quad d_j=S_nw'_j-a w_j,
 \quad N_j\ell_j=-d_j.                                      \tag{B2.1}
\]
Necessity follows by the normalized-base-change lift already proved. Conversely the norm equation has a rational solution $r_j^*$ of
$(I-A_j)r_j^*=d_j/m_j+\ell_j$, since $\ker N_j$ over $\mathbb Q$ equals the image of $I-A_j$. We can choose $\gamma(r_j^*)=0$: the columns $\widehat u,\widehat w$ already span that image over $\mathbb Q$. Then $B_j=\Pi_cr_j^*$ is holomorphic on the whole root disc and solves the exact equation. Thus the norm condition is sufficient for a local isomorphism, not only for conjugacy on the central real torus.

\item The computed norm images translate (B2.1) into
\[
 r'_1\equiv ar_1\pmod3,\quad r'_2\equiv ar_2\pmod4,
 \quad ak_2-k'_2-nr'_2\equiv0\pmod2.                       \tag{B2.2}
\]
There is no condition on $k_1$. Since $r'_2$ is odd, the last condition is the displayed parity. On the finite collar the middle translation is
$A_j=B_j+(\log s_j)/(2\pi i)\,\Pi_cd_j$ and has integral defect $\ell_j$. Its $\gamma$ defect is
$\alpha_j=(ar_j-r'_j)/m_j$.

\item We prove the precise cusp necessity used next. A holomorphic isomorphism between cusp neighbourhoods with fibre-linear part $aI_2$ has, after lifting near a generic smooth central point to the toric covers, the form
\[
 (x,t_c)\longmapsto(t_c^q u(t_c)x^a,t_c),\qquad q\in\mathbb Z^2,
 \quad u(0)\in(\mathbb C^*)^2.                           \tag{B2.3}
\]
Here the assertion concerns this local lift near a generic central point. To prove it, take a source chart near the generic point of its component and a target chart near its image. The etale toric quotient permits a local lift over a simply connected neighbourhood. The torus characters in a chart at the height-one ray $v$ are $x_i=t_c^{v_i}X_i$, with $X_i$ invertible. On $t_c\ne0$ the affine fibre equation gives $x'_i/x_i^a$ independent of the $X$ coordinates. Its expression in the two charts is therefore $t_c^{v'_i-av_i}$ times a holomorphic unit of $t_c$, including zero. This proves (B2.3). Taking logarithms shows that its winding is in $\Lambda_{\rm tor}=\mathbb Z\widehat w+\mathbb Zd$. Changing a lift adds $(I-A_0)\Lambda=\Lambda_{\rm tor}$, so the vanishing of its class in $\bar\Lambda=\Lambda/\Lambda_{\rm tor}$ is intrinsic. This proof supplies the necessary extension test; no assertion about arbitrary punctured translations is used.

\item We retain the actual collar convention. The map from the cusp piece to the middle family is $T_{-\sigma}G_0^{-1}$. Consequently the middle translation induced by (B2.3) is
\[
 A_0=sq+a\sigma-\sigma'+\eta_0,
\]
with $\eta_0$ holomorphic on the cusp disc after a logarithm of the unit. For the clockwise generator $s\mapsto s-1$, the class of its defect in $\bar\Lambda$ is $-a c(\sigma)+c(\sigma')$. On the other hand every middle integral cocycle satisfies
$\ell_0=-A_0(\ell_1+A_1\ell_2)$, whose $\gamma$ coordinate is $-\alpha_1-\alpha_2$. Equality of these coordinates is exactly
\[
 l'_0-a l_0=-(ar_1-r'_1)/3-(ar_2-r'_2)/4,
\]
or $p'=ap$. This proves the remaining necessity. Notice that $p\equiv-r_1\pmod3$ and $p\equiv r_2\pmod2$, so $p$ is nonzero and coprime to six.

\item For sufficiency, $p'=ap$ implies the first two congruences of (B2.2) by reduction modulo three and four. The parity supplies the remaining one. Choose integral solutions $\ell_1,\ell_2$ of (B2.1). Their $\gamma$ coordinates already have the required values. Adding any element of $\ker N_1=(I-A_1)\Lambda$ preserves (B2.1). In particular adding an arbitrary integer multiple of $e_1=(0,1,-2,0)^t$ changes the $\widehat u$ coordinate of $\ell_0$ by an arbitrary integer while leaving its $\gamma$ coordinate fixed. We can therefore arrange
$[\ell_0]=-a c(\sigma)+c(\sigma')$ in the entire rank-two group $\bar\Lambda$. Choose lifts of $\sigma,\sigma'$ on the cusp half-plane. Their remaining defect differs from $\ell_0$ by an actual toric lattice vector; an integral $q$ makes $sq+a\sigma-\sigma'$ have exactly defect $\ell_0$.

\item The finite translations and this cusp reference now define a locally nonempty affine torsor with linear difference sheaf precisely $\mathcal V$ of \ref{fam:src:A.3.1}. The terms $a\sigma-\sigma'$ are retained, even when their chosen lifts have essential singularities at the puncture; they belong to the specified collar reference, while differences of allowed extensions are holomorphic $\eta_0$. The proof of $H^1(B,\mathcal V)=0$ gives a global middle translation. The finite maps descend by (B2.1), and the cusp map is $x\mapsto t_c^qu(t_c)x^a$. The fan is preserved by all integer translations and by $y\mapsto-y$; exponentials of the integer period shift agree. Hence this cusp map and its inverse extend over every stratum. The finite overlap is the equality defining $A_j$, and the cusp overlap is the equality defining $A_0$. Thus the maps glue globally.

\item Finally parity is equivalent to equality of the two displayed exponentials. Passing from $p$ to $-p$ is effected by $a=-1$ and the same proof; for fixed sign the congruences have just been shown sufficient. Every positive integer coprime to six is obtained by choosing $r_1\in\{1,2\}$ and odd $r_2$ in the appropriate residue classes modulo three and four, and then taking $l_0=(p+4r_1+3r_2)/12$. This also proves that the extra collar data have been accounted for, rather than omitted from the statement.
\end{enumerate}
\end{proof}

\begin{proposition}\label{fam:src:B.2.2}\label{fam:st:B.2.2}
\StatementRecord{fam:src:B.2.2}{proved}{Historical events: Claude Fable 5.2 (Anthropic), effort high, 12 November 2026: targeted proof check. Claude Fable 5.2 (Anthropic), effort high, 16 November 2026: used and rechecked in this run. GPT-6 Astra (OpenAI), ChatGPT Pro mode, effort Pro, S6\_09\_PAPERS, 17 November 2026: targeted proof check. Also checked in this writing. Pre-continuation writing events: GPT-6 Astra (OpenAI), ChatGPT Pro mode, effort Pro, S6\_10\_PAPERS, 18 November 2026: computation or formal certificate; targeted proof check. Continuation metadata events: GPT-6 Astra (OpenAI), ChatGPT Pro mode, effort Pro, S6\_10\_PAPERS continuation, 19 November 2026: used and rechecked in this run; status-vocabulary reconciliation; metadata reconciliation only; no new mathematical check. \emph{Pre-continuation recorded proof scope.} Read complete source proof and re-derived the displayed centralizer, finite norm, cusp winding, cohomological gluing, period or reflection argument as applicable; external input locators are separately being verified. PASS exact rational and polynomial verification of matrix orders, centralizer, norm images, even-shift solutions, cusp equation and conjugate-family matrix: code/I/verify\_integral\_classification.py and output JSON. The computation checks the recorded algebra, not analytic extension by itself.}{the hypotheses, notation and definitions stated here; \cref{fam:src:A.5.1}; \cref{fam:src:B.2.1}; \cref{fam:src:A.2.1}}

For multiplicities $(3,4)$ and $|p|=1$, every admissible discrete member is biholomorphic to one member of \ref{fam:src:A.5.1}, uniquely modulo $2\mathbb Z$. At fixed nonzero admissible Engel parameter $u$ there are exactly two isomorphism classes, exchanged by a positive circuit of $u$ about zero. In the convention that Engel's root generator is the inverse of our clockwise one and his circle class is $d$, the identity of the middle quotient sends
\[
 X_{c+1}\longleftrightarrow X(u;a_0-d/3,\ a_1+d/4).
\]
Reversing both declared root generators reverses these two displayed signs.
\end{proposition}


\begin{proof}
\begin{enumerate}[label=\arabic*.,leftmargin=*]
\item Take the canonical target data $(r'_1,r'_2,k'_1,k'_2,l'_0)=(1,-1,0,0,0)$, for which $p'=-1$. Choose $a$ so that $p'=ap$, and an integer $n\equiv k_2\pmod2$. \ref{fam:src:B.2.1} supplies the isomorphism to $X_{c+n}$. The uniqueness is \ref{fam:src:A.5.1}. At fixed $u=e^{-2\pi ic}$, its logarithms differ by integers; the invariant $e^{-\pi i(c+k_2)}$ takes exactly the two square roots of $u$. Both occur, already as $X_c,X_{c+1}$, and \ref{fam:src:A.2.1} distinguishes them.

\item We spell out the integral comparison to Engel's common basis $B_0=(e_1,e_2,e_3,e_4)$, retaining his $C f_4=e_1+e_4$:
\[
 P_{\rm E}=\begin{pmatrix}0&-2&-1&-2\\2&-1&0&-1\\0&2&1&3\\-1&0&0&0\end{pmatrix},\qquad
 T_0P_{\rm E}=P_{\rm E}A_1,\quad T_1P_{\rm E}=P_{\rm E}A_2.
\]
Its determinant is $-1$, $P_{\rm E}d=\delta_0$, $\psi P_{\rm E}=-\gamma$, $P_{\rm E}\varepsilon=-e_4$ and $P_{\rm E}\varepsilon'=\delta_0-Cf_4$. These are exact lattice homomorphism identities, verified independently through Viaclovsky's intermediate basis in \texttt{ab\_geometry.py}. They also agree with the complex periods: at $z_1$ the complex-linear map
$(\zeta_1,\zeta_2)\mapsto(-\zeta_1/(\rho+1),\zeta_2+\zeta_1/3)$ carries $\Pi_cP_{\rm E}^{-1}$ to Engel's columns $(1,0),(\rho-1,0),((1+\rho)/3,1/3),(0,-\kappa_1)$. At $z_2$ his $Cf_4$ has preimage $d-\varepsilon'$ and fixed period $1-\kappa_2$. Thus the comparison fixes the positive circle class and the actual local complex structures, in addition to their monodromy matrices.

\item For the sign assertion, inverting the concrete clockwise affine generator changes its invariant translation $\Pi_cv_j/m_j$ to $-\Pi_cv_j/m_j$. Thus our dictionary is $a_{j-1}=-P_{\rm E}v_j/m_j$ in Engel's lattice marking (and its period image in the torus). Integer shift gives $\Pi_{c+1}v_j=\Pi_c(v_j+\gamma(v_j)d)$, with $\gamma(v_1)=1,\gamma(v_2)=-1$, proving the two signs. A positive $u$ circuit changes $\log u$ by $2\pi i$ and hence changes $c$ by $-1$, exchanging the square roots. This gives the claimed permutation with its orientation.
\end{enumerate}
\end{proof}

\begin{proposition}[counterexample, with corrected family]\label{fam:src:B.2.3}\label{fam:st:B.2.3}
\StatementRecord{fam:src:B.2.3}{proved}{Historical events: Claude Fable 5.2 (Anthropic), effort high, 12 November 2026: targeted proof check. Claude Fable 5.2 (Anthropic), effort high, 16 November 2026: used and rechecked in this run. Claude Fable 5.2 (Anthropic), effort high, 17 November 2026: carried as background without fault found. GPT-6 Astra (OpenAI), ChatGPT Pro mode, effort Pro, S6\_09\_PAPERS, 17 November 2026: targeted proof check. Also checked in this writing. Pre-continuation writing events: GPT-6 Astra (OpenAI), ChatGPT Pro mode, effort Pro, S6\_10\_PAPERS, 18 November 2026: computation or formal certificate; targeted proof check. Continuation metadata events: GPT-6 Astra (OpenAI), ChatGPT Pro mode, effort Pro, S6\_10\_PAPERS continuation, 19 November 2026: used and rechecked in this run; status-vocabulary reconciliation; metadata reconciliation only; no new mathematical check. \emph{Pre-continuation recorded proof scope.} Read complete source proof and re-derived the displayed centralizer, finite norm, cusp winding, cohomological gluing, period or reflection argument as applicable; external input locators are separately being verified. PASS exact rational and polynomial verification of matrix orders, centralizer, norm images, even-shift solutions, cusp equation and conjugate-family matrix: code/I/verify\_integral\_classification.py and output JSON. The computation checks the recorded algebra, not analytic extension by itself.}{the hypotheses, notation and definitions stated here; \cite[Remark 8.1]{Engel2026}; \cref{fam:src:A.3.1}; \cref{fam:src:A.2.1}; \cref{fam:src:B.2.1}; \cref{fam:src:A.4.1}}

The assertion in \cite[Remark 8.1]{Engel2026} of a global proper family over the original punctured $u$-disc for fixed sphere logarithmic data requires a double-cover qualification. The retained family with such data does not descend through $u=w^2$. It does give a proper holomorphic family on each simply connected $u$-open set, and a global proper family over the punctured $w$-disc as in \ref{fam:src:A.3.1}.
\end{proposition}


\begin{proof}
\begin{enumerate}[label=\arabic*.,leftmargin=*]
\item Descent through $u=w^2$ would furnish an isomorphism between the two pulled-back fibres at $w$ and $-w$, namely $X_c$ and $X_{c+1}$. \ref{fam:src:A.2.1} excludes that isomorphism by the nonzero class of $d$ in $\Lambda^{A_2}/N_2\Lambda$. This is an exact order-two obstruction on the normalized finite cover. It applies to all sphere data by \ref{fam:src:B.2.1} because $r_2$ is odd.

\item There is also a direct analytic witness near zero. The relative Albanese invariant of the multiplicity-four reduced fibre in a holomorphic descended family would be a single-valued holomorphic function of $u$. After pullback it would be even in $w$. But \ref{fam:src:A.4.1} gives $j_2(w)=C^{-1}w^{-1}+744+O(w)$, where $C=e^{-\pi i\kappa_2^{\rm part}}\ne0$; its Laurent coefficient of $w^{-1}$ is nonzero. Thus it is not even on any punctured neighbourhood. This verifies the obstruction independently of a choice of isomorphism lift. Local inverse logarithms of $u$ give the proper local families already constructed; coherent even shifts give their global double-cover family. The untwisted middle quotient itself has the integer shear identification, so this refutation concerns the fixed sphere logarithmic data in the stated remark.
\end{enumerate}
\end{proof}

\subsection{Higher multiplicities}\label{fam:section:B.3}

\begin{proposition}\label{fam:src:B.3.1}\label{fam:st:B.3.1}
\StatementRecord{fam:src:B.3.1}{proved}{Historical events: Claude Fable 5.2 (Anthropic), effort high, 12 November 2026: targeted proof check. Claude Fable 5.2 (Anthropic), effort high, 16 November 2026: used and rechecked in this run. Claude Fable 5.2 (Anthropic), effort high, 17 November 2026: carried as background without fault found. GPT-6 Astra (OpenAI), ChatGPT Pro mode, effort Pro, S6\_09\_PAPERS, 17 November 2026: targeted proof check. Also checked in this writing. Pre-continuation writing events: GPT-6 Astra (OpenAI), ChatGPT Pro mode, effort Pro, S6\_10\_PAPERS, 18 November 2026: targeted proof check. Continuation metadata events: GPT-6 Astra (OpenAI), ChatGPT Pro mode, effort Pro, S6\_10\_PAPERS continuation, 19 November 2026: used and rechecked in this run; status-vocabulary reconciliation; metadata reconciliation only; no new mathematical check. \emph{Pre-continuation recorded proof scope.} Read complete source proof and re-derived the displayed centralizer, finite norm, cusp winding, cohomological gluing, period or reflection argument as applicable; external input locators are separately being verified.}{the hypotheses, notation and definitions stated here; \cref{fam:src:A.1.2}; \cref{fam:src:A.2.1}; \cref{fam:src:B.2.1}}

Let $k,l$ be positive integers with $\gcd(3k,4l)=1$, and put $m_1=3k,m_2=4l$. Consider Viaclovsky's sphere members expressed in the same clockwise marking, with $w_j=r_j\varepsilon_j+s_jd$, $\gcd(r_j,m_j)=1$, cusp winding $l_0$, and $p=m_1m_2l_0-m_2r_1-m_1r_2=\pm1$, where $\varepsilon_1=\varepsilon,\varepsilon_2=\varepsilon'$. Two such members are biholomorphic precisely when their multiplicity pair agrees and, for some $a\in\{1,-1\}$ and integer $n=c'-c$,
\[
 p'=ap,\quad r'_j\equiv ar_j\pmod{m_j},\quad
 nr'_1\equiv as_1-s'_1\pmod{k},\quad
 nr'_2\equiv as_2-s'_2\pmod{2l}.
\]
The congruences on $r'_j$ already follow from $p'=ap$. For fixed discrete vectors and cusp collar the identification subgroup is $2kl\mathbb Z$. For fixed admissible $u$ and this multiplicity pair, all the stated sphere translation data give exactly $2kl$ isomorphism classes, permuted in one cycle by a circuit of $u$ about zero.
\end{proposition}


\begin{proof}
\begin{enumerate}[label=\arabic*.,leftmargin=*]
\item We first recover the base for these sphere members without using the degree-two canonical pencil. A biholomorphism $F:X\to X'$ gives $g=f'F$. Since $H^2(S^6,\mathbb R)=0$, the closed form $g^*\omega_{\rm FS}$ is exact. Restrict to a regular torus $T$ and wedge with a flat Kahler form on $T$. Stokes' theorem gives integral zero, while the integrand is nonnegative and strictly positive wherever $dg|_T\ne0$. Thus $g$ is constant on each regular fibre. The image of $(f,g)$ is a compact analytic curve whose projection to $B$ has degree one; its normalization gives a holomorphic $h:B\to B$ with $g=hf$. Apply the same construction to $F^{-1}$ to see that $h$ is an automorphism. The three homological local monodromy orders are three, four and infinity, independently of $k,l$; their intertwining maps under $F$ therefore fix each special point. Hence $h=1$. The nonreduced fibre multiplicities now force equal $k,l$. This is the full reduction in the read Proposition 4.7 of Viaclovsky's paper, written here with its exact maps.

\item The root model of degree $m_1$ is the pullback of the degree-three model by $s\mapsto s^k$, and likewise at the other point. Its linear action remains $A_1$ or $A_2$, while its root generator has order $m_j$. The primitive invariant coordinate $r_j/m_j$ proves freeness of every nonidentity power. The normalization and deck-commutation arguments of \ref{fam:src:A.1.2} and \ref{fam:src:A.2.1} apply to these root covers without changing a generator. Their norms are
\[
 N_{3k}=kN_1,\qquad N_{4l}=lN_2,
\]
so the fixed-coordinate norm images are $3k\mathbb Z\varepsilon+k\mathbb Zd$ and $4l\mathbb Z\varepsilon'+2l\mathbb Zd$. The norm equation (B2.1), with $m_j$ replaced by these actual covering degrees, gives exactly the displayed congruences. The centralizer and first-column comparison remain unchanged.

\item For completeness the integer $p$ records the primitive sphere relation in this marking. The common torus coinvariants are detected by $\gamma$, since $(A_1-I)\Lambda+(A_2-I)\Lambda=\ker\gamma$ by their explicit saturated bases. The finite meridian relations are $x^{m_1}=c^{r_1}$, $y^{m_2}=c^{r_2}$, and the cusp relation is $xy=c^{l_0}$ with $c$ central. Eliminating $y$ gives an abelian relation matrix with determinant $-p$ and primitive first elementary divisor because $\gcd(m_1,m_2)=1$. Thus the sphere data have $p=\pm1$. Reducing $p'=ap$ modulo $m_1,m_2$ gives the two congruences for $r_j$.

\item To prove sufficiency, use the local rational solutions of the norm equation and the identical cusp comparison in \ref{fam:src:B.2.1}. The first-coordinate extension on the new root discs has vanishing orders $2k$ and $l$, precisely the orders of the pulled-back $\Theta$. Dividing by it again gives $\mathcal O_B(-p_0)$, and finite averaging again gives the same rank-two extension sheaf. Thus $H^1(B,\mathcal V)=0$ still removes the analytic cocycle. The images of $I-A_j$ and their arbitrary $\widehat u$ adjustments are unchanged, so $p'=ap$ is exactly the remaining cusp winding condition. Every local map and overlap is the one already written in \ref{fam:src:B.2.1}, with the new denominator $m_j$; this proves global sufficiency.

\item For identical data, $p\ne0$ excludes $a=-1$. The congruences reduce to $k\mid nr_1$ and $2l\mid nr_2$. The two primitivity assumptions imply $k\mid n$ and $2l\mid n$. Since $\gcd(k,2l)=1$, their least positive common multiple is $2kl$. For different sphere translation data, $r'_1$ is invertible modulo $k$, $r'_2$ modulo $2l$, and the Chinese remainder theorem gives exactly one integer shift class modulo $2kl$. The values of $s_1$ modulo $k$ and $s_2$ modulo $2l$ all occur without affecting freeness or $p$. Thus there are exactly $2kl$ classes at fixed $u$. The change $c\mapsto c-1$ on a positive $u$ circuit acts by a generator of this cyclic residue group. Program \texttt{ab\_exact.py} verifies the repeated norms and shift steps for five coprime pairs; the preceding symbolic identities prove all pairs.
\end{enumerate}
\end{proof}

\subsection{Automorphisms and the remaining real-structure calculation}\label{fam:section:B.4}

\begin{proposition}\label{fam:src:B.4.1}\label{fam:st:B.4.1}
\StatementRecord{fam:src:B.4.1}{proved}{Historical events: Claude Fable 5.2 (Anthropic), effort high, 12 November 2026: targeted proof check. Claude Fable 5.2 (Anthropic), effort high, 16 November 2026: used and rechecked in this run. Claude Fable 5.2 (Anthropic), effort high, 17 November 2026: carried as background without fault found. GPT-6 Astra (OpenAI), ChatGPT Pro mode, effort Pro, S6\_09\_PAPERS, 17 November 2026: targeted proof check. Also checked in this writing. Pre-continuation writing events: GPT-6 Astra (OpenAI), ChatGPT Pro mode, effort Pro, S6\_10\_PAPERS, 18 November 2026: computation or formal certificate; targeted proof check; used and rechecked in this run. Continuation metadata events: GPT-6 Astra (OpenAI), ChatGPT Pro mode, effort Pro, S6\_10\_PAPERS continuation, 19 November 2026: used and rechecked in this run; status-vocabulary reconciliation; metadata reconciliation only; no new mathematical check. \emph{Pre-continuation recorded proof scope.} Read complete source proof and re-derived the displayed centralizer, finite norm, cusp winding, cohomological gluing, period or reflection argument as applicable; external input locators are separately being verified. PASS exact rational and polynomial verification of matrix orders, centralizer, norm images, even-shift solutions, cusp equation and conjugate-family matrix: code/I/verify\_integral\_classification.py and output JSON. The computation checks the recorded algebra, not analytic extension by itself. Opened DGKXv1 Theorem7.14 and complete proof p95, including its exact intrinsic-fibration and vertical-translation inputs. The current source locator is corrected to that theorem and stable label; no whole-proof check of all source dependencies is inferred from this bounded reading.}{the hypotheses, notation and definitions stated here; \citet[Theorem~7.14, p.~95, label \nolinkurl{p2:thm:full-automorphisms}]{DuEtAl2026v1}; \cref{app:rederive:fam-src-B.4.1}}

For each member of the retained $(3,4)$ family, $\operatorname{Aut}(X_c)$, including all its components, is $\mathbb C/\mathbb Z\simeq\mathbb C^*$. Its action on the ordinary fibres is $\zeta\mapsto\zeta+be$, and the exponential group isomorphism is $b\mapsto e^{2\pi ib}$.
\end{proposition}
\Uses{the hypotheses, notation and definitions stated here; \citet[Theorem~1.4, p.~3]{DuEtAl2026v1}; \cref{app:rederive:fam-src-B.4.1}}


The full automorphism group is \citet[Theorem~1.4, p.~3]{DuEtAl2026v1}. A complete re-derivation is given in Appendix~\ref{app:rederive:fam-src-B.4.1}.



\begin{proposition}\label{fam:src:B.4.2}\label{fam:st:B.4.2}
\StatementRecord{fam:src:B.4.2}{proved}{Historical events: Claude Fable 5.2 (Anthropic), effort high, 12 November 2026: targeted proof check. Claude Fable 5.2 (Anthropic), effort high, 16 November 2026: used and rechecked in this run. GPT-6 Astra (OpenAI), ChatGPT Pro mode, effort Pro, S6\_09\_PAPERS, 17 November 2026: targeted proof check. Also checked in this writing. Pre-continuation writing events: GPT-6 Astra (OpenAI), ChatGPT Pro mode, effort Pro, S6\_10\_PAPERS, 18 November 2026: computation or formal certificate; targeted proof check. Continuation metadata events: GPT-6 Astra (OpenAI), ChatGPT Pro mode, effort Pro, S6\_10\_PAPERS continuation, 19 November 2026: used and rechecked in this run; status-vocabulary reconciliation; metadata reconciliation only; no new mathematical check. \emph{Pre-continuation recorded proof scope.} Read complete source proof and re-derived the displayed centralizer, finite norm, cusp winding, cohomological gluing, period or reflection argument as applicable; external input locators are separately being verified. PASS exact rational and polynomial verification of matrix orders, centralizer, norm images, even-shift solutions, cusp equation and conjugate-family matrix: code/I/verify\_integral\_classification.py and output JSON. The computation checks the recorded algebra, not analytic extension by itself.}{the hypotheses, notation and definitions stated here; \cref{fam:src:B.4.1}; \cref{fam:src:A.3.1}; \cref{fam:src:B.2.1}; \cref{fam:src:A.5.1}}

Use the normalization $(\beta^{\rm part}+\tau)(p_0)=0$. There is an antiholomorphic diffeomorphism $X_c\to X_{c'}$ exactly when
\[
 c'-(1-\bar c)\in2\mathbb Z,
 \qquad\text{equivalently }w'=-\bar w.
\]
Consequently $\overline{X_c}\simeq X_{1-\bar c}$. The members admitting antiholomorphic self-maps are exactly those with $\operatorname{Re}c\in\tfrac12+\mathbb Z$, equivalently $w=-\bar w$. On each such member the group of holomorphic and antiholomorphic diffeomorphisms is
\[
 \operatorname{Aut}^{\pm}(X_c)\simeq\mathbb C^*\rtimes\mathbb Z/2,
 \qquad z\longmapsto1/\bar z
\]
for the action of the second factor. Its antiholomorphic maps form the nonidentity coset. For all other members the full group is the holomorphic group of \ref{fam:src:B.4.1}.
\end{proposition}


\begin{proof}
\begin{enumerate}[label=\arabic*.,leftmargin=*]
\item We first construct the comparison before returning its affine data to the canonical choice. Lift conjugation of the real orbifold base to the antiholomorphic involution $\iota$ of its universal cover whose cusp branch has $\tau(\iota z)=-\overline{\tau(z)}$. A lift exists by the universal covering property; it extends at the finite orbifold points, whose local equations are $s_j^{m_j}=t_j$. Choose the usual paths from a real point $t>1$, so reflection fixes the second centre and sends the first chosen lift to $g_2^{-1}z_1$. The induced involution of the deck group is
\[
 \phi(g_1)=g_2^{-1}g_1^{-1}g_2,\qquad
 \phi(g_2)=g_2^{-1},\qquad \phi(g_0)=g_0^{-1}.
\]
These identities also follow by reflecting the lifted peripheral paths; their product and their squares agree with $g_1g_2g_0=1$. The cusp branch fixes the lift uniquely.

\item All reflection identities for the periods can be checked before invoking uniqueness. Denote their affine transformation formulas by $G_1,G_2$ on the triple $(\tau,\mu,\beta)$ and put $K(\tau,\mu,\beta)=(-\bar\tau,1-\bar\mu,-\bar\beta)$. Direct rational substitution gives
\[
 KG_1K=G_2^{-1}G_1^{-1}G_2,\qquad KG_2K=G_2^{-1}.
\]
For example $G_2^{-1}(\tau,\mu,\beta)=(-1/\tau,(1-\mu)/\tau,\beta+3-6(1-\mu)^2/\tau)$ and
$G_1^{-1}(\tau,\mu,\beta)=(1/(1-\tau),1-\mu/(1-\tau),\beta-2+6\mu^2/(1-\tau))$. These formulas verify every constant in the two equalities. The reflected $\mu$ is cusp bounded and has the same affine laws, so their difference is a global section of $\mathcal O_B(-p_0)$ and vanishes. After this equality, the difference between the reflected and original $\beta^{\rm part}$ is a holomorphic invariant base function, hence constant; its cusp value is zero under the stated normalization. Therefore
\[
 \mu(\iota z)=1-\bar\mu(z),\qquad
 \beta^{\rm part}(\iota z)=-\overline{\beta^{\rm part}(z)}.
\]
The original $\tau$ identity follows from its $j$ equation and its specified cusp branch. In particular $D_{-\bar c}(\iota z)=D_c(z)$, so this comparison stays in the exact admissible domain.

\item Let
\[
 H=\begin{pmatrix}1&0&0&0\\0&1&0&0\\-6&0&-1&0\\0&-1&0&-1\end{pmatrix}.
\]
Then $H^2=I$, $HA_2H=A_2^{-1}$, $HA_1H=A_2^{-1}A_1^{-1}A_2$, and
\[
 \Pi_{-\bar c}(\iota z)H=-\overline{\Pi_c(z)}.
\]
Thus $(z,\zeta)\mapsto(\iota z,-\bar\zeta)$ is a well-defined antiholomorphic torus map. It also descends under the middle deck action. Indeed covariance and the displayed lattice conjugacies imply $R_{\phi(g)}(\iota z)=\overline{R_g(z)}$ by multiplying the full real-invertible period matrices. This proves exact compatibility, not merely an equality of fibre lattices.

\item At the first finite centre compose this map with the target deck identification $g_2$ to return the centre $g_2^{-1}z_1$ to $z_1$. Its lattice map becomes $A_2H$ and it conjugates the clockwise generator to the inverse generator. At the second centre its lattice map is $H$ with the same reversal. Consequently the target affine vectors are exactly
\[
 w'_1=-A_2H\varepsilon=-\varepsilon+d,\qquad
 w'_2=-H(-\varepsilon')=\varepsilon'-3d.
\]
The minus signs arise because the inverse of a generator with invariant translation $w'_j/m_j$ has translation $-w'_j/m_j$. This calculation verifies the full affine actions on the root covers. For the overlaps choose the root coordinates of the real local base coordinates $t$ and $t-1$ compatible with these local reflections. This choice has an explicit comparison with any original root coordinate: their ratio is a holomorphic unit, so changing it changes the collar by $(2\pi i)^{-1}\log h\,\Pi w_j$, a translation extending on the entire root family. Its inverse gives a holomorphic isomorphism of the two glued descriptions. Thus the extension property and final classification transfer to the original coordinates. Then $\log s'_j=\overline{\log s_j}$ on paired branches. Under $\zeta\mapsto-\bar\zeta$, conjugation of $1/(2\pi i)$ supplies the compensating sign, and the identity of the two logarithmic collar translations follows exactly. Two roots of the same fixed real coordinate differ by an $m_j$th root of unity; this changes the comparison only by a constant period-coordinate translation, and the additive correction of \ref{fam:src:A.3.1} glues those constants. The preceding holomorphic-unit comparison handles a different original local base coordinate. Its inverse has the same construction.

\item At the cusp choose $s=(2\pi i)^{-1}\log t_c$, so $s\mapsto-\bar s$. The full holomorphic unit matrix satisfies
\[
 C_{-\bar c}(\bar t_c)=-\overline{C_c(t_c)}+
       \begin{pmatrix}6&0\\0&1\end{pmatrix}.
\]
The integral extra matrix has trivial exponential. Hence $(x,t_c)\mapsto(\bar x,\bar t_c)$ intertwines the full cusp deck groups, with the same integral index $\lambda$. The toric fan is defined over $\mathbb Z$, so this map extends antiholomorphically on every toric chart and stratum. It preserves the untwisted cusp collar. We have therefore constructed an antiholomorphic map to the actual target datum with parameter $-\bar c$, $r'_1=-1,k'_1=1,r'_2=1,k'_2=-3,l'_0=0$, whose $p'$ is $+1$.

\item Apply \ref{fam:src:B.2.1} to normalize this target back to the canonical $p=-1$ datum. Take $a=-1$ and the odd shift $n=1$. This gives a holomorphic map to $X_{1-\bar c}$. Every other antiholomorphic map, composed with the inverse of this one, is a holomorphic comparison between canonical members. \ref{fam:src:A.5.1} gives exactly the stated congruence. Exponentiating it gives $w'=-\bar w$.

\item It remains to verify the group extension rather than infer its square from $H^2$. For a real member, choose any resulting antiholomorphic self-map $J$. Its derivative on ordinary fibre coordinates is $+\text{complex conjugation}$: the initial reflection has derivative $-\text{complex conjugation}$, while the normalization just used has linear part $-I_2$. Hence it conjugates translation by $b\in\mathbb C/\mathbb Z$ to translation by $\bar b$. Its square is a holomorphic translation $d$. Commuting $J$ with $J^2$ shows $d=\bar d$ in $\mathbb C/\mathbb Z$. This equality forces the imaginary part of a representative of $d$ to vanish, so take $d\in\mathbb R$. Replacing $J$ by the composition of translation by $-d/2$ with $J$ changes its square to $d-d/2-d/2=0$. We have an antiholomorphic involution. Under $b\mapsto e^{2\pi ib}$, its conjugation action is $z\mapsto1/\bar z$. This proves the split group and specifies its antiholomorphic coset.
\end{enumerate}
\end{proof}



\section{Classification and deformation questions}\label{push:section:E}
\subsection{Higher logarithmic transformations}\label{push:section:E.1}
We distinguish Viaclovsky's indices $0,1$ from the indices $1,2$ of
the question.  Put $m=3k$, $n=4l$, and assume $\gcd(m,n)=1$.
The lattice coordinates in this subsection are integral homology
coordinates, and matrices act on columns.

\begin{proposition}[the integral dictionary]\label{push:src:E.1.1}
\StatementRecord{push:src:E.1.1}{proved}{Historical events: Claude Fable 5.2 (Anthropic), effort high, 16 November 2026: carried as background without fault found; used and rechecked in this run. Claude Fable 5.2 (Anthropic), effort high, 17 November 2026: carried as background without fault found. GPT-6 Astra (OpenAI), ChatGPT Pro mode, effort Pro, S6\_09\_PAPERS, 17 November 2026: targeted proof check. Also checked in this writing. Pre-continuation writing events: GPT-6 Astra (OpenAI), ChatGPT Pro mode, effort Pro, S6\_10\_PAPERS, 18 November 2026: targeted proof check. Continuation metadata events: GPT-6 Astra (OpenAI), ChatGPT Pro mode, effort Pro, S6\_10\_PAPERS continuation, 19 November 2026: used and rechecked in this run; status-vocabulary reconciliation; metadata reconciliation only; no new mathematical check. \emph{Pre-continuation recorded proof scope.} Re-derived each generator/base/period convention morphism, including the orientation-reversing base matrix, inverse deck generator, changed cusp nilpotent coefficient and translation period. Exact Viaclovsky source equations remain separately tracked.}{the hypotheses, notation and definitions stated here; \cite[Section 2, equations (2.1), (2.2), (2.7), (2.11)--(2.12)]{ViaclovskyLog2026}}

Let $(\delta,a,b,d)$ be the marking in
\cite[Section 2, equations (2.1), (2.2), (2.7), (2.11)--(2.12)]{ViaclovskyLog2026},
and put $u_0=d+2a+2b$, $u_1=d+3a+3b$.
The homomorphism from the marking of the question to this marking is
\[
 P_V=\begin{pmatrix}
 -2&1&0&1\\-6&0&-1&0\\0&-1&0&0\\-1&0&0&0
 \end{pmatrix},\qquad \det P_V=-1.
\]
It intertwines the two finite monodromies and satisfies
\[
 P_V\widehat\delta=\delta,\qquad
 P_V\varepsilon=-u_0,\qquad
 P_V\varepsilon'=\delta-u_1.
\]
With the counterclockwise deck generators of that paper, its data
$m_0=3k$, $m_1=4l$ and
$a_i=[v_i/m_i]$ with $v_i=\rho_i\delta+N_i u_i$ correspond to
\[
 w_1=N_0\varepsilon-\rho_0\widehat\delta,
 \qquad
 w_2=N_1\varepsilon'-(\rho_1+N_1)\widehat\delta.
\]
For the section-preserving cusp comparison used there, $l_0=0$ and
$p=-(nN_0+mN_1)$.  Thus its sphere condition is
$b:=nN_0+mN_1=\pm1$.
The brackets in $a_i$ mean the torsion point in the smooth
root torus, using its section-preserving origin and the
transported period vector.  Replacing $v_i$ by an integral
multiple $m_i$ times a period does not change that torsion
point, although its chosen lift fixes the collar bookkeeping.
\end{proposition}


\begin{proof}
\begin{enumerate}[label=\arabic*.,leftmargin=*]
\item In Engel's marking the four columns of the change to Viaclovsky's
basis are
\[
 Q_E=\begin{pmatrix}-2&1&0&-2\\-1&0&0&0\\3&-1&1&0\\0&0&0&1\end{pmatrix}.
\]
Direct integral multiplication gives $P_V=Q_E^{-1}P_E$ and
$V_0P_V=P_VA_1$, $V_1P_V=P_VA_2$, where $V_0,V_1$ are the
matrices in that paper's Lemma 2.1.  Multiplication by the three
specified invariant vectors gives the displayed identities.
Inverting a clockwise affine generator replaces its invariant
translation $w_j/m_j$ by $-w_j/m_j$.  Consequently the actual
counterclockwise translation vector is $-P_Vw_j$, which gives the
two formulas, including the otherwise easily missed $N_1$ in the
second coefficient.  Viaclovsky's reference section extends through
the smooth locus of the cusp; its comparison therefore has zero
class in the quotient lattice of toric directions.  This is precisely
$l_0=0$.  Substitution into the definition of $p$ proves its formula.
These matrix identities are independently executed in
\texttt{code/e\_push\_exact.py}, with their output attached.
\end{enumerate}
\end{proof}

For a general cusp collar $\sigma$, write $c(\sigma)$ for
its residual translation defect in
$\overline\Lambda=\Lambda/(\mathbb Z\widehat w+
\mathbb Z\widehat\delta)$, using the convention
$\sigma(g_0z)-R_0\sigma(z)=\Pi_c(g_0z)\ell(\sigma)$
and $c(\sigma)=[\ell(\sigma)]$.
Set $l_0=\gamma(c(\sigma))$.  In the question's notation
$w_j=r_j\varepsilon_j+s_j\widehat\delta$, the integer is
$p=mn l_0-nr_1-mr_2$.
The section-preserving source data of \ref{push:src:E.1.1} have $c(\sigma)=0$.

\begin{theorem}[classification in Viaclovsky's coordinates]\label{push:src:E.1.2}
\StatementRecord{push:src:E.1.2}{proved}{Historical events: Claude Fable 5.2 (Anthropic), effort high, 16 November 2026: carried as background without fault found; used and rechecked in this run. Claude Fable 5.2 (Anthropic), effort high, 17 November 2026: carried as background without fault found. GPT-6 Astra (OpenAI), ChatGPT Pro mode, effort Pro, S6\_09\_PAPERS, 17 November 2026: targeted proof check. Also checked in this writing. Pre-continuation writing events: GPT-6 Astra (OpenAI), ChatGPT Pro mode, effort Pro, S6\_10\_PAPERS, 18 November 2026: targeted proof check. Continuation metadata events: GPT-6 Astra (OpenAI), ChatGPT Pro mode, effort Pro, S6\_10\_PAPERS continuation, 19 November 2026: used and rechecked in this run; status-vocabulary reconciliation; metadata reconciliation only; no new mathematical check. \emph{Pre-continuation recorded proof scope.} Re-derived the full higher-multiplicity affine classification: finite norm images, local affine intertwiners, free adjustment of the cusp defect, global difference-sheaf gluing, necessity through the base-factor map and integral centralizer, and the CRT period 2klZ. All sufficiency maps and local fan extensions were checked.}{the hypotheses, notation and definitions stated here; \cref{push:src:E.1.1}}

Use a local logarithm $c$ of the transported linearization,
$u=e^{-2\pi i c}$, and the same reference section and cusp comparison
on both sides.  Two sphere members with data
$(m,\rho_0,N_0;n,\rho_1,N_1;c)$ and primed data are biholomorphic
if and only if their multiplicities agree and there are
$\alpha\in\{1,-1\}$ and $h=c'-c\in\mathbb Z$ such that
\[
 b'=\alpha b,\qquad
 hN'_0\equiv\rho'_0-\alpha\rho_0\pmod{k},\qquad
 hN'_1\equiv\rho'_1-\alpha\rho_1\pmod{2l}.
\]
Every such biholomorphism covers the identity of $B$.
For fixed discrete data the period identification subgroup is
$2kl\mathbb Z$.  At fixed $u$ all sphere translation data give
exactly $2kl$ classes.  A positive $u$-circuit acts on them by a
single cycle.
\end{theorem}


\begin{proof}
\begin{enumerate}[label=\arabic*.,leftmargin=*]
\item We first give the complete sufficiency argument in the stated problem's
clockwise coordinates, writing $w_j=r_j\varepsilon_j+s_j\widehat\delta$.
Set $S_h=I+hE$ and $d_j=S_hw'_j-\alpha w_j$.
The repeated norm maps have images
\[
 N_m\Lambda=m\mathbb Z\varepsilon+k\mathbb Z\widehat\delta,
 \qquad
 N_n\Lambda=n\mathbb Z\varepsilon'+2l\mathbb Z\widehat\delta.
\]
This follows directly from
\[
 N_3v=3\gamma(v)\varepsilon+(2v_u+v_w+3v_\delta)\widehat\delta,
 \quad
 N_4v=4\gamma(v)\varepsilon'+(2v_u+2v_w+4v_\delta)\widehat\delta.
\]
Suppose $p'=\alpha p$ and
$hr'_1\equiv\alpha s_1-s'_1\pmod k$,
$hr'_2\equiv\alpha s_2-s'_2\pmod{2l}$.
Reduction of $p'=\alpha p$ modulo $m,n$ gives
$r'_j\equiv\alpha r_j\pmod{m_j}$.
We may choose integral solutions of $N_{m_j}\ell_j=-d_j$ explicitly:
\[
 \begin{split}
 \ell_1&=\left(\frac{\alpha r_1-r'_1}{m},0,
 -\frac{s'_1+hr'_1-\alpha s_1}{k},0\right)^t,\\
 \ell_2&=\left(\frac{\alpha r_2-r'_2}{n},
 -\frac{s'_2+hr'_2-\alpha s_2}{2l},0,0\right)^t.
 \end{split}
\]
Put $y_j=d_j/m_j+\ell_j$; these vectors belong to the rational
kernel of $N_{m_j}$ and have first coordinate zero.
The following rational vectors solve $(I-A_j)b_j=y_j$:
\[
 \begin{split}
 b_1&=(0,(2y_{1u}+y_{1w})/3,(y_{1w}-y_{1u})/3,0)^t,\\
 b_2&=(0,(y_{2u}-y_{2w})/2,(y_{2u}+y_{2w})/2,0)^t.
 \end{split}
\]
The last coordinates of the equations follow from the displayed
norm-kernel equations.  The maps
\[
 (s,\zeta)\longmapsto(s,\alpha\zeta+\Pi_c(s)b_j)
\]
are holomorphic maps of the normalized torus families over the
degree-$m_j$ root discs.  Their inverses are holomorphic, and period
covariance and $(I-A_j)b_j=d_j/m_j+\ell_j$ give exact commutation
with the prescribed affine deck generators, modulo the integral
period $\ell_j$.  They therefore descend to biholomorphisms of the
finite fillings, including the full multiple fibres.

\item On the punctured root collar the corresponding middle translation is
\[
 A_j=\Pi_cb_j+\frac{\log s_j}{2\pi i}\Pi_cd_j.
\]
Its integral defect is $\ell_j$.  The middle cocycle relation fixes
$\ell_0=-A_0(\ell_1+A_1\ell_2)$.  In the quotient
$\overline\Lambda=\Lambda/(\mathbb Z\widehat w+\mathbb Z\widehat\delta)$,
the first-coordinate condition is exactly $p'=\alpha p$.
Adding an integral multiple of $(0,1,-2,0)^t\in\ker N_m$ to
$\ell_1$ changes the other quotient coordinate by a generator,
without changing this first coordinate.  We can consequently arrange
\[
 [\ell_0]=-\alpha c(\sigma)+c(\sigma')\quad\hbox{in }\overline\Lambda.
\]
An integral toric vector then supplies $q\in\mathbb Z^2$ such that
the cusp reference $sq+\alpha\sigma-\sigma'$ has precisely defect
$\ell_0$.  Here $s=(2\pi i)^{-1}\log t_c$ is the cusp logarithm;
the maps of cusp toric covers are
\[
 (x,t_c)\longmapsto(t_c^q a(t_c)x^\alpha,t_c),
 \qquad a(0)\in(\mathbb C^*)^2.
\]
Integer translations of the cocharacter lattice, its inversion,
and multiplication by a holomorphic unit preserve the specified
hexagonal fan and all its strata.  These maps and their inverses
extend holomorphically across the three double curves and both
triple points.

\item We justify the global analytic correction.  Differences of allowed
finite extensions are the equivariant holomorphic vertical vectors
on the complete root discs.  The extra subgroup of order $k$ at
the first point acts trivially on the vertical vector coefficients
and rotates $s$; its invariant holomorphic coefficients are exactly
those in $s^k$.  The remaining order-three equivariance is the
original one.  The argument at the second point uses $s^l$ and
order four.  Thus the coarse difference sheaf is literally the same
sheaf $\mathcal V_B$ as for multiplicities $(3,4)$, not an assumed
analogue.  Its first-coordinate quotient is
$\mathcal O_B(-p_0)$: the modular frame has vanishing orders $2k,l$
on the two enlarged roots and a simple cusp pole.  Its fixed
second-coordinate subline is $\mathcal O_B$.  Finite averaging
provides local lifts, hence
\[
 0\longrightarrow\mathcal O_B\longrightarrow\mathcal V_B
 \longrightarrow\mathcal O_B(-p_0)\longrightarrow0.
\]
Both end terms have vanishing $H^1$ on $\mathbb P^1$, so
$H^1(B,\mathcal V_B)=0$.  The permitted local translations form a
torsor under this sheaf: on each overlap their defects agree, so
their difference is a single-valued section of the indicated
difference sheaf.  Vanishing of its \v Cech class gives a global
middle translation.  Its overlap equalities are exactly the two
finite formulas and the cusp formula just displayed.  Consequently
the local biholomorphisms glue, with their inverses, over $B$.

\item For necessity the base is intrinsic even when its anticanonical
pencil changes.  If $F:X\to X'$ is a biholomorphism, the pullback
of the Fubini--Study form by $f'F$ is exact because
$H^2(S^6,\mathbb R)=0$.  Its restriction to a regular complex torus,
wedged with a flat K\"ahler form, has zero integral.  Positivity
forces $f'F$ to be constant on that torus.  The proper analytic
image of $(f,f'F)$ is therefore a curve of degree one over $B$;
normalization and the inverse argument give a base automorphism.
Its three monodromy orders $3,4,\infty$ are distinct, so it fixes
all three marked points and is the identity.  Multiplicities agree.
On ordinary fibres the lift is affine.  The simultaneous rational
centralizer of $A_1,A_2$ is $\mathbb Q I+\mathbb Q E$, as follows
by solving their entrywise commuting equations.  Comparison of the
constant last two columns of $\Pi_c$ and unimodularity then gives
fibre-linear part $\alpha I$, integer difference $h$, and lattice
map $\alpha S_h^{-1}$.  The normalized root base change is
functorial, so the lifted map commutes with the actual deck
generator.  Summing its defect around that generator gives the
two norm conditions above.  A cusp map has an integral toric
winding: at a generic point of a toric component, its ratio of
characters is $t_c^q$ times a holomorphic unit.  Its class in
$\overline\Lambda$ vanishes.  The middle cocycle relation then
gives $p'=\alpha p$.  This proves all the necessary conditions.

\item Substitute \ref{push:src:E.1.1}.  In the second congruence the extra term
$N'_1-\alpha N_1$ is divisible by $4l$, and hence drops out
modulo $2l$.  We obtain exactly the statement's two congruences.
For fixed data $b\ne0$ excludes $\alpha=-1$.
The invertibility of $N_0$ modulo $k$, of $N_1$ modulo $2l$, and
$\gcd(k,2l)=1$ give $h\in2kl\mathbb Z$.
For varying data the Chinese remainder theorem gives exactly one
shift class modulo $2kl$, and all residue classes of
$\rho_0\bmod k$, $\rho_1\bmod2l$ occur without changing $b$ or
freeness.  Finally a positive $u$-circuit sends $c$ to $c-1$.
Its translation on $\mathbb Z/(2kl)$ is a generator.
\end{enumerate}
\end{proof}

The preceding theorem answers Question 3 of
\cite[Section 1.1]{ViaclovskyLog2026}.  The original return asserted the
same arithmetic classification in \cref{fam:src:B.3.1}; the contribution here is
the explicit dictionary and the complete enlarged-root and collar
argument.

\begin{corollary}[holomorphic dependence of the comparisons]\label{push:src:E.1.3}
\StatementRecord{push:src:E.1.3}{proved}{Historical events: Claude Fable 5.2 (Anthropic), effort high, 16 November 2026: carried as background without fault found. GPT-6 Astra (OpenAI), ChatGPT Pro mode, effort Pro, S6\_09\_PAPERS, 17 November 2026: targeted proof check. Also checked in this writing. Pre-continuation writing events: GPT-6 Astra (OpenAI), ChatGPT Pro mode, effort Pro, S6\_10\_PAPERS, 18 November 2026: targeted proof check. Continuation metadata events: GPT-6 Astra (OpenAI), ChatGPT Pro mode, effort Pro, S6\_10\_PAPERS continuation, 19 November 2026: used and rechecked in this run; status-vocabulary reconciliation; metadata reconciliation only; no new mathematical check. \emph{Pre-continuation recorded proof scope.} Re-derived parameter-holomorphic gluing via the product difference bundle and its H1 vanishing, the precise sign W=(-1)\textasciicircum{}s2 w, and the connected automorphism action on smooth monodromy. The automorphism input is credited separately.}{the hypotheses, notation and definitions stated here; \cref{push:src:E.1.2}; \cref{fam:src:B.4.1}}

For fixed discrete data in \ref{push:src:E.1.2} and a fixed admissible shift $h$,
the constructed comparisons are holomorphic isomorphisms of families
over a sufficiently small parameter disc.  In particular, for sphere
data with multiplicities $(3,4)$ the identification with the standard
family uses the parameter
\[
 W=e^{-\pi i(c+s_2)}=(-1)^{s_2}w.
\]
The smooth monodromy conjugacy class around the positive parameter
circle is unchanged by this identification of twist data.
\end{corollary}


\begin{proof}
\begin{enumerate}[label=\arabic*.,leftmargin=*]
\item In \ref{push:src:E.1.2} the rational vectors $b_j$, integral defects $\ell_j$,
and toric exponents are independent of $c$.  All period matrices,
root comparisons, and cusp unit factors are holomorphic in $c$.
On a small Stein parameter disc $D$, the correction sheaf is
$\mathcal V_B\boxtimes\mathcal O_D$.  Its first cohomology on
$B\times D$ is zero: project to $D$, use the same two-chart
complex and $H^1(B,\mathcal V_B)=0$, and then Stein acyclicity.
Thus the affine correction can be chosen holomorphically in $c$.
For $(k,l)=(1,1)$ the only nonvacuous finite congruence is
$h\equiv s_2-s_2'\pmod2$, because the admissible $r_2'$ is odd.
Taking $s_2'=0$ gives the displayed parameter identification.
Two choices of these family comparisons differ fibrewise by an
automorphism of the standard member.  The full automorphism
calculation re-derived in the credited appendix at \ref{fam:src:B.4.1} identifies that group with $\mathbb C^*$;
it is connected.  On overlapping parameter arcs the transitions
are therefore isotopic to the identity.  Composing around a circle
changes the resulting diffeomorphism only by such an isotopy and
by conjugation at the chosen base fibre.  Multiplication of the
parameter by $(-1)^{s_2}$ preserves the positive orientation.
\end{enumerate}
\end{proof}


\section{The monodromy group and its period domain}\label{fam:section:F}

We use the clockwise root deck generators and the lattice basis specified
on the first page.  In this section a group acts properly discontinuously when, for every
compact subset $K$ of its space, only finitely many group elements $g$ satisfy
$gK\cap K\ne\varnothing$.  This definition permits finite stabilizers.  We
write $e=\widehat\delta$, $E=e\otimes\gamma$, and
\[
 J_2=\begin{pmatrix}0&1\\-1&0\end{pmatrix},\qquad
 \Omega=\begin{pmatrix}
 0&0&0&6\\0&0&1&0\\0&-1&0&0\\-6&0&0&0
 \end{pmatrix}.
\]
Thus the alternating form is
$\psi=u\wedge w+6\gamma\wedge\delta$, an integral alternating form of
type $(1,6)$ on $\Lambda$.  We use $\Omega$ for its matrix and retain $J$ for
a real complex structure.  The domain $\mathcal D$ consists of the real
endomorphisms $J$ of $\Lambda_{\mathbb R}$ satisfying
$J^2=-I$, $J^{\mathsf t}\Omega J=\Omega$, and
$\operatorname{sign}(\Omega J)=(2,2)$.  It is the stated
$\operatorname{Sp}_4(\mathbb R)/U(1,1)$, with its complex structure as an
open subset of the Lagrangian Grassmannian.  An integral symplectic matrix
acts by $J\mapsto gJg^{-1}$.

\subsection{An exact arithmetic description}\label{fam:section:F.1}

\begin{theorem}\label{fam:src:F.1.1}\label{fam:st:F.1.1}\label{cdp:st:D.1.14}
\StatementRecord{fam:src:F.1.1}{proved}{Historical events: Claude Fable 5.2 (Anthropic), effort high, 12 November 2026: targeted proof check. Claude Fable 5.2 (Anthropic), effort high, 16 November 2026: carried as background without fault found; used and rechecked in this run. Claude Fable 5.2 (Anthropic), effort high, 17 November 2026: carried as background without fault found. Claude Opus 5.5 (Anthropic), effort max, sweep 1, scope s6, 18 November 2026: used and rechecked in this run. GPT-6 Astra (OpenAI), ChatGPT Pro mode, effort Pro, S6\_09\_PAPERS, 17 November 2026: targeted proof check. Also checked in this writing. Pre-continuation writing events: GPT-6 Astra (OpenAI), ChatGPT Pro mode, effort Pro, S6\_10\_PAPERS, 18 November 2026: targeted proof check. Continuation metadata events: GPT-6 Astra (OpenAI), ChatGPT Pro mode, effort Pro, S6\_10\_PAPERS continuation, 19 November 2026: used and rechecked in this run; status-vocabulary reconciliation; metadata reconciliation only; no new mathematical check. \emph{Pre-continuation recorded proof scope.} Read full proof and re-derived block multiplication, character modulo12, Heisenberg kernel generated by U,V,Z, center and Zariski-closure argument. New exact computer replay forthcoming.}{the hypotheses, notation and definitions stated here; \cref{fam:eq:F-law}; \cref{fam:eq:F-SLpresentation}}

Let $P_{e,\mathbb Z}$ be the integral symplectic group fixing $e$ as a
vector, and let $P_{\mathbb Z}$ be the integral symplectic group preserving
$\mathbb Ze\subset\ker\gamma\subset\Lambda$.  Every element of
$P_{e,\mathbb Z}$ has a unique expression
\begin{equation}\label{fam:eq:F-g}
 g(B,x,t)=
 \begin{pmatrix}
 1&0&0\\6x&B&0\\t&-x^{\mathsf t}J_2B&1
 \end{pmatrix},\qquad
 B\in\operatorname{SL}_2(\mathbb Z),\quad x\in\mathbb Z^2,\quad t\in\mathbb Z,
\end{equation}
where the block sizes are $1,2,1$.  Put
\[
 S=\begin{pmatrix}0&-1\\1&0\end{pmatrix},\quad
 T=\begin{pmatrix}1&1\\0&1\end{pmatrix},\quad
 a=\begin{pmatrix}0&1\\-1&-1\end{pmatrix}.
\]
There is a homomorphism
$\alpha:\operatorname{SL}_2(\mathbb Z)\to\mathbb Z/12$ determined by
$\alpha(S)=9$ and $\alpha(T)=1$; in particular $\alpha(a)=4$.
With $Q(x)=x_1+x_2+x_1x_2$, the actual monodromy group is
\begin{equation}\label{fam:eq:F-Gamma}
 \Gamma=\left\{g(B,x,t):t+6Q(x)-\alpha(B)=0\ \text{in }\mathbb Z/12\right\}.
\end{equation}
Consequently $[P_{e,\mathbb Z}:\Gamma]=12$ and
$P_{\mathbb Z}=\{\pm I\}\times P_{e,\mathbb Z}$, so
$[P_{\mathbb Z}:\Gamma]=24$.  Its image in the Levi quotient is exactly
$\operatorname{SL}_2(\mathbb Z)$, with scalar on $e$ equal to $1$.
\end{theorem}



\begin{proof}
\begin{enumerate}[label=\arabic*.,leftmargin=*]
\item The matrices in the stated problem, in the fixed basis, are
\[
 A_1=\begin{pmatrix}1&0&0&0\\6&0&1&0\\-6&-1&-1&0\\-2&1&0&1\end{pmatrix},
 \qquad
 A_2=\begin{pmatrix}1&0&0&0\\0&0&-1&0\\-6&1&0&0\\3&0&1&1\end{pmatrix}.
\]
Multiplication gives $A_1^3=A_2^4=I$ and
$A_j^{\mathsf t}\Omega A_j=\Omega$.  In particular
$A_1=g(a,(1,-1)^{\mathsf t},-2)$ and
$A_2=g(S,(0,-1)^{\mathsf t},3)$.  Their images $a,S$ generate
$\operatorname{SL}_2(\mathbb Z)$ because $Sa=T$ and $S,T$ perform the
Euclidean algorithm on primitive integral columns.

\item A matrix fixing $e$ and preserving $\Omega$ preserves
$e^{\perp}=\ker\gamma$.  Its first row is $(1,0,0,0)$, its last column is
$e$, its middle block $B$ is symplectic on
$e^{\perp}/\mathbb Re$, and its remaining middle column, say $v$, satisfies
$6r=-v^{\mathsf t}J_2B$.  Integrality of $r$ and unimodularity of $B$
give $v=6x$ with $x\in\mathbb Z^2$.  There is no further condition on
the lower left entry $t$.  This proves \eqref{fam:eq:F-g}.  Direct multiplication
gives
\begin{equation}\label{fam:eq:F-law}
 g(B,x,t)g(C,y,s)
 =g(BC,x+By,t+s-6x^{\mathsf t}J_2By).
\end{equation}

\item For completeness, the presentation of the Levi group needed here follows
from an elementary action on the real projective line.  The images of $a$
and $S$ in $\operatorname{PSL}_2(\mathbb Z)$ have orders $3$ and $2$.
On positive real numbers, $a(z)=-1/(z+1)$ has image $(-1,0)$ and
$a^2(z)=-(z+1)/z$ has image $(-\infty,-1)$; $S$ carries negative real
numbers to positive real numbers.  A cyclically reduced alternating word
starting with a nonidentity power of $a$ and ending with $S$ therefore
carries $(-\infty,0)$ into one of these two proper subintervals.  It cannot
be the identity.  Every reduced word is conjugate to such a word or to a
nonidentity element of one factor.  Together with the Euclidean algorithm,
this proves the presentation
$\operatorname{PSL}_2(\mathbb Z)=\langle\bar a,\bar S\mid
\bar a^3=\bar S^2=1\rangle$.  It follows that
\begin{equation}\label{fam:eq:F-SLpresentation}
 \operatorname{SL}_2(\mathbb Z)=
 \langle a,S\mid a^3=S^4=1,\ [a,S^2]=1\rangle.
\end{equation}
Indeed the group on the right modulo its central element $S^2$ is the
just-proved free product, and $S^2$ maps to the nonidentity matrix $-I$.
Its kernel in the displayed matrix group is consequently trivial.
The values $4,9$ respect all these relations, proving the assertion about
$\alpha$; $T=Sa$ has value $1$.

\item The quadratic function $Q$ modulo $2$ is zero at $(0,0)$ and one at all
three other points of $\mathbb F_2^2$.  Hence $Q(Bx)=Q(x)\pmod2$ for every
$B\in\operatorname{SL}_2(\mathbb Z)$.  Also
\[
 Q(x+y)-Q(x)-Q(y)\equiv x^{\mathsf t}J_2y\pmod2.
\]
Equation \eqref{fam:eq:F-law} now shows that
$h(g(B,x,t))=t+6Q(x)\pmod{12}$ is a homomorphism.
The homomorphism $\chi=h-\alpha\circ B$ is surjective, since
$\chi(g(I,0,1))=1$.  Its values on $A_1,A_2$ are zero.

\item We verify equality with its kernel; containment alone would not determine
the group.  Let $N_{\mathbb Z}$ be the kernel of the map to the middle
block.  The kernel of
$\langle g_1,g_2\mid g_1^3=g_2^4=1\rangle\to
\operatorname{SL}_2(\mathbb Z)$ is normally generated by
$[g_1,g_2^2]$, by \eqref{fam:eq:F-SLpresentation}.  Its image in $\Gamma$ is
therefore normally generated by
\[
 K=[A_1,A_2^2]=g(I,(0,-1)^{\mathsf t},6).
\]
Put
\[
 U=A_2KA_2^{-1}=g(I,(1,0)^{\mathsf t},-6),\quad
 V=K^{-1}=g(I,(0,1)^{\mathsf t},-6),\quad Z=I+12E.
\]
Equation \eqref{fam:eq:F-law} gives $[U,V]=Z^{-1}$.  Thus $U,V,Z$ generate
\begin{equation}\label{fam:eq:F-H}
 H=\{g(I,x,t):t\in6\mathbb Z,\quad t/6\equiv Q(x)\pmod2\}.
\end{equation}
More explicitly,
$U^rV^sZ^k=g(I,(r,s)^{\mathsf t},-6(r+s+rs)+12k)$, so these expressions
give every member of the set.  Conjugation by $g(B,x,t)$ sends
$g(I,y,s)$ to $g(I,By,s-12x^{\mathsf t}J_2By)$; the invariance of $Q$
modulo $2$ proves that $H$ is normal in $P_{e,\mathbb Z}$.  The normal
closure of $K$ is contained in $H$ and contains $U,V,Z$, so it is exactly
$H$.  We have proved $\Gamma\cap N_{\mathbb Z}=H$.  Both $\Gamma$ and
$\ker\chi$ surject onto the same Levi group and have intersection $H$
with $N_{\mathbb Z}$.  Given $g\in\ker\chi$, choose $\gamma\in\Gamma$
with the same middle block.  Then $g\gamma^{-1}\in H$, and $g\in\Gamma$.
This proves \eqref{fam:eq:F-Gamma} and the index $12$.

\item Finally an integral automorphism preserving the primitive line
$\mathbb Ze$ sends $e$ to $e$ or $-e$.  Multiplication by the central
matrix $-I$ converts the latter case to the former.  The intersection of
$\{\pm I\}$ with $P_{e,\mathbb Z}$ is trivial.  This proves the direct
product and the index $24$.
\end{enumerate}
\end{proof}

\begin{proposition}\label{fam:src:F.1.2}\label{fam:st:F.1.2}
\StatementRecord{fam:src:F.1.2}{proved}{Historical events: Claude Fable 5.2 (Anthropic), effort high, 12 November 2026: targeted proof check. Claude Fable 5.2 (Anthropic), effort high, 16 November 2026: used and rechecked in this run. GPT-6 Astra (OpenAI), ChatGPT Pro mode, effort Pro, S6\_09\_PAPERS, 17 November 2026: targeted proof check. Also checked in this writing. Pre-continuation writing events: GPT-6 Astra (OpenAI), ChatGPT Pro mode, effort Pro, S6\_10\_PAPERS, 18 November 2026: targeted proof check. Continuation metadata events: GPT-6 Astra (OpenAI), ChatGPT Pro mode, effort Pro, S6\_10\_PAPERS continuation, 19 November 2026: used and rechecked in this run; status-vocabulary reconciliation; metadata reconciliation only; no new mathematical check. \emph{Pre-continuation recorded proof scope.} Read full proof and re-derived block multiplication, character modulo12, Heisenberg kernel generated by U,V,Z, center and Zariski-closure argument. New exact computer replay forthcoming.}{the hypotheses, notation and definitions stated here; \cref{fam:eq:F-Gamma}; \cref{fam:eq:F-law}}

The group in \eqref{fam:eq:F-Gamma} has center
$\langle Z\rangle\cong\mathbb Z$ and is the following explicit central
extension of $\mathbb Z^2\rtimes\operatorname{SL}_2(\mathbb Z)$ by
$\mathbb Z$.  Choose the representative $a(B)\in\{0,\ldots,11\}$ of
$\alpha(B)$.  Its underlying set consists of triples $(B,x,k)$, with
product
\[
 (B,x,k)(C,y,l)=(BC,x+By,k+l+\kappa((B,x),(C,y))),
\]
where
\begin{align*}
12\kappa((B,x),(C,y))={}&a(B)+a(C)-a(BC)-6Q(x)-6Q(y)\\
 &-6x^{\mathsf t}J_2By+6Q(x+By).
\end{align*}
Its real and rational Zariski closures are respectively the full real
and rational symplectic stabilizers of the vector $e$.  They have
dimension $6$, with three-dimensional Heisenberg unipotent radical and
Levi quotient $\operatorname{SL}_2$.
\end{proposition}



\begin{proof}
\begin{enumerate}[label=\arabic*.,leftmargin=*]
\item The numerator defining $\kappa$ is divisible by $12$, by the two
congruences proved above.  The bijection
\[
 (B,x,k)\longmapsto g(B,x,a(B)-6Q(x)+12k)
\]
intertwines the displayed multiplication with matrix multiplication by
\eqref{fam:eq:F-law}.  This proves associativity as well as the claimed abstract
description, without leaving an unspecified extension class.  The kernel
of the homomorphism to $\mathbb Z^2\rtimes\operatorname{SL}_2(\mathbb Z)$
is $\langle Z\rangle$.  That affine group has trivial center: commuting
with all translations forces its linear part to be $I$, and commuting
with $S,T$ forces its translation part to vanish.  Thus the center of
$\Gamma$ is exactly the stated kernel.

\item The logarithms $U-I,V-I,Z-I$ are linearly independent elements of the
three-dimensional Heisenberg Lie algebra.  The square of each is zero.
A polynomial vanishing on all integral powers $I+n(U-I)$ vanishes on
the whole one-parameter line; the same holds for $V$ and $Z$.  Hence the
Zariski closure contains the full unipotent radical.  Its Levi image
contains $S,T$; the powers of $T$ and of $STS^{-1}$ similarly supply the
two opposite unipotent subgroups of $\operatorname{SL}_2$.  Their Lie
algebras and their bracket span $\mathfrak{sl}_2$.  They generate the
connected algebraic group $\operatorname{SL}_2$.  A subgroup containing
the entire unipotent kernel and surjecting onto this Levi quotient is
the entire vector stabilizer.  Its dimension is $3+3=6$.  This reasoning
uses equations over $\mathbb Q$, so proves the rational assertion as
well as the real assertion.
\end{enumerate}
\end{proof}

\subsection{Properness and the maximal open set}\label{fam:section:F.2}

Define the real-analytic function
\[
 q_e:\mathcal D\longrightarrow\mathbb R,\qquad q_e(J)=\psi(e,Je),
 \qquad\mathcal O=\{J:q_e(J)\ne0\}.
\]
The function $q_e$ is invariant under $P_{\mathbb Z}$: replacing $e$ by
$-e$ changes both of its arguments.  Write $\mathcal O_+$ and
$\mathcal O_-$ for its two sign regions.

\begin{proposition}[counterexample]\label{fam:src:F.2.1}\label{fam:st:F.2.1}
\StatementRecord{fam:src:F.2.1}{proved}{Historical events: Claude Fable 5.2 (Anthropic), effort high, 12 November 2026: targeted proof check. Claude Fable 5.2 (Anthropic), effort high, 16 November 2026: used and rechecked in this run. GPT-6 Astra (OpenAI), ChatGPT Pro mode, effort Pro, S6\_09\_PAPERS, 17 November 2026: targeted proof check. Also checked in this writing. Pre-continuation writing events: GPT-6 Astra (OpenAI), ChatGPT Pro mode, effort Pro, S6\_10\_PAPERS, 18 November 2026: targeted proof check. Continuation metadata events: GPT-6 Astra (OpenAI), ChatGPT Pro mode, effort Pro, S6\_10\_PAPERS continuation, 19 November 2026: used and rechecked in this run; status-vocabulary reconciliation; metadata reconciliation only; no new mathematical check. \emph{Pre-continuation recorded proof scope.} Re-derived infinite unipotent stabilizer; null-locus rational-density argument; compact-stabilizer properness; unique normalized chart and q\_e=-6/D; j\_Gamma/1728 left inverse, extension over finite root discs, and proper embedding argument.}{the hypotheses, notation and definitions stated here; \cref{fam:eq:F-Gamma}}

The assertion that $\Gamma$ acts properly discontinuously on all of
$\mathcal D$ is false.  The following point of $\mathcal D$ has an
explicit infinite stabilizer:
\[
 J_0=\begin{pmatrix}
 0&0&-1&0\\0&0&0&-6\\1&0&0&0\\0&1/6&0&0
 \end{pmatrix},\qquad
 R=\begin{pmatrix}
 1&0&0&0\\0&1&72&0\\0&0&1&0\\-12&0&0&1
 \end{pmatrix}\in\Gamma.
\]
Every power of $R$ fixes $J_0$, and all these powers are distinct.
\end{proposition}



\begin{proof}
\begin{enumerate}[label=\arabic*.,leftmargin=*]
\item Direct multiplication gives $J_0^2=-I$ and
$J_0^{\mathsf t}\Omega J_0=\Omega$.  Moreover
\[
 \Omega J_0=\begin{pmatrix}
 0&1&0&0\\1&0&0&0\\0&0&0&6\\0&0&6&0
 \end{pmatrix}
\]
has eigenvalues $1,-1,6,-6$, so $J_0\in\mathcal D$ and $q_e(J_0)=0$.
The matrix $R$ equals $g(T^{72},0,-12)$ and belongs to \eqref{fam:eq:F-Gamma},
because $\alpha(T^{72})=0\pmod{12}$.  Its commutator with $J_0$ is zero.
Writing $R=I+N$ gives $N\ne0$ and $N^2=0$, so $R^m=I+mN$ and the
powers are distinct.  The compact singleton $\{J_0\}$ therefore violates
the stated definition of proper discontinuity.  These identities are
checked both with SymPy and with a separate rational matrix
implementation in \texttt{code/f\_exact.py}.
\end{enumerate}
\end{proof}

\begin{theorem}\label{fam:src:F.2.2}\label{fam:st:F.2.2}
\StatementRecord{fam:src:F.2.2}{proved}{Historical events: Claude Fable 5.2 (Anthropic), effort high, 12 November 2026: targeted proof check. Claude Fable 5.2 (Anthropic), effort high, 16 November 2026: used and rechecked in this run. GPT-6 Astra (OpenAI), ChatGPT Pro mode, effort Pro, S6\_09\_PAPERS, 17 November 2026: targeted proof check. Also checked in this writing. Pre-continuation writing events: GPT-6 Astra (OpenAI), ChatGPT Pro mode, effort Pro, S6\_10\_PAPERS, 18 November 2026: targeted proof check. Continuation metadata events: GPT-6 Astra (OpenAI), ChatGPT Pro mode, effort Pro, S6\_10\_PAPERS continuation, 19 November 2026: used and rechecked in this run; status-vocabulary reconciliation; metadata reconciliation only; no new mathematical check. \emph{Pre-continuation recorded proof scope.} Re-derived infinite unipotent stabilizer; null-locus rational-density argument; compact-stabilizer properness; unique normalized chart and q\_e=-6/D; j\_Gamma/1728 left inverse, extension over finite root discs, and proper embedding argument.}{the hypotheses, notation and definitions stated here; \cref{fam:eq:F-g}}

The set $\mathcal O$ is the largest $\Gamma$-invariant open subset of
$\mathcal D$ on which the action is properly discontinuous.  It is also
the largest such open subset for $P_{\mathbb Z}$.  Each sign region is
a complex manifold in the analytic topology, and its quotient by either
discrete group is a Hausdorff complex analytic orbifold.
\end{theorem}



\begin{proof}
\begin{enumerate}[label=\arabic*.,leftmargin=*]
\item Let $P_{\mathbb R}$ be the real symplectic group preserving the real line
$\mathbb Re$.  We first prove properness on each sign region.  For a fixed
$J$, the form
$h_J(v,w)=\psi(v,Jw)+i\psi(v,w)$ is a nondegenerate Hermitian form of
signature $(1,1)$ on the complex vector space $(\Lambda_{\mathbb R},J)$,
with the convention conjugate linear in the first entry.  When
$q_e(J)\ne0$, the complex line $\mathbb Ce$ is nondegenerate and its
Hermitian orthogonal complement is a complex line of the opposite sign.
Normalize $e$ by a positive real scalar and choose a unit vector in this
complement.  Comparing these two Hermitian bases for two complex
structures with the same sign gives a real symplectic map preserving
$\mathbb Re$ and conjugating one complex structure to the other.  Thus
$P_{\mathbb R}$ is transitive on each sign region.

\item An element of the stabilizer of $J$ inside $P_{\mathbb R}$ is complex
linear and sends $e$ to $r e$ for a real nonzero $r$.  Preservation of
$q_e(J)$ gives $r^2=1$.  On the complementary complex line it acts by an
element of $U(1)$.  The stabilizer is consequently the compact group
$\{\pm1\}\times U(1)$.  Each sign region is therefore a homogeneous
space of $P_{\mathbb R}$ with compact stabilizer $L$.  Local sections of
$P_{\mathbb R}\to P_{\mathbb R}/L$, and a finite cover of a compact
subset $K$, give a compact subset $C$ of $P_{\mathbb R}$ mapping onto
$K$.  Every element carrying a point of $K$ back into $K$ belongs to
$CLC^{-1}$, which is compact.  This proves properness.  Restricting a
proper action to a closed discrete subgroup gives
the required finite-intersection condition.  Both $\Gamma$ and
$P_{\mathbb Z}$ are closed discrete matrix groups.  The two sign regions
are disjoint invariant open-and-closed subsets of $\mathcal O$, so the
same finite-intersection condition holds on their union.  Proper actions
by discrete groups have finite stabilizers and Hausdorff quotients; their
holomorphic local finite-group quotient charts supply the asserted
analytic orbifolds.

\item It remains to prove maximality, which requires more than one bad point.
The real vector stabilizer $P_{e,\mathbb R}$ acts transitively on
$\mathcal D_0=\{J:q_e(J)=0\}$.  Here is a direct comparison map.  In the
Hermitian space just described, $e$ is a nonzero null vector.  Choose
$v$ with $h_J(e,v)=1$, and replace $v$ by
$v-\tfrac12h_J(v,v)e$.  This preserves $h_J(e,v)=1$ and makes $v$ null.
The ordered complex basis $(e,v)$ has Gram matrix
$\bigl(\begin{smallmatrix}0&1\\1&0\end{smallmatrix}\bigr)$.
The complex-linear map between two such bases fixes $e$ and preserves
$h$, hence its imaginary part $\psi$.  As a real map it belongs to
$P_{e,\mathbb R}$ and conjugates the two complex structures.

\item The subgroup $P_{e,\mathbb Q}$ is dense in $P_{e,\mathbb R}$.
For the unipotent coordinates in \eqref{fam:eq:F-g} this is density of
$\mathbb Q^3$ in $\mathbb R^3$.  For $\operatorname{SL}_2$ one may
approximate $a,b,c$ by rational numbers in the chart $a\ne0$ and set
$d=(1+bc)/a$; translates of that chart cover the group.  Hence
\[
 \{rJ_0r^{-1}:r\in P_{e,\mathbb Q}\}
\]
is dense in $\mathcal D_0$.
Put $R_0=g(T^6,0,-1)$, so $R_0^{12}=R$ and
$(R_0-I)^2=0$.  At $rJ_0r^{-1}$ the nonzero rational nilpotent
$L=r(R_0-I)r^{-1}$ commutes with the complex structure and satisfies
$L^2=0$.  Choose a positive integer $d$ clearing its denominators.
The matrix $I+dL$ and its inverse $I-dL$ are integral, symplectic, and
fix $e$; this follows either by conjugating the identity
$(I+s(R_0-I))^{\mathsf t}\Omega(I+s(R_0-I))=\Omega$ or by expanding it.
The twelfth power of $I+dL$ lies in $\Gamma$, because
$P_{e,\mathbb Z}/\Gamma\cong\mathbb Z/12$.  It has infinite order and
fixes $rJ_0r^{-1}$.  Thus points with infinite $\Gamma$-stabilizer are
dense in $\mathcal D_0$.

\item Every open neighborhood in $\mathcal D$ of any point of $\mathcal D_0$
therefore contains a point with infinite stabilizer.  An open set on
which the action is proper cannot meet $\mathcal D_0$.  This proves the
largest-open-set assertion for $\Gamma$.  The same points have infinite
$P_{\mathbb Z}$-stabilizer, so it proves that assertion for the full
parabolic as well.
\end{enumerate}
\end{proof}

\begin{proposition}\label{fam:src:F.2.3}\label{fam:st:F.2.3}
\StatementRecord{fam:src:F.2.3}{proved}{Historical events: Claude Fable 5.2 (Anthropic), effort high, 12 November 2026: targeted proof check. Claude Fable 5.2 (Anthropic), effort high, 16 November 2026: used and rechecked in this run. GPT-6 Astra (OpenAI), ChatGPT Pro mode, effort Pro, S6\_09\_PAPERS, 17 November 2026: targeted proof check. Also checked in this writing. Pre-continuation writing events: GPT-6 Astra (OpenAI), ChatGPT Pro mode, effort Pro, S6\_10\_PAPERS, 18 November 2026: targeted proof check. Continuation metadata events: GPT-6 Astra (OpenAI), ChatGPT Pro mode, effort Pro, S6\_10\_PAPERS continuation, 19 November 2026: used and rechecked in this run; status-vocabulary reconciliation; metadata reconciliation only; no new mathematical check. \emph{Pre-continuation recorded proof scope.} Re-derived infinite unipotent stabilizer; null-locus rational-density argument; compact-stabilizer properness; unique normalized chart and q\_e=-6/D; j\_Gamma/1728 left inverse, extension over finite root discs, and proper embedding argument.}{the hypotheses, notation and definitions stated here; the explicit constructions and calculations in the proof}

The positive sign region admits the global period chart
\[
 \mathcal O_+=\left\{
 (\tau,\mu,\beta)\in\mathbb C^3:
 \operatorname{Im}\tau>0,\quad
 \operatorname{Im}\beta-\frac{6(\operatorname{Im}\mu)^2}
 {\operatorname{Im}\tau}<0\right\},
\]
represented by the matrix $\Pi$ in the stated problem.  In this chart,
\begin{equation}\label{fam:eq:F-qD}
 q_e(J)=-\frac6D.
\end{equation}
Every admissible period $\Pi_c(z)$ belongs to $\mathcal O_+$, for every
admissible $c$.  Thus this single maximal proper-action open set contains
all the period images in question.
\end{proposition}



\begin{proof}
\begin{enumerate}[label=\arabic*.,leftmargin=*]
\item Let $J\in\mathcal O_+$.  The two vectors $\widehat w,e$ are complex-linearly
independent in $(\Lambda_{\mathbb R},J)$.  Otherwise
$\widehat w=r e+sJe$ for real $r,s$.  Pairing with $e$ gives
$0=sq_e(J)$; hence $s=0$, contradicting the real linear independence
of $\widehat w,e$.  There is consequently a unique complex-linear
isomorphism $(\Lambda_{\mathbb R},J)\to\mathbb C^2$ sending these two
vectors to $(1,0)$ and $(0,1)$.  The corresponding period matrix has
right block $I$.  The equation
$\Pi\Omega^{-1}\Pi^{\mathsf t}=0$ says that its upper left entry is
six times its lower middle entry, giving exactly the displayed matrix.
This construction is the ordinary holomorphic Grassmannian chart,
so its coordinates and inverse are holomorphic.

\item In the real symplectic basis
$(\widehat u,\widehat\gamma/6,\widehat w,e)$ the normalized symmetric
period matrix is
$\bigl(\begin{smallmatrix}\tau&\mu\\\mu&\beta/6\end{smallmatrix}\bigr)$.
Its imaginary part has signature $(1,1)$, equivalently
$6(\operatorname{Im}\mu)^2-
\operatorname{Im}\tau\operatorname{Im}\beta>0$.
Solving $\Pi(Je)=i\Pi e=(0,i)$ in real coefficients gives
\[
 \gamma(Je)=
 \frac{\operatorname{Im}\tau}
 {\operatorname{Im}\tau\operatorname{Im}\beta-
 6(\operatorname{Im}\mu)^2}.
\]
Since $\psi(e,v)=-6\gamma(v)$, positivity of $q_e$ now gives
$\operatorname{Im}\tau>0$ and $D<0$, and proves \eqref{fam:eq:F-qD}.
Conversely these inequalities make the imaginary period matrix
invertible with the stated signature and give $q_e>0$.  The admissibility
inequality in the stated problem is $D_c<0$ everywhere, which proves the last
assertion by this exact comparison map.
\end{enumerate}
\end{proof}

\subsection{The curve in the proper quotient}\label{fam:section:F.3}

\begin{theorem}\label{fam:src:F.3.1}\label{fam:st:F.3.1}
\StatementRecord{fam:src:F.3.1}{proved}{Historical events: Claude Fable 5.2 (Anthropic), effort high, 12 November 2026: targeted proof check. Claude Fable 5.2 (Anthropic), effort high, 16 November 2026: used and rechecked in this run. GPT-6 Astra (OpenAI), ChatGPT Pro mode, effort Pro, S6\_09\_PAPERS, 17 November 2026: targeted proof check. Also checked in this writing. Pre-continuation writing events: GPT-6 Astra (OpenAI), ChatGPT Pro mode, effort Pro, S6\_10\_PAPERS, 18 November 2026: targeted proof check. Continuation metadata events: GPT-6 Astra (OpenAI), ChatGPT Pro mode, effort Pro, S6\_10\_PAPERS continuation, 19 November 2026: used and rechecked in this run; status-vocabulary reconciliation; metadata reconciliation only; no new mathematical check. \emph{Pre-continuation recorded proof scope.} Re-derived infinite unipotent stabilizer; null-locus rational-density argument; compact-stabilizer properness; unique normalized chart and q\_e=-6/D; j\_Gamma/1728 left inverse, extension over finite root discs, and proper embedding argument.}{the hypotheses, notation and definitions stated here; \cref{fam:eq:F-j}; \cref{fam:src:F.2.2}}

For every admissible $c$, the marked period map on the covering of
$B^\circ$ is a holomorphic immersion into $\mathcal O_+$.  The induced
map
\[
 p_c:B^\circ\longrightarrow\Gamma\backslash\mathcal O_+
\]
is injective and a holomorphic immersion.  It extends over $t=0,1$ to a
proper holomorphic embedding
\[
 \bar p_c:\mathbb C=B\setminus\{\infty\}
 \lhook\joinrel\longrightarrow\Gamma\backslash\mathcal O_+.
\]
The closure of $p_c(B^\circ)$ in this quotient is exactly
$\bar p_c(\mathbb C)$.
\end{theorem}



\begin{proof}
\begin{enumerate}[label=\arabic*.,leftmargin=*]
\item For $B=\bigl(\begin{smallmatrix}a&b\\c&d\end{smallmatrix}\bigr)$, changing
the marking by $g(B,x,t)$ replaces $\Pi$ by $\Pi g(B,x,t)^{-1}$ and then
renormalizes its right block.  On the quotient by the complex line of
$e$, the two periods are the row $(\tau,1)$.  Right multiplication by
$B^{-1}$ sends it to $(d\tau-c,a-b\tau)$.  Hence
\begin{equation}\label{fam:eq:F-tauaction}
 \tau\longmapsto\frac{d\tau-c}{a-b\tau}.
\end{equation}
The denominator never vanishes in the upper half plane.  The function
$j(\tau)$ is therefore a holomorphic invariant of this action and
descends to a holomorphic function $j_\Gamma$ on the proper quotient.
Here $j$ is the elliptic invariant, normalized by
$j(E)=1728\,4A^3/(4A^3+27B^2)$ for $y^2=x^3+Ax+B$; changing a lattice
basis preserves the isomorphism of elliptic curves and hence this
function.

\item We derive the required value in the present family.  The elliptic curve
whose period is $\tau$ has the equation stipulated in the stated problem,
\[
 y^2=x^3-3t^3(t-1)x+2t^4(t-1)^2.
\]
Its discriminant is $1728t^8(t-1)^3$ and direct cancellation gives
\begin{equation}\label{fam:eq:F-j}
 j(\tau)=1728t\qquad(t\in B^\circ).
\end{equation}
Equivalently this is the defining lift of the elliptic modular map used
for $\tau$ in the period construction.  The computation is reproduced
exactly in \texttt{f\_exact.py}.  In a local base coordinate $t$, the
derivative of the composite $j\circ\tau$ equals $1728$.  Thus the
derivative of the marked period map is nonzero.  The same identity after
passage to the quotient shows both injectivity and injectivity of the
tangent map of $p_c$: $j_\Gamma/1728$ is a holomorphic left inverse.

\item At the two elliptic points, the period functions are holomorphic in the
cyclic covering coordinates $s_1,s_2$, with $t=s_1^3$ and
$t-1=s_2^4$ after choosing local coordinates as in the stated problem.  Their
values are nondegenerate because $D_c<0$ there.  The maps from these
discs to $\mathcal O_+$ are equivariant for $A_1,A_2$.  Composing with
the quotient makes them invariant, so they descend to holomorphic maps
from the base discs.  This is the finite-quotient comparison map: a
holomorphic invariant germ on a disc is a convergent series in $s_j^{m_j}$,
which proves descent on local function rings.  The extensions agree
with $p_c$ off their centers and therefore with each other.  Identity
\eqref{fam:eq:F-j} extends by holomorphy, so the resulting $\bar p_c$ still
has holomorphic left inverse $j_\Gamma/1728$.

\item The quotient is Hausdorff by \ref{fam:src:F.2.2}.  For a compact subset $K$ of it,
$\bar p_c^{-1}(K)$ is closed in $\mathbb C$ and lies in the compact set
$j_\Gamma(K)/1728$.  Thus $\bar p_c$ is proper.  A holomorphic section of
a holomorphic map is locally a closed immersion: in local embeddings
into polydiscs, its coordinate along the base is $t$ and the remaining
coordinates are its holomorphic component functions, whose graph
defines the image.  Applying this to the left inverse shows that
$\bar p_c$ is a holomorphic embedding; properness makes its image
closed.  Since $\mathbb C\setminus\{0,1\}$ is dense in $\mathbb C$, its
image has the asserted closure.
\end{enumerate}
\end{proof}

\begin{proposition}\label{fam:src:F.3.2}\label{fam:st:F.3.2}
\StatementRecord{fam:src:F.3.2}{proved}{Historical events: Claude Fable 5.2 (Anthropic), effort high, 12 November 2026: targeted proof check. Claude Fable 5.2 (Anthropic), effort high, 16 November 2026: used and rechecked in this run. GPT-6 Astra (OpenAI), ChatGPT Pro mode, effort Pro, S6\_09\_PAPERS, 17 November 2026: targeted proof check. Also checked in this writing. Pre-continuation writing events: GPT-6 Astra (OpenAI), ChatGPT Pro mode, effort Pro, S6\_10\_PAPERS, 18 November 2026: targeted proof check. Continuation metadata events: GPT-6 Astra (OpenAI), ChatGPT Pro mode, effort Pro, S6\_10\_PAPERS continuation, 19 November 2026: used and rechecked in this run; status-vocabulary reconciliation; metadata reconciliation only; no new mathematical check. \emph{Pre-continuation recorded proof scope.} Re-derived all finite special-fibre period limits, the cusp j-limit and the limiting period row-span; computed the two-dimensional intersection with its complex conjugate from the explicit real limiting plane.}{the hypotheses, notation and definitions stated here; \cref{fam:src:F.3.1}}

At the two finite punctures the period map has the holomorphic
finite-cover limits
\[
 \Pi_1=
 \begin{pmatrix}2(2-\rho)&\rho&1&0\\
 \beta_c(z_1)&(2-\rho)/3&0&1\end{pmatrix},\qquad
 \Pi_2=
 \begin{pmatrix}3(1-i)&i&1&0\\
 \beta_c(z_2)&(1-i)/2&0&1\end{pmatrix},
 \quad \rho=e^{\pi i/3}.
\]
Their monodromies are the given order-three and order-four matrices.
At the cusp the curve leaves every compact subset of
$\Gamma\backslash\mathcal O_+$.  On the prescribed cusp branch its
Hodge plane in the compact dual tends to
$\mathbb C\gamma+\mathbb Cu$.  The latter plane has complex dimension
two in common with its conjugate and therefore does not represent a
point of $\mathcal D$.
\end{proposition}



\begin{proof}
\begin{enumerate}[label=\arabic*.,leftmargin=*]
\item The finite values of $\tau,\mu$ follow by evaluating their transformation
laws at the fixed points: $\tau(z_1)=\rho$,
$\mu(z_1)=(2-\rho)/3$, $\tau(z_2)=i$,
$\mu(z_2)=(1-i)/2$.  Substitution in $\Pi$ gives the two matrices.
The finite covering comparison and its descent were proved in \ref{fam:src:F.3.1}.
The unipotent cusp transformation is
$\tau\mapsto\tau-1$, $\mu\mapsto\mu$, $\beta\mapsto\beta+1$.
On the cusp branch, $\operatorname{Im}\tau\to+\infty$,
$\mu\to\mu_0$, and $\beta_c+\tau\to b_0+c$, from the holomorphic
cusp expansions.  Consequently
\[
 D_c=-\operatorname{Im}\tau+O(1),\qquad
 q_e(\Pi_c)\longrightarrow0^+,
 \qquad j_\Gamma(p_c(t))=1728t\longrightarrow\infty.
\]
The last equality alone excludes a convergent subnet in any compact
subset of the proper quotient; the first two identities give the
position relative to the maximal proper-action open set.
The rows of $\Pi_c$, regarded as covectors, are
$6\mu\gamma+\tau u+w$ and $\beta_c\gamma+\mu u+\delta$.
Dividing each by $\tau$ gives limits $u$ and $-\gamma$.
Their exterior product has a nonzero limit, so their span converges
in the Grassmannian to the stated plane.  It is spanned by real
covectors, hence equals its conjugate.  A complex structure has
$H^{1,0}\cap H^{0,1}=0$; here the intersection is the entire
two-dimensional plane.  This explicit comparison proves the assertion
about the compact-dual limit.
\end{enumerate}
\end{proof}

\begin{proposition}\label{fam:src:F.3.3}\label{fam:st:F.3.3}
\StatementRecord{fam:src:F.3.3}{proved}{Historical events: Claude Fable 5.2 (Anthropic), effort high, 12 November 2026: targeted proof check. Claude Fable 5.2 (Anthropic), effort high, 16 November 2026: used and rechecked in this run. GPT-6 Astra (OpenAI), ChatGPT Pro mode, effort Pro, S6\_09\_PAPERS, 17 November 2026: targeted proof check. Also checked in this writing. Pre-continuation writing events: GPT-6 Astra (OpenAI), ChatGPT Pro mode, effort Pro, S6\_10\_PAPERS, 18 November 2026: targeted proof check. Continuation metadata events: GPT-6 Astra (OpenAI), ChatGPT Pro mode, effort Pro, S6\_10\_PAPERS continuation, 19 November 2026: used and rechecked in this run; status-vocabulary reconciliation; metadata reconciliation only; no new mathematical check. \emph{Pre-continuation recorded proof scope.} Re-derived the cusp coordinate Jacobian and the dependent derivative columns at the finite special point. Checked nonzero derivative of the period map through j and the exact matrix map whose determinant vanishes.}{the hypotheses, notation and definitions stated here; \cref{fam:src:F.3.1}; \cref{fam:eq:F-j}}

The determinant used in the workbench is generically nonzero, with
the stated zero at $z_1$:
\[
 \det\dot Z=6\dot\mu^2-\dot\tau\dot\beta,\qquad
 \det\dot Z(z_1)=0.
\]
The immersion established in \ref{fam:src:F.3.1} holds at every point of $B^\circ$
without requiring this determinant to be everywhere nonzero.
\end{proposition}



\begin{proof}
\begin{enumerate}[label=\arabic*.,leftmargin=*]
\item In the cusp coordinate $s$, the expansion
$Z=s\bigl(\begin{smallmatrix}0&1\\-1&0\end{smallmatrix}\bigr)+C(q)$,
$q=e^{2\pi is}$, gives
$\det(dZ/ds)=1+O(q)$, so the determinant is not the zero function.
At $z_1$ the fixed-point transformation law of $\mu$ gives
$\dot\mu=-\dot\tau/3$.  One may obtain this by differentiating
$\mu\circ g_1=(1-\mu)/\tau$, using
$g_1'(z_1)=-\rho$ and $\rho^2-\rho+1=0$.
The law
$\beta\circ g_1=\beta+2-6(1-\mu)^2/\tau$ then gives
$\dot\beta=2\dot\tau/3$.  Substitution yields
$6\dot\tau^2/9-2\dot\tau^2/3=0$.
The derivative of a map from a curve to the three-dimensional period
chart is a vector of three coordinate derivatives.  Identity
\eqref{fam:eq:F-j} proves that this vector is nonzero over $B^\circ$;
the determinant tests the rank of the associated two-by-two
Kodaira--Spencer matrix.  The maps compared here are explicitly the
coordinate derivative and its determinant, and the displayed computation
states exactly what each supplies.
\end{enumerate}
\end{proof}



\section{Exceptional fibres, Neron--Severi classes, and compact subspaces}\label{fam:section:D}
\subsection{Elliptic quotients and the torsion locus}\label{fam:section:D.1}
\begin{proposition}\label{fam:src:D.1.1}\label{fam:st:D.1.1}
\StatementRecord{fam:src:D.1.1}{proved}{Historical events: Claude Fable 5.2 (Anthropic), effort high, 12 November 2026: targeted proof check. Claude Fable 5.2 (Anthropic), effort high, 16 November 2026: used and rechecked in this run. GPT-6 Astra (OpenAI), ChatGPT Pro mode, effort Pro, S6\_09\_PAPERS, 17 November 2026: targeted proof check. Also checked in this writing. Pre-continuation writing events: GPT-6 Astra (OpenAI), ChatGPT Pro mode, effort Pro, S6\_10\_PAPERS, 18 November 2026: targeted proof check. Continuation metadata events: GPT-6 Astra (OpenAI), ChatGPT Pro mode, effort Pro, S6\_10\_PAPERS continuation, 19 November 2026: used and rechecked in this run; status-vocabulary reconciliation; metadata reconciliation only; no new mathematical check. \emph{Pre-continuation recorded proof scope.} Re-derived projection descent and connected kernel, order-four counterexample lattice images, wedge equation for NS, exact complex-orientation sign, monodromy finite-difference independence and explicit projective sequence. The Weierstrass scaling direction is carried by the explicit map I:weierstrass-comparison. Read the complete source proof and re-derived the rank-three lattice/kernel and closure, rational product covering, dense cusp natural-boundary obstruction, exact toric gluing/dihedral action, quotient and Pfaffian criteria, explicit non-torsion algebraic-dimension-one locus, and finite-normalization curve maps as applicable. Corrected the Weierstrass map direction and sign in D.1.1 by the already proved Phi; made the rational product identification in C.3.2 explicit. Exact source passages and finite symbolic replays retain their separately recorded scopes.}{the hypotheses, notation and definitions stated here; the explicit constructions and calculations in the proof}

The specified linear projection \(\operatorname{pr}_1:\mathbb C^2\to
\mathbb C\) induces a holomorphic group homomorphism
\(F_t\to E_t=\mathbb C/(\mathbb Z\tau+\mathbb Z)\) exactly on
\(6\mu=A\tau+B\), \(A,B\in\mathbb Z\). It is then surjective with
connected kernel
\(\mathbb C/(\mathbb Z+(A\mu-\beta)\mathbb Z)\).
On the smooth base the two values of \(t\) are precisely
\[
 t_\pm=\frac{64(47\pm27\sqrt3)}{11}.
\]
\end{proposition}


\begin{proof}
\begin{enumerate}[label=\arabic*.,leftmargin=*]
\item The images of the four lattice generators under the linear projection
are \(6\mu,\tau,1,0\). This proves descent in both directions.
The integral basis change replacing \(\widehat\gamma\) with
\(\widehat\gamma-A\widehat u-B\widehat w\) changes its period to
\((0,\beta-A\mu)\); the remaining vertical generator is \((0,1)\).
The other two periods map onto the entire target lattice, so the kernel
is connected. Its second modulus has imaginary part
\(A\operatorname{Im}\mu-\operatorname{Im}\beta=-D>0\).

\item The geometric model and its marked section are transported by
$\Phi^{-1}(X,Y)=(X/2,-Y/(2\sqrt2))$ to the arithmetic model
below. The inverse map is
$\Phi(x,y)=(2x,-2\sqrt2y)$; substitution in both cubic
equations proves the isomorphism and its effect on the marked point
(\cref{I:weierstrass-comparison}). Thus the arithmetic equation is
\[
 Y^2=X^3-\frac34t^3(t-1)X+\frac14t^4(t-1)^2,
 \qquad P=(0,t^2(t-1)/2).
\]
Its sixth division polynomial at \(P\) is
\[
 \psi_6(P)=\frac9{65536}t^{24}(t-1)^9
 (11t^2-6016t+8192).
\]
For a smooth elliptic fibre, \(\psi_n(P)=0\) is exactly the condition
\(nP=O\), since the multiplication rational map has these and only these
poles in its affine coordinate. The two roots displayed above are simple
and avoid \(0,1\). We computed the same equation independently by the
chord-and-tangent addition algorithm, obtaining
\[
 \begin{aligned}
 x(6P)&=\frac{(89t^2+4736t-4096)Q(t)}
 {1296(11t^2-6016t+8192)^2},\\
 Q(t)&=881t^4+4745344t^3-12079104t^2\\
 &\hspace{2em}+24641536t-16777216.
 \end{aligned}
\]
The division-polynomial recurrence and the addition algorithm are in
separate supplied programs; all polynomial operations use exact rational
arithmetic. Thus the denominator has no unverified numerical root or
cancellation. The decimal values \(1.36510950191\) and
\(545.543981407\) are only convenient approximations to the exact roots.
\end{enumerate}
\end{proof}

\begin{proposition}\label{fam:src:D.1.2}\label{fam:st:D.1.2}
\StatementRecord{fam:src:D.1.2}{proved}{Historical events: Claude Fable 5.2 (Anthropic), effort high, 12 November 2026: targeted proof check. Claude Fable 5.2 (Anthropic), effort high, 16 November 2026: used and rechecked in this run. GPT-6 Astra (OpenAI), ChatGPT Pro mode, effort Pro, S6\_09\_PAPERS, 17 November 2026: targeted proof check. Also checked in this writing. Pre-continuation writing events: GPT-6 Astra (OpenAI), ChatGPT Pro mode, effort Pro, S6\_10\_PAPERS, 18 November 2026: targeted proof check. Continuation metadata events: GPT-6 Astra (OpenAI), ChatGPT Pro mode, effort Pro, S6\_10\_PAPERS continuation, 19 November 2026: used and rechecked in this run; status-vocabulary reconciliation; metadata reconciliation only; no new mathematical check. \emph{Pre-continuation recorded proof scope.} Read complete original appendix proof and independently derived the Betti Jacobian, density/countability, explicit index-N\textasciicircum{}3 sublattice and product isogeny, and 3k\textasciicircum{}2-1 height degree. Full specialized proof moved beside the statement; exact Engel source locator and division replay tracked separately.}{\cref{fam:src:C.2.1}, \cref{fam:src:D.1.1}; the explicit lattice map and section intersection computation below; \cite[Lemma 2.1]{Engel2026}; \cite[Theorem 4.4, p.~3]{Schnell19}}

The locus on \(B^\circ\) where \(P(t)\) is torsion is countable and dense.
Every fibre over this locus is an abelian surface. The divisor on the
smooth base cut out by \(6kP=O\) has degree \(3k^2-1\), including local
intersection multiplicity, for every positive integer \(k\).
\end{proposition}

\begin{proof}
\begin{enumerate}[label=\arabic*.,leftmargin=*]
\item \label{app:rederive:fam-src-D.1.2}
On a simply connected open subset write uniquely
\(\mu=a_0\tau+b_0\), \(a_0,b_0\) real. The section is torsion precisely
when \((a_0,b_0)\in\mathbb Q^2\). In a local complex parameter \(z\), direct
differentiation gives
\[
 \det\frac{\partial(a_0,b_0)}{\partial(\operatorname{Re}z,
 \operatorname{Im}z)}
 =-\frac{|\mu'(z)-a_0\tau'(z)|^2}{\operatorname{Im}\tau}.
\]
This determinant vanishes on no open set. On an open set of vanishing,
the two real differential equations imply \(da_0=db_0=0\), hence
\(\mu=a_0\tau+b_0\) with constant real coefficients. Analytic continuation
would give this identity on the cover. The cusp law
\((\tau,\mu)\circ g_0=(\tau-1,\mu)\) then forces \(a_0=0\), and the
\(g_2\)-law makes a constant \(\mu\) impossible. Thus every open subset
contains a point where this real map is a local diffeomorphism, and its
image contains rational pairs. This proves density. For each integer
triple \((n,A,B)\), \(n\ne0\), the function \(n\mu-A\tau-B\) is
holomorphic and not identically zero by the same argument. Its zeros are
discrete; a countable union proves countability.

\item For \(N\mu=A\tau+B\), \(A,B,N\in\mathbb Z\), \(N>0\), use the linear
coordinate change \(\zeta'_2=\zeta_2-(A/N)\zeta_1\).
The four integral vectors
\[
 N\widehat u-B\widehat\delta,\quad
 N\widehat w+A\widehat\delta,\quad
 N\widehat\gamma-6A\widehat u-6B\widehat w,\quad
 \widehat\delta
\]
span a sublattice of index \(N^3\). Its first two periods lie on
\(\zeta'_2=0\), and its last two lie on \(\zeta_1=0\).
Both intersections are genuine rank-two lattices because \(D<0\).
Their product quotient is therefore a product of elliptic curves, and
inclusion of this sublattice induces a finite surjective holomorphic
homomorphism onto \(F_t\). A positive integral form on the product can
be summed over its finite translation kernel and multiplied by the
index to descend an integral positive form to \(F_t\). The Riemann
polarization criterion proves projectivity.

\item For the degree statement we use the height computation actually read in
\cite[Lemma 2.1 and proof]{Engel2026}: \(\langle P,P\rangle=1/6\),
\(\chi(\mathcal O_S)=1\), the component groups at the two reducible fibres
have orders three and two, and the local corrections for \(P\) are
\(4/3\) and \(1/2\). The application to \(6kP\) is as follows.
The corrected divisor class of a section is its class minus \(O\) and
a fibre multiple, plus the unique rational combination of nonidentity
fibre components that makes it perpendicular to those components.
On restriction to the generic fibre this assignment is additive because
the divisor class \([Q]-[O]\) is the group law in \(\operatorname{Pic}^0\).
The perpendicular correction makes its self-intersection the negative of
the height, so the height is quadratic. The element \(6kP\) is in the
identity component at both reducible fibres; its correction terms vanish.
Hence
\[
 6k^2=(6k)^2\langle P,P\rangle
 =2+2((6kP)\mathbin{\cdot}O),
\]
which gives the formula. There are no missing intersection points at the
three special fibres. The computed rational functions for \(6P\) show
that it has finite affine coordinates at \(0,1\), so it is a nonzero point
of each additive identity component; multiplication by \(k\) remains
nonzero in \(\mathbb C\). At infinity its multiplicative coordinate is
\((-a)^{6k}=a^{6k}\ne1\). Thus all intersections counted by the height
occur on \(B^\circ\). This derives the cited height input in precisely
the form needed here; the first five degrees and the correction matrices
are checked exactly in \texttt{cd\_torsion\_independent.py}.
\end{enumerate}
\end{proof}

\Uses{the hypotheses, notation and definitions stated here; \cref{app:rederive:fam-src-D.1.2}}


This specializes the elliptic-surface height and specialization results cited in the re-derivation. A complete re-derivation is given in Appendix~\ref{app:rederive:fam-src-D.1.2}.


\begin{theorem}[counterexample]\label{fam:src:D.1.3}\label{fam:st:D.1.3}
\StatementRecord{fam:src:D.1.3}{proved}{Historical events: Claude Fable 5.2 (Anthropic), effort high, 12 November 2026: targeted proof check. Claude Fable 5.2 (Anthropic), effort high, 16 November 2026: used and rechecked in this run. GPT-6 Astra (OpenAI), ChatGPT Pro mode, effort Pro, S6\_09\_PAPERS, 17 November 2026: targeted proof check. Also checked in this writing. Pre-continuation writing events: GPT-6 Astra (OpenAI), ChatGPT Pro mode, effort Pro, S6\_10\_PAPERS, 18 November 2026: targeted proof check. Continuation metadata events: GPT-6 Astra (OpenAI), ChatGPT Pro mode, effort Pro, S6\_10\_PAPERS continuation, 19 November 2026: used and rechecked in this run; status-vocabulary reconciliation; metadata reconciliation only; no new mathematical check. \emph{Pre-continuation recorded proof scope.} Re-derived projection descent and connected kernel, order-four counterexample lattice images, wedge equation for NS, exact complex-orientation sign, monodromy finite-difference independence and explicit projective sequence. The Weierstrass scaling direction is carried by the explicit map I:weierstrass-comparison. Read the complete source proof and re-derived the rank-three lattice/kernel and closure, rational product covering, dense cusp natural-boundary obstruction, exact toric gluing/dihedral action, quotient and Pfaffian criteria, explicit non-torsion algebraic-dimension-one locus, and finite-normalization curve maps as applicable. Corrected the Weierstrass map direction and sign in D.1.1 by the already proved Phi; made the rational product identification in C.3.2 explicit. Exact source passages and finite symbolic replays retain their separately recorded scopes.}{the hypotheses, notation and definitions stated here; \cref{fam:src:D.1.1}; \cref{fam:src:C.2.1}}

The proposed assertion that \emph{existence of a surjection onto the
specified \(E_t\) with elliptic kernel} is equivalent to \(6P(t)=O\)
is false. It remains true for the specified projection in \ref{fam:src:D.1.1}.
\end{theorem}


\begin{proof}
\begin{enumerate}[label=\arabic*.,leftmargin=*]
\item Set \(t=32/5\). The exact fourth division polynomial is
\[
 \psi_4(P)=-\frac1{32}t^{10}(t-1)^4(5t-32),
\]
whereas \(\psi_2(P)=t^2(t-1)\ne0\). Thus \(P\) has exact order four.
On the rectangular branch write \(\tau=iT\), \(T>1\), with
\(j(iT)=1728\cdot32/5\). The chosen zero has real part \(1/2\), negative
imaginary part of absolute value less than \(T/2\), and order four;
therefore \(\mu=(2-\tau)/4\). Choose the concrete parameter
\[
 c=\left(\frac98-\frac{101}{4}\right)\tau-\frac94
       -\beta^{\rm part}(t).
\]
It is globally admissible. Indeed the positive-real Jacobi product at this
\((\mu,\tau)\) gives
\(\operatorname{Im}\beta^{\rm part}\geq-T-(3/\pi)\log3\).
Consequently
\(\operatorname{Im}c\leq-185T/8+(3/\pi)\log3<-22<-M\), by \ref{fam:src:C.2.1}.
This inequality proves admissibility without using the numerical estimate
\(c=-2.25-34.02109724218209472308\,i\).

\item The holomorphic linear map \(L:\mathbb C^2\to\mathbb C\),
\(L(\zeta)=3\zeta_1+4\zeta_2\), takes the four periods to
\[
 -101\tau,\qquad2\tau+2,\qquad3,\qquad4.
\]
These generate exactly \(\mathbb Z\tau+\mathbb Z\): the last two generate
\(1\), and the first two then generate \(\tau\), since
\(\gcd(101,2)=1\). Thus \(L\) descends to a surjective holomorphic group
homomorphism \(F_{32/5}\to E_{32/5}\) with connected one-dimensional
kernel. The kernel is compact as a closed subset of a compact torus,
hence an elliptic curve. Nevertheless \(6P=2P\ne O\).
The program \path{cd_torsion_independent.py} checks the displayed
division polynomials and four lattice images by exact arithmetic; its
command and retained output are in \CompendiumRef{I}{app:I:programs}.
The admissibility inequality is proved analytically above.
\end{enumerate}
\end{proof}

\begin{proposition}\label{fam:src:D.1.4}\label{fam:st:D.1.4}
\StatementRecord{fam:src:D.1.4}{proved}{Historical events: Claude Fable 5.2 (Anthropic), effort high, 12 November 2026: targeted proof check. Claude Fable 5.2 (Anthropic), effort high, 16 November 2026: used and rechecked in this run. GPT-6 Astra (OpenAI), ChatGPT Pro mode, effort Pro, S6\_09\_PAPERS, 17 November 2026: targeted proof check. Also checked in this writing. Pre-continuation writing events: GPT-6 Astra (OpenAI), ChatGPT Pro mode, effort Pro, S6\_10\_PAPERS, 18 November 2026: targeted proof check. Continuation metadata events: GPT-6 Astra (OpenAI), ChatGPT Pro mode, effort Pro, S6\_10\_PAPERS continuation, 19 November 2026: used and rechecked in this run; status-vocabulary reconciliation; metadata reconciliation only; no new mathematical check. \emph{Pre-continuation recorded proof scope.} Re-derived projection descent and connected kernel, order-four counterexample lattice images, wedge equation for NS, exact complex-orientation sign, monodromy finite-difference independence and explicit projective sequence. The Weierstrass scaling direction is carried by the explicit map I:weierstrass-comparison. Read the complete source proof and re-derived the rank-three lattice/kernel and closure, rational product covering, dense cusp natural-boundary obstruction, exact toric gluing/dihedral action, quotient and Pfaffian criteria, explicit non-torsion algebraic-dimension-one locus, and finite-normalization curve maps as applicable. Corrected the Weierstrass map direction and sign in D.1.1 by the already proved Phi; made the rational product identification in C.3.2 explicit. Exact source passages and finite symbolic replays retain their separately recorded scopes.}{the hypotheses, notation and definitions stated here; \cref{fam:src:C.3.1}; \cref{fam:src:D.1.6}}

For every smooth fibre and every admissible \(c\), its Neron--Severi lattice
is the following explicit subgroup of \(\mathbb Z^6\):
\begin{equation}\label{fam:eq:cd-NS}
 \left\{E:\ E_{12}-E_{13}\tau+(6E_{23}-E_{14})\mu+
 E_{24}\beta+E_{34}(6\mu^2-\tau\beta)=0\right\}.
\end{equation}
The complex orientation has intersection form
\(E^2=-2\operatorname{Pf}(E)\), where
\(\operatorname{Pf}(E)=E_{12}E_{34}-E_{13}E_{24}+E_{14}E_{23}\).
In particular the invariant class \(\eta=6\gamma\wedge\delta+u\wedge w\)
has \(\int_{F_t}\eta^2=-12\). Its Hermitian form is indefinite.
The fibres with lattice exactly \(\mathbb Z\eta\) have no compact curves
and algebraic dimension zero.
\end{proposition}


\begin{proof}
\begin{enumerate}[label=\arabic*.,leftmargin=*]
\item On the universal cover the holomorphic two-form is the wedge of the two
rows of \(\Pi\). Its six coefficients in the integral coframe are
\(6\mu^2-\tau\beta,-\beta,6\mu,-\mu,\tau,1\).
An integral real two-form is of type \((1,1)\) exactly when its wedge
with this form vanishes: its \((0,2)\) part is then zero, and reality
also annihilates its \((2,0)\) part. Expanding the wedge gives
\eqref{fam:eq:cd-NS}. The exponential sequence identifies these integral
\((1,1)\) classes with Chern classes of holomorphic line bundles, proving
the lattice assertion.

\item The real period matrix in \ref{fam:src:C.3.1} has determinant \(TD<0\). Thus the
integral orientation \((\widehat\gamma,\widehat u,\widehat w,
\widehat\delta)\) is opposite to the complex orientation. Algebraically
\(E\wedge E=2\operatorname{Pf}(E)\gamma\wedge u\wedge w\wedge\delta\),
which proves the sign in the displayed intersection formula. In
particular the expression ``\(\eta^2=12\)'' in AC26, Lemma 9.4 is correct
as its displayed coefficient of \(\gamma\wedge u\wedge w\wedge\delta\),
and has the opposite sign as a complex-oriented numerical
self-intersection. We retain the coefficient identity and record the
orientation comparison explicitly.

\item For the precise Hermitian convention $H_E$ defined in \ref{fam:src:D.1.6}, direct
substitution of the four periods gives $H_\eta=-2G^{-1}$, where
\[
 G=\begin{pmatrix}2T&2U\\2U&\operatorname{Im}\beta/3\end{pmatrix},
 \qquad \det G=2TD/3<0.
\]
Thus $H_\eta$ has one positive and one negative eigenvalue, and every
nonzero multiple remains indefinite. The integration
current of an effective divisor, averaged under translations of the
compact torus, is a nonzero semipositive constant \((1,1)\)-form in its
Chern class. It cannot be an indefinite multiple of \(\eta\).
Thus no curve exists when \(\operatorname{NS}=\mathbb Z\eta\).
A nonconstant meromorphic function would have a nonempty polar divisor
(after subtracting a constant if necessary), so the algebraic dimension
is zero.
\end{enumerate}
\end{proof}

\begin{proposition}\label{fam:src:D.1.5}\label{fam:st:D.1.5}
\StatementRecord{fam:src:D.1.5}{proved}{Historical events: Claude Fable 5.2 (Anthropic), effort high, 12 November 2026: targeted proof check. Claude Fable 5.2 (Anthropic), effort high, 16 November 2026: used and rechecked in this run. GPT-6 Astra (OpenAI), ChatGPT Pro mode, effort Pro, S6\_09\_PAPERS, 17 November 2026: targeted proof check. Also checked in this writing. Pre-continuation writing events: GPT-6 Astra (OpenAI), ChatGPT Pro mode, effort Pro, S6\_10\_PAPERS, 18 November 2026: targeted proof check. Continuation metadata events: GPT-6 Astra (OpenAI), ChatGPT Pro mode, effort Pro, S6\_10\_PAPERS continuation, 19 November 2026: used and rechecked in this run; status-vocabulary reconciliation; metadata reconciliation only; no new mathematical check. \emph{Pre-continuation recorded proof scope.} Re-derived projection descent and connected kernel, order-four counterexample lattice images, wedge equation for NS, exact complex-orientation sign, monodromy finite-difference independence and explicit projective sequence. The Weierstrass scaling direction is carried by the explicit map I:weierstrass-comparison. Read the complete source proof and re-derived the rank-three lattice/kernel and closure, rational product covering, dense cusp natural-boundary obstruction, exact toric gluing/dihedral action, quotient and Pfaffian criteria, explicit non-torsion algebraic-dimension-one locus, and finite-normalization curve maps as applicable. Corrected the Weierstrass map direction and sign in D.1.1 by the already proved Phi; made the rational product identification in C.3.2 explicit. Exact source passages and finite symbolic replays retain their separately recorded scopes.}{the hypotheses, notation and definitions stated here; \cref{fam:eq:cd-NS}}

For a fixed \(c\), all but countably many smooth fibres have
\(\operatorname{NS}=\mathbb Z\eta\). Conversely, for every fixed smooth
\(t\), there are admissible parameters tending to the cusp of the
parameter disc for which \(F_t\) is an abelian surface, irrespective of
the order of \(P(t)\).
\end{proposition}


\begin{proof}
\begin{enumerate}[label=\arabic*.,leftmargin=*]
\item We verify the required independence of five holomorphic functions.
Suppose \(A+B\tau+C\mu+D\beta+E(6\mu^2-\tau\beta)=0\).
Apply \(g_0\), using \((\tau,\mu,\beta)\circ g_0=(\tau-1,\mu,\beta+1)\),
and subtract. The result is \(-B+D+E(\beta-\tau+1)=0\).
A second cusp difference gives \(2E=0\), hence \(E=0,D=B\).
Applying \(g_2\) to the remaining relation and multiplying the difference
by \(-\tau\) gives
\[
 6B\mu^2+C(\tau-1)\mu+B(\tau^2+3\tau+1)-C\tau=0.
\]
Its cusp difference gives \(C\mu+2B(\tau+1)-C=0\), whose next cusp
difference forces \(B=0\). Then \(C\mu=C\); the \(g_2\)-law forces
\(C=0\), and finally \(A=D=0\). Thus an integral class outside
\(\mathbb Z\eta\) has a nonidentically-zero holomorphic left side in
\eqref{fam:eq:cd-NS}; its zeros are discrete. Countability of \(\mathbb Z^6\)
proves the first assertion, with explicit verification of its reduction.

\item For the second assertion put
\(\beta_n=-3n^2\tau+12n\mu\),
\(c_n=\beta_n-\beta^{\rm part}(t)\). Its imaginary part tends to
minus infinity, so the already proved finite bound on \(M\) makes every
sufficiently large \(n\) admissible. The integral alternating matrix
with nonzero upper entries
\[
 E_{13}=-3n^2,\quad E_{14}=6n,\quad E_{23}=-n,\quad E_{24}=1
\]
satisfies \eqref{fam:eq:cd-NS} exactly at \(\beta_n\), and
\(\operatorname{Pf}(E)=-3n^2<0\). Its complex-oriented square is positive.
A nondegenerate Hermitian form in complex dimension two with positive
determinant is definite. Choosing its sign gives a positive integral
Riemann form, so \(F_t\) is projective. This supplies an explicit family
of projective fibres even at any chosen non-torsion value of the section.
\end{enumerate}
\end{proof}

\begin{theorem}\label{fam:src:D.1.6}\label{fam:st:D.1.6}
\StatementRecord{fam:src:D.1.6}{proved}{Historical events: Claude Fable 5.2 (Anthropic), effort high, 12 November 2026: targeted proof check. Claude Fable 5.2 (Anthropic), effort high, 16 November 2026: used and rechecked in this run. GPT-6 Astra (OpenAI), ChatGPT Pro mode, effort Pro, S6\_09\_PAPERS, 17 November 2026: targeted proof check. Also checked in this writing. Pre-continuation writing events: GPT-6 Astra (OpenAI), ChatGPT Pro mode, effort Pro, S6\_10\_PAPERS, 18 November 2026: targeted proof check. Continuation metadata events: GPT-6 Astra (OpenAI), ChatGPT Pro mode, effort Pro, S6\_10\_PAPERS continuation, 19 November 2026: used and rechecked in this run; status-vocabulary reconciliation; metadata reconciliation only; no new mathematical check. \emph{Pre-continuation recorded proof scope.} Read complete proof; re-derived H\_E=E(.,J.)-iE, rational saturated null lattice, actual inclusion/quotient morphisms, Pfaffian sign, effective-divisor positivity and algebraic dimensions. Opened Schnell Theorem4.1 p1, equation4.3 p2, Theorem4.4 p3; full family application is printed.}{\cref{fam:eq:cd-NS}, \cref{fam:src:D.1.4}; \cite[Theorem 4.1, p.~1, and Theorem 4.4, p.~3]{Schnell19}}

Fix any admissible $c$. For each integral alternating matrix $E$, let
$Z_E(c)$ be the zero set on the marking cover of the explicit holomorphic
function in \eqref{fam:eq:cd-NS}, with $\beta=\beta^{\rm part}+c$ given by
\eqref{fam:eq:cd-beta}. Let $\pi$ be the cover map to $B^\circ$, and define
\[
 \mathcal E_c=\pi\!\left(\bigcup_{E\ne0,\ \operatorname{Pf}(E)=0}Z_E(c)\right),
 \qquad
 \mathcal A_c=\pi\!\left(\bigcup_{\operatorname{Pf}(E)<0}Z_E(c)\right).
\]
These give the exact loci of all elliptic quotients and all abelian
surfaces, respectively. The algebraic dimension of every smooth fibre is
\[
 a(F_t)=
 \begin{cases}
 2,&t\in\mathcal A_c,\\
 1,&t\in\mathcal E_c\setminus\mathcal A_c,\\
 0,&t\notin\mathcal E_c\cup\mathcal A_c.
 \end{cases}
\]
Together with \eqref{fam:eq:cd-NS}, these are an exact classification for every
$t$, by countable unions of explicitly given analytic zero sets; they do
not assert that the loci have a finite algebraic description. They include
all quotient maps, without requiring preservation of the first coordinate.
\end{theorem}

\begin{proof}
\begin{enumerate}[label=\arabic*.,leftmargin=*]
\item \label{app:rederive:fam-src-D.1.6}
We apply the Appell--Humbert correspondence and the Riemann projectivity
criterion in \cite[Theorem 4.1, p.~1, and Theorem 4.4, p.~3]{Schnell19}; the reduction to the
present explicit periods is as follows. The real vector-space isomorphism
$\Pi:\mathbb R^4\to\mathbb C^2$ transfers $E$ to the actual alternating
form $\widetilde E(v,w)=E(\Pi^{-1}v,\Pi^{-1}w)$ on the tangent space.
Equation \eqref{fam:eq:cd-NS} gives $J$-invariance. The explicit comparison is
\[
 H_E(v,w)=\widetilde E(v,Jw)-i\widetilde E(v,w).
\]
Direct substitution gives $H_E(Jv,w)=iH_E(v,w)$ and
$H_E(w,v)=\overline{H_E(v,w)}$. Its associated real Chern form is
$-\operatorname{Im}H_E=\widetilde E$, since
$(i/2)(H_E(v,w)-H_E(w,v))=\widetilde E(v,w)$.
Thus the Appell--Humbert integrality condition is exactly integrality
of the given $E$, with its sign fully specified. The complex-oriented square computed in \ref{fam:src:D.1.4} shows that
$\operatorname{Pf}(E)<0$ is equivalent to definiteness of this Hermitian
form. Its sign can be chosen positive. The cited projectivity criterion
then gives precisely the asserted abelian locus, in both directions.

\item For nonzero $E$ with Pfaffian zero, its real rank is two. Its nullspace
$K_{\mathbb R}\subset\mathbb R^4$ is a rational plane because it is the
kernel of an integral matrix. $J$-invariance makes $\Pi K_{\mathbb R}$ a
complex line, and $K=K_{\mathbb R}\cap\mathbb Z^4$ is saturated. The maps
\[
 (\Pi K_{\mathbb R})/\Pi K\ \hookrightarrow F_t
 \ \longrightarrow\
 \mathbb C^2/(\Pi K_{\mathbb R}+\Pi\mathbb Z^4)
\]
are holomorphic group morphisms; the first is a closed embedded elliptic
curve and the second is its quotient onto an elliptic curve, with
connected kernel. A basis of the saturated $K$ extends to an integral basis of $\mathbb Z^4$.
The other two periods project to real-linearly independent vectors in the
one-dimensional complex quotient, because $\Pi$ is a real isomorphism.
They therefore form a discrete lattice there. Compactness follows by
taking the image of its compact real fundamental parallelogram. Conversely any elliptic quotient
pulls an integral positive area form back to a nonzero rank-two integral
$(1,1)$ form. This proves the characterization of $\mathcal E_c$.
Equivalently all its kernels can be found by saturated rank-two lattices
with a basis $v,w$ satisfying
$\det_{\mathbb C}(\Pi v,\Pi w)=0$ and real independence. This is a second
exact integer test, with explicit inclusion and quotient morphisms.

\item For the algebraic-dimension assertion, an effective divisor on a torus
has a semipositive Hermitian form: averaging its integration current under
translations gives that constant form. A nonzero semipositive form has
rank one or two. Rank two gives the abelian case. Rank one gives the
elliptic quotient just constructed. In the absence of a definite form,
there cannot be two different such kernel subtori. Indeed addition from
their product onto $F_t$ would be a surjective holomorphic homomorphism
with finite kernel, and the product polarization would descend after
multiplication by an integer, giving a definite form. Thus there is a
single kernel $K_t$ in this case. Every effective divisor is a sum of
its translates: its intersection number with $K_t$ is zero, whereas a
curve mapping nonconstantly to the elliptic quotient has strictly positive
intersection with a general quotient fibre. Every meromorphic function
therefore factors through this quotient. One can check the last statement
on a fibre avoiding its polar divisor, where it is a holomorphic function
on a compact elliptic curve and hence constant; analytic continuation
then gives the factorization. The quotient supplies nonconstant
meromorphic functions, so the algebraic dimension is exactly one.
If neither locus occurs, a nonconstant meromorphic function would have
a polar divisor of one of the two forbidden ranks. This proves dimension
zero. Finally monodromy acts by integral determinant-one changes of basis,
preserving \eqref{fam:eq:cd-NS} and its Pfaffian. Thus the stated unions descend
without any marking ambiguity. Each individual matrix and every locus
is specified; no expected identification is left in these formulas.
\end{enumerate}
\end{proof}

\Uses{the hypotheses, notation and definitions stated here; \cref{fam:eq:cd-NS}; \cref{fam:eq:cd-beta}; \cref{app:rederive:fam-src-D.1.6}}


This is the Appell--Humbert and Riemann criterion in the present lattice. A complete re-derivation is given in Appendix~\ref{app:rederive:fam-src-D.1.6}.


\begin{proposition}\label{fam:src:D.1.7}\label{fam:st:D.1.7}
\StatementRecord{fam:src:D.1.7}{proved}{Historical events: Claude Fable 5.2 (Anthropic), effort high, 12 November 2026: targeted proof check. Claude Fable 5.2 (Anthropic), effort high, 16 November 2026: used and rechecked in this run. GPT-6 Astra (OpenAI), ChatGPT Pro mode, effort Pro, S6\_09\_PAPERS, 17 November 2026: targeted proof check. Also checked in this writing. Pre-continuation writing events: GPT-6 Astra (OpenAI), ChatGPT Pro mode, effort Pro, S6\_10\_PAPERS, 18 November 2026: targeted proof check. Continuation metadata events: GPT-6 Astra (OpenAI), ChatGPT Pro mode, effort Pro, S6\_10\_PAPERS continuation, 19 November 2026: used and rechecked in this run; status-vocabulary reconciliation; metadata reconciliation only; no new mathematical check. \emph{Pre-continuation recorded proof scope.} Read the complete source proof and re-derived the rank-three lattice/kernel and closure, rational product covering, dense cusp natural-boundary obstruction, exact toric gluing/dihedral action, quotient and Pfaffian criteria, explicit non-torsion algebraic-dimension-one locus, and finite-normalization curve maps as applicable. Corrected the Weierstrass map direction and sign in D.1.1 by the already proved Phi; made the rational product identification in C.3.2 explicit. Exact background citations and the indicated finite symbolic programs remain separate gates.}{the hypotheses, notation and definitions stated here; \cref{fam:eq:cd-NS}; \cref{fam:src:D.1.6}; \cref{fam:src:C.2.1}; \cref{fam:src:D.1.2}}

There is an explicit nonempty locus of admissible members with ordinary
fibres of algebraic dimension one. On a sufficiently small punctured
cusp neighbourhood set
\[
 c(z)=-6\tau(z)-12\mu(z)-\beta^{\rm part}(z).
\]
Outside a countable set of $z$, the fibre of $X_{c(z)}$ over $\pi(z)$ has
Neron--Severi lattice $\mathbb Z\eta+\mathbb ZE$, where the nonzero upper
entries of $E$ are
$E_{13}=-6$, $E_{14}=-6$, $E_{23}=1$, $E_{24}=1$.
Its algebraic dimension is one. The section $P(\pi(z))$ can simultaneously
be required to be non-torsion.
\end{proposition}


\begin{proof}
\begin{enumerate}[label=\arabic*.,leftmargin=*]
\item At the specified parameter $\beta=-6\tau-12\mu$. Both $E$ and $\eta$
satisfy \eqref{fam:eq:cd-NS}, and exact multiplication gives
$\operatorname{Pf}(a\eta+bE)=6a^2$.
For any further integral class $F$, substitution in its period equation
gives a linear relation among $1,\tau,\mu,(\mu+\tau)^2$.
These four functions are linearly independent: applying $g_0$ to a
relation $A+B\tau+C\mu+D(\mu+\tau)^2=0$ and subtracting yields
$-B+D(1-2(\mu+\tau))=0$. A further cusp difference forces $D=0$,
and then $B=0$; nonconstancy of $\mu$ forces $A=C=0$.
Comparison of coefficients now gives
$F_{34}=F_{12}=0$, $F_{13}=-6F_{24}$, and
$6F_{23}-F_{14}=12F_{24}$, exactly
$F\in\mathbb Z\eta+\mathbb ZE$. Every other integral class therefore
has a discrete zero set along this period locus. Excluding their
countable union gives the asserted exact Neron--Severi lattice.
It contains the nonzero isotropic $E$ and no matrix of negative
Pfaffian; \ref{fam:src:D.1.6} gives algebraic dimension one.
Finally
$c=-5\tau-12\mu-(\beta^{\rm part}+\tau)$ has imaginary part tending
to minus infinity at the cusp. A sufficiently small cusp neighbourhood
is therefore admissible by \ref{fam:src:C.2.1}. The torsion locus of $P$ is countable
by \ref{fam:src:D.1.2}, so its exclusion still leaves such points in every open subset
of that neighbourhood. This supplies actual members and fibres, with
all their parameter functions specified.
\end{enumerate}
\end{proof}

\subsection{Compact curves and surfaces}\label{fam:section:D.2}
\begin{proposition}\label{fam:src:D.2.1}\label{fam:st:D.2.1}
\StatementRecord{fam:src:D.2.1}{proved}{Historical events: Claude Fable 5.2 (Anthropic), effort high, 12 November 2026: targeted proof check. Claude Fable 5.2 (Anthropic), effort high, 16 November 2026: used and rechecked in this run. GPT-6 Astra (OpenAI), ChatGPT Pro mode, effort Pro, S6\_09\_PAPERS, 17 November 2026: targeted proof check. Also checked in this writing. Pre-continuation writing events: GPT-6 Astra (OpenAI), ChatGPT Pro mode, effort Pro, S6\_10\_PAPERS, 18 November 2026: targeted proof check. Continuation metadata events: GPT-6 Astra (OpenAI), ChatGPT Pro mode, effort Pro, S6\_10\_PAPERS continuation, 19 November 2026: used and rechecked in this run; status-vocabulary reconciliation; metadata reconciliation only; no new mathematical check. \emph{Pre-continuation recorded proof scope.} Read the complete source proof and re-derived the rank-three lattice/kernel and closure, rational product covering, dense cusp natural-boundary obstruction, exact toric gluing/dihedral action, quotient and Pfaffian criteria, explicit non-torsion algebraic-dimension-one locus, and finite-normalization curve maps as applicable. Corrected the Weierstrass map direction and sign in D.1.1 by the already proved Phi; made the rational product identification in C.3.2 explicit. Exact background citations and the indicated finite symbolic programs remain separate gates.}{the hypotheses, notation and definitions stated here; the explicit constructions and calculations in the proof}

For a standard-twist member \(X_c\) with underlying manifold \(S^6\),
every compact irreducible complex curve lies in one fibre of \(f\), and
every compact irreducible complex surface is one of the reduced fibres.
\end{proposition}


\begin{proof}
\begin{enumerate}[label=\arabic*.,leftmargin=*]
\item The given diffeomorphism with \(S^6\) implies \(H^2(X_c;\mathbb Z)=0\).
For a compact irreducible curve \(C\), let
\(\nu:\widetilde C\to C\hookrightarrow X_c\) be its normalization map.
The holomorphic map \(f\nu\) has
\(\deg (f\nu)^*\mathcal O_{\mathbb P^1}(1)=0\), since its Chern class is
pulled back from the zero class on \(X_c\). A nonconstant holomorphic map
of compact Riemann surfaces has positive degree: a regular value has a
positive number of preimages, all with positive local multiplicity.
Therefore \(f\nu\), and hence \(f|_C\), is constant.

\item An irreducible surface \(Z\) in the smooth threefold is an effective
Cartier divisor. If its image were all of \(B\), a general smooth fibre
would meet it in a nonempty effective divisor \(D\) whose line bundle
is \(\mathcal O_{X_c}(Z)|_{F_t}\). Its integral Chern class is zero.
Choose a flat K\"ahler form \(\omega\) on the torus \(F_t\). The equality
\(\int_{F_t}c_1(D)\wedge\omega=\int_D\omega\) has left side zero and right
side positive, a contradiction. Thus \(Z\) lies in a single fibre.
Each reduced fibre is irreducible of dimension two, so containment forces
equality. This uses the full map into \(X_c\), not a transfer of a
vanishing claim from its normalization.
\end{enumerate}
\end{proof}

\begin{proposition}\label{fam:src:D.2.2}\label{fam:st:D.2.2}
\StatementRecord{fam:src:D.2.2}{proved}{Historical events: Claude Fable 5.2 (Anthropic), effort high, 12 November 2026: targeted proof check. Claude Fable 5.2 (Anthropic), effort high, 16 November 2026: used and rechecked in this run. GPT-6 Astra (OpenAI), ChatGPT Pro mode, effort Pro, S6\_09\_PAPERS, 17 November 2026: targeted proof check. Also checked in this writing. Pre-continuation writing events: GPT-6 Astra (OpenAI), ChatGPT Pro mode, effort Pro, S6\_10\_PAPERS, 18 November 2026: targeted proof check. Continuation metadata events: GPT-6 Astra (OpenAI), ChatGPT Pro mode, effort Pro, S6\_10\_PAPERS continuation, 19 November 2026: used and rechecked in this run; status-vocabulary reconciliation; metadata reconciliation only; no new mathematical check. \emph{Pre-continuation recorded proof scope.} Read the complete source proof and re-derived the rank-three lattice/kernel and closure, rational product covering, dense cusp natural-boundary obstruction, exact toric gluing/dihedral action, quotient and Pfaffian criteria, explicit non-torsion algebraic-dimension-one locus, and finite-normalization curve maps as applicable. Corrected the Weierstrass map direction and sign in D.1.1 by the already proved Phi; made the rational product identification in C.3.2 explicit. Exact background citations and the indicated finite symbolic programs remain separate gates.}{the hypotheses, notation and definitions stated here; \cref{fam:src:D.2.1}; \cref{fam:src:D.1.1}; \cref{fam:src:C.4.1}}

There is no rational curve in \(X_c\setminus W_c\). There are elliptic
curves outside all three special fibres, in particular in each of the
two smooth fibres \(F_{t_\pm}\). All curves in \(W_c\) are precisely the
images of irreducible algebraic curves in \(\mathrm{dP}_6\) under its
finite normalization map. Distinct curves have the same image only for
the three paired boundary lines. All curves in a bielliptic fibre
\(S_j\) are images of irreducible curves in its finite product cover.
\end{proposition}


\begin{proof}
\begin{enumerate}[label=\arabic*.,leftmargin=*]
\item By \ref{fam:src:D.2.1} every curve is vertical. A holomorphic map from
\(\mathbb P^1\) into a torus lifts to its vector-space universal cover,
since \(\mathbb P^1\) is simply connected, and its coordinate functions
are constant by compactness. Thus a smooth torus fibre contains no
rational curve. A map from \(\mathbb P^1\) into a bielliptic surface
lifts through its finite unramified torus cover by the same covering
criterion; the lifted map is constant. These arguments prove the first
assertion, including possibly singular rational images.
\ref{fam:src:D.1.1} gives an embedded connected elliptic kernel in each \(F_{t_\pm}\),
proving the second.

\item For the third assertion let \(\eta:\mathrm{dP}_6\to W_c\) be the finite
normalization morphism. Each component of the inverse image of a compact
curve is compact and one-dimensional and maps onto that curve, because
finite maps have zero-dimensional fibres. It is algebraic in the
projective surface \(\mathrm{dP}_6\). Conversely its image is a compact
analytic curve by properness. Away from the conductor, \(\eta\) is a
biholomorphism, so two distinct curves can have the same image only when
both lie in the conductor; its identifications are exactly the three
opposite pairs computed in \ref{fam:src:C.4.1}. This is also an explicit algebraic
parametrization: the curves of \(\mathrm{dP}_6\) are the three exceptional
lines and strict transforms of irreducible plane curves under the blow-up
of the three coordinate points. Their homogeneous equations and their
three multiplicities specify the strict transforms. The finite-product
cover statement for \(S_j\) follows from the identical proper-image and
finite-preimage argument applied to its unramified morphism
\(E_j\times F_j\to S_j\); curves in the product are algebraic.
\end{enumerate}
\end{proof}

\begin{theorem}\label{fam:src:D.2.3}\label{fam:st:D.2.3}
\StatementRecord{fam:src:D.2.3}{proved}{Historical events: Claude Fable 5.2 (Anthropic), effort high, 12 November 2026: targeted proof check. Claude Fable 5.2 (Anthropic), effort high, 16 November 2026: used and rechecked in this run. GPT-6 Astra (OpenAI), ChatGPT Pro mode, effort Pro, S6\_09\_PAPERS, 17 November 2026: targeted proof check. Also checked in this writing. Pre-continuation writing events: GPT-6 Astra (OpenAI), ChatGPT Pro mode, effort Pro, S6\_10\_PAPERS, 18 November 2026: targeted proof check. Continuation metadata events: GPT-6 Astra (OpenAI), ChatGPT Pro mode, effort Pro, S6\_10\_PAPERS continuation, 19 November 2026: used and rechecked in this run; status-vocabulary reconciliation; metadata reconciliation only; no new mathematical check. \emph{Pre-continuation recorded proof scope.} Read and re-derived complete curve classification: theta automorphy and cocycle on actual lattice, rank-one elliptic fibres, finite product-cover inverse images, and cusp normalization images. Matched Schnell Theorem4.1 and equation4.3 pp1-2; full application printed.}{\cref{fam:src:D.1.6}, \cref{fam:src:D.2.1}, \cref{fam:src:D.2.2}; \cite[Theorem 4.1 and equation (4.3), pp.~1--2]{Schnell19}}

The following descriptions give all compact irreducible curves of $X_c$.
In $W_c$ and $S_j$ they are exactly the images described in \ref{fam:src:D.2.2}.
In a smooth fibre of algebraic dimension zero there are none. In a
smooth fibre of algebraic dimension one they are exactly the translates
of its unique elliptic kernel $K_t$. In an abelian fibre they are exactly
the translates of elliptic subtori, together with the irreducible zero
divisors of the positive Appell--Humbert theta sections described below.
These descriptions, \ref{fam:src:D.1.6}, and \ref{fam:src:D.2.1} constitute a classification including
every smooth exceptional fibre and every compact irreducible surface.
\end{theorem}

\begin{proof}
\begin{enumerate}[label=\arabic*.,leftmargin=*]
\item \label{app:rederive:fam-src-D.2.3}
All curves are vertical by \ref{fam:src:D.2.1}, so the list of fibre types is exhaustive.
The dimension-zero and dimension-one assertions, including uniqueness,
were proved in \ref{fam:src:D.1.6}. Let now $F_t=\mathbb C^2/\Pi\Lambda$ be abelian.
For every positive definite Hermitian form $H$ whose imaginary part takes
integer values on $\Pi\Lambda$, and every function
$\alpha:\Pi\Lambda\to U(1)$ satisfying
\[
 \alpha(\lambda+\nu)=\alpha(\lambda)\alpha(\nu)
                 e^{\pi i\operatorname{Im}H(\lambda,\nu)},
\]
take the nonzero entire solutions $\theta:\mathbb C^2\to\mathbb C$ of
\begin{equation}\label{fam:eq:cd-theta-sections}
 \theta(\zeta+\lambda)=\alpha(\lambda)
 e^{\pi H(\zeta,\lambda)+\pi H(\lambda,\lambda)/2}\theta(\zeta).
\end{equation}
The cocycle identity follows by multiplying the displayed factors; its
integral imaginary cross-term has exponential one. It defines a
holomorphic line bundle by quotienting $\mathbb C^2\times\mathbb C$.
The invariant zero set of $\theta$ descends to a compact divisor on
$F_t$. Retain precisely those zero divisors that are reduced and
irreducible. This is an exact parametrization, with all periods and
multiplier equations given; multiplying $\theta$ by a nonzero scalar
does not change its divisor.

\item The correspondence \cite[Theorem 4.1, p.~1, and equation (4.3), p.~2]{Schnell19} proves
that every holomorphic line bundle and every section on this actual
quotient has this form. For the cohomology matrix the exact comparison is
$H_E=\widetilde E(\,\cdot\,,J\,\cdot\,)-i\widetilde E$ from \ref{fam:src:D.1.6};
we choose the sign of $E$ making $H_E$ positive. Its imaginary part is
$-\widetilde E$, which gives the integral semicharacter equation displayed
above. A divisor class with positive Hermitian rank two gives exactly
these positive theta sections. A divisor of Hermitian rank one is pulled
back from its elliptic quotient and, when irreducible and reduced, is a
single fibre. Thus it is a translate of an elliptic subtorus. The null
lattice is computed by the explicit integer kernel in \ref{fam:src:D.1.6}. No other
rank is possible for an effective nonempty divisor, by positivity of its
averaged current. This proves exhaustiveness on the abelian fibre.

\item The same equations give the curves in the product covers of $S_j$:
compute them on its two-elliptic-curve lattice and take images under the
specified finite morphism. Conversely each curve downstairs is the image
of a component of its finite inverse image, so that this procedure omits
none. In $W_c$ the analogous procedure uses the homogeneous polynomial
parametrization of the normalization in \ref{fam:src:D.2.2}, with the three boundary
pairings already identified. The maps used throughout are finite or
closed embeddings, so they preserve compactness and give one-dimensional
images; no assertion is transferred from a normalization in an invalid
direction.
\end{enumerate}
\end{proof}

\Uses{the hypotheses, notation and definitions stated here; \cref{fam:src:D.2.2}; \cref{fam:src:D.1.6}; \cref{fam:src:D.2.1}; \cref{app:rederive:fam-src-D.2.3}}


This is the effective line-bundle criterion for compact complex tori in the present lattice. A complete re-derivation is given in Appendix~\ref{app:rederive:fam-src-D.2.3}.




\subsection{The rank-three period quotient}\label{fam:section:C.3}
\begin{proposition}\label{fam:src:C.3.1}\label{fam:st:C.3.1}
\StatementRecord{fam:src:C.3.1}{proved}{Historical events: Claude Fable 5.2 (Anthropic), effort high, 12 November 2026: targeted proof check. GPT-6 Astra (OpenAI), ChatGPT Pro mode, effort Pro, S6\_09\_PAPERS, 17 November 2026: targeted proof check. Also checked in this writing. Pre-continuation writing events: GPT-6 Astra (OpenAI), ChatGPT Pro mode, effort Pro, S6\_10\_PAPERS, 18 November 2026: targeted proof check. Continuation metadata events: GPT-6 Astra (OpenAI), ChatGPT Pro mode, effort Pro, S6\_10\_PAPERS continuation, 19 November 2026: used and rechecked in this run; status-vocabulary reconciliation; metadata reconciliation only; no new mathematical check. \emph{Pre-continuation recorded proof scope.} Read the complete source proof and re-derived the rank-three lattice/kernel and closure, rational product covering, dense cusp natural-boundary obstruction, exact toric gluing/dihedral action, quotient and Pfaffian criteria, explicit non-torsion algebraic-dimension-one locus, and finite-normalization curve maps as applicable. Corrected the Weierstrass map direction and sign in D.1.1 by the already proved Phi; made the rational product identification in C.3.2 explicit. Exact background citations and the indicated finite symbolic programs remain separate gates.}{the hypotheses, notation and definitions stated here; the explicit constructions and calculations in the proof}

At every point with \(D=0\), \(T>0\), the real-linear map
\(\Pi:\mathbb R^4\to\mathbb C^2\) has rank three. Write
\(\mu=a_0\tau+b_0\), \(a_0,b_0\in\mathbb R\), and put
\[
 k=(1,-6a_0,-6b_0,6a_0\operatorname{Re}\mu-\operatorname{Re}\beta).
\]
Its real kernel is \(\mathbb R k\). The subgroup \(\Pi\mathbb Z^4\) is
discrete exactly when \(k\in\mathbb Q^4\). In that case it is a
rank-three lattice in a real hyperplane and its quotient in \(\mathbb C^2\)
is a noncompact complex Lie group. In every other case the topological
quotient \(\mathbb C^2/\Pi\mathbb Z^4\) is non-Hausdorff.
\end{proposition}


\begin{proof}
\begin{enumerate}[label=\arabic*.,leftmargin=*]
\item The real matrix, with output order
\((\operatorname{Re}\zeta_1,\operatorname{Im}\zeta_1,
\operatorname{Re}\zeta_2,\operatorname{Im}\zeta_2)\), is
\[
 \begin{pmatrix}6\operatorname{Re}\mu&\operatorname{Re}\tau&1&0\\
 6\operatorname{Im}\mu&T&0&0\\
 \operatorname{Re}\beta&\operatorname{Re}\mu&0&1\\
 \operatorname{Im}\beta&\operatorname{Im}\mu&0&0\end{pmatrix}.
\]
Its last three columns are real-linearly independent because \(T>0\).
Its determinant is \(TD\), so at zero determinant the rank is exactly
three. Substitution gives the displayed kernel, including its signs.
If the image subgroup is discrete, it is a lattice of rank at most three
in its real span. The homomorphism \(\mathbb Z^4\to\Pi\mathbb Z^4\)
therefore has a nonzero kernel, forcing the displayed kernel vector with
first coordinate one to be rational. Conversely a rational line has a
primitive integral generator extendable to a basis of \(\mathbb Z^4\).
The remaining three basis images are a real basis of the image hyperplane
and hence generate a discrete lattice. Its quotient leaves one unbounded
real direction. Finally a non-discrete finitely generated subgroup of
this real hyperplane cannot be closed: a closed subgroup is a product of
a real vector group and a discrete group, and a nonzero real vector group
is uncountable. A quotient topological group is Hausdorff exactly when its
identity subgroup is closed. This proves the last assertion.
\end{enumerate}
\end{proof}

\begin{proposition}\label{fam:src:C.3.2}\label{fam:st:C.3.2}
\StatementRecord{fam:src:C.3.2}{proved}{Historical events: Claude Fable 5.2 (Anthropic), effort high, 12 November 2026: targeted proof check. GPT-6 Astra (OpenAI), ChatGPT Pro mode, effort Pro, S6\_09\_PAPERS, 17 November 2026: targeted proof check. Also checked in this writing. Pre-continuation writing events: GPT-6 Astra (OpenAI), ChatGPT Pro mode, effort Pro, S6\_10\_PAPERS, 18 November 2026: targeted proof check. Continuation metadata events: GPT-6 Astra (OpenAI), ChatGPT Pro mode, effort Pro, S6\_10\_PAPERS continuation, 19 November 2026: used and rechecked in this run; status-vocabulary reconciliation; metadata reconciliation only; no new mathematical check. \emph{Pre-continuation recorded proof scope.} Read the complete source proof and re-derived the rank-three lattice/kernel and closure, rational product covering, dense cusp natural-boundary obstruction, exact toric gluing/dihedral action, quotient and Pfaffian criteria, explicit non-torsion algebraic-dimension-one locus, and finite-normalization curve maps as applicable. Corrected the Weierstrass map direction and sign in D.1.1 by the already proved Phi; made the rational product identification in C.3.2 explicit. Exact background citations and the indicated finite symbolic programs remain separate gates.}{the hypotheses, notation and definitions stated here; \cref{fam:src:C.1.1}; \cref{fam:src:C.2.3}}

The unique degenerate base point at the actual admissibility boundary is
\(t=1\). The boundary constants
\(c=x+i(3/\pi)\log6\) give
\[
 \ker_{\mathbb R}\Pi=\mathbb R(1,3,-3,-x-3/2).
\]
For rational $x+3/2=m/n$ in lowest terms, the quotient is a
finite holomorphic quotient of $E_i\times\mathbb C^*$, with a covering
of degree four. For irrational \(x+3/2\), the closure of the period subgroup is generated
by the real line \(\mathbb R(0,1)\) and the two periods
\((i,(1-i)/2)\), \((1,0)\). Its Hausdorff quotient is a real manifold
homeomorphic to \(\mathbb R\times(S^1)^2\).
\end{proposition}


\begin{proof}
\begin{enumerate}[label=\arabic*.,leftmargin=*]
\item \ref{fam:src:C.1.1} gives \(\operatorname{Im}\beta=3/2\) at this value of \(c\).
The vector \(\widehat\gamma+3\widehat u-3\widehat w\) maps to
\((0,x+3/2)\), while \(\widehat\delta\) maps to \((0,1)\).
For $x+3/2=m/n$ in lowest terms, change coordinates by
$(\zeta_1,\zeta_2)\mapsto(\zeta_1,\zeta_2+\zeta_1/2)$.
The vectors $2\widehat u-\widehat\delta$ and
$2\widehat w-\widehat\delta$ then have periods $(2i,0)$ and $(2,0)$.
Bezout's identity applied to $\widehat\gamma+3\widehat u-3\widehat w$
and $\widehat\delta$ gives the period $(0,1/n)$.
These three periods generate a sublattice of index four in the rank-three
image: its two remaining generators are represented by
$(i,1/2)$ and $(1,1/2)$, each of order two modulo this sublattice.
The product identification is the explicit map
$([\zeta_1],[\zeta_2])\mapsto([\zeta_1/2],e^{2\pi in\zeta_2})$
from $\mathbb C^2$ modulo these three periods to $E_i\times\mathbb C^*$. Lattice inclusion gives the claimed degree-four map.
For irrational $x+3/2$, its additive subgroup with $1$ is dense in the
indicated real line. The other two displayed periods are independent
modulo that line. This proves the closure and the quotient description.
By \ref{fam:src:C.2.3} this is the only degenerate torus cover at the admissibility
boundary; every other finite-base period matrix retains real rank four.
We have described the torus cover of the order-four filling, including
the actual quotient and its topology.
\end{enumerate}
\end{proof}


\paragraph{Local boundary obstruction}\label{fam:src:C.3.3}\label{fam:st:C.3.3}
\StatementRecord{fam:src:C.3.3}{proved}{Historical events: Claude Fable 5.2 (Anthropic), effort high, 12 November 2026: targeted proof check. Claude Fable 5.2 (Anthropic), effort high, 17 November 2026: carried as background without fault found. GPT-6 Astra (OpenAI), ChatGPT Pro mode, effort Pro, S6\_09\_PAPERS, 17 November 2026: targeted proof check. Also checked in this writing. Pre-continuation writing events: GPT-6 Astra (OpenAI), ChatGPT Pro mode, effort Pro, S6\_10\_PAPERS, 18 November 2026: targeted proof check. Continuation metadata events: GPT-6 Astra (OpenAI), ChatGPT Pro mode, effort Pro, S6\_10\_PAPERS continuation, 19 November 2026: used and rechecked in this run; status-vocabulary reconciliation; metadata reconciliation only; no new mathematical check. \emph{Pre-continuation recorded proof scope.} Read the complete source proof and re-derived the rank-three lattice/kernel and closure, rational product covering, dense cusp natural-boundary obstruction, exact toric gluing/dihedral action, quotient and Pfaffian criteria, explicit non-torsion algebraic-dimension-one locus, and finite-normalization curve maps as applicable. Corrected the Weierstrass map direction and sign in D.1.1 by the already proved Phi; made the rational product identification in C.3.2 explicit. Exact background citations and the indicated finite symbolic programs remain separate gates.}{the hypotheses, notation and definitions stated here; \cref{fam:src:C.3.2}; \cref{fam:src:C.3.4}; \cref{push:src:D.1.2}, \cref{push:src:D.1.3}}

The local compactification obstruction at the actual boundary is localized at the order-four fibre over $t=1$. The proposed extension would require a compact complex-space replacement of the rank-three
quotient in \ref{fam:src:C.3.2}, over a neighbourhood of one boundary parameter,
with a proof that it extends the original cyclic filling and its overlap
with the ordinary torus family. Both rational and irrational values of
$x+3/2$ are included in the exact quotient description. The naive
quotients cannot be central fibres of a proper complex-space family:
for rational $x+3/2$ they are noncompact, and for irrational values
they are non-Hausdorff. \ref{fam:src:C.3.4} supplies a further obstruction to preserving
a holomorphic Weierstrass Albanese model in this parameter.
These prove failure of the stated specific replacements. The unique base-fibre location and exact local period obstruction are thereby determined. The general proper holomorphic extension is excluded by \cref{push:src:D.1.2}, \cref{push:src:D.1.3}, which recovers the intrinsic divisor in a hypothetical extension.





\begin{proposition}\label{fam:src:C.3.4}\label{fam:st:C.3.4}
\StatementRecord{fam:src:C.3.4}{proved}{Historical events: Claude Fable 5.2 (Anthropic), effort high, 12 November 2026: targeted proof check. GPT-6 Astra (OpenAI), ChatGPT Pro mode, effort Pro, S6\_09\_PAPERS, 17 November 2026: targeted proof check. Also checked in this writing. Pre-continuation writing events: GPT-6 Astra (OpenAI), ChatGPT Pro mode, effort Pro, S6\_10\_PAPERS, 18 November 2026: targeted proof check. Continuation metadata events: GPT-6 Astra (OpenAI), ChatGPT Pro mode, effort Pro, S6\_10\_PAPERS continuation, 19 November 2026: used and rechecked in this run; status-vocabulary reconciliation; metadata reconciliation only; no new mathematical check. \emph{Pre-continuation recorded proof scope.} Read the complete source proof and re-derived the rank-three lattice/kernel and closure, rational product covering, dense cusp natural-boundary obstruction, exact toric gluing/dihedral action, quotient and Pfaffian criteria, explicit non-torsion algebraic-dimension-one locus, and finite-normalization curve maps as applicable. Corrected the Weierstrass map direction and sign in D.1.1 by the already proved Phi; made the rational product identification in C.3.2 explicit. Exact background citations and the indicated finite symbolic programs remain separate gates.}{the hypotheses, notation and definitions stated here; \cref{fam:section:A.4}; \cref{fam:src:C.1.1}}

On its half-plane of definition, the $j$-invariant of
$\operatorname{Alb}(S_2)$ has no meromorphic continuation across any
point of the line
$\operatorname{Im}c=(3/\pi)\log6$. In particular the rank-three quotient
at $t=1$ cannot be replaced by a Weierstrass family of elliptic Albanese
curves whose coefficients are holomorphic in the same parameter across
that line.
\end{proposition}


\begin{proof}
\begin{enumerate}[label=\arabic*.,leftmargin=*]
\item The explicit Albanese lattice in \ref{fam:section:A.4} is
$\mathbb Z+(\kappa_2/2)\mathbb Z$, with
$\kappa_2=c+3/2-i(3/\pi)\log6$ by \ref{fam:src:C.1.1}. A positive-imaginary modulus is
$z=-\kappa_2/2$, so the invariant is $j(z)$. At a reduced rational
$r=a/b$ on the real boundary, choose an integral determinant-one modular
matrix with bottom row $(b,-a)$. Its value at $r+iy$ has imaginary part
$1/(b^2y)$. The convergent Fourier expansion $j=q^{-1}+744+O(q)$ therefore
gives growth of order $\exp(2\pi/(b^2y))$ as $y\downarrow0$.
A meromorphic function at $r$ is bounded by a fixed power of $1/y$,
so continuation there is impossible. A meromorphic continuation near an
irrational real point would also continue the same function near some
rational real point in that neighbourhood, which has just been excluded.
The affine change from $c$ to $z$ transfers this obstruction to the
stated line. A Weierstrass family with holomorphic coefficients has a
meromorphic $j$-invariant, expressed as the quotient of its cubic
invariants, so it cannot supply the indicated continuation. This
obstruction concerns the stated Albanese construction; it does not
silently assume that every possible compact threefold filling preserves
an Albanese model of its special surface.
\end{enumerate}
\end{proof}

\subsection{Classification of the non-normal cusp surface}\label{fam:section:C.4}
\begin{theorem}\label{fam:src:C.4.1}\label{fam:st:C.4.1}
\StatementRecord{fam:src:C.4.1}{proved}{Historical events: Claude Fable 5.2 (Anthropic), effort high, 12 November 2026: targeted proof check. GPT-6 Astra (OpenAI), ChatGPT Pro mode, effort Pro, S6\_09\_PAPERS, 17 November 2026: targeted proof check. Also checked in this writing. Pre-continuation writing events: GPT-6 Astra (OpenAI), ChatGPT Pro mode, effort Pro, S6\_10\_PAPERS, 18 November 2026: targeted proof check. Continuation metadata events: GPT-6 Astra (OpenAI), ChatGPT Pro mode, effort Pro, S6\_10\_PAPERS continuation, 19 November 2026: used and rechecked in this run; status-vocabulary reconciliation; metadata reconciliation only; no new mathematical check. \emph{Pre-continuation recorded proof scope.} Read the complete source proof and re-derived the rank-three lattice/kernel and closure, rational product covering, dense cusp natural-boundary obstruction, exact toric gluing/dihedral action, quotient and Pfaffian criteria, explicit non-torsion algebraic-dimension-one locus, and finite-normalization curve maps as applicable. Corrected the Weierstrass map direction and sign in D.1.1 by the already proved Phi; made the rational product identification in C.3.2 explicit. Exact background citations and the indicated finite symbolic programs remain separate gates.}{the hypotheses, notation and definitions stated here; \cref{fam:src:C.1.1}}

For the normalization of \ref{fam:src:C.1.1}, the three opposite-side identification
constants of \(W_c\), in the characters specified below, are
\[
 (g_1,g_2,g_3)=(-a^{-1},a^6,-a^{-7}e^{-2\pi ic}).
\]
Two such compact complex surfaces are biholomorphic exactly when
\(c'-c\in\mathbb Z\). The vertical \(\mathbb C^*\)-action inherited from
translation in the \(\widehat\delta\)-direction fixes exactly the double
curve represented by the pair of edges of direction \(e_2\).
\end{theorem}


\begin{proof}
\begin{enumerate}[label=\arabic*.,leftmargin=*]
\item Keep the lattice of toric cocharacters \(\mathbb Z^2\), with
\(e_1=\widehat w\), \(e_2=\widehat\delta\), and cusp logarithm
\(s=(2\pi i)^{-1}\log t_c\). The cusp period matrix is
\([sB_0+C(t_c)\mid I]\), where
\[
 B_0=\begin{pmatrix}0&1\\-1&0\end{pmatrix},\qquad
 C(0)=\begin{pmatrix}6\mu_0&h_0\\b_0-h_0&\mu_0\end{pmatrix},
 \quad b_0=c.
\]
For \(v=e_1,e_2,e_1-e_2\), the gluing map from the side of direction
\(v\) to its opposite is the restriction of the toric action with lattice
label \(B_0v\). A character row \(r\) with \(rv=0\) pulls back by this
map to itself multiplied by
\(\exp(2\pi i\,rC(0)B_0v)\): the remaining power of \(t_c\) is
\(rB_0^2v=-rv=0\). Taking respectively
\(r=(0,1),(1,0),(1,1)\) gives exponents
\(-\mu_0,6\mu_0,-7\mu_0-c\). This proves the three constants, including
the cancellation of \(h_0\). These are morphisms of the whole projective
boundary lines, since multiplication by a nonzero constant extends over
both endpoints. The restrictions agree at both triple points, as the
ambient toric action supplies the same vertex identifications.

\item A biholomorphism of the two surfaces lifts uniquely to a biholomorphism
of their normalizations: integral closure of the reduced local rings is
functorial under an isomorphism. It preserves the inverse image of the
singular locus, the hexagon on \(\mathrm{dP}_6\). Its permutation of the
six components preserves their cyclic incidence, hence lies in the
order-twelve dihedral group. After composition with the corresponding
toric lattice automorphism, it fixes each boundary component. Each torus
character has its divisor supported on those components, so its pullback
divided by itself has divisor zero on the compact normalization and is
constant. Such an automorphism is therefore multiplication by an element
of \((\mathbb C^*)^2\). These torus multipliers cancel in the opposite-side
comparison because the same character is used on both sides.

\item The remaining six effective actions on the triple, with the central
lattice automorphism \(-I\) acting trivially, are described by the equations
\(g'_{k_j}{}^{\epsilon_j}=g_j\), where
\[
 ((k_1,\epsilon_1),(k_2,\epsilon_2),(k_3,\epsilon_3))
 \in\left\{
\begin{array}{lll}
((1,+),(2,+),(3,+)),&((1,-),(3,+),(2,+)),\\
((2,-),(3,-),(1,+)),&((2,+),(1,+),(3,-)),\\
((3,-),(2,-),(1,-)),&((3,+),(1,-),(2,-)).
\end{array}\right.
\]
This finite calculation is reproduced by \texttt{cd\_exact.py}, enumerating
all integral automorphisms of the six rays and transporting the characters.
Every nonidentity action requires one of
\(g_1^{-1}=g_1\), \(g_2^{-1}=g_2\), or equality of a positive real number
\(g_2^{\pm1}\) with a negative real number \(g_1^{\pm1}\).
All fail for \(a>1\). Therefore \(g'_3=g_3\), precisely
\(e^{-2\pi ic'}=e^{-2\pi ic}\). Conversely equality of the three constants
gives the identity gluing morphism on the normalization and hence a
biholomorphism of the quotients. This proves the unmarked assertion.

\item For the action of the subtorus with cocharacter \(e_2\), a one-dimensional
toric orbit on the normalization is pointwise fixed exactly when its
stabilizer line contains \(e_2\). Among the six rays this occurs exactly
for \(\pm e_2\). Their images form one double curve; all zero-dimensional
orbits map to its two endpoints. No dense two-dimensional orbit is fixed.
The quotient action has the same fixed set: a continuous orbit cannot
jump between the finitely many branches of a normalization fibre.
\end{enumerate}
\end{proof}


\section{Real structures and their fixed three-spheres}\label{mapsI:real-forms}

We use the family marking
$\Lambda=\mathbb Z\langle\widehat\gamma,\widehat u,\widehat w,
\widehat\delta\rangle$ and its dual covector $\gamma$.
In this section $d=\widehat\delta$, $g=\widehat\gamma$,
$u=\widehat u$, and $w_{\Lambda}=\widehat w$; an unadorned $w$ in a
lattice expression denotes $w_{\Lambda}$, whereas the scalar parameter is
$w=\exp(-\pi i c)$. Put
\[
 e=\Pi_c d=(0,1)^t,\qquad E=d\otimes\gamma,
 \qquad \Pi_c=
 \begin{pmatrix}6\mu&\tau&1&0\\\beta^{\rm part}+c&\mu&0&1\end{pmatrix},
\]
\[
 D_c=\operatorname{Im}(\beta^{\rm part}+c)
       -6(\operatorname{Im}\mu)^2/\operatorname{Im}\tau<0.
\]
Thus $d$ is a lattice vector and $e$ is its image in the complex translation
space. The circle $\rho(s)$ is translation by $se$, $s\in\mathbb R/\mathbb Z$.
The clockwise root generators and the affine vectors are
\[
 A_1=\begin{pmatrix}1&0&0&0\\6&0&1&0\\-6&-1&-1&0\\-2&1&0&1\end{pmatrix},\quad
 A_2=\begin{pmatrix}1&0&0&0\\0&0&-1&0\\-6&1&0&0\\3&0&1&1\end{pmatrix},
\]
\[
 M_0=(A_1A_2)^{-1},\qquad
 v_1=(1,2,-4,0)^t,\qquad v_2=(-1,-3,3,0)^t.
\]
These are the same conventions as in the family construction. In cusp
coordinates $q=1/t$ and $s=(2\pi i)^{-1}\log q$, the bounded matrix $C_c(q)$
is defined by $Z_c=sB_0+C_c(q)$, with
$B_0=\left(\begin{smallmatrix}0&1\\-1&0\end{smallmatrix}\right)$.
In particular $\Pi_cg=(6\mu,-s+C_{c,21})^t$.
We write $D_1,D_2,D_{12}$ for the three cusp double curves belonging to
edges in directions $e_1,e_2,e_1-e_2$, respectively. Thus $D_2$ is the
fixed sphere of the circle action $\rho$, whereas $D_1$ is the double
curve along which the literal character has a limit. This distinction
is essential in the real calculation as well.

\begin{proposition}[Two conjugacy classes of real structures]\label{mapsI:src:A.3.1}
\StatementRecord{mapsI:src:A.3.1}{proved}{Historical events: Claude Fable 5.2 (Anthropic), effort high, 16 November 2026: targeted proof check; used and rechecked in this run. GPT-6 Astra (OpenAI), ChatGPT Pro mode, effort Pro, S6\_09\_PAPERS, 17 November 2026: targeted proof check. Also checked in this writing. Pre-continuation writing events: GPT-6 Astra (OpenAI), ChatGPT Pro mode, effort Pro, S6\_10\_PAPERS, 18 November 2026: computation or formal certificate; targeted proof check; used and rechecked in this run. Continuation metadata events: GPT-6 Astra (OpenAI), ChatGPT Pro mode, effort Pro, S6\_10\_PAPERS continuation, 19 November 2026: used and rechecked in this run; status-vocabulary reconciliation; metadata reconciliation only; no new mathematical check. \emph{Pre-continuation recorded proof scope.} Read complete proof and re-derived invariant character, finite solid-torus coordinates, cusp nonextension paths, exact lens meridians, relative H1 extension, affine real involutions/cusp units, real cap-meridian determinant and orbit-space recognition. Full finite-congruence replays and Smith/triangulation external locators remain separate gates. Reran real\_involutions.py and real\_finite\_four.py from copied source. Exact translation/Hermite and Smith calculations and all denominator-24 affine congruences pass, including all fixed-lift parities and the two-torus/empty outcomes. This supplements the separately recorded full geometric proof review. Opened Illman Theorem3.1 printedp48 visually and read the following proof opening; read Knudson Theorem4.1 and full proof p10 and Theorem4.2 and full proof p11. The actual smooth compact finite-involution manifold has finite-dimensional equivariant cells and mod2 coefficients required by Smith. These named external gates close; no entire Illman proof review is claimed.}{the hypotheses, notation and definitions stated here; \cref{fam:src:B.4.2}; \cref{fam:src:B.4.1}; \cite[Theorem~3.1, p.~48]{Illman1972}; \cite[Theorem 4.1(1)]{Knudson2004}}

For a real member the antiholomorphic involutions have exactly two
conjugacy classes under holomorphic automorphisms. For any fixed
involutive splitting $\sigma$ they are represented by $\sigma$ and
$\rho(1/2)\sigma$. Every nonempty fixed set is a closed connected
orientable totally real three-manifold with the mod-two homology of
$S^3$. The circle $\rho$ preserves each such fixed set and acts
semifreely on it.
\end{proposition}



\begin{proof}
\begin{enumerate}[label=\arabic*.,leftmargin=*]
\item The real comparison of Proposition~\ref{fam:src:B.4.2} has derivative complex conjugation in the
second period coordinate. It therefore conjugates translation by
$b\in\C/\Z$ to translation by $\bar b$. In multiplicative notation its
action is $z\mapsto1/\bar z$. The square of $z\sigma$ is $z/\bar z$,
so it is an involution exactly for $z\in\R^*$. Conjugation by $y\in\C^*$
changes this coefficient to $|y|^2z$, giving exactly its two signs.
The automorphism computation used here is Proposition~\ref{fam:src:B.4.1}: the integral centralizer has linear part $aI$, the cusp and
finite norm equations force $a=1$ for a self-map, and the global
translation sheaf has $H^0=\C e$ with integral periods $\Z e$.
These are precisely the three steps needed to identify the full group
with $\C^*$, rather than only its identity component.

\item At a fixed point an antiholomorphic involution has a real three-dimensional
fixed tangent space, and its complexification is the full tangent space.
An invariant metric and its exponential map identify its fixed locus
locally with this real subspace. It is therefore a closed totally real
three-manifold. The smooth $\mathbb Z/2$ action admits an equivariant
triangulation by \cite[Theorem~3.1, p.~48]{Illman1972}; compactness and finite
dimension give the finite-dimensional $G$-CW hypothesis. Smith's theorem,
in the form of
\cite[Theorem 4.1(1)]{Knudson2004}, applies to the smooth finite action on
the sphere and gives the stated mod-two homology. A nonempty
three-dimensional fixed set is consequently connected, and its first
mod-two cohomology vanishes; its orientation class $w_1$ therefore
vanishes. On the unit circle $1/\bar z=z$, so $\sigma$ commutes with
$\rho$. Its stabilizers are inherited from the checked semifree action
on $X$ and hence are either trivial or the whole circle.
\end{enumerate}
\end{proof}

We now fix a representative of each sign class, instead of leaving the
symbol $\sigma$ unspecified. Put
\[
 L_F=(2\pi i)^{-1}\log F,\qquad L_G=(2\pi i)^{-1}\log G,
 \qquad F^3=t,\quad G^4=t-1.
 \label{mapsI:real-eq-1}
\]
The roots are positive over $t>1$. We continue from that interval
through the upper half $t$-plane. Their logarithms are taken on the
orbifold uniformizing cover. Write $\iota$ for its reflection lift.
Every point of $-i\R_{>0}$ has a representative $c=1/2+iy$, and every
point of $i\R_{>0}$ has a representative $c=-1/2+iy$, with $y$ in the
admissible range. Set
\[
\begin{array}{c|c|c}
 &\alpha&S(t)\\\hline
 c=1/2+iy&1/3&2L_F-4L_G-1/6\\
 c=-1/2+iy&1/2&-2L_G+1/4
\end{array}
\qquad b_c(z)=\alpha\Pi_c(z)g+S(t(z))e.
\label{mapsI:real-eq-2}
\]

\begin{theorem}[Explicit real involutions and all fixed fibres]\label{mapsI:src:A.3.2}
\StatementRecord{mapsI:src:A.3.2}{proved}{Historical events: Claude Fable 5.2 (Anthropic), effort high, 16 November 2026: carried as background without fault found; used and rechecked in this run. Claude Fable 5.2 (Anthropic), effort high, 17 November 2026: carried as background without fault found. GPT-6 Astra (OpenAI), ChatGPT Pro mode, effort Pro, S6\_09\_PAPERS, 17 November 2026: targeted proof check. Also checked in this writing. Pre-continuation writing events: GPT-6 Astra (OpenAI), ChatGPT Pro mode, effort Pro, S6\_10\_PAPERS, 18 November 2026: computation or formal certificate; targeted proof check. Continuation metadata events: GPT-6 Astra (OpenAI), ChatGPT Pro mode, effort Pro, S6\_10\_PAPERS continuation, 19 November 2026: used and rechecked in this run; status-vocabulary reconciliation; metadata reconciliation only; no new mathematical check. \emph{Pre-continuation recorded proof scope.} Read complete proof and re-derived invariant character, finite solid-torus coordinates, cusp nonextension paths, exact lens meridians, relative H1 extension, affine real involutions/cusp units, real cap-meridian determinant and orbit-space recognition. Full finite-congruence replays and Smith/triangulation external locators remain separate gates. Reran real\_involutions.py and real\_finite\_four.py from copied source. Exact translation/Hermite and Smith calculations and all denominator-24 affine congruences pass, including all fixed-lift parities and the two-torus/empty outcomes. This supplements the separately recorded full geometric proof review.}{the hypotheses, notation and definitions stated here; \cref{fam:src:B.4.2}; \cref{mapsI:real-eq-4}; \cref{mapsI:real-eq-2}; \cref{fam:src:A.3.1}; \cref{mapsI:real-eq-1}}

On the ordinary period charts the formulas
\[
 \sigma_c[z,\zeta]=[\iota z,\bar\zeta+b_c(\iota z)],
 \qquad \sigma_c^- =\rho(1/2)\sigma_c
 \label{mapsI:real-eq-3}
\]
extend to antiholomorphic involutions of the compact threefold.
With these representatives their complete fixed-set data are as follows.
Here $2T^2$ means the disjoint union of two real two-tori; every displayed
circle in a double curve avoids its two triple-point endpoints.
\begin{center}\small
\begin{tabular}{cclccc}
\toprule
$w$&involution&image in $\RP^1$&on $S_1$&on $S_2$&on $W$\\
\midrule
$-ir$, $r>0$&$\sigma_c$&all $\RP^1$&$2T^2$&$2T^2$&$T^2\sqcup S^1_{D_2}\sqcup S^1_{D_{12}}$\\
$-ir$, $r>0$&$\sigma_c^-$&$[1,\infty]$&$\varnothing$&$T^2$&$S^1_{D_1}\sqcup S^1_{D_2}$\\
$ir$, $r>0$&$\sigma_c$&$\varnothing$&$\varnothing$&$\varnothing$&$\varnothing$\\
$ir$, $r>0$&$\sigma_c^-$&$[\infty,1]$ through $0$&$2T^2$&$T^2$&$S^1_{D_1}\sqcup S^1_{D_{12}}$\\
\bottomrule
\end{tabular}
\end{center}
The first interval is the closure of $t>1$; the second is the closure
of $t<1$. On every open real arc contained in the indicated image the
real regular fibre is $2T^2$. No triple point is real for these structures.
\end{theorem}



\begin{proof}
\begin{enumerate}[label=\arabic*.,leftmargin=*]
\item We first verify the global map. Use the lattice reflection from Proposition~\ref{fam:src:B.4.2},
\[
 H=\begin{pmatrix}1&0&0&0\\0&1&0&0\\-6&0&-1&0\\0&-1&0&-1\end{pmatrix},
 \quad E=d\otimes\gamma,
 \quad\Pi_{-\bar c}(\iota z)H=-\overline{\Pi_c(z)}.
\]
The initial map $\zeta\mapsto-\bar\zeta$ has target affine vectors
$-\varepsilon+d$ and $\varepsilon'-3d$. The normalization from this
target at parameter $c_0=-\bar c$ to the canonical datum at $c_0+1$
has linear part $-I$ and translation
\[
 b^{\rm norm}_{c_0}=\tfrac13\Pi_{c_0}g+(2L_F-4L_G)e.
 \label{mapsI:real-eq-4}
\]
Its period defects at the two finite generators are
$(0,-2,2,0)^t$ and $(0,0,2,0)^t$. To check them, the logarithmic increments
are $-2/3$ and $1$, and one multiplies the two displayed matrices $A_j$
by $g/3$. The finite difference vectors in the complex translation space are $2e$ and $-4e$, so subtracting
their logarithmic collars leaves holomorphic translations on both full
root discs. At the cusp
$2L_F-4L_G=s/3-(2\pi i)^{-1}\log(1-q)$ and
$\Pi g=(6\mu,-s+C_{21})^t$. The logarithmic terms cancel and give the
holomorphic toric unit
$(2\mu,C_{21}/3-(2\pi i)^{-1}\log(1-q))$.
Thus \eqref{mapsI:real-eq-4} extends over all three fillings, with its inverse.

\item For $c=1/2+iy$ compose this normalization with the initial reflection
and then with the constant translation $e/6$. This gives exactly the
first line of \eqref{mapsI:real-eq-2}. For $c=-1/2+iy$, additionally close the parameter
shift by the following explicit even-shift formula with $h=-2$,
\[
 b_{h,c_0}=-\frac h{12}\Pi_{c_0}g+
 \left(h(L_F-L_G)-\frac{h^2}{24}\right)e.
\]
A final constant translation $-e/4$ gives the second line of \eqref{mapsI:real-eq-2}.
The displayed translation has integral affine defects on the ordinary cover, holomorphic finite-chart corrections, and a nonvanishing toric-unit expression at the cusp; it represents the even shift of Proposition~\ref{fam:src:A.3.1}. Its cusp expression is checked immediately below. Indeed put $\ell=L_F-L_G=-s/12+\kappa_0(q)$ and
$C_{c,21}(q)=c+C_{0,21}(q)$, and set
$D_0(q)=C_{0,21}(q)-12\kappa_0(q)$.
The explicit even-shift translation becomes
\[
 b_{h,c}=
 \left(-h\mu/2,-h(c+D_0)/12-h^2/24\right)^t.
\]
The logarithmic terms have cancelled exactly. Exponentiating both
coordinates gives holomorphic nonzero torus multipliers on the whole cusp
disc; multiplication by these units extends over every toric stratum.
Compatibility with the period deck group is the integral identity
$\exp(2\pi i h\lambda_1)=1$ for $h\in2\mathbb Z$ and
$\lambda_1\in\mathbb Z$. The finite-chart defects are checked by the
same displayed $A_j$ and root logarithms used for the normalization.

\item Over $t>1$ we have
$\Re\mu=1/2$, $\Re\tau=\Re\beta^{\rm part}=0$ and the logarithms in
\eqref{mapsI:real-eq-1} are purely imaginary. The two squares are respectively the
integral period translations $2\widehat w$ and $3\widehat w$.
The square is a holomorphic automorphism and is consequently the
identity everywhere. This proves the involution claim.

\item We next give exact equations for all real fibres. Put $a=\Re c$ and
$R=-H+2aE$. In the chosen branch the reflection identifications on
$t>1$, $0<t<1$, $t<0$ are respectively $1,g_2,g_0^{-1}$.
They follow by reflecting the chosen upper-half-plane paths; equivalently
the root arguments are $(0,0)$, $(0,\pi/4)$ and $(\pi/3,\pi/4)$.
The three fibre involution matrices are
\[
 R_>=R,\qquad R_m=A_2^{-1}R,\qquad R_<=M_0R.
 \label{mapsI:real-eq-5}
\]
In marked real coordinates let $r_c(t)=\Pi_c(t)^{-1}b_c(t)$,
where $\Pi_c^{-1}$ is the real inverse. The fixed tori are exactly
\[
 \{x\in\R^4/\Lambda:(I-R_\bullet)x=
                       r_c(t)+\epsilon d/2\pmod\Lambda\},
 \qquad\epsilon=0,1.
 \label{mapsI:real-eq-6}
\]
This specifies the cosets, including their parameter dependence. Their
identity subgroups are $T_{\langle w,d\rangle}$ on both outer arcs and
$T_{\langle u+w,d\rangle}$ on the middle arc. When nonempty there are
two cosets. Relative to any one solution their offsets are
$(g+u)/2$ on the right arc and $g/2$ on the other two arcs.

\item For completeness, the finite arithmetic deciding existence can be
performed without approximating the periods. An affine involution
$x\mapsto R_\bullet x+\beta$ has a fixed point if and only if its
square translation $(I+R_\bullet)\beta$ belongs to
$(I+R_\bullet)\Lambda$. In \eqref{mapsI:real-eq-2} the square translation is
\[
 (I+R_\bullet)\alpha g+2\Re S\,d.
\]
For $c=1/2+iy$ the values $2\Re S$ on the three arcs are
$-1/3,-4/3,-2/3$; for $c=-1/2+iy$ they are $1/2,0,0$.
The resulting integral square vectors are respectively
$(2w,0,2w)$ and $(3w,d,3w)$. Adding $d/2$ to the involution adds $d$
to each square vector. Taking column Hermite normal forms of
$I+R_\bullet$ gives exactly the nonempty entries in the table.
All three $I-R_\bullet$ have nonzero Smith factors $1,2$.
This proves the coset counts and offsets; the complete exact computation
is in \texttt{code/maps\_audit/MAPS/code/real\_involutions.py}. In particular the lift on
the left arc is $g_0^{-1}$, not $g_0$; the latter would give a nonintegral
square vector and is the wrong continuation convention.

\item At $t=0$ the root degree is odd. Conjugation by the degree-three deck
group acts transitively on the three reflection lifts; the centralizer
of a reflection in that deck group is trivial.
Thus the quotient has exactly the same two real tori as the adjacent
nonempty real regular fibres. In the empty cases there are none. This
also follows by solving the limiting affine equations, and checks the
full multiplicity-three quotient rather than treating it as a reduced
unramified fibre.

\item At $t=1$ write the affine deck map as
$G(x)=A_2x+v_2/4$. The real root involution has matrix $R_>$ and
translation
\[
 b_-=(1/3,0,0,-1/6)^t,
 \qquad b_+=(1/2,0,0,1/4)^t
\]
for the two parameter signs. The real locus of the quotient is the
quotient of the fixed sets of $G^kh$, $0\leq k<4$. Its even and odd
lift parities are identified internally by $G$, while $G^2$ preserves
each parity. Direct integer congruences show that the nonempty lift
parities are: both for negative $w$ and $\epsilon=0$; only even for
negative $w$ and $\epsilon=1$; neither for positive $w$ and $\epsilon=0$;
only odd for positive $w$ and $\epsilon=1$. Every nonempty lift has
two tori, exchanged by $G^2$. Each parity therefore contributes one
torus downstairs, not a Klein bottle. The exhaustive congruence check
\texttt{code/maps\_audit/MAPS/code/real\_finite\_four.py} uses denominator $24$: translation
denominators divide $12$, and the nonzero Smith factors are $1,2$, so
every component has such a representative. It checks the component
permutation, not just the dimension of the fixed eigenspace. The normal
real line has fourth power of sign positive for the even parity and
negative for the odd parity. This establishes the two one-sided
endpoints at $t=1$ recorded in the table.

\item It remains to compute the cusp, which also excludes hidden real
components confined to a singular fibre. Put
$A=C_{12}(0)$; the cusp normalization gives
\[
 \Re\mu(0)=1/2,\quad A\in i\R,\quad C_{21}(0)=c-A,
 \quad C_{11}(0)=6\mu(0),\quad C_{22}(0)=\mu(0).
\]
In a toric lift the involution is
$x\mapsto U/\bar x$, with
\[
\begin{array}{c|cc}
 &b_1(0)&b_2(0)\\\hline
 w=-ir&2\mu(0)&C_{21}(0)/3-1/6\\
 w=ir&3\mu(0)&C_{21}(0)/2+1/4
\end{array}
\qquad U_j=e^{2\pi i b_j(0)}.
\label{mapsI:real-eq-7}
\]
Thus $U_2>0$ in both rows, whereas $U_1>0$ in the first and $U_1<0$
in the second. Replacing $\sigma_c$ by $\sigma_c^-$ negates $U_2$.
The smooth central orbit has a real torus precisely when both units
are positive, with equations $|x_j|^2=U_j$.

\item For an edge $v$, set $n=(-v_2,v_1)^t$. The period translate bringing
the inverted edge back is $B_0n=v$. Its curve coordinate $z=x^n$ has
involution
\[
 z\longmapsto K_v/\bar z,
 \qquad K_v=U^n e^{2\pi i n^tC(0)n}.
 \label{mapsI:real-eq-8}
\]
The signs for $(D_1,D_2,D_{12})$ are $(-,+,+)$ for negative $w$
and $(-,-,-)$ for positive $w$. Negating $U_2$ reverses the first and
third signs, but not the second. A positive sign gives the single
circle $|z|^2=K_v$; a negative sign gives no fixed point. Central
inversion exchanges the two triangle classes of the periodic fan, so
it exchanges the two triple points. These calculations prove every
cusp entry. Finally a nonempty fixed three-manifold cannot be confined
to the special fibres: a holomorphic function constant on an open
piece of a totally real submanifold is locally constant. This also
proves that the positive-$w$ representative with no regular real
fibres is entirely fixed-point free.
\end{enumerate}
\end{proof}

\begin{theorem}[The real three-manifolds]\label{mapsI:src:A.3.3}
\StatementRecord{mapsI:src:A.3.3}{proved}{Historical events: Claude Fable 5.2 (Anthropic), effort high, 16 November 2026: carried as background without fault found; used and rechecked in this run. GPT-6 Astra (OpenAI), ChatGPT Pro mode, effort Pro, S6\_09\_PAPERS, 17 November 2026: targeted proof check. Also checked in this writing. Pre-continuation writing events: GPT-6 Astra (OpenAI), ChatGPT Pro mode, effort Pro, S6\_10\_PAPERS, 18 November 2026: computation or formal certificate; targeted proof check. Continuation metadata events: GPT-6 Astra (OpenAI), ChatGPT Pro mode, effort Pro, S6\_10\_PAPERS continuation, 19 November 2026: used and rechecked in this run; status-vocabulary reconciliation; metadata reconciliation only; no new mathematical check. \emph{Pre-continuation recorded proof scope.} Read complete proof and re-derived invariant character, finite solid-torus coordinates, cusp nonextension paths, exact lens meridians, relative H1 extension, affine real involutions/cusp units, real cap-meridian determinant and orbit-space recognition. Full finite-congruence replays and Smith/triangulation external locators remain separate gates. Reran real\_involutions.py and real\_finite\_four.py from copied source. Exact translation/Hermite and Smith calculations and all denominator-24 affine congruences pass, including all fixed-lift parities and the two-torus/empty outcomes. This supplements the separately recorded full geometric proof review.}{the hypotheses, notation and definitions stated here; \cref{mapsI:src:A.3.2}}

Every nonempty fixed set in Theorem~\ref{mapsI:src:A.3.2} is diffeomorphic to $S^3$.
Its fundamental group is trivial and its integral homology is
$(\Z,0,0,\Z)$. For negative imaginary $w$, the restricted circle action
is the standard semifree action on $S^3$ with fixed circle
$D_2^{\sigma_c}$. For positive imaginary $w$, the nonempty real form
has a free circle action isomorphic, up to circle and orientation
conventions, to the Hopf action.
\end{theorem}



\begin{proof}
\begin{enumerate}[label=\arabic*.,leftmargin=*]
\item For negative imaginary $w$ both fixed manifolds meet $D_2$ in one circle.
Their circle action is semifree, so the orbit space is a compact
orientable surface with boundary. The local quotient at a fixed circle
is the quotient of $S^1\times\C$ by scalar rotation on $\C$, hence is
a boundary collar. Over the surface interior the principal circle bundle
is trivial because the surface has boundary and $H^2=0$. Changes of
boundary trivialization extend over the attached normal discs by
rotation. Thus the three-manifold is
\[
 (B\times S^1)\ \cup_{\partial B\times S^1}\ (\partial B\times D^2).
\]
If $B$ has genus $g_B$ and $b_B$ boundary components, van Kampen or its
surface handle decomposition gives first homology
$\Z^{2g_B+b_B-1}$. Smith's theorem makes this rank zero. Since $b_B>0$,
we have $g_B=0,b_B=1$. The displayed manifold is then
$\partial(D^2\times D^2)=S^3$, with the indicated circle action. This
derives the smooth recognition from an explicit surface decomposition.

\item For positive imaginary $w$ the nonempty real manifold is free under
$\rho$. The two real tori on $t<1$ join at the single real torus over
$1$. At $0$ the odd root comparison is a torus product. At infinity
they are capped by the real neighbourhoods of $D_1$ and $D_{12}$.
These are solid tori: the primitive circle in the edge direction
collapses and the remaining real circle is their core. In the outer
real-fibre lattice $\Z\langle w,d\rangle$ the two cap meridians are
$w$ and $w-d$.

\item We check the intervening torus identification. Put
$B_1=A_2^{-1}A_1^{-1}A_2$ and $C=B_1^2$. The two real-arc involutions
satisfy $R_<R_m=B_1$, $R_mB_1R_m=B_1^{-1}$ and $B_1^3=I$.
Moving between the real root lifts requires
$B_1^kR_mB_1^{-k}=B_1R_m$, so $2k\equiv1\pmod3$ and $k=2$.
Thus the actual odd-root transport is $C$, not merely an arbitrary
matrix conjugating the two reflections. It satisfies $R_<=CR_mC^{-1}$, and
$C(u+w)=-w+d$. The exchange across the multiplicity-four real torus
is $A_2^2$; after this odd-root transport it is
$CA_2^2C^{-1}$, whose restriction to $\langle w,d\rangle$ is
$w\mapsto-w$, $d\mapsto d$. Consequently the meridians, compared on
one boundary torus, have determinant of absolute value one. The
resulting genus-one Heegaard splitting is $S^3$: primitive meridians
intersecting once can be taken to the two coordinate circles by an
integral torus automorphism, giving the usual two solid tori of $S^3$.
The free circle $d$ then has Euler number of absolute value one, so
the displayed bundle over $S^2$ is the Hopf bundle. This proves the
integral assertions, beyond the mod-two information.
\end{enumerate}
\end{proof}

\begin{proposition}[Restrictions of the repaired circle quotient]\label{mapsI:src:A.3.4}
\StatementRecord{mapsI:src:A.3.4}{proved}{Historical events: Claude Fable 5.2 (Anthropic), effort high, 16 November 2026: carried as background without fault found. Claude Fable 5.2 (Anthropic), effort high, 17 November 2026: carried as background without fault found. GPT-6 Astra (OpenAI), ChatGPT Pro mode, effort Pro, S6\_09\_PAPERS, 17 November 2026: targeted proof check. Also checked in this writing. Pre-continuation writing events: GPT-6 Astra (OpenAI), ChatGPT Pro mode, effort Pro, S6\_10\_PAPERS, 18 November 2026: targeted proof check. Continuation metadata events: GPT-6 Astra (OpenAI), ChatGPT Pro mode, effort Pro, S6\_10\_PAPERS continuation, 19 November 2026: used and rechecked in this run; status-vocabulary reconciliation; metadata reconciliation only; no new mathematical check. \emph{Pre-continuation recorded proof scope.} Read complete proof and re-derived invariant character, finite solid-torus coordinates, cusp nonextension paths, exact lens meridians, relative H1 extension, affine real involutions/cusp units, real cap-meridian determinant and orbit-space recognition. Full finite-congruence replays and Smith/triangulation external locators remain separate gates.}{the hypotheses, notation and definitions stated here; \cref{mapsI:src:A.2.1}; \cref{mapsI:real-eq-2}}

The repaired map $e_{\rm q}$ of Theorem~\ref{mapsI:src:A.2.1} may be chosen $\rho$-invariant. Its restriction to
every nonempty real form is null-homotopic. On an ordinary real fibre
its image consists of the two quotient-circle values
\[
 \theta=\tfrac12\gamma(r_c(t)),\qquad
 \theta=\tfrac12\gamma(r_c(t))+\tfrac12\pmod\Z,
 \quad
 \gamma(r_c(t))=\alpha+\frac{\Im S(t)}{D_c(t)},
 \label{mapsI:real-eq-9}
\]
with $\alpha,S$ from \eqref{mapsI:real-eq-2}. The circle action is along each real
two-torus and is annihilated by $e_{\rm q}$.
\end{proposition}



\begin{proof}
\begin{enumerate}[label=\arabic*.,leftmargin=*]
\item The literal quotient is circle invariant. In a cusp extension
$\theta':N_0\to S^1$, each orbit circle has degree zero: its class is
$d$, which vanishes in $\pi_1(N_0)$. The quotients
$\theta'(\rho(s)x)/\theta'(x)$ therefore have a continuous real lift
vanishing at $s=0$ and periodic in $s$. Averaging this real lift over
$s\in\R/\Z$ changes $\theta'$ to a circle-invariant map. The cocycle
identity proves invariance, and on an invariant outer collar the lift
is zero, so the replacement agrees exactly with the prescribed map.
The same construction is smooth. Thus $e_{\rm q}|_{\operatorname{Fix}\sigma}$
factors through its circle orbit space, which is $D^2$ or $S^2$ by the
preceding proof. Both have only two-dimensional homotopy type and
$S^3$ is two-connected, so this restriction is null-homotopic.
The first row of $I-R_\bullet$ is $(2,0,0,0)$, giving the two values
in \eqref{mapsI:real-eq-9}. The real inverse of the period matrix sends $S e$ to
first coordinate $\Im S/D_c$; this follows by subtracting
$(\Im\mu/\Im\tau)$ times its first imaginary-coordinate equation
from the second. Both real-fibre identity subgroups have zero
$\gamma$-coordinate, proving the remaining assertions.
\end{enumerate}
\end{proof}

The formulas specify the embedded fixed sets by their torus cosets and
their actual finite and toric attachments. No ambient-isotopy
classification or value of the Haefliger knot invariant of these
embedded three-spheres is asserted here.


\begin{cited}[The coherent closing translation]\label{push:src:C.3.5}
\StatementRecord{push:src:C.3.5}{proved}{Historical events: Claude Fable 5.2 (Anthropic), effort high, 16 November 2026: carried as background without fault found. GPT-6 Astra (OpenAI), ChatGPT Pro mode, effort Pro, S6\_09\_PAPERS, 17 November 2026: targeted proof check. Also checked in this writing. Pre-continuation writing events: GPT-6 Astra (OpenAI), ChatGPT Pro mode, effort Pro, S6\_10\_PAPERS, 18 November 2026: used and rechecked in this run. Continuation metadata events: GPT-6 Astra (OpenAI), ChatGPT Pro mode, effort Pro, S6\_10\_PAPERS continuation, 19 November 2026: used and rechecked in this run; status-vocabulary reconciliation; metadata reconciliation only; no new mathematical check. \emph{Pre-continuation recorded proof scope.} Read DGKXv1 Proposition7.8 and full proof pp91–92 and PropositionsC.1–C.2 with complete proofs p102. Rechecked the displayed defect -kl/12 and quadratic correction -k\textasciicircum{}2/24 and matched the fixed roots and analytic markings. This closes the exact source-locator gate for the credited short application.}{the hypotheses, notation and definitions stated here; \cref{fam:src:A.3.1}; \citet[Proposition~7.8, pp.~91--92, and Proposition~C.2, p.~102]{DuEtAl2026v1}; \CompendiumRef{I}{app:rederive:push-src-C.3.5}}

Fix the standard twists and the root and cusp markings of Proposition~\ref{fam:src:A.3.1}.
There is a coherent choice of the even-shift maps
\(\widetilde F_{h,c}:X_c\to X_{c+h}\), \(h\in2\mathbb Z\),
with the following explicit ordinary-chart translation. Choose global
roots \(F^3=t\), \(G^4=t-1\) on the uniformizing cover and put
\(\ell=(2\pi i)^{-1}\log(F/G)\pmod{\mathbb Z}\), \(e=(0,1)^t\).
Then the translation is
\begin{equation}\label{mapsI:closing:eq:C-explicit-closing}
 b_{h,c}=-\frac h{12}\Pi_c\widehat\gamma
              +\left(h\ell-\frac{h^2}{24}\right)e.
\end{equation}
For \(h=12\), it is the scalar multiplication
\begin{equation}\label{mapsI:closing:eq:C-twelve-scalar}
 K(t)=\frac{t^4}{(t-1)^3}
\end{equation}
on the ordinary \(\mathbb C^*\)-coordinate, after removing an integral
period. On the finite root charts its scalar representatives are
\(K(t)s_1^{-12}\) and \(K(t)s_2^{12}\), respectively, both units.
In a cusp coordinate \(q=t_c\), the scalar representative is
\(q^{-1}H(q)\), with \(H\) a holomorphic nonvanishing function.
The same map has a holomorphic \((\mathbb C^*)^2\)-unit representative
on the entire cusp filling. These two representatives differ by
exactly the \(\widehat\gamma\) cusp deck transformation.

For six positive \(w\)-loops the closing map is
\(\widetilde F_{12,c_0-12}\); thus the stated cusp exponent is \(-1\)
with $w=e^{-\pi i c}$ and positive loops in the $w$-plane. This computes the closing
cocycle, not the mapping class of its composition with parallel
transport. In particular it does not imply \(6[\Phi]=0\).
\end{cited}
The expanded source derivation is retained in the proof compendium, Paper~I, label \nolinkurl{push:src:C.3.5}.




The raw even-shift formula and its coherent correction are due to \citet[Proposition~7.8, pp.~91--92, and Proposition~C.2, p.~102]{DuEtAl2026v1}. A complete re-derivation is given in \CompendiumRef{I}{app:rederive:push-src-C.3.5}.




\begin{figure}[htbp]
\centering
\begin{tikzpicture}[x=1cm,y=1cm,>=stealth,font=\small]
\draw[thick] (-5.3,0) .. controls (-1.8,-.35) and (1.8,-.35) .. (5.3,0);
\node at (0,-1.05) {$B=\mathbb P^1_t$};
\foreach \x/\lab in {-4/0,0/1,4/\infty}{
 \fill (\x,-.1) circle (2pt);
 \node[below=5pt] at (\x,-.1) {$\lab$};
 \draw[->] (\x,1.15)--(\x,.15);
}
\node[draw,rounded corners,align=center,text width=2.8cm,minimum height=1.2cm] at (-4,1.75)
 {$3S_1$\\free order-three\\root quotient};
\node[draw,rounded corners,align=center,text width=2.8cm,minimum height=1.2cm] at (0,1.75)
 {$4S_2$\\free order-four\\root quotient};
\node[draw,rounded corners,align=center,text width=2.8cm,minimum height=1.2cm] at (4,1.75)
 {$W$\\normalization $\mathrm{dP}_6$};
\draw[->] (-2,1.15)--(-2,.15);
\node[align=center] at (-2,2.85) {ordinary fibre\\$\mathbb C^2/\Pi_c\mathbb Z^4$};
\node[align=center] at (4,3.05) {three paired boundary curves\\two triple points};
\end{tikzpicture}
\caption{The fixed $(3,4,\infty)$ base data. The arrows denote the
restriction of $f$ to each indicated fibre, with its scheme
multiplicity retained. The ordinary period quotient, finite affine
quotients and the paired hexagon are the three actual local models
in the common conventions; their overlap maps are fixed there.}
\label{fig:I:orbifold}
\end{figure}

\begin{figure}[htbp]
\centering
\begin{tikzpicture}[x=1cm,y=1cm,>=stealth,font=\small]
\fill[blue!4] (0,0) circle (2.6);
\draw[thick,dashed] (0,0) circle (2.6);
\draw[->,gray] (-3.1,0)--(3.2,0) node[right] {$\Re w$};
\draw[->,gray] (0,-3.15)--(0,3.2) node[above] {$\Im w$};
\draw[very thick,blue!65!black] (0,-2.56)--(0,-.13);
\draw[very thick,blue!65!black] (0,.13)--(0,2.56);
\fill[white] (0,0) circle (.11);
\draw[thick] (0,0) circle (.11);
\draw[->] (.23,-.18)--(1.1,-.7);
\node[anchor=west,align=left] at (1.05,-.88) {puncture $w=0$\\corner winding $-pn$};
\node[align=center] at (-1.45,.8) {$0<|w|<216$};
\node[anchor=west,align=left] at (3.15,1.65)
 {outer circle $|w|=216$\\$|q_2|=1$\\$q_2=iw/216$};
\draw[->] (3.05,1.55)--(2.15,1.43);
\node[anchor=east,align=right] at (-.25,1.8)
 {real structures\\$w\in i\mathbb R$};
\node[anchor=west,align=left] at (3.15,-2.1)
 {all admissible members\\retain $3S_1,4S_2,W$};
\draw[->] (3.05,-1.95)--(1.85,-1.7);
\end{tikzpicture}
\caption{The exact parameter disc for the standard fixed data.
The coordinate map is $w=e^{-\pi i c}$ and the coarser coordinate
map is $u=w^2$. The imaginary diameter is the fixed locus of
$w\mapsto-\bar w$, which is the reflected-family map.
The two boundary obstructions use the meridian comparison at the
puncture and the intrinsic Albanese nome at the outer circle.
The special fibres are over $t=0,1,\infty$ in every member, as
indicated in Figure~\ref{fig:I:orbifold}.}
\label{fig:I:parameter-disc}
\end{figure}

\begin{figure}[htbp]
\centering
\begin{tikzpicture}[x=1cm,y=1cm,>=stealth,font=\small]
\begin{scope}[shift={(-3.3,0)}]
 \fill[blue!5] (-1,0)--(-1,2.8)--(1,2.8)--(1,0)
  arc[start angle=30,end angle=150,radius=1.1547];
 \draw[thick] (-1,0)--(-1,2.8);
 \draw[thick] (1,0)--(1,2.8);
 \draw[thick] (1,0) arc[start angle=30,end angle=150,radius=1.1547];
 \node at (0,1.3) {$\mathcal F$};
 \node[align=center] at (0,-1.05)
 {$|\Re\tau|\le\tfrac12,\ |\tau|\ge1$\\$0\le\theta<\pi$};
 \draw[->] (0,2.1)--(0,3.1) node[above] {cusp};
\end{scope}
\begin{scope}[shift={(1.0,.2)},x={(1.8cm,0cm)},y={(.75cm,.5cm)},z={(0cm,2.2cm)}]
 \fill[blue!4] (0,0,0)--(1,0,0)--(1,0,1)--(0,0,1)--cycle;
 \draw[thick] (0,0,0)--(1,0,0)--(1,1,0)--(1,1,1)--(0,1,1)--(0,0,1)--cycle;
 \draw[thick] (1,0,0)--(1,0,1)--(0,0,1);
 \draw[thick] (1,0,1)--(1,1,1);
 \draw[dashed] (0,0,0)--(0,1,0)--(1,1,0);
 \draw[dashed] (0,1,0)--(0,1,1);
 \node[below] at (.5,0,0) {$0\le x_1<1$};
 \node[right] at (1,.5,0) {$0\le x_2<1$};
 \node[left] at (0,0,.5) {$0\le t<12$};
\end{scope}
\node[align=center] at (1.95,-1.55)
 {$U:\ t\mapsto t-6-6x_2$\qquad
  $V:\ t\mapsto t-6+6x_1$\\
  $Z:\ t\mapsto t+12$};
\end{tikzpicture}
\caption{A fundamental set for the index-twelve monodromy kernel
in the real parabolic stabilizer of $d$.
The base is a modular fundamental set with its rotation coordinate;
the fibre is the indicated Heisenberg parallelepiped. The displayed
face maps include the central shears. The representative map and
the boundary identifications are proved in
\ref{I:monodromy-fundamental-set}; the index twelve occurs in the
central fibre coordinate.}
\label{fig:I:monodromy-domain}
\end{figure}

\begin{proposition}[The fundamental set of the monodromy kernel]
\label{I:monodromy-fundamental-set}
\StatementRecord{I:monodromy-fundamental-set}{proved}{No earlier separately identified check is recorded. Checked only in this writing. Pre-continuation writing events: GPT-6 Astra (OpenAI), ChatGPT Pro mode, effort Pro, S6\_10\_PAPERS, 18 November 2026: targeted proof check. Continuation metadata events: GPT-6 Astra (OpenAI), ChatGPT Pro mode, effort Pro, S6\_10\_PAPERS continuation, 19 November 2026: used and rechecked in this run; status-vocabulary reconciliation; metadata reconciliation only; no new mathematical check. \emph{Pre-continuation recorded proof scope.} checked only in this writing by GPT-6 Astra (OpenAI), ChatGPT Pro mode, effort Pro, 18 November 2026}{\cref{fam:src:F.1.1}, \cref{fam:src:F.1.2}; the displayed parabolic group law}
In the real coordinates $g(B,x,t)$ of
\ref{fam:src:F.1.1}, a measurable fundamental set for
$\Gamma\backslash P_d(\mathbb R)$ consists of the matrices
\[
 g(B,x,t),\qquad
 B=n(\Re\tau)a(\Im\tau)k(\theta),\quad
 \tau\in\mathcal F,\quad 0\le\theta<\pi,
 \quad x_1,x_2\in[0,1),\quad t\in[0,12).
\]
Here
$n(u)=\left(\begin{smallmatrix}1&u\\0&1\end{smallmatrix}\right)$,
$a(v)=\operatorname{diag}(\sqrt v,1/\sqrt v)$,
and $k(\theta)$ is the rotation matrix. Paired modular edges
are assigned to one side, and their lower-dimensional
stabilizer identifications use the corresponding lifts in
$\Gamma$: for
$S=\left(\begin{smallmatrix}0&-1\\1&0\end{smallmatrix}\right)$
and $T=\left(\begin{smallmatrix}1&1\\0&1\end{smallmatrix}\right)$,
we use $\widetilde S=g(S,0,9)$,
$\widetilde T=g(T,0,1)$ and
$\widetilde{-I}=g(-I,0,6)$.
Their maps on the three coordinates are
$(B,x,t)\mapsto(AB,Ax,t+\alpha(A))$, with
$\alpha(S)=9$, $\alpha(T)=1$ and $\alpha(-I)=6$.
The Heisenberg face identifications are the
maps printed in Figure~\ref{fig:I:monodromy-domain}.
\end{proposition}
\begin{proof}
\begin{enumerate}[label=\arabic*.,leftmargin=*]
\item The projection $g(B,x,t)\mapsto B$ maps $\Gamma$
onto $\operatorname{SL}_2(\mathbb Z)$ by
\ref{fam:src:F.1.1}. For any $B\in\operatorname{SL}_2(\mathbb R)$,
the usual modular reduction gives an integral matrix $A$
with $AB=n(u)a(v)k(\theta)$ in the stated set. Choose an
actual lift $g(A,y,s)\in\Gamma$ and multiply on the left.
For the displayed edge generators, the character formula
$\chi(g(A,0,t))=t-\alpha(A)\pmod {12}$ proves that the
three specified lifts lie in $\Gamma$. Substitution in
the group law gives each displayed coordinate map.
This is the specified map from a general parabolic element
to one over the modular set.
\item Two choices of lift over the same $A$ differ by
$\Gamma\cap H$. The group calculation in
\ref{fam:src:F.1.1} identifies this kernel with
$\langle U,V,Z\rangle$, where
$U=g(I,(1,0),-6)$, $V=g(I,(0,1),-6)$ and
$Z=g(I,0,12)$.
\item The multiplication law sends $(x,t)$ under $U$
to $(x+(1,0),t-6-6x_2)$, and under $V$ to
$(x+(0,1),t-6+6x_1)$. Powers of these maps first put
$x_1$ and $x_2$ in $[0,1)$. A power of $Z$ then puts
$t$ in $[0,12)$. All three are actual left translations
in $\Gamma$, so they retain the original coset.
\item Within the interiors of these intervals equality
of two reduced representatives forces their kernel
element to have zero $x$-translation and then zero
central translation. Modular uniqueness gives the same
conclusion in the base interior. The usual edge and
elliptic identifications, lifted by the specified
surjection, account for the measure-zero boundary.
This proves the fundamental-set assertion and all
displayed Heisenberg identifications.
\end{enumerate}
\end{proof}

\section{The limit at the puncture}\label{fam:section:E}
\subsection{The untwisted pinching and its local smoothings}\label{fam:section:E.1}

In this section $S$ is the smooth minimal rational elliptic surface with
Weierstrass equation
\[
 y^2=x^3-3t^3(t-1)x+2t^4(t-1)^2,
\]
$\pi:S\to B$ is its elliptic fibration, $P$ is the section used in the
period dictionary, $T=t_{6P}:S\to S$, and $L=M_S=\OO_S(P-O)$.
We use the convention that $\mathbb P_S(\OO\oplus L)$ contains the total
space of $L$, with fibre coordinate $a$, zero section $s_0$, and infinity
section $s_\infty$. The two normal bundles are respectively $L$ and
$L^{-1}$. We write $R=\mathbb P_S(\OO\oplus L)$.

\begin{proposition}\label{fam:src:E.1.1}\label{fam:st:E.1.1}
\StatementRecord{fam:src:E.1.1}{proved}{Historical events: Claude Fable 5.2 (Anthropic), effort high, 12 November 2026: targeted proof check. GPT-6 Astra (OpenAI), ChatGPT Pro mode, effort Pro, S6\_09\_PAPERS, 17 November 2026: targeted proof check. Also checked in this writing. Pre-continuation writing events: GPT-6 Astra (OpenAI), ChatGPT Pro mode, effort Pro, S6\_10\_PAPERS, 18 November 2026: targeted proof check. Continuation metadata events: GPT-6 Astra (OpenAI), ChatGPT Pro mode, effort Pro, S6\_10\_PAPERS continuation, 19 November 2026: used and rechecked in this run; status-vocabulary reconciliation; metadata reconciliation only; no new mathematical check. \emph{Pre-continuation recorded proof scope.} Read complete source proof and re-derived its explicit pushout/plumbing rings, collar gluing, flatness or finite quotient, and integral lattice/monodromy or fan maps as applicable. No proof credit is inferred from compilation; exact algebraic replays and cited source matching are recorded separately.}{the hypotheses, notation and definitions stated here; \cite[Lemma 2.1]{Engel2026}}

The quotient
\[
 Q=R/\{s_0(x)\sim s_\infty(Tx):x\in S\}
\]
has a compact, reduced complex-space structure. Its normalization is
the finite holomorphic map $\nu:R\to Q$. Its singular locus is a copy
of $S$, and its local analytic equation along that locus is $ab=0$,
with two further coordinates on $S$. The induced map $Q\to B$ is
proper. Its sheaf of first-order local smoothing parameters is
\[
 \mathcal T^1_Q=\mathcal Ext^1_{\OO_Q}(\Omega_Q,\OO_Q)
       \simeq L\otimes T^*L^{-1}\simeq\pi^*\OO_B(1)
\]
on the double locus.
\end{proposition}



\begin{proof}
\begin{enumerate}[label=\arabic*.,leftmargin=*]
\item The two sections of $R$ are disjoint closed submanifolds. The
equivalence relation consists of their two graphs and the diagonal;
it is closed in $R\times R$. A quotient of a compact Hausdorff space
by a closed equivalence relation is Hausdorff: for distinct equivalence
classes, compactness gives disjoint saturated neighbourhoods by taking
finite subcovers of the complements of the relation. Thus the displayed
quotient is compact and Hausdorff.

\item For a coordinate neighbourhood $U\subset S$ and frames of $L|_U$ and
$L^{-1}|_{T(U)}$, the two tubular coordinate charts are $(x,a)$ at
$s_0(U)$ and $(Tx,b)$ at $s_\infty(T(U))$. We give their identification
the ring
\[
 \{(f(x,a),g(x,b)):f(x,0)=g(x,0)\}
   =\C\{x_1,x_2,a,b\}/(ab).
\]
Changes of frame are holomorphic and invertible on each branch and
have the same change of $x$ on their intersection. They therefore glue
these hypersurface charts to the inherited charts off the sections.
This constructs the asserted complex space. The direct sum of the two
regular branch rings is the integral closure: an integral element in
the total ring of fractions is integral on each smooth branch, and
conversely the two branch idempotents are integral. The resulting map
$\nu$ is proper with at most two points in each fibre, and its local
ring extension is finite. The map to $B$ is well defined because
$\pi T=\pi$; its properness follows from compactness of $Q$.

\item The conormal presentation for the hypersurface $ab=0$ gives
$\mathcal Ext^1(\Omega,\OO)=\OO/(a,b)$. A change of normal frames
$a'=g(x)a$, $b'=h(x)b$ changes the equation and its smoothing parameter
by $gh$. Consequently these local modules glue as
$N_{s_0/R}\otimes T^*N_{s_\infty/R}=L\otimes T^*L^{-1}$.
For this particular surface, Engel's Lemma 2.1 states
$\operatorname{Hom}(T^*L,L)=\pi^*\OO_B(1)$ \cite[Lemma 2.1]{Engel2026}.
Its identification applies here with exactly $T=t_{6P}$ and $L=\OO(P-O)$;
there is no replacement of the translation by its inverse in this
normal-bundle calculation. One can recover the degree in that lemma
directly: the bundle has multidegree zero on every fibre, since $T$
preserves every component; hence it is $\pi^*\OO_B(n)$. Its degree on
$O$ is
\[
 n=(O-P)\cdot 6P+(P-O)\cdot O=(2-2)+(0-(-1))=1.
\]
The section intersections used here are the explicit computations in
the proof of that lemma. This derives the displayed sheaf in the
normal-bundle convention of our pushout.
\end{enumerate}
\end{proof}

\begin{proposition}\label{fam:src:E.1.2}\label{fam:st:E.1.2}
\StatementRecord{fam:src:E.1.2}{proved}{Historical events: Claude Fable 5.2 (Anthropic), effort high, 12 November 2026: targeted proof check. GPT-6 Astra (OpenAI), ChatGPT Pro mode, effort Pro, S6\_09\_PAPERS, 17 November 2026: targeted proof check. Also checked in this writing. Pre-continuation writing events: GPT-6 Astra (OpenAI), ChatGPT Pro mode, effort Pro, S6\_10\_PAPERS, 18 November 2026: targeted proof check. Continuation metadata events: GPT-6 Astra (OpenAI), ChatGPT Pro mode, effort Pro, S6\_10\_PAPERS continuation, 19 November 2026: used and rechecked in this run; status-vocabulary reconciliation; metadata reconciliation only; no new mathematical check. \emph{Pre-continuation recorded proof scope.} Read complete source proof and re-derived its explicit pushout/plumbing rings, collar gluing, flatness or finite quotient, and integral lattice/monodromy or fan maps as applicable. No proof credit is inferred from compilation; exact algebraic replays and cited source matching are recorded separately.}{the hypotheses, notation and definitions stated here; \cref{fam:src:E.1.1}; \cite{Engel2026}}

Over $S^\circ=S\setminus S_\infty$, the pushout in \ref{fam:src:E.1.1} is the central
fibre of a holomorphic plumbing with parameter $u$. On every relatively
compact open disc $U\Subset B\setminus\{\infty\}$ this plumbing is
proper and flat over $U\times\Delta_\epsilon$, for a sufficiently small
positive $\epsilon$. Its fibre at $u\ne0$ is, over $U$, Engel's untwisted
quotient by the contracting generator $u\phi$ covering $T^{-1}$.
This assertion includes discs containing $0$ or $1$; it is made before
the logarithmic transformations at those points.
\end{proposition}



\begin{proof}
\begin{enumerate}[label=\arabic*.,leftmargin=*]
\item Let $\phi$ be the section of $L\otimes T^*L^{-1}$ with divisor
$S_\infty$, as in Corollary 2.2 of \cite{Engel2026}. On $S^\circ$ it is
nowhere zero. At the pair $s_0(x),s_\infty(Tx)$ we use the node chart
\[
       a\otimes b=u\phi_x,
       \qquad a\in L_x,\quad b\in L^{-1}_{Tx}.                 \tag{E1}
\]
This equation is invariant under changes of the two line-bundle
frames, so it defines a hypersurface chart over $S^\circ\times\Delta$.
We describe the gluing, including the required quotient comparison.
Choose a Hermitian metric on $L$ over $\pi^{-1}(\overline U)$, and
$0<r<1/4$. On this compact set the norm of $\phi$ has a finite upper
bound. Shrink $\epsilon$ so that $|u|\|\phi_x\|<r^2$.
Retain the core $r<\|a\|<r^{-1}$ of each projective fibre. Attach the
node chart with $\|a\|<2r$ and $\|b\|<2r$, using the core coordinate
$a$ on its first collar and the coordinate $1/b$ over $Tx$ on its
second collar. The metric used for $b$ is the dual metric over $Tx$.
The two attachment collars in a node chart are disjoint because
$\|a\|\|b\|<r^2$. The transition functions and their inverses are
holomorphic on the overlaps.

\item These charts are Hausdorff after gluing. Indeed, the smaller closed
core $r\leq\|a\|\leq r^{-1}$ and the closed subcharts
$\|a\|,\|b\|\leq 3r/2$ are compact over a compact subset of
$U\times\Delta_\epsilon$ and have precisely the prescribed closed
identifications on their collars. They cover after a harmless
enlargement of the open cores. The finite closed-chart description
separates distinct points, as in \ref{fam:src:E.1.1}, and proves properness over
$U\times\Delta_\epsilon$.

\item For $u\ne0$, solving (E1) gives
\[
              a_{Tx}=1/b=a_x/(u\phi_x).
\]
This is exactly the inverse of the contracting bundle map
$L_{Tx}\to L_x$, $a_{Tx}\mapsto u\phi_xa_{Tx}$. Thus the map from
the annular charts to $L^\times/\langle u\phi\rangle$ is holomorphic
and identifies exactly the orbits of this generator. Every orbit
meets the core or an attached collar by contraction, and two such
representatives are related by the displayed chart identification.
The comparison map is consequently a biholomorphism. At $u=0$,
(E1) is the pinching chart $ab=0$.

\item The total node chart is smooth because the derivative of its defining
equation with respect to $u$ is $-\phi_x\ne0$. The projection to $u$
is flat: its local rings are torsion free over the convergent
one-variable power-series ring. The joint map to $(t,u)$ is also
flat. On the core it is a product of $\pi$ with the parameter disc;
$\pi$ is flat since a nonconstant map from a smooth connected surface
to a smooth curve has torsion-free local rings over the curve. On a
node chart, $t-t_0$ and $u-u_0$ form a regular sequence in the regular
total local ring: after setting $t=t_0$, the equation (E1) has no
component contained in a level set $u=u_0$. The two elements cut
dimension by two. The local flatness criterion for a regular base
(equivalently its elementary regular-sequence proof) gives flatness
over $(t,u)$. This argument uses smoothness of the surface $S$, not
smoothness of the Kodaira fibre; it therefore applies at every point
of the fibres $IV^*$ and $III$.
\end{enumerate}
\end{proof}

\begin{proposition}\label{fam:src:E.1.3}\label{fam:st:E.1.3}
\StatementRecord{fam:src:E.1.3}{proved}{Historical events: Claude Fable 5.2 (Anthropic), effort high, 12 November 2026: targeted proof check. GPT-6 Astra (OpenAI), ChatGPT Pro mode, effort Pro, S6\_09\_PAPERS, 17 November 2026: targeted proof check. Also checked in this writing. Pre-continuation writing events: GPT-6 Astra (OpenAI), ChatGPT Pro mode, effort Pro, S6\_10\_PAPERS, 18 November 2026: targeted proof check. Continuation metadata events: GPT-6 Astra (OpenAI), ChatGPT Pro mode, effort Pro, S6\_10\_PAPERS continuation, 19 November 2026: used and rechecked in this run; status-vocabulary reconciliation; metadata reconciliation only; no new mathematical check. \emph{Pre-continuation recorded proof scope.} Read complete source proof and re-derived its explicit pushout/plumbing rings, collar gluing, flatness or finite quotient, and integral lattice/monodromy or fan maps as applicable. No proof credit is inferred from compilation; exact algebraic replays and cited source matching are recorded separately.}{the hypotheses, notation and definitions stated here; the explicit constructions and calculations in the proof}

The first simultaneous local plumbing at the nodal point of $S_\infty$
has equation
\[
                 xy=u zv,\qquad q=zv,                         \tag{E2}
\]
after a holomorphic change of one normal coordinate, where $q$ is a
local coordinate on $B$ at infinity. This hypersurface is flat over
the $u$-disc. For each $u\ne0$, its three-dimensional fibre has an
ordinary double point at $x=y=z=v=0$. The projective modification
defined by the blow-up of $(x,z)$ resolves that point for $u\ne0$.
Its two strict-transform charts are
\[
       x=zr,\quad ry=uv;
       \qquad z=xa,\quad y=uav.                               \tag{E3}
\]
The modified family remains flat over $u$. At $u=0$, the normalization
of its central fibre contains the blow-up of one branch of (E2) along
the smooth curve $x=z=0$; hence this particular modification changes
the normalization of the unmodified pinching.
\end{proposition}



\begin{proof}
\begin{enumerate}[label=\arabic*.,leftmargin=*]
\item The surface $S$ is smooth and its nodal elliptic fibre has local
equation $q=zv$ at its node. The section $\phi$ has a simple zero
along that fibre. Thus $\phi=zv$ times an invertible holomorphic
function. Absorbing the function in one normal coordinate gives (E2).
The convergent power-series ring $\C\{x,y,z,v,u\}$ is factorial,
and $u$ is coprime to $xy-uzv$. Hence $uh\in(xy-uzv)$ implies
$h\in(xy-uzv)$. The quotient has no $u$-torsion. At a nonzero
parameter the identical argument uses $u-u_0$. This proves flatness
over the smooth parameter curve.
For fixed $u\ne0$ the four first derivatives vanish only at the
origin, and the quadratic form $xy-uzv$ is nondegenerate. This proves
the ordinary-double-point assertion by its defining equation.

\item The two standard affine charts of the blow-up give (E3), after
removing the common factors $z$ and $x$, respectively. For $u\ne0$
the first chart solves for $v$ and the second for $y$, so both are
smooth. The map is an isomorphism away from the origin in the fixed
$u$-fibre: where the ideal fails to be visibly principal, the
hypersurface equation and one of $y,v$ make it principal. The fibre
over the origin is $\mathbb P^1$. The modification is projective
over (E2), since it is a closed subspace of its product with
$\mathbb P^1$. The first strict-transform equation is coprime to
$u-u_0$ in its convergent power-series ring for every $u_0$; the
second chart is smooth with free parameter $u$. Thus flatness over
$u$ persists.

\item At $u=0$ the first chart has $ry=0$. The branch $r=0$ maps
isomorphically to $x=0$ in the original central fibre. The branch
$y=0$, together with the second chart, is the ordinary blow-up of the
branch $y=0$ in coordinates $(x,z,v)$ along $(x,z)$. This description
follows from the two affine substitutions defining that blow-up and
is an identification of normalizations, not just a birational map.
For $u\ne0$, the base coordinate in the first chart is
$q=zry/u$ and in the second it is $q=xav$. Thus the modification also
exhibits the triple-coordinate strata that the cusp filling requires.
The substitutions and Jacobian are independently checked in
\texttt{code/e\_degeneration.py}.
\end{enumerate}
\end{proof}

\begin{proposition}\label{fam:src:E.1.4}\label{fam:st:E.1.4}
\StatementRecord{fam:src:E.1.4}{proved}{Historical events: Claude Fable 5.2 (Anthropic), effort high, 12 November 2026: targeted proof check. Claude Opus 5.5 (Anthropic), effort max, sweep 1, scope s6, 18 November 2026: computation or formal certificate. GPT-6 Astra (OpenAI), ChatGPT Pro mode, effort Pro, S6\_09\_PAPERS, 17 November 2026: targeted proof check. Also checked in this writing. Pre-continuation writing events: GPT-6 Astra (OpenAI), ChatGPT Pro mode, effort Pro, S6\_10\_PAPERS, 18 November 2026: targeted proof check. Continuation metadata events: GPT-6 Astra (OpenAI), ChatGPT Pro mode, effort Pro, S6\_10\_PAPERS continuation, 19 November 2026: used and rechecked in this run; status-vocabulary reconciliation; metadata reconciliation only; no new mathematical check. \emph{Pre-continuation recorded proof scope.} Read complete source proof and re-derived its explicit pushout/plumbing rings, collar gluing, flatness or finite quotient, and integral lattice/monodromy or fan maps as applicable. No proof credit is inferred from compilation; exact algebraic replays and cited source matching are recorded separately.}{the hypotheses, notation and definitions stated here; \cref{fam:src:E.1.1}; \cite[Lemma 2.1]{Engel2026}}

The compact candidate $Q$ of \ref{fam:src:E.1.1} has
\[
 \pi_1(Q)=\Z,\qquad
 (H_i(Q;\Z))_{i=0}^6
   =(\Z,\Z,\Z^8,\Z^7,\Z^{10},0,\Z),\qquad\chi(Q)=12.
\]
These are invariants of the stated candidate; they do not identify it
as a central fibre of the prescribed compact family.
\end{proposition}



\begin{proof}
\begin{enumerate}[label=\arabic*.,leftmargin=*]
\item The two section embeddings are cofibrations: a tubular neighbourhood
retracts onto each section, and the two neighbourhoods are disjoint.
The ordinary pushout is therefore homotopy equivalent to the space
formed from $R$ by attaching $S\times[0,1]$ along
$s_0$ and $s_\infty T$. This comparison follows by collapsing the
intervals in the mapping cylinder; the tubular retractions construct
its homotopy inverse. Both $S$ and $R$ are simply connected. Van
Kampen applied to this mapping-cylinder model gives one free loop
and no additional generator or relation, so $\pi_1(Q)=\Z$.
The same model has the exact sequence
\[
 H_i(S)\xrightarrow{d_i=s_{0*}-s_{\infty*}T_*}H_i(R)
 \longrightarrow H_i(Q)\longrightarrow H_{i-1}(S).             \tag{E4}
\]
We compute its integral maps, including primitivity.

\item Put $F=[S_t]$, $m=P-O$, and
$v_0=[6P]-[O]-3F\in H_2(S;\Z)$. The intersection computations in
\cite[Lemma 2.1]{Engel2026} give
\[
 F^2=Fv_0=Fm=0,\quad v_0^2=-6,\quad mv_0=-1,\quad m^2=-2,
 \quad OF=1,\quad Ov_0=0,\quad Om=1.                         \tag{E5}
\]
The translation preserves all components of the two reducible
Kodaira fibres. Their root classes span a negative definite space of
rank seven. Its orthogonal complement is generated over $\Q$ by
$F,O,v_0$. Since $T_*O=O+v_0+3F$, preservation of the intersection
form gives, on this complement and hence on the entire lattice,
\[
 T_*a=a+(aF)v_0-\bigl(av_0-3aF\bigr)F.                       \tag{E6}
\]
For clarity, its value on $v_0$ is $v_0+6F$: it has no $O$-component
and has zero intersections with the rank-seven root space, so it is
$\lambda v_0+aF$. The already established identity $T^*m=m-F$ gives
$mT_*v_0=(m-F)v_0=-1$, forcing $\lambda=1$. Orthogonality to
$T_*O$ then forces $a=6$. These
values on $F,O,v_0$ and the root space verify (E6) on a rational basis,
and both sides preserve the actual integral lattice. In particular
$T^*m=m-F$, in agreement with \ref{fam:src:E.1.1}.

\item The rational elliptic surface is the blow-up of $\mathbb P^2$ at nine
points, so its homology in even degrees is $\Z,\Z^{10},\Z$ and
vanishes in odd degrees; its intersection lattice is unimodular.
The projective-bundle theorem follows here by local trivialization
and the fibre hyperplane class: $H_2(R)=H_2(S)\oplus\Z f$ and
$H_4(R)=\Z\oplus H_2(S)$, with $f$ the projective-line fibre.
The hyperplane class that vanishes on $s_0$ restricts to $-m$ on
$s_\infty$. Therefore
\[
 d_2(a)=(1-T_*)a+(mT_*a)f
       =-(aF)v_0+(av_0-3aF)F+(am-aF)f.                       \tag{E7}
\]
The span of $F,v_0,m$ is primitive. To check the only nontrivial
part of this assertion, split off the unimodular plane spanned by
$F,O$; $m-F$ and $v_0$ lie in its negative definite orthogonal
complement. Their Gram determinant is
$\det\left(\begin{smallmatrix}-2&-1\\-1&-6\end{smallmatrix}\right)=11$.
The square of an index of a larger integral lattice would divide
eleven, so the index is one. The span of $F,v_0$ is also primitive:
pairing with $O,m$ gives determinant $-1$. Unimodularity consequently
makes $a\mapsto(aF,av_0,am)$ surjective onto $\Z^3$. The coefficient
matrix in (E7) has determinant $-1$. Thus $d_2$ has primitive image
$\Z v_0\oplus\Z F\oplus\Z f$, rank three, kernel $\Z^7$,
and cokernel $\Z^8$.

\item The map $d_4$ sends $[S]$ to $s_0[S]-s_\infty[S]$, whose Poincare
dual is $\pi_R^*m$. This is primitive and nonzero since $mO=1$.
Thus $d_4$ is injective with cokernel $\Z^{10}$. The map $d_0$ is
zero. Substitution in (E4) proves every asserted homology group.
The Euler characteristic is independently obtained by stratification:
$\chi(R)-\chi(S)=24-12=12$. The Smith computations and both Euler
computations are recorded in \texttt{code/e\_degeneration.out}.
\end{enumerate}
\end{proof}


\paragraph{Local obstruction and the precise gluing requirement}\label{fam:src:E.1.5}\label{fam:st:E.1.5}
\StatementRecord{fam:src:E.1.5}{proved}{Historical events: Claude Fable 5.2 (Anthropic), effort high, 12 November 2026: targeted proof check. Claude Fable 5.2 (Anthropic), effort high, 17 November 2026: carried as background without fault found. GPT-6 Astra (OpenAI), ChatGPT Pro mode, effort Pro, S6\_09\_PAPERS, 17 November 2026: targeted proof check. Also checked in this writing. Pre-continuation writing events: GPT-6 Astra (OpenAI), ChatGPT Pro mode, effort Pro, S6\_10\_PAPERS, 18 November 2026: targeted proof check. Continuation metadata events: GPT-6 Astra (OpenAI), ChatGPT Pro mode, effort Pro, S6\_10\_PAPERS continuation, 19 November 2026: used and rechecked in this run; status-vocabulary reconciliation; metadata reconciliation only; no new mathematical check. \emph{Pre-continuation recorded proof scope.} Re-derived the precise local-to-global scope through the actual xy=uzv charts, blow-up normalization and central-cycle obstruction; supplied a separate numbered receiving proof.}{\cref{fam:src:E.1.1}, \cref{fam:src:E.1.2}, \cref{fam:src:E.1.3}; \cref{push:src:A.1.1}, \cref{push:src:A.2.1}, \cref{push:src:A.2.2}, \cref{push:src:A.3.1}, \cref{push:src:A.3.2}}

The following local construction isolates the obstruction to a flat proper filling of the prescribed compact untwisted family with central fibre exactly $Q$. The full nonexistence result is proved in \cref{push:src:A.3.1}, \cref{push:src:A.3.2}. \ref{fam:src:E.1.1} constructs the compact candidate as an actual complex
space; \ref{fam:src:E.1.2} constructs its smoothing away from infinity, including
the two other special fibres; \ref{fam:src:E.1.3} computes the two-parameter germ and a concrete flat modification of it.

The local requirement identified by this argument is the following analytic gluing
problem at one lift of the node of $S_\infty$. The local model
$xy=uzv$ must be replaced by charts whose restriction to every
$u\ne0$ is the chosen small resolution of the rank-two Mumford
filling, and these charts must be equivariant under both period
generators, with proper quotient and with central normalization $R$.
The explicit blow-up in \ref{fam:src:E.1.3} supplies the required smooth local
general fibre but changes the normalization at $u=0$. This particular blow-up does not settle whether another gluing can preserve that normalization. Thus the local computation alone is an obstruction to the unmodified plumbing argument. The global virtual-Euler argument in \cref{push:src:A.3.1}, \cref{push:src:A.3.2} supplies the exclusion. A modified compact limit
also requires an equivariant periodic fan and verification of its
quotient near the crossing $(u,q)=(0,0)$; the complete fan and quotient data are supplied in \cref{push:src:A.1.1}, \cref{push:src:A.2.1}, \cref{push:src:A.2.2}.



\begin{proof}
\begin{enumerate}[label=\arabic*.,leftmargin=*]
\item The charts of \ref{fam:src:E.1.1} construct $Q$ and its normalization
$R$. On the complement of the crossing, \ref{fam:src:E.1.2}
constructs the smoothing and its finite-fibre gluing maps.
At the crossing, substitution in those charts gives $xy=uzv$.
The blow-up charts in \ref{fam:src:E.1.3} restrict to the chosen
small resolution for $u\ne0$ and have the explicitly computed
changed normalization at $u=0$. They therefore settle precisely
that local modification, while the preservation of $R$ is an
additional gluing requirement on any proposed replacement.

\item For the proposed compact centre exactly equal to $Q$, the central
cycle pushforward and virtual Euler calculation in
\ref{push:src:A.3.1}, \ref{push:src:A.3.2} give the stated global
nonexistence. The equivariant fan construction in
\ref{push:src:A.1.1}, \ref{push:src:A.2.1}, \ref{push:src:A.2.2}
supplies the explicit different central model and its proper
quotient. These exact restrictions of the charts prove each
local and global assertion of the remark.
\end{enumerate}
\end{proof}

\Uses{the hypotheses, notation and definitions stated here; \cref{push:src:A.3.1}, \cref{push:src:A.3.2}; \cref{fam:src:E.1.1}; \cref{fam:src:E.1.2}; \cref{fam:src:E.1.3}; \cref{push:src:A.1.1}, \cref{push:src:A.2.1}, \cref{push:src:A.2.2}}


\subsection{Finite bielliptic limits and the smooth monodromy}\label{fam:section:E.2}

\begin{proposition}\label{fam:src:E.2.1}\label{fam:st:E.2.1}
\StatementRecord{fam:src:E.2.1}{proved}{Historical events: Claude Fable 5.2 (Anthropic), effort high, 12 November 2026: targeted proof check. GPT-6 Astra (OpenAI), ChatGPT Pro mode, effort Pro, S6\_09\_PAPERS, 17 November 2026: targeted proof check. Also checked in this writing. Pre-continuation writing events: GPT-6 Astra (OpenAI), ChatGPT Pro mode, effort Pro, S6\_10\_PAPERS, 18 November 2026: targeted proof check. Continuation metadata events: GPT-6 Astra (OpenAI), ChatGPT Pro mode, effort Pro, S6\_10\_PAPERS continuation, 19 November 2026: used and rechecked in this run; status-vocabulary reconciliation; metadata reconciliation only; no new mathematical check. \emph{Pre-continuation recorded proof scope.} Read complete source proof and re-derived its explicit pushout/plumbing rings, collar gluing, flatness or finite quotient, and integral lattice/monodromy or fan maps as applicable. No proof credit is inferred from compilation; exact algebraic replays and cited source matching are recorded separately.}{the hypotheses, notation and definitions stated here; the explicit constructions and calculations in the proof}

After the auxiliary base change $w=s^6$, each of the two families of
\emph{reduced} special surfaces $S_j(w)$ has a proper flat
semistable filling by explicit finite free quotients. The central
surface for $j=1$ is a cycle of four ruled surfaces over an elliptic
curve; that for $j=2$ is a cycle of three. They have
\[
\begin{array}{c|ccccc|c}
 &H_0&H_1&H_2&H_3&H_4&\chi\\\hline
 j=1&\Z&\Z\oplus\Z/3&\Z^5&\Z^9&\Z^4&0\\
 j=2&\Z&\Z\oplus\Z/2&\Z^4&\Z^7&\Z^3&0
\end{array}
\]
with integral coefficients. Their fundamental groups are
$\Z^2\rtimes_{R_3}\Z$ and $\Z^2\rtimes_{R_4}\Z$, respectively,
where
\[
 R_3=\begin{pmatrix}0&-1\\1&-1\end{pmatrix},\qquad
 R_4=\begin{pmatrix}0&-1\\1&0\end{pmatrix}.
\]
\end{proposition}



\begin{proof}
\begin{enumerate}[label=\arabic*.,leftmargin=*]
\item We retain the actual product presentations, which also specify the
comparison with the surfaces in the stated problem. Let
$\rho=e^{\pi i/3}$,
\[
 E_j=\C/(\Z+\kappa_j\Z),\quad
 F_1=\C/(\Z(\rho-2)+\Z(-3)),\quad
 F_2=\C/(\Z(i-1)+\Z(-2)),
\]
where $\kappa_1=\beta_c(z_1)+2\mu(z_1)$ and
$\kappa_2=\beta_c(z_2)+3\mu(z_2)$ have negative imaginary parts.
The holomorphic maps to the central tori are induced by
$L_1(\xi,\eta)=(\eta,\xi-\eta/3)$ and
$L_2(\xi,\eta)=(\eta,\xi-\eta/2)$. Their kernels are generated by
\[
 k_1=(2/3,\rho),\qquad k_2=(1/2,i).
\]
The residual affine generators lift to
\[
 h_1(\xi,\eta)=(\xi+\kappa_1/3,(\rho-1)\eta),\qquad
 h_2(\xi,\eta)=(\xi-\kappa_2/4,-i\eta).                     \tag{E8}
\]
These formulas are verified by substituting the four product periods
in $L_j$; the lattice indices are three and two. Moreover
$(\rho-1)\rho-\rho$ and $(-i)i-i$ are periods of the displayed
$F_j$, so the generators commute. They give free groups
$G_1=C_3\times C_3$ and $G_2=C_2\times C_4$. Indeed, equality of
the first coordinates of $k_j^a h_j^b(\xi,\eta)$ and $(\xi,\eta)$
forces $b=0$ from the coefficient of the nonreal period $\kappa_j$,
and then $a=0$ from the real period. Thus the composite quotient
maps are exactly the degree-nine and degree-eight covers of $S_j$,
with the given signs, rather than abstract surface isogenies.

\item Use the multiplicative coordinate $v=e^{-2\pi i\xi}$ on $E_j$.
Its Tate parameter is
\[
       q_j=e^{-2\pi i\kappa_j}=A_jw^2=A_js^{12},
       \qquad A_j=e^{-2\pi i\kappa_j^{\mathrm{part}}}\ne0.
\]
Here $A_j$ denotes a scalar only within this proof; it is not a
monodromy matrix. Choose specifically
$A_j^{1/12}=\exp(-2\pi i\kappa_j^{\mathrm{part}}/12)$ and put
$r_j=A_j^{1/12}s$. This choice preserves the fractional period
translations in (E8) exactly. The regular Tate model for $q_j=r_j^{12}$
has a central twelve-gon. We recall its charts to make the extension
explicit. An infinite chain has node charts
$x_i y_i=r_j$; on a branch overlap one uses
$x_{i+1}=y_i^{-1}$ and $y_{i+1}=r_jy_i$. The coordinate
$r_j^i x_i$ on a general fibre is unchanged on overlaps.
The shift $i\mapsto i+12$ acts on that coordinate by $r_j^{12}$.
Quotienting a tubular neighbourhood of the chain by this shift gives
the proper twelve-gon family, and on $r_j\ne0$ it is
$\C^*/q_j^\Z$. Compactness follows from twelve projective components
and a compact annular fundamental domain on a general fibre;
the node charts give flatness and smoothness of its total space.

\item Multiplication by $q_1^{1/3}=r_1^4$ shifts the component index by four.
Multiplication by $q_2^{-1/4}=r_2^{-3}$ shifts it by minus three.
Hence the maps $h_j$ in (E8) extend, together with their rotations
on $F_j$. The maps $k_j$ multiply $v$ by roots of unity and translate
$F_j$, and extend to every projective component. Their relations
hold on the dense set $s\ne0$ and consequently on all charts.
The action is free on the central fibre as well. A nontrivial power
of $h_j$ moves every component and every node. A nontrivial power
of $k_j$ translates $F_j$ by a nonzero torsion point, including
above the nodes where the multiplicative action alone fixes a point.
The finite free quotient is therefore a smooth total space, proper
over $s$, with a reduced normal-crossing central surface. Flatness
descends along the finite etale quotient. The component orbits have
cardinality three and four, leaving four and three components.
Quotienting an individual $\mathbb P^1\times F_j$ by $k_j$ is a
ruled surface over $F_j/\langle\rho\rangle$ or
$F_j/\langle i\rangle$, respectively. Its two distinguished
sections have the same integral homology class: the corresponding
line bundle has finite order and degree zero on the elliptic base.

\item For the topology, unwrap the component cycle into an infinite chain
and then pass to the universal cover $\C$ of its elliptic base.
The resulting union of simply connected ruled pieces, with
simply connected intersections and a tree as incidence graph, is
simply connected by van Kampen. Its deck group is generated by
the enlarged lattice of translations on $F_j$ and one step round
the component cycle. The latter conjugates the lattice by the
order-three or order-four rotation. This gives the asserted
semidirect products; inverses of the displayed rotation matrices
give isomorphic groups by reversing the cycle generator.

\item For homology, the normalization has $n=4$ or $3$ components, each
with ranks $(1,2,2,2,1)$, and $n$ elliptic double curves, each with
ranks $(1,2,1)$. There are no triple points. The mapping-cylinder
sequence for the normalization and the double curves applies
integrally. In degree one the incidence matrix is a twisted cycle:
its cokernel is $\operatorname{coker}(R_j-I)$ and its kernel is
zero, because $R_j$ has no eigenvalue one. The degree-zero incidence
matrix contributes an additional $\Z$ to $H_1$. In degree two the
two section classes on a component agree, so the incidence matrix
is the ordinary cycle matrix on their classes. It has kernel
$\Z$ and primitive image of rank $n-1$. Thus $H_2=\Z^{n+1}$,
$H_3=\Z^{2n+1}$, and $H_4=\Z^n$. Finally
$\operatorname{coker}(R_3-I)=\Z/3$ and
$\operatorname{coker}(R_4-I)=\Z/2$. The exact Smith forms of both
full cycle matrices and the independent $2\times2$ determinants
are in \texttt{code/e\_degeneration.out}.
\end{enumerate}
\end{proof}








\section{The untwisted limit}\label{push:section:A}
\subsection{The nodal plumbing}\label{push:section:A.1}

We write $L=\mathcal O_S(P-O)$, $T=t_{6P}$,
$J=L\otimes T^*L^{-1}=\pi^*\mathcal O_B(1)$, and
$R=\mathbb P_S(\mathcal O\oplus L)$. The convention for this projective
bundle is that its zero and infinity sections have normal bundles
$L$ and $L^{-1}$. The identification defining $Q$ is
$s_0(x)\sim s_\infty(Tx)$. The equality defining $J$ is the equality
in Engel's Lemma 2.1, with this direction of $T$; its intersection
calculation uses $P\cdot O=0$, $P^2=O^2=-1$,
$6P\cdot O=6P\cdot P=2$. Thus the degree on $O$ is
$(O-P)\cdot6P+(P-O)\cdot O=1$. We use the section
$\phi\in H^0(S,J)$ whose divisor is $S_\infty$, as in his
Corollary 2.2. We use the stated version of
\cite{Engel2026}, as is shown by its explicit local equations.

\begin{proposition}\label{push:src:A.1.1}\label{push:A11}
\StatementRecord{push:src:A.1.1}{proved}{Historical events: Claude Fable 5.2 (Anthropic), effort high, 16 November 2026: carried as background without fault found; used and rechecked in this run. GPT-6 Astra (OpenAI), ChatGPT Pro mode, effort Pro, S6\_09\_PAPERS, 17 November 2026: targeted proof check. Also checked in this writing. Pre-continuation writing events: GPT-6 Astra (OpenAI), ChatGPT Pro mode, effort Pro, S6\_10\_PAPERS, 18 November 2026: targeted proof check. Continuation metadata events: GPT-6 Astra (OpenAI), ChatGPT Pro mode, effort Pro, S6\_10\_PAPERS continuation, 19 November 2026: used and rechecked in this run; status-vocabulary reconciliation; metadata reconciliation only; no new mathematical check. \emph{Pre-continuation recorded proof scope.} Read complete source proof and re-derived its explicit pushout/plumbing rings, collar gluing, flatness or finite quotient, and integral lattice/monodromy or fan maps as applicable. No proof credit is inferred from compilation; exact algebraic replays and cited source matching are recorded separately.}{the hypotheses, notation and definitions stated here; the explicit constructions and calculations in the proof}

For $k,l\in\mathbb Z$, let $P_{k,l}$ be the cone defined by $h_q,h_u\geq0$, $kh_q\leq y_1\leq(k+1)h_q$ and $l(h_q+h_u)\leq y_2\leq(l+1)(h_q+h_u)$. These cones and their faces form a fan. Its two
period automorphisms on the cocharacter lattice are
\[
 G_1(y_1,y_2,h_q,h_u)=(y_1+h_q,y_2,h_q,h_u),\qquad
 G_2(y_1,y_2,h_q,h_u)=(y_1,y_2+h_q+h_u,h_q,h_u).
\]
The affine chart of $P_{0,0}$ is
\[
 U_{0,0}=\operatorname{Spec}\mathbb C[x,y,z,v,u]/(xy-uzv),
 \qquad q=zv.
\]
In particular the unsubdivided chart is not a smooth simultaneous
model of the cusp.
\end{proposition}


\begin{proof}
\begin{enumerate}[label=\arabic*.,leftmargin=*]
\item At fixed $h_q,h_u\geq0$, the cone $P_{k,l}$ cuts out the rectangle
\[
 kh_q\leq y_1\leq(k+1)h_q,\qquad
 l(h_q+h_u)\leq y_2\leq(l+1)(h_q+h_u).
\]
When $h_q=0$ the first interval is a point. These rectangles tile
their support, and their intersections are common faces, including
this collapsed boundary. This proves the fan assertion. The displayed
integral transformations preserve the fan and have integral inverses.

\item For a character $(a,b,c,d)$, the six inequalities of the dual cone
are $c,a+c,b+c,a+b+c,d,b+d\geq0$. The characters
\[
 x=(0,1,0,0),\quad y=(0,-1,1,1),\quad
 z=(1,0,0,0),\quad v=(-1,0,1,0),\quad u=(0,0,0,1)
\]
generate its integral monoid. Indeed if $b\geq0$ we subtract $b x$;
if $b<0$ we subtract $(-b)y$. The remaining character has second
coordinate zero and is a nonnegative integral combination of $z,v,u$.
The sole relation is $x+y=u+z+v$, and the character of the first
base coordinate is $z+v$. This proves the coordinate-ring assertion,
including normality. The calculation is checked exactly in
\texttt{code/ab\_push\_local\_checks.py}.
\end{enumerate}
\end{proof}

\begin{theorem}\label{push:src:A.1.2}\label{push:A12}
\StatementRecord{push:src:A.1.2}{proved}{Historical events: Claude Fable 5.2 (Anthropic), effort high, 16 November 2026: carried as background without fault found; used and rechecked in this run. GPT-6 Astra (OpenAI), ChatGPT Pro mode, effort Pro, S6\_09\_PAPERS, 17 November 2026: targeted proof check. Also checked in this writing. Pre-continuation writing events: GPT-6 Astra (OpenAI), ChatGPT Pro mode, effort Pro, S6\_10\_PAPERS, 18 November 2026: targeted proof check. Continuation metadata events: GPT-6 Astra (OpenAI), ChatGPT Pro mode, effort Pro, S6\_10\_PAPERS continuation, 19 November 2026: used and rechecked in this run; status-vocabulary reconciliation; metadata reconciliation only; no new mathematical check. \emph{Pre-continuation recorded proof scope.} Read complete source proof and re-derived its explicit pushout/plumbing rings, collar gluing, flatness or finite quotient, and integral lattice/monodromy or fan maps as applicable. No proof credit is inferred from compilation; exact algebraic replays and cited source matching are recorded separately.}{the hypotheses, notation and definitions stated here; \cref{push:A11}; \cite{EGFS25}}

There is a reduced complex space $\mathcal X^\flat$ and a proper flat
holomorphic map
\[
 (f,\varpi):\mathcal X^\flat\longrightarrow B\times\Delta_u
\]
whose central threefold is $Q$. For $u\ne0$ its fibre is the
contraction of Engel's chosen $X(u)$ at one rational curve of the
Mumford fibre. This curve is the double curve corresponding to the
diagonal of the square in the chosen $A_2$ triangulation; in the
return's coordinates its direction is $e_1-e_2$. Its normal bundle
in $X(u)$ is $\mathcal O_{\mathbb P^1}(-1)^{\oplus2}$. Its image is
the unique ordinary double point of $\mathcal X^\flat_u$.
\end{theorem}


\begin{proof}
\begin{enumerate}[label=\arabic*.,leftmargin=*]
\item Here is a construction on the whole compact surface $S$. Choose a
Hermitian metric on $L$, use the dual metric on $L^{-1}$, and choose
$0<r<1/4$. For sufficiently small $\epsilon>0$ we have
$|u|\|\phi_x\|<r^2$ for all $x\in S$ and $|u|<\epsilon$.
Take the core
\[
 C=\{(x,a,u):a\in L_x,\ r<\|a\|<r^{-1},\ |u|<\epsilon\}
\]
and the node space
\[
 N=\{(x,a,b,u):a\in L_x,\ b\in L^{-1}_{Tx},\
 a\otimes b=u\phi_x,\ \|a\|,\|b\|<2r,\ |u|<\epsilon\}.
\]
On $r<\|a\|<2r$ identify a point of $N$ with $(x,a,u)\in C$.
On $r<\|b\|<2r$ identify it with $(Tx,b^{-1},u)\in C$.
Both identifications are biholomorphisms between open complex spaces;
in the first collar $b=u\phi_x/a$, and in the second
$a=u\phi_xb^{-1}$. Their domains are disjoint because
$\|a\|\|b\|<r^2$. Changes of frames multiply the equation by
units and give the stated bundle-valued equation, so no choice of
frame enters the gluing.

\item This gluing is Hausdorff. At either end of a collar, exactly one
of its two charts contains the limiting circle: the core omits
$\|a\|=r$, and the node chart omits $\|a\|=2r$; the infinity
collar has the identical property in $b$. Thus the graph of the
identification has no unpaired pair of limiting points. The two
collars in each chart are separated, so there are no additional
chains of identifications. These observations give disjoint
saturated neighbourhoods of any two inequivalent points. Properness
is equally concrete. Over a compact subset of $B\times\Delta_\epsilon$,
the closed core $3r/2\leq\|a\|\leq2/(3r)$ together with the
closed node subspace $\|a\|,\|b\|\leq7r/4$ is compact and maps
onto the glued space. To see surjectivity, a point outside the smaller
core lies on a node collar, and a node point outside the smaller
node subspace lies on the smaller core. Compactness of $S$ supplies
the uniform bound used above. This proves properness of the joint
map to $(t,u)$.

\item The total node rings are hypersurface rings over the regular
two-dimensional local rings of $S$. They are Cohen--Macaulay.
Every fibre of the joint map has pure dimension two: first fix
$t$, and then impose $xy=u\phi$ on the two line coordinates over
the one-dimensional fibre of $S$. Consequently $t-t_0,u-u_0$ is a
regular sequence in every total local ring. The regular-sequence
flatness criterion over the smooth two-dimensional base proves
flatness. On the core the same conclusion is immediate from
flatness of $S\to B$. This includes the nonreduced Kodaira fibres
at $0,1$; smoothness of $S$, rather than smoothness of those fibres,
is what the argument uses. At $u=0$ the node equation is $ab=0$,
and the two collar maps recover precisely the specified pinching
of the two sections of $R$.

\item For $u\ne0$ and $t\ne\infty$, the equation gives
$a_{Tx}=a_x/(u\phi_x)$. Hence the resulting quotient is exactly
$L^\times/\langle u\phi\rangle$, where the contracting generator
covers $T^{-1}$. Every orbit meets a core or a collar, by the
uniform contraction estimate; the displayed comparison identifies
exactly its orbit relation. It is therefore a biholomorphism,
not just a meromorphic comparison.

\item At the nodal point of $S_\infty$, write $q=zv$. Absorbing a unit
in one line coordinate gives $xy=uzv$. At a smooth point of
$S_\infty$ the equation is $xy=uq$. For $u\ne0$ the latter
threefold is smooth; the former has one nondegenerate quadratic
singularity at the origin. There are no other singularities of a
nonzero $u$ fibre.

\item For completeness, the relation to the infinite fan quotient is
also explicit. In a Tate covering of a neighbourhood of the $I_1$
fibre use the torus coordinates $(Z,W)$ of Engel's Proposition 2.7.
The two periods act as
\[
 (Z,W,q,u)\longmapsto(c_1qZ,c_2W,q,u),\qquad
 (Z,W,q,u)\longmapsto(c_3Z,c'_4uqW,q,u),
\]
with holomorphic units $c_i$. Their monomial parts are $G_1,G_2$.
The chart indexed by $(k,l)$ has
$z=Zq^{-k}$, $v=q^{k+1}/Z$, $x=W(qu)^{-l}$,
$y=qu/x$, with the harmless unit changes supplied by the line
frames. It is the chart of Proposition~\ref{push:A11}.
The Tate cover unfolds the node of $S_\infty$ into its infinite
chain, and unfolding the two plumbing collars gives the additional
integer index $l$. The foregoing collar identifications are exactly
the two deck identifications. Thus, on an invariant neighbourhood
over a sufficiently small bidisc, the fan space maps by a
holomorphic covering to the plumbing, with deck group $\mathbb Z^2$.
Each core or node coordinate neighbourhood has disjoint lifts
indexed by the two unfolded integers; this proves freeness and
proper discontinuity, including the $\rho_l$ orbits. The first
period stabilizes such a ray but on its orbit translates the Tate
coordinate by $c_1q$, with $q\ne0$ and $|c_1q|<1$; it does not
fix a point. Properness of its quotient was proved by the compact
core and node representatives. This supplies the analytic
comparison needed here. Construction 3.8 of \cite{EGFS25} is a
related multivariable polarized construction; its symmetry
hypothesis on period units is not being silently imposed on these
generally unpolarized tori.

\item Subdividing the square at $h_u=0$ by the chosen diagonal gives
the usual small resolution of $xy=uzv$ for $u\ne0$. Its exceptional
curve has the two affine charts of the total space of
$\mathcal O_{\mathbb P^1}(-1)\oplus
\mathcal O_{\mathbb P^1}(-1)$, as follows by setting $x=za$ on
one chart and $z=xb$ on the other. The two normal coordinates
transform by the reciprocal affine coordinate. There is one
diagonal edge modulo the periods, so there is exactly one such
curve. Removing that edge contracts the paired opposite
$(-1)$-curves in the normalization $dP_6$ to the square model
$\mathbb P^1\times\mathbb P^1$. This identifies the contracted
curve intrinsically in the selected toric model. It is not the
$e_2$-direction curve fixed by $\operatorname{Aut}^0$.
\end{enumerate}
\end{proof}

\subsection{Two minimal toroidal fillings and further modifications}\label{push:section:A.2}

\begin{proposition}\label{push:src:A.2.1}\label{push:A21}
\StatementRecord{push:src:A.2.1}{proved}{Historical events: Claude Fable 5.2 (Anthropic), effort high, 16 November 2026: carried as background without fault found; used and rechecked in this run. GPT-6 Astra (OpenAI), ChatGPT Pro mode, effort Pro, S6\_09\_PAPERS, 17 November 2026: targeted proof check. Also checked in this writing. Pre-continuation writing events: GPT-6 Astra (OpenAI), ChatGPT Pro mode, effort Pro, S6\_10\_PAPERS, 18 November 2026: targeted proof check. Continuation metadata events: GPT-6 Astra (OpenAI), ChatGPT Pro mode, effort Pro, S6\_10\_PAPERS continuation, 19 November 2026: used and rechecked in this run; status-vocabulary reconciliation; metadata reconciliation only; no new mathematical check. \emph{Pre-continuation recorded proof scope.} Re-derived all six prism triangulations, the two periodic bottom-BC choices, the exclusion of mixed rays under joint flatness, and the arbitrary-length disjoint point-star family. Fully read exact ab\_push\_fan.py and reran successfully; all six triangulations reproduced.}{the hypotheses, notation and definitions stated here; \cref{push:src:A.1.1}}

Among subdivisions with no new rays, the fan in \ref{push:src:A.1.1} has exactly
two smooth period-invariant subdivisions restricting to the chosen
$A_2$ triangulation at $h_u=0$. These are also all subdivisions
which are flat over the original $(q,u)$ bidisc, have reduced
$u$-central divisor, and retain that boundary triangulation.
Without the joint-flatness restriction there are infinitely many
smooth invariant subdivisions retaining that triangulation and
having reduced $u$-central divisor.
\end{proposition}


\begin{proof}
\begin{enumerate}[label=\arabic*.,leftmargin=*]
\item Write the six rays of $P_{0,0}$ as
$A=(0,0,1,0)$, $B=(1,0,1,0)$, $C=(0,1,1,0)$,
$D=(1,1,1,0)$, $E=(0,0,0,1)$, $F=(0,1,0,1)$.
The return's $A_2$ diagonal is $BC$. A triangular prism has six
triangulations without new vertices. With this bottom diagonal
the three possibilities are
\[
\begin{array}{c|c|c}
\text{maximal tetrahedra}&\text{left diagonal}&\text{right diagonal}\\\hline
ABCE,\ BCDE,\ CDEF&CE&DE\\
ABCE,\ BCDF,\ BCEF&CE&BF\\
ABCF,\ ABEF,\ BCDF&AF&BF.
\end{array}
\]
Periodicity identifies the left and right faces by $A\mapsto B$,
$C\mapsto D$, $E\mapsto E$, $F\mapsto F$; it excludes exactly
the middle row. Every listed maximal cone has determinant $\pm1$.
There are no missed triangulations: the prism has normalized volume
three, all its nondegenerate four-vertex tetrahedra have normalized
volume one, and enumeration of their triples with pairwise
face intersections gives exactly these six possibilities before
fixing the bottom. The program \texttt{ab\_push\_fan.py} performs
this enumeration in rational arithmetic; intersections are tested
by all vertices of the combined barycentric halfspaces.

\item For the stronger assertion, a new ray with $h_q,h_u>0$ would
give a divisor mapping to the origin of the two-dimensional base.
Its fibre dimension would be three, contradicting joint flatness.
A ray with $h_u=0$ is excluded by the retained boundary
triangulation. On $h_q=0$, the fan support has $y_1=0$.
Reducedness requires the primitive ray to have $h_u=1$, so it
is already one of the $\rho_l$. Thus there are no new rays under
the stated three restrictions. Conversely the two triangulations
have all heights zero or one, and every cone maps onto a face of
the quadrant with saturated image. Their local monomial maps
are flat and their $u$ divisors reduced.

\item If joint flatness is omitted, star-subdivide the orbit of the
maximal cone $ABCE$ by its interior ray
$v_1=A+B+C+E=(1,1,3,1)$. Distinct translates have disjoint
interiors and distinct closed torus orbits. These are simultaneous
ordinary blow-ups at disjoint points, so the operation is
unambiguously period-invariant. Repeat in the cone $ABCv_{m-1}$,
where $v_0=E$. The new ray is
$v_m=v_{m-1}+A+B+C=(m,m,3m,1)$. Replacing any basis vector
by their sum preserves determinant one, and every new ray has
$h_u=1$. The boundary triangulations are unchanged. On the
compact quotient this blows up one smooth central point away
from the conductor at each step. There are therefore infinitely
many invariant smooth refinements, all reduced over $u=0$, with
increasing numbers of rays modulo the period group. They are
not jointly flat. An earlier attempted witness using an orbit
of two-faces did not justify compatibility of intersecting
centres; its correction is retained in
\path{evidence/ab_push_correction_orbit_stars.md}.
\end{enumerate}
\end{proof}

\begin{theorem}\label{push:src:A.2.2}\label{push:A22}
\StatementRecord{push:src:A.2.2}{proved}{Historical events: Claude Fable 5.2 (Anthropic), effort high, 16 November 2026: carried as background without fault found; used and rechecked in this run. GPT-6 Astra (OpenAI), ChatGPT Pro mode, effort Pro, S6\_09\_PAPERS, 17 November 2026: targeted proof check. Also checked in this writing. Pre-continuation writing events: GPT-6 Astra (OpenAI), ChatGPT Pro mode, effort Pro, S6\_10\_PAPERS, 18 November 2026: targeted proof check. Continuation metadata events: GPT-6 Astra (OpenAI), ChatGPT Pro mode, effort Pro, S6\_10\_PAPERS continuation, 19 November 2026: used and rechecked in this run; status-vocabulary reconciliation; metadata reconciliation only; no new mathematical check. \emph{Pre-continuation recorded proof scope.} Re-derived the actual charts for the blow-up of (X,zv), conifold small resolution, normal-product J(-S\_infty), reduced two-branch total equation, and proper descent from the period covering.}{the hypotheses, notation and definitions stated here; \cref{push:src:A.2.1}; \cref{push:A12}; \cref{push:src:A.1.2}}

The two minimal subdivisions of \ref{push:src:A.2.1} give proper flat holomorphic
families $\mathcal X^+,\mathcal X^-\to\Delta_u$ with smooth total
spaces. Their joint maps to $B\times\Delta_u$ are flat. Over
$\Delta_u^*$ they are isomorphic as families to the specified
untwisted family $X(u)$. Each central fibre $Q^\pm$ is reduced,
irreducible, and has normal crossings, with smooth double locus
$S$ and no triple points.

Its normalization $R^\pm$ is obtained by blowing up $R$ along
$S_\infty$ in, respectively, one of its two sections and then
taking the small resolution at the unique ordinary double point
of that blow-up prescribed by the fixed bottom diagonal.
The two proper transforms of the sections are still copies of
$S$, glued by $T$. Their normal-bundle product is trivial;
in particular $\mathcal T^1_{Q^\pm}\simeq\mathcal O_S$.
\end{theorem}


\begin{proof}
\begin{enumerate}[label=\arabic*.,leftmargin=*]
\item Apply either fan subdivision over the cusp neighbourhood in
Theorem~\ref{push:A12}, and use the identity modification elsewhere.
Toric modifications are proper, and a subdivision is an
isomorphism off the strata it subdivides. Its restriction over
$q\ne0$ is the identity, so this gluing has no additional
transition choice. The period actions preserve the subdivisions,
and the covering description in \ref{push:src:A.1.2} descends the proper
modifications to the quotient. All maximal cones are regular.
Every local ring of the resulting fourfold is therefore regular.
The coordinate $u$ is the product of those affine coordinates
whose rays have $h_u=1$; there are at most two such rays in a
maximal cone. This proves reduced normal crossings and the absence
of triple points for the $u$-central fibre. The regular-sequence
argument, or the saturated face images in \ref{push:src:A.2.1}, proves joint
flatness. Properness follows from properness of the toric
modification and of $\mathcal X^\flat$.

\item We give the normalization calculation, including the singular
centre. At a node of $S_\infty$, write the original normalization
as coordinates $(X,z,v)$, with section $X=0$ and centre
$(X,zv)$. Blowing up that ideal has charts
\[
 X=zvA;\qquad zv=XB.
\]
The first is smooth, and the second has one conifold at its
origin. A small resolution of the latter is, for one of the
bottom-diagonal conventions,
\[
 X=za,\quad v=aB;\qquad z=Xd,\quad B=dv.
\]
Interchanging $z,v$ gives the other diagonal convention. These
are holomorphic charts of a proper small modification: they are
the two usual charts of the blow-up of the appropriate rank-one
ideal on the conifold. The strict transform of $X=0$ is $A=0$
in the first chart and never passes through the conifold. It is
therefore unchanged as a surface. The other section is disjoint
from the centre and is unchanged as well. Over a smooth point
of $S_\infty$ the exceptional fibre is one $\mathbb P^1$;
over its node it is a chain of two $\mathbb P^1$'s, one being
the exceptional curve of the small resolution. This identifies
the normalization morphism, including its order of operations.

\item For example the corresponding total fourfold modification of
$xy=uzv$ can be written by first taking $x=ua$, $ay=zv$ and
then resolving the conifold in $(a,y,z,v)$, together with the
other chart $u=xb$, $y=bzv$. It has three smooth affine charts.
The assignments in \texttt{ab\_push\_local\_checks.py} verify
the equations and the map to $q=zv$ exactly. Exchanging the
chosen zero and infinity branches gives the other model.

\item The normal bundle of the section containing the centre changes
by tensoring with $\mathcal O_S(-S_\infty)$; the second normal
bundle is unchanged. Since the original product is $J$, the
new product is $J(-S_\infty)=\mathcal O_S$. The local
formula for $\mathcal T^1$ is the product of the two normal
bundles, with $T$ used on the infinity branch. This proves
d-semistability in this precise sense. The $\rho_l$ form one
orbit, so the central fibre is irreducible and its conductor is
exactly the two stated sections. Over $u\ne0$ both subdivisions
restrict to the fixed bottom triangulation, giving precisely
Engel's Mumford filling, simultaneously in $u$.

\item The birational comparison $\mathcal X^+\dashrightarrow
\mathcal X^-$ is the identity over the punctured $u$ disc and
off the exceptional loci over $(q,u)=(0,0)$. Locally it replaces
the two compatible side-face diagonals by the other two.
All its rays lie in the Gorenstein hyperplane $h_q+h_u=1$,
so both small modifications are crepant over
$\mathcal X^\flat$. Over a smooth point of the nodal curve
this is the ordinary Atiyah flop times a curve; at the node its
extension is exactly the displayed pair of prism subdivisions.
This specifies the birational map without asserting a uniqueness
theorem for arbitrary non-toroidal fillings.
\end{enumerate}
\end{proof}


\paragraph{counterexample; the literal blow-up shape}\label{push:src:A.2.3}\label{push:A23}
\StatementRecord{push:src:A.2.3}{proved}{Historical events: Claude Fable 5.2 (Anthropic), effort high, 16 November 2026: carried as background without fault found; used and rechecked in this run. Claude Fable 5.2 (Anthropic), effort high, 17 November 2026: carried as background without fault found. GPT-6 Astra (OpenAI), ChatGPT Pro mode, effort Pro, S6\_09\_PAPERS, 17 November 2026: targeted proof check. Also checked in this writing. Pre-continuation writing events: GPT-6 Astra (OpenAI), ChatGPT Pro mode, effort Pro, S6\_10\_PAPERS, 18 November 2026: targeted proof check. Continuation metadata events: GPT-6 Astra (OpenAI), ChatGPT Pro mode, effort Pro, S6\_10\_PAPERS continuation, 19 November 2026: used and rechecked in this run; status-vocabulary reconciliation; metadata reconciliation only; no new mathematical check. \emph{Pre-continuation recorded proof scope.} Re-derived the explicit P3 star counterexample and degree-one pullback injection, including b2(R)=10+1=11 from the actual projective bundle; wrote the complete receiving proof.}{\cref{push:src:A.2.1}, \cref{push:src:A.2.4}, \cref{push:src:A.2.5}; \cref{fam:src:E.1.1}}

The unrestricted request that every normalized component
be a blow-up of $R$ is false. The first maximal-cone
star already contributes a $\mathbb P^3$ component,
which admits no proper bimeromorphic map to $R$, since
such a map would inject a rank-eleven rational $H^2$
into a rank-one group. \ref{push:src:A.2.4} gives every normalization
by explicit star charts; \ref{push:src:A.2.5} supplies the terminating
blow-up and blow-down algorithm with centres, order,
exceptional bundles and gluing maps. Thus the literal
shape is refuted and its general algorithmic
replacement is proved.



\begin{proof}
\begin{enumerate}[label=\arabic*.,leftmargin=*]
\item The first maximal-cone star of \ref{push:src:A.2.1} is the
ordinary point star in a smooth four-dimensional toric chart.
Its new central component is the exceptional divisor
$\mathbb P^3$. The specified normalization $R$ is the
projective-line bundle over the rational elliptic surface $S$
of \ref{fam:src:E.1.1}. Its integral projective-bundle basis
gives $b_2(R)=b_2(S)+1=10+1=11$, while
$b_2(\mathbb P^3)=1$.

\item A proper bimeromorphic map $f:\mathbb P^3\to R$ would have
topological degree one. For $a\in H^2(R;\Q)$ and
$b\in H^4(R;\Q)$, the projection formula gives
$\int_{\mathbb P^3}f^*a\,f^*b=\int_Rab$.
Poincare duality makes $f^*$ injective in degree two, contradicting
these ranks. Thus this explicit component refutes the proposed
blow-up-only presentation. The star charts and the terminating
modification sequence in \ref{push:src:A.2.4}, \ref{push:src:A.2.5}
give the stated replacement for every allowed fan.
\end{enumerate}
\end{proof}

\Uses{the hypotheses, notation and definitions stated here; \cref{push:src:A.2.4}; \cref{push:src:A.2.5}}


\begin{theorem}[exhaustive finite fan data]\label{push:src:A.2.4}\label{push:A24}
\StatementRecord{push:src:A.2.4}{proved}{Historical events: Claude Fable 5.2 (Anthropic), effort high, 16 November 2026: carried as background without fault found; used and rechecked in this run. GPT-6 Astra (OpenAI), ChatGPT Pro mode, effort Pro, S6\_09\_PAPERS, 17 November 2026: targeted proof check. Also checked in this writing. Pre-continuation writing events: GPT-6 Astra (OpenAI), ChatGPT Pro mode, effort Pro, S6\_10\_PAPERS, 18 November 2026: targeted proof check. Continuation metadata events: GPT-6 Astra (OpenAI), ChatGPT Pro mode, effort Pro, S6\_10\_PAPERS continuation, 19 November 2026: used and rechecked in this run; status-vocabulary reconciliation; metadata reconciliation only; no new mathematical check. \emph{Pre-continuation recorded proof scope.} Re-derived the exhaustive finite-height ray search, volume(n+2)/6, exact common-face test, finite-orbit necessity, actual equivariant toric modification, normalized star charts, multiplicities and principal-divisor deformation line. Opened DMM2016v2 Section2p6 properness criterion; bounded count replays remain separately tracked.}{the hypotheses, notation and definitions stated here; \cref{push:src:A.2.1}; \cref{push:src:B.2.1}; \cref{push:src:A.1.2}; \cite[Section 2, p.6]{DMM16}}

The following finite-data procedure parametrizes all smooth
period-invariant subdivisions of the cusp fan after
$u=r^n$, retaining the prescribed bottom triangulation and
having finitely many cones modulo the periods. Here $n=1$
gives Part A, and $n=2,12$ give the two cases of B.2.
Each output defines a proper flat family over $\Delta_r$
with smooth total space, whose punctured restriction is the
given untwisted family pulled back by $u=r^n$.

For a bound $H\geq1$, retain the four bottom rays
$A,B,C,D$ and consider all primitive vectors
\[
 \begin{gathered}
 v=(a,b,k,\ell),\qquad
 \ell\geq1,\quad k\geq0,\quad k+\ell\leq H,\\
 0\leq a\leq k,\qquad0\leq b\leq k+n\ell.
 \end{gathered}
 \tag{A.24a}
\]
Take any collection of four-element subsets of this finite
set satisfying the following exact tests:
\begin{enumerate}
\item Each four-column determinant has absolute value one.
\item After dividing each ray by $k+\ell$, the resulting
tetrahedra intersect pairwise in common faces, and their
volumes sum to $(n+2)/6$.
\item Their bottom boundary is the two triangles $ABC,BCD$.
Their boundary subdivisions on $a=0$ and $a=k$ agree under
$G_1$, and those on $b=0$ and $b=k+n\ell$ agree under $G_2$.
\end{enumerate}
All these are finite rational-arithmetic tests. Extend a
passing collection by the two periods, and discard duplicates
as $H$ increases. This is exhaustive. For a reduced central
divisor, retain only $\ell=1$ in (A.24a).

The central prime components are indexed by period orbits of
the selected rays with $\ell>0$, and the multiplicity of the
component indexed by $v$ is exactly $\ell$. The normalization
charts of that component are the star fan of $v$ in
$N/\mathbb Zv$, with the inherited bundle units and the
stabilizer action. For $k>0$ the stabilizer is trivial and this
normalization is a smooth complete toric threefold. All
multiple loci, including possible triple and quadruple loci,
are specified by the positive-$\ell$ rays belonging to each
cone. In the reduced case
$\mathcal T^1\simeq\mathcal O_{\operatorname{Sing}D}$,
so every output is d-semistable.
\end{theorem}


\begin{proof}
\begin{enumerate}[label=\arabic*.,leftmargin=*]
\item Write the period shears as
\begin{align*}
 G_1(a,b,k,\ell)&=(a+k,b,k,\ell),\\
 G_2(a,b,k,\ell)&=(a,b+k+n\ell,k,\ell).
\end{align*}
One closed fundamental cone is
\[
 P_n=\{k,\ell\geq0,\ 0\leq a\leq k,
                    \ 0\leq b\leq k+n\ell\}.
\]
Its bottom is the original square; its top is the edge from
$(0,0,0,1)$ to $(0,n,0,1)$. The slice $k+\ell=1$, in
coordinates $(a,b,k)$, is the compact polytope
\[
 0\leq k\leq1,\qquad0\leq a\leq k,
 \qquad0\leq b\leq n-(n-1)k.
\]
Its volume is
$\int_0^1 k(n-(n-1)k)\,dk=(n+2)/6$.
A tetrahedron with ray columns $v_i$ has slice volume
\[
 \frac{|\det(v_1,v_2,v_3,v_4)|}
      {6\prod_{i=1}^4(k_i+\ell_i)}.
 \tag{A.24b}
\]
This follows by applying the determinant-one coordinate
change from $(a,b,k,\ell)$ to $(a,b,k,k+\ell)$ and then
normalizing the last coordinate to one.

\item A common-face test is also completely explicit: write each
tetrahedron in barycentric halfspaces, intersect the two
finite lists of halfspaces, and enumerate the possible
vertices by solving triples of bounding equations. The
intersection is a common face precisely when all its vertices
are common tetrahedron vertices. Since the candidate
tetrahedra lie in $P_n$, pairwise face intersections and the
volume equality imply that they cover the slice. Otherwise
their finite closed union would omit a nonempty relatively
open set, which has positive volume. The boundary tests make
the translates agree on every shared face. This proves that
every passing collection gives the claimed invariant fan.

\item Conversely, intersect any invariant smooth refinement with
$P_n$. A smooth maximal cone has four primitive generators
with determinant $\pm1$, and its slice is one of the stated
tetrahedra. Finitely many cone orbits give finitely many rays
in this prism, hence some maximum height $H$. No additional
bottom rays are permitted by the boundary condition; every
other ray has the form (A.24a). The resulting collection
passes all three tests. Thus the procedure is a finite,
exhaustive search at every height, rather than a restriction
to regular triangulations or to new rays in a chosen orbit.
There is no projectivity condition in the tests.

\item The finite-orbit condition is necessary for a proper quotient
of such a refinement, not an additional restriction on the
requested proper models. Each full-dimensional cone has a
closed torus orbit consisting of one point with $q=r=0$.
These points form a closed discrete set: an affine chart
of a full-dimensional cone contains exactly its own closed
torus orbit. After the properly discontinuous period
quotient the points are indexed by cone orbits and remain
closed and discrete. Infinitely many such orbits would
therefore give an infinite closed discrete subset of a
compact central fibre, which is impossible.

\item The fundamental orbit representatives of mixed rays are
\[
 (a,b,k,\ell),\quad k,\ell>0,
 \quad0\leq a<k,\quad0\leq b<k+n\ell,
 \quad\gcd(a,b,k,\ell)=1.
 \tag{A.24c}
\]
For pure top rays they are $(0,b,0,\ell)$ with
$0\leq b<n\ell$ and $\gcd(b,\ell)=1$. A smooth top
subdivision necessarily includes each $(0,j,0,1)$,
$0\leq j\leq n$. Indeed an omitted such vector would lie
inside a regular two-cone; writing it in that cone's integral
basis would express its height one as the sum of two
positive integral heights, which is impossible.
In the reduced search there are exactly $n$ pure top orbits
and $\sum_{k=1}^{H-1}k(k+n)$ available mixed orbits.

\item The program \path{code/ab_push_all_refinements.py} enumerates
all candidate rays for $n=1,2,12$ through $H=3$, and all
four-column determinant tests at the bounds reported here:
\[
\begin{array}{c|c|c|r|r}
n&H&\text{reduced}&\text{rays in }P_n&
       \text{unimodular four-subsets}\\\hline
1&1&\text{yes}&6&12\\
1&2&\text{yes}&12&267\\
1&2&\text{no}&13&357\\
2&2&\text{yes}&15&544\\
2&2&\text{no}&17&784\\
12&1&\text{yes}&17&100
\end{array}
\]
These are candidate counts, not a claim that every candidate
four-subset or collection passes the other tests. At height
one, the complete collections are the two and $n+1$ fans
already proved in \ref{push:src:A.2.1} and \ref{push:src:B.2.1}. The saved program also
certifies the first eight members of the explicit infinite
family in \ref{push:src:A.2.1}.

\item For existence of the corresponding compact family, start
with the proper nodal plumbing of \ref{push:src:A.1.2} and base-change it
by $u=r^n$. On the cusp cover apply the proper toric
modification defined by the refinement. Its support is
unchanged, which is the toric properness criterion
\cite[Section 2, p.6]{DMM16}, and each
original affine chart is replaced by a finite fan. On
$q\ne0$ use the subdivision of the top two-dimensional fan
to modify the chart $ab=r^n\varphi$ fibrewise over $S$.
Changes of trivialization multiply $a$ and $b$ by the given
nonvanishing line-bundle units; a torus action extends to
every toric fan chart, so these modifications glue. The same
top fan occurs on the overlap with the cusp construction.
The core outside the plumbing charts is unchanged. This defines
an actual proper modification of the base-changed compact
nodal family, not a gluing inferred from a vanishing theorem.

\item The period maps lift because the fan is invariant. Their
lifted action is free: a fixed point would project to a fixed
point in the nodal covering. It is properly discontinuous:
compact sets project to compact sets, and any translating
element producing an intersection upstairs produces one
downstairs, where \ref{push:src:A.1.2} gives finiteness. Thus the quotient
exists as a Hausdorff complex space, and the equivariant
proper modification descends to a proper modification of
the compact nodal quotient. In a unimodular chart its
coordinates $z_i$ satisfy
\[
 r=\text{unit}\cdot\prod_{i=1}^4 z_i^{\ell_i},
 \qquad q=\text{unit}\cdot\prod_{i=1}^4 z_i^{k_i}.
 \tag{A.24d}
\]
The total space is smooth. The nonconstant coordinate $r$
is a nonzerodivisor, so the morphism to the disc is flat.
Its central multiplicities are the indicated $\ell_i$;
it is reduced exactly when all selected positive heights
are one. The bottom fan gives the specified smooth
Mumford model over $r\ne0$, including its original period
identification. Thus the punctured-family assertion holds
simultaneously in the parameter.

\item For completeness the normalization description is an
explicit chart prescription. For a maximal cone with basis
$(v,v_2,v_3,v_4)$, let $(m_1,m_2,m_3,m_4)$ be the dual
integral basis. The chart of the normalized component is
$z_1=0$ with coordinates $z_i=\chi^{m_i}$, $i=2,3,4$.
On overlaps the transition matrix is the product of the two
integer basis matrices, and the original bundle units are
raised to those same integer exponents. This is precisely
the star fan in $N/\mathbb Zv$. When $k>0$, (A.24c) shows
that the period stabilizer of the ray is trivial. The ray
lies in the interior of the full support, so its star is
complete. Its normalized component is therefore the
asserted complete smooth toric threefold, lying over the
cusp. When $k=0$, the stabilizer is generated by $G_1$.
Away from $S_\infty$ the resulting component is the ruled
bundle $\mathbb P_S(\mathcal O\oplus L^\ell)$: a character
on its open ruled fibre is $W^\ell/r^b$, which transforms
by the $\ell$th power of the transition of $L$. The stated
star charts describe its replacement over $S_\infty$.
This also specifies all gluing maps, rather than only their
combinatorial incidence: they are the original period maps
on those same integer characters.

\item Intersections are read from common cones, with separate
local branches retained before the period quotient. In a
chart (A.24d), the number of branches of the reduced
central divisor is the number of positive $\ell_i$.
When all these heights are one, the hypersurface
calculation gives
$\mathcal T^1_D=\mathcal O_D(D)|_{\operatorname{Sing}D}$.
The divisor is globally principal, cut out by $r$, so this
line on the singular locus is trivial. Finally, any two
output fans have an explicit common refinement by all
intersections of their cones. Its normal toric modification
gives a proper common model, and the resulting birational
map is the identity on the fixed lattice and on the
punctured family. No meromorphic extension of an arbitrary
punctured isomorphism is needed here.
\end{enumerate}
\end{proof}

\begin{proposition}[explicit normalization algorithm and its necessary scope]\label{push:src:A.2.5}\label{push:A25}
\StatementRecord{push:src:A.2.5}{proved}{Historical events: Claude Fable 5.2 (Anthropic), effort high, 16 November 2026: carried as background without fault found; used and rechecked in this run. GPT-6 Astra (OpenAI), ChatGPT Pro mode, effort Pro, S6\_09\_PAPERS, 17 November 2026: targeted proof check. Also checked in this writing. Pre-continuation writing events: GPT-6 Astra (OpenAI), ChatGPT Pro mode, effort Pro, S6\_10\_PAPERS, 18 November 2026: targeted proof check. Continuation metadata events: GPT-6 Astra (OpenAI), ChatGPT Pro mode, effort Pro, S6\_10\_PAPERS continuation, 19 November 2026: used and rechecked in this run; status-vocabulary reconciliation; metadata reconciliation only; no new mathematical check. \emph{Pre-continuation recorded proof scope.} Re-derived the explicit P3 counterexample, blow-up/down search, all chart changes and central multiplicities, and the cotangent-complex Jacobian quotient for nonreduced divisors. Opened AKMW2002 Thms1.6.1/2.7.1 pp543,551-552 and analytic Sections5.5-5.6 pp564-566; opened Matsuki\textquoteright s corrected Weak Factorization Theorem for Toroidal Birational Morphisms, Section~2, p.~2, for smooth-stratum blowups and blowdowns over the target.}{the hypotheses, notation and definitions stated here; \cref{push:src:A.2.4}; \cref{push:src:A.4.6}; \cref{push:src:A.4.8}; \cref{push:src:A.2.2}}

The requirement that every normalized component of every
smooth refinement be a blow-up of
$R=\mathbb P_S(\mathcal O\oplus L)$ is false, already
for a reduced refinement with $n=1$.
The following is a uniform replacement for any finite
fan certificate $\Sigma$ in \ref{push:src:A.2.4}.

The star charts and integer transition matrices of \ref{push:src:A.2.4}
give every normalized component, its gluing maps, and
all multiple loci directly. There is also a terminating
finite-data algorithm that expresses its comparison
with a reference model as an explicit sequence of
blow-ups and blow-downs along specified smooth strata.
At each blow-up the algorithm gives the normalized
strict transforms, the exceptional projective bundle,
its multiplicity, and every new intersection.
This algorithm allows blow-downs; no toric
strong-factorization assertion is used.

For a component of multiplicity $\ell$, the local
nonreduced central equation is the monomial in A.24d.
More generally, in a unimodular chart with central
equation $F=\prod_i z_i^{\ell_i}$, the cotangent-complex
first-order deformation sheaf is
\[
 \mathcal T^1_D
  =\mathcal O_D\Big/
    \left(\ell_i z_i^{\ell_i-1}
          \prod_{j\ne i}z_j^{\ell_j}:\ell_i>0\right).
 \tag{A.25a}
\]
Here $\mathcal T^1_D$ means
$\mathcal Ext^1(\mathbb L_D,\mathcal O_D)$; in the
reduced case it also equals
$\mathcal Ext^1(\Omega_D,\mathcal O_D)$.
The reduced case is the trivial line on the singular
scheme, as in \ref{push:src:A.2.4}. For a nonreduced centre the displayed
Jacobian quotient is the appropriate deformation sheaf;
the reduced normal-crossings notion of d-semistability
is not asserted for it.
\end{proposition}


\begin{proof}
\begin{enumerate}[label=\arabic*.,leftmargin=*]
\item The first maximal-cone star in \ref{push:src:A.4.6} produces a normalized
component $\mathbb P^3$. There cannot be a proper
bimeromorphic holomorphic map $\mathbb P^3\to R$:
between smooth compact oriented manifolds such a map
has degree one, and proper pushforward gives a left
inverse to pullback on rational cohomology. But
$b_2(R)=11$ and $b_2(\mathbb P^3)=1$.
Thus this component is not a blow-up of $R$, regardless
of the proposed centres or their order. The actual
normalization in this example is
$\operatorname{Bl}_pR^+\sqcup\mathbb P^3$,
with its specified plane intersection. This is an
explicit counterexample to the universal shape, not
a failure to find a blow-up sequence.

\item Here is the promised algorithm. Retain all marking
data from \ref{push:src:A.2.4}: the lattice, the two period matrices,
the original line-bundle units, and the common
unsubdivided fan. Strictify the finite quotient cone
complex by the decreasing-dimensional barycentric
stars used in \ref{push:src:A.4.8}. Record these operations and do
the same for the reference model. Form their common
intersection complex. A regular common subdivision
can be found either by toric desingularization or by
enumerating rational subdivisions with increasing
primitive-ray bounds and testing their determinants
and common faces. A regular subdivision exists, so
the latter enumeration terminates.

\item On a finite strict regular complex there are finitely
many possible ordinary star blow-ups: choose a cone
of dimension at least two and insert the sum of its
primitive basis rays. There are also finitely many
possible inverse moves: choose a ray and test whether
its star is exactly the star created by such a
subdivision of a regular cone. Both tests are exact
integer calculations. Require that every intermediate
complex remain a subdivision of the marked original
complex. Record face identifications; equality of two
marked complexes is a finite rational comparison,
using the period orbit representatives in A.24c and
their pure top and pure bottom analogues.

\item Enumerate this locally finite graph of moves in
increasing path length, discarding repeated marked
complexes. Applied separately to the two proper
subdivision morphisms from the common regular
complex, the finite-complex weak-factorization theorem
used in \ref{push:src:A.4.8} guarantees a path, all of whose
intermediates are over the corresponding target.
Thus the search terminates. Concatenate the two paths
and the recorded strictification paths. This gives a
fully specified finite sequence once the input fan is
given. Temporary pure bottom subdivisions may occur;
they are allowed in a presentation of the birational
comparison, and are removed before the prescribed
endpoint is reached.

\item We describe explicitly what the recorded move means
geometrically. If a cone has primitive basis vectors
$v_1,\ldots,v_c$, its centre is the closed stratum
\[
 C=\{z_1=\cdots=z_c=0\}
\]
in the marked toroidal chart. Strictness makes its
global closure smooth; the period translates defining
one centre are disjoint. Blow-up chart $i$ has
coordinates
\[
 z_i=t_i,\qquad z_j=t_i w_j\quad
 (1\leq j\leq c,\ j\ne i),
 \qquad z_j\quad(j>c).
\]
The changes between charts are these integral
monomials, with the original units raised to the
same powers. The exceptional divisor is
$\mathbb P_C(N_{C/Y})$, where $Y$ is the smooth
intermediate total space. If $C$ is central, its
central multiplicity is
\[
 \ell_E=\sum_{j=1}^c\ell_j.
\]
If all $\ell_j$ vanish, it is a horizontal exceptional
divisor and is not a new central component.

\item For an old normalized component, take the inverse
image of this centre on its marked branch and the
ordinary strict transform in these charts. It is
the blow-up along that smooth inverse-image centre;
when that centre is Cartier the blow-up is the
identity. Its intersections with the new exceptional
component are the appropriate coordinate-hyperplane
projective subbundles. Intersections of several old
branches are given by imposing the corresponding
coordinate equations simultaneously. Thus the same
chart prescription records double, triple and
quadruple loci, their normal lines, and their actual
gluing maps. An inverse star is the inverse
contraction of this recorded configuration. This
describes centres and order, not merely a claim that
some birational map exists.

\item The algorithm starts with the explicit reference
normalizations of \ref{push:src:A.2.2} and the $n$-cycle reference
in \ref{push:src:A.4.8}. It is therefore a uniform geometric
presentation of every output by specified blow-ups,
blow-downs and projective bundles. Independently,
the direct star charts avoid running this search
when only a normalization chart or an intersection
is required. The mixed-ray components are complete
smooth toric threefolds; a pure top component
$(0,b,0,\ell)$ is bimeromorphic to
$\mathbb P_S(\mathcal O\oplus L^\ell)$.
There is no assertion that all these distinct
components are modifications of a single copy of $R$.

\item Finally $D$ is globally the principal Cartier
divisor of the parameter $r$ in a smooth total
space. Its two-term cotangent complex is
$[\mathcal O_D\xrightarrow{dF}\Omega_Y|_D]$,
so its first-order deformation module is the
cokernel of the gradient. Equivalently, a lifted
hypersurface equation $F+\epsilon G$ is changed by
an infinitesimal coordinate change precisely by a
Jacobian term. This argument applies to nonreduced
monomials as well; an ordinary conormal sequence
need not be left-exact there. Absorbing the local unit, or reducing its
derivative modulo $F$, gives (A.25a). A global
parameter makes the normal line trivial, so these
Jacobian quotients glue by the canonical
hypersurface presentation. In the reduced case
the Jacobian ideal defines the singular scheme,
giving the asserted trivial deformation line there.

\item The exact certificate
\path{code/ab_push_normalization_algorithm.py}
executes the star search on the first two
maximal-cone examples, verifies both inverse
contractions, and checks all local multiplicity
substitutions with exponents zero, one and two.
The finite-complex theorem supplies termination for
arbitrary input; these tests check the explicit
local implementation.
\end{enumerate}
\end{proof}


\subsection{The virtual Euler obstruction and the conjecture}\label{push:section:A.3}

\begin{lemma}\label{push:src:A.3.1}\label{push:A31}
\StatementRecord{push:src:A.3.1}{proved}{Historical events: Claude Fable 5.2 (Anthropic), effort high, 16 November 2026: carried as background without fault found; targeted proof check; used and rechecked in this run. Claude Fable 5.2 (Anthropic), effort high, 17 November 2026: carried as background without fault found. GPT-6 Astra (OpenAI), ChatGPT Pro mode, effort Pro, S6\_09\_PAPERS, 17 November 2026: targeted proof check. Also checked in this writing. Pre-continuation writing events: GPT-6 Astra (OpenAI), ChatGPT Pro mode, effort Pro, S6\_10\_PAPERS, 18 November 2026: targeted proof check. Continuation metadata events: GPT-6 Astra (OpenAI), ChatGPT Pro mode, effort Pro, S6\_10\_PAPERS continuation, 19 November 2026: used and rechecked in this run; status-vocabulary reconciliation; metadata reconciliation only; no new mathematical check. \emph{Pre-continuation recorded proof scope.} Re-derived the relative perfect tangent K-class and radial fundamental-cycle comparison; opened Levy08 Lemmas3.1-3.2 Section3.1 pp588-590 and Lojasiewicz64 Section3 Thms1-2 pp463-464. Triangulation over a singular analytic total space, beyond the printed ambient-manifold theorem, remains an explicit source/application gate. Closed the singular-ambient triangulation application gate: replaced it by a proper analytic resolution, triangulated the radial inverse image inside its smooth real analytic total space, computed boundary multiplicities by u z\textasciicircum{}m, and pushed the chain down. BM95 Section0.2(2)p3,Thm1.6p5,Remark1.16p11 and Loj64 Thms1-2pp463-464 opened at exact scopes; full replacement written in receiving revision.}{the hypotheses, notation and definitions stated here; \cite[Lemmas 3.1--3.2 and Section 3.1]{Levy08}; \cite[Section~3, Theorems~1--2, pp.~463--464]{Loj64}; \cite[Theorem~1.6, p.~5; Remark~1.16, p.~11]{BM95AB}; \cite{FG10}}

Let $Z$ be a compact reduced local-complete-intersection complex
space of pure dimension $n$. In any proper flat holomorphic
family over a disc with central fibre $Z$ and smooth nearby
fibres, their Euler number is
\[
 e_{\mathrm{vir}}(Z)=\langle c_n(T_Z^{\mathrm{vir}}),[Z]\rangle.
\]
This statement needs no projectivity or K\"ahler hypothesis.
\end{lemma}


\begin{proof}
\begin{enumerate}[label=\arabic*.,leftmargin=*]
\item Flat deformation preserves the local-complete-intersection
property relative to the disc near $Z$. Locally this follows by
lifting an embedding and its regular sequence of equations;
Nakayama's lemma and flatness show that their lifts generate
the ideal and remain a regular sequence. Properness allows one
to shrink the disc so that these neighbourhoods cover the total
space. The relative cotangent complex is consequently perfect
of amplitude $[-1,0]$. Its dual defines a topological complex
$K$-class $V$ on that total space; local finite vector-bundle
complexes define the class by descent, or by stabilization and
partitions of unity; the analytic-to-topological comparison is
given in \cite[Lemmas 3.1--3.2 and Section 3.1]{Levy08}.
Its restriction to a smooth fibre is its
tangent bundle, and its restriction to $Z$ is its intrinsic
virtual tangent class.

\item We compare the fundamental cycles on a smooth proper modification,
so that the ambient-manifold hypothesis of the triangulation theorem
is satisfied. Analytic resolution, in the complex-analytic scope of
\cite[Section~0.2(2), p.~3; Theorem~1.6 and its noncompact extension,
p.~5; Remark~1.16, p.~11]{BM95AB}, gives
\[
 \rho:\widetilde{\mathcal X}\longrightarrow\mathcal X
\]
with smooth total space, isomorphic over the smooth locus. Only a
neighbourhood of the compact central fibre is needed; the locally
finite construction is finite there after shrinking the base.
The total space is reduced: a nilpotent local section lies in every
power of the base parameter by reducedness of the central fibre and
flatness, and hence vanishes by Krull intersection. Thus the stated
resolution applies to the actual dominating space.

\item Write the divisor of the pulled-back parameter as
$\sum_j m_jD_j$. The inverse image of a short radial segment
$[0,\epsilon]$ is a compact semianalytic subset of the smooth real
analytic manifold $\widetilde{\mathcal X}$. We triangulate it,
compatibly with its two endpoint fibres and the singular loci of the
$D_j$, by \cite[Section~3, Theorems~1--2, pp.~463--464]{Loj64}.
Over the interior of the segment its top-dimensional part is an
oriented real $(2n+1)$-manifold. Near a general point of $D_j$ the
parameter is $u z^{m_j}$, with $u$ a holomorphic unit; choosing a
local root of $u$ exhibits $m_j$ radial half-sheets, each with the
complex boundary orientation. The sum of their oriented top
simplices therefore has boundary
\[
 [\mathcal X_\epsilon]-\sum_jm_j[D_j].
\]
The singular and crossing sets have real dimension at most $2n-2$
and contribute no boundary simplex of dimension $2n$.

\item At every smooth point of the reduced central fibre the original
map is a submersion, by flatness and smoothness of that fibre germ.
The modification is consequently an isomorphism there and the
strict-transform multiplicity is one. Every other $D_j$ maps into
$\operatorname{Sing}Z$, of complex dimension at most $n-1$.
Its fundamental class has zero pushforward in dimension $2n$:
its map factors through that lower-dimensional analytic set.
Thus pushing the displayed chain forward gives exactly
$[\mathcal X_\epsilon]-[Z]$. This proves the required homology
identity in the original total space. Pairing the globally defined
$c_n(V)$ with these homologous cycles gives the equality of
Chern numbers. Gauss--Bonnet on the smooth compact complex
fibre gives its Euler number. This proves the analytic assertion
directly. The algebraic counterpart is Theorem 5.8 and
Remark 5.9 of \cite{FG10}; those scheme statements alone would
not justify omitting the analytic comparison above.
\end{enumerate}
\end{proof}

\begin{theorem}[Conjecture 8.2 refuted]\label{push:src:A.3.2}\label{push:A32}
\StatementRecord{push:src:A.3.2}{proved}{Historical events: Claude Fable 5.2 (Anthropic), effort high, 16 November 2026: carried as background without fault found; targeted proof check; used and rechecked in this run. Claude Fable 5.2 (Anthropic), effort high, 17 November 2026: carried as background without fault found. GPT-6 Astra (OpenAI), ChatGPT Pro mode, effort Pro, S6\_09\_PAPERS, 17 November 2026: targeted proof check. Also checked in this writing. Pre-continuation writing events: GPT-6 Astra (OpenAI), ChatGPT Pro mode, effort Pro, S6\_10\_PAPERS, 18 November 2026: targeted proof check. Continuation metadata events: GPT-6 Astra (OpenAI), ChatGPT Pro mode, effort Pro, S6\_10\_PAPERS continuation, 19 November 2026: used and rechecked in this run; status-vocabulary reconciliation; metadata reconciliation only; no new mathematical check. \emph{Pre-continuation recorded proof scope.} Re-derived normalization pullback of the virtual tangent class, both normal-line contributions and virtual Euler24-24=0, the actual untwisted Euler2, and the distinct smooth-zero-section smoothing. Exact symbolic virtual-c3 replay passed; Engel attached-source edition Conjecture8.2p22 and Lemma2.1pp3-4 opened.}{the hypotheses, notation and definitions stated here; \cite{Engel2026}; \cref{push:src:A.1.2}; \cref{push:src:A.2.2}; \cref{push:src:A.4.3}; \cref{push:A31}}

For the specified pinching $Q$ one has $e_{\mathrm{vir}}(Q)=0$.
Every smooth proper-flat smoothing of $Q$ therefore has Euler
number zero. In particular there is no proper flat holomorphic
filling of Engel's untwisted family with central fibre exactly
$Q$, even if no map to $B$ is required. Thus Conjecture 8.2
of the attached \cite{Engel2026} is false as worded. The proper nodal
filling of \ref{push:src:A.1.2} and the two natural semistable fillings of
\ref{push:src:A.2.2} are precise replacements, with different central fibres.
\end{theorem}


\begin{proof}
\begin{enumerate}[label=\arabic*.,leftmargin=*]
\item Let $D_0,D_\infty\subset R$ be the conductor sections and
$\nu:R\to Q$ the normalization. On either branch of the
hypersurface $xy=0$, the derivative comparison with its
normalization gives in $K^0(R)$, where pullback of the
perfect tangent class means derived pullback,
\[
 \nu^*T_Q^{\mathrm{vir}}
 =T_R-i_{0*}(J|_{D_0})-i_{\infty*}(J|_{D_\infty}).
\]
Indeed in a smooth local ambient space the normal line $A$ of
the branch and the normal line of the whole hypersurface
differ by $D$: the latter is $A(D)$. Their difference in
$K$-theory is $-[i_*(A(D)|_D)]$, and $A(D)|_D$ is exactly
the product of the two branch normals, namely $\mathcal T^1_Q$.
These local comparisons are intrinsic and agree on overlaps.
On the infinity branch $J$ means its $T^{-1}$-pullback
from the double locus, canonically the same line bundle
because it comes from $B$. Thus $J$ extends from both sections by the pullback from $B$,
so $[i_*J]=[J]-[J(-D)]$.

\item Put $j=c_1(J)=\pi_R^*F$, $d_i=[D_i]$. Since the sections are
disjoint, $d_0d_\infty=0$, and since $F$ is a fibre, $j^2=0$.
Consequently
\[
 c(\nu^*T_Q^{\mathrm{vir}})
 =c(T_R)\frac{(1+j-d_0)(1+j-d_\infty)}{(1+j)^2}
 =c(T_R)\{1-d_0-d_\infty+j(d_0+d_\infty)\}.
\]
The degree-three part has integral
\[
 e(R)-\sum_i\int_{D_i}c_2(T_R|_{D_i})
       +\sum_i\int_{D_i}j\,c_1(T_R|_{D_i}).
\]
We have $e(R)=2e(S)=24$. On a section with normal class
$n_i$, $c_2(T_R|_{D_i})=c_2(S)+c_1(S)n_i$ and
$c_1(T_R|_{D_i})=c_1(S)+n_i$. The two normal classes on
the unglued normalization are $m,-m$, where $m=c_1(L)$.
Their sum is zero; $c_1(S)=F$ and $F^2=0$. The displayed
number is therefore $24-24+0=0$. Because
$\nu_*[R]=[Q]$, this is the intrinsic virtual number on $Q$.
The projective-bundle calculation is independently checked
in \texttt{ab\_push\_local\_checks.py}, using
$H(H+m)=0$, $d_0=H+m$, $d_\infty=H$.

\item The untwisted quotient away from its cusp has Euler number zero:
scalar multiplication supplies a free circle action, including
the two additive places before logarithmic transformation. Indeed
a circle scalar identified with an iterate of $u\phi$ must be
the identity, since every positive iterate strictly contracts
the Hermitian norm and every negative iterate expands it.
Compact fibres with a free circle action have Euler number zero.
The chosen Mumford fibre has two zero-dimensional torus orbits
and all other torus strata have Euler number zero, so
$e(X(u))=2$. The same number follows independently from the
semistable nearby-cycle count in \ref{push:src:A.4.3} below. Lemma~\ref{push:A31}
therefore excludes such a smoothing of $Q$. No requirement
on its leading smoothing section, and no extension of $f$,
entered this argument.

\item There is also a direct check that zero is the correct smooth
number. Replace $\phi$ in the global construction of \ref{push:src:A.1.2}
by a section of $J$ vanishing on a smooth elliptic fibre of
$S$. The compact collar construction still applies. At its
zero divisor the fixed nonzero-parameter equation is $xy=uz$
with $dz\ne0$, so every nearby threefold is smooth. The
double-locus nearby Euler function is one on that elliptic
curve and zero off it; its integral is zero. This is an actual
proper-flat smoothing of $Q$, but is not Engel's family.
\end{enumerate}
\end{proof}

\begin{proposition}[counterexample and corrected]\label{push:src:A.3.3}\label{push:A33}
\StatementRecord{push:src:A.3.3}{proved}{Historical events: Claude Fable 5.2 (Anthropic), effort high, 16 November 2026: carried as background without fault found. GPT-6 Astra (OpenAI), ChatGPT Pro mode, effort Pro, S6\_09\_PAPERS, 17 November 2026: targeted proof check. Also checked in this writing. Pre-continuation writing events: GPT-6 Astra (OpenAI), ChatGPT Pro mode, effort Pro, S6\_10\_PAPERS, 18 November 2026: targeted proof check. Continuation metadata events: GPT-6 Astra (OpenAI), ChatGPT Pro mode, effort Pro, S6\_10\_PAPERS continuation, 19 November 2026: used and rechecked in this run; status-vocabulary reconciliation; metadata reconciliation only; no new mathematical check. \emph{Pre-continuation recorded proof scope.} Re-derived 2P coordinates, its disjointness from P/O on every special fibre, global graph centre, blow-up charts over(t,u), and repeatable O plus O(-1) normal sequence. Exact doubling substitution replay passed.}{the hypotheses, notation and definitions stated here; \cref{push:section:A.3}; \cref{push:src:A.2.1}}

The assertion that the two natural
fillings exhaust smooth-total-space reduced-normal-crossings
fillings is false, even when joint flatness over
$B\times\Delta$ is imposed. There are infinitely many such
fillings of the same punctured untwisted family. The
two-model classification in \ref{push:src:A.2.1} is a classification of the
specified minimal toroidal refinements, not of all fillings.
\end{proposition}


\begin{proof}
\begin{enumerate}[label=\arabic*.,leftmargin=*]
\item We give centres that preserve joint flatness, rather than
blowing up a point over a point of the two-dimensional base.
The section $2P$ is disjoint from $P$ and $O$. In the
Weierstrass equation its coordinates are
\[
 x(2P)=\frac{9t^2}{8},\qquad
 y(2P)=\frac{\sqrt2\,t^2(5t-32)}{32}.
\]
At $0$ it meets the nonidentity component different from
that met by $P$; at $1$ the displayed coordinates are finite
and different from those of $P,O$; at infinity the minimal
coordinates are finite and different as well. This also
checks directly the zero intersection numbers
$P\cdot2P=O\cdot2P=0$. The rational section of
$\mathcal O(P-O)$ therefore restricts to a nowhere-vanishing
holomorphic section on $2P\simeq B$. Its graph is a compact
curve $C\subset L^\times\subset R^\pm$ disjoint from both
conductor sections and from the modification centres.
It gives a curve in the smooth locus of $Q^\pm$ mapping
isomorphically to $B$.

\item Blow up $C$ in the smooth total fourfold $\mathcal X^\pm$.
Locally the coordinates are $(t,u,a,b)$ with
$C=(u,a,b)=0$. The blow-up charts replace $u$ by $u$ itself
or by $aU$ or $bU$. Thus the new central divisor is reduced
normal crossings, and the total space is smooth. All fibres
over $(t,u)$ still have pure dimension two; the regular base
criterion proves joint flatness. The morphism is proper and
is an isomorphism for $u\ne0$. It adds a
$\mathbb P^2$-bundle over $C$, hence a new central irreducible
component. To justify repetition, $N_{C/Q^\pm}$ is an extension
of $\mathcal O_B(-1)$ by $\mathcal O_B$ and splits, and
$H^1(N_{C/Q^\pm})=0$. The normal sequence in the total family
therefore splits off the trivial parameter direction. Its
section in the exceptional $\mathbb P^2$-bundle is disjoint
from the intersection with the strict central transform and
has the same normal sequence. We may repeat indefinitely.
Their differing component counts preclude
isomorphism as fillings. The local substitutions and the
section-doubling calculation are exact computations, recorded
with the other local checks.
\end{enumerate}
\end{proof}

\subsection{Integral topology and nearby cycles}\label{push:section:A.4}

\begin{proposition}\label{push:src:A.4.1}\label{push:A41}
\StatementRecord{push:src:A.4.1}{proved}{Historical events: Claude Fable 5.2 (Anthropic), effort high, 16 November 2026: carried as background without fault found. GPT-6 Astra (OpenAI), ChatGPT Pro mode, effort Pro, S6\_09\_PAPERS, 17 November 2026: targeted proof check. Also checked in this writing. Pre-continuation writing events: GPT-6 Astra (OpenAI), ChatGPT Pro mode, effort Pro, S6\_10\_PAPERS, 18 November 2026: targeted proof check. Continuation metadata events: GPT-6 Astra (OpenAI), ChatGPT Pro mode, effort Pro, S6\_10\_PAPERS continuation, 19 November 2026: used and rechecked in this run; status-vocabulary reconciliation; metadata reconciliation only; no new mathematical check. \emph{Pre-continuation recorded proof scope.} Re-derived primitive normalization maps, determinant11 lattice saturation, exceptional-node generization matrix[[1,0],[1,1]], Leray and mapping-cylinder groups of Rplus/Qplus, and both Euler counts. Current e\_degeneration and ab\_push\_local\_checks exact replays passed.}{the hypotheses, notation and definitions stated here; \cref{fam:src:E.1.1}, \cref{fam:src:E.1.2}, \cref{fam:src:E.1.3}, \cref{fam:src:E.1.4}; \cref{push:src:A.2.2}}

With integral coefficients, the spaces constructed above have
\[
\begin{array}{c|rrrrrrr|r}
&b_0&b_1&b_2&b_3&b_4&b_5&b_6&\chi\\\hline
R&1&0&11&0&11&0&1&24\\
R^\pm&1&0&12&0&12&0&1&26\\
Q&1&1&8&7&10&0&1&12\\
Q^\pm&1&1&9&7&11&0&1&14
\end{array}
\]
All these homology groups are torsion free. The fundamental
groups of $R,R^\pm$ are trivial and those of $Q,Q^\pm$ are
$\mathbb Z$. The dual complex of either $Q^\pm$ is a circle.
\end{proposition}


\begin{proof}
\begin{enumerate}[label=\arabic*.,leftmargin=*]
\item We check the lattice step of \cref{fam:src:E.1.1}, \cref{fam:src:E.1.2}, \cref{fam:src:E.1.3}, \cref{fam:src:E.1.4} on which the calculation depends.
Let $\mathcal L=H_2(S;\mathbb Z)$, with its unimodular intersection
pairing, and put
\[
 F=[S_t],\qquad m=P-O,\qquad v=[6P]-[O]-3F.
\]
The intersections are
\[
 F^2=Fm=Fv=0,\quad m^2=-2,\quad v^2=-6,\quad mv=-1,
 \quad OF=Om=1,\quad Ov=0.
\]
Translation is the integral isometry
\[
 T_*a=a+(aF)v-(av-3aF)F.
\]
This is checked on $F,O,v$ and on the seven root classes of
the two reducible Kodaira fibres, which translation fixes.
It also gives $T^*m=m-F$. The span of $F,v,m$ is primitive:
split off the unimodular plane generated by $F,O$; the Gram
determinant of $m-F,v$ in its orthogonal complement is eleven,
so an index larger than one would have its square divide eleven.
Unimodularity therefore makes
$a\mapsto(aF,av,am)$ surjective onto $\mathbb Z^3$.

\item The normalization mapping cylinder gives
\[
 H_i(S)\xrightarrow{d_i}H_i(R)\longrightarrow H_i(Q)
 \longrightarrow H_{i-1}(S),\qquad
 d_i=s_{0*}-s_{\infty*}T_*.
\]
The maps are induced by actual embeddings, and the mapping
cylinder is homotopy equivalent to the pinching because the
two section embeddings are disjoint cofibrations. With $f$ the
projective-line fibre, the degree-two map is
\[
 d_2(a)=-(aF)v+(av-3aF)F+(am-aF)f.
\]
Its three coefficient columns have determinant $-1$, so the
image is primitive of rank three, the kernel has rank seven,
and the cokernel has rank eight. In degree four the image is
the primitive class Poincar\'e dual to $m$, since $mO=1$.
In degree zero the difference map is zero. The projective-bundle
formula and this exact sequence give the asserted groups of
$R,Q$. Van Kampen on the same mapping cylinder gives the
single stable loop and no other generator, hence
$\pi_1(Q)=\mathbb Z$.

\item Now let $h:R^+\to R$ be the modification constructed in \ref{push:src:A.2.2}.
Its fibres are points off the nodal curve $C=S_\infty$, one
$\mathbb P^1$ over $C_{\mathrm{reg}}$, and a chain of two
$\mathbb P^1$'s over the node. The chart calculation identifies
the sheaf $R^2h_*\mathbb Z$ with
$\eta_*\mathbb Z_{\mathbb P^1}$, where
$\eta:\mathbb P^1\to C$ is the normalization. Here is the
integral check at the node. If $(a,b)$ are the cohomological
degrees on the two exceptional curves, generization to the
two branches is $a$ and $a+b$: on one branch the generic
exceptional fibre specializes to the first curve, on the
other it specializes to their sum. The matrix
$\left(\begin{smallmatrix}1&0\\1&1\end{smallmatrix}\right)$
is unimodular. These generizations follow by substituting
$X=za$, $v=aB$ in the blow-up charts, or by taking the two
limits of their exceptional fibres. Proper local neighbourhoods
retract onto these fibres by the positive toric contraction,
so the fibre calculation computes the higher direct images.
There is no other positive-degree direct image.

\item The Leray spectral sequence has only the row $H^*(R,\mathbb Z)$
and the row $H^{*-2}(\mathbb P^1,\mathbb Z)$; both are free
and in even degrees. It degenerates integrally for degree
reasons, giving the stated groups of $R^+$. The same local
neighbourhood retractions and Van Kampen show that $h$ induces
an isomorphism on fundamental groups, so $R^+$ is simply
connected. The case $R^-$ interchanges the branches.

\item For $Q^+$ use the normalization mapping cylinder once again.
The new degree-two difference map is the old $d_2$ with
one extra coordinate $aF$: the exceptional divisor restricts
to $F$ on the blown-up section and to zero on the other.
It has the same primitive rank-three image and the same
kernel; its cokernel rank increases by one. The degree-four
difference remains primitive, since its pushforward to $R$
already was primitive, and its cokernel likewise increases
by one. This proves the groups of $Q^+$. Interchanging the
sections changes the extra coordinate by a sign and does
not change any group. Smith normal forms are recorded in
\texttt{ab\_push\_local\_checks.out}. The dual complex has
one vertex and one self-edge, with no higher cells, and
Van Kampen again gives $\mathbb Z$.

\item As independent Euler checks, stratification gives
$\chi(Q)=\chi(R)-\chi(S)=24-12$. In the modification $h$,
$C_{\mathrm{reg}}\simeq\mathbb C^*$ contributes zero and
the exceptional chain above the node increases Euler number
by $3-1=2$. Thus $\chi(R^\pm)=26$ and
$\chi(Q^\pm)=26-12=14$.
\end{enumerate}
\end{proof}

\begin{proposition}\label{push:src:A.4.2}\label{push:A42}
\StatementRecord{push:src:A.4.2}{proved}{Historical events: Claude Fable 5.2 (Anthropic), effort high, 16 November 2026: carried as background without fault found. Claude Opus 5.5 (Anthropic), effort max, sweep 1, scope s6, 18 November 2026: computation or formal certificate. GPT-6 Astra (OpenAI), ChatGPT Pro mode, effort Pro, S6\_09\_PAPERS, 17 November 2026: targeted proof check. Also checked in this writing. Pre-continuation writing events: GPT-6 Astra (OpenAI), ChatGPT Pro mode, effort Pro, S6\_10\_PAPERS, 18 November 2026: targeted proof check. Continuation metadata events: GPT-6 Astra (OpenAI), ChatGPT Pro mode, effort Pro, S6\_10\_PAPERS continuation, 19 November 2026: used and rechecked in this run; status-vocabulary reconciliation; metadata reconciliation only; no new mathematical check. \emph{Pre-continuation recorded proof scope.} Re-derived the Thom and Gysin maps, every integral Mayer-Vietoris rank, the oriented positive-u neck tube map K to L/<F,v,m>, its square zero, rank6 and discriminant11 Smith factors. Full matrix replay remains separately tracked.}{the hypotheses, notation and definitions stated here; \cref{push:src:A.2.2}; \cref{push:src:A.4.1}}

The untwisted smooth fibre has torsion-free integral homology
with Betti numbers
\[
 (b_0,b_1,b_2,b_3,b_4,b_5,b_6)=(1,1,8,14,8,1,1).
\]
The positive-$u$ monodromy is the identity outside degree three.
In degree three there is an integral exact sequence
\[
 0\longrightarrow \mathcal L/\langle F,v,m\rangle
 \longrightarrow H_3(X(u);\mathbb Z)
 \longrightarrow K\longrightarrow0,\qquad
 K=\{a\in\mathcal L:aF=av=am=0\},
\]
in which $T_u-I$ is zero on the left term and, on the quotient,
is the map
\[
 K\longrightarrow \mathcal L/\langle F,v,m\rangle,\qquad
 a\longmapsto a\bmod\langle F,v,m\rangle.
\]
The sign convention orients the neck circle by increasing
$\arg x$ in $xy=u$; the boundary convention for a transported
chain is chosen so that a positive $u$ loop adds this circle
times its boundary cycle. With these conventions
$(T_u-I)^2=0$, its rank is six, and its nonzero Smith factors
are $1,1,1,1,1,11$.
\end{proposition}


\begin{proof}
\begin{enumerate}[label=\arabic*.,leftmargin=*]
\item We use the actual semistable model \ref{push:src:A.2.2}. Remove disjoint
open tubular neighbourhoods of the two conductor sections
from $R^+$, and denote the remaining compact manifold by
$Z$. Its two boundary circle bundles are identified by the
plumbing; let $E$ denote the first. The Euler class of
$E\to S$ is $n=m-F$. It is primitive. The Thom sequence
for $(R^+,Z)$ gives
\[
 H_5(Z)=\mathbb Z,\quad H_4(Z)=0,\quad
 H_3(Z)=\mathcal L/\langle F,m\rangle\simeq\mathbb Z^8,
 \quad H_2(Z)\simeq\mathbb Z^{10},\quad H_1(Z)=0.
\]
We spell out the rank and saturation assertions. If $H$ is
the projective-bundle class restricting to zero on $s_0$
and to $-m$ on $s_\infty$, and $A$ is the exceptional
divisor class, the two restriction maps on $H^2(R^+)$ are
\[
 (\alpha,h,a)\longmapsto\alpha+aF,\qquad
 (\alpha,h,a)\longmapsto\alpha-hm.
\]
Their combined image is primitive of rank twelve, and its
cokernel is $\mathcal L/\langle F,m\rangle$. In degree two,
intersection with the two sections is onto $\mathbb Z^2$:
a generic projective fibre meets both once, and an exceptional
projective fibre meets just the section containing the centre
once. In degree six the map is the diagonal
$\mathbb Z\to\mathbb Z^2$. These are precisely the Thom
maps giving the asserted groups.

\item The Gysin sequence of the primitive circle bundle gives
\[
 H_2(E)=n^\perp\subset\mathcal L,\qquad
 H_3(E)=\mathcal L/\mathbb Z n,\qquad
 H_1(E)=H_4(E)=0,\quad H_5(E)=\mathbb Z.
\]
The smooth fibre is obtained from $Z$ by the boundary
mapping cylinder. In its Mayer--Vietoris sequence the
degree-two difference map is the normalization difference
$d_2^+$ restricted to $n^\perp$. In the coefficient
coordinates $(s,t,h)=(aF,av,am)$ the constraint is $h=s$,
so this map has primitive rank two and kernel $K$.
The degree-three difference map is, modulo $F,m$,
$(T_*-I)a=(aF)v$, hence has primitive image $\mathbb Z v$.
These statements use the surjectivity of the three
intersection coordinates proved in \ref{push:src:A.4.1}, so no integral
index is omitted. The degree-five difference is zero:
the plumbing reverses the normal circle orientation, and
the two boundary fundamental classes sum to zero in $Z$.
The degree-zero difference is zero as well. Exactness now
gives the claimed groups and the displayed degree-three
extension. Poincar\'e duality, or the same integral exact
sequence, confirms the absence of torsion in complementary
degrees.

\item We also compute monodromy directly from this decomposition.
In the neck $xy=u$, a positive circuit of $u$ rotates one
boundary normal coordinate once relative to the other.
A cycle in $Z$ is unchanged. A three-cycle whose cut-open
boundary represents $a\in K$ acquires the tube obtained
by sweeping $a$ once around that normal circle. The Gysin
description sends this tube to the class of $a$ in
$H_3(E)=\mathcal L/\mathbb Z n$, and then to its class in
$\mathcal L/\langle F,m,v\rangle$. This is exactly the
map in the statement. The same description in the other
degrees is zero, using the groups and difference maps just
computed. A tube has no cut-open boundary, which proves
$(T_u-I)^2=0$ without an appeal to a monodromy theorem.

\item Finally $K\cap\langle F,m,v\rangle=\mathbb ZF$, since
the Gram matrix of $m-F,v$ has determinant eleven.
Hence the rank is six. In the unimodular $E_8$ summand,
the primitive rank-two lattice spanned by $m-F,v$ has
discriminant eleven; its orthogonal complement maps into
the dual quotient with cokernel of order eleven. Thus the
nonzero Smith factors are five ones and eleven. The fully
integral check uses an explicit $E_8$ basis in
\texttt{ab\_push\_monodromy.py}; its output records the
kernel basis, quotient projection, and all entries of the
seven-by-seven monodromy block. This supplies a second,
exact lattice verification of the discriminant argument.
\end{enumerate}
\end{proof}

\begin{proposition}[the pointwise nearby Euler function]\label{push:src:A.4.3}\label{push:A43}
\StatementRecord{push:src:A.4.3}{proved}{Historical events: Claude Fable 5.2 (Anthropic), effort high, 16 November 2026: carried as background without fault found. GPT-6 Astra (OpenAI), ChatGPT Pro mode, effort Pro, S6\_09\_PAPERS, 17 November 2026: targeted proof check. Also checked in this writing. Pre-continuation writing events: GPT-6 Astra (OpenAI), ChatGPT Pro mode, effort Pro, S6\_10\_PAPERS, 18 November 2026: targeted proof check. Continuation metadata events: GPT-6 Astra (OpenAI), ChatGPT Pro mode, effort Pro, S6\_10\_PAPERS continuation, 19 November 2026: used and rechecked in this run; status-vocabulary reconciliation; metadata reconciliation only; no new mathematical check. \emph{Pre-continuation recorded proof scope.} Re-derived actual local nearby Euler values from xy=u,xy=uz,xy=uzv and the small resolution; central integral1 for nodal plumbing and nearby2 for semistable resolution.}{the hypotheses, notation and definitions stated here; \cref{push:src:A.3.2}}

For $\mathcal X^\flat\to\Delta_u$ the nearby-fibre Euler
function on the double locus of $Q$ is zero off $S_\infty$
and one at every point of $S_\infty$, including its node.
Its integral is one. In particular the assertion that the
pointwise value on the smooth part of $S_\infty$ is zero
is false; its integrated contribution is zero.
For either semistable family $\mathcal X^\pm$ the nearby
Euler function is one off the double locus and zero on
it, giving $\chi(X(u))=26-2\cdot12=2$.
\end{proposition}


\begin{proof}
\begin{enumerate}[label=\arabic*.,leftmargin=*]
\item Off the zero divisor of $\phi$, the transverse nearby
fibre of $xy=u\phi$ is an annulus and has Euler number
zero. At a smooth zero, the local equation is $xy=uz$;
for fixed $u\ne0$ one solves for $z$, and the small
local nearby space is contractible. Its Euler number
is one. At the node it is the quadratic cone
$xy=uzv$, also contractible, with Euler number one.
The smooth part of the nodal elliptic curve is
$\mathbb C^*$ and has Euler number zero; the node
contributes one. The core $Q\setminus S$ is a
$\mathbb C^*$-bundle and contributes zero.
This proves the first count.

\item In the resolved family the local map on the double locus
is simply $ab=u$ times two transverse coordinates,
including above the old node. Its nearby annulus has
Euler number zero everywhere on that locus. Away from
it, the family is a submersion and the local value is
one. Integrating gives the second count. If the nodal
threefold in the first family is smoothed rather than
small-resolved, its ordinary-double-point Milnor fibre
is homotopy equivalent to $S^3$ and has Euler number
zero; the replacement changes $1$ to $0$, whereas
the small resolution changes $1$ to $2$. These agree
with the intrinsic virtual number in \ref{push:src:A.3.2}.
\end{enumerate}
\end{proof}

\begin{proposition}[limitation of Clemens--Schmid]\label{push:src:A.4.4}\label{push:A44}
\StatementRecord{push:src:A.4.4}{proved}{Historical events: Claude Fable 5.2 (Anthropic), effort high, 16 November 2026: carried as background without fault found. GPT-6 Astra (OpenAI), ChatGPT Pro mode, effort Pro, S6\_09\_PAPERS, 17 November 2026: targeted proof check. Also checked in this writing. Pre-continuation writing events: GPT-6 Astra (OpenAI), ChatGPT Pro mode, effort Pro, S6\_10\_PAPERS, 18 November 2026: targeted proof check. Continuation metadata events: GPT-6 Astra (OpenAI), ChatGPT Pro mode, effort Pro, S6\_10\_PAPERS continuation, 19 November 2026: used and rechecked in this run; status-vocabulary reconciliation; metadata reconciliation only; no new mathematical check. \emph{Pre-continuation recorded proof scope.} Re-derived odd b1, null small-exceptional curve from the unimodular generization matrix, and monodromy filtration graded ranks6,2,6. No limiting polarized mixed-Hodge identification inferred.}{the hypotheses, notation and definitions stated here; \cref{push:src:A.4.1}}

The classical K\"ahler or projective Clemens--Schmid
comparison cannot be applied to these untwisted models.
Their smooth fibres have $b_1=1$ and are not K\"ahler.
Moreover the exceptional small-resolution curve in
$R^\pm$ is null-homologous, so that normalization is
not K\"ahler either. The rational monodromy weight
filtration, defined algebraically from $N=T_u-I$ and
centred at degree three, has graded ranks $6,2,6$ in
weights $2,3,4$. This is a filtration of the computed
monodromy representation; we do not identify it with
a polarized limiting mixed Hodge structure.
\end{proposition}


\begin{proof}
\begin{enumerate}[label=\arabic*.,leftmargin=*]
\item The odd first Betti number excludes a K\"ahler metric.
In the generization matrix of \ref{push:src:A.4.1} a global section of
$R^2h_*\mathbb Z$ has equal branch values, hence degree
zero on the small exceptional curve. All classes pulled
back from $R$ also have degree zero on it. The Leray
calculation exhausts $H^2(R^\pm)$, so the curve has
zero homology class; it cannot have positive area for
a K\"ahler form. Finally $N^2=0$, rank $N=6$, and
$\dim H_3=14$, so the centred filtration is
$W_2=\operatorname{im}N$, $W_3=\ker N$, $W_4=H_3$.
Its graded dimensions are the asserted ones. These
topological computations replace, rather than assume,
the inapplicable K\"ahler comparison.
\end{enumerate}
\end{proof}

\begin{corollary}[unrestricted toroidal topology]\label{push:src:A.4.5}\label{push:A45}
\StatementRecord{push:src:A.4.5}{proved}{Historical events: Claude Fable 5.2 (Anthropic), effort high, 16 November 2026: carried as background without fault found. GPT-6 Astra (OpenAI), ChatGPT Pro mode, effort Pro, S6\_09\_PAPERS, 17 November 2026: targeted proof check. Also checked in this writing. Pre-continuation writing events: GPT-6 Astra (OpenAI), ChatGPT Pro mode, effort Pro, S6\_10\_PAPERS, 18 November 2026: targeted proof check. Continuation metadata events: GPT-6 Astra (OpenAI), ChatGPT Pro mode, effort Pro, S6\_10\_PAPERS continuation, 19 November 2026: used and rechecked in this run; status-vocabulary reconciliation; metadata reconciliation only; no new mathematical check. \emph{Pre-continuation recorded proof scope.} Re-derived the finite-data deduction from the exact k,m,t4 homology table, primitive integral blow-up splitting and monodromy Smith factors; supplied a numbered proof.}{\cref{push:src:A.2.4}, \cref{push:src:A.4.4}, \cref{push:src:A.4.7}, \cref{push:src:A.4.8}, \cref{push:src:A.4.9}}

For every output of \ref{push:src:A.2.4}, \ref{push:src:A.4.7}--\ref{push:src:A.4.9} determine the
Euler number, fundamental group, all integral homology
groups with torsion, nearby Euler function and marked
monodromy. The integral normalization computation is
given by the finite algorithm and the primitive
central pushforward matrices of \ref{push:src:A.4.9}; the independent
ambient computation is \ref{push:src:A.4.8}, and the torus-stratum
Euler count is \ref{push:src:A.4.7}. The non-K\"ahler limitation on
classical Clemens--Schmid is the explicit one in
\ref{push:src:A.4.4}; no polarized limiting Hodge structure is
asserted.
\end{corollary}

\begin{proof}
\begin{enumerate}[label=\arabic*.,leftmargin=*]
\item An output of \ref{push:src:A.2.4} has finitely many ray and
cone orbits with specified integral maps. Count its pure top
rays $k$, mixed rays $m$ and maximal four-cones $t_4$.
Theorem~\ref{push:src:A.4.7} gives $\chi=11k+t_4$ and
$\pi_1=\Z$, and evaluates the nearby Euler function on
each branch stratum; its integral is two.

\item Theorem~\ref{push:src:A.4.8} gives the full free integral
homology table and the marked degree-three monodromy $I+nN$,
with six nonzero Smith factors $n,n,n,n,n,11n$. The integral
splitting of \ref{push:src:A.4.9} computes the same groups
directly along the finite sequence of normalization modifications;
its exceptional-fibre diagonal entries are units, so it also
determines primitivity and torsion. Thus every output named in
the statement is determined by the indicated finite data.
The exact failure of the particular polarized comparison is
already the calculation in \ref{push:src:A.4.4}; no additional
Hodge-theoretic conclusion is used here.
\end{enumerate}
\end{proof}

\Uses{the hypotheses, notation and definitions stated here; \cref{push:src:A.2.4}; \cref{push:src:A.4.7}; \cref{push:src:A.4.9}; \cref{push:src:A.4.8}; \cref{push:src:A.4.4}}


\begin{proposition}[an infinite explicit subclass]\label{push:src:A.4.6}\label{push:A46}
\StatementRecord{push:src:A.4.6}{proved}{Historical events: Claude Fable 5.2 (Anthropic), effort high, 16 November 2026: carried as background without fault found. GPT-6 Astra (OpenAI), ChatGPT Pro mode, effort Pro, S6\_09\_PAPERS, 17 November 2026: targeted proof check. Also checked in this writing. Pre-continuation writing events: GPT-6 Astra (OpenAI), ChatGPT Pro mode, effort Pro, S6\_10\_PAPERS, 18 November 2026: targeted proof check. Continuation metadata events: GPT-6 Astra (OpenAI), ChatGPT Pro mode, effort Pro, S6\_10\_PAPERS continuation, 19 November 2026: used and rechecked in this run; status-vocabulary reconciliation; metadata reconciliation only; no new mathematical check. \emph{Pre-continuation recorded proof scope.} Re-derived every point-star normalization, the chain of P3 components and P2 conductor maps, primitive integral Mayer-Vietoris induction adding one in degrees2,4,6, and nearby Euler2 for all m.}{the hypotheses, notation and definitions stated here; \cref{push:src:A.2.1}; \cref{push:src:A.4.2}; \cref{push:src:A.4.1}}

Let $D_m$ be the central fibre after the $m$ maximal-cone
star subdivisions in \ref{push:src:A.2.1}, starting with $Q^+$, where
$m\geq0$. Then
\[
 \begin{gathered}
 \pi_1(D_m)=\mathbb Z,\qquad
 H_*(D_m;\mathbb Z)\text{ is free},\\
 (b_0,\ldots,b_6)=(1,1,9+m,7,11+m,0,1+m).
 \end{gathered}
\]
and $\chi(D_m)=14+3m$. For $m\geq1$ its normalization is
the disjoint union of $\operatorname{Bl}_pR^+$, $m-1$
copies of $\operatorname{Bl}_{p_i}\mathbb P^3$, and a final
$\mathbb P^3$. The original two sections are still glued
by $T$. The new components form a chain attached to the
original component. Consecutive components meet in one
$\mathbb P^2$, with no triple intersections. The dual
complex is the original circle with a chain of $m$ edges
attached. The nearby Euler function is one off the double
locus and zero on it; its integral remains two. The nearby
integral monodromy remains exactly that of \ref{push:src:A.4.2}.
\end{proposition}


\begin{proof}
\begin{enumerate}[label=\arabic*.,leftmargin=*]
\item The torus orbit of the cone $ABCE$ is a point on a single
local central branch, because precisely one of its four
rays has positive $u$-height. Hence it is a smooth point
of the central fibre, away from the conductor. A star at
the sum of the four basis rays is the ordinary blow-up
of that point in the smooth fourfold. The central divisor
is the union of its strict transform and the exceptional
$\mathbb P^3$, both with multiplicity one. Their
intersection is the exceptional $\mathbb P^2$ in the
blow-up of the central threefold, identified with a
hyperplane in $\mathbb P^3$. Their two normal lines there
are $\mathcal O(-1)$ and $\mathcal O(1)$.

\item The next distinguished cone $ABCv_1$ has only the new
central ray $v_1$. Its point lies on the new $\mathbb P^3$
away from its attaching hyperplane. Repeating therefore
blows up that component at a point away from the preceding
double locus and attaches the next $\mathbb P^3$. The
incoming hyperplane and outgoing exceptional plane are
disjoint in each middle $\operatorname{Bl}_{p_i}\mathbb P^3$.
This proves the complete normalization and gluing
description, the absence of triples and the dual complex.

\item For the integral calculation, blowing up a smooth point
of a compact complex threefold adds one free generator
in each of $H_2$ and $H_4$ and changes no other homology.
This also holds when the threefold has singularities away
from that point: remove a small ball and use the
Mayer--Vietoris sequence with boundary $S^5$; the
replacement retracts to $\mathbb P^2$. Now attach
$\mathbb P^3$ along that $\mathbb P^2$. In degrees
$0,2,4$ the inclusion into $\mathbb P^3$ is an integral
isomorphism, so the Mayer--Vietoris boundary map is
injective and primitive. It contributes no torsion or
odd homology. The new component contributes one free
$H_6$ generator, since $H_6(\mathbb P^2)=0$. Each step
therefore adds precisely one to $b_2,b_4,b_6$.
Van Kampen gives no change of fundamental group: the
blow-up ball replacement, the attaching plane and the
new projective component are simply connected.
Induction from \ref{push:src:A.4.1} proves the asserted groups.

\item Independently, the strict-transform blow-up adds two to
the central Euler number, the new $\mathbb P^3$ adds
four, and its intersection $\mathbb P^2$ subtracts
three. Thus each step adds three, agreeing with the
homology computation. The double locus is the disjoint
union of the original $S$ and $m$ planes, of Euler number
$12+3m$. Since the local equation is $u=x$ off this locus
and $u=xy$ on it, the nearby Euler integral is
$(14+3m)-(12+3m)=2$. All modifications are the identity
over $u\ne0$, so they do not alter the marked punctured
family or its homological monodromy. The exact saved
certificates check all determinants, boundary faces,
slice volumes and these integer formulas for $0\leq m\leq8$;
the displayed induction proves them for every $m$.
\end{enumerate}
\end{proof}

\begin{theorem}[Euler number and fundamental group for every toroidal output]\label{push:src:A.4.7}\label{push:A47}
\StatementRecord{push:src:A.4.7}{proved}{Historical events: Claude Fable 5.2 (Anthropic), effort high, 16 November 2026: carried as background without fault found. GPT-6 Astra (OpenAI), ChatGPT Pro mode, effort Pro, S6\_09\_PAPERS, 17 November 2026: targeted proof check. Also checked in this writing. Pre-continuation writing events: GPT-6 Astra (OpenAI), ChatGPT Pro mode, effort Pro, S6\_10\_PAPERS, 18 November 2026: targeted proof check. Continuation metadata events: GPT-6 Astra (OpenAI), ChatGPT Pro mode, effort Pro, S6\_10\_PAPERS continuation, 19 November 2026: used and rechecked in this run; status-vocabulary reconciliation; metadata reconciliation only; no new mathematical check. \emph{Pre-continuation recorded proof scope.} Re-derived torus-orbit stabilizers, Euler11k+t4, the two bottom-triangle contributions to nearby Euler2, common-modification fundamental group and monomial gradient exponent. BM95 input already matched separately; universal gradient argument read and independently derived.}{the hypotheses, notation and definitions stated here; \cref{push:src:A.2.4}; \cite[Theorem 1.6 and its following paragraph]{BM95AB}; \cref{push:src:B.4.2}}

Let $\Sigma$ be any output of \ref{push:src:A.2.4} for $u=r^n$, and let
$D_\Sigma$ denote the underlying reduced central space.
Let $k(\Sigma)$ be the number of selected pure top ray
orbits, and let $t_4(\Sigma)$ be the number of maximal
four-cones in one fundamental prism. Then
\[
 \chi(D_\Sigma)=11k(\Sigma)+t_4(\Sigma),
 \qquad \pi_1(D_\Sigma)=\mathbb Z,
 \qquad H_1(D_\Sigma;\mathbb Z)=\mathbb Z.
\]
These assertions include nonreduced central divisors.
At a point with precisely one local central branch of
multiplicity $\ell$, the nearby Euler function is $\ell$;
at a point with at least two local branches it is zero.
Its integral is always two. For every reduced output the
pure top count is $k(\Sigma)=n$, so
$\chi(D_\Sigma)=11n+t_4(\Sigma)$.
The higher integral homology of every arbitrary output is
not asserted by this theorem.
\end{theorem}


\begin{proof}
\begin{enumerate}[label=\arabic*.,leftmargin=*]
\item Over $B\setminus\{\infty\}$ the top fan gives a cycle
of $k(\Sigma)$ ruled bundles over
$S\setminus S_\infty$, glued along their zero and
infinity sections with the specified translations.
There are no triple intersections there. Its Euler
number is therefore
\[
 2k(\Sigma)\chi(S\setminus S_\infty)
       -k(\Sigma)\chi(S\setminus S_\infty)
 =11k(\Sigma),
\]
because $\chi(S)=12$ and $\chi(S_\infty)=1$.
This computation concerns the support of the divisor,
so the multiplicities do not enter.

\item At $q=r=0$ the strata are the torus orbits of cones whose
generators collectively have positive $q$-height and
positive $r$-height. Each such cone has trivial period
stabilizer. Indeed, an element $G_1^aG_2^b$ preserving
its finite set of primitive rays has a power fixing them
all. A ray of positive $q$-height then forces $a=0$;
any ray with positive $q$- or $r$-height forces $b=0$.
Consequently a stratum of cone dimension $d$ is an actual
$(\mathbb C^*)^{4-d}$, with no additional torus quotient.
The induced gluing units only identify different period
translates of the stratum. Its compactly supported Euler
number is zero unless $d=4$, in which case it is one.
Every maximal cone has both kinds of height, and each has
one representative in the interior of the fundamental
prism. Thus the cusp contributes exactly $t_4(\Sigma)$.
Additivity proves the displayed central Euler formula.

\item The local total space is smooth, and its parameter is
$r=z_1^{\ell_1}\cdots z_s^{\ell_s}$ times a holomorphic
unit, with $\ell_i>0$. Its Milnor fibre has the homotopy
type of a finite disjoint union of $(S^1)^{s-1}$.
For $s=1$ it consists of $\ell_1$ contractible pieces;
for $s\geq2$ its Euler number is zero. This proves the
pointwise assertion. Off the cusp the one-branch parts
are $\mathbb C^*$-bundles over the appropriate parts of
$S\setminus S_\infty$ and have Euler number zero; the
double parts have nearby Euler value zero. At the cusp,
positive-dimensional torus strata again contribute zero.
A maximal cone with exactly one positive-$r$ ray must
have three bottom rays. Its bottom face is consequently
one of the two prescribed triangles $ABC$ and $BCD$.
Each triangle belongs to exactly one such maximal cone.
Its determinant is the $r$-height $\ell$ of the fourth
ray, because the bottom triangle is unimodular.
Smoothness forces $\ell=1$. Exactly these two points
contribute one each. Thus the nearby Euler integral is
two for every output, even when some other components
have multiplicity greater than one.

\item For the fundamental group, first choose the reference
model obtained by base-changing $\mathcal X^+$ by
$u=r^n$ and resolving the surface equation $xy=r^n$
along its smooth conductor $S$. Its normalized central
components are $R^+$ and $n-1$ projective-line bundles
over $S$. All are simply connected, and their double
surfaces are copies of the simply connected $S$.
The gluing graph is an $n$-cycle, interpreted as the
self-gluing circle when $n=1$. Van Kampen therefore
gives fundamental group $\mathbb Z$ for this reference
centre.

\item Any two fans of \ref{push:src:A.2.4} have the explicit proper common
normal modification formed by intersecting their cones.
Resolve this analytic space to a smooth common model;
analytic resolution, including the noncompact version
near the compact central set, is supplied by
\cite[Theorem 1.6 and its following paragraph]{BM95AB}.
This use of a common modification is legitimate here
because the common fan has already constructed it.
It is not an assumed extension of an arbitrary punctured
isomorphism.

\item A proper bimeromorphic map $F:Y\to X$ between smooth
complex manifolds induces an isomorphism of fundamental
groups. To see this directly, choose an analytic
$A\subset X$ of complex codimension at least two such
that $F$ is an isomorphism outside $A$. General position
gives $\pi_1(X\setminus A)\simeq\pi_1(X)$, and gives
surjectivity of
$\pi_1(Y\setminus F^{-1}A)\to\pi_1(Y)$.
Under the common-open identification, composition with
$F_*$ is the displayed isomorphism to $\pi_1(X)$.
Thus the first map is also injective, and $F_*$ is an
isomorphism. Apply this argument to both maps from the
smooth common model, after restricting to a sufficiently
small parameter disc.

\item Finally, each smooth total space retracts to the support
of its central divisor, including in the nonreduced case.
The gradient proof in \ref{push:src:B.4.2} applies with the explicit
bound for $\rho=|r|^2$
\[
 \|\nabla\rho\|\geq c\rho^{1-1/(2L)},
 \qquad L=\ell_1+\cdots+\ell_s.
\]
On the complement of the divisor it follows from the
weighted arithmetic--geometric mean inequality applied
to $\sum_i\ell_i^2/|z_i|^2$. A holomorphic unit can be
absorbed into one coordinate locally. Compactness of
the central set gives a uniform exponent below one,
and the resulting finite remaining path-length bound
gives the continuous gradient retraction. Hence all
the centres and their smooth total neighbourhoods have
the same fundamental group as the reference centre.
Abelianization gives the stated $H_1$.

\item The program \path{code/ab_push_all_fan_topology.py}
checks the Euler and two-point nearby formulas for every
minimal bridge at $n=1,2,12$, eight point-star
modifications, and five nonreduced conductor-star
modifications. In the latter family the added pure top
rays are $(0,1,0,j)$ for $2\leq j\leq m+1$; the formula
gives $\chi(D_m)=14+12m$, independently agreeing with
successive blow-ups of the conductor surface of Euler
number twelve. The finite checks are recorded separately
from the proof for every fan.
\end{enumerate}
\end{proof}

\begin{theorem}[integral homology of every toroidal output]\label{push:src:A.4.8}\label{push:A48}
\StatementRecord{push:src:A.4.8}{proved}{Historical events: Claude Fable 5.2 (Anthropic), effort high, 16 November 2026: carried as background without fault found. Claude Opus 5.5 (Anthropic), effort max, sweep 1, scope s6, 18 November 2026: computation or formal certificate. GPT-6 Astra (OpenAI), ChatGPT Pro mode, effort Pro, S6\_09\_PAPERS, 17 November 2026: targeted proof check. Also checked in this writing. Pre-continuation writing events: GPT-6 Astra (OpenAI), ChatGPT Pro mode, effort Pro, S6\_10\_PAPERS, 18 November 2026: targeted proof check. Continuation metadata events: GPT-6 Astra (OpenAI), ChatGPT Pro mode, effort Pro, S6\_10\_PAPERS continuation, 19 November 2026: used and rechecked in this run; status-vocabulary reconciliation; metadata reconciliation only; no new mathematical check. \emph{Pre-continuation recorded proof scope.} Re-derived reference n-cycle integral elimination, precise finite-cone-complex factorization, all compact and horizontal centre types, integral blow-up splitting, invariants b2-(boundary rays), b6-(central rays), full Betti tuple and I+nN monodromy. AKMW locators opened; Matsuki correction source gate and bounded implementation replays remain separately tracked.}{the hypotheses, notation and definitions stated here; \cref{push:src:A.2.4}; \cref{push:src:A.4.2}; \cref{push:src:A.4.1}; \cite[Section 5.6, pp.565--566]{AKMW02AB}; \cite[Theorems 1.6.1 and 2.7.1]{AKMW02AB}; \cite[Section 2, p.2]{Matsuki00ErrAB}; \cite[Section 5.5]{AKMW02AB}; \cref{push:src:A.4.7}}

Let $D_\Sigma$ be the reduced central space of an output of
\ref{push:src:A.2.4} for $u=r^n$. Write $k$ for its number of pure top
ray orbits, $m$ for its number of mixed ray orbits, and
$t_4$ for its number of maximal four-cone orbits. Then
every integral homology group is free, and
\[
 \begin{split}
 (b_0,b_1,b_2,b_3,b_4,b_5,b_6)
   ={}&(1,1,8+k+m,7,\\
       &\hspace{20mm}9k+t_4-2m-1,0,k+m).
 \end{split}
\]
This includes nonreduced central divisors: the homology
is that of their underlying support. The marked nearby
fibre is $X(r^n)$, and its integral monodromy is
$I+nN$ in degree three, where $N$ is the map in \ref{push:src:A.4.2};
it is the identity in every other degree. In particular
the six nonzero Smith factors of this degree-three
monodromy minus the identity are
$n,n,n,n,n,11n$.
\end{theorem}


\begin{proof}
\begin{enumerate}[label=\arabic*.,leftmargin=*]
\item We first compute a reference centre by its normalization
mapping cylinder. Base-change $\mathcal X^+$ by $u=r^n$
and resolve $xy=r^n$ along its smooth conductor $S$.
Put $N_0=\mathcal O_S(m-F)$, using the lattice notation
of \ref{push:src:A.4.1}. The normalization consists of $R^+$ and
$n-1$ copies of $\mathbb P_S(\mathcal O\oplus N_0)$,
in a cycle. The latter assertion follows directly on
the surface-resolution charts: if $x$ has transition
function $g$, the coordinates $x/r^i$ have the same
transition function, while their complementary
coordinates have transition function $g^{-1}$.
Only the closing identification is $T$.

\item In the integral mapping-cylinder sequence the degree-two
map has domain $\mathcal L^n$ and target of rank $11n+1$.
Project each intermediate ruled component to $S$.
The resulting $n-1$ difference blocks are integral
identity matrices. Eliminate them by integral row and
column operations. The remaining map from $\mathcal L$
is $d_2^+$ from \ref{push:src:A.4.1}, together with $n-1$ coordinates
$a\mapsto a(m-F)$. Those coordinates are integral
combinations of the three primitive coordinates of
$d_2^+$. Thus the original degree-two map has primitive
image of rank $10(n-1)+3$, kernel $\mathbb Z^7$, and
cokernel $\mathbb Z^{n+8}$.

\item In degree four, the same elimination uses the $n-1$
identity blocks on the fundamental classes of the bases.
The remaining map from $\mathbb Z$ contains $d_4^+$,
which is primitive by \ref{push:src:A.4.1}; the other entries are the
section differences $m-F$. Hence the degree-four map
is injective with primitive image of rank $n$, and
cokernel of rank $10n+1$. In degree zero it is the
incidence map of an $n$-cycle. In degree six the
normalization contributes one independent fundamental
class per component. Consequently the reference centre
has free integral homology and Betti tuple
\[
 (1,1,n+8,7,10n+1,0,n).
\]
This is also a direct integral calculation of all the
maps needed below; it uses no degeneration theorem.

\item We next justify the precise factorization used to compare
an arbitrary output with this reference. On each total
space take the toroidal boundary
\[
 E=\bigl(\{r=0\}\cup f^{-1}(\infty)\bigr)_{\rm red}.
\]
In the cusp charts it is the toric boundary, and outside
the cusp its nontrivial toroidal charts are the plumbing
charts along $S$. The unit transition functions preserve
monomial ideals. \ref{push:src:A.2.4} gives a finite cone complex modulo
the two periods, with the possible self-identifications
of its faces explicitly retained.

\item Make this complex strict by barycentric star subdivisions,
in decreasing cone dimension. At the moment a stratum
is blown up, all of its deeper self-intersections have
already been separated. Its closure is therefore smooth,
and different translates in the chosen orbit are disjoint.
The star at the sum of the primitive basis generators is
the ordinary blow-up of that closure. Repeating over the
finitely many cone orbits gives smooth strict toroidal
embeddings. This is the geometric barycentric subdivision
used in \cite[Section 5.6, pp.565--566]{AKMW02AB}; its
analytic realization by the same monomial ideals is
described in Section 5.5 of that paper. Pure $q$-strata
are allowed in this preparatory step. Thus we do not
assume that strictification preserves the punctured
$r$-family.

\item Intersect the two subdivided fans, and take a regular
subdivision, to obtain their explicit common strict
toroidal modification. Apply toroidal weak factorization
to each of its two subdivision maps. The needed theorem
is the finite polyhedral-complex factorization by
nonsingular star subdivisions and inverse subdivisions,
\cite[Theorems 1.6.1 and 2.7.1]{AKMW02AB}. The corrected
form for a morphism states that its centres are smooth
irreducible closed strata and that all intermediate
complexes map to the target
\cite[Section 2, p.2]{Matsuki00ErrAB}.

\item There is no invocation here of a general relative
analytic weak-factorization theorem. The complexes in
question are finite. The cited combinatorial
factorization applies to them; the analytic
subdivision--modification correspondence of
\cite[Section 5.5]{AKMW02AB} realizes each of its finitely
many stars by the indicated monomial blow-up.
The correspondence respects the face identifications
and the bundle units, so the local analytic blow-ups
glue. All intermediate morphisms are proper over the
target of the subdivision map, hence over a sufficiently
small $r$-disc. They remain smooth, and both parameters
remain monomials times units. This gives the required
finite chain of actual analytic blow-ups and blow-downs
between the total neighbourhoods.

\item We classify the connected centres in this chain. A
nontrivial centre has codimension at least two. For a
cone having a positive $q$-height ray, its period
stabilizer is trivial: a power of a stabilizing element
fixes each of its finitely many primitive rays; the
positive $q$-height first forces the exponent of $G_1$
to vanish and then the exponent of $G_2$ to vanish.
The same argument holds for pure bottom cones. Thus
there is no elliptic quotient hidden in those centres.
The normalization of the closure of their torus stratum
is the toric variety of its star. If the closed stratum
is smooth, this finite birational normalization map is
an isomorphism.

\item For a mixed cone the centre is contained in $q=r=0$;
its star is complete, and it is a compact smooth toric
point, curve, or surface. It has free integral
cohomology in even degrees only. For a pure top centre,
the cone must have dimension two. Its stabilizer is
precisely $\langle G_1\rangle$, and the chart quotient
and the plumbing identification show that its dense
open part is the corresponding open part of $S$.
The compact smooth centre is therefore a surface
birational to the rational surface $S$. Equivalently,
its proper map to the unmodified top stratum is a
birational map of smooth surfaces. It is a smooth
rational surface and again has free even cohomology.
It need not remain isomorphic to $S$ after earlier
mixed modifications.

\item For a pure bottom centre the cone has dimension two
or three. Its star gives a smooth toric surface or
curve proper over the $r$-disc. The induced character
$r$ is primitive on its quotient lattice, since it
was primitive before quotienting by a smooth cone
on which it vanishes. Its generic fibre is therefore
connected, and is respectively a complete toric
$\mathbb P^1$ or a point. In the surface case the
ordered two-dimensional fan gives a chain of rational
curves as the central support, possibly with positive
multiplicities. In the curve case the support is one
point. Monomial gradient retraction, as in \ref{push:src:A.4.7},
shows that the ordinary integral cohomology of the
centre is free and even, with $H^0=\mathbb Z$ and
$H^4=0$. These centres are not compact; this is why
they must be distinguished from the first two types.

\item For any closed smooth complex centre $C\subset Y$ of
complex codimension $c$, compact or not, the integral
blow-up splitting is
\[
 H^j(\operatorname{Bl}_C Y;\mathbb Z)
 \simeq H^j(Y;\mathbb Z)\oplus
       \bigoplus_{a=1}^{c-1}H^{j-2a}(C;\mathbb Z).
\]
For clarity, this statement requires no K\"ahler
hypothesis. Tubular excision identifies the relative
groups by the integral Thom class; the projective-bundle
theorem for the exceptional divisor supplies the powers
of its hyperplane class. Their exceptional Gysin maps,
together with pullback from $Y$, give the displayed
map. The relation that restriction of the exceptional
class is minus the hyperplane class makes the relative
matrix triangular with unit diagonal. Proper
blow-down pushforward, defined by Borel--Moore duality,
splits the pullback summand. This proves the splitting
also for the noncompact centres just listed.

\item It follows that all odd cohomology groups and all
torsion are unchanged along the factorization. Each
connected centre adds one to $b_2$ and one to the number
of prime boundary components. Each compact centre
adds one to $b_6$ and one to the number of central
prime components: in degree six the only new group
is its top cohomology. A pure bottom centre adds
neither, because its $H^4$ (and, for a curve, its
$H^2$) vanishes. The reverse operations subtract
the same contributions.

\item Consequently $b_2$ minus the number of all boundary
ray orbits, and $b_6$ minus the number of central ray
orbits, are invariants of this comparison. Both
endpoints have exactly one pure bottom ray orbit.
The reference has $n$ central components and
$b_2=n+8$, so the arbitrary output has
\[
 b_2=8+k+m,\qquad b_6=k+m,
 \qquad (b_1,b_3,b_5)=(1,7,0).
\]
All groups are free by the integral splitting in both
directions. The total neighbourhoods retract onto
their central supports by \ref{push:src:A.4.7}, and those compact
analytic supports have finite triangulations. The
universal coefficient theorem therefore transfers
the same freeness and ranks to integral homology.
Finally the independently computed torus-stratum
Euler number $11k+t_4$ of \ref{push:src:A.4.7} gives
$b_4=9k+t_4-2m-1$.

\item Every original output, unlike temporary intermediate
factorizations, agrees with the marked family
$X(r^n)$ off $r=0$. A positive $r$-loop makes $n$
positive $u$-loops. Since $N^2=0$, its monodromy is
$(I+N)^n=I+nN$. Multiplication of the already computed
integral Smith matrix by $n$ gives the stated factors.

\item The exact program \path{code/ab_push_all_fan_homology.py}
checks the full reference normalization maps and their
Smith forms at $n=1,2,3,4,6,12$, the blow-up rank
increments for every centre type used here, and the
resulting formula on the certified fan families of
\ref{push:src:A.4.7}. These finite checks supplement the integral
row-reduction and factorization argument for every
fan; they are not a bounded search substituted for it.
\end{enumerate}
\end{proof}

\begin{proposition}[a second integral central-space computation]\label{push:src:A.4.9}\label{push:A49}
\StatementRecord{push:src:A.4.9}{proved}{Historical events: Claude Fable 5.2 (Anthropic), effort high, 16 November 2026: carried as background without fault found. GPT-6 Astra (OpenAI), ChatGPT Pro mode, effort Pro, S6\_09\_PAPERS, 17 November 2026: targeted proof check. Also checked in this writing. Pre-continuation writing events: GPT-6 Astra (OpenAI), ChatGPT Pro mode, effort Pro, S6\_10\_PAPERS, 18 November 2026: targeted proof check. Continuation metadata events: GPT-6 Astra (OpenAI), ChatGPT Pro mode, effort Pro, S6\_10\_PAPERS continuation, 19 November 2026: used and rechecked in this run; status-vocabulary reconciliation; metadata reconciliation only; no new mathematical check. \emph{Pre-continuation recorded proof scope.} Re-derived actual integral derived-sheaf splitting using exceptional Gysin maps with stalk basis1,-xi,..., proper base change to central support, and simplex contraction of normalized local-branch complex. This is a sheaf map calculation, not a rank-only inference.}{the hypotheses, notation and definitions stated here; \cref{push:src:A.4.8}; \cref{push:src:A.2.5}; \cref{push:src:A.4.7}}

The homology formula of \ref{push:src:A.4.8} can also be computed
on the central spaces, by their normalization
resolutions and the strata of each recorded star.
More precisely, let $Y'\to Y$ be one ordinary
blow-up in the algorithm of \ref{push:src:A.2.5}, with closed
smooth centre $C$ of complex codimension $c$.
Let $D,D'$ be the reduced central supports and
$C_0=C\cap D$. In the derived category of integral
sheaves on $D$ there is an isomorphism
\[
 R(\pi_D)_*\mathbb Z_{D'}\simeq
 \mathbb Z_D\oplus
 \bigoplus_{a=1}^{c-1}(i_0)_*\mathbb Z_{C_0}[-2a],
 \tag{A.49a}
\]
where $\pi_D:D'\to D$ and $i_0:C_0\hookrightarrow D$
are the induced continuous maps. In particular
\[
 H^j(D';\mathbb Z)\simeq H^j(D;\mathbb Z)
 \oplus\bigoplus_{a=1}^{c-1}H^{j-2a}(C_0;\mathbb Z).
 \tag{A.49b}
\]
The maps in this splitting are integral, with unit
diagonal on the exceptional-fibre basis. Thus it
also determines the primitive boundary maps and
the absence of torsion in the normalization
mapping-cylinder computation along the entire
finite algorithm, not just the Euler characteristic.
\end{proposition}


\begin{proof}
\begin{enumerate}[label=\arabic*.,leftmargin=*]
\item We construct the sheaf maps rather than infer a
splitting from fibre ranks. Let $E=\mathbb P_C(N_{C/Y})$
be the exceptional divisor, let $p:E\to C$, and put
$\xi=c_1(\mathcal O_E(1))$. Pullback from $Y$ gives
the degree-zero summand. For $1\leq a\leq c-1$,
pullback by $p$, multiplication by $\xi^{a-1}$,
and the exceptional Gysin morphism give a derived
map
\[
 i_*\mathbb Z_C[-2a]\longrightarrow
 R\pi_*\mathbb Z_{Y'},\qquad i:C\hookrightarrow Y.
\]
The Gysin map exists integrally for a closed complex
divisor, by its oriented tubular Thom class; neither
the centre nor the ambient space must be compact.
At a point of $C$, proper base change identifies
the stalk of the target with
$H^*(\mathbb P^{c-1};\mathbb Z)$.
Restriction of the exceptional normal line is
$\mathcal O_E(-1)$, so the displayed maps send
the generators to
\[
 1,-\xi,-\xi^2,\ldots,-\xi^{c-1}.
\]
This is an integral basis. Off $C$ the map is
the identity. It is therefore a stalkwise
quasi-isomorphism
\[
 \mathbb Z_Y\oplus
 \bigoplus_{a=1}^{c-1}i_*\mathbb Z_C[-2a]
 \xrightarrow{\ \simeq\ }R\pi_*\mathbb Z_{Y'}.
\]

\item Now restrict this actual sheaf morphism to $D$.
The central support upstairs is the whole inverse
image $\pi^{-1}D$, including when the centre is
horizontal. Proper base change applies to this
Cartesian square. The inverse-image functor for
integral sheaves is exact, and the inverse image
of $i_*\mathbb Z_C$ is
$(i_0)_*\mathbb Z_{C_0}$. This proves (A.49a).
Taking derived global sections proves (A.49b).
This calculation concerns the stratified map of
central spaces. It does not first calculate the
cohomology of either ambient total neighbourhood.

\item Here is its explicit relation to the requested
normalization mapping cylinders. Resolve
$\mathbb Z_{D'}$ by the alternating ordered-local-branch
complex: its first term is the constant sheaf on
the normalization, its next term is on the normalized
double-branch space, and so on. Its augmentation is
an integral quasi-isomorphism, because at a point
with $s$ branches it is the augmented simplex
complex on those branches. No division by $s!$ is
used; one uses alternating integral cochains, or
equivalently the orientation local system of each
unordered branch set. Proper pushforward of this
resolution computes the left side of (A.49a).
The same resolution for $D$, together with the
centre complex, computes its right side.

\item The local matrices can be read directly from the
blow-up charts in \ref{push:src:A.2.5}. If the centre lies on $k$
central branches among its $c$ defining coordinates,
the exceptional fibre is $\mathbb P^{c-1}$. The new
exceptional branch supplies that whole projective
space; the strict transform of old branch $i$
supplies its coordinate hyperplane $H_i$.
Intersections supply the corresponding coordinate
projective subspaces. In each even cohomological
degree the normalization differential is the
alternating restriction matrix with entries
$0,1,-1$. The cover consisting of the whole
projective space and its selected hyperplanes
has the elementary simplex contraction, leaving
one generator $\xi^a$ in each degree $2a$,
$0\leq a<c$. For a horizontal centre, every
central branch meeting its generic central
stratum supplies the whole exceptional fibre;
the same simplex contraction applies to these
copies. Higher boundary strata are treated by
the same coordinate restrictions. The global
Gysin maps just constructed glue the surviving
classes with unit coefficient, so this local
cancellation cannot hide an extension or an
integral index.

\item Starting from the explicitly computed reference
normalization matrices in \ref{push:src:A.4.8}, (A.49a) therefore
computes the normalization complex along every
recorded blow-up or inverse blow-up. A compact
centre has $C_0=C$. A horizontal toric surface
centre retracts to its rational-chain $C_0$;
a horizontal curve centre has $C_0$ a point.
Their cohomology is free and even. Consequently
the central computation preserves odd cohomology
and torsion, adds one to $b_2$ for every new
boundary ray, and adds one to $b_6$ exactly for
a new central ray. It gives the same
$(b_1,b_3,b_5)=(1,7,0)$,
$b_2=8+k+m$ and $b_6=k+m$ as \ref{push:src:A.4.8}.

\item Independently, stratification of $\pi_D$ into
$D\setminus C_0$ and $C_0$, with fibres a point
and $\mathbb P^{c-1}$, gives
\[
 \chi(D')-\chi(D)=(c-1)\chi(C_0).
\]
The absolute torus-stratum count of \ref{push:src:A.4.7} is
$\chi(D)=11k+t_4$. Together with the central
group calculation this gives again
$b_4=9k+t_4-2m-1$. Thus the two computations
agree integrally and in Euler number.

\item The program \path{code/ab_push_central_pushforward.py}
constructs every local alternating projective-subspace
matrix for $2\leq c\leq4$ and $1\leq k\leq c$,
as well as the possible horizontal cases. Its
Smith calculations check the unit cancellations
in every local dimension occurring in the
fourfolds. The sheaf construction proves their
global compatibility for arbitrary input fans.
\end{enumerate}
\end{proof}


\section{Local limits and the global component obstruction}\label{push:section:B}
\subsection{The finite root neighbourhoods}\label{push:section:B.1}

\begin{proposition}[counterexample and corrected]\label{push:src:B.1.1}\label{push:B11}
\StatementRecord{push:src:B.1.1}{proved}{Historical events: Claude Fable 5.2 (Anthropic), effort high, 16 November 2026: carried as background without fault found; targeted proof check. Claude Opus 5.5 (Anthropic), effort max, 16 November 2026: carried as background without fault found. GPT-6 Astra (OpenAI), ChatGPT Pro mode, effort Pro, S6\_09\_PAPERS, 17 November 2026: targeted proof check. Also checked in this writing. Pre-continuation writing events: GPT-6 Astra (OpenAI), ChatGPT Pro mode, effort Pro, S6\_10\_PAPERS, 18 November 2026: targeted proof check. Continuation metadata events: GPT-6 Astra (OpenAI), ChatGPT Pro mode, effort Pro, S6\_10\_PAPERS continuation, 19 November 2026: used and rechecked in this run; status-vocabulary reconciliation; metadata reconciliation only; no new mathematical check. \emph{Pre-continuation recorded proof scope.} Re-derived root-disc C-star bundle, twelvefold base change xy=r\textasciicircum{}12 and regular chain; proved generic nonproduct ruled components by the Picard class and finite-isogeny norm map.}{the hypotheses, notation and definitions stated here; \cref{fam:src:B.1.1}, \cref{fam:src:B.1.4}; \cref{push:src:A.1.2}; \cref{fam:src:E.2.1}}

The assertion that the filling is the product of the smooth elliptic
family with a twelve-gon Tate model is false on the whole root disc. There is, however, a proper flat holomorphic filling of
the marked root family over
$\Delta_{s_j}\times\Delta_r$, where $u=r^{12}$, with smooth
total space and reduced normal-crossings $r$-central divisor.
Its central surface over an ordinary value of $s_j$ is a cycle
of twelve ruled surfaces over the smooth elliptic curve
$E'_{j,s_j}$. Their normalizations are projective bundles
associated with the degree-zero line bundle represented by
$P_{s_j}-O_{s_j}$, with the appropriate pullbacks by translations.
They are not generically products with $\mathbb P^1$.
\end{proposition}


\begin{proof}
\begin{enumerate}[label=\arabic*.,leftmargin=*]
\item On the root disc the elliptic family $E'_j\to\Delta_{s_j}$ is
smooth. In the additive period coordinates $(\xi,\zeta_2)$,
put $W=\exp(-2\pi i\zeta_2)$. Quotienting by
$\ker\gamma\subset\Lambda$ gives a holomorphic
$\mathbb C^*$-bundle over this elliptic family. The period
$\widehat\gamma$ acts by
\[
 (\xi,W)\longmapsto
 (\xi+6\mu(s),a(s)uW),\qquad
 a(s)=\exp(-2\pi i\beta^{\rm part}(s)).
\]
This is the bundle and translation in \cref{fam:src:B.1.1}, \cref{fam:src:B.1.4}, with inverse fibre
coordinate; the period identities check its transition
functions directly. The factor $a$ is a holomorphic unit on
the entire root disc. Apply the compact-core plumbing of
\ref{push:src:A.1.2} relatively over a closed smaller root disc. Here the
smoothing section never vanishes, so the only local equation
at the conductor is $xy=u$ after a change of frame. This
gives a smooth total space before the base change.

\item After $u=r^{12}$, the local equation is $xy=r^{12}$,
times the two smooth coordinates on $E'_j$. Resolve its
$A_{11}$ surface singularity by its minimal toric resolution.
In lattice coordinates this inserts each primitive height-one
ray between the two boundary rays, at unit intervals. The
construction is canonical for changes $x\mapsto a x$,
$y\mapsto a^{-1}y$ by a holomorphic unit, so these local
resolutions glue. Its affine charts have the equation
$r=XY$, or $r=X$ at a smooth point of the central divisor,
times the elliptic and root coordinates. Thus the total space
is smooth, the divisor is reduced with normal crossings, and
the joint map to $(s_j,r)$ is flat. The resolution is proper,
and the plumbing was proper, so the joint map is proper.
The comparison over $r\ne0$ is the original period quotient,
since the resolution is an isomorphism there. There are
twelve components in the cycle because one application of
$\widehat\gamma$ shifts the unfolded component index by
twelve.

\item The ruled normalization of a component is obtained by adding
the zero and infinity sections to the indicated
$\mathbb C^*$-bundle. For a degree-zero line bundle $L_0$
on an elliptic curve, $\mathbb P(\mathcal O\oplus L_0)$
is a product over that curve only if $L_0$ is trivial.
Indeed an isomorphism would give two independent nowhere
vanishing quotient sections of a common twist of
$\mathcal O\oplus L_0$; degree zero forces that twist and
$L_0$ to be trivial. Equivalently, the two distinguished
disjoint sections have normal bundles $L_0,L_0^{-1}$.
A finite isogeny can trivialize $L_0$ only when $L_0$ is
torsion, because the norm of its pullback is a positive
tensor power of $L_0$.

\item The section $P$ of the rational elliptic surface is not
torsion, as is checked by its positive Mordell--Weil height
in Engel's Lemma 2.1. If its restriction to the root disc
were identically killed by a positive integer, analytic
continuation on the smooth part of $S\to B$ would give the
same identity globally. Each nonzero torsion equation
therefore has a discrete zero set. Outside a countable union
of such sets, $P_s-O_s$ is non-torsion. This disproves the
product assertion, even after an arbitrary fixed finite
isogeny of the elliptic factor. The special torsion
isogenies at $s_j=0$ used in \cref{fam:src:E.2.1} remain valid.
\end{enumerate}
\end{proof}

\begin{theorem}\label{push:src:B.1.2}\label{push:B12}
\StatementRecord{push:src:B.1.2}{proved}{Historical events: Claude Fable 5.2 (Anthropic), effort high, 16 November 2026: carried as background without fault found; targeted proof check. Claude Opus 5.5 (Anthropic), effort max, 16 November 2026: carried as background without fault found. GPT-6 Astra (OpenAI), ChatGPT Pro mode, effort Pro, S6\_09\_PAPERS, 17 November 2026: targeted proof check. Also checked in this writing. Pre-continuation writing events: GPT-6 Astra (OpenAI), ChatGPT Pro mode, effort Pro, S6\_10\_PAPERS, 18 November 2026: targeted proof check. Continuation metadata events: GPT-6 Astra (OpenAI), ChatGPT Pro mode, effort Pro, S6\_10\_PAPERS continuation, 19 November 2026: used and rechecked in this run; status-vocabulary reconciliation; metadata reconciliation only; no new mathematical check. \emph{Pre-continuation recorded proof scope.} Re-derived full affine root deck maps and nonconstant unit factors, component shifts4/-3, absence of fixed points at nodes, quotient flatness and t-multiplicity, and explicit marked product-cover identifications.}{the hypotheses, notation and definitions stated here; \cref{push:src:B.1.1}; \cref{fam:src:E.2.1}; \cref{fam:src:A.5.1}}

The twisted finite deck generators extend to the models of
\ref{push:src:B.1.1}. On the twelve component labels their actions are
$i\mapsto i+4$ at $j=1$ and $i\mapsto i-3$ at $j=2$,
with labels in $\mathbb Z/12$. Every nonidentity group
element acts freely, including at every node of the central
cycle. The quotient is a smooth complex fourfold, proper
and flat over $(t,r)=(s_j^{m_j},r)$, with a reduced
normal-crossings $r$-central divisor. Its fibre divisor
over $t=p_j$ has multiplicity exactly $m_j$.

The reduction over $s_j=r=0$ is the central surface of \cref{fam:src:E.2.1}:
a cycle of four ruled surfaces at $j=1$, and a cycle of
three at $j=2$. In the root presentation the normal
character to $s_j=0$ is $e^{-2\pi i/m_j}$ and the
character on $r$ is trivial.
\end{theorem}


\begin{proof}
\begin{enumerate}[label=\arabic*.,leftmargin=*]
\item We retain source and target arguments. Write primes for
evaluation at $g_js$, put $\zeta_j=e^{-2\pi i/m_j}$,
and set
\[
 R_1=\begin{pmatrix}-1/\tau&0\\(1-\mu)/\tau&1\end{pmatrix},
 \qquad
 R_2=\begin{pmatrix}1/\tau&0\\-\mu/\tau&1\end{pmatrix}.
\]
The identity $R_j\Pi_c=\Pi'_c A_j$ and the fixed-vector
identities $A_jv_j=v_j$ are checked symbolically in
\texttt{ab\_push\_finite\_decks.py}. In the coordinates
$(s,\xi,W,r)$ the two affine deck maps become
\begin{align*}
 s'&=\zeta_1s,\qquad
 \xi'=-\xi/\tau+(6\mu'+2\tau'-4)/3,\\
 W'&=r^4 U_1(s,\xi)W,\\
 U_1&=\exp\!\left[-2\pi i
  \left(\frac{1-\mu}{\tau}\xi+
        \frac{\beta^{{\rm part}\prime}+2\mu'}3\right)\right],
 \end{align*}
and
\begin{align*}
 s'&=\zeta_2s,\qquad
 \xi'=\xi/\tau+(-6\mu'-3\tau'+3)/4,\\
 W'&=r^{-3} U_2(s,\xi)W,\\
 U_2&=\exp\!\left[-2\pi i
  \left(-\frac{\mu}{\tau}\xi-
        \frac{\beta^{{\rm part}\prime}+3\mu'}4\right)\right].
 \end{align*}
Each $U_j$ is a nowhere vanishing holomorphic function on
the additive elliptic cover. Period covariance says
exactly that it defines the required bundle isomorphism
after quotient by $\ker\gamma$; it is not being treated
as a scalar function on the elliptic curve when it is not
one. Multiplication by a unit preserves every toric chart.
Multiplication by $r^a$ sends the chart of component index
$i$ to the chart of index $i+a$. Therefore the displayed
maps and their inverses extend holomorphically to the
resolved infinite chain and then to its twelve-component
quotient. The identities $g_1^3=g_2^4=1$ on the quotient
hold on the dense punctured family, and hence everywhere.

\item A nonidentity power has component shift $4,8$ in the first
case and $-3,-6,-9$ in the second. None is zero modulo
twelve. The same shifts act on the edge labels of the
cycle, so no node can be fixed either. At $s\ne0$ the
rotation of $s$ already excludes fixed points. The finite
quotient map is consequently a holomorphic local
isomorphism. Its source is smooth, so its target is smooth.
Properness follows from the proper root model and the
finite base map. Flatness can be checked after the
surjective local-isomorphism cover: there $s^{m_j}$ and
$r$ are a regular sequence, and their fibres have pure
dimension two. The local central equation remains
$r=XY$ or $r=X$. Finally $t=s^{m_j}$ shows that the
Cartier divisor over $p_j$ is exactly $m_j$ times its
reduction, including where this divisor meets a central
node. This also proves the asserted normal characters.

\item At $s=0$ use the explicitly marked isogenies
\[
 L_1(\xi,\eta)=(\eta,\xi-\eta/3),\qquad
 L_2(\xi,\eta)=(\eta,\xi-\eta/2).
\]
Their kernel generators are $(2/3,\rho)$ and $(1/2,i)$,
and the finite deck maps lift to
\[
 (\xi,\eta)\mapsto(\xi+\kappa_1/3,(\rho-1)\eta),
 \qquad
 (\xi,\eta)\mapsto(\xi-\kappa_2/4,-i\eta).
\]
In the Tate coordinate $v=e^{-2\pi i\xi}$ these are the
shifts $4,-3$ just computed. The kernel generators act
by nonzero translations on the smooth elliptic factor
$F_j$, hence freely even at the two fixed points of each
Tate component. Every other element with nonzero deck
power has a nonzero component shift. The full product
groups are therefore free on the product semistable
models. Their component orbits have cardinalities three
and four, giving precisely four and three ruled
components. The original lattice isogenies intertwine
the maps on the dense tori and their unit-interval fans,
so the quotient identifications extend across all nodes.
This verifies the particular comparison with \cref{fam:src:E.2.1} that
we use.

\item For the punctured family obtained by descending the
$c$-family, the holomorphic correction $\eta_j(s)$ in
\cref{fam:src:A.5.1} causes no additional local obstruction here. At the
shift $c\mapsto c+12$, the reference translations
$-4\Pi_c(\widehat u+\widehat w)$ and
$3\Pi_c(\widehat u+\widehat w)$ are periods.
The remaining correction is equivariant and independent
of $c$. Translation by $-c\eta_j(s)/12$ on the additive
cover conjugates that descent to the period descent.
It commutes with the finite deck action by equivariance.
Thus the constructed finite models fill the actual local
pullback by $w=r^6$ after a specified holomorphic
identification over $r\ne0$. This local gauge is not
asserted to extend across $r=0$ on overlaps with other
base neighbourhoods.
\end{enumerate}
\end{proof}

\subsection{The base-changed cusp fan}\label{push:section:B.2}

\begin{proposition}\label{push:src:B.2.1}\label{push:B21}
\StatementRecord{push:src:B.2.1}{proved}{Historical events: Claude Fable 5.2 (Anthropic), effort high, 16 November 2026: carried as background without fault found. GPT-6 Astra (OpenAI), ChatGPT Pro mode, effort Pro, S6\_09\_PAPERS, 17 November 2026: targeted proof check. Also checked in this writing. Pre-continuation writing events: GPT-6 Astra (OpenAI), ChatGPT Pro mode, effort Pro, S6\_10\_PAPERS, 18 November 2026: targeted proof check. Continuation metadata events: GPT-6 Astra (OpenAI), ChatGPT Pro mode, effort Pro, S6\_10\_PAPERS continuation, 19 November 2026: used and rechecked in this run; status-vocabulary reconciliation; metadata reconciliation only; no new mathematical check. \emph{Pre-continuation recorded proof scope.} Re-derived complete n+1 bridge classification, forced top rays, boundary periodicity, disjoint-monomial joint flatness, n component labels and trivial cusp component monodromy. Bounded fan replay remains separately tracked.}{the hypotheses, notation and definitions stated here; \cref{push:src:A.2.1}; \cref{push:src:A.1.2}; \cref{push:src:A.2.2}}

For the plain base change $u=r^n$ of the untwisted cusp
model, the smooth period-invariant refinements which
retain the prescribed bottom triangulation, are jointly
flat over $(q,r)$, and have reduced $r$-central divisor
are indexed by $k=0,\ldots,n$. In a fundamental prism
write its four bottom rays as $A,B,C,D$ as in \ref{push:src:A.2.1},
and put $\rho_i=(0,i,0,1)$ for $0\leq i\leq n$.
The maximal cones for the index $k$ are
\[
 ABC\rho_k,\quad BCD\rho_k,\quad
 AB\rho_i\rho_{i+1}\ (i<k),\quad
 CD\rho_i\rho_{i+1}\ (i\geq k).
\]
There are thirteen such models for $n=12$ and three
for $n=2$. Each has smooth total space and a reduced
normal-crossings central divisor with $n$ components
and no triple locus. Its component labels over the
punctured $q$-disc are globally $\mathbb Z/n$; in
particular their monodromy around $q=0$ is the identity.
\end{proposition}


\begin{proof}
\begin{enumerate}[label=\arabic*.,leftmargin=*]
\item The map of cocharacter lattices for the base change is
\[
 (y_1,y_2,h_q,h_r)\longmapsto
 (y_1,y_2,h_q,nh_r).
\]
The old top ray indexed by $l$ pulls back to
$(0,nl,0,1)$. Smoothness on the two-dimensional top
face forces insertion of every primitive ray
$\rho_{nl+i}$, $0<i<n$: a skipped interval of length
$d>1$ would retain lattice index $d$ on that face.
Joint flatness excludes any new ray with both base
heights positive, and the fixed bottom excludes new
rays of $h_r=0$. Reducedness permits only top rays
of $h_r=1$. These are consequently all the rays.

\item On the left face $AC\rho_0\rho_n$, a triangulation
with these collinear top vertices has a unique triangle
containing the bottom edge $AC$. Write its third
vertex as $\rho_a$. The triangles on either side
of that bridge are then forced. The analogous right
face has bridge $BD\rho_b$. The first period takes
$A,C$ to $B,D$ and fixes every top ray, so periodicity
requires $a=b=k$. The fixed bottom triangles $ABC$
and $BCD$ must then cone to $\rho_k$. The remaining
successive top intervals force exactly the displayed
cones, leaving no choice. Conversely those cones
have determinant $\pm1$ and common-face
intersections. Their normalized volumes sum to
$n+2$, the volume of the prism, so they cover it.
The determinant, intersection and volume assertions
for both requested exponents are certified in
\texttt{ab\_push\_basechange\_fan.py} using rational
arithmetic.

\item The period actions are
\[
 G_1:(y_1,y_2)\mapsto(y_1+h_q,y_2),\qquad
 G_2:(y_1,y_2)\mapsto(y_1,y_2+h_q+nh_r).
\]
Thus the choices propagate consistently from the
fundamental prism. The proper toric subdivisions
descend by the covering construction of \ref{push:src:A.1.2}.
Every ray has pure $q$-height or pure $r$-height, so
in a regular affine chart $q$ and $r$ are monomials
in disjoint sets of coordinates. This proves joint
flatness. All nonzero $r$-heights equal one, and
there are at most two top rays in any maximal cone;
this proves reduced normal crossings and absence of
triple points. Properness follows by descent from
the proper toric modification of the base-changed
proper cusp model.

\item The central components and all their intersection
maps are specified by the stars of the $\rho_i$.
Modulo the periods their indices are $i\bmod n$;
adjacent indices meet in the conductor elliptic
surface. Away from the cusp their normalizations
are the ruled bundles of the expanded plumbing.
At the unique bridge class the star has both bottom
triangles, and gives the nodal-curve blow-up and
small resolution described in \ref{push:src:A.2.2}, on the
corresponding ruled bundle. This star description
also fixes the gluing when the bridge is on an
inserted component. The endpoint choices $k=0,n$
are the canonical resolutions of the base changes
of the two original models. Intermediate values
are related to them by the successive side flops.

\item Finally, $G_1$ fixes each index and $G_2$ changes it
by $n$. The units in the cusp periods are
single-valued holomorphic functions of $q$ and do
not change a ray label. The toric charts themselves
are defined over $q$, not a root of $q$. Hence
these labels remain global on the punctured
$q$-disc, proving the monodromy assertion. Further
smooth periodic subdivisions with mixed rays give
infinitely many models flat over $r$ alone, as in
\ref{push:src:A.2.1}; they are excluded by the stated joint-flatness
condition.
\end{enumerate}
\end{proof}

\begin{corollary}[general plain-cusp normalization]\label{push:src:B.2.2}\label{push:B22}
\StatementRecord{push:src:B.2.2}{proved}{Historical events: Claude Fable 5.2 (Anthropic), effort high, 16 November 2026: carried as background without fault found. GPT-6 Astra (OpenAI), ChatGPT Pro mode, effort Pro, S6\_09\_PAPERS, 17 November 2026: targeted proof check. Also checked in this writing. Pre-continuation writing events: GPT-6 Astra (OpenAI), ChatGPT Pro mode, effort Pro, S6\_10\_PAPERS, 18 November 2026: targeted proof check. Continuation metadata events: GPT-6 Astra (OpenAI), ChatGPT Pro mode, effort Pro, S6\_10\_PAPERS continuation, 19 November 2026: used and rechecked in this run; status-vocabulary reconciliation; metadata reconciliation only; no new mathematical check. \emph{Pre-continuation recorded proof scope.} Re-derived the n=2/12 specialization through the actual parameter characteristic map, ray-height multiplicities and finite fan classification; supplied the complete short application proof.}{\cref{push:src:A.2.4}, \cref{push:src:A.2.5}, \cref{push:src:A.4.7}, \cref{push:src:A.4.8}, \cref{push:src:A.4.9}}

For $n=2$ and $n=12$, \ref{push:src:A.2.4}--\ref{push:src:A.2.5} give all
finite-orbit smooth invariant plain-cusp refinements,
their proper flat families, normalization charts,
finite blow-up and blow-down presentations,
multiplicities and all multiple loci. \ref{push:src:A.4.7}--\ref{push:src:A.4.9}
give their complete integral topology and nearby
Euler function. These statements concern the plain
untwisted cusp. They assert no compatibility with
the finite logarithmic-transform pieces of B.1.
\end{corollary}

\begin{proof}
\begin{enumerate}[label=\arabic*.,leftmargin=*]
\item The base maps $u=w^2$ and $u=r^{12}$ induce multiplication by
$2$ and $12$ on the integral parameter characteristic. Insert
$n=2$ and $n=12$ respectively into the periodic lattice action
and fan construction of \ref{push:src:A.2.4}. Its converse
recovers the invariant subdivision from each such plain-cusp
model, so the parametrization includes every stated output.

\item For each primitive ray $v$, the local parameter is a monomial
whose exponent on its divisor is its integral height $\ell(v)$.
Thus the central divisor is reduced exactly when every such
height equals one. Smoothness, flatness, properness, the
normalization charts and all intersection maps are those proved
in \ref{push:src:A.2.4}; \ref{push:src:A.2.5} gives the finite
blow-up and blow-down presentation. Finally substitute the
resulting orbit counts and maps into
\ref{push:src:A.4.7}, \ref{push:src:A.4.8}, \ref{push:src:A.4.9} for the
claimed topology. Every step is on the specified plain-cusp
charts, so the conclusion has exactly the scope in the statement.
\end{enumerate}
\end{proof}

\Uses{the hypotheses, notation and definitions stated here; \cref{push:src:A.2.4}; \cref{push:src:A.2.5}; \cref{push:src:A.4.7}; \cref{push:src:A.4.9}}


\begin{corollary}[unrestricted toroidal parametrization]\label{push:src:B.2.3}\label{push:B23}
\StatementRecord{push:src:B.2.3}{proved}{Historical events: Claude Fable 5.2 (Anthropic), effort high, 16 November 2026: carried as background without fault found. GPT-6 Astra (OpenAI), ChatGPT Pro mode, effort Pro, S6\_09\_PAPERS, 17 November 2026: targeted proof check. Also checked in this writing. Pre-continuation writing events: GPT-6 Astra (OpenAI), ChatGPT Pro mode, effort Pro, S6\_10\_PAPERS, 18 November 2026: targeted proof check. Continuation metadata events: GPT-6 Astra (OpenAI), ChatGPT Pro mode, effort Pro, S6\_10\_PAPERS continuation, 19 November 2026: used and rechecked in this run; status-vocabulary reconciliation; metadata reconciliation only; no new mathematical check. \emph{Pre-continuation recorded proof scope.} Re-derived the n=2/12 specialization through the actual parameter characteristic map, ray-height multiplicities and finite fan classification; supplied the complete short application proof.}{\cref{push:src:A.2.4}, \cref{push:src:A.2.5}, \cref{push:src:A.4.7}, \cref{push:src:A.4.8}, \cref{push:src:A.4.9}}

For $u=r^{12}$ and $u=w^2$, all finite-orbit smooth invariant
subdivisions of the plain cusp fan are obtained by \ref{push:src:A.2.4}
with $n=12$ and $n=2$, respectively. The construction there
gives their proper flat smooth-total-space families,
central multiplicities, normalization charts and all
multiple loci. Reducedness is exactly the restriction
$\ell=1$. This is a classification of the plain untwisted
cusp models; it does not supply the missing transition to
the finite logarithmic-transform models in B.3.
\end{corollary}

\begin{proof}
\begin{enumerate}[label=\arabic*.,leftmargin=*]
\item The base maps $u=w^2$ and $u=r^{12}$ induce multiplication by
$2$ and $12$ on the integral parameter characteristic. Insert
$n=2$ and $n=12$ respectively into the periodic lattice action
and fan construction of \ref{push:src:A.2.4}. Its converse
recovers the invariant subdivision from each such plain-cusp
model, so the parametrization includes every stated output.

\item For each primitive ray $v$, the local parameter is a monomial
whose exponent on its divisor is its integral height $\ell(v)$.
Thus the central divisor is reduced exactly when every such
height equals one. Smoothness, flatness, properness, the
normalization charts and all intersection maps are those proved
in \ref{push:src:A.2.4}; \ref{push:src:A.2.5} gives the finite
blow-up and blow-down presentation. Finally substitute the
resulting orbit counts and maps into
\ref{push:src:A.4.7}, \ref{push:src:A.4.8}, \ref{push:src:A.4.9} for the
claimed topology. Every step is on the specified plain-cusp
charts, so the conclusion has exactly the scope in the statement.
\end{enumerate}
\end{proof}

\Uses{the hypotheses, notation and definitions stated here; \cref{push:src:A.2.4}}


\subsection{The prescribed pieces do not glue}\label{push:section:B.3}

\begin{theorem}[counterexample to simultaneous gluing]\label{push:src:B.3.1}\label{push:B31}
\StatementRecord{push:src:B.3.1}{proved}{Historical events: Claude Fable 5.2 (Anthropic), effort high, 16 November 2026: carried as background without fault found. GPT-6 Astra (OpenAI), ChatGPT Pro mode, effort Pro, S6\_09\_PAPERS, 17 November 2026: targeted proof check. Also checked in this writing. Pre-continuation writing events: GPT-6 Astra (OpenAI), ChatGPT Pro mode, effort Pro, S6\_10\_PAPERS, 18 November 2026: targeted proof check. Continuation metadata events: GPT-6 Astra (OpenAI), ChatGPT Pro mode, effort Pro, S6\_10\_PAPERS continuation, 19 November 2026: used and rechecked in this run; status-vocabulary reconciliation; metadata reconciliation only; no new mathematical check. \emph{Pre-continuation recorded proof scope.} Re-derived the actual Stein component cover from proper normalized ruled fibres, dihedral transport, and peripheral contradiction residues7,1,11,5 modulo12.}{the hypotheses, notation and definitions stated here; \cref{push:section:B.1}; \cref{push:section:B.2}; \cref{push:src:B.1.2}; \cref{push:src:B.2.1}}

The finite models of \ref{push:section:B.1} and any of the plain
$u=r^{12}$ cusp models of \ref{push:section:B.2} cannot be glued to a
proper complex space over $B\times\Delta_r$ with
the prescribed central fibres on all three base
neighbourhoods. In particular the proposed simultaneous gluing is false. This conclusion does
not rule out a different cusp model of the same
punctured $S^6$ family.
\end{theorem}


\begin{proof}
\begin{enumerate}[label=\arabic*.,leftmargin=*]
\item If a further plain cusp refinement changes the top fan,
its central fibre over an ordinary $q\ne0$ has additional
components or different multiplicities and already cannot
agree with the twelve-cycle model of \ref{push:section:B.1} on the overlap.
We may therefore assume its top fan is the standard reduced
twelve-cycle. Additional mixed-height rays lie over $q=0$
and do not change the following argument over the punctured
base.

\item Suppose such a gluing exists and put
$B^\circ=B\setminus\{0,1,\infty\}$. Normalize
its reduced central fibre over $B^\circ$. In
each of the prescribed plumbing charts this is
a proper smooth family whose fibres are the
disjoint union of twelve connected ruled
surfaces. The charts and their transition
isomorphisms therefore give a locally constant
twelve-element set of connected components.
Equivalently, the Stein factor of this
normalization is an unramified degree-twelve
cover of $B^\circ$: the algebra of fibrewise
holomorphic functions is locally $\mathcal O^{12}$,
since each proper connected ruled surface has
only constant holomorphic functions. This
description constructs the cover and requires
no conjecture about components at the missing
three points.

\item The double curves give the adjacency graph of
a twelve-cycle. Thus its transition
permutations lie in the dihedral group
$D_{12}$, the automorphism group of that graph.
We record a peripheral permutation by its
deck-identification map. Holonomy defined by
lifting a clockwise path and returning its
endpoint to the initial root slice uses the
inverse permutations throughout; the relation
and the following contradiction are unchanged.
At $0$ the monodromy is a rotation by four
steps; at $1$ it is a rotation by minus three
steps, by the explicit root deck maps in
\ref{push:src:B.1.2}. Relabelling a cycle can reverse either
orientation, but changes these only to
$\pm4$ and $\pm3$. At the cusp the
monodromy is the identity by \ref{push:src:B.2.1}.

\item The peripheral relation in
$\pi_1(B^\circ)$ is $g_1g_2g_0=1$ with
the consistently chosen clockwise loops.
In any common labelling, the product of the
two finite rotations would therefore have
to be the identity. Its possible rotation
numbers are
\[
 4+3,\quad4-3,\quad-4+3,\quad-4-3
 \equiv7,1,11,5\pmod {12}.
\]
None is zero. This contradicts the peripheral
relation. The four residues and the finite
component cycles are checked exactly by
\texttt{ab\_push\_finite\_decks.py}.
\end{enumerate}
\end{proof}

\begin{proposition}[the marked genus-three component cover]\label{push:src:B.3.2}\label{push:B32}
\StatementRecord{push:src:B.3.2}{proved}{Historical events: Claude Fable 5.2 (Anthropic), effort high, 16 November 2026: carried as background without fault found. GPT-6 Astra (OpenAI), ChatGPT Pro mode, effort Pro, S6\_09\_PAPERS, 17 November 2026: targeted proof check. Also checked in this writing. Pre-continuation writing events: GPT-6 Astra (OpenAI), ChatGPT Pro mode, effort Pro, S6\_10\_PAPERS, 18 November 2026: targeted proof check. Continuation metadata events: GPT-6 Astra (OpenAI), ChatGPT Pro mode, effort Pro, S6\_10\_PAPERS continuation, 19 November 2026: used and rechecked in this run; status-vocabulary reconciliation; metadata reconciliation only; no new mathematical check. \emph{Pre-continuation recorded proof scope.} Re-derived the connected cyclic representation4,-3,-1, its Kummer function-field cover, ramification3/4/12 and Riemann-Hurwitz genus3. Receiving version will state this actual cover directly without a hypothetical threefold construction. Receiving theorem now states the actual marked cyclic cover directly; full explicit covering morphism, normal local branches, peripheral action and genus-three proof written. The hypothetical-central-threefold wording is superseded by the actual-cover formulation.}{\cref{push:src:B.3.1}; the explicit function-field equation and its local normalizations}

Let $B^\circ=\mathbb P^1\setminus\{0,1,\infty\}$ and use the
clockwise deck-identification convention of \ref{push:src:B.3.1}.
The peripheral rotations $4,-3,-1$ in $\mathbb Z/12$ define a
connected degree-twelve cover of $B^\circ$. Its compact
normalization is the smooth genus-three curve obtained from
\[
 y^{12}=\frac{t^4}{(t-1)^3}.
\]
The ramification indices over $0,1,\infty$ are $3,4,12$,
respectively. The peripheral relation determines the last rotation
from the first two, so this is the marked cyclic component cover
used in the partial filling of \ref{push:src:B.4.19}.
\end{proposition}

\begin{proof}
\begin{enumerate}[label=\arabic*.,leftmargin=*]
\item The presentation $g_1g_2g_0=1$ sends the third peripheral generator
to $-4+3=-1$. The first two images generate $\mathbb Z/12$, so
the resulting action on the twelve sheets is transitive. This
constructs the connected cover by its monodromy representation.

\item The rational function $t^4/(t-1)^3$ has orders $4,-3,-1$ at the
three punctures. Its twelfth-root equation has the same marked
deck representation. The covering-space identification is the map
which sends a chosen twelfth root at one base point to its chosen
sheet and continues both along the same paths. Its equality of
monodromy makes that map single-valued; its local analytic root
coordinates make it holomorphic. Normalization extends this
identification to the compact covers.

\item Locally $y^{12}=z^a u(z)$ has $\gcd(12,a)$ branches of ramification
index $12/\gcd(12,a)$, after absorbing the holomorphic unit in a
local twelfth root. Thus the three contributions to ramification
are $4(3-1)$, $3(4-1)$, and $12-1$. Their sum is $28$.
Riemann--Hurwitz gives $2-2g=2\cdot12-28=-4$, so $g=3$.
Reversing the deck convention is realized by $y\mapsto y^{-1}$,
with the reciprocal defining equation; the compact curve and the
ramification calculation are unchanged.
\end{enumerate}
\end{proof}


\begin{proposition}[obstruction to a specified repair]\label{push:src:B.3.3}\label{push:B33}
\StatementRecord{push:src:B.3.3}{proved}{Historical events: Claude Fable 5.2 (Anthropic), effort high, 16 November 2026: carried as background without fault found. GPT-6 Astra (OpenAI), ChatGPT Pro mode, effort Pro, S6\_09\_PAPERS, 17 November 2026: targeted proof check. Also checked in this writing. Pre-continuation writing events: GPT-6 Astra (OpenAI), ChatGPT Pro mode, effort Pro, S6\_10\_PAPERS, 18 November 2026: targeted proof check. Continuation metadata events: GPT-6 Astra (OpenAI), ChatGPT Pro mode, effort Pro, S6\_10\_PAPERS continuation, 19 November 2026: used and rechecked in this run; status-vocabulary reconciliation; metadata reconciliation only; no new mathematical check. \emph{Pre-continuation recorded proof scope.} Re-derived component translation order d/gcd(d,a,b), generic normal isotropy and resulting multiplicity d/h0, including the multiplicity-at-least12 obstruction for a primitive shift.}{the hypotheses, notation and definitions stated here; \cref{push:src:B.3.1}; \cref{push:src:A.1.2}}

A cyclic cusp-root construction using a
fractional translation by the two cusp
periods cannot repair \ref{push:src:B.3.1} while retaining
the reduced Mumford fibre for $r\ne0$.
More precisely, let $q=s^d$, $u=r^n$,
and suppose a cyclic deck generator has
translation $d^{-1}(aG_1+bG_2)$ in the
additive notation for the two periods.
Assume its $r$-component shift is integral.
Then its $r$-component shift is
$nb/d\bmod n$, whereas the multiplicity
of every generic quotient $q$-component
is a positive multiple of
\[
 \frac{d}{\gcd(d,a,b)}.
\]
A nonzero $r$-component shift therefore
forces a nonreduced $q$-fibre. In
particular a shift $-1\bmod12$ in this
class of constructions forces multiplicity
at least twelve.
\end{proposition}


\begin{proof}
\begin{enumerate}[label=\arabic*.,leftmargin=*]
\item In the coordinates of \ref{push:src:A.1.2} the two
periods have $q$-valuation vectors
$(1,0)$ and $(0,1)$; their $r$-valuation
vectors are $(0,0)$ and $(0,n)$.
After $q=s^d$, the period lattice on
the $q$-central component indices is
$d\mathbb Z^2$. A regular toroidal
resolution has indices in
$(\mathbb Z/d)^2$. The proposed deck
generator translates them by $(a,b)$.
The order of that translation is
$d/\gcd(d,a,b)$, so the subgroup
preserving a component has order
$h=\gcd(d,a,b)$. Its subgroup fixing
a generic point has some order
$h_0$ dividing $h$; unit factors may
make $h_0$ smaller. That isotropy
acts on the normal coordinate $s$
by a primitive $h_0$th root of
unity. The quotient normal
coordinate is $s^{h_0}$, and
$q=(s^{h_0})^{d/h_0}$ has
multiplicity $d/h_0$, a positive
multiple of $d/h$. Unit factors
cannot decrease this lower bound.

\item Reducedness would require $h_0=d$,
and therefore $h=d$,
so both $a,b$ are divisible by $d$.
Then $nb/d$ is zero modulo $n$.
If $nb/d$ is relatively prime to
$n$, its order in $\mathbb Z/n$
is $n$, which divides the order
$d/h$ of the full component
translation; this gives the final
lower bound. The residue calculations
and examples for the requested
exponents are independently recorded
in \texttt{ab\_push\_cusp\_root\_repair.py}.
This rules out the specified repair,
not a cusp model with a different
geometry or a different comparison
to the punctured family.
\end{enumerate}
\end{proof}

\begin{theorem}[split-normalization trace obstruction]\label{push:src:B.3.4}\label{push:B34}
\StatementRecord{push:src:B.3.4}{proved}{Historical events: Claude Fable 5.2 (Anthropic), effort high, 16 November 2026: carried as background without fault found. GPT-6 Astra (OpenAI), ChatGPT Pro mode, effort Pro, S6\_09\_PAPERS, 17 November 2026: targeted proof check. Also checked in this writing. Pre-continuation writing events: GPT-6 Astra (OpenAI), ChatGPT Pro mode, effort Pro, S6\_10\_PAPERS, 18 November 2026: targeted proof check. Continuation metadata events: GPT-6 Astra (OpenAI), ChatGPT Pro mode, effort Pro, S6\_10\_PAPERS continuation, 19 November 2026: used and rechecked in this run; status-vocabulary reconciliation; metadata reconciliation only; no new mathematical check. \emph{Pre-continuation recorded proof scope.} Re-derived explicit b\textasciicircum{}3=a\textasciicircum{}4+1 cover and Weierstrass scaling, all local section intersections/component4 expansions, height Gram matrix, rank3 and index dividing4, and integral-trace obstruction12ZP via half-integral congruences. Exact Shioda and localization inputs tracked separately. Closed Shioda primary-source height/rank/trace input gates by reading Theorems1.3/4.1/8.4/8.6,table8.16,Proposition8.12 in the author-uploaded 37-page edition; exact edition-specific pages recorded.}{the hypotheses, notation and definitions stated here; \cref{push:src:B.3.2}; \cite[Theorems 8.4 and 8.6, table (8.16)]{Shioda90}; \cite[Theorem 4.1]{Shioda90}; \cite[Theorem 1.3]{Shioda90}}

Let $C$ be the genus-three component cover
in \ref{push:src:B.3.2}, and let $\mathscr E\to C$ be
the relatively minimal elliptic surface
obtained by pulling back $S\to B$.
Write $\sigma$ for the order-twelve deck
transformation, with its natural lift to
$\mathscr E$. Then
\[
 \operatorname{Tr}_{\sigma}
   \mathscr E(\mathbb C(C))
 :=\left\{\sum_{i=0}^{11}\sigma^i Q:
                 Q\in\mathscr E(\mathbb C(C))\right\}
 =12\mathbb Z P.
\]
In particular there is no section $Q$
for which $(t_Q\circ\sigma)^{12}=t_{6P}$.
Thus the proposed repair that takes the
normalization over $C$ to be ruled over
this split elliptic surface, with adjacent
conductor map $t_Q\circ\sigma$, is
impossible.

This conclusion is restricted to the
stated split normalization. A ruled
normalization over a nontrivial genus-one
torsor need not admit a section or such
an expression for its affine deck lift;
it is not excluded by this theorem.
\end{theorem}


\begin{proof}
\begin{enumerate}[label=\arabic*.,leftmargin=*]
\item There is a convenient model of the cover
and its elliptic surface:
\[
 C:\ b^3=a^4+1,\qquad
 t=b^3,\quad y=b/a,\qquad
 \mathscr E:\ Y^2=X^3-3bX+2a^2.
\]
Indeed $y^{12}=t^4/(t-1)^3$.
The substitution from the original
Weierstrass coordinates is
$x=a^2b^4X$, $y_{\rm old}=a^3b^6Y$.
The affine discriminant of the new
equation is $1728$; it has no singular
fibre on the affine curve. The curve
$C$ has one point at infinity, with
pole orders $3,4$ for $a,b$.
Put $z=a/b$ and $c=z^3a=z^4b$ there.
Then $z$ is a uniformizer and
\[
 c^3(1-c)=z^{12},\qquad c(0)=1.
\]
In the integral Weierstrass coordinates
$X_0=z^2X$, $Y_0=z^3Y$, the equation is
\[
 Y_0^2=X_0^3-3cX_0+2c^2.
\]
Its discriminant has order twelve and
its $c_4$ is a unit. The unique singular
fibre of the minimal regular surface
is therefore $I_{12}$. The fundamental
line bundle is $\mathcal O_C(\infty)$,
so $\chi(\mathcal O_{\mathscr E})=1$.
This is a projective elliptic surface.

\item Let $\omega=e^{2\pi i/3}$ and
$\lambda=e^{2\pi i/12}$. The deck lift is
\[
 \sigma(a,b,X,Y)=(ia,\omega b,\lambda^2X,\lambda^3Y).
\]
It fixes the section $P=(0,-\sqrt2a)$.
Two further sections are
\[
 R=\left(\frac{a^2}{b-1},
                 \frac{\sqrt3a}{b-1}\right),
 \qquad
 R'=\sigma R=
   \left(\frac{a^2}{b-\omega},
                 \frac{\sqrt3\omega a}{b-\omega}\right).
\]
Direct substitution verifies their
equations. The three orbit points
$R,\sigma R,\sigma^2R$ lie on the line
\[
 Y=\frac{\sqrt3b}{a}X-\sqrt3a.
\]
Thus $R+\sigma R+\sigma^2R=O$ by the
elliptic chord law, and $\sigma^3R=R$.
The ansatz that produced these sections,
its other solutions, and this line
identity are checked symbolically in
\texttt{ab\_push\_trace\_sections.py}.

\item We give the intersection computation
needed for the height lattice; this
also fixes the component indices used
below. The only pole of $R$ on the affine
curve is at $(a,b)=(0,1)$. There $a$ is
a uniformizer and $b-1=a^4/3+O(a^8)$,
so $X=3a^{-2}+O(a^2)$ and
$Y=3\sqrt3a^{-3}+O(a)$.
Consequently $R\cdot O=1$.
The same computation at $(0,\omega)$
gives $R'\cdot O=1$, while $P\cdot O=0$.
At the other two points with $a=0$,
$P$ and $R$ meet simply: the fibre is
smooth at $(X,Y)=(0,0)$, and the
difference of their $Y$ coordinates
has a nonzero linear term in $a$.
Thus $P\cdot R=2$, and likewise
$P\cdot R'=2$. The sections $R,R'$
meet simply at $(a,b)=(0,\omega^2)$.
They cannot meet elsewhere on the
affine curve, since their $X$
coordinates are distinct for $a\ne0$.

\item To check infinity, for the section
$\sigma^jR$ put $k=\omega^j$. Its
integral coordinates are
\[
 X_0=\frac{c^2}{c-kz^4},\qquad
 Y_0=\frac{\sqrt3kc z^4}{c-kz^4}.
\]
At the node use analytic coordinates
\[
 U=Y_0-(X_0-1)\sqrt{X_0+2},\qquad
 V=Y_0+(X_0-1)\sqrt{X_0+2}.
\]
They satisfy
\[
 UV=(1-c)\{3X_0-2(c+1)\},
\]
whose right side is a unit times $z^{12}$.
The section expansions are
\[
 U=-\frac{\sqrt3k^2}{6}z^8+O(z^{12}),
 \qquad V=2\sqrt3kz^4+O(z^8).
\]
Hence $R,R'$ both meet component $4$
of the resolved $I_{12}$ cycle, with
distinct coordinates on that component.
Reversing the cycle numbering replaces
both indices by $8$ and has no effect
on the argument. The section $P$ meets
the identity component, at the smooth
point $(X_0,Y_0)=(0,-\sqrt2)$.
Thus no further intersections occur at
infinity, and $R\cdot R'=1$.

\item The height pairing formula and its
positivity are those of
\cite[Theorems 8.4 and 8.6, table (8.16)]{Shioda90}.
For an $I_{12}$ fibre its correction is
$\operatorname{contr}(i,j)=\min(i,j)-ij/12$.
The preceding intersections give, in
the ordered basis $P,R,R'$,
\[
 G=\begin{pmatrix}
 2&0&0\\
 0&4/3&-2/3\\
 0&-2/3&4/3
 \end{pmatrix},\qquad \det G=8/3.
\]
In particular these are three independent
sections. The Leray sequence and
$R^1\pi_*\mathcal O_{\mathscr E}
=\mathcal O_C(-\infty)$ show that its
irregularity equals $g(C)=3$; the
right-hand sheaf has no global section.
This also agrees with the Picard-variety
statement in \cite[Theorem 4.1]{Shioda90}.
The equalities $\chi=1$ and
$e(\mathscr E)=12$ give $p_g=3$,
$b_2=22$ and $h^{1,1}=16$.
The zero section, a fibre and the eleven
nonidentity $I_{12}$ components generate
a rank-thirteen trivial lattice.
The identification of sections with
the N\'eron--Severi quotient
\cite[Theorem 1.3]{Shioda90}, together
with $\rho\le h^{1,1}$, now shows that
the Mordell--Weil rank is exactly three
and $\rho=16$.

\item There is no nonzero torsion section.
Indeed its height would be zero, so
for some $0\le j<12$ the height formula
would give
\[
 0=2+2(Q\cdot O)-\frac{j(12-j)}{12}.
\]
The correction is at most three and
$Q\cdot O$ is a nonnegative integer.
This forces $Q\cdot O=0$ and
$j(12-j)=24$, which has no integral
solution.

\item Let $L=\mathbb ZP+\mathbb ZR+\mathbb ZR'$.
The integral N\'eron--Severi sublattice
generated by the trivial lattice and
these three sections has determinant
$-32$. This follows either by eliminating
the trivial lattice with the displayed
height matrix, or from the integral
intersection matrix recorded in
\texttt{ab\_push\_trace\_obstruction.py}.
All its entries have just been computed:
sections have square $-1$, their mutual
intersections are $2,2,1$, and $R,R'$
meet precisely component $4$.
If $d=[\mathscr E(\mathbb C(C)):L]$,
the section--N\'eron--Severi quotient
identification gives the same index
for this sublattice in
$\operatorname{NS}(\mathscr E)$.
Its integral determinant implies
$d^2\mid32$, hence $d\in\{1,2,4\}$.
In particular $4Q\in L$ for every section.

\item In the rational basis $P,R,R'$ the
deck action and its trace are
\[
 M=\begin{pmatrix}1&0&0\\0&0&-1\\0&1&-1\end{pmatrix},
 \qquad \sum_{i=0}^{11}M^i=
       \operatorname{diag}(12,0,0).
\]
Suppose first that the trace of $Q$ is
$6P$. By the index bound we may write
\[
 Q=\frac12P+\frac14(uR+vR'),\qquad u,v\in\mathbb Z.
\]
Let $j\in\mathbb Z/12$ be its component
index at infinity. Since both $R,R'$
meet component four, the height formula
modulo integers gives
\[
 \frac{2u-v}{6}\equiv\frac j3\pmod{\mathbb Z},
 \qquad
 \frac{-u+2v}{6}\equiv\frac j3\pmod{\mathbb Z}.
\]
Equivalently $2u-v-2j$ and $-u+2v-2j$
are divisible by six. Reducing modulo
two shows that both $u,v$ are even.
Write $u=2u'$, $v=2v'$. The component
map is a group homomorphism, so
\[
 2j\equiv4(u'+v')\pmod{12}.
\]
Therefore $j$ is even. The geometric
height formula then lies in
$\tfrac13\mathbb Z$. In contrast the
height from $G$ is
\[
 \langle Q,Q\rangle
 =\frac12+\frac{u'^2-u'v'+v'^2}{3},
\]
which lies in the disjoint coset
$\tfrac12+\tfrac13\mathbb Z$.
This is a contradiction.

\item The same argument determines the full
trace image. If $Q=\alpha P+\beta R+\gamma R'$
in the rational span, pairing with the
narrow section $P$ gives
$2\alpha=\langle Q,P\rangle\in\mathbb Z$.
Every odd half-integral value of $\alpha$
leads to the identical congruences and
height contradiction, with $1/2$
replaced by an odd square divided by
two. Thus $\alpha\in\mathbb Z$ and
the trace is in $12\mathbb ZP$.
The section $P$ realizes the generator.
The exact matrices, local series,
torsion test, and all relevant residue
classes are verified in
\texttt{ab\_push\_trace\_obstruction.py}.
Finally the identity
$(t_Q\circ\sigma)^{12}=t_{\operatorname{Tr}_\sigma Q}$
proves the claimed obstruction, with
the scope stated in the theorem.
\end{enumerate}
\end{proof}

\begin{corollary}[exact Mordell--Weil lattice]\label{push:src:B.3.5}\label{push:B35}
\StatementRecord{push:src:B.3.5}{proved}{Historical events: Claude Fable 5.2 (Anthropic), effort high, 16 November 2026: carried as background without fault found. GPT-6 Astra (OpenAI), ChatGPT Pro mode, effort Pro, S6\_09\_PAPERS, 17 November 2026: targeted proof check. Also checked in this writing. Pre-continuation writing events: GPT-6 Astra (OpenAI), ChatGPT Pro mode, effort Pro, S6\_10\_PAPERS, 18 November 2026: targeted proof check. Continuation metadata events: GPT-6 Astra (OpenAI), ChatGPT Pro mode, effort Pro, S6\_10\_PAPERS continuation, 19 November 2026: used and rechecked in this run; status-vocabulary reconciliation; metadata reconciliation only; no new mathematical check. \emph{Pre-continuation recorded proof scope.} Re-derived pole semigroup3i+4j, the only zero-intersection sections plus/minusP, height lower bound1 for all others, and exclusion of every index2/4 overlattice by the three height1/3 half-classes. Closed Shioda primary height-input gate; exact positive height formula and local correction table opened in author-uploaded edition.}{the hypotheses, notation and definitions stated here; \cref{push:src:B.3.4}}

In the notation of \ref{push:src:B.3.4},
\[
 \mathscr E(\mathbb C(C))
   =\mathbb ZP\oplus\mathbb ZR\oplus\mathbb ZR',
 \qquad \det\operatorname{NS}(\mathscr E)=-32.
\]
\end{corollary}


\begin{proof}
\begin{enumerate}[label=\arabic*.,leftmargin=*]
\item The affine coordinate ring of $C$ is
$\mathbb C[a,b]/(b^3-a^4-1)$, with
unique representatives of degree at most
two in $b$. The pole orders $3i+4j$ of
$a^ib^j$, $0\le j\le2$, are distinct,
so cancellation cannot create functions
with pole order one or two at infinity.
Consequently
\[
 H^0(C,\mathcal O_C(2\infty))=\mathbb C,
 \qquad
 H^0(C,\mathcal O_C(3\infty))
       =\mathbb C\oplus\mathbb Ca.
\]
A section disjoint from $O$ therefore
has $X=x_0$, $Y=y_0+y_1a$ for constants
$x_0,y_0,y_1$. Substitution in the
Weierstrass equation forces $x_0=y_0=0$
and $y_1^2=2$. Thus the only nonzero
sections with intersection zero with
$O$ are $P,-P$.
Every other nonzero section has
intersection at least one with $O$,
and hence height at least $4-3=1$.

\item \ref{push:src:B.3.4} proves that every section has
integral coefficient along $P$ in the
rational height decomposition. The
remaining lattice can have index only
one, two or four over $\mathbb ZR+\mathbb ZR'$.
Any proper such overlattice contains a
nonzero element of order two modulo
this sublattice. After subtracting an
integral combination of $R,R'$ it is
one of $R/2$, $R'/2$, $(R+R')/2$.
Each has height $1/3$, contradicting
the preceding lower bound. Thus the
index is one. The determinant already
computed in \ref{push:src:B.3.4} is the full
N\'eron--Severi determinant.
The pole-basis calculation, coefficient
equations and three half-class heights
are included in the exact certificate
\texttt{ab\_push\_trace\_obstruction.py}.
\end{enumerate}
\end{proof}

\begin{proposition}[an unramified torsor alternative]\label{push:src:B.3.6}\label{push:B36}
\StatementRecord{push:src:B.3.6}{proved}{Historical events: Claude Fable 5.2 (Anthropic), effort high, 16 November 2026: carried as background without fault found. GPT-6 Astra (OpenAI), ChatGPT Pro mode, effort Pro, S6\_09\_PAPERS, 17 November 2026: targeted proof check. Also checked in this writing. Pre-continuation writing events: GPT-6 Astra (OpenAI), ChatGPT Pro mode, effort Pro, S6\_10\_PAPERS, 18 November 2026: targeted proof check. Continuation metadata events: GPT-6 Astra (OpenAI), ChatGPT Pro mode, effort Pro, S6\_10\_PAPERS continuation, 19 November 2026: used and rechecked in this run; status-vocabulary reconciliation; metadata reconciliation only; no new mathematical check. \emph{Pre-continuation recorded proof scope.} Read full proof and re-derived the cyclic norm quotient, local norm of6P by henselian division, subtraction of localized torsor class, and carry cocycle producing actual affine A\textasciicircum{}12=t6P. Closed HS53/CDL24/Xiao91 primary input gates: five-term transgression pp129-130; H2 field vanishing and surjective localization CDLp445; Xiao exact sequence4p21 and translating-twist theorem3.4p14 opened. Scope includes the actual nontrivial Jacobian hypotheses.}{the hypotheses, notation and definitions stated here; \cref{push:src:B.3.4}; \cite[Chapter III, Section 4, Theorem 2
and the following remark, pp.129--130]{HS53}; \cite[Lemma 4.7.8]{CDL24}; \cite{CDL24}; \cite[Section 4, equation (4)]{Xiao91}; \cite[Sections 3--4]{Xiao91}; \cref{push:src:B.4.1}}

Put $K=\mathbb C(t)$, $L=\mathbb C(C)$,
$G=\operatorname{Gal}(L/K)=\langle\sigma\rangle$,
and let $E/K$ be the generic elliptic
curve of $S\to B$. There is a nonzero
class
\[
 \xi\in\operatorname{Sha}(E_L/L)^G
 \quad\hbox{whose transgression in}\quad
 H^2(G,E(L))=\mathbb ZP/12\mathbb ZP
 \quad\hbox{is }6P.
\]
Here $\operatorname{Sha}$ denotes the
subgroup of genus-one torsors trivial
over the henselian local fields of all
points of $C$.
In particular the trace obstruction of
\ref{push:src:B.3.4} does not obstruct all genus-one
torsors. Such a class has a relatively
minimal projective genus-one surface
$V\to C$, locally isomorphic to
$\mathscr E\to C$, and an affine lift
$A:V\to V$ of $\sigma$ with
\[
 A^{12}=t_{6P}.
\]
This is an existence statement for a
possible normalization base, not a
construction of a threefold central
fibre or of its smoothing.
\end{proposition}


\begin{proof}
\begin{enumerate}[label=\arabic*.,leftmargin=*]
\item We use Galois cohomology with its
continuous cochains. The five-term
Hochschild--Serre sequence gives
\[
 H^1(K,E)\longrightarrow H^1(L,E)^G
 \xrightarrow{\ d_2\ }H^2(G,E(L))
 \longrightarrow H^2(K,E).
\]
The transgression sequence is
\cite[Chapter III, Section 4, Theorem 2
and the following remark, pp.129--130]{HS53};
its continuous form follows by passing
to finite Galois quotients for the
finitely many values of each continuous
cochain. For this complex function field,
$H^2(K,E)=0$, as in
\cite[Lemma 4.7.8]{CDL24}.
One can also see this from
$\operatorname{cd}K=1$: the torsion
submodule of $E(\overline K)$ has no
second cohomology, and its quotient is
a rational vector space, whose positive
continuous cohomology vanishes by
averaging on finite quotients.

\item Cyclic group cohomology gives
$H^2(G,E(L))=E(K)/\operatorname{Tr}_\sigma E(L)$.
By \ref{push:src:B.3.4} this is $\mathbb Z/12$, with
generator the class of $P$. Hence there
is a $G$-invariant torsor class
$\xi_0\in H^1(L,E)$ with $d_2\xi_0=6P$.
We next remove all its local ramification
without changing this transgression.

\item At a place $v$ of $B$ choose a point
$w$ of $C$ above it and its decomposition
subgroup $G_w$. The only nontrivial
decomposition groups have orders
$3,4,12$, at $0,1,\infty$ respectively.
The component groups of the original
fibres are $\mathbb Z/3$, $\mathbb Z/2$
and zero. Thus $6P$ belongs to the
identity component of the local N\'eron
group at every ramified place.
Multiplication by the ramification
index is surjective on that identity
group over the henselian local field:
on the residue group it is surjective
on $\mathbb G_a(\mathbb C)$ or
$\mathbb G_m(\mathbb C)$, and its
derivative is invertible; Hensel's
lemma lifts a division point.
Choose such a division point over
$K_v^h$. Viewed over $L_w^h$, it is
invariant, so its local norm is $6P$.
The restriction of the transgression
class to $H^2(G_w,E(L_w^h))$ is
therefore zero. For unramified places
the decomposition group is trivial,
so the same assertion is automatic.

\item Naturality and exactness of the local
five-term sequences now show that
$(\xi_0)_w$ is the restriction of some
$\eta_v\in H^1(K_v^h,E)$.
Only finitely many of these classes
need be nonzero. The original elliptic
surface is nontrivial, since
$j=1728t$. Ogg--Shafarevich localization
is therefore surjective:
\[
 H^1(K,E)\longrightarrow
 \bigoplus_{v\in B}H^1(K_v^h,E)
 \longrightarrow0.
\]
This is the combination of exact
sequence (4.7.4) and Theorem 4.7.9
of \cite{CDL24}; the same surjectivity
and its geometric fibre-twist
interpretation are given in
\cite[Section 4, equation (4)]{Xiao91}.
Take a global $\eta$ with these local
values and put
$\xi=\xi_0-\operatorname{res}_{L/K}\eta$.
It is locally trivial at every point
of $C$, remains $G$-invariant, and has
$d_2\xi=6P$. It is nonzero because its
transgression is nonzero.

\item The standard interpretation of the
Weil--Ch\^atelet group represents
$\xi$ by a smooth projective genus-one
curve over $L$; taking its relatively
minimal regular model over $C$ gives
$V$. Local triviality identifies each
henselian model with that of $E_L$.
Over $\mathbb C$ the corresponding
\'etale identifications yield
holomorphic identifications over small
base discs. Equivalently, $V$ is obtained
by gluing copies of $\mathscr E$ over
these discs by translations by local
sections. These translations extend
over the regular $I_{12}$ model.
Thus $V$ is locally isomorphic to
$\mathscr E$, has smooth total surface,
no multiple fibres, and the same single
$I_{12}$ fibre. Properness and flatness
hold locally over the base and hence
globally. This use of torsors and local
twists is the one in
\cite[Sections 3--4]{Xiao91}; it does
not assert an explicit equation for
the selected torsor.

\item The $G$-invariance of $\xi$ gives an
affine isomorphism of its generic fibre
covering $\sigma$. It extends to the
relatively minimal regular surface:
in this particular situation the
extension can be checked on a common
resolution. If an exceptional
$(-1)$-curve for one projection is not
contracted by the other, its image is
a vertical rational curve and has
self-intersection at least $-1$,
since blowing down only increases
self-intersection. On either locally
trivial model all vertical rational
curves are components of the $I_{12}$
fibre, of square $-2$. Thus every such
$(-1)$-curve is contracted by both
projections and can be removed from
the common resolution. Repeating this
proves that the map is an isomorphism.
The same argument gives the local
comparison with the Jacobian model
used above. Transporting the base by
$\sigma$ changes none of these facts.
Write $A_0$ for this lift. Its twelfth
power is translation by a section
$N\in E(L)^G=\mathbb ZP$.
Choosing the powers $1,A_0,\ldots,A_0^{11}$
as lifts of $G$ gives the cyclic
factor set
$N\lfloor(i+j)/12\rfloor$.
Its cohomology class is precisely the
descent obstruction $d_2\xi$, so
$N=(6+12m)P$ for an integer $m$.
Replacing $A_0$ by $A=t_{-mP}\circ A_0$
gives $A^{12}=t_{6P}$.
The finite carry-cocycle identities,
its exact order two modulo twelve,
and the three local component checks
are recorded in
\texttt{ab\_push\_torsor\_cocycles.py}.

\item Neither this cohomological construction
nor local triviality specifies the
normal line bundles, the threefold
cusp charts, or a smoothing with the
given punctured-family identification.
Those remain requirements of \ref{push:src:B.4.1}.
\end{enumerate}
\end{proof}

\begin{proposition}[an auxiliary family with the wrong cusp]\label{push:src:B.3.7}\label{push:B37}
\StatementRecord{push:src:B.3.7}{proved}{Historical events: Claude Fable 5.2 (Anthropic), effort high, 16 November 2026: carried as background without fault found. GPT-6 Astra (OpenAI), ChatGPT Pro mode, effort Pro, S6\_09\_PAPERS, 17 November 2026: targeted proof check. Also checked in this writing. Pre-continuation writing events: GPT-6 Astra (OpenAI), ChatGPT Pro mode, effort Pro, S6\_10\_PAPERS, 18 November 2026: targeted proof check. Continuation metadata events: GPT-6 Astra (OpenAI), ChatGPT Pro mode, effort Pro, S6\_10\_PAPERS continuation, 19 November 2026: used and rechecked in this run; status-vocabulary reconciliation; metadata reconciliation only; no new mathematical check. \emph{Pre-continuation recorded proof scope.} Re-derived the Picard-degree calculation, degree-one coboundary correction from Phi12Phi6, exact compact auxiliary plumbing and twelve ordinary double points, Euler12/24 and valuation-divisibility obstruction. Closed relative-Picard/Brauer and height-scaling source gates: CDL2024 Prop0.9.2p143,Section4.2p370; Shioda Prop8.12/eq8.25p24 in author edition opened.}{the hypotheses, notation and definitions stated here; \cref{push:src:B.3.6}; \cite[Proposition 0.9.2 and the opening
of Section 4.2, p.370]{CDL24}; \cite[Proposition 8.12, equation (8.25)]{Shioda90}; \cref{push:src:A.1.2}; \cref{push:src:A.2.2}}

For the torsor and affine lift $(V,A)$
of \ref{push:src:B.3.6}, there is a line bundle $M$ on
$V$ representing the narrow Jacobian
section $P$, with
\[
 M^2=-2,\qquad
 M\otimes(A^*M)^{-1}
       \simeq\pi^*\mathcal O_C(\infty).
\]
These data give a proper flat nodal
plumbing $\mathcal Z^\flat\to\Delta_h$
and two uniform smooth-total-space
semistable modifications
$\mathcal Z^\pm\to\Delta_h$, together
with holomorphic maps to $B$.
For $h\ne0$, the nodal fibre has
twelve ordinary double points and
Euler number twelve; each small-resolved
fibre has Euler number twenty-four.
The cusp fibre of the latter map to
$B$ has multiplicity twelve.
Consequently this auxiliary construction
does not fill the prescribed $S^6$
family. No proper birational modification
over $B$ can turn its multiple cusp
into a reduced cusp.
\end{proposition}


\begin{proof}
\begin{enumerate}[label=\arabic*.,leftmargin=*]
\item The section $P$ determines a relative
Picard class of multidegree zero on
every fibre component of $V\to C$.
It lifts to a line bundle $M$ because
$\operatorname{Br}(C)=0$; this is the
relative Picard consequence of
\cite[Proposition 0.9.2 and the opening
of Section 4.2, p.370]{CDL24}.
In local trivializations of the torsor,
translation acts trivially on the
relative degree-zero Picard class.
Thus the local Poincar\'e lines glue
up to base lines, with precisely that
Brauer obstruction. In particular we
may and do retain multidegree zero
on all twelve fibre components.

\item Choose a finite multisection cover of
degree $d$ which trivializes the torsor
generically, and resolve its base change.
There the pulled-back $M$ becomes the
narrow Poincar\'e line for the pulled-back
$P$, up to a base line. Its square is
$-2d$, by height base change
\cite[Proposition 8.12, equation (8.25)]{Shioda90}.
Base lines do not change the square
of a fibre-degree-zero line. The
intersection pullback formula therefore
gives $M^2=-2$.

\item Since $\sigma P=P$, the line
$J=M\otimes(A^*M)^{-1}$ is trivial
on every fibre, including the
$I_{12}$ fibre. Grauert base change
and the evaluation map give
$J=\pi^*J_C$ for a line on $C$.
The line $M\otimes(t_{6P}^*M)^{-1}$
has base degree twelve. To check the
sign and degree, pass to the same
trivializing cover. On a Jacobian,
the degree of the restriction of
$\mathcal O(P-O)$ to a section $Q$
minus its degree on $O$ is
$-\langle P,Q\rangle$, by the height
formula and the narrowness of $P$.
For $Q=6P$ this is $-12d$.
Descending the degree factor $d$
proves the assertion. Iterating $A^*$
twelve times now gives
$12\deg J_C=12$, so $\deg J_C=1$.

\item The forms $da/b^2$, $a\,da/b^2$ and
$da/b$ form a basis of $H^0(C,K_C)$.
Their $\sigma$-eigenvalues are
$\lambda^7,\lambda^{10},\lambda^{11}$.
Thus the characteristic polynomial
on $H^1(C,\mathbb Z)$ is
\[
 p(x)=\Phi_{12}(x)\Phi_6(x)
     =x^6-x^5+x^3-x+1,
 \qquad
 p(x)+(1-x)(x^5+x^2+x)=1.
\]
It follows that $1-\sigma^*$ is an
automorphism of $\operatorname{Pic}^0(C)$,
with inverse $(\sigma^*)^5+(\sigma^*)^2+\sigma^*$.
Twisting $M$ by a base line changes
$J_C$ by $N\otimes(\sigma^*N)^{-1}$.
The displayed inverse therefore supplies
a base twist making
$J_C=\mathcal O_C(\infty)$.
Let $\phi$ be its section with divisor
the reduced $I_{12}$ fibre $F_\infty$.

\item We spell out the comparison with \ref{push:src:A.1.2}
to identify the family actually obtained.
Set $R_V=\mathbb P_V(\mathcal O\oplus M)$,
and identify its zero and infinity
sections by $s_0(x)\sim s_\infty(Ax)$.
On the node neighbourhood over $V$ use
\[
 a\in M_x,\quad b\in M_{Ax}^{-1},
 \qquad a\otimes b=h\phi_x.
\]
The $a$ collar is glued to the annular
core at $(x,a)$; the $b$ collar is glued
at $(Ax,b^{-1})$. Choose the Hermitian
metric and collar radii as in \ref{push:src:A.1.2}.
Compactness of $V$ and the bounded norm
of $\phi$ give disjoint collars,
Hausdorffness and compact representatives
by exactly the same inequalities.
These are holomorphic gluing maps of
fourfolds over $\Delta_h$.
Their local hypersurface rings are flat
over $h$, and the compact representatives
prove properness. The map to $B$ is
well defined because $t\circ\sigma=t$.
No identification with the requested
punctured $S^6$ family is being made.

\item At a smooth point of $F_\infty$ the
equation is $ab=hz_C$; at one of its
twelve nodes, $z_C=uv$ and it is
$ab=huv$. For $h\ne0$ the latter
has one ordinary double point.
These twelve centres are distinct
even if $A$ rotates the fibre
components. A centre has $a=b=0$
and is in neither annular collar;
no overlap identifies it with the
centre indexed by $Ax$.
The double locus is $V$, not $V/A$.
The two uniform prism subdivisions
of \ref{push:src:A.2.2} extend over these charts:
the affine action preserves the
orientation of the $I_{12}$ cycle,
and all remaining changes are units.
They give proper small modifications,
smooth total spaces and reduced
normal-crossings $h$-central fibres.
Equivalently their normalizations are
the blow-up of $R_V$ along $F_\infty$
in one section, followed by the twelve
chosen small resolutions. The normal
product becomes
$\pi^*\mathcal O_C(\infty)(-F_\infty)=\mathcal O_V$.

\item The nearby function of the nodal
plumbing is one on $F_\infty$ and
zero on its complement in the conductor.
Since $e(F_\infty)=12$, its integral
is twelve. Each small resolution
replaces a point by $\mathbb P^1$
and adds one, giving twenty-four.
Independently $e(R_V)=24$; above each
of the twelve nodes the modified
normalization replaces a point by
a chain of two projective lines,
of Euler number three. The open
parts of the fibre are copies of
$\mathbb C^*$ and contribute zero.
The modified normalization therefore
has Euler number $48$, its pinching
has Euler number $36$, and the
semistable nearby count is
$48-2e(V)=24$.

\item Finally $z=a/b$ is a coordinate on
$C$ near infinity and
$q_B=1/t=z^{12}/c^3$, with $c$ a unit.
At a generic point of the auxiliary
cusp fibre this gives multiplicity
twelve. In fact the map $q_B$ is
locally a twelfth power throughout
the smooth nearby threefold; a
holomorphic root of the unit can
be chosen locally. On any common
resolution of a proper birational
comparison over $B$, every divisor
over that fibre consequently has
$q_B$-valuation divisible by twelve.
A reduced fibre on another smooth
model would have a strict transform
of valuation one, a contradiction.
The finite Euler, polynomial and
multiplicity checks are in
\texttt{ab\_push\_auxiliary\_plumbing.py}.
\end{enumerate}
\end{proof}

\begin{proposition}[analytic unramified torsors do not evade the Hodge obstruction]\label{push:src:B.3.8}\label{push:B38}
\StatementRecord{push:src:B.3.8}{proved}{Historical events: Claude Fable 5.2 (Anthropic), effort high, 16 November 2026: carried as background without fault found. GPT-6 Astra (OpenAI), ChatGPT Pro mode, effort Pro, S6\_09\_PAPERS, 17 November 2026: targeted proof check. Also checked in this writing. Pre-continuation writing events: GPT-6 Astra (OpenAI), ChatGPT Pro mode, effort Pro, S6\_10\_PAPERS, 18 November 2026: targeted proof check. Continuation metadata events: GPT-6 Astra (OpenAI), ChatGPT Pro mode, effort Pro, S6\_10\_PAPERS continuation, 19 November 2026: used and rechecked in this run; status-vocabulary reconciliation; metadata reconciliation only; no new mathematical check. \emph{Pre-continuation recorded proof scope.} Re-derived whole-base R1Q comparison through translation action on the I12 cycle, b1=6 by Leray, and explicit positive curvature form on projectivization. Buchdahl and ddbar bimeromorphic input matching tracked separately.}{the hypotheses, notation and definitions stated here; \cref{push:src:B.3.4}; \cref{push:src:B.4.5}; \cref{push:src:B.3.6}; \cite[Theorem 11, p.297]{Buchdahl1999}; \cite[Section 2.1, p.8]{Chen2026HodgeRoot}}

Let $V\to C$ be any compact smooth analytic genus-one surface
locally isomorphic, as a torsor over every base disc, to the
Jacobian surface $\mathscr E\to C$ of \ref{push:src:B.3.4}. Algebraicity
or finite order of this analytic torsor class is not assumed.
Then $b_1(V)=6$, and $V$ is K\"ahler. For every holomorphic
line bundle $M$ on $V$, the projective bundle
$\mathbb P_V(\mathcal O\oplus M)$ is K\"ahler. Every smooth
compact threefold bimeromorphic to such a bundle satisfies
the $\partial\bar\partial$-lemma.

Consequently the necessary non-$\partial\bar\partial$
normalization component in \ref{push:src:B.4.5} cannot be obtained simply
by replacing the projective torsor of \ref{push:src:B.3.6} by a nonalgebraic
locally trivial analytic torsor, followed by a ruled bundle
and bimeromorphic modifications. This excludes that specified
repair; it does not assert that every possible cusp
normalization has this form.
\end{proposition}


\begin{proof}
\begin{enumerate}[label=\arabic*.,leftmargin=*]
\item Write $p:\mathscr E\to C$ for the projective Jacobian.
The calculation in \ref{push:src:B.3.4} gives
$R^1p_*\mathcal O_{\mathscr E}=\mathcal O_C(-\infty)$,
and $g(C)=3$. The coherent Leray sequence therefore gives
$h^1(\mathscr E,\mathcal O_{\mathscr E})=3$. Projectivity
and Hodge decomposition imply $b_1(\mathscr E)=6$.
The topological Leray sequence with rational coefficients is
\[
 0\longrightarrow H^1(C;\mathbb Q)
 \longrightarrow H^1(\mathscr E;\mathbb Q)
 \longrightarrow H^0(C,R^1p_*\mathbb Q)
 \longrightarrow H^2(C;\mathbb Q)
 \longrightarrow H^2(\mathscr E;\mathbb Q).
\]
The last arrow is injective because $p$ has a section.
The first arrow is an isomorphism of six-dimensional
vector spaces. Exactness thus implies
$H^0(C,R^1p_*\mathbb Q)=0$.

\item The local torsor identifications differ by translations.
On a smooth elliptic fibre translations act trivially on
$H^1$. On the $I_{12}$ fibre they preserve the orientation
of the component polygon and therefore act trivially on
its one-dimensional $H^1$ as well. These identifications
commute with the local specialization maps. Hence, for
$\pi:V\to C$, they glue to
$R^1\pi_*\mathbb Q\simeq R^1p_*\mathbb Q$ on the whole
curve, including infinity. The same Leray edge sequence
for $V$ now gives
$H^1(V;\mathbb Q)\simeq H^1(C;\mathbb Q)$.
Thus $b_1(V)=6$ is even. Buchdahl's theorem
\cite[Theorem 11, p.297]{Buchdahl1999} proves that $V$
is K\"ahler without an algebraicity hypothesis.

\item Choose a Hermitian metric on $\mathcal O\oplus M$.
The induced metric on the tautological relatively ample
line of its projectivization has a closed real $(1,1)$
curvature form $\theta$ positive in the vertical direction.
If $\omega$ is a K\"ahler form on $V$, compactness implies
that $\theta+A\pi^*\omega$ is positive in every direction
for a sufficiently large constant $A$. It is closed, so
it is a K\"ahler form on the projective bundle. A smooth
compact bimeromorphic model is therefore in Fujiki class
$\mathcal C$ and satisfies the $\partial\bar\partial$-lemma;
the domination statement is recalled, with its precise
reference to the $\partial\bar\partial$ theorem, in
\cite[Section 2.1, p.8]{Chen2026HodgeRoot}. Apply \ref{push:src:B.4.5}
for the final obstruction.
\end{enumerate}
\end{proof}


\subsection{Nonexistence of a proper analytic filling}

\begin{corollary}[counterexample; the proposed global filling]\label{push:src:B.4.1}\label{push:B41}
\StatementRecord{push:src:B.4.1}{proved}{Historical events: Claude Fable 5.2 (Anthropic), effort high, 15 November 2026: carried as background without fault found. Claude Fable 5.2 (Anthropic), effort high, 16 November 2026: carried as background without fault found. GPT-6 Astra (OpenAI), ChatGPT Pro mode, effort Pro, S6\_09\_PAPERS, 17 November 2026: targeted proof check. Also checked in this writing. Pre-continuation writing events: GPT-6 Astra (OpenAI), ChatGPT Pro mode, effort Pro, S6\_10\_PAPERS, 18 November 2026: targeted proof check. Continuation metadata events: GPT-6 Astra (OpenAI), ChatGPT Pro mode, effort Pro, S6\_10\_PAPERS continuation, 19 November 2026: used and rechecked in this run; status-vocabulary reconciliation; metadata reconciliation only; no new mathematical check. \emph{Pre-continuation recorded proof scope.} Re-derived the no-filling corollary from the already checked arbitrary-singularity exclusion, including preservation of properness by further finite base change and the actual punctured-circle monodromy scope.}{the hypotheses, notation and definitions stated here; \cref{push:src:B.4.33}; \cref{push:section:B.1}; \cref{push:section:B.2}; \cref{push:src:B.4.19}}

The actual punctured $S^6$ family has no proper
Hausdorff complex-analytic filling after any finite
cover of $w=0$. In particular, a proper flat filling
with smooth total space and reduced normal-crossings
centre does not exist. There is no positive least
semistable exponent, no equivariant analytic filling
to descend, and no central multiplicity data for
such a descent.
\end{corollary}


\begin{proof}
\begin{enumerate}[label=\arabic*.,leftmargin=*]
\item Theorem \ref{push:src:B.4.33} below proves the stronger assertion
allowing arbitrary singular central spaces and
multiplicities. A hypothetical filling after any
initial cover remains proper after a further cover;
the theorem excludes that further cover as well.
The local constructions in \ref{push:section:B.1} and \ref{push:section:B.2} and the
partial filling \ref{push:src:B.4.19} therefore cannot be completed
to the requested global object.
\end{enumerate}
\end{proof}

The obstruction cannot be removed by citing
$H^1(B,\mathcal V)=0$ from \cref{fam:src:A.5.1}. That
vanishing solves an additive translation
cocycle for a fixed nonzero parameter and
its integer shifts. On a finite overlap
the translation contains
\[
 \frac{\log s_j}{2\pi i}\,\Pi_c v_j,
 \qquad c=-\frac{12}{2\pi i}\log r.
\]
Its vertical exponential contains a product
of the two logarithms. No assertion in
the cited vanishing controls its growth
at $r=0$ or makes it a unit in the central
toric charts. \ref{push:src:B.3.1} proves that at least
one of those extension requirements must
fail for the prescribed pieces. That local
failure alone would not rule out a different cap.
\ref{push:src:B.4.30}--\ref{push:src:B.4.33} supply the additional global argument
which excludes every proper analytic filling of
the actual family, including such alternatives.

\begin{theorem}[fundamental-group obstruction]\label{push:src:B.4.2}\label{push:B42}
\StatementRecord{push:src:B.4.2}{proved}{Historical events: Claude Fable 5.2 (Anthropic), effort high, 16 November 2026: carried as background without fault found. GPT-6 Astra (OpenAI), ChatGPT Pro mode, effort Pro, S6\_09\_PAPERS, 17 November 2026: targeted proof check. Also checked in this writing. Pre-continuation writing events: GPT-6 Astra (OpenAI), ChatGPT Pro mode, effort Pro, S6\_10\_PAPERS, 18 November 2026: targeted proof check. Continuation metadata events: GPT-6 Astra (OpenAI), ChatGPT Pro mode, effort Pro, S6\_10\_PAPERS continuation, 19 November 2026: used and rechecked in this run; status-vocabulary reconciliation; metadata reconciliation only; no new mathematical check. \emph{Pre-continuation recorded proof scope.} Re-derived both connected-fibre covering-space surjectivity and monomial-gradient/meridian proof, with actual winding-one map from a one-pair pinching. Opened Nakayama--Ogus Theorem~5.1, Corollary~5.2 and Proposition~5.3, p.~2226, and matched the proper separated exact relatively smooth log-analytic hypotheses for the oriented fibre-bundle route. The second route is derived directly.}{the hypotheses, notation and definitions stated here; \cite[Theorem 5.1, p.2226]{NO10}}

Let $f:\mathcal Y\to\Delta$ be a
proper holomorphic map from a smooth
complex manifold, with smooth connected
fibres over $\Delta^*$ and reduced
normal-crossings central fibre $D$.
There is a continuous specialization
map from a nearby fibre $Y_t$ to $D$
which induces a surjection
\[
 \pi_1(Y_t)\longrightarrow\pi_1(D).
\]
No K\"ahler hypothesis is needed.
Hence the central fibre of any such
semistable filling of $S^6$ is simply
connected. In particular it cannot
be a simple pinching of a connected
normalization along one pair of
disjoint nonempty conductor sections,
whether its elliptic normalization
base is split or is a genus-one torsor.
\end{theorem}


\begin{proof}
\begin{enumerate}[label=\arabic*.,leftmargin=*]
\item Equip $\mathcal Y$ and $\Delta$ with
their divisorial logarithmic structures
from $D$ and $0$. Locally, after
absorbing a unit, the map is
$t=z_1\cdots z_k$. Its monoid chart
is the diagonal
$\mathbb N\to\mathbb N^k$; it is
fine, saturated, exact and log smooth.
Thus \cite[Theorem 5.1, p.2226]{NO10}
applies in the complex analytic
category. The induced map of real
oriented logarithmic blow-ups is a
proper topological fibre bundle over
$[0,\epsilon)\times S^1$.
Fix $\theta=\arg t$. Trivializing
this bundle along the radial interval
identifies $Y_t$ homeomorphically with
the fibre $D^{\log}_\theta$ over
$(0,\theta)$. Its natural projection
\[
 p_\theta:D^{\log}_\theta\longrightarrow D
\]
is proper and surjective. Above a
point lying on $k$ local branches,
its fibre consists of the phases
$(\theta_1,\ldots,\theta_k)$ whose
product is the fixed phase $\theta$,
including the harmless phase of the
local unit. It is a torus
$(S^1)^{k-1}$, in particular connected.

\item Here is a covering-space proof that
$p_\theta$ is surjective on fundamental
groups; path lifting for the collapse
is not assumed. The normal-crossings
space $D$ is locally contractible
and has a connected universal cover
$\widetilde D$. Pull it back along
$p_\theta$. The map from the pullback
to $\widetilde D$ is still proper,
surjective and has connected fibres.
Its total space is connected: a
nontrivial separation into open and
closed sets would, by properness and
connected fibres, give such a
separation of $\widetilde D$.
On the other hand its connected
components are indexed by the cosets
of the image of
$\pi_1(D^{\log}_\theta)$ in $\pi_1(D)$.
There is therefore only one coset.
Composing $p_\theta$ with the chosen
radial homeomorphism proves the claim.

\item We give an independent elementary
surjectivity argument. Put
$W=f^{-1}(\overline\Delta_\epsilon)$
and $W^*=W\setminus D$. The negative
gradient of $|f|^2$, for a Hermitian
metric, retracts $W$ onto $D$.
For completeness, in a product chart
$f=z_1\cdots z_k$ the arithmetic--
geometric mean inequality gives
\[
 \|\nabla |f|^2\|
   \ge c\,|f|^{2-1/k}.
\]
A finite cover of compact $D$ gives
a common exponent $\alpha<1$ in
$\|\nabla\rho\|\ge c\rho^\alpha$,
where $\rho=|f|^2$; for a fourfold
one may take $\alpha=7/8$.
The remaining trajectory length is
bounded by
$c^{-1}\rho^{1-\alpha}/(1-\alpha)$.
Hence trajectories converge uniformly
to $D$, and parametrization by the
level of $\rho$ gives a continuous
strong deformation retraction fixing
$D$. General position moves every
loop of $W$ off the real-codimension-two
normal-crossings strata, so
$\pi_1(W^*)\to\pi_1(W)$ is onto.
The smooth proper fibration over the
punctured disc gives
\[
 \pi_1(Y_t)\longrightarrow\pi_1(W^*)
       \longrightarrow\mathbb Z\longrightarrow0.
\]
A small transverse meridian at a
smooth point of a reduced component
maps to the generator of $\mathbb Z$
and bounds a disc in $W$.
Thus the image of $\pi_1(Y_t)$ alone
generates $\pi_1(W)=\pi_1(D)$.

\item For the last assertion, let $R$ be
the connected normalization and
$D_0,D_\infty$ its two disjoint
conductor sections. Choose a
continuous function $h:R\to[0,1]$
equal to zero on $D_0$ and one on
$D_\infty$. The function
$\exp(2\pi i h)$ descends to the
pinching. A path from a point of
$D_0$ to its paired point in
$D_\infty$ becomes a loop whose
image has winding number one.
Consequently the pinching has a
surjection of its fundamental group
onto $\mathbb Z$. This contradicts
the preceding necessary condition
when $Y_t\simeq S^6$.
\end{enumerate}
\end{proof}

\subsection{4 supplement. A Hodge obstruction to the central strata}

\begin{cited}[consequences with explicit derivation]\label{push:src:B.4.3}
\StatementRecord{push:src:B.4.3}{proved}{Historical events: Claude Fable 5.2 (Anthropic), effort high, 16 November 2026: carried as background without fault found. GPT-6 Astra (OpenAI), ChatGPT Pro mode, effort Pro, S6\_09\_PAPERS, 17 November 2026: targeted proof check. Also checked in this writing. Pre-continuation writing events: GPT-6 Astra (OpenAI), ChatGPT Pro mode, effort Pro, S6\_10\_PAPERS, 18 November 2026: used and rechecked in this run. Continuation metadata events: GPT-6 Astra (OpenAI), ChatGPT Pro mode, effort Pro, S6\_10\_PAPERS continuation, 19 November 2026: used and rechecked in this run; status-vocabulary reconciliation; metadata reconciliation only; no new mathematical check. \emph{Pre-continuation recorded proof scope.} Opened Nishiguchi MPI/88-48 Theorem1.4 printedp7 and Example2.1 printedp13, with their exact K3 and Hopf-component scope. The receiving citations now use these precise loci. Istrati source checks remain separately recorded; no new broad degeneration theorem is inferred.}{the hypotheses, notation and definitions stated here; \CompendiumRef{I}{app:rederive:push-src-B.4.3}; \cref{push:src:B.4.33}; \citet[Theorem 1.4, preprint p.~7, and Example 2.1, p.~13]{Nishiguchi1988IIRoot}; \citet[Theorem 7.1]{IstratiEtAl2022Root}}

Let $g:\mathcal Z\to\Delta_s$ be a proper holomorphic map from a
smooth complex fourfold, smooth over $\Delta_s^*$, with reduced
normal-crossings central fibre $D$. Suppose its nearby fibres are
biholomorphic to the six-sphere members of the specified family,
after any finite base change $w=s^N$. Then
\[
 H^1(D;\C)=0,\qquad h^1(D,\OO_D)\geq1.
\]
Consequently the natural map
$H^1(D;\C)\to H^1(D,\OO_D)$ is not surjective. In particular,
$D$ cannot be the analytification of a proper complex algebraic
variety with normal-crossings singularities.

If $D$ has \emph{simple} normal crossings, at least one of its
nonempty closed strata fails the $\partial\bar\partial$-lemma.
In particular, no such simple normal-crossings filling can have
all irreducible components K\"ahler or all irreducible components
Moishezon. These obstructions hold for every $N\geq1$.
\end{cited}
The expanded source derivation is retained in the proof compendium, Paper~I, label \nolinkurl{push:src:B.4.3}.




This is the stated application of the classical Euler and normal-crossing cohomology results; all hypotheses and the independent derivation are retained in the appendix. A complete re-derivation is given in \CompendiumRef{I}{app:rederive:push-src-B.4.3}.


This is an application of existing Hodge-theoretic results, not
a claim of a new degeneration theorem. Searches on 15 November
2026 for six-sphere semistable degenerations, normal-crossings
limits, Moishezon central fibres, and Conjecture 8.2 with Hodge
theory did not locate this specific application. We opened
Chen's version 2 and checked Definition 2.3
and Theorem 2.4(4), including the condition inserted before
parts (3)--(4); omitting that condition would make the
statement false. We also opened the original Graf--Kov\'acs
article and checked the proper Du Bois assertion in Section
1.D. The separate source audit records the searches.

The theorem prevents a repair made solely from compact
Moishezon components in a simple normal-crossings model.
It does not prove that an arbitrary analytic filling must have
such components. In particular, no argument here makes an
unknown three-dimensional component over the cusp Moishezon.
These Hodge arguments alone give conditional restrictions, not
unrestricted nonexistence. The latter is proved in \ref{push:src:B.4.33} for
every finite parameter cover, using the actual affine clutching.
The Hopf-component example of
\citet[Theorem 1.4, preprint p.~7, and Example 2.1, p.~13]{Nishiguchi1988IIRoot}
shows why an analytic degeneration cannot simply be assumed to
have class-$\mathcal C$ strata. The toric-Kato construction of
\citet[Theorem 7.1]{IstratiEtAl2022Root} is another relevant
non-K\"ahler analogue. Neither source supplies the required
cusp comparison for the six-sphere family.

\begin{proposition}[extension of the marked fibration]\label{push:src:B.4.4}
\StatementRecord{push:src:B.4.4}{proved}{Historical events: Claude Fable 5.2 (Anthropic), effort high, 16 November 2026: carried as background without fault found; targeted proof check. GPT-6 Astra (OpenAI), ChatGPT Pro mode, effort Pro, S6\_09\_PAPERS, 17 November 2026: targeted proof check. Also checked in this writing. Pre-continuation writing events: GPT-6 Astra (OpenAI), ChatGPT Pro mode, effort Pro, S6\_10\_PAPERS, 18 November 2026: targeted proof check. Continuation metadata events: GPT-6 Astra (OpenAI), ChatGPT Pro mode, effort Pro, S6\_10\_PAPERS continuation, 19 November 2026: used and rechecked in this run; status-vocabulary reconciliation; metadata reconciliation only; no new mathematical check. \emph{Pre-continuation recorded proof scope.} Read complete receiving proof and re-derived displayed sheaf/topological argument. The external logarithmic comparison applications in B4.30/B4.32/B4.33 were audited separately against IKN and its erratum; the independent direct commuting-meridian proof is night:P2.7. Boundary proof checked using intrinsic anticanonical divisor, coherence and proper-image exceptional loci.}{the hypotheses, notation and definitions stated here; \cref{fam:src:A.1.1}, \cref{fam:src:A.1.2}; \citep[Section 6, Hauptsatz I]{Grauert1960D}}

Suppose $g:\mathcal Z\to\Delta_s$ is a proper holomorphic map
from a smooth connected complex fourfold, and its restriction
over $\Delta_s^*$ is
isomorphic to a finite base change of the specified six-sphere
family. The given marked holomorphic map
\[
 f^*:\mathcal Z|_{\Delta_s^*}\longrightarrow\PP^1_t
\]
extends to a meromorphic map $\mathcal Z\dashrightarrow\PP^1_t$.
Its graph gives a proper projective modification
$\mathcal Z'\to\mathcal Z$, an isomorphism over
$\Delta_s^*$, on which the marked map is holomorphic.
No assertion that $\mathcal Z'$ is smooth or semistable is
included in this statement.
\end{proposition}


\begin{proof}
\begin{enumerate}[label=\arabic*.,leftmargin=*]
\item We shrink the disc whenever necessary. Put
$L=K_{\mathcal Z/\Delta_s}^{-2}$. The relative canonical
bundle here is the line bundle
$K_{\mathcal Z}\otimes g^*K_{\Delta_s}^{-1}$; its
restriction to a nonzero fibre is that fibre's canonical bundle.
The precise \cref{fam:src:A.1.1}, \cref{fam:src:A.1.2} input is
$K_{X_s}^{-2}\simeq f_s^*\mathcal O_{\PP^1}(1)$,
with connected fibres and no further sections. Therefore
$h^0(X_s,L|_{X_s})=2$, and its complete linear system is
the given map $f_s$, up to a projective change of coordinate.
The three critical values are intrinsically distinguished by
their fibres: multiplicity three, multiplicity four, and reduced
normal crossings, respectively.

\item Grauert's proper direct-image theorem
\citep[Section 6, Hauptsatz I]{Grauert1960D} makes
$\mathcal F=g_*L$ coherent. It is torsion-free: multiplication
by a nonzero base germ is injective on $L$, because $g$ is
nonconstant on the smooth connected total space, and hence is injective on
its direct image. A coherent torsion-free sheaf on the smooth
disc is locally free. Its rank is two. On the punctured disc,
constant $h^0=2$ and the analytic base-change theorem checked
in E.3 identify its fibres with the complete spaces of sections.
Choose a frame $(\sigma_0,\sigma_1)$ of $\mathcal F$.
Evaluation gives two global holomorphic sections of $L$ on
$\mathcal Z$, and thus a meromorphic map
$F_0=[\sigma_0:\sigma_1]$ to $\PP^1$.

\item For completeness its graph is analytic and proper without an
appeal to an unproved extension of a punctured graph. In a local
trivialization of $L$, write the sections as $a,b$. Remove their
common divisorial factor in the regular local ring of the smooth
space $\mathcal Z$. The equation
$v a-u b=0$ in $\mathcal Z\times\PP^1_{[u:v]}$
then has a unique component dominating $\mathcal Z$; it is
the closure of the graph where $(a,b)\ne(0,0)$. These local
components glue. Denote the resulting reduced graph by
$\mathcal G$. The projection $\mathcal G\to\mathcal Z$
is proper and projective, and is an isomorphism over
$\Delta_s^*$. The map
\[
 H:\mathcal G\longrightarrow\PP^1\times\Delta_s,
 \qquad H=(F_0,g),
\]
is proper. Locally $\mathcal G$ is a hypersurface in the smooth
space $\mathcal Z\times\PP^1$. The vanishing of the maximal
Jacobian minors for its equation and the two target coordinates
defines the closed analytic locus $C$ where $H$ is not a
submersion; singular points of $\mathcal G$ are included.

\item The image $H(C)$ is a closed analytic subset of the target.
One may deduce this precise proper-image assertion from the same
direct-image theorem: the coherent sheaf
$H_*\mathcal O_C$ has support exactly $H(C)$, since the
constant section $1$ is nonzero at every nonempty proper fibre.
Over $\Delta_s^*$ this image consists exactly of three distinct
critical-value sections. It cannot have a two-dimensional
component, since such an analytic component would contain an
open set of the connected two-dimensional target and hence meet
the punctured region in dimension two. Remove any component
contained in $\PP^1\times\{0\}$. The remaining reduced
analytic curve is proper with zero-dimensional fibres over the
disc, and is therefore finite over it.

\item Each of the three labelled sections on the punctured disc belongs
to a degree-one branch of that finite curve. Indeed their labels
are intrinsic and cannot be permuted by continuation. The
normalization of a degree-$d$ branch over a punctured disc has
cyclic monodromy of length $d$; a single-valued labelled section
therefore has $d=1$. The normalization extends it holomorphically
to a map from the full disc to $\PP^1$. Let these three maps be
$c_1(s),c_2(s),c_0(s)$, in the multiplicity-three,
multiplicity-four and cusp order. They may collide at $s=0$.

\item For $s\ne0$ the unique projective transformation sending
$(c_1,c_2,c_0)$ to $(0,1,\infty)$ has meromorphic coefficients
at $0$. In a common affine chart its value at $z$ is
\[
 A_s(z)=
 \frac{(z-c_1(s))(c_2(s)-c_0(s))}
      {(z-c_0(s))(c_2(s)-c_1(s))};
\]
homogeneous brackets give the same formula if a value is infinite.
The denominators are not identically zero, since the three values
are distinct on the punctured disc. Thus $A_s\circ F_0$ is
meromorphic on $\mathcal Z$. On every nonzero fibre it agrees
with the actual marked $f_s$, because both are the complete
bi-anticanonical map and have the same three labelled values.
This equality is an equality of maps over the punctured disc,
not just a fibrewise biholomorphism statement.

\item Finally the graph of $A_s\circ F_0$ is constructed by the same
local two-section argument after clearing powers of the base
coordinate in its meromorphic coefficients. Its projection to
$\mathcal Z$ is the asserted proper projective modification,
and its second projection is the required holomorphic marked map.
\end{enumerate}
\end{proof}

No condition on the central divisor was used: a nonzero base
germ is a nonzero divisor on the smooth dominating total space.
In particular the proposition applies after resolving the
reduced main component of any proper analytic filling, even
when the resulting central fibre has multiplicities.

The proposition removes a possible essential projective-coordinate
ambiguity. It does not extend a punctured isomorphism to any of
the auxiliary models of \ref{push:section:B.3}, and does not make a component over
$t=\infty,s=0$ Moishezon. Those would be additional statements.

\begin{theorem}[the obstruction on normalization components]\label{push:src:B.4.5}
\StatementRecord{push:src:B.4.5}{proved}{Historical events: Claude Fable 5.2 (Anthropic), effort high, 16 November 2026: carried as background without fault found. GPT-6 Astra (OpenAI), ChatGPT Pro mode, effort Pro, S6\_09\_PAPERS, 17 November 2026: targeted proof check. Also checked in this writing. Pre-continuation writing events: GPT-6 Astra (OpenAI), ChatGPT Pro mode, effort Pro, S6\_10\_PAPERS, 18 November 2026: targeted proof check. Continuation metadata events: GPT-6 Astra (OpenAI), ChatGPT Pro mode, effort Pro, S6\_10\_PAPERS continuation, 19 November 2026: used and rechecked in this run; status-vocabulary reconciliation; metadata reconciliation only; no new mathematical check. \emph{Pre-continuation recorded proof scope.} Re-derived ordered local-branch resolutions for constants and convergent functions, the exact Dolbeault-to-de Rham comparison, correction alpha by dbar h, and compact-pluriharmonic constancy on double branch spaces. Obtained degree-one surjectivity using ddbar only on normalization components.}{the hypotheses, notation and definitions stated here; \cref{push:src:B.4.2}; \cref{push:src:B.4.3}}

Let $D$ be a compact reduced complex space with normal-crossings
singularities, allowing self-intersections of global components.
Suppose that every connected component of its smooth normalization
satisfies the $\partial\bar\partial$-lemma. Then the natural map
\[
 H^1(D;\C)\longrightarrow H^1(D,\OO_D)
\]
is surjective.

Consequently every smooth-total-space reduced normal-crossings
filling of the specified six-sphere family, after any finite base
change, has a normalization component which fails the
$\partial\bar\partial$-lemma. In particular its normalization
cannot consist entirely of compact K\"ahler or Moishezon manifolds.
This conclusion includes normal crossings which are not simple.
\end{theorem}


\begin{proof}
\begin{enumerate}[label=\arabic*.,leftmargin=*]
\item We give a degree-one argument directly on $D$; no nearby
Hodge-de Rham degeneration theorem is needed. Let
$D^{[p]}$ be the smooth compact branch space whose points are
ordered $(p+1)$-tuples of distinct local branches through one
point of $D$. Equivalently, we normalize the loci where these
branches meet, retaining their ordering. These are distinct
\emph{local branches}, even when they belong to one global
irreducible component. The finite holomorphic maps
$\pi_p:D^{[p]}\to D$ and their face maps are specified by
forgetting one branch. The symmetric group acts by reordering;
we take the alternating summand, with the sign character, in
each degree. Averaging over this finite group is exact over
$\C$. Write $\delta$ for the alternating sum of the face
pullbacks.

\item Both augmented complexes
\[
 \C_D\longrightarrow
 \bigl((\pi_p)_*\C_{D^{[p]}}\bigr)_{p\geq0}^{\mathrm{alt}},
 \qquad
 \OO_D\longrightarrow
 \bigl((\pi_p)_*\OO_{D^{[p]}}\bigr)_{p\geq0}^{\mathrm{alt}}
\]
are resolutions. This can be checked locally without a global
ordering. In a chart $z_1\cdots z_k=0$, the constant complex is
the augmented cochain complex of the simplex on the $k$ branches.
For the holomorphic complex, fix a monomial in the convergent
power series ring. It survives on exactly the intersections
indexed by subsets of those $i$ for which its $z_i$-exponent is
zero. If there are no such indices it already vanishes on $D$;
otherwise its coefficient complex is again an augmented simplex
complex. The usual contraction is a finite combination of
coordinate restrictions, so preserves convergence. This proves
exactness for analytic functions, not just for formal series.
The descriptions agree under permutations of the local branches.
The finite simplex and permutation identities for the at most
four branches in our threefold application are checked exactly
in \path{code/b_nc_degree_one.py}.

\item Resolve the holomorphic terms by smooth Dolbeault forms and
the constant terms by smooth de Rham forms on the branch spaces.
These compute the respective hypercohomologies. Indeed finite
pushforward is exact here, and partitions of unity from local
ambient smooth functions pull back to make the resulting sheaves
fine. The ordinary de Rham lemma and Dolbeault lemma apply on
each smooth branch space. Projection to the $(0,\bullet)$
part defines the natural comparison between the two total
complexes and induces $H^1(D;\C)\to H^1(D,\OO_D)$.

\item Represent a coherent degree-one class by a total cocycle
$(\alpha,f)$, where $\alpha$ is a smooth $(0,1)$-form on
$D^{[0]}$ and $f$ is an alternating smooth function on
$D^{[1]}$. Choose the total-complex signs so that its equations
are
\[
 \bar\partial\alpha=0,\qquad
 \delta\alpha=\bar\partial f,\qquad \delta f=0.
\]
On a normalization component, $\partial\alpha$ is of type
$(1,1)$, is $\partial$-exact and is $\bar\partial$-closed.
The $\partial\bar\partial$-lemma gives a smooth function $h$
with $\partial\alpha=\partial\bar\partial h$.
Do this on every component, and put
\[
 \alpha'=\alpha-\bar\partial h,\qquad
 f'=f-\delta h.
\]
The pair represents the same coherent class and still satisfies
the displayed cocycle equations; now $d\alpha'=0$.
Holomorphic face maps commute with both differential operators,
so
\[
 \partial\bar\partial f'
       =\partial\delta\alpha'=\delta\partial\alpha'=0.
\]
A smooth complex-valued pluriharmonic function on a compact
connected complex manifold is constant: apply the maximum
principle to its real and imaginary parts in local complex
coordinates. Therefore $f'$ is constant on each connected
component of $D^{[1]}$. It follows that
$df'=0$ and $\delta\alpha'=0$.

\item Thus the same pair $(\alpha',f')$, with $\alpha'$ viewed as
a complex-valued smooth one-form, is a cocycle of the total
de Rham branch complex. Its image is the original coherent
class. This proves surjectivity. Notice that no
$\partial\bar\partial$ hypothesis on $D^{[1]}$ or higher
branch spaces was used; compactness suffices for the only
function argument there.

\item For the asserted application, \ref{push:src:B.4.2} and \ref{push:src:B.4.3} give
$H^1(D;\C)=0$ and $h^1(D,\OO_D)\geq1$.
Surjectivity is impossible, so some normalization component
fails the lemma. Compact K\"ahler and smooth Moishezon
manifolds satisfy it, as in the primary comparison checked
in \ref{push:src:B.4.3}. The same argument applies after every finite base
change, since its smooth fibres are still the specified
six-spheres.
\end{enumerate}
\end{proof}

This strengthens the simple-normal-crossings obstruction in
\ref{push:src:B.4.3} by using only degree one. It also avoids an unverified
extension of an algebraic Du Bois theorem to arbitrary analytic
spaces. Searches for normal-crossings cohomology with K\"ahler
normalization, the normalization Dolbeault complex, and
degree-one surjectivity located the broader Hodge-theoretic
setting discussed by \citet{Chen2026HodgeRoot} and the proper
algebraic statement of \citet{GrafKovacs2014Root}. The proof above
states and verifies the particular analytic comparison required
here. It still does not prove that all possible analytic
normalization components satisfy its hypothesis.

\begin{proposition}[a restricted non-K\"ahler cap obstruction]\label{push:src:B.4.6}\label{push:B46}
\StatementRecord{push:src:B.4.6}{proved}{Historical events: Claude Fable 5.2 (Anthropic), effort high, 16 November 2026: carried as background without fault found. GPT-6 Astra (OpenAI), ChatGPT Pro mode, effort Pro, S6\_09\_PAPERS, 17 November 2026: targeted proof check. Also checked in this writing. Pre-continuation writing events: GPT-6 Astra (OpenAI), ChatGPT Pro mode, effort Pro, S6\_10\_PAPERS, 18 November 2026: targeted proof check; used and rechecked in this run. Continuation metadata events: GPT-6 Astra (OpenAI), ChatGPT Pro mode, effort Pro, S6\_10\_PAPERS continuation, 19 November 2026: used and rechecked in this run; status-vocabulary reconciliation; metadata reconciliation only; no new mathematical check. \emph{Pre-continuation recorded proof scope.} Re-derived closed-block simplex resolutions, the five-term comparison and its exact surjectivity/injectivity hypotheses; computed horizontal conductor action Phi12 Phi6 with p(1)=1. Exact Kato-source input remains a source-match gate. Matched the exact toric Kato construction and cohomology inputs: Istrati et al Section2pp1351-1353,Proposition8.4p1377,Theorem8.5p1378. The hyperbolic scope of H1O=C is retained. Opened Schroeer–Siebert Definition2.1 p4 and the Section3 conductor/normal-line propositions pp6–10, together with the previously read Kamenova and KKN hypotheses. Rechecked the receiving explicit cohomology comparison: these source models are not asserted to identify an arbitrary cap. The only supported comparisons are the displayed actual maps and the specific hypotheses stated there.}{the hypotheses, notation and definitions stated here; \cref{push:src:B.3.8}; \cref{push:src:B.3.7}; \cite[Section 2, p.1352]{IstratiEtAl2022Root}; \cite[Proposition 8.4 and Theorem 8.5]{IstratiEtAl2022Root}; \cite[proof of Theorem 8.8, p.1380]{IstratiEtAl2022Root}; \cite[Theorem 5.2]{IstratiEtAl2022Root}; \cref{push:src:B.4.5}; \cref{push:src:B.4.2}; \cref{push:src:B.4.3}}

Let $D$ be a compact reduced normal-crossings complex
space. Suppose its irreducible components have been
partitioned into closed unions $D_0,\ldots,D_s$, with
their reduced induced structures. Assume that
\[
 H^1(D_i,\mathbb C)\longrightarrow H^1(D_i,\mathcal O)
 \quad\hbox{is surjective for every }i,
 \tag{B.46a}
\]
and that the analogous map is injective for every
pairwise intersection $D_i\cap D_j$. Then
\[
 H^1(D,\mathbb C)\longrightarrow H^1(D,\mathcal O_D)
 \quad\hbox{is surjective}.
 \tag{B.46b}
\]

In particular these hypotheses hold for the following
specified cap assembly. The horizontal block $D_0$ is
a two-section pinching of a ruled bundle over an
analytic locally trivial torsor $V\to C$ of \ref{push:src:B.3.8},
with conductor map $A$ covering $\sigma$ and
$A^{12}=t_{6P}$. The normalization may undergo the
modifications of \ref{push:src:B.3.7}, or any smooth bimeromorphic
modifications for which the two conductor restriction
maps on $H^1$ are the same pullbacks as before.
The other blocks are smooth hyperbolic toric Kato
threefolds or compact smooth spaces satisfying
(B.46a), for example class-$\mathcal C$ spaces.
Assume every pairwise overlap is either a reduced
union of invariant divisors of a hyperbolic toric
Kato threefold, or a compact normal-crossings surface
whose normalized components and normalized double
curves are rational. No such assembly can be a
semistable central fibre of the $S^6$ family.

This excludes the stated cap construction even when
an invariant Kato surface in an overlap kills a
fundamental-group generator of a Kato component.
It is not a classification of arbitrary analytic
cusp components or of their possible overlaps.
\end{proposition}


\begin{proof}
\begin{enumerate}[label=\arabic*.,leftmargin=*]
\item Write $D^{[p]}$ for the disjoint union of intersections
of $p+1$ distinct blocks. The augmented closed-cover
complexes resolve both $\mathbb C_D$ and $\mathcal O_D$.
For the coherent assertion this can be checked in a
normal-crossings coordinate chart. The local equation
of $D$ is a product of distinct coordinates, and each
block groups a disjoint subset of those factors.
The corresponding ideals are generated by products
in disjoint sets of variables. The ordinary alternating
intersection complex is exact, as one sees on each
monomial, where it is the augmented complex of a
simplex. The same contraction acts coefficientwise
on convergent power series. For constant sheaves,
small intersections are contractible unions of
coordinate subspaces, and the same stalkwise simplex
complex is exact. This proves both resolutions also
when the global blocks are self-crossing.

\item Their hypercohomology spectral sequences have
\[
 E_1^{p,q}(\mathcal F)=H^q(D^{[p]},\mathcal F),
 \qquad \mathcal F=\mathbb C\ \hbox{or}\ \mathcal O.
\]
Every connected compact reduced branch intersection
has only constant holomorphic functions. Thus the
comparison maps in the entire $q=0$ row are
isomorphisms. They induce isomorphisms on
$E_2^{1,0}$ and $E_2^{2,0}$. Surjectivity in (B.46a)
and injectivity on pairwise intersections imply
surjectivity on $E_2^{0,1}$: lift a coherent
degree-one class on the blocks, then its topological
restriction difference vanishes because its coherent
image vanishes and the pairwise comparison is
injective.

\item Now use the natural five-term exact sequences
\[
 0\longrightarrow E_2^{1,0}
 \longrightarrow H^1(D,\mathcal F)
 \longrightarrow E_2^{0,1}
 \xrightarrow{d_2} E_2^{2,0}.
\]
The isomorphism on the last term makes a lift of a
coherent $d_2$-closed element topologically
$d_2$-closed. The isomorphism on the first term then
gives (B.46b). If the block comparisons are
isomorphisms, the same argument also proves
injectivity, provided the map on $E_2^{0,1}$ is
injective.

\item We verify the horizontal-block assertion rather than
apply the pairwise criterion directly to its two
conductor branches. By \ref{push:src:B.3.8},
$H^1(V,\mathbb C)=H^1(C,\mathbb C)$ and
$H^1(V,\mathcal O)=H^1(C,\mathcal O_C)$.
The first equality is the Leray calculation there;
the second also follows from
$R^1\pi_*\mathcal O_V=\mathcal O_C(-\infty)$,
which has no global section. The action of $A$ on
these groups is $\sigma^*$, since it covers $\sigma$
and the displayed maps are pullbacks from the base.
The characteristic polynomial
\[
 p(x)=\Phi_{12}(x)\Phi_6(x),\qquad p(1)=1,
\]
shows that $1-A^*$ is invertible on both groups.
The normalization sequence of the two-section
pinching, for constants and for holomorphic
functions, therefore gives
\[
 H^1(D_0,\mathbb C)=\mathbb C,
 \qquad H^1(D_0,\mathcal O)=\mathbb C.
\]
Both classes come from the same difference map on
constant functions of the conductor. Their natural
comparison is an isomorphism. Projective-bundle
pullback identifies the normalization groups with
those of $V$. The stipulated modifications preserve
these groups and both restriction maps; in particular
the explicit modifications of \ref{push:src:B.3.7} have this property.
This grouped computation is essential: on an individual
conductor $V$, the comparison has dimensions six and
three and is not injective.

\item For a hyperbolic toric Kato manifold $Z$, its
fundamental group is $\mathbb Z$, as in
\cite[Section 2, p.1352]{IstratiEtAl2022Root};
\cite[Proposition 8.4 and Theorem 8.5]{IstratiEtAl2022Root}
give $H^0(Z,\Omega_Z^1)=0$ and
$H^1(Z,\mathcal O_Z)=\mathbb C$.
The exact sequence
$0\to\mathbb C\to\mathcal O\xrightarrow{d}d\mathcal O\to0$
shows that $H^1(Z,\mathbb C)\to H^1(Z,\mathcal O)$
is injective, because a global section of
$d\mathcal O$ is a closed holomorphic one-form.
The two dimensions are one, so it is an isomorphism.
The same reasoning applies to hyperbolic toric Kato
surfaces. Compact smooth rational varieties have
both groups zero. Class-$\mathcal C$ blocks satisfy
(B.46a) by their $\partial\bar\partial$ property.

\item Finally let $F$ be a reduced union of invariant
divisors in a hyperbolic toric Kato threefold.
Its normalized components are compact toric surfaces
or hyperbolic toric Kato surfaces
\cite[proof of Theorem 8.8, p.1380]{IstratiEtAl2022Root}.
Its normalized double curves are rational, since
every invariant curve is rational in the hyperbolic
case \cite[Theorem 5.2]{IstratiEtAl2022Root}.
Use the ordered-local-branch resolution from \ref{push:src:B.4.5}.
The comparisons on component $H^1$ are isomorphisms,
and the groups on normalized double curves vanish.
The same five-term argument shows that
$H^1(F,\mathbb C)\to H^1(F,\mathcal O_F)$ is an
isomorphism. For an overlap with rational normalized
components and double curves, the argument is
identical with zero component $H^1$. Thus every
pairwise overlap in the statement has the required
injective comparison, proving (B.46b) for the
specified assembly.

\item For an actual semistable $S^6$ centre, \ref{push:src:B.4.2} gives
$H^1(D,\mathbb C)=0$, while upper semicontinuity
and the nearby value $h^1(\mathcal O)=1$ in \ref{push:src:B.4.3}
give $h^1(D,\mathcal O_D)\geq1$. This contradicts
(B.46b). The obstruction uses the cohomology maps
on the overlaps; it does not assume that every
Kato component's loop survives in the union.
\end{enumerate}
\end{proof}


\paragraph{Why automatic analytic comparison is not used.}
We also examined whether general results on degenerations of tori
would force an arbitrary filling into the excluded toroidal class.
The distinction between a construction theorem and such a comparison
theorem is essential. \citet[Theorem 5.1.3 and Sections 5.2--5.4]
{KKN2008Root} construct proper analytic models of log complex tori
without assuming a polarization. Their input is a specified log
complex torus and an admissible pairing; the theorem does not assert
that every proper analytic filling is one of these models.
\ref{push:src:B.4.24} constructs the actual local Jacobian log input and isolates
its affine torsor obstruction. \ref{push:src:B.4.33} derives an associated torsor
from traces in a hypothetical cap; it does not identify that cap
with a proper fan model.

The classification in \citet[Theorem 3.1]{KamenovaToriRoot} assumes
a K\"ahler total space and trivial canonical bundle. Those hypotheses
have not been proved for a slice of an arbitrary filling here.
\citet[Proposition 3.1, preprint pp.~6--7]{SchroeerSiebert2002Root} explicitly
study normal-crossings degenerations with Hopf normalization, but
their smooth fibres are primary Kodaira surfaces. This is a source
of alternative analytic models to test, not a comparison with our
punctured family.

Even extending the marked map in \ref{push:src:B.4.4} does not extend a chosen
punctured isomorphism to a second model. Properness alone cannot
justify that inference: the automorphism
$(s,[z_0:z_1])\mapsto(s,[e^{1/s}z_0:z_1])$ of
$\Delta^*\times\PP^1$ has no meromorphic extension across zero,
although both product families are proper over the disc. Its
restriction to $[1:1]$ has an essential singularity. Thus an
unrestricted proof needs additional information beyond a
punctured isomorphism. \ref{push:src:B.4.30}--\ref{push:src:B.4.33} supply exactly such information:
bounded differences of extending sections and a global radial-loop
invariant. Their contradiction proves nonexistence without extending
an individual punctured isomorphism meromorphically.


\begin{proposition}[a parabolic analytic model and its comparison failure]\label{push:src:B.4.7}
\StatementRecord{push:src:B.4.7}{proved}{Historical events: Claude Fable 5.2 (Anthropic), effort high, 16 November 2026: carried as background without fault found. GPT-6 Astra (OpenAI), ChatGPT Pro mode, effort Pro, S6\_09\_PAPERS, 17 November 2026: targeted proof check. Also checked in this writing. Pre-continuation writing events: GPT-6 Astra (OpenAI), ChatGPT Pro mode, effort Pro, S6\_10\_PAPERS, 18 November 2026: targeted proof check; used and rechecked in this run. Continuation metadata events: GPT-6 Astra (OpenAI), ChatGPT Pro mode, effort Pro, S6\_10\_PAPERS continuation, 19 November 2026: used and rechecked in this run; status-vocabulary reconciliation; metadata reconciliation only; no new mathematical check. \emph{Pre-continuation recorded proof scope.} Re-derived the actual E-bundle over the blown-up parameter surface, Chern vectors 1,i,1+i and resulting Hopf multipliers, proper smooth gluing, unimodular pi1/H1 maps, all central homology groups and coherent comparison. Re-derived proper-graph obstruction by fundamental groups. Read Doi Section5.3 through Theorem5.5 pp881–883 and Example5.3 pp880–881; read Nishiguchi Theorem1.4 and Example2.1 and Schroeer–Siebert Section3. Matched smooth-map versus requested holomorphic-map data explicitly and retained the receiving complete proper product-cap construction. The SS Theorem4.4 quotient-generator exponent issue is recorded in the report and is unused by this receiving proof.}{the hypotheses, notation and definitions stated here; \citep[Definition 1.1, pp.~85--86]{Fujimoto1990AB}; \cref{push:src:B.4.6}; \citet[Section 2.2]{AbashevaDeev2024AB}; \cref{push:src:B.4.2}; \cref{push:src:A.4.7}}

There is a proper semistable analytic family of compact complex
surfaces $p:\mathcal H\to\Delta$ with smooth total space such that:
\begin{enumerate}
\item its punctured restriction is a holomorphic product with the
scalar primary Hopf surface
\[
 H_-=(\C^2\setminus\{0\})/\langle -e^{-\pi}\id\rangle;
\]
\item its central fibre is $Y=H_1\cup_E H_2$, where both $H_i$
are isomorphic to
\[
 H_+=(\C^2\setminus\{0\})/\langle e^{-2\pi}\id\rangle,
 \qquad E=\C/(\Z+i\Z);
\]
\item $N_{E/H_1}\simeq N_{E/H_2}\simeq\mathcal O_E$, and
\[
 \pi_1(Y)=0,\qquad h^1(Y,\mathcal O_Y)=1,
 \qquad H_*(Y;\Z)=(\Z,0,0,\Z^3,\Z^2);
\]
\item no punctured product isomorphism has a meromorphic extension
whose graph is proper over the two total spaces near the central
fibres.
\end{enumerate}
Consequently $\mathcal H\times\PP^1\to\Delta$ is an actual
semistable family of compact threefolds with simply connected centre
and $h^1(\mathcal O)=1$, whose normalized components fail the
$\partial\bar\partial$ lemma.  Its smooth fibres have Euler number
zero and fundamental group $\Z$.  It is not a filling of the
six-sphere family.
\end{proposition}


\begin{proof}
\begin{enumerate}[label=\arabic*.,leftmargin=*]
\item We construct the family by an explicit instance of Fujimoto's
generalized logarithmic transformation
\citep[Definition 1.1, pp.~85--86]{Fujimoto1990AB}.
Let
\[
 b:\mathcal B=\operatorname{Bl}_{(p,0)}(\PP^1\times\Delta)
       \longrightarrow\PP^1\times\Delta.
\]
The central curve is $F_1\cup F_2$, where $F_1$ is the strict
transform of $\PP^1\times\{0\}$ and $F_2$ is the exceptional
projective line.  They meet transversely in one point.  The two
line bundles
\[
 L=b^*\operatorname{pr}_1^*\mathcal O_{\PP^1}(1),
 \qquad A=\mathcal O_{\mathcal B}(F_2)
\]
have degree-restriction matrix
\[
 \begin{array}{c|rr}
       &L&A\\\hline
 F_1&1&1\\
 F_2&0&-1
 \end{array}.
\]
In particular this matrix is unimodular.  Take nonzero holomorphic
transition functions $l_{ab},a_{ab}$ for these two line bundles.
On small intersections choose logarithms and put
\[
 t_{ab}=\left[
       \frac{(1+i)\log l_{ab}-i\log a_{ab}}{2\pi i}
                         \right]\in E.
\]
This is independent of the logarithm choices, since $1+i,-i$
belong to the lattice of $E$.  On triple intersections the sum
is again a lattice element.  The $t_{ab}$ therefore form a
holomorphic cocycle with values in the translation group $E$.
Glue $U_a\times E$ by these translations, obtaining a
holomorphic principal $E$-bundle $\mathcal H\to\mathcal B$.
It is a proper submersion, since it is locally a product with a
compact elliptic curve.  Its total space is smooth and Hausdorff.
Compose with $\mathcal B\to\Delta$.

\item The Chern vector of the bundle, in the lattice $\Z+i\Z$, is
\[
 \kappa=(1+i)c_1(L)-i c_1(A).
\]
Its restrictions to $F_1,F_2$ are respectively $1$ and $i$.
On a nonzero fibre its restriction is $1+i$.  These statements
also follow directly from the displayed transition functions.
Locally at the crossing of $F_1$ and $F_2$ the base map is
$xy=u$, and the principal elliptic bundle is locally trivial.
The composed family is thus locally $(x,y,z)\mapsto xy$,
with $z$ a coordinate on $E$.  It is flat, has smooth total
space and reduced simple normal crossings centre.  It is proper
because both maps in the composition are proper.

\item We identify each surface explicitly.  For a primitive vector
$c\in\Z+i\Z$, choose $d$ so that $(c,d)$ is an oriented
lattice basis.  The homomorphism
\[
 \C^*\longrightarrow E,\qquad
 z\longmapsto\left[\frac{c\log z}{2\pi i}\right]
\]
has kernel generated by
$\alpha_c=\exp(2\pi i d/c)$, with $|\alpha_c|<1$.
The principal $E$-bundle on $\PP^1$ associated to
$\mathcal O(1)^\times$ by this homomorphism is the quotient
of $\mathcal O(1)^\times$ by that cyclic kernel.  Fibrewise
inversion identifies $\mathcal O(1)^\times$ with
$\mathcal O(-1)^\times\simeq\C^2\setminus\{0\}$;
inverting the cyclic generator does not change its group.
The resulting surface is the scalar Hopf surface $H_{\alpha_c}$.
For the three vectors in the construction we may take
\[
 (c,d)=(1,i),\quad(i,-1),\quad(1+i,i),
\]
giving respectively
\[
 \alpha_1=\alpha_i=e^{-2\pi},\qquad
 \alpha_{1+i}=\exp\bigl(2\pi i(1+i)/2\bigr)=-e^{-\pi}.
\]
On $\mathcal B|_{\Delta^*}=\PP^1\times\Delta^*$,
$A$ has the canonical nowhere-zero section of its exceptional
divisor.  Its cocycle is therefore a coboundary there, while
$L$ is the constant $\mathcal O(1)$ bundle.  The construction
gives an actual holomorphic product $H_-\times\Delta^*$,
not merely an isotrivial set of fibres.

\item Each central $H_i\to F_i$ is a principal elliptic bundle.
Its intersection with the other component is a smooth fibre,
so a local trivialization over a base disc gives
$N_{E/H_i}=\mathcal O_E$ exactly.  Thus the product of the
two normal lines is trivial, with the smoothing section supplied
by the coordinate $u=xy$.

\item Here are the integral topology maps.  A scalar Hopf surface is
diffeomorphic to $S^1\times S^3$: divide the radial coordinate
by multiplication by $|\alpha|$ and untwist its constant phase
rotation along the radial interval.  Its fibre inclusion has
fundamental-group map
\[
 \Lambda=\Z^2\longrightarrow\Lambda/\Z c_i.
\]
For $c_1=(1,0)$ and $c_2=(0,1)$, these are the two coordinate
quotients.  The van Kampen pushout kills both generators of
$\Lambda$, and hence $\pi_1(Y)=0$.  In the normalization
Mayer--Vietoris sequence the degree-one map, with the second
inclusion negated, is
\[
 H_1(E;\Z)\longrightarrow H_1(H_1;\Z)\oplus H_1(H_2;\Z),
 \qquad \begin{pmatrix}0&1\\-1&0\end{pmatrix}.
\]
It is unimodular.  The degree-zero map is the primitive column
$(1,-1)^t$, the degree-two target is zero, and $H_3(E)=0$.
The exact sequence now gives
\[
 H_1Y=H_2Y=0,\quad
 0\longrightarrow\Z^2\longrightarrow H_3Y
       \longrightarrow\Z\longrightarrow0,
 \quad H_4Y=\Z^2.
\]
All groups displayed in the proposition are therefore free.

\item For coherent cohomology, translations act trivially on
$H^1(E,\mathcal O_E)$, so for either central principal bundle
the sheaf $R^1\pi_*\mathcal O$ on $\PP^1$ is the constant
line bundle.  Its Leray sequence shows that restriction
\[
 H^1(H_i,\mathcal O_{H_i})\longrightarrow H^1(E,\mathcal O_E)
\]
is an isomorphism.  The normalization exact sequence therefore
has degree-one map $(s,t)\mapsto s-t$ from $\C^2$ to $\C$.
The preceding degree-zero map is surjective.  Its kernel has
dimension one, proving $h^1(Y,\mathcal O_Y)=1$.
This also displays exactly why the overlap-injectivity hypothesis
in \ref{push:src:B.4.6} fails: $H^1(E;\C)\to H^1(E,\mathcal O_E)$ has a
one-dimensional kernel.

\item For completeness, the two scalar Hopf surfaces used above are
not biholomorphic.  An isomorphism lifts to an automorphism of
$\C^2\setminus\{0\}$ that intertwines the two cyclic generators
or their inverses.  Hartogs extension applies to the lift and
its inverse; both fix zero and have invertible derivative there.
Differentiating the intertwining identity gives
$e^{-2\pi}=(-e^{-\pi})^{\pm1}$, which is false.  They have
no rational curves, since a map $\PP^1\to H_\alpha$ lifts to
$\C^2\setminus\{0\}$ and is constant.  The usual resolution
of a bimeromorphic map of smooth compact surfaces therefore
also shows that these two Hopf surfaces are not bimeromorphic.
This agrees with the explicit scalar case of
\citet[Section 2.2]{AbashevaDeev2024AB}.

\item There is an even shorter argument excluding a meromorphic
extension of the punctured product comparison.  A sufficiently
small closed semistable neighbourhood in $\mathcal H$ retracts
onto $Y$, by the monomial-gradient argument of \ref{push:src:B.4.2}.  Its
fundamental group is zero.  The product neighbourhood
$H_-\times\Delta$ has fundamental group $\Z$.  If the
punctured comparison had a proper meromorphic graph, resolving
that graph would give a common smooth proper modification of
the two neighbourhoods.  A proper bimeromorphic map between
smooth complex manifolds preserves $\pi_1$, by the
codimension-two avoidance argument in \ref{push:src:A.4.7}.  This is a
contradiction.  Both models already have the indicated
holomorphic map to the same $\PP^1$; the example therefore
retains an extended base map while the total-space comparison
fails.

\item Finally take the product with $\PP^1$.  Its central integral
Betti numbers are
\[
 (1,0,1,3,2,3,2),
\]
and $h^1(\mathcal O)$ remains one.  The nearby fibre is
$H_-\times\PP^1$, of fundamental group $\Z$ and Euler
number zero.  These explicit invariants prevent this example
from being mistaken for a filling of $S^6$.
\end{enumerate}
\end{proof}


\paragraph{Primary literature and the scope of the auxiliary cap}
The Cartier-divisor cocycle above is established machinery, not
a new logarithmic-transformation construction.  Fujimoto's
Definition~1.1 permits an arbitrary Cartier divisor; his
Remark~1.4 attributes non-K\"ahler torus degenerations to Ueno.
His Proposition~1.14 extends principal elliptic bundles across
subsets of codimension at least three.  Example~1.15 concerns
failure of extension as a principal elliptic bundle in
codimension two and explicitly leaves general relative
compactification open there; it supplies no unrestricted
nonexistence theorem for the present cusp
\citep[pp.~85--86, 93--94]{Fujimoto1990AB}.

Hopf components in analytic semistable degenerations are already
present in the original work of
\citet[Theorem 1.4, preprint p.~7, and Example 2.1, p.~13]{Nishiguchi1988IIRoot}.
The explicit diagonal Hopf quotients, elliptic conductor gluing
and normal-line calculation in
\citet[Propositions 3.1--3.4, preprint pp.~6--9]{SS00}
give a closely related analytic model.  Those are comparisons,
not smoothing theorems for the present threefold.  The construction
above supplies its own complete proper family and does not invoke
such a theorem by analogy.

There are two concrete obstacles to directly inserting the product
cap into B.3.  The divisor $E\times\PP^1$ in
$H_i\times\PP^1$ has normal degree zero along the $\PP^1$
ruling.  The exceptional divisor obtained by blowing up an
elliptic curve in a horizontal normalization has normal degree
$-1$ on its exceptional $\PP^1$ fibres.  For a two-component
attachment with no triple divisor on this interface, these cannot
be dual normal lines.  In the general normal-crossings formula
a triple divisor must also be included: a section
$E\times\{p\}$ has degree $+1$ on the ruling and could cancel
this particular mismatch.  No obstruction to such a higher-intersection
attachment is asserted.  Moreover the actual generic two-tori
do not possess the elliptic fibration of a principal elliptic
bundle, as the following exact calculation shows.

Fix an ordinary base point for which
$\mu\notin\Q\tau+\Q$.  Such points exist: otherwise the
ordinary base would be the countable union of the analytic loci
where one rational relation $\mu=r\tau+s$ holds; Baire's
theorem would force one such relation identically, contradicting
that $P$ is a nontorsion section.  For integral vectors
$v=(a,b,c,d)$ and $v'=(A,B,C,D)$, write
\[
 \Pi v=(6a\mu+b\tau+c,\ a\beta+b\mu+d).
\]
The coefficient of $\beta$ in
$\det(\Pi v,\Pi v')$ is
\[
 (Ab-aB)\tau+(Ac-aC).
\]
If it vanishes and one of $a,A$ is nonzero, the first three
coordinates of $v,v'$ are rationally proportional.  The
remaining determinant is
$(6a\mu+b\tau+c)(D-(A/a)d)$ after choosing $a\ne0$;
the first factor is nonzero, so a zero determinant would make
the two vectors dependent.  The case $a=0,A\ne0$ is the
same statement with the vectors interchanged.  If $a=A=0$,
the determinant is
\[
 (bD-Bd)\tau+(cB-Cb)\mu+(cD-Cd),
\]
and the rational independence of $1,\tau,\mu$ again forces
all the minors to vanish.  Thus every independent integral
pair spans a complex line for at most one value of $\beta$.
There are countably many pairs.  For all other admissible
parameters the two-torus is simple and contains no elliptic
subtorus.  In particular the distinguished line $\C\delta$
meets the period lattice in $\Z\delta$, so its complexified
orbit is $\C^*$, rather than an elliptic curve.

A principal elliptic-bundle construction may therefore describe
a central cap, but cannot by itself equal the entire prescribed
punctured cusp family.  It would need an additional affine
identification that destroys its generic elliptic fibration,
together with the actual nodal $q$-interface and the finite-place
component rotations. The construction here establishes none of
those additional comparisons. \ref{push:src:B.4.33} separately rules out a proper
analytic filling of the actual global family, without denying the
auxiliary Hopf degeneration constructed above.


\Uses{the hypotheses, notation and definitions stated here; \citep[pp.~85--86, 93--94]{Fujimoto1990AB}; \citet[Theorem 1.4, preprint p.~7, and Example 2.1, p.~13]{Nishiguchi1988IIRoot}; \citet[Propositions 3.1--3.4, preprint pp.~6--9]{SS00}; \cref{push:src:B.4.33}}



\paragraph{Limits of automatic analytic comparison}
The abstract identity $\Aut(X)=\C^*$ does not supply a
meromorphic compactification of the action. On an ordinary
two-torus the distinguished orbit is a nonconstant holomorphic
map from $\C^*$. If it extended meromorphically across both
ends to $\PP^1$, it would extend holomorphically: a
meromorphic map from a smooth curve to a compact complex
manifold has a removable indeterminacy. Every holomorphic
map $\PP^1\to T^2$ is constant, since it lifts to $\C^2$.
This is a contradiction. Thus an algebraic structure on the
abstract automorphism group is insufficient here.

The compactification theorem of
\citet[Theorem 1.2]{BiliottiGhigi2021Root} assumes Fujiki
class $\mathcal C$; that hypothesis is unavailable.
The local isomorphism theorem of
\citet[Theorem 3]{Meersseman2011Root} assumes pointwise
isomorphism on the whole base, including its central point,
and constant vector-field dimension. It cannot identify the
unknown central fibre. The Hopf compactification example of
\citet[Example 1]{KatzarkovEtAl2025Root} likewise illustrates
that an analytic open part need not determine a bimeromorphic
compactification. Their definability-based comparison is not
used: no definability of our punctured comparison has been
proved. \ref{push:src:B.4.7} supplies a proper family exhibiting the relevant
failure directly. None of these automatic-comparison results
proves the conclusion of \ref{push:src:B.4.33}, whose argument instead uses
extending-section differences and the global affine loop invariant.

Non-K\"ahler modifications of elliptic degenerations are
established. The product case in
\citet[Section 3, Example 3.4(1)]{Fujimoto1988Root}, using
the preceding construction, already gives such degenerations
with elliptic abelian nearby surfaces. Its elliptic-bundle
structure prevents a direct identification with our generic
simple two-tori. For the related codimension-two extension
question, \citet[Theorem 7.1.3]{Nakayama2002Root} proves an
extension criterion for locally bimeromorphically projective
elliptic fibrations over a punctured polydisc: the fibration
must be bimeromorphically projective over that punctured base.
It rules out the sectionless elliptic example discussed there.
Our regular fibres have dimension two and no asserted elliptic
quotient, so that theorem does not establish a nonexistence
result for the present filling.

The differential-geometric smoothing theorem of
\citet[Theorem 5.5, pp.~882--883]{Doi2009Root} allows
non-K\"ahler components with dual normal bundles,
anticanonical double divisors and compatible residues.
Its total space is a smooth real six-manifold and its
map is smooth; the complex structures on the surface
fibres vary continuously in the stated sense.
This is not a holomorphic total-space filling theorem.
Example 5.3 identifies the degree-zero elliptic ruled
self-gluing with a torus and the nonzero-degree cases
with Kodaira surfaces. Neither result supplies the
prescribed holomorphic two-parameter comparison at our
missing origin.


\Uses{the hypotheses, notation and definitions stated here; \citet[Theorem 1.2]{BiliottiGhigi2021Root}; \citet[Theorem 3]{Meersseman2011Root}; \citet[Example 1]{KatzarkovEtAl2025Root}; \cref{push:src:B.4.7}; \cref{push:src:B.4.33}; \citet[Section 3, Example 3.4(1)]{Fujimoto1988Root}; \citet[Theorem 7.1.3]{Nakayama2002Root}; \citet[Theorem 5.5, pp.~882--883]{Doi2009Root}}


\begin{proposition}[a normal-line obstruction on every semistable good-reduction interface]\label{push:src:B.4.8}
\StatementRecord{push:src:B.4.8}{proved}{Historical events: Claude Fable 5.2 (Anthropic), effort high, 16 November 2026: carried as background without fault found. GPT-6 Astra (OpenAI), ChatGPT Pro mode, effort Pro, S6\_09\_PAPERS, 17 November 2026: targeted proof check. Also checked in this writing. Pre-continuation writing events: GPT-6 Astra (OpenAI), ChatGPT Pro mode, effort Pro, S6\_10\_PAPERS, 18 November 2026: targeted proof check. Continuation metadata events: GPT-6 Astra (OpenAI), ChatGPT Pro mode, effort Pro, S6\_10\_PAPERS continuation, 19 November 2026: used and rechecked in this run; status-vocabulary reconciliation; metadata reconciliation only; no new mathematical check. \emph{Pre-continuation recorded proof scope.} Re-derived mixed Chern clutching covector (0,mN), semistable genus-one dual-graph alternatives and their explicit Dehn-twist monodromy, two bounding solid tori, translation invariance modulo fibre degree, and exact actual-cusp multipliers -(5 plus/minus 2sqrt6).}{the hypotheses, notation and definitions stated here; \citet[Section 2.2]{AbashevaDeev2024AB}}

Let $a,b$ be holomorphic units on a disc $\Delta_q$, with
$0<|a(q)|<1$, and let $m\ne0$ be an integer.  Put
\[
 \mathcal E_a=(\C^*\times\Delta_q)/
       \bigl((z,q)\sim(a(q)z,q)\bigr).
\]
This is a smooth proper elliptic family over the full disc.
Over $\Delta_q^*$ define a line bundle $\mathcal L_m$ by
\[
 (z,v,q)\sim\bigl(a(q)z,\ q^m b(q)v,\ q\bigr).
\]
Its restriction to each elliptic fibre has degree zero.
After any finite base change $q=s^N\varepsilon(s)$, where
$N>0$ and $\varepsilon$ is a holomorphic unit, the following
extension is impossible: a proper semistable family
$p:\mathcal C\to\Delta_s$ of connected curves, with smooth
total surface, an isomorphism
$\mathcal C|_{\Delta_s^*}\simeq\mathcal E_a|_{\Delta_s^*}$,
and a line bundle on $\mathcal C$ restricting to
$\mathcal L_m$ under that isomorphism.

The assertion allows any finite collection of rational tails
in the central curve and arbitrary translations in the
punctured isomorphism.  It remains true after tensoring
$\mathcal L_m$ by the restriction of a line bundle extending
on the proposed interface.  It does not assert an obstruction
for arbitrary singular, non-flat, or higher-dimensional
interface networks.

In particular the axis $z_1=0$ in the diagonal Hopf surfaces
\[
 H_{q,a(q)}=(\C^2\setminus\{0\})/
       \langle(z_1,z_2)\mapsto(qz_1,a(q)z_2)\rangle
\]
cannot acquire such a semistable completion together with
its normal line, after any finite base change.
\end{proposition}


\begin{proof}
\begin{enumerate}[label=\arabic*.,leftmargin=*]
\item Choose a logarithm of $a$ on the disc and set
$\tau_a=\log a/(2\pi i)$.  Then $\Im\tau_a>0$, and
\[
 \mathcal E_a\simeq
 (\C\times\Delta_q)/(\Z+\tau_a(q)\Z).
\]
The real coordinates $z=\exp(2\pi i(x+\tau_a(q)y))$
give a smooth trivialization by the marked torus
$(x,y)\in\R^2/\Z^2$.  In particular its monodromy on
$H_1$ is the identity.  The line factors in this marking
are $1$ along the $x$ cycle and $q^m b(q)$ along the
$y$ cycle.  On a positive parameter circle, the latter
factor has winding $m$.  Equivalently, the clutching
calculation gives
\[
 \langle c_1(\mathcal L_m),[T_x]\rangle=0,
 \qquad
 \langle c_1(\mathcal L_m),[T_y]\rangle=m,
 \qquad
 \langle c_1(\mathcal L_m),[\mathcal E_{a,q}]\rangle=0.
\]
Here $T_x,T_y$ sweep the indicated cycles around the
parameter circle; the orientation of $T_y$ is chosen to
agree with the clutching convention in the displayed
line factor.  Reversing that convention changes both
sides' sign and has no effect on nonvanishing.  One can
also compute the integer on the square
$[0,1]_y\times S^1_q$: the two $y$ edges are glued by
$q^m b(q)$, whose degree is $m$.

\item Under $q=s^N\varepsilon(s)$ the degree becomes $mN$.
The unit $\varepsilon(s)^m b(q(s))$ has zero winding.
Thus the mixed Chern class is the nonzero covector
\[
  (0,mN)\in
 H^1(\mathcal E_{a,q};\Z)\otimes H^1(S_s^1;\Z).
\]
This is an integral calculation; no limiting logarithm
or asymptotic estimate is involved.

\item Suppose an extension $\mathcal C$ existed.  Its central
fibre is connected.  Flatness and properness give arithmetic
genus one, and the normalization exact sequence gives
\[
 1=\sum_i g(C_i)+b_1(\Gamma),
\]
where $C_i$ are the smooth normalization components and
$\Gamma$ is the dual graph, including a loop for a
self-node.  Consequently either there is a unique
elliptic component and $\Gamma$ is a tree, or all
components are rational and $\Gamma$ has one cycle.

\item The second case has nontrivial monodromy.  This can be
seen directly in the local plumbing $uv=s$.  Transport
around $s=0$ is a positive Dehn twist on its neck annulus.
An edge outside the unique graph cycle gives a separating
neck and acts trivially on $H_1$.  Each of the $k\ge1$
edges in the graph cycle gives a nonseparating neck.
These disjoint necks represent the same primitive class
up to sign on the smoothed genus-one surface.  The
transvection is unchanged by reversing that class.
Their product therefore has, in a symplectic basis, the
matrix
\[
  \begin{pmatrix}1&k\\0&1\end{pmatrix}
\]
or its inverse under the opposite monodromy convention.
It is not the identity.  The given punctured elliptic
family has identity monodromy, so this case is impossible.
This elementary plumbing argument is all that is needed
here; no algebraic stable-reduction or comparison theorem
for arbitrary analytic total spaces is being assumed.

\item There is therefore an elliptic core $C_0$ with rational
trees attached at finitely many points.  Choose two
oriented loops on $C_0$ avoiding those points and with
intersection number one.  The map $p$ is a submersion
near their union.  Compactness of the loops and a smooth
horizontal distribution provide, after shrinking the
disc, families of loops over the full disc.  Each swept
torus on its boundary bounds the resulting solid torus
in $\mathcal C$.  Hence the first Chern class of every
line bundle on $\mathcal C$ evaluates to zero on both
such swept tori.

\item The two transported loops form a basis of the nearby
elliptic fibre's $H_1$: cutting off the rational trees
adds only genus-zero pieces along separating necks,
and their intersection number remains one.  Under the
given punctured isomorphism, their swept tori therefore
project to a basis of the mixed group
$H_1(\mathcal E_{a,q};\Z)\otimes H_1(S_s^1;\Z)$.
They may differ from the marked swept tori by integral
fibre classes.  This accounts in particular for an
arbitrary loop of fibre translations in the isomorphism.
The fibre degree of $\mathcal L_m$ is zero, so these
additional fibre classes do not change either pairing.
The nonzero covector $(0,mN)$ cannot vanish on both
basis vectors.  This contradicts the bounding-solid-torus
calculation and proves the assertion.

\item Finally an extending tensor factor has zero pairing on
both bounding tori by the same argument.  It cannot
cancel the obstruction.  Vertical divisors and any
Cartier triple divisor on such an extended interface
are included among these harmless tensor factors.

\item For the stated Hopf consequence, the axis is precisely
$\mathcal E_{a,q}$.
The derivative of the deck transformation in the $z_1$
direction multiplies a normal vector by $q$; its action
on the axis coordinate is multiplication by $a(q)$.
These are exactly the two factors defining
$\mathcal L_1$.  This is also the explicit diagonal-Hopf
normal-line convention of
\citet[Section 2.2]{AbashevaDeev2024AB}.
The assertion just proved applies with $m=1$ and $b=1$.
It also applies to the first multiplier $q^m b(q)$ for
any positive $m$.  At the actual cusp the nonzero unit
$e^{2\pi i\mu(q)}$ has absolute value different from one
at $q=0$.  Indeed the nodal cubic is
$Y^2=(X-1)^2(X+2)$; its normalization
$X=h^2-2$, $Y=h(h^2-3)$ and the multiplicative coordinate
$(h-\sqrt3)/(h+\sqrt3)$ show that the marked section at
$(X,Y)=(0,\pm\sqrt2)$ has multiplier one of
$-(5\pm2\sqrt6)$.
Choose that multiplier or its inverse to obtain
such a contracting $a(q)$.
This explains why the candidate axis has good reduction.
No claim is made that an arbitrary filling is obliged
to contain this particular axis or normal factor.
\end{enumerate}
\end{proof}


\paragraph{What a corrected cap must change}
The obstruction concerns an individual normal line,
not just the tensor product of two normal lines in
the $T^1$ condition.  Opposite windings can cancel in
that tensor product while neither line extends on its
own proposed smooth branch.  Thus cancelling the
scalar $q$ by $q^{-1}$ is insufficient for this
good-reduction interface.  A successful correction
would have to change the interface or its punctured
normal bundle outside the class specified above.
The constant-elliptic example \ref{push:src:B.4.7} has no such
winding and remains consistent with \ref{push:src:B.4.8}.


\Uses{the hypotheses, notation and definitions stated here; \cref{push:src:B.4.7}; \cref{push:src:B.4.8}}


\begin{proposition}[ordinary logarithmic torus cocycles and the remaining descent classes]\label{push:src:B.4.9}
\StatementRecord{push:src:B.4.9}{proved}{Historical events: Claude Fable 5.2 (Anthropic), effort high, 16 November 2026: carried as background without fault found. Claude Fable 5.2 (Anthropic), effort high, 17 November 2026: carried as background without fault found. GPT-6 Astra (OpenAI), ChatGPT Pro mode, effort Pro, S6\_09\_PAPERS, 17 November 2026: targeted proof check. Also checked in this writing. Pre-continuation writing events: GPT-6 Astra (OpenAI), ChatGPT Pro mode, effort Pro, S6\_10\_PAPERS, 18 November 2026: targeted proof check. Continuation metadata events: GPT-6 Astra (OpenAI), ChatGPT Pro mode, effort Pro, S6\_10\_PAPERS continuation, 19 November 2026: used and rechecked in this run; status-vocabulary reconciliation; metadata reconciliation only; no new mathematical check. \emph{Pre-continuation recorded proof scope.} Re-derived C-star homomorphism-to-flat-cycle correspondence, exact commuting monodromy kernels and primitive Koszul cokernel Z2, Cartier class zero map, and actual quadratic affine section sigma(q,c) with sigma(c+h)-sigma(c)=b\_h,c. Cartan-B exact primary locator still tracked separately. Closed the CartanB source gate by opening Expose19 Section5pp19-14--19-15, including the degree-one convergence/gluing argument.}{the hypotheses, notation and definitions stated here; \cref{push:src:B.3.1}; \cref{push:src:C.3.5}; \citep[Definition 1.1]{Fujimoto1990AB}; \citep[Section 5]{Cartan1952B_AB}}

Let $\mathcal T\to U$, $U=\Delta_q^*\times\Delta_r^*$,
be the Jacobian two-torus family of the marked cusp,
after $w=r^6$.  A holomorphic homomorphism of group
families
\[
 \C^*\times U\longrightarrow\mathcal T
\]
is determined by a flat integral cycle.  For this family
all such cycles are
\[
 c=a\widehat w+b\widehat\delta,\qquad a,b\in\Z.
\]
In the exponential cusp coordinates its homomorphism
is $z\mapsto(z^a,z^b)$.  Thus the ordinary Cartier
logarithm construction along a globally specified
cocharacter has transition maps
\[
 (Z,W)\longmapsto(g_{ij}^{\,a}Z,g_{ij}^{\,b}W).
\]
If these units extend over the prescribed fan charts,
the maps extend there and preserve each toric
boundary component.  Changing only these transition
units cannot repair the component permutations
$+4,-3,0$ refuted in \ref{push:src:B.3.1}.

This is not a classification of all torus-torsor
descent.  In fact, for the sheaf $\mathscr T$ of
holomorphic sections of $\mathcal T$,
\[
 H^1(U,\mathscr T)\simeq H^2(U,\Lambda)\simeq\Z^2.
\]
The last group is represented by the coinvariant
classes of $\widehat\gamma,\widehat u$.  The ordinary
fixed-cocharacter Cartier constructions have zero
image in this group.  The actual cusp family with
the coherent closing translations of \ref{push:src:C.3.5} also
has zero torsor class, by the explicit holomorphic
section below.  Thus nonzero classes in this
$\Z^2$ do not have its prescribed marked punctured
cusp germ.  No compactification of an arbitrary
torsor is asserted.
\end{proposition}


\begin{proof}
\begin{enumerate}[label=\arabic*.,leftmargin=*]
\item First consider any smooth holomorphic family of
complex tori with integral period local system
$\Lambda$ and relative Lie bundle $\mathcal V$.
Lift a group homomorphism from $\C^*$ along the
universal covering $z=\exp(2\pi i t)$.
The lifted homomorphism from the additive group
$\C$ to the universal covering vector space of
each torus is $t\mapsto t v$ for a holomorphic
section $v$ of $\mathcal V$.  Periodicity under
$t\mapsto t+1$ says precisely that $v$ belongs
to the period lattice in every fibre.  In any
local marking its integral coefficients are
continuous, and hence locally constant.  It
therefore defines a flat global section of
$\Lambda$.  Conversely any such section gives
the homomorphism by this formula.  The image
need not be a closed elliptic subgroup; this
argument applies without a polarization and
also to simple two-tori.

\item In the ordered basis
$(\widehat\gamma,\widehat u,\widehat w,\widehat\delta)$,
the actual commuting linear monodromies have
nilpotent parts
\[
 N_q=A_0-I=
 \begin{pmatrix}
 0&0&0&0\\0&0&0&0\\0&1&0&0\\-1&0&0&0
 \end{pmatrix},\qquad
 N_r=-12\widehat\delta\otimes\gamma.
\]
The sign of $N_r$ is the positive-loop closing
convention of \ref{push:src:C.3.5}; reversing it leaves all
groups below unchanged.  Their common integral
kernel is exactly $\Z\widehat w\oplus\Z\widehat\delta$.
The two corresponding columns of $\Pi_c$ are
$(1,0)^t,(0,1)^t$.  Exponentiating the homomorphism
therefore gives $(z^a,z^b)$, as asserted.

\item Fujimoto's logarithmic cocycle formula
\citep[Definition 1.1]{Fujimoto1990AB} works in
the same way for a complex torus of any dimension:
apply the displayed homomorphism to the transition
functions $g_{ij}$ of a line bundle.  Equivalently
use $c\log(g_{ij})/(2\pi i)$ modulo the period
lattice.  A different logarithm adds an integral
cycle, and the triple-overlap defect is an
integral cycle as well.  With the present common
invariant cycles this is exactly multiplication
by $(g_{ij}^a,g_{ij}^b)$.  Torus multiplication
extends to every toric variety of a fixed fan
and preserves every torus orbit and boundary
ray.  For an extending unit it is an automorphism
on the full chart.  Consequently it contributes
no permutation of the central component labels.
The product-of-permutations obstruction of \ref{push:src:B.3.1}
is unchanged.  This conclusion keeps the fan,
period action, and local component models fixed.

\item To exhibit the scope left outside that calculation,
retract $U$ to its real two-torus.  The cellular
cochain complex with coefficients in the local
system is the two-generator Koszul complex.
Its last differential has image
$N_q\Lambda+N_r\Lambda$, and hence
\[
 H^2(U,\Lambda)
   =\Lambda/(N_q\Lambda+N_r\Lambda)
   =\Lambda/(\Z\widehat w+\Z\widehat\delta)
   \simeq\Z\widehat\gamma\oplus\Z\widehat u.
\]
The relation lattice is primitive: $N_q$ already
has image $\Z\widehat w+\Z\widehat\delta$.
The exact matrix and Smith calculation is saved
in \path{code/ab_push_hopf_winding.py}.

\item The analytic exponential sequence is
\[
 0\longrightarrow\Lambda\longrightarrow
 \mathcal O(\mathcal V)\longrightarrow
 \mathscr T\longrightarrow0.
\]
The product of two punctured discs is Stein,
and $\mathcal O(\mathcal V)$ is coherent.
Cartan's theorem B therefore kills its positive
cohomology
\citep[Section 5]{Cartan1952B_AB}.
The long exact sequence gives the stated
isomorphism $H^1(U,\mathscr T)\simeq H^2(U,\Lambda)$.
These classes do define holomorphic torsors by
translation gluing, but this calculation contains
no compactness assertion at either missing axis.

\item For the last assertion, naturality of the two
exponential sequences sends the first Chern class
of a Cartier line bundle to its multiple of the
flat cycle $c$ in $H^2(U,\Lambda)$.  Both common
invariant cycles lie in $N_q\Lambda$, so this
multiple is zero.  A holomorphic trivializing
gauge over $U$ need not extend at the missing
axes.  The proof thus neither identifies
arbitrary descent data with monomial cocycles
nor asserts a meromorphic comparison of
arbitrary analytic fillings.

\item We finally compute the actual affine class rather
than infer it from these linear monodromies.  Retain
the cusp functions of \ref{push:src:C.3.5},
\[
 q=e^{2\pi i s},\qquad
 C_0(q)=\beta^{\rm part}+s,\qquad
 \ell=-s/12+\kappa_0(q),
 \qquad D(q)=C_0(q)-12\kappa_0(q).
\]
The functions $\mu,C_0,\kappa_0$ are single-valued
and holomorphic at $q=0$.  The actual cusp unit
representative of the even-shift translation is
\[
 \bar b_{h,c}(q)=
 \left(-\frac h2\mu(q),\,
       -\frac h{12}(c+D(q))-\frac{h^2}{24}\right).
\]
On the logarithmic parameter cover define
\[
 \sigma(q,c)=
 \left(-\frac c2\mu(q),\,
       -\frac{c^2}{24}-\frac{cD(q)}{12}\right).
\]
Direct subtraction gives the exact identity
\[
 \sigma(q,c+h)-\sigma(q,c)=\bar b_{h,c}(q).
\]
In particular it is equivariant for the deck
transformation $c\mapsto c-12$ corresponding to
a positive $r$ loop, and it is single-valued in
$q$.  The original cusp has no extra affine
$q$ translation ($l_0=0$).  Thus $\sigma$ descends
to a holomorphic section of the actual affine
torus family.  Equivalently its $q$ and $r$
affine lifts commute exactly: the $r$ translation
$\bar b_{-12,c}$ is independent of the logarithm
branch $s$.  The resulting integral commutator
representing the Chern class is zero already in
$\Lambda$, before passage to coinvariants.

\item Multiplication by the exponential of $\sigma$
is a torus-unit automorphism of every cusp fan
chart for a fixed nonzero parameter.  Since its
coefficients are holomorphic in the full $q$
disc, this comparison includes the prescribed
Mumford cusp fibre for $r\ne0$.  The quadratic
dependence on $c=-(12/(2\pi i))\log r$ is retained:
it supplies no extension across $r=0$ and no
global comparison with the finite-place models.
Changing coherent closings by a scalar depending
only on the parameter has zero $q$ commutator
and leaves this Chern-class conclusion unchanged.
The explicit polynomial identity is checked in
\path{code/ab_push_hopf_winding.py}.
\end{enumerate}
\end{proof}

\begin{proposition}[an explicit parabolic family and its nonextending exceptional correction]\label{push:src:B.4.10}
\StatementRecord{push:src:B.4.10}{proved}{Historical events: Claude Fable 5.2 (Anthropic), effort high, 16 November 2026: carried as background without fault found. GPT-6 Astra (OpenAI), ChatGPT Pro mode, effort Pro, S6\_09\_PAPERS, 17 November 2026: targeted proof check. Also checked in this writing. Pre-continuation writing events: GPT-6 Astra (OpenAI), ChatGPT Pro mode, effort Pro, S6\_10\_PAPERS, 18 November 2026: targeted proof check; used and rechecked in this run. Continuation metadata events: GPT-6 Astra (OpenAI), ChatGPT Pro mode, effort Pro, S6\_10\_PAPERS continuation, 19 November 2026: used and rechecked in this run; status-vocabulary reconciliation; metadata reconciliation only; no new mathematical check. \emph{Pre-continuation recorded proof scope.} Re-derived moving Kato embedding and uniform compact containment, both normal frames, convergent regular gauge H and developing product G, weighted n[-s] exceptional correction and its nonextension, and degree-eight quadric/parabolic normal product with positive triple section. No compact third component or global sphere filling inferred. Closed exact Kato/normal-line source gates: Istrati et al Section2 and Proposition5.6pp1369-1370, Fujiki Section2equations1-8p4. The proof already supplies its own relative proper collar construction. Opened Teleman Introduction pp1750–1751, including the exact modification and spherical-shell gluing description. This closes the named background source gate; the moving Kato embedding, uniform collar containment and exceptional correction are checked by the previously recorded receiving proof.}{the hypotheses, notation and definitions stated here; \cref{push:src:B.4.8}; \citet[Section 2, pp.1351--1353]{IstratiEtAl2022Root}; \citet[Introduction, pp.1750--1751]{Teleman2010AB}; \citet[Example 9.10, p.1386]{IstratiEtAl2022Root}; \citet[Proposition 5.6, pp.1369--1370]{IstratiEtAl2022Root}; \citet[Section 1.2]{DlousskyTeleman2012AB}}

Fix an integer $n\ge1$ and holomorphic units $a(s),\kappa(s)$ on a
disc, with $0<|a(s)|<1$. After shrinking the disc there is a proper
holomorphic submersion $\mathcal K\to\Delta_s$, with an embedded smooth
proper elliptic family
\[
 \mathcal E=(\C^*\times\Delta_s)/
       ((v,s)\sim(a(s)v,s)).
\]
Its central surface is a parabolic Inoue surface $S_n$, with
$\pi_1(S_n)=\Z$, $b_2(S_n)=n$, and elliptic axis of square $-n$.
For $s\ne0$, the surface $\mathcal K_s$ is obtained from the diagonal
Hopf surface $H_{\kappa(s)s^n,a(s)}$ by $n$ successive blow-ups at one
axis point and its indicated infinitely near strict-axis points.
In a fixed annulus marking of $\mathcal E$, its normal line has factor
\[
 g_s(v)=\kappa(s)(a(s)v+s)^n.
\]
This line extends across $s=0$ and is isomorphic to the holomorphic
family with factor $\kappa(s)a(s)^n v^n$.
For constant $a,\kappa$ it is a constant family of degree $-n$ lines.

The relation with the normal line of the unmodified Hopf axis is
\[
 N_{\mathcal E_s/\mathcal K_s}
   =L_{\kappa(s)s^n}\otimes\OO_{\mathcal E_s}(-n[-s]),
       \qquad s\ne0.
\]
Here $n[-s]$ is the weighted total exceptional correction, not the
intersection with the reduced final exceptional locus. It has no
Cartier extension on $\mathcal E$, even after a finite base change or
the addition of a vertical component. Thus this construction does
not contradict \ref{push:src:B.4.8}: the marked divisor used to change the normal
line fails to extend.

For $n=8$, the degree $-8$ axis normal can be chosen to be the inverse
of $\OO_E(2)$ on any chosen smooth elliptic complete intersection
$E=(2,2)\subset\PP^3$. This supplies the normal-line matching with
$\operatorname{Bl}_E\PP^3$ described below, including one positive
triple section. It is an auxiliary construction and first-order
matching result. No attachment to the prescribed punctured
$S^6$ family, compact third component, or global smoothing is asserted.
\end{proposition}


\begin{proof}
\begin{enumerate}[label=\arabic*.,leftmargin=*]
\item We give the gluing domains as well as the local formulas. Let
$B_R\subset\C^2$ be a ball, and let
$\pi_n:\widehat B_R\to B_R$ be the following modification.
Blow up the origin, then successively the intersection of the newest
exceptional curve with the strict transform of the axis $x=0$.
In the last chart,
\[
 \pi_n(z,w)=(zw^n,w).
\]
Define the parameter-dependent embedding in that chart by
\[
 \sigma_s(z,w)=
       \bigl(\kappa(s)z,\ a(s)w+(1-a(s))s\bigr).
\]
Choose constants $A<\rho<1$ and $K>0$ with
$|a(s)|\le A$ and $|\kappa(s)|\le K$ on a closed smaller disc.
Shrink it further so that
$|a(s)w+(1-a(s))s|\le\rho R$ on $\overline B_R$.
Then
\[
 |\pi_n\sigma_s(z,w)|^2
 \le K^2R^2(\rho R)^{2n}+\rho^2R^2<R^2
\]
when $R$ is small enough. The last-chart coordinates of these images
are bounded, so this is compact containment of
$\sigma_s(\overline B_R)$ in $\widehat B_R$, uniformly in $s$.
Notice that no special smallness condition on the Picard parameter
$\kappa$ is required.

\item Remove the embedded ball and identify the outer and inner collars by
$\sigma_s\pi_n$. The outer collar avoids the exceptional locus, and
$\pi_n$ is a biholomorphism there; thus the collar identification is
a biholomorphism. Carry out the same construction over the parameter
disc. The embeddings and the identifications are holomorphic and
preserve $s$. Uniformly disjoint collars give a Hausdorff smooth
threefold, and the projection is a holomorphic submersion in every
chart. Uniform compact fundamental regions over compact parameter
sets prove properness. This is the usual Kato construction, with
the explicit moving embedding above; see
\citet[Section 2, pp.1351--1353]{IstratiEtAl2022Root}.
The general occurrence of deformations from Kato surfaces to
blown-up Hopf surfaces is classical; see
\citet[Introduction, pp.1750--1751]{Teleman2010AB}.
We use explicit formulas here to retain the elliptic marking and
its exceptional correction.

\item The contracted map $\pi_n\sigma_s$ fixes $(0,s)$.
In the coordinate $v=w-s$ it becomes
\[
 f_s(x,v)=
    \bigl(\kappa(s)x(a(s)v+s)^n,\ a(s)v\bigr).
\]
At $s=0$ this is the parabolic monomial germ with exponent matrix
$\left(\begin{smallmatrix}1&n\\0&1\end{smallmatrix}\right)$.
The chosen modification is exactly the $n$-fold modification in
\citet[Example 9.10, p.1386]{IstratiEtAl2022Root}.
Consequently the central surface is the stated parabolic Inoue
surface; its elliptic axis is disjoint from the rational cycle.
The construction gives $\pi_1=\Z$ and $b_2=n$, or equivalently one
can use the cited description of this surface. These are invariants
of the auxiliary surface, not of a proposed $S^6$ central fibre.

\item The strict axis is preserved in the collar gluing. Its coordinate
$v$ is multiplied by $a(s)$. To see its full quotient, the axis disc
in the blowdown coordinates is
\[
 D_s=\{v:|v+s|<R\}.
\]
For small $s$, it contains $0$ and satisfies $a(s)D_s\Subset D_s$.
Its dilates have union $\C$ and shrinking intersection $\{0\}$.
Thus $D_s\setminus a(s)D_s$, with the two boundary collars glued,
is a fundamental annulus for $\C^*/a(s)^\Z$.
This proves the exact elliptic marking throughout the parameter disc.

\item There are two normal frames, whose distinction is essential.
In the last blow-up chart the collar map $\sigma_s\pi_n$ has factor
\[
 g_z(v)=\kappa(s)(v+s)^n.
\]
On a fixed annulus avoiding $v=0$, shrink $s$ so that both
$v+s$ and $a(s)v+s$ are nonzero. The blowdown normal coordinate is
$x=z(v+s)^n$, and in that frame the factor is
\[
 g_x(v)=\kappa(s)(a(s)v+s)^n.
\]
The change between the frames is a nonzero gauge on these annuli.
The central factor is $\kappa(0)a(0)^n v^n$ and has degree $-n$,
with the factor-of-automorphy convention of
\citet[Proposition 5.6, pp.1369--1370]{IstratiEtAl2022Root}.

\item An explicit regular gauge verifies the whole family. On a fixed
annulus and its overlap with its $a(s)$-translate put
\[
 H_s(v)=\prod_{j\ge0}\left(1+\frac{a(s)^j s}{v}\right)^n.
\]
The product converges locally uniformly in $(v,s)$ and is nonzero
after shrinking the parameter disc. It satisfies
\[
 \frac{H_s(a(s)v)}{H_s(v)}
      =\left(1+\frac{s}{a(s)v}\right)^n.
\]
Replacing $x$ by $x/H_s(v)$ changes $g_x$ to
$\kappa(s)a(s)^n v^n$.
The collar maps are linear in the normal coordinate, so the same
calculation identifies an actual tubular line neighbourhood, not
only its first infinitesimal normal bundle.

\item For completeness, the generic Hopf blowdown also has an explicit
developing map. Fix $s\ne0$ and set
\[
 G_s(v)=\prod_{j\ge1}\left(1+\frac{a(s)^j v}{s}\right)^n,
       \qquad
 \Psi_s(x,v)=(xG_s(v),v).
\]
This product is entire in $v$, converges locally uniformly, and obeys
\[
 G_s(a(s)v)=\frac{G_s(v)}{(1+a(s)v/s)^n},\qquad
 \Psi_s f_s=\operatorname{diag}(\kappa(s)s^n,a(s))\Psi_s.
\]
For sufficiently small nonzero $s$, both diagonal multipliers have
absolute value less than one. Compose $\Psi_s$ with $\pi_n$ on a
Kato fundamental region. The displayed equivariance glues it to a
holomorphic map to $H_{\kappa(s)s^n,a(s)}$.
It never reaches the excluded origin: that would require
$x=v=0$, which is the fixed point $(0,s)$ removed with the inner ball.
All zero-factor curves of $G_s$ have
$v=-s/a(s)^j$, so they lie in the one axis orbit $[-s]$.
The original modification also contracts over $v=-s$.
At that representative,
\[
 G_s(-s)=\prod_{j\ge1}(1-a(s)^j)^n\ne0.
\]
Hence locally the map is precisely $\pi_n$ followed by a coordinate
unit. Away from that orbit it is a local biholomorphism of degree
one. Compactness makes it a proper modification, and the displayed
local model identifies it as the stated $n$ successive blow-ups.
This is a punctured-parameter construction: no meromorphic extension
of $G_s$ through $s=0$ is being claimed.

\item Each blow-up of a point on the strict elliptic curve subtracts that
point from its normal line. Following all $n$ stages gives
\[
 N_s=L_{\kappa(s)s^n}\otimes\OO_{\mathcal E_s}(-n[-s]).
\]
The multiplicity $n$ records total exceptional transforms.
The reduced final exceptional locus meets the strict axis only once.
One may check the formula directly: the standard theta factor
$\theta_a(av)=-v^{-1}\theta_a(v)$ shows that $\OO([p])$
has factor $-p/v$. At $p=-s$ the tensor product above has factor
\[
 \kappa(s)s^n(v/s)^n=\kappa(s)v^n,
\]
which is equivalent to $\kappa(s)a(s)^n v^n$ by a nonzero monomial
gauge. This checks the signs and degrees independently of the
successive-blow-up formula.

\item The divisor $n[-s]$ cannot extend. For constant $a$, choose two
distinct classes $[v_1],[v_2]\in\C^*/a^\Z$.
The sequences $s_j=-a^jv_i$ tend to zero, whereas the corresponding
divisors are constantly $n[v_i]$ and have different limits in
$\operatorname{Sym}^n(E_a)$.
For varying $a(s)$, choose the holomorphic logarithm
$\tau_a(s)=\log a(s)/(2\pi i)$ on the disc.
The Abel--Jacobi coordinate of $n[-s]$ is
$n\log(-s)/(2\pi i)$ modulo $\Z+\tau_a(s)\Z$.
On a positive parameter loop it winds $n$ times around a primitive
cycle in the smooth elliptic family. An extending section over the
disc would have zero winding. A finite base change of order $N$
multiplies this nonzero integer by $N$.
Adding a vertical elliptic component does not help: remove that
Cartier component, and the remaining proper horizontal divisor of
degree $n$ would be finite over the disc and define a continuous
section of the relative symmetric power, which has just been ruled
out. Nonproper exceptional divisors in deformations of Kato surfaces
are a classical phenomenon, studied by
\citet[Section 1.2]{DlousskyTeleman2012AB}; the calculation here
identifies the precise weighted divisor in this marked family.

\item This also explains the scope of ordinary elementary modifications.
If a degree-zero line $L_m$ from \ref{push:src:B.4.8} is changed by point blow-ups,
then its new normal is $L_m\otimes\OO(-D_{\rm exc})$.
If both that normal and the weighted Cartier correction extended on
a semistable good-reduction interface, their tensor product would
extend $L_m$, contradicting \ref{push:src:B.4.8}. Drifting translations cannot
remove both requirements. The present Kato family escapes exactly
because $D_{\rm exc}=n[-s]$ has no such extension.

\item Finally let $E=(2,2)\subset\PP^3$ be smooth and set
$L=\OO_E(2)$, of degree eight. Its normal in $\PP^3$ is
$L\oplus L$. For $A=\operatorname{Bl}_E\PP^3$, the exceptional
divisor is $D=E\times\PP^1$ and, in the lines convention,
\[
 N_{D/A}=L\boxtimes\OO_{\PP^1}(-1).
\]
A bare Hopf-product cap has normal $M\boxtimes\OO$ with
$\deg M=0$. The required triple-divisor line would then be
$(L\otimes M)^{-1}\boxtimes\OO(1)$, which has no nonzero section
because its elliptic degree is $-8$. Thus no effective triple divisor
repairs that particular bare interface.

\item For the parabolic surface with $n=8$, changing $\kappa$ varies the
flat factor through all of $\operatorname{Pic}^0(E)$:
constant factors are parametrized by $\C^*/a^\Z$.
Choose its axis normal to equal $L^{-1}$.
In $B=S_8\times\PP^1$ the divisor $D=E\times\PP^1$ then has normal
$L^{-1}\boxtimes\OO$. With triple section $T=E\times\{p\}$,
\[
 N_{D/A}\otimes N_{D/B}\otimes\OO_D(T)
   =(L\boxtimes\OO(-1))\otimes
     (L^{-1}\boxtimes\OO)\otimes(\OO\boxtimes\OO(1))
   \simeq\OO_D.
\]
This proves precisely the first-order matching claimed.
\end{enumerate}
\end{proof}


\paragraph{Auxiliary normal matching and the actual filling obstruction}
Choose a smooth quadric $Q\subset\PP^3$ containing $E$.
Its strict transform in $A$ is still $Q\simeq\PP^1\times\PP^1$:
blowing up its Cartier divisor $E$ does not change $Q$.
Its normal is
$\OO_Q(2,2)(-E)\simeq\OO_Q$, and it meets $D$ in a section.
The matching surface $S_8\times\{p\}$ in $B$ also has trivial
normal. A third threefold would therefore need divisors $Q$ and
$S_8$ meeting along $E$, with normals
$\OO_Q(-E)$ and $\OO_{S_8}(-E)$ respectively.
These are exactly the normal conditions for the central surface
$Q\cup_E S_8$ in a semistable surface smoothing.
The surface smoothing is constructed in \ref{push:src:B.4.11}--\ref{push:src:B.4.12}.
Its topological Euler number is $4+8-2\chi(E)=12$ and its
holomorphic Euler number is one. Its total space over an open
disc is not a compact third threefold, so that construction
does not supply a compact threefold assembly.

A full assembly must also preserve the required coherent cohomology
from \ref{push:src:B.4.5}; a lone parabolic component attached to rational components
can lose its $H^1(\OO)$ through the elliptic overlap.
Additional conductor surfaces meeting $E$ would contribute further
triple divisors and must be included in the normal equation.
Nor is a local change of the cusp parameter automatically a permitted
global base change of the prescribed family.
These auxiliary matching statements do not produce a filling
of the actual family. \ref{push:src:B.4.33} proves that no such proper analytic
filling exists after any finite parameter cover.


\Uses{the hypotheses, notation and definitions stated here; \cref{push:src:B.4.11}; \cref{push:src:B.4.12}; \cref{push:src:B.4.5}; \cref{push:src:B.4.33}}


\begin{proposition}[globalization at the contracted parabolic elliptic axis]\label{push:src:B.4.11}
\StatementRecord{push:src:B.4.11}{proved}{Historical events: Claude Fable 5.2 (Anthropic), effort high, 16 November 2026: carried as background without fault found. GPT-6 Astra (OpenAI), ChatGPT Pro mode, effort Pro, S6\_09\_PAPERS, 17 November 2026: targeted proof check. Also checked in this writing. Pre-continuation writing events: GPT-6 Astra (OpenAI), ChatGPT Pro mode, effort Pro, S6\_10\_PAPERS, 18 November 2026: targeted proof check. Continuation metadata events: GPT-6 Astra (OpenAI), ChatGPT Pro mode, effort Pro, S6\_10\_PAPERS continuation, 19 November 2026: used and rechecked in this run; status-vocabulary reconciliation; metadata reconciliation only; no new mathematical check. \emph{Pre-continuation recorded proof scope.} Re-derived exact degree-eight elliptic cone contraction, logarithmic tangent Atiyah extension and formal-neighbourhood H1, equality of pushed tangent sheaves via lifted flows, nonzero period Kodaira-Spencer map, Leray H2-vanishing and local-to-global Cech obstruction calculation. Grauert formal-functions and Manetti analytic-effectivity source gates remain separately tracked. Closed Manetti analytic-effectivity gate by visually reading the original Lemma1p93 for compact reduced spaces with isolated singularities; opened Grauert1960 formal-functions HauptsatzIIap59 and DT2020 Proposition3.3p11 at exact scopes. The receiving proof explicitly verifies H2(ThetaY)=0 and isolatedness.}{the hypotheses, notation and definitions stated here; \cref{push:src:B.4.10}; \citep[Section~6, Hauptsatz~IIa, p.~59]{Grauert1960D}; \citet[Lemma~1, p.~93]{Manetti1991D}; \citet[Section~3.2, Proposition~3.3 and the discussion preceding
Proposition~3.4]{DlousskyTeleman2020D}}

Let $S=S_8(a,\kappa)$ be the central parabolic surface in \ref{push:src:B.4.10},
let $E\subset S$ be its elliptic axis, and put
$N=N_{E/S}$, so $\deg N=-8$.
The contraction $\pi:S\to Y$ of $E$ is a compact normal surface
with a single isolated singular point $p$.
Its germ at $p$ is the vertex of the affine cone over
$(E,N^{-1})$.  In particular,
\[
 H^2(Y,\Theta_Y)=0.
\]
The restriction morphism from a semiuniversal deformation of $Y$
to a semiuniversal deformation of $(Y,p)$ is regular.
Consequently every prescribed analytic one-parameter deformation
of this isolated singularity is induced, after shrinking its
parameter disc, by a proper deformation of the compact surface $Y$.
If the prescribed local deformation smooths the singularity, the
global deformation has smooth fibres away from the central value.

This concerns globalization of a specified local deformation of
an auxiliary compact surface.  It does not assert an attachment to
the punctured $S^6$ family or a compactification of the resulting
threefold over the open parameter disc.
\end{proposition}


\begin{proof}
\begin{enumerate}[label=\arabic*.,leftmargin=*]
\item We first retain the analytic neighbourhood, not merely its normal
line.  At the central parameter the Kato gluing of \ref{push:src:B.4.10} is
\[
 (x,v)\longmapsto(\kappa a^8xv^8,av).
\]
It is exactly linear in $x$.  Thus a neighbourhood of $E$ is a
neighbourhood of the zero section in the negative line bundle
$N$, with the indicated factor of automorphy.
Its contraction is the cone associated with the section ring
\[
 R(E,N^{-1})=\bigoplus_{m\ge0}H^0(E,N^{-m}).
\]
A line bundle of degree eight on an elliptic curve is very ample
and normally generated.  The cone is normal, its only singularity
is its vertex, and the blow-up of the vertex maximal ideal is
$\operatorname{Tot}(N)$ near its zero section.
One can see the last assertion directly: the degree-one sections
generate the section ring and have no common zero on $E$;
their ratios on the blow-up recover the projective embedding of
$E$, and their common radial coordinate is the coordinate in $N$.
In particular, with $I=\mathcal O_S(-E)$ and $\mathfrak m_p$ the
vertex maximal ideal,
\[
 \mathfrak m_p\mathcal O_S=I.
\]
The local contraction glues to the identity on $S\setminus E$,
giving the stated proper contraction of the compact surface.

\item Write $C$ for the disjoint rational cycle of the parabolic surface.
The meromorphic two-form
$d\log x\wedge d\log v$ is invariant under the displayed
monomial gluing.  In the toric blow-up charts it has simple poles
along precisely the axis and the rational cycle, and no zeros.
Consequently
\[
 K_S=\mathcal O_S(-E-C).
\]
Also $b_1(S)=1$, by \ref{push:src:B.4.10}.
Put $\mathcal F=\Theta_S(-\log E)$, the locally free sheaf of
vector fields tangent to $E$.
Serre duality and the displayed canonical divisor give
\[
 H^2(S,\mathcal F)^*
 =H^0\bigl(S,\Omega^1_S(\log E)\otimes K_S\bigr)=0.
\]
Indeed there are sheaf injections
\[
 \Omega^1_S(\log E)(-E-C)
       \lhook\joinrel\longrightarrow\Omega^1_S(-C)
       \lhook\joinrel\longrightarrow\Omega^1_S.
\]
Locally the first sends the logarithmic generators
$dx/x,dy$, multiplied by $x$, to the ordinary forms $dx,xdy$.
Finally $H^0(S,\Omega^1_S)=0$ follows already from $b_1(S)=1$.
For completeness, a holomorphic one-form $\alpha$ on a compact
complex surface is closed: Stokes' theorem applied to
$\alpha\wedge\overline{d\alpha}$ makes the integral of the
positive volume form associated with
$d\alpha\wedge\overline{d\alpha}$ vanish.
If $\alpha\ne0$, its real and imaginary parts give independent
real de Rham classes.  An exact real linear combination would
be the differential of a global real pluriharmonic function,
which is constant on the compact surface.  This would give
$b_1(S)\ge2$, a contradiction.

\item We compute the one possible local higher direct image explicitly.
The restriction of $\mathcal F$ to $E$ is the Atiyah extension
of $N$:
\[
 0\longrightarrow\mathcal O_E
 \longrightarrow\mathcal A_N:=\mathcal F|_E
 \longrightarrow\Theta_E\longrightarrow0.
\]
In normal coordinates related by $x_j=g_{ij}(v)x_i$, the
vertical Euler field is $x_i\partial_{x_i}$, and the changes of
horizontal lifts have cocycle $d\log g_{ij}$.
Thus this extension has class $c_1(N)$ in
$H^1(E,\Omega^1_E)$, which is nonzero because $\deg N=-8$.
Since $\Theta_E\simeq\mathcal O_E$, the connecting map
$H^0(E,\Theta_E)\to H^1(E,\mathcal O_E)$ is an isomorphism.
The same cohomology sequence therefore gives the natural
isomorphism
\[
 H^1(E,\mathcal A_N)\xrightarrow{\ \sim\ }
 H^1(E,\Theta_E)\simeq\C.
\]

\item The successive quotients for the formal neighbourhood of $E$ are
\[
 I^m\mathcal F/I^{m+1}\mathcal F
     =\mathcal A_N\otimes N^{-m}\qquad(m\ge0).
\]
For $m\ge1$ both line bundles in the induced extension have
degree $8m>0$ on the elliptic curve, so their first cohomology
vanishes.  It follows that all restriction maps
\[
 H^1\bigl(E,\mathcal F/I^{m+1}\mathcal F\bigr)
 \longrightarrow H^1(E,\mathcal F/I^m\mathcal F)
 \quad(m\ge1)
\]
are isomorphisms.  Their common value is
$H^1(E,\mathcal A_N)\simeq\C$.
The analytic theorem on formal functions
\citep[Section~6, Hauptsatz~IIa, p.~59]{Grauert1960D}
identifies this inverse limit with the completed stalk of
$R^1\pi_*\mathcal F$ at $p$.
This coherent sheaf is supported at the single point $p$,
so its stalk is a finite-length module and equals its completion.
We have proved
\[
 R^1\pi_*\mathcal F\simeq\C_p,
 \qquad
 H^0(Y,R^1\pi_*\mathcal F)
    \xrightarrow{\ \sim\ }H^1(E,\Theta_E).
\]
The second map is the natural restriction followed by the
Atiyah-extension map just computed, not a dimension count alone.

\item There is also an equality of sheaves
\[
 \pi_*\mathcal F=\Theta_Y.
\]
To check it at $p$, a holomorphic derivation of the cone
integrates locally to a holomorphic flow: embed the cone in a
polydisc, lift its coefficient functions, and use the ordinary
holomorphic vector-field flow, which preserves the defining ideal.
The flow preserves the unique singular point, hence its maximal
ideal.  Functoriality of the blow-up of that ideal lifts the flow
to the resolution, preserving the exceptional curve.
Its derivative is therefore a section of $\mathcal F$.
Conversely a section of $\mathcal F$ differentiates pullbacks
of holomorphic functions on $Y$ and descends to a derivation,
because $\pi_*\mathcal O_S=\mathcal O_Y$.
Both operations agree off $E$ and are inverse there; uniqueness
extends the equality over $E$.

\item The global-to-local cohomology map is surjective for a concrete
reason.  Write $a=e^{2\pi i\tau}$ with $\operatorname{Im}\tau>0$,
and vary $\tau$ to $\tau+u$ in the explicit Kato construction,
keeping its translation parameter zero.
The same uniform collar argument as in \ref{push:src:B.4.10} gives a proper
holomorphic deformation of the pair $(S,E)$.
Its Kodaira--Spencer class belongs to $H^1(S,\mathcal F)$.
Under the natural composite
\[
 H^1(S,\mathcal F)\longrightarrow
 H^0(Y,R^1\pi_*\mathcal F)
 \xrightarrow{\ \sim\ }H^1(E,\Theta_E)
\]
it is the Kodaira--Spencer class of
$\C/(\Z+(\tau+u)\Z)$, which is nonzero.
For example, differentiating its second deck transformation
$z\mapsto z+\tau+u$ gives the usual nonzero period cocycle.
This argument works also at the elliptic curves with extra
automorphisms; it does not differentiate the coarse $j$-invariant.
The target is one dimensional, so the first arrow is onto.

\item The degree-two part of the Leray sequence is
\[
 H^1(S,\mathcal F)\longrightarrow
 H^0(Y,R^1\pi_*\mathcal F)\longrightarrow
 H^2(Y,\pi_*\mathcal F)\longrightarrow H^2(S,\mathcal F).
\]
The first map is surjective and the last group is zero.
Using $\pi_*\mathcal F=\Theta_Y$ proves
$H^2(Y,\Theta_Y)=0$.

\item We spell out the local-to-global criterion used in the final
step, since the rational-cycle contraction theorem does not
apply to this elliptic contraction.
For a compact complex space with isolated singularities,
$H^2(\Theta)=0$ makes the restriction from global to local
deformations formally smooth.
Here is the obstruction argument for the present single
singularity.  Take a finite Stein cover with one small
neighbourhood of $p$ and all other members smooth.
For a small extension of Artinian bases, prescribe the lifted
deformation on the singular member and take smooth local
lifts on the other members.  Their pairwise identifications
lift on the smooth Stein intersections.  The failure of the
cocycle identity on triple intersections is a Cech two-cocycle
in $\Theta_Y$ tensored with the square-zero kernel.
Its class vanishes by the proved $H^2$-vanishing; correcting
the identifications by a one-cochain glues the desired lift.
This proves the relative infinitesimal lifting property.

\item The analytic semiuniversal deformation theory turns this
property into regularity of the induced morphism of analytic
base germs.  The general criterion is
\citet[Lemma~1, p.~93]{Manetti1991D}, stated for reduced compact
complex spaces with isolated singularities and proved by the
local-to-global spectral sequence of the cotangent complex.
It is also used in
\citet[Section~3.2, Proposition~3.3 and the discussion preceding
Proposition~3.4]{DlousskyTeleman2020D}.
There is no hypothesis about the kind of exceptional curve
in this criterion.  A regular morphism of analytic
germs is locally a projection from a product with a smooth
germ.  Consequently a specified analytic disc in the local
deformation base lifts to the global base, without replacing
the prescribed local deformation by a different smoothing
component.
The global family is proper by the compact deformation
construction.  Away from $p$ its central fibre is smooth,
and smoothness is open; compactness permits one common
shrinking of the disc.  A local smoothing at $p$ therefore
gives smooth global noncentral fibres, as asserted.
\end{enumerate}
\end{proof}

\begin{proposition}[a proper quadric--parabolic surface smoothing]\label{push:src:B.4.12}
\StatementRecord{push:src:B.4.12}{proved}{Historical events: Claude Fable 5.2 (Anthropic), effort high, 16 November 2026: carried as background without fault found. GPT-6 Astra (OpenAI), ChatGPT Pro mode, effort Pro, S6\_09\_PAPERS, 17 November 2026: targeted proof check. Also checked in this writing. Pre-continuation writing events: GPT-6 Astra (OpenAI), ChatGPT Pro mode, effort Pro, S6\_10\_PAPERS, 18 November 2026: targeted proof check. Continuation metadata events: GPT-6 Astra (OpenAI), ChatGPT Pro mode, effort Pro, S6\_10\_PAPERS continuation, 19 November 2026: used and rechecked in this run; status-vocabulary reconciliation; metadata reconciliation only; no new mathematical check. \emph{Pre-continuation recorded proof scope.} Re-derived elliptic (2,2) embedding with prescribed degree-eight normal, full cone section ring and flat linear smoothing, blow-up equation t=xi e, compact globalization using B4.11, all three normal products and Thom map (a,b) to 2a+2b. Preserved the exact scope: proper surface smoothing over an open disc. Matched the precise degree-eight cone-smoothing construction against Pinkham Remark7.6(iii)p46 and Section9pp63-65; the actual anticanonical-cone family and normal bundles are derived in the receiving proof.}{the hypotheses, notation and definitions stated here; \cref{push:src:B.4.11}; \cref{push:src:B.4.10}; \citep[Remark~7.6(iii), p.~46; Section~9, pp.~63--65]{Pinkham1974Root}}

Let $S=S_8(a,\kappa)$ be as in \ref{push:src:B.4.11}, with elliptic axis $E$
and normal line $N$ of degree $-8$.
There is an embedding $E\hookrightarrow Q=\PP^1\times\PP^1$
as a smooth divisor of bidegree $(2,2)$ with
$N_{E/Q}\simeq N^{-1}$.
There is a proper flat map from a smooth complex threefold
\[
 \varpi:\mathcal V\longrightarrow\Delta
\]
with smooth noncentral fibres and reduced simple normal-crossings
central fibre exactly $S\cup_E Q$.  Its two normal lines are
\[
 N_{S/\mathcal V}=\OO_S(-E),\qquad
 N_{Q/\mathcal V}=K_Q=\OO_Q(-2,-2).
\]
The family is obtained by globalizing a specified cone smoothing
and blowing up its vertex; no finite base change is needed.

For the elliptic-quartic configuration of \ref{push:src:B.4.10}, the three
pairwise normal-line products, including the triple-curve factors,
can consequently all be made trivial using
\[
 A=\operatorname{Bl}_E\PP^3,\qquad
 B=S\times\PP^1,\qquad C=\mathcal V.
\]
Here $C$ is a proper family over an open disc and is not a compact
threefold.  This assertion supplies the indicated surface smoothing
and normal-line products only.  It does not construct a compact
three-component central fibre, a compatible global smoothing of
such a fibre, or a comparison with the prescribed punctured
$S^6$ family.
\end{proposition}


\begin{proof}
\begin{enumerate}[label=\arabic*.,leftmargin=*]
\item Put $L=N^{-1}$, of degree eight.  Multiplication by two on the
elliptic Picard variety is surjective, so choose a degree-four
line bundle $M$ with $M^{\otimes2}=L$.
Choose degree-two line bundles $M_1,M_2$ with
$M_1\otimes M_2=M$ and $M_1\not\simeq M_2$.
Each complete linear system is a basepoint-free double cover of
$\PP^1$.  The product of the two maps is an embedding in $Q$.
Indeed two distinct points identified by both maps would give
$M_1\simeq\OO_E(p+q)\simeq M_2$.
A common ramification point would similarly give
$M_1\simeq\OO_E(2p)\simeq M_2$.
Thus the product is injective and immersive; properness makes it a
closed embedding.  Its image has bidegree $(2,2)$, and
\[
 \OO_Q(E)|_E=(M_1\otimes M_2)^{\otimes2}=L.
\]
This also allows us to start with the fixed quadric and elliptic
quartic in \ref{push:src:B.4.10}: the normal line there was chosen to be the
inverse of this same $L$.

\item Use the anticanonical embedding of $Q$ in $\PP^8$ and write
\[
 R_Q=\bigoplus_{d\ge0}H^0(Q,\OO_Q(2d,2d)),\qquad
 Z=\operatorname{Spec}R_Q.
\]
The ring is generated in degree one and is normal.
One direct verification uses the bihomogeneous monomials:
every monomial of bidegree $(2d,2d)$ is a product of $d$
monomials of bidegree $(2,2)$, and the associated monoid is the
intersection of a rational polyhedral cone with its full lattice.
The punctured cone is the complement of the zero section in
$\OO_Q(-2,-2)$, hence is smooth.
Choose the section $e\in H^0(Q,\OO_Q(2,2))$ cutting out $E$.
The degree-one function $t=e$ on $Z$ defines a flat map
$Z\to\mathbb A^1$: its coordinate ring is torsion-free over
$\C[t]$, a principal ideal domain.
For every $d\ge0$, the sequence
\[
 0\longrightarrow\OO_Q(2d-2,2d-2)
 \xrightarrow{\ e\ }\OO_Q(2d,2d)
 \longrightarrow L^{\otimes d}\longrightarrow0
\]
has vanishing first cohomology in its left term.
For $d=0$ this follows from the Kunneth formula for
$\OO_Q(-2,-2)$; for $d\ge1$ it follows from the same formula
and nonnegative degrees.
Taking sections identifies
\[
 R_Q/(e)=\bigoplus_{d\ge0}H^0(E,L^{\otimes d}).
\]
The central surface is therefore exactly the affine cone over
$(E,L)$, including its scheme structure.
For $t\ne0$, projection from the cone vertex identifies the
fibre with $Q\setminus E$: over each point where $e$ is nonzero
there is exactly one radial coordinate on which $e$ has value $t$.
In particular this is an explicit smoothing of the isolated
singularity of \ref{push:src:B.4.11}.
It is the affine version of Pinkham's cone construction
\citep[Remark~7.6(iii), p.~46; Section~9, pp.~63--65]{Pinkham1974Root}.
In degree eight the quadric component is distinguished from the
other smoothing component; no replacement of it by that other
component is made here.

\item Blowing up the vertex of $Z$ gives
$\operatorname{Tot}(\OO_Q(-2,-2))=\operatorname{Tot}(K_Q)$.
This can be checked on the blow-up charts given by ratios of the
degree-one homogeneous coordinates; these recover a point of
$Q$, with the remaining common factor the tautological fibre
coordinate $\xi$.
The map to the parameter line is precisely
\[
 t=\xi e.
\]
Its central divisor is the union of the zero section $Q$ and
the restriction $\operatorname{Tot}(K_Q|_E)$, each with
multiplicity one.  Near $E$ the equation is $t=\xi z$, where
$z$ is a local defining function of $E$ in $Q$.
Thus the total space is smooth and the central divisor has
simple normal crossings.

\item Now contract $E$ in the compact surface $S$ as in \ref{push:src:B.4.11}.
The regular global-to-local deformation map in that proposition
lifts the specified analytic germ $Z\to\mathbb A^1$ to a proper
deformation $\mathcal Y\to\Delta$ of $Y=S/E$.
The lifted deformation agrees analytically with this cone
smoothing in a neighbourhood of its singular point.
Away from that neighbourhood the central surface is smooth;
openness of smoothness and properness allow a common shrinking
of $\Delta$ for which the family there is a submersion.
Blow up the single cone vertex in $\mathcal Y$.
The resulting map $\mathcal V\to\Delta$ remains proper.
Near the exceptional divisor its local model is the smooth
space just computed; elsewhere it is unchanged and smooth.
Its central strict transform is exactly $S$, because the
maximal-ideal blow-up of the degree-eight elliptic cone is its
given resolution $\operatorname{Tot}(N)$ near $E$.
The other component is $Q$, and their intersection is the original
marked elliptic curve.  The local equation $t=\xi z$ proves
flatness and reducedness near the intersection; at a smooth
point of either component $t$ is a local coordinate.
The noncentral fibres are the smooth compact fibres of
$\mathcal Y$.
This proves every asserted properness and smoothness property
for the surface family.

\item The exceptional zero section has normal $K_Q$ in its local
total-space model.  Since $S+Q$ is the principal divisor of $t$,
restriction to $S$ gives
\[
 \OO_{\mathcal V}(S)|_S
   =\OO_{\mathcal V}(-Q)|_S=\OO_S(-E).
\]
These are the two asserted normal lines.  They also show directly
that the two branch normals along $E$ multiply to the trivial line.

\item We record the precise extent of the three-piece matching.
Realize $Q$ as a smooth quadric in $\PP^3$ and $E$ as its
intersection with a second quadric.  Put $L=\OO_E(2)$.
Then $N_{E/\PP^3}=L\oplus L$, so the exceptional divisor
of $A=\operatorname{Bl}_E\PP^3$ is
$D_{AB}=E\times\PP^1$, with normal
$L\boxtimes\OO_{\PP^1}(-1)$.
The strict transform of $Q$ is still $Q$, because the centre is
Cartier on it, and its normal in $A$ is
\[
 N_{Q/A}=\OO_Q(2)\otimes\OO_Q(-E)=\OO_Q.
\]
It meets $D_{AB}$ in the constant section $E\times\{p\}$
corresponding to that member of the pencil of quadrics.
In $B=S\times\PP^1$, use the divisors
$D_{BA}=E\times\PP^1$ and $D_{BC}=S\times\{p\}$.
Their normals are $L^{-1}\boxtimes\OO$ and $\OO_S$.
In $C=\mathcal V$, use its central divisors $Q$ and $S$.
The three normal products with the triple contribution are then
\[
\begin{array}{c|c}
 \text{double surface}&
 N_{D/\text{first branch}}\otimes
 N_{D/\text{second branch}}\otimes\OO_D(\text{triple curve})\\ \hline
 E\times\PP^1&
 (L\boxtimes\OO(-1))\otimes(L^{-1}\boxtimes\OO)
       \otimes(\OO_E\boxtimes\OO(1))=\OO\\
 Q&\OO_Q\otimes K_Q\otimes\OO_Q(E)=\OO_Q\\
 S&\OO_S\otimes\OO_S(-E)\otimes\OO_S(E)=\OO_S.
\end{array}
\]
These are pairwise line-bundle identities.  They do not supply
a compact third component or a compatible global smoothing
section for a compact normal-crossings threefold.
In fact the third space constructed here maps surjectively and
properly to an open disc, so it cannot itself be compact.

\item For a useful topological check on the local model,
the affine smooth fibre $Q\setminus E$ has
\[
 H_1(Q\setminus E;\Z)=\Z/2,\qquad
 H_2(Q\setminus E;\Z)=\Z^3.
\]
Indeed the Thom sequence of the pair $(Q,Q\setminus E)$ has
intersection map $H_2(Q;\Z)\to H_0(E;\Z)$ given by
$(a,b)\mapsto2a+2b$.  Its cokernel is $\Z/2$ and its
kernel is $\Z$; the preceding term $H_1(E;\Z)=\Z^2$
injects into $H_2(Q\setminus E;\Z)$.
This exact extension of free groups splits abstractly.
The map $H_4(Q;\Z)\to H_2(E;\Z)$ is an isomorphism,
so $H_3(Q\setminus E;\Z)=0$, and its Euler number is four.
The Euler number is also obtained independently by projection
to either ruling of $Q$: away from the four branch values of
$E\to\PP^1$, the fibre is $\C^*$ and has Euler number zero;
at each branch value it is $\C$ and has Euler number one.
In particular no simply connected local Milnor fibre is being
assumed.  The compact noncentral surface has Euler number
$\chi(S)+\chi(Q)-2\chi(E)=8+4=12$ by the local
$\xi z=t$ retraction, not the Euler number of an $S^6$ fibre.
Exact polynomial, Hilbert-function, normal-product and Smith-form
checks appear in \texttt{code/root\_quadric\_smoothing.py}.
\end{enumerate}
\end{proof}

\begin{proposition}[the second period does not lift to a finite marked Hopf modification]\label{push:src:B.4.13}
\StatementRecord{push:src:B.4.13}{proved}{Historical events: Claude Fable 5.2 (Anthropic), effort high, 16 November 2026: carried as background without fault found. GPT-6 Astra (OpenAI), ChatGPT Pro mode, effort Pro, S6\_09\_PAPERS, 17 November 2026: targeted proof check. Also checked in this writing. Pre-continuation writing events: GPT-6 Astra (OpenAI), ChatGPT Pro mode, effort Pro, S6\_10\_PAPERS, 18 November 2026: targeted proof check. Continuation metadata events: GPT-6 Astra (OpenAI), ChatGPT Pro mode, effort Pro, S6\_10\_PAPERS continuation, 19 November 2026: used and rechecked in this run; status-vocabulary reconciliation; metadata reconciliation only; no new mathematical check. \emph{Pre-continuation recorded proof scope.} Re-derived finite exceptional-support invariance under the marked period and exact normal Picard morphism b to -d b; checked actual second cusp period has nonzero parameter valuation, so the generic translation has infinite order after each finite base change.}{the hypotheses, notation and definitions stated here; \cref{push:src:B.4.10}; \cref{push:src:A.1.2}; \cref{push:src:B.4.8}}

Fix $0<|\alpha|,|a|<1$, and put
\[
 H=(\C^2\setminus\{0\})/
       \langle(x,v)\mapsto(\alpha x,av)\rangle,
       \qquad E=\{x=0\}\simeq\C^*/a^\Z.
\]
Let $h$ be the automorphism induced by
$(x,v)\mapsto(c x,bv)$, where $b,c\in\C^*$.
Suppose $p:Y\to H$ is a finite composition of point blow-ups,
allowing infinitely near centres, and at least one point of $E$
is an image of its exceptional locus. If $[b]\in\C^*/a^\Z$ has
infinite order, there is no holomorphic automorphism of $Y$
which agrees, on the dense unmodified core, with the marked
automorphism $h$.

The same obstruction holds for a finite collection of such marked
modifications if the proposed period permutes the collection and
acts as $h$ on each marked core.
More generally, if a compact smooth surface contains a preserved
marked elliptic curve $E$ of nonzero normal degree $d$, every
holomorphic automorphism restricting to translation by $b$ on $E$
must satisfy $d\,b=0$ in $E$.
The latter assertion also excludes a finite cycle of such preserved
marked cap axes when the cycle composite is an infinite-order
translation. These assertions apply to the direct attempt
to lift the second actual cusp period to the finite Hopf
modifications in \ref{push:src:B.4.10}, for generic nonzero parameters, after every
finite base change. It concerns a holomorphic period action on
these full modifications. It does not exclude a meromorphic
correspondence, a different connected unfolding made from
noncompact pieces, or a cap present only over the central parameter.
\end{proposition}


\begin{proof}
\begin{enumerate}[label=\arabic*.,leftmargin=*]
\item Let $S\subset H$ be the image of the exceptional curves of $p$.
It is finite. Repeated infinitely near blow-ups change the curves
over a point but do not make $S$ infinite.
If $\widetilde h$ were a lift, the maps
\[
 p\widetilde h,\ hp:Y\longrightarrow H
\]
would agree on the dense core and hence everywhere, by the identity
principle for holomorphic maps to the separated complex surface $H$.
The image under $\widetilde h$ of a curve contracted by $p$ is
therefore again contracted by $p$. Applying the same argument to
$\widetilde h^{-1}$ proves that the total exceptional locus is
preserved. Thus $h(S)=S$.

\item The automorphism $h$ preserves $E$, where it is translation by
$[b]$. Consequently the nonempty finite set $S\cap E$ would be
invariant under that translation. The orbit of any one of its
points is finite, so for some $k>0$,
\[
 [b^k v]=[v],\qquad\text{equivalently } b^k\in a^\Z.
\]
This says exactly that $[b]$ is torsion, a contradiction.
In particular, no choice of multiplicities or infinitely near
centres changes the obstruction: only their nonempty finite
support on $E$ is needed.

\item For finitely many modifications $p_i:Y_i\to H$, a proposed deck
map has biholomorphisms
$Y_i\to Y_{\sigma(i)}$ satisfying the same core comparison.
The identity principle again gives equivariance everywhere.
The union of the finitely many exceptional image sets is
$h$-invariant; its intersection with $E$ gives the same contradiction.
Equivalently, around a cycle of length $\ell$ the map $h^\ell$
would lift to a single modification, while $[b^\ell]$ still has
infinite order.

\item For the more general assertion no blowdown is needed.
If an automorphism preserves $E$, its derivative in the normal
direction gives
\[
 t_b^*N_{E/S}\simeq N_{E/S}.
\]
For a line bundle $N$ of degree $d$ on an elliptic curve, the
homomorphism $b\mapsto t_b^*N\otimes N^{-1}$ is multiplication
by $-d$ under the usual identification $\operatorname{Pic}^0(E)=E$.
Indeed write $N=\OO_E(\sum_j n_j[p_j])$ with
$\sum_j n_j=d$; the Abel--Jacobi class of its translated divisor
minus the original one is $-d b$.
Thus preservation of $N$ forces $d b=0$.
Taking a power proves the finite-cycle version, since its composite
preserves one axis and its nonzero-degree normal line.
In particular a nonzero-degree parabolic axis cannot carry the
marked generic period translation even if no finite Hopf
modification is specified. The original finite-support argument
also covers modifications where that normal-degree test alone
would give no information.

\item We check the actual cusp rather than assume that torsion fails.
The periods of \ref{push:src:A.1.2} have the form
\[
 g_1:(Z,W)\longmapsto(q c_1 Z,c_2W),\qquad
 g_2:(Z,W)\longmapsto(c_3Z,r^{12}q c_4 W),
\]
where the $c_i$ are nonvanishing holomorphic units near the cusp.
Choose $W$ or $W^{-1}$ so that
$a=c_2^{\varepsilon}$ has modulus less than one,
$\varepsilon\in\{1,-1\}$.
This is possible at the actual marked cusp because its limiting
unit has modulus different from one, as computed in \ref{push:src:B.4.8}.
For $q\ne0$ small, the partial $g_1$ quotient of $(\C^*)^2$
is the dense core of the diagonal Hopf surface with multipliers
$(q c_1,a)$. The second period induces the above $h$, with
\[
 b(q,r)=(r^{12}q c_4(q,r))^{\varepsilon}.
\]
Fix a sufficiently small nonzero $q$. For each $k>0$ and
$j\in\Z$, the equation $b(q,r)^k=a(q)^j$ has a discrete
solution set in the punctured $r$-disc: its left-hand side has
a zero or pole of nonzero order at $r=0$, while the right-hand
side is a nonzero constant in $r$.
Hence outside a countable union of discrete sets, $[b(q,r)]$
has infinite order. Under a finite base change in $r$ the
nonzero order is multiplied by its ramification index and
the same conclusion holds. A finite base change in $q$ does
not affect this fixed-nonzero-$q$ argument.

\item To compare with \ref{push:src:B.4.10}, write $q c_1=\kappa(s)s^n$ after
the chosen finite cover in the cusp variable. Its generic Kato
fibre is precisely a finite Hopf modification with nonempty
axis centre $[-s]$. A holomorphic extension of the actual second
period to that whole fibre would therefore contradict the result
just proved. Allowing finitely many translated copies does not
remove the contradiction, by the finite-collection assertion.
\end{enumerate}
\end{proof}


\paragraph{What an infinite unfolding must additionally prove}
There is no corresponding obstruction to merely indexing translated
modifications by all integers and letting the deck map shift the
index. That construction alone has the wrong core. Indeed suppose
$p_j:Y_j\to H$ and
$\widetilde h_j:Y_j\to Y_{j+1}$ satisfy
$p_{j+1}\widetilde h_j=hp_j$.
On the disjoint union of their unmodified cores, the quotient by
the shift is identified with the single core $H^\circ$ by
\[
 [(j,y)]\longmapsto h^{-j}p_j(y).
\]
It is not the required quotient $H^\circ/\langle h\rangle$.
Thus a useful unfolding must have additional connected overlaps
and must recover the original second-period identifications there.

Nor can one simply glue several full compact Hopf modifications
along the same dense core and assume a Hausdorff space results.
If two embedded compact spaces in a Hausdorff space share an open
subset dense in both, their images, being closed, both equal its
closure. They coincide. A construction by genuinely distinct
compact copies with only that common core consequently cannot be
such a Hausdorff gluing.
These observations alone do not classify connected unfoldings
of noncompact pieces with a new boundary geometry. Analytic
transition maps, proper discontinuity and the retained punctured
quotient are not supplied by the auxiliary compact family in
\ref{push:src:B.4.10}. Independently of that restricted classification question,
\ref{push:src:B.4.33} rules out any such construction giving a proper filling
of the actual global family.


\Uses{the hypotheses, notation and definitions stated here; \cref{push:src:B.4.10}; \cref{push:src:B.4.33}}


\begin{proposition}[tangent weights and the compact-bundle obstruction reduction]\label{push:src:B.4.14}
\StatementRecord{push:src:B.4.14}{proved}{Historical events: Claude Fable 5.2 (Anthropic), effort high, 16 November 2026: carried as background without fault found. GPT-6 Astra (OpenAI), ChatGPT Pro mode, effort Pro, S6\_09\_PAPERS, 17 November 2026: targeted proof check. Also checked in this writing. Pre-continuation writing events: GPT-6 Astra (OpenAI), ChatGPT Pro mode, effort Pro, S6\_10\_PAPERS, 18 November 2026: targeted proof check. Continuation metadata events: GPT-6 Astra (OpenAI), ChatGPT Pro mode, effort Pro, S6\_10\_PAPERS continuation, 19 November 2026: used and rechecked in this run; status-vocabulary reconciliation; metadata reconciliation only; no new mathematical check. \emph{Pre-continuation recorded proof scope.} Re-derived vector-field functional equations for the Kato germ, all node and elliptic fixed-point contributions and character -8, equivariant Leray passage from nine to eight tangent H1 classes, and associated-bundle direct images/twisted H2 vanishing. Atiyah-Singer III original fixed-point theorem and Fujiki exact source gates remain separately tracked. Closed the finite-action Lefschetz source gate: original Atiyah-SingerIII scan Theorem4.6p566 and equations4.4-4.5p565 visually read. Verified that the proof applies it only to roots of unity and then identifies the finite Laurent character; Fujiki Theorem1.1pp1-2/equations1-8p4 also opened.}{the hypotheses, notation and definitions stated here; \cref{push:src:B.4.11}; \citet[Section~2, equations~(1)--(8)]{Fujiki2009D}; \citet[Theorem~1.1]{Fujiki2009D}; \citet[Theorem~(4.6), p.~566]{AtiyahSinger1968IIIE}}

For the explicit parabolic surface $S=S_8(a,\kappa)$ and its
elliptic contraction $\pi:S\to Y$ of \ref{push:src:B.4.11}, the genuine action
$\lambda\cdot(x,v)=(\lambda x,v)$ induces trivial representations
on the following cohomology groups:
\[
 H^0(S,\Theta_S(-\log E))=\C,\qquad
 H^1(S,\Theta_S(-\log E))=\C^9,
\]
\[
 H^0(Y,\Theta_Y)=\C,\qquad H^1(Y,\Theta_Y)=\C^8.
\]
The two second cohomology groups vanish.

Let $p:Z_L\to\PP^1$ be the compact holomorphic $Y$-bundle
associated to any holomorphic line bundle $L$ and this
$\C^*$ action.  Put
$\mathcal T=\Theta_{Z_L/\PP^1}$.
For every $b\in\PP^1$ one has
\[
 R^0p_*\mathcal T=\mathcal O_{\PP^1},\qquad
 R^1p_*\mathcal T=\mathcal O_{\PP^1}^{\oplus8},\qquad
 R^2p_*\mathcal T=0,
\]
and
\[
 H^2(Z_L,\mathcal T)=0,
 \qquad
 H^2\bigl(Z_L,\mathcal T\otimes
                   p^*\mathcal O_{\PP^1}(-b)\bigr)=0.
\]
These are cohomological reductions.  They do not, by themselves,
globalize a prescribed deformation along the nonisolated
singular section of $Z_L$ or construct a compact third component.
\end{proposition}


\begin{proof}
\begin{enumerate}[label=\arabic*.,leftmargin=*]
\item The action commutes with the central Kato gluing and hence
acts on the whole compact surface and its contraction.
Its toric coordinate description is also given in
\citet[Section~2, equations~(1)--(8)]{Fujiki2009D}.
For any $\lambda\ne1$, its fixed locus consists of the elliptic
curve $E$ and the eight nodes of the rational cycle.
To check that there are no further points, on the dense torus
an equality with a Kato deck iterate would require $a^jv=v$;
since $0<|a|<1$, this forces $j=0$, and then $\lambda=1$.
The action on each rational component is nontrivial and fixes
only its two endpoints.  In the unfolded fan the consecutive
rays are $(j,1)$ and $(j+1,1)$.
Their dual characters evaluate on the scaling cocharacter
$(1,0)$ as $-1,+1$.
Thus the two tangent weights at every cycle node are
$\lambda^{-1},\lambda$.

\item We also record a direct check of the space of global vector
fields, agreeing with \citet[Theorem~1.1]{Fujiki2009D}.
Write the contracted Kato germ as
\[
 F(x,v)=(Kxv^8,av),\qquad K=\kappa a^8.
\]
A global vector field restricts on the outer spherical annulus
to $P(x,v)\partial_x+Q(x,v)\partial_v$.
Hartogs extension gives holomorphic coefficients on the ball.
Compatibility with the collar gluing, followed by analytic
continuation, gives the exact germ identity
\[
 V(F)=DF\cdot V.
\]
Expand $Q(x,v)=\sum_{i\ge0}x^iQ_i(v)$.
The second component says
\[
 K^iv^{8i}Q_i(av)=aQ_i(v).
\]
For $i>0$ the lowest nonzero $v$-order on the two sides would
differ by $8i$, so $Q_i=0$.
For $i=0$, power-series coefficients give $Q_0(v)=cv$.
The constant-in-$x$ coefficient of the first component also
vanishes by the same lowest-order argument.
Its linear coefficient satisfies
$P_1(av)=P_1(v)+8c$.
Evaluation at zero gives $c=0$, and then $P_1$ is constant.
For $i>1$ the remaining equations are
\[
 P_i(v)=K^{i-1}v^{8(i-1)}P_i(av),
\]
so again $P_i=0$.
Consequently $H^0(S,\Theta_S)=\C\,x\partial_x$.
This Euler field is tangent to $E$ and has weight zero.

\item Set $\mathcal F=\Theta_S(-\log E)$.
\ref{push:src:B.4.11} proved $H^2(S,\mathcal F)=0$.
Apply the holomorphic Lefschetz formula to a root of unity
$\lambda\ne1$, using its finite cyclic group action.
The applicable original statement is
\citet[Theorem~(4.6), p.~566]{AtiyahSinger1968IIIE}; it requires
a compact complex manifold, not a K\"ahler manifold.
On $E$, the Atiyah extension $\mathcal F|_E$ has two
weight-zero terms and Chern character $2$; the normal line
has weight $\lambda$ and degree $-8$.
Since the tangent bundle of the elliptic curve has trivial
Todd class, the contribution of $E$ is
\[
 \int_E\frac{2}{1-\lambda^{-1}\exp(-c_1(N))}
   =-\frac{2\lambda^{-1}\deg N}
           {(1-\lambda^{-1})^2}
   =\frac{16\lambda}{(\lambda-1)^2}.
\]
At a cycle node $\mathcal F$ equals the ordinary tangent
bundle.  Its contribution is
\[
 \frac{\lambda+\lambda^{-1}}
      {(1-\lambda)(1-\lambda^{-1})}
   =-\frac{\lambda^2+1}{(\lambda-1)^2}.
\]
Adding the eight nodes gives the constant character
\[
 \sum_q(-1)^q\operatorname{Tr}
       (\lambda\mid H^q(S,\mathcal F))
 =\frac{16\lambda-8(\lambda^2+1)}{(\lambda-1)^2}
 =-8.
\]
Since $H^0$ is one weight-zero copy and $H^2=0$,
$H^1(S,\mathcal F)$ consists of nine trivial copies.
Equality on infinitely many roots of unity determines the
finite Laurent-polynomial character of the $\C^*$ representation.

\item Every map in the Leray sequence of \ref{push:src:B.4.11} is equivariant.
Its surjection
\[
 H^1(S,\mathcal F)\longrightarrow
 H^0(Y,R^1\pi_*\mathcal F)
       \simeq H^1(E,\Theta_E)=\C
\]
has a trivial one-dimensional target.
The kernel is $H^1(Y,\Theta_Y)$, so it is eight-dimensional
and trivial.  The equality in degree zero follows from
$\pi_*\mathcal F=\Theta_Y$; the second cohomology vanishing
was already proved in \ref{push:src:B.4.11}.

\item Finally form $Z_L$ by choosing line-bundle transition functions
and gluing products $Y\times U_i$ by their action on $Y$.
This is a holomorphic locally trivial bundle with compact
fibre over the compact curve $\PP^1$, hence is proper and
compact.  It has the section of isolated cone points as its
singular locus; that locus is one dimensional in the total
space.  On each product chart, coherent cohomology and base
change identify $R^qp_*\mathcal T$ with
$H^q(Y,\Theta_Y)\otimes\mathcal O_{U_i}$.
The gluing is exactly by the representations computed above.
Their triviality proves the displayed direct-image formulas.
The projection formula and the Leray sequence then give
\[
 H^2\bigl(Z_L,\mathcal T\otimes p^*M\bigr)
      \simeq H^1(\PP^1,M^{\oplus8})
\]
for $M=\mathcal O$ or $\mathcal O(-b)$, both of which have
zero first cohomology on $\PP^1$.
This proves the asserted reduction, including the twist
appropriate to retaining one whole fibre.

\item The singular locus here is not isolated, so the isolated
criterion of \ref{push:src:B.4.11} is not applicable to $Z_L$.
In particular its $H^2$-vanishing is not a proof of analytic
gluing of a prescribed local smoothing along the whole
singular section.  Compatibility of local lifts and their
analytic effectivity remain separate requirements.
The exact character and ray calculations are recorded in
\texttt{code/ab\_push\_kato\_weights.py}.
\end{enumerate}
\end{proof}

\subsection{A local insertion into the prescribed simple-torus family}

\begin{proposition}\label{push:src:B.4.15}\label{push:B415}
\StatementRecord{push:src:B.4.15}{proved}{Historical events: Claude Fable 5.2 (Anthropic), effort high, 16 November 2026: carried as background without fault found. GPT-6 Astra (OpenAI), ChatGPT Pro mode, effort Pro, S6\_09\_PAPERS, 17 November 2026: targeted proof check. Also checked in this writing. Pre-continuation writing events: GPT-6 Astra (OpenAI), ChatGPT Pro mode, effort Pro, S6\_10\_PAPERS, 18 November 2026: targeted proof check. Continuation metadata events: GPT-6 Astra (OpenAI), ChatGPT Pro mode, effort Pro, S6\_10\_PAPERS continuation, 19 November 2026: used and rechecked in this run; status-vocabulary reconciliation; metadata reconciliation only; no new mathematical check. \emph{Pre-continuation recorded proof scope.} Re-derived the actual semiabelian character action through all twelve boundary components, the two small-resolution charts z=xQ,y=Qs and x=zX,s=Xy, and the associated O(F) twist. Checked properness, flatness of the one-parameter map, component normal geometry and the exact comparison away from F. Joint flatness is not asserted.}{the hypotheses, notation and definitions stated here; \cref{push:src:B.3.1}; \citep{Cartan1952B_AB}; \citet[Definition~1.1]{Fujimoto1990AB}}

Let $q_0$ be a regular point of the elliptic base, and let
$\Delta_z$, $z=q-q_0$, be a simply connected disc avoiding
$0,1,\infty$. After $w=r^6$, take the standard twelve-cycle
model
\[
 \mathcal X\longrightarrow\Delta_z\times\Delta_r
\]
of the prescribed local two-torus family. There is a proper
replacement $\mathcal X'\to\Delta_z\times\Delta_r$, with
a specified isomorphism to $\mathcal X$ off $(z,r)=(0,0)$,
whose total fourfold is smooth and whose divisor $r=0$ is
reduced normal crossings. The composite map to $\Delta_r$
is flat.

The replacement has twelve new compact vertical normalization
components. Six are $\PP^1$-bundles over a scalar Hopf
surface; the other six are the blow-ups of those bundles
along two disjoint elliptic curves. Every new component
has fundamental group $\Z$, and therefore fails the
$\partial\bar\partial$ lemma. The punctured family is the
actual prescribed family, so its generic simplicity is
preserved.

This is a statement about a proper model over a bidisc,
flat over the degeneration parameter. It does not assert
joint flatness over the bidisc, a one-parameter degeneration
of compact surfaces with Hopf components, or a global
$S^6$ filling. All twelve component labels, and hence the
obstruction of \ref{push:src:B.3.1}, are unchanged.
\end{proposition}



\begin{proof}
\begin{enumerate}[label=\arabic*.,leftmargin=*]
\item \emph{The actual action.}
Use the rank-three period subsystem
\[
 \mathbf u(z)=(\tau(z),\mu(z)),\qquad
 \mathbf w=(1,0),\qquad \boldsymbol\delta=(0,1).
\]
It defines the analytic semiabelian group
\[
 G_z=\C^2/(\Z\mathbf u(z)+\Z\mathbf w+
                    \Z\boldsymbol\delta),
 \qquad
 1\longrightarrow\C^*_{\delta}\longrightarrow G_z
       \longrightarrow E_z\longrightarrow0,
 \quad E_z=\C/(\Z+\tau(z)\Z).
\]
This is the punctured degree-zero line bundle determined
by the actual $\mu(z)$. We make neither a product assumption
nor a torsion assumption on that line bundle. The formula
\[
 \chi_z(t)=
 \left[\frac{\log t}{2\pi i}\,\mathbf u(z)\right]
 \quad(t\in\C^*)
 \tag{B.4.15.1}
\]
is a single-valued holomorphic homomorphism: a logarithm
change adds an actual lattice period. It varies
holomorphically with $z$. The $r$-monodromy fixes
$\mathbf u$, and the $z$-disc has no period monodromy.

\item Here the action extends to the entire twelve-cycle model,
not just to its smooth fibres. Each unfolded normalization
is
\[
 R_z=G_z\mathbin{\times^{\C^*_{\delta}}}\PP^1
       =\PP_{E_z}(\OO\oplus L_z).
\]
Translation in $G_z$ preserves both conductor sections.
More explicitly, a class represented by $(a,b)$ acts
before the closing quotient by
\[
  (\xi,W)\longmapsto
            (\xi+a,e^{-2\pi i b}W).
 \tag{B.4.15.2}
\]
It commutes exactly with the closing period, which is
another translation of the same commutative group and
whose scalar factor contains $r^{12}$. The adjacent line
directions are paired dually in the local equation $xy=r$;
the two scalar factors cancel. Scalar multiplication also
preserves the one-dimensional toric subdivision producing
the twelve components. Thus (\ref{push:src:B.4.15}.2) extends over every
node, preserves each branch, and changes a branch equation
$x$ by a unit times $x$. This checks the required action
using the actual semiabelian marking of B.1.

\item The affine family itself admits the necessary local
punctured comparison. On
$U=\Delta_z\times\Delta_r^*$ its lattice local system
$\Lambda$ has $H^2(U,\Lambda)=0$, since $U$ retracts onto
a circle, also for local coefficients. Its relative Lie
algebra $\mathcal V$ is a coherent vector bundle. The
exponential sequence
\[
 0\longrightarrow\Lambda\longrightarrow\mathcal V
       \longrightarrow\mathcal T\longrightarrow0
\]
and Cartan's theorem B on the Stein base give
$H^1(U,\mathcal T)=0$ \citep{Cartan1952B_AB}.
Consequently the actual affine torsor has a holomorphic
section on $U$, which supplies the stated punctured
comparison with its Jacobian. No extension of this
section across $r=0$ is asserted.

\item \emph{The base blow-up and small resolutions.}
Blow up the base origin:
\[
 \beta:\widetilde B=\operatorname{Bl}_{(0,0)}
            (\Delta_z\times\Delta_r)
       \longrightarrow\Delta_z\times\Delta_r,
 \qquad F=\beta^{-1}(0,0)\cong\PP^1.
\]
In the chart $r=zs$, the pulled-back node is
$xy=zs$, with one further smooth elliptic coordinate.
Label the original components by $i\in\Z/12$.
At each node choose $x=0$ as the even-indexed branch,
and blow up $(x,z)$. This is globally the sum of the
pulled-back invertible ideal of the chosen central
component and the ideal $\OO_{\widetilde B}(-F)$.
The ideals are $G$-invariant. Their non-Cartier loci are
the twelve disjoint elliptic node loci, so the choices
at the six even components are compatible.

\item The two strict-transform charts are
\[
 \begin{aligned}
 z&=xQ,&y&=Qs,&r&=xQs,\\
 x&=zX,&s&=Xy,&r&=zXy.
 \end{aligned}
 \tag{B.4.15.3}
\]
They are smooth, and all three factors of $r$ occur once.
In the other base chart $z=r\zeta$, the node is $xy=r$
and the same ideal is $(x,r)=(x)$, so no change occurs.
Away from the indicated node loci it is also invertible.
The resulting total space is smooth and its $r$-central
divisor is reduced normal crossings.

\item For the component identifications, before resolution put
\[
 \begin{array}{ll}
 V_x=\{x=z=0\},&V_y=\{y=z=0\},\\
 H_x=\{x=s=0\},&H_y=\{y=s=0\}.
 \end{array}
\]
The $V$ branches are vertical over $F$, and the $H$
branches lie over the strict transform of $r=0$.
On $V_x$, the two maps in (\ref{push:src:B.4.15}.3) are
\[
 (Q,s)\longmapsto(y,s)=(Qs,s),\qquad
 (X,y)\longmapsto(y,s)=(y,Xy).
\]
These are exactly the blow-up charts of $(y,s)$.
Thus $V_x$ is blown up along its elliptic curve over
$y=s=0$; $V_y$ is unchanged. Similarly $H_y$ is blown
up along the corresponding elliptic curve and $H_x$
is unchanged. Each even vertical component is therefore
blown up at its two disjoint conductor curves, whereas
each odd vertical component is unchanged. The horizontal
parities are reversed.

\item \emph{The Cartier logarithm and properness.}
Set $M=\OO_{\widetilde B}(F)$. On a trivializing cover
write its transition units as $g_{ij}$. Glue the resolved
pullbacks by $\chi_z(g_{ij})$. These are genuine
automorphisms satisfying the cocycle identity.
This is the Cartier-logarithm construction of
\citet[Definition~1.1]{Fujimoto1990AB}; the explicit
semiabelian action above verifies its applicability
here, without assuming that the smooth two-tori are
principal elliptic bundles. The invariant ideals show
that the action lifts to the resolutions; equivalently
one may twist first and then resolve the descended ideals.

\item Over every base chart the glued space is a copy of the
proper resolved pullback. It is Hausdorff: two points
over the same base point are separated in such a chart,
and points over distinct base points are separated in
the base. Properness is local on the target, and
composition with the proper $\beta$ proves properness
over the original bidisc. The canonical divisor section
of $M$ is nonzero off $F$ and trivializes the twist there.
The base blow-up and the resolutions are also
isomorphisms off the origin. This gives the specified
isomorphism with the original family there, in particular
over every $r\ne0$.

\item The construction can be inserted into any larger model
which is standard near $q_0$: the preceding
trivialization supplies the overlap map outside the
origin. The cocharacter $\mathbf u$ is needed only on
this disc and is not extended around the special base
points. Since the local rings of the total space are
regular and $r$ is a nonzero divisor, they are
torsion-free, hence flat, over the one-dimensional
regular parameter ring. This proves flatness over $r$.
The new vertical threefolds rule out joint flatness over
the bidisc. Properness over the $r$-disc alone follows
only when this local replacement is inserted into a
model with compact elliptic base.

\item \emph{The compact vertical components.}
Before resolving, each vertical normalization is the
twist of $R_{q_0}\times F$ by
$M|_F=\OO_{\PP^1}(-1)$, namely
\[
 V_i^0=\OO_{\PP^1}(-1)^\times
                    \mathbin{\times^{\C^*}}R_{q_0}.
\]
Its elliptic quotient is a principal $E_{q_0}$-bundle
$H\to\PP^1$ with primitive Chern cycle
$\pm\mathbf u$. It is explicitly a scalar Hopf surface.
Indeed put
\[
 \tau_0=\tau(q_0),\qquad
 \eta=-1/\tau_0,\qquad a'=e^{2\pi i\eta},
 \qquad |a'|<1.
\]
After dividing the elliptic coordinate by $\tau_0$,
the projection of (\ref{push:src:B.4.15}.1) is
$[\log t/(2\pi i)]$ in $\C/(\Z+\eta\Z)$.
Its kernel is $(a')^\Z$. Consequently
\[
 H=\OO_{\PP^1}(-1)^\times/(a')^\Z
   =(\C^2\setminus\{0\})/\langle a'\id\rangle.
 \tag{B.4.15.4}
\]
The principal $G$-bundle
$\OO(-1)^\times\mathbin{\times^{\C^*}}G_{q_0}$
maps to $H$ with fibre $\C^*_{\delta}$.
Compactifying this principal $\C^*$-bundle gives an
actual projective bundle
\[
 V_i^0=\PP_H(\OO\oplus M_H)
\]
with its two global conductor sections. On each
elliptic fibre of $H\to\PP^1$, the line $M_H$ retains
the original degree-zero line $L_{q_0}$.

\item The two conductor sections over the distinguished point
$F\cap\operatorname{strict}(r=0)$ are disjoint elliptic
curves. The component calculation in (\ref{push:src:B.4.15}.3) gives
exactly the six unchanged bundles and six two-curve
blow-ups in the statement. A $\PP^1$-bundle has the
fundamental group of its base, and blowing up a smooth
complex codimension-two centre preserves fundamental
group. Hence every new normalization has
$\pi_1=\pi_1(H)=\Z$. A compact
$\partial\bar\partial$ manifold has even first Betti
number, so these normalizations fail that lemma.

\item Finally, the connected translation action fixes every
component label. The rotations $+4,-3,0$ of the
incompatible global assembly are unchanged. This proves
the precise local assertion and its stated limitation.
\end{enumerate}
\end{proof}


\paragraph{Scope}
For a generic arc $z=kr$ through the blown-up base
origin, the strict transform meets $F$ at an ordinary
point. Its central surface remains a cycle of ruled
surfaces over the elliptic curve. The non-K\"ahler
geometry in (\ref{push:src:B.4.15}.4) is carried by the variation
around the whole exceptional $\PP^1$. Therefore a
theorem only about one-parameter limits of simple
two-tori would not, by itself, exclude these vertical
threefold components.


\Uses{the hypotheses, notation and definitions stated here; \cref{push:src:B.4.15}}


\begin{proposition}[scope of the insertion operations]\label{push:src:B.4.16}
\label{push:B416}
\StatementRecord{push:src:B.4.16}{proved}{Historical events: Claude Fable 5.2 (Anthropic), effort high, 16 November 2026: carried as background without fault found. GPT-6 Astra (OpenAI), ChatGPT Pro mode, effort Pro, S6\_09\_PAPERS, 17 November 2026: targeted proof check. Also checked in this writing. Pre-continuation writing events: GPT-6 Astra (OpenAI), ChatGPT Pro mode, effort Pro, S6\_10\_PAPERS, 18 November 2026: targeted proof check. Continuation metadata events: GPT-6 Astra (OpenAI), ChatGPT Pro mode, effort Pro, S6\_10\_PAPERS continuation, 19 November 2026: used and rechecked in this run; status-vocabulary reconciliation; metadata reconciliation only; no new mathematical check. \emph{Pre-continuation recorded proof scope.} Re-derived the C12 times interval dual cell complex (24 vertices,48 edges,24 faces), the relative collapse sequence and the degree-one edge-map injection into H1 of the fibre. The finite collapse companion program remains a separate computation check.}{the hypotheses, notation and definitions stated here; \cref{push:src:B.4.15}; \cref{push:src:B.3.1}; \cref{push:src:B.4.2}}

The local operation \ref{push:src:B.4.15} attaches a triangulated
cylinder to the original twelve-cycle in the dual
complex and admits an elementary collapse back to that
cycle. In a compact normal-crossings model, performing
it at a regular place consequently preserves the
homotopy type of the global dual complex.

The same conclusion holds for finite sequences of these
insertions, twists by a connected holomorphic group
action preserving the prescribed normal-crossings
model, and toroidal subdivisions of its central cone
complex. In particular, starting with a model whose
dual complex has a nonzero first integral cohomology
class, these operations cannot produce a semistable
$S^6$ centre. They also cannot change the component
rotations in \ref{push:src:B.3.1}.

This statement concerns the specified generated class
of operations. It makes no assertion that arbitrary
analytic caps or arbitrary changes of the punctured
comparison belong to that class.
\end{proposition}



\begin{proof}
\begin{enumerate}[label=\arabic*.,leftmargin=*]
\item Write $H_i$, $i\in\Z/12$, for the original component
vertices and $V_i$ for the new vertical ones. At a node
between an even component $i$ and its odd neighbour
$j$, the two charts (\ref{push:src:B.4.15}.3) have triple strata
\[
 \{V_i,V_j,H_j\},\qquad
 \{V_i,H_i,H_j\}.
 \tag{B.4.16.1}
\]
The square with horizontal edges $H_iH_j$ and $V_iV_j$
is thus divided by the diagonal $V_iH_j$. Over all
twelve nodes, (\ref{push:src:B.4.16}.1) is exactly a triangulation of
$C_{12}\times[0,1]$, attached along the old horizontal
cycle $C_{12}\times\{0\}$.

\item There is a literal elementary collapse relative to the
old cycle. First remove each triangle containing a top
edge $V_iV_j$, paired with that free edge. Next remove
each remaining triangle, paired with its now free
diagonal $V_iH_j$. Finally remove each vertical edge
$H_iV_i$, paired with its free vertex $V_i$. What remains
is precisely the original twelve-cycle. These are
$36$ elementary collapses. Equivalently, the retraction
sends $V_i$ to $H_i$ and is the identity on all old
vertices and edges.

\item This description is valid inside a larger dual complex
$K$. The only new strata lie over the base origin,
so the new complex is the mapping-cylinder attachment
\[
 K\ \cup_{\,f:C_{12}\to K}\
           (C_{12}\times[0,1]).
\]
It retracts onto $K$, regardless of possible global
identifications among the local horizontal branches.
For self-crossings one uses the ordered local-branch
dual complex, so no identification of distinct local
branches is silently discarded.

\item A connected group acts trivially on the finite sets
of component labels and connected components of every
stratum: its permutation is locally constant and is
the identity at the identity element. Cartier twists
by such an action therefore keep this same incidence
complex. Toroidal subdivisions subdivide the
corresponding central polyhedral complex and hence
preserve its underlying homotopy type. These arguments
apply successively, proving the assertion for finite
compositions. Around each regular insertion point the
relative component monodromy is still the identity;
no peripheral rotation is added to \ref{push:src:B.3.1}.

\item For completeness, a compact normal-crossings space $D$
has a natural injection
\[
 H^1(\Delta(D),\C)\ \lhook\joinrel\longrightarrow\
 H^1(D,\C),
 \tag{B.4.16.2}
\]
where $\Delta(D)$ denotes its local-branch dual
complex with connected strata. Use the ordered finite
branch spaces $D^{[p]}$ and take the alternating
summand for the branch-permutation action. Over $\C$
this is an exact direct-summand operation. Resolve the
constant sheaf by this augmented alternating branch
complex. At a point its row is the augmented simplex
complex, so it is exact, including for self-crossings.
The associated first-quadrant spectral sequence has
\[
 E_1^{p,q}=H^q(D^{[p]},\C)^{\mathrm{alt}},
 \qquad
 E_2^{1,0}=H^1(\Delta(D),\C).
\]
No higher differential enters or leaves $(1,0)$.
Its filtration piece in total degree one gives
(\ref{push:src:B.4.16}.2). Equivalently one can use the orientation
local system on each unordered branch stratum; an
ordinary constant sheaf there would be incorrect if
branches permute. A nonzero integral first cohomology
class of the finite dual complex remains nonzero over
$\C$, because first integral cohomology is torsion-free.
Thus a surviving circle class forces $H^1(D,\C)\ne0$.
By \ref{push:src:B.4.2} every semistable centre of
a simply connected nearby $S^6$ is simply connected,
which is incompatible with this conclusion.

\item The finite incidence and collapse certificates themselves
are integral. They are
saved in \texttt{code/ab\_push\_regular\_cap\_dual.py}
and its output. For the twelve-cycle the chain ranks
are $24,48,24$, the nonzero Smith entries of the two
boundary maps are respectively $1^{23}$ and $1^{24}$,
and the integral homology is $(\Z,\Z,0)$.
The explicit retraction pulls the integral angular
cocycle back with winding one. These exact checks
certify the finite combinatorics; the preceding chart
calculation supplies its analytic interpretation.
\end{enumerate}
\end{proof}

\begin{proposition}[trivial monodromy does not force equality of first Betti numbers]\label{push:src:B.4.17}
\StatementRecord{push:src:B.4.17}{proved}{Historical events: Claude Fable 5.2 (Anthropic), effort high, 16 November 2026: carried as background without fault found. GPT-6 Astra (OpenAI), ChatGPT Pro mode, effort Pro, S6\_09\_PAPERS, 17 November 2026: targeted proof check. Also checked in this writing. Pre-continuation writing events: GPT-6 Astra (OpenAI), ChatGPT Pro mode, effort Pro, S6\_10\_PAPERS, 18 November 2026: targeted proof check. Continuation metadata events: GPT-6 Astra (OpenAI), ChatGPT Pro mode, effort Pro, S6\_10\_PAPERS continuation, 19 November 2026: used and rechecked in this run; status-vocabulary reconciliation; metadata reconciliation only; no new mathematical check. \emph{Pre-continuation recorded proof scope.} Re-derived the compact principal-elliptic-bundle counterexample: central b1=0, nearby b1=1, trivial monodromy, and b2(total)=1 from the explicit two Chern classes. Matched Honda-Viaclovsky v1 Theorems1.2/3.2 and equation3.9: the counterexample refutes the converse while preserving the inequality and the equality-to-monodromy implication.}{the hypotheses, notation and definitions stated here; \cref{push:src:B.4.7}}

There is a compact connected complex threefold $Z$ and a proper
holomorphic map $f:Z\to\PP^1$ with connected fibres, with precisely
one singular fibre $Y$, such that
\[
 f^{-1}(\PP^1\setminus\{0\})
 \simeq H_-\times(\PP^1\setminus\{0\}),\qquad
 b_1(Y)=0<1=b_1(H_-).
\]
Here $H_-$ is the scalar Hopf surface with multiplier $-e^{-\pi}$,
and $Y=H_+\cup_E H_+$ is a reduced simple-normal-crossings union
with $E=\C/(\Z+i\Z)$ and $H_+$ of multiplier $e^{-2\pi}$.
In particular, trivial local monodromy on $H^1$ does not imply
$b_1(Y)=b_1(H_-)$, even with compact smooth total space.
\end{proposition}


\begin{proof}
\begin{enumerate}[label=\arabic*.,leftmargin=*]
\item This is the compact-base version of the explicit Cartier cocycle
of \ref{push:src:B.4.7}; no extension of a meromorphic comparison is required.
Put
\[
 b:\overline B=\operatorname{Bl}_{(p,0)}(\PP^1\times\PP^1)
      \longrightarrow\PP^1\times\PP^1,
 \quad L=b^*\operatorname{pr}_1^*\mathcal O(1),
 \quad A=\mathcal O(F_2),
\]
where $F_2$ is the exceptional curve and $F_1$ the strict transform
of $\PP^1\times\{0\}$.  Given transition units $l_{ab},a_{ab}$,
the formula
\[
 t_{ab}=\left[
 \frac{(1+i)\log l_{ab}-i\log a_{ab}}{2\pi i}
 \right]\in E
\]
is independent of logarithm choices and satisfies the cocycle
identity.  Glue the products $U_a\times E$ by these translations.
This gives a holomorphic principal elliptic bundle
$\pi:Z\to\overline B$.  Its total space is smooth and Hausdorff;
it is compact because $\overline B$ and the fibre are compact and
$\pi$ is locally a proper product.  Let
$f=\operatorname{pr}_2\circ b\circ\pi$.

\item The restrictions of the Chern vector
$(1+i)c_1(L)-i c_1(A)$ to $F_1,F_2$ and a smooth ruling fibre
are $1,i,1+i$, respectively.  As checked directly in \ref{push:src:B.4.7},
the corresponding principal elliptic bundles over $\PP^1$
are the two copies of $H_+$ and the surface $H_-$.  Locally
at $F_1\cap F_2$, the base map is $xy=t$ and the elliptic
bundle is a product.  Thus $f$ has smooth total space and just
the one reduced normal-crossings singular fibre $Y$.
The canonical section of $A$ trivializes $A$ outside $F_2$;
there the blow-up is an isomorphism and $L$ is pulled back
from the first factor.  Consequently the asserted punctured
product is an actual holomorphic product over the whole
$\PP^1\setminus\{0\}$.  In particular every local monodromy
operator is the identity.

\item For completeness the fibre inclusions on fundamental groups
are the two quotient maps
\[
 \Lambda=\Z+i\Z\longrightarrow\Lambda/\Z,
 \qquad
 \Lambda\longrightarrow\Lambda/\Z i.
\]
Their van Kampen pushout kills both lattice generators, so
$\pi_1(Y)=0$.  In Mayer--Vietoris the degree-one map, with
the second inclusion negated, has matrix
\[
 \begin{pmatrix}0&1\\-1&0\end{pmatrix}.
\]
It is unimodular.  Since a scalar Hopf surface is diffeomorphic
to $S^1\times S^3$, the sequence gives
\[
 H_*(Y;\Z)=(\Z,0,0,\Z^3,\Z^2),\qquad b_1(H_-)=1.
\]
This proves the asserted strict inequality despite identity
monodromy.  This compact threefold is not $S^6$: for example,
the principal-bundle spectral sequence gives $b_2(Z)=1$.
Indeed the two Chern classes are independent in
$H^2(\overline B;\Z)\simeq\Z^3$, and the differential on
$\bigwedge^2H^1(E)$ is injective, leaving exactly one
degree-two class.
\end{enumerate}
\end{proof}


\paragraph{Source audit and the valid remaining implication}
The converse just disproved is asserted in
\citet[Theorems 1.2 and 3.2, arXiv version~1]{HondaViaclovsky2025D}.
The published primary PDF excerpt, printed p.~553,
retains it as Theorem~1.1.  The two versions and the precise
access record are distinguished in
\texttt{B\_monodromy\_source\_audit.md}.
In the version read in full, equation~(3.9) establishes only
\[
 b_1(Y)\leq\dim H^1(F;\R)^{T}\leq b_1(F).
\]
Identity of $T$ gives equality in the second inequality;
it does not give equality in the first.  The inequality and
the implication from equality of the endpoint Betti numbers
to identity monodromy survive this counterexample.
The irreducible-fibre argument in its Corollary~3.7 also
survives: the top-dimensional homology of an irreducible
central fibre has rank one, forcing equality with the
invariant subspace in the displayed relative-homology
sequence.  No irreducibility is present in the counterexample.
This is a source correction relevant to possible
non-K\"ahler completions, not a nonextension theorem for
the prescribed $S^6$ family.


\Uses{the hypotheses, notation and definitions stated here; \citet[Theorems 1.2 and 3.2, arXiv version~1]{HondaViaclovsky2025D}}


\subsection{The corrected cusp on the punctured bidisc}

\begin{proposition}[a partial cusp filling]\label{push:src:B.4.18}
\label{push:B418}
\StatementRecord{push:src:B.4.18}{proved}{Historical events: Claude Fable 5.2 (Anthropic), effort high, 16 November 2026: carried as background without fault found. Claude Opus 5.5 (Anthropic), effort max, 16 November 2026: carried as background without fault found. GPT-6 Astra (OpenAI), ChatGPT Pro mode, effort Pro, S6\_09\_PAPERS, 17 November 2026: targeted proof check. Also checked in this writing. Pre-continuation writing events: GPT-6 Astra (OpenAI), ChatGPT Pro mode, effort Pro, S6\_10\_PAPERS, 18 November 2026: targeted proof check. Continuation metadata events: GPT-6 Astra (OpenAI), ChatGPT Pro mode, effort Pro, S6\_10\_PAPERS continuation, 19 November 2026: used and rechecked in this run; status-vocabulary reconciliation; metadata reconciliation only; no new mathematical check. \emph{Pre-continuation recorded proof scope.} Re-derived the affine corner cocycle and its explicit primitive F; checked both finite differences, the twelve-fold root shears, the phase cancellation, the free rotating action on the cusp fan, and the actual isomorphism gluing the two partial models away from the missing origin. No origin extension is asserted.}{the hypotheses, notation and definitions stated here; \cref{push:src:B.3.1}; \cref{push:src:B.3.2}; \cref{push:src:B.4.9}; \cref{push:src:A.2.1}; \cref{push:src:A.1.2}; \cref{push:src:B.3.3}}

Put
\[
 B^\times=(\Delta_q\times\Delta_r)\setminus\{(0,0)\}.
\]
There is a proper Hausdorff analytic family over
$B^\times$ with smooth total fourfold, flat over $r$,
whose restriction to $r\ne0$ is the actual prescribed
cusp family after $w=r^6$, and whose reduced
normal-crossings $r$-central fibre over $q\ne0$ has
component deck permutation the shift $-1$ on $\Z/12$
for the clockwise cusp generator $g_0$ of \ref{push:src:B.3.1}.
Equivalently, the positive $q$ loop has shift $+1$.
In particular its $q$-fibre for every $r\ne0$ is the
actual reduced Mumford fibre, while its other axis has
the component transport required by \ref{push:src:B.3.2}.

The overlap is an explicit holomorphic translation,
including the retained affine marking. This proposition
does not assert a proper completion over the missing
origin, nor the global finite-place overlap maps.
\end{proposition}



\begin{proof}
\begin{enumerate}[label=\arabic*.,leftmargin=*]
\item \emph{The fractional deck and its exact comparison.}
On the logarithmic cover of
$U=\Delta_q^*\times\Delta_r^*$ write
\[
 s=\frac{\log q}{2\pi i},\qquad
 t=\frac{\log r}{2\pi i},\qquad
 C(q)=\beta^{\rm part}(q)+s.
\]
The functions $\mu(q)$ and $C(q)$ are single-valued
and holomorphic at $q=0$. After the affine gauge of
\ref{push:src:B.4.9} the actual family is its Jacobian, with period
\[
 \boldsymbol\gamma(s,t)
       =(6\mu(q),C(q)-s-12t),\qquad
 \boldsymbol\delta=(0,1).
 \tag{B.4.18.1}
\]
The other periods are
$\mathbf u=(\tau,\mu)$ and $\mathbf w=(1,0)$.
Under $s\mapsto s+1$, they change by
$\boldsymbol\gamma\mapsto\boldsymbol\gamma-
\boldsymbol\delta$ and $\mathbf u\mapsto\mathbf u+\mathbf w$.
Under $t\mapsto t+1$, the period $\boldsymbol\gamma$
changes by $-12\boldsymbol\delta$.
These are integral changes of the period basis; the
ambient additive $\C^2$ coordinate is unchanged.

\item Set
\[
 b(s,t)=\frac1{12}\boldsymbol\gamma(s,t)
                         -\frac1{24}\boldsymbol\delta.
\]
The two exact identities
\[
 \sum_{k=0}^{11}b(s+k,t)
       =\boldsymbol\gamma(s,t)-6\boldsymbol\delta,
 \qquad
 b(s,t+1)-b(s,t)=-\boldsymbol\delta
 \tag{B.4.18.2}
\]
show that this gives a genuine order-twelve affine
$q$-deck action on the root cover $q=z^{12}$,
well defined around the $r$ loop. The first right-hand
side is an actual period. The scalar correction
$-\boldsymbol\delta/24$ is retained: the bare fraction
$\boldsymbol\gamma/12$ would not have the required
twelfth iterate.

\item There is an explicit holomorphic trivialization on $U$.
Put $L=s+12t$ and
\[
 F(s,t)=
 \left(\frac{\mu(q)L}{2},\,
       \frac{C(q)L}{12}-\frac{L^2}{24}+6t\right).
 \tag{B.4.18.3}
\]
Direct subtraction gives
\[
 F(s+1,t)-F(s,t)=b(s,t),\qquad
 F(s,t+1)-F(s,t)=\boldsymbol\gamma(s,t).
 \tag{B.4.18.4}
\]
The second difference is a period, not zero; that is
exactly what is needed for descent to the torus.
Thus $\zeta\mapsto\zeta+F$ is a holomorphic
isomorphism from the untwisted Jacobian to the
$+b$ affine torus torsor over $U$; its inverse
$\zeta\mapsto\zeta-F$ trivializes that torsor.
Replacing $b,F$ by $-b,-F$ gives the opposite
component rotation. If $\sigma$ denotes the actual
affine section of \ref{push:src:B.4.9}, the map from the actual
family to the $+b$ torsor is explicitly
$\zeta\mapsto\zeta-\sigma+F$.
This retains the affine cusp marking, rather
than only identifying its fibrewise Jacobians.

\item \emph{A rotating compactification before changing the
$q$-boundary.}
We describe the root-cover model to verify the finite
action and properness. In cocharacter coordinates
$(y_1,y_2,h_z,h_r)$, its two period shears are
\[
 \begin{aligned}
 G_1 &: (y_1,y_2)\longmapsto(y_1+12h_z,y_2),\\
 G_2 &: (y_1,y_2)\longmapsto
                         (y_1,y_2+12h_z+12h_r).
 \end{aligned}
\]
The positive fractional deck has the shear
\[
 H:(y_1,y_2)\longmapsto(y_1,y_2+h_z+h_r),
 \qquad H^{12}=G_2,
 \tag{B.4.18.5}
\]
and rotates $z$ by $e^{2\pi i/12}$.
It commutes with $G_1$. Its unit factors are holomorphic
functions of $q$. For example, in the convention
$W=e^{-2\pi i\zeta_2}$, it multiplies $W$ by
\[
 r z\,e^{-2\pi i C(q)/12}e^{\pi i/12}
\]
and the first multiplicative coordinate by
$e^{\pi i\mu(q)}$. The twelve successive root rotations
contribute the factor
$e^{2\pi i(0+\cdots+11)/12}=-1$, and the scalar
correction contributes $e^{\pi i}=-1$.
They cancel. The twelfth iterate is consequently
the actual second period, including its unit, as
also follows additively from (\ref{push:src:B.4.18}.2).

\item Use the unit-square cusp triangulation, with twelve
consecutive bottom cells in the $y_1$ direction and
the same bridge choice in every unit prism. Its
translates have bottom cell indices $(i,j)$ and top
interval $[j,j+1]$. Thus $H$ increases the bottom
second index and the top index by one, while $G_1$
shifts the first index by twelve. The fan is invariant
under both, and therefore under $G_2$.
For each bottom cell its three maximal cones are
the determinant-one cones of \ref{push:src:A.2.1}. The usual
compact-core construction of \ref{push:src:A.1.2}, now on the
$I_{12}$ elliptic root model and with twelve
plumbing intervals before the $G_2$ quotient,
gives a proper model on the root cover. Its
unfolded charts are precisely this fan. Thus
properness and proper discontinuity are supplied
by the compact core and collars, not inferred
from determinants alone.

\item The affine deck extends to this model because its
unit multiplication preserves every toric chart
and its shear preserves the fan. It is free.
On the $r$-central divisor a point lies on at most
two adjacent component labels in $\Z/12$; no
nonidentity translation can preserve that set.
This covers their intersection over the common
base origin as well. On the $q$-central divisor
with $r\ne0$, the bottom triangle has second
component indices contained in $\{j,j+1\}$,
again not preserved by a nonzero translation
modulo twelve. Off both axes the root-base action
is free. Its finite quotient is therefore smooth,
proper over $(q,r)$, flat over $r$, and has a
reduced normal-crossings $r$-central divisor.

\item For completeness, for the negative rotation keep the same positive
root-base rotation and use the inverse translation
$-b$. The resulting valuation shear is $H^{-1}$;
its twelfth iterate is $G_2^{-1}$. It preserves
the same fan and has the same freeness properties.
This is the opposite convention and is not the model
used below. In the retained convention the clockwise
cusp generator is $g_0:s\mapsto s-1$. Thus the
positive fractional deck $b$ has shift $+1$ for
$s\mapsto s+1$ and shift $-1$ for $g_0$, exactly
as required by \ref{push:src:B.3.2}. The finite shifts $4,-3$
and this cusp shift all refer to the same clockwise
deck-identification convention.

\item Denote this rotating model by $\mathcal R$.
Its $q$-fibre for $r\ne0$ has multiplicity twelve.
Indeed the original $z$-central divisor is reduced
and the finite action has trivial generic isotropy
there, while $q=z^{12}$. This is exactly the
limitation established in \ref{push:src:B.3.3}. We do not identify
this full $q$-boundary with the prescribed cusp.

\item \emph{The actual two-open gluing.}
On the base open set $q\ne0$ use $\mathcal R$ with
the positive fractional deck $b$. On the base open set $r\ne0$
use the actual reduced Mumford cusp family.
Their common base is $U$, where (\ref{push:src:B.4.18}.3)--
(\ref{push:src:B.4.18}.4), composed with \ref{push:src:B.4.9}, give the specified
holomorphic isomorphism. These two open sets cover
$B^\times$, and there is no additional triple
overlap. Glue by that isomorphism.

\item The result is Hausdorff: points over one base
point are separated in one of the two given
Hausdorff models, and points with different
base images are separated in the base.
Properness is local on the target, so follows
from the properness of those two restrictions.
Smoothness and flatness over $r$ are also local.
The restriction over $r\ne0$ is, by construction,
the actual prescribed cusp including its reduced
$q$-fibre. Over $r=0,q\ne0$, the construction has
the required twelve-cycle and rotation $-1$.
This proves every asserted property.

\item The exact symbolic differences, finite norm,
scalar phase, commuting shears and unit-prism
determinants are saved in
\texttt{code/ab\_push\_corrected\_cusp\_overlap.py}
and its output. No completion at the origin
follows from these finite calculations.
\end{enumerate}
\end{proof}


\paragraph{Scope}
This construction fills the punctured bidisc, not its
origin. \ref{push:src:B.4.30} proves that its prescribed partial family
has no proper analytic completion at that point.
The comparison (\ref{push:src:B.4.18}.3) contains quadratic
logarithms. Neither a meromorphic graph closure
nor a common toroidal modification is thereby
available. For fixed $r\ne0$, the multiple
$q$-boundary of $\mathcal R$ cannot be substituted
for the actual reduced boundary: the valuation
argument of \ref{push:src:B.3.7} excludes a proper birational
comparison over $q$ between them. \ref{push:src:B.4.33} further
excludes every proper filling of the actual global family,
even if its other central axial model is changed. Thus
the missing origin is not left as a construction claim.


\Uses{the hypotheses, notation and definitions stated here; \cref{push:src:B.4.30}; \cref{push:src:B.4.18}; \cref{push:src:B.3.7}; \cref{push:src:B.4.33}}


\subsection{The actual global family with one base point removed}

\begin{proposition}[the global partial filling]\label{push:src:B.4.19}
\label{push:B419}
\StatementRecord{push:src:B.4.19}{proved}{Historical events: Claude Fable 5.2 (Anthropic), effort high, 16 November 2026: carried as background without fault found; targeted proof check. GPT-6 Astra (OpenAI), ChatGPT Pro mode, effort Pro, S6\_09\_PAPERS, 17 November 2026: targeted proof check. Also checked in this writing. Pre-continuation writing events: GPT-6 Astra (OpenAI), ChatGPT Pro mode, effort Pro, S6\_10\_PAPERS, 18 November 2026: targeted proof check. Continuation metadata events: GPT-6 Astra (OpenAI), ChatGPT Pro mode, effort Pro, S6\_10\_PAPERS continuation, 19 November 2026: used and rechecked in this run; status-vocabulary reconciliation; metadata reconciliation only; no new mathematical check. \emph{Pre-continuation recorded proof scope.} Re-derived global punctured-base descent through the explicit r-power component coordinate maps, cusp exponent cancellation, connector root actions and integral fan translation. Checked the overlapping gauges agree and preserve the proper local models.}{the hypotheses, notation and definitions stated here; \cref{push:src:B.1.2}; \cref{push:src:B.4.18}; \cref{push:src:B.4.30}; \cref{push:src:B.4.33}; \cref{push:src:C.3.5}; \cref{push:src:B.1.1}; \cref{push:src:B.3.2}; \cref{push:src:B.4.9}; \cref{push:src:A.1.2}}

Let
\[
 \Omega=(\mathbb P^1\times\Delta_r)
                  \setminus\{(\infty,0)\}.
\]
There is a proper Hausdorff holomorphic map
$\mathcal Y^\circ\to\Omega$ with smooth total fourfold,
flat over $r$, whose restriction to $r\ne0$ is the
actual family pulled back by $w=r^6$. Its reduced
normal-crossings central divisor has the finite-place
models of \ref{push:src:B.1.2} and the corrected punctured cusp of
\ref{push:src:B.4.18}, with their retained markings. Over
$B^\circ=\mathbb P^1\setminus\{0,1,\infty\}$ its
relative component cover is
\[
 y^{12}=K(x),\qquad K(x)=\frac{x^4}{(x-1)^3}.
\]
Thus the global finite-place overlap maps and their
cocycle are constructed. The sole missing base point
is $(\infty,0)$. This proposition constructs only the
partial family: \ref{push:src:B.4.30} proves its noncompletion, and
\ref{push:src:B.4.33} excludes every proper filling of the actual
family over the whole $r$ disc.
\end{proposition}



\begin{proof}
\begin{enumerate}[label=\arabic*.,leftmargin=*]
\item We distinguish the ordinary base coordinate $x$ from
the logarithmic parameter
\[
 v=\frac{\log r}{2\pi i},\qquad c=-12v,
 \qquad e=(0,1)^t.
\]
Use the coherent closing of \ref{push:src:C.3.5}. Its additive formula
for a positive $r$ loop, that is $c\mapsto c-12$, is
\[
 b_{-12,c}=\Pi_c\widehat\gamma-(k+6)e,
 \qquad k=12\ell=\frac{\log K}{2\pi i}.
 \tag{B.4.19.1}
\]
The first summand and $6e$ are actual periods.
Consequently its torus translation has the scalar
representative $-ke$. This use of a period changes
the representative, not the actual quotient map.

\item The coherent closing in \ref{push:src:C.3.5} may be used for the
specified punctured family. More explicitly, two
coherent choices differing by a scalar automorphism
give a principal $\mathbb C^*$ torsor over the punctured
parameter disc. That torsor is holomorphically trivial:
the exponential sequence and Cartan's theorem B give
$\operatorname{Pic}(\Delta_r^*)=H^2(\Delta_r^*,\mathbb Z)=0$.
The resulting scalar gauge is a holomorphic
isomorphism of the punctured families. It is not
asserted to extend at $r=0$.

\item \emph{The ordinary charts and their exact cocycle.}
Choose a good analytic cover of $B^\circ$, refining
period-frame charts, and on each member choose a
holomorphic logarithm $k_i=(2\pi i)^{-1}\log K$.
On the parameter logarithm cover put
\[
 \widetilde\zeta_i=\zeta+k_i v e.
 \tag{B.4.19.2}
\]
Under the positive $r$ deck, $\zeta$ changes by
$-k_i e$ modulo periods and $v$ changes by one.
Thus (\ref{push:src:B.4.19}.2) has the ordinary period descent.
It identifies the actual local punctured family
with the standard twelve-cycle plumbing of \ref{push:src:B.1.1},
now over an ordinary base chart.

\item On an overlap let $n_{ij}=k_j-k_i\in\mathbb Z$.
The invariant vector $e$ is fixed by every period-frame
change. In the multiplicative coordinate
$W_i=\exp(-2\pi i\widetilde\zeta_{i,2})$, the extra
transition is exactly
\[
 W_j=r^{-n_{ij}}W_i.
 \tag{B.4.19.3}
\]
Every integral power in (\ref{push:src:B.4.19}.3) extends to an
automorphism of the whole twelve-cycle plumbing.
For clarity, on the unfolded node chart write
\[
 x_j=W/r^j,\qquad y_j=r^{j+1}/W,
 \qquad x_jy_j=r.
\]
Multiplication $W\mapsto r^{-n}W$ sends the chart
index $j$ to $j-n$ and leaves both displayed chart
coordinates unchanged. The same statement holds
in bundle frames, whose additional factors are
holomorphic units. It commutes with the closing
period, which shifts the index by twelve and
translates the elliptic base by $6P$. In particular
a shift of twelve need not be the identity on a
central ruled component; no such identity is used.
The usual semiabelian translations and period-frame
changes preserve the two sections, hence also extend.

\item The integers satisfy
$n_{ik}=n_{ij}+n_{jk}$. Since the scalar translations
commute with each other and with the period-frame
changes, (\ref{push:src:B.4.19}.3) is an exact cocycle on the entire
plumbing. This constructs the family over
$B^\circ\times\Delta_r$, with the original
punctured-family identification (\ref{push:src:B.4.19}.2).

\item \emph{The finite-place cancellations.}
At the two root charts put
\[
 w_j=\frac{\log s_j}{2\pi i},\qquad
 \epsilon_1=1,\quad\epsilon_2=-1,
 \qquad k_j=12\kappa_j.
\]
The notation is that of \ref{push:src:C.3.5}: $k_j$ is a
holomorphic logarithm of the unit
$K(x)s_j^{-12\epsilon_j}$ on the full root disc.
Its invariance under the finite deck is exact.
Indeed the unit is invariant since $m_j\mid12$;
the difference of two logarithms is an integer
constant, and evaluating at the fixed root centre
makes that constant zero. The scalar representative
of the local positive $r$ closing is $-k_j e$.
Apply the local gauge
$\widetilde\zeta_j=\zeta_j+k_jve$ and use the
smooth twelve-cycle root model of \ref{push:src:B.1.1}, with its
free finite deck action from \ref{push:src:B.1.2}. Because $e$ is
invariant and $k_j$ is invariant, this gauge does
not change that finite action.

\item On a sector of the punctured root disc choose
compatible logarithms. The actual root-to-ordinary
collar is
\[
 \zeta=\zeta_j+\sigma_{j,c},\qquad
 \sigma_{j,c}=w_j\Pi_c v_j,
 \qquad \gamma(v_j)=\epsilon_j.
\]
Moreover $k-k_j=12\epsilon_jw_j$. Therefore its
expression in the new gauges is
\[
 \begin{aligned}
 \widetilde\zeta-\widetilde\zeta_j
 &=w_j\Pi_c v_j+12\epsilon_jw_jve\\
 &=w_j\Pi_{\rm part}v_j,
 \qquad \Pi_{\rm part}=\Pi_c-ce\gamma.
 \end{aligned}
 \tag{B.4.19.4}
\]
Here $c=-12v$ is essential: the parameter logarithms
cancel exactly. The last translation is independent
of $r$. On each sector it is a holomorphic section
of the semiabelian group and therefore extends to
the ruled normalization and all node charts.

\item Changing a logarithm adds an integral transition
of the form (\ref{push:src:B.4.19}.3), together with an actual
period-frame identification. Those maps already
extend. Their deck and overlap identities agree
on the dense locus $r\ne0$ with the actual collar
identities, hence agree everywhere by the identity
principle. Thus the sector maps descend and glue
exactly to the finite quotients. This constructs
the specified family over
$(\mathbb P^1\setminus\{\infty\})\times\Delta_r$.
It retains the multiplicities three and four and
the reductions already proved in \ref{push:src:B.1.2}.

\item The component cover can also be read without
guessing a global normalization. If $i_i$ is the
unfolded component index modulo twelve, then
(\ref{push:src:B.4.19}.3) gives $i_j=i_i-n_{ij}$. Hence
\[
 y=\exp(2\pi i k_i/12)\exp(2\pi i i_i/12)
\]
is independent of the chart and satisfies $y^{12}=K$.
The clockwise deck shifts are consequently $4,-3,-1$.
The local finite quotients complete this cover at
$0,1$ exactly as in \ref{push:src:B.3.2}. No claim about the missing
cusp point of the central divisor is implicit here.

\item \emph{The corrected cusp overlap, including $r=0$.}
Retain $s=\log q/(2\pi i)$ and the holomorphic
functions $C,\mu,\kappa_0$ of \ref{push:src:B.4.18} and \ref{push:src:C.3.5}.
Then $k=-s+12\kappa_0$. The actual-to-rotating
comparison of the corrected positive-deck model is
$\zeta\mapsto\zeta-\sigma+F$. Substitution in the
displayed formulas of \ref{push:src:B.4.9} and \ref{push:src:B.4.18} gives
\[
 F-\sigma=kve+G(s)+6ve,
 \qquad
 G(s)=\left(\frac{\mu s}{2},\,
                    \frac{Cs}{12}-\frac{s^2}{24}\right).
 \tag{B.4.19.5}
\]
Thus its comparison with the ordinary gauge
(\ref{push:src:B.4.19}.2) is the translation $G(s)+6ve$.
On a sector of $q\ne0$, $G$ is independent of $r$
and gives a holomorphic semiabelian translation.
The other term is the integral fan shift
$W\mapsto r^{-6}W$. Both extend, with their inverses,
over the full twelve-cycle $r$-central divisor.
For completeness the branch identities are
\[
 G(s+1)-G(s)=b(s,v)+ve,
 \qquad k(s+1)-k(s)=-1.
\]
The ordinary coordinate deck contributes $-ve$;
their sum is exactly the positive fractional deck
$b$. A positive $r$ loop changes $6ve$ by the
integral period $6e$. These identities prove
descent of the extended comparison to the full
punctured $q$ disc. In particular the common
clockwise cusp shift is $-1$, as required.

\item Glue the just constructed global model on
$x\ne\infty$ to \ref{push:src:B.4.18} on the cusp neighbourhood
minus its origin, using (\ref{push:src:B.4.19}.5). The overlap
contains $r=0$ and the map extends there by the
preceding chart calculation. All identities
retain the actual punctured marking; there is no
unresolved scalar or translation cocycle.

\item Each chart used is a proper model over its base
open set: this was proved by compact-core plumbing
in \ref{push:src:A.1.2} and by the finite free quotient in \ref{push:src:B.1.2};
\ref{push:src:B.4.18} is proper on its stated punctured bidisc.
Properness is local on the target, so the glued
map to $\Omega$ is proper. The same local argument
gives Hausdorffness, smoothness of the total space,
flatness over $r$, and the reduced normal-crossings
central divisor. For $r\ne0$ the base is the entire
$\mathbb P^1$, and the specified comparison is
precisely the actual compact family. At $r=0$ the
base omits $\infty$; properness over the full
parameter disc does not follow and is not claimed.

\item The exact scalar, finite-root, cusp and node-chart
identities are verified in
\path{code/ab_push_global_partial.py} and its
saved output. They do not supply an origin cap; the
nonexistence statements are proved in \ref{push:src:B.4.30} and \ref{push:src:B.4.33}.
\end{enumerate}
\end{proof}

\subsection{An embedded Tate interface at the missing corner}

\begin{proposition}[the other Kato axis]\label{push:src:B.4.20}
\label{push:B420}
\StatementRecord{push:src:B.4.20}{proved}{Historical events: Claude Fable 5.2 (Anthropic), effort high, 16 November 2026: carried as background without fault found. GPT-6 Astra (OpenAI), ChatGPT Pro mode, effort Pro, S6\_09\_PAPERS, 17 November 2026: targeted proof check. Also checked in this writing. Pre-continuation writing events: GPT-6 Astra (OpenAI), ChatGPT Pro mode, effort Pro, S6\_10\_PAPERS, 18 November 2026: targeted proof check. Continuation metadata events: GPT-6 Astra (OpenAI), ChatGPT Pro mode, effort Pro, S6\_10\_PAPERS continuation, 19 November 2026: used and rechecked in this run; status-vocabulary reconciliation; metadata reconciliation only; no new mathematical check. \emph{Pre-continuation recorded proof scope.} Re-derived the Cartier Kato degeneration with h transformed to a h, its cycle normals and actual Tate interface. Computed the ambient critical-locus map s=uv-h and the distinct collar map s=uv explicitly; persistent rational curves exclude a torus identification for the complete Kato nearby fibre.}{the hypotheses, notation and definitions stated here; \cref{push:src:B.4.10}; \cite[Example 9.10]{IstratiEtAl2022Root}; \cref{push:src:B.4.33}}

In the proper smooth Kato family $\mathcal K\to\Delta_s$
of \ref{push:src:B.4.10} there is a second relative effective Cartier
divisor $\mathcal D\subset\mathcal K$, disjoint from
the persistent elliptic divisor $\mathcal E$.
The surface $\mathcal D$ is smooth and proper over
$\Delta_s$. Its general fibre and central fibre are
\[
 \mathcal D_s=\mathbb C^*/(\kappa(s)s^n)^\mathbb Z
 \quad(s\ne0),\qquad
 \mathcal D_0=\hbox{the reduced rational }n\hbox{-cycle}.
\]
At its nodes the relative local equation is $uv=s$.
Its normal line in $\mathcal K$ is the multidegree-zero
Tate line whose factor on the closing collar is $a(s)$.
In particular it has degree zero on every component
of the central cycle. Scaling the first blowdown
coordinate gives a holomorphic $\mathbb C^*$ action
on the pair $(\mathcal K,\mathcal D)$ over the base.

For $n=12$ this constructs an embedded proper copy of
the actual cusp's elliptic root interface, with its
degree-zero normal line, after choosing the units
$\kappa,a$ to be the actual Tate and normal multipliers.
In the additive period convention $\mathbf u=(\tau,\mu)$,
$W=e^{-2\pi i\zeta_2}$, these multipliers are
$\kappa(s)s^{12}=e^{2\pi i\tau}$ and
$a(s)=e^{-2\pi i\mu}$. The elliptic translation by
$6P$ is realized by first-coordinate scaling
$a(s)^{-6}$.

This is an embedded Cartier and normal-line comparison,
not an identification of relative tubular neighbourhoods.
No lift of the whole second period, attachment to the
global partial filling, or origin completion is asserted.
\end{proposition}



\begin{proof}
\begin{enumerate}[label=\arabic*.,leftmargin=*]
\item Use the modification, embedding and compact collar
domains of \ref{push:src:B.4.10}. Write $w$ for the second blowdown
coordinate and set, on the modified fundamental region,
\[
 h_s=\pi_n^*w-s.
\]
The collar identification is $\sigma_s\pi_n$, where
$\sigma_s(z,w)=(\kappa(s)z,a(s)w+(1-a(s))s)$ in
the last chart. Consequently
\[
 h_s\circ(\sigma_s\pi_n)=a(s)h_s.
 \tag{B.4.20.1}
\]
Its zeros therefore glue to a relative effective
Cartier divisor, and (\ref{push:src:B.4.20}.1) is simultaneously the
normal-line transition. The divisor is closed in the
proper family $\mathcal K\to\Delta_s$, so it is proper.

\item For $s\ne0$, the line $w=s$ misses every blow-up centre.
In the coordinate $v=w-s$ the contracting germ is
\[
 (x,v)\longmapsto
       (\kappa(s)x(a(s)v+s)^n,a(s)v).
\]
Its restriction to $v=0$ is
$x\mapsto\kappa(s)s^n x$. The generic developing
map of \ref{push:src:B.4.10} restricts to the identity in this
coordinate, since its product $G_s(v)$ has
$G_s(0)=1$. Thus this is precisely the unmodified
other Hopf axis, with the stated Tate multiplier
and normal multiplier $a(s)$.

\item The divisor of $\pi_n^*w$ is reduced. After the first
blow-up, both the strict horizontal axis and the
exceptional curve have coefficient one. Every later
centre is the intersection of the newest exceptional
curve with the strict vertical axis. It lies on no
other component of the divisor of $\pi_n^*w$.
The new exceptional coefficient is therefore again
one. Before the end collars are identified, the
divisor is a chain comprising the strict horizontal
axis and the $n$ exceptional pieces. The end annuli
join the two endpoint pieces, leaving the familiar
rational $n$-cycle of the parabolic surface.
It is disjoint from its elliptic curve, as also
described in the parabolic model of
\cite[Example 9.10]{IstratiEtAl2022Root}.

\item At a crossing of two adjacent pieces, ordinary
blow-up coordinates give $\pi_n^*w=uv$ up to a
holomorphic unit, which can be absorbed in $u$.
Hence $\mathcal D$ has equation $uv=s$. It is
smooth, and its central fibre is reduced nodal.
Away from these nodes its defining equation is
a coordinate minus $s$. This proves smoothness
of the whole divisor and semistability over the disc.
The monomial transition charts are those of the
unit-interval Tate model, and the closing on the
nonzero axis is exactly $x\mapsto\kappa(s)s^n x$.
They identify the marked family with that Tate
model, not only its individual smooth fibres.

\item The same local defining functions trivialize the
normal line. Their only nontrivial closing factor
is (\ref{push:src:B.4.20}.1), namely the holomorphic unit $a(s)$.
On each rational component the line has degree
zero; equivalently its two adjacent-node degree
corrections cancel its self-intersection $-2$.
For $n=1$ the same statement uses the two branches
of the self-node. The factor $a(s)$ gives the
retained gluing parameter around the cycle.
This is exactly the multidegree-zero extension
of the line with factor $a(s)$ on the smooth
Tate curves.

\item The first-coordinate scaling commutes with the
modification and the collar maps and preserves
$h_s$. One may equivalently descend the vector
field $x\partial_x$ and its local flow. Properness
over the base makes that flow complete on compact
fibres, uniformly on compact parameter sets;
its $2\pi i$ period yields the stated holomorphic
$\mathbb C^*$ action. On the Tate axis it is
ordinary multiplication of $x$, so a holomorphic
unit $a(s)^{-6}$ realizes the marked translation
by $6P$.

\item For the actual cusp, the first lattice period
acts in multiplicative coordinates by
\[
 (Z,W)\longmapsto
       (e^{2\pi i\tau}Z,e^{-2\pi i\mu}W).
\]
After the twelfth root change, its first multiplier
is $\kappa(s)s^{12}$ for a holomorphic unit.
The second multiplier tends to the contracting
inverse of $-(5+2\sqrt6)$, so after shrinking
$0<|a(s)|<1$ as required in \ref{push:src:B.4.10}.
The preceding exact Tate charts and the normal
factor identify this interface and its line with
the actual ones. The two section normals in the
ruled model are inverse; one chooses the matching
branch accordingly. More precisely, in the total
line-coordinate convention, this $W$ normal is
$M^{-1}$ for $M=\mathcal O(P-O)$: the theta ratio
with divisor $P-O$ has multiplier $e^{2\pi i\mu}$,
the inverse of $a$. The opposite ruled section
uses $M$. The separate first-coordinate factor
$a^{-6}$ is the translation by $6P$.
These marked quantities are not conflated.

\item We record the neighbourhood limitation explicitly.
At a central node, the smooth Kato total space
has coordinates $(u,v,s)$ with divisor
$h=uv-s=0$. In coordinates $(u,v,h)$ its map to
the base is
\[
 s=uv-h.
 \tag{B.4.20.2}
\]
It is a submersion. The total space of the normal
line, with the inherited map from $\mathcal D$,
would instead have $s=uv$, with a critical normal
line over that node. Therefore there is no
relative linear tubular-neighbourhood identification
across the nodes. This is a literal local
singularity distinction, not only a possible
global obstruction to linearization.

\item An inverse-normal plumbing variable $f$ would
give the local equation
\[
 r=fh=f(uv-s).
 \tag{B.4.20.3}
\]
The total fourfold in these coordinates is
smooth and its $r$-central divisor is reduced
normal crossing. For $r\ne0$ one has $f\ne0$
and $s=uv-r/f$, whose derivative in $f$ is
$r/f^2\ne0$. Thus this local plumbing changes
the behaviour at the Tate node; (\ref{push:src:B.4.20}.3) by
itself does not identify the prescribed two-period
Mumford cusp. This auxiliary proposition supplies neither
global plumbing data nor that period comparison. \ref{push:src:B.4.33}
subsequently excludes any proper filling of the actual family.

\item There is also a concrete obstruction to the most
direct global use of this plumbing. For fixed
$s\ne0$, all the $n$ blow-up centres lie on the
first Hopf axis $x=0$, at $[v=-s]$ and its
infinitely near strict-axis points. The other
axis $\mathcal D_s$ is $v=0$, $x\ne0$.
The axes are disjoint, and hence every exceptional
rational curve misses $\mathcal D_s$. A plumbing
that leaves the complement of a small tubular
neighbourhood of $\mathcal D_s$ unchanged retains
these compact curves in its $r\ne0$ fibres over
the same $s$. A complex torus has no rational
curve: a holomorphic map from $\mathbb P^1$
lifts to its vector-space universal cover and
is constant. Thus this fixed-base localized
plumbing of the entire Kato surface cannot
produce the prescribed ordinary two-tori.
Further deformations or a different map to the
original base are outside this obstruction.

\item The Cartier factor, Tate multiplier, reduced
divisor multiplicities, scaling action and the
two local base-map comparisons are checked in
\path{code/ab_push_tate_interface.py} and
its saved output. Properness rests on the
explicit compact domains of \ref{push:src:B.4.10}.
\end{enumerate}
\end{proof}

\begin{proposition}[a degree-one Kato model with a linear twelve-gon collar]\label{push:src:B.4.21}
\StatementRecord{push:src:B.4.21}{proved}{Historical events: Claude Fable 5.2 (Anthropic), effort high, 16 November 2026: carried as background without fault found. GPT-6 Astra (OpenAI), ChatGPT Pro mode, effort Pro, S6\_09\_PAPERS, 17 November 2026: targeted proof check. Also checked in this writing. Pre-continuation writing events: GPT-6 Astra (OpenAI), ChatGPT Pro mode, effort Pro, S6\_10\_PAPERS, 18 November 2026: targeted proof check. Continuation metadata events: GPT-6 Astra (OpenAI), ChatGPT Pro mode, effort Pro, S6\_10\_PAPERS continuation, 19 November 2026: used and rechecked in this run; status-vocabulary reconciliation; metadata reconciliation only; no new mathematical check. \emph{Pre-continuation recorded proof scope.} Re-derived the six singularity blowups and multiplicity sequence 2 min(i,12-i), the 72 normal-regularizing curve blowups, all normal twists, and the uniformly convergent developing product that conjugates the return map to (kappa q Z,a H). Checked persistence of the marked negative elliptic normal line.}{the hypotheses, notation and definitions stated here; \cref{push:src:B.4.10}; \cref{push:src:B.4.19}; \cref{push:src:B.4.33}; \cref{push:src:B.4.20}}

Start with the $n=1$ Kato family of \ref{push:src:B.4.10} over a disc
$\Delta_q$, with holomorphic units $a(q),\kappa(q)$ and
$0<|a(q)|<1$.  Pull it back by $q=s^{12}$.
There is a proper modification of this pulled-back threefold,
obtained by six point blow-ups and seventy-two smooth curve
blow-ups supported on its central rational cycle, with the
following properties.  Its total space $\mathcal K^*$ is smooth,
it is proper and flat over $\Delta_s$, and its restriction to
$s\ne0$ is the original pulled-back Kato family.

The strict transform $\mathcal D^*$ of the other Cartier axis
is the smooth semistable Tate surface with multiplier
$\kappa(s^{12})s^{12}$ and central fibre $I_{12}$.  A neighbourhood
of $\mathcal D^*$ in $\mathcal K^*$ is biholomorphic over
$\Delta_s$ to a neighbourhood of the zero section in the
multidegree-zero Tate line with closing factor $a(s^{12})$.
The comparison retains the explicit transverse coordinate
\[
 H=\frac{h}{q},\qquad h=\pi_1^*w-q,
 \tag{B.4.21.1}
\]
on the original punctured family.  Thus its normal comparison,
relative to the original punctured normal frame, has factor
$q^{-1}=s^{-12}$; its inverse has factor $q$.

The persistent elliptic axis has an unchanged neighbourhood.
In particular its normal in the strict transform of the original
central Kato surface still has degree $-1$, rather than $-12$.
The construction supplies an actual linear collar with this
marking.  It does not extend the whole second period, construct
an origin cap for \ref{push:src:B.4.19}, or identify the modified Kato family
with the prescribed torus family. Such a proper filling of
the actual family is excluded by \ref{push:src:B.4.33}; the present
statement retains its auxiliary embedded-collar scope.
\end{proposition}


\begin{proof}
\begin{enumerate}[label=\arabic*.,leftmargin=*]
\item We distinguish the original parameter $q$ from its twelfth root
$s$.  The pulled-back Kato total space is smooth: the original
map to $q$ is a submersion, so locally its pullback has smooth
coordinates consisting of two fibre coordinates and $s$.
The other axis of \ref{push:src:B.4.20} is Cartier, with closing factor $a(q)$
for its defining function $h$.  At its unique central node its
local equation after this base change is
\[
 h=uv-s^{12}.
 \tag{B.4.21.2}
\]
Away from that point the axis surface is smooth.  The persistent
elliptic axis is disjoint from it throughout the family.

\item \emph{The six point blow-ups.}
Blow up the singular point of the strict transform of the
Cartier surface, six times in succession.  All these are points
of the smooth ambient threefold.  If its current equation is
$uv=s^m$, with $m\ge2$, the $s$ chart of the ambient point
blow-up has
\[
 u=su_1,\qquad v=sv_1,
 \qquad u_1v_1=s^{m-2}.
 \tag{B.4.21.3}
\]
The $u$ chart has $v=uv_1$, $s=us_1$, and strict-transform
equation
\[
 v_1=u^{m-2}s_1^m;
 \tag{B.4.21.4}
\]
the $v$ chart is symmetric.  The latter charts are smooth.
For $m>2$ the only remaining singular point is the origin in
(\ref{push:src:B.4.21}.3); for $m=2$ that chart has $u_1v_1=1$.
The exponents are $12,10,8,6,4,2$, so precisely six point
blow-ups resolve the surface.

\item Write $D$ for its smooth strict transform.  Its central fibre is
\[
 F=C_0+C_1+\cdots+C_{11}=\operatorname{div}_D(s).
 \tag{B.4.21.5}
\]
Here $C_0$ is the strict transform of the old nodal rational
curve.  Each of the first five point blow-ups contributes two
rational exceptional curves, and the sixth contributes the
smooth conic in its exceptional plane.  In cyclic order the
two chains are $C_1,\ldots,C_5$ and $C_{11},\ldots,C_7$, with
the conic $C_6$ between them.  Equations (\ref{push:src:B.4.21}.3)--
(\ref{push:src:B.4.21}.4) show that every component has multiplicity one,
that consecutive components meet transversely once, and that
there are no other intersections.  Thus the fibre is $I_{12}$.
Since $F$ is principal, $C_i\cdot F=0$ and $C_i^2=-2$.
The same charts identify $D$ with the minimal regular Tate
model of the given multiplier, including the original closing
unit $\kappa(q)$.

\item Let $M$ be the pullback to $D$ of the original Cartier normal
line.  It is the Tate line with factor $a(q)$ and multidegree
zero.  A component contracted to the original node has degree
zero for this pullback, and its restriction to $C_0$ is the
pullback of the degree-zero line on the nodal curve.
In the iterated $s$ chart after $j$ point blow-ups the original
function is
\[
 h=s^{2j}(u_jv_j-s^{12-2j}).
\]
The $j$th exceptional ambient divisor therefore has coefficient
$2j$ in the total transform of $\operatorname{div}(h)$.
For $j<6$ its intersection with $D$ consists of
$C_j+C_{12-j}$, and for $j=6$ it is $C_6$.
Consequently the new normal line is
\[
 N_{D/\mathcal K^{(6)}}=M\otimes\mathcal O_D(-Z),
 \quad Z=\sum_{i=0}^{11}z_iC_i,
 \quad z_i=2\min(i,12-i).
 \tag{B.4.21.6}
\]
The list of coefficients is
$(0,2,4,6,8,10,12,10,8,6,4,2)$.
In particular its degree vector is $-4$ on $C_0$, $+4$ on
$C_6$, and zero on the other ten curves.

\item \emph{The seventy-two curve blow-ups and the actual function.}
For each $i$, blow up the ambient threefold along the current
copy of $C_i$ exactly
\[
 b_i=12-z_i
\]
times, in increasing order of $i$; the entry $b_6=0$ requires
no blow-up.  At every step $C_i$ is a smooth embedded compact
curve.  Its ideal on $D$ is invertible, so the strict transform
of $D$ is canonically isomorphic to $D$ itself.  The normal
line is tensored by $\mathcal O_D(-C_i)$ at that step.
Thus the final normal is
\[
 M\otimes\mathcal O_D(-Z-\textstyle\sum_i b_iC_i)
   =M\otimes\mathcal O_D(-12F)\simeq M.
 \tag{B.4.21.7}
\]
The sum of the $b_i$ is $72$.
The last isomorphism divides the local generator $s^{12}$
of $\mathcal O_D(-12F)$ by $s^{12}$.  Relative to the
original punctured normal identification it is multiplication
by $s^{-12}$, as asserted.

\item We need more than the line isomorphism (\ref{push:src:B.4.21}.7).
The point-blow-up charts give the following exact local
description near $D$.  At a node $C_i\cap C_{i+1}$ one can
choose ambient coordinates $(x,y,e)$ such that
\[
 D=(e=0),\qquad s=xy,\qquad
 h=e x^{z_i}y^{z_{i+1}}\epsilon,
 \tag{B.4.21.8}
\]
where $\epsilon$ is a holomorphic unit.  At a general point
of $C_i$, omit $y$ and use $s=x$.  These formulas follow
directly from (\ref{push:src:B.4.21}.4): the strict transform is a coordinate
minus a monomial, and the two factors of $s$ are the two
vertical divisors.  They also show transversality of $D$ to
their intersection.  Absorbing a unit in the second base
coordinate preserves this description if different original
coordinate frames are used.

\item Blowing up $C_i$ is locally the blow-up of $(e,x)$.
In the chart containing the strict transform of $D$, put
$e=xe'$.  Formula (\ref{push:src:B.4.21}.8) keeps $s=xy$ and increases
the exponent of $x$ in $h$ by exactly one.  The old vertical
divisor in this direction no longer meets the strict transform
of $D$.  At an adjacent node the exponent of $y$ is unchanged.
This calculation applies at every repetition and for the
other component as well.  At the end, every vertical ambient
divisor meeting $D$ has $s$-order one and $h$-order twelve.
Hence (\ref{push:src:B.4.21}.8) becomes
\[
 h=s^{12}e\epsilon,
 \qquad H=h/q=e\epsilon.
 \tag{B.4.21.9}
\]
In particular $H$ is holomorphic near all of $D$, has divisor
exactly $D$, and has nonzero normal derivative there.  On the
closing collar its transition is exactly $H\mapsto a(q)H$.
The construction has regularized the actual meromorphic
function, rather than just its normal restriction.

\item \emph{The relative product charts.}
There is an equivalent useful local target for the preceding
modifications.  The ideal $(h,q)$ is globally defined, since
its changes of generators multiply $h$ by the unit $a(q)$.
Near $D$ the modification maps to its $q$ blow-up chart,
because (\ref{push:src:B.4.21}.9) makes the pullback ideal principal.
At the old node this chart has
\[
 uv=s^{12}(1+H).
\]
Restrict to $|H|<\eta<1$ in a local closing frame and put
$V=v/(1+H)$.  The equation becomes
\[
 uV=s^{12},
\]
times the independent coordinate $H$.
The preceding smooth modification near $D$ is its minimal
resolution times that coordinate.  We verify this assertion
in characters, rather than infer it from the normal line.

\item Label the exceptional chain so that the orders of $u,V,s$
on $C_i$ are $i,12-i,1$, respectively, with the two end
branches at indices $0,12$.  These are the values in the
explicit surface resolution above.  Along any ambient vertical
divisor meeting $C_i$, the orders of $u$ and $V$ sum to twelve,
because $uV=s^{12}$ and $s$ has order one.  Restriction to
$D$ can only increase those orders.  Since the restricted
orders already sum to twelve, the ambient orders are exactly
$i,12-i$.  No other zero or pole of either character occurs
near $D$, because their product is $s^{12}$.

\item On the chart of the two adjacent rays $i,i+1$, the functions
\[
 X_i=u/s^i,\qquad Y_i=s^{i+1}/u,
 \qquad X_iY_i=s
 \tag{B.4.21.10}
\]
therefore extend holomorphically.  A meromorphic function on
the smooth ambient space with no codimension-one pole is
holomorphic.  Their restrictions are the two standard minimal
resolution coordinates on $D$.  Together with $H$, whose
normal derivative is nonzero, they have invertible differential
along $D$.  Thus (\ref{push:src:B.4.21}.10) and $H$ are actual biholomorphic
product charts near $D$.  Their character identities give the
standard transitions of the minimal resolution of $A_{11}$.
In particular there is no additional exceptional divisor
meeting $D$ hidden in this product comparison.

\item \emph{The closing and the uniform linear collar.}
The remaining issue is the global closing, which the Kato
formula determines exactly.  In its blowdown coordinate $x$
and $h=w-q$, that formula is
\[
 (x,h)\longmapsto
 \bigl(\kappa(q)x(a(q)h+q),a(q)h\bigr).
\]
On the regularized collar this reads
\[
 (x,H)\longmapsto
 \bigl(\kappa(q)q x(1+a(q)H),a(q)H\bigr).
 \tag{B.4.21.11}
\]
Shrink the parameter disc so $|a(q)|\le A<1$.
For $|H|<\eta<1$ the product
\[
 G(q,H)=\prod_{j\ge1}(1+a(q)^jH)
 \tag{B.4.21.12}
\]
converges locally uniformly in both variables, is holomorphic
and nonzero, and satisfies
\[
 G(q,a(q)H)=\frac{G(q,H)}{1+a(q)H}.
\]
Uniform convergence follows from
$\sum_{j\ge1}|a(q)^jH|\le\eta A/(1-A)$; the factors have
no zero in the chosen neighbourhood, and the convergent
logarithm series gives a nonzero product and holomorphic inverse.
In the new character $Z=xG(q,H)$, (\ref{push:src:B.4.21}.11) becomes
\[
 (Z,H)\longmapsto
       (\kappa(q)qZ,a(q)H).
 \tag{B.4.21.13}
\]
Multiplication of a toric character by this unit extends on
the product charts (\ref{push:src:B.4.21}.10): it multiplies $X_i$ by the
unit and $Y_i$ by its inverse.  It therefore gives actual
isomorphisms at all twelve nodes as well as on the annular
core.  The exact closing (\ref{push:src:B.4.21}.13) identifies these charts
with the Tate line of factor $a(q)$, over the Tate family
of multiplier $\kappa(q)q$.  Their inverse maps obey the
same identities.  Compactness of the central divisor and
a finite cover by core and node charts allow one common
shrinking of the base and the collar neighbourhood.
This proves the claimed neighbourhood isomorphism.

\item Every ambient modification used above is a blow-up of a smooth
closed centre and is proper.  The total space remains smooth,
and its map to the parameter disc is flat because a nonzero
base function is a non-zero-divisor in its smooth local rings.
All centres lie over $s=0$, so the punctured family is unchanged.
All centres are disjoint from the persistent elliptic axis,
so its entire neighbourhood is unchanged; in particular its
central normal degree remains $-1$.
We do not assert that the modified ambient map is a submersion
at $s=0$.

\item The exact exponents, intersection vector, blow-up count and
local substitutions are checked by
\path{code/d_push_degree_one_repair.py} and its saved
output.  The construction stops at the collar.  It does not
extend the second-period action to the complementary Kato
region or remove the exceptional-curve obstruction of \ref{push:src:B.4.20}
to the naive fixed-base plumbing of the entire family.
\end{enumerate}
\end{proof}

\subsection{A local resolution retaining the cusp factor}

\begin{proposition}[the repaired local plumbing model]\label{push:src:B.4.22}
\label{push:B422}
\StatementRecord{push:src:B.4.22}{proved}{Historical events: Claude Fable 5.2 (Anthropic), effort high, 16 November 2026: carried as background without fault found. GPT-6 Astra (OpenAI), ChatGPT Pro mode, effort Pro, S6\_09\_PAPERS, 17 November 2026: targeted proof check. Also checked in this writing. Pre-continuation writing events: GPT-6 Astra (OpenAI), ChatGPT Pro mode, effort Pro, S6\_10\_PAPERS, 18 November 2026: targeted proof check. Continuation metadata events: GPT-6 Astra (OpenAI), ChatGPT Pro mode, effort Pro, S6\_10\_PAPERS continuation, 19 November 2026: used and rechecked in this run; status-vocabulary reconciliation; metadata reconciliation only; no new mathematical check. \emph{Pre-continuation recorded proof scope.} Re-derived the three smooth small-resolution charts for fuv=gs, then the regular height-one subdivision after s to s\textasciicircum{}n. Checked gluing, the r-map and normal-crossing fibres on every chart. This is the stated local geometric construction.}{the hypotheses, notation and definitions stated here; \cref{push:src:B.4.21}; \cref{push:src:B.4.19}; \cref{push:src:B.4.20}; \cref{push:src:B.4.30}; \cref{push:src:B.4.33}}

Let
\[
 X_n=\{fuv=g s^n\}\subset\mathbb C^5,
 \qquad r=g-f,\qquad q=s^n.
\]
There is a proper bimeromorphic local modification
$\widetilde X_n\to X_n$ with smooth total fourfold
and reduced normal-crossings divisor $r=0$.
For fixed $r\ne0$ it resolves the cyclic base change
of a three-branch $q$-normal-crossings chart.
This is a local statement and does not identify its
global smooth fibres with the prescribed tori.
The divisor $s=0$ after the displayed cyclic base change
is reduced; the divisor of the original function $q=s^n$
has multiplicity $n$. No descent removing that multiplicity
is asserted.
\end{proposition}



\begin{proof}
\begin{enumerate}[label=\arabic*.,leftmargin=*]
\item For $n=1$ first blow up $(u,q)$ and then the pullback
of $(f,g)$. The charts of the resulting modification are
\[
\begin{array}{c|c|c}
 \text{coordinates}&\text{map to }X_1&r\\\hline
 (q,U,v,f)&u=qU,\ g=fUv&f(Uv-1)\\
 (u,Q,G,f)&q=uQ,\ v=GQ,\ g=fG&f(G-1)\\
 (u,v,F,g)&q=uvF,\ f=gF&g(1-F).
\end{array}
\]
To derive them, the $q$ chart of the first blow-up
has $u=qU$, hence $g=fUv$; the ideal $(f,g)$ is
already principal there. Its other chart has
$q=uQ$ and equation $fv=gQ$. Blowing up $(f,g)$
gives the last two charts. These are the actual
blow-up charts, with the common factors cancelled
on the closure of the dense torus; no additional
component of the naive total transform is included.
They are smooth. In the first chart $Uv-1$ is a
smooth coordinate along its zero set, and in the
others $G-1$ or $1-F$ is a coordinate. Thus the
displayed products prove the reduced NC assertion
for $r$.

\item For arbitrary $n$, pull back this proper modification
by $q=s^n$. The first chart remains smooth. The
second is $uQ=s^n$ times the two coordinates $G,f$.
Resolve its $A_{n-1}$ factor by the unit-interval fan.
The third chart is
\[
 uvF=s^n
\]
times the independent coordinate $g$. Its cone lattice
may be written as $N=\mathbb Z^3$ with coordinates
$(a,b,d)$, where the valuations of $u,v,F,s$ are
$a,b,nd-a-b,d$. The cone is the cone over
$n\Delta_2$. Insert all rays $(i,j,1)$ with
$i,j\ge0$ and $i+j\le n$, and triangulate by the
unit triangles
\[
\begin{aligned}
 &(i,j),(i+1,j),(i,j+1) &&(i,j\ge0,\ i+j\le n-1),\\
 &(i+1,j),(i,j+1),(i+1,j+1)
                       &&(i,j\ge0,\ i+j\le n-2).
\end{aligned}
\]
Every determinant is one in absolute value; the
triangles partition the large triangle and give a
proper smooth toric subdivision. Their number is
$n^2$. All new rays have $s$-height one.

\item The open set $F\ne0$ is the toric open corresponding
to $nd-a-b=0$. Its fan is the unit subdivision of
the edge $i+j=n$, hence the minimal $A_{n-1}$
resolution times $\mathbb C^*$. This is precisely
the second chart's resolution under $Q=vF$.
Where that chart meets the first, the $A$-type
resolution is an isomorphism. These facts glue
the local proper resolutions to a proper
modification of the whole local $X_n$.

\item The function $F$ is a toric monomial. Above $F=1$
it is submersive: in a regular toric chart a
nonzero monomial has nonzero derivative in a
coordinate with positive exponent. Equivalently,
on $F\ne0$ it is the displayed torus projection.
Thus $r=g(1-F)$ has two smooth transverse
branches in the third chart after resolution.
The second chart has $r=f(G-1)$ independently
of the resolved surface coordinates, and the
first was already checked. This proves smoothness
and reduced normal crossings over $r$ for all $n$.

\item For $r\ne0$ in the third chart,
$g=r/(1-F)$, so the local fibre is the resolved
threefold $uvF=s^n$ with $F\ne1$.
For $n=1$, its $q$ divisor is $uvF=0$, a literal
three-branch normal-crossings chart. For $n=12$
the standard triangulation gives 144 unimodular
charts above this one triple point. This is the
same local cyclic-cover count used earlier for
each of the two actual cusp triple points; it
is not a new global Euler count. The other
potential singularity, at $f=-r$, is included
in this construction: the first blow-up small-resolves it
for $n=1$, and the second is an isomorphism there
because $f$ is nonzero.

\item The equation arises from the actual repaired
neighbourhood of \ref{push:src:B.4.21}. There $h=uv-q$ in the
old node chart and $H=h/q$ is holomorphic in the
repaired collar. The unit inverse-normal plumbing
$Hf=r$ is exactly
\[
 f(uv-q)=rq,
 \qquad fuv=(f+r)q.
\]
Putting $g=f+r$ gives $X_n$. Thus the retained
$q$ factor has been derived from the actual
meromorphic function regularized by the
modification, rather than inferred solely from
an abstract normal-line isomorphism.
More precisely, the actual repaired collar, where $H$ is small,
identifies with the third chart near $F=1$ by
\[
 F=(1+H)^{-1},\qquad g=f(1+H),\qquad r=fH.
\]
There $uvF=s^n$ is exactly the resolved equation
$uv=s^n(1+H)$ from \ref{push:src:B.4.21}. The charts near $F=0$
are additional local completion charts; their attachment
to a complementary global Kato piece has not been established.

\item This local equation and its resolution do not
supply the complementary Kato charts, their
second-period identification, the proper closure
of the surviving exceptional-curve families, or
an origin cap for \ref{push:src:B.4.19}. In particular the
fixed-base obstruction in \ref{push:src:B.4.20} is not removed
by recognizing a local Mumford triple chart. \ref{push:src:B.4.30} and
\ref{push:src:B.4.33} exclude the respective actual origin completion
and global proper filling; the local model above remains
a proved auxiliary construction.
The exact chart identities, determinants, edge
incidences and fan counts are saved in
\path{code/ab_push_corner_resolution.py} and its
output.

\item The full period action is a separate condition. The established
first collar action is $H\mapsto aH$ and $f\mapsto a^{-1}f$;
the product $r=fH$ is fixed. The product $G(q,H)$ from \ref{push:src:B.4.21}
conjugates this on a small collar, but this does not extend that
conjugacy past its zeros $H=-a^{-j}$. If an actual marked period
were to scale the test coordinate $W=1+H=1+r/f$ by $b$, its
required action on $f$ would be
\[
 f\longmapsto\frac{rf}{(b-1)f+br},
\]
not scalar multiplication. The test coordinate has not been
identified with a marked period coordinate, so this identity
records a compatibility requirement, not an obstruction.
\end{enumerate}
\end{proof}

\begin{proposition}[exceptional-curve nonclosure and persistence in a semistable cap]\label{push:src:B.4.23}
\StatementRecord{push:src:B.4.23}{proved}{Historical events: Claude Fable 5.2 (Anthropic), effort high, 16 November 2026: carried as background without fault found. GPT-6 Astra (OpenAI), ChatGPT Pro mode, effort Pro, S6\_09\_PAPERS, 17 November 2026: targeted proof check. Also checked in this writing. Pre-continuation writing events: GPT-6 Astra (OpenAI), ChatGPT Pro mode, effort Pro, S6\_10\_PAPERS, 18 November 2026: targeted proof check. Continuation metadata events: GPT-6 Astra (OpenAI), ChatGPT Pro mode, effort Pro, S6\_10\_PAPERS continuation, 19 November 2026: used and rechecked in this run; status-vocabulary reconciliation; metadata reconciliation only; no new mathematical check. \emph{Pre-continuation recorded proof scope.} Re-derived the exceptional intersection winding (12,0), its multiplication under finite base change, and the proper-divisor non-extension argument. Checked the O(-1) plus O plus O normal bundle and its embedded deformation map. Exact Dloussky-Teleman proposition is a separately tracked source gate. Closed the persistence input gate by reading Dloussky-Teleman2012 Proposition2.3pp6-7 and its complete relative normal-sequence/Douady proof.}{the hypotheses, notation and definitions stated here; \cref{push:src:B.4.21}; \cref{push:src:B.4.10}; \citet[Section 1.2]{DlousskyTeleman2012AB}; \citet[Proposition 2.3 and its proof, pp.~6--7 of the
original preprint]{DlousskyTeleman2012AB}; \citet{Merkulov1995D}}

Let $\mathcal K^*\to\Delta_s$ be the repaired degree-one
Kato family of \ref{push:src:B.4.21}, with $q=s^{12}$, and let
$\mathcal E\to\Delta_s$ be its unchanged persistent elliptic
axis. For $s\ne0$ let $C_s\simeq\PP^1$ be the exceptional
curve of the single Hopf blow-up.

The union of the $C_s$ has no proper analytic divisor
closure in $\mathcal K^*$. Further finite proper modifications
cannot give these curves a proper analytic divisor closure:
any such closure would push forward to a forbidden one.
In particular, the generic Hopf blowdown does not extend
as a relative proper holomorphic modification, or as a
relative meromorphic contraction represented by a proper
analytic graph, across $s=0$.

There is also a deformation obstruction which does not
assume a particular plumbing. Suppose a smooth complex
fourfold $\mathcal Y$ has holomorphic functions $(Q,r)$,
its divisor $r=0$ is reduced normal crossings, and one
component $K$ contains a neighbourhood of some $C_{s_0}$,
$s_0\ne0$, agreeing with the complementary Kato region
and its ordinary base coordinate $Q$. Suppose this curve
is disjoint from every conductor branch of $r=0$.
Then every sufficiently nearby fibre $(Q,r)=(Q_0,r_0)$
contains an embedded rational curve. Consequently such a
cap cannot have the prescribed ordinary two-tori as its
$r\ne0$ surface fibres. The conclusion applies to the
whole complementary Kato component with these hypotheses;
it makes no assertion about a different central component,
a curve meeting the conductor, or a replacement changing
the holomorphic base map.
\end{proposition}


\begin{proof}
\begin{enumerate}[label=\arabic*.,leftmargin=*]
\item We first use the exact marked blowdown of \ref{push:src:B.4.10}.
The exceptional curve of the single generic Hopf blow-up
meets the persistent elliptic axis transversely at
\[
 [v=-q]\in E_{a(q)}=\C^*/a(q)^{\Z}.
\]
All modifications in \ref{push:src:B.4.21} are disjoint from that axis.
The intersection section in the repaired family is therefore
still
\[
 p(s)=[-s^{12}].                                      \tag{B.4.23.1}
\]
The units $a(q)$ obey $0<|a(q)|<1$. Choose a holomorphic
logarithm on the full parameter disc and put
\[
 \tau(s)=\frac{\log a(s^{12})}{2\pi i}.
\]
Then $\mathcal E_s=\C/(\Z+\tau(s)\Z)$ and
$\operatorname{Im}\tau(s)>0$. The maps
$(x,y,s)\mapsto[x+y\tau(s)]$, with $(x,y)\in\R^2/\Z^2$,
give a smooth trivialization of the elliptic family.
A logarithmic lift of (\ref{push:src:B.4.23}.1) is
\[
 z(s)=\frac{\log(-s^{12})}{2\pi i}.
\]
On a positive loop around zero, $\tau$ returns to itself
and $z$ changes by $12$. Thus the loop of the section has
class $(12,0)\ne0$ in the integral first homology of this
trivialized family. A continuous section over the disc
would fill that loop. Hence $p$ has no continuous, and
therefore no holomorphic, extension. A meromorphic section
from a nonsingular curve would extend holomorphically:
its proper graph is generically degree one over the curve,
and normalization gives a finite degree-one map to that
normal base. This is equally impossible. A finite base
change of degree $d>0$ multiplies the winding to $(12d,0)$
and does not change the conclusion.

\item Suppose an analytic divisor $Z$ in $\mathcal K^*$, proper
over the disc, restricted over the punctured disc to the
union of the $C_s$. Its horizontal irreducible part does
not contain $\mathcal E$, since the generic intersection
is the one transverse point (\ref{push:src:B.4.23}.1). Both are Cartier
locally in the smooth threefold, so their intersection
is an analytic curve, proper over the disc. Its reduced
horizontal part has generic degree one. An irreducible
horizontal analytic curve cannot have a positive-dimensional
fibre without being contained in that fibre. Hence its
normalization is finite of degree one over the disc and
is isomorphic to the disc. It extends $p(s)$, a
contradiction. Vertical components of the intersection
have no effect on this horizontal component.

\item If an additional finite proper modification admitted a
proper analytic closure, its proper image would give
such a $Z$, by the proper mapping theorem; the image has
dimension two and the same generic exceptional curves.
The horizontal exceptional locus of a relative proper
holomorphic blowdown likewise gives $Z$. For a meromorphic
contraction with proper analytic graph, resolve the graph
and apply the same argument to its horizontal exceptional
locus and its proper image in $\mathcal K^*$.
This is a direct marked-intersection proof of a familiar
nonproperness phenomenon for exceptional curves in GSS
surface deformations; see
\citet[Section 1.2]{DlousskyTeleman2012AB} for the primary
context. No classification of all central curves is used.

\item We next prove the deformation assertion. Write
$C=C_{s_0}$ and $S=K\cap\{Q=Q_0\}$. Near this curve,
$S$ is a smooth surface and $C$ is its exceptional curve,
so
\[
 N_{C/S}\simeq\OO_{\PP^1}(-1).
\]
The base coordinate $Q$ is submersive on $K$ near $C$.
Because $C$ misses every other central branch, the reduced
function $r$ is a local defining function for $K$ near
$C$ and its normal derivative is nowhere zero.
Consequently the differential of $(Q,r)$ gives the exact
normal sequence
\[
 0\longrightarrow\OO_{\PP^1}(-1)
 \longrightarrow N_{C/\mathcal Y}
 \xrightarrow{\ (dQ,dr)\ }
 \OO_{\PP^1}^{\oplus2}\longrightarrow0.
                                                        \tag{B.4.23.2}
\]
Its extension class vanishes because
$H^1(\PP^1,\OO(-1))=0$. Hence
\[
 N_{C/\mathcal Y}\simeq
 \OO(-1)\oplus\OO^{\oplus2},\qquad
 H^1(C,N_{C/\mathcal Y})=0,\qquad
 h^0(C,N_{C/\mathcal Y})=2.                             \tag{B.4.23.3}
\]
Moreover (\ref{push:src:B.4.23}.2) induces an isomorphism on global
sections, since both $H^0(\OO(-1))$ and $H^1(\OO(-1))$
vanish.

\item Kodaira's embedded-deformation theorem therefore gives a
smooth two-dimensional local parameter space $T$ of
compact embedded deformations of $C$, with tangent space
$H^0(C,N_{C/\mathcal Y})$. In precisely this relative
surface setting one may apply the directly opened
\citet[Proposition 2.3 and its proof, pp.~6--7 of the
original preprint]{DlousskyTeleman2012AB}: after shrinking
the neighbourhood of $C$, $(Q,r)$ is a holomorphic
submersion of relative dimension two, and $C$ is a
compact simply exceptional divisor of its fibre.
That proposition does not require the submersion to be
proper. The general completeness theorem is due to Kodaira; we use the exact statement in the opening of \citet{Merkulov1995D}.
Shrink this family so its curves
remain in the chosen neighbourhood. Their genus remains
zero, so all are isomorphic to $\PP^1$.
The functions $Q,r$ are constant on every compact
connected curve and thus define a holomorphic map
\[
 T\longrightarrow\C^2.
\]
Its derivative at $[C]$ is precisely the isomorphism on
sections induced by (\ref{push:src:B.4.23}.2). The holomorphic inverse
function theorem shows that the map covers a neighbourhood
of $(Q_0,0)$. In particular the nearby fixed pairs with
$r\ne0$ contain rational curves. A holomorphic map from
$\PP^1$ to a complex torus lifts to its vector-space
universal cover and is constant. Thus those fibres cannot
be the prescribed ordinary two-tori.

\item The conductor condition is essential. For example, if a
further central component meets $K$ along the persistent
elliptic axis, then its intersection number with $C$ is
one and $N_{K/\mathcal Y}|_C\simeq\OO(-1)$.
The normal sequence need no longer supply the $r$ direction
in (\ref{push:src:B.4.23}.2). The proposition therefore does not exclude
such a changed central arrangement; its proper construction
and comparison with the prescribed family remain separate
requirements.
\end{enumerate}
\end{proof}

\subsection{The associated logarithmic torsor at the missing origin}

\begin{proposition}[obstruction in the fixed log-torus category]\label{push:src:B.4.24}
\label{push:B424}
\StatementRecord{push:src:B.4.24}{proved}{Historical events: Claude Fable 5.2 (Anthropic), effort high, 16 November 2026: carried as background without fault found. Claude Fable 5.2 (Anthropic), effort high, 17 November 2026: carried as background without fault found. Claude Fable 5.2 (Anthropic), effort max, 16 November 2026: targeted proof check. GPT-6 Astra (OpenAI), ChatGPT Pro mode, effort Pro, S6\_09\_PAPERS, 17 November 2026: targeted proof check. Also checked in this writing. Pre-continuation writing events: GPT-6 Astra (OpenAI), ChatGPT Pro mode, effort Pro, S6\_10\_PAPERS, 18 November 2026: targeted proof check. Continuation metadata events: GPT-6 Astra (OpenAI), ChatGPT Pro mode, effort Pro, S6\_10\_PAPERS continuation, 19 November 2026: used and rechecked in this run; status-vocabulary reconciliation; metadata reconciliation only; no new mathematical check. \emph{Pre-continuation recorded proof scope.} Re-derived the explicit logarithmic period pairing, facewise positivity, two actual local origins and their transition F, and the injective valuation map that detects the nonzero gamma class. The obstruction remains nonzero after finite roots. Kato-Kajiwara-Nakayama source/application scopes remain separate.}{the hypotheses, notation and definitions stated here; \citet{KKN2008Root}; \cref{push:src:B.4.18}; \cref{push:src:B.4.19}; \citep[Definition 1.2.3, pp.~77--78]{KKN2008Root}; \citep[Definition 1.2.6, p.~79]{KKN2008Root}; \cref{push:src:B.4.9}; \citep[Sections 1.3.1--1.3.4, pp.~81--82]{KKN2008Root}; \citet[Sections 7--9]{KKN2008PicardE}; \cref{push:src:B.4.30}; \cref{push:src:B.4.33}}

The actual local Jacobian admits a nondegenerate logarithmic
complex-torus extension over the full $(q,r)$ bidisc and hence
has proper analytic models in the sense of
\citet{KKN2008Root}. After the normalized base change
$q=Q^{12}$, the two marked axial fan presentations of
\ref{push:src:B.4.18}--\ref{push:src:B.4.19} determine a torsor under its associated
log-torus sheaf on the punctured bidisc. This torsor does
not extend across the origin as a torsor locally trivial
for the usual analytic topology. The obstruction persists
after every further finite coordinate-root base change.

Here the torsor is formed from the \emph{associated
log-torus sheaves}, not from the fixed proper fan models
themselves, which need not be group objects. The conclusion
excludes an origin cap obtained from a locally trivial
torsor of this specified kind with the given axial
identifications. No implication from an arbitrary proper
analytic cap with vertical threefold components to this
torsor category is asserted.
\end{proposition}



\begin{proof}
\begin{enumerate}[label=\arabic*.,leftmargin=*]
\item \emph{The actual Jacobian pairing.}
Retain the period columns and logarithms of \ref{push:src:B.4.18}:
\[
 \gamma=(6\mu,C-s-12v)^t,\quad
 u=(\tau,\mu)^t,\quad w=(1,0)^t,\quad\delta=(0,1)^t,
 \qquad
 s=\frac{\log q}{2\pi i},\quad
 v=\frac{\log r}{2\pi i}.
\]
Here $\tau-s,\mu,C$ are holomorphic at $q=0$.
Use multiplicative characters
$Z=\exp(2\pi i\zeta_1)$ and
$W=\exp(-2\pi i\zeta_2)$, with character lattice
$X=\Z e_1\oplus\Z e_2$ and period lattice
$Y=\Z u\oplus\Z\gamma$. Put
\[
 a_0=e^{2\pi i(\tau-s)},\qquad
 a_1=e^{2\pi i\mu},\qquad a_2=e^{-2\pi iC}.
\]
These are holomorphic units on the full $q$ disc.
The multiplicative pairing is
\[
 \begin{array}{c|cc}
       &u&\gamma\\\hline
 e_1&q a_0&a_1^6\\
 e_2&a_1^{-1}&q r^{12}a_2.
 \end{array}
 \tag{B.4.24.1}
\]
Give the base the fine saturated log structure from
$qr=0$. In its sharp monoid $\N^2$, written additively,
the pairing modulo units is $\operatorname{diag}(q,q+12r)$.
We check the exact admissibility definition
\citep[Definition 1.2.3, pp.~77--78]{KKN2008Root}.
On the full face a nonnegative functional
$N(q)=h,N(r)=k$ gives
$\operatorname{diag}(h,h+12k)$. If $h>0$ it is
nondegenerate. If $h=0<k$, its kernel on each lattice
is the first basis line, exactly the face lattice
for the face generated by $q$; the quotient pairing
is multiplication by $12k$. For the $q$ face the
nonzero face lattices are these first basis lines
and the nontrivial quotient pairing is multiplication
by $h$. The $r$ face and the zero face have zero
face lattices. Zero functionals give zero-dimensional
quotients. These are all the face tests.

\item The residual unit condition is also satisfied
\citep[Definition 1.2.6, p.~79]{KKN2008Root}.
At $q=0$ the residual lattices vanish. At $r=0$,
$q\ne0$, their pairing is $\log|q a_0|\ne0$ after
shrinking the disc. At ordinary points the real
logarithmic matrix has determinant
\[
 \log|q a_0|\log|q r^{12}a_2|
                 +6\bigl(\log|a_1|\bigr)^2>0,
\]
after shrinking so both of the first two logarithms
are negative. Thus (\ref{push:src:B.4.24}.1) defines a log complex
torus $\mathcal A$ by Definition 1.3.4, pp.~81--82.
Theorems 5.2.1 and 5.3.2, pp.~171--172 and 180--181,
respectively supply a complete invariant fan with
finitely many $Y$-orbits and its proper analytic
model. This use of their proper-model theorem does
not require polarization.

\item \emph{The two extending origins.}
Set
\[
 B=\Delta_Q\times\Delta_r,\quad q=Q^{12},\qquad
 B_q=\Delta_Q^*\times\Delta_r,\quad
 B_r=\Delta_Q\times\Delta_r^*,\quad U=B_q\cap B_r.
\]
On $B_r$, \ref{push:src:B.4.9} supplies the actual affine section
\[
 \sigma=
 \left(6\mu v,\,-6v^2+v(C-12\kappa_0)\right)^t.
 \tag{B.4.24.2}
\]
For every local branch of $\log r$, both of its
multiplicative coordinates are holomorphic units
in the full $q$ disc. Its exact $r$-deck difference
is the bounded actual cusp closing translation
of \ref{push:src:B.4.9}. Hence it descends as a global section
of the actual affine family, including $q=0$.
Translation by $-\sigma$ identifies the prescribed
cusp fan presentation with its group presentation.
The normalized toric base change $q=Q^{12}$ preserves
this identification and its extending origin.

\item On $B_q$, the rotating model of \ref{push:src:B.4.18} is the
free cyclic quotient of its root fan model. Selecting
$Q$ recovers that model; more generally the normalized
base change recovers it also at the ramification
divisor. The twelfth fractional deck is
$\gamma-6\delta$, an actual period, so the
multiplicative period action is exactly the second
column of (\ref{push:src:B.4.24}.1). The point $(Z,W)=(1,1)$
extends in the root fan at $r=0$ and gives the
group origin. Thus both axial presentations have
actual extending origins in their associated
$\mathcal A$-sheaves. Refining a fan does not
replace the pairing or its associated quotient
sheaf; we use the explicit presentations, without
claiming that an arbitrary proper model determines
such a presentation.

\item In these two origins the exact actual comparison
$\zeta\mapsto\zeta-\sigma+F$ of \ref{push:src:B.4.18} becomes
translation by
\[
 F(s,v)=\left(
 \frac{\mu(s+12v)}2,
 \frac{C(s+12v)}{12}-\frac{(s+12v)^2}{24}+6v
 \right)^t.
 \tag{B.4.24.3}
\]
This defines the associated $\mathcal A$-torsor
by gluing trivial torsors on $B_q$ and $B_r$.
The differences are exactly
\[
 F(s+12,v)-F(s,v)=\gamma-6\delta,
 \qquad F(s,v+1)-F(s,v)=\gamma.
 \tag{B.4.24.4}
\]
They are periods, so $F$ is an $\mathcal A$-section
on $U$.

\item \emph{The ghost-valuation obstruction.}
By the quotient-sheaf definition
\citep[Sections 1.3.1--1.3.4, pp.~81--82]{KKN2008Root},
on the underlying analytic site there is an exact sequence
\[
 0\longrightarrow Y\longrightarrow
 \mathcal T_{\log}^{(Y)}\longrightarrow
 \mathcal A\longrightarrow0.
 \tag{B.4.24.5}
\]
The middle term is the sheaf of bounded multiplicative
log characters. On $U$ it is the ordinary
two-dimensional multiplicative torus sheaf.
In particular (\ref{push:src:B.4.24}.4) says that the connecting
class of $F$ in $H^1(U,Y)$ has $r$-loop component
$\gamma$. The terms in $\Z w+\Z\delta$ disappear
on exponentiation, and must not be included in $Y$.

\item Let $i:D_Q=\{Q=0\}\hookrightarrow B_r$.
Taking the $Q$-order of each character gives
\[
 \operatorname{ord}_Q:
 \mathcal T_{\log}^{(Y)}\longrightarrow i_*\Z^2,
 \qquad
 u\longmapsto(12,0),\quad
 \gamma\longmapsto(0,12)
 \tag{B.4.24.6}
\]
on its period subsheaf. If $h$ is any
$\mathcal A$-section on $B_r$, local lifts in
the middle sheaf of (\ref{push:src:B.4.24}.5) represent its
connecting class by a $Y$-cocycle. Applying
(\ref{push:src:B.4.24}.6) makes that cocycle a coboundary,
the difference of the local ghost vectors.
But restriction identifies $H^1(B_r,Y)$ with
$H^1(D_Q,Y)$, since $B_r$ retracts onto $D_Q$;
the induced map to $H^1(D_Q,\Z^2)$ is the
injective diagonal map $(12,12)$. Therefore
the connecting class of $h$ is zero. In
particular its $r$-loop component is zero.

\item For a section on $B_q$ that component is zero
as well, because the $r$ loop contracts in $B_q$.
Consequently $F$ cannot be the difference of
sections from $B_q$ and $B_r$. The glued torsor
is nontrivial on $B\setminus\{0\}$. An extension
locally trivial near the origin would have a
section there; after shrinking the bidisc its
two coordinate expressions would give precisely
the forbidden splitting. This proves nonextension.

\item For $q=Q^{12m}$ and $r=R^n$, with $m,n>0$,
the $R$-loop difference is
\[
 F(s,v+n)-F(s,v)
    =n\gamma-6n(n-1)\delta.
\]
Its $Y$-component is $n\gamma\ne0$,
whereas the $Q$-ghost map on $Y$ is
$\operatorname{diag}(12m,12m)$ and remains
injective. The same proof applies after every
such cover. Finite Kummer covers of the standard
coordinate log structure are locally dominated
by coordinate-root covers, so allowing them
does not remove this obstruction either.

\item \emph{The full gauge and the scope of the conclusion.}
There is no contradiction with \ref{push:src:B.4.19}'s identity
\[
 F-\sigma-kv\delta=G(s)+6v\delta,
 \quad k=-s+12\kappa_0,
 \quad G(s)=\left(\frac{\mu s}2,
                  \frac{Cs}{12}-\frac{s^2}{24}\right)^t.
 \tag{B.4.24.7}
\]
The ordinary gauge $kv\delta$ is not an extending
log section at $q=0$. Equation (\ref{push:src:B.4.24}.7) gives
the correct raw overlap and its fan extension
on $q\ne0$; it is not an allowed change of
both axial log-group origins. Changing either
origin by a section that does extend along its
own axis changes the clutching section by one
of the two terms already excluded by the
ghost calculation.

\item Finally, the cited proper-model theorem is a
construction from a specified log torus, not
a classification of arbitrary analytic caps.
The log Picard comparison in
\citet[Sections 7--9]{KKN2008PicardE} requires
explicit log Hodge and invariant-cycle hypotheses;
they have not been established for an arbitrary
non-K\"ahler cap here. Nor does (\ref{push:src:B.4.24}.5)
assert that every such cap is a locally trivial
torsor. Thus the obstruction has exactly the
scope stated. The additional bridge from hypothetical
arbitrary caps is proved later in \ref{push:src:B.4.30}--\ref{push:src:B.4.33},
using traces and a global radial-loop invariant.
The period differences and complete
gauge cancellation are checked in
\path{code/e_push_completion_diagnostics.py}.
\end{enumerate}
\end{proof}

\begin{proposition}[a one-eigenvalue Hopf boundary has no proper analytic completion]\label{push:src:B.4.25}
\label{push:B425}
\StatementRecord{push:src:B.4.25}{proved}{Historical events: Claude Fable 5.2 (Anthropic), effort high, 16 November 2026: carried as background without fault found. GPT-6 Astra (OpenAI), ChatGPT Pro mode, effort Pro, S6\_09\_PAPERS, 17 November 2026: targeted proof check. Also checked in this writing. Pre-continuation writing events: GPT-6 Astra (OpenAI), ChatGPT Pro mode, effort Pro, S6\_10\_PAPERS, 18 November 2026: targeted proof check. Continuation metadata events: GPT-6 Astra (OpenAI), ChatGPT Pro mode, effort Pro, S6\_10\_PAPERS continuation, 19 November 2026: used and rechecked in this run; status-vocabulary reconciliation; metadata reconciliation only; no new mathematical check. \emph{Pre-continuation recorded proof scope.} Re-derived anticanonical automorphy monomials, the resonance equations alpha=beta\textasciicircum{}N, the Rouche zeros approaching the cusp, and the coherent-direct-image base-change contradiction. Checked the proper-resolution reduction and persistence after finite covers. Matched Grauert1960 HauptsatzIp59 coherence and Section7pp62-64 base-change hypotheses; read1963erratumpp35-36. The direct-image proof does not use the corrected exceptional-set closedness assertion.}{the hypotheses, notation and definitions stated here; \cref{push:src:B.4.21}; \cref{push:src:B.4.19}; \citep[Section 6, Hauptsatz I, p.~59]{Grauert1960D}; \citet[Theorem 1.6 and the noncompact analytic
extension immediately following it]{BM95AB}}

Let $\alpha(q)=q^m\kappa(q)$, where $m\ge1$ and
$\kappa$ is a holomorphic unit on a disc, and let $\beta$
be holomorphic with $0<|\beta(0)|<1$. Shrink the disc so
$\beta$ is a unit and $0<|\alpha(q)|<|\beta(q)|<1$ for
$q\ne0$. Consider the diagonal primary Hopf surfaces
\[
 H_q=(\C^2\setminus\{0\})/
       \langle(z,w)\mapsto(\alpha(q)z,\beta(q)w)\rangle.
                                                        \tag{B.4.25.1}
\]
Their holomorphic family over the punctured disc has no
proper complex-analytic completion over the full disc,
even allowing singular or nonreduced central fibres and
a singular total space. This remains true after every
finite base change ramified at zero.

In particular one cannot replace the blown-up Hopf
fibres in \ref{push:src:B.4.21} by the unblown family with one multiplier
$\kappa(q)q\to0$ and the other tending to a nonzero
contraction. This concerns that specified auxiliary
Hopf family. It is not a nonexistence theorem for a
completion of the actual two-torus family in \ref{push:src:B.4.19}.
\end{proposition}


\begin{proof}
\begin{enumerate}[label=\arabic*.,leftmargin=*]
\item \emph{Anticanonical sections of a diagonal Hopf surface.}
A section of $K_{H_q}^{-1}$ lifts to
$f(z,w)\,\partial_z\wedge\partial_w$ on
$\C^2\setminus\{0\}$. Hartogs extension makes $f$
entire on $\C^2$. Deck equivariance is exactly
\[
 f(\alpha z,\beta w)=\alpha\beta f(z,w).
                                                        \tag{B.4.25.2}
\]
For the Taylor coefficient of $z^i w^j$, with $i,j\ge0$,
this says
\[
 \alpha^{i-1}\beta^{j-1}=1.                             \tag{B.4.25.3}
\]
When $i,j\ge1$, the contraction inequalities force
$i=j=1$. When $j=0$, (\ref{push:src:B.4.25}.3) would give
$\beta=\alpha^{i-1}$, which is impossible under
$|\alpha|<|\beta|<1$ (and also for $i\le1$).
When $i=0$, the only possibility is
\[
 \alpha=\beta^N,\qquad j=N+1,\qquad N\ge2.
                                                        \tag{B.4.25.4}
\]
At a fixed parameter there is at most one such $N$,
since $0<|\beta|<1$. Thus
\[
 h^0(H_q,K_{H_q}^{-1})=
 \begin{cases}
 2,&\alpha(q)=\beta(q)^N\text{ for some }N\ge2,\\
 1,&\text{otherwise}.
 \end{cases}                                           \tag{B.4.25.5}
\]
The sections in the first case are represented by
$zw$ and $w^{N+1}$. The computation uses actual
bivector automorphy, with no choice of Picard convention.

\item \emph{Infinitely many accumulating resonances.}
Put $b=|\beta(0)|\in(0,1)$ and $k=|\kappa(0)|>0$.
For all sufficiently small $q$ we may arrange
$|\kappa(q)|\ge3k/4$ and
\[
 |\beta(q)|\le b\exp(C|q|)
\]
for some fixed $C$, by a holomorphic logarithm of
$\beta(q)/\beta(0)$. Set
\[
 R_N=\left(\frac{2b^N}{k}\right)^{1/m}.
\]
Then $R_N\to0$ and $N R_N\to0$. On $|q|=R_N$, for
all sufficiently large $N$,
\[
 |q^m\kappa(q)|\ge\frac32 b^N,
 \qquad |\beta(q)^N|\le b^N\exp(CN R_N)<\frac54 b^N.
\]
Rouch\'e's theorem shows that
$q^m\kappa(q)-\beta(q)^N$ has $m$ zeros counted with
multiplicity in $|q|<R_N$. None is zero. Choose one,
$q_N$. Different $N$ give different points, since a
nonzero contraction cannot have two equal distinct
positive powers. Hence there are infinitely many
resonances $q_N\to0$ with the first value in
(\ref{push:src:B.4.25}.5).

\item For each $N$ the resonance function is nonzero at
$q=0$, so is not identically zero. Its zeros in the
punctured disc are discrete. The countable union of
these discrete sets cannot contain an open disc.
There are therefore also nonresonant points arbitrarily
near zero, with the second value in (\ref{push:src:B.4.25}.5).

\item \emph{Coherent direct images forbid this pattern.}
Suppose first that a proper completion
$\pi:\mathcal X\to\Delta$ has a smooth connected
total threefold and agrees with (\ref{push:src:B.4.25}.1) over
$\Delta^*$. Put $L=K_{\mathcal X}^{-1}$.
On each nonzero fibre, adjunction identifies
$L|_{\mathcal X_q}$ with $K_{H_q}^{-1}$, because
the normal direction is the trivial base direction.
Grauert's proper direct-image theorem gives coherent
sheaves
\[
 M=\pi_*L,\qquad N=R^1\pi_*L
\]
on the disc
\citep[Section 6, Hauptsatz I, p.~59]{Grauert1960D}.
Multiplication by a nonzero base function is injective
on $L$, and hence on $M$, since the total space is
smooth and the map is nonconstant. Thus $M$ is
locally free of a fixed rank on the disc. The torsion
subsheaf of $N$ is coherent with locally finite
support. After shrinking, it has no support in
$\Delta^*$.

\item For any $a\in\Delta^*$, apply the direct-image long
exact sequence to
\[
 0\longrightarrow L\xrightarrow{\ q-a\ }L
 \longrightarrow L|_{\mathcal X_a}\longrightarrow0.
\]
At $a$ its beginning yields
\[
 0\longrightarrow M_a/(q-a)M_a
 \longrightarrow H^0(\mathcal X_a,L|_{\mathcal X_a})
 \longrightarrow \ker(q-a:N_a\to N_a)
 \longrightarrow0.                                   \tag{B.4.25.6}
\]
The last term vanishes. Therefore
$h^0(H_a,K_{H_a}^{-1})=\operatorname{rank}M$ for
\emph{every} sufficiently small $a\ne0$. This
contradicts (\ref{push:src:B.4.25}.5) and the accumulating resonances.
This argument uses coherence and a direct exact
sequence; it does not use the corrected closedness
assertion discussed in the Grauert source audit.

\item Now suppose an arbitrary proper analytic completion
$\mathcal Z\to\Delta$ exists. Pass to its reduction
and retain the unique irreducible component containing
the connected smooth punctured total space. This is a
closed analytic component, remains proper over the
disc, and dominates the disc. Its restriction over
$\Delta^*$ is exactly the given family. Resolve this
component by a proper analytic resolution chosen to
be an isomorphism on its already nonsingular locus;
see \citet[Theorem 1.6 and the noncompact analytic
extension immediately following it]{BM95AB}.
All singular points lie over zero, so the resolution
is an isomorphism over the entire punctured family.
Properness of the composition is preserved. Near the
compact central fibre, local finiteness of the
analytic resolution permits one shrinking of the
disc; its dominating smooth component is a smooth
proper total threefold as in the preceding paragraph.
Its map to the disc is nonconstant, so the base
function is a non-zero-divisor. No vertical component
was retained which could invalidate that assertion.
The preceding contradiction applies. Reduction also
shows that nilpotent structure supported at zero
cannot evade it.

\item Finally, a finite base change replaces $\alpha$ by
a holomorphic function of some positive vanishing
order and replaces $\beta$ by a unit whose value
at zero still has modulus strictly between zero and
one. The same proof applies verbatim.
\end{enumerate}
\end{proof}


\paragraph{Scope and comparison with the classical toric Hopf degeneration}
There is no conflict with the degeneration in
\citet[Section 13, pp.~116--118]{Oda1978D}.
There the boundary rays are generated by $n$ and
$-n+a n'$ with $a>0$, and the period is translation
by $\lambda^{n'}$. In the universal-cover basis
\[
 v_1=n,\qquad v_2=-n+a n',\qquad
 n'=\frac{v_1+v_2}{a},
\]
the base change $\lambda=t^a$ makes this period
act with eigenvalues $(t,t)$, up to fixed torus
units in the indicated variant. Both eigenvalues
tend to zero. \ref{push:src:B.4.25} assumes instead a nonzero
limit for the second eigenvalue. The distinction
is forced by the exact period marking of the
candidate Tate interface, and cannot be removed
by citing the existence of some Hopf degeneration.
The blown-up Kato family of \ref{push:src:B.4.10} is also consistent
with \ref{push:src:B.4.25}: its generic fibres are blow-ups of
$H_q$, so the anticanonical sections must satisfy
the extra vanishing conditions at the marked
blow-up centre and are not the spaces in
(\ref{push:src:B.4.25}.5).


\Uses{the hypotheses, notation and definitions stated here; \citet[Section 13, pp.~116--118]{Oda1978D}; \cref{push:src:B.4.25}; \cref{push:src:B.4.10}}


\subsection{A necessary vertical fibre over the missing origin}

\begin{proposition}[every origin component is a threefold]\label{push:src:B.4.26}
\label{push:B426}
\StatementRecord{push:src:B.4.26}{proved}{Historical events: Claude Fable 5.2 (Anthropic), effort high, 16 November 2026: carried as background without fault found. Claude Fable 5.2 (Anthropic), effort max, 16 November 2026: targeted proof check. GPT-6 Astra (OpenAI), ChatGPT Pro mode, effort Pro, S6\_09\_PAPERS, 17 November 2026: targeted proof check. Also checked in this writing. Pre-continuation writing events: GPT-6 Astra (OpenAI), ChatGPT Pro mode, effort Pro, S6\_10\_PAPERS, 18 November 2026: targeted proof check. Continuation metadata events: GPT-6 Astra (OpenAI), ChatGPT Pro mode, effort Pro, S6\_10\_PAPERS continuation, 19 November 2026: used and rechecked in this run; status-vocabulary reconciliation; metadata reconciliation only; no new mathematical check. \emph{Pre-continuation recorded proof scope.} Re-derived the finite multisection trace including normalized ramification and explicit multiplicative norms, the induced nonzero degree-times-gamma obstruction, and the Weierstrass finite-cut argument. Log-torsor extension inputs remain separately tracked; the receiving formulation will state the actual exclusion directly.}{the hypotheses, notation and definitions stated here; \cref{push:src:B.4.19}; \cref{push:src:B.4.30}; \cref{push:src:B.4.33}; \cref{push:src:B.4.24}; \citet[Chapter II, Sections 7--8]{Demailly2012E}}

Suppose the actual partial family of \ref{push:src:B.4.19} has a
proper analytic completion $f:\mathcal Z\to\Delta_q\times\Delta_r$
with smooth pure four-dimensional total space, with
$r\circ f$ flat, and with its prescribed identification
over the punctured bidisc. The map to this fixed bidisc
is assumed holomorphic. Then every irreducible component
of $f^{-1}(0,0)$ has dimension three. In particular,
there is no equidimensional completion, and no completion
having even one two-dimensional component of the origin
fibre.

The conclusion still permits the vertical threefold
components allowed in B.4. It is not an assertion that
an arbitrary $r$-proper filling initially has a holomorphic
extension of $q$; if $q$ is only meromorphic, a graph
modification must first be supplied and its resulting
geometry retained. \ref{push:src:B.4.30} subsequently excludes even
the permitted vertical-component completions of the
fixed partial family, and \ref{push:src:B.4.33} excludes every proper
filling of the actual global family.
\end{proposition}



\begin{proof}
\begin{enumerate}[label=\arabic*.,leftmargin=*]
\item We first prove the finite-multisection consequence of
the actual log-torsor obstruction \ref{push:src:B.4.24}, without
assuming that the unknown cap is itself a log torsor.
Make the normalized finite base change $q=Q^{12}$.
The resulting space $\mathcal Z'$ is normal and
pure four-dimensional. Its restrictions to the two
punctured axes have the explicit normalized fan
presentations used in \ref{push:src:B.4.24}. Write
\[
 B=\Delta_Q\times\Delta_r,\quad
 B_q=\Delta_Q^*\times\Delta_r,\quad
 B_r=\Delta_Q\times\Delta_r^*,\quad
 U=B_q\cap B_r.
\]

\item Assume there is a finite surjective map
$p:V\to B$ of positive generic degree $d$, together
with a holomorphic map $V\to\mathcal Z'$ over $B$.
We may normalize $V$ and retain any collection of
components dominating $B$; its degree remains a
positive integer. A normal surface is Cohen--Macaulay;
finite over the smooth bidisc, it is flat. Thus the
finite holomorphic trace and the degree-$d$ sum with
multiplicities are available also at branch points.
On the ordinary locus, sum the
$d$ points of this multisection, with multiplicity,
in each of the two actual axial group origins of
\ref{push:src:B.4.24}. This gives two local expressions for the
trace in the Jacobian group. We show that they
extend as log-torus sections on $B_q$ and $B_r$.

\item \emph{Ramification away from the axes.}
Over the ordinary locus the torus family is a smooth
holomorphic group family. Near a branch point, choose
small neighbourhoods of the finitely many preimages
in $V$ and local Lie lifts of their maps to the torus.
The trace of these holomorphic functions under the
finite map is holomorphic, and its exponential is
the ordinary sum. Changes of lifts add periods,
which disappear in the group quotient. Thus the
ordinary trace extends across the ordinary branch
locus. Equivalently this is the usual addition map
on the relative symmetric power, applied to the
finite multisection.

\item \emph{The explicit norm at an axis.}
After shrinking $B$, every component of the finite
cover's branch divisor which is not an axis meets
each punctured axis nowhere: its intersection with
either axis is an isolated analytic germ at the
origin. At a point of a punctured axis choose
coordinates $(x,y)$, with $x$ the normal coordinate
to that axis. Locally the normal finite cover is
a disjoint union of covers
\[
 x=t^e,\qquad y=y,
 \tag{B.4.26.1}
\]
including $e=1$ sheets. Indeed the complement of
the axis in a small polydisc has cyclic fundamental
group; each connected unramified finite cover is
the displayed root cover there, and its normal
finite extension is unique. No assumption that the
whole finite cover is a coordinate-root cover is
being made.

\item The map from such a root polydisc into the
prescribed fan model gives a section of its
associated log-torus sheaf. It has a local bounded
multiplicative log lift. This can also be seen
directly in a toric chart: the pulled-back characters
are meromorphic functions with zeros and poles
supported on $t=0$, hence are $t^k u(t,y)$ with
$k\in\Z$ and $u$ a holomorphic unit. The quotient
by the period lattice causes only the period
ambiguity already retained in \ref{push:src:B.4.24}.

\item The multiplicative norm of a character on (\ref{push:src:B.4.26}.1)
is explicitly
\[
 \prod_{j=0}^{e-1}(\zeta_e^j t)^k
              u(\zeta_e^j t,y)
 =(-1)^{k(e-1)}x^k
              \prod_{j=0}^{e-1}u(\zeta_e^j t,y).
 \tag{B.4.26.2}
\]
The last product is a holomorphic unit in $(x,y)$.
Thus the norm has an integral ghost vector and
is an actual log character on the base. It is
bounded in the required sense: on the $r$ axis
the first character has zero normal order and
the second is unrestricted; on the $Q$ axis
both are unrestricted. These are precisely the
bounded-character conditions for (\ref{push:src:B.4.24}.1).
Multiplication of lifts corresponds to addition
in the log torus. Changing a lift by a period
changes its norm by an integral period, so the
resulting quotient section is independent of
the lift. Summing over the finitely many covers
(\ref{push:src:B.4.26}.1), whose degrees add to $d$, proves the
claimed extension along that axis.

\item The local extensions agree on overlaps because
they agree on the ordinary locus. To see uniqueness
without a compactification assumption, lift locally
to log characters; agreement there is an equality
up to a locally constant $Y$-period, and that
meromorphic equality extends across the axis.
We have therefore obtained actual sections
$h_q$ of $\mathcal A$ on $B_q$ and $h_r$ on
$B_r$. The change between the two origins for
each point is $F$, so on $U$
\[
 h_q-h_r=dF
 \tag{B.4.26.3}
\]
up to the harmless overall choice of subtraction
convention. The ghost proof of \ref{push:src:B.4.24} forbids
(\ref{push:src:B.4.26}.3): its $r$-loop $Y$-component is
$d\gamma\ne0$. Thus no such finite multisection
can exist.

\item \emph{A two-dimensional fibre germ would supply one.}
Suppose an origin point of $\mathcal Z$ has local
fibre dimension two. Finite base change and
normalization are finite maps and preserve that
local fibre dimension at points lying above it.
At a resulting point of $\mathcal Z'$, choose
two generic local linear sections in an analytic
embedding. Their intersection with the fibre is
isolated. Intersecting with $\mathcal Z'$ gives
a two-dimensional analytic germ whose map to
$B$ is finite. This is the local finite-projection
criterion from Weierstrass preparation; see also
\citet[Chapter II, Sections 7--8]{Demailly2012E}.
Choose a component through the point. Its finite
image has dimension two and is therefore the
whole base germ. A finite representative of this
germ, after shrinking the bidisc, is finite and
surjective over the entire smaller bidisc. Its
normalization is still finite and maps to
$\mathcal Z'$, giving exactly the forbidden
multisection of positive degree. Properness of
this finite representative prevents loss of
sheets at its boundary.

\item Every fibre component of a dominant map from a
pure fourfold to a smooth surface has dimension
at least two. If a two-dimensional irreducible
component existed, a generic point on it would
have local fibre dimension two, which we have
excluded. Flatness over $r$ makes $r=0$ a
Cartier divisor of pure dimension three, so
the origin fibre has no component of dimension
greater than three. All its irreducible
components consequently have dimension three,
as asserted.
\end{enumerate}
\end{proof}

\subsection{The obstruction survives modifications of the base}

\begin{proposition}[no log-torsor repair by point blowups]\label{push:src:B.4.27}
\label{push:B427}
\StatementRecord{push:src:B.4.27}{proved}{Historical events: Claude Fable 5.2 (Anthropic), effort high, 16 November 2026: carried as background without fault found; targeted proof check. Claude Fable 5.2 (Anthropic), effort max, 16 November 2026: carried as background without fault found. GPT-6 Astra (OpenAI), ChatGPT Pro mode, effort Pro, S6\_09\_PAPERS, 17 November 2026: targeted proof check. Also checked in this writing. Pre-continuation writing events: GPT-6 Astra (OpenAI), ChatGPT Pro mode, effort Pro, S6\_10\_PAPERS, 18 November 2026: targeted proof check. Continuation metadata events: GPT-6 Astra (OpenAI), ChatGPT Pro mode, effort Pro, S6\_10\_PAPERS continuation, 19 November 2026: used and rechecked in this run; status-vocabulary reconciliation; metadata reconciliation only; no new mathematical check. \emph{Pre-continuation recorded proof scope.} Read complete receiving proof and re-derived displayed sheaf/topological argument. The external logarithmic comparison applications in B4.30/B4.32/B4.33 were audited separately against IKN and its erratum; the independent direct commuting-meridian proof is night:P2.7. Boundary proof checked using intrinsic anticanonical divisor, coherence and proper-image exceptional loci.}{the hypotheses, notation and definitions stated here; \cref{push:src:B.4.24}; \citep[Lemma 1.2.9, pp.~80--81]{KKN2008Root}; \cref{push:src:B.4.30}; \cref{push:src:B.4.33}}

Let $B=\Delta_Q\times\Delta_R$, with
$q=Q^{12m}$ and $r=R^n$, where $m,n$ are positive
integers. Let $\beta:S\to B$ be any finite sequence
of point blowups whose centres lie over the origin.
Give $S$ the divisorial log structure of the reduced
inverse image of $QR=0$, and let $\mathcal A_S$ be
the pullback of the actual Jacobian log torus of
\ref{push:src:B.4.24}. Then
\[
 H^1(S,\mathcal A_S)=0
 \tag{B.4.27.1}
\]
for the usual analytic topology. In particular the
actual associated axial log-torus torsor on the
punctured bidisc has no locally trivial
$\mathcal A_S$-torsor extension after these
modifications.

This result concerns the specified associated
log-torus sheaf. It does not assert that flattening
or resolving an arbitrary analytic origin cap
produces a torsor under that sheaf.
\end{proposition}



\begin{proof}
\begin{enumerate}[label=\arabic*.,leftmargin=*]
\item Write $E_1,\ldots,E_N$ for the exceptional curves,
and $D_Q,D_R$ for the strict transforms of the
two coordinate axes. The reduced total boundary
is normal crossing. Every $E_i$ is a rational
curve, and the exceptional divisor is a connected
tree if $N>0$. Its two orders
\[
 a_i=\operatorname{ord}_{E_i}(Q),\qquad
 b_i=\operatorname{ord}_{E_i}(R)
 \tag{B.4.27.2}
\]
are positive: the image of $E_i$ is the origin.
Set
\[
 D_1=D_Q+\sum_i E_i,\qquad
 D_2=D_Q+D_R+\sum_i E_i.
\]
For a reduced normal-crossing divisor $D$, denote
by $\mathcal M_D^{\mathrm{gp}}$ the sheaf of
meromorphic units whose divisors are supported
on $D$.

\item \emph{The exact bounded-character subsheaf.}
The period pairing remains (\ref{push:src:B.4.24}.1), pulled
back to $S$. By the base-change lemma
\citep[Lemma 1.2.9, pp.~80--81]{KKN2008Root}
it defines $\mathcal A_S$. More explicitly its
bounded multiplicative-character sheaf is
\[
 \mathcal T_{\log}^{(Y)}
   =\mathcal M_{D_1}^{\mathrm{gp}}
          \times\mathcal M_{D_2}^{\mathrm{gp}}.
 \tag{B.4.27.3}
\]
Indeed the first character is bounded by powers
of $q$, and therefore has zero order on $D_R$.
It may have arbitrary integer order on $D_Q$
and on each $E_i$, where $q$ has positive order.
The second character is bounded by powers of
$q r^{12}$, whose order is positive on every
boundary component, so all its boundary orders
are permitted. These individual conditions
also suffice for every character combination:
if its second coefficient is nonzero, a suitable
signed multiple of the second period bounds
both local normal orders; otherwise powers of
the first period do so on the permitted support.
At a crossing one chooses the bounding powers
simultaneously for its two integer orders.
This proves (\ref{push:src:B.4.27}.3), including at crossings
and at intersections with $D_R$; replacing the
bounded sheaf by the full log-character sheaf
in both coordinates would be incorrect.

\item \emph{Cohomology of the blown-up bidisc.}
We record the elementary ingredients to fix
the analytic scope. For a point blowup $p$ of
a smooth surface,
$p_*\mathcal O=\mathcal O$ and
$R^j p_*\mathcal O=0$ for $j>0$.
The first identity follows from Hartogs.
For the second, the two standard blowup charts
give a local Laurent calculation: on their
intersection a holomorphic function expands
as a convergent sum of $u^k z^\ell$ with
$k\ge0$. Terms with $\ell\ge0$ extend in
the first chart. Terms with $\ell<0$ extend
in the second, because in coordinates
$z'=1/z$, $u'=uz$ they are
$(u')^k(z')^{k-\ell}$. Cauchy estimates on
smaller coordinate polydiscs preserve convergence
of the two sums. The charts and their
intersection are Stein, so this kills the
first cohomology and there is no higher
\v Cech cohomology. Iterating and using the
Leray sequence over the Stein bidisc gives
\[
 H^j(S,\mathcal O_S)=0\qquad(j>0).
 \tag{B.4.27.4}
\]

\item The exceptional tree has one $H^2(\ ,\Z)$
generator per rational component and no positive
cohomology in other degrees. Proper base change
and Leray for $\beta$ therefore give
\[
 H^2(S,\Z)\xrightarrow{\ \sim\ }
       \bigoplus_i H^2(E_i,\Z)\simeq\Z^N,
 \qquad H^j(S,\Z)=0\quad(j=1,3,4).
 \tag{B.4.27.5}
\]
Equivalently these groups follow from the usual
tubular plumbing of the exceptional tree.
The exceptional intersection matrix is
unimodular: at each blowup the total-transform
basis adds a vector of square $-1$ orthogonal
to the preceding basis. Passing to the strict
transform basis is an integral unimodular
change of basis. Thus the Chern classes of
$\mathcal O(E_i)$ form a basis of
$H^2(S,\Z)$. The exponential sequence and
(\ref{push:src:B.4.27}.4)--(\ref{push:src:B.4.27}.5) give
\[
 \operatorname{Pic}(S)\simeq\Z^N
       =\langle\mathcal O(E_i)\rangle,
 \qquad
 H^2(S,\mathcal O_S^*)=H^3(S,\mathcal O_S^*)=0.
 \tag{B.4.27.6}
\]

\item For either $D=D_1$ or $D_2$ the divisor-order
sequence is exact:
\[
 0\longrightarrow\mathcal O_S^*
 \longrightarrow\mathcal M_D^{\mathrm{gp}}
 \longrightarrow
   \bigoplus_{C\subset D}i_{C*}\Z
 \longrightarrow0.
 \tag{B.4.27.7}
\]
The map from global integer orders to the Picard
group sends a component to its divisor line
bundle. It is surjective because both divisors
contain every $E_i$. Every component has
$H^1(C,\Z)=0$: the strict transforms of the
axes are discs and the exceptional curves are
rational. Their only nonzero $H^2$ contributions
are the $N$ exceptional curves. Consequently
the long exact sequence of (\ref{push:src:B.4.27}.7) gives
\[
 H^1(S,\mathcal M_{D_j}^{\mathrm{gp}})=0,
 \qquad
 H^2(S,\mathcal M_{D_j}^{\mathrm{gp}})
       \simeq\bigoplus_i H^2(E_i,\Z)
       \simeq\Z^N\quad(j=1,2).
 \tag{B.4.27.8}
\]
The second isomorphism is induced by the actual
divisor-order map, a fact needed for the next
calculation.

\item \emph{Injectivity of the period map in degree two.}
The quotient definition of $\mathcal A_S$ gives
\[
 0\longrightarrow Y\longrightarrow
 \mathcal T_{\log}^{(Y)}\longrightarrow
 \mathcal A_S\longrightarrow0,
 \qquad Y=\Z u\oplus\Z\gamma.
 \tag{B.4.27.9}
\]
By (\ref{push:src:B.4.27}.3) and (\ref{push:src:B.4.27}.8), its middle
sheaf has zero $H^1$ and $H^2\simeq\Z^{2N}$.
Using restriction to $E_i$ as the coordinates
in (\ref{push:src:B.4.27}.5), the map
\[
 H^2(S,Y)\longrightarrow
 H^2(S,\mathcal T_{\log}^{(Y)})
\]
is block diagonal. Its block at $E_i$ is the
valuation matrix of the two period columns,
\[
 \begin{pmatrix}
 12m a_i&0\\0&12m a_i+12n b_i
 \end{pmatrix}.
 \tag{B.4.27.10}
\]
Both diagonal entries are positive. Thus this
map is injective over $\Z$, not merely over
$\Q$. The long exact sequence of
(\ref{push:src:B.4.27}.9) now gives (\ref{push:src:B.4.27}.1).

\item Put $U_S=\beta^{-1}(B\setminus\{0\})$.
The actual torsor on $U_S\simeq B\setminus\{0\}$
is nontrivial by \ref{push:src:B.4.24}, including after the
stated finite coordinate roots. An extending
torsor on $S$ would be trivial by (\ref{push:src:B.4.27}.1),
and its restriction would be trivial as well,
a contradiction. For $N>0$, $U_S$ is the
complement of the exceptional union; for
$N=0$ it is the punctured bidisc. Thus the
case $N=0$ is included as well.

\item No step supplies a log-group presentation of
an arbitrary analytic cap. In particular a
flattening of its map to the bidisc changes
the base, but does not by itself produce
(\ref{push:src:B.4.27}.9) for that unknown space. The result
therefore removes a specific logarithmic
repair by base blowups. \ref{push:src:B.4.30}--\ref{push:src:B.4.33} supply the
additional extending-section and global-loop arguments
that exclude an arbitrary proper filling; no such
comparison is assumed in this proposition.
The valuation blocks and the unimodular
intersection matrices for the first 124 marked
point-blowup sequences are independently checked
in \path{code/e_push_log_blowup.py}. This finite
audit supplements the all-length argument above.
\end{enumerate}
\end{proof}

\begin{proposition}[maximal monodromy with trivial central dualizing bundle]\label{push:src:B.4.28}
\StatementRecord{push:src:B.4.28}{proved}{Historical events: Claude Fable 5.2 (Anthropic), effort high, 16 November 2026: carried as background without fault found. GPT-6 Astra (OpenAI), ChatGPT Pro mode, effort Pro, S6\_09\_PAPERS, 17 November 2026: targeted proof check. Also checked in this writing. Pre-continuation writing events: GPT-6 Astra (OpenAI), ChatGPT Pro mode, effort Pro, S6\_10\_PAPERS, 18 November 2026: targeted proof check. Continuation metadata events: GPT-6 Astra (OpenAI), ChatGPT Pro mode, effort Pro, S6\_10\_PAPERS continuation, 19 November 2026: used and rechecked in this run; status-vocabulary reconciliation; metadata reconciliation only; no new mathematical check. \emph{Pre-continuation recorded proof scope.} Re-derived the anticanonical component equation, Euler-charge reduction, normal-degree balance, rank-one bound for the elliptic-conductor graph, and hyperbolic monodromy boundary homology obstruction. Checked the nonzero residue integral on the Clifford torus and the remaining orientable toric dual complex. Classification and charge source locators remain separately tracked. Closed the charge, NAC coefficient-one classification, half-Inoue exclusion and hyperbolic complement input locators: FriedmanLemmas1.2/2.2/2.7pp5/11/14;DlousskyCor2.25/Lem2.26/Prop2.27pp16-19;Fujiki-PontecorvoSection3pp20-22 and Lem3.2proofp68;OdaTheorem15.1andfollowingpp130-131. General surface classification and canonical-base-change scopes remain in the source inventory.}{the hypotheses, notation and definitions stated here; \citet[Definition 1.1 and Lemmas 1.2, 2.2, 2.7]{Friedman2016AB}; \citet[Corollary 2.25, Lemma 2.26 and Proposition
2.27]{Dloussky2005AB}; \citet[Sections 3.2--3.3]{FujikiPontecorvo2010AB}; \citet[Section 15, Theorem 15.1 and the following
paragraph, pp.~130--131]{Oda1978D}; \citet[Chapter III, Theorem 4.12(ii)]{BS76}; \citep[Chapter III, Theorem 4.12(iii)]{BS76}}

Let $p:\mathcal X\to\Delta$ be a proper holomorphic map
from a smooth complex threefold, smooth over $\Delta^*$,
whose nonzero fibres are complex two-tori. Suppose that
$D=p^{-1}(0)$ is a reduced normal-crossing divisor and
\emph{assume in addition} that $\omega_D\simeq\OO_D$.
If its monodromy on $H_1(\mathcal X_t,\Q)$ satisfies
\[
 \operatorname{rank}(T-1)=2,                         \tag{B.4.28.1}
\]
then every normalization component of $D$ is a smooth
rational toric surface, its conductor is its toric
boundary, and the dual complex of $D$ is a topological
two-torus. In particular all normalized components are
K\"ahler.

No K\"ahler hypothesis on $\mathcal X$ is used. The
additional triviality of $\omega_D$ is not inferred from
the neighbouring two-tori. The proposition does not
produce a modification with this property from an
arbitrary analytic completion.
\end{proposition}


\begin{proof}
\begin{enumerate}[label=\arabic*.,leftmargin=*]
\item Write $V_i$ for the compact smooth normalization
components and $C_i$ for their reduced conductor divisors.
Adjunction gives
\[
 K_{V_i}+C_i\sim0.                                  \tag{B.4.28.2}
\]
If $D$ is smooth, the map is a proper submersion near
its compact central fibre and its monodromy is trivial.
Otherwise every $C_i$ is nonempty, by connectedness.
Adjunction on a conductor branch says that it is either
a smooth elliptic curve disjoint from the other branches,
or a rational curve with two boundary points, counted
on its normalization. The rational branches form cycles,
including the nodal irreducible and two-component cases.

\item \emph{Nonnegative Euler contributions.}
A nonempty effective anticanonical divisor makes every
pluricanonical system empty: multiplying a nonzero
pluricanonical section by the appropriate power of the
anticanonical section would give a nonzero holomorphic
function with a nonempty zero divisor. Thus the surface
has Kodaira dimension $-\infty$. Its minimal model is
rational, ruled, or of class VII. We spell out the
consequences for the reduced anticanonical pairs here.

\item On a rational surface, the conductor is connected, by
$0\to K_V\to\OO_V\to\OO_C\to0$ and $H^1(K_V)=0$.
For a rational cycle, the charge identity and positivity
are
\[
 e(V\setminus C)=12-C^2-\#\{\text{cycle components}\}
 \ge0,\qquad
 e(V\setminus C)=0\ \Longleftrightarrow\ (V,C)
 \text{ is toric}.                                  \tag{B.4.28.3}
\]
Here $e$ denotes compactly supported Euler characteristic.
This includes the nodal irreducible convention:
\citet[Definition 1.1 and Lemmas 1.2, 2.2, 2.7]{Friedman2016AB}.
For a smooth elliptic anticanonical curve,
$e(V\setminus C)=2+b_2(V)>0$ instead.

\item On a nonrational ruled minimal model, a conductor
component dominating the base has degree at most two
and genus at most one. The base is therefore elliptic.
The minimal conductor is two disjoint elliptic sections
or one unramified elliptic bisection. Its Euler
contribution is zero. The two section squares sum to
zero, and the bisection has square zero, also directly
from $K_V+C\sim0$ and $K_V^2=0$.

\item A minimal class-VII surface with $b_2=0$ and a curve is
a Hopf surface. Its elliptic curves have square zero,
and a reduced anticanonical divisor has at most two
elliptic components. For this last assertion the exact
sequence $0\to K_V\to\OO_V\to\OO_C\to0$ gives
$h^0(\OO_C)\le1+h^1(K_V)=2$.
The Euler contribution is zero. This also covers
secondary Hopf surfaces; no primary marking is required.

\item For a minimal class-VII surface with $b_2>0$, the
reduced anticanonical divisor is a numerically
anticanonical divisor whose coefficients are all one.
\citet[Corollary 2.25, Lemma 2.26 and Proposition
2.27]{Dloussky2005AB} imply that it is parabolic Inoue
or Inoue--Hirzebruch. Half-Inoue has only a nontrivially
twisted anticanonical divisor and cannot satisfy
(\ref{push:src:B.4.28}.2). The remaining possibilities are hyperbolic
Inoue with two disjoint rational cycles, or parabolic
Inoue with one rational cycle and a disjoint elliptic
curve. If $n=b_2>0$, there are $n$ rational cycle
components and $e(V)=n$, so the contribution is zero.
In the parabolic case the elliptic curve has square
$-n$. The curve configurations, anticanonical identities
and underlying $(S^1\times S^3)\#n\overline{\PP}^2$
topology are recorded in
\citet[Sections 3.2--3.3]{FujikiPontecorvo2010AB}.

\item These assertions extend to nonminimal pairs. An
exceptional curve outside $C$ meets it once, transversely
at a smooth point, since $-K_V\cdot E=1$; its contraction
reverses an interior blow-up. An exceptional curve in
$C$ meets the rest at two points and its contraction
reverses a corner blow-up. The reduced anticanonical
condition is preserved. Interior blow-up increases
$e(V\setminus C)$ by one, whereas corner blow-up leaves
it unchanged. Thus every contribution is nonnegative.
Equality permits only toric rational pairs, minimal
elliptic ruled or Hopf pairs, and corner blow-ups of
hyperbolic or parabolic Inoue pairs. For these Inoue
pairs the cycle count and $b_2$ increase together, and
the fundamental group stays $\Z$. Corner blow-ups do
not touch the parabolic elliptic curve, whose square
remains strictly negative.

\item The local nearby fibre at a multiplicity-$k$ stratum
of a reduced normal-crossing divisor retracts to
$(S^1)^{k-1}$. Euler integration consequently gives
\[
 0=e(\mathcal X_t)=\sum_i e(V_i\setminus C_i).          \tag{B.4.28.4}
\]
Every summand is nonnegative and therefore zero.
Rational components with a smooth elliptic conductor,
and all interior blow-ups, are excluded.

\item \emph{Elliptic normal-degree balance.}
An elliptic double curve contains no triple point.
The local model $p=xy$ gives opposite degrees to its
two normal lines. Sum these degrees over all elliptic
conductor branches. The total is zero. An elliptic
ruled or Hopf component contributes zero; a parabolic
Inoue component contributes the strictly negative
degree of its sole elliptic branch. No rational
component with elliptic conductor remains. Hence there
are no parabolic Inoue components.

\item By connectedness, all conductors are now elliptic or
all are rational: no remaining component has both
types. In the elliptic case there are no triple points
and every vertex of the dual graph has valence at most
two. The graph is a segment or circle, allowing a loop
or multiple edges. In this case
\[
 \operatorname{rank}(T-1)\le b_1(\text{dual graph})
 \le1.                                               \tag{B.4.28.5}
\]
To see this without an invariant-cycle theorem, remove
small tubes about the disjoint double curves. The
remaining compact cores extend as smooth submersions
over the whole disc and can be trivialized there.
Monodromy is then a product of the annular twists
$xy=t$. A passage of a loop across one neck adds its
oriented normal-circle class, with coefficient the
signed crossing number of that edge. Core loops have
zero change. Crossing numbers of a closed loop form a
cycle of the dual graph. Thus $T-1$ factors through its
rational cycle space. The normal-circle rotation exists
on the circle bundle also for a nontrivial normal line;
a product normal bundle is not assumed. This proves
(\ref{push:src:B.4.28}.5) and excludes the elliptic case.

\item \emph{Hyperbolic components contradict a nonzero nearby period.}
The remaining possibilities are toric rational pairs
and corner blow-ups of hyperbolic Inoue pairs. Suppose
one component is of the latter kind. A boundary of a
small neighborhood of either rational conductor cycle,
lying in $V_i\setminus C_i$, is a torus mapping torus
$M_A$, whose monodromy on its fibre lattice is
\[
 A\in\operatorname{SL}_2(\Z),\qquad
 |\operatorname{tr}A|>2.                             \tag{B.4.28.6}
\]
Here is the exact marking needed. In
\citet[Section 15, Theorem 15.1 and the following
paragraph, pp.~130--131]{Oda1978D}, the universal
toric construction has lattice $N$ and deck action
given by multiplication by a nontrivial totally
positive unit of a real quadratic field. This is
the hyperbolic matrix $A$. The complement of the
two cycles is $\R$ times the compact angular-torus
bundle over a circle. A path in the logarithmic
absolute-value space from a point to its image by
$A$ realizes this mapping torus. The point can be
chosen in a node cone arbitrarily near a conductor
node. Its two adjacent primitive rays form a lattice
basis, so the angular fibre is exactly the Clifford
torus in the corresponding local coordinates
$(x,y)$. Corner blow-ups are toric star subdivisions;
they alter neither the dense torus nor $A$, only
this unimodular coordinate basis. Thus the same
marked mapping torus exists for every corner
blow-up allowed above.

\item Choose a compact such $M_A$ in the smooth open stratum.
A neighborhood of this compact set extends as a
smooth submersion over the full parameter disc.
Transport gives an embedding $M_A\to\mathcal X_t$ for
all small $t$, with the same marked fibre $L_t\simeq
(S^1)^2$. In the abelian group
$H_1(\mathcal X_t,\Q)$, the mapping-torus relations give
\[
 (A-1)H_1(L_t,\Q)\longmapsto0.
\]
Since $A-1$ is invertible over $\Q$, both angular
generators map to zero. The cohomology ring of a
complex two-torus is the exterior algebra on $H^1$.
It follows that the map $H^2(\mathcal X_t,\Q)\to
H^2(L_t,\Q)$ is zero and hence
\[
 [L_t]=0\quad\hbox{in }H_2(\mathcal X_t,\Q).           \tag{B.4.28.7}
\]

\item On the other hand, $h^0(D,\omega_D)=1$, since $D$ is
compact and connected, and the corresponding dimension
on every nearby two-torus is also one. Grauert's
cohomology and base-change theorem applied to the
invertible relative dualizing sheaf supplies a relative
section extending the chosen nonzero central generator;
see \citet[Chapter III, Theorem 4.12(ii)]{BS76}.
Evaluation is
nonzero at every point of $D$. Compactness and shrinking
make this section nonvanishing near the entire fibre.
Write it as $\Omega$. At the conductor node just used,
the total family has coordinates $xyz=t$ and
\[
 \Omega_t=U(x,y,z)\frac{dx}{x}\wedge\frac{dy}{y},
 \qquad U(0,0,0)\ne0.                                \tag{B.4.28.8}
\]
The angular torus can be transported specifically as
$|x|=\epsilon_1$, $|y|=\epsilon_2$, $z=t/(xy)$.
This agrees, up to isotopy, with its compact-core
transport. If $U=\sum a_{ijk}x^iy^jz^k$ in a small
polydisc, direct integration of the convergent series
gives
\[
 \int_{L_t}\Omega_t
   =(2\pi i)^2\sum_{k\ge0}a_{kkk}t^k
   \longrightarrow(2\pi i)^2U(0,0,0)\ne0.            \tag{B.4.28.9}
\]
For small $t$ this contradicts (\ref{push:src:B.4.28}.7), because
the holomorphic two-form $\Omega_t$ is closed. This
excludes every hyperbolic Inoue component. The
argument uses the nearby torus, not a free-product
description of the central fundamental group.

\item \emph{The dual complex.}
All normalization pairs are now rational toric pairs.
Every vertex link is one circle and every edge has
two incident faces, because each rational double
curve has two triple points. Thus the finite dual
complex $K$ is a connected closed topological surface.
It is oriented: the iterated residues of a
trivialization of $\omega_D$ are nonzero at every
triple point, and the residue theorem on each rational
double curve makes its two coefficients opposite,
with the gluing sign for the two branches. These
coefficients form a nowhere-zero two-cycle on $K$.
Constancy of coherent Euler characteristic in the
proper flat analytic family
\citep[Chapter III, Theorem 4.12(iii)]{BS76}, and the alternating
normalization resolution, give
\[
 0=\chi(\OO_D)=\chi(K).                              \tag{B.4.28.10}
\]
The surface therefore has genus one.

\item For self-crossings, use the finite local-branch
incidence cell structure; the link of a normalized
toric component is its circular boundary. Over the
normalization $\PP^1$ of a double curve, the two-branch
cover extends unramified over its triple-point ends
by local normal-crossing coordinates and hence splits.
There is no hidden edge-reflection monodromy. Subdivision
displays the same closed surface and the residue signs.
The alternating normalization resolution uses these
branch signs and gives the same Euler count. The
hyperbolic-link argument took place in the open
normalization stratum, which embeds in $D$, and is
unchanged by self-identification of boundary branches.
\end{enumerate}
\end{proof}


\paragraph{The separate canonical-model reduction}
This conditional proposition is not the noncompletion
argument for \ref{push:src:B.4.19}. Suppose a relative holomorphic two-form
on a hypothetical transverse punctured slice has first
been proved to extend meromorphically. Rescale it so
that its vertical divisor is $\sum_i a_iV_i$ with
$\min_i a_i=0$. On each component its restriction gives
\[
 K_{V_i}+C_i\sim\sum_{j\ne i}(a_j-a_i)C_{ij}.          \tag{B.4.28.11}
\]
Here $\OO_{\mathcal X}(V_i)|_{V_i}\sim
-\sum_{j\ne i}C_{ij}$ since the full fibre is principal.
The minimal coefficient does not make this right side
zero: neighbouring higher-coefficient components remain.
Contracting them, while retaining semistability and
properness in an arbitrary non-K\"ahler threefold, is
not proved here. Nor is arbitrary meromorphic extension
of a punctured comparison assumed. The K\"ahler,
$K$-trivial torus classification of
\citet[Theorem 3.1]{KamenovaToriRoot} does not supply
that missing reduction. The unrestricted nonexistence
of the actual sphere filling in \ref{push:src:B.4.33} does not require
such a canonical model; it uses extending-section traces
and the global affine loop invariant instead.


\Uses{the hypotheses, notation and definitions stated here; \cref{push:src:B.4.19}; \citet[Theorem 3.1]{KamenovaToriRoot}; \cref{push:src:B.4.33}}



\subsection{What invariant forms do not control}

\begin{proposition}[full-rank periods do not bound an affine comparison]\label{push:src:B.4.29}
\label{push:B429}
\StatementRecord{push:src:B.4.29}{proved}{Historical events: Claude Fable 5.2 (Anthropic), effort high, 16 November 2026: carried as background without fault found. GPT-6 Astra (OpenAI), ChatGPT Pro mode, effort Pro, S6\_09\_PAPERS, 17 November 2026: targeted proof check. Also checked in this writing. Pre-continuation writing events: GPT-6 Astra (OpenAI), ChatGPT Pro mode, effort Pro, S6\_10\_PAPERS, 18 November 2026: targeted proof check. Continuation metadata events: GPT-6 Astra (OpenAI), ChatGPT Pro mode, effort Pro, S6\_10\_PAPERS continuation, 19 November 2026: used and rechecked in this run; status-vocabulary reconciliation; metadata reconciliation only; no new mathematical check. \emph{Pre-continuation recorded proof scope.} Re-derived the explicit Tate-product small resolution and rank-two monodromy, the two translated-section sequences with distinct central limits, and the proper-graph non-extension argument. Checked the absolute three-form in local charts and its exact invariance under the translation.}{the hypotheses, notation and definitions stated here; \citet[Sections 1.1.1--1.1.3, pp.~73--74]{KKN2008Root}; \cref{push:src:B.4.30}; \cref{push:src:B.4.31}; \cref{push:src:B.4.33}}

There is a proper semistable family of two-dimensional
complex tori over a disc, with smooth total threefold
and rank-two unipotent monodromy on $H^1$, and a
fibre-preserving automorphism over the punctured disc
with the following properties. It acts identically
on all relative invariant one- and two-forms and
on the Gauss--Manin variation. Their coherent
extensions supplied by the proper family therefore
give the same lattices under this identification.
An absolute holomorphic canonical form extending
across the centre is also unchanged. Nevertheless,
the automorphism has no meromorphic extension to
the total family.

Thus these invariant-form and connection data do
not alone control an affine origin or a marked
total-space comparison, even in the presence of
a proper smooth semistable cap. This statement
does not refute a separate extension theorem
for the difference of two sections which both
extend holomorphically into the same cap.
\end{proposition}



\begin{proof}
\begin{enumerate}[label=\arabic*.,leftmargin=*]
\item \emph{A proper full-rank example.}
Let $\mathcal E\to\Delta_t$ be the standard Tate
family with ordinary fibre
\[
 E_t=\C^*/t^{\Z}\qquad(t\ne0)
\]
and reduced nodal rational central fibre. Its
proper analytic model and node chart $xy=t$
are the explicit example in
\citet[Sections 1.1.1--1.1.3, pp.~73--74]{KKN2008Root}.
Form $\mathcal E\times_{\Delta}\mathcal E$.
Its only total-space singularity is the product
of the two nodes, with local equations
\[
 xy=t,\qquad uv=t,\qquad xy=uv.
\]
Blow up the ideal $(x,u)$ in this conifold and
take its strict-transform charts. They are
\[
 u=xA,\quad y=Av,\quad t=xAv
\]
with smooth coordinates $(x,A,v)$, and
\[
 x=uB,\quad v=By,\quad t=uBy
\]
with smooth coordinates $(u,B,y)$. This is the
standard proper small resolution, an isomorphism
away from the conifold point. Glue it to the
identity on the complement of that point. The
resulting $p:\mathcal Y\to\Delta$ is proper,
has smooth total threefold, and has reduced
normal-crossing centre: its displayed equations
are products of distinct coordinates. The
resolution does not alter the ordinary fibres
$E_t\times E_t$.

\item For each elliptic factor choose period cycles
$w_i,u_i$ with $u_i\mapsto u_i+w_i$ around
$t=0$. On the direct sum the logarithm of
monodromy has rank two and square zero.
The same rank holds on dual $H^1$. Thus this
example has the full-rank weight-one
unipotent degeneration at issue, rather than
a constant elliptic factor.

\item \emph{The nonextending translation.}
On the punctured disc use multiplicative
coordinates $(Z,W)$ for the two factors and put
\[
 T_h(Z,W,t)=\bigl(e^{2\pi i h(t)}Z,W,t\bigr),
 \qquad h(t)=1/t.
 \tag{B.4.29.1}
\]
This is a holomorphic automorphism of the
punctured family. The zero section
$e(t)=([1],[1])$ extends to the smooth torus
locus of the central normalization. For
positive integers $j$ tending to infinity,
\[
 t_j=1/j,\qquad t'_j=1/(j+1/2)
\]
lie in the disc and tend to zero, whereas
\[
 T_h(e(t_j))=([1],[1]),\qquad
 T_h(e(t'_j))=([-1],[1]).
\]
In the Tate model these approach the distinct
central smooth points with multiplicative
coordinates $(1,1)$ and $(-1,1)$.
The small resolution is an isomorphism near
both points. Hence $T_h\circ e$ has no
continuous extension to the centre.

\item This also excludes a meromorphic total
extension, rather than only a holomorphic
one. If such an extension existed, let
$\Gamma\to\mathcal Y$ be its proper
meromorphic graph and pull it back along
$e:\Delta\to\mathcal Y$. The resulting
analytic space is proper over the disc and
has a single point in each noncentral fibre.
Its unique component dominating the disc
is a proper analytic curve, generically of
degree one. After normalization it is finite
of degree one over the normal disc, hence
is the disc itself. Its map to the target
$\mathcal Y$ would extend $T_h\circ e$
holomorphically. This contradicts the two
distinct limits just found. Vertical
components in the pulled-back graph do
not affect this dominating curve.

\item \emph{Relative forms and their coherent extensions.}
On the ordinary locus the invariant relative
one-forms are
\[
 \omega_1=dZ/Z,\qquad \omega_2=dW/W
 \quad\pmod{dt}.
\]
They and $\omega_1\wedge\omega_2$ are
unchanged by every $T_h$. The holomorphic
forms of a compact torus are invariant, so
this accounts for all relative holomorphic
one- and two-forms on the ordinary fibres.
The relative log forms extend in the Tate
model. Ordinary coherent relative forms
also give finite-order extensions: for
example $t\omega_1,t\omega_2$, and
$t^2\omega_1\wedge\omega_2$ extend as
ordinary relative forms. In the first
resolution chart the first two are explicitly
\[
 t\omega_1=Av\,dx,\qquad
 t\omega_2=Av\,dx+xv\,dA;
\]
the other chart gives the analogous regular
expressions. At an isolated Tate node the
same identity is $t\,dx/x=y\,dx$.

\item More generally, take the torsion-free image
of each coherent direct image of relative
forms in its bundle on the punctured disc.
Under the two punctured identifications
$\mathrm{id}$ and $T_h$, these are identical
coherent lattices, since the induced map on
the entire ordinary relative-form bundle
is the identity. Coherence therefore gives
no bound on $h$ in this example.

\item \emph{Gauss--Manin and an absolute form.}
Every fibre translation is homotopic to the
identity and acts identically on integral
cohomology. Consequently $T_h$ induces the
identity on the flat Gauss--Manin bundle,
its connection, and its Hodge subbundle.
The logarithmic period growth and regular
singularity of that connection are unchanged.
They concern the linear cohomology variation,
not the translated origin.

\item There is even an unchanged absolute form
\[
 \Omega=dt\wedge\frac{dZ}{Z}\wedge\frac{dW}{W}.
 \tag{B.4.29.2}
\]
It descends under both Tate period actions,
since their extra differentials are multiples
of $dt$. In the first small-resolution chart
it is
\[
 d(xAv)\wedge\frac{dx}{x}
          \wedge\frac{d(xA)}{xA}
       =dv\wedge dx\wedge dA,
\]
and in the other chart it is
$dy\wedge dB\wedge du$.
Thus $\Omega$ extends as a nonvanishing
holomorphic canonical form on the smooth
total space. Translation adds
$2\pi i h'(t)dt$ to $dZ/Z$, which vanishes
in (\ref{push:src:B.4.29}.2), so $T_h^*\Omega=\Omega$.

\item \emph{What a total differential would have to add.}
On a local multiplicative chart, an absolute
lift of the first relative form transforms by
\[
 T_h^*(dZ/Z)=dZ/Z+2\pi i h'(t)dt.
 \tag{B.4.29.3}
\]
This horizontal term is erased both on passing
to relative forms and on wedging with the base
differential. It is not controlled by the
linear Gauss--Manin connection. The stronger
choice $h(t)=e^{1/t}$ makes that term
essential, while leaving all the preceding
relative and canonical calculations unchanged.
For this choice the two sequences
$t_j=1/\log j$ and
$t'_j=1/\log(j+1/2)$ give the same two
distinct central limits. A condition imposing
a regular extending affine connection or
specified absolute lifts could detect such
growth, but it is additional data; properness
and the invariant forms above do not supply it.

\item The example compares two markings of the same
proper family. It does not identify either
with the actual six-sphere cap. Nor does it
show that differences of two sections both
extending into a single cap fail to extend:
a common essential translation cancels in
such a difference. The difference-extension bridge
is proved in \ref{push:src:B.4.30}; \ref{push:src:B.4.31}--\ref{push:src:B.4.33} add the global
meridian and loop control needed for the actual family.
They do not contradict this counterexample. Exact chart,
monodromy, and translation identities are
checked in
\path{code/e_push_forms_translation.py}.
\end{enumerate}
\end{proof}

\subsection{Noncompletion of the prescribed partial filling}

\begin{proposition}[the fixed partial filling cannot be completed]\label{push:src:B.4.30}
\label{push:B430}
\StatementRecord{push:src:B.4.30}{proved}{Historical events: Claude Fable 5.2 (Anthropic), effort high, 16 November 2026: carried as background without fault found; targeted proof check. Claude Fable 5.2 (Anthropic), effort max, 16 November 2026: targeted proof check. GPT-6 Astra (OpenAI), ChatGPT Pro mode, effort Pro, S6\_09\_PAPERS, 17 November 2026: targeted proof check. Also checked in this writing. Pre-continuation writing events: GPT-6 Astra (OpenAI), ChatGPT Pro mode, effort Pro, S6\_10\_PAPERS, 18 November 2026: targeted proof check; used and rechecked in this run. Continuation metadata events: GPT-6 Astra (OpenAI), ChatGPT Pro mode, effort Pro, S6\_10\_PAPERS continuation, 19 November 2026: used and rechecked in this run; status-vocabulary reconciliation; metadata reconciliation only; no new mathematical check. \emph{Pre-continuation recorded proof scope.} Read complete receiving proof and re-derived displayed sheaf/topological argument. The external logarithmic comparison applications in B4.30/B4.32/B4.33 were audited separately against IKN and its erratum; the independent direct commuting-meridian proof is night:P2.7. Boundary proof checked using intrinsic anticanonical divisor, coherence and proper-image exceptional loci. Opened IKN Proposition6.1 and complete Theorems6.2–6.3 proofs, pp30–34, the adjacent remarks and the entire mathematical erratum pp113–115. Rechecked the receiving proper separated log-smooth fs setup, torsion-free characteristic cokernels and unipotent-coefficient application. Also read the exact strict-analytic toroidal Lemma2.14 and full proof, and Parusinski Corollary1.2 with complete proof pp3–6. This closes these named source gates; it adds no second whole-proof credit.}{the hypotheses, notation and definitions stated here; \cref{push:src:B.4.19}; \cref{push:src:B.4.4}; \cref{push:src:B.4.31}; \cref{push:src:B.4.33}; \cref{push:src:B.4.27}; \citet[Theorems 6.2 and 6.3,
pp.~32 and 34]{IKN2005E}; \citet{IKN2007ErratumE}; \citet[Definition 7.1(i) and Section 9.2,
pp.~164--165 and 173]{KKN2008PicardE}; \cref{push:src:B.4.29}; \cref{push:src:B.4.24}; \citet[Corollary 1.2 and its proof,
manuscript pp.~3--6]{Parusinski1990E}; \citet[Theorems 1.6 and 1.10]{BM95AB}; \citet[Lemma 2.14, pp.~6--7]{EH2024E}; \cref{push:src:B.4.26}; \cref{push:src:B.4.32}}

The partial filling of \ref{push:src:B.4.19} has no proper
Hausdorff complex-analytic completion over its
missing base point. This remains true after finite
coordinate roots of the two base coordinates and
allows arbitrary additional three-dimensional
components of the origin fibre. No K\"ahler or
projective hypothesis is imposed on a putative
completion.

This is a statement about the prescribed partial
filling, including its two punctured axial models.
This proposition alone does not exclude every filling
over the $r$ disc: a different completion along
$r=0$, $q\ne0$ need not have the axial model of
\ref{push:src:B.4.19}. Resolving the graph in \ref{push:src:B.4.4} does not by
itself supply that identification. \ref{push:src:B.4.31}--\ref{push:src:B.4.33}
remove the need for it by controlling meridians and
the global affine loop class, and prove unrestricted
nonexistence for the actual family.
\end{proposition}



\begin{proof}
\begin{enumerate}[label=\arabic*.,leftmargin=*]
\item We first prove the local extension statement
needed for traces. We then make a hypothetical
completion equidimensional by an explicit
coefficient-ideal construction and glue multiplied
local traces into the logarithmic torsor
forbidden by \ref{push:src:B.4.27}.

\item \emph{The analytic logarithmic de Rham input.}
We use \citet[Theorems 6.2 and 6.3,
pp.~32 and 34]{IKN2005E}. For a proper separated
log-smooth analytic map with log-smooth base and
torsion-free cokernels of characteristic groups,
its ordinary relative log-de Rham cohomology is
locally free and, with its Gauss--Manin
connection, is the nilpotent logarithmic
Riemann--Hilbert extension of the log Betti
system. These statements have no K\"ahler or
projectivity assumption. The correction in
\citet{IKN2007ErratumE} concerns the algebraic
Hodge-degeneration Theorem~7.1, which is not
used here. The same distinction is explicit in
\citet[Definition 7.1(i) and Section 9.2,
pp.~164--165 and 173]{KKN2008PicardE}.

\item \emph{Closed logarithmic forms extend by saturation.}
Consider a proper semistable family of surfaces
\[
 f:Z\longrightarrow\Delta_t\times V,
 \tag{B.4.30.1}
\]
where $V$ is a disc or a point. Its noncentral
fibres are the specified two-dimensional torus
torsors. The semistable charts, with the
transverse parameter retained, are
\[
 t=x_1\cdots x_k,\qquad u=u,\qquad k\le3.
 \tag{B.4.30.2}
\]
Write $\Lambda$ for the integral first-homology
lattice and assume that $T=1+N$, $N^2=0$, and
\[
 \ker(N:\Lambda\longrightarrow\Lambda)
       =\Z w_1\oplus\Z w_2=:W.
 \tag{B.4.30.3}
\]
Here $W$ is the primitive cocharacter lattice
of the specified multiplicative Jacobian
presentation. Suppose its normalized invariant
one-forms $\omega_1,\omega_2$, with periods
$2\pi i\delta_{ij}$ on $w_j$, span a rank-two
subbundle $F_A$ of the canonical extension
$\mathcal H$ of $H^1$.

\item These assumptions hold at every generic
exceptional divisor in our application. Indeed
$q=t^a$ times a unit and $r=t^b$ times a unit,
with $a,b>0$. The two positive valuation entries
of \ref{push:src:B.4.27} give (\ref{push:src:B.4.30}.3), with $w_1=w$ and
$w_2=\delta$, after the fixed integral character
change. In the actual period matrix, subtracting
the logarithmic $N$ columns leaves holomorphic
unit periods. The $w,\delta$ columns remain
the identity. Hence their normalized form span
extends as the asserted subbundle $F_A$.

\item Let $K^\bullet=\Omega^\bullet_{Z/(\Delta\times V)}
(\log)$ and
\[
 \mathcal F=\ker\bigl(f_*\Omega^1(\log)
                   \longrightarrow f_*\Omega^2(\log)\bigr).
\]
This is coherent: the displayed differential
is linear over the base. The triangle
$K^{\ge1}\to K^\bullet\to\OO_Z$ gives the
low-degree exact sequence
\[
 0\longrightarrow\mathcal F
 \longrightarrow\mathcal H
 \longrightarrow R^1f_*\OO_Z.
 \tag{B.4.30.4}
\]
The preceding connecting map is zero since
$f_*\OO_Z=\OO_{\Delta\times V}$ and the
differential is relative. Thus no degeneration
of a Hodge spectral sequence is assumed.

\item The reduced connected compact central fibres
have only constant holomorphic functions.
For $D=f^{-1}(t=0)$ this gives
$f_*\OO_D=\OO_V$. Applying $f_*$ to
\[
 0\longrightarrow\OO_Z\xrightarrow{\ t\ }
 \OO_Z\longrightarrow\OO_D\longrightarrow0
\]
shows that $R^1f_*\OO_Z$ has no $t$-torsion:
the map $\OO_{\Delta\times V}\to\OO_V$ is
surjective. Consequently $\mathcal F$ is
$t$-saturated in $\mathcal H$. The subbundle
$F_A$ is also $t$-saturated, and the two agree
on $t\ne0$, where all holomorphic one-forms
on a torus are invariant and closed. The
coherent quotients of $\mathcal F+F_A$ by
either summand have support in $t=0$ and
are locally killed by a power of $t$. This
is the localized equality being used; it
does not identify arbitrary punctured
holomorphic sections with meromorphic ones.
Therefore
\[
 \mathcal F=\mathcal H\cap F_A[t^{-1}]=F_A.
 \tag{B.4.30.5}
\]
In particular the normalized forms are actual
closed relative logarithmic forms on $Z$,
holomorphic also in $u$.

\item \emph{Integral local residues and convergent primitives.}
A local fibre in (\ref{push:src:B.4.30}.2) retracts to
$(S^1)^{k-1}$. Its angular classes are
monodromy-invariant, and their inclusion into
the global fibre intertwines monodromy. By
(\ref{push:src:B.4.30}.3) they belong to $W$. Thus all periods
of $\omega_j/(2\pi i)$ on these angular loops
are constant integers, for all $(t,u)$.

\item We record the elementary convergent calculation.
Modulo $d\log t$, write a closed form as
\[
 \omega=\sum_{i<k}a_i(x,z,u)d\log x_i
                   +\sum_\ell b_\ell(x,z,u)dz_\ell.
\]
First integrate the ordinary $z$ variables by
the holomorphic homotopy formula. For a
remaining monomial $x^\alpha$, its relative
weight is
\[
 (\alpha_1-\alpha_k,\ldots,\alpha_{k-1}-\alpha_k).
\]
Closedness makes its coefficient vector
proportional to this weight. If the weight is
nonzero, divide by a nonzero integer entry to
integrate the monomial. The denominator has
absolute value at least one, so the series
converges on smaller polydiscs. Zero-weight
monomials are powers of $t$. Their coefficients
are exactly the local angular periods, and
therefore reduce to constant integers in the
normalized forms. We obtain
\[
 \omega_j=d_{\rm rel}h_j+
               \sum_{i<k}n_{ij}d\log x_i,
 \qquad n_{ij}\in\Z,
 \tag{B.4.30.6}
\]
with $h_j$ holomorphic in every variable.
Hence
\[
 \chi_j=e^{h_j}\prod_{i<k}x_i^{n_{ij}}
 \tag{B.4.30.7}
\]
is a meromorphic character whose divisor is
supported on the central coordinate branches.

\item \emph{Two sections extending into the same cap.}
Let $a,b$ be holomorphic sections of (\ref{push:src:B.4.30}.1).
A finite chain of semistable charts through the
connected compact central fibre joins their
limits. Shrinking $V$ makes the chain uniform
in $u$. Every nonempty open intersection with
a reduced NC fibre contains a smooth point of
that fibre. At such a point $f$ is a
submersion, so choose a holomorphic transverse
auxiliary section in each successive overlap.

\item In each chart, evaluate the two characters
(\ref{push:src:B.4.30}.7) at the entering and leaving sections
and take their ratios. A pulled-back $x_i$
is a power of $t$ times a holomorphic unit,
since it vanishes only over $t=0$. Thus these
ratios are logarithmic characters on the base.
Multiply them along the finite chain. On each
ordinary fibre, the result is the exponential
of the integrals of $(\omega_1,\omega_2)$ from
$a$ to $b$, so its class modulo the specified
period lattice is exactly $b-a$. A different
path changes it by an integral period. Every
integral ghost is bounded along this full-rank
positive degeneration. Consequently $b-a$
extends as a section of the associated
logarithmic torus $A$.

\item This construction is holomorphic in the
transverse variable. It is not an inference
from separate one-dimensional restrictions.
It concerns two sections extending into one
cap, and is compatible with \ref{push:src:B.4.29}, where an
essential change of reference origin need not
extend into that cap.

\item \emph{Equidimensionalization without a projective hypothesis.}
Suppose a proper completion exists. Retain its
reduced main component, normalize it, and make
the finite coordinate roots used in \ref{push:src:B.4.24}.
Its ordinary locus has dimension four; it is
equidimensional of relative dimension two
over the punctured bidisc. Since $r$ is
nonzero on the main component, its origin
fibre has local dimension at most three.

\item We use the coefficient-ideal argument in
\citet[Corollary 1.2 and its proof,
manuscript pp.~3--6]{Parusinski1990E}, and give
its excess-one specialization. At a point of
the origin fibre with local fibre dimension
three, choose three generic vertical linear
functions $z=(z_1,z_2,z_3)$ so that the germ
\[
 (f,z): (Z,p)\longrightarrow(B\times\C^3,(0,0))
\]
is finite. Its reduced four-dimensional image
is a hypersurface, with an equation
\[
 F(y,z)=\sum_\alpha a_\alpha(y)z^\alpha.
 \tag{B.4.30.8}
\]
The coefficient ideal $I=(a_\alpha)$ is finitely
generated. Moreover $F\in I\OO_{B\times\C^3}$:
reduce modulo $I$ and use injectivity into
the $z$-adic completion, or Krull intersection.
Thus the division below is convergent and
holomorphic, not merely coefficientwise formal.
After shrinking, $I$ is a unit off the origin. Otherwise the image fibre over
some $y\ne0$ would contain an open subset of
$\C^3$; finiteness would force a
three-dimensional source fibre, contrary to
the prescribed partial family.

\item Blow up $I$. Where its pullback is generated
by $a_{\alpha_0}$, divide (\ref{push:src:B.4.30}.8) by that
generator. The divided equation has a
coefficient equal to one, so is nonzero on
every base fibre. The strict transform lies
in its zero set, and the finite map is
preserved by base change. Its fibres near
this chart therefore have dimension at most
two. A finite number of these charts covers
the compact origin fibre. Blowing up the
product of their coefficient ideals makes
each factor invertible and performs the same
reduction simultaneously. Resolve the
modified surface base. It is dominated by a
finite sequence of point blowups over the
origin. Further base change preserves the
fibre-dimension bound, and a strict transform
cannot increase it. Dominance gives the
opposite bound, so every fibre now has
dimension exactly two.

\item We only need this equidimensionality, not
flatness of the whole transformed space.
At every point of the modified base it gives
a finite local multisection by a transverse
analytic slice. Normalize that surface. A
normal surface is Cohen--Macaulay, so its
finite map to the smooth surface base is
flat. Thus its positive degree and
multiplicity trace are well defined.

\item \emph{Uniform semistable reduction along a generic exceptional divisor.}
Let $S\to B$ be the preceding point-blown-up
base and $E$ its compact exceptional tree.
Apply analytic embedded resolution near
the compact exceptional set, in the scope of
\citet[Theorems 1.6 and 1.10]{BM95AB}, to
the total space and the inverse image of $E$
so the latter is \emph{simple} normal crossings. Its components give global
markings of all divisor rays. For each closed
intersection of these components, use its
proper map to a component of $E$. Delete the
critical values of every dominating such map
and the images of nondominating intersections.
These sets are finite. Near any remaining
point, with $u$ a coordinate along $E$, the
map has the form
\[
 t=\text{unit}\prod_{i=1}^k x_i^{m_i},
 \qquad u=u,\qquad k\le3.
 \tag{B.4.30.9}
\]
The last bound follows because the transverse
fibres are surfaces; a fourfold intersection
cannot dominate the curve $E$.

\item Choose $M$ divisible by the finitely many
multiplicities. Also let $L$ be a common
multiple of $1,\ldots,d_{\max}$ for the
finite multisections to be used. Make the
simultaneous cover $t=\tau^{ML}$ and
normalize. Replace $M$ by $ML$ in the
following lattice calculation. In the lattice
\[
 \sum_i m_i a_i=M b,
\]
each extreme ray has primitive generator
$a=(M/m_i)e_i$, $b=1$. Thus the height-one
slice is an integral polygon of dimension
at most two. Insert all its lattice points
by star subdivisions, using one order on the
finite marked cone complex. Shared faces
receive the same subdivisions. The final
empty lattice triangles are unimodular, so
the toric modification is smooth and its
$\tau$ fibre is reduced NC. The monomial
ideals glue under coordinate changes by
units. This is precisely the analytic
subdivision construction of
\citet[Lemma 2.14, pp.~6--7]{EH2024E}; their
projective main theorem is not used.
The construction retains the transverse
parameter $u$ throughout.

\item The finite multisections split on this
simultaneous cyclic cover at generic points
of $E$.
Indeed, away from the remaining branch-curve
intersections, their finite normal covers
are the cyclic covers of a punctured disc,
times the $u$ disc. A proper modification
gives a meromorphic lift of each resulting
section; on its normal surface base the
indeterminacy is isolated. Delete those
finitely many additional values of $u$.
The sections then lift holomorphically into
the semistable model. The preceding
saturation and primitive argument applies.

\item If a difference extends only after
$t=\tau^e$, take its multiplicative norm.
For each character the formula is
\[
 \prod_{j=0}^{e-1}(\zeta_e^j\tau)^k
                      U(\zeta_e^j\tau,u)
   =(-1)^{k(e-1)}t^k\,\text{unit}.
 \tag{B.4.30.10}
\]
A change of lift adds an integral period,
whose norm remains a period. The resulting
section downstairs is $e$ times the original
difference. This is an explicit norm on the
usual analytic site, not an assumption of
Kummer descent for a quotient sheaf. A
common multiplier suffices: there are
finitely many divisor multiplicities and
multisections, and ramification indices of
a finite multisection are bounded by its
degree. The product of the multiplicity
common multiple and
$\operatorname{lcm}(1,\ldots,d_{\max})$
bounds the simultaneous root and norm.
The toric subdivision is performed after
this combined root, so no further base
change silently destroys semistability.

\item \emph{Hartogs for logarithmic sections, including the period ambiguity.}
Let $V_p$ be a small bidisc about an isolated
remaining point on $S$, and put
$V_p^*=V_p\setminus\{p\}$. The period sheaf
$Y=\Z^2$ is constant and
\[
 H^1(V_p^*,Y)=H^1(S^3,\Z^2)=0.
\]
The usual-site exact sequence
\[
 0\longrightarrow Y\longrightarrow
 T_{\log}^{(Y)}\longrightarrow A\longrightarrow0
 \tag{B.4.30.11}
\]
therefore lifts every $A$-section on $V_p^*$
to bounded logarithmic characters. By
\ref{push:src:B.4.27} the middle sheaf is exactly
\[
 M_{D_Q+E}^{\rm gp}\times
 M_{D_Q+D_R+E}^{\rm gp}.
\]
For either character divide by the local
boundary monomial carrying its integer
orders. The resulting function and its
inverse are holomorphic on the punctured
bidisc. Hartogs extends both, and their
product remains one. Restore the monomial.
The forbidden strict-$r$ order of the first
character stays zero, so the extension is
still bounded.

\item The extension is unique. Locally lift the
difference of two extensions to the middle
sheaf. On the connected punctured bidisc
it is an integral period, hence one constant
element of $Y$. Subtract that period and
apply the identity theorem to its two
characters. This proves uniqueness and
preserves all cocycle identities. It does
not assert Hartogs extension for $A$-torsors.

\item \emph{Multiplying local traces and gluing the forbidden torsor.}
Choose a finite cover $U_i$ of a
neighbourhood of the compact exceptional
tree and a finite local multisection of
degree $d_i>0$ on each member. Let $d$ be
a common positive multiple of these degrees.
On $U_i\setminus E$, its trace multiplied
by $d/d_i$ is a section $t_i$ of the
$d$-fold contracted product $\mathcal P^d$
of the actual punctured $A$-torsor.
On the two original punctured axes these
are genuine logarithmic trace sections:
a holomorphic map into either prescribed
fan model pulls its characters back to
powers of the local root coordinate times
units, and their norms are exactly those
in \ref{push:src:B.4.26}. A proper modification projecting
to that model does not change this argument.

\item On overlaps, $t_i-t_j$ is an $A$-section
away from $E$. The uniform argument above
gives a positive multiple extending across
the generic part of every exceptional
component. Choose one common multiplier
$m$ for the finite set of curves and
multisections and replace $t_i$ by $mt_i$.
These are sections of $\mathcal P^{dm}$.
Their differences now extend away from
finitely many points of $E$; the Hartogs
argument extends them across those points
as well, including crossings. Uniqueness
preserves the cocycle equations.

\item Glue trivial $A$-torsors on the $U_i$ by
these extended differences and glue them
to $\mathcal P^{dm}$ off $E$ using the
trace identifications. This produces an
$A$-torsor on $S$ restricting to
$\mathcal P^{dm}$ on
$\beta^{-1}(B\setminus\{0\})$.
\ref{push:src:B.4.27} says that $H^1(S,A)=0$.
But \ref{push:src:B.4.24} says that every positive
multiple of this actual punctured torsor
is nontrivial: its $r$-loop connecting
class is a nonzero multiple of $\gamma$.
This is the required contradiction.

\item Finally, all period and cohomology arguments
in \ref{push:src:B.4.24} and \ref{push:src:B.4.27} allow further finite
coordinate roots. The proof is therefore
unchanged by them. Its use of the two
\emph{prescribed} punctured axial models
is essential. Along the strict $r$ axis the linear monodromy
condition alone leaves invariant elliptic directions.
The proof here does not control their geometric
residues or identify an unknown cap with \ref{push:src:B.4.19}.
That control is supplied by the global argument in
\ref{push:src:B.4.31} and used, together with \ref{push:src:B.4.32}, in \ref{push:src:B.4.33}.
\end{enumerate}
\end{proof}

\subsection{Global elliptic monodromy constrains every central meridian}

\begin{proposition}[the global meridians have only the prescribed toric direction]\label{push:src:B.4.31}
\StatementRecord{push:src:B.4.31}{proved}{Historical events: Claude Fable 5.2 (Anthropic), effort high, 16 November 2026: carried as background without fault found; targeted proof check. Claude Fable 5.2 (Anthropic), effort max, 16 November 2026: carried as background without fault found. GPT-6 Astra (OpenAI), ChatGPT Pro mode, effort Pro, S6\_09\_PAPERS, 17 November 2026: targeted proof check. Also checked in this writing. Pre-continuation writing events: GPT-6 Astra (OpenAI), ChatGPT Pro mode, effort Pro, S6\_10\_PAPERS, 18 November 2026: targeted proof check. Continuation metadata events: GPT-6 Astra (OpenAI), ChatGPT Pro mode, effort Pro, S6\_10\_PAPERS continuation, 19 November 2026: used and rechecked in this run; status-vocabulary reconciliation; metadata reconciliation only; no new mathematical check. \emph{Pre-continuation recorded proof scope.} Completed the previously unread angular-circle and finite-orbit paragraphs. Re-derived primitive monomial circles with their gcd multiplicities, normal-bundle transport and conjugation invariance in H1, finite component permutation under parameter monodromy, kerE reduction, and actual Gauss--Manin finite-orbit transport across a finite product-chart cover. Exact analytic resolution locator remains a citation gate.}{the hypotheses, notation and definitions stated here; \cref{push:src:B.4.4}; \cref{push:src:B.4.19}; \citet[Theorems 1.6 and 1.10 and the paragraph following
Theorem 1.6]{BM95AB}; \cref{push:src:B.4.9}}

Suppose $g:\mathcal Z\to\Delta_r$ is any proper complex-analytic
filling of a finite parameter cover of the actual six-sphere
family. Retain its reduced main component, normalize and resolve
it, extend its intrinsically marked map to $\PP^1$ by \ref{push:src:B.4.4},
resolve the graph, and resolve the support of its central divisor. Denote
the resulting smooth proper modification by $\widetilde{\mathcal Z}$,
and its holomorphic marked map by $h$. All modifications can be
chosen to be isomorphisms over $r\ne0$. The central support has
normal crossings; its multiplicities need not be one.

There is a finite subset $\Sigma\subset\PP^1$, containing
$0,1,\infty$, with the following property. For every
$q\in B^\circ:=\PP^1\setminus\Sigma$, every local angular
vanishing circle of the $r$-central normal-crossings divisor,
viewed in the homology of the nearby ordinary two-torus, lies in
\[
             \Z\widehat\delta\subset H_1(X_{q,r};\Z).
 \tag{B.4.31.1}
\]
The same conclusion holds for the image of the entire local
angular torus at a higher crossing over $B^\circ$.
Here $\widehat\delta$ is the primitive cycle in the audited
period marking. The conclusion is unaffected by any further
finite cover of the degeneration parameter.

This is a constraint on every hypothetical global filling. It
does not assert that the unknown rank-one axial completion is
the prescribed model of \ref{push:src:B.4.19}, or that an affine origin in one
model extends into the other.
\end{proposition}


\begin{proof}
\begin{enumerate}[label=\arabic*.,leftmargin=*]
\item \emph{The global model and the finite exceptional set.}
The punctured total family is connected and smooth, hence
irreducible. Its closure is the unique reduced main component
dominating the parameter disc. Retain this component, normalize
it, and resolve it properly, without changing its smooth
punctured part. The resulting total space is a smooth connected
fourfold, proper over the disc. The parameter is a nonzero
divisor, so its map to the disc is flat. The proof of \ref{push:src:B.4.4}
requires no reducedness or normal-crossings condition at the
centre: its torsion-free coherent bi-anticanonical pencil and
labelled critical-value argument apply to this smooth model.
Its marked graph is therefore a proper modification.
Further analytic resolution and principalization of the central
ideal give the stated smooth
model and normal-crossings support; the complex-analytic scope,
including locally finite resolution, is explicit in
\citet[Theorems 1.6 and 1.10 and the paragraph following
Theorem 1.6]{BM95AB}. Properness over a smaller parameter disc
allows us to work over a compact central fibre throughout.
The modifications have their centres over $r=0$, so neither
the actual punctured family nor its marked map changes.
Flatness over $r$ persists because $r$ is a nonzero divisor in
the smooth local rings. We do not infer that its central
multiplicities remain one.

\item For possible self-crossings use the finite closed spaces of
ordered local branch intersections, with each local component
normalized. Locally these are the coordinate intersections of
the central branches. In particular, the double-branch spaces
are compact smooth surfaces and the triple-branch spaces are
compact smooth curves, even when two branches belong to one
global component. There are finitely many such components.

\item Add to $\Sigma$ the images of all branch-intersection components
on which $h$ is constant and the critical values of $h$ on every
remaining branch-intersection component, including the
normalizations of the single branches. These sets are finite.
Indeed the maps are proper; on a component mapping nontrivially
to the curve $\PP^1$ their critical image is a proper analytic
subset of that curve. Zero-dimensional deeper strata contribute
only finitely many further values.

\item Thus, over $B^\circ$, the closed double-branch spaces form
proper smooth families of compact curves. The connected
components of their fibres form finite covering spaces of
$B^\circ$: this follows either from proper submersion
trivializations, or from their finite Stein factorizations.
Removing the finitely many triple points from one connected
smooth curve leaves it connected. Branch orderings, central
multiplicities, and these connected fibre components therefore
constitute finite discrete data over $B^\circ$.

\item \emph{The local angular circles, including multiplicities.}
At a point of an open double-branch stratum, the submersion
condition for $h$ gives coordinates in which $q=h$ is a
tangential coordinate and
\[
                 r=x^a y^b,\qquad a,b>0.
 \tag{B.4.31.2}
\]
A holomorphic unit has been absorbed in one coordinate on this
small chart. Put $d=\gcd(a,b)$. Each connected component of the
local nonzero fibre has the primitive angular circle
\[
 \theta\longmapsto
 \bigl(x_0e^{i(b/d)\theta},\,
       y_0e^{-i(a/d)\theta}\bigr),
       \qquad 0\le\theta\le2\pi,
 \tag{B.4.31.3}
\]
where $x_0^a y_0^b=r$. There are at most $d$ choices of the
connected monomial-fibre component. Interchanging the ordered
branches can reverse the circle's orientation. These are only
finite choices.

\item Move the chosen base point along a path in the same open
double-curve fibre. Local product charts and the normal line
bundles transport the annular circle continuously. Changes of
normal frame rotate its parameter and do not add a tangential
one-cycle. Equivalently, this is the fibre circle of the normal
circle bundle along that path. If a based loop is used, changing
the path changes it by conjugation, which does not affect its
class in $H_1$. Consequently, for fixed small nonzero $r$ and
fixed $q$, each of the finitely many branch-component data
determines only finitely many homology classes, including signs.

\item The same local formula proves monodromy invariance in the
$r$ direction. Continuing around $r=0$ permutes the finitely
many connected components of $x^a y^b=r$; after a positive
power it carries (\ref{push:src:B.4.31}.3) to the same oriented circle class.
In the actual marking, after a finite parameter cover the
$r$-monodromy is
\[
 T_r=I-mE,\qquad
 E=\widehat\delta\otimes\gamma,\qquad m\in\Z_{>0},
 \quad E^2=0.
 \tag{B.4.31.4}
\]
For example, $w=r^6$ gives $m=12$ with the positive-loop
convention used in \ref{push:src:B.4.9}. If $T_r^k v=v$ for one of the
angular classes $v$, then $kmEv=0$. The homology lattice is
torsion-free, so
\[
 v\in\ker E
    =\Z\widehat u+\Z\widehat w+\Z\widehat\delta.
 \tag{B.4.31.5}
\]
No assertion of reducedness after graph resolution and no
semistable-reduction theorem is needed for this step.

\item \emph{Why global continuation gives a finite orbit.}
Choose a base point $q_*\in B^\circ$. Let $\mathscr S$ be the
finite set of angular classes furnished at $q_*$ by all the
connected double-curve fibres, their ordered branches, their
monomial-fibre components, and both orientations.
Take generators of the finitely generated group
$\pi_1(B^\circ,q_*)$. A generator can be lifted to a path in
the appropriate family of double curves; it ends in one of
the finitely many connected fibre components. Avoid the
triple points during this lift. Properness over the compact
generator paths and finitely many product charts give one
sufficiently small $r_*$ for which all their annular circles
can be transported in the actual smooth fibres
$X_{q,r_*}$.

\item The transported circles are a continuous family of cycles,
so they realize the actual Gauss--Manin transport of $H_1$.
At the end, moving the base point within the resulting
connected punctured double curve does not change its $H_1$
class, by the previous paragraph. Hence each generator
permutes $\mathscr S$. The same holds for all words in the
generators. This proves a finite orbit in the actual lattice,
not merely a finite orbit of abstract stratum labels.
Its projection to
\[
 \ker E/\Z\widehat\delta
       \simeq \Z\widehat u+\Z\widehat w
 \tag{B.4.31.6}
\]
therefore has finite orbit as well.

\item \emph{The exact elliptic-monodromy calculation.}
In the fixed marking the two finite peripheral
matrices are
\[
 A_1=\begin{pmatrix}
 1&0&0&0\\6&0&1&0\\-6&-1&-1&0\\-2&1&0&1
 \end{pmatrix},\qquad
 A_2=\begin{pmatrix}
 1&0&0&0\\0&0&-1&0\\-6&1&0&0\\3&0&1&1
 \end{pmatrix}.
 \tag{B.4.31.7}
\]
They fix $\widehat\delta$ and $\gamma$, commute with $E$,
and induce on (\ref{push:src:B.4.31}.6)
\[
 M_1=\begin{pmatrix}0&1\\-1&-1\end{pmatrix},\qquad
 M_2=\begin{pmatrix}0&-1\\1&0\end{pmatrix}.
\]
Deleting additional ordinary values from the elliptic base
does not decrease this monodromy image: the map on fundamental
groups to $\PP^1\setminus\{0,1,\infty\}$ is surjective.
The image therefore contains
\[
 U=M_1M_2=\begin{pmatrix}1&0\\-1&1\end{pmatrix},\qquad
 V=M_2UM_2^{-1}=\begin{pmatrix}1&1\\0&1\end{pmatrix}.
 \tag{B.4.31.8}
\]
If a vector $\bar v$ has finite orbit, some positive powers
$U^k$ and $V^\ell$ both fix it. Their fixed lines are,
respectively, $\Q(0,1)$ and $\Q(1,0)$, whose intersection
is zero. Thus $\bar v=0$. Applying this to the projected
classes in $\mathscr S$ proves that each angular class lies
in $\Q\widehat\delta$. Since $\widehat\delta$ is a member
of an integral basis of $H_1(T^4;\Z)$, it is primitive;
an integral class in this rational line belongs to
$\Z\widehat\delta$. This proves (\ref{push:src:B.4.31}.1).

\item At a crossing with local equation
$r=\prod_{j=1}^k x_j^{a_j}$, the angular torus has
rational first homology
\[
 \ker\bigl((a_1,\ldots,a_k):\Q^k\longrightarrow\Q\bigr).
\]
It is spanned by the pairwise angular vectors
$a_j e_i-a_i e_j$. Each is a rational multiple of the
primitive circle of the corresponding double branch,
which can be moved to an open double-curve point by
keeping the other coordinates small and nonzero.
The same calculation therefore places the entire angular
image in $\Q\widehat\delta$, and primitivity again gives
the integral conclusion. This also handles self-crossings,
because the ordered local branches were retained throughout.
\end{enumerate}
\end{proof}


\paragraph{Scope of the affine comparison}
The finite-orbit mechanism uses the global compact marked
base and its actual elliptic monodromy. A rank-one degeneration
over an isolated ordinary base disc has no such monodromy
constraint. Nor does (\ref{push:src:B.4.31}.1) distinguish an affine origin:
even if differences of two sections extending into the same
cap have controlled logarithmic characters, a common
nonextending translation cancels from those differences.
Controlling its affine loop relative to the particular
origin of \ref{push:src:B.4.19} is an additional step, proved in \ref{push:src:B.4.32}.
The invariant-form counterexample \ref{push:src:B.4.29} is consistent with
this limitation. \ref{push:src:B.4.33} uses the loop class, rather than
an extending comparison of individual origins, to prove
unrestricted nonexistence for Part B.


\Uses{the hypotheses, notation and definitions stated here; \cref{push:src:B.4.31}; \cref{push:src:B.4.19}; \cref{push:src:B.4.32}; \cref{push:src:B.4.29}; \cref{push:src:B.4.33}}


The finite matrix calculation, including arbitrary positive
powers of the two unipotents, is reproduced in
\path{code/ab_push_global_meridians.py}; its saved output and
JSON record all checks. The selected fixed-space minor is
$-k\ell\ne0$. An initial sign-only assertion mismatch is
retained with its exact script and output; it did not affect
the nonvanishing argument.

\subsection{The radial-loop class of every global cap}

\begin{proposition}[vanishing of the global affine loop defect]\label{push:src:B.4.32}
\label{push:B432}
\StatementRecord{push:src:B.4.32}{proved}{Historical events: Claude Fable 5.2 (Anthropic), effort high, 16 November 2026: carried as background without fault found; targeted proof check. GPT-6 Astra (OpenAI), ChatGPT Pro mode, effort Pro, S6\_09\_PAPERS, 17 November 2026: targeted proof check. Also checked in this writing. Pre-continuation writing events: GPT-6 Astra (OpenAI), ChatGPT Pro mode, effort Pro, S6\_10\_PAPERS, 18 November 2026: targeted proof check; used and rechecked in this run. Continuation metadata events: GPT-6 Astra (OpenAI), ChatGPT Pro mode, effort Pro, S6\_10\_PAPERS continuation, 19 November 2026: used and rechecked in this run; status-vocabulary reconciliation; metadata reconciliation only; no new mathematical check. \emph{Pre-continuation recorded proof scope.} Read complete receiving proof and re-derived displayed sheaf/topological argument. The external logarithmic comparison applications in B4.30/B4.32/B4.33 were audited separately against IKN and its erratum; the independent direct commuting-meridian proof is night:P2.7. Boundary proof checked using intrinsic anticanonical divisor, coherence and proper-image exceptional loci. Matched the exact IKN Theorem6.3 and full proof p34, with the hypotheses carried by B.4.30 and the separately read erratum. This closes the named logarithmic comparison citation gate; the already recorded receiving proof re-derivation retains its own scope.}{the hypotheses, notation and definitions stated here; \cref{push:src:B.4.31}; \cref{push:src:B.4.19}; \cref{push:src:B.4.4}; \cref{push:src:B.1.2}; \cref{push:src:B.4.30}; \citet[Theorem 6.3]{IKN2005E}}

Suppose a proper complex-analytic filling of a finite parameter
cover of the actual six-sphere family exists. Make the
modifications in \ref{push:src:B.4.31} and delete its finite exceptional set
from the marked base. Use the ordinary local origins of
\ref{push:src:B.4.19} as reference. There is a well-defined normalized
radial-loop defect of a local extending cap section,
\[
 \lambda\in
 V:=\bigl(\Lambda\otimes\Q\bigr)/\Q\widehat\delta .
 \tag{B.4.32.1}
\]
It is independent of the section, of the finite parameter root
needed to obtain that section, and of the modifications used.
It patches to a flat section of the actual local system $V$
over the remaining ordinary base, and this section is zero.

In particular, let $t$ be the trace of an extending finite
multisection of degree $d$, and let $o$ be the reference ordinary
origin. The radial-loop defect of $t-do$ in $V$ is zero.
The same holds for an origin obtained from multiplied such
traces by bounded logarithmic differences along $r=0$.
This conclusion controls a loop class; it does not assert
that $t-do$ is bounded or extends holomorphically at $r=0$.
\end{proposition}



\begin{proof}
\begin{enumerate}[label=\arabic*.,leftmargin=*]
\item After a common further parameter cover we may compare with
$w=r^6$ in \ref{push:src:B.4.19}; further positive powers are allowed throughout.
Retain and resolve the reduced main component of the putative
filling. The extension of the marked map uses \ref{push:src:B.4.4}, whose
smooth fourfold proof requires no condition on the central
divisor. \ref{push:src:B.4.31} supplies a proper smooth marked graph with
normal-crossings central support, possibly with multiplicities,
and the integral meridian conclusion. These operations are
the identity on the punctured family.

\item \emph{The connecting class and its normalization.}
On a simply connected ordinary base chart, write
\[
 T_r=I-mE,\qquad
 E=\widehat\delta\otimes\gamma,\qquad m>0,\qquad E^2=0.
 \tag{B.4.32.2}
\]
For a holomorphic section $a$ of the Jacobian on the punctured
parameter disc, its connecting class for the Lie-bundle
exponential sequence belongs to
\[
 H^1(\Delta_r^*,\Lambda)
       =\Lambda/(T_r-I)\Lambda .
\]
Equivalently, continue a local additive lift around the positive
loop and record its endpoint period. A change of lift alters
the period label by $(T_r-I)\ell$. Since this lies in
$\Z\widehat\delta$, the image
\[
 \nu_r(a)\in V
 \tag{B.4.32.3}
\]
is well defined. No single-valued additive lift or meromorphic
extension of $a$ has been assumed. In a holomorphic family of
sections this class is locally constant in the ordinary base:
the integral period label is continuous and discrete.

\item Under $r=R^e$, restriction of the connecting class sends a
period label $v$ to
\[
 \sum_{j=0}^{e-1}T_r^jv
   =ev-\frac{me(e-1)}2Ev
   \equiv ev\pmod{\Q\widehat\delta}.
 \tag{B.4.32.4}
\]
If $b$ is a cap section on this root cover, define
\[
 \lambda(b;o)=\frac1e\,\nu_R\bigl(b-o(R^e)\bigr).
 \tag{B.4.32.5}
\]
Formula (\ref{push:src:B.4.32}.4) proves invariance under further roots and
fixes the normalization used below for finite traces.

\item \emph{The actual affine changes of reference origin.}
This step uses the specified origins, not arbitrary zeros of
abstract isomorphic tori. Retain
\[
 v=\frac{\log r}{2\pi i},\qquad c=-12v,\qquad
 k_i=\frac{\log K}{2\pi i},\qquad
 K(x)=\frac{x^4}{(x-1)^3}.
\]
The ordinary gauge of \ref{push:src:B.4.19} is
\[
 \widetilde\zeta_i=\zeta+k_i v e,\qquad
 e=(0,1)^t=\Pi_c\widehat\delta .
 \tag{B.4.32.6}
\]
Its zero is the local reference origin. On an ordinary overlap
$k_j-k_i=n_{ij}\in\Z$. The extra translation is $n_{ij}ve$;
its radial-loop label is $n_{ij}\widehat\delta$, hence zero
in $V$.

\item The linear Lie-frame changes are independent of $c$ and $r$.
For the two generating finite peripheral changes they are
\[
 R_1=
 \begin{pmatrix}-1/\tau&0\\(1-\mu)/\tau&1\end{pmatrix},
 \qquad
 R_2=
 \begin{pmatrix}1/\tau&0\\-\mu/\tau&1\end{pmatrix}.
 \tag{B.4.32.7}
\]
They satisfy $R_je=e$ and $R_j\Pi_c=\Pi'_cA_j$, with
source and target arguments retained as in \ref{push:src:B.1.2}. Their
compositions therefore act on (\ref{push:src:B.4.32}.3) by the actual lattice
transport and have no additional parameter-dependent
Lie-frame contribution.

\item The finite-root connector is the potentially nonconstant
affine term. Its full cancellation is \ref{push:src:B.4.19}.4:
\[
 \begin{split}
  \sigma_{j,c}&=w_j\Pi_c v_j,\qquad
   k-k_j=12\epsilon_jw_j,\qquad
   \gamma(v_j)=\epsilon_j,\\
  \sigma_{j,c}+(k-k_j)ve
     &=w_j\Pi_{\rm part}v_j .
 \end{split}
 \tag{B.4.32.8}
\]
The last expression is independent of $r$, since $c=-12v$.
It has zero radial-loop defect. The remaining finite deck
maps have the multiplicative factors $r^4$ and $r^{-3}$
in \ref{push:src:B.1.2}, which add only integral $\widehat\delta$ labels.
Changing a sector's period representative adds
$(T_r-I)\ell$, again in $\Q\widehat\delta$.

\item On a punctured cusp sector the comparison is, by \ref{push:src:B.4.19}.5,
\[
 G(s)+6ve,\qquad
 G(s)=\left(\frac{\mu s}{2},
                 \frac{Cs}{12}-\frac{s^2}{24}\right)^t .
 \tag{B.4.32.9}
\]
Here $G$ is independent of $r$. Its additional radial-loop
label is $6\widehat\delta$. A shift by all twelve components
need not be the identity on the ruled normalization; it may
translate the elliptic base by $6P$. As an actual period
change its radial-loop contribution is nevertheless
$(T_r-I)\widehat\gamma$, which is in the same $\delta$ line.

\item These equations prove that the reference origins determine one
globally consistent loop class in $V$. They do not assert that
the origins themselves glue. A different coherent closing
related by a global scalar $\C^*$ gauge has only an integral
$\widehat\delta$ winding and gives the same class. A further
power of the parameter multiplies the displayed loop labels
and leaves all these conclusions unchanged.

\item \emph{Comparison of sections extending into one cap.}
Choose an ordinary base disc avoiding the finite set in
\ref{push:src:B.4.31}. Its closed branch strata map submersively to the
disc. The local finite root and compatible toric subdivision
argument of \ref{push:src:B.4.30} gives a proper semistable model with this
ordinary base parameter retained. Only a model over this
small disc is required. On an overlap of two choices use a
common further root and common refinement. No global
semistable-reduction theorem over the entire punctured
$\PP^1$ is used.

\item The central surface fibres of this model are connected.
Indeed the proper map from its normal dominating total space
to the normal surface base has a finite Stein factor.
It is generically degree one because the ordinary torus
fibres are connected. A finite birational map to a normal
base is an isomorphism.

\item The saturation argument (\ref{push:src:B.4.30}.4)--(\ref{push:src:B.4.30}.5) still applies.
It does not require rank-two monodromy: it only requires the
known normalized rank-two form subbundle of the canonical
log-de Rham extension. The ordinary period matrix supplies
that subbundle, since the $r$ logarithm is solely in the
$\widehat\delta\otimes\gamma$ term and the $w,\delta$
columns stay the identity. The analytic comparison theorem
used here is precisely
\citet[Theorem 6.3]{IKN2005E}, with the hypotheses and erratum
scope already checked in \ref{push:src:B.4.30}.

\item Rank two in \ref{push:src:B.4.30} supplied the integral local residues.
Here \ref{push:src:B.4.31} supplies what is needed more directly: every
local angular class is an integral multiple of
$\widehat\delta$. Thus the first normalized form has residue
zero and the second has integral residue. The convergent
primitive calculation (\ref{push:src:B.4.30}.6) gives exponentials in
\[
              \OO^*\times M_r^{\rm gp}.
 \tag{B.4.32.10}
\]
This is the exact bounded character sheaf along the strict
$r$ axis. In particular the first character is a holomorphic
unit across $r=0$ and has zero radial winding. The second
is a power of $r$ times a unit and contributes only an
integral $\widehat\delta$ winding.

\item Given two extending sections, join their limits by the finite
chart chain of \ref{push:src:B.4.30}. The central connectedness just proved
allows different components as endpoints. The resulting
character comparison lies in (\ref{push:src:B.4.32}.10), so
\[
               \nu_R(b-b')=0\quad\hbox{in }V.
 \tag{B.4.32.11}
\]
Different roots are compared on a common further root.
Proper lifts through refinements are meromorphic on the
two-dimensional section base and holomorphic away from
isolated points. Comparing on the nonempty open set where
the lifts are holomorphic suffices: both loop classes are
locally constant and hence the equality holds throughout
the ordinary overlap. Thus the normalized class is
independent of all the choices.

\item Smooth points of a reduced central fibre have transverse
holomorphic sections, so this common class exists locally.
If an extending finite multisection of degree $d$ is used
instead, split it after a finite cyclic root at a generic
ordinary value. Its $d$ branches have the same normalized
class. Additivity of the connecting homomorphism and
(\ref{push:src:B.4.32}.4) give
\[
            \nu_r(t-do)=d\,\lambda .
 \tag{B.4.32.12}
\]
This also proves independence from finite trace
presentations without presuming an extending chosen origin.

\item \emph{Global transport and its fixed space.}
On one ordinary base chart choose one extending section.
Its normalized integral loop label is locally constant.
On overlaps, (\ref{push:src:B.4.32}.11) removes the choice of cap section,
and (\ref{push:src:B.4.32}.6)--(\ref{push:src:B.4.32}.9) remove the reference-origin
ambiguity in $V$. The local classes therefore patch under
the actual homology transport. Hence $\lambda$ is a flat
section of $V$ over the ordinary base with finitely many
additional points removed.

\item In the quotient basis
$(\widehat\gamma,\widehat u,\widehat w)$, the matrices are
\[
 \overline A_1=
 \begin{pmatrix}1&0&0\\6&0&1\\-6&-1&-1\end{pmatrix},
 \qquad
 \overline A_2=
 \begin{pmatrix}1&0&0\\0&0&-1\\-6&1&0\end{pmatrix}.
 \tag{B.4.32.13}
\]
Their rational fixed spaces are respectively
\[
 \Q(1,2,-4),\qquad \Q(1,3,-3),
 \tag{B.4.32.14}
\]
whose intersection is zero. Removing additional ordinary
base values does not decrease the monodromy image: the
map on fundamental groups to the three-punctured sphere
is surjective. Therefore $\lambda=0$.

\item Even if one retained only a finite orbit of trace branches,
the conclusion would be the same. The covector $\gamma$
is invariant. Averaging a finite orbit would give a common
invariant vector with its original $\gamma$ coefficient,
so that coefficient is zero. The remaining orbit lies in
the elliptic two-dimensional subspace, where the two
unipotents in \ref{push:src:B.4.31} have distinct fixed lines and no
nonzero finite orbit.

\item Finally, a bounded logarithmic difference on the strict
$r$ axis has the character form (\ref{push:src:B.4.32}.10), hence its
loop label lies in $\Q\widehat\delta$. Adding such a
difference to a multiplied finite trace preserves its
zero defect. This proves the last assertion and in
particular gives zero after applying the invariant
covector $\gamma$.
\end{enumerate}
\end{proof}


\paragraph{Scope of the loop calculation}
No assertion of a bounded, meromorphic or holomorphic
comparison with the reference origin has been made.
Essential translations with zero loop defect remain
possible. The conclusion uses the global marked
$\PP^1$ base and its two peripheral matrices; an isolated
rank-one degeneration over an ordinary base disc has
no such constraint. Arbitrary singular proper fillings
are covered by resolving their reduced main component
before \ref{push:src:B.4.4} and \ref{push:src:B.4.31}; their central multiplicities
are retained until the local finite-root calculations.


\Uses{the hypotheses, notation and definitions stated here; \cref{push:src:B.4.4}; \cref{push:src:B.4.31}}


The exact quotient fixed spaces, the root-transfer identity,
the finite connector cancellation and the cusp loop label
are checked in
\path{code/d_push_global_loop_defect.py}, with its saved
output and JSON record. The proof of the global geometric
comparison is the argument above.

\subsection{No proper analytic filling of the actual punctured family}

\begin{theorem}[nonexistence after every finite parameter cover]\label{push:src:B.4.33}
\label{push:B433}
\StatementRecord{push:src:B.4.33}{proved}{Historical events: Claude Fable 5.2 (Anthropic), effort high, 16 November 2026: carried as background without fault found; targeted proof check; used and rechecked in this run. Claude Fable 5.2 (Anthropic), effort high, 17 November 2026: carried as background without fault found. Claude Fable 5.2 (Anthropic), effort max, 16 November 2026: targeted proof check. GPT-6 Astra (OpenAI), ChatGPT Pro mode, effort Pro, S6\_09\_PAPERS, 17 November 2026: targeted proof check. Also checked in this writing. Pre-continuation writing events: GPT-6 Astra (OpenAI), ChatGPT Pro mode, effort Pro, S6\_10\_PAPERS, 18 November 2026: targeted proof check; used and rechecked in this run. Continuation metadata events: GPT-6 Astra (OpenAI), ChatGPT Pro mode, effort Pro, S6\_10\_PAPERS continuation, 19 November 2026: used and rechecked in this run; status-vocabulary reconciliation; metadata reconciliation only; no new mathematical check. \emph{Pre-continuation recorded proof scope.} Read complete receiving proof and re-derived displayed sheaf/topological argument. The external logarithmic comparison applications in B4.30/B4.32/B4.33 were audited separately against IKN and its erratum; the independent direct commuting-meridian proof is night:P2.7. Boundary proof checked using intrinsic anticanonical divisor, coherence and proper-image exceptional loci. Read and matched IKN Theorems6.2–6.3 and the full mathematical erratum at the exact scope recorded in current\_source\_matches.json. This supplements the existing complete receiving-proof review and closes the named IKN source gate only.}{the hypotheses, notation and definitions stated here; \cref{push:src:B.4.29}; \cref{push:src:B.4.32}; \cref{push:src:B.4.31}; \citet[Theorems 1.6 and 1.10]{BM95AB}; \cref{push:src:B.4.4}; \cref{push:src:B.4.24}; \cref{push:src:B.4.30}; \cref{push:src:B.4.27}; \cref{push:src:B.4.19}; \cref{push:src:B.4.26}}

The actual punctured six-sphere family has no proper
Hausdorff complex-analytic filling after any finite
cover of its degeneration parameter. In particular,
there is no positive integer $N$ for which the cover
$w=r^N$ has a proper filling with smooth total space
and reduced normal-crossings central fibre.

The assertion permits an arbitrary singular central
space, arbitrary additional components, and arbitrary
central multiplicities. It does not impose a
K\"ahler, projective, or logarithmic hypothesis on
the putative filling. Its proof derives the specific
logarithmic torsor needed for the contradiction from
differences of sections extending into that filling.
\end{theorem}



\begin{proof}
\begin{enumerate}[label=\arabic*.,leftmargin=*]
\item We distinguish an individual affine origin from its
radial-loop class throughout. \ref{push:src:B.4.29} prevents an
argument based only on extending invariant forms
from controlling an individual origin. \ref{push:src:B.4.32} is
the additional global input that controls the loop
class without making such an assertion.

\item \emph{Reduction of an arbitrary filling to a marked proper map.}
Suppose a filling exists. Its punctured part is a
connected smooth fourfold. Retain its unique reduced
main component, normalize, and resolve it properly,
with every centre over the central parameter. The
result is a smooth connected fourfold proper over
the parameter disc and unchanged over its puncture.
The parameter is a nonzero divisor, so its map to
the disc is flat. The complex-analytic resolution
used here has the scope recorded in \ref{push:src:B.4.31} and
\citet[Theorems 1.6 and 1.10]{BM95AB}.

\item \ref{push:src:B.4.4} extends the intrinsically marked map to
$\PP^1$ meromorphically. Resolve its proper graph
and the support of the central divisor. We obtain
a proper holomorphic map
\[
 (f,r):Z\longrightarrow\PP^1\times\Delta_r
 \tag{B.4.33.1}
\]
with $Z$ smooth, whose restriction over $r\ne0$
is exactly the marked actual family. The central
support is simple normal crossing after further
embedded resolution, but its multiplicities are
retained. \ref{push:src:B.4.31} and \ref{push:src:B.4.32} apply in this scope.
There is no use of a semistable-reduction theorem
for threefolds.

\item A hypothetical filling remains proper after any
further finite base change, normalization, and
retention of its main component. A finite map of
disc germs is a coordinate power after absorbing
a holomorphic unit into the source coordinate.
We may therefore make a further cover so that
the actual parameter is $w=r^6$, followed, if
necessary, by $r=R^n$ with $n>0$. This notation
includes every original positive cover degree:
pulling $w=\rho^N$ back by $\rho=R^6$ gives
$w=R^{6N}$ and hence $r=R^N$.

\item Put $q=1/f$ near the marked cusp and make the
normalized base change $q=Q^{12m}$, $m>0$.
Write
\[
 B=\Delta_Q\times\Delta_R,\qquad
 B^\times=\Delta_Q^*\times\Delta_R^*,\qquad
 B_q=\Delta_Q^*\times\Delta_R,\quad
 B_r=\Delta_Q\times\Delta_R^*.
 \tag{B.4.33.2}
\]
The proper normalized main component above $B$
has the actual torus torsor $\mathscr P$ on
$B^\times$. Its Jacobian is the restriction of
the specific log torus $\mathcal A_B$ of
\ref{push:src:B.4.24}. The completion on $B_q$ is unknown;
we do not identify it with the rotating model.
On $B_r$ it is the prescribed actual cusp,
up to the proper modifications just made.

\item \emph{Only the origin can have excess fibre dimension.}
There are finitely many irreducible components
of the compact central divisor in (\ref{push:src:B.4.33}.1).
On each component the marked map is either
constant or has two-dimensional level sets.
Thus a three-dimensional fibre of $(f,r)$
can occur only at one of the finitely many
constant images. Shrink the cusp disc to
exclude all these images except its centre.
Also exclude the finitely many noncusp values
removed in \ref{push:src:B.4.31}--\ref{push:src:B.4.32}. The map is now
equidimensional of relative dimension two
on $B\setminus\{0\}$. Finite base change
and normalization preserve this assertion.

\item Apply the coefficient-ideal construction of
\ref{push:src:B.4.30} to its possible three-dimensional
origin fibre. It supplies a finite sequence
of point blowups
\[
 \beta:S\longrightarrow B
 \tag{B.4.33.3}
\]
with centres over the origin, and a proper
normalized main transform whose every fibre
has dimension two. We use precisely the
convergent finite-projection argument proved
there, so this is not an invocation of a
projective flattening theorem. Put
$K=\beta^{-1}(0)$. When no blowup is needed,
$K$ is a point; otherwise it is the compact
exceptional tree. Write $E$ for the exceptional
divisor, possibly empty, and $D_Q,D_R$ for
the strict transforms of the axes.

\item At every point of $K$, a transverse analytic
slice gives a finite local multisection.
Normalize it; a normal surface is
Cohen--Macaulay, so its finite map to the
smooth surface base is flat. Choose a finite
open cover $U_i$ of a neighbourhood of $K$
and such multisections of degrees $d_i>0$.
Properness of $\beta$ implies that this
neighbourhood contains the inverse image
of a smaller bidisc. After shrinking, the
$U_i$ cover all of $S$. Consequently all
degree, resolution, and norm choices below
can be made from one finite collection.

\item Let $d$ be a common multiple of the $d_i$.
On the ordinary locus of $U_i$, the trace
of its multisection multiplied by $d/d_i$
is a section $t_i$ of the contracted
$d$-fold torsor $\mathscr P^d$. A trace
uses the group law of the Jacobian and
does not require an affine origin. Thus
$t_i-t_j$ is an ordinary Jacobian section
on every overlap.

\item \emph{Extending the differences at every boundary divisor.}
Let $\mathcal A_S=\beta^*\mathcal A_B$ in
the associated-log-torus sense of \ref{push:src:B.4.27}.
Its exact bounded-character sheaf is
\[
 \mathcal T_{\log}^{(Y)}
  =\mathcal M_{D_Q+E}^{\rm gp}
        \times\mathcal M_{D_Q+D_R+E}^{\rm gp},
 \qquad Y=\Z\widehat u\oplus\Z\widehat\gamma.
 \tag{B.4.33.4}
\]
We verify this specific target for the
extended differences, rather than assigning
a log-group structure to the unknown cap.

\item At a generic exceptional component, the
orders of both $Q$ and $R$ are positive.
The actual monodromy has full rank two;
its primitive invariant lattice is
$\Z\widehat w\oplus\Z\widehat\delta$.
The saturation, integral-residue, and
convergent-primitive proof of \ref{push:src:B.4.30}
therefore extends differences of two
extending sections to $\mathcal A_S$.
Every integer character order is bounded
at such a divisor. The same full-rank
argument applies at a generic point of
$D_Q$; here the underlying punctured-axis
family is additionally the actual given
cusp model.

\item At a generic point of $D_R$, the monodromy
has rank one. The saturation step of
\ref{push:src:B.4.30} still applies to the actual
rank-two subbundle of normalized invariant
one-forms:
the logarithmic term is solely
$\widehat\delta\otimes\gamma$, and the
$\widehat w,\widehat\delta$ columns
normalize a holomorphic subbundle of the
canonical log-de Rham extension. In this
case we do not replace its three-dimensional
monodromy kernel by a two-dimensional one.
Instead \ref{push:src:B.4.31} supplies the needed stronger
geometric fact: every local angular
meridian of the cap is an integral multiple
of $\widehat\delta$. After any local
semistable cover and toric subdivision,
the same holds for the new angular
meridians: their rational images are
combinations of the previous angular
images, and the primitive cycle
$\widehat\delta$ makes an integral image
an integral multiple of it.

\item Consequently the first normalized form
has zero local residues, and the second
has integral residues. The explicit
relative primitives of \ref{push:src:B.4.30} exponentiate
to
\[
             \mathcal O^*\times
                    \mathcal M_{D_R}^{\rm gp}.
 \tag{B.4.33.5}
\]
This is exactly the restriction of
(\ref{push:src:B.4.33}.4): the first character has zero
order at the strict $R$ axis. The finite
chart-chain comparison works because the
proper central surface fibres are
connected. Indeed the normal proper map
to the normal surface base has trivial
finite Stein factor, since its ordinary
fibres are connected tori. Normalization
and retention of the main component after
a root have the same property.

\item These statements concern local reductions
near a generic point of a boundary curve.
\ref{push:src:B.4.30} gives them holomorphically in the
transverse parameter, with a finite
common root before the toric subdivision.
The root simultaneously accounts for
central multiplicities and splitting of
the chosen finite multisections. Norms
then show that one positive multiple of
each trace difference extends downstairs.
Its first order at $D_R$ remains zero:
the norm of a holomorphic unit is a unit.
Choose one common positive multiplier
for all the finitely many curves and
traces, and replace $d,t_i$ accordingly.
Denote their final common degree by $D>0$.
For precision, resolve the inverse image
of the full boundary over a relatively
compact neighbourhood of $K$; properness
then supplies finitely many divisor
multiplicities and finitely many trace
degrees. Their multiplicity common multiple
times $\operatorname{lcm}(1,\ldots,d_{\max})$
works for every local transverse root;
no global root along a boundary curve
is required, because norms are taken
locally and the unique extended sections
glue. Other branch curves of the finite
multisections meet a boundary component
only at isolated points unless they
equal it; after shrinking around the
compact set these intersections are
finite. The split sections lift through
proper modifications meromorphically from
their smooth surface bases, with only
codimension-two indeterminacy, and these
isolated points are treated by the
Hartogs argument below.

\item The only remaining exceptional points of
the boundary are isolated. The Hartogs
argument in \ref{push:src:B.4.30} applies to sections
of exactly the sheaf (\ref{push:src:B.4.33}.4). On a
punctured small bidisc the constant
period sheaf $Y$ has zero $H^1$, so a
section lifts to bounded characters.
After division by their boundary monomials,
both units and inverses extend by ordinary
Hartogs. Restoring the monomials preserves
the zero first-character order on $D_R$.
The extension is unique modulo $Y$.
Thus every $t_i-t_j$ extends to an
$\mathcal A_S$-section on $U_i\cap U_j$,
and uniqueness preserves the cocycle
identities.

\item \emph{The cap-derived torsor has a global origin.}
Glue the trivial $\mathcal A_S$-torsors
on $U_i$ by these extended differences.
The result is a torsor $\mathcal Q$ on
\emph{all} of $S$. Its ordinary restriction
is canonically $\mathscr P^D$, using
the actual trace sections as the local
identifications. It has not been required
to match the prescribed rotating model
on $D_R$. \ref{push:src:B.4.27} gives
\[
 H^1(S,\mathcal A_S)=0.
 \tag{B.4.33.6}
\]
Choose a global section $h$ of $\mathcal Q$.
On the ordinary locus it is a section
of the actual $D$-fold affine torsor.
Near a generic point of $D_R$ its
difference from a multiplied local cap
trace is an $\mathcal A_S$-section.

\item \emph{The strict-$R$ comparison controls only the loop class.}
For an ordinary Jacobian section $a$
write $\nu_R(a)$ for its radial connecting
class in
\[
 V=(\Lambda\otimes\Q)/\Q\widehat\delta,
 \tag{B.4.33.7}
\]
as defined in \ref{push:src:B.4.32}. Let $\ell_\gamma$
be the integral covector taking value one
on $\widehat\gamma$ and zero on
$\widehat u,\widehat w,\widehat\delta$.
It descends to $V$ and is invariant under
the actual monodromy.

\item Let $o_{\rm ord}$ be a \ref{push:src:B.4.19} ordinary
local reference origin. \ref{push:src:B.4.32}, including
its degree-normalized trace and root
formula, gives
\[
 \nu_R(t_i-D o_{\rm ord})=0.
 \tag{B.4.33.8}
\]
An $\mathcal A_S$-section near $D_R$
has a local lift (\ref{push:src:B.4.33}.5). Its radial
loop can therefore add only a multiple
of $\widehat\delta$. It follows that
$\nu_R(h-D o_{\rm ord})=0$ as well.

\item The rotating origin $o$ of \ref{push:src:B.4.24} on
$B_q$ has the same radial-loop class
as $o_{\rm ord}$. This is the exact
comparison (\ref{push:src:B.4.19}.5): their coordinate
change is $G(s)+6v\widehat\delta$,
where $G$ is independent of $r$.
After $r=R^n$ its radial defect is still
purely $\widehat\delta$. Therefore
\[
 \ell_\gamma\nu_R(h-D o)=0.
 \tag{B.4.33.9}
\]
No boundedness or meromorphic extension
of $h-D o$ is asserted or needed.

\item \emph{The actual cusp origin gives the second zero.}
The section $\sigma$ of (\ref{push:src:B.4.24}.2)
is an actual extending origin on $B_r$.
Its characters are the actual unit
characters of the known cusp fan.
Over $B_r$ the cap still projects properly
to that prescribed fan presentation,
so comparing a finite trace with
$d_i\sigma$ is exactly the character-norm
calculation of \ref{push:src:B.4.26}. This does not
require a global lift of $\sigma$
through every modification; alternatively
such a lift is meromorphic and is
holomorphic away from isolated points,
where the same Hartogs argument applies.
Include these bounded-degree norm
multipliers in the choice of $D$ above.
Hence $D\sigma$ is a
section of $\mathcal Q|_{B_r}$ and
\[
                  h-D\sigma\in
                         \mathcal A_B(B_r).
 \tag{B.4.33.10}
\]
Here $S$ is identified with $B$ away
from its origin, so there is no base
modification in this restriction.

\item The $Q$-ghost map on the period sheaf
along the cusp is
\[
 Y\longrightarrow\Z^2,\qquad
 \widehat u\longmapsto(12m,0),\quad
 \widehat\gamma\longmapsto(0,12m).
 \tag{B.4.33.11}
\]
It is injective. The ghost-cocycle
argument of \ref{push:src:B.4.24} therefore makes the
radial $Y$-connecting class of every
$\mathcal A_B(B_r)$-section zero.
In particular
\[
             \ell_\gamma\nu_R(h-D\sigma)=0.
 \tag{B.4.33.12}
\]
The quotient by the cocharacter cycles
$\widehat w,\widehat\delta$ in that
argument does not change this evaluation.

\item \emph{The exact actual clutching gives a contradiction.}
On $B^\times$, the coordinates based at
$\sigma$ and $o$ differ by the exact
translation $F$ of (\ref{push:src:B.4.24}.3):
\[
 \zeta\longmapsto\zeta-\sigma+F.
\]
Thus $\sigma-o=F$ in the common Jacobian,
and
\[
 (h-D o)-(h-D\sigma)=D F.
 \tag{B.4.33.13}
\]
With $s=(2\pi i)^{-1}\log q$ and
$v=(2\pi i)^{-1}\log r$, the exact
period difference is
\[
 F(s,v+n)-F(s,v)
       =n\gamma-6n(n-1)\delta.
 \tag{B.4.33.14}
\]
The corresponding integral period label
has $\ell_\gamma$ value $n$.
Applying $\ell_\gamma\nu_R$ to
(\ref{push:src:B.4.33}.13), equations (\ref{push:src:B.4.33}.9) and
(\ref{push:src:B.4.33}.12) give zero on the left,
whereas (\ref{push:src:B.4.33}.14) gives $Dn>0$ on
the right. This is impossible.

\item All auxiliary coordinate roots, toric
norms, and trace multipliers have positive
degree and merely multiply this nonzero
integer. The contradiction excludes the
initial arbitrary proper analytic filling.
In particular the smooth reduced-NC
filling requested in Part B does not
exist for any positive cover exponent.
\end{enumerate}
\end{proof}

The final step uses two separate facts: a
cap-derived origin has zero global radial
defect by \ref{push:src:B.4.32}, and a section extending
along the actual cusp has zero radial
ghost class by \ref{push:src:B.4.24}. Neither fact
identifies the unknown central model with
\ref{push:src:B.4.19}. A common essential translation
with zero loop class, as allowed by
\ref{push:src:B.4.29}, is harmless in this comparison.
The untwisted Engel family provides a
useful scope check: its identity closing
has zero affine origins on both axial
fan models, hence $F=0$. The nonzero
clutching class used in (\ref{push:src:B.4.33}.14)
is specific to the actual twisted
family; the proof does not exclude
the untwisted family's known filling.


\begin{corollary}[counterexample; the premise of a global central fibre]\label{push:src:B.5.1}\label{push:B51}
\StatementRecord{push:src:B.5.1}{proved}{Historical events: Claude Fable 5.2 (Anthropic), effort high, 16 November 2026: carried as background without fault found. GPT-6 Astra (OpenAI), ChatGPT Pro mode, effort Pro, S6\_09\_PAPERS, 17 November 2026: targeted proof check. Also checked in this writing. Pre-continuation writing events: GPT-6 Astra (OpenAI), ChatGPT Pro mode, effort Pro, S6\_10\_PAPERS, 18 November 2026: targeted proof check. Continuation metadata events: GPT-6 Astra (OpenAI), ChatGPT Pro mode, effort Pro, S6\_10\_PAPERS continuation, 19 November 2026: used and rechecked in this run; status-vocabulary reconciliation; metadata reconciliation only; no new mathematical check. \emph{Pre-continuation recorded proof scope.} Re-derived the no-filling corollary from the already checked arbitrary-singularity exclusion, including preservation of properness by further finite base change and the actual punctured-circle monodromy scope.}{the hypotheses, notation and definitions stated here; \cref{push:src:B.4.33}}

There is no central fibre or descent resolution
of the requested proper analytic filling of the
actual $S^6$ family. Its proposed integral homology,
fundamental group, central Euler number and
nearby-cycle comparison are therefore undefined.
\end{corollary}


\begin{proof}
\begin{enumerate}[label=\arabic*.,leftmargin=*]
\item Theorem \ref{push:src:B.4.33} excludes the underlying proper
analytic space after every finite parameter cover.
An invariant of that nonexistent central fibre
cannot be computed by assigning it the topology
of a local or auxiliary model. The local topologies
and necessary degeneration conditions proved above
retain exactly their stated hypotheses. Restriction of the actual punctured family to a parameter circle
gives its smooth mapping torus; that monodromy is computed in Part C.
\end{enumerate}
\end{proof}



\section{The boundary of the parameter disc}\label{push:section:D}

We use the standard twists throughout this section. Put
\(H_*=(3/\pi)\log6\), fix a boundary parameter
\(c_*=x+iH_*\), and set \(a=x+3/2\). The distinction between a
single degenerate fibre and the total quotient over its root disc is
essential below. All topological group quotients have their quotient
topology; a Hausdorff reflection means a continuous quotient universal
for continuous maps to Hausdorff spaces.

\subsection{The intrinsic obstruction to extension}\label{push:section:D.1}

\begin{proposition}\label{push:src:D.1.1}
\StatementRecord{push:src:D.1.1}{proved}{Historical events: Claude Fable 5.2 (Anthropic), effort high, 16 November 2026: carried as background without fault found; targeted proof check; used and rechecked in this run. GPT-6 Astra (OpenAI), ChatGPT Pro mode, effort Pro, S6\_09\_PAPERS, 17 November 2026: targeted proof check. Also checked in this writing. Pre-continuation writing events: GPT-6 Astra (OpenAI), ChatGPT Pro mode, effort Pro, S6\_10\_PAPERS, 18 November 2026: targeted proof check. Continuation metadata events: GPT-6 Astra (OpenAI), ChatGPT Pro mode, effort Pro, S6\_10\_PAPERS continuation, 19 November 2026: used and rechecked in this run; status-vocabulary reconciliation; metadata reconciliation only; no new mathematical check. \emph{Pre-continuation recorded proof scope.} Read complete receiving proof and re-derived displayed sheaf/topological argument. The external logarithmic comparison applications in B4.30/B4.32/B4.33 were audited separately against IKN and its erratum; the independent direct commuting-meridian proof is night:P2.7. Boundary proof checked using intrinsic anticanonical divisor, coherence and proper-image exceptional loci.}{the hypotheses, notation and definitions stated here; \cref{fam:src:A.4.1}; \cref{fam:src:C.1.1}, \cref{fam:src:C.2.3}; \cref{fam:src:C.3.4}}

For \(0<|w|<216\), the reduced multiplicity-four surface has
\[
 j(\operatorname{Alb}S_2(w))=J(iw/216),\qquad
 J(q)=q^{-1}+744+196884q+\cdots.
\]
This holomorphic function has no meromorphic continuation through any
point of \(|w|=216\).
\end{proposition}


\begin{proof}
\begin{enumerate}[label=\arabic*.,leftmargin=*]
\item The fixed-coordinate periods on the product cover, including the
isogeny kernel and the affine deck translation, generate
\(\frac12\mathbb Z+\frac{\kappa_2}{4}\mathbb Z\).
Multiplication by two therefore identifies the Albanese lattice with
\(\mathbb Z+(\kappa_2/2)\mathbb Z\). This verifies the precise step of
\cref{fam:src:A.4.1} used here; taking only the invariant eigenspace would omit the
half-period. The values in \cref{fam:src:C.1.1}, \cref{fam:src:C.2.3} and \(\mu(z_2)=(1-i)/2\) give
\[
 \kappa_2=c+\frac32-iH_*,\qquad
 e^{-\pi i\kappa_2}=e^{-3\pi i/2}e^{-\pi H_*}e^{-\pi ic}
 =\frac{iw}{216}.
\]
The positive-imaginary representative is \(-\kappa_2/2\), proving the
identity with the stated Fourier normalization.

\item For completeness we reproduce the obstruction in \cref{fam:src:C.3.4}, rather than regard a numerical Fourier plot as evidence. At a
reduced rational \(r=b/d\) choose a matrix in
\(\operatorname{SL}_2(\mathbb Z)\) with bottom row \((d,-b)\).
For \(z=r+iy\), its transformed imaginary part is \(1/(d^2y)\).
Modular invariance and the convergent Fourier series imply
\(\log|j(r+iy)|=2\pi/(d^2y)+o(1)\).
A meromorphic function near \(r\) grows at most polynomially in
\(1/y\), so no continuation exists there. Every real interval contains
such a rational, excluding continuation at irrational points as well.
A local logarithm of \(q=iw/216\) transfers this result to the circle.
This is a derivation of the earlier natural-boundary result, with the
normalizing constant checked exactly.
\end{enumerate}
\end{proof}

\begin{lemma}\label{push:src:D.1.2}
\StatementRecord{push:src:D.1.2}{proved}{Historical events: Claude Fable 5.2 (Anthropic), effort high, 16 November 2026: carried as background without fault found; targeted proof check; used and rechecked in this run. Claude Fable 5.2 (Anthropic), effort high, 17 November 2026: carried as background without fault found. GPT-6 Astra (OpenAI), ChatGPT Pro mode, effort Pro, S6\_09\_PAPERS, 17 November 2026: targeted proof check. Also checked in this writing. Pre-continuation writing events: GPT-6 Astra (OpenAI), ChatGPT Pro mode, effort Pro, S6\_10\_PAPERS, 18 November 2026: targeted proof check. Continuation metadata events: GPT-6 Astra (OpenAI), ChatGPT Pro mode, effort Pro, S6\_10\_PAPERS continuation, 19 November 2026: used and rechecked in this run; status-vocabulary reconciliation; metadata reconciliation only; no new mathematical check. \emph{Pre-continuation recorded proof scope.} Read complete receiving proof and re-derived displayed sheaf/topological argument. The external logarithmic comparison applications in B4.30/B4.32/B4.33 were audited separately against IKN and its erratum; the independent direct commuting-meridian proof is night:P2.7. Boundary proof checked using intrinsic anticanonical divisor, coherence and proper-image exceptional loci.}{the hypotheses, notation and definitions stated here; the explicit constructions and calculations in the proof}

Let \(V\) be a connected Riemann surface and
\(\pi:\mathcal X\to V\) a flat proper holomorphic map of reduced
complex spaces. Suppose that over a nonempty open subset \(V_0\), the
fibres are smooth compact threefolds whose unique anticanonical divisor
is \(2S_v\), where \(S_v\) is a smooth connected surface with
\(b_1(S_v)=2\). Then there is an open subset \(V^{\rm good}\subset V\), with at
most countable complement, on which the reduced support of the
relative anticanonical divisor is a smooth proper surface family. Its
Albanese elliptic invariant is holomorphic.
\end{lemma}


\begin{proof}
\begin{enumerate}[label=\arabic*.,leftmargin=*]
\item Flatness and the smoothness of the fibres over \(V_0\) imply that
\(\pi\) is a submersion there: the local analytic smoothness criterion
applies because the base is smooth. The relative nonsmooth locus is a
closed analytic subset, by the Jacobian criterion. Its proper image is
analytic by Grauert (1960, \S7, Satz~1), the proper mapping theorem.
It misses \(V_0\), hence is discrete. We remove it and work over the
connected remaining curve \(V'\). Removing a discrete subset from a
connected Riemann surface does not disconnect it. Thus
\(\mathcal X'\to V'\) is a smooth proper map.

\item Write \(L=K_{\mathcal X'/V'}^{-1}\). Grauert's direct-image theorem
(1960, \S6, Hauptsatz~I) makes \(F=\pi_*L\) and \(R^1\pi_*L\)
coherent. We prove the needed base change directly over the smooth
curve; we do not use the originally stated \S7, Satz~2, whose claimed
closed exceptional locus was corrected in Grauert's 1963 erratum.
For a local base parameter \(t\) vanishing at \(v\), smoothness and
invertibility of \(L\) give an exact sequence
\(0\to L\xrightarrow{t}L\to i_*L_v\to0\), where
\(i:X_v\hookrightarrow\mathcal X'\) is the closed embedding.
Its direct-image sequence gives, at \(v\),
\[
 0\longrightarrow F_v/tF_v\longrightarrow H^0(X_v,L_v)
 \longrightarrow (R^1\pi_*L)_v[t]\longrightarrow0.
\]
The last term denotes elements annihilated by \(t\). The support of
the torsion subsheaf of the coherent module \(R^1\pi_*L\) is closed
discrete, by the structure theorem for finitely generated modules
over the local discrete valuation rings of the smooth base curve.
Outside that support the canonical base-change map is an isomorphism.
Also \(F\) is torsion-free: multiplication by a nonzero base function
is injective on sections of \(L\). It is therefore locally free.
Its rank is one by choosing an interior point outside the discrete
torsion support. These statements concern the canonical map, not
merely equal dimensions.
Locally a frame gives a relative section \(\sigma\) of \(L\),
nowhere identically zero on a fibre. Its divisor is independent of the
frame. Let \(D\) be its reduced support. On \(V_0\), the divisor is
\(2S_v\) and \(D\) is exactly \(S_v\).

\item The map \(D\to V'\) is proper, since \(D\) is closed in a proper
space. There is no vertical divisor component: a vertical divisor in
the smooth total space would contain a whole irreducible fibre,
contrary to the nonzero restriction of \(\sigma\). Every horizontal
component has proper image all of the connected base. Since the
interior fibre is irreducible, there is exactly one horizontal
component, and its multiplicity in \(\operatorname{div}\sigma\) is
two. Thus \(\operatorname{div}\sigma=2D\) on this base. In the smooth
ambient space \(D\) is Cartier. It is flat over the curve because
its local rings are torsion-free over the base discrete valuation
ring. Restriction of the Cartier identity to an interior fibre gives
\(2D_v=2S_v\), whence \(D_v=S_v\) with its reduced structure.
Flatness and the smoothness of these fibres make \(D\to V'\) smooth
over \(V_0\). Its analytic nonsmooth locus therefore misses \(V_0\);
the proper mapping theorem once more makes its image discrete.
After deleting that image, \(D\to V^{\rm good}\) is a smooth proper
family of connected surfaces. Each deletion is closed discrete in the
then-current open curve. Their finite union is at most countable,
and the final good locus is open; we do not assert that successive
exceptional sets cannot accumulate at an earlier removed point.

\item Ehresmann's theorem identifies the smooth fibres along every path;
in particular their first Betti number is two. Buchdahl (1999,
Theorem~11) then gives a K\"ahler metric on every fibre. Hence
\(h^{1,0}=1\), every holomorphic one-form is closed, and the integral
period map has a rank-two lattice image. Constant \(h^{1,0}\),
properness, and analytic base change supply a local relative
holomorphic one-form. Locally on a simply connected
base disc, a relative one-form and a basis of the free part of
\(H_1\) give two holomorphic periods. Their ratio lies in one
half-plane after orienting the basis. To see holomorphicity directly,
choose local primitives of the closed relative form on simply
connected fibre charts; their differences are holomorphic functions
of the base, and the finite sum around a transported cycle is its
period. The modular invariant of the ratio is independent of all
choices and is holomorphic. This proves the last assertion without
requiring a global choice of section of the Albanese family.
\end{enumerate}
\end{proof}

\begin{theorem}[nonextension across the parameter boundary]\label{push:src:D.1.3}
\StatementRecord{push:src:D.1.3}{proved}{Historical events: Claude Fable 5.2 (Anthropic), effort high, 16 November 2026: carried as background without fault found; used and rechecked in this run. Claude Fable 5.2 (Anthropic), effort high, 17 November 2026: carried as background without fault found. GPT-6 Astra (OpenAI), ChatGPT Pro mode, effort Pro, S6\_09\_PAPERS, 17 November 2026: targeted proof check. Also checked in this writing. Pre-continuation writing events: GPT-6 Astra (OpenAI), ChatGPT Pro mode, effort Pro, S6\_10\_PAPERS, 18 November 2026: targeted proof check. Continuation metadata events: GPT-6 Astra (OpenAI), ChatGPT Pro mode, effort Pro, S6\_10\_PAPERS continuation, 19 November 2026: used and rechecked in this run; status-vocabulary reconciliation; metadata reconciliation only; no new mathematical check. \emph{Pre-continuation recorded proof scope.} Read complete receiving proof and re-derived displayed sheaf/topological argument. The external logarithmic comparison applications in B4.30/B4.32/B4.33 were audited separately against IKN and its erratum; the independent direct commuting-meridian proof is night:P2.7. Boundary proof checked using intrinsic anticanonical divisor, coherence and proper-image exceptional loci.}{the hypotheses, notation and definitions stated here; \cref{fam:src:A.1.1}; \cref{fam:src:A.1.1}, \cref{fam:src:A.1.2}; \cref{fam:src:A.4.1}; \cref{push:src:D.1.2}; \cref{push:src:D.1.1}}

For every \(w_0\) with \(|w_0|=216\) and every connected open
neighbourhood \(U\) of \(w_0\), there is no flat proper holomorphic
map \(\mathcal X\to U\) from a reduced complex space whose fibres
over \(U\cap\{0<|w|<216\}\) are biholomorphic, fibre by fibre, to
\(X_w\). In particular an extension isomorphic to the given family
over that intersection does not exist.
\end{theorem}


\begin{proof}
\begin{enumerate}[label=\arabic*.,leftmargin=*]
\item We may restrict a putative extension to a small disc meeting the
circle, with closure away from zero. This is a holomorphic base
restriction; properness and flatness are preserved. The calculation
of the canonical evaluation in \cref{fam:src:A.1.1}, gives
\(K_{X_w}=\mathcal O(-2S_2(w))\) and
\(h^0(K_{X_w}^{-1})=1\). Its sole divisor is \(2S_2(w)\), so it is
intrinsic under every fibrewise biholomorphism. The surface has
\(b_1=2\): on its finite product cover the rotation kills the moving
factor in invariant first cohomology and fixes the elliptic factor.
These are the exact \cref{fam:src:A.1.1}, \cref{fam:src:A.1.2} and \cref{fam:src:A.4.1} steps required in Lemma~\ref{push:src:D.1.2}.

\item The at most countable exceptional set in that lemma cannot contain
an arc of the circle. Choose a boundary point outside it and then a smaller disc
on which the reduced divisor is a smooth proper surface family. Its
holomorphic Albanese invariant agrees, pointwise on the inner part,
with \(J(iw/216)\). It is therefore an analytic continuation through
that point, contrary to Proposition~\ref{push:src:D.1.1}. No choice of the
fibrewise biholomorphisms was used. Thus the weaker hypothesis already
leads to the contradiction.
\end{enumerate}
\end{proof}

\subsection{The local quotient and its Hausdorff reflection}\label{push:section:D.2}

\begin{lemma}\label{push:src:D.2.1}
\StatementRecord{push:src:D.2.1}{proved}{Historical events: Claude Fable 5.2 (Anthropic), effort high, 16 November 2026: carried as background without fault found. GPT-6 Astra (OpenAI), ChatGPT Pro mode, effort Pro, S6\_09\_PAPERS, 17 November 2026: targeted proof check. Also checked in this writing. Pre-continuation writing events: GPT-6 Astra (OpenAI), ChatGPT Pro mode, effort Pro, S6\_10\_PAPERS, 18 November 2026: computation or formal certificate; targeted proof check. Continuation metadata events: GPT-6 Astra (OpenAI), ChatGPT Pro mode, effort Pro, S6\_10\_PAPERS continuation, 19 November 2026: used and rechecked in this run; status-vocabulary reconciliation; metadata reconciliation only; no new mathematical check. \emph{Pre-continuation recorded proof scope.} Read complete proof and re-derived the exact period basis change/dual inverse, limiting covector graph, primitive-survivor realization, central Hausdorff quotient, explicit nonseparated sequence, properness obstruction, theta gradient or equianharmonic coefficient as applicable. Cornulier duality locator remains separate. Reran d\_push\_exact.py and d\_push\_equianharmonic.py. Exact period inverse, quadratic limit vector, central gcd, theta heat coefficients, the degree-three equianharmonic mean identity, and coarse derivative 7/(64 pi) pass.}{the hypotheses, notation and definitions stated here; \cref{fam:src:C.2.2}}

Near the order-four root centre there is a holomorphic coordinate
\[
 h=\mu-(1-\tau)/2
\]
with nonzero derivative in \(s_2\). In the fixed complex-linear fibre
coordinates \(\eta=\zeta_1\), \(\xi=\zeta_2+\zeta_1/2\), and the
integral ordered basis
\((\widehat w,\widehat u,\widehat\delta,
\widehat\gamma+3\widehat u-3\widehat w)\), the periods at \(c_*\) are
\begin{equation}\label{eq:d-period-normal}
 (1,\tfrac12),\quad(\tau,\tfrac12+h),\quad(0,1),\quad(6h,a+B(h)),
\end{equation}
where
\[
 \tau=i-\frac{3\pi i}{2}h^2+O(h^4),\qquad
 B(h)=-3ih^2+O(h^4),\qquad
 D_{c_*}=-3|h|^2+O(|h|^4).
\]
For \(c\) near \(c_*\), only the last \(\xi\)-period changes, by
\(c-c_*\).
\end{lemma}


\begin{proof}
\begin{enumerate}[label=\arabic*.,leftmargin=*]
\item The basis change has determinant of absolute value one. Substitution
into the original period matrix gives \eqref{eq:d-period-normal}
exactly, with
\(B=\beta+6h+3(1-\tau)/2-a\).
At the centre \(B(0)=0\). Put
\(G(h,\tau)=\wp_\tau((1-\tau)/2+h)\) and
\(H(h,\tau)=\log(\vartheta_3(h,\tau)/\eta_D(\tau))\).
The half-period identities and theta heat equation give
\[
 G(0,i)=0,\quad G_{h^2}(0,i)=-g_2(i)/4,\quad
 G_\tau(0,i)=\frac{i g_2(i)}{6\pi},
\]
\[
 H(0,i)=\log\sqrt2,\qquad H_\tau(0,i)=0,\qquad
 H_{h^2}(0,i)=-\pi/2.
\]
Here \(G_{h^2}\) is the power-series coefficient. For the derivative
constant, write \(A=\vartheta_2(0,i)^4=\vartheta_4(0,i)^4\).
Then \(g_2(i)=4\pi^4A^2\), while differentiating
\(e_2=(\pi^2/3)(\vartheta_2^4-\vartheta_4^4)\) with the heat
identities gives \(e_2'(i)=2\pi^3iA^2/3\).
The logarithmic derivatives used in this differentiation are
\[
 \frac{\partial_\tau\vartheta_2}{\vartheta_2}
  =\frac{\pi i}{12}(E_2+\vartheta_3^4+\vartheta_4^4),\qquad
 \frac{\partial_\tau\vartheta_4}{\vartheta_4}
  =\frac{\pi i}{12}(E_2-\vartheta_2^4-\vartheta_3^4).
\]
They follow by applying the theta heat equation to the half-period
identities. The coefficient arithmetic is independently checked in
\texttt{d\_push\_exact.py}. In particular \(G_\tau(0,i)\ne0\).
The holomorphic implicit-function theorem and evenness of \(G\) give
\(\tau=\tau(h)\), even in \(h\), with the displayed quadratic
coefficient. Since \(\tau-i\) has order two in the root coordinate,
\(h\) has order one.

\item The identity verified in \cref{fam:src:C.2.2} is
\[
 \log\mathcal N=\operatorname{Re}H-
       \frac{\pi(\operatorname{Im}h)^2}{\operatorname{Im}\tau}.
\]
The displayed coefficients, \(\tau-i=O(h^2)\), and evenness give
\(\log\mathcal N-\log\sqrt2=-(\pi/2)|h|^2+O(|h|^4)\).
Multiplication by \(6/\pi\), using the boundary normalization, gives
the formula for \(D_{c_*}\). Directly from the definition,
\[
 D_{c_*}=\operatorname{Im}B-
             6(\operatorname{Im}h)^2/\operatorname{Im}\tau.
\]
Consequently \(\operatorname{Im}B=-3\operatorname{Re}(h^2)
+O(|h|^4)\). Since \(B\) is holomorphic and \(B(0)=0\), its Taylor
coefficients force \(B=-3ih^2+O(h^4)\). This argument checks the
quadratic coefficient omitted by the proposed first-order approximation.
\end{enumerate}
\end{proof}

\begin{proposition}[nonseparation beyond the rank-drop line]\label{push:src:D.2.2}
\StatementRecord{push:src:D.2.2}{proved}{Historical events: Claude Fable 5.2 (Anthropic), effort high, 16 November 2026: carried as background without fault found. GPT-6 Astra (OpenAI), ChatGPT Pro mode, effort Pro, S6\_09\_PAPERS, 17 November 2026: targeted proof check. Also checked in this writing. Pre-continuation writing events: GPT-6 Astra (OpenAI), ChatGPT Pro mode, effort Pro, S6\_10\_PAPERS, 18 November 2026: computation or formal certificate; targeted proof check. Continuation metadata events: GPT-6 Astra (OpenAI), ChatGPT Pro mode, effort Pro, S6\_10\_PAPERS continuation, 19 November 2026: used and rechecked in this run; status-vocabulary reconciliation; metadata reconciliation only; no new mathematical check. \emph{Pre-continuation recorded proof scope.} Read complete proof and re-derived the exact period basis change/dual inverse, limiting covector graph, primitive-survivor realization, central Hausdorff quotient, explicit nonseparated sequence, properness obstruction, theta gradient or equianharmonic coefficient as applicable. Cornulier duality locator remains separate. Reran d\_push\_exact.py and d\_push\_equianharmonic.py. Exact period inverse, quadratic limit vector, central gcd, theta heat coefficients, the degree-three equianharmonic mean identity, and coarse derivative 7/(64 pi) pass.}{the hypotheses, notation and definitions stated here; \cref{push:src:D.2.1}}

The assertion that total-space nonseparation is confined to
\(\mathbb C(2,-1)^t\), modulo central periods, is false. At the
boundary value \(a=0\), the central points with \((\eta,\xi)=(0,0)\)
and \((0,-3i)\) cannot be separated in the root-family quotient.
They are not related by that complex line and the central periods.
\end{proposition}


\begin{proof}
\begin{enumerate}[label=\arabic*.,leftmargin=*]
\item Choose \(h_n=1/n\), which defines actual root points by Lemma~\ref{push:src:D.2.1},
and take the integral lattice vector
\[
 \lambda_n=n^2(\widehat\gamma+3\widehat u-3\widehat w)
            -6n\widehat w+3n\widehat\delta.
\]
Its period is exactly \((0,n^2B(1/n))\), and this tends to
\((0,-3i)\). Thus the orbit of \((h_n,0)\) also has a representative
converging to \((0,0,-3i)\). The two central quotient points are
nonseparated. In these coordinates the asserted complex line is
\(\{\xi=0\}\), and every central period has real \(\xi\)-coordinate;
therefore their difference is not in the asserted relation. The
quadratic model calculation and its integer coefficients are checked
exactly in \texttt{d\_push\_exact.py}; the analytic error is
\(n^2O(n^{-4})\to0\), so the counterexample is for the actual family.
\end{enumerate}
\end{proof}

\begin{theorem}\label{push:src:D.2.3}
\StatementRecord{push:src:D.2.3}{proved}{Historical events: Claude Fable 5.2 (Anthropic), effort high, 16 November 2026: carried as background without fault found. GPT-6 Astra (OpenAI), ChatGPT Pro mode, effort Pro, S6\_09\_PAPERS, 17 November 2026: targeted proof check. Also checked in this writing. Pre-continuation writing events: GPT-6 Astra (OpenAI), ChatGPT Pro mode, effort Pro, S6\_10\_PAPERS, 18 November 2026: targeted proof check. Continuation metadata events: GPT-6 Astra (OpenAI), ChatGPT Pro mode, effort Pro, S6\_10\_PAPERS continuation, 19 November 2026: used and rechecked in this run; status-vocabulary reconciliation; metadata reconciliation only; no new mathematical check. \emph{Pre-continuation recorded proof scope.} Read complete proof and re-derived the exact period basis change/dual inverse, limiting covector graph, primitive-survivor realization, central Hausdorff quotient, explicit nonseparated sequence, properness obstruction, theta gradient or equianharmonic coefficient as applicable. Cornulier duality locator remains separate.}{the hypotheses, notation and definitions stated here; \cref{push:src:D.3.1}; \cref{eq:d-dual-exact}; \cref{push:src:D.2.2}}

Let \(\mathcal T=(\Delta_h\times\mathbb C^2)/\Lambda\) at the fixed
boundary parameter \(c_*\), after shrinking \(\Delta_h\). Two
central points represented by \((\eta,\xi)\) and
\((\eta',\xi')\) are nonseparated precisely when
\[
 \operatorname{Re}(\xi'-\xi)\in G_a
       :=\tfrac12\mathbb Z+a\mathbb Z\quad(a\in\mathbb Q),
\]
and every pair is nonseparated when \(a\notin\mathbb Q\).
There is no nonseparation between distinct base fibres, or between
distinct points of a noncentral fibre. The Hausdorff reflection
\(\mathcal T^{\mathrm H}\to\Delta_h\) has the original complex tori
over \(h\ne0\), and its central fibre is
\[
 \mathbb R/G_a\simeq S^1\quad(a\in\mathbb Q),\qquad
 \{*\}\quad(a\notin\mathbb Q).
\]
If \(a=m/n\) is in lowest terms, \(n>0\), then
\(G_a=\operatorname{lcm}(2,n)^{-1}\mathbb Z\).
\end{theorem}


\begin{proof}
\begin{enumerate}[label=\arabic*.,leftmargin=*]
\item We give the full limiting-group argument in \ref{push:src:D.3.1}; its formula also
applies to the fixed boundary family with \(c=c_*\), \(h\ne0\).
Put \(d_0=\operatorname{lcm}(2,n)\) in the rational case. In the dual
lattice, the fixed integral covector whose central limit is
\(d_0\operatorname{Re}\xi\) has
\((p,q,r)=(0,0,d_0)\). The exact dual formula
\eqref{eq:d-dual-exact} below shows that it extends continuously, as a
real-linear functional \(\ell_h\) on \(\mathbb C^2\), to that limit:
its last coefficient is \(O(h^2)\), its other moving coefficients
are \(O(h)\). Therefore
\[
 (h,z)\longmapsto\exp(2\pi i\ell_h(z))
\]
is continuous, invariant under every period, and restricts at the
centre to \(\exp(2\pi i d_0\operatorname{Re}\xi)\).
Points with different indicated central classes are separated by this
map to the circle.

\item For the converse take a ray \(h=\epsilon\omega\),
\(\epsilon\downarrow0\), whose real and imaginary parts have
irrational ratio. More precisely choose \(\omega\) so that
\(p\operatorname{Re}\omega+q\operatorname{Im}\omega\ne0\)
for every nonzero pair \((p,q)\) in the countable half-integral
lattice in \ref{push:src:D.3.1}. The asymptotic dual formula shows that all such
covectors escape to infinity. For rational \(a\), the surviving dual
group is exactly \(\mathbb Z d_0\operatorname{Re}\xi\); for
irrational \(a\), it is zero. Pontryagin--Chabauty duality
(Cornulier, Theorem~3.1) now gives the period-group limits
\[
 \{z:d_0\operatorname{Re}\xi(z)\in\mathbb Z\},\qquad
 \mathbb C^2,
\]
respectively. Every element of the indicated limit is the limit of
actual period vectors. Such a sequence supplies intersecting images
of every two prescribed central neighbourhoods, proving
nonseparation. This includes the quadratic-scale differences of
\ref{push:src:D.2.2}, not only first-order differences.

\item Collapse exactly these central classes and leave every other fibre
unchanged. The resulting quotient is Hausdorff. The base coordinate
separates different fibres. Within a noncentral fibre, continuous
functions supported over a compact subdisc away from zero separate
distinct points, because there the original action is properly
discontinuous and gives a smooth torus bundle. In the rational central
fibre the displayed circle-valued function separates different
classes. These functions descend through the proposed quotient, so
it is Hausdorff. Every continuous map from \(\mathcal T\) to a
Hausdorff space identifies nonseparated pairs and therefore factors
uniquely and continuously through this quotient. This proves the
universal property, including its topology. The elementary greatest
common divisor calculation gives the stated value of \(G_a\).
\end{enumerate}
\end{proof}

\begin{proposition}\label{push:src:D.2.4}
\StatementRecord{push:src:D.2.4}{proved}{Historical events: Claude Fable 5.2 (Anthropic), effort high, 16 November 2026: carried as background without fault found. GPT-6 Astra (OpenAI), ChatGPT Pro mode, effort Pro, S6\_09\_PAPERS, 17 November 2026: targeted proof check. Also checked in this writing. Pre-continuation writing events: GPT-6 Astra (OpenAI), ChatGPT Pro mode, effort Pro, S6\_10\_PAPERS, 18 November 2026: targeted proof check. Continuation metadata events: GPT-6 Astra (OpenAI), ChatGPT Pro mode, effort Pro, S6\_10\_PAPERS continuation, 19 November 2026: used and rechecked in this run; status-vocabulary reconciliation; metadata reconciliation only; no new mathematical check. \emph{Pre-continuation recorded proof scope.} Read complete proof and re-derived the exact period basis change/dual inverse, limiting covector graph, primitive-survivor realization, central Hausdorff quotient, explicit nonseparated sequence, properness obstruction, theta gradient or equianharmonic coefficient as applicable. Cornulier duality locator remains separate.}{the hypotheses, notation and definitions stated here; \cref{fam:src:C.1.1}, \cref{fam:src:C.2.3}; \cref{fam:src:C.3.1}, \cref{fam:src:C.3.2}}

At the boundary, the construction over \(B\setminus\{1\}\) is a
complex manifold, since \cref{fam:src:C.1.1}, \cref{fam:src:C.2.3} makes every remaining finite period map
invertible and the cusp completion is unchanged.
The order-four deck action descends continuously to
\(\mathcal T^{\mathrm H}\). On its rational central circle it is
translation by \(-a/4\); on the irrational central point it is
trivial. Its finite quotient is Hausdorff with central fibre a circle
or a point, respectively. Neither quotient is a complex space over
the base compatible with the given holomorphic family off the centre.
The linear foliation by \(\mathbb C(2,-1)^t\) on the original central
torus cover has compact elliptic leaves. For rational \(a\) the
leaves are closed; for irrational \(a\) the closure of a leaf is the
set with its fixed value of \(\operatorname{Im}\xi\).
\end{proposition}


\begin{proof}
\begin{enumerate}[label=\arabic*.,leftmargin=*]
\item The deck transformation is a biholomorphism of the original quotient,
so the universal property extends it uniquely to a homeomorphism of
the Hausdorff reflection. The action on the fixed product coordinate
at the centre is \(\xi\mapsto\xi-\kappa_2/4=\xi-a/4\).
A finite group quotient of a Hausdorff space is Hausdorff: finite
orbits are separated by intersecting finitely many translated open
neighbourhoods. A finite rotation quotient of a circle is a circle.

\item We verify properness of the reflected map, rather than assume it.
On a closed small base disc every noncentral torus point has a
representative in the image of \([0,1]^4\) under its period matrix.
These representatives belong to a single compact subset of
\(\mathbb C^2\). At the centre their images meet every class of the
circle or the point. The image of that compact parameterized
parallelepiped is therefore the full inverse image of the closed
disc in the reflection. The reflected map is proper; the same holds
after the finite quotient. If it were holomorphic on a compatible
complex space, the central fibre would be analytic. A circle cannot
be a complex analytic space: a positive-dimensional analytic germ
has real dimension at least two, and a compact zero-dimensional
analytic space is finite. A central point is also impossible here:
upper semicontinuity of fibre dimension, or the local dimension
inequality at a point in the closure of the three-dimensional smooth
locus, forces a fibre dimension of at least two for a holomorphic
map to a curve. Thus this Hausdorff reflection cannot be glued as a
complex-space filling to the remaining constructed manifold.

\item For the leaf assertion, the moving line is \(\{\xi=0\}\).
At the centre the projection of the period group to the \(\xi\)-line
is \(G_a\subset\mathbb R\). Its intersection with the moving line
is a rank-two lattice in \(\mathbb C_\eta\): the periods are
\(p+qi\) with \((p+q)/2+ka+\ell=0\) for integers
\(p,q,k,\ell\). For irrational \(a\) this is the index-two lattice
\(p+q\equiv0\pmod2\); for rational \(a\) it is either that lattice
or \(\mathbb Z[i]\). Both give intrinsic elliptic leaf manifolds isomorphic to
\(E_i\). In the irrational case their maps into the non-Hausdorff
central quotient need not be topological embeddings; ``leaf'' here
means the image with this intrinsic immersed topology. When \(G_a\)
is discrete, the saturated leaf is closed.
When it is dense in \(\mathbb R\), its saturation has closure
\(\mathbb C_\eta\times(\operatorname{Re}\xi+\mathbb R)\), with
\(\operatorname{Im}\xi\) fixed. This is a statement in the original,
possibly non-Hausdorff central fibre. Its own Hausdorffization is the
rank-three quotient described in \cref{fam:src:C.3.1}, \cref{fam:src:C.3.2}; it is not the central fibre of
\(\mathcal T^{\mathrm H}\).
\end{enumerate}
\end{proof}

\subsection{All local limiting period groups}\label{push:section:D.3}

\begin{theorem}\label{push:src:D.3.1}
\StatementRecord{push:src:D.3.1}{proved}{Historical events: Claude Fable 5.2 (Anthropic), effort high, 16 November 2026: carried as background without fault found. GPT-6 Astra (OpenAI), ChatGPT Pro mode, effort Pro, S6\_09\_PAPERS, 17 November 2026: targeted proof check. Also checked in this writing. Pre-continuation writing events: GPT-6 Astra (OpenAI), ChatGPT Pro mode, effort Pro, S6\_10\_PAPERS, 18 November 2026: computation or formal certificate; targeted proof check. Continuation metadata events: GPT-6 Astra (OpenAI), ChatGPT Pro mode, effort Pro, S6\_10\_PAPERS continuation, 19 November 2026: used and rechecked in this run; status-vocabulary reconciliation; metadata reconciliation only; no new mathematical check. \emph{Pre-continuation recorded proof scope.} Read complete proof and re-derived the exact period basis change/dual inverse, limiting covector graph, primitive-survivor realization, central Hausdorff quotient, explicit nonseparated sequence, properness obstruction, theta gradient or equianharmonic coefficient as applicable. Cornulier duality locator remains separate. Reran d\_push\_exact.py and d\_push\_equianharmonic.py. Exact period inverse, quadratic limit vector, central gcd, theta heat coefficients, the degree-three equianharmonic mean identity, and coarse derivative 7/(64 pi) pass.}{the hypotheses, notation and definitions stated here; \cref{push:src:D.2.1}; \cref{eq:d-dual-exact}}

Let \(c_n\to c_*\) with \(\operatorname{Im}c_n<H_*\), and let
\(h_n\to0\). Set
\[
 \delta_n=-D_{c_n}(h_n)>0,\qquad
 V_n=\left(\frac{6\operatorname{Re}h_n}{\delta_n},
           \frac{6\operatorname{Im}h_n}{\delta_n},
           \frac{\operatorname{Re}(c_n-c_*)}{\delta_n}\right).
\]
Define the discrete group
\[
 L_a=\{(p,q,r)\in\mathbb R^3:
    r\in\mathbb Z,\ p+r/2\in\mathbb Z,
    \ q+r/2\in\mathbb Z,\ ar\in\mathbb Z\}.
\]
It has rank three for rational \(a\), and is
\(\mathbb Z^2\times\{0\}\) for irrational \(a\).
Every convergent subsequence of the closed lattices
\(\Gamma_n=\Pi_{c_n}(h_n)\mathbb Z^4\subset\mathbb C^2\) has the
following form. There is a primitive subgroup \(K\subset L_a\) and
a homomorphism \(\ell:K\to\mathbb R\) such that
\begin{equation}\label{eq:d-limit-group}
 \Gamma_\infty=
 \{(\eta,\xi):p\operatorname{Re}\eta+q\operatorname{Im}\eta
       +r\operatorname{Re}\xi+\ell(p,q,r)\operatorname{Im}\xi
       \in\mathbb Z\text{ for every }(p,q,r)\in K\}.
\end{equation}
Here
\[
 K=\{\lambda\in L_a:\lambda\cdot V_n
                      \text{ has a finite limit}\},\qquad
 \ell(\lambda)=\lim_n\lambda\cdot V_n,
\]
and for \(\lambda\notin K\) the absolute value tends to infinity.
Conversely every primitive \(K\subset L_a\) and every such
\(\ell\) occur for sequences with \(\operatorname{Im}c_n<H_*\).
Thus the possible pointed additive quotient limits are real tori of
any dimension from zero to three for rational \(a\), and from zero
to two for irrational \(a\), with the marked groups given explicitly
by \eqref{eq:d-limit-group}.
\end{theorem}


\begin{proof}
\begin{enumerate}[label=\arabic*.,leftmargin=*]
\item We use the fixed complex-linear coordinates of \ref{push:src:D.2.1}. Write
\(h=X+iY\), \(\tau=R+iT\), and
\(z=B(h)+c-c_*\). The real period matrix, with the basis of \ref{push:src:D.2.1},
is
\[
 A=\begin{pmatrix}
 1&R&0&6X\\
 0&T&0&6Y\\
 1/2&1/2+X&1&a+\operatorname{Re}z\\
 0&Y&0&\operatorname{Im}z
 \end{pmatrix},\qquad
 \det A=T D_c.
\]
For an integral dual label \((m_w,m_u,m_d,m_g)\), put
\(r=m_d\), \(p=m_w-r/2\), \(q=m_u-r/2\).
The real-linear covector \(A^{-t}(m_w,m_u,m_d,m_g)^t\) is
\begin{equation}\label{eq:d-dual-exact}
 \left(p,\frac{q-pR-rX-tY}{T},r,t\right),\quad
 t=\frac{m_g-ra-6pX-6Y(q-pR-rX)/T-r\operatorname{Re}z}{D_c}.
\end{equation}
Multiplication by \(A^t\) proves the formula; the entire equation is
checked symbolically in \texttt{d\_push\_exact.py}.

\item A bounded sequence of dual vectors has bounded integral labels,
since \(A_n\) is bounded. After subselection its label is constant.
Its limit must satisfy \(m_g=ra\), because
\(A_0^t\) has that relation. The remaining labels are exactly
\(L_a\). Conversely consider a fixed label from \(L_a\).
Lemma~\ref{push:src:D.2.1} gives
\[
 T=1+O(|h|^2),\quad R=O(|h|^2),\quad
 \operatorname{Re}B=6XY+O(|h|^4),
\]
\[
 \delta_n=-\operatorname{Im}(c_n-c_*)
                       +3|h_n|^2+O(|h_n|^4).
\]
In particular \(\delta_n\ge C|h_n|^2\) for a fixed \(C>0\)
after shrinking the disc. The numerator in
\eqref{eq:d-dual-exact}, with \(m_g=ra\), is
\[
 -6pX-6qY-r\operatorname{Re}(c-c_*)
       +O(|h|^3)(|p|+|q|+|r|).
\]
The cancellation of the quadratic real term is essential. Division
by \(D_c=-\delta_n\) makes the error \(O(|h_n|)\). Thus
\[
 t_n=(p,q,r)\cdot V_n+o(1).
\]
Whenever this is bounded, the other three coefficients converge to
\((p,q,r)\). This proves that the dual limit, if it exists, is the
discrete graph
\(\{(p,q,r,\ell(p,q,r)):(p,q,r)\in K\}\).
There can be no further elements from varying labels, because a
bounded sequence of labels lies in a finite set. A convergent dual
group cannot contain two different limits with the same label: their
difference would have label zero, whose sole preimage under every
\(A_n^t\) is zero. Nor can a surviving label escape along a
subsequence, since an element of a Chabauty limit must be approximable
along the whole convergent sequence. Consequently every label either
has the indicated finite limit or escapes to infinity. Addition and
integer division show that \(K\) is a subgroup and is primitive.

\item Cornulier's Theorem~3.1 identifies convergence of closed subgroups of
a finite-dimensional real vector space with convergence of their
Pontryagin annihilators. Applying it to the dual graph gives exactly
\eqref{eq:d-limit-group}. This theorem's hypotheses are met by the
real vector space underlying \(\mathbb C^2\) and by the closed full
lattices \(\Gamma_n\); no nonclosed group is used as a Chabauty
point.

\item To prove realizability, extend \(\ell\) linearly to a covector
represented by a vector \(v_0\in\mathbb R^3\). Since \(K\) is
primitive, \(L_a\cap\operatorname{span}_{\mathbb R}K=K\).
Choose \(v\) orthogonal to that span and not orthogonal to any member
of \(L_a\setminus K\); the excluded proper hyperplanes are
countably many, and cannot cover the orthogonal space. If
\(K=L_a\), take \(v=0\). Put \(V_n=nv+v_0=(A_n',B_n',C_n')\).
Choose positive \(d_n\to0\) with
\(d_n(1+\|V_n\|^2)\to0\), and define
\[
 h_n=d_n(A_n'+iB_n')/6,\quad x_n=d_nC_n',\quad
 \epsilon_n=d_n+D_{c_*}(h_n)>0,
 \qquad c_n=c_*+x_n-i\epsilon_n.
\]
The inequality follows from \(D_{c_*}(h)=-3|h|^2+O(|h|^4)\).
These choices give \(-D_{c_n}(h_n)=d_n\) exactly, so their vectors
are the prescribed \(V_n\). On \(K\) the scalar products converge
to \(\ell\), and off \(K\) they escape. This constructs every
claimed limit inside the admissible parameter domain.

\item Finally the dual graph is a discrete free group of rank
\(\operatorname{rank}K\). Its annihilator quotient is a compact
connected real torus of that dimension. In the fixed Euclidean metric,
\(d_n([z],[z'])=\operatorname{dist}(z-z',\Gamma_n)\).
Chabauty convergence gives local uniform convergence of these
functions: minimizers relevant on a bounded set can be restricted to
a bounded set, and the two defining Chabauty approximation properties
then give both inequalities. This identifies the stated pointed
quotient limits. They are assertions about real topological and
metric groups, not complex tori of a varying complex dimension.
\end{enumerate}
\end{proof}

\begin{corollary}\label{push:src:D.3.2}
\StatementRecord{push:src:D.3.2}{proved}{Historical events: Claude Fable 5.2 (Anthropic), effort high, 16 November 2026: carried as background without fault found. GPT-6 Astra (OpenAI), ChatGPT Pro mode, effort Pro, S6\_09\_PAPERS, 17 November 2026: targeted proof check. Also checked in this writing. Pre-continuation writing events: GPT-6 Astra (OpenAI), ChatGPT Pro mode, effort Pro, S6\_10\_PAPERS, 18 November 2026: targeted proof check. Continuation metadata events: GPT-6 Astra (OpenAI), ChatGPT Pro mode, effort Pro, S6\_10\_PAPERS continuation, 19 November 2026: used and rechecked in this run; status-vocabulary reconciliation; metadata reconciliation only; no new mathematical check. \emph{Pre-continuation recorded proof scope.} Read complete proof and re-derived the exact period basis change/dual inverse, limiting covector graph, primitive-survivor realization, central Hausdorff quotient, explicit nonseparated sequence, properness obstruction, theta gradient or equianharmonic coefficient as applicable. Cornulier duality locator remains separate.}{the hypotheses, notation and definitions stated here; \cref{push:src:D.2.3}; \cref{push:src:D.3.1}; \cref{push:src:D.1.3}}

There is no path-independent limiting group at the boundary even
when \(c_n\to c_*\) from the admissible side. At \(h_n=0\) and
\(c_n=c_*-i\epsilon_n\), all of \(L_a\) survives, giving a
three-dimensional real torus if \(a\) is rational and a
two-dimensional real torus if \(a\) is irrational. Other admissible
paths give a point. The fixed-boundary Hausdorff reflection of \ref{push:src:D.2.3}
therefore cannot be identified with all pointed fibre limits.
\end{corollary}


\begin{proof}
\begin{enumerate}[label=\arabic*.,leftmargin=*]
\item The first sequence has \(V_n=0\), hence \(K=L_a\) and \(\ell=0\).
For the second, take \(K=0\) in the explicit realization in \ref{push:src:D.3.1}.
The limits have different real dimensions. Away from \(t=1\), the
periods stay invertible and the local complex families converge to
their ordinary value at \(c_*\); the ambiguity is at the specified
root centre. No statement of convergence of compact complex
threefolds follows from a real-group limit. In particular \ref{push:src:D.1.3}
already excludes the proposed holomorphic proper extension.
\end{enumerate}
\end{proof}

\subsection{The supercritical locus}\label{push:section:D.4}

\begin{proposition}\label{push:src:D.4.1}
\StatementRecord{push:src:D.4.1}{proved}{Historical events: Claude Fable 5.2 (Anthropic), effort high, 16 November 2026: carried as background without fault found. GPT-6 Astra (OpenAI), ChatGPT Pro mode, effort Pro, S6\_09\_PAPERS, 17 November 2026: targeted proof check. Also checked in this writing. Pre-continuation writing events: GPT-6 Astra (OpenAI), ChatGPT Pro mode, effort Pro, S6\_10\_PAPERS, 18 November 2026: targeted proof check. Continuation metadata events: GPT-6 Astra (OpenAI), ChatGPT Pro mode, effort Pro, S6\_10\_PAPERS continuation, 19 November 2026: used and rechecked in this run; status-vocabulary reconciliation; metadata reconciliation only; no new mathematical check. \emph{Pre-continuation recorded proof scope.} Read complete proof and re-derived the exact period basis change/dual inverse, limiting covector graph, primitive-survivor realization, central Hausdorff quotient, explicit nonseparated sequence, properness obstruction, theta gradient or equianharmonic coefficient as applicable. Cornulier duality locator remains separate.}{the hypotheses, notation and definitions stated here; \cref{push:src:D.2.1}; \cref{fam:src:C.1.1}, \cref{fam:src:C.2.3}}

For \(\operatorname{Im}c=H_*+\epsilon\), with
\(0<\epsilon\ll1\), the zero locus \(D_c=0\) is one real-analytic
simple closed curve surrounding \(t=1\). In the coordinate of \ref{push:src:D.2.1},
\[
 |h|^2=\epsilon/3+O(\epsilon^2).
\]
In the fixed coarse coordinate it has the leading shape
\[
 |t-1|=\frac{\pi^4E_4(i)}4\epsilon^2
                         \bigl(1+O(\epsilon)\bigr).
\]
The error is uniform along the curve. The region containing \(t=1\)
has \(D_c>0\); the region containing the cusp has \(D_c<0\).
\end{proposition}


\begin{proof}
\begin{enumerate}[label=\arabic*.,leftmargin=*]
\item The Hessian of \(D_{c_*}(h)=-3|h|^2+O(|h|^4)\) is negative
at zero. Its radial derivative is \(-6|h|+O(|h|^3)\), uniformly
in angle. Consequently each sufficiently small positive level of
\(-D_{c_*}\) has exactly one radius in each direction and is a
real-analytic simple closed curve, with the displayed expansion.
The order-four deck action preserves it, acts freely on it, and the
coarse map \(t\) is a four-fold covering of its image. Since the
punctured root quotient is a punctured disc, that image is again a
simple closed curve. The strict global maximum in \cref{fam:src:C.1.1}, \cref{fam:src:C.2.3}, including the
independent certificate in Part E.4, gives a uniform gap on the
complement of any fixed root neighbourhood; the cusp tends to minus
infinity. Thus no other zero-locus component exists for small enough
\(\epsilon\).

\item Ramanujan's differential identities give
\(j''(i)=-2\pi^2j(i)E_4(i)\): indeed
\(j'=-2\pi i jE_6/E_4\) and
\(E_6'(i)=-\pi iE_4(i)^2\). Since \(j=1728t\), Lemma~\ref{push:src:D.2.1} gives
\[
 t-1=\frac{9\pi^4E_4(i)}4h^4+O(h^6).
\]
Insertion of the radial expansion proves the coarse-coordinate
formula, including its order of error. The sign statement follows
from the strict radial decrease near zero and from the cusp sign.
\end{enumerate}
\end{proof}

\begin{proposition}\label{push:src:D.4.2}
\StatementRecord{push:src:D.4.2}{proved}{Historical events: Claude Fable 5.2 (Anthropic), effort high, 16 November 2026: carried as background without fault found. GPT-6 Astra (OpenAI), ChatGPT Pro mode, effort Pro, S6\_09\_PAPERS, 17 November 2026: targeted proof check. Also checked in this writing. Pre-continuation writing events: GPT-6 Astra (OpenAI), ChatGPT Pro mode, effort Pro, S6\_10\_PAPERS, 18 November 2026: targeted proof check. Continuation metadata events: GPT-6 Astra (OpenAI), ChatGPT Pro mode, effort Pro, S6\_10\_PAPERS continuation, 19 November 2026: used and rechecked in this run; status-vocabulary reconciliation; metadata reconciliation only; no new mathematical check. \emph{Pre-continuation recorded proof scope.} Read complete proof and re-derived the exact period basis change/dual inverse, limiting covector graph, primitive-survivor realization, central Hausdorff quotient, explicit nonseparated sequence, properness obstruction, theta gradient or equianharmonic coefficient as applicable. Cornulier duality locator remains separate.}{the hypotheses, notation and definitions stated here; \cref{fam:src:C.3.1}, \cref{fam:src:C.3.2}; \cref{fam:src:C.1.1}, \cref{fam:src:C.2.3}; \cref{fam:src:A.1.1}, \cref{fam:src:A.1.2}}

For any height \(H=\operatorname{Im}c>H_*\), the zero locus is the
well-defined modular-norm level set
\[
 Z_H=\left\{t:\mathcal N(\mu(t),\tau(t))
                    =2\sqrt3\,e^{-\pi H/6}\right\},\qquad
 \wp_{\tau(t)}(\mu(t))=0.
\]
It is compact away from the cusp, and is a real-analytic set in the
finite root charts. At every regular point it is a smooth real curve.
The construction gives a complex manifold precisely over
\(B\setminus Z_H\), retaining the original Mumford fibre at the cusp.
The open sublevel set \(\{D_c<0\}\), with the cusp adjoined, is
connected. At every point of \(Z_H\), the period quotient is the
rank-three quotient of \cref{fam:src:C.3.1}, \cref{fam:src:C.3.2}; its exact Hausdorffization is described
below.
\end{proposition}


\begin{proof}
\begin{enumerate}[label=\arabic*.,leftmargin=*]
\item The formula for \(D\) in \cref{fam:src:C.1.1}, \cref{fam:src:C.2.3} gives the equation, and the theta
transformation law shows that it is independent of the modular and
period representatives. Compactness away from the cusp follows from
\(D_{\rm part}\to-\infty\). The analytic and regular-level
assertions follow in each root chart from the real-analytic implicit
function theorem. On the complement, the real period matrix is
invertible. On compact subsets of that complement its inverse is
bounded, which proves uniform discreteness and proper discontinuity
of the lattice action. The quotient is therefore a complex manifold
with compact torus fibres. The finite deck action is free for the
same invariant-coordinate reason as \cref{fam:src:A.1.1}, \cref{fam:src:A.1.2}: a nontrivial power translates
that circle by a noninteger fraction, regardless of the sign of
\(D\). Near the cusp \(D<0\), so its existing holomorphic Mumford
completion and its overlaps remain available.
At a zero with irrational kernel line, the fibre quotient is already
non-Hausdorff. At a zero with rational kernel choose a primitive
integral kernel vector. Its holomorphic period vanishes at that root
point but is not identically zero nearby. Its first nonzero Taylor
coefficient, multiplied by arbitrarily large integer multiples and
approached along suitable root directions, yields nonseparated
central points. Thus no neighbourhood of a zero in the naive quotient
is a Hausdorff complex manifold. This justifies the asserted maximal
open base locus.

\item For connectedness, in any holomorphic coordinate \(z\) with
\(T=\operatorname{Im}\tau>0\), differentiation gives
\[
 \partial_z\partial_{\bar z}D_{\rm part}
 =-\frac3T\left|\mu'(z)-
         \frac{\operatorname{Im}\mu}{T}\tau'(z)\right|^2\le0.
\]
Thus the function is superharmonic on the root charts. This property
and the minimum principle descend under the finite quotient; one
can equally apply the principle on a root cover of a putative
minimum. A component of the strict sublevel set not containing the
cusp would have compact closure with boundary at level \(-H\), and
would attain a smaller interior minimum. The strong minimum
principle excludes this. A sufficiently small cusp neighbourhood is
contained in the sublevel set, so all components reaching the cusp
are the same one. Hence the sublevel set is connected.
\end{enumerate}
\end{proof}

\begin{proposition}\label{push:src:D.4.3}
\StatementRecord{push:src:D.4.3}{proved}{Historical events: Claude Fable 5.2 (Anthropic), effort high, 16 November 2026: carried as background without fault found. GPT-6 Astra (OpenAI), ChatGPT Pro mode, effort Pro, S6\_09\_PAPERS, 17 November 2026: targeted proof check. Also checked in this writing. Pre-continuation writing events: GPT-6 Astra (OpenAI), ChatGPT Pro mode, effort Pro, S6\_10\_PAPERS, 18 November 2026: targeted proof check. Continuation metadata events: GPT-6 Astra (OpenAI), ChatGPT Pro mode, effort Pro, S6\_10\_PAPERS continuation, 19 November 2026: used and rechecked in this run; status-vocabulary reconciliation; metadata reconciliation only; no new mathematical check. \emph{Pre-continuation recorded proof scope.} Read complete proof and re-derived the exact period basis change/dual inverse, limiting covector graph, primitive-survivor realization, central Hausdorff quotient, explicit nonseparated sequence, properness obstruction, theta gradient or equianharmonic coefficient as applicable. Cornulier duality locator remains separate.}{the hypotheses, notation and definitions stated here; \cref{fam:src:C.3.1}, \cref{fam:src:C.3.2}}

At a zero of \(D_c\), let \(A:\mathbb R^4\to\mathbb C^2\) be
the real period map and \(\ker A=\mathbb R k\), with the kernel
vector stated in \cref{fam:src:C.3.1}, \cref{fam:src:C.3.2}. Define
\[
 L_k=\{m\in\mathbb Z^4:m\cdot k=0\},\qquad r=\operatorname{rank}L_k,
 \qquad V_k=\bigcap_{m\in L_k}\ker m.
\]
Then
\[
 \overline{A\mathbb Z^4}=A(\mathbb Z^4+V_k),\qquad
 \mathbb C^2/\overline{A\mathbb Z^4}\simeq
                         \mathbb R\times(S^1)^r
\]
as real topological groups. The original quotient is a complex Lie
group precisely when \(k\) is rational, in which case \(r=3\).
Otherwise it is non-Hausdorff. Every Chabauty limit of nearby full
lattices contains \(\overline{A\mathbb Z^4}\); the quotient topology
at the fixed zero and a limit of full lattices need not coincide.
\end{proposition}


\begin{proof}
\begin{enumerate}[label=\arabic*.,leftmargin=*]
\item The last three period columns are real-linearly independent, so the
rank is exactly three. This confirms the linear algebra of \cref{fam:src:C.3.1}, \cref{fam:src:C.3.2}
without a positivity assumption. The annihilator characters of
\(A\mathbb Z^4\) are the real covectors \(\lambda\) with
\(A^t\lambda\in\mathbb Z^4\). Their integral labels belong exactly
to \(L_k\). The smallest rational subspace containing \(k\) is
\(V_k\); the closure of the image of \(\mathbb Z^4\) in
\(\mathbb R^4/\mathbb R k\) is its image after adding \(V_k\).
For an elementary verification, the two groups have the same
annihilator characters, and characters separate closed subgroups of
a real vector space by putting such a subgroup into the form
\(\mathbb R^d\times\mathbb Z^e\). This proves the closure formula.
Inside the three-dimensional image hyperplane its connected part
has dimension \(3-r\), and its discrete quotient has rank \(r\).
The one transverse real direction gives the displayed group.

\item The image is discrete exactly when the kernel line has a primitive
integral generator: extend such a generator to an integral basis for
the converse; for necessity, a discrete subgroup of a real
three-space has rank at most three, forcing a nonzero integral
kernel. A nondiscrete finitely generated subgroup cannot be closed,
since a nondiscrete closed subgroup contains an uncountable vector
group. Finally each fixed integral period converges to its limiting
period; a Chabauty limit is closed and therefore contains their
closure. This does not identify it with that closure.
\end{enumerate}
\end{proof}

\begin{proposition}\label{push:src:D.4.4}
\StatementRecord{push:src:D.4.4}{proved}{Historical events: Claude Fable 5.2 (Anthropic), effort high, 16 November 2026: carried as background without fault found. GPT-6 Astra (OpenAI), ChatGPT Pro mode, effort Pro, S6\_09\_PAPERS, 17 November 2026: targeted proof check. Also checked in this writing. Pre-continuation writing events: GPT-6 Astra (OpenAI), ChatGPT Pro mode, effort Pro, S6\_10\_PAPERS, 18 November 2026: computation or formal certificate; targeted proof check. Continuation metadata events: GPT-6 Astra (OpenAI), ChatGPT Pro mode, effort Pro, S6\_10\_PAPERS continuation, 19 November 2026: used and rechecked in this run; status-vocabulary reconciliation; metadata reconciliation only; no new mathematical check. \emph{Pre-continuation recorded proof scope.} Read complete proof and re-derived the exact period basis change/dual inverse, limiting covector graph, primitive-survivor realization, central Hausdorff quotient, explicit nonseparated sequence, properness obstruction, theta gradient or equianharmonic coefficient as applicable. Cornulier duality locator remains separate. Reran d\_push\_exact.py and d\_push\_equianharmonic.py. Exact period inverse, quadratic limit vector, central gcd, theta heat coefficients, the degree-three equianharmonic mean identity, and coarse derivative 7/(64 pi) pass.}{the hypotheses, notation and definitions stated here; the explicit constructions and calculations in the proof}

On an ordinary modulus chart where \(P=\wp'_\tau(\mu)\ne0\), put
\[
 L=\partial_\mu\log\vartheta_1(\mu,\tau),\qquad
 Z=L+2\pi i\frac{\operatorname{Im}\mu}{\operatorname{Im}\tau}.
\]
Then the exact real-analytic derivative is
\[
 \partial_\tau D_{\rm part}
       =\frac{iZ}{4\pi^2P}(3ZP-2g_2).
\]
Away also from the ramification of \(j\), the critical-point equation
on the coarse base is therefore \(Z=0\) or \(3ZP=2g_2\).
\end{proposition}


\begin{proof}
\begin{enumerate}[label=\arabic*.,leftmargin=*]
\item Put \(E=\pi^2E_2/3\). The theta heat equation,
\(\wp=-\partial_\mu L-E\), the differential equation
\(\wp''=6\wp^2-g_2/2\), and
\(4\pi i E'=-2E^2+g_2/6\) give, with \(\mu\) held fixed,
\[
 \partial_\tau\wp=
 \frac{2\wp^2-2E\wp+L\wp'-g_2/3}{2\pi i}.
\]
Indeed differentiating the heat equation twice in \(\mu\) gives
\((\wp''-2(\wp+E)^2+2L\wp')/(4\pi i)-E'\), which simplifies to
the displayed formula. Implicit differentiation of \(\wp(\mu)=0\)
therefore yields
\(\mu'=iL/(2\pi)-ig_2/(6\pi P)\).
At a zero of \(\wp\), the same heat equation and the eta derivative
give
\[
 \frac{d}{d\tau}\log(\vartheta_1(\mu,\tau)/\eta_D(\tau))
       =L\mu'+\frac{L^2}{4\pi i}.
\]
Differentiating \(D=\operatorname{Im}\beta-
6(\operatorname{Im}\mu)^2/\operatorname{Im}\tau\), and writing
\(b=\operatorname{Im}\mu/\operatorname{Im}\tau\), now gives
\[
 \partial_\tau D=\frac3\pi Z(\mu'-b)
                         -\frac{3i}{4\pi^2}Z^2,
 \qquad \mu'-b=\frac{iZ}{2\pi}-\frac{ig_2}{6\pi P}.
\]
Substitution proves the formula. Where \(dt/d\tau\ne0\), the chain
rule transfers vanishing of this derivative to the coarse base.
The program \texttt{d\_push\_gradient.py} cross-checks the identity
against a separate symmetric difference of the theta-product norm;
that numerical check is not needed for the algebraic proof.
\end{enumerate}
\end{proof}

\begin{proposition}\label{push:src:D.4.5}
\StatementRecord{push:src:D.4.5}{proved}{Historical events: Claude Fable 5.2 (Anthropic), effort high, 16 November 2026: carried as background without fault found. GPT-6 Astra (OpenAI), ChatGPT Pro mode, effort Pro, S6\_09\_PAPERS, 17 November 2026: targeted proof check. Also checked in this writing. Pre-continuation writing events: GPT-6 Astra (OpenAI), ChatGPT Pro mode, effort Pro, S6\_10\_PAPERS, 18 November 2026: computation or formal certificate; targeted proof check. Continuation metadata events: GPT-6 Astra (OpenAI), ChatGPT Pro mode, effort Pro, S6\_10\_PAPERS continuation, 19 November 2026: used and rechecked in this run; status-vocabulary reconciliation; metadata reconciliation only; no new mathematical check. \emph{Pre-continuation recorded proof scope.} Read complete proof and re-derived the exact period basis change/dual inverse, limiting covector graph, primitive-survivor realization, central Hausdorff quotient, explicit nonseparated sequence, properness obstruction, theta gradient or equianharmonic coefficient as applicable. Cornulier duality locator remains separate. Reran d\_push\_exact.py and d\_push\_equianharmonic.py. Exact period inverse, quadratic limit vector, central gcd, theta heat coefficients, the degree-three equianharmonic mean identity, and coarse derivative 7/(64 pi) pass.}{the hypotheses, notation and definitions stated here; \cref{push:section:D.4}; \cref{fam:src:C.1.1}, \cref{fam:src:C.2.3}; \cref{push:src:D.4.4}}

At the equianharmonic point \(\tau=\rho=e^{\pi i/3}\), with the
original period representative \(\mu=(2-\rho)/3\), one has
\[
 \mathcal N(\mu,\rho)=3^{1/4},\qquad
 D_{\rm part}(z_1)=-\frac3{2\pi}\log48.
\]
Consequently the zero locus passes through the coarse point \(t=0\)
exactly at the height
\[
 H_0=\frac3{2\pi}\log48
 =H_*+\frac3{2\pi}\log\frac43.
\]
It has \(D_c(0)<0\) below this height and \(D_c(0)>0\) above it.
In the fixed coarse coordinate the function is \(C^1\) at zero and
\[
 D_{\rm part}(t)=D_{\rm part}(0)
       +\frac7{32\pi}\operatorname{Re}t+O(|t|^{4/3}),\qquad
 \partial_tD_{\rm part}(t)=\frac7{64\pi}+O(|t|^{1/3}).
\]
Thus at \(H=H_0\) the level set is a regular \(C^1\) arc through
zero, tangent to the imaginary axis. This proves the crossing value
and its local regularity requested in \ref{push:section:D.4}, without using an unproved
closed form for \(\beta^{\rm part}(z_1)\).
\end{proposition}


\begin{proof}
\begin{enumerate}[label=\arabic*.,leftmargin=*]
\item For a fixed modulus write \(T=\operatorname{Im}\tau\) and
\[
 u(z)=\log\mathcal N(z,\tau)
 =\log\left|\frac{\vartheta_1(z,\tau)}{\eta_D(\tau)}\right|
       -\frac{\pi(\operatorname{Im}z)^2}{T}.
\]
The theta transformation law makes \(u\) periodic for
\(\Lambda=\mathbb Z+\mathbb Z\tau\). It is locally integrable,
with \(u(z)=\log|z|+c_\tau+o(1)\) at zero, and has no other
singularity on the torus. We first fix its additive normalization
exactly. In the Jacobi product, with
\(Q=e^{2\pi i\tau}\) and \(X=e^{2\pi iz}\), the quotient is
\[
 \frac{\vartheta_1(z,\tau)}{\eta_D(\tau)}
 =iQ^{1/12}X^{-1/2}(1-X)
       \prod_{m\ge1}(1-Q^mX)(1-Q^mX^{-1}),
\]
up to a constant of absolute value one. For
\(z=x+iy\), \(0<y<T\), every product factor has horizontal
mean logarithm zero, by the power series for
\(\log(1-ae^{2\pi ix})\), \(|a|<1\). The series converges
uniformly on closed substrips. Thus
\[
 \int_0^1u(x+iy)\,dx
       =-\frac{\pi T}{6}+\pi y-\frac{\pi y^2}{T}.
\]
The logarithmic singularities at the strip endpoints are
integrable; alternatively one integrates first on closed substrips
and passes to the limit. Integration in \(y\) gives
\(\int_{\mathbb C/\Lambda}u\,dA/T=0\). A tilted fundamental
parallelogram has the same horizontal averages, by period one.
This proves the zero-mean normalization directly.

\item For \(\tau=\rho\) put \(\alpha=1+\rho\). Since
\(\rho^2=\rho-1\), multiplication by \(\alpha\) on
\((1,\rho)\) has integral matrix
\[
 \begin{pmatrix}1&-1\\1&2\end{pmatrix},\qquad \det=3.
\]
Moreover \(\alpha\mu=1\), \(3\mu=2-\rho\in\Lambda\),
and \(\mu\notin\Lambda\). Its kernel on the torus is therefore
\(K=\{0,\mu,-\mu\}\), and \(|\alpha|^2=3\).
The distributional Euclidean Laplacian is
\[
 \Delta u=2\pi\delta_0-\frac{2\pi}{T}.
\]
Indeed the holomorphic theta factor contributes its simple zero,
and the Gaussian contributes the displayed constant. Pulling back
by the covering \(z\mapsto\alpha z\) shows that
\(u(\alpha z)\) and \(\sum_{v\in K}u(z+v)\) have the same
Laplacian: both have one logarithmic mass at every point of \(K\)
and constant term \(-6\pi/T\). Their difference is consequently
a periodic harmonic function. The logarithmic singularities cancel
locally, so the removable-singularity theorem gives an ordinary
harmonic function on the compact torus. It is constant.
Multiplication by \(\alpha\) and translation preserve normalized
Haar measure. The mean just computed therefore makes that constant
zero, proving the exact distribution identity
\[
 \mathcal N(\alpha z,\rho)
       =\prod_{v\in K}\mathcal N(z+v,\rho).
\]
As \(z\to0\), divide by \(\mathcal N(z,\rho)\). The left
ratio tends to \(|\alpha|=\sqrt3\). Since \(\mathcal N\) is
even, the other two factors tend to \(\mathcal N(\mu,\rho)^2\).
Positivity now gives \(\mathcal N(\mu,\rho)=3^{1/4}\).

\item Finally the exact \cref{fam:src:C.1.1}, \cref{fam:src:C.2.3} normalization is
\[
 D_{\rm part}=\frac6\pi
          \log\frac{\mathcal N(\mu,\tau)}{2\sqrt3}.
\]
Substitution gives \(-3\log48/(2\pi)\). The boundary height is
\(H_*=3\log6/\pi\), so the stated difference is
\(3\log(4/3)/(2\pi)>0\). Addition of \(\operatorname{Im}c\)
proves the crossing height and sign assertions.

\item We verify regularity in the coarse coordinate rather than infer it
from the ramified modulus. Multiplication by \(\rho^2\) is an
isometry of the torus and preserves the zero-mean Green norm just
used. Also \((\rho^2-1)\mu=\rho-1\in\Lambda\). Hence its
nontrivial rotation fixes \(\mu\), and the real gradient of
\(\log\mathcal N\) vanishes there. In the notation of \ref{push:src:D.4.4} this
says \(Z(\rho)=0\). Write \(P_0=\wp'_\rho(\mu)\ne0\),
\(E=\pi^2E_2/3\), and \(b=\operatorname{Im}\mu/T\).
The identities of \ref{push:src:D.4.4} give \(\mu'=b\) at \(\rho\), since
\(g_2(\rho)=0\) and \(L=-2\pi ib\). The heat equation, with
\(\mu\) held fixed, gives
\[
 L_\tau=\frac{-P-2EL}{4\pi i},\qquad L_\mu=-E.
\]
When differentiating \(Z\) along \(\mu(\tau)\), its two terms
involving \(E b\) cancel and the derivative of \(b\) vanishes at
\(\rho\). Thus
\[
 \partial_\tau Z(\rho)=\frac{iP_0}{4\pi},\qquad
 \partial_{\bar\tau}Z(\rho)=0.
\]
Ramanujan's identity and the Weierstrass equation yield
\[
 g_2'(\rho)=-\frac{3i}{\pi}g_3(\rho)
            =\frac{3i}{\pi}P_0^2.
\]
For \(v=\tau-\rho\), it follows that
\[
 3ZP-2g_2=-\frac{21iP_0^2}{4\pi}v+O(|v|^2),\qquad
 \partial_\tau D_{\rm part}
       =\frac{21iP_0^2}{64\pi^4}v^2+O(|v|^3).
\]
Here all remainders are real analytic in \(v,\bar v\).
The exact normalization \(t=j/1728\) gives
\[
 t=-\frac{i g_3(\rho)}{\pi^3}v^3+O(v^4)
   =\frac{iP_0^2}{\pi^3}v^3+O(v^4).
\]
For example this follows immediately from
\(t=E_4^3/(E_4^3-E_6^2)\),
\(E_4'(\rho)=-(2\pi i/3)E_6(\rho)\), and
\(g_3=8\pi^6E_6/27\). Dividing the derivative expansions on
\(v\ne0\) proves
\(\partial_tD_{\rm part}=7/(64\pi)+O(|t|^{1/3})\).
The conjugate derivative has the conjugate expansion. The resulting
real Taylor expansion extends the coarse function as a \(C^1\)
function with nonzero real gradient, giving the asserted local arc
by the implicit function theorem. In particular this is not a
singular level at \(t=0\), despite the vanishing modulus derivative
at the order-three ramification point.

\item The degree-three lattice matrix, its kernel representatives and all
rational derivative coefficients are checked exactly in
\path{code/d_push_equianharmonic.py}. Independent theta-derivative
and finite-difference evaluations at 90 digits are retained in
\path{code/d_push_equianharmonic_crossing.py} and its output;
they support, but do not supply, the analytic proof.
\end{enumerate}
\end{proof}

\begin{cited}\label{push:src:D.4.6}
\StatementRecord{push:src:D.4.6}{proved}{Historical events: Claude Fable 5.2 (Anthropic), effort high, 16 November 2026: carried as background without fault found. GPT-6 Astra (OpenAI), ChatGPT Pro mode, effort Pro, S6\_09\_PAPERS, 17 November 2026: targeted proof check. This writing records metadata reconciliation only; no new mathematical check. Continuation metadata events: GPT-6 Astra (OpenAI), ChatGPT Pro mode, effort Pro, S6\_10\_PAPERS continuation, 19 November 2026: used and rechecked in this run; status-vocabulary reconciliation; metadata reconciliation only; no new mathematical check. \emph{Pre-continuation recorded proof scope.} Current-writing verification remains in progress. The earlier checks named above retain their recorded scope.}{the hypotheses, notation and definitions stated here; \cref{push:src:D.4.2}; \cref{push:src:D.4.5}; \cref{push:src:D.4.7}; \cref{push:src:D.4.8}}

This identifier recorded an open step at the earlier checkpoint:
the minimum principle in \ref{push:src:D.4.2} did not exclude additional maxima,
and the exact crossing in \ref{push:src:D.4.5} did not classify all levels. That
open status is superseded by \ref{push:src:D.4.7}--\ref{push:src:D.4.8} below. We retain the
record of the unsuccessful reduction. Lin and Wang (2010,
Theorem~1.3) count critical points of the Green function on a fixed
elliptic curve, where \(Z=0\) is its critical equation. That theorem
does not count zeros along \(\wp_\tau(\mu(\tau))=0\) as the modulus
varies, nor does it address \(3ZP-2g_2\). Our completed argument
therefore supplies a separate nonvanishing certificate for both
factors, with analytic local and cusp estimates.
\end{cited}
The expanded source derivation is retained in the proof compendium, Paper~I, label \nolinkurl{push:src:D.4.6}.




\begin{proposition}\label{push:src:D.4.7}
\StatementRecord{push:src:D.4.7}{proved}{Historical events: Claude Fable 5.2 (Anthropic), effort high, 15 November 2026: carried as background without fault found. Claude Fable 5.2 (Anthropic), effort high, 16 November 2026: carried as background without fault found. GPT-6 Astra (OpenAI), ChatGPT Pro mode, effort Pro, S6\_09\_PAPERS, 17 November 2026: targeted proof check. Also checked in this writing. Pre-continuation writing events: GPT-6 Astra (OpenAI), ChatGPT Pro mode, effort Pro, S6\_10\_PAPERS, 18 November 2026: targeted proof check. Continuation metadata events: GPT-6 Astra (OpenAI), ChatGPT Pro mode, effort Pro, S6\_10\_PAPERS continuation, 19 November 2026: used and rechecked in this run; status-vocabulary reconciliation; metadata reconciliation only; no new mathematical check. \emph{Pre-continuation recorded proof scope.} Independently re-proved all three local/cusp nonvanishing claims from the checked global critical-point classification via the exact affine map Dpart=(6/pi)(L-log(2sqrt3)); computed the coarse derivative7/(64pi). The earlier long interval route remains a second proof in the compendium.}{\cref{ari:C.8.15}; \cref{push:src:D.4.4}, \cref{push:src:D.4.5}; \cref{ari:C.1.1}}

The function $D_{\rm part}$ has no critical point on the punctured
root neighbourhood $0<|\tau-i|\leq3/20000$, no critical point on the
coarse neighbourhood represented by $|\tau-\rho|\leq1/50$, and no
critical point where $\Im\tau\geq2$. At $\rho$ the two factors
$Z$ and $F=3ZP-2g_2$ of \ref{push:src:D.4.4} vanish, but the coarse gradient is
nonzero, as in \ref{push:src:D.4.5}. The earlier interval proof and its 400 local rectangles remain in
the S6\_09 proof compendium under this label. We give the consequence
of the full critical-point classification here.
\end{proposition}

\begin{proof}
\begin{enumerate}[label=\arabic*.,leftmargin=*]
\item The two norms agree by their exact sigma--theta formula in
\ref{ari:C.1.1}. Hence, with $L=\log\mathcal N$,
\[
 D_{\rm part}=\frac6\pi\bigl(L-\log(2\sqrt3)\bigr).
\]
Theorem~\ref{ari:C.8.15} proves that the coarse differential
is nonzero away from $t=1$, including at $t=0$, where
$L_t(0)=7/384$. It gives
$(D_{\rm part})_t(0)=7/(64\pi)$. Its two exact root-critical
orbits are $i$ and $\rho$, with the factor vanishing specified
in \ref{push:src:D.4.4}, \ref{push:src:D.4.5}.

\item The root-to-coarse map is unramified off those elliptic orbits.
The stated small disc about $i$ contains no other elliptic
representative, and removing its centre therefore removes its
only critical point. The coarse disc about $\rho$ has the
nonzero differential just calculated. Finally
\[
 \Im\frac{ai+b}{ci+d}=\frac1{c^2+d^2}\le1,\qquad
 \Im\frac{a\rho+b}{c\rho+d}
 =\frac{\Im\rho}{c^2+cd+d^2}\le1.
\]
Thus no elliptic-orbit critical point lies at height at least two.
These facts prove every claimed domain of nonvanishing.
\end{enumerate}
\end{proof}



\begin{proposition}\label{push:src:D.4.8}
\StatementRecord{push:src:D.4.8}{proved}{Historical events: Claude Fable 5.2 (Anthropic), effort high, 16 November 2026: carried as background without fault found. GPT-6 Astra (OpenAI), ChatGPT Pro mode, effort Pro, S6\_09\_PAPERS, 17 November 2026: targeted proof check. Also checked in this writing. Pre-continuation writing events: GPT-6 Astra (OpenAI), ChatGPT Pro mode, effort Pro, S6\_10\_PAPERS, 18 November 2026: targeted proof check. Continuation metadata events: GPT-6 Astra (OpenAI), ChatGPT Pro mode, effort Pro, S6\_10\_PAPERS continuation, 19 November 2026: used and rechecked in this run; status-vocabulary reconciliation; metadata reconciliation only; no new mathematical check. \emph{Pre-continuation recorded proof scope.} Independently re-proved the entire level-locus and degeneration statement from exact L-to-H level conversion and the checked strict radial foliation; retained the determinant/rank-three quotient argument. The older global interval/transport proof remains in S6\_09 as a second proof.}{\cref{ari:C.8.15}; \cref{push:src:D.4.1}, \cref{push:src:D.4.2}, \cref{push:src:D.4.3}, \cref{push:src:D.4.5}; \cref{ari:C.1.1}}

For every \(H>H_*\), the degeneracy locus \(Z_H\) is exactly one
embedded \(C^1\) Jordan curve in the finite \(t\)-plane. It is
real analytic unless \(H=H_0=3\log48/(2\pi)\); at that height it
is real analytic away from \(t=0\) and has the regular crossing
described in \ref{push:src:D.4.5}. There are no singular coarse levels at any
supercritical height. The bounded complementary disc contains
\(t=1\) and has \(D_c>0\); the complementary disc in
\(\mathbb P^1\), containing the cusp, has \(D_c<0\). The point
\(t=0\) is in the latter disc for \(H<H_0\), on the curve for
\(H=H_0\), and in the former disc for \(H>H_0\).

The bounded discs are strictly nested as \(H\) increases, shrink
to \(\{1\}\) as \(H\downarrow H_*\), and exhaust the finite
plane as \(H\to\infty\). Their small-height boundary shape is
the expansion in \ref{push:src:D.4.1}. The maximal open base locus of the complex
manifold construction consists of these two complementary discs.
Approaching their common boundary, the covolume of the marked
root torus lattice is
\[
 \operatorname{covol}(\Pi_c\mathbb Z^4)
       =\operatorname{Im}\tau\,|D_c|\longrightarrow0.
\]
At the boundary the period map has rank three and the exact quotient
and Hausdorff reflection are those of \ref{push:src:D.4.3}. Thus this statement
classifies the locus and the local degeneration requested in D.4.
\end{proposition}

\begin{proof}
\begin{enumerate}[label=\arabic*.,leftmargin=*]
\item Use the same exact affine relation and write
\[
 f=-D_{\rm part},\qquad
 Z_H=\{L=\log(2\sqrt3)-\pi H/6\}.
\]
The maximum $L(1)=\log\sqrt2$ corresponds exactly to
$H_*=3\log6/\pi$. Thus $H>H_*$ corresponds to a level
strictly below that maximum. Theorem~\ref{ari:C.8.15} proves
that each such level is one strictly star-shaped $C^1$ Jordan
curve about $1$, real analytic away from the possible point
$t=0$. Its nonzero coarse derivative there proves regularity
also on that exceptional level.

\item Since $L(0)=\tfrac14\log3$, its height is
$H_0=3\log48/(2\pi)$. The exact equation $D_c=H-f$ gives
the two signs, the position of $0$, and the bounded component
$\{f<H\}$ containing $1$. Strict radial monotonicity proves
strict nesting. The unique minimum $f(1)=H_*$ and properness
of $f$ off $1$ prove shrinking as $H\downarrow H_*$ and
exhaustion of every compact subset as $H\to\infty$.
The expansion at the minimum is \ref{push:src:D.4.1}.

\item In the ordered real coordinates and marked lattice basis of
\ref{push:src:D.4.3}, direct expansion of the period determinant
gives
\[
 \det_{\R}\Pi_c
   =\Im\tau\,\Im\beta_c-6(\Im\mu)^2
   =\Im\tau\,D_c.
\]
Its absolute value is the covolume and tends to zero at the
level curve. The same matrix retains three independent columns
there, so its rank is exactly three. Apply the exact quotient
and Hausdorff-reflection computation of \ref{push:src:D.4.3}.
The local proper-action criterion of \ref{push:src:D.4.2}
gives complex manifolds on both complementary regions and
excludes every point of their common boundary, including its
discrete fixed-fibre quotients. This proves the maximality and
all the local degeneration assertions.
\end{enumerate}
\end{proof}




\section{The uniform corner obstruction}

\subsection{Period transport and the actual parameter descent}

\begin{theorem}\label{night:P2.1}\label{stmt:P2.1}
\StatementRecord{night:P2.1}{proved}{Historical events: GPT-6 Astra (OpenAI), ChatGPT Pro mode, effort Pro, S6\_06\_BLIND, 16--18 November 2026: targeted proof check. Also checked in this writing. Pre-continuation writing events: GPT-6 Astra (OpenAI), ChatGPT Pro mode, effort Pro, S6\_10\_PAPERS, 18 November 2026: targeted proof check. Continuation metadata events: GPT-6 Astra (OpenAI), ChatGPT Pro mode, effort Pro, S6\_10\_PAPERS continuation, 19 November 2026: used and rechecked in this run; status-vocabulary reconciliation; metadata reconciliation only; no new mathematical check. \emph{Pre-continuation recorded proof scope.} Read complete cp25 proofs, re-derived actual origin/deck differences, quotient fundamental-group presentation, finite-orbit vanishing, connected-corner commuting ratio, higher-multiplicity determinant and rational-base recovery, and untwisted plumbing/virtual Euler comparison. Arithmetic chart replay and exact external locators remain separately recorded gates.}{the definitions and exact calculations in this statement}

For the prescribed $(3,4)$ construction with integral twists $(l_0,l_1,l_2)$, put
\[
 b=4l_1+3l_2,\qquad p=12l_0-4l_1-3l_2,\qquad a=b-12l_0=-p.
\]
On the logarithmic cover of the cusp bidisc the period columns are
\[
 \gamma=\binom{6\mu}{C-s-12v},\quad
 u=\binom{s+h}{\mu},\quad w=\binom10,\quad\delta=\binom01,
\]
where $h,\mu,C$ are holomorphic at $q=0$. Their positive-loop column comparisons are
\[
 K_q=\begin{pmatrix}1&0&0&0\\0&1&0&0\\0&1&1&0\\-1&0&0&1\end{pmatrix},\qquad
 K_r=\begin{pmatrix}1&0&0&0\\0&1&0&0\\0&0&1&0\\-12&0&0&1\end{pmatrix}.
\]
The parameter descent on the ordinary Jacobian chart is
\begin{equation}\label{eq:P2H}
 R:(t,v,\zeta)\longmapsto(t,v+1,\zeta+H(t)\delta),\qquad
 H(t)=-\frac{4l_1\log t+3l_2\log(t-1)}{2\pi i}.
\end{equation}
In the cusp collar it extends across $q=0$ and has coefficient $a$ on $s\delta$.
\end{theorem}
\begin{proof}
\begin{enumerate}[label=\arabic*.,leftmargin=*]
\item We substitute $c=-12v$ in the supplied period matrix and use its cusp expansions $\tau=s+h$, $\beta^{\rm part}=C-s$. Direct subtraction gives the two displayed matrices. On fixed vector coordinates the lattice conjugation matrices are $K_q^{-1}$ and $K_r^{-1}$. In the characters $Z,W$, addition of $u$ and $\gamma$ gives respectively
\[
 (Z,W)\mapsto(qe^{2\pi ih}Z,e^{-2\pi i\mu}W),\qquad
 (Z,W)\mapsto(e^{12\pi i\mu}Z,qr^{12}e^{-2\pi iC}W).
\]
These are holomorphic character identities, including the signs of both meridians.

\item The collar section carrying the primitive $\widehat\gamma$ class has lift
\begin{equation}\label{eq:P2zg}
 z_\gamma(s,v)=\binom{6s\mu}{-s(s+1)/2+s(C-12v)}.
\end{equation}
It satisfies $z_\gamma(s+1,v)-z_\gamma(s,v)=\gamma(s+1,v)$ and $z_\gamma(s,v+1)-z_\gamma(s,v)=-12s\delta$. We therefore use the cusp coordinate $\eta=\zeta+l_0z_\gamma$, retaining the actual collar integer.

\item At finite point $j$, the logarithmic gluing is translation by $(\log z_j)/(2\pi i)\Pi v_j$, with $t_j=z_j^{m_j}$ and $\widehat\gamma$ coefficient of $v_j$ equal to $l_j$. Changing $v$ by one changes this gluing by $-12l_j\log z_j/(2\pi i)\delta$. This is $-4l_1\log t/(2\pi i)\delta$ at $0$ and $-3l_2\log(t-1)/(2\pi i)\delta$ at $1$. The complementary summand of $H$ is holomorphic on each disc, so its translation extends. The finite affine action changes by $-12l_j/m_j$ times the integral period $\delta$, and is unchanged on the torus. Thus \eqref{eq:P2H} is the actual descent compatible with both finite caps. Branch changes of $H$ are integral $\delta$ periods.

\item At infinity,
\[
 H=bs+U(q),\qquad U(q)=-\frac{3l_2\log(1-q)}{2\pi i}.
\]
In $\eta$ the descent is translation by $(as+U)\delta$. Adding the lattice period $a\gamma$ rewrites this vector as $(6a\mu,U+aC-12av)$. It is holomorphic at $q=0$ for each $v$; its exponentials act by units on every toric chart. This proves the extension and its collar coefficient. The identities are recomputed by \texttt{code/part2\_corner.py} and \texttt{code/part2\_origins\_and\_group.py}.
\end{enumerate}
\end{proof}

\begin{theorem}\label{night:P2.2}\label{stmt:P2.2}
\StatementRecord{night:P2.2}{proved}{Historical events: Claude Opus 5.5 (Anthropic), effort max, sweep 1, scope s6, 18 November 2026: used and rechecked in this run. GPT-6 Astra (OpenAI), ChatGPT Pro mode, effort Pro, S6\_06\_BLIND, 16--18 November 2026: targeted proof check. Also checked in this writing. Pre-continuation writing events: GPT-6 Astra (OpenAI), ChatGPT Pro mode, effort Pro, S6\_10\_PAPERS, 18 November 2026: targeted proof check. Continuation metadata events: GPT-6 Astra (OpenAI), ChatGPT Pro mode, effort Pro, S6\_10\_PAPERS continuation, 19 November 2026: used and rechecked in this run; status-vocabulary reconciliation; metadata reconciliation only; no new mathematical check. \emph{Pre-continuation recorded proof scope.} Read complete cp25 proofs, re-derived actual origin/deck differences, quotient fundamental-group presentation, finite-orbit vanishing, connected-corner commuting ratio, higher-multiplicity determinant and rational-base recovery, and untwisted plumbing/virtual Euler comparison. Arithmetic chart replay and exact external locators remain separately recorded gates.}{P2.3}

On $q=Q^{12}$, actual extending origins have lifts
\begin{align*}
 o&=\binom{-b\mu s/2}{bs^2/24-bCs/12+l_0s/2+v(bs+U)},\\
 \sigma&=-l_0z_\gamma+\binom{6a\mu v}{vU+aCv-6av^2+6av}.
\end{align*}
The section $\sigma$ extends across the cusp for $r\ne0$, and $o$ extends across the parameter central fibre of the partial filling. Their difference is
\begin{equation}\label{eq:P2F}
 F=\sigma-o=a\binom{\mu(s+12v)/2}{C(s+12v)/12-(s+12v)^2/24+6v}.
\end{equation}
On $q=Q^{12m}$, $r=R^n$, with $m,n$ positive integers,
\begin{align}
 \Delta_QF&=am\gamma-6am^2\delta,\label{eq:P2FQ}\\
 \Delta_RF&=an\gamma-6an(n-1)\delta.\label{eq:P2FR}
\end{align}
The radial class modulo $\Z\widehat\delta$ is $-pn\widehat\gamma$. The proposed $+pn\widehat\gamma$ has the opposite sign in the fixed conventions.
\end{theorem}
\begin{proof}
\begin{enumerate}[label=\arabic*.,leftmargin=*]
\item The lift descent identities are
\begin{align*}
 \Delta_{s,12}o&=-b\gamma+6(b+l_0)\delta,&\Delta_{v,1}o&=H\delta,\\
 \Delta_{s,12}\sigma&=-12l_0\gamma+78l_0\delta,&\Delta_{v,1}\sigma&=H\delta+a\gamma.
\end{align*}
Every additional term is a lattice period. Both lifts therefore define holomorphic sections of the descended torus torsor. In the cusp chart $\sigma+l_0z_\gamma$ is holomorphic and bounded in $q$ and lies in the smooth toric open orbit at $q=0$. It extends there. Subtracting $o$ makes the vector coordinate invariant under the parameter deck. The zero coordinate extends in the rank-one charts of P2.3 below. This constructs the required second origin.

\item Subtracting gives \eqref{eq:P2F}. Integral shifts of $s$ leave $\mu,C,U$ unchanged; expansion of the square gives \eqref{eq:P2FQ}--\eqref{eq:P2FR}. The symbolic program verifies each coefficient. For $(l_0,l_1,l_2)=(0,1,-1)$, $p=-1$, $a=1$, and \eqref{eq:P2F} is the function printed in the prompt. Its radial class is $+n\widehat\gamma$. For every admissible twist $p\ne0$, and the difference between the two proposed classes is $2pn\widehat\gamma\ne0$ in the free quotient lattice.
\end{enumerate}
\end{proof}

\subsection{The partial parameter filling}
\begin{theorem}\label{night:P2.3}\label{stmt:P2.3}
\StatementRecord{night:P2.3}{proved}{Historical events: GPT-6 Astra (OpenAI), ChatGPT Pro mode, effort Pro, S6\_06\_BLIND, 16--18 November 2026: targeted proof check. Also checked in this writing. Pre-continuation writing events: GPT-6 Astra (OpenAI), ChatGPT Pro mode, effort Pro, S6\_10\_PAPERS, 18 November 2026: targeted proof check. Continuation metadata events: GPT-6 Astra (OpenAI), ChatGPT Pro mode, effort Pro, S6\_10\_PAPERS continuation, 19 November 2026: used and rechecked in this run; status-vocabulary reconciliation; metadata reconciliation only; no new mathematical check. \emph{Pre-continuation recorded proof scope.} Read complete cp25 proofs, re-derived actual origin/deck differences, quotient fundamental-group presentation, finite-orbit vanishing, connected-corner commuting ratio, higher-multiplicity determinant and rational-base recovery, and untwisted plumbing/virtual Euler comparison. Arithmetic chart replay and exact external locators remain separately recorded gates.}{P2.2}

There is a proper holomorphic partial filling over $(\Pp^1\times\Delta_r)\setminus\{(\infty,0)\}$ with smooth total space and reduced normal-crossing parameter central divisor. Over a regular finite base point its central fibre is a cycle of twelve ruled surfaces. Its component rotations at $0,1,\infty$ are $4l_1,3l_2,-b$ modulo $12$, with component cover
\[
 y^{12}=t^{4l_1}(t-1)^{3l_2}.
\]
\end{theorem}
\begin{proof}
\begin{enumerate}[label=\arabic*.,leftmargin=*]
\item On a simply connected finite-base chart subtract a local origin with $v$-difference $H\delta$, quotient by $\delta$, and put $W=e^{-2\pi i\eta_2}$, $z=\eta_1$. The remaining lattice transformations are
\begin{align*}
 (z,W)&\mapsto(z+1,W),\\
 (z,W)&\mapsto(z+\tau,e^{-2\pi i\mu}W),\\
 (z,W)&\mapsto(z+6\mu,d(t)r^{12}W),
\end{align*}
with $d$ a nonzero holomorphic unit. The first two give the complement of the zero section of a degree-zero line bundle over $\C/(\Z+\tau\Z)$. Use the infinite chain of node charts
\begin{equation}\label{eq:P2chain}
 X_i=W/r^i,\quad Y_i=r^{i+1}/W,\quad X_iY_i=r\qquad(i\in\Z).
\end{equation}
Adjacent charts are identified on $Y_i\ne0$ by $X_{i+1}=Y_i^{-1}$, $Y_{i+1}=rY_i$. The charts and their inverse overlaps give an infinite projective-line chain at $r=0$. Unit multiplication on $W$ multiplies $X_i,Y_i$ inversely, so all elliptic line-bundle transitions extend. The third period shifts the index by twelve, translates the elliptic coordinate by $6\mu$, and multiplies by $d$. Its quotient has the stated twelve-cycle.

\item The total node chart is smooth by eliminating $r=X_iY_i$, and its central divisor is reduced NC. On a compact smaller finite-base chart the elliptic lattices have uniformly bounded compact fundamental parallelograms. Choose middle annuli and node discs using upper and lower bounds for $d$ and the finitely many transition units. Twelve consecutive indices give a compact fundamental cover over a closed smaller parameter disc. The overlap formulas have closed graphs on that cover, and the positive index translation separates nonadjacent charts. Thus the quotient is Hausdorff and proper. At $r\ne0$ it is exactly the original torus quotient.

\item At root point $j$, addition of $v_j/m_j$ multiplies $W$ by a factor containing $r^{12l_j/m_j}$ and rotates the cycle by $12l_j/m_j$. Every nonidentity power moves the cycle, because $3\nmid l_1$ and $l_2$ is odd. The full finite action is free on the central divisor, including its intersection strata. Its quotient is smooth and preserves the reduced NC parameter divisor. Logarithmic gluing translations differ on overlaps by holomorphic units and integral chain shifts, which extend by the same formulas. The rotations give the cyclic cover; the third is their negative sum and equals $p$ modulo $12$.

\item We glue this finite-base filling to the prescribed cusp filling over $r\ne0$ by the actual translation $\sigma-o$. P2.2 identifies the overlap with the original family. Properness and the Hausdorff property are local on the displayed target, so the glued map has both. The local equations prove flatness over the parameter disc.
\end{enumerate}
\end{proof}

\subsection{Extending the marked base map}
\begin{theorem}\label{night:P2.4}\label{stmt:P2.4}
\StatementRecord{night:P2.4}{proved}{Historical events: GPT-6 Astra (OpenAI), ChatGPT Pro mode, effort Pro, S6\_06\_BLIND, 16--18 November 2026: targeted proof check. Also checked in this writing. Pre-continuation writing events: GPT-6 Astra (OpenAI), ChatGPT Pro mode, effort Pro, S6\_10\_PAPERS, 18 November 2026: targeted proof check. Continuation metadata events: GPT-6 Astra (OpenAI), ChatGPT Pro mode, effort Pro, S6\_10\_PAPERS continuation, 19 November 2026: used and rechecked in this run; status-vocabulary reconciliation; metadata reconciliation only; no new mathematical check. \emph{Pre-continuation recorded proof scope.} Read complete cp25 proofs, re-derived actual origin/deck differences, quotient fundamental-group presentation, finite-orbit vanishing, connected-corner commuting ratio, higher-multiplicity determinant and rational-base recovery, and untwisted plumbing/virtual Euler comparison. Arithmetic chart replay and exact external locators remain separately recorded gates.}{the definitions and exact calculations in this statement}

For every admissible $(3,4)$ twist, the marked map $f:X_c\to\Pp^1$ extends meromorphically over every hypothetical proper Hausdorff analytic filling of a finite parameter cover. A proper modification carries a holomorphic marked map to $\Pp^1$, and hence a proper map to $\Pp^1\times\Delta$.
\end{theorem}
\begin{proof}
\begin{enumerate}[label=\arabic*.,leftmargin=*]
\item Take the reduced irreducible component containing the punctured family. It is closed and proper; discarding other central components preserves that family. Resolve it isomorphically over the smooth punctured family. We use analytic resolution \cite[Theorem 1.6 and the following analytic-neighborhood paragraph, preprint p.~5]{BM97}: reduced complex analytic spaces admit proper modifications by locally finite sequences of smooth-centre blowups with smooth result. It can be chosen unchanged over the nonsingular locus. Properness makes the neighborhoods of the central fibre relatively compact, so the stated analytic hypotheses hold. Write $p:Z\to\Delta$ for the smooth result.

\item The determinant calculation in Part~5 gives $K_{X_c}^{-2}=f^*\OO(1)$. Connected compact fibres give $f_*\OO_X=\OO_{\Pp^1}$, so $h^0(K_{X_c}^{-2})=2$. Let $L=\omega_{Z/\Delta}^{-2}$. Grauert's coherence theorem \cite[Hauptsatz I]{Grauert60} makes $p_*L$ coherent. Multiplication by the parameter is injective on $L$ and on $p_*L$. A finitely generated torsion-free module over $\C\{r\}$ is free, so this direct image has rank two. On the punctured disc, $p$ is a submersion, $L$ is flat and $h^0=2$ is constant. Grauert's base-change theorem \cite[\S7, Satz 5]{Grauert60} identifies its frame with the complete fibrewise linear system. Evaluation gives a meromorphic map $f':Z\dashrightarrow\Pp^1$ whose restriction is $f$ followed by a projective linear transformation on each noncentral fibre.

\item Blowing up the base ideal and resolving its graph gives a proper holomorphic map $(f',p)$ to $\Pp^1\times\Delta$. The rank-drop locus is analytic. Its image is analytic by the proper-image theorem \cite[\S7, Satz 1]{Grauert60}. Away from the central parameter the critical image has exactly three points, labelled by the multiplicity-three fibre, multiplicity-four fibre and cusp. Each traces a nonvertical analytic curve of degree one over the punctured disc. The normalization of its closure gives a meromorphic section of $\Pp^1$ over the disc. The limits may coincide, but the three meromorphic sections remain distinct.

\item The unique cross-ratio transformation taking these sections to $0,1,\infty$ has meromorphic coefficients. Composing with $f'$ gives the meromorphic marked map. Its coefficients have finite-order poles. Resolving its graph gives the asserted holomorphic map; the modifications are proper and preserve the punctured family. This proves the needed extension without making a graph-closure assertion for an arbitrary map from a punctured disc.
\end{enumerate}
\end{proof}

\begin{theorem}\label{night:P2.5}\label{stmt:P2.5}
\StatementRecord{night:P2.5}{proved}{Historical events: GPT-6 Astra (OpenAI), ChatGPT Pro mode, effort Pro, S6\_06\_BLIND, 16--18 November 2026: targeted proof check. Also checked in this writing. Pre-continuation writing events: GPT-6 Astra (OpenAI), ChatGPT Pro mode, effort Pro, S6\_10\_PAPERS, 18 November 2026: targeted proof check. Continuation metadata events: GPT-6 Astra (OpenAI), ChatGPT Pro mode, effort Pro, S6\_10\_PAPERS continuation, 19 November 2026: used and rechecked in this run; status-vocabulary reconciliation; metadata reconciliation only; no new mathematical check. \emph{Pre-continuation recorded proof scope.} Read complete cp25 proofs, re-derived actual origin/deck differences, quotient fundamental-group presentation, finite-orbit vanishing, connected-corner commuting ratio, higher-multiplicity determinant and rational-base recovery, and untwisted plumbing/virtual Euler comparison. Arithmetic chart replay and exact external locators remain separately recorded gates.}{the definitions and exact calculations in this statement}

The global monodromy on $\Lambda_\Q/\Q\widehat\delta$ has no nonzero vector with finite orbit.
\end{theorem}
\begin{proof}
\begin{enumerate}[label=\arabic*.,leftmargin=*]
\item In coordinates $(g,u,w)$, a cusp unipotent has nilpotent part $N=E_{32}$. The two finite homology generators are
\[
 A_1=\begin{pmatrix}1&0&0&0\\6&0&1&0\\-6&-1&-1&0\\-2&1&0&1\end{pmatrix},\qquad
 A_2=\begin{pmatrix}1&0&0&0\\0&0&-1&0\\-6&1&0&0\\3&0&1&1\end{pmatrix}.
\]
A vector with finite orbit is fixed by a positive power of every conjugate of that unipotent. Since $(I+N)^j=I+jN$ in characteristic zero, it belongs to $\ker N$, $\ker(NA_1)$ and $\ker(NA_2)$. The equations are $u=0$, $6g+w=0$, $-w=0$, forcing zero. The exact kernel calculation is in \texttt{part2\_origins\_and\_group.py}.
\end{enumerate}
\end{proof}

\begin{theorem}\label{night:P2.6}\label{stmt:P2.6}
\StatementRecord{night:P2.6}{proved}{Historical events: GPT-6 Astra (OpenAI), ChatGPT Pro mode, effort Pro, S6\_06\_BLIND, 16--18 November 2026: targeted proof check. Also checked in this writing. Pre-continuation writing events: GPT-6 Astra (OpenAI), ChatGPT Pro mode, effort Pro, S6\_10\_PAPERS, 18 November 2026: targeted proof check. Continuation metadata events: GPT-6 Astra (OpenAI), ChatGPT Pro mode, effort Pro, S6\_10\_PAPERS continuation, 19 November 2026: used and rechecked in this run; status-vocabulary reconciliation; metadata reconciliation only; no new mathematical check. \emph{Pre-continuation recorded proof scope.} Read complete cp25 proofs, re-derived actual origin/deck differences, quotient fundamental-group presentation, finite-orbit vanishing, connected-corner commuting ratio, higher-multiplicity determinant and rational-base recovery, and untwisted plumbing/virtual Euler comparison. Arithmetic chart replay and exact external locators remain separately recorded gates.}{P2.5}

In a smooth resolved proper filling carrying the marked base map, every parameter-central component dominating $\Pp^1$ has zero normal-meridian lattice class modulo $\Z\widehat\delta$, relative to the ordinary parameter origin.
\end{theorem}
\begin{proof}
\begin{enumerate}[label=\arabic*.,leftmargin=*]
\item Put $B^\circ=\Pp^1\setminus\{0,1,\infty\}$. After $r=R^n$, the quotient fundamental group of the punctured family over $B^\circ$ by the normal subgroup generated by $\widehat\delta$ is
\begin{equation}\label{eq:P2split}
 \bigl((\Lambda/\Z\widehat\delta)\rtimes\pi_1(B^\circ)\bigr)\times\Z_R.
\end{equation}
The finite Jacobian linear changes are independent of $c$ and fix $\widehat\delta$. The $R$ descent is the $\delta$ translation $H$, all its branch changes are integral $\delta$ periods, and its lattice conjugation differs from the identity by a $\delta$ multiple. The finite gluing changes under $c\mapsto c-12n$ also lie in $\Z\delta$. Thus every cross-cocycle disappears in exactly the quotient \eqref{eq:P2split}; this proves the presentation from the actual chart action.

\item Let $D$ dominate $B$. Its positive normal circle has parameter degree $e=\ord_D(R)>0$, and in \eqref{eq:P2split} is $(\lambda,e)$ with trivial base-loop coordinate. Changing a basing path by a fibre loop changes $\lambda$ by $(I-T_R^e)$ of that loop, which vanishes modulo $\delta$. Transport along $D$ therefore acts on $\lambda$ only by base monodromy.

\item The proper map $D\to B$ has Stein factorization, in the precise analytic form of \cite[Theorem 8.5]{Movasati24}: a proper holomorphic map factors into a map with connected fibres and a finite map. Here $D$ is compact smooth and $B$ a smooth curve, so the hypotheses hold. Remove finitely many branch and critical values. The first map is then a proper submersion with connected fibres; path lifting and a path in a fibre show surjectivity on fundamental groups. The finite curve map has finite-index image in the punctured base group. Removing the other divisors from $D$ preserves this image, since loops in a smooth complex manifold can be moved away from a proper complex analytic subset of real codimension at least two. Hence a finite-index subgroup is represented by paths that transport the normal meridian.

\item The orbit of $\lambda$ is finite. P2.5 kills it rationally, and the quotient lattice is free, so it vanishes integrally. This includes every multiplicity $e$. After a local finite root in the base, each horizontal component inherits this zero class from its global preimage; we do not assume that the local component itself is global.
\end{enumerate}
\end{proof}

\subsection{The commuting-meridian obstruction}
\begin{theorem}\label{night:P2.7}\label{stmt:P2.7}
\StatementRecord{night:P2.7}{proved}{Historical events: GPT-6 Astra (OpenAI), ChatGPT Pro mode, effort Pro, S6\_06\_BLIND, 16--18 November 2026: targeted proof check. Also checked in this writing. Pre-continuation writing events: GPT-6 Astra (OpenAI), ChatGPT Pro mode, effort Pro, S6\_10\_PAPERS, 18 November 2026: targeted proof check. Continuation metadata events: GPT-6 Astra (OpenAI), ChatGPT Pro mode, effort Pro, S6\_10\_PAPERS continuation, 19 November 2026: used and rechecked in this run; status-vocabulary reconciliation; metadata reconciliation only; no new mathematical check. \emph{Pre-continuation recorded proof scope.} Read complete cp25 proofs, re-derived actual origin/deck differences, quotient fundamental-group presentation, finite-orbit vanishing, connected-corner commuting ratio, higher-multiplicity determinant and rational-base recovery, and untwisted plumbing/virtual Euler comparison. Arithmetic chart replay and exact external locators remain separately recorded gates.}{P2.2; P2.4; P2.6}

For every admissible $(3,4)$ twist, the punctured family has no proper Hausdorff complex-analytic filling at $w=0$ after any finite parameter cover. This includes every sphere member and permits arbitrary singularities, multiplicities and extra central components.
\end{theorem}
\begin{proof}
\begin{enumerate}[label=\arabic*.,leftmargin=*]
\item A proposed filling pulls back properly to a further cover $w=R^{6n}$. By P2.4 we resolve its dominating component and marked-map graph. Near infinity take the finite base change $q=Q^{12m}$, normalize the dominating component, and resolve it. Principalize the inverse image of the two axes using \cite[Theorem 1.10, preprint p.~7, with the analytic category in (0.2), p.~3]{BM97}. That theorem turns a coherent ideal on a smooth analytic space into a locally monomial ideal after proper smooth-centre blowups. The axis ideal is coherent and the relevant inverse image is compact over a smaller closed bidisc, so its hypotheses hold. The smooth torus family over $Q^*\times R^*$ is unchanged.

\item The proper map to the bidisc has connected fibres. Its Stein finite factor is generically one sheet because its general fibre is a connected torus. A finite birational map to the normal bidisc is an isomorphism: its finite local algebra is integral over, and contained in the fraction field of, the integrally closed local ring of that bidisc. In particular the corner fibre is connected.

\item In the axis complement, the original logarithmic deck lifts have
\begin{align*}
 A_Q&=I-12mE_{32}+12mE_{41},&
 A_R&=I+12nE_{41},&
 [Q,R]&=-12mnb\widehat\delta.
\end{align*}
The affine relation is the difference $H(s+12m)-H(s)=12mb$ through $n$ parameter turns. The lifts
\[
 Q_o=T_{-mb\widehat\gamma}Q,\qquad R_o=R
\]
commute. They give a specific semidirect product $\Lambda\rtimes\Z^2$; central $\delta$ corrections do not affect the following coefficients.

\item For every irreducible axis component $D_i$, put
\[
 \alpha_i=\ord_{D_i}(Q),\quad\beta_i=\ord_{D_i}(R),\quad
 \kappa_i=\text{the }\widehat\gamma\text{ coefficient of its lattice meridian}.
\]
Its positive meridian is $(\lambda_i,\alpha_i,\beta_i)$, with nonnegative orders and positive sum. Conjugation changes its lattice vector by a vector killed by the first-coordinate covector; thus $\kappa_i$ is independent of its basing path in the complement. Meeting components have commuting meridians, represented by coordinate circles in a smooth NC chart. The $\widehat\delta$ coordinate of their commutator gives
\begin{equation}\label{eq:P2ratio}
 (m\alpha_i+n\beta_i)\kappa_j=(m\alpha_j+n\beta_j)\kappa_i.
\end{equation}
Both nilpotents have square zero and their products vanish, so there are no higher terms. The $E_{32}$ contribution lies in the $\widehat w$ coordinate; $E_{41}$ gives exactly \eqref{eq:P2ratio}. The ratio $\kappa_i/(m\alpha_i+n\beta_i)$ is therefore constant along every intersection edge. Its denominator is positive for every axis component, including vertical exceptional components.

\item A parameter-central component dominating the original base exists, since the proper map onto $\Pp^1\times\Delta$ is surjective and finitely many vertical components cannot cover its central base. Its strict transform meets the corner by properness. Its ratio is zero by P2.6.

\item The strict cusp component met by $\sigma$ also meets the corner. Its meridian in the smooth open orbit is the loop traced by the extending section. P2.2 gives $\widehat\gamma$ coefficient $-12ml_0$ in the original coordinates. Relative to $Q_o$ this becomes $-12ml_0+mb=ma$. Its orders are $(1,0)$, so its ratio is $a$. Further ramification multiplies the order and coefficient by the same positive factor, and proper transforms preserve a smooth open part of the branch.

\item The compact connected corner fibre is covered by its intersections with finitely many axis components. The incidence graph of this closed cover is connected: a partition with no edge would partition the fibre into disjoint nonempty closed subsets, both also open. A chain therefore joins the horizontal central component to the strict cusp component. Equation \eqref{eq:P2ratio} forces $a=0$.

\item Admissibility gives $3\nmid l_1$ and $l_2$ odd, whence $p$ is coprime to $12$ and nonzero. Since $a=-p$, this is a contradiction. Each comparison was proper and preserved the punctured family, so it applies to the original filling and every finite parameter cover.
\end{enumerate}
\end{proof}

\begin{theorem}\label{night:P2.8}\label{stmt:P2.8}
\StatementRecord{night:P2.8}{proved}{Historical events: GPT-6 Astra (OpenAI), ChatGPT Pro mode, effort Pro, S6\_06\_BLIND, 16--18 November 2026: targeted proof check. Also checked in this writing. Pre-continuation writing events: GPT-6 Astra (OpenAI), ChatGPT Pro mode, effort Pro, S6\_10\_PAPERS, 18 November 2026: targeted proof check. Continuation metadata events: GPT-6 Astra (OpenAI), ChatGPT Pro mode, effort Pro, S6\_10\_PAPERS continuation, 19 November 2026: used and rechecked in this run; status-vocabulary reconciliation; metadata reconciliation only; no new mathematical check. \emph{Pre-continuation recorded proof scope.} Read complete cp25 proofs, re-derived actual origin/deck differences, quotient fundamental-group presentation, finite-orbit vanishing, connected-corner commuting ratio, higher-multiplicity determinant and rational-base recovery, and untwisted plumbing/virtual Euler comparison. Arithmetic chart replay and exact external locators remain separately recorded gates.}{P2.1; P2.3; P2.7}

For $(l_0,l_1,l_2)=(0,2,-1)$ we have $p=-5$, $a=b=5$, component cover $y^{12}=t^8/(t-1)^3$, and rotations $8,-3,-5$ modulo $12$. The descent and origins are
\begin{align*}
 H&=5s+U,& U&=\frac{3\log(1-q)}{2\pi i},\\
 o&=\binom{-5\mu s/2}{5s^2/24-5Cs/12+v(5s+U)},\\
 \sigma&=\binom{30\mu v}{vU+5Cv-30v^2+30v},& F&=5F_0,
\end{align*}
where $F_0$ is the vector in \eqref{eq:P2F} without $a$. Its two differences are $5m\gamma-30m^2\delta$ and $5n\gamma-30n(n-1)\delta$. No finite parameter cover admits a proper filling.
\end{theorem}
\begin{proof}
\begin{enumerate}[label=\arabic*.,leftmargin=*]
\item Substitution in P2.1--P2.3 gives every displayed formula. The lattice and character laws remain those of P2.1 and the charts those of \eqref{eq:P2chain}. The corner has $[Q,R]=-60mn\widehat\delta$ and $Q_o=T_{-5m\widehat\gamma}Q$. Its horizontal ratio is zero and strict cusp ratio five. Thus P2.7 gives the contradiction for this member itself. The symbolic program checks the substitutions and differences.
\end{enumerate}
\end{proof}

\subsection{Higher coprime multiplicities}
\begin{theorem}\label{night:P2.9}\label{stmt:P2.9}
\StatementRecord{night:P2.9}{proved}{Historical events: GPT-6 Astra (OpenAI), ChatGPT Pro mode, effort Pro, S6\_06\_BLIND, 16--18 November 2026: targeted proof check. Also checked in this writing. Pre-continuation writing events: GPT-6 Astra (OpenAI), ChatGPT Pro mode, effort Pro, S6\_10\_PAPERS, 18 November 2026: targeted proof check. Continuation metadata events: GPT-6 Astra (OpenAI), ChatGPT Pro mode, effort Pro, S6\_10\_PAPERS continuation, 19 November 2026: used and rechecked in this run; status-vocabulary reconciliation; metadata reconciliation only; no new mathematical check. \emph{Pre-continuation recorded proof scope.} Read complete cp25 proofs, re-derived actual origin/deck differences, quotient fundamental-group presentation, finite-orbit vanishing, connected-corner commuting ratio, higher-multiplicity determinant and rational-base recovery, and untwisted plumbing/virtual Euler comparison. Arithmetic chart replay and exact external locators remain separately recorded gates.}{P2.4; P2.6}

Every admissible logarithmic datum in the coprime $(3k,4l)$ family of the supplied Viaclovsky construction has no proper Hausdorff analytic filling at parameter zero after any finite cover.
\end{theorem}
\begin{proof}
\begin{enumerate}[label=\arabic*.,leftmargin=*]
\item Put $m_1=3k$, $m_2=4l$, $N=m_1m_2=12kl$, with $\gcd(m_1,m_2)=1$. Let $l_j$ be the $\widehat\gamma$ coefficient of the invariant finite twist. Admissibility is $\gcd(l_j,m_j)=1$. Retain $l_0$ and define $b=m_2l_1+m_1l_2$, $a=b-Nl_0$. On $w=r^{N/2}$, $c=-Nv$, and the gamma period is $(6\mu,C-s-Nv)$. The descent is
\[
 H=-\frac{m_2l_1\log t+m_1l_2\log(t-1)}{2\pi i}=bs+U,
 \qquad U=-\frac{m_1l_2\log(1-q)}{2\pi i}.
\]
The finite actions are unchanged by $c\mapsto c-N$, since $N/m_j$ is integral. The collar coefficient is $a$; adding $a\gamma$ gives a unit action on cusp charts. The rank-one filling has $N$ components with rotations $m_2l_1,m_1l_2,-b$. Admissibility makes the finite actions free on the cycle.

\item On $q=Q^N$ the origins are
\begin{align*}
 o&=\binom{-6b\mu s/N}{bs^2/(2N)-bCs/N+l_0s/2+v(bs+U)},\\
 \sigma&=-l_0z_\gamma+\binom{6a\mu v}{vU+aCv-Nav^2/2+Nav/2},
\end{align*}
with $z_\gamma=(6s\mu,-s(s+1)/2+s(C-Nv))$. Their difference is
\[
 aF_N=a\binom{6\mu(s+Nv)/N}{C(s+Nv)/N-(s+Nv)^2/(2N)+Nv/2}.
\]
On $q=Q^{Nm}$, $r=R^n$,
\begin{align*}
 \Delta_Q(aF_N)&=am\gamma-\tfrac N2am^2\delta,\\
 \Delta_R(aF_N)&=an\gamma-\tfrac N2an(n-1)\delta.
\end{align*}
All coefficients are integral because $12\mid N$. The exact symbolic program \texttt{part2\_general.py} verifies them.

\item We also prove extension of the intrinsic base map. Pulling the order-three and order-four tangent root charts through $k$ and $l$ powers changes their vertical determinant orders from $2,1$ to $2k,l$. The vertical determinant is therefore $f^*\OO(-1)\otimes\OO(2kS_1+lS_2)$. The determinant of the image of $df$ is $f^*\OO(2)\otimes\OO(-(3k-1)S_1-(4l-1)S_2)$; the cusp singular ideal has codimension two. Hence
\begin{align*}
 K_X^{-1}&=f^*\OO(1)\otimes\OO((1-k)S_1+(1-3l)S_2),\\
 K_X^{-N}&=f^*\OO(d),& d&=4l+3k-kl.
\end{align*}
The integer $d$ is nonzero. Otherwise $k\mid4l$, while $\gcd(k,4l)=1$, so $k=1$, giving $d=3l+3>0$.

\item Choose $K^{-N}$ when $d>0$ and $K^N$ when $d<0$; its restriction is $f^*\OO(|d|)$. On a resolved hypothetical filling the corresponding relative direct image is coherent torsion free of rank $|d|+1$. Evaluation and graph resolution give a proper pluricanonical image. Over the field $K$ of meromorphic parameter germs, its image is a smooth rational normal curve. More explicitly, the kernels of multiplication on the direct-image modules, tensored with $K$, give its quadratic equations; on the punctured disc these are the Veronese equations and generate its ideal. The three labelled critical-value curves give distinct $K$-rational points. The coprime multiplicities are unequal, and the third point is the cusp.

\item Over an algebraic closure of $K$ there is a unique rational parameter taking those points to $0,1,\infty$. Uniqueness makes it invariant under every field automorphism over $K$, so its coefficients belong to $K$. They are meromorphic germs with finite principal parts. Composing and resolving the graph yields the marked base map required in P2.4. This explicitly recovers a degree-one base parameter.

\item Modulo $\delta$, the global lattice and finite-orbit calculation are unchanged. The cross-cocycles still lie in $\Z\delta$, so P2.6 applies. At the corner,
\[
 A_Q=I-NmE_{32}+NmE_{41},\quad A_R=I+NnE_{41},\quad [Q,R]=-Nmnb\widehat\delta.
\]
The commuting lifts $Q_o=T_{-mb\widehat\gamma}Q$, $R_o=R$ give \eqref{eq:P2ratio}. Connectedness forces equality of the horizontal ratio zero and cusp ratio $a$. But $a\equiv m_2l_1\pmod{m_1}$ and $a\equiv m_1l_2\pmod{m_2}$, so $\gcd(a,N)=1$. This contradiction proves nonexistence. Proper pullback to the further cover preserves the punctured family and gives the implication back to every proposed cover.
\end{enumerate}
\end{proof}

\subsection{The untwisted control}
\begin{theorem}\label{night:P2.10}\label{stmt:P2.10}
\StatementRecord{night:P2.10}{proved}{Historical events: GPT-6 Astra (OpenAI), ChatGPT Pro mode, effort Pro, S6\_06\_BLIND, 16--18 November 2026: targeted proof check. Also checked in this writing. Pre-continuation writing events: GPT-6 Astra (OpenAI), ChatGPT Pro mode, effort Pro, S6\_10\_PAPERS, 18 November 2026: targeted proof check. Continuation metadata events: GPT-6 Astra (OpenAI), ChatGPT Pro mode, effort Pro, S6\_10\_PAPERS continuation, 19 November 2026: used and rechecked in this run; status-vocabulary reconciliation; metadata reconciliation only; no new mathematical check. \emph{Pre-continuation recorded proof scope.} Read complete cp25 proofs, re-derived actual origin/deck differences, quotient fundamental-group presentation, finite-orbit vanishing, connected-corner commuting ratio, higher-multiplicity determinant and rational-base recovery, and untwisted plumbing/virtual Euler comparison. Arithmetic chart replay and exact external locators remain separately recorded gates.}{the definitions and exact calculations in this statement}

The untwisted family $X(u)$ has a proper flat Hausdorff analytic filling at $u=0$. Its corner winding is zero. The filling is obtained from two projective-line bundles by explicit plumbing, a proper modification and a finite quotient.
\end{theorem}
\begin{proof}
\begin{enumerate}[label=\arabic*.,leftmargin=*]
\item Let $S$ be the supplied rational elliptic surface, $M=\OO_S(P-O)$, $h=t_{6P}$, and $L=\Hom(h^*M,M)=\OO_S(F_\infty)$, with section $\phi$ vanishing on $F_\infty$. The absence of logarithmic twists gives $b=a=0$ in the corner character calculation. Thus the origins agree and their winding is zero. We now construct the filling with the actual untwisted finite fibres.

\item Take two copies $Y_0,Y_1$ of $\Pp_S(\OO\oplus M)$. The finite coordinate $z_i$ lies in $M$, its inverse $w_i$ in $M^{-1}$. The first neck, over $p\in S$, joins zero in $Y_0$ to infinity in $Y_1$ over $h(p)$ by $z_0w_1=r\phi_p$. The second joins zero in $Y_1$ to infinity in $Y_0$ over the same point by $z_1w_0=r$. These equations are typed in $M\otimes h^*M^{-1}=L$ and $M\otimes M^{-1}=\OO$.

\item Choose a Hermitian metric and $0<\epsilon<1/4$. Use middle annuli $\epsilon/2<|z_i|<2/\epsilon$ in $Y_i$, times the parameter disc, and neck hypersurfaces with $|z|,|w|<\epsilon$. Glue their ends on $\epsilon/2<|z|<\epsilon$ or $\epsilon/2<|w|<\epsilon$ by coordinate projection, using $h$ at the first infinity end. In holomorphic frames these are holomorphic maps, and they commute with bundle transitions because the product lines are exactly the ones displayed above. Choose $\rho$ so that $\rho\max(1,\|\phi\|)<\epsilon^2/16$. The product equation makes the two ends of a neck disjoint; the two ends in each middle chart are disjoint as well.

\item For $|r|\le\rho'<\rho$, a compact fundamental cover consists of closed middle annuli $3\epsilon/4\le|z_i|\le4/(3\epsilon)$ and closed neck pieces $|z|,|w|\le3\epsilon/4$. Every omitted neck end is represented in a middle annulus. The gluing relation on this finite union of compact bundles is the union of diagonals and closed graphs of the coordinate identifications. No additional graph accumulates, because the two neck ends are disjoint. The relation is closed, so its quotient is compact and Hausdorff. The open charts give the analytic structure. This proves properness over the smaller disc and the global Hausdorff property.

\item Over $S\setminus F_\infty$, traversing the two necks gives
\[
 (p,z_0)\longmapsto\left(h(p),\frac{z_0}{r^2\phi_p}\right).
\]
The inverse covers $h^{-1}$ and is multiplied by $u\phi$, with $u=r^2$. It is the precise generator of Engel's quotient on that open set.

\item At a smooth point of $F_\infty$ the first neck is $zw=rq$; at its node it is $zw=rxy$, with $q=xy$. The ideal of one globally distinguished central component is locally $(z,r)$. Because the two global components are distinct, this is a global coherent ideal. Blow it up. At the node its two charts are
\begin{align*}
 r=za,\quad w=axy;&\qquad (z,a,x,y)\text{ are smooth coordinates},\\
 z=rb,\quad bw=xy;&\qquad (r,b,w,x,y)\text{ with the displayed relation}.
\end{align*}
The remaining singularity is an ordinary threefold node times the parameter disc. Its singular section does not meet the first chart. Resolve $(b,x)$ in the second chart. The two charts are
\begin{align*}
 b=xk,\quad y=kw;&\qquad (r,x,k,w),\\
 x=bk,\quad w=ky;&\qquad (r,b,k,y).
\end{align*}
They are smooth, with overlap $k\mapsto k^{-1}$. Off the singular section the ideal is invertible, so the proper small resolution glues to the identity. The analogous blowup charts at smooth points of $F_\infty$ are already smooth. The resulting total space is smooth, and its parameter central divisor is reduced NC: in the first chart $r=za$, in the small-resolution charts $r$ is a coordinate. Parameter torsion freeness proves flatness. The substitutions are checked exactly in \texttt{part2\_plumbing\_charts.py}.

\item We identify the compactified generic fibre, including the cusp. For fixed $r\ne0$ the first blowup is an isomorphism. After multiplying by a unit the first neck is $zw=u xy$. On its dense torus use $(x,z,q)$, with $y=q/x$, $w=uq/z$. The four affine monomials have exponents
\[
 (1,0,0),\quad(-1,0,1),\quad(0,1,0),\quad(0,-1,1).
\]
The dual cone is $0\le a,b\le c$, the cone over the unit square at height one. The resolution of $(z,x)$, equal to that of $(b,x)$ when $r$ is a unit, divides it along $a=b$. The two cones over the main-diagonal triangles are unimodular. In both resolved charts $q$ is the product of three coordinates. The Tate and linearization periods shift the two lattice directions by one, while their unit factors act holomorphically on the charts. Thus this is exactly the periodic triangular Mumford filling in Engel's Proposition 2.8. Choosing $x$ in the Tate direction fixes the diagonal and the comparison. Together with the quotient comparison this identifies the entire compact generic fibre with $X(u)$.

\item The involution $r\mapsto-r$ fixes the $Y_0$ fibre coordinate and negates the $Y_1$ coordinate. It preserves both equations and the global component ideal. In the nodal blowup chart it acts as $(b,w,x,y,r)\mapsto(-b,-w,x,y,-r)$, preserving $(b,x)$ and lifting holomorphically to the small resolution. Its finite quotient is proper and Hausdorff over $u=r^2$. Its invariant local rings are torsion free over the discrete valuation ring $\C\{u\}$ and hence flat. Away from zero it exchanges the two roots of $u$; the comparison on $Y_0$ is the identity, so the quotient fibre is exactly one copy of $X(u)$. This is a filling in the original parameter.
\end{enumerate}
\end{proof}

\begin{theorem}\label{night:P2.11}\label{stmt:P2.11}
\StatementRecord{night:P2.11}{false as stated}{Historical events: GPT-6 Astra (OpenAI), ChatGPT Pro mode, effort Pro, S6\_06\_BLIND, 16--18 November 2026: targeted proof check. Also checked in this writing. Pre-continuation writing events: GPT-6 Astra (OpenAI), ChatGPT Pro mode, effort Pro, S6\_10\_PAPERS, 18 November 2026: targeted proof check. Continuation metadata events: GPT-6 Astra (OpenAI), ChatGPT Pro mode, effort Pro, S6\_10\_PAPERS continuation, 19 November 2026: used and rechecked in this run; status-vocabulary reconciliation; metadata reconciliation only; no new mathematical check. \emph{Pre-continuation recorded proof scope.} Read complete cp25 proofs, re-derived actual origin/deck differences, quotient fundamental-group presentation, finite-orbit vanishing, connected-corner commuting ratio, higher-multiplicity determinant and rational-base recovery, and untwisted plumbing/virtual Euler comparison. Arithmetic chart replay and exact external locators remain separately recorded gates.}{P2.10}

The one-component central space specified in Engel's Conjecture 8.2 cannot be the central fibre of a proper flat smoothing with smooth fibres $X(u)$. Its intrinsic virtual Euler number is zero, whereas $e(X(u))=2$.
\end{theorem}
\begin{proof}
\begin{enumerate}[label=\arabic*.,leftmargin=*]
\item Let $Q$ denote that proposed space. Its normalization is $Y=\Pp_S(\OO\oplus M)$ and the two sections are identified by $h$. Locally $Q$ is the lci hypersurface $xy=0$. The double locus is $S$ and its first-order smoothing line is $L=M\otimes h^*M^{-1}=\OO(F_\infty)$.

\item We compute its virtual tangent Chern number on the normalization. Along one labelled branch let $a$ be the normal Chern class of the double surface in that branch, and $b$ that of the other branch. A local smooth hypersurface presentation gives
\[
 c(T_Q^{\vir})|_Y=c(T_Y)\frac{1+B}{1+B+D},\qquad B|_D=b,\quad D|_D=a.
\]
The degree-three correction, localized to $D=S$, is
\[
 -c_2(S)+c_1(S)b-ab-b^2.
\]
Indeed expanding the fraction gives the terms $-Dc_2(Y)+D(B+D)c_1(Y)-D(B+D)^2$, and $c(T_Y)|_D=c(T_S)(1+a)$ yields the expression. Adding the other branch gives $-2c_2(S)+c_1(S)(a+b)-(a+b)^2$. It is symmetric and $a+b=c_1(L)$, so it descends through the actual normalization gluing. Consequently
\[
 e_{\vir}(Q)=e(Y)-2e(S)+\int_S\bigl(c_1(S)c_1(L)-c_1(L)^2\bigr).
\]
Since $e(S)=12$, $e(Y)=24$, $c_1(S)=F_\infty$ and $F_\infty^2=0$, this is zero.

\item We verify that this intrinsic value applies to an arbitrary flat smoothing. Embed a local total-space germ in a smooth polydisc over the parameter disc, with ideal $I$ and quotient $A$. Flatness gives $\operatorname{Tor}_1^{\C\{r\}}(A,\C)=0$, identifying $I/rI$ with the central ideal. Lift a regular sequence generating that central ideal. Nakayama's lemma gives generators of $I$; the sequence consisting of $r$ and these lifts is regular in the ambient regular local ring, so the lifted equations define a relative lci morphism. Its perfect relative cotangent complex restricts to the intrinsic complex of $Q$. The virtual tangent Chern class therefore specializes to that of $Q$. Compact flat fibre cycles have the same oriented local regular-sequence specialization class: their difference bounds the oriented chain over a parameter interval, constructed in the local presentations and glued by their complex orientations. The degree-three virtual pairing is constant. On a smooth fibre it equals the ordinary tangent Euler number. Thus every smooth fibre of such a smoothing would have Euler number zero.

\item For the actual $X(u)$, the $M^\times$ scaling circle is free over all finite base strata, whose Euler contributions vanish. The index-one triangular cusp has two zero-dimensional torus strata; all positive-dimensional torus strata have Euler number zero. Euler additivity gives $e(X(u))=2$. This contradiction excludes the proposed central space. The explicit proper flat filling in P2.10 has a different central space.
\end{enumerate}
\end{proof}
\section{The cotangent direct images and the decisive period scalar}
\label{inv:sec:partA}

\subsection{Local logarithmic bundles and a duality calculation}

We use the marked threefold, with its arbitrary admissible
constant $c$.  Coefficients of coherent cohomology are complex.  In this
section $\omega_f=K_X\otimes f^*K_B^{-1}$,
$\mathcal L=\Omega_B^1(\log\Sigma)=\OO_B(1)$,
$\mathcal A=\Omega_X^1(\log D)$, and
$\mathcal E=\mathcal A/f^*\mathcal L$.  We write $\gamma,u$ also for
their images in $R^1f_*\OO_X$; this convention concerns cohomology
classes and does not identify them with differential-form representatives.

\begin{cited}[The structure-sheaf input in the precise form used below]\label{inv:A1.0}\label{wb:thm:structure-direct-images-audit}\label{wb:lem:cusp-structure-sheaf-audit}
\StatementRecord{inv:A1.0}{proved}{Historical events: Claude Fable 5.2 (Anthropic), effort high, 12 November 2026: targeted proof check. Claude Fable 5.2 (Anthropic), effort high, 16 November 2026: carried as background without fault found; used and rechecked in this run. Claude Fable 5.2 (Anthropic), effort high, 17 November 2026: carried as background without fault found. Claude Fable 5.2 (Anthropic), effort max, 16 November 2026 14:25 UTC: carried as background without fault found. GPT-6 Astra (OpenAI), ChatGPT Pro mode, effort Pro, S6\_09\_PAPERS, 17 November 2026: targeted proof check. Also checked in this writing. Pre-continuation writing events: GPT-6 Astra (OpenAI), ChatGPT Pro mode, effort Pro, S6\_10\_PAPERS, 18 November 2026: used and rechecked in this run. Continuation metadata events: GPT-6 Astra (OpenAI), ChatGPT Pro mode, effort Pro, S6\_10\_PAPERS continuation, 19 November 2026: used and rechecked in this run; status-vocabulary reconciliation; metadata reconciliation only; no new mathematical check. \emph{Pre-continuation recorded proof scope.} Read AC26 original PDF pp72–76 and p92: matched the entire structure-sheaf statement to Lemma9.12 and Proposition9.13, including cyclic exponents, gamma subbundle, cusp cup pairing and full scheme-fibre base-change maps. Its expanded re-derivation remains a compendium proof of a known input.}{\cite[Lemma~9.12, pp.~72--74; Proposition~9.13, pp.~74--76; p.~92]{AC2026}}

This is the structure-sheaf calculation of Claude and Alp\"oge: \cite[Lemma~9.12, pp.~72--74; Proposition~9.13, pp.~74--76, and the base-change clarification on p.~92]{AC2026}. The direct images $R^qf_*\OO_X$, for $q=0,1,2$, are
$\OO$, $G=\OO\gamma\oplus\OO(-1)$, and $\OO(-1)$.
The inclusion $\OO\gamma\subset G$ is a subbundle.  At the cusp,
$(\gamma,u)$ is a holomorphic frame, and cup product identifies
$\bigwedge^2G$ with $R^2f_*\OO_X$.  In the cyclic covering coordinate
$s$, the action on $G$ has characters $(1,\zeta^{-a})$, where
$(m,a)=(3,2)$ or $(4,1)$ and $\zeta=e^{-2\pi i/m}$.
These assertions concern full scheme fibres, including $3S_1$ and $4S_2$.
At the cusp, the natural map $H^1(W;\mathbb C)\to H^1(W,\mathcal O_W)$ is an isomorphism. In degree two the comparison is surjective and is injective on $\ker[H^2(W;\mathbb C)\to H^2(\mathrm{dP}_6;\mathbb C)]$. The cup product $\bigwedge^2H^1(W,\mathcal O_W)\to H^2(W,\mathcal O_W)$ is an isomorphism. Every base-change map $(R^qf_*\mathcal O_X)\otimes k(b)\to H^q(X_b,\mathcal O_{X_b})$ is an isomorphism, including at the nonreduced fibres.
\end{cited}


The expanded source derivation is retained in the proof compendium, Paper~I, label \nolinkurl{inv:A1.0}.

\Uses{the hypotheses, notation and definitions stated here; \citet[Lemma~9.12 and Proposition~9.13]{AC2026}; \CompendiumRef{I}{app:rederive:inv-A1.0}}


The structure-sheaf and cusp statements are due to \citet[Lemma~9.12 and Proposition~9.13]{AC2026}; we retain the independent root-line derivation. A complete re-derivation is given in \CompendiumRef{I}{app:rederive:inv-A1.0}.


\begin{proposition}[Local freeness and the determinant line]\label{inv:A1.1}
\StatementRecord{inv:A1.1}{proved}{Historical events: Claude Fable 5.2 (Anthropic), effort high, 16 November 2026: carried as background without fault found; used and rechecked in this run. Claude Fable 5.2 (Anthropic), effort high, 17 November 2026: carried as background without fault found. Claude Fable 5.2 (Anthropic), effort max, 16 November 2026 14:25 UTC: carried as background without fault found. GPT-6 Astra (OpenAI), ChatGPT Pro mode, effort Pro, S6\_09\_PAPERS, 17 November 2026: targeted proof check. Also checked in this writing. Pre-continuation writing events: GPT-6 Astra (OpenAI), ChatGPT Pro mode, effort Pro, S6\_10\_PAPERS, 18 November 2026: targeted proof check. Continuation metadata events: GPT-6 Astra (OpenAI), ChatGPT Pro mode, effort Pro, S6\_10\_PAPERS continuation, 19 November 2026: used and rechecked in this run; status-vocabulary reconciliation; metadata reconciliation only; no new mathematical check. \emph{Pre-continuation recorded proof scope.} Read full source proofs and re-derived logarithmic module, root-character/duality and residue cancellation maps, modular determinant and selected scalar, elliptic-section comparison, Leray table/lifts, and ordinary-relative differential syzygies. Opened AC original pp5,80–81 to verify precise scope and correction wording; external classical citation locators under separate review.}{the hypotheses, notation and definitions stated here; the explicit constructions and calculations in the proof}

Both $\mathcal A$ and $\mathcal E$ are locally free, of ranks three and
two.  The vertical vector $\xi=\partial/\partial\zeta_2$ is nowhere
zero as a section of $\mathcal E^\vee$.  There is an exact sequence
of holomorphic vector bundles
\begin{equation}\label{inv:eq:A-E-line}
0\longrightarrow N\longrightarrow\mathcal E
 \xrightarrow{\langle\cdot,\xi\rangle}\OO_X\longrightarrow0,
\qquad N=\det\mathcal E=\omega_f(-2S_1-3S_2).
\end{equation}
\end{proposition}


\begin{proof}
\begin{enumerate}[label=\arabic*.,leftmargin=*]
\item At a smooth point of a reduced cusp branch the assertion is the usual
coordinate calculation.  At a double point write $t_c=xy$ with third
coordinate $z$; then
$\mathcal A=R\,d\log x\oplus R\,d\log y\oplus R\,dz$ and its
base subline is $R(d\log x+d\log y)$.  At either triple point write
$t_c=xyz$; then $\mathcal A$ has the three logarithmic coordinate
basis vectors and its base subline is their sum.  In each case that
subline is a direct summand.  These calculations apply on every one
of the three double curves, including their endpoints; no simplicity
of the global divisor is assumed.

\item At $S_j$ we use the actual free cyclic quotient
$q:J'_j\to N_j$, which is finite and etale on total spaces, and
$t_j=s^{m_j}$ on the base.  On $J'_j$ the logarithmic conormal is
$m_j\,ds/s$, and its coefficient is a unit over $\C$.  It is again a
direct summand of the logarithmic cotangent bundle.  The comparison
$q^*\mathcal A=\Omega_{J'_j}^1(\log(s=0))$ identifies the quotient
bundles and their exact sequence; faithful flat descent proves the
claim downstairs.  This is not a replacement of $m_jS_j$ by $S_j$.

\item On the marked torus family $\xi$ is the constant second vertical
vector.  The linear monodromies fix it, and translations preserve it.
On a cusp toric chart it is a nonzero constant vector in the relative
logarithmic tangent lattice of rank two.  A toric coordinate change
$x_i\mapsto c_i(t_c)t_c^{b_i}x_i$ changes $d\log x_i$ by a multiple
of $d\log t_c$, so this constant relative vector descends and is
nowhere zero in that bundle.  It can vanish as an ordinary tangent
vector at a toric stratum; this does not imply vanishing in the
logarithmic bundle.  Hence evaluation is surjective and its kernel
is the determinant line.  Taking determinants of
$0\to f^*\mathcal L\to\mathcal A\to\mathcal E\to0$ gives
\[
\det\mathcal E=K_X(D)\otimes f^*\mathcal L^{-1}
=\omega_f\bigl(D-f^*\Sigma\bigr)
=\omega_f(-2S_1-3S_2),
\]
which proves \eqref{inv:eq:A-E-line}.
\end{enumerate}
\end{proof}

\begin{proposition}[The logarithmic direct images by vertical twisting and duality]\label{inv:A1.2}\label{wb:prop:log-relative-direct-images}
\StatementRecord{inv:A1.2}{proved}{Historical events: Claude Fable 5.2 (Anthropic), effort high, 12 November 2026: targeted proof check. Claude Fable 5.2 (Anthropic), effort high, 16 November 2026: carried as background without fault found; used and rechecked in this run. Claude Fable 5.2 (Anthropic), effort max, 16 November 2026 14:25 UTC: carried as background without fault found. GPT-6 Astra (OpenAI), ChatGPT Pro mode, effort Pro, S6\_09\_PAPERS, 17 November 2026: targeted proof check. Also checked in this writing. Pre-continuation writing events: GPT-6 Astra (OpenAI), ChatGPT Pro mode, effort Pro, S6\_10\_PAPERS, 18 November 2026: targeted proof check. Continuation metadata events: GPT-6 Astra (OpenAI), ChatGPT Pro mode, effort Pro, S6\_10\_PAPERS continuation, 19 November 2026: used and rechecked in this run; status-vocabulary reconciliation; metadata reconciliation only; no new mathematical check. \emph{Pre-continuation recorded proof scope.} Read full source proofs and re-derived logarithmic module, root-character/duality and residue cancellation maps, modular determinant and selected scalar, elliptic-section comparison, Leray table/lifts, and ordinary-relative differential syzygies. Opened AC original pp5,80–81 to verify precise scope and correction wording; external classical citation locators under separate review.}{the hypotheses, notation and definitions stated here; \cref{inv:eq:A-H-images}; \cref{inv:eq:A-E-line}; \cref{inv:eq:A-N-images}; \cref{inv:A1.0}}

One has, as holomorphic vector bundles on $B$,
\[
\begin{array}{c|ccc}
q&0&1&2\\\hline
R^qf_*\mathcal E&\OO\oplus\OO(-1)&
 \OO^2\oplus\OO(-1)^2&\OO\oplus\OO(-1).
\end{array}
\]
They commute with base change to every full scheme fibre.
\end{proposition}


\begin{proof}
\begin{enumerate}[label=\arabic*.,leftmargin=*]
\item Put $H=2S_1+3S_2$.  We first compute the line bundle $\OO_X(H)$,
rather than a tensor product of two separately descended lattices.
On the cover at $p_j$, its sections allow poles of order at most $m_j-1$
in $s$.  A fixed character has invariant exponent zero and gains no
pole, whereas the variable character $\zeta^{-a}$ has least allowed
invariant exponent $a-m$.  Thus the inclusion
$R^1f_*\OO\to R^1f_*\OO(H)$ is the identity on $\OO\gamma$ and
changes $s^ab_1$ to $s^{a-m}b_1=t_j^{-1}s^ab_1$ on its quotient.
It makes exactly one upper modification at each finite point.
The same computation applies to $R^2$.  No modification is made at
the cusp or a regular fibre.  The result is
\begin{equation}\label{inv:eq:A-H-images}
f_*\OO(H)=\OO,\qquad
R^1f_*\OO(H)=\OO\oplus\OO(1),\qquad
R^2f_*\OO(H)=\OO(1).
\end{equation}
Here the middle extension splits because $H^1(\OO(-1))=0$.
The local invariant modules are free.  Tensoring them with
$\C\{t_j\}/(t_j)$ gives invariants of
$s^{-(m_j-1)}\C\{s\}/(s^m)$, with the indicated surviving character
basis vectors.  The Reynolds idempotent is linear over the base,
so this calculation also proves the comparison with the nonreduced
fibre.  At the cusp the line bundle is trivial.  Cohomology and base
change, applied successively from the top degree, proves the claimed
local freeness and comparison for \eqref{inv:eq:A-H-images}.

\item For clarity the relative duality map we use is the cup product followed
by the analytic trace, giving a morphism of bundles
\[
R^qf_*(\omega_f\otimes\OO(-H))\longrightarrow
\bigl(R^{2-q}f_*\OO(H)\bigr)^\vee.
\]
Every fibre is a compact Gorenstein surface: it is the Cartier fibre
in a smooth threefold.  Serre duality on that scheme makes the map
an isomorphism on its residue field, also for $m_jS_j$.  The constant
dimensions just computed imply local freeness on both sides, so
Nakayama's lemma makes the bundle map an isomorphism.  Therefore
\begin{equation}\label{inv:eq:A-N-images}
(f_*N,R^1f_*N,R^2f_*N)
=(\OO(-1),\OO\oplus\OO(-1),\OO).
\end{equation}
This is the explicit reduction of analytic relative duality needed
here; the scheme fibre is retained throughout.

\item The connecting maps of \eqref{inv:eq:A-E-line} are zero on a regular
fibre: its relative cotangent is a trivial bundle and evaluation on
a constant nonzero vector splits.  Source and target direct images
in \eqref{inv:eq:A-N-images} and Statement \ref{inv:A1.0} are vector bundles,
so these holomorphic connecting maps, zero on a dense open set, are
zero everywhere.  We obtain extensions in each degree of the listed
$R^qN$ by $R^q\OO$.  Every resulting pair of line-bundle degrees is
zero or minus one; their extension groups are
$H^1(\OO(-1)),H^1(\OO),H^1(\OO(1))$, all zero.  These extensions
give exactly the three asserted bundles.  Applying the already proved
fibre comparisons to the short exact sequences proves base change
for $R^q\mathcal E$ as well.  This argument uses the vertical twist
$\OO(H)$ and duality; it does not use the workbench's filtration of
$f_*\mathcal E\otimes G$.
\end{enumerate}
\end{proof}

\begin{proposition}[All absolute cotangent direct images]\label{inv:A1.3}\label{wb:prop:absolute-cotangent-direct-images}
\StatementRecord{inv:A1.3}{proved}{Historical events: Claude Fable 5.2 (Anthropic), effort high, 12 November 2026: targeted proof check. Claude Fable 5.2 (Anthropic), effort high, 16 November 2026: used and rechecked in this run. Claude Fable 5.2 (Anthropic), effort high, 16 November 2026 13:58 UTC: targeted proof check. Claude Fable 5.2 (Anthropic), effort max, 16 November 2026 14:25 UTC: carried as background without fault found. GPT-6 Astra (OpenAI), ChatGPT Pro mode, effort Pro, S6\_09\_PAPERS, 17 November 2026: targeted proof check. Also checked in this writing. Pre-continuation writing events: GPT-6 Astra (OpenAI), ChatGPT Pro mode, effort Pro, S6\_10\_PAPERS, 18 November 2026: targeted proof check. Continuation metadata events: GPT-6 Astra (OpenAI), ChatGPT Pro mode, effort Pro, S6\_10\_PAPERS continuation, 19 November 2026: used and rechecked in this run; status-vocabulary reconciliation; metadata reconciliation only; no new mathematical check. \emph{Pre-continuation recorded proof scope.} Read full source proofs and re-derived logarithmic module, root-character/duality and residue cancellation maps, modular determinant and selected scalar, elliptic-section comparison, Leray table/lifts, and ordinary-relative differential syzygies. Opened AC original pp5,80–81 to verify precise scope and correction wording; external classical citation locators under separate review.}{the hypotheses, notation and definitions stated here; \cref{inv:A2.1}; \cref{inv:A2.2}; \cref{inv:A2.4}; \cref{inv:A1.1}}

For every admissible $c$, the six direct images are
\begin{equation}\label{inv:eq:A-absolute-table}
\begin{array}{c|cc}
q&R^qf_*\Omega_X^1&R^qf_*\Omega_X^2\\\hline
0&\OO(-2)&\OO(-2)\oplus\OO(-1)\\
1&\OO\oplus\OO(-1)^2&\OO(-2)\oplus\OO(-1)^2\\
2&\OO\oplus\OO(-1)&\OO.
\end{array}
\end{equation}
In particular every listed sheaf is torsion-free.
\end{proposition}


\begin{proof}
\begin{enumerate}[label=\arabic*.,leftmargin=*]
\item We use Statements \ref{inv:A2.1}, \ref{inv:A2.2} and \ref{inv:A2.4} below, whose proofs do not use
\eqref{inv:eq:A-absolute-table}.  They give
$\coker\kappa_0=\C_{p_1}\oplus\C_{p_2}$,
surjectivity of $\kappa_1$, and rank one on its global sections.
Write $\mathcal K=\ker\kappa_1$.  The logarithmic conormal sequence
then gives
\[
f_*\mathcal A=\OO(1),\quad
0\to T\to R^1f_*\mathcal A\to\mathcal K\to0,\quad
R^2f_*\mathcal A=\OO\oplus\OO(-1),
\quad T=\C_{p_1}\oplus\C_{p_2}.
\]
Let $\nu:\widetilde D=\mathrm{dP}_6\sqcup S_1\sqcup S_2\to D$
be normalization.  The local bases in Statement \ref{inv:A1.1} give the exact
residue sequence
\[
0\to\Omega_X^1\to\mathcal A\to\nu_*\OO_{\widetilde D}\to0.
\]
Its three residue direct images are
$\C_{p_0}\oplus\C_{p_1}\oplus\C_{p_2}$,
$\C_{p_1}\oplus\C_{p_2}$, and zero: the normalization of $W$ is
rational, and a bielliptic surface has $(h^0,h^1,h^2)(\OO)=(1,1,0)$.
In degree zero the map from $\OO(1)$ is local evaluation with the
nonzero multipliers $1,3,4$.  Its kernel is $\OO(-2)$.

\item We check the torsion cancellation, rather than discarding it.  At
$p_j$, in the invariant frame $(b_0,s^ab_1)$ of $G$, residue of
$(dt_j/t_j)(x b_0+y s^ab_1)$ in $H^1(S_j,\OO)$ is $m_jx[b_0]$.
The residue factors through the base logarithmic line on the covering
chart.  Naturality of the connecting map consequently gives
$\operatorname{res}_1\kappa_0=0$.  The determinant of $\kappa_0$
has a simple zero, so its reduction has rank one, and its image is
exactly the one-dimensional kernel $\C s^ab_1$ of this residue.
Hence $\coker\kappa_0\to H^1(S_j,\OO)$ is an isomorphism.
Both torsion summands cancel, giving the canonical identification
$R^1f_*\Omega_X^1=\mathcal K$ and the asserted $R^2$.
The rank-one global map
$\OO^2\oplus\OO(-1)^2\to\OO$ has a nonzero constant component
on a trivial summand.  Eliminating all other components by bundle
automorphisms gives
$\mathcal K=\OO\oplus\OO(-1)^2$.

\item Finally $\Omega_X^2=(\Omega_X^1)^\vee\otimes K_X$.  Relative
duality, with the locally free direct images now established, gives
\[
R^qf_*\Omega_X^2
=\bigl(R^{2-q}f_*\Omega_X^1\bigr)^\vee\otimes K_B.
\]
Substitution of $K_B=\OO(-2)$ gives every entry of the second column.
No Euler-characteristic identity is used to select a Hodge number.
\end{enumerate}
\end{proof}

\subsection{An exact modular determinant and two numerical routes}

The scalar depends on the fourth root chosen in $t-1=s^4$.  We make
that ambiguity explicit.  We choose the local branch for which, on
the positive $t>1$ arc leading to the distinguished cusp,
$\tau=iy$ with $y>1$, and
\begin{equation}\label{inv:eq:A-s-branch}
\tau=i+\frac{i}{\pi\sqrt{E_4(i)}}s^2+O(s^4),\qquad
\operatorname{Im}\mu_s(0)>0.
\end{equation}
The nonvanishing proved below makes the second condition meaningful;
it fixes the remaining sign of $s$.  Equivalently it is the lower
cusp-regular zero branch of the Weierstrass function in Statement \ref{inv:A2.3}.
Replacing $s$ by $i^js$ multiplies the reported derivative by $i^{-j}$.
The workbench instead normalizes $ds/dz(z_2)=1$; its number is obtained
by the chain rule and is not numerically the same unspecified coordinate.

\begin{proposition}[The complete raw period determinant]\label{inv:A2.1}\label{wb:prop:log-ks-divisor}\label{wb:prop:period-derivative-determinant}
\StatementRecord{inv:A2.1}{proved}{Historical events: Claude Fable 5.2 (Anthropic), effort high, 12 November 2026: targeted proof check. Claude Fable 5.2 (Anthropic), effort high, 16 November 2026: used and rechecked in this run. Claude Fable 5.2 (Anthropic), effort high, 16 November 2026 13:58 UTC: targeted proof check. Claude Fable 5.2 (Anthropic), effort max, 16 November 2026: targeted proof check. GPT-6 Astra (OpenAI), ChatGPT Pro mode, effort Pro, S6\_09\_PAPERS, 17 November 2026: targeted proof check. Also checked in this writing. Pre-continuation writing events: GPT-6 Astra (OpenAI), ChatGPT Pro mode, effort Pro, S6\_10\_PAPERS, 18 November 2026: targeted proof check. Continuation metadata events: GPT-6 Astra (OpenAI), ChatGPT Pro mode, effort Pro, S6\_10\_PAPERS continuation, 19 November 2026: used and rechecked in this run; status-vocabulary reconciliation; metadata reconciliation only; no new mathematical check. \emph{Pre-continuation recorded proof scope.} Read full source proofs and re-derived logarithmic module, root-character/duality and residue cancellation maps, modular determinant and selected scalar, elliptic-section comparison, Leray table/lifts, and ordinary-relative differential syzygies. Opened AC original pp5,80–81 to verify precise scope and correction wording; external classical citation locators under separate review.}{the hypotheses, notation and definitions stated here; the explicit constructions and calculations in the proof}

On the marked upper half-plane,
\begin{equation}\label{inv:eq:A-modular-K}
6\mu_z^2-\tau_z\beta_z
=\tau_z^2\frac{E_4(\tau)^2}{E_6(\tau)}.
\end{equation}
The right side extends holomorphically through the apparent poles.
Its exact orders are two at $z_1$ and zero at $z_2$; it has no other
zero away from their translates.  The logarithmic connecting map
$\kappa_0:f_*\mathcal E\to\mathcal L\otimes G$ has divisor
$\operatorname{divisor}(\det\kappa_0)=p_1+p_2$.
\end{proposition}


\begin{proof}
\begin{enumerate}[label=\arabic*.,leftmargin=*]
\item Consider the holomorphic quadratic differential
$\mathscr Q=6(d\mu)^2-d\tau\,d\beta$.  Direct differentiation of
the two affine transformation laws in AC26 Definition 3.1 gives
\[
g_j^*\mathscr Q=\tau^{-2}\mathscr Q\quad(j=1,2).
\]
In particular $g_2^2$ fixes $\tau$ and sends
\[
\mu\longmapsto1-\tau-\mu,\qquad
\beta\longmapsto\beta+6\tau+12\mu-6,
\]
which leaves $\mathscr Q$ exactly unchanged.  The symbolic residuals
of all three identities are zero in \texttt{a\_calculations.py}.
The kernel of
$\Delta=\Z/3*\Z/4\to\Z/3*\Z/2=\operatorname{PSL}_2(\Z)$
is the normal subgroup generated by $g_2^2$.  The squared multiplier
of $\mathscr Q$ is the ordinary modular multiplier
$(c\tau+d)^{-2}$, so it is one on every kernel element.  Two generic
points of the marked cover with the same $\tau$ have the same $t$
and differ by an element of $\Delta$; its modular image fixes a
generic point, hence is the identity.  It follows that
$k=\mathscr Q/(d\tau)^2$ is a single-valued meromorphic function of
$\tau$, transforming as a modular form of weight two.  This kernel
argument is needed: mere equivariance of the separate period functions
would not prove descent.

\item At $z_1$, write a linearizing parameter $r$ with
$g_1r=\zeta r$, $\zeta=e^{-2\pi i/3}$.  If $\mathscr Q=K_rdr^2$
has first nonzero coefficient of order $n$, its transformation law
at the fixed point reads $\zeta^{n+2}=\rho^{-2}=\zeta$.
Thus $n\equiv2\pmod3$, and $K_r$ vanishes to order at least two.
The map $\tau$ is unramified there by AC26 Theorem 3.4(i).  At
$z_2$, $\tau-i$ has order two and $K_s$ is holomorphic; therefore
$k$ has at most a simple pole at $i$.  Away from these points
$\tau$ is unramified.  At the cusp $\mu$ and $\beta+\tau$ are
holomorphic in $q$, so $k=1+O(q)$.

\item The modular function $k E_6/E_4^2$ is therefore holomorphic at
$\rho$, at $i$, and at the cusp, as well as everywhere else.
The usual zero statements for $E_4,E_6$ (simple at $\rho,i$,
respectively) are the modular facts used in AC26 Section 3.1.
This function descends to the compact modular curve $X(1)=\Pone$,
so it is constant.  Its cusp value is one.  This proves
\eqref{inv:eq:A-modular-K} and the exact two local orders, without the
workbench's degree count for the coarse determinant line.

\item For completeness we translate to the actual coarse lattices.  Choose
cover eigenframes $e_0,e_1$ of $\mathcal E$ with characters
$1,\zeta^a$, and $b_0,b_1$ of $G$ with characters
$1,\zeta^{-a}$.  Their invariant bases are
$e_0,s^{m-a}e_1$ and $b_0,s^ab_1$.  The logarithmic base frame is
$\ell=dt_j/t_j=m\,ds/s$.  Changing determinants from these frames
to the marked ones gives
\[
\det\kappa_0=\epsilon(t_j)m^{-2}s^{m+2-2a}K_s,
\qquad \epsilon(0)\ne0.
\]
This equality follows by taking the two factors $s/m$ from replacing
$ds$ by $\ell$, the factor $s^{m-a}$ in the source, and the factor
$s^{-a}$ in the target.  Thus at $(m,a)=(3,2)$ the total order in
$s$ is $1+2=3$, and at $(4,1)$ it is $4+0=4$.
Both coarse zeros are simple.  At the cusp, the frame change has
no elementary divisor and $dZ/d\tau$ tends to
$\left(\begin{smallmatrix}0&1\\-1&0\end{smallmatrix}\right)$,
with determinant one.  Formula \eqref{inv:eq:A-modular-K} excludes any
further zero, proving the asserted divisor.
\end{enumerate}
\end{proof}

\begin{proposition}[The evaluated scalar and the two local leading terms]\label{inv:A2.2}
\StatementRecord{inv:A2.2}{proved}{Historical events: Claude Fable 5.2 (Anthropic), effort high, 16 November 2026: used and rechecked in this run. Claude Fable 5.2 (Anthropic), effort high, 16 November 2026 13:58 UTC: targeted proof check. Claude Fable 5.2 (Anthropic), effort max, 16 November 2026: targeted proof check. GPT-6 Astra (OpenAI), ChatGPT Pro mode, effort Pro, S6\_09\_PAPERS, 17 November 2026: targeted proof check. Also checked in this writing. Pre-continuation writing events: GPT-6 Astra (OpenAI), ChatGPT Pro mode, effort Pro, S6\_10\_PAPERS, 18 November 2026: targeted proof check. Continuation metadata events: GPT-6 Astra (OpenAI), ChatGPT Pro mode, effort Pro, S6\_10\_PAPERS continuation, 19 November 2026: used and rechecked in this run; status-vocabulary reconciliation; metadata reconciliation only; no new mathematical check. \emph{Pre-continuation recorded proof scope.} Read full source proofs and re-derived logarithmic module, root-character/duality and residue cancellation maps, modular determinant and selected scalar, elliptic-section comparison, Leray table/lifts, and ordinary-relative differential syzygies. Opened AC original pp5,80–81 to verify precise scope and correction wording; external classical citation locators under separate review.}{the hypotheses, notation and definitions stated here; \cref{inv:eq:A-s-branch}; \cref{inv:eq:A-modular-K}; \cref{inv:A2.3}}

In the coordinate \eqref{inv:eq:A-s-branch},
\begin{equation}\label{inv:eq:A-mu-value}
\mu_s(0)=i\,C_4,
\qquad
C_4=\frac{\sqrt{2/3}}{\pi E_4(i)^{1/4}}
=\frac{4\sqrt\pi}{3^{3/4}\Gamma(1/4)^2}>0.
\end{equation}
Numerically,
\[
C_4=0.2366093140807852246877051134504496062943051340858243372192243072\ldots.
\]
For $r^3=t$, with $r>0$ on the arc $0<t<1$ and $\tau$ travelling
from $\rho$ to $i$ along the unit circle, set
\[
\lambda=\frac{3e^{5\pi i/6}}{2\pi E_6(\rho)^{1/3}}.
\]
Then $\tau=\rho+\lambda r+O(r^2)$ and
\[
K_r=-\frac{3i\lambda}{2\pi}r^2+O(r^3).
\]
The leading coefficient has real and imaginary parts
\[
\begin{aligned}
\operatorname{Re}K_2&=0.08010221108753827140763650309934451243\ldots,\\
\operatorname{Im}K_2&=0.13874109940222333968402480721280640441\ldots.
\end{aligned}
\]
\end{proposition}


\begin{proof}
\begin{enumerate}[label=\arabic*.,leftmargin=*]
\item Ramanujan's differential relations give
$E_6'(i)=-\pi i E_4(i)^2$ and
$E_4'(\rho)=-(2\pi i/3)E_6(\rho)$.
Using $j=1728E_4^3/(E_4^3-E_6^2)$, we obtain
\[
j(\tau)-1728=-1728\pi^2E_4(i)(\tau-i)^2+O((\tau-i)^3).
\]
Since $j=1728(1+s^4)$, the chosen branch is precisely
\eqref{inv:eq:A-s-branch}.  Formula \eqref{inv:eq:A-modular-K} now gives
\[
K_s(0)=-\frac4{\pi^2\sqrt{E_4(i)}}=6\mu_s(0)^2,
\]
because $\tau_s(0)=0$.  This proves nonvanishing and the first
formula in \eqref{inv:eq:A-mu-value}; the sign is fixed by the stated
coordinate convention.  Similarly
$j(\tau)=-512i\pi^3E_6(\rho)(\tau-\rho)^3+cdots$.
This gives the stated $\lambda$, and substitution into
\eqref{inv:eq:A-modular-K} gives the coefficient $-3i\lambda/(2\pi)$.

\item There is also a geometric nonvanishing proof independent of this
modular determinant.  Using the section identification proved in
\ref{inv:A2.3}, set $x=s^2X,y=s^3Y,t=1+s^4$.  The smooth family becomes
\[
Y^2=X^3-3(1+s^4)^3X+2(1+s^4)^4s^2,
\qquad P=(0,-\sqrt2(1+s^4)^2s)
\]
on the chosen branch.  At $s=0$ its equation is
$Y^2=X^3-3X$ and the point is $(0,0)$.  The first variation of the
curve itself is zero.  In the local coordinate $Y$ at that point,
$dX/(2Y)=dY/(3X^2-3)$, so the Abel integral of $P$ has derivative
$\sqrt2/3$.  The marked period is $-i\Omega_R$, rather than the
positive real period $\Omega_R$: the positive imaginary branch
$\tau=iy$, $y>1$, has elliptic scale $F/s=i/\Omega_R$ at zero.
Its period has zero first derivative.  Division gives
$\mu_s(0)=i\sqrt2/(3\Omega_R)\ne0$.
Using $g_2(\Z+i\Z)=(4\pi^4/3)E_4(i)$ and
$g_2(\Omega_R(\Z+i\Z))=12$ gives
$\Omega_R=\pi E_4(i)^{1/4}/\sqrt3$, which identifies this value
with \eqref{inv:eq:A-mu-value}.  This also explains why choosing the
positive real period with a positive real $s$ instead would select
the reciprocal branch $\operatorname{Im}\tau<1$.

\item Here are the numerical methods actually run.  The first evaluates
$E_4(i)=1+240\sum_{n=1}^{100}\sigma_3(n)e^{-2\pi n}$ at 100
decimal working digits and inserts it in \eqref{inv:eq:A-mu-value}.
The second integrates
\[
\Omega_R=\frac{\sqrt2}{3^{1/4}}
\int_0^{\pi/2}(1-\tfrac12\sin^2\theta)^{-1/2}\,d\theta,
\qquad C_4=\frac{\sqrt2}{3\Omega_R}.
\]
This is the real period of $Y^2=X^3-3X$ for $dX/(2Y)$; substituting
the three roots $\sqrt3,0,-\sqrt3$ in the elementary elliptic-integral
change of variable proves the formula.  The beta-integral evaluation
of the displayed complete elliptic integral gives the gamma
expression in \eqref{inv:eq:A-mu-value}.  The Fourier and quadrature
answers agree with absolute discrepancy $1.79\cdot10^{-102}$.
As a third check, Newton iteration on the Weierstrass Fourier series
with $j(\tau)=1728(1+s^4)$, using the two roots adjacent to
$(1-\tau)/2$, gives their symmetric derivative.  At $s=10^{-15}$ it
is
\[
i\,0.23660931408078522468770511345048081721971218\ldots,
\]
with discrepancy $3.13\cdot10^{-32}$ and maximum residual below
$2.9\cdot10^{-101}$.  The order-three coefficient was independently
evaluated using the divisor-sum expansion of $E_6$ and its Jacobi
theta-constant formula, with agreement to the working precision.
The code, all output, and precision are supplied.  These are numerical
corroborations, not interval certificates; nonvanishing and the exact
orders were proved analytically above.  For example the elementary
bound $\sigma_3(n)\le n^4$ and $e^{-2\pi}<1/500$ bounds the omitted
Fourier tail after $N$ terms by
$250(N+1)^4 500^{-(N+1)}$; roundoff is separately controlled by
the working precision and the independent calculation.
\end{enumerate}
\end{proof}

\begin{proposition}[Identification of the elliptic logarithm and the verified zero formula]\label{inv:A2.3}
\StatementRecord{inv:A2.3}{proved}{Historical events: Claude Fable 5.2 (Anthropic), effort high, 16 November 2026 13:58 UTC: carried as background without fault found. GPT-6 Astra (OpenAI), ChatGPT Pro mode, effort Pro, S6\_09\_PAPERS, 17 November 2026: targeted proof check. Also checked in this writing. Pre-continuation writing events: GPT-6 Astra (OpenAI), ChatGPT Pro mode, effort Pro, S6\_10\_PAPERS, 18 November 2026: targeted proof check. Continuation metadata events: GPT-6 Astra (OpenAI), ChatGPT Pro mode, effort Pro, S6\_10\_PAPERS continuation, 19 November 2026: used and rechecked in this run; status-vocabulary reconciliation; metadata reconciliation only; no new mathematical check. \emph{Pre-continuation recorded proof scope.} Read full source proofs and re-derived logarithmic module, root-character/duality and residue cancellation maps, modular determinant and selected scalar, elliptic-section comparison, Leray table/lifts, and ordinary-relative differential syzygies. Opened AC original pp5,80–81 to verify precise scope and correction wording; external classical citation locators under separate review.}{the hypotheses, notation and definitions stated here; \cref{inv:eq:A-s-branch}; \citet{EichlerZagier1982}; \citet{DukeImamoglu2008}; \cref{inv:eq:A-mu-value}}

For the Weierstrass equation
\[
y^2=x^3-3t^3(t-1)x+2t^4(t-1)^2,
\]
the section $[\mu]$ is a generator of its rational section group and
equals one of $P_*=(0,\sqrt2t^2(t-1))$ and $-P_*$.
In particular $\wp(\mu,\tau)=0$ on every regular fibre.
Along the branch in \eqref{inv:eq:A-s-branch}, with
$x_0=1-1/t$, the exact branch is
\begin{equation}\label{inv:eq:A-DI-branch}
\mu=\frac{1-\tau}{2}
 +\frac{i\sqrt6}{3\pi}
 \frac{x_0^{1/4}\,{}_3F_2(1/3,2/3,1;3/4,5/4;x_0)}
 {{}_2F_1(1/12,5/12;1;1-x_0)}.
\end{equation}
Consequently $\mu-(2-\tau)/3$ has exact order two at the
order-three point in $r^3=t$.
\end{proposition}


\begin{proof}
\begin{enumerate}[label=\arabic*.,leftmargin=*]
\item We include the comparison, since a matching $j$-map alone would not
identify a section.  Write $F$ for the elliptic scale, not the
threefold projection.  It is specified by
\[
F^2=-\frac{3t(t-1)E_4(\tau)}{\pi^2E_6(\tau)},\qquad
F/t\longrightarrow i\sqrt3/\pi
\]
at the distinguished cusp.  The local orders are two at $r=0$ and
one at $s=0$; hence the square root exists on the simply connected
marked cover.  Its transformation laws are
$F(g_1z)=-F(z)/\tau(z)$ and $F(g_2z)=F(z)/\tau(z)$.
These follow first after squaring by modularity, and the signs are
fixed by those local orders.  Direct substitution in the cubic
identifies its regular fibres with $\C/(\Z+\tau\Z)$ by
\[
x=F^2\wp(\nu,\tau),\qquad y=\tfrac12F^3\wp'(\nu,\tau).
\]
The transformation laws of $\mu$ make this a single-valued section
$P$.  Its coordinates are meromorphic at the two finite special
points and at the cusp, so are rational functions on $\Pone$.

\item Every rational section $Q$ admits a global holomorphic elliptic
logarithm $\nu$ on the marked cover.  Indeed the exponential map
from the trivial complex line bundle to this smooth elliptic family
is a covering map with locally uniformly discrete lattice; pullback
over a simply connected base lifts the section.  The apparent poles
in its projective Weierstrass coordinates at a finite point are
meromorphic poles of finite order: the coordinates of $Q$ are
rational in $t$, while $F$ is holomorphic with the finite orders
just computed.  Multiplying the three homogeneous coordinates by
a common power of the local uniformizer and cancelling their
common zero gives holomorphic coordinates not all zero at the
centre.  Their limiting point lies on the closed projective
Weierstrass family.  This is the valuative criterion for properness
over the local analytic discrete valuation ring; it is not an
assertion that an arbitrary punctured-disc holomorphic map to a
compact space has a removable singularity.
We can choose this logarithm bounded at the distinguished cusp.
To verify this last assertion, put $v=1/t$, $X=v^2x$, $Y=v^3y$.
The equation becomes
$Y^2=X^3-3(1-v)X+2(1-v)^2$.  A section cannot meet its nodal point
$(X,Y)=(1,0)$ at $v=0$: substituting $X=1+O(v),Y=O(v)$ leaves
nonzero coefficient one in the defining equation at order $v$.
On the smooth nodal locus the coordinate
$u=\exp(2\pi i\nu)$ is a nonzero finite coordinate, as follows
directly from
\[
X=1+\frac{12u}{(u-1)^2},\qquad
Y=\frac{12\sqrt3u(u+1)}{(u-1)^3}.
\]
The normally convergent Weierstrass Fourier series extends this
coordinate to a nearby smoothing.  The inverse function theorem
therefore writes the section as a holomorphic unit $u(v)$, also
near $u=1$, the point at infinity.  A logarithm of that unit fixes
the bounded lift.  Any two such lifts differ by an integer, since
an additional nonzero multiple of $\tau$ is unbounded.

\item For this lift, write the integral transformation errors as
\[
\nu(g_1z)=-\nu/\tau+\alpha_1+\beta_1(\tau-1)/\tau,
\quad
\nu(g_2z)=\nu/\tau+\alpha_2-\beta_2/\tau.
\]
Bounded cusp regularity and $g_1g_2=g_0^{-1}$ imply
$\alpha_2=-\beta_1$, $\beta_2=\alpha_1+\beta_1$.
Putting $n=-\alpha_1-2\beta_1$, $a=\alpha_1+\beta_1$, we find that
$\nu-(n\mu+a)$ satisfies the homogeneous law of AC26 Lemma 3.10
and is cusp-regular.  We recall a direct verification of the degree
used here.  The function
$F_h=E_4(\tau)^2E_6(\tau)^{1/2}/\Delta(\tau)$ on the marked cover
has orders two and one at $z_1,z_2$, and has a simple pole at the
cusp.  The square root exists because its zeros on this cover have
even order.  Modularity and these local orders fix the two signs:
$F_h(g_1z)=-F_h/\tau$, $F_h(g_2z)=F_h/\tau$.
For any holomorphic solution of the same homogeneous law, the
quotient by $F_h$ is an invariant meromorphic function of $t$.
A possible pole at the order-three point would have order at most
$2/3$ in $t$, and one at the order-four point at most $1/4$; its
integral order is therefore nonnegative.  At the cusp regularity
requires this quotient to vanish.  Thus the solution sheaf is
exactly $\OO_B(-p_0)$, with no nonzero global sections.
We obtain $\nu=n\mu+a$ and hence $Q=nP$.  The two distinct
finite values of $\mu$ exclude $nP=0$ for $n\ne0$ by cusp boundedness;
there is no torsion.  This proves that $P$ generates the section group.

\item We identify $P_*$ with a primitive generator by an intersection
calculation.  The rational elliptic surface and its fibre types
$\mathrm{IV}^*,\mathrm{III},I_1$ are those in Engel Section 2.
At $t=1$, the blowup chart $x=(t-1)X,y=(t-1)Y$ gives
$Y^2=(t-1)X^3-3t^3X+2t^4$; $P_*$ meets its exceptional
multiplicity-one component at $(X,Y)=(0,\sqrt2)$, away from the
identity component.  At $t=0$, two successive charts
$x=t^2X,y=t^2Y$ give
$Y^2=t^2X^3-3t(t-1)X+2(t-1)^2$; $P_*$ meets the smooth exceptional
component $Y=-\sqrt2$ at $X=0$.  These are the two nonidentity
simple components of the $E_6$ configuration.  The section is
disjoint from $O$ in every chart, including infinity.

\item Here the height calculation can be made directly, without assuming
primitivity.  Orthogonally project $Q-O$ in
$\operatorname{NS}(S)\otimes\Q$ off $O$, the fibre $F_S$, and the
nonidentity components.  The inverse intersection matrices are
$-C_{E_6}^{-1}$ and $-1/2$.  Their diagonal entries for the two
nonidentity simple components are $-4/3$ and $-1/2$.
Here we may number the nonidentity components so that
\[
C_{E_6}=\begin{pmatrix}
2&-1&0&0&0&0\\
-1&2&-1&0&-1&0\\
0&-1&2&-1&0&0\\
0&0&-1&2&0&0\\
0&-1&0&0&2&-1\\
0&0&0&0&-1&2
\end{pmatrix},\qquad
(C_{E_6}^{-1})_{44}=(C_{E_6}^{-1})_{66}=4/3.
\]
Since $K_S=-F_S$, adjunction gives $Q^2=O^2=-1$.  Expanding the
square of this projection gives
\[
h(Q)=2+2(Q\cdot O)-c_0(Q)-c_1(Q),\quad
c_0(Q)\in\{0,4/3\},\ c_1(Q)\in\{0,1/2\}.
\]
The projection is additive: the chord-and-tangent law makes
$Q+R-(Q+R)_{\rm section}-O$ linearly equivalent to a vertical
divisor, which the projection kills.  Its negative square therefore
satisfies $h(nQ)=n^2h(Q)$.  A section meets a single simple
component, since its base parameter has order one; consequently
the listed local corrections exhaust the possibilities.
For $Q\ne O$ we have $Q\cdot O\ge0$, so $h(Q)\ge1/6$.
For $P_*$ the explicit incidences give $h(P_*)=1/6$.
Writing $P_*=nP$ now gives $1/6=n^2h(P)\ge n^2/6$, and hence
$n=\pm1$.  Inverting the displayed $E_6$ Cartan matrix and the
chord-and-tangent formulas is included in the exact computation log.
This proves $\wp(\mu,\tau)=0$.

\item We checked the original theorem on page 399 of
\citet{EichlerZagier1982}, as well as equation (6) of
\citet{DukeImamoglu2008}.  With the normalized $\Delta$, their
coefficient is indeed $144\pi i\sqrt6$, their constant is
$\log(5+2\sqrt6)/(2\pi i)$, and the integrand is
$(\sigma-\tau)\Delta(\sigma)/E_6(\sigma)^{3/2}$ on the vertical
path with its continued branch.  Thus the displayed Eichler--Zagier integral formula is
correct with its indicated sign and period ambiguity.  Theorem 1
of Duke--Imamoglu gives \eqref{inv:eq:A-DI-branch}: we take one minus
their $z_0$, rather than $z_0-\tau$, because the latter is unbounded
at the distinguished cusp.  The chosen branch approaches
$(1-i)/2$ at $\tau=i$ and has lower imaginary cusp value.
These conditions identify the original bounded lift.  In particular
its linear coefficient is
$i\sqrt6/(3\pi\,{}_2F_1(1/12,5/12;1;1))$, agreeing with
\eqref{inv:eq:A-mu-value}.

\item Finally exact elliptic addition gives
\[
x(2P_*)=9t^2/8,\quad
y(2P_*)=\frac{t^2(32-5t)}{16\sqrt2},\quad
x(3P_*)=\frac{16}{81}(t-1)(5t-32).
\]
The last expression has nonzero limit $512/81$ at $t=0$.
If $\nu=\mu-(2-\tau)/3$, then $3\nu$ is the local logarithm of
$3P$ modulo its exact period.  Since $F$ has order two and
$\wp(3\nu)=(3\nu)^{-2}+O(\nu^2)$, the equality
$x(3P)=F^2\wp(3\nu)$ forces $\nu$ to have exact order two.
It also gives its leading coefficient
$\nu_2=\pm3F_2/(16\sqrt2)$, where $F_2=\lim_{r\to0}F/r^2$;
the sign is the explicitly stated elliptic-logarithm branch.
\end{enumerate}
\end{proof}

\begin{proposition}[The global rank and the decisive scalar]\label{inv:A2.4}\label{wb:thm:h12-scalar}
\StatementRecord{inv:A2.4}{proved}{Historical events: Claude Fable 5.2 (Anthropic), effort high, 12 November 2026: targeted proof check. Claude Fable 5.2 (Anthropic), effort high, 16 November 2026: used and rechecked in this run. Claude Fable 5.2 (Anthropic), effort high, 16 November 2026 13:58 UTC: targeted proof check. Claude Fable 5.2 (Anthropic), effort max, 16 November 2026: targeted proof check. GPT-6 Astra (OpenAI), ChatGPT Pro mode, effort Pro, S6\_09\_PAPERS, 17 November 2026: targeted proof check. Also checked in this writing. Pre-continuation writing events: GPT-6 Astra (OpenAI), ChatGPT Pro mode, effort Pro, S6\_10\_PAPERS, 18 November 2026: targeted proof check. Continuation metadata events: GPT-6 Astra (OpenAI), ChatGPT Pro mode, effort Pro, S6\_10\_PAPERS continuation, 19 November 2026: used and rechecked in this run; status-vocabulary reconciliation; metadata reconciliation only; no new mathematical check. \emph{Pre-continuation recorded proof scope.} Read full source proofs and re-derived logarithmic module, root-character/duality and residue cancellation maps, modular determinant and selected scalar, elliptic-section comparison, Leray table/lifts, and ordinary-relative differential syzygies. Opened AC original pp5,80–81 to verify precise scope and correction wording; external classical citation locators under separate review.}{the hypotheses, notation and definitions stated here; \cref{inv:A2.1}; \cref{inv:A1.3}}

The map
$\kappa_1:R^1f_*\mathcal E\to\mathcal L\otimes R^2f_*\OO=\OO$
is surjective, and $H^0(\kappa_1):\C^2\to\C$ has rank one.
Let $h=a\,d\zeta_1+d\zeta_2$ be the unique global section of
$f_*\mathcal E$ with $h(\xi)=1$.  In a genuine invariant target
frame normalized as below,
\[
\kappa_1(h\gamma)(p_2)=-\tfrac14\mu_s(0)=-\tfrac i4 C_4.
\]
The rank-one condition is equivalent to
$a\tau_z+\mu_z\not\equiv0$.  Under the same direct-image and
sheaf-surjectivity hypotheses a rank-zero global map would instead
give $h^{1,2}=2$, $h^{1,1}=3$; that alternative does not occur here.
\end{proposition}


\begin{proof}
\begin{enumerate}[label=\arabic*.,leftmargin=*]
\item Let $d_{ij}(s)$ be the cover Kodaira--Spencer coefficients in the
eigenframes of Statement \ref{inv:A2.1} after removing $ds$.
The invariant coarse matrix and its cup-product row have the forms
\[
[\kappa_0]=\begin{pmatrix}A&t_jB\\C&D_0\end{pmatrix},\qquad
[\kappa_1]=(-C,A,-D_0,B),
\]
in the ordered basis
$e_0b_0,s^ae_0b_1,s^{m-a}e_1b_0,e_1b_1$.  Specifically
$A=sd_{00}/m$, $B=s^{1-a}d_{01}/m$,
$C=s^{1-a}d_{10}/m$, $D_0=s^{m+1-2a}d_{11}/m$.
These are holomorphic invariant functions; the row follows by
wedging the connecting class with $b_0,b_1$.
If the row had no unit at a finite point, every entry would vanish
modulo $t_j$ and $\det\kappa_0$ would vanish to order at least two,
contradicting \ref{inv:A2.1}.  Away from the finite points, including the cusp,
$\kappa_0$ is invertible, which gives a unit in the cup-product row.
Thus $\kappa_1$ is surjective.

\item The lift $h$ exists and is unique because $f_*\mathcal E$ is an
extension of $\OO$ by $\OO(-1)$.  Its invariance under $g_2$ gives
$a(g_2z)=\tau a+\mu$, hence $a(z_2)=1/2$.
In the quotient frame $(\gamma,u)$ of $G$ the connecting map is
\[
\kappa_0(h)=\frac{s}{4}
\bigl((6a\mu_s+\beta_s)\gamma+(a\tau_s+\mu_s)u\bigr)\ell.
\]
Choose the variable eigenframe $b_1$ by Reynolds averaging so that
$\gamma\wedge b_1=v(s)\gamma\wedge u$, $v(0)=1$.
Then $s\,v(s)(\gamma\wedge u)\ell$ is the actual invariant
target frame; the uncorrected expression with $v=1$ need not be
invariant away from zero.  Cup product with $\gamma$ gives at the
central point $-(a\tau_s+\mu_s)/4=-\mu_s/4$, proving rank one.

\item To check the stated equivalence, there is a second global section
\[
q_E=\sigma_1\otimes u+6\sigma_2\otimes\gamma
\quad\text{of }R^1f_*\mathcal E.
\]
It is the negative of the relative $(1,1)$ image of the invariant
flat class $\eta=u\wedge w+6\gamma\wedge\delta$.
It extends at the cusp because the relative logarithmic forms and
$\gamma,u$ extend there.  At a finite point it is an invariant
holomorphic tensor on the smooth cover, hence belongs to the four
explicit invariant generators above.  Its cup-product row is zero,
since the two terms are $6\mu_z$ and $-6\mu_z$.
It is independent of $h\gamma$ on every regular fibre (its
$\sigma_1\otimes u$ coefficient is nonzero), so these two sections
form a basis of the two-dimensional global section space.
Consequently the global map has rank one precisely when its value
on $h\gamma$, equivalently $a\tau_z+\mu_z$, is not identically
zero.  If its rank were zero, the long exact cohomology sequence of
$0\to\mathcal K\to\OO^2\oplus\OO(-1)^2\to\OO\to0$
would give $h^0(\mathcal K)=2$, $h^1(\mathcal K)=1$.
The residue calculation in \ref{inv:A1.3} and Leray would then give exactly
the two counterfactual Hodge numbers stated in the claim.
\end{enumerate}
\end{proof}

\subsection{The Hodge table and its Leray generators}

\begin{proposition}[The complete Hodge table for every admissible parameter]\label{inv:A3.1}\label{wb:cor:exact-hodge-diamond}\label{wb:thm:parametric-hodge-diamond}\label{wb:prop:workbench-holomorphic-form-vanishing}
\StatementRecord{inv:A3.1}{proved}{Historical events: Claude Fable 5.2 (Anthropic), effort high, 12 November 2026: targeted proof check. Claude Fable 5.2 (Anthropic), effort high, 16 November 2026: used and rechecked in this run. Claude Fable 5.2 (Anthropic), effort high, 16 November 2026 13:58 UTC: targeted proof check. Claude Fable 5.2 (Anthropic), effort high, 17 November 2026: carried as background without fault found. Claude Fable 5.2 (Anthropic), effort max, 16 November 2026: targeted proof check. Claude Opus 5.5 (Anthropic), effort max, sweep 1, scope s6, 18 November 2026: computation or formal certificate. GPT-6 Astra (OpenAI), ChatGPT Pro mode, effort Pro, S6\_09\_PAPERS, 17 November 2026: targeted proof check. Also checked in this writing. Pre-continuation writing events: GPT-6 Astra (OpenAI), ChatGPT Pro mode, effort Pro, S6\_10\_PAPERS, 18 November 2026: targeted proof check. Continuation metadata events: GPT-6 Astra (OpenAI), ChatGPT Pro mode, effort Pro, S6\_10\_PAPERS continuation, 19 November 2026: used and rechecked in this run; status-vocabulary reconciliation; metadata reconciliation only; no new mathematical check. \emph{Pre-continuation recorded proof scope.} Read full source proofs and re-derived logarithmic module, root-character/duality and residue cancellation maps, modular determinant and selected scalar, elliptic-section comparison, Leray table/lifts, and ordinary-relative differential syzygies. Opened AC original pp5,80–81 to verify precise scope and correction wording; external classical citation locators under separate review.}{the hypotheses, notation and definitions stated here; \cref{inv:eq:A-absolute-table}; \cref{inv:A1.0}}

For every $\operatorname{Im}c<-M$, the Hodge numbers are
\[
\begin{array}{c|rrrr}
q\backslash p&0&1&2&3\\\hline
3&0&0&0&1\\
2&0&1&2&1\\
1&1&2&1&0\\
0&1&0&0&0.
\end{array}
\]
\end{proposition}



The vanishing of holomorphic forms and the Hodge-diamond
pattern with one undetermined integer are due to
\citet[Proposition~9.20, p.~80; Theorem~9.1(vi), p.~65]{AC2026}.
The value of that integer is determined here by the nonzero direct-image
scalar. The full calculations of both the pattern and its value are retained.

\begin{proof}
\begin{enumerate}[label=\arabic*.,leftmargin=*]
\item Since $B$ is a curve, coherent Leray has no differential and gives
for $F=\Omega_X^p$ the exact sequences
$0\to H^1(B,R^{q-1}f_*F)\to H^q(X,F)\to H^0(B,R^qf_*F)\to0$.
Substitution of \eqref{inv:eq:A-absolute-table}, with
$h^0(\OO)=h^1(\OO(-2))=1$ and all cohomology of $\OO(-1)$ zero,
gives columns $p=1,2$ directly.  Statement \ref{inv:A1.0} gives column zero,
and Serre duality gives column three.  Thus the numbers are outputs
of the sheaf computation.  Replacing $\beta$ by $\beta+c$ leaves
every derivative, finite eigencharacter, invariant lattice and
logarithmic residue used above unchanged.  The toric gluing units
change, but their relative logarithmic derivatives still differ only
by base logarithmic forms.  Hence the same computations apply at
each admissible $c$, including every point arbitrarily close to the
strict boundary; no manifold at the excluded boundary is asserted.
\end{enumerate}
\end{proof}

\begin{proposition}[Specified associated-graded generators and a global Dolbeault lift]\label{inv:A3.2}\label{wb:lem:bcpr-eta-extension}
\StatementRecord{inv:A3.2}{proved}{Historical events: Claude Fable 5.2 (Anthropic), effort high, 12 November 2026: targeted proof check. Claude Fable 5.2 (Anthropic), effort high, 16 November 2026: carried as background without fault found. GPT-6 Astra (OpenAI), ChatGPT Pro mode, effort Pro, S6\_09\_PAPERS, 17 November 2026: targeted proof check. Also checked in this writing. Pre-continuation writing events: GPT-6 Astra (OpenAI), ChatGPT Pro mode, effort Pro, S6\_10\_PAPERS, 18 November 2026: targeted proof check. Continuation metadata events: GPT-6 Astra (OpenAI), ChatGPT Pro mode, effort Pro, S6\_10\_PAPERS continuation, 19 November 2026: used and rechecked in this run; status-vocabulary reconciliation; metadata reconciliation only; no new mathematical check. \emph{Pre-continuation recorded proof scope.} Read full source proofs and re-derived logarithmic module, root-character/duality and residue cancellation maps, modular determinant and selected scalar, elliptic-section comparison, Leray table/lifts, and ordinary-relative differential syzygies. Opened AC original pp5,80–81 to verify precise scope and correction wording; external classical citation locators under separate review.}{the hypotheses, notation and definitions stated here; \cref{inv:A2.4}; \cref{inv:eq:A-E-line}; \cref{inv:A1.3}}

The Leray associated-graded spaces for the four middle groups are
\[
\begin{array}{c|cc}
&E_2^{1,q-1}&E_2^{0,q}\\\hline
H^{1,1}&H^1(\OO(-2))=\C\alpha&\C\eta_{\rm rel}\\
H^{1,2}&0&\C\upsilon\\
H^{2,1}&\C\upsilon^\vee&0\\
H^{2,2}&\C\eta_{\rm rel}^\vee&\C\alpha^\vee.
\end{array}
\]
Here $\alpha$ is the Cech class of $dt/t$ on the two standard affine
base charts, $\eta_{\rm rel}=-q_E$ is the invariant relative class
from \ref{inv:A2.4}, and $\upsilon$ is the image of the trace-dual of
$1\in f_*\OO(H)$ under $R^2f_*N\to R^2f_*\mathcal E$ and the
isomorphism $R^2f_*\Omega_X^1=R^2f_*\mathcal E$.
The superscript $\vee$ denotes the bases dual under the relative
trace and the base Serre pairing.  There exists a Dolbeault class
$Q\in H^1(X,\Omega_X^1)$ whose restriction to every regular fibre
is $\eta=u\wedge w+6\gamma\wedge\delta$.
\end{proposition}


\begin{proof}
\begin{enumerate}[label=\arabic*.,leftmargin=*]
\item The isomorphism $f_*\Omega_X^1=\Omega_B^1$ in the residue sequence
is the actual pullback map, so $\alpha$ gives its stated Leray
class; it can be represented by a nonzero normalization of
$f^*\omega_{\rm FS}$ in Dolbeault cohomology.  The section
$\eta_{\rm rel}$ belongs to the canonical kernel identification
$R^1f_*\Omega_X^1=\ker\kappa_1$ and spans its global sections.
Leray is surjective onto this space, so any preimage defines $Q$
with the asserted fibre restriction.  Its ambiguity is exactly
$\C\alpha$; the statement does not claim a canonical global form.
The injection $R^2f_*N=\OO\to R^2f_*\mathcal E$ from
\eqref{inv:eq:A-E-line} spans its global sections, since the quotient
is $\OO(-1)$.  This defines $\upsilon$ with no choice of a negative
degree summand.  Duality used in \ref{inv:A1.3} gives the remaining entries
and their specified dual pairings.  Choosing a lift of
$\eta_{\rm rel}$ and of $\alpha^\vee$ gives bases of the whole
middle cohomology groups.  The associated-graded descriptions are
canonical with the named choices; the splitting of Leray vector
spaces is not asserted canonical.
\end{enumerate}
\end{proof}

\subsection{Ordinary relative differentials and the correction at the cusp}

\begin{proposition}[All local modules needed to distinguish degrees one and two]\label{inv:A4.1}
\StatementRecord{inv:A4.1}{proved}{Historical events: Claude Fable 5.2 (Anthropic), effort high, 12 November 2026: targeted proof check. Claude Fable 5.2 (Anthropic), effort high, 16 November 2026: used and rechecked in this run. GPT-6 Astra (OpenAI), ChatGPT Pro mode, effort Pro, S6\_09\_PAPERS, 17 November 2026: targeted proof check. Also checked in this writing. Pre-continuation writing events: GPT-6 Astra (OpenAI), ChatGPT Pro mode, effort Pro, S6\_10\_PAPERS, 18 November 2026: targeted proof check. Continuation metadata events: GPT-6 Astra (OpenAI), ChatGPT Pro mode, effort Pro, S6\_10\_PAPERS continuation, 19 November 2026: used and rechecked in this run; status-vocabulary reconciliation; metadata reconciliation only; no new mathematical check. \emph{Pre-continuation recorded proof scope.} Read full source proofs and re-derived logarithmic module, root-character/duality and residue cancellation maps, modular determinant and selected scalar, elliptic-section comparison, Leray table/lifts, and ordinary-relative differential syzygies. Opened AC original pp5,80–81 to verify precise scope and correction wording; external classical citation locators under separate review.}{the hypotheses, notation and definitions stated here; the explicit constructions and calculations in the proof}

Let $R=\C\{x,y,z\}$.  For $t=xy$,
\[
\Omega_{R/\C\{t\}}^1=R^3/R(y,x,0)
\simeq(x,y)\oplus R,
\]
and
\[
\Omega_{R/\C\{t\}}^2\simeq
(x,y)\oplus R/(x,y),\qquad
\Omega_{R/\C\{t\}}^3\simeq R/(x,y).
\]
For $t=xyz$,
$\Omega^1=R^3/R(yz,xz,xy)$ is torsion-free of rank two and is
not free at the triple point.  If $I=(xy,xz,yz)$, there is an exact
sequence
\[
0\to T\to\Omega^2\xrightarrow{\wedge dt} I\to0,
\qquad
T=\frac{R u\oplus R v}{(xu,\ y(u-v),\ zv)},
\quad IT=0,
\]
and $\Omega^3=R/I$.  Finally, for $t=s^m$ in
$R=\C\{s,y_1,y_2\}$,
\[
\begin{aligned}
\Omega^1&=R^2\oplus R/(s^{m-1})\,ds,\\
\Omega^2&=R\,dy_1\wedge dy_2
 \oplus\bigl(R/(s^{m-1})\bigr)^2,\\
\Omega^3&=R/(s^{m-1})\,ds\wedge dy_1\wedge dy_2.
\end{aligned}
\]
\end{proposition}


\begin{proof}
\begin{enumerate}[label=\arabic*.,leftmargin=*]
\item The degree-one presentation is the quotient by $dt$; higher degrees
are its exterior powers.  At a double point the transverse quotient
$R^2/R(y,x)$ maps isomorphically to $(x,y)$ by
$(a,b)\mapsto xa-yb$ after an immaterial change of sign of one
basis vector.  More explicitly one may use
$(a,b)\mapsto xa-yb$ with relation $(y,x)$; its kernel is precisely
that relation, since $x,y$ are relatively prime.  In degree two,
with basis $A=dx\wedge dy,B=dx\wedge dz,C=dy\wedge dz$, the
relations are $xA=yA=0$, $yB+xC=0$.  This gives the displayed
torsion summand and the ideal $(x,y)$; wedging once more gives
$R/(x,y)$.

\item At the triple point the map from ordinary to logarithmic relative
one-forms sends $(a,b,c)$ to $(xa-zc,yb-zc)$.
Its kernel is precisely $R(yz,xz,xy)$: the two equations imply
$c=xyh$, $a=yzh$, $b=xzh$ by relative primality of the coordinates.
Thus it injects into $R^2$, proving torsion-freeness.  Its fibre at
the maximal ideal has dimension three although its rank is two,
so it is not free.  The same minimal-generator argument works at
the double locus.

\item In degree two, wedging with $dt$ is the row
$(xy,-xz,yz)$ in the basis $(A,B,C)$.
Its syzygies are generated freely by
$u=(z,y,0)$ and $v=(0,y,x)$.
To check generation, an equation $xyA-xzB+yzC=0$ implies
$x\mid C$, say $C=xC_1$; then $yA=z(B-yC_1)$ implies
$A=zA_1$, $B=y(A_1+C_1)$.  Thus the syzygy is $A_1u+C_1v$.
The three relations from $dt\wedge(dx,dy,dz)$ are
$-xu,y(u-v),zv$, proving the stated presentation of the kernel.
The ideal $I$ kills it by those relations.  Conversely the quotient
$I$ is torsion-free, so this kernel is exactly the torsion of
$\Omega^2$.  The presentation immediately gives $\Omega^3=R/I$.
All matrix identities have also been checked over $\Z[x,y,z]$ by
the supplied exact program.  For $t=s^m$ the sole degree-one
relation is $ms^{m-1}ds=0$; $m$ is a unit and taking exterior
powers gives the three formulas.
\end{enumerate}
\end{proof}

\begin{proposition}[Refuted: The claimed degree-one torsion along the cusp]\label{inv:A4.2}
\StatementRecord{inv:A4.2}{proved}{Historical events: Claude Fable 5.2 (Anthropic), effort high, 16 November 2026: carried as background without fault found; used and rechecked in this run. Claude Fable 5.2 (Anthropic), effort high, 17 November 2026: carried as background without fault found. GPT-6 Astra (OpenAI), ChatGPT Pro mode, effort Pro, S6\_09\_PAPERS, 17 November 2026: targeted proof check. Also checked in this writing. Pre-continuation writing events: GPT-6 Astra (OpenAI), ChatGPT Pro mode, effort Pro, S6\_10\_PAPERS, 18 November 2026: targeted proof check. Continuation metadata events: GPT-6 Astra (OpenAI), ChatGPT Pro mode, effort Pro, S6\_10\_PAPERS continuation, 19 November 2026: used and rechecked in this run; status-vocabulary reconciliation; metadata reconciliation only; no new mathematical check. \emph{Pre-continuation recorded proof scope.} Read full source proofs and re-derived logarithmic module, root-character/duality and residue cancellation maps, modular determinant and selected scalar, elliptic-section comparison, Leray table/lifts, and ordinary-relative differential syzygies. Opened AC original pp5,80–81 to verify precise scope and correction wording; external classical citation locators under separate review.}{the hypotheses, notation and definitions stated here; \cref{inv:A4.1}}

The assertion in AC26 Remark 9.21, PDF page 81, that
$\Omega_{X/B}^1$ has torsion along $W\cup S_1\cup S_2$ is false.
Its torsion is supported precisely on $S_1\cup S_2$; its degree-one
cusp modules are torsion-free but not locally free on the double
locus.  The plural statement in Remark 1.3, PDF page 5, must be
qualified by degree: $\Omega^2_{X/B}$ does have torsion on the cusp
double locus.  The workbench's degree-one correction is correct;
it must not be extended to degree two.
\end{proposition}


\begin{proof}
\begin{enumerate}[label=\arabic*.,leftmargin=*]
\item Both disputed sentences were checked on rendered pages of the
original PDF, so this is not a transcription correction.  The
double and triple charts in \ref{inv:A4.1} exhaust the singular strata of
$W$.  They prove absence of degree-one torsion at the cusp and
exhibit nonzero degree-two torsion there.  At a multiple fibre the
finite etale comparison
$q^*\Omega^1_{N_j/B}=\Omega^1_{J'_j/B}$ is an isomorphism,
and its last summand in \ref{inv:A4.1} is nonzero for $m=3,4$.
Faithful flatness descends both its existence and its support.
There are no other critical fibres and the ordinary relative
cotangent is locally free on their complement.  Thus the stated
support is exact.  The erroneous sentence is at transcription
line 4330 (the final paragraph of Remark 9.21); the workbench's
correct degree-one replacement occupies lines 980--1006 of
\texttt{hodge\_frolicher.tex}.  No numerical Hodge value in that
workbench is refuted by this calculation.
\end{enumerate}
\end{proof}



\section{Fr\"olicher, Bott--Chern, Aeppli, and Hermitian metrics}
\label{inv:sec:partC}

\subsection{The differentials and their dependence on the Hodge number}

In this section $X$ is the standard marked $(3,4)$ threefold and
$h=h^{1,2}(X)=1$ by Part A.  We retain $h$ in the formal calculations
to specify exactly which conclusions depend on its evaluated value.
Every assertion involving $h$ is valid for every value
consistent with the Hodge table
\begin{equation}\label{inv:eq:C-hodge-input}
\begin{array}{c|rrrr}
q\backslash p&0&1&2&3\\\hline
3&0&0&0&1\\
2&0&h&h+1&1\\
1&1&h+1&h&0\\
0&1&0&0&0.
\end{array}
\end{equation}
Thus no conclusion below silently uses $h=1$.  The geometric inputs for
this part are the structure-sheaf calculation, the vanishing of
holomorphic forms, and the sphere topology, which give
\eqref{inv:eq:C-hodge-input}; Statements \ref{inv:A1.3} and \ref{inv:A3.1} determine $h$.
We use $d=\partial+\bar\partial$ on complex-valued
smooth forms and the complex orientation for integration.

\begin{proposition}[The complete Fr\"olicher differential in terms of $h$]\label{inv:C1.1}\label{wb:thm:frolicher-e2-candidate}
\StatementRecord{inv:C1.1}{proved}{Historical events: Claude Fable 5.2 (Anthropic), effort high, 12 November 2026: targeted proof check. Claude Fable 5.2 (Anthropic), effort high, 16 November 2026: used and rechecked in this run. Claude Fable 5.2 (Anthropic), effort high, 16 November 2026 13:58 UTC: targeted proof check. GPT-6 Astra (OpenAI), ChatGPT Pro mode, effort Pro, S6\_09\_PAPERS, 17 November 2026: targeted proof check. Also checked in this writing. Pre-continuation writing events: GPT-6 Astra (OpenAI), ChatGPT Pro mode, effort Pro, S6\_10\_PAPERS, 18 November 2026: targeted proof check. Continuation metadata events: GPT-6 Astra (OpenAI), ChatGPT Pro mode, effort Pro, S6\_10\_PAPERS continuation, 19 November 2026: used and rechecked in this run; status-vocabulary reconciliation; metadata reconciliation only; no new mathematical check. \emph{Pre-continuation recorded proof scope.} Re-derived spectral-sequence rank exhaustion, arrow/square counts, signed collar factor12, Bott–Chern product vanishing using explicit primitives and pairings, metric/Stokes obstruction and Aeppli restriction to the indefinite fibre class. This does not award a reading of the externally cited double-complex source yet.}{the hypotheses, notation and definitions stated here; \cref{inv:eq:C-hodge-input}}

For a compact complex threefold with \eqref{inv:eq:C-hodge-input} and the
complex de Rham cohomology of $S^6$, the Fr\"olicher spectral sequence
satisfies $E_1\ne E_2=E_\infty$.  Its only nonzero $d_1$ ranks are
\[
\rank d_1^{0,1}=1,\quad \rank d_1^{1,1}=h,\quad
\rank d_1^{1,2}=h,\quad \rank d_1^{2,2}=1.
\]
In particular its two middle rows are exact.
\end{proposition}


\begin{proof}
\begin{enumerate}[label=\arabic*.,leftmargin=*]
\item Here $E_1^{p,q}=H^q(X,\Omega_X^p)$ and
$d_r:E_r^{p,q}\longrightarrow E_r^{p+r,q-r+1}$ are complex-linear.
For $r=2$, any source in the two middle rows has target in one of
$(2,0),(3,0),(3,1)$ or outside the table, except the source $(0,2)$,
which is zero.  Sources in row $q=3$ outside $(3,3)$ are zero, and
$(3,3)$ has no target.  For $r=3$ the only possible targets inside the
table have sources $(0,2)$ or $(0,3)$, both zero.  For $r>3$ there is
no arrow inside the table.  A zero term remains zero on subsequent
pages.  Hence $E_2=E_\infty$.

\item The filtration on $H^k_{\mathrm{dR}}(X;\C)$ has no nonzero associated
graded term for $1\le k\le5$.  Every middle term of $E_2$ therefore
vanishes.  Consequently
\[
0\longrightarrow\C\longrightarrow\C^{h+1}
\longrightarrow\C^h\longrightarrow0,
\qquad
0\longrightarrow\C^h\longrightarrow\C^{h+1}
\longrightarrow\C\longrightarrow0
\]
are exact, giving the stated ranks.  The first rank is one, so $E_1$
cannot equal $E_2$.
\end{enumerate}
\end{proof}

\begin{proposition}[A basis on which all first differentials are specified]\label{inv:C1.2}
\StatementRecord{inv:C1.2}{proved}{Historical events: Claude Fable 5.2 (Anthropic), effort high, 16 November 2026: carried as background without fault found. GPT-6 Astra (OpenAI), ChatGPT Pro mode, effort Pro, S6\_09\_PAPERS, 17 November 2026: targeted proof check. Also checked in this writing. Pre-continuation writing events: GPT-6 Astra (OpenAI), ChatGPT Pro mode, effort Pro, S6\_10\_PAPERS, 18 November 2026: targeted proof check. Continuation metadata events: GPT-6 Astra (OpenAI), ChatGPT Pro mode, effort Pro, S6\_10\_PAPERS continuation, 19 November 2026: used and rechecked in this run; status-vocabulary reconciliation; metadata reconciliation only; no new mathematical check. \emph{Pre-continuation recorded proof scope.} Re-derived spectral-sequence rank exhaustion, arrow/square counts, signed collar factor12, Bott–Chern product vanishing using explicit primitives and pairings, metric/Stokes obstruction and Aeppli restriction to the indefinite fibre class. This does not award a reading of the externally cited double-complex source yet.}{the hypotheses, notation and definitions stated here; \cref{inv:A1.3}; \cref{inv:A3.2}; \cref{inv:C1.1}; \cref{inv:C3.7}}

Choose a nonzero $g\in H^{0,1}_{\bar\partial}(X)$, a basis
$q_1,\ldots,q_h$ of a complement to $\C d_1g$ in
$H^{1,1}_{\bar\partial}(X)$, a basis $v_1,\ldots,v_h$ of
$H^{1,2}_{\bar\partial}(X)$, and a class $e$ complementary to the
image of $d_1:H^{1,2}\to H^{2,2}$.  With
$a=d_1g$, $c_j=d_1q_j$, $d_j=d_1v_j$, and $t=d_1e$, the first
differential is
\[
g\mapsto a\mapsto0,\quad q_j\mapsto c_j\mapsto0,
\qquad v_j\mapsto d_j\mapsto0,\quad e\mapsto t\mapsto0.
\]
These are bases in the corresponding bidegrees.  If $h=1$ and a
class $q$ restricts to the invariant fibre class $\eta$ described below,
we may take $q_1=q$, $v_1=gq$, and $e=q^2$, after scaling $g$ if
necessary.  Statements \ref{inv:A1.3} and \ref{inv:A3.2} prove these hypotheses for $X$,
so this geometric choice applies to the threefold in the stated problem.
\end{proposition}


\begin{proof}
\begin{enumerate}[label=\arabic*.,leftmargin=*]
\item The first assertions follow from the injectivity and surjectivity
proved in \ref{inv:C1.1}.  For the last assertion write
$\eta=u\wedge w+6\gamma\wedge\delta$ in the marked integral
cohomology of a regular fibre.  Its complex-oriented square has
integral $-12$.  Its product with the $(0,1)$ image of $\gamma$ is
nonzero: the explicit period calculation in \ref{inv:C3.7} gives its nonzero
pairing with $d\zeta_2$.  Thus $gq$ and $q^2$ have nonzero fibre
restrictions.  The images under $d_1$ restrict to zero on the fibre,
because the fibre is a complex torus and its invariant representatives
are $d$-closed.  Therefore $q$ complements $d_1g$, $gq$ spans the
one-dimensional $H^{1,2}$, and $q^2$ complements $d_1H^{1,2}$.
The extension used in this last choice is the actual Leray lift in
Statement \ref{inv:A3.2}.  This argument uses that extension; it does not infer
it from a calculation on the smooth fibre alone.
\end{enumerate}
\end{proof}

\subsection{A computation from the two definitions}

\begin{proposition}[Bott--Chern and Aeppli tables as functions of $h$]\label{inv:C2.1}\label{wb:thm:exact-bott-chern-aeppli}
\StatementRecord{inv:C2.1}{proved}{Historical events: Claude Fable 5.2 (Anthropic), effort high, 12 November 2026: targeted proof check. Claude Fable 5.2 (Anthropic), effort high, 16 November 2026: used and rechecked in this run. Claude Fable 5.2 (Anthropic), effort high, 16 November 2026 13:58 UTC: targeted proof check. GPT-6 Astra (OpenAI), ChatGPT Pro mode, effort Pro, S6\_09\_PAPERS, 17 November 2026: targeted proof check. Also checked in this writing. Pre-continuation writing events: GPT-6 Astra (OpenAI), ChatGPT Pro mode, effort Pro, S6\_10\_PAPERS, 18 November 2026: computation or formal certificate; targeted proof check. GPT-6 Astra (OpenAI), ChatGPT Pro mode, effort Pro, S6\_10\_PAPERS, 18 November 2026 04:30:14.617321+00:00: used and rechecked in this run. Continuation metadata events: GPT-6 Astra (OpenAI), ChatGPT Pro mode, effort Pro, S6\_10\_PAPERS continuation, 19 November 2026: used and rechecked in this run; status-vocabulary reconciliation; metadata reconciliation only; no new mathematical check. \emph{Pre-continuation recorded proof scope.} Re-derived spectral-sequence rank exhaustion, arrow/square counts, signed collar factor12, Bott–Chern product vanishing using explicit primitives and pairings, metric/Stokes obstruction and Aeppli restriction to the indefinite fibre class. This does not award a reading of the externally cited double-complex source yet. PASS exact reruns c\_double\_complex.py (matrix kernels/images versus arrow counts for h=0..5) and c\_sign\_and\_fibre.py (collar orientation, lift lattice, Hermitian determinant and real-period orientation). Programs and stdout retained under code/I/replay\_invariants.}{the hypotheses, notation and definitions stated here; \cref{inv:A3.1}; \citet[Theorem 3]{Stelzig2021}; \cref{inv:C1.1}; \citet[Proposition 6(1)]{Stelzig2021}; \cref{inv:eq:C-arrows}}

With
\[
H^{p,q}_{BC}=
\frac{\ker\partial\cap\ker\bar\partial\subset A^{p,q}}
{\im\partial\bar\partial},\qquad
H^{p,q}_{A}=
\frac{\ker\partial\bar\partial\subset A^{p,q}}
{\im\partial+\im\bar\partial},
\]
the complete tables are
\begin{equation}\label{inv:eq:C-BC-A}
\begin{array}{c|rrrr}
q\backslash p&0&1&2&3\\\hline
3&0&0&1&1\\
2&0&h&2h&1\\
1&0&2&h&0\\
0&1&0&0&0
\end{array}_{BC}
\qquad
\begin{array}{c|rrrr}
q\backslash p&0&1&2&3\\\hline
3&0&0&0&1\\
2&0&h&2&0\\
1&1&2h&h&0\\
0&1&1&0&0.
\end{array}_{A}
\end{equation}
In particular, substitution of the value $h=1$ proved in \ref{inv:A3.1} gives
the workbench's numerical tables.  The totals are
$(1,0,2,2,2,2,1)$ and $(1,2,2,2,2,0,1)$, respectively.
\end{proposition}


\begin{proof}
\begin{enumerate}[label=\arabic*.,leftmargin=*]
\item We apply the structure theorem for bounded double complexes over a
field, \citet[Theorem 3]{Stelzig2021}, to the actual algebraic double
complex $(A^{\bullet,\bullet}(X),\partial,\bar\partial)$.
The word bounded refers to the bidegrees, not to the dimensions of the
spaces of smooth forms, so that theorem applies here.  It gives a
direct-sum decomposition by isomorphisms of double complexes into
squares and zigzags.  An odd zigzag has one-dimensional total
cohomology.  All odd zigzags in intermediate total degree are therefore
excluded by $b_1=\cdots=b_5=0$.  An odd zigzag of length greater than
one would have at least two vertices of its cohomological total degree.
The extreme diagonals $p+q=0$ and $p+q=6$ contain just one bidegree
each, so the remaining odd zigzags are exactly one dot at $(0,0)$ and
one at $(3,3)$.

\item An even zigzag of length $2r$ contributes a nonzero $d_r$ to one of
the two Fr\"olicher spectral sequences.  \ref{inv:C1.1} and complex conjugation
exclude $r\ge2$ in both sequences.  Thus the remaining summands are
arrows.  The horizontal arrows, with their multiplicities, are
\begin{equation}\label{inv:eq:C-arrows}
\begin{array}{c|c|r}
\text{source}&\text{target}&\text{multiplicity}\\\hline
(0,1)&(1,1)&1\\
(1,1)&(2,1)&h\\
(1,2)&(2,2)&h\\
(2,2)&(3,2)&1.
\end{array}
\end{equation}
The vertical arrows are their reflections in $p=q$, with the same
multiplicities.  The correspondence with spectral-sequence ranks is
\citet[Proposition 6(1)]{Stelzig2021}; alternatively one computes
column cohomology on each horizontal arrow directly.

\item We now evaluate the displayed definitions, rather than use a formula
for the dimensions.  On an arrow, the source is killed by neither
the outgoing differential nor its total differential, whereas its
target is closed under both.  There is no composite
$\partial\bar\partial$.  Hence the target contributes one dimension
to Bott--Chern cohomology, and the source contributes one dimension
to Aeppli cohomology.  On a square, only the top right corner is closed
under both operators, and it is the image of
$\partial\bar\partial$ from the bottom left corner.  The other
three corners of the square are images of $\partial$ or
$\bar\partial$, while the bottom left corner is not in
$\ker\partial\bar\partial$.  Thus a square contributes to neither
cohomology.  A dot contributes once to each.  These computations
commute with direct sums, including the possibly infinite sum of
squares.  Counting the sources and targets in \eqref{inv:eq:C-arrows}
and its reflection gives \eqref{inv:eq:C-BC-A}.

\item The program \texttt{code/c\_double\_complex.py} performs a separate
kernel-and-image computation over $\Q$ and a source-and-target
count.  It includes an acyclic square with its signs, checks
$\partial\bar\partial+\bar\partial\partial=0$, and verifies the
two answers for $h=0,1,2,3,4,5$.  Its output is
\texttt{outputs/c\_double\_complex.txt}.  This is exact linear
algebra, not a floating-point check and not a determination of $h$.
\end{enumerate}
\end{proof}

\begin{proposition}[Comparison with McHugh and Angella]\label{inv:C2.2}
\StatementRecord{inv:C2.2}{proved}{Historical events: Claude Fable 5.2 (Anthropic), effort high, 16 November 2026: carried as background without fault found. GPT-6 Astra (OpenAI), ChatGPT Pro mode, effort Pro, S6\_09\_PAPERS, 17 November 2026: targeted proof check. Also checked in this writing. Pre-continuation writing events: GPT-6 Astra (OpenAI), ChatGPT Pro mode, effort Pro, S6\_10\_PAPERS, 18 November 2026: targeted proof check. GPT-6 Astra (OpenAI), ChatGPT Pro mode, effort Pro, S6\_10\_PAPERS, 18 November 2026 04:30:14.617321+00:00: used and rechecked in this run. Continuation metadata events: GPT-6 Astra (OpenAI), ChatGPT Pro mode, effort Pro, S6\_10\_PAPERS continuation, 19 November 2026: used and rechecked in this run; status-vocabulary reconciliation; metadata reconciliation only; no new mathematical check. \emph{Pre-continuation recorded proof scope.} Re-derived the McHugh and Angella dimension identities in the stated double-complex dimensions, including HBC22=2h. Exact primary formula locators remain separately tracked; no source-verification credit is implied.}{the hypotheses, notation and definitions stated here; \cref{inv:eq:C-hodge-input}; \cref{inv:C2.1}; \cref{inv:C1.1}; \citet[Theorem 6.3, Proposition 6.4, (7), pp. 20--25]{McHugh2018}; \cref{inv:eq:C-BC-A}; \citet[Section 4, pp. 13--16]{Angella2018}}

The specialization of McHugh's Theorem 6.3, Proposition 6.4, equation
(7), and Section 7.1 to \eqref{inv:eq:C-hodge-input} agrees with \ref{inv:C2.1}.
The corresponding relations and Figure 5 in Angella's Section 4 also
agree.  No correction to these specific formulas is required.
\end{proposition}


\begin{proof}
\begin{enumerate}[label=\arabic*.,leftmargin=*]
\item McHugh's quotient $K^{p,q}$ is zero in intermediate total degrees
on every arrow and square just used: a closed corner of an arrow is
the image of the incoming differential and a closed corner of a square
is in both images.  Thus $k^{p,q}=0$ for $0<p+q<6$.
Also $E_2^{3,1}=0$ by \ref{inv:C1.1} and $H^{1,1}(X;\R)=0$ because
$H^2(X;\R)=0$.  Reading the formulas in
\citet[Theorem 6.3, Proposition 6.4, (7), pp. 20--25]{McHugh2018}
therefore gives
\[
h_{BC}^{1,1}=2h^{0,1}=2,\qquad
h_{BC}^{1,2}=h^{1,2}=h,\qquad
h_{BC}^{2,2}=-1+(h+1)+h=2h.
\]
The remaining entries in his table are precisely the boundary entries
of \eqref{inv:eq:C-BC-A}.  In
\citet[Section 4, pp. 13--16]{Angella2018}, the relation
\[
h_{BC}^{2,2}=2h_{BC}^{2,1}-2h^{0,1}+2
\]
gives $2h$, and the boundary values reduce to those above because
$h^{2,0}=h^{0,2}=0$.  This comparison concerns exactly the formulas
used here; it does not certify every assertion in either article.
\end{enumerate}
\end{proof}

\subsection{The algebra, line bundles, and metric consequences}

\begin{proposition}[The real Bott--Chern Chern-class map and a global $G$]\label{inv:C3.1}\label{wb:cor:picard-line-audit}
\StatementRecord{inv:C3.1}{proved}{Historical events: Claude Fable 5.2 (Anthropic), effort high, 12 November 2026: targeted proof check. Claude Fable 5.2 (Anthropic), effort high, 16 November 2026: carried as background without fault found. GPT-6 Astra (OpenAI), ChatGPT Pro mode, effort Pro, S6\_09\_PAPERS, 17 November 2026: targeted proof check. Also checked in this writing. Pre-continuation writing events: GPT-6 Astra (OpenAI), ChatGPT Pro mode, effort Pro, S6\_10\_PAPERS, 18 November 2026: targeted proof check. Continuation metadata events: GPT-6 Astra (OpenAI), ChatGPT Pro mode, effort Pro, S6\_10\_PAPERS continuation, 19 November 2026: used and rechecked in this run; status-vocabulary reconciliation; metadata reconciliation only; no new mathematical check. \emph{Pre-continuation recorded proof scope.} Re-derived the explicit Cech exponential-to-unitary-connection morphism, curvature sign and real BC image; checked the Green identity, nonvanishing via compact plurisubharmonic maximum principle and exact f*O(1) identification. Read AC26 Th10.5(b) Steps1/4 p88 for the separately credited Picard/fibre restriction input and included its short exact-sequence application.}{\cref{inv:C1.1}, \cref{inv:C2.1}, \cref{inv:A1.0}; \cite[Theorem~10.5(b), proof Steps~1 and~4, p.~88]{AC2026}}

The Picard-group identification and fibrewise injectivity below are due to Claude and Alp\"oge \cite[Theorem~10.5(b), proof Steps~1 and~4, p.~88]{AC2026}. We compute the real Chern map and a specified Green-operator representative in that coordinate. Let $G$ be a $\bar\partial$-closed smooth $(0,1)$ representative of
a nonzero $g\in H^1(X,\OO_X)$.  Set
\[
\mathsf A=[\partial G]_{BC},\qquad
\mathsf B=[\bar\partial\bar G]_{BC}.
\]
These classes form a complex basis of $H^{1,1}_{BC}(X)$.
With the exponential $z\mapsto e^{2\pi iz}$ and the Chern convention
$c_{1,BC}=[iF/(2\pi)]$, the map of real Lie groups is
\begin{equation}\label{inv:eq:C-Pic-map}
\Pic(X)=\{\exp(zg):z\in\C\}\longrightarrow H^{1,1}_{BC}(X),
\qquad \exp(zg)\longmapsto z\mathsf A+\bar z\mathsf B.
\end{equation}
It is injective and its image is the real subspace fixed by conjugation.

There is an operator-explicit global representative.  Normalize the
Fubini--Study form by $\int_{\Pone}\omega_{FS}=1$, put
$\theta=f^*\omega_{FS}$, choose any smooth Riemannian metric on $X$,
and let $\mathcal G_d$ be its Green operator for the real de Rham
Laplacian.  Then
\begin{equation}\label{inv:eq:C-G-Green}
\alpha=d^*\mathcal G_d\theta,\qquad G_0=\alpha^{0,1}
\end{equation}
satisfy $d\alpha=\theta$, $\bar\partial G_0=0$,
$[G_0]\ne0$, and $\exp([G_0])=f^*\OO_{\Pone}(1)$.
Restriction $H^1(X,\mathcal O_X)\to H^1(X_b,\mathcal O_{X_b})$ is injective for every scheme-theoretic fibre; on a smooth fibre its image is $\mathbb C\gamma^{0,1}$.
\end{proposition}

\begin{proof}
\begin{enumerate}[label=\arabic*.,leftmargin=*]
\item By \ref{inv:C1.1}, the Dolbeault image of $\mathsf A$ is
$d_1g\ne0$. The class $\mathsf B=\overline{\mathsf A}$ is nonzero,
whereas its Dolbeault image is zero. Thus these two classes are
independent; the dimension calculation of \ref{inv:C2.1} makes
them a basis. In this basis conjugation sends $(a,b)$ to
$(\bar b,\bar a)$, so its fixed real subspace consists precisely
of $(z,\bar z)$.

\item The exponential sequence and $H^1(X;\Z)=H^2(X;\Z)=0$ give the
cited identification $\Pic(X)=H^1(X,\OO_X)=\C g$.
On a smooth trivial line with unit frame $s$, put
$\bar\partial_zs=-2\pi izG\otimes s$. Integrability follows from
$\bar\partial G=0$. On a cover with $G=\bar\partial\phi_i$,
the holomorphic frames $e_i=e^{2\pi iz\phi_i}s$ satisfy
$e_j=e^{2\pi iz(\phi_j-\phi_i)}e_i$; hence this is exactly the
line represented by the exponential of the Cech class $zg$.
Its unitary Chern connection and curvature are
\[
 \nabla_z=d-2\pi i(zG+\bar z\bar G),\qquad
 F_z=-2\pi i(z\partial G+\bar z\bar\partial\bar G).
\]
Multiplication by $i/(2\pi)$ proves \eqref{inv:eq:C-Pic-map}.
The two independent basis vectors prove injectivity, additivity
as real Lie groups and the assertion about the image.

\item The closed real form $\theta$ is exact since $H^2(X;\R)=0$.
It has zero harmonic projection. De Rham Hodge decomposition,
with $d\mathcal G_d=\mathcal G_dd$, therefore gives
$dd^*\mathcal G_d\theta=\theta$. The Green operator and $d^*$
preserve reality, so $\alpha$ is real and $d\alpha=\theta$.
The $(0,2)$ component yields $\bar\partial G_0=0$.
To prove nonvanishing, an equation $G_0=\bar\partial u$ would give
\[
 \theta=\partial\bar\partial(u-\bar u)
       =2i\partial\bar\partial\operatorname{Im}u.
\]
Semipositivity of $\theta$ would make $\operatorname{Im}u$
plurisubharmonic. The maximum principle on compact connected $X$
would make it constant, contradicting the nonzero pullback of
$\omega_{FS}$ on the regular locus of $f$.

\item Since $\alpha=G_0+\bar G_0$, the Chern formula already proved
sends $\exp([G_0])$ to $[\theta]_{BC}$, which is also
$c_{1,BC}(f^*\OO(1))$. Its injectivity proves the stated
line-bundle equality, with exactly the declared $2\pi i$
exponential and curvature conventions.

\item For the last assertion, the exact sequence of a full scheme
fibre is
\[
 0\longrightarrow f^*\OO(-1)\longrightarrow\OO_X
   \longrightarrow\OO_{X_b}\longrightarrow0.
\]
The cited direct images in \ref{inv:A1.0} give
$f_*f^*\OO(-1)=\OO(-1)$ and
$R^1f_*f^*\OO(-1)=\OO(-1)\oplus\OO(-2)$.
Both terms bounding $H^1(X,f^*\OO(-1))$ in the degree-one Leray
edge sequence vanish, so this group vanishes. Restriction on
$H^1(\OO)$ is injective. On a smooth fibre, the same edge map is
the evaluation of $H^0(B,G)=\C\gamma^{0,1}$ in $G_b$.
This identifies its image as claimed and reproduces exactly
the application of \cite[Theorem~10.5(b), Step~4, p.~88]{AC2026}.
\end{enumerate}
\end{proof}

\Uses{the hypotheses, notation and definitions stated here; \citet[Theorem~10.5]{AC2026}; \CompendiumRef{I}{app:rederive:inv-C3.1}}


The Picard and fibre-restriction statement is due to \citet[Theorem~10.5]{AC2026}. A complete re-derivation is given in \CompendiumRef{I}{app:rederive:inv-C3.1}.


\begin{proposition}[The factor $12$ with a fixed collar orientation]\label{inv:C3.2}\label{wb:lem:bcpr-degree-one-input}
\StatementRecord{inv:C3.2}{proved}{Historical events: Claude Fable 5.2 (Anthropic), effort high, 12 November 2026: targeted proof check. Claude Fable 5.2 (Anthropic), effort high, 16 November 2026: carried as background without fault found. GPT-6 Astra (OpenAI), ChatGPT Pro mode, effort Pro, S6\_09\_PAPERS, 17 November 2026: targeted proof check. Also checked in this writing. Pre-continuation writing events: GPT-6 Astra (OpenAI), ChatGPT Pro mode, effort Pro, S6\_10\_PAPERS, 18 November 2026: computation or formal certificate; targeted proof check. Continuation metadata events: GPT-6 Astra (OpenAI), ChatGPT Pro mode, effort Pro, S6\_10\_PAPERS continuation, 19 November 2026: used and rechecked in this run; status-vocabulary reconciliation; metadata reconciliation only; no new mathematical check. \emph{Pre-continuation recorded proof scope.} Re-derived spectral-sequence rank exhaustion, arrow/square counts, signed collar factor12, Bott–Chern product vanishing using explicit primitives and pairings, metric/Stokes obstruction and Aeppli restriction to the indefinite fibre class. This does not award a reading of the externally cited double-complex source yet. PASS exact reruns c\_double\_complex.py (matrix kernels/images versus arrow counts for h=0..5) and c\_sign\_and\_fibre.py (collar orientation, lift lattice, Hermitian determinant and real-period orientation). Programs and stdout retained under code/I/replay\_invariants.}{the hypotheses, notation and definitions stated here; \cref{inv:eq:C-collar-sum}; \cref{inv:C3.1}}

Use the marked cohomology class $\gamma$ dual to
$\hat\gamma$, the analytic peripheral generators and the lift relation of
[AC26, Section 7.5], and put
$p=12l_0-4l_1-3l_2=\pm1$.  Let $g_\gamma$ be the unique coherent
Leray class whose fibre image is $\gamma^{0,1}$.  Then
\begin{equation}\label{inv:eq:C-factor12}
\exp(12g_\gamma)=f^*\OO_{\Pone}(-p).
\end{equation}
For the standard data $(l_0,l_1,l_2)=(0,1,-1)$, this is
$\exp(12g_\gamma)=f^*\OO(1)$.  For the opposite value $p=1$,
the exponent representing $f^*\OO(1)$ is $-12g_\gamma$.
\end{proposition}


\begin{proof}
\begin{enumerate}[label=\arabic*.,leftmargin=*]
\item We include the comparison that fixes the sign.  A degree-one integral
class $k\gamma$ on a regular fibre extends over the cyclic filling of
order $m_j$ precisely when the affine generator $\hat g_j$ can be
assigned an integer value $n_j$ satisfying
$m_j n_j=k\gamma(v_j)=kl_j$.  Admissibility gives
$\gcd(l_j,m_j)=1$, so this is precisely $m_j\mid k$.
At the cusp $\gamma$ extends, as follows from its actual specialization
to $H^1(W;\Z)$ in [AC26, Theorem B.1 and Proposition 7.12].
The global monodromy invariants have rank one generated by $\gamma$.
It follows, by intersection of these three local images in the marked
local system, that
$H^0(B,R^1f_*\Z)=12\Z\gamma$.  This recovers the degree-one
part of [AC26, Proposition 7.26], without identifying the reduced and
non-reduced fibre: this is topological cohomology, unchanged by
nilpotents, whereas the coherent comparison uses the full fibre.

\item Let $U=\bigsqcup_iD_i$ be small discs around the special points and
let $V$ be their pair-of-pants complement with collars.  Lift
$12\gamma$ to integral classes $\Theta_i$ on $f^{-1}(D_i)$ and
$\Theta_V$ on $f^{-1}(V)$.  Existence over $V$ follows from
$H^2(V;\Z)=0$ and the low-degree Leray sequence; over each disc it
follows from the local images just computed.  Write
$a_j=\Theta_V(\hat\rho_j)$.  On the two finite analytic peripheral generators
the differences $\Theta_i-\Theta_V$ have periods
\[
n_1=4l_1-a_1,\qquad n_2=3l_2-a_2.
\]
At the cusp, the specified lift relation
$\widetilde a_0=t_\mu\hat\rho_2^{-1}\hat\rho_1^{-1}$,
with $\gamma(\mu)=l_0$, and the null-homotopy of
$\widetilde a_0$ in the cusp filling give
$n_0=-12l_0+a_1+a_2$.  Hence
\begin{equation}\label{inv:eq:C-collar-sum}
n_0+n_1+n_2=-p.
\end{equation}
These are the explicitly oriented boundary evaluations in [AC26,
Proposition 7.27, pp. 62--63], not an unspecified sign of a spectral
sequence differential.

\item Represent the lifts by real closed one-forms.  On a collar their
difference is the pullback of a real closed base form $b_i$ with
clockwise period $n_i$, modulo an exact form.  This follows from the
degree-one Leray sequence over the annulus, since the difference has
zero fibre restriction and the annulus has no second cohomology.
Subtract exact terms on smaller collars, extending their primitives
with cutoffs, so the equality of forms holds wherever a partition of
unity changes.  For a real base cutoff $\rho_i$ equal to one toward
$D_i$ and zero toward $V$, the resulting glued real one-form $\alpha$
on $X$ has
\[
d\alpha=f^*\vartheta,\qquad
\vartheta=\sum_i d\rho_i\wedge b_i.
\]
The form $\vartheta$ is a smooth real base two-form; thus
$\bar\partial\alpha^{0,1}=0$.  Its fibre restriction is the
$(0,1)$ image of $12\gamma$, so the coherent Leray edge isomorphism
$H^1(X,\OO_X)\simeq H^0(B,R^1f_*\OO_X)$ identifies
$[\alpha^{0,1}]=12g_\gamma$.

\item In a positively oriented polar coordinate $(r,\varphi)$ on $D_i$,
$b_i=-n_i\,d\varphi/(2\pi)$ modulo an exact form.
Because $\rho_i$ decreases from one to zero as $r$ increases,
\[
\int_{D_i\cap V}d\rho_i\wedge b_i
=-n_i\int_0^1d\rho_i=n_i.
\]
The exact part does not change this integral, since $d\rho_i$ is
supported inside the collar.  Therefore
$\int_B\vartheta=-p$ by \eqref{inv:eq:C-collar-sum}.  \ref{inv:C3.1} identifies
the Bott--Chern class of $\exp(12g_\gamma)$ with $[f^*\vartheta]$.
On $\Pone$ two real $(1,1)$ forms of the same integral differ by
$i\partial\bar\partial$ of a real function (solve the ordinary
scalar Laplace equation).  Hence this class is the Chern class of
$f^*\OO(-p)$, and injectivity in \ref{inv:C3.1} proves
\eqref{inv:eq:C-factor12}.  The program
\texttt{code/c\_sign\_and\_fibre.py} verifies the integer
intersection, \eqref{inv:eq:C-collar-sum}, and the collar integral exactly.
\end{enumerate}
\end{proof}

\begin{proposition}[The complete Bott--Chern algebra, for every $h$]\label{inv:C3.3}\label{wb:thm:bcpr-square-zero-algebra}\label{wb:prop:bcpr-generators-are-bases}
\StatementRecord{inv:C3.3}{proved}{Historical events: Claude Fable 5.2 (Anthropic), effort high, 12 November 2026: targeted proof check. Claude Fable 5.2 (Anthropic), effort high, 16 November 2026: used and rechecked in this run. Claude Fable 5.2 (Anthropic), effort high, 16 November 2026 13:58 UTC: targeted proof check. Claude Fable 5.2 (Anthropic), effort high, 17 November 2026: carried as background without fault found. GPT-6 Astra (OpenAI), ChatGPT Pro mode, effort Pro, S6\_09\_PAPERS, 17 November 2026: targeted proof check. Also checked in this writing. Pre-continuation writing events: GPT-6 Astra (OpenAI), ChatGPT Pro mode, effort Pro, S6\_10\_PAPERS, 18 November 2026: targeted proof check. Continuation metadata events: GPT-6 Astra (OpenAI), ChatGPT Pro mode, effort Pro, S6\_10\_PAPERS continuation, 19 November 2026: used and rechecked in this run; status-vocabulary reconciliation; metadata reconciliation only; no new mathematical check. \emph{Pre-continuation recorded proof scope.} Re-derived spectral-sequence rank exhaustion, arrow/square counts, signed collar factor12, Bott–Chern product vanishing using explicit primitives and pairings, metric/Stokes obstruction and Aeppli restriction to the indefinite fibre class. This does not award a reading of the externally cited double-complex source yet.}{the hypotheses, notation and definitions stated here; \cref{inv:C2.1}; \CompendiumRef{I}{wb:eq:bcpr-generators}; \cref{inv:C1.2}; \cref{inv:C3.1}; \cref{inv:C1.1}}

The positive-total-degree ideal $N\subset H^{\bullet,\bullet}_{BC}(X)$
satisfies
\[
H^{\bullet,\bullet}_{BC}(X)=\C\cdot1\oplus N,\qquad N^2=0.
\]
Its nonzero bidegrees and dimensions are those of \ref{inv:C2.1}.  Thus the
square-zero conclusion is independent of the evaluated value of $h$.
For $h=1$, the nine classes specified in \CompendiumRef{I}{wb:eq:bcpr-generators}, together with the unit, give bases of every nonzero Bott--Chern group.
\end{proposition}


\begin{proof}
\begin{enumerate}[label=\arabic*.,leftmargin=*]
\item Take global $\bar\partial$-closed representatives $G,Q_j,V_j,E$
of the classes $g,q_j,v_j,e$ of \ref{inv:C1.2}.  By the arrow computation in
\ref{inv:C2.1}, bases for the positive Bott--Chern groups consist of
\[
\partial G,\ \bar\partial\bar G;\qquad
\partial Q_j,\ \bar\partial\bar Q_j;\qquad
\partial V_j,\ \bar\partial\bar V_j;\qquad
\partial E,\ \bar\partial\bar E;\qquad \nu,
\]
where the pairs have bidegrees $(1,1)$, $(2,1),(1,2)$,
$(2,2)$, $(3,2),(2,3)$, respectively, and
$\int_X\nu=1$.  In particular the natural map
$H^{3,2}_{BC}\to H^{3,2}_{\bar\partial}$ is an isomorphism:
on its unique horizontal arrow the closed target maps nontrivially
to column cohomology.

\item Write $a=\partial G$, $b=\bar\partial\bar G$ and
$c_j=\partial Q_j$.  The graded Leibniz rule gives
\[
\partial\bar\partial(G\wedge\bar G)=-a\wedge b,
\quad
\partial\bar\partial(G\wedge\bar Q_j)
=-a\wedge\bar c_j,
\quad
\partial\bar\partial(\bar G\wedge Q_j)=b\wedge c_j.
\]
Consequently $\mathsf A\mathsf B=0$,
$\mathsf A\overline{\mathsf C_j}=0$, and
$\mathsf B\mathsf C_j=0$.

\item Let $f^*\OO(1)=\exp(\lambda g)$, where $\lambda\ne0$.
Its curvature may be chosen to be the pullback of a two-form on $B$,
whose square vanishes pointwise.  By \ref{inv:C3.1},
$\lambda^2\mathsf A^2+\bar\lambda^2\mathsf B^2=0$.
The Dolbeault image of $\mathsf B^2$ is zero.  Hence
$[a^2]_{\bar\partial}=0$.  But
$a^2=\partial(G\wedge a)$, and \ref{inv:C1.1} says that
$d_1:H^{1,2}\to H^{2,2}$ is injective.  Thus
$[G\wedge a]_{\bar\partial}=0$.  Choose a smooth $(1,1)$ form
$\Phi$ satisfying $G\wedge a=\bar\partial\Phi$.  Applying
$\partial$ proves $a^2=\partial\bar\partial\Phi$.
Conjugation gives $\mathsf B^2=0$ as well.

\item To treat $\mathsf A\mathsf C_j$, pair its Dolbeault image with the
generator $g\in H^{0,1}_{\bar\partial}$.  Since
$\bar\partial c_j=0$ and $G\wedge a=\bar\partial\Phi$,
\[
\int_X a\wedge c_j\wedge G
=-\int_X\bar\partial(\Phi\wedge c_j)=0.
\]
Dolbeault--Serre duality and the isomorphism
$H^{3,2}_{BC}\to H^{3,2}_{\bar\partial}$ give
$\mathsf A\mathsf C_j=0$.  Complex conjugation gives
$\mathsf B\overline{\mathsf C_j}=0$.

\item Every remaining product whose bidegree lies in $[0,3]^2$ lands in
$(3,3)$: it is a product of a $(1,1)$ class with a $(2,2)$ class,
or of a $(2,1)$ class with a $(1,2)$ class.  Each non-top generator
displayed above is $d$-exact, since its chosen primitive is
$\bar\partial$-closed, or is the conjugate of such a form.
The products therefore have integral zero.  Integration is an
isomorphism $H^{3,3}_{BC}\to\C$ by \ref{inv:C2.1} and the class of $\nu$,
so all these products vanish.  All other products exceed the complex
dimension in at least one bidegree.  This exhausts the possibilities.
\end{enumerate}
\end{proof}

\begin{proposition}[Non-K\"ahler degrees and the precise degreewise statement]\label{inv:C3.4}
\StatementRecord{inv:C3.4}{proved}{Historical events: Claude Fable 5.2 (Anthropic), effort high, 16 November 2026: carried as background without fault found. Claude Opus 5.5 (Anthropic), effort max, sweep 1, scope s6, 18 November 2026: computation or formal certificate. GPT-6 Astra (OpenAI), ChatGPT Pro mode, effort Pro, S6\_09\_PAPERS, 17 November 2026: targeted proof check. Also checked in this writing. Pre-continuation writing events: GPT-6 Astra (OpenAI), ChatGPT Pro mode, effort Pro, S6\_10\_PAPERS, 18 November 2026: targeted proof check. Continuation metadata events: GPT-6 Astra (OpenAI), ChatGPT Pro mode, effort Pro, S6\_10\_PAPERS continuation, 19 November 2026: used and rechecked in this run; status-vocabulary reconciliation; metadata reconciliation only; no new mathematical check. \emph{Pre-continuation recorded proof scope.} Re-derived spectral-sequence rank exhaustion, arrow/square counts, signed collar factor12, Bott–Chern product vanishing using explicit primitives and pairings, metric/Stokes obstruction and Aeppli restriction to the indefinite fibre class. This does not award a reading of the externally cited double-complex source yet.}{the hypotheses, notation and definitions stated here; \cref{inv:C2.1}}

For $k=0,\ldots,6$ the non-K\"ahler degrees are
\begin{equation}\label{inv:eq:C-delta}
(\Delta^k)=(0,2,2+2h,4h,2+2h,2,0).
\end{equation}
If $h\ge1$, the two maps
$\bigoplus_{p+q=k}H^{p,q}_{BC}\to H^k_{\mathrm{dR}}(X;\C)
\to\bigoplus_{p+q=k}H^{p,q}_A$
are both isomorphisms exactly for $k=0,6$.
For the weaker assertion that a $d$-exact pure $d$-closed form of
degree $k$ is $\partial\bar\partial$-exact, the valid degrees are
$k=0,1,6$, and also $k=3,4$ if $h=0$.
\end{proposition}


\begin{proof}
\begin{enumerate}[label=\arabic*.,leftmargin=*]
\item The totals from \ref{inv:C2.1} are
\[
(h^k_{BC})=(1,0,2,2h,2h,2,1),\qquad
(h^k_A)=(1,2,2h,2h,2,0,1).
\]
Subtracting twice the Betti numbers gives \eqref{inv:eq:C-delta}.
On the two endpoint dots both comparison maps are isomorphisms;
all intermediate de Rham groups vanish.  The first assertion follows
by inspecting when both outside groups vanish.  The weaker assertion
is exactly injectivity of the first comparison map.  In intermediate
degree this is equivalent to vanishing of the Bott--Chern group alone,
which gives the stated list.  In degree zero there is no nonzero
exact form, and in degree six integration identifies Bott--Chern and
de Rham cohomology.  This distinguishes the two meanings of a
``$\partial\bar\partial$-lemma in a degree''; a vanishing Bott--Chern
group in degree one does not make $\Delta^1$ vanish.
\end{enumerate}
\end{proof}

\begin{proposition}[Metric obstructions that do not require $h$]\label{inv:C3.5}
\StatementRecord{inv:C3.5}{proved}{Historical events: Claude Fable 5.2 (Anthropic), effort high, 16 November 2026: carried as background without fault found; used and rechecked in this run. GPT-6 Astra (OpenAI), ChatGPT Pro mode, effort Pro, S6\_09\_PAPERS, 17 November 2026: targeted proof check. Also checked in this writing. Pre-continuation writing events: GPT-6 Astra (OpenAI), ChatGPT Pro mode, effort Pro, S6\_10\_PAPERS, 18 November 2026: targeted proof check. Continuation metadata events: GPT-6 Astra (OpenAI), ChatGPT Pro mode, effort Pro, S6\_10\_PAPERS continuation, 19 November 2026: used and rechecked in this run; status-vocabulary reconciliation; metadata reconciliation only; no new mathematical check. \emph{Pre-continuation recorded proof scope.} Re-derived spectral-sequence rank exhaustion, arrow/square counts, signed collar factor12, Bott–Chern product vanishing using explicit primitives and pairings, metric/Stokes obstruction and Aeppli restriction to the indefinite fibre class. This does not award a reading of the externally cited double-complex source yet.}{the hypotheses, notation and definitions stated here; the explicit constructions and calculations in the proof}

The threefold $X$ admits no balanced, strongly Gauduchon, K\"ahler,
locally conformally K\"ahler, or Hermitian-symplectic metric.
\end{proposition}


\begin{proof}
\begin{enumerate}[label=\arabic*.,leftmargin=*]
\item Let $\theta=f^*\omega_{FS}$.  It is a nonzero semipositive exact
real $(1,1)$ form.  If $\omega$ were balanced, the real four-form
$\omega^2$ would be closed and
$\int_X\theta\wedge\omega^2>0$, because the integrand is positive
on the regular locus.  Exactness of $\theta$ and Stokes' theorem
give zero, a contradiction.

\item More generally, if $\omega$ were strongly Gauduchon, write
$\partial\omega^2=\bar\partial u$ with $u\in A^{3,1}$.
The real four-form $\Omega=\omega^2-u-\bar u$ is closed.
Bidegrees give $\theta\wedge\Omega=\theta\wedge\omega^2$,
so the same argument applies.  K\"ahler metrics are balanced.
If $d\omega=\ell\wedge\omega$ with a closed real Lee form
$\ell$, then $b_1=0$ gives $\ell=dv$, and
$d(e^{-v}\omega)=0$, excluding locally conformally K\"ahler metrics.
Finally any symplectic two-form on $X$ is exact because $b_2=0$.
Its cube has nonzero integral for its symplectic orientation, whereas
an exact cube has zero integral.  Thus $X$ is not symplectic and in
particular admits no Hermitian-symplectic metric.
\end{enumerate}
\end{proof}

\begin{cited}[Gauduchon metrics with specified Chern--Ricci form]\label{inv:C3.6}
\StatementRecord{inv:C3.6}{proved}{Historical events: Claude Fable 5.2 (Anthropic), effort high, 16 November 2026: carried as background without fault found. GPT-6 Astra (OpenAI), ChatGPT Pro mode, effort Pro, S6\_09\_PAPERS, 17 November 2026: targeted proof check. Also checked in this writing. Pre-continuation writing events: GPT-6 Astra (OpenAI), ChatGPT Pro mode, effort Pro, S6\_10\_PAPERS, 18 November 2026 04:30:14.617321+00:00: used and rechecked in this run. Continuation metadata events: GPT-6 Astra (OpenAI), ChatGPT Pro mode, effort Pro, S6\_10\_PAPERS continuation, 19 November 2026: used and rechecked in this run; status-vocabulary reconciliation; metadata reconciliation only; no new mathematical check. \emph{Pre-continuation recorded proof scope.} Current-writing verification remains in progress. The earlier checks named above retain their recorded scope.}{the hypotheses, notation and definitions stated here; \citet[Theorem~2]{Gauduchon1985}; \citet[Theorem~1.3]{STW2017}; \CompendiumRef{I}{app:rederive:inv-C3.6}}

There exists a Gauduchon metric on $X$.  More precisely, for every
Gauduchon metric $\omega_0$ there is another Gauduchon metric
$\omega$ with $[\omega^2]_A=[\omega_0^2]_A$ and
\begin{equation}\label{inv:eq:C-Ricci}
\operatorname{Ric}^{\mathrm{Ch}}(\omega)=\pi f^*\omega_{FS}.
\end{equation}
This Ricci form is semipositive and has rank one at regular points
of $f$.  No Chern--Ricci-flat Hermitian metric exists on $X$.
\end{cited}
The expanded source derivation is retained in the proof compendium, Paper~I, label \nolinkurl{inv:C3.6}.




This is the application of \citet[Theorem~2]{Gauduchon1985} and \citet[Theorem~1.3]{STW2017} with the computed Bott--Chern class. A complete re-derivation is given in \CompendiumRef{I}{app:rederive:inv-C3.6}.


\begin{proposition}[The fibre obstruction to SKT and astheno-K\"ahler metrics]\label{inv:C3.7}
\StatementRecord{inv:C3.7}{proved}{Historical events: Claude Fable 5.2 (Anthropic), effort high, 12 November 2026: targeted proof check. Claude Fable 5.2 (Anthropic), effort high, 16 November 2026: carried as background without fault found; used and rechecked in this run. GPT-6 Astra (OpenAI), ChatGPT Pro mode, effort Pro, S6\_09\_PAPERS, 17 November 2026: targeted proof check. Also checked in this writing. Pre-continuation writing events: GPT-6 Astra (OpenAI), ChatGPT Pro mode, effort Pro, S6\_10\_PAPERS, 18 November 2026: computation or formal certificate; targeted proof check. GPT-6 Astra (OpenAI), ChatGPT Pro mode, effort Pro, S6\_10\_PAPERS, 18 November 2026 04:30:14.617321+00:00: used and rechecked in this run. Continuation metadata events: GPT-6 Astra (OpenAI), ChatGPT Pro mode, effort Pro, S6\_10\_PAPERS continuation, 19 November 2026: used and rechecked in this run; status-vocabulary reconciliation; metadata reconciliation only; no new mathematical check. \emph{Pre-continuation recorded proof scope.} Re-derived spectral-sequence rank exhaustion, arrow/square counts, signed collar factor12, Bott–Chern product vanishing using explicit primitives and pairings, metric/Stokes obstruction and Aeppli restriction to the indefinite fibre class. This does not award a reading of the externally cited double-complex source yet. PASS exact reruns c\_double\_complex.py (matrix kernels/images versus arrow counts for h=0..5) and c\_sign\_and\_fibre.py (collar orientation, lift lattice, Hermitian determinant and real-period orientation). Programs and stdout retained under code/I/replay\_invariants.}{the hypotheses, notation and definitions stated here; \cref{inv:A3.2}; \cref{inv:C1.2}; \cref{inv:C2.1}; \cref{inv:eq:C-eta-negative}}

The threefold $X$ admits no pluriclosed (SKT) metric and no
astheno-K\"ahler metric.  The geometric input is the value $h=1$
and the global Dolbeault class $q$ in \ref{inv:A3.2}, whose restriction to a
regular fibre is $\eta=u\wedge w+6\gamma\wedge\delta$.
More generally, it is enough that the real image of
$H_A^{1,1}(X;\R)\to H_A^{1,1}(F;\R)$ be contained in
$\R\eta$ for one regular fibre $F$.
\end{proposition}


\begin{proof}
\begin{enumerate}[label=\arabic*.,leftmargin=*]
\item We first compute the sign that matters.  Put
$y=\operatorname{Im}\tau>0$, $m=\operatorname{Im}\mu$, and
$\Delta=\operatorname{Im}\beta-6m^2/y<0$.
For the marked period rows
$\sigma_1=6\mu\gamma+\tau u+w$ and
$\sigma_2=\beta\gamma+\mu u+\delta$ we have
$\eta=u\wedge\sigma_1+6\gamma\wedge\sigma_2$.
This proves that the real class $\eta$ has type $(1,1)$.
The determinant of the four real period columns, with rows ordered
as $\operatorname{Re}\zeta_1,\operatorname{Im}\zeta_1,
\operatorname{Re}\zeta_2,\operatorname{Im}\zeta_2$, is
$y\Delta<0$.  Since $\eta^2=12\gamma\wedge u\wedge w\wedge\delta$,
\begin{equation}\label{inv:eq:C-eta-negative}
\int_{F,\mathrm{or}_{\C}}\eta^2=-12.
\end{equation}
Equivalently the invariant representative of $\eta$ has Hermitian
coefficient matrix
\[
\frac1{y\Delta}
\begin{pmatrix}\operatorname{Im}\beta&-6m\\-6m&6y\end{pmatrix},
\qquad \det=\frac6{y\Delta}<0,
\]
so it has one positive and one negative eigenvalue.  The exact
symbolic determinant and an exact rational example are checked in
\texttt{code/c\_sign\_and\_fibre.py} by both eigenvalues and
diagonal elimination.

\item The explicit invariant fibre representative of $\gamma^{0,1}$ is
\[
\Gamma=\frac{i}{2\Delta}
\left(d\bar\zeta_2-\frac myd\bar\zeta_1\right).
\]
Writing $Q_F$ for the invariant $(1,1)$ representative of $\eta$,
direct exterior multiplication gives
\[
Q_F\wedge\Gamma=
\frac1{4y\Delta}
d\zeta_1\wedge d\bar\zeta_1\wedge d\bar\zeta_2,
\quad
\int_FQ_F\wedge\Gamma\wedge d\zeta_2=-1.
\]
This also proves the nonvanishing used in \ref{inv:C1.2}.

\item When $h=1$, \ref{inv:C2.1} identifies $H_A^{1,1}(X)$ with the sum of the
one horizontal source and its vertical conjugate.  A representative
$Q$ of $q$ is nonzero in the horizontal source, since its fibre
restriction is nonzero whereas $d_1g$ restricts to zero.  Thus
$[Q]_A,[\bar Q]_A$ form a basis.  Their restrictions to $F$ both
equal the real class $\eta$.  The real image of the restriction map
is therefore $\R\eta$.

\item If $\omega$ were a pluriclosed positive real $(1,1)$ form, its
restriction to $F$ would be pluriclosed and positive.  Average this
restriction under translations of the compact complex torus $F$.
The average $\omega_{\mathrm{av}}$ is invariant, closed, and
strictly positive.  Translation preserves the Aeppli class: on a
torus the $\partial\bar\partial$-lemma identifies Aeppli
cohomology with the invariant $(p,q)$ forms, so the average is the
unique invariant representative of that class.  This identification
can also be checked Fourier mode by Fourier mode: each nonzero mode
is a Koszul square for two independent type components of its real
nonzero covector and contributes zero Aeppli cohomology.  Thus
$[\omega_{\mathrm{av}}]=r\eta$ for some real $r$.
But positivity gives $\int_F\omega_{\mathrm{av}}^2>0$, whereas
\eqref{inv:eq:C-eta-negative} gives $-12r^2\le0$, a contradiction.
In dimension three the astheno-K\"ahler equation
$\partial\bar\partial\omega^{n-2}=0$ is exactly the SKT equation.
\end{enumerate}
\end{proof}





\section{Twisted tangent cohomology}
\subsection{All integral twists and the generic nonintegral twist}

\begin{proposition}[Integral Picard twists, including their torsion]\label{inv:D4.1}
\StatementRecord{inv:D4.1}{proved}{Historical events: Claude Fable 5.2 (Anthropic), effort high, 16 November 2026: carried as background without fault found. Claude Fable 5.2 (Anthropic), effort high, 17 November 2026: carried as background without fault found. GPT-6 Astra (OpenAI), ChatGPT Pro mode, effort Pro, S6\_09\_PAPERS, 17 November 2026: targeted proof check. Also checked in this writing. Pre-continuation writing events: GPT-6 Astra (OpenAI), ChatGPT Pro mode, effort Pro, S6\_10\_PAPERS, 18 November 2026: computation or formal certificate; targeted proof check. Continuation metadata events: GPT-6 Astra (OpenAI), ChatGPT Pro mode, effort Pro, S6\_10\_PAPERS continuation, 19 November 2026: used and rechecked in this run; status-vocabulary reconciliation; metadata reconciliation only; no new mathematical check. \emph{Pre-continuation recorded proof scope.} Read full root-base, logarithmic tangent, direct-image and twisted-Dolbeault proofs. Re-derived coarse degree function, local normalization quotient, resonance equation, Fourier contraction comparison, joint order-three coefficient 123ik/(64pi), stated local modules, and jump-locus finiteness. Exact Koszul and twist-program computations retain their separately recorded replay scopes. PASS: unchanged d\_exact.py and d\_koszul.py replays. Two independent weighted cohomology counts for 481 integers, tangent Euler characteristics, root-character invariants, polynomial Smith forms and joint order-three coefficient 123i/(64pi). These are finite algebra checks; analytic comparison is in the targeted proof event.}{the hypotheses, notation and definitions stated here; the explicit constructions and calculations in the proof}

For $M_n=\OO_X(n(S_1-S_2))$, define
$u(n)=\max\{d(n)+1,0\}$, $v(n)=\max\{-d(n)-1,0\}$ and
$\epsilon_n=0$ if $3\mid n$, $\epsilon_n=1$ otherwise. Then
\[
 f_*(T_X\otimes M_n)=\OO(d(n))\oplus\OO(d(n-1)),
\]
\[
 R^1f_*(T_X\otimes M_n)=
 \OO(d(n))\oplus\OO(d(n-1))\oplus\OO(d(n+2)),           \tag{D.12}
\]
\[
 R^2f_*(T_X\otimes M_n)=
 \OO(d(n+6))\oplus\C_{p_1}^{\epsilon_n}.
\]
In particular its cusp stalk is free of rank one. The four global
dimensions, for every integer $n$, are
\begin{align*}
 h^0&=u(n)+u(n-1),\\
 h^1&=v(n)+v(n-1)+u(n)+u(n-1)+u(n+2),\\
 h^2&=v(n)+v(n-1)+v(n+2)+u(n+6)+\epsilon_n,\\
 h^3&=v(n+6).                                          \tag{D.13}
\end{align*}
\end{proposition}


\begin{proof}
\begin{enumerate}[label=\arabic*.,leftmargin=*]
\item The defining roots give $\bar f^*R=\OO_X(S_1-S_2)$, so the projection
formula applies to (D.11), with its actual coherent torsion summand.
Tensoring (D.8) changes its generator character from $-1$ to $n-1$.
There is exactly one invariant coefficient among exponents $0,1$ if
$n-1$ is $0$ or $1$ modulo three, and none if it is $2$. Multiplication
by the coarse parameter $t=s^3$ is zero on this module. Its invariant
direct image is therefore exactly the displayed skyscraper, of
length $\epsilon_n$, not a higher-order thickening. The two other
special places contribute no torsion. Formula (D.1) proves (D.12),
and the two-column Leray sequence proves (D.13). The code evaluates
these formulas for $-240\le n\le240$ using two independent integer
cohomology counts; their alternating sums are all one. That last
check is supplementary to the sheaf proof.
\end{enumerate}
\end{proof}

\begin{proposition}[The cusp length and the generic twist]\label{inv:D4.2}
\StatementRecord{inv:D4.2}{proved}{Historical events: Claude Fable 5.2 (Anthropic), effort high, 16 November 2026: carried as background without fault found. GPT-6 Astra (OpenAI), ChatGPT Pro mode, effort Pro, S6\_09\_PAPERS, 17 November 2026: targeted proof check. Also checked in this writing. Pre-continuation writing events: GPT-6 Astra (OpenAI), ChatGPT Pro mode, effort Pro, S6\_10\_PAPERS, 18 November 2026: targeted proof check. Continuation metadata events: GPT-6 Astra (OpenAI), ChatGPT Pro mode, effort Pro, S6\_10\_PAPERS continuation, 19 November 2026: used and rechecked in this run; status-vocabulary reconciliation; metadata reconciliation only; no new mathematical check. \emph{Pre-continuation recorded proof scope.} Read full root-base, logarithmic tangent, direct-image and twisted-Dolbeault proofs. Re-derived coarse degree function, local normalization quotient, resonance equation, Fourier contraction comparison, joint order-three coefficient 123ik/(64pi), stated local modules, and jump-locus finiteness. Exact Koszul and twist-program computations retain their separately recorded replay scopes.}{the hypotheses, notation and definitions stated here; \cref{inv:C3.2}; \cref{inv:D4.3}; \cref{inv:D1.2}}

Use the coordinate $M_z=\exp(z\mathfrak g)$ of \ref{inv:C3.2}. For every
$z\notin\Z$ the cusp contribution is exactly
\[
 (R^0f_*(T_X\otimes M_z))_{p_0}=0,\quad
 (R^1f_*(T_X\otimes M_z))_{p_0}=0,\quad
 (R^2f_*(T_X\otimes M_z))_{p_0}=\C\{q\}/(q).           \tag{D.14}
\]
Outside the countable exceptional set explicitly defined in \ref{inv:D4.3},
these are the only nonzero direct-image stalks. Hence the generic
twist has $(h^0,h^1,h^2,h^3)=(0,0,1,0)$, and the cusp nilpotence exponent $r$ equals one there.
\end{proposition}


\begin{proof}
\begin{enumerate}[label=\arabic*.,leftmargin=*]
\item The map $H^1(W,\C)\to H^1(W,\OO_W)$ is an isomorphism by the
normalization complex in \ref{inv:D1.2}. The restriction of $M_z$ has trivial
pullback to $P$ and descent character
$(\exp(2\pi iz),1)$ on the two generators of the dual torus. The
normalization complex with these coefficients is the cellular
cochain complex of that rank-one local system. It may be computed
on the product cell structure as
\[
 \C\xrightarrow{(a-1,0)^t}\C^2
 \xrightarrow{(0,a-1)}\C,
 \qquad a=\exp(2\pi iz).
\]
For $z\notin\Z$ this is acyclic. All normalization components remain
rational with trivial pulled-back twist, so this calculates the
actual coherent cohomology $H^q(W,M_z|_W)$, not just a topological
analogy. Analytic semicontinuity and base change make
$R^q\bar f_*M_z$ zero on a neighbourhood of $p_0$ for all $q$.
The projection formula and (D.3) imply the same vanishing for
$\mathcal T\otimes M_z$. Tensoring (D.10) remains exact, and
$\nu^*M_z=\OO_P$ because $P$ is rational and the topological first
Chern class of $M_z$ is zero. Thus the degree-two direct image is
exactly the reduced skyscraper $H^2(P,K_P)\otimes\C_{p_0}$.
This proves (D.14), including its annihilator and not only its residue
dimension.

\item If the twist is nontrivial on every smooth fibre and on both finite
root-cover central tori, all cohomology of its fibre restriction
vanishes at those places. Indeed a nontrivial topologically trivial
line bundle on a complex torus has zero cohomology: in a unitary flat
presentation its Dolbeault Laplacian has no zero Fourier mode.
Tensoring the tangent extension by that line preserves acyclicity.
Base change on the root charts and exact invariants then give zero
direct images away from $p_0$. \ref{inv:D4.3} proves that the excluded twists
form a countable set. Leray now gives the asserted four dimensions.
\end{enumerate}
\end{proof}

\begin{proposition}[The exceptional nonintegral twists cannot be omitted]\label{inv:D4.3}
\StatementRecord{inv:D4.3}{proved}{Historical events: GPT-6 Astra (OpenAI), ChatGPT Pro mode, effort Pro, S6\_09\_PAPERS, 17 November 2026: targeted proof check. Also checked in this writing. Pre-continuation writing events: GPT-6 Astra (OpenAI), ChatGPT Pro mode, effort Pro, S6\_10\_PAPERS, 18 November 2026: targeted proof check. Continuation metadata events: GPT-6 Astra (OpenAI), ChatGPT Pro mode, effort Pro, S6\_10\_PAPERS continuation, 19 November 2026: used and rechecked in this run; status-vocabulary reconciliation; metadata reconciliation only; no new mathematical check. \emph{Pre-continuation recorded proof scope.} Read full root-base, logarithmic tangent, direct-image and twisted-Dolbeault proofs. Re-derived coarse degree function, local normalization quotient, resonance equation, Fourier contraction comparison, joint order-three coefficient 123ik/(64pi), stated local modules, and jump-locus finiteness. Exact Koszul and twist-program computations retain their separately recorded replay scopes.}{the hypotheses, notation and definitions stated here; \cref{inv:D4.2}}

On a regular marked fibre, $M_z$ is trivial if and only if there are
integers $n,p,q,r$ such that
\[
 q\tau(t)+r\mu(t)=p,\qquad
 z=n-6q\mu(t)-r\beta_c(t).                             \tag{D.15}
\]
For $z\notin\Z$ these resonances occur in a countable set of parameter
values $z$; every resonant $z$ is nonreal. On the finite root covers
the exceptional lattices are
\[
 \Lambda_1=\Z+(3\beta_c(z_1)+4-2\rho)\Z,\qquad
 \Lambda_2=\Z+(2\beta_c(z_2)+3-3i)\Z.                  \tag{D.16}
\]
Take the union of $\Z$, these two lattices, and all values from (D.15)
over regular fibres; this is the countable exceptional set used in
\ref{inv:D4.2}. The set is larger than $\Z$: the regular value $t=32/5$
supplies nonintegral twists with $h^1(X,T_X\otimes M_z)\ge2$.
\end{proposition}


\begin{proof}
\begin{enumerate}[label=\arabic*.,leftmargin=*]
\item The exponential sequence on a torus identifies trivial line bundles
with integral constant cocycles modulo the two holomorphic period
rows. Write
$z\gamma=(n,p,q,r)+a(6\mu,\tau,1,0)+b(\beta_c,\mu,0,1)$.
The third and fourth coordinates give $a=-q,b=-r$; the first two
give exactly (D.15). If $r=0$, the positivity of $\operatorname{Im}\tau$
forces $q=p=0$, leaving $z=n$. Otherwise the imaginary parts of
(D.15) give
\[
 \operatorname{Im}\mu=-\frac qr\operatorname{Im}\tau,
 \qquad\operatorname{Im}z=-rD\ne0.                    \tag{D.17}
\]
For fixed $(p,q,r)$ with $(q,r)\ne(0,0)$, the first equation of
(D.15) has discrete zeros, since $1,\tau,\mu$ are linearly
independent. A discrete subset of this second-countable Riemann
surface is countable. There are countably many integer tuples and
markings, so the union of the resulting $z$ is countable.

\item At $z_1$, insert $\tau=\rho$, $\mu=(2-\rho)/3$ in the first
equation. Its solutions have $(q,r,p)=(k,3k,2k)$.
At $z_2$ they have $(q,r,p)=(k,2k,k)$.
The second equation gives (D.16). These are discrete lattices because
their nonintegral generators have imaginary parts $3D$ and $2D$.

\item For the explicit regular resonance, use the section
$P=(0,\sqrt2t^2(t-1))$ of the stated elliptic surface. The
chord-and-tangent law gives
\[
 x(2P)=9t^2/8,\qquad
 y(2P)=t^2(32-5t)/(16\sqrt2).
\]
At $t=32/5$, the discriminant
$1728t^8(t-1)^3$ is nonzero, $P$ is not two-torsion, and $2P$ is a
nonzero two-torsion point. Thus $P$ has exact order four. The period
identification in Part A supplies a nonzero relation
$q\tau+4\mu=p$ with integers $p,q$. Equation (D.15), with $r=4$,
defines an actual twist for which the regular torus restriction is
trivial; (D.17) shows that it is nonintegral. The period derivative
there has nonzero determinant, so the tangent restriction has
$h^0(T_X|_F)=2$. The degree-zero base-change exact sequence gives
\[
 H^0(F,(T_X\otimes M_z)|_F)
   \simeq (R^1f_*(T_X\otimes M_z))_t[t],
\]
because $f_*(T_X\otimes M_z)=0$ for a twist nontrivial on the general
fibre. Its right side has dimension two. The degree-one direct image
is a finite-support sheaf, so its global sections have dimension at
least two, proving the bound on $h^1$. The exact elliptic calculation
is checked by both duplication and substitution in
\texttt{d\_exact.py}; its output records the nonzero rational
discriminant. The choice of integers $p,q$ depends on the marked
elliptic logarithm, and no unverified numerical marking is substituted.
\end{enumerate}
\end{proof}

\begin{proposition}[Local complexes and every noncusp torsion length]\label{inv:D4.4}
\StatementRecord{inv:D4.4}{proved}{Historical events: GPT-6 Astra (OpenAI), ChatGPT Pro mode, effort Pro, S6\_09\_PAPERS, 17 November 2026: targeted proof check. Also checked in this writing. Pre-continuation writing events: GPT-6 Astra (OpenAI), ChatGPT Pro mode, effort Pro, S6\_10\_PAPERS, 18 November 2026: computation or formal certificate; targeted proof check. Continuation metadata events: GPT-6 Astra (OpenAI), ChatGPT Pro mode, effort Pro, S6\_10\_PAPERS continuation, 19 November 2026: used and rechecked in this run; status-vocabulary reconciliation; metadata reconciliation only; no new mathematical check. \emph{Pre-continuation recorded proof scope.} Read full root-base, logarithmic tangent, direct-image and twisted-Dolbeault proofs. Re-derived coarse degree function, local normalization quotient, resonance equation, Fourier contraction comparison, joint order-three coefficient 123ik/(64pi), stated local modules, and jump-locus finiteness. Exact Koszul and twist-program computations retain their separately recorded replay scopes. PASS: unchanged d\_exact.py and d\_koszul.py replays. Two independent weighted cohomology counts for 481 integers, tangent Euler characteristics, root-character invariants, polynomial Smith forms and joint order-three coefficient 123i/(64pi). These are finite algebra checks; analytic comparison is in the targeted proof event.}{the hypotheses, notation and definitions stated here; \cref{inv:D4.1}; \cref{inv:A2.3}}

Let $z\notin\Z$. At a regular resonance $t_*$ in (D.15), with
$A=\C\{t-t_*\}$, the two stalks are
\[
 (R^1f_*(T_X\otimes M_z))_{t_*}=(A/(t-t_*))^2,
 \qquad (R^2f_*(T_X\otimes M_z))_{t_*}=A/((t-t_*)^2).    \tag{D.18}
\]
At $p_1$, write $z=n+k(3\beta_1+4-2\rho)$, $k\ne0$, when this is
possible. Both coarse stalks are killed by $t$ and have length
$1+\epsilon_n$, with $\epsilon_n$ as in \ref{inv:D4.1}.
At $p_2$, write $z=n+k(2\beta_2+3-3i)$, $k\ne0$, when possible.
Both coarse stalks are killed by $t-1$ and have length one for
$n\equiv1,2\pmod4$ and length zero for $n\equiv0,3\pmod4$.
If the relevant lattice representation does not exist, both stalks
at that finite fibre are zero.
\end{proposition}


\begin{proof}
\begin{enumerate}[label=\arabic*.,leftmargin=*]
\item We first give the comparison that turns the local calculation into
a computation of the sheaves. On a simply connected disc of smooth
torus fibres, including a finite root cover, retain the real marking
and the holomorphic period matrix. Translation by $\Pi(s)v$, for
fixed $v\in\R^4$, is holomorphic on the total family. After
subtracting the integer cocycle of a resonance, the line bundle has
constant lattice multipliers. Multiplying the lifted real translation
by its constant factor of automorphy makes the lattice act as the
identity and gives a genuine compact real torus action on the line
bundle. The same action acts on the tangent bundle by differentiation.
It commutes with the full Dolbeault differential.

\item Averaging over this compact torus projects the relative Dolbeault
complex onto its invariant part. Choose the lift so that the central
trivial line has trivial character. On every nonzero Fourier character,
the relative Dolbeault symbol is exterior multiplication by a
nonzero covector; contraction by a dual covector is a contracting
homotopy. After shrinking the disc, these nonzero covectors have a
uniform positive lower bound for bounded Fourier frequency and grow
linearly for unbounded frequency, because the real period matrices
and their inverses stay bounded. The homotopies therefore preserve
smooth Fourier series, with all derivatives, and are uniform on
compact subsets of the disc. The tangent extension is triangular
with respect to its vertical and horizontal parts; its nilpotent
off-diagonal term is treated by the finite homological perturbation
series, so every nonzero Fourier complex is still contractible.
Thus averaging is a quasi-isomorphism, not merely a fibrewise
dimension comparison.

\item The chain-level comparison can be written without an unspecified
holomorphic gauge. Write the normalized period matrix as $[Z(s),I]$
and use the real coordinates $\zeta=Z(s)x+y$, with
$x,y\in\R^2/\Z^2$. Set
\[
 G_j=(dx_j)^{0,1},\qquad e_i=\partial_{\zeta_i},\qquad
 h=\partial_s+\sum_i(Z'(s)x)_i e_i.
\]
These are invariant under the real translations. Since each $dx_j$
is closed on the total real product, $\bar\partial G_j=0$.
The forms $G_j$ frame the relative invariant $(0,1)$ forms because
$\operatorname{Im}Z$ is invertible, and
\[
 \bar\partial h=\sum_{i,j}Z'_{ij}(s)e_iG_j,
 \qquad (dy)^{0,1}=-Z(s)(dx)^{0,1}.
\]
For the smooth line frame defined by its flat character, the last
identity gives $\bar\partial\ell=2\pi i(a_1G_1+a_2G_2)\ell$,
with $a_1,a_2$ in (D.20); there is no additional base term.
All displayed coefficients are holomorphic in $s$. The full
invariant complex is thus the ordinary base Dolbeault resolution
of the following finite complex of holomorphic bundles:
\[
 C^j=(V\oplus L)\otimes\bigwedge^jG,\quad0\le j\le2,
 \qquad
 d(v\otimes\eta)=v\otimes a\wedge\eta,
\]
\[
 d(h\otimes\eta)=h\otimes a\wedge\eta+\kappa\wedge\eta. \tag{D.19}
\]
Here $V$ is the rank-two vertical bundle, $L$ the horizontal line,
$G=R^1\OO$ on the root chart, and $\kappa\in V\otimes G\otimes L^*$
is the actual tangent-extension cocycle. The matrix is obtained
by differentiating the period translations. In the raw period and
constant-cocycle frames it is
\[
 \kappa=\begin{pmatrix}6\mu'&\tau'\\\beta_c'&\mu'\end{pmatrix},
 \qquad
 a=(z-n+6q\mu+r\beta_c,-p+q\tau+r\mu).               \tag{D.20}
\]
The conventional factor $2\pi i$ in the Dolbeault representative of
the exponential class multiplies both components of $a$ by the
same nonzero constant; absorbing it in the frame has no effect on
any module below. The displayed formula also follows directly by
subtracting the holomorphic period rows from the constant lattice
cocycle. This identifies the chain map and its tangent extension;
no unspecified map between fibre cohomology groups is being used.
For the nonzero Fourier characters, use the filtration by the number
of base $(0,1)$ factors: the relative contractions make the associated
graded complexes acyclic, and the two-step filtration then proves
acyclicity of the full complexes. Thus the contraction need not
commute with the base Dolbeault operator. The finite cyclic group
normalizes the real translation group, so averaging and this
comparison are equivariant; exact invariants compute the actual
multiple-fibre stalks.

\item At a regular resonance, $a'=\kappa^t(q,r)^t$ and
$(q,r)\ne(0,0)$. The determinant of $\kappa$ is nonzero there by
Part A. Therefore $a$ has a simple zero. Choose its nonzero leading
direction as the second frame of $G$ and use invertibility of
$\kappa$ to choose the two vertical frames. Holomorphic invertible
changes reduce (D.19) to
\[
 a=(0,s),\qquad \kappa=\begin{pmatrix}0&1\\1&0\end{pmatrix}. \tag{D.21}
\]
The second cohomology is generated by $h\gamma\delta$ and is killed
by $s^2$; the relations eliminate the two vertical generators.
The first cohomology is generated by $e_1\delta,e_2\delta$, each
killed by $s$. These are equalities of modules over $\C\{s\}$
and prove (D.18).

\item There is an additional joint coefficient to check at the order-three
point. Put $t=s^3$, $\nu=\mu-(2-\tau)/3$ and
$b=\beta_c+4\mu+(2/3)\tau-b(0)$. The order-three transformation
laws give $\nu=\nu_2s^2+O(s^3)$ and $b=b_3s^3+O(s^4)$.
The elliptic calculation in \ref{inv:A2.3} gives, with elliptic scale
$f(s)=F s^2+O(s^3)$,
\[
 \nu_2^2=9F^2/512.
\]
For clarity we evaluate the other coefficients with the original
Weierstrass normalization. Since $x=f^2\wp$,
$y=f^3\wp'/2$, one has
$f^4g_2=12t^3(t-1)$ and $f^6g_3=-8t^4(t-1)^2$.
The identities $g_2=4\pi^4E_4/3$, $g_3=8\pi^6E_6/27$ and
$E_4'(\rho)=-2\pi iE_6(\rho)/3$ give
\[
 E_6(\rho)=-\frac{27}{\pi^6F^6},\qquad
 \tau(s)=\rho+\tau_1s+O(s^2),\quad \tau_1=\frac{i\pi F^2}{2}.
\]
The exact period determinant identity in Part A gives
\[
 K(s)=6\mu'^2-\tau'\beta_c'=K_2s^2+O(s^3),
 \quad K_2=\frac{3F^2}{4}.
\]
Equivalently $K=6\nu'^2-\tau'b'$, so
$K_2=24\nu_2^2-3\tau_1b_3$ and $b_3=7i/(32\pi)$.

\item For a nonintegral resonance the relation is
$(q,r,p)=\pm(k,3k,2k)$, $k\ne0$; the sign only changes the
following nonzero scalar. After a constant change of the two
cocycle frames, its Picard section is $3k(\nu,b)$, and the
transpose period derivative becomes
\[
 \begin{pmatrix}\tau'&\nu'\\6\nu'&b'\end{pmatrix}.
\]
Eliminating the unit $\tau'$ leaves the joint transverse component
$3k(b-6\nu\nu'/\tau')$. Its coefficient of $s^3$ is
\[
 -\frac{k(K_2+12\nu_2^2)}{\tau_1}
       =\frac{123ik}{64\pi}\ne0.                      \tag{D.22}
\]
Thus neither the order of the Picard section nor the determinant
order alone is used to select a Smith form. The joint coefficient
has been evaluated exactly.

\item Choose the second cocycle frame to be the Picard section divided by
$s^2$. Equivariance and (D.22) give a matrix with a unit in its
lower-right entry, first-order off-diagonal terms with nonzero
coefficients, and a second-order upper-left term. Subtracting a
holomorphic multiple of the second vertical row removes the last
term; the nonzero determinant coefficient and (D.22) keep both
off-diagonal coefficients units. Rescaling the equivariant frames
therefore reduces (D.19) to
\[
 a=(0,s^2),\qquad \kappa=\begin{pmatrix}0&s\\s&1\end{pmatrix}. \tag{D.23}
\]
Its cohomology is
\[
 H^1=\C\{s\}/(s^2)\oplus\C\{s\}/(s^3),\qquad
 H^2=\C\{s\}/(s^2)\oplus\C\{s\}/(s^3).             \tag{D.24}
\]
For example in degree one put
$u=e_1\delta$, $v=e_2\gamma+s h\delta$; the remaining cycle
$e_2\delta$ is eliminated by the boundary
$s u+e_2\delta+s v$. The two surviving relations are
$s^2u=0,s^3v=0$. In degree two eliminate $e_2\gamma\delta$;
the remaining generators $e_1\gamma\delta,h\gamma\delta$
have annihilators $s^2,s^3$ respectively.

\item We specify the character convention to avoid a sign ambiguity.
A frame has invariant exponent $w$ when its coefficient satisfies
$g(\zeta s)=\zeta^w g(s)$; its pullback-frame weight is the negative
of $w$. For the trivialized twist in (D.16), these exponents are
$n$ at $p_1$ and $3n$ at $p_2$. This agrees with $M_n$ in \ref{inv:D4.1}
and follows by evaluating its integer cocycle on $v_1,v_2$.
At $p_1$ the two generators in $H^1$ of (D.24) have invariant
exponents $n+2,n+2$, and the two generators in $H^2$ have exponents
$n+2,n$. The length-three summand has one invariant coefficient
for every $n$; the length-two summand has one precisely when
$n\not\equiv0\pmod3$. The coarse parameter is $s^3$, so all these
invariants are killed by it. This proves the asserted order-three
stalks with the full nilpotent layers included.

\item At $p_2$ the determinant is a unit. The Picard section has a simple
zero, so (D.21) applies on the order-four cover. The invariant
exponents of the two degree-one generators are $3n+1,3n+2$;
the cyclic degree-two module of length two has generator exponent
$3n+2$. In both degrees its invariant dimension is one exactly for
$n\equiv1,2\pmod4$. The coarse parameter $s^4$ kills the modules.
This proves the stated order-four calculation. The two independent
exact checks in \texttt{d\_koszul.py} are Smith normal forms and
ranks of truncated jet complexes; the latter are also restricted to
every stabilizer character. Their outputs agree with (D.21)--(D.24)
and with every invariant length just computed.
\end{enumerate}
\end{proof}

\begin{proposition}[The full twisted cohomology and its jump loci]\label{inv:D4.5}
\StatementRecord{inv:D4.5}{proved}{Historical events: GPT-6 Astra (OpenAI), ChatGPT Pro mode, effort Pro, S6\_09\_PAPERS, 17 November 2026: targeted proof check. Also checked in this writing. Pre-continuation writing events: GPT-6 Astra (OpenAI), ChatGPT Pro mode, effort Pro, S6\_10\_PAPERS, 18 November 2026: targeted proof check. Continuation metadata events: GPT-6 Astra (OpenAI), ChatGPT Pro mode, effort Pro, S6\_10\_PAPERS continuation, 19 November 2026: used and rechecked in this run; status-vocabulary reconciliation; metadata reconciliation only; no new mathematical check. \emph{Pre-continuation recorded proof scope.} Read full root-base, logarithmic tangent, direct-image and twisted-Dolbeault proofs. Re-derived coarse degree function, local normalization quotient, resonance equation, Fourier contraction comparison, joint order-three coefficient 123ik/(64pi), stated local modules, and jump-locus finiteness. Exact Koszul and twist-program computations retain their separately recorded replay scopes.}{the hypotheses, notation and definitions stated here; \cref{inv:D4.1}; \cref{inv:D4.2}; \cref{inv:D4.4}; \cref{inv:C3.8}}

For $z\notin\Z$, let $\mathcal R_z$ be the set of regular base
points satisfying (D.15), counted as points, not as marked integer
tuples. It is finite. Put $N_z=|\mathcal R_z|$. Define
$\ell_1(z)=0$ off $\Lambda_1$, and
$\ell_1(z)=1+\epsilon_n$ if
$z=n+k(3\beta_1+4-2\rho)$, $k\ne0$.
Define $\ell_2(z)=1$ if
$z=n+k(2\beta_2+3-3i)$, $k\ne0$, with
$n\equiv1,2\pmod4$, and $\ell_2(z)=0$ otherwise.
Then, setting $J_z=2N_z+\ell_1(z)+\ell_2(z)$,
\[
 \boxed{(h^0,h^1,h^2,h^3)(T_X\otimes M_z)=(0,J_z,1+J_z,0).} \tag{D.25}
\]
More precisely, with $\mathfrak m_t$ the maximal ideal at $t$,
\[
 R^1f_*(T_X\otimes M_z)=
   \bigoplus_{t\in\mathcal R_z}\C_t^2
      \oplus\C_{p_1}^{\ell_1(z)}\oplus\C_{p_2}^{\ell_2(z)},
\]
\[
 R^2f_*(T_X\otimes M_z)=\C_{p_0}\oplus
   \bigoplus_{t\in\mathcal R_z}\OO_{B,t}/\mathfrak m_t^2
      \oplus\C_{p_1}^{\ell_1(z)}\oplus\C_{p_2}^{\ell_2(z)}. \tag{D.26}
\]
Together with \ref{inv:D4.1} these formulas cover every element of
$\operatorname{Pic}(X)=\C$. The actual cohomology jump loci are
locally discrete subsets of this complex line.
\end{proposition}


\begin{proof}
\begin{enumerate}[label=\arabic*.,leftmargin=*]
\item The support assertions and all stalks were proved in \ref{inv:D4.2}--\ref{inv:D4.4}.
The direct image in degree zero is zero for a nonintegral twist:
it has generic rank zero and is torsion-free, since multiplication
by a nonzero base function is injective on the bundle on the smooth
integral total space. The support of either higher coherent direct
image is a proper closed analytic subset of the compact curve,
and hence is finite. This proves finiteness of $\mathcal R_z$ as
well as the absence of an accumulation at the special values.
If two integer tuples at a marked fibre represented the same
trivialized line, their difference would be an integral real class
in $H^{1,0}$, which is zero. Thus there is no multiplicity from
duplicate tuples. Formula (D.26) follows by the decomposition of a
finite-support sheaf into its local modules. Its global sections
have dimension equal to its length and its $H^1$ is zero, so Leray
gives (D.25).

\item To clarify the word generic, the global closed form $G_\gamma$ in
\ref{inv:C3.8} defines a holomorphic line bundle on $X\times\C$ by the
integrable operators $\bar\partial-2\pi i zG_\gamma$. Proper
coherent base change for the projection to $\C$ makes each
cohomology-dimension threshold an analytic subset. A real noninteger
$z$ has no resonance by (D.17) and (D.16), so the generic values are
attained. All possible exceptional values were shown to form a
countable set. Thus each actual threshold exceeding a generic
dimension is a proper analytic subset of a complex curve and is
locally discrete. Some indexed resonances, namely the two invisible
order-four congruences, produce no jump; (D.25) distinguishes them.

\item This is an explicit countable description of the jump loci in the
Picard coordinate, rather than a generic assertion. It can also be
indexed algebraically on the elliptic surface. If a regular fibre
has torsion section $P$ of exact order $r_t$, choose its marked
relation $q_t\tau+r_t\mu=p_t$. Its full resonance lattice is
\[
 \Z+\bigl(6q_t\mu(t)+r_t\beta_c(t)\bigr)\Z.
\]
Every regular resonance arises this way. The underlying values
$t$ are exactly the non-special roots of the division polynomials
$\psi_N(P)$, $N\ge2$. The nonintegral lattice points contribute two
to $J_z$ for each distinct such fibre. The finite-point lattices
have the congruence weights already stated. The script
\texttt{d\_division\_polynomials.py} evaluates the polynomials
exactly for $2\le N\le12$; for example the first nonconstant
regular factors are $5t-32$, $89t^2+4736t-4096$ and
$11t^2-6016t+8192$. These computations specify an executable
enumeration of the resonance loci without inventing decimal values
for the unspecified parameter $c$.

\item As an independent index check, $c_1(X)=c_2(X)=c_1(M)=0$ in integral
cohomology, so the degree-six Riemann--Roch integrand for the actual
twisted tangent bundle is $c_3/2$. Its integral is one. Both (D.13)
and (D.25) have alternating sum one. The index is used here to
verify the completed sheaf calculation, not to select its lengths.
\end{enumerate}
\end{proof}

\paragraph{Source and re-derivation of the tangent theorem.}
The tangent cohomology $(1,1,1,0)$ and the smooth one-dimensional
Kuranishi germ are theorems of Du, Guo, Kusakabe, and Xie,
arXiv:2610.00250v1, Section~6 \citep{DuEtAl2026v1}.
The invariant run re-derived these results after reading that paper.
The present proof prints the root-base calculation, its local
comparisons, and the additional coarse direct images and all Picard
twists. Version~2 \citep{DuEtAl2026v2} does not contain the deformation
section; this does not alter the credit for the version~1 theorems.



\subsection{Tangent cohomology for arbitrary coprime multiplicities}\label{push:section:E.2}
\begin{theorem}[tangent cohomology and local miniversality]\label{push:src:E.2.1}
\StatementRecord{push:src:E.2.1}{proved}{Historical events: Claude Fable 5.2 (Anthropic), effort high, 15 November 2026: carried as background without fault found. Claude Fable 5.2 (Anthropic), effort high, 16 November 2026: carried as background without fault found; used and rechecked in this run. GPT-6 Astra (OpenAI), ChatGPT Pro mode, effort Pro, S6\_09\_PAPERS, 17 November 2026: targeted proof check. Also checked in this writing. Pre-continuation writing events: GPT-6 Astra (OpenAI), ChatGPT Pro mode, effort Pro, S6\_10\_PAPERS, 18 November 2026: targeted proof check. Continuation metadata events: GPT-6 Astra (OpenAI), ChatGPT Pro mode, effort Pro, S6\_10\_PAPERS continuation, 19 November 2026: used and rechecked in this run; status-vocabulary reconciliation; metadata reconciliation only; no new mathematical check. \emph{Pre-continuation recorded proof scope.} Re-derived the weighted-stack Picard presentation and root-forgetting map, all Leray direct images and weighted Laurent counts, the logarithmic tangent sequences and residue connecting rank, and the resulting tangent tuple (1,1,1,0). Checked the explicit nonzero groupoid Kodaira-Spencer class and the one-variable Kuranishi-ring argument.}{the hypotheses, notation and definitions stated here; \cref{push:section:E.1}; \cref{push:src:E.3.1}}

For every sphere member with multiplicities $(3k,4l)$ and the
admissible logarithmic data of \ref{push:section:E.1}, put $z=kl$ and $d=3k+4l$.
On $\mathfrak B_{k,l}=\mathbb P(3k,4l)$ the root lift
$\bar f:X\to\mathfrak B_{k,l}$ satisfies
\[
 \begin{split}
 \mathcal E&=\mathcal O\oplus\mathcal O(-z),\\
 R^q\bar f_*\mathcal O_X&=
 \mathcal O,\quad\mathcal O\oplus\mathcal O(-z),\quad\mathcal O(-z),\\
 K_X&=\bar f^*\mathcal O(z-d).
 \end{split}
\]
The tangent cohomology is $(h^0,h^1,h^2,h^3)(X,T_X)=(1,1,1,0)$.
With the groupoid derivative convention
$\kappa(\Phi)=(\partial_c\Phi)\Phi^{-1}$, the class of
$\partial/\partial c$ is $e\otimes[\gamma]$.  The class of
$\partial/\partial u$ is
$-(2\pi i u)^{-1}e\otimes[\gamma]$.  Both are nonzero.
The Kuranishi germ is a smooth reduced line, and the fixed-data
period family is locally miniversal.
\end{theorem}


\begin{proof}
\begin{enumerate}[label=\arabic*.,leftmargin=*]
\item The root object is defined by the actual divisor identities
$f^*(0)=mS_1$ and $f^*(1)=nS_2$, with $m=3k$, $n=4l$.
It gives a holomorphic map to the analytic root stack.  Its base
changes to the two root discs are the actual smooth torus
families; the quotient actions are free.  At the cusp its charts
are $q=z_0z_1z_2$, so the source is smooth and the map is flat,
of relative dimension two, with Gorenstein fibres.  Properness is
checked on the root atlas, where it is properness of the compact
torus families and the original compact cusp filling.

\item Use weighted coordinates $x,y$ of weights $m,n$, with coarse
coordinate $t=y^m/(y^m-x^n)$.
The two root divisors have bundles $\mathcal O(n)$ and
$\mathcal O(m)$, and the cusp has bundle $\mathcal O(mn)$.
The root-forgetting morphism
\[
 r:\mathfrak B_{k,l}\longrightarrow\mathfrak B_{1,1},
 \qquad (x,y)\longmapsto(x^l,y^k),
\]
uses the group homomorphism $\lambda\mapsto\lambda^{kl}$ on
the weighted quotient.  It is a holomorphic morphism of analytic
stacks over $B$, with local maps $s\mapsto s^k$ and
$s\mapsto s^l$.  In particular $r^*\mathcal O(1)=\mathcal O(z)$.
For completeness, the analytic Picard group of this root curve has
the presentation
$\langle H,R_1,R_2\mid mR_1=H,\ nR_2=H\rangle$:
after subtracting the two stabilizer characters, a line bundle
descends to the coarse projective line.  Its only relations are
the two defining root identities.  The Smith form and
$\gcd(m,n)=1$ identify this group with $\mathbb Z$, with
$H=mn\mathcal O(1)$, $R_1=n\mathcal O(1)$, and
$R_2=m\mathcal O(1)$.  This proves the analytic torsion-free
Picard assertion used below.
Translations have identity vertical derivative and act trivially
on structure-sheaf cohomology.  Thus the vertical bundle computed
from the linear period factors is $r^*(\mathcal O\oplus
\mathcal O(-1))$.  The original determinant calculation can also
be repeated: its quotient line $L$ has
$L^{12}=\mathcal O(-p_0)=\mathcal O(-12z)$, so
$L=\mathcal O(-z)$ because the Picard group is torsion-free.
The fixed vector gives a subline $\mathcal O$; its extension
splits because $H^1(\mathcal O(z))=0$.

\item The cusp normalization resolution has terms
$\mathcal O_{\mathrm{dP}_6}$, three copies of
$\mathcal O_{\mathbb P^1}$, and two point sheaves.
Its global differential is $(a,b,c)\mapsto
(a-b+c,a-b+c)$, giving cohomology $(1,2,1)$; its alternating cup
pairing in degree one is nondegenerate.  The degree-one cocycle
$\gamma$ extends across the cusp because it vanishes on the toric
period lattice.  Its edge values $(0,-1,-1)$ have nonzero cup
product with $(1,0,-1)$.  On finite root charts its class is the
nonzero constant lattice cocycle $\gamma$, invariant under every
deck transformation.  These are explicit comparisons of the
same class, so it gives a nowhere-zero subline in
$G=R^1\bar f_*\mathcal O_X$.
The base-change and relative-duality argument \ref{push:src:E.3.1} applies on
these charts.  It gives $R^2\bar f_*\mathcal O_X=L$ and
$\det G=L$.  The extension of $L$ by the specified
$[\gamma]$-subline splits by the same $H^1(\mathcal O(z))=0$.
This proves all direct-image formulas.

\item The logarithmic tangent sequence is
\[
 0\longrightarrow\mathcal V=\bar f^*\mathcal E
 \longrightarrow\mathcal T=T_X(-\log W)
 \longrightarrow\mathcal H=\bar f^*\mathcal O(d-12z)
 \longrightarrow0.
\]
Its last map is the relative differential on smooth root charts
and sends $\sum a_i z_i\partial_{z_i}$ to
$(\sum a_i)q\partial_q$ at the cusp.  It is surjective there.
Since $K_{\mathfrak B}=\mathcal O(-d)$, determinants give
$K_X=\bar f^*\mathcal O(z-d)$.

\item We give the weighted cohomology calculation explicitly; simply
replacing $1$ by $kl$ in an acyclicity assertion would be incorrect.
The two quotient charts $x\ne0$, $y\ne0$ are finite quotients
of affine lines, and their overlap is a finite quotient of a
punctured line.  Finite averaging and Stein acyclicity compute
coherent cohomology by this two-chart complex.  Its homogeneous
Laurent monomials give
\[
 \begin{split}
 h^0\mathcal O(a)&=\#\{(i,j)\in\mathbb Z_{\ge0}^2:mi+nj=a\},\\
 h^1\mathcal O(a)&=h^0\mathcal O(-a-d).
 \end{split}
\]
The second equality identifies the remaining monomials
$x^{-i-1}y^{-j-1}$; it does not require a positivity theorem.
The bundles $\mathcal O(-z),\mathcal O(-2z)$ are acyclic.
Indeed a possible degree-one monomial would give
$m(i+1)+n(j+1)=az$, with $a=1$ or $2$.
Coprimality forces $i+1=lA$, $j+1=kB$, hence
$3A+4B=a$ with $A,B\ge1$, which is impossible.
Positive degree has no degree-one monomials, so
$H^1\mathcal O(z)=0$ as used above.

\item Put $b=d-12z$.  Both $b$ and $b-z$ are negative.
Degree-one monomials for $\mathcal O(b)$ would satisfy
$3A+4B=12$, with $A,B\ge1$, and none exist.
For $\mathcal O(b-z)$ the equation is $3A+4B=13$;
its only positive solution is $(A,B)=(3,1)$, producing
$x^{3l-2}y^{k-2}$ in the dual description.  It is allowed
exactly when $k>1$; write $e_k=1$ in that case and $e_k=0$
otherwise.  Projection formula and the two-column Leray sequence
now give
\[
 h^\bullet(\mathcal V)=(1,1,0,0),\qquad
 h^\bullet(\mathcal H)=(0,0,e_k,e_k),\qquad
 h^\bullet(\mathcal T)=(1,1,e_k,e_k).
\]
The program \texttt{e\_push\_exact.py} independently checks these
counts for all $1965$ coprime pairs with $k,l\le80$; the monomial
argument proves them for all pairs.

\item On the normalization $\nu:E=\mathrm{dP}_6\to W\subset X$,
the branchwise normal line is $K_E$: the invariant cusp form
$d\log x_1\wedge d\log x_2\wedge dq$ trivializes $K_X$
on the entire cusp neighbourhood.  Local normal coefficients
therefore give the exact sequence
\[
 0\longrightarrow\mathcal T\longrightarrow T_X
 \longrightarrow\nu_*K_E\longrightarrow0,
 \qquad h^\bullet(K_E)=(0,0,1).
\]
We must determine its connecting map when $k>1$.
The dual logarithmic sequence, tensored with $K_X$, has horizontal
subline $\bar f^*\mathcal O(13z-2d)$ and vertical quotient
$\bar f^*(\mathcal O(z-d)\oplus\mathcal O(2z-d))$.
The latter has no global sections, by the same equations
$3A+4B=1,2$ above.  Consequently
\[
 H^0(\Omega_X^1(\log W)\otimes K_X)
 =H^0(\mathfrak B,\mathcal O(13z-2d)).
\]
For $k>1$ this is generated by the monomial
$x^{3l-2}y^{k-2}$ multiplying the horizontal logarithmic
differential.  Both $x$ and $y$ are nonzero over $p_0$, so its
residue on $E$ is nonzero.  The ordinary logarithmic residue
sequence, tensored with $K_X$, identifies this residue map with
the Serre dual of
$H^2(E,K_E)\to H^3(X,\mathcal T)$.
One can check this identification directly in each chart
$q=z_0z_1z_2$: the residues of $dq/q$ are one on each branch,
and contraction with the invariant three-form is the normal
coefficient pairing.  Thus that connecting map has rank $e_k$.
The long exact sequence proves $h^2(T_X)=1$ and $h^3(T_X)=0$;
degrees zero and one are unchanged.

\item Finally $\partial_c\Pi=e\otimes\gamma$ gives the constant
groupoid cocycle $\lambda\mapsto\gamma(\lambda)e$ on every
regular torus.  On the enlarged finite roots the derivative of the
affine generator is $(r_j/m_j)e$, and the derivative of its
punctured connector is $(\log s_j)/(2\pi i)\,r_je$.
The identity $m_j(r_j/m_j)=\gamma(w_j)$ verifies the power
relation, while $\gamma A_j=\gamma$ verifies conjugation with
every lattice arrow.  Changing $\log s_j$ by $2\pi i N$
adds the derivative of the integral arrow $Nw_j$; the connector
formula is used only on the punctured overlap.  Its central
root-chart representative is the constant $(r_j/m_j)e$,
which is holomorphic.  At the cusp $\gamma$ kills the toric periods,
so the constant deck cocycle is exactly its extending cocycle
above.  Its image under the two isomorphisms
$H^1(\mathcal V)\simeq H^1(\mathcal T)\simeq H^1(T_X)$
is $e\otimes[\gamma]\ne0$.  The chain rule gives the asserted
$u$-class.  The Kuranishi local algebra is
$\mathbb C\{t\}/I$ since $h^1=1$.  The existing family induces
$t\mapsto g(c)$ with $g'(c_0)\ne0$.  The analytic inverse
function theorem makes this substitution injective; every
element of $I$ maps to zero, so $I=0$.  Its classifying map is
then an isomorphism of germs, proving local miniversality.
\end{enumerate}
\end{proof}

\begin{proposition}[local separation and canonical dimension]\label{push:src:E.2.2}
\StatementRecord{push:src:E.2.2}{proved}{Historical events: Claude Fable 5.2 (Anthropic), effort high, 16 November 2026: carried as background without fault found. GPT-6 Astra (OpenAI), ChatGPT Pro mode, effort Pro, S6\_09\_PAPERS, 17 November 2026: targeted proof check. Also checked in this writing. Pre-continuation writing events: GPT-6 Astra (OpenAI), ChatGPT Pro mode, effort Pro, S6\_10\_PAPERS, 18 November 2026: targeted proof check. Continuation metadata events: GPT-6 Astra (OpenAI), ChatGPT Pro mode, effort Pro, S6\_10\_PAPERS continuation, 19 November 2026: used and rechecked in this run; status-vocabulary reconciliation; metadata reconciliation only; no new mathematical check. \emph{Pre-continuation recorded proof scope.} Re-derived deformation constancy of the marked multiplicities and the canonical Kodaira-dimension sign; checked the exact divisibility argument excluding the equality case for the stated coprime parameters.}{the hypotheses, notation and definitions stated here; \cref{push:src:E.1.2}; \cref{push:src:E.2.1}}

The multiplicity pair $(3k,4l)$ is constant on every sufficiently
small deformation of one of the preceding members.  Distinct
pairs therefore determine disjoint open loci among these
isomorphism classes.  Moreover
\[
 \kappa(X)=\begin{cases}-\infty,&kl<3k+4l,\\1,&kl>3k+4l.
 \end{cases}
\]
Equality $kl=3k+4l$ never occurs under $\gcd(3k,4l)=1$.
\end{proposition}


\begin{proof}
\begin{enumerate}[label=\arabic*.,leftmargin=*]
\item Local miniversality identifies every small deformation with the
fixed-data period family; \ref{push:src:E.1.2} makes its multiplicity pair an
intrinsic biholomorphic invariant.  The canonical formula of
\ref{push:src:E.2.1} and connected fibres identify its pluricanonical sections
with those of $\mathcal O(r(kl-3k-4l))$ on the root curve.
Negative weight has no sections.  Positive weight has linearly
growing section counts along multiples of $12kl$, giving
Kodaira dimension one.  Equality would imply $k\mid4l$;
coprimality forces $k=1$, and then $l=3+4l$ is impossible.
\end{enumerate}
\end{proof}

\begin{proposition}[separation in irreducible analytic families]\label{push:src:E.2.3}
\StatementRecord{push:src:E.2.3}{proved}{Historical events: Claude Fable 5.2 (Anthropic), effort high, 16 November 2026: carried as background without fault found. GPT-6 Astra (OpenAI), ChatGPT Pro mode, effort Pro, S6\_09\_PAPERS, 17 November 2026: targeted proof check. Also checked in this writing. Pre-continuation writing events: GPT-6 Astra (OpenAI), ChatGPT Pro mode, effort Pro, S6\_10\_PAPERS, 18 November 2026: targeted proof check. GPT-6 Astra (OpenAI), ChatGPT Pro mode, effort Pro, S6\_10\_PAPERS, 18 November 2026 04:30:14.617321+00:00: used and rechecked in this run. Continuation metadata events: GPT-6 Astra (OpenAI), ChatGPT Pro mode, effort Pro, S6\_10\_PAPERS continuation, 19 November 2026: used and rechecked in this run; status-vocabulary reconciliation; metadata reconciliation only; no new mathematical check. \emph{Pre-continuation recorded proof scope.} Re-derived the pluricanonical growth difference and recovery of the ordered (3k,4l) pair from the reduced rational coefficient; checked the local universal-family and coherent-jump arguments that make the pair constant on the connected deformation family.}{the hypotheses, notation and definitions stated here; \cref{push:src:E.2.1}; \cite[Chapter III, Theorem 4.1]{BS76}}

Suppose a proper holomorphic submersion over a connected
irreducible reduced complex base has two fibres of the type
in \ref{push:src:E.2.1}.  Then their multiplicity pairs agree, even if some
other fibres have not been constructed by this method.
For these members the asymptotic number
\[
 s(X)=\lim_{r\to\infty}
 \frac{h^0(X,K_X^r)-h^0(X,K_X^{-r})}{r}
 =\frac1{12}-\frac1{3k}-\frac1{4l}
\]
recovers the ordered multiplicity pair uniquely.
\end{proposition}


\begin{proof}
\begin{enumerate}[label=\arabic*.,leftmargin=*]
\item For coprime $m,n$, the number of nonnegative solutions of
$mi+nj=N$ is $N/(mn)+O(1)$ as $N\to\infty$:
there is one permitted residue class of $j$ modulo $m$, and
its representatives lie between $0$ and $N/n$.
The canonical formula in \ref{push:src:E.2.1} proves the displayed limit,
with a negative value when the anticanonical counts grow.
Write $1/12-s(X)=a/b$ in lowest positive terms.
Since $\gcd(m+n,mn)=1$, we have $a=m+n$, $b=mn$.
Thus $m,n$ are the two roots of $T^2-aT+b$.
Their order is determined because only $m$ is divisible by
three and only $n$ is divisible by four.

\item For every fixed integer $r$, apply the finite-free-complex
theorem of \cite[Chapter III, Theorem 4.1]{BS76} to
the relative canonical power, which is an invertible flat
sheaf for a holomorphic submersion.  Its zeroth cohomology
dimension has analytic jumping loci.  On an irreducible
base any value constant on a nonempty open set is its
generic value: a proper analytic jumping locus has empty
interior.  Local miniversality in \ref{push:src:E.2.1} makes these dimensions
constant near either of the two given fibres, since all
nearby fibres have the same multiplicities and their
canonical formula is independent of the period parameter.
Consequently the two fibres have equal dimensions for
every positive and negative $r$.  Their limits $s(X)$
are equal, and the preceding arithmetic recovers the same
pair.  This proof uses relative cohomology only; it makes
no general invariance assertion for Kodaira dimension.
\end{enumerate}
\end{proof}

\subsection{Analytic foundations and the version 1}\label{push:section:E.3}
\begin{cited}[base change and duality]\label{push:src:E.3.1}
\StatementRecord{push:src:E.3.1}{proved}{Historical events: Claude Fable 5.2 (Anthropic), effort high, 16 November 2026: carried as background without fault found. GPT-6 Astra (OpenAI), ChatGPT Pro mode, effort Pro, S6\_09\_PAPERS, 17 November 2026: targeted proof check. Also checked in this writing. Pre-continuation writing events: GPT-6 Astra (OpenAI), ChatGPT Pro mode, effort Pro, S6\_10\_PAPERS, 18 November 2026 04:30:14.617321+00:00: used and rechecked in this run. Continuation metadata events: GPT-6 Astra (OpenAI), ChatGPT Pro mode, effort Pro, S6\_10\_PAPERS continuation, 19 November 2026: used and rechecked in this run; status-vocabulary reconciliation; metadata reconciliation only; no new mathematical check. \emph{Pre-continuation recorded proof scope.} Current-writing verification remains in progress. The earlier checks named above retain their recorded scope.}{the hypotheses, notation and definitions stated here; \cref{push:src:E.2.1}; \cref{fam:src:A.5.2}; \cite{DuEtAl2026v1}; \CompendiumRef{I}{app:rederive:push-src-E.3.1}; \cite[Section 6, Hauptsatz I, and Section 7, Satz 5]{Grauert1960D}; \cite{Grauert1963ErratumD}; \cite[Chapter III, Theorem 4.1]{BS76}}

For the proper flat root-curve morphisms in \ref{push:src:E.2.1}, the canonical
geometric-point comparison
\[
 (R^q\bar f_*\mathcal O_X)\otimes k(b)
 \longrightarrow H^q(X_b,\mathcal O_{X_b})
\]
is an isomorphism for $q=0,1,2$, and the direct images are locally
free.  The trace induces the canonical isomorphism
\[
 (R^2\bar f_*\mathcal O_X)^\vee
 \simeq\bar f_*\omega_{X/\mathfrak B}.
\]
Both comparisons descend on the analytic root-curve groupoid.
These conclusions, and hence the use made of them in \cref{fam:src:A.5.2}, require
no result specific to the version 1 of
\cite{DuEtAl2026v1}.
\end{cited}
The expanded source derivation is retained in the proof compendium, Paper~I, label \nolinkurl{push:src:E.3.1}.




This is analytic base change and duality, with the source hypotheses specified in the re-derivation. A complete re-derivation is given in \CompendiumRef{I}{app:rederive:push-src-E.3.1}.


Grauert's original direct-image theorem
\cite[Section 6, Hauptsatz I, and Section 7, Satz 5]{Grauert1960D}
was also opened together with its correction \cite{Grauert1963ErratumD}.
The correction explicitly warns that the exceptional locus in
Section 7, Satz 2 need not be closed.  \ref{push:src:E.3.1} uses the
constant-cohomology theorem on every chart and the later
base-change formulation, so it does not make that erroneous
closedness inference.  For a generic argument over a smooth
curve one can instead use the finite free complex of
\cite[Chapter III, Theorem 4.1]{BS76}; its rank-jump minors
give the required local analytic exceptional sets.


% Compendium-only push:src:E.3.2: apparatus: historical/current dependency audit of the five credited DGKX version1 theorems; complete mathematical applications retained at their actual receiving results and in the credited-result appendix



% Compendium-only push:src:E.3.3: apparatus: exact scope comparison for the five DGKX version1 standard(3,4) results; the source-specific receiving results are separately retained



\section{Other logarithmic transformations and finite covers}

\subsection{Comparison under logarithmic transformation}

\begin{proposition}[The marked geometric dictionary on the first page]\label{inv:B1.0}
\StatementRecord{inv:B1.0}{proved}{Historical events: Claude Fable 5.2 (Anthropic), effort high, 17 November 2026: carried as background without fault found. GPT-6 Astra (OpenAI), ChatGPT Pro mode, effort Pro, S6\_09\_PAPERS, 17 November 2026: targeted proof check. Also checked in this writing. Pre-continuation writing events: GPT-6 Astra (OpenAI), ChatGPT Pro mode, effort Pro, S6\_10\_PAPERS, 18 November 2026: targeted proof check. Continuation metadata events: GPT-6 Astra (OpenAI), ChatGPT Pro mode, effort Pro, S6\_10\_PAPERS continuation, 19 November 2026: used and rechecked in this run; status-vocabulary reconciliation; metadata reconciliation only; no new mathematical check. \emph{Pre-continuation recorded proof scope.} Read complete proofs and re-derived the marked theta-line transport, finite-root base changes and canonical characters, weighted section multiplication, integral fundamental-group quotient, cyclic-cover residue sequence, equivariant extra cohomology, and the rational Mayer--Vietoris ranks. Exact external theorem locators remain a separate citation gate.}{the hypotheses, notation and definitions stated here; the explicit constructions and calculations in the proof}

The theta quotient identifies the period family with the line-bundle
quotient in Engel's construction. In the marking and reference
linearization fixed in the conventions, $u=w^2$. Replacing that
reference multiplies $u$ by the corresponding nonzero constant.
\end{proposition}


\begin{proof}
\begin{enumerate}[label=\arabic*.,leftmargin=*]
\item First quotient the period presentation by its last three columns,
using $z=\zeta_1$ and $v=e^{2\pi i\zeta_2}$. The resulting principal
$\C^*$-bundle over $\C/(\Z+\tau\Z)$ has factors
\[
 (z,v)\mapsto(z+1,v),\qquad
 (z,v)\mapsto(z+\tau,e^{2\pi i\mu}v).
\]
For an odd theta function $\vartheta$, the meromorphic function
$\vartheta(z-\mu)/\vartheta(z)$ has precisely these factors and
divisor $(\mu)-(0)$. Thus this is the complement of the zero section
of $L_\mu=\OO(P-O)$, with $P=\mu$ on the elliptic curve. The first
period column acts by
$(z,v)\mapsto(z+6\mu,e^{2\pi i\beta}v)$; its inverse covers $-6P$,
as does Engel's contracting generator. Replacing $\beta^{\rm part}$
by $\beta^{\rm part}+c$ multiplies this inverse lift by
$e^{-2\pi i c}$. A change of reference lift multiplies its scale by
a single constant $\lambda_0\in\C^*$, proving the asserted parameter
comparison and naming its entire ambiguity.

\item For the integral marking fixed in the conventions, use
\[
 P_{\rm mark}=\begin{pmatrix}0&-2&-1&-2\\2&-1&0&-1\\0&2&1&3\\-1&0&0&0\end{pmatrix}.
\]
With the monodromy matrices in the construction, direct multiplication gives
$\det P_{\rm mark}=-1$, $T_0P_{\rm mark}=P_{\rm mark}A_1$ and
$T_1P_{\rm mark}=P_{\rm mark}A_2$. A positive base meridian ends at the
inverse root-deck translate, and its endpoint identification applies the
clockwise root generator; thus its homology action is $A_j$.
Because $A_jv_j=v_j$, the inverse affine generator has translation
$-v_j/m_j$. Its Engel coordinates are therefore
$a_0=-P_{\rm mark}v_1/3$ and $a_1=-P_{\rm mark}v_2/4$.
The chosen reference inverse lift is the transported lift at $c=0$.
Consequently $\lambda_0=1$ and $u=e^{-2\pi i c}=w^2$.

\item For concordance, the reversed-meridian intertwiner in the invariant run was
\[
 Q=\begin{pmatrix}2&-1&0&-2\\2&-1&0&-1\\0&2&1&3\\1&0&0&0\end{pmatrix}
   =P_{\rm mark}
 \begin{pmatrix}-1&0&0&0\\-6&0&-1&0\\6&-1&0&0\\2&1&1&1\end{pmatrix}.
\]
It satisfies $T_0Q=QA_1^{-1}$ and $T_1Q=QA_2^{-1}$, hence uses a
different meridian convention. The displayed transition matrix is part
of the notation translation; $Q$ is not identified with $P_{\rm mark}$.
The theta-quotient argument above identifies the marked holomorphic
families in the convention fixed here and retains all integral collar
translations.
\end{enumerate}
\end{proof}

\begin{proposition}[The relative cohomology comparison]\label{inv:B1.1}
\StatementRecord{inv:B1.1}{proved}{Historical events: Claude Fable 5.2 (Anthropic), effort high, 16 November 2026: carried as background without fault found; used and rechecked in this run. GPT-6 Astra (OpenAI), ChatGPT Pro mode, effort Pro, S6\_09\_PAPERS, 17 November 2026: targeted proof check. Also checked in this writing. Pre-continuation writing events: GPT-6 Astra (OpenAI), ChatGPT Pro mode, effort Pro, S6\_10\_PAPERS, 18 November 2026: targeted proof check. Continuation metadata events: GPT-6 Astra (OpenAI), ChatGPT Pro mode, effort Pro, S6\_10\_PAPERS continuation, 19 November 2026: used and rechecked in this run; status-vocabulary reconciliation; metadata reconciliation only; no new mathematical check. \emph{Pre-continuation recorded proof scope.} Read complete proofs and re-derived the marked theta-line transport, finite-root base changes and canonical characters, weighted section multiplication, integral fundamental-group quotient, cyclic-cover residue sequence, equivariant extra cohomology, and the rational Mayer--Vietoris ranks. Exact external theorem locators remain a separate citation gate.}{the hypotheses, notation and definitions stated here; the explicit constructions and calculations in the proof}

Fix the period family and cusp filling. Let $X'$ be obtained by any
admissible translation twists and multiplicities $m_1=3k$, $m_2=4l$,
where $k,l$ are positive integers. Write $D'=W+S'_1+S'_2$ and
\[
E'=\Omega^1_{X'}(\log D')/f'^*\Omega^1_B(\log\Sigma).
\]
The identity on the common smooth family induces isomorphisms on $B$
\[
R^qf_*\OO_X\simeq R^qf'_*\OO_{X'},\qquad
R^qf_*E\simeq R^qf'_*E'\quad(0\le q\le2).
\]
These isomorphisms identify the logarithmic connecting maps $\kappa_q$.
They also hold when the cusp collar is changed by an admissible translation.
\end{proposition}


\begin{proof}
\begin{enumerate}[label=\arabic*.,leftmargin=*]
\item We construct the comparison at every point. Over $B\setminus\Sigma$
the two families are identified, and a translation acts trivially on
$H^q(F,\OO_F)$ and $H^q(F,\Omega^1_F)$ for a complex torus $F$.
For the second assertion, translation fixes the invariant holomorphic
one-forms, and the first follows either from invariant Dolbeault forms
or the exterior algebra $H^\bullet(F,\OO_F)=\bigwedge^\bullet H^1(F,\OO_F)$.
These are statements about the induced maps on cohomology, including
their holomorphic variation.

\item At a point with original order $d=3$ or $4$, let
$\widetilde f:T\to\Delta_z$ be the smooth proper torus model with
coarse coordinate $q=z^d$. The new model is its pullback by $z=s^e$,
where $m=de$, followed by the free cyclic quotient of order $m$.
Smooth proper base change identifies the higher direct images upstairs
with the pullbacks of those for $T$. The relative logarithmic cotangent
bundle upstairs is precisely $\Omega^1_{T/\Delta_z}$ pulled back:
the base logarithmic direction is killed in the quotient defining $E'$.
The derivative of a translation contributes only a multiple of $ds$;
it consequently acts on this quotient as ordinary translation on relative forms.

\item Average over the cyclic action and choose holomorphic eigenframes.
If a frame for $H^q(\OO)$ or $H^q(\Omega^1)$ has character
$\zeta_d^{-a}$, $0\le a<d$, its original invariant extension is
$z^a\OO_{\Delta_q}$. On the order-$m$ cover its character is
$\zeta_m^{-ea}$, and its invariant extension is
$s^{ea}\OO_{\Delta_q}=z^a\OO_{\Delta_q}$.
Taking invariants is exact in characteristic zero; hence this equality
is an equality of the extensions inside the common punctured-disc
bundle, not only an equality of their ranks or degrees. A translation
twist induces the identity on the cohomology just used and changes none
of these lattices. This proves the comparison at both finite points,
including the full nonreduced fibres.

\item At the cusp the relative logarithmic forms have the global toric frame
$d\log x_1,d\log x_2$, modulo $d\log q$.
A toric gluing has the form
$x_i\mapsto c_i(q)q^{b_i}x_i$ and therefore changes $d\log x_i$ by
$(q c_i'/c_i+b_i)d\log q$. It induces the identity on the relative
logarithmic quotient. The structure-sheaf comparison is also the identity
on the normalization cochain complex. This verifies the cusp extension
and the assertion about translated collars.

\item The logarithmic cotangent exact sequences are isomorphic on the smooth
base, so their connecting maps agree there under these comparisons.
Their source and target sheaves are locally free by the local calculations
in Part A and the eigenframe calculation above. A holomorphic map between
vector bundles on a reduced curve which is zero on a dense open set is
zero everywhere. Thus the connecting maps agree globally.
\end{enumerate}
\end{proof}

\begin{proposition}[All admissible twists have the same Hodge numbers]\label{inv:B1.2}
\StatementRecord{inv:B1.2}{proved}{Historical events: Claude Fable 5.2 (Anthropic), effort high, 16 November 2026: carried as background without fault found; used and rechecked in this run. Claude Fable 5.2 (Anthropic), effort high, 17 November 2026: carried as background without fault found. GPT-6 Astra (OpenAI), ChatGPT Pro mode, effort Pro, S6\_09\_PAPERS, 17 November 2026: targeted proof check. Also checked in this writing. Pre-continuation writing events: GPT-6 Astra (OpenAI), ChatGPT Pro mode, effort Pro, S6\_10\_PAPERS, 18 November 2026: targeted proof check. Continuation metadata events: GPT-6 Astra (OpenAI), ChatGPT Pro mode, effort Pro, S6\_10\_PAPERS continuation, 19 November 2026: used and rechecked in this run; status-vocabulary reconciliation; metadata reconciliation only; no new mathematical check. \emph{Pre-continuation recorded proof scope.} Read complete proofs and re-derived the marked theta-line transport, finite-root base changes and canonical characters, weighted section multiplication, integral fundamental-group quotient, cyclic-cover residue sequence, equivariant extra cohomology, and the rational Mayer--Vietoris ranks. Exact external theorem locators remain a separate citation gate.}{the hypotheses, notation and definitions stated here; \cref{inv:B1.1}}

For every family in \ref{inv:B1.1} whose cusp is the one-component hexagonal filling,
the direct images of the absolute cotangent sheaves are
\[
\begin{array}{c|ccc}
 &R^0f'_*&R^1f'_*&R^2f'_*\\\hline
\Omega^1_{X'}&\OO(-2)&\OO\oplus\OO(-1)^2&\OO\oplus\OO(-1)\\
\Omega^2_{X'}&\OO(-2)\oplus\OO(-1)&\OO(-2)\oplus\OO(-1)^2&\OO
\end{array}
\]
and $R^qf'_*\OO_{X'}=(\OO,\OO\oplus\OO(-1),\OO(-1))$.
In particular every admissible $(3,4)$ member with
$|12l_0-4l_1-3l_2|=1$ has the Hodge diamond of Part A.
\end{proposition}


\begin{proof}
\begin{enumerate}[label=\arabic*.,leftmargin=*]
\item \ref{inv:B1.1} identifies the logarithmic direct images and their connecting maps.
In Part A the cokernel of $\kappa_0$ is
$\C_{p_1}\oplus\C_{p_2}$, and $\kappa_1$ is surjective with kernel
$\OO\oplus\OO(-1)^2$. These statements therefore hold for $X'$.
The residue sheaf of its normalized divisor has higher direct images
\[
R^0f'_*\OO_{\widetilde D'}=\C_{p_0}\oplus\C_{p_1}\oplus\C_{p_2},
\quad R^1f'_*\OO_{\widetilde D'}=\C_{p_1}\oplus\C_{p_2},
\quad R^2=0.
\]
Indeed the finite reductions are bielliptic, and the cusp normalization
is rational. The base logarithmic residue is evaluation with weights
$(1,m_1,m_2)$; its kernel is $\OO(1)(-\Sigma)=\OO(-2)$.
On $H^1$ the residue of the invariant class $\gamma^{0,1}$ is
$m_j\gamma^{0,1}|_{S'_j}\ne0$. The other eigenline has zero residue.
The local rank-one image of $\kappa_0$ is its kernel, as follows either
from the residue diagram of Part A or from the eigenframe calculation
in \ref{inv:B1.1}. Thus the two finite skyscrapers cancel exactly. This proves
the first row, with no surviving torsion.

\item The proper flat map is Gorenstein of relative dimension two. Relative
duality for $F=\Omega^1_{X'}$, using
$F^*\otimes\omega_{X'/B}=\Omega^2_{X'}\otimes f'^*\OO(2)$,
gives the second row because all sheaves in the first row are locally
free. Coherent Leray on $\Pone$ gives every Hodge number. In particular
$h^{p,0}=0$ for $p>0$ follows from these direct images and duality,
so no separate invariance assumption about absolute forms is needed.
\end{enumerate}
\end{proof}

\subsection{Higher multiplicities and the anticanonical ring}

\begin{proposition}[Canonical bundles for all coprime multiplicities]\label{inv:B2.1}
\StatementRecord{inv:B2.1}{proved}{Historical events: Claude Fable 5.2 (Anthropic), effort high, 16 November 2026: used and rechecked in this run. Claude Fable 5.2 (Anthropic), effort high, 16 November 2026 13:58 UTC: targeted proof check. GPT-6 Astra (OpenAI), ChatGPT Pro mode, effort Pro, S6\_09\_PAPERS, 17 November 2026: targeted proof check. Also checked in this writing. Pre-continuation writing events: GPT-6 Astra (OpenAI), ChatGPT Pro mode, effort Pro, S6\_10\_PAPERS, 18 November 2026: targeted proof check. Continuation metadata events: GPT-6 Astra (OpenAI), ChatGPT Pro mode, effort Pro, S6\_10\_PAPERS continuation, 19 November 2026: used and rechecked in this run; status-vocabulary reconciliation; metadata reconciliation only; no new mathematical check. \emph{Pre-continuation recorded proof scope.} Read complete proofs and re-derived the marked theta-line transport, finite-root base changes and canonical characters, weighted section multiplication, integral fundamental-group quotient, cyclic-cover residue sequence, equivariant extra cohomology, and the rational Mayer--Vietoris ranks. Exact external theorem locators remain a separate citation gate.}{the hypotheses, notation and definitions stated here; \citet[Theorem 1.1 and Proposition 2.3]{ViaclovskyLog2026}; \cref{inv:B1.1}; \cref{inv:B1.2}}

Let $X=X_{3k,4l}(u)$ with admissible translation data as in
\citet[Theorem 1.1 and Proposition 2.3]{ViaclovskyLog2026},
$\gcd(3k,4l)=1$, and any nonzero $u$ in its construction domain.
Then
\[
K_X=f^*\OO_B(-1)\otimes
\OO_X\bigl((k-1)S_1+(3l-1)S_2\bigr).
\]
Moreover $(h^{0,q})_{q=0}^3=(1,1,0,0)$, $h^{p,0}=0$ for $p>0$,
and $(h^{1,1},h^{1,2})=(2,1)$ for every such pair.
Thus these Hodge numbers do not distinguish the multiplicity families.
\end{proposition}


\begin{proof}
\begin{enumerate}[label=\arabic*.,leftmargin=*]
\item The cited proposition supplies exactly the free cyclic quotients used
in \ref{inv:B1.1}; its Theorem 1.1 supplies the stated sphere topology for its
admissible choices. \ref{inv:B1.2} computes the cohomology, with no use of that
topological conclusion.
For the canonical bundle we check the natural evaluation
$f^*f_*\omega_{X/B}\to\omega_{X/B}$.
The invariant relative determinant has weights $a_1=2k$ and $a_2=l$
under the clockwise generators of orders $3k$ and $4l$. The pullback
of a canonical form is $F(s)\,ds\wedge e$, where the character of
$ds\wedge e$ is $\zeta_m^{1+a}$. Therefore its invariant coefficients
are $s^{m-1-a}$ times an invariant function and a nowhere vanishing
unit. Since the quotient of total spaces is etale, the evaluation
vanishes to order $m-1-a$ on the reduced fibre downstairs. These orders
are $k-1$ and $3l-1$. At the cusp the form
$d\log x_1\wedge d\log x_2\wedge dq$ is a nowhere vanishing canonical
frame in every smooth toric chart, so no cusp component occurs.

\item By \ref{inv:B1.1} and relative duality,
$f_*\omega_{X/B}=(R^2f_*\OO_X)^*=\OO(1)$.
The evaluation divisor just computed therefore gives
$\omega_{X/B}=f^*\OO(1)\otimes
\OO((k-1)S_1+(3l-1)S_2)$. Tensoring by $f^*K_B=f^*\OO(-2)$
proves the formula.
\end{enumerate}
\end{proof}

\begin{proposition}[The complete graded anticanonical ring]\label{inv:B2.2}
\StatementRecord{inv:B2.2}{proved}{Historical events: Claude Fable 5.2 (Anthropic), effort high, 16 November 2026: carried as background without fault found. GPT-6 Astra (OpenAI), ChatGPT Pro mode, effort Pro, S6\_09\_PAPERS, 17 November 2026: targeted proof check. Also checked in this writing. Pre-continuation writing events: GPT-6 Astra (OpenAI), ChatGPT Pro mode, effort Pro, S6\_10\_PAPERS, 18 November 2026: targeted proof check. Continuation metadata events: GPT-6 Astra (OpenAI), ChatGPT Pro mode, effort Pro, S6\_10\_PAPERS continuation, 19 November 2026: used and rechecked in this run; status-vocabulary reconciliation; metadata reconciliation only; no new mathematical check. \emph{Pre-continuation recorded proof scope.} Read complete proofs and re-derived the marked theta-line transport, finite-root base changes and canonical characters, weighted section multiplication, integral fundamental-group quotient, cyclic-cover residue sequence, equivariant extra cohomology, and the rational Mayer--Vietoris ranks. Exact external theorem locators remain a separate citation gate.}{the hypotheses, notation and definitions stated here; \cref{inv:B2.1}}

Put $m_1=3k$, $m_2=4l$, and $d=4l+3k-kl$.
Choose integers $r,s$ with $m_2r+m_1s=1$ and set
$\mathcal L=\OO_X(rS_1+sS_2)$.
There is a graded algebra isomorphism
\[
\bigoplus_{n\ge0}H^0(X,-nK_X)
\simeq\bigoplus_{n\ge0}\C[x,y]_{dn},
\qquad\deg x=m_2,\quad\deg y=m_1,
\]
where the right-hand summand is assigned degree $n$.
For negative weights the polynomial space is zero. Explicitly
\[
h^0(X,-nK_X)=
\max\left\{1+n-\left\lceil\frac{n(k-1)}{3k}\right\rceil
-\left\lceil\frac{n(3l-1)}{4l}\right\rceil,\,0\right\}.
\]
The integer $d$ is nonzero for admissible pairs. If $d<0$ the
anticanonical ring is $\C$; if $d>0$ it is the indicated weighted
Veronese ring. For $(k,l)=(1,1)$ it is
$\C[U,V]$, $\deg U=1$, $\deg V=2$.
\end{proposition}


\begin{proof}
\begin{enumerate}[label=\arabic*.,leftmargin=*]
\item The divisor of $f^*(t/(t-1))$ is $m_1S_1-m_2S_2$.
It follows, retaining this rational function as the comparison, that
$\mathcal L^{m_2}\simeq\OO(S_1)$,
$\mathcal L^{m_1}\simeq\OO(S_2)$ and
$\mathcal L^{m_1m_2}\simeq f^*\OO(1)$.
\ref{inv:B2.1} then gives $-K_X=\mathcal L^d$.

\item We first establish the required section algebra for every integer
power of $\mathcal L$. If a meromorphic function represents a section
of $\OO(aS_1+bS_2)$, it is holomorphic on every regular compact fibre
and is constant there. It descends to the base on the smooth locus.
On a cyclic cover $q=s^m$, the finite pole bound gives a Laurent
expansion bounded below, and invariance leaves only powers divisible
by $m$. Thus descent extends meromorphically across the finite point,
with the exact valuation inequality $m\operatorname{ord}_q(h)+a\ge0$.
The identical argument at the reduced cusp has $m=1$.
Conversely each base function satisfying these bounds pulls back to
an allowed section. Consequently, as subsheaves of meromorphic functions,
\[
f_*\OO(aS_1+bS_2)=
\OO_B\left(\left\lfloor a/m_1\right\rfloor p_1+
\left\lfloor b/m_2\right\rfloor p_2\right).
\]
This proves the displayed dimension formula using \ref{inv:B2.1} and
$h^0(\Pone,\OO(e))=\max\{e+1,0\}$.

\item For the algebra, let $x,y$ be the canonical sections of
$\OO(S_1),\OO(S_2)$ under the chosen identifications with powers of
$\mathcal L$. Their powers $x^{m_1},y^{m_2}$ pull back two independent
sections of $\OO_B(1)$, with zeros at the two distinct points.
For any integer $j$, choose the unique residues
$0\le a_0<m_1$, $0\le b_0<m_2$ for which
$j-m_2a_0-m_1b_0=e m_1m_2$ is integral.
The valuation formula identifies $H^0(\mathcal L^j)$ with
$x^{a_0}y^{b_0}H^0(\OO_B(e))$. When $e\ge0$ a basis is
\[
x^{a_0+m_1i}y^{b_0+m_2(e-i)},\qquad0\le i\le e;
\]
when $e<0$ there are no sections. These are exactly the monomials of
weighted degree $j$. Their multiplication is ordinary polynomial
multiplication, proving the algebra assertion in every degree.
When $k=l=1$, the weight-$6n$ monomials are generated by $y^2$ and
$x^3$, of regraded degrees $1,2$.
Finally, $d=0$ would imply $k\mid4l$ and $l\mid3k$; coprimality
forces $k=l=1$, which instead gives $d=6$.
\end{enumerate}
\end{proof}

For completeness the requested small pairs give the following exact data.
Here $r_1,r_2$ are the coefficients in \ref{inv:B2.1} and $P_n^-=h^0(-nK_X)$.
\[
\begin{array}{c|c|c|c|c|ccc}
(k,l)&(m_1,m_2)&(r_1,r_2)&(\deg x,\deg y)&d&P^-_1&P^-_2&P^-_3\\\hline
(1,1)&(3,4)&(0,2)&(4,3)&6&1&2&2\\
(1,2)&(3,8)&(0,5)&(8,3)&9&1&1&2\\
(3,1)&(9,4)&(2,2)&(4,9)&10&0&1&1\\
(1,4)&(3,16)&(0,11)&(16,3)&15&1&1&1\\
(5,1)&(15,4)&(4,2)&(4,15)&14&0&1&1\\
(1,5)&(3,20)&(0,14)&(20,3)&18&1&1&1
\end{array}
\]
The program \texttt{b\_families.py} enumerates all coprime pairs with
$m_1m_2\le60$ and checks the floor formula against exact weighted
monomial counting for $0\le n\le12$. These checks accompany the
all-degree proof above.

\subsection{Finite fundamental groups and universal covers}

\begin{proposition}[Connected cyclic covers and the cusp comparison]\label{inv:B3.1}
\StatementRecord{inv:B3.1}{proved}{Historical events: Claude Fable 5.2 (Anthropic), effort high, 16 November 2026: carried as background without fault found; used and rechecked in this run. GPT-6 Astra (OpenAI), ChatGPT Pro mode, effort Pro, S6\_09\_PAPERS, 17 November 2026: targeted proof check. Also checked in this writing. Pre-continuation writing events: GPT-6 Astra (OpenAI), ChatGPT Pro mode, effort Pro, S6\_10\_PAPERS, 18 November 2026: targeted proof check. Continuation metadata events: GPT-6 Astra (OpenAI), ChatGPT Pro mode, effort Pro, S6\_10\_PAPERS continuation, 19 November 2026: used and rechecked in this run; status-vocabulary reconciliation; metadata reconciliation only; no new mathematical check. \emph{Pre-continuation recorded proof scope.} Read complete proofs and re-derived the marked theta-line transport, finite-root base changes and canonical characters, weighted section multiplication, integral fundamental-group quotient, cyclic-cover residue sequence, equivariant extra cohomology, and the rational Mayer--Vietoris ranks. Exact external theorem locators remain a separate citation gate.}{the hypotheses, notation and definitions stated here; \cref{inv:B3.3}; \cref{inv:B1.1}}

For an admissible $(3,4)$ member with
$n=|12l_0-4l_1-3l_2|>1$, let $\pi:Y\to X$ be the universal
cover and $g=f\pi$. It is a finite etale holomorphic map of degree $n$.
Every fibre of $g$ is connected; its finite reductions are bielliptic,
its multiplicities are $3,4$, and its cusp $W_n$ normalizes to a
disjoint union of $n$ copies of $\mathrm{dP}_6$.
The pullback maps induce isomorphisms
\[
R^qf_*\OO_X\simeq R^qg_*\OO_Y,\qquad
R^qf_*E_X\simeq R^qg_*E_Y.
\]
\end{proposition}


\begin{proof}
\begin{enumerate}[label=\arabic*.,leftmargin=*]
\item We first derive the fundamental-group presentation from the fillings.
The smooth family over the punctured base has zero section, so its
fundamental group is $\Lambda\rtimes\langle x,y\rangle$, where
$x,y$ are the analytic peripheral generators fixed by the root-deck convention and act by $A_1,A_2$.
On a finite filling the lifted affine generator has $m_j$th power
translation by $v_j$: its $m_j$ translations add to $v_j$ because
$A_jv_j=v_j$. Thus filling the normal disc imposes
$x^3=t_{v_1}$ and $y^4=t_{v_2}$, with the signs fixed by the
clockwise convention on the first page.

\item The toric covering space of the cusp tube is simply connected.
Indeed its dense torus has cocharacter lattice $\Z^3$; a loop may
be perturbed into that torus, and the normal discs of the three rays
$(0,1),(e_1,1),(e_2,1)$ kill a basis of this lattice.
This argument applies also to the infinite fan, since a compact loop
and a homotopy meet only finitely many of its locally finite charts.
The same conclusion holds for the small tube: if a loop in
$|t_c|<\epsilon$ bounds a disc in the toric space, the torus action
$(x_1,x_2,t_c)\mapsto(x_1,x_2,a t_c)$ for a sufficiently small
positive $a$ moves that compact disc into the tube, and the scaling
homotopy of its boundary stays in the tube.
The cusp quotient therefore has fundamental group
$\Lambda/\Lambda_{\rm tor}$, and its toric meridian bounds the
normal disc at fixed $x_1,x_2$. A translated collar makes this
meridian $t_\mu y^{-1}x^{-1}$, where $\gamma(\mu)=l_0$.
Van Kampen consequently kills $\Lambda_{\rm tor}$ and imposes
$xy=t_\mu$.

\item The normal closure of $\Lambda_{\rm tor}$ in the fibre lattice is
$\ker\gamma$: it contains $\hat w,\hat\delta$ and
$A_1\hat w+\hat w=\hat u$, and all these relations lie in
$\ker\gamma$. The remaining generator $c$, the image of
$\hat\gamma$, is central because $\gamma A_j=\gamma$.
The presentation is thus
$\langle c,x,y\mid c\text{ central},\ xy=c^{l_0},\
x^3=c^{l_1},\ y^4=c^{l_2}\rangle$.
Eliminating $y$ makes it abelian, with relation matrix
\[
 \begin{pmatrix}3&-l_1\\4&l_2-4l_0\end{pmatrix},
 \qquad\det=-p.
\]
Its first Smith divisor is one because its entries include three
and four; hence the group is $\Z/|p|$.
After imposing $c=1$ the quotient has $x^3=x^4=1$, so it is trivial.
Thus the image of the regular-fibre lattice, containing $c$, generates
the entire group $\Z/n$. The same is true of each reduced special
fibre. At a finite point the cyclic model is a smooth torus family
over a disc, with a free affine finite-group action. An equivariant
smooth trivialization, obtained by averaging a horizontal connection,
contracts the disc to its centre and descends to a deformation
retraction of the local total space onto $S_j$. Thus the image of
the regular-fibre lattice in the local fundamental group is contained
in the image of $\pi_1(S_j)$; the same holds after mapping to
$\pi_1(X)$. At the cusp $\hat\gamma$ survives in
$\Lambda/\Lambda_{\rm tor}$. The explicit normal-crossing contraction
of the tube onto $W$ given in \ref{inv:B3.3} identifies their fundamental groups,
so the same generator is in the image of $\pi_1(W)$.
The covering-space criterion therefore makes the inverse image of
every fibre connected. Admissibility gives $3\nmid l_1$ and odd $l_2$,
so $\gcd(n,12)=1$. The etale cover preserves fibre multiplicities.
On a finite reduction it cannot kill the nontrivial order-$3$ or
order-$4$ linear holonomy, since the covering degree is prime to it;
the resulting finite torus quotient is again bielliptic.

\item Normalization commutes with the finite etale base change $Y\to X$.
The pullback of the simply connected normalization $\mathrm{dP}_6$
is its trivial degree-$n$ cover. The double curves and triple points
also pull back etale, and their incidence maps are the lifts of the
original incidence maps. Consequently the dual complex is an actual
degree-$n$ covering of the original real two-torus, rather than merely
a cell complex with the same cell counts. It is connected because
$W_n$ is connected. The normalization resolution for $\OO_{W_n}$
has acyclic components, double curves and point strata. Its global
cochain complex is consequently the cellular cochain complex of this
dual torus, with $n,3n,2n$ cells. Thus its coherent cohomology is
$(1,2,1)$, and pullback from $W$ is an isomorphism over $\C$:
on a finite cover of real tori it is induced by a full-rank map of
their integral first cohomology lattices.

\item Modulo $d\log q$, the two toric forms $d\log x_i$ are invariant
under all cusp gluings as computed in \ref{inv:B1.1}. Hence
$E_X|_W=\OO_W^2$ and $E_Y|_{W_n}=\OO_{W_n}^2$,
and the same pullback argument applies to their cohomology.
Over a regular fibre, an isogeny of complex tori induces isomorphisms
on $H^q(\OO)$ and $H^q(\Omega^1)$ over $\C$.
At a finite fibre, make the cyclic base change. The map is an isogeny
of the smooth torus families and induces isomorphisms of their
cohomology bundles. Taking the same cyclic invariant eigenspaces
descends these isomorphisms, including all nilpotent layers of the
scheme fibre. Grauert base change at the cusp, and this explicit
calculation at the finite points, prove the asserted bundle maps are
isomorphisms on all of $B$.
\end{enumerate}
\end{proof}

\begin{proposition}[Cotangent direct images and Hodge numbers of the covers]\label{inv:B3.2}
\StatementRecord{inv:B3.2}{proved}{Historical events: Claude Fable 5.2 (Anthropic), effort high, 16 November 2026: used and rechecked in this run. Claude Fable 5.2 (Anthropic), effort high, 16 November 2026 13:58 UTC: targeted proof check. Claude Fable 5.2 (Anthropic), effort high, 17 November 2026: carried as background without fault found. GPT-6 Astra (OpenAI), ChatGPT Pro mode, effort Pro, S6\_09\_PAPERS, 17 November 2026: targeted proof check. Also checked in this writing. Pre-continuation writing events: GPT-6 Astra (OpenAI), ChatGPT Pro mode, effort Pro, S6\_10\_PAPERS, 18 November 2026: targeted proof check. Continuation metadata events: GPT-6 Astra (OpenAI), ChatGPT Pro mode, effort Pro, S6\_10\_PAPERS continuation, 19 November 2026: used and rechecked in this run; status-vocabulary reconciliation; metadata reconciliation only; no new mathematical check. \emph{Pre-continuation recorded proof scope.} Read complete proofs and re-derived the marked theta-line transport, finite-root base changes and canonical characters, weighted section multiplication, integral fundamental-group quotient, cyclic-cover residue sequence, equivariant extra cohomology, and the rational Mayer--Vietoris ranks. Exact external theorem locators remain a separate citation gate.}{the hypotheses, notation and definitions stated here; \cref{inv:B3.1}; \cref{inv:B1.2}}

Write $T_n=\C_{p_0}^{\,n-1}$. Then
\[
\begin{array}{c|ccc}
 &R^0g_*&R^1g_*&R^2g_*\\\hline
\Omega^1_Y&\OO(-2)&\OO\oplus\OO(-1)^2\oplus T_n&\OO\oplus\OO(-1)\\
\Omega^2_Y&\OO(-2)\oplus\OO(-1)&\OO(-2)\oplus\OO(-1)^2&\OO\oplus T_n
\end{array}
\]
The splittings involving $T_n$ are not canonical.
The only nonzero Hodge numbers of $Y$ are
\[
h^{0,0}=h^{0,1}=h^{3,2}=h^{3,3}=1,\quad
h^{1,2}=h^{2,1}=1,\quad h^{1,1}=h^{2,2}=n+1.
\]
The original finite-fundamental-group threefold $X$ itself has the
same Hodge diamond as the simply connected $(3,4)$ member.
\end{proposition}


\begin{proof}
\begin{enumerate}[label=\arabic*.,leftmargin=*]
\item \ref{inv:B3.1} identifies the logarithmic cotangent sequence on cohomology,
and therefore its connecting maps and direct images, with those of
the original family. The normalized residue sheaf differs only at
the cusp: its degree-zero cohomology there is now $\C^n$.
The residue of a base logarithmic form is the diagonal vector
$(1,\ldots,1)$, since each cusp component has multiplicity one.
Its cokernel is $T_n$. The degree-one residues at the finite fibres
cancel the same two skyscrapers as in \ref{inv:B1.2}. The long exact sequence
therefore gives
\[
0\longrightarrow T_n\longrightarrow R^1g_*\Omega^1_Y
\longrightarrow\OO\oplus\OO(-1)^2\longrightarrow0
\]
and the remaining entries of the first row. The extension splits
because $\operatorname{Ext}^1(V,T_n)=H^1(V^*\otimes T_n)=0$
for a vector bundle $V$ on $B$.

\item For clarity relative duality must now be used in its derived form.
For a sheaf $F$ locally free on the total space, it gives
\[
Rg_*(F^*\otimes\omega_{Y/B})
\simeq R\mathcal Hom_B(Rg_*F,\OO_B)[-2].
\]
For $F=\Omega^1_Y$, the torsion in $R^1g_*F$ contributes
$\mathcal Ext^1_B(T_n,\OO_B)$ in degree two of the right side.
This is a length-$(n-1)$ skyscraper at $p_0$. After tensoring by
$\OO(-2)$ it still has that length. The locally free parts dualize
in complementary degrees, giving exactly the second row; the
extension by the locally free quotient splits for the same reason
as above. Coherent Leray computes the displayed Hodge numbers.
For $X$ with one-component cusp, \ref{inv:B1.2} already applies independently
of its fundamental group.
\end{enumerate}
\end{proof}

\begin{proposition}[Canonical rings and rational topology of the covers]\label{inv:B3.3}
\StatementRecord{inv:B3.3}{proved}{Historical events: Claude Fable 5.2 (Anthropic), effort high, 16 November 2026: used and rechecked in this run. Claude Fable 5.2 (Anthropic), effort high, 16 November 2026 13:58 UTC: targeted proof check. GPT-6 Astra (OpenAI), ChatGPT Pro mode, effort Pro, S6\_09\_PAPERS, 17 November 2026: targeted proof check. Also checked in this writing. Pre-continuation writing events: GPT-6 Astra (OpenAI), ChatGPT Pro mode, effort Pro, S6\_10\_PAPERS, 18 November 2026: targeted proof check. Continuation metadata events: GPT-6 Astra (OpenAI), ChatGPT Pro mode, effort Pro, S6\_10\_PAPERS continuation, 19 November 2026: used and rechecked in this run; status-vocabulary reconciliation; metadata reconciliation only; no new mathematical check. \emph{Pre-continuation recorded proof scope.} Read complete proofs and re-derived the marked theta-line transport, finite-root base changes and canonical characters, weighted section multiplication, integral fundamental-group quotient, cyclic-cover residue sequence, equivariant extra cohomology, and the rational Mayer--Vietoris ranks. Exact external theorem locators remain a separate citation gate.}{the hypotheses, notation and definitions stated here; \cref{inv:B3.1}; \cref{inv:B2.2}; \cref{inv:B3.2}}

For $X$ and $Y$ as above,
\[
K_Y=\pi^*K_X=\OO_Y(-2S'_2),\quad
\bigoplus_{m\ge0}H^0(Y,-mK_Y)=\C[U,V],\quad
\deg(U,V)=(1,2).
\]
The same ring holds for $X$.
Over $\C$, the Betti numbers of $Y$ are
$(1,0,n-1,0,n-1,0,1)$; its Fr\"olicher sequence degenerates at $E_2$.
For every nontrivial character of the deck group, the character
part of Dolbeault cohomology consists of one class in $(1,1)$ and
one in $(2,2)$ and no others.
\end{proposition}


\begin{proof}
\begin{enumerate}[label=\arabic*.,leftmargin=*]
\item The canonical formula pulls back through an etale map, and
$4S'_2=g^*p_2$. \ref{inv:B3.1} gives $g_*\OO_Y=\OO_B$.
The valuation argument in \ref{inv:B2.2} therefore applies unchanged and proves
the ring statement with all multiplication maps.

\item The comparisons of \ref{inv:B3.1} are pullback maps and identify the logarithmic
cohomology bundles with the invariant part; the deck group acts
trivially on those bundles. At the cusp its action on the $n$
normalization components is the regular permutation representation.
The quotient by the diagonal has each nontrivial character once.
Thus the first-row additional skyscraper in \ref{inv:B3.2} has precisely these
characters, and duality gives the conjugate characters in its
second-row additional skyscraper. All other cohomology is invariant
and equal to that of $X$. Coherent Leray now proves the assertion
about the Dolbeault character spaces.

\item The de Rham complex and its Fr\"olicher sequence split under the
idempotents of the finite deck group. A nontrivial character has just
the two stated Dolbeault entries. No differential can connect them
or connect either to another nonzero entry, so both survive.
The invariant summand is the de Rham complex of $X$ by finite
etale descent and averaging. We compute its rational topology here
from the actual attachment maps; the integral homology assertion
of AC26 is not used as an input. In the remainder of this proof
all homology and cohomology coefficients are $\Q$.

\item Let $N$ be a small closed cusp tube, $M=\partial N$, and
$U=X\setminus\operatorname{int}N$. The space $U$ is covered by
two pieces retracting onto $S_1,S_2$, whose intersection retracts
onto a regular fibre $F$. The retractions are obtained from the
equivariant smooth torus trivializations used in \ref{inv:B3.1}. The maps
$F\to S_j$ are the central finite covering maps; translation of
their origins induces the identity on their rational cohomology.
Finite-cover transfer identifies their pullbacks with the invariant
subspaces for $A_j$. Put
\[
 \psi_1=2u+w+3\delta,\qquad \psi_2=u+w+2\delta,
 \qquad \eta=u\wedge w+6\gamma\wedge\delta.
\]
Direct substitution in $A_1,A_2$ shows that the invariant
degree-one spaces are $\langle\gamma,\psi_j\rangle$ and the
invariant degree-two spaces are
$\langle\gamma\wedge\psi_j,\eta\rangle$.
Their dimensions are two: the nontrivial eigenvalues of $A_j$ on
the complementary plane are the inverse primitive third or fourth
roots, so the exterior-square fixed space has exactly the two
displayed generators.
In the transfer-normalized dual homology bases, the two
Mayer--Vietoris maps are therefore the following actual matrices:
\[
 \alpha_1=\begin{pmatrix}
 1&0&0&0\\0&2&1&3\\-1&0&0&0\\0&-1&-1&-2
 \end{pmatrix},\qquad
 \alpha_2=\begin{pmatrix}
 2&1&3&0&0&0\\0&0&6&1&0&0\\
 -1&-1&-2&0&0&0\\0&0&-6&-1&0&0
 \end{pmatrix}.
\]
The second matrix uses the ordered wedge basis
$(\hat\gamma\hat u,\hat\gamma\hat w,\hat\gamma\hat\delta,
\hat u\hat w,\hat u\hat\delta,\hat w\hat\delta)$.
Both matrices have rank three, and
$\ker\alpha_1=\Q\nu_0$, with
$\nu_0=-\hat u-\hat w+\hat\delta$.
Consequently
\[
 H_1(U)=\Q,\qquad
 0\longrightarrow\operatorname{coker}\alpha_2=\Q
 \longrightarrow H_2(U)
 \xrightarrow{\partial}\Q\nu_0\longrightarrow0.       \tag{B.1}
\]
The generator $\Theta\in H^1(U)$ is normalized by
$\Theta|_F=\gamma$. The finite affine generators have
$3\bar g_1=l_1c$, $4\bar g_2=l_2c$ in rational homology, where
$c$ is the image of $\hat\gamma$ and $\Theta(c)=1$.
For the toric meridian $a_0=t_\mu y^{-1}x^{-1}$ this gives the
explicit collar value
\[
 \Theta(j_*a_0)=l_0-l_1/3-l_2/4=p/12.                 \tag{B.2}
\]

\item We justify the comparison with $N$ without assuming a collapse
theorem from AC26. In each normal-crossing chart
$t_c=z_1\cdots z_r$, $1\le r\le3$, choose the real vector field
whose $z_i$ component is
$-z_i|z_i|^{-2}/\sum_j|z_j|^{-2}$ off the divisor.
It satisfies $\dot t_c=-t_c$ and has norm at most
$C|t_c|^{1/r}$. A finite partition of unity preserves the first
identity and gives a global field with norm at most
$C'|t_c|^{1/3}$ after shrinking the tube. Extend it by zero on $W$.
Its flow exists for all positive time, and its tail length is
bounded by $3C'\epsilon^{1/3}e^{-T/3}$ after time $T$.
Thus the flow converges uniformly to $W$ and gives a continuous
strong deformation retraction fixing $W$. The same estimate proves
continuity at the divisor and at the endpoint of the retraction.

\item The constant-sheaf normalization resolution of $W$ has strata
$\mathrm{dP}_6$, three copies of $\Pone$, and two points.
Its row of degree-zero cohomology is
\[
 \Q\xrightarrow{0}\Q^3
 \xrightarrow{(a,b,c)\mapsto(a-b+c,a-b+c)}\Q^2.
\]
The strata have no degree-one cohomology. This resolution therefore
gives $H_1(N)=H_1(W)=\Q^2$, $H_4(N)=\Q[W]$ and $H_5(N)=0$;
it also gives $e(W)=6-3\cdot2+2=2$.
Poincare--Lefschetz duality for the compact smooth manifold with
boundary $(N,M)$ now shows that $H_1(M)\to H_1(N)$ and
$H_2(M)\to H_2(N)$ are surjective. For degree one the next
relative group is $H_1(N,M)=H^5(N)=0$.
For degree two the map to $H_2(N,M)=H^4(N)$ is the intersection
pairing with $[W]$, which is zero because $W$ is the divisor of
the global function $t_c$ on $N$. The Wang sequence for the
boundary mapping torus, with $\rank(M_0-I)=2$, gives
$\dim H_1(M)=3$. Its kernel to $H_1(N)$ is exactly
$\Q a_0$: this meridian is nonzero under projection to the base
circle and bounds its normal disc in $N$.

\item We exhibit two classes in
$K_2=\ker(H_2(M)\to H_2(N))$.
For $\nu=\hat w,\hat\delta$, let $\xi_\nu$ be the two-torus
obtained by translating the toric meridian around the angular
circle $S^1_\nu$. In the chart $t_c=z_0z_1z_2$ corresponding to
the rays $(0,1),(e_1,1),(e_2,1)$, these are the boundaries of
$D^2_{z_0}\times S^1_{z_i}$ with the other coordinate fixed
and the radii chosen so that $|t_c|=\epsilon$ on the boundary.
They therefore bound in $N$ with the literal toric meridian;
no choice of an unrecorded collar lift is being made.
Cutting the meridian into the two arcs of the cover defining
(B.1), the two cylinders have end circles differing by
$(I-A_1^{-1})\nu$. With this orientation of the Mayer--Vietoris
boundary, matrix multiplication gives
\[
 \partial j_*\xi_{\hat w}=-\nu_0,\qquad
 \partial j_*\xi_{\hat\delta}=0.                       \tag{B.3}
\]
The affine translation of either endpoint has no effect on its
homology class, which explains why this boundary formula is the
same for every translation twist.

\item Choose $\mathcal A\in H^2(U)$ restricting to $\eta$ on $F$;
it exists by Mayer--Vietoris because the two transfer-normalized
classes $\eta$ have equal restrictions. Fibrewise translation by
$S^1_{\hat\delta}$ defines a global smooth action on $U$:
$A_j\hat\delta=\hat\delta$, and all affine translations commute
with this circle. Slant product of its action map with the circle
defines $\iota:H^2(U)\to H^1(U)$. Naturality on $F$ gives
$\iota\mathcal A|_F=\iota_{\hat\delta}\eta=-6\gamma$.
The restriction $H^1(U)\to H^1(F)$ is injective by (B.1), so
$\iota\mathcal A=-6\Theta$. By the definition of slant product,
\[
 \langle\mathcal A,j_*\xi_{\hat\delta}\rangle
 =\langle\iota\mathcal A,j_*a_0\rangle=-p/2.           \tag{B.4}
\]
We may add the class pulled back from $(\Q\nu_0)^*$ under
the boundary map in (B.1) so that
$\langle\mathcal A,j_*\xi_{\hat w}\rangle=0$;
this does not change its restriction to $F$ or (B.4).
In the coordinates on $H_2(U)$ given by $\partial/\nu_0$ and
this normalized pairing, the two actual attachment cycles have
matrix
\[
 \bigl[j_*\xi_{\hat w}\ \ j_*\xi_{\hat\delta}\bigr]
 =\begin{pmatrix}-1&0\\0&-p/2\end{pmatrix}.             \tag{B.5}
\]
It is invertible over $\Q$ because admissibility implies
$p\ne0$. Thus the image of $K_2$ already fills $H_2(U)$.
Likewise (B.2) says the image of $\ker(H_1(M)\to H_1(N))$
is nonzero in $H_1(U)$. Using the surjectivity to $N$ just proved,
Mayer--Vietoris for $X=U\cup N$ now gives $H_1(X)=H_2(X)=0$.
Every other fibre has Euler number zero, so additivity over the
three special fibres and the smooth torus bundle gives $e(X)=2$.
Poincare duality then gives $b_4=b_5=0$ and $b_3=0$.
This proves the asserted rational homology of $X$ for all admissible
twists; the only changing determinant in (B.2)--(B.5) is a nonzero
constant multiple of $p$. The integer matrices and all exterior-power
identities were checked exactly in \texttt{code/b\_topology.py},
with output \texttt{outputs/b\_topology.txt}.

\item Part C's argument using this rational topology and the same Hodge
table now gives degeneration of the invariant summand at $E_2$.
Summing it with the nontrivial character spaces proves the Betti
numbers and the degeneration statement for $Y$. In particular
$e(Y)=2n$, as also follows from the covering degree.
\end{enumerate}
\end{proof}



\section{The integral circle quotient and its cusp replacement}\label{mapsI:circle-quotient}

We write $g,u,w,d$ for the ordered homology basis
$(\widehat\gamma,\widehat u,\widehat w,\widehat\delta)$ in this
section, and $\gamma$ for its first dual covector.  Thus $d$ is a
vector, whereas $\gamma$ is a covector.  The cusp coordinate is
$q=1/t$.  Positive target orientations are the complex orientation
on the base followed by the positive $\gamma$-circle orientation.
The complex orientation on $X$ fixes the source orientation.

\begin{proposition}[The primitive integral quotient and its finite fillings]\label{mapsI:src:A.1.1}
\StatementRecord{mapsI:src:A.1.1}{proved}{Historical events: Claude Fable 5.2 (Anthropic), effort high, 16 November 2026: carried as background without fault found. GPT-6 Astra (OpenAI), ChatGPT Pro mode, effort Pro, S6\_09\_PAPERS, 17 November 2026: targeted proof check. Also checked in this writing. Pre-continuation writing events: GPT-6 Astra (OpenAI), ChatGPT Pro mode, effort Pro, S6\_10\_PAPERS, 18 November 2026: targeted proof check. Continuation metadata events: GPT-6 Astra (OpenAI), ChatGPT Pro mode, effort Pro, S6\_10\_PAPERS continuation, 19 November 2026: used and rechecked in this run; status-vocabulary reconciliation; metadata reconciliation only; no new mathematical check. \emph{Pre-continuation recorded proof scope.} Read complete proof and re-derived invariant character, finite solid-torus coordinates, cusp nonextension paths, exact lens meridians, relative H1 extension, affine real involutions/cusp units, real cap-meridian determinant and orbit-space recognition. Full finite-congruence replays and Smith/triangulation external locators remain separate gates.}{the hypotheses, notation and definitions stated here; \cite{AC2026}}

For the markings and the clockwise root generators of the family construction,
the homological matrices and twisting vectors are
\[
 A_1=\begin{pmatrix}1&0&0&0\\6&0&1&0\\-6&-1&-1&0\\-2&1&0&1\end{pmatrix},
 \qquad
 A_2=\begin{pmatrix}1&0&0&0\\0&0&-1&0\\-6&1&0&0\\3&0&1&1\end{pmatrix},
\]
\[
 M_0=\begin{pmatrix}1&0&0&0\\0&1&0&0\\0&1&1&0\\-1&0&0&1\end{pmatrix},
 \qquad v_1=g+2u-4w,\qquad v_2=-g-3u+3w.
\]
The common invariant vectors and covectors are $\mathbb Zd$ and
$\mathbb Z\gamma$, respectively.  The smallest invariant primitive
sublattice containing the cusp vanishing lattice is
$L=\mathbb Z\langle u,w,d\rangle=\ker\gamma$.
The formula
\[
 [z,\Pi(z)x]\longmapsto([z],\gamma(x)\bmod\mathbb Z)
 \label{mapsI:quotient-eq-1}
\]
defines a real-analytic fibrewise group epimorphism
$J\longrightarrow B^\circ\times S^1$ with connected kernel $T_L$.
At the finite fillings it descends to a smooth proper $T^3$-bundle
over a Seifert solid torus, of multiplicities $3$ and $4$ and
translation residues $1$ and $-1$.
\end{proposition}



\begin{proof}
\begin{enumerate}[label=\arabic*.,leftmargin=*]
\item The first two displayed matrices are the inverse transposes of
the matrices printed on page 2 of \cite{AC2026}; their product
has inverse $M_0$.  Direct integral multiplication gives
$A_1^3=A_2^4=A_1A_2M_0=I$ and $A_jv_j=v_j$.
Solving the two fixed-vector equations gives $(0,0,0,d_0)$.
For a row vector $(a,b,c,e)$, the first fixed-covector equation
gives $b=2c,e=3c$ and the second gives $b=c,e=2c$;
therefore $b=c=e=0$.  Also
\[
 (M_0-I)(a,b,c,e)=(0,0,b,-a),\qquad
 A_1w=u-w.
\]
These identities prove the assertions about $L$ and the cusp
lattice $K=\mathbb Z\langle w,d\rangle$.  Their primitivity follows
from the displayed bases.  The flat real marking identifies the
deck action on a fibre with $x\mapsto A_jx$.  Since
$\gamma A_j=\gamma$, \eqref{mapsI:quotient-eq-1} descends.  Its kernel is connected
because $\gamma$ is primitive.

\item On a root disc the affine generator is
\[
 (s,x)\longmapsto
 (e^{-2\pi i/m}s,A_jx+v_j/m),\qquad (m,\kappa)=(3,1),(4,-1).
\]
It induces $(s,\theta)\mapsto
(e^{-2\pi i/m}s,\theta+\kappa/m)$ on the character quotient.
This action is free since $\gcd(\kappa,m)=1$.  Writing
$a=e^{2\pi i\theta}$ and choosing $r$ with $r\kappa\equiv1\pmod m$,
the invariant coordinates $s a^r$ and $a^m$ identify its quotient
with $D^2\times S^1$.  The base coordinate becomes
$(s a^r)^m(a^m)^{-r}$, and hence has multiplicity $m$ at its
central circle.  Before the finite quotient the map is a trivial
$T_L$-bundle.  The free actions on source and target are
equivariant for that map, so it descends locally to a smooth
$T_L$-bundle.  Properness follows from compactness of $T_L$.
The logarithmic translation used to glue the root chart to $J$
changes the character by its $\gamma$-component.  Applying
$\gamma$ to the actual gluing translation therefore gives the
target gluing, proving agreement on every overlap.
The exact checks, including the exhaustive mod-two fixed-vector
calculation used below, are in \texttt{code/maps\_audit/MAPS/code/maps\_exact.py}.
\end{enumerate}
\end{proof}


Engel's primitive invariant covector is named $\psi$ in his
Lemma~3.8 \cite[Lemma~3.8]{Engel2026}.  It is not named $\varepsilon$ there.  In his
common marking $\mathcal B_0$, the integral comparison is
\[
 P_E=\begin{pmatrix}0&-2&-1&-2\\2&-1&0&-1\\0&2&1&3\\-1&0&0&0\end{pmatrix},
 \qquad \psi P_E=-\gamma.
 \label{mapsI:quotient-eq-2}
\]
Thus the quotient characters agree after reversing the target
circle orientation.  With the inverse root generators the
translation vectors are $a_{j-1}=-P_Ev_j/m_j$.
This is the precise sign convention in Proposition~\ref{fam:src:B.2.2}; it explains why equality of the unnamed primitive
covectors without a sign convention would be insufficient.

\begin{proposition}[Pointwise cusp obstruction and the corrected double curve]\label{mapsI:src:A.1.2}
\StatementRecord{mapsI:src:A.1.2}{proved}{Historical events: Claude Fable 5.2 (Anthropic), effort high, 15 November 2026: carried as background without fault found. Claude Fable 5.2 (Anthropic), effort high, 16 November 2026: carried as background without fault found. GPT-6 Astra (OpenAI), ChatGPT Pro mode, effort Pro, S6\_09\_PAPERS, 17 November 2026: targeted proof check. Also checked in this writing. Pre-continuation writing events: GPT-6 Astra (OpenAI), ChatGPT Pro mode, effort Pro, S6\_10\_PAPERS, 18 November 2026: targeted proof check. Continuation metadata events: GPT-6 Astra (OpenAI), ChatGPT Pro mode, effort Pro, S6\_10\_PAPERS continuation, 19 November 2026: used and rechecked in this run; status-vocabulary reconciliation; metadata reconciliation only; no new mathematical check. \emph{Pre-continuation recorded proof scope.} Read complete proof and re-derived invariant character, finite solid-torus coordinates, cusp nonextension paths, exact lens meridians, relative H1 extension, affine real involutions/cusp units, real cap-meridian determinant and orbit-space recognition. Full finite-congruence replays and Smith/triangulation external locators remain separate gates.}{the hypotheses, notation and definitions stated here; \cref{mapsI:quotient-eq-4}}

The literal character has no continuous extension at either
triple point or at generic points of two of the three double
curves.  It has the constant limit $0$ on the smooth locus of
$W$ and on the third double curve away from the triple points.
In the chart $q=z_0z_1z_2$, $x_1=z_1$, $x_2=z_2$, the extendible
double curve is $z_0=z_1=0$, $z_2\ne0$.  The fixed curve of the
$\mathbb C^*$-action in direction $d=e_2$ is instead
$z_0=z_2=0$, $z_1\ne0$, and is one of the nonextendible curves.
Consequently the extendible double curve and the fixed double curve are distinct.
\end{proposition}



\begin{proof}
\begin{enumerate}[label=\arabic*.,leftmargin=*]
\item Put $q=e^{-2\pi T}$ on a positive radial ray and use the exact
cusp periods
\[
 Z_T=iTB_0+C(q),\qquad
 B_0=\begin{pmatrix}0&1\\-1&0\end{pmatrix},
\]
where $C$ is bounded and holomorphic at $0$.  For a fixed
$a\in\mathbb R^2$ the representative
$x_T(a)=\exp(2\pi iZ_Ta)$ has character $a_1\bmod\mathbb Z$ and
logarithmic position
\[
 \frac{\log|x_T(a)|}{\log|q|}
 =B_0a+T^{-1}\operatorname{Im}C(q)a.
 \label{mapsI:quotient-eq-3}
\]
The choices $a=(-1/4,1/4)$ and $a'=(-1/2,1/4)$ have logarithmic
positions tending to $(1/4,1/4)$ and $(1/4,1/2)$.
Both are interior to the triangle with vertices $0,e_1,e_2$,
so all three chart coordinates tend to $0$.  Their character
values are $3/4$ and $1/2$.  This proves nonextension at the
first triple point.  The other triangle type has vertices
$e_1,e_2,e_1+e_2$.  Its coordinates can be taken as
$x_1x_2/q,q/x_2,q/x_1$.  Positions $(3/4,3/4)$ and
$(3/4,1/2)$ lie in its interior, give distinct values of $-y_2$
modulo one, and converge to its origin.  Translation by the
lattice preserves triangle type; these are precisely the two
triple points of $W$.

\item For arbitrary convergent paths in the first chart, inversion
of the exact logarithmic period matrix gives
\[
 \theta_\gamma=-\frac{\log|z_2|}{\log|q|}+o(1)\pmod{\mathbb Z}.
 \label{mapsI:quotient-eq-4}
\]
Here the error tends to zero whenever the logarithmic ratios
stay bounded; they do stay bounded near a fixed point of a
normal-crossing stratum because $q=z_0z_1z_2$ and the three
coordinates are bounded.  If $z_0,z_1\to0$ and $z_2\to c\ne0$,
the right side tends to zero.  If $z_0,z_2\to0$ and $z_1=c\ne0$,
the paths $z_2=q^a,z_0=q^{1-a}/c$, $0<a<1$, have limits
$-a\pmod1$ at the same double point.  The same construction
with $z_0=c$ proves nonextension on $z_1=z_2=0$.
At a smooth point of a component the ratio in \eqref{mapsI:quotient-eq-4} tends
to $0$ or $1$, giving zero modulo one.  Integer translates
of these charts give the stated assertion on all of $W$.

\item The torus action in direction $e_2=d$ is
\[
 \lambda\cdot(z_0,z_1,z_2)
       =(\lambda^{-1}z_0,z_1,\lambda z_2).
\]
Its fixed curve is $z_0=z_2=0$.  This directly locates the
incorrect identification of the family construction.  Edge direction,
torus stabilizer, and coordinate on the quotient polygon must
not be interchanged.
\end{enumerate}
\end{proof}

The nonextension above is local and pointwise.  It does not
contradict extension of the cohomology class over a cusp
neighbourhood.  The toric covering has rays $(v,1)$ for all
$v\in\mathbb Z^2$.  Meridians about its toric divisors kill
the whole cocharacter lattice: $(0,1)$ is a ray and
$(e_i,1)-(0,1)=(e_i,0)$.  The toric cover is therefore simply
connected.  Its free period deck group is $\overline\Lambda$
and hence
\[
 \pi_1(N_0)=\overline\Lambda=\Lambda/K.
 \label{mapsI:quotient-eq-5}
\]
For the punctured cusp the map to this group kills $K$ and
the base meridian.  The latter bounds the disc given by a
point of the dense orbit on a central component.  Both
long-period covectors extend in integral $H^1$, since they
kill precisely these relations.  In contrast, an extension
of $\theta_\gamma$ to all of $X\simeq S^6$ would induce a
nonzero map $H_1(F)\to\mathbb Z$ factoring through
$H_1(X)=0$, which is impossible.  This last argument is the
global cohomological obstruction; it says nothing about
pointwise cusp limits.

\begin{proposition}[The filled three-dimensional target]\label{mapsI:src:A.1.3}
\StatementRecord{mapsI:src:A.1.3}{proved}{Historical events: Claude Fable 5.2 (Anthropic), effort high, 16 November 2026: carried as background without fault found. GPT-6 Astra (OpenAI), ChatGPT Pro mode, effort Pro, S6\_09\_PAPERS, 17 November 2026: targeted proof check. Also checked in this writing. Pre-continuation writing events: GPT-6 Astra (OpenAI), ChatGPT Pro mode, effort Pro, S6\_10\_PAPERS, 18 November 2026: targeted proof check. Continuation metadata events: GPT-6 Astra (OpenAI), ChatGPT Pro mode, effort Pro, S6\_10\_PAPERS continuation, 19 November 2026: used and rechecked in this run; status-vocabulary reconciliation; metadata reconciliation only; no new mathematical check. \emph{Pre-continuation recorded proof scope.} Read complete proof and re-derived invariant character, finite solid-torus coordinates, cusp nonextension paths, exact lens meridians, relative H1 extension, affine real involutions/cusp units, real cap-meridian determinant and orbit-space recognition. Full finite-congruence replays and Smith/triangulation external locators remain separate gates.}{the hypotheses, notation and definitions stated here; the explicit constructions and calculations in the proof}

For the original data the quotient over the complement of
the cusp disc is the exterior of the $(3,4)$ torus knot,
with the specified Seifert structure.  For general admissible
$(l_0,l_1,l_2)$ the filled quotient is the lens space determined
by the following meridian slopes in one oriented integral
torus basis $(\rho_1,h)$:
\[
 \mu_1=(3,-l_1),\qquad
 \mu_2=(-4,4l_0-l_2),\qquad
 \det(\mu_1,\mu_2)=p=12l_0-4l_1-3l_2.
 \label{mapsI:quotient-eq-6}
\]
In the lens-space convention that the second meridian is
$-q_0\mu_1+p\lambda_1$, where
$\det(\mu_1,\lambda_1)=1$, the lens parameter satisfies
$3q_0\equiv4\pmod p$.  In particular the target depends only
on the signed integer $p$, and is a three-sphere exactly
when $|p|=1$.
\end{proposition}



\begin{proof}
\begin{enumerate}[label=\arabic*.,leftmargin=*]
\item Let $P$ be the three-holed sphere, with oriented boundary
meridians $\rho_0+\rho_1+\rho_2=0$ and positive fibre $h$.
The two finite meridian relations are
$3\rho_1-l_1h=0$ and $4\rho_2-l_2h=0$; the cusp solid-torus
meridian is $\rho_0+l_0h$.  Eliminating $\rho_0$ gives
$\rho_2=l_0h-\rho_1$ on the annulus joining the finite solid
tori, which gives exactly \eqref{mapsI:quotient-eq-6}.  Since both slopes are
primitive, this is a genus-one Heegaard splitting.  Choose
$a,b\in\mathbb Z$ with $al_1+3b=1$ and put
$\lambda_1=(a,b)$.  Direct multiplication gives
\[
 \mu_2=-q_0\mu_1+p\lambda_1,
 \quad q_0=a(4l_0-l_2)+4b,
 \quad 3q_0-ap=4.
 \label{mapsI:quotient-eq-7}
\]
This proves the claimed lens parameter in a convention fixed
by actual meridians, including orientations.  The fundamental
group of this gluing is $\mathbb Z/|p|$.  Admissibility implies
$\gcd(p,6)=1$, so the modular inverse of $3$ used above exists.

\item For $(l_0,l_1,l_2)=(0,1,-1)$ the two exceptional-fibre data
are $(3,1),(4,-1)$.  They are realized by the circle action
\[
 s\cdot(z_1,z_2)=(e^{4is}z_1,e^{3is}z_2)
 \quad\text{on } S^3\subset\mathbb C^2.
\]
Its invariant map is $[z_1^3:z_2^4]$; the local circle
translation residues are $4^{-1}=1\pmod3$ and
$3^{-1}=-1\pmod4$.  Removing a regular orbit gives the
same two Seifert solid tori and their product annulus as
above.  A regular orbit is the $(3,4)$ torus knot, up to
the harmless exchange of the two coordinates.  This
identifies the manifold with its peripheral structure;
the abstract presentation $\langle x,y\mid x^3=y^{-4}\rangle$
alone would not have justified that identification.
\end{enumerate}
\end{proof}

\subsection{The repaired map and its relative uniqueness}

\begin{theorem}[The repaired map and its relative homotopy class]\label{mapsI:src:A.2.1}
\StatementRecord{mapsI:src:A.2.1}{proved}{Historical events: Claude Fable 5.2 (Anthropic), effort high, 16 November 2026: carried as background without fault found. GPT-6 Astra (OpenAI), ChatGPT Pro mode, effort Pro, S6\_09\_PAPERS, 17 November 2026: targeted proof check. Also checked in this writing. Pre-continuation writing events: GPT-6 Astra (OpenAI), ChatGPT Pro mode, effort Pro, S6\_10\_PAPERS, 18 November 2026: targeted proof check. Continuation metadata events: GPT-6 Astra (OpenAI), ChatGPT Pro mode, effort Pro, S6\_10\_PAPERS continuation, 19 November 2026: used and rechecked in this run; status-vocabulary reconciliation; metadata reconciliation only; no new mathematical check. \emph{Pre-continuation recorded proof scope.} Read complete proof and re-derived invariant character, finite solid-torus coordinates, cusp nonextension paths, exact lens meridians, relative H1 extension, affine real involutions/cusp units, real cap-meridian determinant and orbit-space recognition. Full finite-congruence replays and Smith/triangulation external locators remain separate gates.}{the hypotheses, notation and definitions stated here; \cref{mapsI:quotient-eq-5}}

For every admissible member, there is a smooth map
$e_{\rm q}:X\to Q$ over $B$ whose restriction outside an arbitrarily
small cusp neighbourhood is the literal circle quotient.
For fixed quotient gluing, its homotopy class is independent
of the cusp replacement, provided the replacement maps the
cusp neighbourhood into its designated solid torus and is
kept fixed on its outer collar.  In particular there is a
well-defined class for the original member $X\simeq S^6$.
\end{theorem}



\begin{proof}
\begin{enumerate}[label=\arabic*.,leftmargin=*]
\item On the cusp boundary use the product coordinates
$(q,\theta_\gamma)$ after the specified collar shear.
By \eqref{mapsI:quotient-eq-5}, the boundary circle class extends to $N_0$.
Since $S^1=K(\mathbb Z,1)$, it has a continuous representative
$\theta':N_0\to S^1$.  Its restriction to the boundary is
homotopic to the given one: the two maps induce the same
homomorphism on $\pi_1$ of the boundary mapping torus.
The homotopy-extension property for a manifold with boundary
changes $\theta'$ on a collar so that it agrees exactly with
$\theta_\gamma$ there.  Relative smooth approximation gives
a smooth representative, still equal on a smaller collar.
Then $(f,\theta')$ is the required map into $D^2\times S^1$,
and agrees with the finite-filling maps on the overlap.
This proves smoothness even at the critical points of $f$:
the product map is smooth without any submersion assertion.

\item The inclusion $\partial N_0\to N_0$ surjects on $\pi_1$,
as \eqref{mapsI:quotient-eq-5} explicitly shows.  It follows that
$H^1(N_0;\mathbb Z)\to H^1(\partial N_0;\mathbb Z)$
is injective.  Since both spaces are connected, the long
exact sequence gives $H^1(N_0,\partial N_0;\mathbb Z)=0$.
Two circle replacements agreeing on the boundary therefore
have quotient whose relative $H^1$ class is zero, hence
have a relative homotopy.  Keeping the base coordinate
fixed gives the required homotopy of maps into the solid
torus.  Different choices of the collar size are compared
on a common smaller collar.  This proves the assertion.
\end{enumerate}
\end{proof}




\subsection{The Albanese coordinates and real structures}
\begin{proposition}[exact nomes and continuation radius]\label{ari:C.2.1}
\StatementRecord{ari:C.2.1}{proved}{Historical events: Claude Fable 5.2 (Anthropic), effort high, 16 November 2026: used and rechecked in this run. Claude Fable 5.2 (Anthropic), effort high, 17 November 2026: carried as background without fault found. GPT-6 Astra (OpenAI), ChatGPT Pro mode, effort Pro, S6\_09\_PAPERS, 17 November 2026: targeted proof check. Also checked in this writing. Pre-continuation writing events: GPT-6 Astra (OpenAI), ChatGPT Pro mode, effort Pro, S6\_10\_PAPERS, 18 November 2026: computation or formal certificate; targeted proof check. Continuation mathematical and source checks: GPT-6 Astra (OpenAI), ChatGPT Pro mode, effort Pro, S6\_10\_PAPERS continuation, 19 November 2026: computation or formal certificate; independent re-proof; used and rechecked in this run. Continuation metadata events: GPT-6 Astra (OpenAI), ChatGPT Pro mode, effort Pro, S6\_10\_PAPERS continuation, 19 November 2026: used and rechecked in this run; status-vocabulary reconciliation; metadata reconciliation only; no new mathematical check. \emph{Continuation recorded proof scope.} Complete new mathematical derivation written before opening the target proof, followed by comparison of every new step. The fresh program \path{code/continuation_root/ari_C2_C3_fresh.py} passed 34 exact identities by a differential-polynomial and real-polynomial method. Ramanujan, theta and period inputs were opened with their exact normalizations; the renewed credit search and positive antecedents are recorded in the continuation comparison.}{the hypotheses, notation and definitions stated here; \cref{fam:src:A.4.1}; \cref{ari:C.1.2}; \cref{fam:src:C.2.3}, \cref{fam:src:A.5.1}; \cite{DuGuoKusakabeXie2026v1}}

For $w=e^{-\pi ic}$, the fixed branches give
\begin{align*}
 \kappa_1&=c+\tfrac12-\frac{3i}{2\pi}\log48,&
 q_1&=-w^2/48^3,\\
 \kappa_2&=c+\tfrac32-\frac{3i}{\pi}\log6,&
 q_2&=iw/216,\qquad q_1=(27/64)q_2^2.
\end{align*}
The family is parameterized by $0<|q_2|<1$ and exists precisely in
that domain. Writing $J(q)=q^{-1}+744+\cdots$, one has
\[
 j(\operatorname{Alb}S_1)=J(-w^2/48^3).
\]
This function is holomorphic on $0<|w|<48^{3/2}$, has a pole of order
two at zero, and has a natural boundary at $|w|=48^{3/2}$.
In particular it continues holomorphically through every point of the
family boundary $|w|=216$.
\end{proposition}


\begin{proof}
\begin{enumerate}[label=\arabic*.,leftmargin=*]
\item The precise Albanese lattices in \cref{fam:src:A.4.1} are verified by projecting all
translations on the product covers. At the order-three fibre the
projection is $(1/3)(\mathbb Z+\kappa_1\mathbb Z)$; multiplication by
three gives the first stated lattice. At the order-four fibre it is
$\tfrac12\mathbb Z+\tfrac14\kappa_2\mathbb Z$; multiplication by two
gives $\mathbb Z+(\kappa_2/2)\mathbb Z$. Thus the positive-imaginary
period representatives are $-\kappa_1$ and $-\kappa_2/2$.
Adding $2\mu_1=1-i/\sqrt3$ and
$3\mu_2=3/2-3i/2$ to \ref{ari:C.1.2} proves the two $\kappa$ identities,
including their real parts. Exponentiation gives the signs in $q_1,q_2$;
the ratio is $216^2/48^3=27/64$.

\item The identity $q_2=iw/216$ in Proposition~\ref{push:src:D.1.1} has
factor $1/2$ precisely because of the Albanese projection just computed.
Its natural-boundary proof uses modular invariance at rational boundary
points. The existence criterion and exact $2\mathbb Z$
identification are \cref{fam:src:C.2.3}, \cref{fam:src:A.5.1}; the latter is also Theorem~1.4
(label \texttt{p2:thm:exact-period}) of
\cite{DuGuoKusakabeXie2026v1}. Its hypotheses specify precisely the
retained period family, affine finite twists and cusp collar. They
coincide with the marked family $X_c$; the normalization differs
only by the fixed choice of the particular solution, which does not
alter parameter differences. Version 2 drops that theorem; we infer
no mathematical conclusion from this omission.

\item The Fourier expansion for $J$ converges for $0<|q|<1$ and has a simple
pole at $q=0$. The inverse image under $q=-w^2/48^3$ is the claimed
punctured disc and the pole has order two. The inequality
$216<48^{3/2}$ is equivalent to $27<64$. At every nonzero point on the
larger boundary the map to $q$ is locally biholomorphic. A continuation
there would continue $J$ through the unit circle, contradicting the
rational-cusp argument proved in Proposition~\ref{push:src:D.1.1}. The phrase ``up to the
radius'' therefore means the open punctured disc, not holomorphy on
its bounding circle or at its central pole.
\end{enumerate}
\end{proof}

\begin{proposition}[the real locus]\label{ari:C.2.2}
\StatementRecord{ari:C.2.2}{proved}{Historical events: Claude Fable 5.2 (Anthropic), effort high, 15 November 2026: carried as background without fault found. Claude Fable 5.2 (Anthropic), effort high, 16 November 2026: carried as background without fault found; used and rechecked in this run. GPT-6 Astra (OpenAI), ChatGPT Pro mode, effort Pro, S6\_09\_PAPERS, 17 November 2026: targeted proof check. Also checked in this writing. Pre-continuation writing events: GPT-6 Astra (OpenAI), ChatGPT Pro mode, effort Pro, S6\_10\_PAPERS, 18 November 2026: computation or formal certificate; targeted proof check. Continuation mathematical and source checks: GPT-6 Astra (OpenAI), ChatGPT Pro mode, effort Pro, S6\_10\_PAPERS continuation, 19 November 2026: targeted proof check. Continuation metadata events: GPT-6 Astra (OpenAI), ChatGPT Pro mode, effort Pro, S6\_10\_PAPERS continuation, 19 November 2026: used and rechecked in this run; status-vocabulary reconciliation; metadata reconciliation only; no new mathematical check. \emph{Pre-continuation recorded proof scope.} Re-derived the exact reflection intertwiners, primitive norm vectors, Laurent gluing, conjugation action on translation C/Z and the real square correction. The full-automorphism source pointer is being corrected to the credited local proof. Fully read and freshly replayed unchanged Cspecial\_reflection.py; exit zero. Output and SHA256 are retained in code/I/replay\_special/replay\_manifest.json. Numerical outputs retain numerical scope.}{the hypotheses, notation and definitions stated here; \cref{fam:src:B.4.2}; \cref{fam:src:A.5.1}; \cite{DuGuoKusakabeXie2026v1}}

An antiholomorphic equivalence $X_c\to X_{c'}$ exists exactly when
$c'-(1-\bar c)\in2\mathbb Z$. Consequently $X_c$ admits an
antiholomorphic involution exactly when $q_2\in\mathbb R\setminus\{0\}$.
\end{proposition}


\begin{proof}
\begin{enumerate}[label=\arabic*.,leftmargin=*]
\item We check the additional reflection argument, rather than infer a real
structure from the reality of one elliptic invariant. On the period
cover, reflection acts by
\[
 (\tau,\mu,\beta)\longmapsto(-\bar\tau,1-\bar\mu,-\bar\beta).
\]
The theta product and the cusp normalization prove this equality of
continued period functions. For the integral monodromy matrices
$A_1,A_2$ of the retained construction, the lattice comparison is
\[
 H=\begin{pmatrix}1&0&0&0\\0&1&0&0\\-6&0&-1&0\\0&-1&0&-1\end{pmatrix},
 \quad H^2=I,\quad HA_2H=A_2^{-1},\quad
 HA_1H=A_2^{-1}A_1^{-1}A_2,
\]
\[
 \Pi_{-\bar c}(\iota z)H=-\overline{\Pi_c(z)}.
\]
Thus $(z,\zeta)\mapsto(\iota z,-\bar\zeta)$ is an actual torus-family
map. Write $d=(0,0,0,1)^t$, $\varepsilon=(1,2,-4,0)^t$ and
$\varepsilon'=(1,3,-3,0)^t$. Reversing the finite deck generators
changes the target affine vectors to
\[
 v'_1=-\varepsilon+d,\qquad v'_2=\varepsilon'-3d.
\]
The second equality includes the sign of the inverse affine action.
The file \path{Cspecial_reflection.py} checks these matrix identities
over the integers.

\item Normalize this reflected datum to the retained vectors
$(\varepsilon,-\varepsilon')$. The linear fibre map is $-I_2$, the
parameter shift is $1$, and the lattice map is $-I+E$,
$E=d(1,0,0,0)$. The differences between target and transported affine
vectors are $2d$ and $-4d$. They lie in the images of the respective
norm operators. Explicit integral correcting cocycles are
$\ell_1=(0,0,-2,0)^t$ and $\ell_2=(0,2,0,0)^t$; the norm equations
are $N_1\ell_1=-2d$ and $N_2\ell_2=4d$. Rational solutions of
$(I-A_j)r_j=(v_j-(-I+E)v'_j)/m_j+\ell_j$ can be chosen with first
coordinate zero: $r_1=(0,-2/3,-2/3,0)^t$ and
$r_2=(0,1,1,0)^t$. These give translations holomorphic across both
root discs. Their cusp cocycle is
$\ell_0=-A_0(\ell_1+A_1\ell_2)=(0,0,4,-2)^t$; its coordinates are
integral, so the corresponding shear and inversion preserve the
periodic triangular cusp fan. They therefore extend across its toric
charts. The differences of the local corrections form a cocycle in
the rank-two vector bundle whose invariant line and quotient are
$\mathcal O$ and $\mathcal O(-1)$ on the base sphere. Both have zero
first cohomology, so the local maps glue. Equivalently one splits the
Laurent series of the cocycle into its nonnegative and negative
parts in these two line-bundle frames. This checks the precise
gluing step of Proposition~\ref{fam:src:B.4.2}, including the finite affine translations
and cusp extension. The resulting map has target $X_{1-\bar c}$.

\item Every other antiholomorphic comparison composed with its inverse is
holomorphic, so \cref{fam:src:A.5.1} gives exactly the congruence in the statement.
A self-map exists iff $\operatorname{Re}c\in1/2+\mathbb Z$, iff
$w=-\bar w$, iff $q_2=iw/216$ is real. Its square is a holomorphic
translation $d_0\in\mathbb C/\mathbb Z$, since the automorphism group
is the $\mathbb C^*$ computed in the credited re-derivation
of Proposition~\ref{fam:src:B.4.1}, Appendix~\ref{app:rederive:fam-src-B.4.1}. The constructed reflection acts on
this additive translation coordinate by conjugation. Commutation
with its square gives $d_0=\bar d_0$ in $\mathbb C/\mathbb Z$, hence
$d_0$ has a real representative. Composing the reflection with
translation by $-d_0/2$ makes its square the identity. Necessity and
sufficiency for an involution follow.
\end{enumerate}
\end{proof}

\subsection{The next local terms}
Put $G=g_2(i)=\Gamma(1/4)^8/(16\pi^2)$ and
$D_*=-3\log6/\pi$.
Our theta convention is
$\vartheta_j(z,\tau)=\theta_j(\pi z\mid\tau)$, with lattice
$\mathbb Z+\tau\mathbb Z$.  The heat identity and modular
transformations are \cite[Equations 20.13.1, 20.7.26, 20.7.30
and 23.18.5]{DLMFTheta2026}.  The classical differential system,
in the notation $P=E_2$, $Q=E_4$, $R=E_6$, is
\[
 xP_x=(P^2-Q)/12,\qquad
 xQ_x=(PQ-R)/3,\qquad
 xR_x=(PR-Q^2)/2;
\]
see \cite[Equations (1.1)--(1.4), p.~1]{ErbanVanGorder2021}.
Here $x=e^{2\pi i\tau}$, so
$d/d\tau=2\pi i x\,d/dx$.
For $E=\pi^2E_2/3$, $g_2=4\pi^4E_4/3$ and
$g_3=8\pi^6E_6/27$, this gives the exact normalization
\[
 \begin{split}
 E'&=\frac{i}{2\pi}(E^2-g_2/12),\\
 g_2'&=\frac{i}{\pi}(2Eg_2-3g_3),\qquad
 g_3'=\frac{i}{\pi}(3Eg_3-g_2^2/6).
 \end{split}
\]

\begin{proposition}[fourth order at the maximum]\label{ari:C.3.1}
\StatementRecord{ari:C.3.1}{proved}{Historical events: Claude Fable 5.2 (Anthropic), effort high, 16 November 2026: carried as background without fault found; used and rechecked in this run. GPT-6 Astra (OpenAI), ChatGPT Pro mode, effort Pro, S6\_09\_PAPERS, 17 November 2026: targeted proof check. Also checked in this writing. Pre-continuation writing events: GPT-6 Astra (OpenAI), ChatGPT Pro mode, effort Pro, S6\_10\_PAPERS, 18 November 2026: targeted proof check. Continuation mathematical and source checks: GPT-6 Astra (OpenAI), ChatGPT Pro mode, effort Pro, S6\_10\_PAPERS continuation, 19 November 2026: computation or formal certificate; independent re-proof; used and rechecked in this run. Continuation metadata events: GPT-6 Astra (OpenAI), ChatGPT Pro mode, effort Pro, S6\_10\_PAPERS continuation, 19 November 2026: used and rechecked in this run; status-vocabulary reconciliation; metadata reconciliation only; no new mathematical check. \emph{Continuation recorded proof scope.} Complete new mathematical derivation written before opening the target proof, followed by comparison of every new step. The fresh program \path{code/continuation_root/ari_C2_C3_fresh.py} passed 34 exact identities by a differential-polynomial and real-polynomial method. Ramanujan, theta and period inputs were opened with their exact normalizations; the renewed credit search and positive antecedents are recorded in the continuation comparison.}{the hypotheses, notation and definitions stated here; the explicit constructions and calculations in the proof}

On either root branch at $\tau=i$, let
$h=\mu-(1-\tau)/2$ and $\delta=\tau-i$. Then
\begin{align*}
 \tau&=i-\frac{3\pi i}{2}h^2-\frac{9\pi^2i}{8}h^4+O(h^6),\\
 h^2&=\frac{2i}{3\pi}\delta+\frac1{3\pi}\delta^2+O(\delta^3),\\
 \mu&=\frac{1-\tau}{2}
  \mathbin{\pm}\left(\frac{2i\delta}{3\pi}\right)^{1/2}
          \left(1-\frac{i}{4}\delta+O(\delta^2)\right).
\end{align*}
The square root is taken holomorphically on the chosen root chart.
Furthermore
\begin{align}\label{ari:eq:cs-Li}
 \log N={}&\tfrac12\log2-\frac\pi2|h|^2
   +\frac{3\pi^2}{8}|h|^4
   -\frac{3\pi^2}{4}|h|^2\operatorname{Re}(h^2)
   -\frac{5G}{192}\operatorname{Re}(h^4)+O(|h|^6),\\
 D_{\rm part}={}&D_* -\frac2\pi|\delta|
  +\frac1\pi\bigl(|\delta|^2+|\delta|\operatorname{Im}\delta\bigr)
  +\frac{5G}{72\pi^3}\operatorname{Re}(\delta^2)
  +O(|\delta|^3).\nonumber
\end{align}
In particular the leading terms agree with the boundary expansion in Lemma~\ref{push:src:D.2.1}.
\end{proposition}


\begin{proof}
\begin{enumerate}[label=\arabic*.,leftmargin=*]
\item Write $e(\tau)=\wp_\tau((1-\tau)/2)$ and
$E=\pi^2E_2/3$. The half-period Taylor series is
\[
 \wp_\tau((1-\tau)/2+h)
 =e+(3e^2-g_2/4)h^2+(3e^3-g_2e/4)h^4+O(h^6).
\]
The exact half-period coefficients used in Lemma~\ref{push:src:D.2.1} are
$e(i)=0$, $e'(i)=iG/(6\pi)$. This proves the claimed identity
$3e'(i)=(2i/\pi)G/4$. For clarity its closed form follows from
$e=\pi^2(\vartheta_2^4-\vartheta_4^4)/3$ and the heat equation:
if $A=\vartheta_2(i)^4=\vartheta_4(i)^4$, then
$G=4\pi^4A^2$ and $e'(i)=2\pi^3iA^2/3$.
For the classical closed value, \cite[Equation 23.5.3]{DLMFTheta2026}
gives $e_1=\Gamma(1/4)^4/(32\pi\omega_1^2)$,
$e_2=0$ and $e_3=-e_1$.  The lattice identification is
$2\omega_1=1$, $2\omega_3=i$; hence
$G=-4(e_1e_2+e_2e_3+e_3e_1)=4e_1^2
=\Gamma(1/4)^8/(16\pi^2)$.

\item Ramanujan's identities give $g_2'(i)=2iG$. Differentiating
$e(-1/\tau)=\tau^2 e(\tau)$ twice at $i$ gives
$e''(i)=3i e'(i)=-G/(2\pi)$. Inserting
$\delta=A_2h^2+A_4h^4+O(h^6)$ into the displayed half-period series
therefore gives
\[
 A_2=-3\pi i/2,\qquad
 A_4=\frac{g_2'A_2/4-e''A_2^2/2}{e'}=-9\pi^2i/8.
\]
The implicit-function theorem applies because $e'(i)\ne0$;
evenness in $h$ gives the stated remainder. Series inversion gives
the formulas for $h^2$ and $\mu$.

\item Let $H(h,\tau)=\log(\vartheta_3(h,\tau)/\eta(\tau))$. The exact
identity from the half-period translation is
$\log N=\operatorname{Re}H-\pi(\operatorname{Im}h)^2/
\operatorname{Im}\tau$. At $(0,i)$ its relevant Taylor data are
\[
 \begin{gathered}
 H_0=\tfrac12\log2,\quad H_\tau=0,\quad
 [h^2]H=-\pi/2,\quad [h^4]H=G/48,\\
 \partial_\tau[h^2]H=-\pi i/4-iG/(16\pi),\qquad
 H_{\tau\tau}=-G/(24\pi^2).
 \end{gathered}
\]
Indeed $\partial_h^2\log\vartheta_3=-\wp(h+(1-\tau)/2)-E$;
the heat equation at $h=0$ gives $H_\tau=ie/(4\pi)$, and
$E'(i)=\pi i/2-iG/(24\pi)$. These identities prove every entry.
After substitution $\tau=\tau(h)$, the pure $h^4$ coefficient of
$H$ is $-3\pi^2/8-5G/192$. Expanding
$\operatorname{Im}\tau=1-(3\pi/2)\operatorname{Re}(h^2)+O(|h|^4)$
in the Gaussian gives \eqref{ari:eq:cs-Li}. Finally
\[
 |h|^2=\frac{2}{3\pi}|\delta|
          \left(1+\tfrac12\operatorname{Im}\delta+O(|\delta|^2)\right)
\]
and substitution into the quartic expression, followed by
multiplication by $6/\pi$, gives the formula for $D_{\rm part}$.
The fresh exact differential-polynomial calculation in
\path{code/continuation_root/ari_C2_C3_fresh.py}
checks all these coefficients and the real quartic polynomial.
Its exact output is retained beside the program; the earlier theta
evaluation diagnostics remain in the computation record.
\end{enumerate}
\end{proof}

\begin{proposition}[the equianharmonic point is neither an extremum]\label{ari:C.3.2}
\StatementRecord{ari:C.3.2}{proved}{Historical events: Claude Fable 5.2 (Anthropic), effort high, 15 November 2026: carried as background without fault found. GPT-6 Astra (OpenAI), ChatGPT Pro mode, effort Pro, S6\_09\_PAPERS, 17 November 2026: targeted proof check. Also checked in this writing. Pre-continuation writing events: GPT-6 Astra (OpenAI), ChatGPT Pro mode, effort Pro, S6\_10\_PAPERS, 18 November 2026: targeted proof check. Continuation mathematical and source checks: GPT-6 Astra (OpenAI), ChatGPT Pro mode, effort Pro, S6\_10\_PAPERS continuation, 19 November 2026: computation or formal certificate; independent re-proof; used and rechecked in this run. Continuation metadata events: GPT-6 Astra (OpenAI), ChatGPT Pro mode, effort Pro, S6\_10\_PAPERS continuation, 19 November 2026: used and rechecked in this run; status-vocabulary reconciliation; metadata reconciliation only; no new mathematical check. \emph{Continuation recorded proof scope.} Complete new mathematical derivation written before opening the target proof, followed by comparison of every new step. The fresh program \path{code/continuation_root/ari_C2_C3_fresh.py} passed 34 exact identities by a differential-polynomial and real-polynomial method. Ramanujan, theta and period inputs were opened with their exact normalizations; the renewed credit search and positive antecedents are recorded in the continuation comparison.}{the hypotheses, notation and definitions stated here; the explicit constructions and calculations in the proof}

Set $T_0=\sqrt3/2$, $\mu_0=(2-\rho)/3$, and
$p=\wp'_\rho(\mu_0)$. Then $p\ne0$, $p^2=-g_3(\rho)$, and, for
$\delta=\tau-\rho$ on the order-three root chart,
\begin{align*}
 \mu&=\frac{2-\tau}{3}
      +\frac{3p}{16\pi^2}\delta^2
      +\frac{3ip}{32\pi^2T_0}\delta^3+O(\delta^4),\\
 t&=\frac{ip^2}{\pi^3}\delta^3
        \left(1+\frac{3i}{2T_0}\delta+O(\delta^2)\right),\\
 \log N&=\tfrac14\log3+\frac7{192}\operatorname{Re}t
       -\frac{9\pi}{512T_0|p|^{2/3}}|t|^{4/3}
       +O(|t|^{5/3}).
\end{align*}
Thus
\[
 D_{\rm part}(t)=-\frac3{2\pi}\log48
       +\frac7{32\pi}\operatorname{Re}t
       -\frac{27}{256T_0|p|^{2/3}}|t|^{4/3}
       +O(|t|^{5/3}).
\]
The point $t=0$ is neither a local maximum nor a local minimum.
\end{proposition}


\begin{proof}
\begin{enumerate}[label=\arabic*.,leftmargin=*]
\item We rederive the differentiation in Proposition~\ref{push:src:D.4.4}: put
$\ell=\partial_z\log\vartheta_1(\mu,\tau)$,
$a=\operatorname{Im}\mu/T$, $Z=\ell+2\pi ia$, and
$R=\mu'-a$. The theta heat equation and
$\wp=-\partial_z\ell-E$ give on the zero locus
\[
 \mu'=\frac{i\ell}{2\pi}-\frac{ig_2}{6\pi\wp'},\qquad
 R=\frac{iZ}{2\pi}-\frac{ig_2}{6\pi\wp'},\qquad
 \partial_\tau\log N=\frac{Z}{8\pi i}(Z+4\pi iR).
\]
These equations give the factored formula of Proposition~\ref{push:src:D.4.4}
with its stated normalizations. At $\rho$ the unit endomorphism fixing the
nonzero order-three point forces the gradient of $U_\rho$ there to
vanish, so $Z=0$. Explicitly the unit $\rho-1$ fixes $\mu_0$ modulo
the lattice and has no fixed nonzero real covector. Hence
$a=\mu'=-1/3$ and $\ell=2\pi i/3$.

\item Along the zero locus, a further heat-equation differentiation gives
\[
 \ell'=\frac{i\wp'}{4\pi}+\frac{iEg_2}{6\pi\wp'}.
\]
The Ramanujan derivatives at $\rho$ are
\[
 g_2'=\frac{3ip^2}{\pi},\quad
 g_2''=\frac{5i}{T_0}g_2',\quad
 \frac{(\wp')'}{\wp'}=\frac{3i}{2T_0},\quad E=\frac\pi{T_0}.
\]
They give $\mu''=3p/(8\pi^2)$,
$\mu'''=9ip/(16\pi^2T_0)$ and hence the asserted $\mu$ expansion.
Expansion of $t=g_2^3/(g_2^3-27g_3^2)$ gives the displayed $t$
formula. Since $\partial_\tau a=R/(2iT)$, one has
\[
 Z=\frac{ip}{4\pi}\delta+O(|\delta|^2),\qquad
 R=\frac{3p}{8\pi^2}\delta+O(|\delta|^2).
\]
The derivative formula now gives the harmonic cubic
$\operatorname{Re}(7ip^2\delta^3/(192\pi^3))=(7/192)
\operatorname{Re}t+O(|\delta|^4)$.

\item For the next term, direct differentiation of the Gaussian gives
\[
 \partial_\tau\partial_{\bar\tau}\log N
       =-\frac\pi{2T}|R|^2
       =-\frac{9|p|^2}{128\pi^3T_0}|\delta|^2+O(|\delta|^3).
\]
After the harmonic cubic has been written as a function of $t$, the
order-three stabilizer permits only a radial term in degree four:
the actual stabilizer map is $\gamma\tau=1-1/\tau$ and the zero
branch transforms as
$\mu(\gamma\tau)=(1-\mu(\tau))/\tau$.
The theta and eta transformations cited above cancel against the
Gaussian factor, so $N(\gamma\tau)=N(\tau)$.
Thus
in the linearizing coordinate
$(\tau-\rho)/(\tau-\bar\rho)$ invariant monomials have exponents
whose difference is divisible by three. The coefficient of
$|\delta|^4$ is therefore $-9|p|^2/(512\pi^3T_0)$, since
$\partial_\delta\partial_{\bar\delta}|\delta|^4=4|\delta|^2$.
Converting $|\delta|^4$ to $|t|^{4/3}$ proves the expansion, including
the remainder. Its nonzero linear term in $\operatorname{Re}t$
has both signs, which proves the last assertion. In particular
superharmonicity at this root point does not make it a local maximum.
\end{enumerate}
\end{proof}



\subsection{CM and torsion}
\begin{cited}[no smooth CM-torsion parameter]\label{ari:C.5.1}
\StatementRecord{ari:C.5.1}{proved}{Historical events: Claude Opus 5.5 (Anthropic), effort max, sweep 1, scope s6, 18 November 2026: used and rechecked in this run. GPT-6 Astra (OpenAI), ChatGPT Pro mode, effort Pro, S6\_09\_PAPERS, 17 November 2026: targeted proof check. This writing records metadata reconciliation only; no new mathematical check. Continuation metadata events: GPT-6 Astra (OpenAI), ChatGPT Pro mode, effort Pro, S6\_10\_PAPERS continuation, 19 November 2026: used and rechecked in this run; status-vocabulary reconciliation; metadata reconciliation only; no new mathematical check. \emph{Pre-continuation recorded proof scope.} Current-writing verification remains in progress. The earlier checks named above retain their recorded scope.}{the hypotheses, notation and definitions stated here; \citet[Theorem~4.15]{BourdonClarkStankewicz2017}; \citet[Theorem~20.12 and Corollary~20.13]{Sutherland2023CM}; \CompendiumRef{I}{app:rederive:ari-C.5.1}}

There is no $t\in\mathbb C\setminus\{0,1\}$ for which $E_t$ has
complex multiplication and $P_t$ is torsion. On the smooth models
at $t=0,1$ after base change of orders $3,4$, respectively, the
limiting points have orders $3,2$. Thus these are the only CM
torsion points on the completed family.
\end{cited}
The expanded source derivation is retained in the proof compendium, Paper~I, label \nolinkurl{ari:C.5.1}.




This applies \citet[Theorem~4.15]{BourdonClarkStankewicz2017} and \citet[Theorem~20.12 and Corollary~20.13]{Sutherland2023CM}; the exact finite enumeration is retained. A complete re-derivation is given in \CompendiumRef{I}{app:rederive:ari-C.5.1}.



\section{Global heights and the eleven rational CM fibres}

We use the equation
\[
 E_t:Y^2=X^3-\frac34t^3(t-1)X+\frac14t^4(t-1)^2,
 \qquad P_t=(0,t^2(t-1)/2).
\]
Its discriminant is $27t^8(t-1)^3$. Our canonical height is
$\widehat h_O(P)=\frac12\lim_{n\to\infty}n^{-2}h(x(nP))$.
In particular the bilinear Mordell--Weil self-pairing used in the
supplied surface calculation is $2\widehat h_O$.
For a finite place $v$ above the rational prime $p$, we normalize
$|p|_v=p^{-1}$ and set $n_v=[K_v:\Q_p]$. At an archimedean
place $n_v$ is one or two, according as the place is real or complex.

\begin{proposition}[the global formula]\label{ari:C.6.1}
\StatementRecord{ari:C.6.1}{proved}{Historical events: GPT-6 Astra (OpenAI), ChatGPT Pro mode, effort Pro, S6\_09\_PAPERS, 17 November 2026: targeted proof check. Also checked in this writing. Pre-continuation writing events: GPT-6 Astra (OpenAI), ChatGPT Pro mode, effort Pro, S6\_10\_PAPERS, 18 November 2026: targeted proof check; used and rechecked in this run. Continuation metadata events: GPT-6 Astra (OpenAI), ChatGPT Pro mode, effort Pro, S6\_10\_PAPERS continuation, 19 November 2026: used and rechecked in this run; status-vocabulary reconciliation; metadata reconciliation only; no new mathematical check. \emph{Pre-continuation recorded proof scope.} Read complete printed arithmetic proofs and re-derived the exact scaling/local-height equations and the genus-trace square denominator argument. Rank, saturation, directed L-derivative replays and precise external theorem matches remain explicit gates; no proof verdict is awarded to those computational claims by this read. Opened DeMarco–Mavraki1701.07947v2 Section2.1 pp5–6 and matched the canonical/local height normalization, continuity/pole properties and duplication discriminant term1/4 log|Delta|. This closes this source-normalization gate.}{the hypotheses, notation and definitions stated here; \cref{ari:C.7.1}; \citet[Section 2.1]{demarco2017}}

For every number field $K$, every $t\in K\setminus\{0,1\}$,
and the section just specified,
\begin{equation}\label{ari:eq:ch-global}
 \sum_{\sigma:K\hookrightarrow\C}\log N(P_{\sigma(t)})
 =\sum_{v\nmid\infty}n_v\lambda_v(P_t)
          -[K:\Q]\widehat h_O(P_t).
\end{equation}
Here the local heights are invariant under change of Weierstrass
model and satisfy, when $nP\ne O$,
\begin{equation}\label{ari:eq:ch-local-mult}
 \lambda_v(nP)=n^2\lambda_v(P)-\log|\psi_n(P)|_v
                   +\frac{n^2-1}{12}\log|\Delta|_v.
\end{equation}
At a complex embedding $\lambda_v(P)=-\log N(P)$; complex
embeddings are counted separately in \eqref{ari:eq:ch-global}.
For rational $t$, put $A_t=\exp(\sum_p\lambda_p(P_t))$.
Then $A_t$ is a positive real algebraic number and
\[
 N(P_t)=A_t e^{-\widehat h_O(P_t)},\qquad
 D_{\rm part}(t)=\frac6\pi
       \left(\log\frac{A_t}{2\sqrt3}-\widehat h_O(P_t)\right).
\]
Thus $N(P_t)$ is algebraic if and only if
$e^{\widehat h_O(P_t)}$ is algebraic.
\end{proposition}


\begin{proof}
\begin{enumerate}[label=\arabic*.,leftmargin=*]
\item The sigma expression for the archimedean local height is
\[
 \lambda_v(P)=-\log|\sigma(u)|
       +\frac12\Re(u\eta(u))-\frac1{12}\log|\Delta|.
\]
The sigma--theta relation and the Legendre relation identify this
with $-\log N$ exactly, including its constant; this is also the
calculation in Statement \ref{ari:C.7.1}. The division identity
$\sigma(nu)=(-1)^{n-1}\psi_n(P)\sigma(u)^{n^2}$ gives
\eqref{ari:eq:ch-local-mult}. At finite places the same identity is
the defining multiplication law of the canonical N\'eron function.
Its normalization is uniquely determined by continuity away from
$O$, its logarithmic pole at $O$, and the multiplication law.
These are the three properties printed in
\citet[Section 2.1]{demarco2017}; their duplication law includes
the discriminant term $\frac14\log|\Delta|_v$.
Under a coordinate scaling the changes in the sigma or formal
parameter term and in $\frac1{12}\log|\Delta|_v$ cancel.

\item Sum the duplication law over places. The product formula cancels
the division polynomial and discriminant terms, so that the sum
$H(P)=[K:\Q]^{-1}\sum_vn_v\lambda_v(P)$ satisfies
$H(2P)=4H(P)$. The pole normalization implies
$H(P)=\frac12h(x(P))+O_E(1)$. Iterating multiplication by two
therefore identifies $H$ with $\widehat h_O$. Moving the
archimedean terms to the other side gives \eqref{ari:eq:ch-global}.
For a finite place, a multiple of $P$ lies in the identity component
of a minimal model. There the local height is a rational multiple
of $\log p$; the multiplication law shows the same for $P$.
Only finitely many places contribute. This proves the assertion
about $A_t$ and both of its consequences.
\end{enumerate}
\end{proof}

\begin{proposition}[exact finite local-height calculation]\label{ari:C.6.2}
\StatementRecord{ari:C.6.2}{proved}{Historical events: GPT-6 Astra (OpenAI), ChatGPT Pro mode, effort Pro, S6\_09\_PAPERS, 17 November 2026: targeted proof check. Also checked in this writing. Pre-continuation writing events: GPT-6 Astra (OpenAI), ChatGPT Pro mode, effort Pro, S6\_10\_PAPERS, 18 November 2026: targeted proof check. Continuation metadata events: GPT-6 Astra (OpenAI), ChatGPT Pro mode, effort Pro, S6\_10\_PAPERS continuation, 19 November 2026: used and rechecked in this run; status-vocabulary reconciliation; metadata reconciliation only; no new mathematical check. \emph{Pre-continuation recorded proof scope.} Read complete printed arithmetic proofs and re-derived the exact scaling/local-height equations and the genus-trace square denominator argument. Rank, saturation, directed L-derivative replays and precise external theorem matches remain explicit gates; no proof verdict is awarded to those computational claims by this read.}{the hypotheses, notation and definitions stated here; \cref{ari:eq:ch-local-mult}}

Let $E/\Q$ be a global minimal Weierstrass model. If $mP$ is
in the identity component at $p$, and $mP\ne O$, then
\begin{equation}\label{ari:eq:ch-local-exact}
 \frac{\lambda_p(P)}{\log p}
 =\frac{v_p(\Delta)}{12}
   +\frac{\frac12\max(0,-v_p(x(mP)))-v_p(\psi_m(P))}{m^2}.
\end{equation}
For $m=1$ we read $\psi_1=1$. This formula, with exact rational
coordinates and valuations, gives all the finite contributions
in the tables below.
\end{proposition}


\begin{proof}
\begin{enumerate}[label=\arabic*.,leftmargin=*]
\item On the identity component of a minimal model the canonical
N\'eron function is
$\frac12\log\max(1,|x|_p)-\frac1{12}\log|\Delta|_p$.
This follows either from the formal parameter at $O$ followed
by the group law on the smooth locus, or from the uniqueness
properties used above. Substitute it for $\lambda_p(mP)$ in
\eqref{ari:eq:ch-local-mult} and divide by $m^2\log p$.
The contribution $v_p(\Delta)/(12m^2)$ combines with the
discriminant term to give $v_p(\Delta)/12$.
Membership in the identity component is checked by reduction
to the nonsingular locus of the minimal cubic. Thus the test is
exact also at $2$ and $3$; it does not use a formula restricted
to residue characteristic at least five.
\end{enumerate}
\end{proof}

\begin{proposition}[the CM models and their height identities]\label{ari:C.6.3}
\StatementRecord{ari:C.6.3}{proved}{Historical events: GPT-6 Astra (OpenAI), ChatGPT Pro mode, effort Pro, S6\_09\_PAPERS, 17 November 2026: targeted proof check. Also checked in this writing. Pre-continuation writing events: GPT-6 Astra (OpenAI), ChatGPT Pro mode, effort Pro, S6\_10\_PAPERS, 18 November 2026: computation or formal certificate; targeted proof check. Continuation metadata events: GPT-6 Astra (OpenAI), ChatGPT Pro mode, effort Pro, S6\_10\_PAPERS continuation, 19 November 2026: used and rechecked in this run; status-vocabulary reconciliation; metadata reconciliation only; no new mathematical check. \emph{Pre-continuation recorded proof scope.} Read complete printed arithmetic proofs and re-derived the exact scaling/local-height equations and the genus-trace square denominator argument. Rank, saturation, directed L-derivative replays and precise external theorem matches have their separately recorded scopes. Reran C\_cm\_height\_arb.py at 1152 bits with 350 Fourier terms. All eleven CM j intervals certify the printed integers; all prescribed uniformizing root brackets are proved by opposite interval signs, and all theta-height radii are below 1e-236. Reran the exact model/local-height programs in PARI2.17.2 and the unchanged eclib20250122 saturation code. All eleven displayed models and local coefficients match; ranks are exactly2 for rows1,8 and exactly1 for the remaining nine. All nine saturation indices are1, with no remaining primes. The full actual outputs and source hashes are retained. Opened PARI ellrank/ellheight normalization documentation and Cremona--Siksek Theorem1, author-PDFp8, for the algorithmic scope.}{the hypotheses, notation and definitions stated here; \cref{ari:eq:ch-local-exact}}

The eleven rational CM fibres have the global minimal models,
section coordinates, component classes, and finite local heights
in the following tables. On each model $a_1=a_2=0$, and we
write $[a_3,a_4,a_6]$ for
$y^2+a_3y=x^3+a_4x+a_6$. The displayed point is $(0,y_P)$.
The sign of $P$, which is immaterial to all these heights,
is chosen by the recorded change of variables.
The constants $A_t$ in increasing row order are
\[
 \sqrt6,\quad 2^{2/3}3^{1/4},\quad 2^{3/4},\quad
 2,\quad 2,\quad \sqrt3,\quad 1,\quad 3^{5/12},\quad
 1,\quad1,\quad1.
\]
\end{proposition}


\begin{proof}
\begin{enumerate}[label=\arabic*.,leftmargin=*]
\item We substitute the eleven exact values $t=j/1728$, and perform
the integral minimal-model transformations and Tate reduction
algorithm. The transformations and resulting invariants are
printed in \path{code/C_cm_models.out}; the exact valuation
computations are in \path{code/C_cm_local_heights.out}.
For each bad prime the program computes $mP$ by the rational
group law, tests the two partial derivatives of the minimal
cubic, and evaluates $\psi_m$ by its exact recursion. Formula
\eqref{ari:eq:ch-local-exact} then proves each entry. Outside the
listed primes the discriminant is a unit and $P$ is integral,
so the local height is zero. Multiplying $p^{\lambda_p/\log p}$
gives the listed constants $A_t$.

\item We specify the component labels without relying on an
unrecorded ordering of a resolution graph. The label $0$
means the identity component, and $c_P$ is the class represented
by the explicitly displayed section $P$. At an odd prime of
type $I_0^*$ this nonidentity class is the outer component
indexed by the root zero of the residual Tate cubic; the
scaling $x=pX$ identifies that root. When there are two
rational components, $c_P$ is their unique nonzero class.
The only remaining ambiguity is row 5 at $2$, where the
component group is $(\Z/2)^2$. Here let $T=(-14,0)$;
we label it with the basis $(c_P,c_T)$. Indeed
$P+T=(18,64)$ after the corresponding choice of sign,
and $P,T,P+T$ all reduce to the singular point, whereas
$2P=(25/4,99/8)$ reduces to the identity component.
Thus the three nonzero classes are explicitly separated.

\item The independent analytic check uses the exact quadratic CM
period ratio, solves $\wp(\mu)=0$ by its convergent Fourier
series, and evaluates the theta product at 270-digit working
precision. The value of $j$ agrees with its prescribed integer
to absolute error less than $10^{-245}$. For every row the
direct value of $\log N$ agrees with
$\log A_t-\widehat h_O(P_t)$ to more than 245 decimal places.
These numerical comparisons test the exact normalization;
the algebraic local-height entries were already proved by
integer and rational arithmetic. The full precision values
are in \path{records/C_cm_theta_pslq.json}.
\end{enumerate}
\end{proof}

\begingroup\footnotesize\setlength{\tabcolsep}{4pt}
\begin{longtable}{r r r r r r}
\toprule Row & $j$ & $a_3$ & $a_4$ & $a_6$ & $y_P$\\\midrule\endhead
1 & 8000 & 0 & -5880 & 153664 & 392\\
2 & 54000 & 0 & -7260 & 234256 & 484\\
3 & 287496 & 0 & -2156 & 38416 & 196\\
4 & 16581375 & 0 & -7732620 & 8275904784 & 90972\\
5 & -3375 & 0 & -140 & 784 & -28\\
6 & -32768 & 1 & -3234 & 72630 & -270\\
7 & -884736 & 1 & -38 & 90 & -10\\
8 & -12288000 & 1 & -1920270 & 1024288020 & -32005\\
9 & -884736000 & 1 & -379260 & 89898842 & -9482\\
10 & -147197952000 & 1 & -347045930 & 2488447882842 & -1577482\\
11 & -262537412640768000 & 1 & -75065352465446420 & 7916029533155605844266842 & -2813543945482\\
\bottomrule\end{longtable}
\begin{longtable}{r r l c c l}
\toprule Row & $p$ & Type & Component & Order & $\lambda_p/\log p$\\\midrule\endhead
1 & 2 & $III^*$ & $c_P$ & 2 & $\frac{1}{2}$\\
1 & 3 & $I_0^*$ & $0$ & 1 & $\frac{1}{2}$\\
1 & 7 & $I_0^*$ & $c_P$ & 2 & $0$\\
2 & 2 & $I_{4}^*$ & $c_P$ & 2 & $\frac{2}{3}$\\
2 & 3 & $III$ & $0$ & 1 & $\frac{1}{4}$\\
2 & 11 & $I_0^*$ & $c_P$ & 2 & $0$\\
3 & 2 & $I_{5}^*$ & $c_P$ & 2 & $\frac{3}{4}$\\
3 & 7 & $I_0^*$ & $c_P$ & 2 & $0$\\
4 & 2 & $I_{8}^*$ & $c_P$ & 2 & $1$\\
4 & 3 & $I_0^*$ & $c_P$ & 2 & $0$\\
4 & 7 & $III$ & $c_P$ & 2 & $0$\\
4 & 19 & $I_0^*$ & $c_P$ & 2 & $0$\\
5 & 2 & $I_{8}^*$ & $c_P$ & 2 & $1$\\
5 & 7 & $III$ & $c_P$ & 2 & $0$\\
6 & 3 & $I_0^*$ & $0$ & 1 & $\frac{1}{2}$\\
6 & 7 & $I_0^*$ & $c_P$ & 2 & $0$\\
6 & 11 & $III$ & $c_P$ & 2 & $0$\\
7 & 19 & $III$ & $c_P$ & 2 & $0$\\
8 & 3 & $IV$ & $0$ & 1 & $\frac{5}{12}$\\
8 & 11 & $I_0^*$ & $c_P$ & 2 & $0$\\
8 & 23 & $I_0^*$ & $c_P$ & 2 & $0$\\
9 & 3 & $I_0^*$ & $c_P$ & 2 & $0$\\
9 & 7 & $I_0^*$ & $c_P$ & 2 & $0$\\
9 & 43 & $III$ & $c_P$ & 2 & $0$\\
10 & 7 & $I_0^*$ & $c_P$ & 2 & $0$\\
10 & 31 & $I_0^*$ & $c_P$ & 2 & $0$\\
10 & 67 & $III$ & $c_P$ & 2 & $0$\\
11 & 7 & $I_0^*$ & $c_P$ & 2 & $0$\\
11 & 11 & $I_0^*$ & $c_P$ & 2 & $0$\\
11 & 19 & $I_0^*$ & $c_P$ & 2 & $0$\\
11 & 127 & $I_0^*$ & $c_P$ & 2 & $0$\\
11 & 163 & $III$ & $c_P$ & 2 & $0$\\
\bottomrule\end{longtable}
\endgroup
\noindent The next table prints 52 significant digits. The last column is both the directly computed $\log N$ and the independently computed $\log A_t-\widehat h_O(P_t)$; the agreement is more precise than the digits printed.
\begingroup\small
\begin{longtable}{r l}
\toprule Row & Height and archimedean logarithm\\\midrule\endhead
1 & $\widehat h_O=0.68444758838339190721754413638128047178526153154571$\\
 & $\log N=0.2114321462306355931886945428090706645762338145457924$\\[3pt]
2 & $\widehat h_O=0.6417998032172362080432024005106271224631326545214062$\\
 & $\log N=0.09495138932308808775043032302545534908240674117445132$\\[3pt]
3 & $\widehat h_O=0.5335302122611885050129127004864161587766141055417204$\\
 & $\log N=-0.01366982684122952294998860939278373271998900477152897$\\[3pt]
4 & $\widehat h_O=0.9940668685687274300038154462317330832788241518412543$\\
 & $\log N=-0.300919688008782120586583324773556515203324017480999$\\[3pt]
5 & $\widehat h_O=0.4673656771814712800075202012628467836356203437033082$\\
 & $\log N=0.2257815033784740294097119201953297844398797906569471$\\[3pt]
6 & $\widehat h_O=0.4286091723813209970260194687110167110834518087280487$\\
 & $\log N=0.120696971952733848671603149750246141240293470183326$\\[3pt]
7 & $\widehat h_O=0.09064583701789963707894077069188164993051948122321821$\\
 & $\log N=-0.09064583701789963707894077069188164993051948122321821$\\[3pt]
8 & $\widehat h_O=0.7366208256267509672496883390906560396314066510934334$\\
 & $\log N=-0.2788657053483719291683361570396036626949522520006211$\\[3pt]
9 & $\widehat h_O=0.6018296198896332298783588845164397135769872973598109$\\
 & $\log N=-0.6018296198896332298783588845164397135769872973598109$\\[3pt]
10 & $\widehat h_O=1.002649286144114870320961646648539589847474352238484$\\
 & $\log N=-1.002649286144114870320961646648539589847474352238484$\\[3pt]
11 & $\widehat h_O=2.165489246796061077105970244492001901750597685962121$\\
 & $\log N=-2.165489246796061077105970244492001901750597685962121$\\[3pt]
\bottomrule\end{longtable}\endgroup


\begin{proposition}[the maximum as a height inequality]\label{ari:C.6.4}
\StatementRecord{ari:C.6.4}{proved}{Historical events: Claude Fable 5.2 (Anthropic), effort high, 15 November 2026: carried as background without fault found. GPT-6 Astra (OpenAI), ChatGPT Pro mode, effort Pro, S6\_09\_PAPERS, 17 November 2026: targeted proof check. Also checked in this writing. Pre-continuation writing events: GPT-6 Astra (OpenAI), ChatGPT Pro mode, effort Pro, S6\_10\_PAPERS, 18 November 2026: targeted proof check; used and rechecked in this run. Continuation metadata events: GPT-6 Astra (OpenAI), ChatGPT Pro mode, effort Pro, S6\_10\_PAPERS continuation, 19 November 2026: used and rechecked in this run; status-vocabulary reconciliation; metadata reconciliation only; no new mathematical check. \emph{Pre-continuation recorded proof scope.} Read complete printed arithmetic proofs and re-derived the exact scaling/local-height equations and the genus-trace square denominator argument. Rank, saturation, directed L-derivative replays and precise external theorem matches remain explicit gates; no proof verdict is awarded to those computational claims by this read. Opened DeMarco–Mavraki Theorem1.1 p2, Theorems2.2–2.3 pp7–8 and the metric/proof loci pp15–16,23; matched the positive generic height1/12 and the finite-place/continuous-remainder application. Other explicitly recorded CM rank or directed-numerical gates are unchanged. Opened Silverman author slides21–22 (PDF26–27), matching the uniform height bound and specialization limit with their hypotheses. The exact DeMarco–Mavraki application and its bounded local remainder remain the operative proof. The original1983Silverman and Tate articles remain explicitly unverified background records.}{the hypotheses, notation and definitions stated here; \cref{ari:C.8.5}; \cref{ari:eq:ch-global}; \citet{demarco2017}; \citet[pp.~183--184]{callSilverman1993}; \citet{silverman2006heights}}

For $t\in K\setminus\{0,1\}$,
\begin{equation}\label{ari:eq:ch-bound}
 \widehat h_O(P_t)\geq
 \frac1{[K:\Q]}\sum_{v\nmid\infty}n_v\lambda_v(P_t)
                    -\frac12\log2.
\end{equation}
The inequality is strict on this smooth base. Over $\Q$ it
is $\widehat h_O(P_t)\geq\log A_t-\frac12\log2$.
The generic height is $\widehat h_O(P)=1/12$, and the
variation divisor is $D_P=(\infty)/12$. Consequently
\[
 \widehat h_O(P_t)=\frac1{12}h(t)+O(1)
 \qquad(t\in\overline\Q\setminus\{0,1\}),
\]
where the constant is uniform over algebraic $t$ and is for
the fixed family and section.
\end{proposition}


\begin{proof}
\begin{enumerate}[label=\arabic*.,leftmargin=*]
\item Apply the analytic maximum theorem \ref{ari:C.8.5} at every embedding
and substitute into \eqref{ari:eq:ch-global}. Equality would
require $j(\sigma(t))=1728$ for every embedding, hence $t=1$,
which has been excluded.

\item For the generic local heights at $0$ and $1$, the minimal
discriminant orders are $8$ and $3$, and the section meets
the nonidentity component of types $IV^*$ and $III$.
Their local corrections are $2/3$ and $1/4$, respectively,
so their normalized local heights are zero. At infinity,
put $s=1/t$, $X=s^{-2}x$, $Y=s^{-3}y$. The new equation is
\[
 y^2=x^3-\tfrac34(1-s)x+\tfrac14(1-s)^2,
 \qquad P=(0,(1-s)/2),\qquad\Delta=27s(1-s)^3.
\]
At $s=0$ the node is $(1/2,0)$, whereas $P=(0,1/2)$ is
nonsingular. Its height contribution is $1/12$.
At all other function-field places the point is integral
and the discriminant is a unit. This proves
$D_P=(\infty)/12$ and the generic height.

\item Theorem 1.1 of \citet{demarco2017} applies with $K=\Q$,
$B=\mathbb P^1$, the displayed elliptic surface and rational
section, since the generic height is positive. It equips
this divisor with continuous adelic semipositive metrics
whose associated height is exactly $\widehat h_O(P_t)$.
Comparing with the usual height for $(\infty)/12$ gives the
bounded difference. This is the specialization here of the
variation theorem of Tate, with the local refinements of
Silverman reproduced as Theorems 2.2 and 2.3 in that paper:
outside finitely many places the local correction is zero,
and at each remaining place it extends continuously after
subtracting its prescribed logarithmic singularity.
Our inequality is an explicit archimedean estimate with
constant $\frac12\log2$; the general bounded-variation
statement alone does not supply that constant.

\item We spell out the comparison with the opened primary author sources.
Theorem 4.1 of \citet[pp.~183--184]{callSilverman1993} concerns a
smooth projective base curve $T$, a degree-one base height $h_T$,
a section, and a fibrewise morphism whose pullback multiplies the
chosen divisor class by a real number $\alpha>1$. It gives
\[
 \lim_{h_T(t)\to\infty}
 \frac{\widehat h_t(P_t)}{h_T(t)}
       =\widehat h_{\Q(T)}(P).
\]
Here take $T=\mathbb P^1$, the class of $(O)$, and the generic
morphism $[2]$, with $[2]^*(O)\sim4(O)$ and $\alpha=4$.
The smooth fibres are precisely the complement of $0,1,\infty$;
the section extends on the regular projective elliptic surface.
All the hypotheses hold, and the limit is $1/12$.
The unnumbered remark immediately following that theorem states
Tate's stronger elliptic-family formula
$\widehat h_t(P_t)=h_{T,D_P}(t)+O_P(1)$, with
$\deg D_P=\widehat h_{\Q(T)}(P)$. This is an author-verified
reproduction of Tate's result, not a claim to have read the
inaccessible original scan \citep{tate1983variation}.
The same distinction applies to \citet{silverman1983specialization}:
the limit theorem is reproduced by the author on slide~22 of
\citet{silverman2006heights}, while slide~21 states the uniform
canonical-versus-Weil-height bound in terms of the base height.
For our explicit bounded remainder and local normalization, the
fully opened Theorems~1.1, 2.2 and 2.3 of
\citet{demarco2017}, applied above, remain the cited proof.
\end{enumerate}
\end{proof}

\begin{proposition}[rigorous high-precision height enclosures]\label{ari:C.6.5}
\StatementRecord{ari:C.6.5}{proved}{Historical events: GPT-6 Astra (OpenAI), ChatGPT Pro mode, effort Pro, S6\_09\_PAPERS, 17 November 2026: targeted proof check. Also checked in this writing. Pre-continuation writing events: GPT-6 Astra (OpenAI), ChatGPT Pro mode, effort Pro, S6\_10\_PAPERS, 18 November 2026: targeted proof check. Continuation metadata events: GPT-6 Astra (OpenAI), ChatGPT Pro mode, effort Pro, S6\_10\_PAPERS continuation, 19 November 2026: used and rechecked in this run; status-vocabulary reconciliation; metadata reconciliation only; no new mathematical check. \emph{Pre-continuation recorded proof scope.} Read complete printed arithmetic proofs and re-derived the exact scaling/local-height equations and the genus-trace square denominator argument. Rank, saturation, directed L-derivative replays and precise external theorem matches remain explicit gates; no proof verdict is awarded to those computational claims by this read.}{the hypotheses, notation and definitions stated here; \cref{ari:C.6.3}; \cref{ari:C.5.1}}

For every one of the eleven rational CM parameters in \ref{ari:C.6.3},
the files \path{records/C_cm_height_arb.json} and
\path{code/C_cm_height_arb.out} give certified enclosures of
$\log N(P)$ and $\widehat h_O(P)$ of radius less than
$10^{-236}$. These enclosures use no floating-point height
value as a proof step.
\end{proposition}


\begin{proof}
\begin{enumerate}[label=\arabic*.,leftmargin=*]
\item The CM discriminants, in the order of \ref{ari:C.6.3}, are
\[
 -8,-12,-16,-28,-7,-11,-19,-27,-43,-67,-163.
\]
Enumeration of primitive reduced forms
$(a,b,c)$, $|b|\le a\le c$, $b^2-4ac=D$, with the usual
nonnegative boundary convention, gives exactly one form of
each discriminant, always with $a=1$.
The range is finite and proved by $a\le\sqrt{|D|/3}$.
Thus the CM theorem used in \ref{ari:C.5.1} makes the corresponding
singular modulus an integer. The Fourier enclosures below
place it within $1/4$ of the specified $j$, proving its
identity, rather than only numerical closeness. Exact class
polynomials were also computed independently in
\path{code/C_cm_class_polynomials.gp} and its output.

\item Write $\tau=iT$ for the first four entries and
$\tau=1/2+iT$ for the other seven, with
\[
 (T_1,T_2,T_3,T_4)=(\sqrt2,\sqrt3,2,\sqrt7),\qquad
 T_k=\sqrt{|D_k|}/2\quad(k\ge5).
\]
Put $Q=e^{-2\pi T}$, $q=Q$ or $-Q$ accordingly,
$M=350$, and $Y=2/5$. On $0\le y\le Y$,
\begin{equation}\label{ari:eq:ch-wp-bracket}
 \frac{\wp(1/2-iy)}{\pi^2}
 =\frac1{\cosh^2\pi y}-\frac13+
 8\sum_{n\ge1}\frac{nq^n}{1-q^n}
                 \{1-(-1)^n\cosh(2\pi ny)\}.
\end{equation}
For $a=Qe^{2\pi Y}$, the absolute tail after $M$ in
$\wp$ is at most
\[
 \frac{16\pi^2}{1-Q}
 \frac{a^{M+1}\{M+1-Ma\}}{(1-a)^2}<10^{-875}.
\]
The proof is $1+\cosh(2\pi ny)\le2e^{2\pi nY}$ and
$\sum_{n>M}na^n=a^{M+1}(M+1-Ma)/(1-a)^2$.
Each previously computed decimal supplies only a rational
candidate centre $y_0$. At the two rational endpoints
$y_0\pm10^{-240}$ the truncated expression plus this
infinite tail is respectively strictly positive and strictly
negative. The code records all twenty-two strict signs.
Continuity produces a zero within each bracket. Every zero
of $\wp$ is one of the pair $\pm P$, and $N$ agrees on
that pair, so no unproved choice of root is needed.

\item The product formula on this bracket has positive real factors:
\begin{align*}
 \log N={}&\log(2\cosh\pi y)-\frac{\pi T}{6}
                      -\frac{\pi y^2}{T}\\
 &+\sum_{n\ge1}\{\log(1+q^ne^{2\pi y})
                         +\log(1+q^ne^{-2\pi y})\}.
\end{align*}
Since $|\log(1+x)|\le |x|/(1-|x|)$ for real $|x|<1$,
the omitted tail is bounded by
\[
 \frac{2e^{2\pi Y}Q^{M+1}}
 {(1-Q)(1-e^{2\pi Y}Q^{M+1})}<10^{-1260}.
\]
All calculations use Arb enclosures at $1152$ bits; the
radius comparisons and endpoint signs are exact interval
comparisons. The resulting radius is less than $10^{-236}$.
Adding the exact finite local logarithms of \ref{ari:C.6.3} and
subtracting this enclosure gives the claimed height bound.

\item For completeness, the singular-modulus certificate uses
$E_4=1+240\sum n^3q^n/(1-q^n)$ and
$\Delta_0=q\prod(1-q^n)^{24}$. Their respective error
bounds, after truncation at $M$, are
\[
 \frac{240(M+1)^3Q^{M+1}(1+4Q+Q^2)}{(1-Q)^5},
 \qquad
 \frac{24Q^{M+1}}{(1-Q)(1-Q^{M+1})}
\]
for $E_4$ and for the logarithm of the product. The first
uses $(M+1+k)^3\le(M+1)^3(k+1)^3$ and the differentiated
geometric series; the second uses the same logarithm bound.
Thus $E_4^3/\Delta_0$ is rigorously enclosed. The complete
certificate is \path{code/C_cm_height_arb.py}, with its
output and rational bracket centres retained.
\end{enumerate}
\end{proof}


\subsection{Supplement: ranks, saturation, and central derivatives}

We retain the order of the eleven rational CM values in the stated conventions, and
write $E_i$ for the global minimal models recorded in the preceding
tables. The real period in this subsection is
\[
 \Omega_i=\int_{E_i(\R)}|\omega_i|,
\]
including both connected components when $E_i(\R)$ is disconnected.
The canonical height $\widehat h$ is relative to $(O)$; PARI's
\texttt{ellheight} and the regulator convention used in the BSD formula
are relative to $2(O)$, and hence equal $2\widehat h$ in rank one.
This factor of two is essential in every comparison below.

\begin{proposition}[ARI \ref{ari:C.6.R1}; Proved by descent and full saturation]\label{ari:C.6.R1}
\StatementRecord{ari:C.6.R1}{proved}{Historical events: GPT-6 Astra (OpenAI), ChatGPT Pro mode, effort Pro, S6\_09\_PAPERS, 17 November 2026: targeted proof check. Also checked in this writing. Pre-continuation writing events: GPT-6 Astra (OpenAI), ChatGPT Pro mode, effort Pro, S6\_10\_PAPERS, 18 November 2026: computation or formal certificate; targeted proof check. Continuation metadata events: GPT-6 Astra (OpenAI), ChatGPT Pro mode, effort Pro, S6\_10\_PAPERS continuation, 19 November 2026: used and rechecked in this run; status-vocabulary reconciliation; metadata reconciliation only; no new mathematical check. \emph{Pre-continuation recorded proof scope.} Read complete printed arithmetic proofs and re-derived the exact scaling/local-height equations and the genus-trace square denominator argument. Rank, saturation, directed L-derivative replays and precise external theorem matches remain explicit gates; no proof verdict is awarded to those computational claims by this read. Reran the exact model/local-height programs in PARI2.17.2 and the unchanged eclib20250122 saturation code. All eleven displayed models and local coefficients match; ranks are exactly2 for rows1,8 and exactly1 for the remaining nine. All nine saturation indices are1, with no remaining primes. The full actual outputs and source hashes are retained. Opened PARI ellrank/ellheight normalization documentation and Cremona--Siksek Theorem1, author-PDFp8, for the algorithmic scope.}{the hypotheses, notation and definitions stated here; \citet[Theorem 1]{CrCremonaSiksek2006}}

The algebraic and analytic ranks agree for the eleven curves. They are
two for $j=8000$ and $j=-12288000$, and one for all nine other values.
For every rank-one curve the section $P$ is a primitive generator of
$E(\Q)/E(\Q)_{\rm tors}$. Its index is therefore one. For the two
rank-two curves the index of $\Z P$ in this free quotient is infinite.
\[
\begin{array}{r|r|r|r}
i&j&\operatorname{rank}E(\Q)&[E(\Q)/E(\Q)_{\rm tors}:\Z P]\\\hline
1&8000&2&\infty\\
2&54000&1&1\\
3&287496&1&1\\
4&16581375&1&1\\
5&-3375&1&1\\
6&-32768&1&1\\
7&-884736&1&1\\
8&-12288000&2&\infty\\
9&-884736000&1&1\\
10&-147197952000&1&1\\
11&-262537412640768000&1&1
\end{array}
\]
\end{proposition}


\begin{proof}[Arithmetic verification]
\begin{enumerate}[label=\arabic*.,leftmargin=*]
\item Exact $2$-descent in PARI/GP returns equal lower and upper rank bounds
for all eleven curves; its entire output, including independent points,
is in \path{code/C_cm_rank_lvalues.out}. Independent eclib descent
and saturation were completed for the first nine curves. The separate
full-saturation calculation \path{code/Crank_saturation.cpp} treats
all nine rank-one points, using eclib's automatic height-based index
bound and checking every prime up to that bound. All nine returned
index one and an empty list of unsaturated primes.

\item For clarity, the rounded-down bounds for the indices in the subgroup
of points having good reduction everywhere are respectively
\[
 1,1,1,2,2,2,3,5,5
\]
in the rank-one order $i=2,3,4,5,6,7,9,10,11$; the algorithm also
checks the primes dividing component-group orders. The underlying
height lower bound is the one in
\citet[Theorem 1]{CrCremonaSiksek2006}, applied to a minimal
Weierstrass equation and the subgroup
$E(\Q)\cap E^0(\R)\cap\bigcap_pE^0(\Q_p)$.
For the largest curve the lower bound on the nonzero height in this
subgroup is $0.651940632595967949\ldots$, its regulator is
$17.323913974368488616\ldots$, and the resulting index bound is
$5.154885040175890631\ldots$. The finite-field saturation maps for
$2,3,5$ already have full rank at residue primes $17,17,59$.
These data certify primitivity; merely finding $P$ in the output of
\texttt{ellrank} would not do so. The complete height-bound and
saturation transcripts are retained, including the independent
descent run that was stopped after its ninth completed curve.

\item The analytic ranks follow from the explicit derivative sums and bounds
proved below. The two rank-two bases furnished by descent are
$(-7,441),(105,833)$ for $i=1$, and
$(0,32004),(22218,3305456)$ for $i=8$.
\end{enumerate}
\end{proof}

\begin{proposition}[ARI \ref{ari:C.6.R2}; Proved: central derivatives with explicit tails]\label{ari:C.6.R2}
\StatementRecord{ari:C.6.R2}{proved}{Historical events: GPT-6 Astra (OpenAI), ChatGPT Pro mode, effort Pro, S6\_09\_PAPERS, 17 November 2026: targeted proof check. Also checked in this writing. Pre-continuation writing events: GPT-6 Astra (OpenAI), ChatGPT Pro mode, effort Pro, S6\_10\_PAPERS, 18 November 2026: computation or formal certificate; targeted proof check. Continuation metadata events: GPT-6 Astra (OpenAI), ChatGPT Pro mode, effort Pro, S6\_10\_PAPERS continuation, 19 November 2026: used and rechecked in this run; status-vocabulary reconciliation; metadata reconciliation only; no new mathematical check. \emph{Pre-continuation recorded proof scope.} Read complete printed arithmetic proofs and re-derived the exact scaling/local-height equations and the genus-trace square denominator argument. Rank, saturation, directed L-derivative replays and precise external theorem matches remain explicit gates; no proof verdict is awarded to those computational claims by this read. Read the complete retained C++ summation program and replayed both segmented runs unchanged:384bits/radius128/cutoff170 and512bits/radius96/cutoff190. The latter reproduces915818441 coefficients,69802132 nonzero values,sumabs1904681267996,4745134 blocks and the saved derivative exactly. Explicit tail is3.09321800754876143967e-78. Independently enclosed all eleven real periods and fresh derivative ratios with Arb; all displayed fifty-digit values pass. Opened Rubin Proposition5.3 p15 and Theorem10.1/full Section10 pp47–50; the exact CM isogeny and central-vanishing application are now written out.}{the hypotheses, notation and definitions stated here; \cref{ari:eq:Cr-tails}; \citet[Section 3]{CrDarmon2009}; \citet[Theorem 10.1]{CrRubin1999}}

The following are the values of $L'(E_i,1)/\Omega_i$; zeros in rows
one and eight are exact. More digits, the periods, derivatives and
tail bounds are in \path{results_Crank.csv}.
\begin{center}\small
\begin{tabular}{r|l}
$i$&$L'(E_i,1)/\Omega_i$\\\hline
1&$0$\\
2&$5.134398425737889664345619204085016979705061236171249722654226\ldots$\\
3&$2.134120849044754020051650801945664635106456422166881627967598\ldots$\\
4&$31.81013979419927776012209427941545866492237285892013668350005\ldots$\\
5&$1.869462708725885120030080805051387134542481374813232701178304\ldots$\\
6&$3.428873379050567976208155749688133688667614469824389504923957\ldots$\\
7&$0.362583348071598548315763082767526599722077924892872839203426\ldots$\\
8&$0$\\
9&$9.629273918234131678053742152263035417231796757756973886621556\ldots$\\
10&$16.04238857830583792513538634637663343755958963581574670348088\ldots$\\
11&$1247.321806154531180413038860827393095408344267114181914400926\ldots$
\end{tabular}
\end{center}
The two second derivatives, divided by the same real periods, are
\begin{align*}
 L''(E_1,1)/\Omega_1
 &=28.71408722758099118128474692322079946598060512912402172882525\ldots,\\
 L''(E_8,1)/\Omega_8
 &=78.03222041637032223732443782051140882655135939512581188900823\ldots.
\end{align*}
In particular the leading Taylor coefficients for these two curves
are half the displayed second derivatives.
\end{proposition}


\begin{proof}[Derivative formulas and error control]

\begin{enumerate}[label=\arabic*.,leftmargin=*]
\item Let $N$ be the conductor, $c=2\pi/\sqrt N$, and
$L(E,s)=\sum_{n\geq1}a_nn^{-s}$. Put
\[
 E_1(x)=\int_x^\infty\frac{e^{-u}}u\,du,
 \qquad G(x)=\int_x^\infty e^{-u}\frac{\log(u/x)}u\,du.
\]
Splitting the Mellin integral at $1/\sqrt N$ and using the functional
equation gives, for root number $-1$,
\begin{equation}\label{ari:eq:Cr-first}
 L'(E,1)=2\sum_{n\geq1}\frac{a_n}{n}E_1(cn).
\end{equation}
For root number $+1$ and $L(E,1)=0$, it gives
\begin{equation}\label{ari:eq:Cr-second}
 L''(E,1)=4\sum_{n\geq1}\frac{a_n}{n}G(cn).
\end{equation}
One can verify the factors directly by writing the completed Mellin
integral as an integral over $[1,\infty)$ of
$(u^{s-1}+wu^{1-s})f(iu/\sqrt N)$: its first derivative has
the factor $2\log u$ when $w=-1$, and its second derivative
the factor $2\log^2u$ when $w=1$. Integration by parts gives
the two formulas above. The vanishing of the lower derivatives
removes all derivatives of the gamma and conductor factors.

\item Hasse's bound and multiplicativity imply
$|a_n|\leq d(n)\sqrt n\leq2n$. Also
$E_1(x)\leq e^{-x}/x$ and $G(x)\leq e^{-x}/x^2$.
The latter follows on writing $u=x+v$ and using
$\log(1+v/x)\leq v/x$. Consequently the absolute tails after
$n=M$ are at most
\begin{equation}\label{ari:eq:Cr-tails}
 \frac{4e^{-c(M+1)}}{c(M+1)(1-e^{-c})},\qquad
 \frac{8e^{-c(M+1)}}{c^2(M+1)^2(1-e^{-c})},
\end{equation}
respectively. For the first ten curves we used
$M=\lceil180/c\rceil$, exact integer coefficients, and 1024-bit
MPFR arithmetic. For the rank-two kernel we used the convergent series
\[
 G(x)=\frac12(\log x+\gamma)^2+\frac{\pi^2}{12}
       +\sum_{k\geq1}\frac{(-x)^k}{k^2k!}.
\]
After $k>x$ its terms decrease in absolute value. The alternating
remainder was cut below $10^{-140}$. At most 740 terms were needed;
the largest intermediate term was less than $1.4\cdot10^{72}$.
The 1024-bit arithmetic therefore leaves far more than 100 digits
after this cancellation. The largest bound in
\eqref{ari:eq:Cr-tails} was $3.46\cdot10^{-77}$. The source and complete
outputs are \path{code/Crank_coefficients.gp},
\path{code/Crank_derivatives.cpp}, and their corresponding outputs.

\item For $i=11$, a full vector of $819$ million coefficients would be
unnecessarily large. We instead evaluated the CM theta series in
bounded-memory segments. The curve is the quadratic twist, by
$d=185801=7\cdot11\cdot19\cdot127$, of the conductor-$163^2$
curve
\[
 y^2+y=x^3-2174420x+1234136692.
\]
Write $D=163$, $\chi_D(a)=(a/D)$, and $\chi_d(n)=(d/n)$.
The exact coefficient $a_n$ is the sum of
\[
 \chi_d(n)\chi_D(a)a \quad(n=a^2,\ a>0)
\]
and
\[
 \chi_d(n)\chi_D((D+1)u/2)u
 \quad\left(n=\frac{u^2+Dv^2}{4},\ u,v>0,
                           \ u\equiv v\pmod2\right).
\]
This is the Hecke-character theta series for
$\mathcal O_{\Q(\sqrt{-D})}$, grouping conjugate algebraic integers
and then applying the quadratic twist; see
\citet[Section 3]{CrDarmon2009}. The first 10000 coefficients agree
exactly with PARI for both the untwisted and twisted curve. The same
code at $D=67,d=217$ reproduces the independent PARI derivative.

\item To accelerate the sum, in a block $n=n_0+r$, $|r|\leq R$, form the
exact integer moments $M_k=\sum a_nr^k$. For $x=cn_0$ write
\[
 e_k=e^{-x}(-x)^k/k!,\quad g_0=e^{-x},\quad
 g_k=e_k-g_{k-1},\quad
 \phi_0=E_1(x),\quad\phi_k=-\phi_{k-1}-g_{k-1}/k.
\]
Then the block contribution to $L'/2$ is
$\sum_k\phi_k M_k/n_0^{k+1}$. This follows by Taylor expansion
of $E_1(x(1+s))/(1+s)$. On $|s|=1/2$ its modulus is at most
$2E_1(x/2)$, which is less than $14$ throughout the block range.
We take enough terms that $(2R/n_0)^{k+1}<10^{-85}$.
The Cauchy remainder, summed with the actual
$\sum|a_n|<2\cdot10^{12}$ and $n_0>7700$, is less than
$10^{-75}$. The values of $E_1$ at consecutive block centres
are continued by the analogous Taylor expansion and refreshed
from a correctly rounded MPFR evaluation every 1024 blocks.
The continuation bound is $10^{-110}$ per normalized expansion.
All moments are integers, not floating-point sums.

\item The 384-bit run used $R=128$ and cutoff $170$; a separate 512-bit
run used $R=96$ and cutoff $190$. Their values agree beyond 75
decimal places. The latter used $915818441$ coefficients in
segments, with an omitted tail less than $3.10\cdot10^{-78}$,
and gave
\[
 L'(E_{11},1)=
 2.296582198163380608965915774436446055185027904705664673866754\ldots.
\]
The elementary operation-rounding bounds at these precisions are
smaller than $10^{-90}$: in the recurrences above
$|e_k|\leq1$, $|g_k|\leq1$, and $|\phi_k|<20$;
there are fewer than $5\cdot10^6$ blocks and the recurrences are
refreshed before 1024 consecutive continuations. Bounding each
normalized moment by $\sum|a_n|(R/n_0)^k$ gives this estimate
without assuming cancellation in the total sum. The arithmetic
and truncation margins are amply sufficient for every displayed
50-digit value. Both executable builds and both transcripts are
retained as \path{code/Crank_cm_stream*}.

\item The real periods are enclosed from the exact minimal invariants,
independently of the saved decimal period values.  For rows $1,2,3,4$
put respectively $T=\sqrt2,\sqrt3,2,\sqrt7$ and $\tau=iT$;
for rows $5,\ldots,11$ put
$T=\sqrt{7}/2,\sqrt{11}/2,\sqrt{19}/2,\sqrt{27}/2,
\sqrt{43}/2,\sqrt{67}/2,\sqrt{163}/2$ and $\tau=1/2+iT$.
Let $r=e^{-2\pi T}$ and $q=e^{2\pi i\tau}$.  The identity
$E_4=1+240\sum_{n\geq1}\sigma_3(n)q^n$ and
$\sigma_3(n)<2n^3$ give, after $M=100$ terms, the tail bound
\[
480r^{A}\left\{
\frac{A^3}{1-r}+\frac{3A^2r}{(1-r)^2}
+\frac{3Ar(1+r)}{(1-r)^3}
+\frac{r(1+4r+r^2)}{(1-r)^4}\right\},\qquad A=101.
\]
This follows by expanding $(A+j)^3$ and summing each geometric
moment.  All eleven parameters satisfy
$1-1008r/(1-32r)>0$; the bound
$\sigma_5(n)<2n^5$ and $n^5\leq32^{n-1}$ therefore proves
$E_6(\tau)>0$.  The exact Weierstrass invariant relations yield
$\kappa=2\pi(E_4/c_4)^{1/4}$.  For $c_6>0$ the real lattice
generator has length $\kappa$.  For $c_6<0$ the lattice is rotated
by $i$: its real generator has length $T\kappa$ in the rectangular
case and $2T\kappa$ in the other case.  Indeed the real vectors in
$i(\Z+iT\Z)$ are $T\Z$, while those in
$i(\Z+(1/2+iT)\Z)$ are $2T\Z$.  Multiplication by the number
$n_{\R}=2$ for $\Delta>0$ and $n_{\R}=1$ for $\Delta<0$
gives the full real period used in the statement.

The current program \path{code/I/certify_all_period_ratios.py}
performs these operations with 768-bit Arb intervals and verifies
every displayed period and derivative ratio.  For the first ten
derivatives it enlarges the explicit tail by one percent and adds
$10^{-100}$ for arithmetic and kernel-series errors.  This last
bound is conservative: fewer than $5\cdot10^5$ coefficients and
740 series terms occur, the intermediate terms are below
$1.4\cdot10^{72}$, and the 1024-bit accumulated rounding and
$10^{-140}$ kernel remainders are below this allowance.
For the last derivative it adds $1.1\cdot10^{-75}$, covering the
proved block Taylor and rounding bounds above.  Division by the
strictly positive period intervals preserves all fifty displayed
digits.  The resulting intervals and exact invariants are in
\path{code/I/records/all_period_ratios_arb.json}.

\item The exact root numbers are $+1$ for $i=1,8$ and $-1$ otherwise.
All points are nontorsion by Statement C.5. For a CM curve,
the Coates--Wiles theorem, in the form
\citet[Theorem 10.1]{CrRubin1999}, says that a nonzero central
Hecke $L$-value forces the Mordell--Weil group over the CM field
to be finite. Proposition~5.3 of the same source supplies, for a nonmaximal CM order, an isogenous
curve with maximal order: the quotient is by the Galois-stable subgroup attached to the conductor ideal. The isogeny induces an isomorphism on Mordell--Weil groups tensored with $\Q$ and on their rational Tate modules, preserving ranks and local $L$-factors.
Since $L(E/K,s)=L(E/\Q,s)^2$ for the CM field $K$, nontorsion
therefore forces $L(E,1)=0$. This proves the needed central
vanishing also in the two sign-$+1$ cases. The functional equation
then forces $L'(E,1)=0$ for those two curves. The positive first
or second derivatives just evaluated prove exactly the asserted
analytic ranks.
\end{enumerate}

\end{proof}

\begin{proposition}[ARI \ref{ari:C.6.R3}; Proved: the consequences of Gross--Zagier and Kolyvagin]\label{ari:C.6.R3}
\StatementRecord{ari:C.6.R3}{proved}{Historical events: GPT-6 Astra (OpenAI), ChatGPT Pro mode, effort Pro, S6\_09\_PAPERS, 17 November 2026: targeted proof check. Also checked in this writing. Pre-continuation writing events: GPT-6 Astra (OpenAI), ChatGPT Pro mode, effort Pro, S6\_10\_PAPERS, 18 November 2026: targeted proof check. Continuation metadata events: GPT-6 Astra (OpenAI), ChatGPT Pro mode, effort Pro, S6\_10\_PAPERS continuation, 19 November 2026: used and rechecked in this run; status-vocabulary reconciliation; metadata reconciliation only; no new mathematical check. \emph{Pre-continuation recorded proof scope.} Read complete printed arithmetic proofs and re-derived the exact scaling/local-height equations and the genus-trace square denominator argument. Rank, saturation, directed L-derivative replays and precise external theorem matches remain explicit gates; no proof verdict is awarded to those computational claims by this read. Matched Gross–Zagier I.7.3 pp11–12 and final proof p129 visually, Gross Theorem1.3 pp235–236 visually, and Jetchev–Lauter–Stein Theorem1.1 p2. Re-derived the rational differential/real-component period factor and added the explicit restriction–inflation argument carrying finite Sha from the Heegner field to Q. Exact derivative and algebraic-rank inputs were freshly replayed. The claimed nine positive rational factors and finite Sha conclusions are supported at this complete receiving-proof scope.}{the hypotheses, notation and definitions stated here; \citep{CrGrossZagier1986}; \citet[Theorem 1.3]{CrGross1991}; \citet[Theorem 1.1]{CrJetchev2008}}

For each of the nine rank-one curves there is a positive rational
number $q_i$ such that
\begin{equation}\label{ari:eq:Cr-rational-factor}
 \frac{L'(E_i,1)}{\Omega_i}=q_i\widehat h(P_i).
\end{equation}
Their Tate--Shafarevich groups over $\Q$ are finite. Consequently
the corresponding $D_{\rm part}$ is an explicit logarithm minus
$(6/(\pi q_i))L'(E_i,1)/\Omega_i$. For the two rank-two curves, $L'(E_i,1)=0$ while
$\widehat h(P_i)>0$.
\end{proposition}


\begin{proof}

\begin{enumerate}[label=\arabic*.,leftmargin=*]
\item Gross--Zagier's theorem applies because each curve is modular and
$L(E_i,1)=0$. CM modularity follows from its Hecke-character
theta series. Their Theorem I.7.3, in the opened original
preprint \citep{CrGrossZagier1986}, gives a rational point $R$ and
a nonzero rational factor relating $L'(E,1)$ to its canonical
height times the real period. The differentials in the two conventions are rational multiples:
write $\omega'=a\omega_E$ with $a\in\Q^*$.  A primitive real
period for $\omega'$ is $|a|\Omega_i/n_{\R}$, where
$n_{\R}\in\{1,2\}$ is the number of real components.
Thus the change to the full real period multiplies the rational
factor by $|a|/n_{\R}$.
As $L'(E_i,1)>0$, $R$ has infinite order. Primitivity gives
$R=mP_i+T$, $m\ne0$, $T$ torsion; its height is
$m^2\widehat h(P_i)$. This proves
\eqref{ari:eq:Cr-rational-factor}, including positivity.

\item The rank-one Gross--Zagier--Kolyvagin implication asserts rank one
and finite Tate--Shafarevich group when the central zero is simple.
One can use Kolyvagin's Heegner-point theorem as stated in
\citet[Theorem 1.3]{CrGross1991}, together with the nonzero
Gross--Zagier Heegner height; a modern primary research statement
of the rank conclusion is \citet[Theorem 1.1]{CrJetchev2008}.
The hypotheses are satisfied separately by the positive derivative
and exact sign in each of the nine rows.  Explicitly, the auxiliary
imaginary quadratic field $K$ in the proof of Gross--Zagier I.7.3
has split Heegner primes and $L(E_i^D,1)\ne0$; hence
$L'(E_i/K,1)=L'(E_i,1)L(E_i^D,1)\ne0$.  Its Heegner point has
positive height, so Gross's Theorem~1.3 gives finite
$\operatorname{Sha}(E_i/K)$.  The restriction map
\[
\operatorname{Sha}(E_i/\Q)\longrightarrow
\operatorname{Sha}(E_i/K)
\]
has kernel contained in
$H^1(\operatorname{Gal}(K/\Q),E_i(K))$ by restriction--inflation.
This last group is finite: the Mordell--Weil group is finitely
generated, and its degree-one cohomology for this order-two group
is finitely generated and killed by two.  The finite kernel and
finite image prove the claimed finiteness over $\Q$.
For the two other rows the established equations are
$L'(E_i,1)=0$ and $\widehat h_O(P_i)>0$; therefore an equation
$L'(E_i,1)/\Omega_i=q\widehat h_O(P_i)$ with $q>0$ would contradict
these exact values.  Their rank-two regulator, for any free basis
$P,Q$, is the determinant
$\widehat h_O(P)\widehat h_O(Q)-\langle P,Q\rangle^2$.
\end{enumerate}

\end{proof}

\paragraph{What the numerical BSD comparison does and does not determine.}
For a primitive rank-one point, the BSD prediction in these conventions is
\[
 \frac{L'(E,1)}{\Omega_E}
 =\frac{2\#\operatorname{Sha}(E/\Q)\prod_pc_p}
        {\#E(\Q)_{\rm tors}^{\,2}}\,\widehat h(P).
\]
The independently evaluated quotient after dividing by
$2\widehat h(P)\prod_pc_p/\#E(\Q)_{\rm tors}^2$ is one, to
more than 69 digits, for $i=2,3,4,5,6,7,9,10$, and nine for
$i=11$. Thus the numerically predicted rational factors $q_i$ are
\[
 8,4,32,4,8,4,16,16,576,
\]
in that order. This initial numerical comparison alone proves neither
these exact factors nor the orders $1$ or $9$ of $\operatorname{Sha}$.
The exact factors are subsequently proved in \ref{ari:C.6.R4} and \ref{ari:C.6.R5}:
the latter supplies an independent Gross--Zagier denominator bound
and rigorous intervals for all remaining ratios. This determines
the analytic BSD quotients exactly, while the orders of
$\operatorname{Sha}$ are proved here only in the three cases of
\ref{ari:C.6.R4}. In particular primitivity alone does not determine
$\#\operatorname{Sha}$ for the last curve.

\begin{cited}[ARI \ref{ari:C.6.R4}; Known input: three exact leading factors]\label{ari:C.6.R4}
\StatementRecord{ari:C.6.R4}{proved}{Historical events: GPT-6 Astra (OpenAI), ChatGPT Pro mode, effort Pro, S6\_09\_PAPERS, 17 November 2026: targeted proof check. Also checked in this writing. Pre-continuation writing events: GPT-6 Astra (OpenAI), ChatGPT Pro mode, effort Pro, S6\_10\_PAPERS, 18 November 2026: targeted proof check; used and rechecked in this run. Continuation metadata events: GPT-6 Astra (OpenAI), ChatGPT Pro mode, effort Pro, S6\_10\_PAPERS continuation, 19 November 2026: used and rechecked in this run; status-vocabulary reconciliation; metadata reconciliation only; no new mathematical check. \emph{Pre-continuation recorded proof scope.} Read complete printed arithmetic proofs and re-derived the exact scaling/local-height equations and the genus-trace square denominator argument. Rank, saturation, directed L-derivative replays and precise external theorem matches remain explicit gates; no proof verdict is awarded to those computational claims by this read. Read Miller Definition1.1/Theorem1.2 p2, Section7 opening p16 and complete exception Table9 p23. Conductors3136,3136,361 lie outside the exceptions; all have proved analytic rank1. Exact local/torsion data give coefficient4. Current Arb height and derivative/period intervals isolate the resulting positive integer Sha order as1 in all three rows. Expanded known application remains in the unchanged proof compendium.}{the exact models and primitivity in \cref{ari:C.6.R1};
the certified derivatives in \cref{ari:C.6.R2};
\citet[Definition~1.1, Theorem~1.2, and Table~9]{miller2011};
the full application retained in the proof compendium, Paper~I,
label \nolinkurl{app:rederive:ari-C.6.R4}}

For $i=3,5,7$, namely $j=287496,-3375,-884736$, one has
\[
 \operatorname{Sha}(E_i/\Q)=0,\qquad
 \frac{L'(E_i,1)}{\Omega_i}=4\widehat h_O(P_i).
\]
Thus the factors $q_3=q_5=q_7=4$ are determined exactly by the integer bounds in the cited application. For the smallest of these examples,
$E:y^2+y=x^3-38x+90$ and $P=(0,9)$, this gives the exact
reduction
\[
 N(P)=\exp\left(-\frac{L'(E,1)}{4\Omega_E}\right).
\]
\end{cited}

The expanded source derivation is retained in the proof compendium, Paper~I, label \nolinkurl{ari:C.6.R4}.




The three exact cases specialize \citet[Definition~1.1, Theorem~1.2 and Table~9]{miller2011}. The complete application is retained in the proof compendium, Paper~I, label \nolinkurl{app:rederive:ari-C.6.R4}.


\paragraph{The scope of Rubin's theorem.}
The applicable ordinary CM results concern a prime $p$ at which
$E$ has good ordinary reduction, equivalently $p$ splits in the CM
field. One precise opened formulation of Rubin's rank-one formula
is \citet[Theorem 1]{CrBDP2012}: for a curve over $\Q$ with CM
by the maximal order, an ordinary prime prime to the conductor,
and a simple central zero, an appropriate Katz $p$-adic $L$-value
is, up to a nonzero algebraic factor, the square of a formal
logarithm of a nontorsion rational point divided by a $p$-adic
period. Its input is analytic rank one, already verified here.
It is a $p$-adic identity; it does not identify the real exponential
$\exp\widehat h(P)$ as algebraic.

For the $p$-part of the real leading-term formula, a modern precise
version is \citet[Theorem A]{CrCastella2025}. In addition to analytic
rank one it requires full CM, abelian torsion field over the CM
field, $p\nmid6h_K$, splitting of $p$, conductor prime to $p$,
and that the different of $K$ divide the Hecke conductor exactly.
For the maximal odd-discriminant cases here, the latter follows
from the exponent two of the ramified prime in the rational
conductor. The nonmaximal cases of discriminant $-12$ and $-28$
can be passed through a $2$-isogeny at $p>3$. The discriminant
$-16$ case requires care: the direct application of this version
has a different-conductor hypothesis at $2$ that fails, and we
do not claim that this particular formulation proves its full
leading-term formula. Rubin's original ordinary main-conjecture
paper is \citep{CrRubin1991}; its bibliographic record was checked,
but its complete theorem text was not available in the opened
publisher copy, so no additional hypothesis-free claim is based
on it here.

In particular these ordinary-prime statements do not prove that
the $3$-primary part of $\operatorname{Sha}$ has order $9$ for the
last curve: $3$ is inert in $\Q(\sqrt{-163})$ and also excluded
from the stated $p>3$ hypotheses. \ref{ari:C.6.R5} proves the corresponding
analytic quotient to be exactly $9$ by a separate denominator
argument. The identification with the actual Tate--Shafarevich
group remains a separate arithmetic question in that row.


\subsection{Supplement: exact leading factors from a genus-character denominator bound}

We retain the full real period $\Omega_i$ and the absolute canonical
height $\widehat h_O$ of the preceding subsection. The argument below
uses the Gross--Zagier formula to bound a rational denominator; it does
not assume the leading-term Birch--Swinnerton-Dyer formula.

\begin{proposition}[ARI \ref{ari:C.6.R5}; Proved: all remaining rational leading factors]\label{ari:C.6.R5}
\StatementRecord{ari:C.6.R5}{proved}{Historical events: GPT-6 Astra (OpenAI), ChatGPT Pro mode, effort Pro, S6\_09\_PAPERS, 17 November 2026: targeted proof check. Also checked in this writing. Pre-continuation writing events: GPT-6 Astra (OpenAI), ChatGPT Pro mode, effort Pro, S6\_10\_PAPERS, 18 November 2026: computation or formal certificate; targeted proof check; used and rechecked in this run. Continuation metadata events: GPT-6 Astra (OpenAI), ChatGPT Pro mode, effort Pro, S6\_10\_PAPERS continuation, 19 November 2026: used and rechecked in this run; status-vocabulary reconciliation; metadata reconciliation only; no new mathematical check. \emph{Pre-continuation recorded proof scope.} Read complete printed arithmetic proofs and re-derived the exact scaling/local-height equations and the genus-trace square denominator argument. Rank, saturation, directed L-derivative replays and precise external theorem matches remain explicit gates; no proof verdict is awarded to those computational claims by this read. Reran B\_C6\_gz\_exact\_metadata.gp with only the unavailable database-label display omitted. All six exact twist/differential/conductor/coprimality/Heegner-splitting identities pass. All four CM theta coefficient vectors and auxiliary quadratic-twist vectors agree with independently computed elliptic Fourier coefficients through n=10000, beyond each printed Sturm bound. Directed numerical enclosures and the remaining primary theorem normalizations are still being checked. All six current directed MPFR derivative/auxiliary central-value enclosures and Arb exact-factor recognitions passed. Ratios r=64,64,8,16,400,64 and q=8,32,8,16,16,576; genus-trace magnitudes32,64,8,16,80,192. Fresh GP Fourier coefficients and exact eleven-model comparison passed. An initial auxiliary script used PARI reserved variable I and failed; changed to idx and reran all five affected file certificates successfully. No failed output was accepted. Opened original Cai-Shu-Tian2014 Theorem1.1 pp2524-2525 and Belabas-Bernardi-Perrin-Riou2019 AppendixB p102 unnumbered26569a1/a2 block. Verified its Manin constant1 and modular degree59368 visually. Height-normalization and remaining modular-input checks remain separately tracked. Recomputed all first-five small-base Manin constants by installed PARI2.17.2 ellmaninconstant(E); all equal1. Local runtime help explicitly states unconditional algorithm. Initial two-argument call from newer web docs was rejected; preserved failed logs, removed unsupported flag and reran five computations successfully. Opened author-uploaded YZZ2011 Section7.1.1 printedpp218-219 including Theorems7.1.1/7.1.2 and Proposition7.1.3; confirmed base-field-relative height and pairing without half, yielding exact factor4 relative to absolute O-height. Opened Darmon2009 Section3pp4-5 defining canonical CM Hecke theta series. Opened Stein Theorem9.18 and its entire proof pp171–173, Corollaries9.19–9.20 pp173–174, matching the coefficient bounds and coset-product argument. Together with the separately recorded YZZ, Cai–Shu–Tian, Darmon, Manin-constant and six directed numerical checks, this closes the named primary-source gates for the six denominator recognitions.}{the hypotheses, notation and definitions stated here; \cref{ari:C.6.R4}; \citet[Appendix B, p.~102, unnumbered $26569\mathrm a1/26569\mathrm a2$
block]{CgzBelabas2019}; \citet[Theorem 1.1, pp.~2524--2525]{CgzCai2014}; \cref{ari:C.6.R2}, \cref{ari:C.6.R3}; \cref{ari:C.6.R1}; \citet[\S7.1.1, pp.~218--219, Theorem 7.1.2]{CgzYZZ2011}; \cref{ari:eq:Cgz-product}; \citet[\S3, pp.~4--5]{CrDarmon2009}; \citet[Theorem 9.18]{CgzStein2007}; \cref{ari:eq:Cr-first}; \cref{ari:C.6.5}; \cref{ari:eq:Cgz-central-integrality}; \cref{ari:eq:Cgz-denominator}}

For the six rank-one rows $i=2,4,6,9,10,11$, one has
\[
 \frac{L'(E_i,1)}{\Omega_i}
   =q_i\widehat h_O(P_i),\qquad
 (q_2,q_4,q_6,q_9,q_{10},q_{11})=(8,32,8,16,16,576).
\]
Together with \ref{ari:C.6.R4}, this determines every rank-one factor:
\[
\begin{array}{c|rrrrrrrrr}
 i&2&3&4&5&6&7&9&10&11\\\hline
 q_i&8&4&32&4&8&4&16&16&576.
\end{array}
\]
The certified auxiliary ratios and denominator bounds used to identify
the six additional factors are as follows. Here $D$ is the auxiliary
negative fundamental discriminant,
$r_i=L(E_i^D,1)/\Omega_{i,D}$, and the denominator of $q_i$ is at most
$B_i$.
\[
\begin{array}{c|r|r|r|r|r}
 i&D&r_i&B_i&q_i&|a_i|\\\hline
 2&-23&64&256&8&32\\
 4&-31&64&256&32&64\\
 6&-19&8&32&8&8\\
 9&-19&16&64&16&16\\
 10&-3&400&1600&16&80\\
 11&-3&64&256&576&192
\end{array}
\]
The last column is the absolute genus-trace integer in the proof. These
identities determine the analytic BSD quotients exactly: they are one
for $i=2,4,6,9,10$ and nine for $i=11$. They do not by themselves
identify those numbers with the orders of the Tate--Shafarevich groups.
\end{proposition}



\begin{proof}
\begin{enumerate}[label=\arabic*.,leftmargin=*]
\item We give the arithmetic denominator bound first and then the interval
certificate that isolates the listed rational numbers.

\item \paragraph{1. A parametrization under quadratic twisting.}
Let $E_0/\Q$ have conductor $N_0$, and suppose a parametrization
$\phi:X_0(N_0)\to E_0$ sends $\infty$ to $O$ and satisfies
\[
 \phi^*\omega_0=c_0f_0(q)\,dq/q,
\]
where $\omega_0$ is a minimal differential and $f_0$ is the normalized
newform. Let $b$ be an odd fundamental discriminant coprime to $N_0$,
and suppose that the minimal twist has conductor $N_0b^2$ and exact
invariants
\[
 c_4(E_0^b)=b^2c_4(E_0),\qquad
 c_6(E_0^b)=b^3c_6(E_0).
\]
Then $E_0^b$ has a parametrization from $X_0(N_0b^2)$ with the same
Manin factor $c_0$.

\item Indeed, put $m=|b|$ and, using addition on $E_0$, define
\[
 F(z)=\sum_{a\bmod m}\chi_b(a)\phi(z+a/m).
\]
For $\gamma=\left(\begin{smallmatrix}A&B\\C&D_0\end{smallmatrix}\right)$
in $\Gamma_0(N_0m^2)$ choose $a'\equiv aD_0^2\pmod m$.
Direct multiplication gives
$T_{a/m}\gamma T_{a'/m}^{-1}\in\Gamma_0(N_0)$, where $T_t$
denotes translation by $t$. Because $\chi_b(D_0^2)=1$, the sum is
invariant. Each summand is an algebraic morphism on a common
congruence cover; invariance therefore descends the sum to
$X_0(N_0m^2)$ over $\Q(\zeta_m)$. Its expansion at $\infty$
gives $F^{\sigma_u}=\chi_b(u)F$ for
$\sigma_u(\zeta_m)=\zeta_m^u$. Quadratic descent gives the required
map to $E_0^b$ over $\Q$, with $\infty$ mapping to $O$.
On differentials, the Gauss-sum identity gives
\[
 F^*\omega_0=c_0\tau(\chi_b)f_b(q)\,dq/q,
 \qquad \tau(\chi_b)^2=b.
\]
For the twist isomorphism $\iota:E_0^b\to E_0$ the exact invariant
identities imply
$\iota^*\omega_0=\sqrt b\,\omega_b$ up to sign, with no additional
minimal-model scaling. Cancellation proves the assertion. The map
need not be optimal; optimality is not used below.

\item If $c_0=1$, let $t_b=\#E_0^b(\Q)_{\rm tors}$ and
$n_b=\#\pi_0(E_0^b(\R))$. Manin--Drinfeld says that the image of
the rational cuspidal divisor $[0]-[\infty]$ is torsion. Its elliptic
logarithm, along the imaginary axis, is $\pm L(E_0^b,1)$.
Consequently $t_bL(E_0^b,1)$ is an integral multiple of the positive
real lattice period $\Omega_b^+$, and
\begin{equation}\label{ari:eq:Cgz-central-integrality}
 t_bn_b\,\frac{L(E_0^b,1)}{\Omega_b}\in\Z,
 \qquad \Omega_b=n_b\Omega_b^+.
\end{equation}

\item \paragraph{2. The six small base curves and Heegner hypotheses.}
Use the following data; $E_i=E_0^d$ and $K=\Q(\sqrt D)$.
\[
\begin{array}{c|c|r|r|r}
 i&E_0&N_0&d&D\\\hline
 2&576\mathrm e2&576&-11&-23\\
 4&3136\mathrm r2&3136&57&-31\\
 6&121\mathrm b1&121&21&-19\\
 9&1849\mathrm a1&1849&21&-19\\
 10&4489\mathrm a1&4489&217&-3\\
 11&26569\mathrm a1&26569&185801&-3
\end{array}
\]
All six base Manin factors are one. For the first five bases this
is independently returned by PARI's unconditional modular-symbol
algorithm. For $26569\mathrm a1$, the value one and modular degree
$59368$ are explicitly tabulated by
\citet[Appendix B, p.~102, unnumbered $26569\mathrm a1/26569\mathrm a2$
block]{CgzBelabas2019}; their Algorithm 1 describes the computation.
The first five completed computations are in
\path{code/invariants_arithmetic/arithmetic/outputs/writing_B_C6_gz_exact_metadata.gp.txt}. The separate exact audit
\path{code/B_C6_gz_exact_metadata.gp} checks the minimal equations,
both $c_4,c_6$ scaling identities for $b=d,dD$, both conductor
identities, all torsion orders, and all splitting conditions. It
does not recognize any of these identities from floating-point
scaling factors.

\item In every row $d,D,dD$ are fundamental discriminants,
$\gcd(d,D)=\gcd(dD,N_0)=1$, and every prime dividing $N_0$ splits
in $K$. The genus character for $K(\sqrt d)/K$ factors through
$\operatorname{Pic}(\mathcal O_{|d|})$ and has conductor exactly
$|d|$. On ideals prime to $d$ it is
\[
 \chi(\mathfrak a)=\chi_d(N\mathfrak a).
\]
For a unit of the order locally at a prime $p\mid d$, its norm is
a square modulo $p$, so this character is trivial on the order's
units. It is nontrivial on the full local unit group: at a split
prime the norm is multiplication on two copies of $\Q_p^\times$,
and at an inert prime the residue-field norm onto $\mathbb F_p^\times$
is surjective. Since $p$ is odd and unramified in $K$, this also
proves that the conductor contains every $p\mid d$ exactly once.
Thus the primitive ring-class-character hypotheses of
\citet[Theorem 1.1, pp.~2524--2525]{CgzCai2014} hold. Their
$\mu(N_0,D)=0$ and
$u=[\mathcal O_{|d|}^{\times}:\Z^{\times}]=1$; the latter remains
true for $D=-3$ because $|d|>1$.

\item Let $P_{\mathfrak a}$ be the corresponding Heegner points and put
\[
 Y=\sum_{[\mathfrak a]}\chi(\mathfrak a)\phi(P_{\mathfrak a}).
\]
The coefficients are $\pm1$, so this is an actual algebraic point,
not merely an element of a tensor product with $\Q$. It is fixed
by $\ker\chi$ and anti-invariant under the other automorphism of
$K(\sqrt d)/K$. Twisting therefore gives $Y_d\in E_i(K)$.
Artin factorization gives
\[
 L(s,E_0,\chi)=L(E_i,s)L(E_i^D,s).
\]
The intervals below prove $L(E_i^D,1)>0$, hence the auxiliary
curve has rank zero. The rank-zero implication used here is the
same unconditional modular/CM rank implication used in \cref{ari:C.6.R2}, \cref{ari:C.6.R3}.
Thus $E_i(K)$ has rank one. \ref{ari:C.6.R1} proves that $P_i$ is primitive
over $\Q$. Since $Y_d-\overline{Y_d}$ is torsion and
$Y_d+\overline{Y_d}$ is rational, there is an integer $a_i$ with
\[
 2Y_d=a_iP_i\pmod{E_i(K)_{\rm tors}},\qquad
 \widehat h_O(Y_d)=\frac{a_i^2}{4}\widehat h_O(P_i).
\]

\item \paragraph{3. The exact height and period normalization.}
Write $V_0=\operatorname{area}(\C/\Lambda_0)$ for the period area
of the minimal differential. Integrating $\phi^*\omega_0$ gives
\[
 4\pi^2(f_0,f_0)_{\Gamma_0(N_0)}
   =\frac{\deg\phi}{c_0^2}V_0.
\]
The degree in the denominator of CST's Theorem 1.1 cancels, giving
\begin{equation}\label{ari:eq:Cgz-product}
 L'(E_i,1)L(E_i^D,1)
  =\frac{2V_0}{c_0^2|d|\sqrt{|D|}}\,\widehat h_K(Y).
\end{equation}
Here $\widehat h_K$ is the diagonal N\'eron--Tate pairing relative
to $K$. This convention is explicit in
\citet[\S7.1.1, pp.~218--219, Theorem 7.1.2]{CgzYZZ2011}:
their underlying height is normalized by $[L:K]^{-1}$ over a
field of definition $L$, and their pairing is
$\widehat h_{\mathcal L}(x+y)-\widehat h_{\mathcal L}(x)
-\widehat h_{\mathcal L}(y)$.

\item One can also see the factor geometrically. For
$\mathcal L=\mathcal O(O)$, pull the normalized Poincar\'e bundle
back under $1\times\lambda_{\mathcal L}$. The result is
$m^*\mathcal L\otimes p_1^*\mathcal L^{-1}\otimes
p_2^*\mathcal L^{-1}$. On the diagonal its canonical height is
that of $[2]^*\mathcal L\otimes\mathcal L^{-2}\simeq
\mathcal L^2$, namely $2\widehat h_O$ in the absolute convention.
Changing from the absolute convention to relative heights over
$K$ multiplies by $[K:\Q]=2$. Therefore
\[
 \widehat h_K(Y)=4\widehat h_O(Y_d)
                =a_i^2\widehat h_O(P_i).
\]
This factor comes from the Poincar\'e pairing and the base field,
independently of any numerical height routine.

\item A real-adapted basis of $\Lambda_0$ gives
\[
 \Omega_i\Omega_{i,D}=
 \begin{cases}
 2V_0/(|d|\sqrt{|D|}),&\Delta(E_0)<0,\\
 4V_0/(|d|\sqrt{|D|}),&\Delta(E_0)>0.
 \end{cases}
\]
Indeed, a triangular real lattice has one real component and its
purely imaginary primitive period has twice the height of its
nonreal basis vector. A rectangular lattice has two real
components. Twisting by the two opposite signs interchanges the
real and imaginary directions, and the exact minimal differential
scalings in step 2 supply $1/(|d|\sqrt{|D|})$.
Substituting into \eqref{ari:eq:Cgz-product} gives
\begin{equation}\label{ari:eq:Cgz-denominator}
 q_ir_i=
 \begin{cases}
 a_i^2/c_0^2,&\Delta(E_0)<0,\\
 a_i^2/(2c_0^2),&\Delta(E_0)>0.
 \end{cases}
\end{equation}
In particular, if $r_i=A_i/B_i'$ is reduced and $c_0=1$, the
denominator of $q_i$ divides $2|A_i|$. We use the weaker upper
bound $B_i=4|A_i|$ in the interval certificate.

\item \paragraph{4. Global identification of the CM coefficient series.}
For $s=11,43,67,163$, put
$F_s=\Q(\sqrt{-s})$ and $\mathfrak p=(\sqrt{-s})$.
These fields have class number one and units $\{\pm1\}$.
The character
\[
 \psi((\alpha))=
  \left(\frac{\alpha\bmod\mathfrak p}{s}\right)\alpha
\]
is well defined because $(-1/s)=-1$. It has infinity type $(1,0)$
and finite conductor $\mathfrak p$. Hecke's theta construction
places $\sum_{\mathfrak a}\psi(\mathfrak a)q^{N\mathfrak a}$ in
$S_2(\Gamma_0(s^2))$: the nebentypus is
$\chi_{-s}(n)(n/s)=1$. This exact special case is also written
in \citet[\S3, pp.~4--5]{CrDarmon2009}. Grouping conjugate
generators $\alpha=(u+v\sqrt{-s})/2$ gives the integral formula
used by the coefficient stream:
\[
 b_n=\sum_{a^2=n,\ a>0}\left(\frac as\right)a
  +\sum_{\substack{u,v>0,\ u\equiv v\ (2)\\u^2+sv^2=4n}}
       \left(\frac{(s+1)u/2}{s}\right)u.
\]
For all four bases, exact PARI elliptic-curve coefficients equal
$b_n$ for $1\leq n\leq10000$.

\item This finite comparison proves equality of the entire modular forms.
For completeness, if a nonzero form $g$ of weight $k$ on
$\Gamma_0(N)$ has trivial character, take the product of
$g|_k\gamma$ over right coset representatives in
$\Gamma_0(N)\backslash\mathrm{SL}_2(\Z)$. It is a nonzero
holomorphic full-level form of weight $kI$, where
$I=N\prod_{p\mid N}(1+1/p)$. Every factor has nonnegative
order at infinity, so the product's order is at least
$\operatorname{ord}_\infty g$. The full-level valence formula
therefore implies $\operatorname{ord}_\infty g\leq kI/12$.
Equivalently, agreement through $\lfloor kI/12\rfloor$ proves
equality. The stronger arithmetic congruence statement is
\citet[Theorem 9.18]{CgzStein2007}; only the elementary
characteristic-zero argument is needed here. For $k=2,N=s^2$
the four bounds are
\[
 \left\lfloor\frac{s^2+s}{6}\right\rfloor
     =22,315,759,4455,
\]
all below $10000$. Applying the exact quadratic twist character
then identifies every coefficient in both the target and auxiliary
streams, including coefficients past the tested range. The exact
metadata transcript prints the four bounds and their comparisons.

\item \paragraph{5. Directed enclosures of both $L$-sums.}
For $c=2\pi/\sqrt N$, the root-number $+1$ formula is
\[
 L(E,1)=2\sum_{n\geq1}a_n\frac{e^{-cn}}n,
\]
while for root number $-1$ the first derivative is
\eqref{ari:eq:Cr-first}. The exact coefficients in the smaller cases
come from retained PARI files. The largest auxiliary stream uses
$s=163$, signed twist $-557403$, conductor
$8254937936042721$, and $M=867617471$; the target stream uses
$s=163,d=185801,M=289205824$. Both use the cutoff $cM\geq60$
up to rounding of the integer cutoff, with the actual $M$ used in
the tail bound. No precision conclusion depends on the nominal
cutoff.

\item Here are the finite-sum and rounding certificates implemented in
\path{code/B_C6_central_stream.cpp}. All intervals use 256-bit
MPFR with outward rounding at every operation. In a block centred
at $b$ with radius $R=128$, set
$S_k=\sum a_n(n-b)^k\in\Z$. For $g(z)=e^{-cz}/z$, put
\[
 C_0=1,\quad C_k=C_{k-1}c/k,\qquad
 T_0=1/b,\quad T_k=(T_{k-1}+C_k)/b.
\]
The coefficient of $(z-b)^k$ is $(-1)^ke^{-cb}T_k$.
For $h(z)=E_1(cz)/z$, put
\[
 e_0=E_1(cb),\quad e_k=(-1)^ke^{-cb}T_{k-1}/k\ (k\geq1),
 \qquad h_0=e_0/b,\quad h_k=(e_k-h_{k-1})/b.
\]
The block sums are the finite sums of these interval coefficients
times the exact integers $S_k$.

\item On $|z-b|=b/2$, the two kernels have bounds
\[
 |g(z)|\leq\frac{2e^{-cb/2}}b=:U_g,
 \qquad
 |h(z)|\leq\frac{4e^{-cb/2}}{cb^2}=:U_h.
\]
For the second bound, use
$E_1(cz)=\int_1^\infty e^{-czt}\,dt/t$ when $\Re z>0$.
Thus truncation at degree $k$ contributes at most
\[
 U_g\frac{(2R/b)^{k+1}}{1-2R/b}\sum|a_n|
 \quad\hbox{or}\quad
 U_h\frac{(2R/b)^{k+1}}{1-2R/b}\sum|a_n|.
\]
The degree is selected to make the geometric factor small, but
the actual remainder is separately evaluated with upward rounding.
For the derivative stream, the initial $E_1$ in each segment is
evaluated using correctly rounded $-\operatorname{Ei}(-x)$ and
monotonicity. Subsequent centres, spaced by $W=257$, use
\[
 E_1(c(b+W))=E_1(cb)-\int_0^Wg(b+x)\,dx.
\]
Integrating the same Taylor expansion gives the interval for this
advance. Its omitted part is bounded by
$WU_g(2W/b)^{k+1}/(1-2W/b)$ and added to both endpoints.
The first $10280$ terms are summed directly. Segment boundaries
are multiples of $257$, so no coefficient is omitted or repeated.

\item Hasse's bound and multiplicativity give $|a_n|\leq2n$.
The infinite tails, including the factor two, are consequently at
most
\[
 \frac{4e^{-c(M+1)}}{1-e^{-c}},\qquad
 \frac{4e^{-c(M+1)}}{c(M+1)(1-e^{-c})},
\]
respectively. The program evaluates both with the downward bound
for $c$ and upward bounds for the remaining positive expressions.
It finally adds the block remainders, interval arithmetic errors,
and infinite tail. For row eleven the central-value tail is below
$5.07\cdot10^{-19}$, and the derivative tail is below
$2.82\cdot10^{-21}$. These modest precisions already suffice for
rational separation. All twelve completed interval outputs are
retained, as are the exact coefficient files.

\item \paragraph{6. Periods, heights, and exact rational separation.}
The periods are enclosed independently with 512-bit Arb. Let
$\tau=i\sqrt s$ for $i=2,4$ and
$\tau=(1+i\sqrt s)/2$ for $i=6,9,10,11$; here
$s=3,7,11,43,67,163$, respectively. The exact CM identification
and reduced-form class-number check are part of \ref{ari:C.6.5}. Put
$T=\Im\tau$, $q_\tau=e^{2\pi i\tau}$, and
$r=|q_\tau|<10^{-4}$. We enclose
\[
 E_4(\tau)=1+240\sum_{n=1}^{100}\sigma_3(n)q_\tau^n
                   +\text{tail}.
\]
Since $\sigma_3(n)<2n^3$, with $A=101$ the absolute tail is at most
\[
 480r^A\left\{
 \frac{A^3}{1-r}+\frac{3A^2r}{(1-r)^2}
 +\frac{3Ar(1+r)}{(1-r)^3}
 +\frac{r(1+4r+r^2)}{(1-r)^4}\right\}.
\]
The exact invariants $c_4,c_6$ of each minimal equation give
$\kappa=2\pi(E_4(\tau)/c_4)^{1/4}$. The sign of $E_6(\tau)$
is positive: $\sigma_5(n)<2n^5$ and
$n^5\leq32^{n-1}$ give
$|E_6(\tau)-1|\leq1008r/(1-32r)<1$.
Thus the full real period is
\[
 \Omega=n_{\R}\kappa\begin{cases}
 1,&c_6>0,\\
 T,&c_6<0,\ \Delta>0,\\
 2T,&c_6<0,\ \Delta<0,
 \end{cases}
 \qquad n_{\R}=\begin{cases}2,&\Delta>0,\\1,&\Delta<0.\end{cases}
\]
These formulas apply separately to the target and auxiliary curve.
Every operation, the Fourier tail, and the fourth root is enclosed
by Arb. The height intervals are the independent certificates
\path{records/C_cm_height_arb.json} of \ref{ari:C.6.5}.

\item The final computation is
\path{code/B_C6_exact_factors.py}. First, the intervals for
$r_i$ isolate a unique integer after multiplication by the exact
$t_bn_b$ in \eqref{ari:eq:Cgz-central-integrality}. This proves the
six auxiliary rational values, with no BSD input. Next form the
interval $L'(E_i,1)/(\Omega_i\widehat h_O(P_i))$. The enclosures
below are deliberately enlarged decimal summaries of the recorded
Arb intervals:
\[
\begin{array}{c|c|c}
 i&|r_i-r_i^{\rm listed}|\text{ is at most}
   &|q_i-q_i^{\rm listed}|\text{ is at most}\\\hline
 2&4\cdot10^{-22}&2\cdot10^{-53}\\
 4&4\cdot10^{-20}&3\cdot10^{-45}\\
 6&4\cdot10^{-22}&2\cdot10^{-53}\\
 9&5\cdot10^{-21}&4\cdot10^{-53}\\
 10&3\cdot10^{-20}&10^{-23}\\
 11&7\cdot10^{-15}&8\cdot10^{-19}
\end{array}
\]
Each last-column bound is strictly less than $1/B_i$. A rational
number with denominator at most $B_i$ differing from the listed
integer would differ by at least $1/B_i$. Equation
\eqref{ari:eq:Cgz-denominator} therefore makes every listed $q_i$
exact. The same equation gives the listed $|a_i|$ by integer
square-root checks. The JSON certificate records all input
intervals, exact invariants, rational values, and positive separation
margins. Its stdout concludes that all six remaining factors are
exactly certified. This completes the proof.
\end{enumerate}
\end{proof}


\subsection{Supplement: the three specified smooth torsion fibres}

\begin{proposition}[T1; Proved, with exact local and descent computations]\label{ari:C.6.T1}
\StatementRecord{ari:C.6.T1}{proved}{Historical events: GPT-6 Astra (OpenAI), ChatGPT Pro mode, effort Pro, S6\_09\_PAPERS, 17 November 2026: targeted proof check. Also checked in this writing. Pre-continuation writing events: GPT-6 Astra (OpenAI), ChatGPT Pro mode, effort Pro, S6\_10\_PAPERS, 18 November 2026: computation or formal certificate; targeted proof check. Continuation metadata events: GPT-6 Astra (OpenAI), ChatGPT Pro mode, effort Pro, S6\_10\_PAPERS continuation, 19 November 2026: used and rechecked in this run; status-vocabulary reconciliation; metadata reconciliation only; no new mathematical check. \emph{Pre-continuation recorded proof scope.} Read complete printed arithmetic proofs and re-derived the exact scaling/local-height equations and the genus-trace square denominator argument. Rank, saturation, directed L-derivative replays and precise external theorem matches remain explicit gates; no proof verdict is awarded to those computational claims by this read. Read and replayed the unmodified local and final torsion programs in PARI2.17.2. Exact rank0,torsion4,conductor40,discriminant80,component orders2/1 and local-height coefficients1/12 match. Exact plus modular symbol is-1/4; the printed period/L-value comparisons reproduce. All starred final-log entries are stack-growth warnings; complete transcript contains no GP error.}{the hypotheses, notation and definitions stated here; the explicit constructions and calculations in the proof}

At $t=32/5$ a global minimal model and the image of the section are
\[
 E:y^2=x^3-2x+1,\qquad P=(0,1),\qquad \Delta=80.
\]
The conductor is $40$, $E(\Q)=\Z/4\Z$ is generated by $P$,
and $\widehat h(P)=0$. The index of $\Z P$ in the quotient
modulo torsion is one. The complete bad-prime data are
\[
\begin{array}{c|c|c|c|c}
p&\operatorname{ord}_p\Delta&\text{type}&\text{component of }P&\lambda_p(P)\\\hline
2&4&\mathrm{III}&\text{unique nonzero class, order }2&\log2/12\\
5&1&\mathrm I_1&\text{identity}&\log5/12.
\end{array}
\]
All other finite local heights vanish. Consequently
\[
 \log N=\frac{\log10}{12},\qquad
 D_{\rm part}=\frac6\pi\log\frac{10^{1/12}}{2\sqrt3}.
\]
Numerically the two sides of the global height identity equal
\[
 0.191882091082837140334832621223697017300091790719064414669444\ldots.
\]
With the full real period $\Omega=5.9376498936895458116293755517\ldots$,
the central derivative is
\[
 \frac{L'(E,1)}\Omega=
 0.071331625531738774092617964261659879861112925567606085287841\ldots.
\]
The analytic rank is zero; a vanishing first derivative is therefore
not expected even though the point $P$ is torsion.
\end{proposition}


\begin{proof}
\begin{enumerate}[label=\arabic*.,leftmargin=*]
\item The change of variables from the displayed family is
$x=(24/5)^2X$, $y=(24/5)^3Y$. Tate's algorithm gives the
listed minimal model, conductor and local data. Since $P$ has
order four, $3P=-P$. The local division identity consequently
gives
\[
 \lambda_p(P)=\left(\frac{v_p(\Delta)}{12}
                         -\frac{v_p(\psi_3(P))}{8}\right)\log p.
\]
Here $\psi_3(P)=-4$, proving both local coefficients exactly.
The exact descent output is $[0,0,0,[]]$, and the torsion
calculation gives $\Z/4\Z$. These computations, the direct
global equality and the analytic derivative are in
\path{code/Crank_torsion_local.gp} and
\path{code/Crank_torsion_final.gp}, with their outputs.
As a further check, the exact plus modular symbol from $0$ to
$\infty$ is $-1/4$ in the real-period lattice normalization;
the numerical ratio $L(E,1)/\Omega$ is $1/8$. The completed
functional equation then gives the checked expression
$(\gamma+\log(2\pi/\sqrt{40}))/8$ for the displayed derivative
ratio. The local-height proof does not use this $L$-value check.
\end{enumerate}
\end{proof}

\begin{proposition}[T2; Proved, with exact number-field local arithmetic]\label{ari:C.6.T2}
\StatementRecord{ari:C.6.T2}{proved}{Historical events: GPT-6 Astra (OpenAI), ChatGPT Pro mode, effort Pro, S6\_09\_PAPERS, 17 November 2026: targeted proof check. Also checked in this writing. Pre-continuation writing events: GPT-6 Astra (OpenAI), ChatGPT Pro mode, effort Pro, S6\_10\_PAPERS, 18 November 2026: computation or formal certificate; targeted proof check. Continuation metadata events: GPT-6 Astra (OpenAI), ChatGPT Pro mode, effort Pro, S6\_10\_PAPERS continuation, 19 November 2026: used and rechecked in this run; status-vocabulary reconciliation; metadata reconciliation only; no new mathematical check. \emph{Pre-continuation recorded proof scope.} Read complete printed arithmetic proofs and re-derived the exact scaling/local-height equations and the genus-trace square denominator argument. Rank, saturation, directed L-derivative replays and precise external theorem matches remain explicit gates; no proof verdict is awarded to those computational claims by this read. Read and replayed the exact local program and the Simon two-isogeny descent driver with only its library read paths changed. bnfcertify=1 and LIMBIGPRIME=0; the complete new transcript has the displayed surviving points and excluded local coverings, quotient orders2 and4, hence rank1. Torsion6,three local types,valuations and weighted height sum match. Also reproduced the optional100-digit number-field L-value display, which retains numerical scope. No independent review of the entire Simon library is claimed.}{the hypotheses, notation and definitions stated here; \cref{ari:C.4.4}}

Put $s=\sqrt3$ and
$t_\pm=(3008\pm1728s)/11$. Over $K=\Q(s)$, a global minimal
model for the fibre at $t_-$ is
\[
 E_-:y^2=x^3-(18+10s)x+21+12s,\qquad P_-=(0,3+2s).
\]
The model and point for $t_+$ are its conjugates. The
discriminant of this model is $1032912+596352s$. Define
\[
 \mathfrak p_2=(2,1+s),\qquad\mathfrak p_3=(s),\qquad
 \mathfrak p_{11,-}=(11,s+5),\quad
 \mathfrak p_{11,+}=(11,s-5).
\]
The bad-prime data of $E_-$ are
\[
\begin{array}{c|c|c|c|c|c}
v&e_v&v(\Delta)&\text{type}&\text{component order of }P_-&\lambda_v(P_-)\\\hline
\mathfrak p_2&2&8&\mathrm{IV}^*&3&0\\
\mathfrak p_3&2&3&\mathrm{III}&2&0\\
\mathfrak p_{11,-}&1&1&\mathrm I_1&1&\log11/12.
\end{array}
\]
For $E_+$ replace $\mathfrak p_{11,-}$ by
$\mathfrak p_{11,+}$. The nonzero component at
$\mathfrak p_3$ is unique; the two nonzero classes at
$\mathfrak p_2$ are interchanged by $P\mapsto-P$ and the
class of $P$ is a generator. All omitted local heights vanish.
Both points have order six, and
\begin{align*}
 \log N(P_-)&=\frac1{12}\log(17+10\sqrt3),\\
 \log N(P_+)&=\frac1{12}\log(10\sqrt3-17),\\
 \sum_{\sigma:K\hookrightarrow\R}\log N(P_\sigma)
 &=\frac{\log11}{12}
 =\sum_{v<\infty}n_v\lambda_v(P)-2\widehat h(P).
\end{align*}
The common value in the last line is
\[
 0.199824606066530878671828631497094108318475571161451431268214\ldots.
\]
In particular $A_\mp=(17+10\sqrt3)^{1/12}$ and
$(10\sqrt3-17)^{1/12}$ respectively give
$D_{\rm part}(t_\mp)=(6/\pi)\log(A_\mp/(2\sqrt3))$.

The natural Mordell--Weil group here is $E_\pm(K)$; the
notation $E_{t_\pm}(\Q)$ requested in the stated conventions does not
define an elliptic curve over $\Q$, since its $j$-invariant
is irrational. Over $K$, both curves have rank one and torsion
group $\Z/6\Z$. Thus the index of the torsion section in the
free quotient is infinite. For example $(1,-1-s)$ has infinite
order on $E_-$.
\end{proposition}


\begin{proof}
\begin{enumerate}[label=\arabic*.,leftmargin=*]
\item For $t_-$ the change from the original family to this model
has scale $u=48(-64+37s)/11$, with translation parameters zero.
It is obtained by multiplying the first global minimalization
scale by the unit $(2+s)^{-3}$. Hence it preserves all finite
minimal valuations. The conjugate scale gives $E_+$.
Tate's algorithm yields the displayed table. Its discriminant
ideal is
$\mathfrak p_2^8\mathfrak p_3^3\mathfrak p_{11,-}$ and its
conductor ideal is
$\mathfrak p_2^2\mathfrak p_3^2\mathfrak p_{11,-}$.

\item As $5P=-P$, the normalized local heights are computed exactly by
\[
 \frac{\lambda_v(P)}{\log p}
 =\frac1{e_v}\left(\frac{v(\Delta)}{12}
                         -\frac{v(\psi_5(P))}{24}\right).
\]
The valuations of $\psi_5(P)$ at the three listed primes are
$16,6,0$, respectively. This proves the table, including zero
contribution at the two additive primes. Here $f_v=1$ at all
three primes, so weighting by $n_v=e_vf_v$ gives precisely
$\log11/12$. The exact torsion values proved in Proposition~\ref{ari:C.4.4} have product
$(17+10\sqrt3)(10\sqrt3-17)=11$, proving the archimedean side.

\item For the rank, translate the rational two-torsion point to zero.
The resulting equation is
\[
 y^2=x^3-(9+6s)x^2+(45+26s)x.
\]
The $2$-isogeny descent gives
$[E(K):\widehat\phi E'(K)]=2$ and
$[E'(K):\phi E(K)]=4$. Every surviving covering has a
displayed point; every excluded covering has a recorded local
obstruction at $\mathfrak p_2$ or $\mathfrak p_3$.
The quotient formula $2^{r+2}=2\cdot4$ proves $r=1$.
The field arithmetic was certified by \texttt{bnfcertify};
Simon descent was run with \texttt{LIMBIGPRIME=0}, so all local
tests were deterministic. The complete transcript is
\path{code/Crank_torsion_descent_sage.out}. The unmodified
Simon scripts maintained in Sage are retained with their source
URLs; the older upstream scripts and their unsuccessful
compatibility attempts are also retained. The torsion calculation
gives $\Z/6\Z$, generated by $P$, independently of the rank
calculation.
\end{enumerate}
\end{proof}

\paragraph{The corresponding number-field $L$-derivative.}
For completeness, the natural replacement for the nonexistent
rational-curve quantity at $t_\pm$ can also be computed. Put
\[
 \Omega_K=\frac1{\sqrt{12}}
  \prod_{\sigma:K\hookrightarrow\R}
       \int_{E_\sigma(\R)}|\omega_\sigma|.
\]
The Hasse--Weil $L$-function over $K$ has analytic conductor
$12^2\cdot396=57024$, gamma parameters $(0,0,1,1)$, and sign
$-1$. Direct 100-digit evaluation gives
\begin{align*}
 L'(E/K,1)&=2.046708855992864213402012895411307398829979893338038716037006\ldots,\\
 L'(E/K,1)/\Omega_K
 &=0.260797680551692455136263550519918387218960539765929916318910\ldots.
\end{align*}
This is a value for the curve over $K$, common to the two
conjugates. If the period is instead defined as the product of
the two full real periods, omitting the discriminant factor,
the ratio is
$0.075285805535274833467509818428226475613981023404668273560620\ldots$.
Both conventions and all input periods are explicitly recorded
in \path{code/Crank_torsion_final.out}; neither value enters
the exact torsion height proof.



\subsection{Integral and Gaussian descriptions}

\begin{proposition}[the Mahler-measure identities]\label{ari:C.8.6}
\StatementRecord{ari:C.8.6}{proved}{Historical events: Claude Fable 5.2 (Anthropic), effort high, 16 November 2026: used and rechecked in this run. GPT-6 Astra (OpenAI), ChatGPT Pro mode, effort Pro, S6\_09\_PAPERS, 17 November 2026: targeted proof check. Also checked in this writing. Pre-continuation writing events: GPT-6 Astra (OpenAI), ChatGPT Pro mode, effort Pro, S6\_10\_PAPERS, 18 November 2026: targeted proof check. Continuation metadata events: GPT-6 Astra (OpenAI), ChatGPT Pro mode, effort Pro, S6\_10\_PAPERS continuation, 19 November 2026: used and rechecked in this run; status-vocabulary reconciliation; metadata reconciliation only; no new mathematical check. \emph{Pre-continuation recorded proof scope.} Read complete proof and re-derived its Green-function divisor/mean, beta integral, signed Gaussian Fourier, Betti Jacobian/current-density, division-divisor multiplicity, or heat-equation calculation. For C.8.10 the DeMarco-Mavraki external theorem locator remains a separate gate.}{the hypotheses, notation and definitions stated here; \cref{ari:C.8.5}}

Let $\Lambda\subset\C$ be a lattice of area $A$, let $v$ be a zero
of $\wp_\Lambda$, and let $d\mu=A^{-1}dA_z$ be Haar probability
measure on $\C/\Lambda$. Put $U(z)=\log N(z)$ and
$L=U(v)$. Then
\begin{align}
 6L&=\int_{\C/\Lambda}\log\left|\frac{\wp'(z)^2}{\wp(z)^3}\right|d\mu(z),
 \label{ari:eq:ca-mahler}\\
 N(v)^2&=|\Delta(\Lambda)|^{1/6}
      \exp\left(-\int_{\C/\Lambda}\log|\wp(z)|d\mu(z)\right).
 \label{ari:eq:ca-wpmean}
\end{align}
If $C(u)=1-(g_2/4)u^2-(g_3/4)u^3$, the pushforward of $d\mu$
under $u=1/\wp(z)$ is
\begin{equation}\label{ari:eq:ca-nu}
 d\nu(u)=\frac1{2A}\frac{dA_u}{|u|\,|C(u)|}.
\end{equation}
In particular its unnormalized integral is exactly $2A$, and
\[
 N(v)^6=4\exp\left(\int_{\C}\log|C(u)|d\nu(u)\right).
\]
Writing $e_1,e_2,e_3$ for the three branch values and
$p(x)=\int\log|x-\wp(z)|d\mu(z)$, the maximum theorem is equivalently
\[
 \int\log|C|d\nu\leq\log2,\qquad
 p(0)\geq\frac{p(e_1)+p(e_2)+p(e_3)}3-\frac{\log2}3.
\]
Equality in either inequality occurs exactly when $g_3=0$.
\end{proposition}


\begin{proof}
\begin{enumerate}[label=\arabic*.,leftmargin=*]
\item The mean of $U$ is zero by the theta product calculation in C.1,
and near zero
$U(z)=\log|z|+\frac1{12}\log|\Delta|+o(1)$.
The divisor of $\wp$ is $(v)+(-v)-2(O)$, with multiplicity
two when $v=-v$. Subtract its logarithmic divisor from the
corresponding combination of $U$'s. The result is harmonic
on the compact torus and hence constant. Its expansion at
zero gives
\[
 \log|\wp(z)|=U(z-v)+U(z+v)-2U(z)
                  +\frac16\log|\Delta|-2L.
\]
Integration proves \eqref{ari:eq:ca-wpmean}.
Similarly the divisor of $\wp'$ is the sum of the three
nonzero half-periods minus $3(O)$. The multiplication-by-two
identity of C.1 gives the product of $N$ at those half-periods
as $2$. Since $\wp'(z)\sim-2z^{-3}$, comparison at zero yields
\[
 \log|\wp'(z)|=\sum_{w\in E[2]\setminus\{O\}}U(z-w)-3U(z)
                         +\frac14\log|\Delta|.
\]
Its mean is $\frac14\log|\Delta|$. Twice this identity minus
three times the preceding mean proves \eqref{ari:eq:ca-mahler}.
All logarithmic singularities are locally integrable, so these
integrations also cover a zero at a half-period.

\item The map $z\mapsto x=\wp(z)$ has degree two. Its two inverse
sheets each have Jacobian $|\wp'(z)|^{-2}$; branch values form
a measure-zero set. Thus its pushforward density is
$2/(A|4x^3-g_2x-g_3|)$. Under $x=1/u$ the Jacobian is
$|u|^{-4}$ and the denominator is $4|u|^{-3}|C(u)|$.
This proves \eqref{ari:eq:ca-nu} and its normalization. The identity
$\wp'^2/\wp^3=4C(u)$ gives the next formula.

\item Replacing $v$ above by a half-period gives
$p(e_j)=\frac16\log|\Delta|-2U(w_j)$; summing and using
$\sum_jU(w_j)=\log2$ proves the centroid formulation.
The equality characterization is that of \ref{ari:C.8.5},
since $j=1728$ is equivalent to $g_3=0$ for a nonsingular lattice.
\end{enumerate}
\end{proof}

\begin{proposition}[the two complex beta integrals]\label{ari:C.8.7}
\StatementRecord{ari:C.8.7}{proved}{Historical events: Claude Fable 5.2 (Anthropic), effort high, 16 November 2026: used and rechecked in this run. GPT-6 Astra (OpenAI), ChatGPT Pro mode, effort Pro, S6\_09\_PAPERS, 17 November 2026: targeted proof check. Also checked in this writing. Pre-continuation writing events: GPT-6 Astra (OpenAI), ChatGPT Pro mode, effort Pro, S6\_10\_PAPERS, 18 November 2026: targeted proof check. Continuation metadata events: GPT-6 Astra (OpenAI), ChatGPT Pro mode, effort Pro, S6\_10\_PAPERS continuation, 19 November 2026: used and rechecked in this run; status-vocabulary reconciliation; metadata reconciliation only; no new mathematical check. \emph{Pre-continuation recorded proof scope.} Read complete proof and re-derived its Green-function divisor/mean, beta integral, signed Gaussian Fourier, Betti Jacobian/current-density, division-divisor multiplicity, or heat-equation calculation. For C.8.10 the DeMarco-Mavraki external theorem locator remains a separate gate.}{the hypotheses, notation and definitions stated here; \cref{ari:C.8.6}; \cref{ari:eq:ca-mahler}}

For the square and equianharmonic lattices respectively, the
integrals in \ref{ari:C.8.6} are
\[
 \int\log|C|\,d\nu=\log2,\qquad
 \int\log|C|\,d\nu=\frac32\log3-2\log2.
\]
They recover $N^2=2$ and $N^4=3$ with the same normalization.
\end{proposition}


\begin{proof}
\begin{enumerate}[label=\arabic*.,leftmargin=*]
\item By homogeneity it suffices to take $C(u)=1-u^m$, with $m=2$
or $m=3$. For $-1<\Re s<1-1/m$ put
\[
 I_m(s)=\int_{\C} |u|^{-1}|1-u^m|^{s-1}dA_u.
\]
The substitution $w=u^m$, with all $m$ sheets counted, gives
\[
 I_m(s)=\frac1m B_{\C}\left(\frac1{2m},\frac{s+1}2\right),
 \quad
 B_{\C}(a,b)=\pi\frac{\Gamma(a)\Gamma(b)\Gamma(1-a-b)}
                   {\Gamma(1-a)\Gamma(1-b)\Gamma(a+b)}.
\]
We include a proof of the integral evaluation. For
$0<a,b$ and $a+b<1$, insert
\[
 |w|^{2a-2}=\frac1{\Gamma(1-a)}
       \int_0^\infty t^{-a}e^{-t|w|^2}\,dt
\]
and the corresponding formula for $|1-w|^{2b-2}$.
Tonelli's theorem applies to the positive integrand. The
Gaussian integral over $w$ is
$\pi(t+s)^{-1}\exp(-ts/(t+s))$.
Putting $t=rv$, $s=r(1-v)$ reduces the remaining integrals
to $\Gamma(1-a-b)$ and the beta integral $B(a,b)$, proving
the displayed formula. The same integrable bounds justify
differentiation in a neighbourhood of $s=0$.

\item Consequently, with $a=1/(2m)$ and $\psi=\Gamma'/\Gamma$,
\[
 \frac{I_m'(0)}{I_m(0)}
 =\frac12\left(2\psi(1/2)-\psi(1/2-a)-\psi(1/2+a)\right).
\]
The gamma duplication and triplication formulas, differentiated
logarithmically, evaluate this as $\log2$ for $m=2$ and
$\frac32\log3-2\log2$ for $m=3$. These can also be checked
directly from
$\psi(1/2)=-\gamma-2\log2$,
$\psi(1/4)+\psi(3/4)=-2\gamma-6\log2$ and
$\psi(1/3)+\psi(2/3)=-2\gamma-3\log3$.
Equation \eqref{ari:eq:ca-mahler} completes the normalization test.
\end{enumerate}
\end{proof}

\begin{proposition}[the signed Gaussian identity]\label{ari:C.8.8}
\StatementRecord{ari:C.8.8}{proved}{Historical events: Claude Fable 5.2 (Anthropic), effort high, 16 November 2026: used and rechecked in this run. GPT-6 Astra (OpenAI), ChatGPT Pro mode, effort Pro, S6\_09\_PAPERS, 17 November 2026: targeted proof check. Also checked in this writing. Pre-continuation writing events: GPT-6 Astra (OpenAI), ChatGPT Pro mode, effort Pro, S6\_10\_PAPERS, 18 November 2026: targeted proof check. Continuation metadata events: GPT-6 Astra (OpenAI), ChatGPT Pro mode, effort Pro, S6\_10\_PAPERS continuation, 19 November 2026: used and rechecked in this run; status-vocabulary reconciliation; metadata reconciliation only; no new mathematical check. \emph{Pre-continuation recorded proof scope.} Read complete proof and re-derived its Green-function divisor/mean, beta integral, signed Gaussian Fourier, Betti Jacobian/current-density, division-divisor multiplicity, or heat-equation calculation. For C.8.10 the DeMarco-Mavraki external theorem locator remains a separate gate.}{the hypotheses, notation and definitions stated here; the explicit constructions and calculations in the proof}

For $\Lambda=\Z+\tau\Z$, $T=\Im\tau$, and every $z\in\C$,
\[
 \sqrt{2T}\,|\eta(\tau)|^2N(z,\tau)^2
 =\sum_{m\in\frac12\Lambda}\epsilon(m)
                         e^{-2\pi|z-m|^2/T},
\]
where $\epsilon=-1$ on $\Lambda$ and $+1$ on the three other
cosets of $\Lambda$ in $\frac12\Lambda$.
Both sides and all their derivatives converge locally uniformly.
\end{proposition}


\begin{proof}
\begin{enumerate}[label=\arabic*.,leftmargin=*]
\item Write $z=b+a\tau$ and expand in the Fourier characters
$e^{2\pi i(kb+la)}$ of $(a,b)\in(\R/\Z)^2$.
For $G(z)=\sum_{\lambda\in\Lambda}e^{-2\pi|z-\lambda|^2/T}$,
Poisson summation gives the coefficient
\[
 \frac12\exp\left(-\frac{\pi|k\tau-l|^2}{2T}\right).
\]
The right side is the sum of the four half-period translates
of $G$, minus $2G$. Its coefficient is therefore
\[
 (-1)^{k+l+kl}
       \exp\left(-\frac{\pi|k\tau-l|^2}{2T}\right).
\]
To check the sign, it is $+1$ when both $k,l$ are even and
$-1$ otherwise.

\item On the left use
\[
 \vartheta_1(z,\tau)
 =-i\sum_{n\in\Z}(-1)^n
        e^{\pi i(n+1/2)^2\tau+2\pi i(n+1/2)z}.
\]
Multiply by its conjugate and by
$\sqrt{2T}e^{-2\pi(\Im z)^2/T}$.
In the double sum put $k=n-m$ and apply one-dimensional
Poisson summation to the remaining Gaussian in $m$.
Completing the square gives precisely the coefficient just
displayed. This proves equality of the Fourier series.
Positive definiteness of $|k\tau-l|^2$ and Gaussian decay
justify the changes of order, Poisson summation, and every
claimed differentiation.
\end{enumerate}
\end{proof}




\subsection{The Betti measure and torsion equidistribution}

We take $dd^c=(i/\pi)\partial\bar\partial$, so that
$dd^c\log|w|=\delta_0$. For a local complex coordinate $w$,
$dA_w=(i/2)dw\wedge d\bar w$.

\begin{proposition}[Refuted sign and normalization, with the corrected current identity proved]\label{ari:C.8.9}
\StatementRecord{ari:C.8.9}{proved}{Historical events: Claude Fable 5.2 (Anthropic), effort high, 15 November 2026: carried as background without fault found. Claude Fable 5.2 (Anthropic), effort high, 16 November 2026: used and rechecked in this run. GPT-6 Astra (OpenAI), ChatGPT Pro mode, effort Pro, S6\_09\_PAPERS, 17 November 2026: targeted proof check. Also checked in this writing. Pre-continuation writing events: GPT-6 Astra (OpenAI), ChatGPT Pro mode, effort Pro, S6\_10\_PAPERS, 18 November 2026: targeted proof check. Continuation metadata events: GPT-6 Astra (OpenAI), ChatGPT Pro mode, effort Pro, S6\_10\_PAPERS continuation, 19 November 2026: used and rechecked in this run; status-vocabulary reconciliation; metadata reconciliation only; no new mathematical check. \emph{Pre-continuation recorded proof scope.} Read complete proof and re-derived its Green-function divisor/mean, beta integral, signed Gaussian Fourier, Betti Jacobian/current-density, division-divisor multiplicity, or heat-equation calculation. For C.8.10 the DeMarco-Mavraki external theorem locator remains a separate gate.}{the hypotheses, notation and definitions stated here; \cref{ari:C.3.2}; \cref{ari:eq:ca-cuspexp}}

On a local period chart write $\mu=a\tau+b$, with $a,b$ real,
and put $T=\Im\tau$. The positive Betti measure is
\[
 m_P=db\wedge da
 =\frac{|\mu_w-a\tau_w|^2}{T}\,dA_w.
\]
Thus the proposed form $da\wedge db$ has the wrong sign.
The globally single-valued function $L(t)=\log N(P_t)$ satisfies
on $\mathbb P^1(\C)$
\begin{equation}\label{ari:eq:ca-current}
 dd^cL=\frac1{12}\delta_\infty-m_P,\qquad
               \int_{\mathbb P^1}m_P=\frac1{12}.
\end{equation}
There are no atoms of $m_P$ at $0$, $1$ or $\infty$.
Let $T_0=\sqrt3/2$, $\rho=(1+i\sqrt3)/2$, and
$p=\wp'_\rho((2-\rho)/3)$. The density in the ordinary
$t$ coordinate satisfies
\begin{align}
 \frac{dm_P}{dA_t}
 &=\frac{|t|^{-2/3}}{64T_0|p|^{2/3}}
                      \bigl(1+O(|t|^{1/3})\bigr),
                         &&t\longrightarrow0,\label{ari:eq:ca-density0}\\
 \frac{dm_P}{dA_t}
 &=\frac{|t-1|^{-3/2}}{24\pi^2\sqrt{E_4(i)}}
                      \bigl(1+O(|t-1|^{1/2})\bigr),
                         &&t\longrightarrow1.\label{ari:eq:ca-density1}
\end{align}
\end{proposition}


\begin{proof}
\begin{enumerate}[label=\arabic*.,leftmargin=*]
\item On a chart avoiding a zero of the section, the theta quotient
is holomorphic and nonzero. Its log modulus is pluriharmonic.
Only $-\pi(\Im\mu)^2/T$ contributes to the curvature.
Direct differentiation of this real quotient gives
\[
 \partial_w\partial_{\bar w}\frac{(\Im\mu)^2}{T}
       =\frac{|\mu_w-a\tau_w|^2}{2T}.
\]
Also
\[
 a_w=\frac{\mu_w-a\tau_w}{2iT},\qquad
 b_w=-\bar\tau\frac{\mu_w-a\tau_w}{2iT}.
\]
Taking the wedge product proves the formula for $m_P$ and
$dd^cL=-m_P$ on the chart. The transformation of $(a,b)$
under an integral change of period basis has determinant one,
so the measure patches. Translation of $\mu$ by a locally
constant integral period changes neither the formula nor
the measure. The formula remains valid on a root chart
at an elliptic modulus.

\item At $t=0$ the expansion proved in Proposition~\ref{ari:C.3.2} is
\[
 L(t)=\frac14\log3+\frac7{192}\Re t
       -B|t|^{4/3}+O(|t|^{5/3}),\qquad
 B=\frac{9\pi}{512T_0|p|^{2/3}}.
\]
The expansion is a convergent real-analytic expansion on a
root chart, so it can be differentiated termwise with the
indicated orders. The Euclidean identity
$\Delta_{\R^2}|t|^\alpha=\alpha^2|t|^{\alpha-2}$ gives
\eqref{ari:eq:ca-density0}.

\item At $t=1$ put $h=\mu-(1-\tau)/2$. The expansion in Proposition~\ref{ari:C.3.1}
and the modular $j$ expansion give
\[
 t-1=A_1h^4(1+O(h^2)),\qquad
 A_1=\frac94\pi^4E_4(i)>0,\qquad
 L=\log\sqrt2-\frac\pi2|h|^2+O(|h|^4).
\]
Therefore
$L=\log\sqrt2-\pi|t-1|^{1/2}/(2\sqrt{A_1})+O(|t-1|)$.
Termwise differentiation on the root chart proves
\eqref{ari:eq:ca-density1}. Both densities are locally integrable,
and the continuous potentials at $0$ and $1$ have no atoms.

\item For infinity take $u=1/t$, and set $\ell=\log(5+2\sqrt6)$.
The cusp branch is
$\mu=1/2-i\ell/(2\pi)+O(q)$, $q=e^{2\pi i\tau}$.
The theta and eta products, with the nonzero factor
$|2\sin\pi\mu|=2\sqrt3+O(q)$, give
\begin{equation}\label{ari:eq:ca-cuspexp}
 L=\frac1{12}\log|u|+\log(2\sqrt3)-\frac1{12}\log1728
          -\frac{\ell^2}{2\log(1728/|u|)}
          +O\left(\frac{|u|}{|\log|u||}+|u|\right).
\end{equation}
Here $q=u/1728+O(u^2)$; the estimate follows directly by
bounding the convergent products. The nonlogarithmic
correction is continuous at $u=0$ and has no point mass.
Equivalently its Betti density is
\[
 \frac{dm_P}{dA_u}
  =\frac{\ell^2}{2\pi|u|^2\log(1728/|u|)^3}
       \bigl(1+o(1)\bigr),
\]
which is integrable. The single logarithmic term in
\eqref{ari:eq:ca-cuspexp} contributes $\delta_\infty/12$.
This proves the global current identity; pairing it with
the constant function one gives the total mass.
\end{enumerate}
\end{proof}

\begin{proposition}[the correctly normalized equidistribution]\label{ari:C.8.10}
\StatementRecord{ari:C.8.10}{proved}{Historical events: Claude Fable 5.2 (Anthropic), effort high, 15 November 2026: carried as background without fault found. Claude Fable 5.2 (Anthropic), effort high, 16 November 2026: used and rechecked in this run. GPT-6 Astra (OpenAI), ChatGPT Pro mode, effort Pro, S6\_09\_PAPERS, 17 November 2026: targeted proof check. Also checked in this writing. Pre-continuation writing events: GPT-6 Astra (OpenAI), ChatGPT Pro mode, effort Pro, S6\_10\_PAPERS, 18 November 2026: targeted proof check; used and rechecked in this run. Continuation metadata events: GPT-6 Astra (OpenAI), ChatGPT Pro mode, effort Pro, S6\_10\_PAPERS continuation, 19 November 2026: used and rechecked in this run; status-vocabulary reconciliation; metadata reconciliation only; no new mathematical check. \emph{Pre-continuation recorded proof scope.} Read complete proof and re-derived its Green-function divisor/mean, beta integral, signed Gaussian Fourier, Betti Jacobian/current-density, division-divisor multiplicity, or heat-equation calculation. For C.8.10 the DeMarco-Mavraki external theorem locator remains a separate gate. Opened DeMarco–Mavraki Corollary1.2 p2, current pullback description pp13–14 and final equidistribution proof p23; rechecked the receiving passage from distinct Galois orbits to weighted division divisors using bounded fixed-point multiplicities. The probability limit is curvature divided by degree1/12, namely12m\_P; this closes the named source gate.}{the hypotheses, notation and definitions stated here; \citet[Theorem 1.1 and Corollary 1.2]{demarco2017}; \cref{ari:C.6.4}; \cref{ari:C.8.9}}

Let $D_n$ be the divisor, with intersection multiplicity, of
parameters $t\in\mathbb P^1\setminus\{0,1,\infty\}$ for which
$nP_t=O$. Then
\[
 \deg D_n=\frac{n^2}{12}+O(1),\qquad
 \frac1{n^2}[D_n]\ \longrightarrow\ m_P,\qquad
 \frac1{\deg D_n}[D_n]\ \longrightarrow\ 12m_P
\]
weakly as $n\to\infty$ through indices for which $D_n\ne0$.
For any nonrepeating sequence of individual algebraic torsion
parameters, their Galois-orbit probability measures likewise
converge to $12m_P$. Thus the probability limit requires the factor $12$.
\end{proposition}


\begin{proof}
\begin{enumerate}[label=\arabic*.,leftmargin=*]
\item We apply \citet[Theorem 1.1 and Corollary 1.2]{demarco2017}.
Their hypotheses are a number field, a smooth projective
base curve, an elliptic surface and a rational section of
nonzero generic canonical height. Here they are $\Q$,
$\mathbb P^1$, the displayed rational elliptic surface,
and $P$, whose generic height for $(O)$ is $1/12$ by \ref{ari:C.6.4}.
At a torsion specialization the canonical height is zero.
Corollary 1.2 consequently gives convergence of every
nonrepeating sequence of Galois-orbit probability measures.

\item The limiting measure is the curvature of the height metric
divided by the degree of its line bundle. In our normalization,
\ref{ari:C.8.9} computes the curvature as $m_P$, and \ref{ari:C.6.4} computes
that degree as $1/12$. Hence the probability measure is
$12m_P$. This division by the degree is essential: a measure
of mass $1/12$ cannot be a weak limit of probability measures.
The source's informal description of the pullback current
is interpreted with this normalization.

\item For completeness, the passage from individual orbits to
the full division divisor does not require uniformity of
the equidistribution theorem in $n$. The surface height
formula, with our half-self-pairing convention, reads
\[
 \widehat h_O(Q)=\chi(\mathcal O_{\mathcal E})+(Q\cdot O)
                -\frac12\sum_v\operatorname{contr}_v(Q).
\]
Here $\chi=1$, and the finitely many fibre corrections are
bounded as $Q=nP$ varies, because they depend only on its
component. Since $\widehat h_O(nP)=n^2/12$, it follows
that $(nP\cdot O)=n^2/12+O(1)$. Removing the singular fibres
changes this by $O(1)$: the local intersection at each
such fibre is bounded as $n$ varies in characteristic zero.
One can see this in the formal group after a fixed
multiple enters it; multiplication by a nonzero integer
is a local analytic isomorphism over the residue field
$\C$. Thus $\deg D_n=n^2/12+O(1)$.

\item At any fixed smooth torsion parameter $t_0$, of exact order
$d$, the multiplicity in $D_n$ is zero unless $d\mid n$,
and is constant when $d\mid n$. Indeed the formal parameter
of $nP=[n/d](dP)$ has first term $(n/d)$ times that of $dP$,
with nonzero constant coefficient. Fix a continuous test
function and $\varepsilon>0$. Corollary 1.2 implies that
all but finitely many distinct Galois orbits of torsion
parameters have average within $\varepsilon$ of the limiting
integral; otherwise a countersequence would violate it.
The total weighted contribution of the finitely many
exceptional orbits is bounded independently of $n$, by
the multiplicity observation. Divide by $\deg D_n$ and
let $n\to\infty$. This proves the probability convergence
for $D_n$, and the degree asymptotic proves the unnormalized
version.
\end{enumerate}
\end{proof}

\begin{proposition}[the two first-derivative identities]\label{ari:C.8.11}
\StatementRecord{ari:C.8.11}{proved}{Historical events: Claude Fable 5.2 (Anthropic), effort high, 16 November 2026: used and rechecked in this run. GPT-6 Astra (OpenAI), ChatGPT Pro mode, effort Pro, S6\_09\_PAPERS, 17 November 2026: targeted proof check. Also checked in this writing. Pre-continuation writing events: GPT-6 Astra (OpenAI), ChatGPT Pro mode, effort Pro, S6\_10\_PAPERS, 18 November 2026: targeted proof check. Continuation metadata events: GPT-6 Astra (OpenAI), ChatGPT Pro mode, effort Pro, S6\_10\_PAPERS continuation, 19 November 2026: used and rechecked in this run; status-vocabulary reconciliation; metadata reconciliation only; no new mathematical check. \emph{Pre-continuation recorded proof scope.} Read complete proof and re-derived its Green-function divisor/mean, beta integral, signed Gaussian Fourier, Betti Jacobian/current-density, division-divisor multiplicity, or heat-equation calculation. For C.8.10 the DeMarco-Mavraki external theorem locator remains a separate gate.}{the hypotheses, notation and definitions stated here; the explicit constructions and calculations in the proof}

On an ordinary modulus chart put
\[
 a=\frac{\Im\mu}{T},\qquad
 W=\frac{\vartheta_1'(\mu,\tau)}{\vartheta_1(\mu,\tau)}
                         +2\pi ia,\qquad
 R=\frac{d\mu}{d\tau}-a.
\]
Then
\[
 L_\tau=\frac{W}{8\pi i}(W+4\pi iR),\qquad
 \det\frac{d}{d\tau}
 \begin{pmatrix}\tau&\mu\\\mu&\beta^{\rm part}/6\end{pmatrix}
       =\frac{(W+2\pi iR)^2}{4\pi^2}.
\]
Moreover
\[
 R=\frac{g_2}{6\pi i\wp'(\mu)}-\frac{W}{2\pi i},
 \qquad
 L_\tau=\frac{iW(3W\wp'(\mu)-2g_2)}
                         {24\pi\wp'(\mu)}.
\]
Multiplication by $6/\pi$ gives the critical-point factorization for
$D_{\rm part}$ in Proposition~\ref{push:src:D.4.4}.
\end{proposition}


\begin{proof}
\begin{enumerate}[label=\arabic*.,leftmargin=*]
\item Put $f=\vartheta_1'/\vartheta_1$ and
$E=\pi^2E_2/3$. The theta heat equation and
$\eta'/\eta=\pi iE_2/12$ give
\[
 \partial_\tau\log(\vartheta_1/\eta)
       =\frac{f_z+f^2+E}{4\pi i}.
\]
At $\wp(\mu)=0$, the identity $f_z=-\wp-E$ reduces
this to $f^2/(4\pi i)$. Thus, along the section,
\[
 \frac{d}{d\tau}\log(\vartheta_1(\mu,\tau)/\eta)
       =f\mu'+\frac{f^2}{4\pi i}.
\]
Differentiate the Gaussian $-\pi(\Im\mu)^2/T$ using
Wirtinger derivatives, retaining the factor $1/2$ in
the derivative of the real part of a holomorphic logarithm.
Substitution of $f=W-2\pi ia$ and $\mu'=R+a$ gives the
first formula. Since
$\beta^{\rm part}/6=(i/\pi)\log(\vartheta_1/(2\sqrt3\eta))$,
its derivative is $if\mu'/\pi+f^2/(4\pi^2)$.
Subtracting $\mu'^2$ gives
$(f+2\pi i\mu')^2/(4\pi^2)$, the second formula.

\item Finally differentiate $\wp_\tau(\mu(\tau))=0$ by the theta
heat equation. After substituting $\wp(\mu)=0$,
$\wp''(\mu)=-g_2/2$, and $f=W-2\pi ia$, the result is
$6\pi i\wp'R=g_2-3W\wp'$. This is the formula for $R$.
Inserting it into the first identity gives the last one.
These computations independently verify the differential formula used in Proposition~\ref{ari:C.3.2};
its apparent singularity at a half-period is
handled by the root chart rather than division by zero.
\end{enumerate}
\end{proof}



\section{Critical points and the global level curves}

We continue to write $L(t)=\log N(\mu(\tau),\tau)$, where
$t=j(\tau)/1728$ and $\wp_\tau(\mu)=0$. The sign and lattice
translation of $\mu$ do not affect any expression below. We use
Eichler and Zagier's unnumbered theorem on p.~399 of
\cite{EichlerZagier1982}, also stated as equation (6) and the
unnumbered theorem in Section~2 of \cite{DukeImamoglu2008}.
For the cusp branch in this problem their theorem says
\begin{equation}\label{ari:eq:cc-EZ}
 \begin{aligned}
 \mu&=\frac12-i\beta_{\rm cusp}-
       \frac{36i\sqrt6}{\pi}\sum_{n\ge1}\frac{a_nq^n}{n^2},
 &\beta_{\rm cusp}&=\frac{\log(5+2\sqrt6)}{2\pi},\\
 \sum_{n\ge1}a_nq^n&=\frac{\Delta_0}{E_6^{3/2}},
 &\mu''&=144\pi i\sqrt6\frac{\Delta_0}{E_6^{3/2}}.
 \end{aligned}
\end{equation}
The branch of $E_6^{3/2}$ tends to one at the cusp. Thus the
formula applies directly on $\Im\tau>1$ and then by analytic
continuation away from the half-period branch points. Our use at $i$ is on the square-root chart specified explicitly below.
The original theorem and the hypergeometric reformulation concern
the logarithm $\mu$; neither asserts the global inequality for
$L$ proved here. Searches through these papers and their cited
work, and through the torus Green-function literature, did not
locate the radial monotonicity statement below.

\begin{proposition}[a radial derivative identity]\label{ari:C.8.12}
\StatementRecord{ari:C.8.12}{proved}{Historical events: Claude Fable 5.2 (Anthropic), effort high, 15 November 2026: carried as background without fault found. GPT-6 Astra (OpenAI), ChatGPT Pro mode, effort Pro, S6\_09\_PAPERS, 17 November 2026: targeted proof check. Also checked in this writing. Pre-continuation writing events: GPT-6 Astra (OpenAI), ChatGPT Pro mode, effort Pro, S6\_10\_PAPERS, 18 November 2026: targeted proof check. Continuation metadata events: GPT-6 Astra (OpenAI), ChatGPT Pro mode, effort Pro, S6\_10\_PAPERS continuation, 19 November 2026: used and rechecked in this run; status-vocabulary reconciliation; metadata reconciliation only; no new mathematical check. \emph{Pre-continuation recorded proof scope.} Read complete proof. Re-derived critical factor, radial derivative, cusp Fourier bound, Cayley ramification removal, rational-series tail bounds, scalar Bernstein implication and raywise Jordan-curve argument. The precise imported elliptic-function formula remains an external-locator gate until opened.}{the hypotheses, notation and definitions stated here; \cref{ari:C.8.11}; \cref{ari:C.3.2}}

Put $p=\wp'_\tau(\mu)$ and
\[
 X=\frac{3Wp}{g_2},\qquad F=1-X,
 \qquad W=\frac{\vartheta_1'(\mu,\tau)}{\vartheta_1(\mu,\tau)}
              +\frac{2\pi i\Im\mu}{\Im\tau}.
\]
These are single-valued modular invariants wherever $g_2\ne0$.
They have continuous extensions at $t=0,1$, with
$X(0)=1/4$ and $X(1)=0$. At every $t\ne0,1$,
\begin{equation}\label{ari:eq:cc-radial}
 (t-1)L_t=\frac{X(X-2)}{24},\qquad
 r\frac{d}{dr}L(1+re^{i\theta})
       =\frac{\Re(X(X-2))}{12}.
\end{equation}
Also, if $a=\Im\mu/\Im\tau$ and $R=\mu'-a$, then
\begin{equation}\label{ari:eq:cc-FR}
 F=\frac{6\pi ip}{g_2}R.
\end{equation}
Consequently $|F|<1$ implies strict decrease along the ray
from $1$ through $t$.
\end{proposition}


\begin{proof}
\begin{enumerate}[label=\arabic*.,leftmargin=*]
\item A lattice translation leaves $W$ unchanged, and replacing $\mu$
by $-\mu$ changes the signs of both $W$ and $p$. Under a modular
change, $W,p,g_2$ have respective weights $1,3,4$, so $X$ is
invariant. The two identities of \ref{ari:C.8.11} give
\[
 L_\tau=\frac{iW(3Wp-2g_2)}{24\pi p},\qquad
 R=\frac{g_2-3Wp}{6\pi ip}.
\]
Ramanujan's identities and $p^2=-g_3$ give
\[
 t'_\tau=-2\pi it\frac{E_6}{E_4},\qquad
 \frac{t-1}{t'_\tau}=-\frac{E_6}{2\pi iE_4}.
\]
Substitution, using $g_2=4\pi^4E_4/3$ and
$g_3=8\pi^6E_6/27$, proves \eqref{ari:eq:cc-radial} and
\eqref{ari:eq:cc-FR}. Since $X(X-2)=F^2-1$, $|F|<1$ makes
its real part strictly negative.

\item At $i$, $W=p=0$ and $g_2\ne0$, so $X\to0$. At $\rho$, put $\delta=\tau-\rho$ and $p_0=\wp'_\rho(\mu_0)$.
The expansion in \ref{ari:C.3.2} and Ramanujan's identities give
\[
 R=\frac{3p_0}{8\pi^2}\delta+O(|\delta|^2),\qquad
 g_2=\frac{3ip_0^2}{\pi}\delta+O(\delta^2).
\]
Substitution into \eqref{ari:eq:cc-FR} gives $F(\rho)=3/4$.
Equivalently, the two possible limits furnished by
\eqref{ari:eq:cc-radial} and $L_t(0)=7/384$ are $X=1/4,7/4$;
the expansion of $R$ selects $1/4$. The same expansions show
continuity on the coarse $t$-line.
\end{enumerate}
\end{proof}

\begin{proposition}[the cusp exclusion]\label{ari:C.8.13}
\StatementRecord{ari:C.8.13}{proved}{Historical events: GPT-6 Astra (OpenAI), ChatGPT Pro mode, effort Pro, S6\_09\_PAPERS, 17 November 2026: targeted proof check. Also checked in this writing. Pre-continuation writing events: GPT-6 Astra (OpenAI), ChatGPT Pro mode, effort Pro, S6\_10\_PAPERS, 18 November 2026: computation or formal certificate; targeted proof check. Continuation metadata events: GPT-6 Astra (OpenAI), ChatGPT Pro mode, effort Pro, S6\_10\_PAPERS continuation, 19 November 2026: used and rechecked in this run; status-vocabulary reconciliation; metadata reconciliation only; no new mathematical check. \emph{Pre-continuation recorded proof scope.} Read complete proof. Re-derived critical factor, radial derivative, cusp Fourier bound, Cayley ramification removal, rational-series tail bounds, scalar Bernstein implication and raywise Jordan-curve argument. The precise imported elliptic-function formula remains an external-locator gate until opened. Reran Ccritical\_rational\_certificate.py and Ccritical\_scalar.py. All 100 exact ODE coefficients and outward integer bounds were recomputed; all 15 exact Q(sqrt(3)) Bernstein coefficients are positive. Saved stdout and exact output certificates accompany the code.}{the hypotheses, notation and definitions stated here; \cref{ari:eq:cc-EZ}; \cref{ari:eq:ca-exp}}

If $\Im\tau\ge7/5$, then $|F|<17/20$.
\end{proposition}


\begin{proof}
\begin{enumerate}[label=\arabic*.,leftmargin=*]
\item Here $|q|<1/6600$. On the auxiliary circle $|q|=1/2000$,
the Lambert series and eta product give
\[
 |E_6-1|\le
 504r\frac{1+26r+66r^2+26r^3+r^4}{(1-r)^7},\qquad
 |\Delta_0|\le r\exp\!\frac{24r}{(1-r)^2},
 \quad r=\frac1{2000}.
\]
The right sides imply $|\Delta_0/E_6^{3/2}|<1/1250$.
In particular $E_6$ has no zero on this disc and the specified
branch is holomorphic there. By Cauchy's inequality, with
$s=10/33$,
\[
 \sum_{n\ge1}\frac{|a_nq^n|}{n}
 \le\frac{-\log(1-s)}{1250},\qquad
 \sum_{n\ge1}\frac{|a_nq^n|}{n^2}
 \le\frac1{1250}\left(s+\frac{s^2}{4(1-s)}\right).
\]
Equation \eqref{ari:eq:cc-EZ} therefore yields
\[
 |\mu-(1/2-i\beta_{\rm cusp})|<\frac1{125},\qquad
 |\mu'|<\frac{13}{250},\qquad \beta_{\rm cusp}<\frac{365}{1000}.
\]
For $r=1/6600$, the corresponding weight-four and weight-six
Lambert bounds are
\[
 240r\frac{1+4r+r^2}{(1-r)^5}<\frac{37}{1000},\qquad
 504r\frac{1+26r+66r^2+26r^3+r^4}{(1-r)^7}
       <\frac{77}{1000}.
\]
Thus $|E_4|>963/1000$ and $|E_6|^{1/2}<1038/1000$.
Since $6\pi|p/g_2|=\sqrt6|E_6|^{1/2}/|E_4|$,
\begin{equation}\label{ari:eq:cc-cuspbound}
 |F|<\sqrt6\frac{1038}{963}
 \left\{\frac{365/1000+1/125}{7/5}+\frac{13}{250}\right\}
 <\frac{17}{20}.
\end{equation}
All scalar comparisons here follow from geometric series and
the exponential Taylor bounds \eqref{ari:eq:ca-exp}; for the logarithm
use $-\log(1-s)=\sum s^n/n$ with its geometric remainder.
For reference, the upper quantities before the last rounding are
less than $789/10^6$, $755/10^5$, $51/1000$, and $841/1000$
for $|\Delta_0/E_6^{3/2}|$, the error in $\mu$, the bound on
$\mu'$, and the final bound on $|F|$, respectively.
\end{enumerate}
\end{proof}

\begin{proposition}[the compact exclusion by one Cayley disc]\label{ari:C.8.14}
\StatementRecord{ari:C.8.14}{proved}{Historical events: GPT-6 Astra (OpenAI), ChatGPT Pro mode, effort Pro, S6\_09\_PAPERS, 17 November 2026: targeted proof check. Also checked in this writing. Pre-continuation writing events: GPT-6 Astra (OpenAI), ChatGPT Pro mode, effort Pro, S6\_10\_PAPERS, 18 November 2026: computation or formal certificate; targeted proof check. Continuation metadata events: GPT-6 Astra (OpenAI), ChatGPT Pro mode, effort Pro, S6\_10\_PAPERS continuation, 19 November 2026: used and rechecked in this run; status-vocabulary reconciliation; metadata reconciliation only; no new mathematical check. \emph{Pre-continuation recorded proof scope.} Read complete proof. Re-derived critical factor, radial derivative, cusp Fourier bound, Cayley ramification removal, rational-series tail bounds, scalar Bernstein implication and raywise Jordan-curve argument. The precise imported elliptic-function formula remains an external-locator gate until opened. PASS: unchanged Ccritical\_rational\_certificate.py recomputed exact modular ODE coefficients and outward-rounded rational coefficient norms; Ccritical\_scalar.py recomputed all 15 positive Bernstein coefficients over Q(sqrt(3)). Outputs code/I/replay\_arithmetic/. Reran Ccritical\_rational\_certificate.py and Ccritical\_scalar.py. All 100 exact ODE coefficients and outward integer bounds were recomputed; all 15 exact Q(sqrt(3)) Bernstein coefficients are positive. Saved stdout and exact output certificates accompany the code.}{the hypotheses, notation and definitions stated here; \cref{ari:eq:cc-EZ}; \cref{ari:eq:cc-Fseries}; \cref{ari:eq:cc-ode}; \cref{ari:eq:cc-recurrence}; \cref{ari:eq:cc-power}; \cref{ari:eq:cc-finite-table}; \cref{ari:eq:cc-norms}; \cref{ari:eq:cc-Ffactored}; \cref{ari:eq:cc-scalar}; \cref{ari:C.8.12}}

Let $\chi=(\tau-i)/(\tau+i)$. On
$0<|\chi|\le2-\sqrt3$, the invariant $F$ satisfies $|F|<1$.
At the two equianharmonic points on the boundary its continuous
value is $3/4$.
\end{proposition}


\begin{proof}
\begin{enumerate}[label=\arabic*.,leftmargin=*]
\item Set
\[
 \begin{gathered}
 r_0=2-\sqrt3,\qquad A=E_4(i),\qquad k=\pi^2A,\\
 v=\pi\sqrt A\,\chi,\qquad z=v^2=k\chi^2,\\
 z_0=kr_0^2,\qquad c_0=1-\frac1{\sqrt3}.
 \end{gathered}
\]
In this proof $z$ is a coordinate on the modular disc, not a
point of the elliptic curve. Define analytic functions $f_2,f_4,f_6$
by
\[
 \begin{gathered}
 \sqrt A f_2(v)=\frac{E_2(\tau)}{(1-\chi)^2}
                         -\frac3{\pi(1-\chi)},\\
 A f_4(v)=\frac{E_4(\tau)}{(1-\chi)^4},\qquad
 A^{3/2} f_6(v)=\frac{E_6(\tau)}{(1-\chi)^6}.
 \end{gathered}
\]
Ramanujan's differential identities become the rational system
\begin{equation}\label{ari:eq:cc-ode}
 \begin{gathered}
 f_2'=\frac{f_4-f_2^2}{3},\qquad
 f_4'=\frac43(f_6-f_2f_4),\qquad
 f_6'=2(f_4^2-f_2f_6),\\
 (f_2,f_4,f_6)(0)=(0,1,0).
 \end{gathered}
\end{equation}
Uniqueness makes $f_4$ even and $f_2,f_6$ odd. Write
\[
 B(z)=f_4(v),\quad J(z)=\frac{f_6(v)}{2v},\quad
 D(z)=\frac{B(z)^3-4zJ(z)^2}{J(z)^{3/2}},\qquad
 H(z)=\sum_{n\ge0}\frac{-[z^n]D(z)}{16n^2-1}z^n.
\]
The branch has $J(0)^{1/2}=1$. The estimates below establish
its regularity on a larger disc, so these definitions are
unambiguous on $|z|\le z_0$.

\item Write $h=(1-\chi)(\mu-(1-\tau)/2)$. The Möbius chain rule and
\eqref{ari:eq:cc-EZ} give
\[
 h''(\chi)=-\frac{4\mu''(\tau)}{(1-\chi)^3}
       =-\frac{i}{2\sqrt{3\pi}}\chi^{-3/2}D(k\chi^2).
\]
The root expansion at $i$ gives the branch
\begin{equation}\label{ari:eq:cc-h}
 h=\frac{2i}{\sqrt{3\pi}}\sqrt \chi\,H(z).
\end{equation}
Indeed differentiating this series twice gives precisely the
coefficient denominator $16n^2-1$ above. The two possible affine
homogeneous terms are absent because the half-period difference
has a Puiseux expansion consisting only of half-integral powers.
Changing the sign of the square root changes the signs of $h,p,R$
and leaves $F$ unchanged.

\item Put $r=|\chi|$. The definitions of $a$ and $R$ give exactly
\[
 R=\frac{1-\chi}{2i}
       \left(h'+\frac{\overline \chi h+\overline h}{1-r^2}\right),
 \qquad
 \frac p{g_2}=-\frac{i\sqrt \chi}{\sqrt{3\pi}(1-\chi)}
                         \frac{\sqrt J}{B}.
\]
Consequently
\begin{equation}\label{ari:eq:cc-Fseries}
 F=\frac{\sqrt J}{B}
 \left\{H+4zH'+\frac{2r}{1-r^2}(rH-\overline H)\right\}.
\end{equation}
All conjugates in this display are values at the same point;
$H$ has real coefficients, so $\overline{H(z)}=H(\overline z)$.

\item The apparent pole at the equianharmonic points must be removed
before estimating this expression. We have $B(-z_0)=0$.
At these points $R=0$, since the expansion at $\rho$ gives
$\mu'=-1/3=a$. Thus, if
\[
 S=\sqrt J H,\qquad T=\sqrt J(H+4zH'),
\]
then $T(-z_0)=c_0S(-z_0)$. Define the analytic quotients
\[
 C(z)=\frac{B(z)}{1+z/z_0},\qquad
 U(z)=\frac{T(z)-c_0S(z)}{(1+z/z_0)C(z)},\qquad
 V(z)=\frac{\sqrt{J(z)}}{C(z)}.
\]
The estimates below prove $C\ne0$ on the needed larger disc.
For $r<r_0$, equation \eqref{ari:eq:cc-Fseries} is now
\begin{equation}\label{ari:eq:cc-Ffactored}
 F=U+V\frac{dH+q(H-\overline H)}{1+z/z_0},
 \qquad
 d=\frac{2(r_0-r)}{(1+r_0)(1+r)},\quad q=\frac{2r}{1-r^2}.
\end{equation}
Here $q$ is only a real abbreviation within this display and
its estimates, not the Fourier nome.

\item We prove the following four coefficient bounds. For an analytic
power series $K=\sum K_nz^n$, put
$\|K\|_s=\sum|K_n|s^n$. Then, for $0\le s\le z_0$,
\begin{equation}\label{ari:eq:cc-norms}
 \begin{aligned}
 \|U\|_s&\le\frac1{\sqrt3}+\frac{13}{200}s,&
 \|V\|_s&\le1+\frac7{50}s,\\
 \|H\|_s&\le1+\frac{113}{1000}s,&
 \|H'\|_s&\le\frac{13}{100}.
 \end{aligned}
\end{equation}
We give the coefficient arithmetic and all tails, rather than
leaving these inequalities to sampling a circle. In
\eqref{ari:eq:cc-ode}, write $f_2=\sum a_nv^n$,
$f_4=\sum b_nv^n$, $f_6=\sum c_nv^n$. Starting with
$(a_0,b_0,c_0)=(0,1,0)$ (only in this recurrence the symbol
$c_0$ denotes the zero initial coefficient), compute
\begin{equation}\label{ari:eq:cc-recurrence}
 \begin{aligned}
 a_{n+1}&=\frac{b_n-\sum_{j=0}^na_ja_{n-j}}{3(n+1)},\\
 b_{n+1}&=\frac{4(c_n-\sum_{j=0}^na_jb_{n-j})}{3(n+1)},\\
 c_{n+1}&=\frac{2(\sum_{j=0}^nb_jb_{n-j}-
                          \sum_{j=0}^na_jc_{n-j})}{n+1}.
 \end{aligned}
\end{equation}
Thus $B_n=b_{2n}$ and $J_n=c_{2n+1}/2$. For powers of a series
$K$ with $K_0=1$, the equally elementary recurrence is
\begin{equation}\label{ari:eq:cc-power}
 [z^n]K^\alpha=\frac1n\sum_{j=1}^n
                 ((\alpha+1)j-n)K_j[z^{n-j}]K^\alpha.
\end{equation}
Together with $H_n=-D_n/(16n^2-1)$, these determine every
coefficient used. For example,
\[
 \begin{array}{c|rrrr}
 n&0&1&2&3\\ \hline
 B_n&1&10/9&5/27&4/81\\
 J_n&1&14/27&28/135&14/405\\
 H_n&1&13/135&-1999/153090&29299/28146690.
 \end{array}
\]
Division by $1+z/z_0$ uses $C_0=1$ and
$C_n=B_n-C_{n-1}/z_0$; the same recurrence applies to
$(T-c_0S)/(1+z/z_0)$. Division by $C$ uses the usual
convolution recurrence with constant coefficient one.

\item Here are sufficient rational enclosures for the only nonrational
constants in this arithmetic:
\begin{equation}\label{ari:eq:cc-constants}
 \begin{gathered}
 14<k<72/5,\qquad
 1.0315619043558156493568263343177356527388253980203
 <z_0\\
 <1.0315619043558156493568263343177356527388253980205.
 \end{gathered}
\end{equation}
They follow without special-function evaluation from
$\pi=16\arctan(1/5)-4\arctan(1/239)$, the first 100 terms
of the alternating arctangent series, the exponential series
through degree 160, and
\[
 A=1+240\sum_{n=1}^{49}\frac{n^3e^{-2\pi n}}{1-e^{-2\pi n}}
       +\varepsilon_A,\qquad
 0<\varepsilon_A<
 \frac{240(4e^{-2\pi})^{50}}
 {(1-4e^{-2\pi})(1-e^{-2\pi})}.
\]
The last bound uses $n^3\le4^n$. Bounds for $\sqrt3$ follow
by squaring rational endpoints. All ensuing coefficient
comparisons are rational arithmetic using these enclosures.

\item To bound infinite tails, first consider $|z|=12/5$.
Here $|\chi|<5/12$, $\Im\tau>7/17$ and
$|e^{2\pi i\tau}|<1/13$. Lambert series give
\begin{equation}\label{ari:eq:cc-outer}
 |B|<400,\qquad |J|<3000\qquad (|z|=12/5).
\end{equation}
For clarity the respective upper bounds before rounding are
\[
 \left(\frac{12}{7}\right)^4
 \left(1+240u\frac{1+4u+u^2}{(1-u)^5}\right)<400,
 \quad u=1/13,
\]
\[
 \frac13\left(\frac{12}{7}\right)^6
 \left(1+504u\frac{1+26u+66u^2+26u^3+u^4}{(1-u)^7}\right)<3000.
\]
We have used $A>1$ and $2\sqrt{12/5}>3$.
For $R=3/2$, let $\xi=5/8$ and retain coefficients through
degree 59. The tails of $B,J,C$ are respectively at most
\begin{equation}\label{ari:eq:cc-tail60}
 \frac{400\xi^{60}}{1-\xi},\qquad
 \frac{3000\xi^{60}}{1-\xi},\qquad
 \frac{400z_0}{12/5-z_0}\frac{\xi^{60}}{1-\xi}.
\end{equation}
For the last bound, $B(-z_0)=0$ implies
$|C_n|\le400z_0(12/5)^{-n}/(12/5-z_0)$.
The finite sums from \eqref{ari:eq:cc-recurrence}--\eqref{ari:eq:cc-power}
with these tails give
\[
 \begin{array}{c|cccc}
 &\|B\|_{3/2}&\|J\|_{3/2}&\|C-1\|_{3/2}
                   &\min_{|z|\le3/2}\Re J(z)\\ \hline
 \text{bound}&<3276/1000&<2411/1000&<335/1000&>2049/10000.
 \end{array}
\]
For the last column, the lower bound used is exactly
\[
 1-J_2R^2-\frac{J_1^2}{8J_2}
       -\sum_{n=3}^{59}|J_n|R^n-
                   \frac{3000\xi^{60}}{1-\xi}.
\]
This is obtained by minimizing the quadratic
$\Re(1+J_1z+J_2z^2)$ on $|z|=R$; the minimum occurs at
$\cos\arg z=-J_1/(4J_2R)$. In particular, on $|z|\le3/2$,
\[
 |B|<33/10,\quad |J|<5/2,\quad |C|>3/5,\quad
 \Re J>1/5,\quad
 |D|\le\frac{(33/10)^3+6(5/2)^2}{(1/5)^{3/2}}<1000.
\]
Cauchy's inequality and the denominators in $H_n$ now give
\[
 |H|\le1+1000\sum_{n\ge1}\frac1{16n^2-1}
       =1+1000\left(\frac12-\frac\pi8\right)<110
                      \quad (|z|\le3/2).
\]
On $|z|\le7/5$, the same argument gives
\[
 |H+4zH'|\le1+1000\sum_{n\ge1}
                 \frac{(14/15)^n}{4n-1}
       \le1+\frac{1000}{3}\log15<910.
\]
It follows on $|z|=7/5$ that $|U|<10000$ and $|V|<3$;
indeed $|(1+z/z_0)C|>(7/(5z_0)-1)3/5$, and
$|\sqrt J|<8/5$.

\item Finally retain coefficients through degree 99 and put
$\eta=5z_0/7$. The tails in the coefficient norms for
$U,V,H,H'$ at $z_0$ are bounded respectively by
\begin{equation}\label{ari:eq:cc-tail100}
 \frac{10000\eta^{100}}{1-\eta},\quad
 \frac{3\eta^{100}}{1-\eta},\quad
 \frac{110\eta^{100}}{1-\eta},\quad
 \frac{110\eta^{100}}{z_0}
 \left(\frac{100}{1-\eta}+\frac{\eta}{(1-\eta)^2}\right).
\end{equation}
The finite coefficient sums plus precisely these tails give
\begin{equation}\label{ari:eq:cc-finite-table}
 \begin{array}{c|cccc}
 & (\|U\|_{z_0}-1/\sqrt3)/z_0
 & (\|V\|_{z_0}-1)/z_0
 & (\|H\|_{z_0}-1)/z_0 & \|H'\|_{z_0}\\ \hline
 \text{upper bound}&635/10000&1368/10000&1115/10000&1293/10000.
 \end{array}
\end{equation}
The recurrences, cutoffs and remainders above specify each finite
comparison completely. The attached exact integer/rational
script reproduces them with outward rounding on denominator
$10^{100}$ and independently checks every rational ODE
coefficient; no point of the modular domain is enumerated.
Equations \eqref{ari:eq:cc-finite-table} imply
\eqref{ari:eq:cc-norms}, with room in each inequality.

\item We finish with a single polynomial inequality. For $s=|z|=kr^2$,
\[
 |H-\overline H|\le2|\Im z|\,\|H'\|_s,\qquad
 \frac{|\Im z|}{|1+z/z_0|}\le s.
\]
The second inequality follows by writing $z/z_0=ue^{i\phi}$
and using $(\cos\phi+u)^2\ge0$. The useful cancellation is
\[
 |dH+q(H-\overline H)|^2
      =d^2|H|^2+4q(d+q)(\Im H)^2.
\]
Also
\[
 \frac d{|1+z/z_0|}\le
 e(r):=\frac{c_0}{(1+r)(1+r/r_0)},\qquad
 q\le\kappa r,\quad d+q\le c_0+\kappa r^2,
 \quad\kappa=1+\frac2{\sqrt3}.
\]
The last bounds use $(1-r_0^2)^{-1}=1/2+1/\sqrt3$.
It follows from \eqref{ari:eq:cc-Ffactored}--\eqref{ari:eq:cc-norms},
using $k<72/5$, that
\begin{equation}\label{ari:eq:cc-scalar}
 |F|\le U_0(r)+V_0(r)
 \sqrt{e(r)^2H_0(r)^2+
       4\kappa r(c_0+\kappa r^2)(72/5)^2r^4(13/100)^2},
\end{equation}
where
\[
 U_0=\frac1{\sqrt3}+\frac{13}{200}\frac{72}{5}r^2,
 \quad V_0=1+\frac7{50}\frac{72}{5}r^2,
 \quad H_0=1+\frac{113}{1000}\frac{72}{5}r^2.
\]
To verify that the right side is less than one, let
\[
 M(r)=(1+r)^2(1+r/r_0)^2,
\]
\begin{equation}\label{ari:eq:cc-P}
 \begin{aligned}
 P(r)={}&(1-U_0)^2M(r)\\
 &-V_0^2\left\{c_0^2H_0^2+
 4\kappa r(c_0+\kappa r^2)(72/5)^2r^4(13/100)^2M(r)\right\}.
 \end{aligned}
\end{equation}
This polynomial has $P(0)=0$. Write
\[
 \frac{P(r_0x)}{r_0x}
     =\sum_{j=0}^{14}b_j\binom{14}{j}x^j(1-x)^{14-j}.
\]
All coefficients belong to $\Q(\sqrt3)$. Explicit lower bounds,
obtained by expanding the displayed factored polynomial, are
\[
 \begin{array}{c|rrrrrrrr}
 j&0&1&2&3&4&5&6&7\\ \hline
 100b_j>&169&175&181&187&192&197&202&205\\[2pt]
 j&8&9&10&11&12&13&14&\\ \hline
 100b_j>&207&207&204&197&184&160&118&.
 \end{array}
\]
For an exact independent recipe, if $P(r)/r=\sum q_nr^n$,
then
\[
 b_j=\sum_{n=0}^j q_nr_0^n
                  \frac{\binom jn}{\binom{14}n}.
\]
For example $b_0=4-4\sqrt3/3$ and
$b_1=25357/13125-453\sqrt3/4375$. The rational enclosure
\[
 \begin{gathered}
 1.73205080756887729352744634150587236694280525381038<\sqrt3,\\
 \sqrt3<1.73205080756887729352744634150587236694280525381039.
 \end{gathered}
\]
verifies every entry by integer multiplication; its endpoints
are checked by squaring. The complete exact coefficient list is
also retained with the code. Every Bernstein basis function is
nonnegative on $[0,1]$, so $P(r)>0$ for $0<r\le r_0$.
Since $U_0<1$ there, squaring \eqref{ari:eq:cc-scalar} is legitimate
and proves $|F|<1$. Formula \eqref{ari:eq:cc-Ffactored} applies
at $r=r_0$ except at $z=-z_0$ by continuity; at that remaining
point \ref{ari:C.8.12} gives $F=3/4$. This proves the assertion.
\end{enumerate}
\end{proof}

\begin{proposition}[every critical point and every level set]\label{ari:C.8.15}
\StatementRecord{ari:C.8.15}{proved}{Historical events: Claude Fable 5.2 (Anthropic), effort high, 15 November 2026: carried as background without fault found. Claude Fable 5.2 (Anthropic), effort high, 16 November 2026: used and rechecked in this run. GPT-6 Astra (OpenAI), ChatGPT Pro mode, effort Pro, S6\_09\_PAPERS, 17 November 2026: targeted proof check. Also checked in this writing. Pre-continuation writing events: GPT-6 Astra (OpenAI), ChatGPT Pro mode, effort Pro, S6\_10\_PAPERS, 18 November 2026: computation or formal certificate; targeted proof check. Continuation metadata events: GPT-6 Astra (OpenAI), ChatGPT Pro mode, effort Pro, S6\_10\_PAPERS continuation, 19 November 2026: used and rechecked in this run; status-vocabulary reconciliation; metadata reconciliation only; no new mathematical check. \emph{Pre-continuation recorded proof scope.} Read complete proof. Re-derived critical factor, radial derivative, cusp Fourier bound, Cayley ramification removal, rational-series tail bounds, scalar Bernstein implication and raywise Jordan-curve argument. The precise imported elliptic-function formula remains an external-locator gate until opened. Reran Ccritical\_rational\_certificate.py and Ccritical\_scalar.py. All 100 exact ODE coefficients and outward integer bounds were recomputed; all 15 exact Q(sqrt(3)) Bernstein coefficients are positive. Saved stdout and exact output certificates accompany the code.}{the hypotheses, notation and definitions stated here; \cref{ari:C.8.13}; \cref{ari:C.8.14}; \cref{ari:eq:cc-radial}; \cref{ari:eq:cc-EZ}; \cref{fam:src:C.2.3}}

For every zero $\mu$ of $\wp_\tau$, the zero sets of the two
critical factors are exactly
\[
 \begin{aligned}
 W=0&\quad\Longleftrightarrow\quad
       \tau\in\operatorname{SL}_2(\Z)\{i,\rho\},\\
 3W\wp'_\tau(\mu)-2g_2=0&\quad\Longleftrightarrow\quad
       \tau\in\operatorname{SL}_2(\Z)\rho.
 \end{aligned}
\]
On the coarse $t$-plane the function $L$ has no critical point
away from $t=1$. The point $t=0$ is a noncritical $C^1$ point,
with $L_t(0)=7/384$. The point $t=1$ is the unique global
maximum, of value $\log\sqrt2$, and has the fractional-power
local maximum described in C.3. In particular there is no
saddle in $0<t<1$. On the root charts, $h=0$ over $t=1$
is a nondegenerate maximum, while $\delta=0$ over $t=0$
is a degenerate critical point with nonzero harmonic cubic
leading term and three alternating ascending and descending
sectors. These are all critical points on the root charts.

For every angle $\theta$, $r\mapsto L(1+re^{i\theta})$ is
strictly decreasing for $r>0$, from $\log\sqrt2$ to $-\infty$.
For each $c<\log\sqrt2$, the level set $\{L=c\}$ is one
Jordan curve, strictly star-shaped about $1$. Its bounded
complement is $\{L>c\}$; its unbounded complement is
$\{L<c\}$. The curve is real analytic except possibly at
$t=0$ when $c=\frac14\log3$, where it is $C^1$ with nonzero
tangent. The closed superlevel sets form nested topological
discs. At the maximum the level set is the single point $1$,
and above it the level set is empty. The corresponding
statements for $D_{\rm part}$ follow by its increasing affine
change from $L$.
\end{proposition}


\begin{proof}
\begin{enumerate}[label=\arabic*.,leftmargin=*]
\item Take a representative in the closed standard modular fundamental
region. If $T\ge7/5$, \ref{ari:C.8.13} applies. Otherwise
$|\Re\tau|\le1/2$ and $\sqrt3/2\le T\le7/5$. For fixed
$\Re\tau$, the function
\[
 |\chi|^2=\frac{(\Re\tau)^2+(T-1)^2}
                    {(\Re\tau)^2+(T+1)^2}
\]
has its maximum on this interval at an endpoint and increases
with $(\Re\tau)^2$. The two endpoint bounds are $r_0^2$
and $41/601<r_0^2$. Thus \ref{ari:C.8.14} covers the entire remaining
region. We conclude $|1-X|<1$ for every $t\ne1$.
Equations \eqref{ari:eq:cc-radial} now prove strict radial
monotonicity, including $t=0$ by continuity. They also exclude
both critical equations $W=0$ and $3Wp-2g_2=0$ on the
ordinary base. The elliptic points must be checked on their
root charts: at $\rho$, the nonzero coarse derivative already
given rules out a critical point; at $i$, $W=p=0$ but $3Wp-2g_2=-2g_2\ne0$, and C.3 gives
the strict maximum. At $\rho$, $W=g_2=0$ and $p\ne0$,
so both displayed factors vanish. This proves their exact zero
sets, including the elliptic orbits. In particular the vanishing of both factors at
$\rho$ in the $\tau$ coordinate is a ramification effect,
not a saddle on the coarse base. More explicitly, on the root
chart at $\rho$,
\[
 L=L(0)+\frac7{192}\Re\left(\frac{ip_0^2}{\pi^3}\delta^3\right)
                   +O(|\delta|^4),
\]
with $p_0\ne0$, whereas on the chart at $i$,
$L=\log\sqrt2-\pi|h|^2/2+O(|h|^4)$.
These expansions give the asserted root-chart critical types.

\item At the cusp, the product formula and \eqref{ari:eq:cc-EZ} give,
uniformly in $\Re\tau$,
\[
 L=-\frac{\pi T}{6}+\log(2\sqrt3)-
                      \frac{\pi\beta_{\rm cusp}^2}{T}+O(e^{-2\pi T}),
 \qquad |t|=\frac{e^{2\pi T}}{1728}(1+o(1)).
\]
Hence $L\to-\infty$ uniformly as $|t|\to\infty$.
Continuity at $t=1$ supplies the other endpoint of each ray.
The intermediate value theorem and strict monotonicity give
one and only one radius $r_c(\theta)>0$ at each level $c$.
Properness and strict monotonicity imply continuity in $\theta$;
the nonvanishing radial derivative gives the stated local
regularity by the implicit-function theorem. Thus
$\theta\mapsto1+r_c(\theta)e^{i\theta}$ is a continuous
injective parametrization of a Jordan curve and has exactly
the asserted complementary regions.

\item The proposed route in the stated conventions that posits one saddle on
$(0,1)$ is therefore false. The corrected statement is the
strict radial foliation above. It supplies, in addition to
the separate duplication proof, a global proof of Theorem~\ref{fam:src:C.2.3} and
resolves the critical-point and global-level-set question
posed for this normalized norm function.
\end{enumerate}
\end{proof}



\subsection{An independent interval certificate}\label{push:section:E.4}
\label{sec:e4}

We write $\eta_D$ for the Dedekind eta function and set
\[
 h=z-(1-\tau)/2,\qquad
 \mathcal N^2(z,\tau)=
 \left|\frac{\vartheta_3(h,\tau)}{\eta_D(\tau)}\right|^2
 \exp\left(-\frac{2\pi(\Im h)^2}{\Im\tau}\right).
\]
The half-period theta identity identifies this with the norm used in
\cref{fam:src:C.1.1}, \cref{fam:src:C.2.3}. The coefficient of $\Im c$ in $D_c$ is one. Consequently the
identity from the period calculation above which we use is
\[
 D_{\rm part}=\frac6\pi
 \log\left(\frac{\mathcal N}{2\sqrt3}\right).
\]
It follows directly by taking the imaginary part of the theta formula
for $\beta^{\rm part}$; its Gaussian factor removes
$6(\Im\mu)^2/\Im\tau$. Thus no numerical normalization is inferred
from the new computation. The analytic formula and the first
certificate are recorded in the period calculation above, \cref{fam:src:C.1.1}, \cref{fam:src:C.2.3}. The result below supplies independent
arithmetic evidence for that same conclusion, with the independent certificate below.

\begin{theorem}[independent certification of every box]\label{push:src:E.4.1}
\StatementRecord{push:src:E.4.1}{proved}{Historical events: Claude Fable 5.2 (Anthropic), effort high, 16 November 2026: carried as background without fault found; used and rechecked in this run. Claude Fable 5.2 (Anthropic), effort high, 17 November 2026: carried as background without fault found. GPT-6 Astra (OpenAI), ChatGPT Pro mode, effort Pro, S6\_09\_PAPERS, 17 November 2026: targeted proof check. Also checked in this writing. Pre-continuation writing events: GPT-6 Astra (OpenAI), ChatGPT Pro mode, effort Pro, S6\_10\_PAPERS, 18 November 2026: computation or formal certificate; targeted proof check; used and rechecked in this run. Continuation metadata events: GPT-6 Astra (OpenAI), ChatGPT Pro mode, effort Pro, S6\_10\_PAPERS continuation, 19 November 2026: used and rechecked in this run; status-vocabulary reconciliation; metadata reconciliation only; no new mathematical check. \emph{Pre-continuation recorded proof scope.} The earlier attribution of code/I/replay\_interval\_leaves.output.json and replay\_family/f\_interval\_verify.stdout to this independent Arb theorem is withdrawn. Those executions used the earlier mpmath implementation and remain recorded under fam:src:C.2.3. The fresh unchanged Arb execution is separately recorded below. Read Johansson2016 Section2 pp2–4 and complete Section3 pp4–6, and official python-flint0.9.0 constructor/lower/upper/man\_exp/str semantics. These source-reading scopes are retained. The earlier reference to an already replayed Arb certificate was premature: that earlier execution used mpmath. The actual Arb implementation match and execution are recorded in the subsequent events. Read and rederived the full six-step proof and both complete original programs. Checked rational-to-Arb enclosures, exact dyadic Newton choices, the complex sup-norm contraction, finite pole-free Fourier enclosures, both moment tails, theta and logarithmic eta tails, upper-endpoint acceptance and the exact rational sweep. Wrote the explicit half-period theta identity and period/evenness transformations proving that the one-root norm enclosure controls both zeros. The successful contraction forces the selected preconditioner to be nonzero. The proof review is separate from the executable certificates. PASS: executed the original unchanged e4\_arb\_certificate.py with python-flint0.9.0,192bits,all19701boxes,four workers. All19701 strict root contractions and norm bounds succeeded,0failures. The largest upper endpoint is enclosed by1.999877571767410524070536478244687423578792301873142913 plus or minus9.26e-56,below1.999877571767410525; worst index17866. The input SHA256 is87b7c2dc672e98cab926c092ebdc757373a6817bb651c3472235a42956b5c9ab. Only rational endpoints from the earlier certificate are imported. Full outputs and identical source hashes are in code/I/replay\_arb and analysis/I\_arb\_replay\_manifest.json. PASS: executed the original unchanged e4\_cover\_and\_local.py. The exact Fraction sweep checks19911leaves,438vertical edges,area23/40,disjoint interiors and complete coverage;36local rectangles lie in the prescribed disc and174excluded rectangles lie below the unit circle. The separately implemented256bit Arb local-constant inequalities also all passed. Execution of those constants grants no additional whole-proof review of the older local argument.}{the hypotheses, notation and definitions stated here; \citet[Sections 2--3]{Johansson2016}; \citet{PythonFlint2026}}



For all 19701 rational rectangles in the supplied certificate, and for
each $\tau$ in the rectangle, both zeros of $\wp_\tau$ satisfy
$\mathcal N^2<2$. The new largest certified upper bound is less than
\[
 1.999877571767410525.
\]
The 19701 rectangles, 36 local rectangles and 174 excluded rectangles
have disjoint interiors and cover exactly
$[0,1/2]\times[17/20,2]$. Every local rectangle is contained in
$|\tau-i|<3/20000$, and every excluded rectangle satisfies
$|\tau|<1$ throughout.
\end{theorem}



\begin{proof}

\begin{enumerate}[label=\arabic*.,leftmargin=*]
\item The programs \path{code/I/replay_arb/S6_02_PUSH/code/e4_arb_certificate.py} and
\path{code/I/replay_arb/S6_02_PUSH/code/e4_cover_and_local.py} use respectively 192-bit and 256-bit
Arb arithmetic through python-flint 0.9.0. The arithmetic implementation
is independent of mpmath. We import neither mpmath nor any routine
from the previous certificate. The only numerical data taken from that
certificate are its rational rectangle endpoints. We recompute roots,
preconditioners, radii, contraction bounds, function values and tails.
Arb's midpoint-radius representation and its outward error bounds are
described by \citet[Section~2, pp.~2--4, and Section~3, pp.~4--6]{Johansson2016}; the precise
constructor and endpoint semantics used here are documented in
\citet{PythonFlint2026}. The final comparison is made with an upper
endpoint, before conversion to a decimal string.

\item For a rectangle $Q$, let $z_0$ and $A\ne0$ be the exact dyadic complex
numbers selected by the midpoint Newton iteration. That iteration is
only a choice of constants. For $Z=z_0+[-r,r]+i[-r,r]$ we independently
enclose $\wp_\tau(z_0)$ and $\wp'_\tau(Z)$ for $\tau\in Q$.
We accept only if
\[
 \begin{split}
 &\left|-A\wp_\tau(z_0)
 +(1-A\wp'_\tau(Z))(Z-z_0)\right|_\infty<r,\\
 &|\Re(1-A\wp'_\tau(Z))|
       +|\Im(1-A\wp'_\tau(Z))|<1.
 \end{split}
\]
Here absolute values on enclosures mean their outward upper bounds.
The derivative enclosure and integration along a line segment show
that $z\mapsto z-A\wp_\tau(z)$ is a strict contraction of $Z$ into
itself, uniformly for $\tau\in Q$. Thus it has a genuine zero in $Z$.
The finite Fourier enclosures exclude the poles on $Z$, and the strict
Lipschitz inequality forces $A\ne0$. The two zeros of $\wp$ are
negatives modulo the period lattice, counted with multiplicity. The
exact half-period transformation, obtained by replacing the summation
index in the theta series, is
\[
 \vartheta_3\bigl(z-(1-\tau)/2,\tau\bigr)
   =i e^{-\pi i\tau/4-\pi i z}\vartheta_1(z,\tau).
\]
Writing $v=\Im z$ and $y=\Im\tau$, its modulus and Gaussian factor
cancel to give
\[
 \mathcal N(z,\tau)=
 \left|\frac{\vartheta_1(z,\tau)}{\eta_D(\tau)}\right|
 e^{-\pi v^2/y}.
\]
Now $\vartheta_1(-z)=-\vartheta_1(z)$,
$\vartheta_1(z+1)=-\vartheta_1(z)$ and
$\vartheta_1(z+\tau)=-e^{-\pi i\tau-2\pi i z}\vartheta_1(z)$.
In the last transformation the modulus contributes
$e^{\pi y+2\pi v}$, while the Gaussian contributes its reciprocal.
Thus the displayed norm is even and lattice-periodic, proving directly
that the one-root enclosure controls both zeros. No continuity or
success of Newton iteration is used in that conclusion.

\item We use the normally convergent cosine expansion
\[
 \wp_\tau(z)=\frac{\pi^2}{\sin^2\pi z}-\frac{\pi^2}{3}
 +8\pi^2\sum_{n\geq1}\frac{nq^n}{1-q^n}
                      (1-\cos 2\pi nz),\qquad q=e^{2\pi i\tau},
\]
and differentiate it in $z$ for $\wp'$. This identity follows by twice
differentiating the normally convergent Jacobi product and inserting
$E_2=1-24\sum_{n\geq1}nq^n/(1-q^n)$; this also verifies the
Fourier formula used in the period calculation above. The new program sums
32 terms directly in cosine form. Put $r_q\geq|q|$,
$H_z\geq|\Im z|$, and $R_q=r_q e^{2\pi H_z}<1$.
For $j=1,2$, the inequality $(N+k)^j\leq(N+1)^jk^j$ gives
\[
 \begin{aligned}
 U_1(a)&=\frac{(N+1)a^{N+1}}{(1-a)^2},\\
 U_2(a)&=\frac{(N+1)^2a^{N+1}(1+a)}{(1-a)^3},
 \qquad
 \sum_{n>N}n^j a^n\leq U_j(a).
 \end{aligned}
\]
We therefore add complex rectangles enclosing discs of radii
\[
 \frac{8\pi^2}{1-r_q}\bigl(U_1(r_q)+U_1(R_q)\bigr),
 \qquad \frac{16\pi^3}{1-r_q}U_2(R_q)
\]
to the two partial sums. These estimates follow from
$|\cos z|,|\sin z|\leq e^{|\Im z|}$ and
$|1-q^n|\geq1-r_q$. They differ from the moment formulas used in
the old implementation.

\item For $\vartheta_3$ we sum
$1+2\sum_{n=1}^{12}e^{\pi i n^2\tau}\cos(2\pi nh)$.
With $T_*\leq\Im\tau$ and $H_*\geq|\Im h|$, the remaining
absolute value is at most
\[
 \frac{2e^{-169\pi T_*+26\pi H_*}}
 {1-e^{-27\pi T_*+2\pi H_*}}.
\]
The ratio in this expression is tested to be less than one. For eta
we independently sum its logarithm,
\[
 \log\eta_D(\tau)=\frac{\pi i\tau}{12}
          +\sum_{n=1}^{24}\log(1-q^n)+\mathcal E,
 \qquad
 |\mathcal E|\leq
 \frac{r_q^{25}}{(1-r_q)(1-r_q^{25})}.
\]
The logarithms are holomorphic on $|q^n|<1$. The tail follows by
summing $|\log(1-z)|\leq|z|/(1-|z|)$.
We enclose $\mathcal N^2$ as
\[
 ((\Re\vartheta_3)^2+(\Im\vartheta_3)^2)
 \exp\left(-2\Re\log\eta_D-
                   \frac{2\pi(\Im h)^2}{\Im\tau}\right),
\]
and test its upper endpoint against the exact integer 2.

\item All 19701 tests succeeded. The full output records each new enclosure;
the largest is enclosed by
\[
 1.999877571767410524070536478244687423578792301873142913
       \ \pm 9.26\cdot10^{-56}.
\]
This is the enclosure of a certified upper bound, not an estimate of
the actual maximum. Its box has input index 17866, with zero-based
indexing. The file digest of the input is recorded in the output.

\item The cover check is exact rational arithmetic. At all 438 distinct
vertical edges we maintain the active closed vertical intervals. In
every open vertical slab they form an ordered partition of
$[17/20,2]$, each with multiplicity one. The last active set is empty;
all leaves lie in the root rectangle, and the sum of their exact areas
is $23/40$. The closed rectangles also cover their common edges.
The omitted and excluded conditions are checked separately using
rational endpoint comparisons. Finally,
$2(1/10000)^2<(3/20000)^2$ puts the entire omitted square inside the
analytic disc used below. This verifies both the inequalities and
the coverage, rather than just replaying the old list of successful
flags.
\end{enumerate}

\end{proof}

\begin{theorem}[quantitative local maximum]\label{push:src:E.4.2}
\StatementRecord{push:src:E.4.2}{proved}{Historical events: Claude Fable 5.2 (Anthropic), effort high, 16 November 2026: carried as background without fault found. GPT-6 Astra (OpenAI), ChatGPT Pro mode, effort Pro, S6\_09\_PAPERS, 17 November 2026: targeted proof check. Also checked in this writing. Pre-continuation writing events: GPT-6 Astra (OpenAI), ChatGPT Pro mode, effort Pro, S6\_10\_PAPERS, 18 November 2026: targeted proof check. Continuation metadata events: GPT-6 Astra (OpenAI), ChatGPT Pro mode, effort Pro, S6\_10\_PAPERS continuation, 19 November 2026: used and rechecked in this run; status-vocabulary reconciliation; metadata reconciliation only; no new mathematical check. \emph{Pre-continuation recorded proof scope.} Independently re-proved the quantitative local maximum through the explicit sigma-to-Jacobi norm identity and the checked global maximum; the small root-disc contains only the i representative. Earlier constants and interval proof remain in the compendium as a second proof.}{\cref{ari:C.1.1}, \cref{ari:C.8.15}}

If $|\tau-i|\leq3/20000$ and $\wp_\tau(z)=0$, then
$\mathcal N(z,\tau)\leq\sqrt2$, with equality precisely at
$\tau=i$. The explicit numerical second proof is retained in the compendium;
the global norm theorem gives the following short proof.
\end{theorem}

\begin{proof}
\begin{enumerate}[label=\arabic*.,leftmargin=*]
\item For $\Lambda=\Z+\Z\tau$, the normalization identities are
\[
 \Delta_\Lambda=(2\pi)^{12}\eta_D(\tau)^{24},\qquad
 \sigma_\Lambda(z)=
 \frac{\vartheta_1(z,\tau)}{2\pi\eta_D(\tau)^3}
                 e^{\eta_1z^2/2}.
\]
The real-linear quasiperiod formula in \ref{ari:C.1.1} then gives
\[
 |\Delta_\Lambda|^{1/12}|\sigma_\Lambda(z)|
 e^{-\Re(z\eta(z))/2}
 =\left|\frac{\vartheta_1(z,\tau)}{\eta_D(\tau)}\right|
            e^{-\pi(\Im z)^2/\Im\tau}
 =\mathcal N(z,\tau).
\]
For $\wp_\tau(z)=0$, Theorem~\ref{ari:C.8.15} gives the
global maximum $\log\sqrt2$ of its logarithm, exactly on the
orbit of $i$. In the disc $|\tau-i|\le3/20000$ the only such
representative is $i$: a modular image with $c^2+d^2\ge2$
has imaginary part at most $1/2$, and the remaining images
are integral translates of $i$. Exponentiating proves the
stated inequality and its equality locus.
\end{enumerate}
\end{proof}




\clearpage\appendix

\appendix
\section{Credited proofs of parameter equivalence and deformation results}
\subsection{The canonical divisor and the exact anticanonical ring}
\label{wb:sec:canonical-ring-recovery}

The analytic construction admits a further explicit calculation which does
not use its claimed identification with a sphere. Retain the actual proper
map $f:X\to B=\mathbb P^1$, the cusp $p_0$, the finite points $p_1,p_2$,
the reduced finite fibres $S_1,S_2$, and
$f^*(p_1)=3S_1$, $f^*(p_2)=4S_2$. All line-bundle isomorphisms below
refer to this constructed analytic space. No replacement period matrix or
new compactification is introduced.

\begin{theorem}[Canonical bundle and recovery of the original fibration]
\label{wb:thm:canonical-ring-recovery}
\StatementRecord{wb:thm:canonical-ring-recovery}{proved}{Historical events: GPT-6 Astra (OpenAI), ChatGPT Pro mode, effort Pro, S6\_09\_PAPERS, 17 November 2026: targeted proof check. Also checked in this writing. Pre-continuation writing events: GPT-6 Astra (OpenAI), ChatGPT Pro mode, effort Pro, S6\_10\_PAPERS, 18 November 2026: targeted proof check. Continuation metadata events: GPT-6 Astra (OpenAI), ChatGPT Pro mode, effort Pro, S6\_10\_PAPERS continuation, 19 November 2026: used and rechecked in this run; status-vocabulary reconciliation; metadata reconciliation only; no new mathematical check. \emph{Pre-continuation recorded proof scope.} Re-derived the full receiving proof: invariant local evaluation exponents0 and2; actual meromorphic pushforward floor(k/4); every graded basis, the weighted polynomial ring and the intrinsic degree-two pencil, including dependence on scalar trivializations.}{the hypotheses, notation and definitions stated here; \CompendiumRef{I}{wb:lem:hodge-line-audit}}

The evaluation map
\[
 f^*(f_*\omega_{X/B})\longrightarrow\omega_{X/B}
\]
has zero divisor exactly $2S_2$. Consequently
\[
 \omega_X\simeq f^*\mathcal O_B(-p_2)\otimes\mathcal O_X(2S_2)
       \simeq\mathcal O_X(-2S_2),\qquad
 \omega_X^{-2}\simeq f^*\mathcal O_B(p_2).
 \tag{CR1}
\]
For every integer $k$ the actual direct image of the vertical divisor is
\[
 f_*\mathcal O_X(kS_2)
       =\mathcal O_B\bigl(\lfloor k/4\rfloor p_2\bigr)
 \tag{CR2}
\]
as subsheaves of the meromorphic-function sheaf, using pullback to identify
their sections. In particular, for each integer $m\geq0$,
\[
 h^0(X,\omega_X^{-m})=\lfloor m/2\rfloor+1,
 \qquad h^0(X,\omega_X^m)=0\quad(m>0).
 \tag{CR3}
\]
There is a graded algebra isomorphism
\[
 \mathbb C[U,V]\xrightarrow{\sim}
       \bigoplus_{m\geq0}H^0(X,\omega_X^{-m}),
       \qquad\deg U=1,\quad\deg V=2.
 \tag{CR4}
\]
The degree-two complete linear system has no base points and its map is
the original $f$, followed by a specified projective-coordinate map of $B$.
The degree-one system consists of the single divisor $2S_2$.
\end{theorem}



\begin{proof}
\begin{enumerate}[label=\arabic*.,leftmargin=*]
\item Write $\mathfrak K=f_*\omega_{X/B}$. The map is the natural adjunction
evaluation, not an identification inferred from the ranks of its two
terms. On a smooth torus fibre the invariant top form is nowhere zero.
At the cusp, the original form
$(dx_1/x_1)\wedge(dx_2/x_2)\wedge dt_c$ is a nowhere-zero canonical
frame in each of the unimodular toric charts. Its relative version is
the local generator of $\mathfrak K$ used in
\CompendiumRef{I}{wb:lem:hodge-line-audit}, so evaluation has no zero there.

\item For a finite filling retain $t_j=s_j^{m_j}$,
$(m_1,m_2)=(3,4)$ and $(a_1,a_2)=(2,1)$. On its finite \'etale
total-space cover $q_j:J'_j\to N_j$, the actual direct-image generator
pulls back, by (HL.9.5c)--(HL.9.5f), to
\[
 s_j^{m_j-1-a_j}e^{\varphi_j(s_j)}
       (ds_j\wedge e)\otimes(dt_j)^{-1}.
 \tag{CR5}
\]
Here $(dt_j)^{-1}$ is the pulled-back inverse canonical frame on the
base, not division by $m_js_j^{m_j-1}ds_j$. The coefficient
$e^{\varphi_j}$ is a unit. Thus the zero divisor on the cover is
$(m_j-1-a_j)\{s_j=0\}$. Since $q_j$ is \'etale, faithful pullback
of Cartier divisors makes the downstairs divisor
$(m_j-1-a_j)S_j$. The two coefficients are respectively $0$ and $2$.
These neighbourhoods and the smooth locus cover $X$, proving the stated
global zero divisor with no omitted component.

\item The evaluation therefore identifies
\[
 \omega_{X/B}\simeq f^*\mathfrak K\otimes\mathcal O_X(2S_2).
\]
By \CompendiumRef{I}{wb:lem:hodge-line-audit}, $\mathfrak K\simeq\mathcal O_B(p_0)$.
In the original affine base coordinate $t$, with
$t(p_1)=0$, $t(p_2)=1$, $t(p_0)=\infty$, retain the explicit choices
\[
 h=\frac1{t-1},\qquad \eta_B=\frac{dt}{(t-1)^2}.
\]
Their divisors are $p_0-p_2$ and $-2p_2$, respectively: at infinity,
with $w=1/t$, they are $w/(1-w)$ and $-dw/(1-w)^2$.
Thus $\eta_B$ specifies $\omega_B\simeq\mathcal O_B(-2p_2)$.
Tensoring with $f^*\omega_B$ gives
\[
 \omega_X\simeq
 \mathcal O_X\bigl(f^*(p_0)-2f^*(p_2)+2S_2\bigr).
\]
Multiplication by $f^*h$, with its displayed divisor, identifies the
base factor $\mathcal O_B(p_0-2p_2)$ with $\mathcal O_B(-p_2)$.
Using the actual multiplicity $f^*(p_2)=4S_2$ proves (CR1).

\item To prove (CR2), a section of $\mathcal O_X(kS_2)$ is represented by
a meromorphic function on $X$ with the indicated vertical order bound.
Away from $p_2$ it is holomorphic. Its restriction to each connected
compact fibre there is constant; $f_*\mathcal O_X=\mathcal O_B$
therefore descends it uniquely to a holomorphic function on that base
open set. Near $p_2$, pull back to the original cover $J'_2$, including
its zero section. The descended function becomes a function $a(s_2)$,
and $s_2^k a(s_2)$ is holomorphic at zero (with the evident vanishing
interpretation when $k<0$). Thus $a$ is meromorphic and has a Laurent
series bounded below. Invariance under the actual quotient, whose base
coordinate is $s_2\mapsto\zeta_2s_2$, says
$a(\zeta_2s_2)=a(s_2)$. Only exponents divisible by four occur, so
$a(s_2)=b(s_2^4)$ for a meromorphic $b$ on the original base disc.
The order inequality is exactly
\[
 4\operatorname{ord}_{p_2}(b)+k\geq0
 \quad\Longleftrightarrow\quad
 \operatorname{ord}_{p_2}(b)\geq-\lfloor k/4\rfloor.
\]
Conversely every such $b$ pulls back to an allowed meromorphic function.
This proves equality in both directions, as well as compatibility on
overlaps. No assumption about the field of all meromorphic functions
on $X$ is used. Applying (CR2) with $k=2m$ and $k=-2m$ and using
$h^0(\mathbb P^1,\mathcal O(d))=\max\{d+1,0\}$ proves (CR3).

\item Here is the exact algebra map in (CR4). Fix the isomorphism
$\omega_X^{-1}\simeq\mathcal O_X(2S_2)$ obtained above. Let $u$ be
the section corresponding to the canonical divisor section of
$\mathcal O_X(2S_2)$, so $\operatorname{div}(u)=2S_2$.
Its square, under the induced degree-two identification in (CR1),
is $f^*\sigma_{p_2}$, where $\sigma_{p_2}$ is the canonical divisor
section of $\mathcal O_B(p_2)$. Choose a second section $\ell$ of this
line bundle with $\ell(p_2)\ne0$, retaining that choice, and let $v$
correspond to $f^*\ell$ under the same degree-two identification.
Define $U\mapsto u$, $V\mapsto v$.

\item For $m=2r$, the monomials $u^{2r-2j}v^j$, $0\leq j\leq r$,
correspond exactly to
$f^*(\sigma_{p_2}^{r-j}\ell^j)$, a basis of
$H^0(B,\mathcal O_B(rp_2))$. For $m=2r+1$, multiplication by $u$
identifies the same basis with the full space: in the meromorphic
description (CR2), both $4rS_2$ and $(4r+2)S_2$ allow precisely
the base pole orders at most $r$. Thus the displayed map is an
isomorphism in every degree, proving both surjectivity and absence of
polynomial relations. Replacing $\ell$ by $a\ell+b\sigma_{p_2}$,
$a\ne0$, changes the second generator by exactly
$v\mapsto av+bu^2$; no scalar is silently discarded.
Replacing the chosen isomorphism $\omega_X^{-1}\to\mathcal O_X(2S_2)$
by $c$ times itself changes the generators by
$(u,v)\mapsto(c^{-1}u,c^{-2}v)$, with $c\in\mathbb C^*$ retained.

\item The two degree-two generators have no common zero because
$\sigma_{p_2},\ell$ do not. Their map is
$x\mapsto[\sigma_{p_2}(f(x)):\ell(f(x))]$, which is exactly
$f$ followed by that chosen coordinate map $B\to\mathbb P^1$.
Explicitly, writing $\ell=A/(t-1)+B$ with $A\ne0$, it is
$t\mapsto[t-1:A+B(t-1)]$, including the values $[0:A]$ at $p_2$
and $[1:B]$ at $p_0$.
Degree one has basis $u$ and its divisor is $2S_2$, as claimed.
\end{enumerate}
\end{proof}

In the already specified additive Picard coordinate \eqref{inv:eq:C-factor12}, with
$L_{12}=f^*\mathcal O_B(1)$, the identities retain the original constants:
\[
 \mathcal O_X(S_1)=L_4,\qquad
 \mathcal O_X(S_2)=L_3,\qquad
 \omega_X=L_{-6},\qquad \omega_X^{-2}=L_{12}.
\]
This last expression uses the previously proved Picard identification;
the divisor, direct-image and ring proofs above do not depend on it.


\subsection{The retained period constant is an effective deformation parameter}
\label{wb:sec:retained-period-deformation}

Keep the functions, lattice marking, base coordinate and local fillings of
\CompendiumRef{I}{wb:sec:candidate-geometry}. In particular, fix the actual cusp-regular
solution $\beta^{\rm part}$, not a replacement with a cusp pole, and write
\[
 \beta_c=\beta^{\rm part}+c,\qquad
 \mathcal P=\{c\in\mathbb C:\operatorname{Im}c<-M\},
 \quad M=\max_{B\setminus\{p_0\}}D_{\rm part}.
 \tag{PD1}
\]
The maximum and the exact admissible domain were proved in
\CompendiumRef{I}{wb:prop:period-exact-admissible-constants}. The following calculation
uses the analytic construction, not its integral cohomology or sphere
recognition. Its purpose is to determine whether the freedom in $c$
actually changes the constructed complex manifold.

\begin{theorem}[Exact period deformation]
\label{wb:thm:retained-period-deformation}
\StatementRecord{wb:thm:retained-period-deformation}{proved}{Historical events: GPT-6 Astra (OpenAI), ChatGPT Pro mode, effort Pro, S6\_09\_PAPERS, 17 November 2026: targeted proof check. Also checked in this writing. Pre-continuation writing events: GPT-6 Astra (OpenAI), ChatGPT Pro mode, effort Pro, S6\_10\_PAPERS, 18 November 2026: targeted proof check. Continuation metadata events: GPT-6 Astra (OpenAI), ChatGPT Pro mode, effort Pro, S6\_10\_PAPERS continuation, 19 November 2026: used and rechecked in this run; status-vocabulary reconciliation; metadata reconciliation only; no new mathematical check. \emph{Pre-continuation recorded proof scope.} Read and re-derived the full parameter-retained construction, uniform cusp properness, exact real lattice morphism, plus-graph/connecting-class sign conversion, nonzero torus class, relative canonical-pencil base change over dual numbers, scheme-theoretic constancy of all three critical values, and resulting nonzero Kodaira-Spencer class of the full threefold.}{the hypotheses, notation and definitions stated here; \CompendiumRef{I}{wb:lem:middle-period-proper}; \CompendiumRef{I}{wb:lem:toric-chart-control}; \CompendiumRef{I}{wb:thm:candidate-cusp-geometry}; \cref{wb:thm:canonical-ring-recovery}; \cite[Section~67, p.~15]{McMullen1996Complex}}

For every open disc $Q$ with $\overline Q\subset\mathcal P$, the retained
construction gives a proper holomorphic submersion
\[
 \pi:\mathcal X_Q\longrightarrow Q,
 \qquad F:\mathcal X_Q\longrightarrow B\times Q,
 \qquad \pi=\operatorname{pr}_Q\circ F,
 \tag{PD2}
\]
whose fibre over $c$ is the analytically constructed threefold $X_c$ and
whose restriction there is the original $f_c:X_c\to B$.
The three distinguished values remain the original
$t(p_1)=0$, $t(p_2)=1$, $t(p_0)=\infty$.

Fix a regular base point and one of its original lifts $z$, and fix
$c_*\in Q$. With all functions below evaluated at that $z$, set
\[
 T=\operatorname{Im}\tau>0,\quad m=\operatorname{Im}\mu,\quad
 D=\operatorname{Im}\beta_{c_*}-\frac{6m^2}{T}<0,\quad a=\frac mT,
 \qquad P=\begin{pmatrix}0&0\\-a&1\end{pmatrix}.
 \tag{PD3}
\]
For $\delta=c-c_*$, the exact marked real-linear map from the universal
cover of the fibre at $c_*$ to that at $c$ is
\[
 w=A_\delta\zeta+B_\delta\overline\zeta,\qquad
 A_\delta=I+\frac{\delta}{2iD}P,\qquad
 B_\delta=-\frac{\delta}{2iD}P.
 \tag{PD4}
\]
It maps the original period lattice onto the new period lattice, and
therefore descends to a real-analytic marked diffeomorphism of the two
tori. In the explicit convention that the transported $(0,1)$ tangent
space is spanned by
$\partial_{\overline\zeta_k}+\sum_j\nu^j_k\partial_{\zeta_j}$,
its exact coefficient matrix and its first derivative are
\[
 \nu_\delta=\frac{\delta}{2iD+\delta}P,\qquad
 \dot\nu_z=\frac{1}{2iD}\,
 \partial_{\zeta_2}\otimes
 (d\overline\zeta_2-a\,d\overline\zeta_1).
 \tag{PD5}
\]
The class of $\dot\nu_z$ in
$H^{0,1}(f_{c_*}^{-1}(t(z)),T^{1,0})$ is nonzero.
With the tangent-exact-sequence convention
$\operatorname{KS}(\partial_c)=[\bar\partial h]$ for a smooth
$(1,0)$ lift $h$ of $\partial_c$, the actual fibre class is
$\operatorname{KS}_z(\partial_c)=-[\dot\nu_z]$.
Moreover the first-order deformation of the entire $X_{c_*}$ obtained
from (PD2) is nontrivial: its Kodaira--Spencer class in
$H^1(X_{c_*},T_{X_{c_*}})$ is nonzero. In particular
$h^1(X_{c_*},T_{X_{c_*}})\geq1$.
\end{theorem}



\begin{proof}
\begin{enumerate}[label=\arabic*.,leftmargin=*]
\item \emph{Construction with the parameter retained.}
On an ordinary lifted base open set use the original period map
\[
 \Pi(z,c)=
 \begin{pmatrix}
 6\mu(z)&\tau(z)&1&0\\
 \beta^{\rm part}(z)+c&\mu(z)&0&1
 \end{pmatrix}:
 \Lambda\otimes\mathbb R\longrightarrow\mathbb C^2.
 \tag{PD6}
\]
Its real determinant is $\operatorname{Im}\tau\,D_c$, which never
vanishes on $Q$. Thus $(z,c,u)\mapsto[z,c,\Pi(z,c)u]$ identifies the
quotient, as a real-analytic family, with the original marked real torus
over this base open set times $Q$. The lattice action is holomorphic in
$(z,c,\zeta)$; it is free and properly discontinuous locally uniformly
on compact base-parameter sets, since the inverse real period matrices
are uniformly bounded there. The quotient is therefore a complex
manifold. Projection to $Q$ is a submersion because it is so before
the covering quotient. Its map to the lifted base times $Q$ is proper:
over any compact subset the image of that subset times a fixed closed
lattice parallelotope is the full inverse image.

\item The deck matrices are the original integral $M_g$ and the original
matrices $R_g(z)$, with
$\Pi(gz,c)=R_g(z)\Pi(z,c)M_g$. The $R_g$ do not depend on $c$.
Consequently the ordinary base quotient is jointly holomorphic.
Proper discontinuity there follows from proper discontinuity on the
base, exactly as in \CompendiumRef{I}{wb:lem:middle-period-proper}, with compact
parameter sets carried along.

\item At a finite filling retain $m_1=3$, $m_2=4$,
$s_j(g_jz)=e^{-2\pi i/m_j}s_j(z)$, and the actual logarithmic lift
\[
 (z,c,\zeta)\longmapsto
 \left(g_jz,c,R_{g_j}(z)\zeta+
                  \Pi(g_jz,c)\frac{v_j}{m_j}\right).
 \tag{PD7}
\]
In the real marking this is $(z,c,u)\mapsto
(g_jz,c,A_ju+v_j/m_j)$, independent of $c$.
Every nonidentity power remains free: the invariant integral covector
$\gamma$ takes the translation of its $k$-th power to $k/3$ or $-k/4$,
which is not an integer for $0<k<m_j$.
The finite quotient is therefore a local biholomorphism and is a
submersion to $Q$. The same compact-parallelotope argument proves
properness over $D_j\times Q$. Its map is still exactly
$t_j=s_j^{m_j}$. On the punctured overlap the translating section is
\[
 \frac{\log s_j(z)}{2\pi i}\Pi(z,c)v_j.
 \tag{PD8}
\]
A change of logarithm adds an actual period. It is therefore a
branch-independent holomorphic torus section, jointly in $z,c$.
The calculation (2.9e) applies with $c$ retained and proves the
punctured-overlap conjugacy, not merely equality of the monodromy orders.

\item For the cusp retain the same infinite unimodular fan and the full matrix
\[
 C(t_c,c)=
 \begin{pmatrix}6\mu(t_c)&h(t_c)\\
 b^{\rm part}(t_c)-h(t_c)+c&\mu(t_c)
 \end{pmatrix},\qquad
 B_0=\begin{pmatrix}0&1\\-1&0\end{pmatrix}.
 \tag{PD9}
\]
Choose a closed cusp disc of radius $r_1$ on which these functions are
holomorphic. Put
\[
 K_Q=\max_{|t_c|\le r_1,\ c\in\overline Q}
          \|-2\pi\operatorname{Im}C(t_c,c)\|,
 \qquad 0<\epsilon<\min(r_1,1),\quad
          |\log\epsilon|>2K_Q.
 \tag{PD10}
\]
The action on $Y_\epsilon\times Q$ is the original one,
\[
 (x,t_c,c)\longmapsto
 \left(e^{2\pi iC(t_c,c)\lambda}
                  t_c^{B_0\lambda}x,t_c,c\right),
 \quad\lambda\in\mathbb Z^2.
 \tag{PD11}
\]
It is holomorphic across the boundary by the same toric automorphism
and nowhere-zero torus multiplier. The estimates in (2.8) now hold
uniformly on $Q$:
\[
 B_{t_c,c}=B_0+
       \frac{-2\pi\operatorname{Im}C(t_c,c)}{\log|t_c|},\quad
 \|B_{t_c,c}-B_0\|<\tfrac12,\quad
 \|B_{t_c,c}\lambda\|>\tfrac12\|\lambda\|\quad(\lambda\ne0).
 \tag{PD12}
\]
They exclude stabilizers on the dense torus and bound the finite set
of translations that can join each pair of bounded charts, uniformly
in $c$. At $t_c=0$ a nonzero translation moves the bounded cell, so
there is no stabilizer there either. A compact set is contained in finitely
many of the open charts $\mathcal U_\sigma(2,\epsilon)\times Q$ from
\CompendiumRef{I}{wb:lem:toric-chart-control}.
The chart-pair estimate proves proper discontinuity on the whole
product, including sequences tending to its boundary.

\item For clarity properness is also uniform, not only pointwise in $c$.
For $\epsilon_1<\epsilon$, let $K_0$ be the finite union of closed
unit polydiscs associated to triangles meeting the closed radius
$\sqrt2$ ball, intersected with $|t_c|\le\epsilon_1$, exactly as in
\CompendiumRef{I}{wb:thm:candidate-cusp-geometry}. For every $c$ and $t_c\ne0$,
choose $\lambda$ so that
$\|B_{t_c,c}^{-1}y+\lambda\|_\infty\le1/2$.
The translated logarithmic position has norm at most
$3/(2\sqrt2)<\sqrt2$, and hence lies in a chart defining $K_0$.
At $t_c=0$, the six closed bidiscs covering $E_0$ give the same
representatives after the exact translation by $B_0v$, with multiplier
$\exp(2\pi iC(0,c)B_0v)$ retained. Therefore, for compact $K\subset Q$,
the image of the compact $K_0\times K$ is the full inverse image of
$\{|t_c|\le\epsilon_1\}\times K$. Every compact set in the cusp
base times $Q$ is closed in a set of this form. This proves properness.
The quotient is a submersion to $Q$ because its covering charts are
open subsets of $Y_\epsilon\times Q$.

\item Exponentiating the original logarithmic torus coordinates gives the
jointly holomorphic cusp overlap, with period identity
$Z=sB_0+C(t_c,c)$. Glue it and (PD8) to the middle family.
The special base discs are disjoint. Each resulting inverse image of
a member of this base cover is its entire Hausdorff local piece.
Two points with different $(t,c)$ values are separated on the base;
two with the same value lie in one Hausdorff open piece and can be
separated there. This proves Hausdorffness. Local complex charts give
a complex fourfold and the submersion $\pi$.
Properness of $F$ follows by covering any compact set of $B\times Q$
by finitely many relatively compact sets inside these base pieces and
using the just proved local properness. Since $B$ is compact, projection
$B\times Q\to Q$ is proper, proving (PD2). The fibres and overlaps
are precisely those in the original analytic construction.

\item \emph{The exact variation on a regular fibre.}
Write $u=(u_1,u_2,u_3,u_4)^t$ in the original ordered lattice basis and
$\zeta=\Pi(z,c_*)u$. Taking imaginary parts gives the two exact equations
\[
 \operatorname{Im}\zeta_1=6m u_1+T u_2,\qquad
 \operatorname{Im}\zeta_2=(D+6m^2/T)u_1+m u_2.
 \tag{PD13}
\]
Subtracting $a$ times the first equation from the second yields
\[
 u_1=\frac{\operatorname{Im}\zeta_2-a\operatorname{Im}\zeta_1}{D}.
 \tag{PD14}
\]
In the changed periods the same $u$ maps to
$w=\Pi(z,c)u=\zeta+\delta e_2u_1$. Substitution of
$\operatorname{Im}\zeta=(\zeta-\overline\zeta)/(2i)$ proves (PD4).
This also proves its asserted lattice map and real invertibility:
it is exactly $\Pi(z,c)\Pi(z,c_*)^{-1}$ as a real-linear map.
Its real determinant is $(D+\operatorname{Im}\delta)/D>0$.
Neither the nonholomorphic part nor the original marking is dropped.

\item Here $P^2=P$, and
\[
 \det A_\delta=1+\frac\delta{2iD},\qquad
 A_\delta^{-1}=I-\frac{\delta}{2iD+\delta}P.
 \tag{PD15}
\]
The denominator can vanish only for $\delta=-2iD$. At that parameter
$D+\operatorname{Im}\delta=-D>0$, which is outside the admissible
domain. Thus this inverse is defined at every $c\in Q$.
Annihilating $w$ by the declared new $(0,1)$ fields gives
$A_\delta\nu_\delta+B_\delta=0$, which proves (PD5), including
its sign. All its coefficients are constant in the fibre coordinates.
Thus its first derivative is a globally defined $\bar\partial$-closed
tangent-valued $(0,1)$ form on the torus.

\item It cannot be $\bar\partial V$ for a globally defined smooth vector
field $V$. Indeed the invariant holomorphic frame
$\partial_{\zeta_1},\partial_{\zeta_2}$ writes $V$ with periodic
smooth coefficients. The integral of each of their derivatives
$\partial_{\overline\zeta_k}V^j$ against the translation-invariant
probability volume is zero: real constant vector fields preserve this
volume, and integration of their derivatives on the compact torus is
zero. But the coefficient of
$\partial_{\zeta_2}\otimes d\overline\zeta_2$ in $\dot\nu_z$ is
the nonzero constant $1/(2iD)$. Its integral is that same constant.
This contradiction proves nonzero Dolbeault class without a dimension
count or an appeal to numerical period variation.

\item The conversion to the tangent-exact-sequence convention retains a
minus sign. In the moving coordinates the real-marking lift has
$(1,0)$ part $h=\partial_c+u_1\partial_{w_2}$ at $c_*$.
Since (PD14) gives
\[
 \bar\partial u_1=-\frac{d\overline\zeta_2-a\,d\overline\zeta_1}{2iD},
\]
the connecting class $[\bar\partial h]$ is exactly $-[\dot\nu_z]$.
Both conventions therefore give a nonzero class; they are not identified
without this sign conversion.

\item \emph{The intrinsic pencil in the parameter family.}
Let $\mathcal S_2\subset\mathcal X_Q$ be the reduced finite divisor
over $p_2\times Q$, so $F^*(p_2\times Q)=4\mathcal S_2$.
The proof of \cref{wb:thm:canonical-ring-recovery} works for relative
forms over $Q$, as can be checked on each original chart.
On the middle family the frame is $d\zeta_1\wedge d\zeta_2$ relative
to $B\times Q$. At the cusp the relative canonical frame over $Q$ is
$(dx_1/x_1)\wedge(dx_2/x_2)\wedge dt_c$.
Extra $dc$ terms in transition differentials vanish in relative forms.
The finite determinant characters $\chi_j$, their correcting units
$e^{\varphi_j(s_j)}$, and the modular function $\Theta$ used in
(HL.9.4)--(HL.9.6) are independent of $c$, since $R_g,\tau,\mu$ are.
Thus their coefficient identities hold jointly over $Q$.
The evaluation exponents remain $m_j-1-a_j=0,2$.
It follows by the actual same transition functions and divisors that
\[
 \omega_{\mathcal X_Q/Q}\simeq\mathcal O_{\mathcal X_Q}(-2\mathcal S_2),
 \qquad
 \omega_{\mathcal X_Q/Q}^{-2}\simeq
          F^*\operatorname{pr}_B^*\mathcal O_B(p_2).
 \tag{PD16}
\]
Here the base identifications retain $h=1/(t-1)$ and
$\eta_B=dt/(t-1)^2$ from (CR1); these functions also do not depend on $c$.
In particular (PD16) is a relative line-bundle isomorphism, not a
collection of separately chosen isomorphisms of its fibres.

\item We need its complete degree-two system after an infinitesimal base
change, and prove this explicitly. Put
$A=\mathbb C[\epsilon]/(\epsilon^2)$, base-change by
$c=c_*+\epsilon$, write $X_A$ for the resulting deformation and
$L_A=\omega_{X_A/A}^{-2}$. For $L_*=\omega_{X_{c_*}}^{-2}$,
flatness of the deformation and local freeness of $L_A$ give
\[
 0\longrightarrow L_*\xrightarrow{\ \epsilon\ }L_A
        \longrightarrow L_*\longrightarrow0.
 \tag{PD17}
\]
The two sections of $\mathcal O_B(p_2)$ used in (CR4), namely
$\sigma_{p_2}$ and $\ell=A_0/(t-1)+B_0'$ with $A_0\ne0$,
pull back by (PD16) to sections $s_0,s_1$ of $L_A$ whose reductions
are a basis of $H^0(X_{c_*},L_*)$. This basis assertion is exactly
the proved degree-two case of (CR2)--(CR4).
Any section of $L_A$ reduces to a complex linear combination of these
two reductions; subtract that lifted combination. Exactness of (PD17)
then writes the remainder as $\epsilon$ times a section of $L_*$,
again a combination of that basis. Conversely, if an $A$-linear
combination of $s_0,s_1$ is zero, reduce modulo $\epsilon$ and then
use the injective first arrow of (PD17) to see that all four complex
coefficients vanish. Therefore
\[
 H^0(X_A,L_A)=A s_0\oplus A s_1.
 \tag{PD18}
\]
The system is basepoint-free and its morphism is $f_A$ followed by
the fixed coordinate map $t\mapsto[t-1:A_0+B_0'(t-1)]$.
This proves the required base-change statement without presuming that
pointwise dimensions alone imply it.

\item \emph{Nontriviality for the complete threefold.}
Suppose its first-order deformation were trivial, with an isomorphism
$X_A\simeq X_{c_*}\times\operatorname{Spec}A$ reducing to the identity.
Intrinsic relative canonical bundles and their complete spaces of
sections are preserved by that isomorphism. By (PD18), it would
therefore identify $f_A$ with $a_A\circ f_{c_*}$ for a projective
linear change $a_A\in\operatorname{PGL}_2(A)$ reducing to the identity.

\item This change fixes all three distinguished points even to first order.
Here is a scheme-sensitive check. In the finite covering charts
$f_A$ is $t_j=s_j^{m_j}$, so its relative critical ideal is
$(s_j^{m_j-1})$, and its restriction to that critical scheme has
$t_j=0$. The map from the parameter ring $A$ to this critical
scheme is injective; it contains a constant-parameter section given
by $s_j=0$ and the zero section of the covering torus.
At the cusp, in a chart containing a triple point,
$t_c=z_0z_1z_2$ and the critical ideal is
$(z_1z_2,z_0z_2,z_0z_1)$. Again $t_c=0$ in its ring and $A$ injects,
using $z_0=z_1=z_2=0$. At a double chart the same assertion follows
from $t_c=z_0z_1$; a smooth chart has no critical point.
These critical ideals are intrinsic to the differential of the map
over $A$, and the charts are local biholomorphisms onto $X_A$.
Away from the three distinguished base values the map is a
submersion, so it has no critical point.
Consequently the scheme-theoretic critical-value sections are exactly
the three constant sections $p_1,p_2,p_0$. An automorphism reducing
to the identity cannot permute these three distinct closed points.
Its value at each is therefore the same point over $A$, not a
nonzero infinitesimal displacement: such a displacement would be a
nonzero element of $\epsilon\mathbb C$ in an injective copy of $A$
on the just displayed critical chart.

\item For completeness write $a_A(t)=t+\epsilon(b_0+b_1t+b_2t^2)$.
Fixing $0$ gives $b_0=0$, fixing $1$ gives $b_0+b_1+b_2=0$,
and fixing infinity, in $u=1/t$, gives $b_2=0$.
Hence $a_A$ is the identity. The hypothesized trivialization thus
preserves $f$, and restricting it to any regular base value would
trivialize the torus deformation (PD4). Its nonzero class (PD5)
excludes this.

\item Finally, the zero Kodaira--Spencer class is precisely the condition for
such a first-order trivialization. One may see the needed direction
directly: choose local product charts for the holomorphic submersion
$\pi$ and differentiate their transition functions at $c_*$. The
resulting tangent-valued overlap functions form a Cech $1$-cocycle.
If its class vanishes, after refinement it is the coboundary of local
holomorphic vector fields. Changing each chart by that vector field
times $\epsilon$ cancels the first-order overlap terms, producing
the identity-reducing trivialization. This is the usual
Kodaira--Spencer connecting class (for the standard deformation-space
interpretation see \cite[Section~67, p.~15]{McMullen1996Complex});
the explicitly recorded conversion
from the plus-graph convention in (PD5) changes its sign, not its
vanishing. The contradiction proves the final assertion.
\end{enumerate}
\end{proof}

The theorem establishes a genuinely varying analytic family with its
original marked fibration and its intrinsic degree-two canonical pencil.
It does not assert that every two distinct parameters give nonisomorphic
unmarked manifolds, or that this family exhausts their deformation space.
The global integral topology and the CDP argument are not inputs to this
deformation calculation.


The global identifications within this retained family are computed in
\cref{wb:pshift:thm:classification}: they are precisely $c'-c\in2\mathbb Z$,
including the finite-filling obstruction and every cusp overlap.


\subsection{Exact global identifications of the retained period parameter}
\label{wb:pshift:sec:classification}

Keep the original base $B$, coordinate $t$, values
$t(p_1)=0,t(p_2)=1,t(p_0)=\infty$, and the original functions
$\tau,\mu,\beta^{\mathrm{part}}$. Write
\[
 \beta_c=\beta^{\mathrm{part}}+c,\qquad
 \mathcal P=\{c\in\mathbb C:\operatorname{Im}c<-M\},\qquad
 \Pi_c(z)=\begin{pmatrix}6\mu(z)&\tau(z)&1&0\\
                      \beta_c(z)&\mu(z)&0&1\end{pmatrix}.
 \tag{PS1}\label{wb:pshift:eq:periods}
\]
All assertions below concern the analytic pieces and their actual gluing.
No integral cohomology, sphere recognition, or CDP assertion is an input.
In particular, the global parameter classification is not inferred from
an equality between two fibre lattices.

Put $e=(0,1)^{\mathsf t}\in\mathbb C^2$ and retain the ordered lattice
\[
 \Lambda=\mathbb Z\langle\widehat\gamma,\widehat u,
                        \widehat w,\widehat\delta\rangle,
 \qquad E=\widehat\delta\otimes\gamma.
 \tag{PS2}\label{wb:pshift:eq:rank-one-map}
\]
Thus $E$ is the matrix with its only nonzero entry, equal to one, in
row four and column one; $E^2=0$ and $\Pi_cE=e\gamma$.
The finite data, with every sign retained, are
\[
\begin{gathered}
 A_1=\begin{pmatrix}1&0&0&0\\6&0&1&0\\-6&-1&-1&0\\-2&1&0&1\end{pmatrix},
 \qquad
 A_2=\begin{pmatrix}1&0&0&0\\0&0&-1&0\\-6&1&0&0\\3&0&1&1\end{pmatrix},\\
 v_1=(1,2,-4,0)^{\mathsf t},\quad
 v_2=(-1,-3,3,0)^{\mathsf t},\quad
 (m_1,m_2)=(3,4),\quad
 (\epsilon_1,\epsilon_2)=(1,-1),\quad \epsilon_j=\gamma(v_j).
\end{gathered}
\tag{PS3}\label{wb:pshift:eq:finite-data}
\]
The covariance is $R_g\Pi_c=\Pi_c(gz)A_g$, with
\[
 R_{g_1}=\begin{pmatrix}-1/\tau&0\\(1-\mu)/\tau&1\end{pmatrix},\qquad
 R_{g_2}=\begin{pmatrix}1/\tau&0\\-\mu/\tau&1\end{pmatrix},\qquad
 R_{g_0}=I_2.
 \tag{PS4}\label{wb:pshift:eq:covariance}
\]
In particular $R_ge=e$ for every $g$.

We give the full period-shift proof of Theorem~\ref{fam:st:A.5.1}.


\subsubsection{The middle family and its complete rational centralizer}

For an integer $n$, put $S_n=I+nE$. Since $S_n^{-1}=S_{-n}$,
\[
 \Pi_{c+n}=\Pi_c S_n,\qquad
 \Pi_{c+n}\Lambda=\Pi_c\Lambda.
 \tag{PS6}\label{wb:pshift:eq:integer-shear}
\]
Both original matrices fix $\widehat\delta$ and preserve $\gamma$.
Consequently $A_jE=EA_j=E$, and the same commutation holds for every
monodromy word. The identity of $\mathbb C^2$ therefore identifies
the two full middle torus families with their actual $R_g$-action.
Its source-to-target lattice map is $S_{-n}$, since
$\Pi_{c+n}S_{-n}=\Pi_c$.

\begin{lemma}[All middle-family linear parts]
\label{wb:pshift:lem:centralizer}
\StatementRecord{wb:pshift:lem:centralizer}{proved}{Historical events: GPT-6 Astra (OpenAI), ChatGPT Pro mode, effort Pro, S6\_09\_PAPERS, 17 November 2026: targeted proof check. Also checked in this writing. Pre-continuation writing events: GPT-6 Astra (OpenAI), ChatGPT Pro mode, effort Pro, S6\_10\_PAPERS, 18 November 2026: targeted proof check. Continuation metadata events: GPT-6 Astra (OpenAI), ChatGPT Pro mode, effort Pro, S6\_10\_PAPERS continuation, 19 November 2026: used and rechecked in this run; status-vocabulary reconciliation; metadata reconciliation only; no new mathematical check. \emph{Pre-continuation recorded proof scope.} Re-derived the entire simultaneous rational centralizer using both fixed vector and fixed covector planes and the two quotient matrices; solved the extension equations, imposed integrality and determinant, and compared all period columns to obtain the integral shift condition.}{the hypotheses, notation and definitions stated here; the explicit constructions and calculations in the proof}

The simultaneous rational centralizer of $A_1,A_2$ is exactly
\[
 \{aI+bE:a,b\in\mathbb Q\}.
 \tag{PS7}\label{wb:pshift:eq:centralizer}
\]
A biholomorphism between the middle families for $c,c'$ over the
identity of $B^\circ$ forces $c'-c\in\mathbb Z$. Its complex-linear
part on every fibre is either $I_2$ or $-I_2$.
\end{lemma}



\begin{proof}
\begin{enumerate}[label=\arabic*.,leftmargin=*]
\item Solving the two fixed-vector systems from (PS3) gives
\[
 (\Lambda\otimes\mathbb Q)^{A_1}
   =\mathbb Q(1,2,-4,0)^{\mathsf t}+\mathbb Q\widehat\delta,
 \quad
 (\Lambda\otimes\mathbb Q)^{A_2}
   =\mathbb Q(1,3,-3,0)^{\mathsf t}+\mathbb Q\widehat\delta.
\]
Their intersection is $\mathbb Q\widehat\delta$. The fixed covectors
are respectively
\[
 \mathbb Q\gamma+\mathbb Q(2u+w+3\delta),\qquad
 \mathbb Q\gamma+\mathbb Q(u+w+2\delta).
\]
Their intersection is $\mathbb Q\gamma$: comparing the $u$ and $w$
coefficients of a common multiple gives $2x=y$ and $x=y$, hence
$x=y=0$. Every commuting endomorphism therefore preserves
$\mathbb Q\widehat\delta$ and $\ker\gamma$.

\item On $\ker\gamma/\mathbb Q\widehat\delta$, in the classes of
$(\widehat u,\widehat w)$, the two matrices are
\[
 Q_1=\begin{pmatrix}0&1\\-1&-1\end{pmatrix},\qquad
 Q_2=\begin{pmatrix}0&-1\\1&0\end{pmatrix}.
\]
A matrix commuting with $Q_2$ has the form
$\left(\begin{smallmatrix}b&d\\-d&b\end{smallmatrix}\right)$.
Commutation with $Q_1$ gives $d=0$. Thus a commuting endomorphism $H$
has columns
\[
 H\widehat\gamma=(a,r,s,k)^{\mathsf t},\quad
 H\widehat u=b\widehat u+u_0\widehat\delta,\quad
 H\widehat w=b\widehat w+w_0\widehat\delta,\quad
 H\widehat\delta=d_0\widehat\delta.
\]
Use the actual columns $A_1\widehat u=-\widehat w+\widehat\delta$,
$A_1\widehat w=\widehat u-\widehat w$, and
$A_2\widehat u=\widehat w$. Commutation on these columns gives
\[
 d_0-w_0=b+u_0,\qquad u_0=2w_0,\qquad w_0=u_0.
\]
Hence $u_0=w_0=0$ and $d_0=b$. Commutation on
$\widehat\gamma$ with $A_1$ gives
\[
 r=2(a-b),\qquad s=-4(a-b).
\]
Commutation on that column with $A_2$ also gives $r=-s$.
Therefore $a=b$, $r=s=0$, and $H=aI+kE$. Conversely these matrices
commute, by the fixed-vector and fixed-covector identities preceding
(PS6). This proves (PS7).

\item For the geometric assertion, a holomorphic map between two compact
complex tori lifts to their universal covers. The derivative of the
lift is periodic, so its entries descend to holomorphic functions on
the compact source torus and are constant. The lift is consequently
affine, with a complex-linear part $L$. In a holomorphic family these
linear parts vary holomorphically. Their lattice maps are integral
and locally constant, hence give an isomorphism of the two marked
local systems. Because the map is over the identity of the base,
that integral map $H$ commutes with both $A_j$. Thus
$H=aI+bE$, with $a,b\in\mathbb Z$. Its determinant is $a^4$;
invertibility over $\mathbb Z$ gives $a=1$ or $a=-1$.

\item The complete period equation is
\[
                   L(z)\Pi_c(z)=\Pi_{c'}(z)(aI+bE).
 \tag{PS8}\label{wb:pshift:eq:period-morphism}
\]
Its third and fourth columns give $L(z)=aI_2$. Its first column
then gives $a(c'-c)+b=0$. Therefore $c'-c=-b/a\in\mathbb Z$,
and the asserted two linear parts exhaust the possibilities.
\end{enumerate}
\end{proof}

\subsubsection{The exact finite obstruction, including normalization}

The two cyclic norm matrices are, by multiplication of (PS3),
\[
 N_1=I+A_1+A_1^2=
 \begin{pmatrix}3&0&0&0\\6&0&0&0\\-12&0&0&0\\0&2&1&3\end{pmatrix},
\]
\[
 N_2=I+A_2+A_2^2+A_2^3=
 \begin{pmatrix}4&0&0&0\\12&0&0&0\\-12&0&0&0\\0&2&2&4\end{pmatrix}.
 \tag{PS9}\label{wb:pshift:eq:norms}
\]
For example these formulas give
$N_1\widehat w=\widehat\delta$,
$N_2\widehat u=N_2\widehat w=2\widehat\delta$,
$N_1\widehat\delta=3\widehat\delta$, and
$N_2\widehat\delta=4\widehat\delta$.

We first justify why a local isomorphism must satisfy the exact affine
equations on the original finite cover. The actual ramified base change
$t_j=u^{m_j}$ has local equation
\[
 s_j^{m_j}-u^{m_j}
   =\prod_{a=0}^{m_j-1}(s_j-e^{2\pi ia/m_j}u)=0.
\]
Its normalization separates all these smooth branches, each with
coordinates $(z_1,z_2,u)$. The original map
$y\mapsto(q_j(y),s_j(y))$ identifies the original torus family $J'_j$
with this normalization, as proved by the explicit branches in
\CompendiumRef{I}{wb:ffcert:thm:normalization}. A biholomorphism over the identity of
$D_j$ base-changes over the identity of the $u$-disc and lifts uniquely
to the normalization: on local rings it carries integral elements in
the total fraction ring to integral elements, and its inverse does
the same. This lift commutes with the original deck action because
both composites lift the same map $(x,u)\mapsto(F(x),\zeta_j u)$.
In particular one may not replace the generator by an arbitrary
coprime power.

Let $c'=c+n$, where $n$ is an integer, and suppose the linear part of
the lifted map is $aI_2$, $a\in\{1,-1\}$. The lifted map on the torus
family is $[z,\zeta]\mapsto[z,a\zeta+B_j(z)]$. Its translating
section has a holomorphic lift $B_j:\Delta_j\to\mathbb C^2$:
$\Delta_j\times\mathbb C^2\to J'_j$ is a covering map and the disc is
simply connected. Local inverses of this covering show that the lift
is holomorphic. The deck-equivariance identity holds modulo an
actual lattice vector, so, after choosing this lift, it is
\[
 B_j(g_jz)-R_{g_j}(z)B_j(z)
  =\Pi_c(g_jz)\left(
       \frac{(1-a)v_j+n\epsilon_j\widehat\delta}{m_j}
          +\ell_j\right),\qquad \ell_j\in\Lambda.
 \tag{PS10}\label{wb:pshift:eq:finite-equation}
\]
The integer column $\ell_j$ is constant: it is locally constant by
discreteness of the marked lattice and the disc is connected.
Iterating $m_j$ times, with $A_jv_j=v_j$ and
$A_j\widehat\delta=\widehat\delta$, yields the necessary equality
\[
 N_j\ell_j=-(1-a)v_j-n\epsilon_j\widehat\delta.
 \tag{PS11}\label{wb:pshift:eq:finite-norm-obstruction}
\]
Here the left side of the iterated equation is zero because the
base point and the original vector coordinate return after $m_j$
steps; real invertibility of $\Pi_c$ makes the lattice equality exact.

If $a=-1$, applying $\gamma$ to (PS11) gives
\[
                   m_j\gamma(\ell_j)=-2\epsilon_j.
\]
This is impossible already at $j=1$, since $3$ does not divide $2$.
Thus a global isomorphism cannot have negative linear part.

For $a=1$, the first row of (PS11) gives $\gamma(\ell_j)=0$.
Equation (PS9) then proves that the order-three filling imposes no
restriction on $n$, whereas the order-four filling imposes precisely
$n\in2\mathbb Z$. More intrinsically, the finite obstruction is the
class of $-n\epsilon_j\widehat\delta$ in
\[
 \Lambda^{A_j}/N_j\Lambda.
 \tag{PS12}\label{wb:pshift:eq:finite-class}
\]
In the original fixed bases this quotient is
$\mathbb Z/3$ at $j=1$, with $\widehat\delta$ zero, and
$\mathbb Z/4\oplus\mathbb Z/2$ at $j=2$, with
$\widehat\delta$ generating the second summand. Indeed (PS9) gives
\[
 N_1\Lambda=3\mathbb Z v_1\oplus\mathbb Z\widehat\delta,
 \quad
 N_2\Lambda=4\mathbb Z\varepsilon'\oplus
                      2\mathbb Z\widehat\delta,
 \qquad \varepsilon'=(1,3,-3,0)^{\mathsf t}.
\]
This proves the exact parity obstruction rather than merely failure
of the untranslated identity.

\subsubsection{Local maps for every even shift}

Let $n\in2\mathbb Z$. Define the following original rational columns,
integral cocycles, and holomorphic translation vectors:
\[
\begin{aligned}
 r_1&=-\frac n3(\widehat u+\widehat w),&
 \ell_1&=-n\widehat w,&
 b_1(z)&=\Pi_c(z)r_1=-\frac n3(\tau(z)+1,\mu(z))^{\mathsf t},\\
 r_2&= \frac n4(\widehat u+\widehat w),&
 \ell_2&=\frac n2\widehat u,&
 b_2(z)&=\Pi_c(z)r_2=\frac n4(\tau(z)+1,\mu(z))^{\mathsf t}.
\end{aligned}
\tag{PS13}\label{wb:pshift:eq:local-corrections}
\]
Direct multiplication of the complete matrices in (PS3) gives
\[
 (I-A_j)r_j=\frac{n\epsilon_j}{m_j}\widehat\delta+\ell_j.
 \tag{PS14}\label{wb:pshift:eq:local-matrix-identity}
\]
For $j=1$ the four entries are $(0,0,-n,n/3)$, and for $j=2$
they are $(0,n/2,0,-n/4)$. Pairing with the full $\Pi_c(g_jz)$
shows that $b_j$ satisfies (PS10) for $a=1$.
Translation by $b_j$ is therefore already an isomorphism of the
two finite local fillings. The remaining issue is its exact overlap
with the fixed middle identification.

Keep the original linearizing coordinate, radius, and phase
$s_j(g_jz)=e^{-2\pi i/m_j}s_j(z)$. On a logarithm branch put
\[
 w_j(z)=\frac{\log s_j(z)}{2\pi i},\qquad
 a_j(z)=b_j(z)+n\epsilon_j w_j(z)e.
 \tag{PS15}\label{wb:pshift:eq:finite-middle-sections}
\]
Changing the logarithm adds $n\epsilon_j\widehat\delta$ under the
period map, so $a_j$ is branch independent as a torus section.
For the specified negative generator, $w_j(g_jz)=w_j(z)-1/m_j$.
Equations (PS4) and (PS14) consequently give the exact equality
\[
 a_j(g_jz)=R_{g_j}(z)a_j(z)+\Pi_c(g_jz)\ell_j.
 \tag{PS16}\label{wb:pshift:eq:local-affine-cocycle}
\]
This is why the logarithmic term in (PS15) cannot be omitted.
In fact the original finite overlap maps are translations by
$\sigma_{j,c}=w_j\Pi_c v_j$, and
\[
                  \sigma_{j,c+n}-\sigma_{j,c}
                     =n\epsilon_jw_je.
 \tag{PS17}\label{wb:pshift:eq:actual-overlap-difference}
\]
Thus $a_j$ is exactly the middle translation induced by the local
translation $b_j$, with the original gluing maps retained.

At the cusp retain $s=\tau-h$, $e^{2\pi is}=t_c$, $s(g_0z)=s(z)-1$,
and the complete period matrix
\[
 Z_c=sB_0+C_c(t_c),\quad
 B_0=\begin{pmatrix}0&1\\-1&0\end{pmatrix},\quad
 C_c=\begin{pmatrix}6\mu&h\\ b^{\mathrm{part}}-h+c&\mu\end{pmatrix}.
 \tag{PS18}\label{wb:pshift:eq:cusp-periods}
\]
For $\lambda=(\lambda_1,\lambda_2)^{\mathsf t}\in\mathbb Z^2$,
$C_{c+n}\lambda-C_c\lambda=(0,n\lambda_1)^{\mathsf t}$. Therefore
\[
 \exp(2\pi iC_{c+n}\lambda)t_c^{B_0\lambda}
       =\exp(2\pi iC_c\lambda)t_c^{B_0\lambda}.
 \tag{PS19}\label{wb:pshift:eq:cusp-actions-identical}
\]
The entire cusp actions, including their nonconstant holomorphic units,
are identical for an integer shift.

The affine cocycle convention in (PS16) has multiplication law
\[
                    \ell_{gh}=\ell_g+A_g\ell_h.
 \tag{PS20}\label{wb:pshift:eq:cocycle-law}
\]
Since $g_0=(g_1g_2)^{-1}$, it gives
\[
\begin{aligned}
 \ell_1+A_1\ell_2
   &=-\frac{3n}{2}\widehat w+\frac n2\widehat\delta,\\
 \ell_0=-A_{g_0}(\ell_1+A_1\ell_2)
   &=\frac{3n}{2}\widehat w-\frac n2\widehat\delta.
\end{aligned}
\tag{PS21}\label{wb:pshift:eq:cusp-cocycle}
\]
Here $A_{g_0}$ fixes both toric lattice vectors, as is also immediate
from $A_{g_0}=\left(\begin{smallmatrix}I&0\\B_0&I\end{smallmatrix}\right)$.
Put
\[
 q_n=(-3n/2,n/2)^{\mathsf t}\in\mathbb Z^2,
                  \qquad a_0(z)=s(z)q_n.
 \tag{PS22}\label{wb:pshift:eq:cusp-translation}
\]
Then $a_0(g_0z)-a_0(z)=-q_n=\Pi_c(g_0z)\ell_0$, exactly as
required by (PS16). Exponentiating gives the map
$(x,t_c)\mapsto(t_c^{q_n}x,t_c)$ on the dense cusp torus.
It extends holomorphically and invertibly to the original infinite fan:
the integral map
\[
 (y,y_3)\longmapsto(y+y_3q_n,y_3)
\]
translates each height-one triangle by the integer vector $q_n$ and
therefore permutes the actual cones, with inverse given by $-q_n$.
It commutes with the cusp quotient action in (PS19). This proves
extension across every cusp stratum, not only across its dense torus.

\subsubsection{The global additive correction and all overlap identities}

We give the global correction explicitly as an additive gluing
construction and prove its solvability. Let $\widetilde B^\circ$ be
the universal cover of the three-punctured base. Its fundamental group
is free on the original two negative peripheral generators. Pull back
the original $\Pi_c,R$ to this cover. Assign $\ell_1,\ell_2$ from
(PS13) to its free generators and extend by (PS20). Freeness makes
this an actual integral cocycle; its third peripheral value is
precisely (PS21). The maps
\[
 (z,y)\longmapsto
       (gz,R_g(z)y+\Pi_c(gz)\ell_g)
 \tag{PS23}\label{wb:pshift:eq:affine-bundle-action}
\]
form an affine action. To check the action law, covariance gives
$R_g(hz)\Pi_c(hz)=\Pi_c(ghz)A_g$; substituting this equality
in the composition gives exactly (PS20).
The quotient is an affine holomorphic bundle over $B^\circ$, since
the deck action on the base is free and properly discontinuous and
one may use evenly covered base opens as affine charts.

On a finite punctured disc, (PS16) says that $a_j$ is a section of
this affine bundle. On the cusp punctured disc, (PS21)--(PS22) say
the same for $a_0$. Define a sheaf $\mathscr H_n$ on $B$ as follows.
Away from the special points it is the sheaf of sections of (PS23).
Near $p_j$, a section is required to have the form $a_j+\eta_j$,
where $\eta_j:\Delta_j\to\mathbb C^2$ is holomorphic and obeys
$\eta_j(g_jz)=R_{g_j}(z)\eta_j(z)$. Near the cusp it is required
to have the form $a_0+\eta_0$, with $\eta_0(t_c)$ holomorphic on
the full cusp disc. These are compatible definitions on the punctured
overlaps by the equal affine cocycles just proved. They also prove
local nonemptiness, including at all three special points.

The differences of these local sections form a vector sheaf
$\mathscr V$. Its ordinary sections have covariance $\eta(gz)=R_g\eta(z)$;
at the two finite points and the cusp its extensions are precisely
the holomorphic extensions specified above. Projection to the first
coordinate gives an exact sequence on the original compact base
\[
 0\longrightarrow\mathcal O_B e\longrightarrow\mathscr V
        \longrightarrow\mathcal L\longrightarrow0,
                  \qquad \mathcal L\simeq\mathcal O_B(-p_0).
 \tag{PS24}\label{wb:pshift:eq:vector-sheaf}
\]
Here $\mathcal L$ is exactly the homogeneous period sheaf from
\CompendiumRef{I}{wb:prop:period-input-audit}, not a newly chosen extension. Indeed
the first diagonal entries in (PS4) are its factors $-1/\tau$ and
$1/\tau$, and at the cusp both extensions require ordinary
holomorphicity. The original function
\[
 \Theta=\frac{(E_4\circ\tau)^2(E_6\circ\tau)^{1/2}}
                 {\Delta_{\mathrm{mod}}\circ\tau}
\]
has those exact factors, orders two and one in $s_1,s_2$, and a
simple pole in $t_c$. Dividing by $\Theta$ identifies the ordinary
and finite stalks with $\mathcal O_B$, and the cusp stalk with
$\mathfrak m_{p_0}$. This recovers the displayed $\mathcal L$
with its original cusp condition. These transformation and order
calculations are the ones used to construct the retained $\mu$;
no topological conclusion about $X_c$ is involved.

For completeness, exactness on the finite stalks in (PS24) retains
the off-diagonal entry of $R_g$. Given an equivariant first component
$u(z)$, begin with the vector $(u(z),0)^{\mathsf t}$ and average it
under the finite group action on vector-valued holomorphic functions.
The resulting vector is equivariant, and its first component is
still $u$. Thus projection is surjective there. Its kernel is
exactly the invariant second-component functions, which are
holomorphic functions of $s_j^{m_j}=t_j$, multiplied by $e$.
The ordinary and cusp stalks are immediate. This also proves that
$\mathscr V$ is locally free of rank two: lift a local frame of
$\mathcal L$ and join it to $e$.

Both $H^1(B,\mathcal O_B)$ and
$H^1(B,\mathcal O_B(-p_0))$ vanish. One can compute these two
groups without a degree approximation: on the two standard affine
charts of $\mathbb P^1$, split a holomorphic Laurent series on
their overlap into its nonnegative and negative powers. For
$\mathcal O_B$ these are the two coboundary contributions, and
for $\mathcal O_B(-p_0)$ its transition factor shifts the second
contribution by one, still leaving no missing exponent. The positive
part converges on the full first affine chart and the transformed
negative part on the full second affine chart. This proves the two
vanishings, equivalently the usual additive Cousin calculation on
these charts. The exact sequence (PS24) now gives
\[
                             H^1(B,\mathscr V)=0.
 \tag{PS25}\label{wb:pshift:eq:no-analytic-obstruction}
\]

Here is how to choose the correction entering the maps. Choose
local sections $h_\alpha$ of $\mathscr H_n$, taking the displayed
$a_j,a_0$ on sufficiently small special discs and affine charts
elsewhere. Their differences $d_{\alpha\beta}=h_\beta-h_\alpha$
are a Cech $1$-cocycle with values in $\mathscr V$. Project it to
$\mathcal L$ and solve that additive cocycle, using its just proved
$H^1$ vanishing. Lift the resulting local $\mathcal L$ sections to
$\mathscr V$ by the local lifting construction in (PS24). Subtracting
their coboundary from $d$ leaves a scalar $\mathcal O_Be$ cocycle.
Solve that scalar additive cocycle by its $H^1$ vanishing. This
produces sections $v_\alpha$ with
$d_{\alpha\beta}=v_\beta-v_\alpha$. The local sections
$h_\alpha-v_\alpha$ agree and define a global section
$a\in\mathscr H_n(B)$. This is a two-stage construction of the
required correction, not an unproved additional gluing condition.
It gives holomorphic vectors $\eta_j,\eta_0$ with
\[
 a=a_j+\eta_j\quad(j=1,2),\qquad a=a_0+\eta_0
                              \quad\hbox{at the cusp}.
 \tag{PS26}\label{wb:pshift:eq:global-correction}
\]
Any choices made in the additive solutions give maps with the same
asserted properties. Their differences as sections of this affine
bundle lie in $H^0(B,\mathscr V)=\mathbb C e$: this follows from
(PS24), $H^0(B,\mathcal L)=0$, and $H^0(B,\mathcal O_B)=\mathbb C$.

Projecting $a$ modulo $\Pi_c\Lambda$ gives a holomorphic section of
the middle torus family because every $\ell_g$ in (PS23) is integral.
The promised three kinds of local maps are therefore
\[
\begin{aligned}
 F_{\mathrm{mid}}[z,\zeta]&=[z,\zeta+a(z)],\\
 F_j[z,\zeta]&=[z,\zeta+b_j(z)+\eta_j(z)]\quad(j=1,2),\\
 F_0[x,t_c]&=[\exp(2\pi i\eta_0(t_c))t_c^{q_n}x,t_c].
\end{aligned}
\tag{PS27}\label{wb:pshift:eq:full-isomorphism}
\]
The first descends by (PS23) and equality of the integer-shift
lattices. The finite maps descend by (PS10), (PS14), and the exact
covariance of $\eta_j$. The cusp map descends by (PS19); its
holomorphic-unit factor extends by torus multiplication on every
original chart, and its integer factor extends by the fan map in
(PS22). Each map is a biholomorphism, with inverse given by the
negative translation or the inverse torus multiplier and fan map.

On a finite overlap, write the original gluing as
$h_{j,c}(\zeta)=\zeta+\sigma_{j,c}$. Then (PS17), (PS26), and
(PS27) give the equality of actual torus maps
\[
                     h_{j,c+n}F_j=F_{\mathrm{mid}}h_{j,c}.
 \tag{PS28}\label{wb:pshift:eq:finite-overlap-identity}
\]
Both sides add
$\sigma_{j,c}+a=\sigma_{j,c+n}+b_j+\eta_j$.
On the cusp overlap the original map is exponentiation of $\zeta$;
(PS22) and (PS26) give
$\exp(2\pi ia)=t_c^{q_n}\exp(2\pi i\eta_0)$, exactly the
third line of (PS27). Thus that overlap also agrees, including
all nonconstant period units. The maps and their inverses glue to
a global biholomorphism $X_c\to X_{c+n}$ over $B$.

The choices can be made coherently for every parameter and every even
shift. Both initial cocycle columns in (PS13) have $\gamma=0$.
Because every $A_g$ preserves $\gamma$, (PS20) shows that
$\gamma(\ell_g)=0$ for every word. Therefore the affine term
$\Pi_c(gz)\ell_g$ in (PS23) is independent of $c$. The displayed
$a_j,a_0$, and the vector sheaf $\mathscr V$ are also independent
of $c$. Perform the additive construction once for $n=2$, keeping
one resulting section $a^{(2)}$. For $n=2k$ choose
$a^{(2k)}=k a^{(2)}$ and multiply its local correcting vectors by
$k$. All their affine cocycles and local formulas are then exactly
those for $n=2k$. The resulting maps satisfy the full group law
\[
 F_{2h,c+2k}\circ F_{2k,c}=F_{2(h+k),c},\qquad
                 F_{0,c}=\mathrm{id}_{X_c}.
 \tag{PS28a}\label{wb:pshift:eq:coherent-shift-action}
\]
This follows by adding the actual translation vectors in the first
two lines of (PS27), and by multiplying their units and adding the
integer exponents in its third line. The formulas are jointly
holomorphic in $c$, since their corrections are independent of it.

\subsubsection{Intrinsic necessity and the resulting parameter space}

The analytically computed complete linear system
$|\omega_{X_c}^{-2}|$ is the original map $f_c$, by
\cref{wb:thm:canonical-ring-recovery}. A biholomorphism preserves its
canonical bundle, the complete space of sections, and its evaluation
map. Consequently it induces $\alpha\in\operatorname{PGL}_2(\mathbb C)$
with $f_{c'}F=\alpha f_c$. The three special fibres have distinct
intrinsic types: the cusp is reduced and nonnormal, while the other
two have multiplicities three and four. All remaining fibres are
smooth reduced tori. Thus $\alpha$ fixes $p_0,p_1,p_2$ individually.
In the original coordinate, a fractional-linear map fixing
$\infty,0,1$ has zero lower-left and upper-right entries and equal
diagonal entries; hence $\alpha=\mathrm{id}_B$.

Now \cref{wb:pshift:lem:centralizer} forces $c'-c=n\in\mathbb Z$
and linear part $\pm I_2$. Equations (PS10)--(PS11), obtained on
the original normalized finite covers, exclude $-I_2$ and then
exclude odd $n$ for $+I_2$. Conversely (PS27)--(PS28) construct
the global map for every even $n$. This proves
\cref{wb:pshift:thm:classification} without converting a failed
particular map into a nonisomorphism assertion.

As a complex parameter space for the isomorphism classes represented
by this retained family, the quotient $\mathcal P/(2\mathbb Z)$ has
the explicit coordinate
\[
 w=\exp(-\pi ic),\qquad 0<|w|<\exp(-\pi M).
 \tag{PS29}\label{wb:pshift:eq:parameter-disc}
\]
Indeed $|w|=\exp(\pi\operatorname{Im}c)$, and two complex numbers
have the same value of this exponential exactly when they differ
by an even integer. Every point in the displayed punctured disc
has a logarithm giving a $c\in\mathcal P$. The exponential has
nonzero derivative, so these local inverse logarithms give the
claimed quotient complex structure. The integer-period quotient
seen by the middle family instead has coordinate $w^2$; it is
covered two-to-one by (PS29). Its two lifts are distinguished by
the exact order-two summand in (PS12). This statement classifies
the parameters of the retained family; it does not assert that
the family exhausts all deformations or all complex threefolds
with the same singular-fibre types.



\subsection{Re-derivation of Theorem~\ref{fam:src:A.5.1}}\label{app:rederive:fam-src-A.5.1}
\begin{proof}
\ref{fam:src:A.1.1} makes the base intrinsic. \ref{fam:src:A.2.1} gives the necessary parity and excludes the negative linear part. \ref{fam:src:A.3.1} supplies the inverse maps, all collar identities, and descent of the jointly holomorphic family. \ref{fam:src:A.4.2} proves its infinitesimal effectiveness. Finally $|e^{-\pi ic}|=e^{\pi\operatorname{Im}c}$ and equality of exponentials is exactly difference in $2\mathbb Z$, so the coordinate has precisely the stated image and fibres.
\end{proof}

\subsection{Re-derivation of Theorem~\ref{fam:src:A.5.2}}\label{app:rederive:fam-src-A.5.2}
\begin{proof}
We give the calculation, including the singular-fibre maps. The
corresponding calculation appears in the version 1
\cite[Section 6, especially Proposition 6.10, Lemma 6.11,
Theorem 6.12 and Proposition 6.16]{DuEtAl2026v1}; those deformation
statements are omitted from version two. We use the verified maps and
cohomology calculation below, rather than infer their validity from the
earlier theorem label.

\emph{The base and its line bundles.}
Write $X=X_c$ and form the analytic root curve
\[
 \mathfrak B=\sqrt[3]{(B,p_1)}\mathbin{\times_B}\sqrt[4]{(B,p_2)},
 \qquad \rho:\mathfrak B\longrightarrow B.
\]
A root object is a line bundle with section whose indicated power is the
pullback of the section cutting out $p_j$. The actual equalities
$f^*p_1=3S_1$ and $f^*p_2=4S_2$ supply these objects as
$(\mathcal O_X(S_j),1_{S_j})$, and hence supply a morphism
$\bar f:X\to\mathfrak B$ with $\rho\bar f=f$.
The explicit weighted quotient
\[
 \mathfrak B\simeq[(\mathbb C^2\setminus\{0\})/\mathbb C^*],
 \qquad \lambda(x,y)=(\lambda^3x,\lambda^4y),
 \qquad t=\frac{y^3}{y^3-x^4},
\]
identifies the roots with $(\mathcal O(4),y)$ and
$(\mathcal O(3),x)$ and the ordinary point $p_0$ with
$y^3-x^4=0$. In particular $\rho^*\mathcal O_B(1)=\mathcal O(12)$.
This isomorphism follows on each root chart by the displayed root
sections; away from their zeros both sides have the same ordinary coarse
coordinate, so the chart isomorphisms glue.

For completeness the analytic line-bundle computation can be made on
the coarse curve. Removing the two stabilizer characters from any line
bundle makes it descend to $B$: on each root chart, averaging a local
frame whose central value is nonzero gives an invariant frame still
nonzero at the centre, and its invariant coefficients are holomorphic
in $s_j^{m_j}$. These frames give the descended line and its transition
functions. Consequently its Picard group has
generators $H,R_1,R_2$ and relations $3R_1=H=4R_2$.
Thus it is freely generated by $R_1-R_2$, denoted $\mathcal O(1)$;
$H=12\mathcal O(1)$, $R_1=4\mathcal O(1)$ and
$R_2=3\mathcal O(1)$. The chart descent just used is exact because
invariants of a finite group over $\mathbb C$ are an exact functor.
With the clockwise root generators, a local invariant section of
$\mathcal O(k)$ has vanishing exponents
\[
 r_3(k)\in\{0,1,2\},\quad r_3(k)\equiv k\pmod3,
 \qquad
 r_4(k)\in\{0,1,2,3\},\quad r_4(k)\equiv-k\pmod4.
\]
Dividing by those local powers gives the precise coarse direct image
\begin{equation}\label{fam:eq:ab-root-degree}
 \rho_*\mathcal O(k)=\mathcal O_B\bigl(d(k)\bigr),\qquad
 d(k)=\frac{k}{12}-\frac{r_3(k)}3-\frac{r_4(k)}4.
\end{equation}
Indeed $k-4r_3-3r_4=12d(k)$ identifies the line bundle with
$\rho^*\mathcal O_B(d(k))\otimes R_1^{r_3}\otimes R_2^{r_4}$.
The canonical root sections generate the invariant direct image of
the last two factors; division by their local powers is the displayed
line-bundle isomorphism on the coarse charts. Finite invariants also give
$R^q\rho_*=0$ for $q>0$, as seen on each Stein quotient chart.
For $k=0$ the degree is zero; for
$k=1,-1,-2,-5,-6$ it is $-1$. It follows directly from the
cohomology of $\mathbb P^1$ that these five line bundles are acyclic,
whereas $H^0(\mathcal O)=\mathbb C$ and $H^1(\mathcal O)=0$.
The local equation $dt_j=m_js_j^{m_j-1}ds_j$ gives
\[
 K_{\mathfrak B}=\mathcal O(-24+2\cdot4+3\cdot3)
 =\mathcal O(-7),\qquad
 T_{\mathfrak B}(-p_0)=\mathcal O(-5).
\]

\emph{The actual base changes and cusp cohomology.}
On a finite root disc $\Delta_{s_j}$ the local presentations are
$[J'_j/\mu_{m_j}]\to[\Delta_{s_j}/\mu_{m_j}]$.
The fibre-product identity
\[
 [J'_j/\mu_{m_j}]
 \mathbin{\times_{[\Delta_{s_j}/\mu_{m_j}]}}\Delta_{s_j}
 \simeq J'_j
\]
is induced by the quotient torsor and its equivariant map to the disc;
its inverse sends a torsor with a chosen root-chart trivialization to
the corresponding point of $J'_j$. The action on $J'_j$ is free by
\ref{fam:src:A.1.2}. Thus this is an etale atlas base change and is exactly the
smooth torus family, with its given deck generator. Over the cusp the
base change is the original cusp family. Locally it is
$q=z_0z_1z_2$; the homomorphism
$\mathbb C\{q\}\to\mathbb C\{z_0,z_1,z_2\}$ is flat because its
target is torsion-free over the discrete valuation ring.
These charts prove that $\bar f$ is representable, proper and flat,
with connected Gorenstein fibres of dimension two. Properness on the
cusp is the compact representative construction of \ref{fam:src:A.3.1}, and on the
finite charts it is the compact torus quotient. Also
$\bar f^*p_0=W$ as reduced Cartier divisors.

Let $\nu:E\to W$ be the normalization, $E=\mathrm{dP}_6$, let
$C_h,C_v,C_d$ be the three double-curve normalizations, and let
$P_+,P_-$ be the two triple points. Each $C_e$ is $\mathbb P^1$.
Write $\rho_e^0,\rho_e^1:C_e\to E$ for its two paired boundary maps
and $j_e=\nu\rho_e^0=\nu\rho_e^1$. The sheaf sequence is
\begin{equation}\label{fam:eq:ab-cusp-resolution}
 0\longrightarrow\mathcal O_W\longrightarrow\nu_*\mathcal O_E
 \xrightarrow{\delta^0}\bigoplus_{e=h,v,d}j_{e*}\mathcal O_{C_e}
 \xrightarrow{\delta^1}\mathbb C_{P_+}\oplus\mathbb C_{P_-}
 \longrightarrow0,
\end{equation}
where $\delta^0(s)_e=(\rho_e^1)^*s-(\rho_e^0)^*s$ and
$\delta^1$ evaluates $t_h-t_v+t_d$ at the two ordered triple points.
Its exactness is local: for $xyz=0$ it is the alternating restriction
complex of the three coordinate planes, their pairwise lines and their
common point. Decomposing convergent power series by the variables
occurring in each monomial gives the augmented simplex complex, which
is exact. The one- and two-branch cases are the same calculation with
fewer planes. This verifies exactness for the actual boundary maps,
including the torus multipliers in their gluing.

All positive-degree structure-sheaf cohomology of $E,C_e$ vanishes.
Consequently \eqref{fam:eq:ab-cusp-resolution} computes cohomology by
\[
 \mathbb C\xrightarrow{0}\mathbb C^3
 \xrightarrow{\left(\begin{smallmatrix}1&-1&1\\1&-1&1\end{smallmatrix}\right)}
 \mathbb C^2.
\]
It gives $(h^0,h^1,h^2)(\mathcal O_W)=(1,2,1)$.
We also need its multiplication. The branch-intersection resolution
is the augmented semisimplicial sheaf algebra with its restriction face
maps; its local exactness makes the augmentation a multiplicative
quasi-isomorphism. Its Alexander--Whitney product therefore computes
the sheaf-cohomology cup product. In the ordered triangles
$[0,e_1,e_2]$ and $[e_1,e_2,e_1+e_2]$ it is
\[
 \alpha\smile\beta=(\alpha_h\beta_d,\alpha_d\beta_h),
 \qquad [\alpha]\smile[\beta]
 =\alpha_h\beta_d-\alpha_d\beta_h
 \quad\hbox{in }\mathbb C^2/\mathbb C(1,1).
\]
The cocycles $(1,1,0)$ and $(0,1,1)$ pair to one. Thus
$\bigwedge^2H^1(\mathcal O_W)\to H^2(\mathcal O_W)$ is an
isomorphism, including its product rather than just its dimensions.

\emph{Vertical fields and the direct images.}
Set $\mathcal T=T_X(-\log W)$ and
$\mathcal V=\ker(d\bar f^{\log})$.
The vertical transitions on the core are precisely $R_g(z)$.
Finite collar translations have vertical differential the identity,
and at the cusp the frame $\partial_{\zeta_i}$ corresponds to
$2\pi i x_i\partial_{x_i}$. The cusp transformations
$x\mapsto u(q)q^{B_0\lambda}x$ fix the coefficients of these
relative logarithmic fields. These exact identifications construct
a bundle $\mathcal E$ on $\mathfrak B$ with
$\mathcal V=\bar f^*\mathcal E$.
The fixed vector $e=(0,1)^t$ is nonzero in every logarithmic frame,
so there is an exact sequence
\[
 0\longrightarrow\mathcal O\xrightarrow{e}\mathcal E
 \longrightarrow L\longrightarrow0.
\]
The factors of $L$ are $-1/\tau$ and $1/\tau$. Therefore the usual
weight-twelve modular discriminant makes $1/\Delta(\tau)$ a
meromorphic section of $L^{12}$. It has no finite zero or pole,
and its cusp pole is simple: $\Delta(\tau)=q\,u(q)$ with
$u(0)\ne0$. Hence $L^{12}=\mathcal O(-p_0)=\mathcal O(-12)$.
The Picard calculation gives $L=\mathcal O(-1)$, and the extension
splits since $H^1(\mathcal O(1))=0$. Thus
\[
 \mathcal E\simeq\mathcal O\oplus\mathcal O(-1).
\]
There is also the exact sequence
\begin{equation}\label{fam:eq:ab-log-tangent}
 0\longrightarrow\bar f^*\mathcal E\longrightarrow\mathcal T
 \longrightarrow\bar f^*\mathcal O(-5)\longrightarrow0.
\end{equation}
On the smooth root charts it is the relative tangent sequence. At the
cusp its last map sends $\sum a_i z_i\partial_{z_i}$ to
$(\sum a_i)q\partial_q$, so it is onto and its kernel is the
two height-zero lattice directions already used above. Taking its
determinant and using $W=\bar f^*p_0$ gives the actual relative
dualizing line
\begin{equation}\label{fam:eq:ab-relative-dualizing}
 \omega_{X/\mathfrak B}=\bar f^*(\det\mathcal E)^{-1}.
\end{equation}
In particular this computation gives
$K_X=\bar f^*\mathcal O(-6)=\mathcal O_X(-2S_2)$, with the last
identification supplied by $\bar f^*R_2=\mathcal O_X(S_2)$;
it agrees through this actual pullback with \ref{fam:src:A.1.1}.

The fibres of $\bar f$ all have structure-sheaf cohomology dimensions
$(1,2,1)$: on ordinary and finite root charts they are tori, and at
the cusp we have just computed them. Analytic proper flat base change
on these complex-space charts gives locally free $R^q\bar f_*\mathcal O_X$
with these fibres. Its canonical comparison maps agree on the atlas
double overlaps, so they descend to $\mathfrak B$. The unit gives
$R^0\bar f_*\mathcal O_X=\mathcal O$.
The source charts are smooth and the fibres are Cartier divisors,
so relative analytic duality applies in its Gorenstein form. The
relative trace pairing, compatible on the same overlaps, and
\eqref{fam:eq:ab-relative-dualizing} give
\[
 (R^2\bar f_*\mathcal O_X)^\vee
 =\bar f_*\omega_{X/\mathfrak B}
 =(\det\mathcal E)^{-1},
 \qquad R^2\bar f_*\mathcal O_X=\mathcal O(-1).
\]
Here local freeness removes the higher sheaf-Ext terms in the duality
spectral sequence. This is precisely the pointwise base-change and
top-degree duality application spelled out in the cited Lemmas 6.8--6.9.

Put $G=R^1\bar f_*\mathcal O_X$. Cup product and the verified fibre
comparisons show $\det G=R^2\bar f_*\mathcal O_X=\mathcal O(-1)$;
at $W$ this uses the nonzero pairing above.
We construct the required subline of $G$. On an ordinary fibre,
constant lattice cocycles identify
\[
 H^1(F_t,\mathcal O_{F_t})
 =\Lambda_{\mathbb C}^*/\operatorname{rowspan}_{\mathbb C}\Pi_c.
\]
Indeed a coboundary among constant cocycles is given by an entire
function with constant period differences. Its derivatives are
periodic holomorphic functions on the compact torus, hence constant;
the function is affine linear. This gives exactly the displayed
kernel, and both sides then have dimension two. The class of
$\gamma=(1,0,0,0)$ is nonzero, since evaluation on columns
$\widehat w,\widehat\delta$ would force both coefficients of a
putative period-row combination to vanish.
The matrices satisfy $\gamma A_j=\gamma$. Translations act trivially
on $H^1(\mathcal O)$, so this holomorphic class extends on both finite
root charts and respects their gluing.

At the cusp $\gamma$ kills
$\Lambda_{\mathrm{tor}}=\mathbb Z\langle\widehat w,\widehat\delta\rangle$.
For local sections of the free toric quotient with changes
$s_i=\Psi_{\lambda_{ij}}s_j$, the constants
$\gamma(\lambda_{ij})$ therefore form a holomorphic Cech cocycle
on the entire neighbourhood. With
$B_0=\left(\begin{smallmatrix}0&1\\-1&0\end{smallmatrix}\right)$,
the edges $e_1,e_2,e_2-e_1$ have deck vectors
$(0,1),(-1,0),(-1,-1)$. The central cocycle values are consequently
$(0,-1,-1)$, which pair to one with $(1,0,-1)$.
Its value on $W$ is nonzero. On the punctured overlap the exponential
map kills precisely $\Lambda_{\mathrm{tor}}$, so its full-lattice
cocycle is exactly $\lambda\mapsto\gamma(\lambda)$, proving equality
with the preceding class. Over the Stein cusp disc, the Leray edge
map $H^1(\bar f^{-1}D_0,\mathcal O_X)\to H^0(D_0,G)$ is an
isomorphism, since $H^1(D_0,\mathcal O)=H^2(D_0,\mathcal O)=0$.
Thus this is a holomorphic, nonvanishing extension of $[\gamma]$.
It supplies a global subline and exact sequence
\[
 0\longrightarrow\mathcal O\xrightarrow{[\gamma]}G
 \longrightarrow\mathcal O(-1)\longrightarrow0.
\]
Again $H^1(\mathcal O(1))=0$ splits it. We have proved
\begin{equation}\label{fam:eq:ab-direct-images}
 R^q\bar f_*\mathcal O_X=
 \begin{cases}
 \mathcal O,&q=0,\\
 \mathcal O\oplus\mathcal O(-1),&q=1,\\
 \mathcal O(-1),&q=2.
 \end{cases}
\end{equation}

\emph{From logarithmic to ordinary tangent cohomology.}
The projection formula gives, for $q=0,1,2$ respectively,
\[
 R^q\bar f_*\mathcal V:
 \quad\mathcal O\oplus\mathcal O(-1),\quad
 \mathcal O\oplus2\mathcal O(-1)\oplus\mathcal O(-2),\quad
 \mathcal O(-1)\oplus\mathcal O(-2).
\]
For $\bar f^*\mathcal O(-5)$ the three expressions are
$\mathcal O(-5)$, $\mathcal O(-5)\oplus\mathcal O(-6)$ and
$\mathcal O(-6)$. Formula \eqref{fam:eq:ab-root-degree} makes every
nontrivial summand in this list acyclic. Leray, whose base has
cohomological dimension one, therefore gives
$h^\bullet(\mathcal V)=(1,1,0,0)$ and zero cohomology for the
horizontal term. Sequence \eqref{fam:eq:ab-log-tangent} yields
$h^\bullet(\mathcal T)=(1,1,0,0)$.

The normalization map $\nu:E\to X$ is an immersion on each local
branch. Let $N_\nu=\operatorname{coker}(T_E\to\nu^*T_X)$.
The exact normal-derivative map is
\begin{equation}\label{fam:eq:ab-normal-quotient}
 0\longrightarrow\mathcal T\longrightarrow T_X
 \longrightarrow\nu_*N_\nu\longrightarrow0.
\end{equation}
For a local equation $z_0\cdots z_{r-1}=0$, it sends
$\sum b_i\partial_{z_i}$ to the $r$ independent normal coefficients
$b_i|_{z_i=0}$. These can be lifted independently, proving
surjectivity; the kernel is exactly the logarithmic tangent module.
This also specifies the global map under changes of branch charts.
The cusp three-form
\[
 \Omega=d\log x_1\wedge d\log x_2\wedge dq
\]
is invariant under $x\mapsto u(q)q^{B_0\lambda}x$ because the
extra logarithmic differentials are multiples of $dq$. In any
unimodular maximal cone the height-one equation is $q=z_0z_1z_2$,
and the determinant-one character change gives
$\Omega=\pm dz_0\wedge dz_1\wedge dz_2$.
It therefore trivializes $K_X$ on the entire cusp neighbourhood,
including every branch intersection. Taking determinants of the
immersion tangent sequence now gives the explicit line identification
$N_\nu=K_E\otimes(\nu^*K_X)^{-1}\simeq K_E$.
Since $E$ is the rational projective surface $\mathrm{dP}_6$,
Serre duality gives $h^\bullet(E,K_E)=(0,0,1)$.
The cohomology sequence of \eqref{fam:eq:ab-normal-quotient} proves
$h^\bullet(X,T_X)=(1,1,1,0)$, as asserted.

\emph{The reduced deformation germ.}
\ref{fam:src:A.4.2} gives a nonzero Kodaira--Spencer vector at this same $c$;
the just computed $h^1=1$ makes its map an isomorphism.
Embed the actual Kuranishi germ in its one-dimensional tangent space
and write its local analytic algebra as $\mathbb C\{t\}/I$.
The existing proper family has a local classifying map, whose pullback
sends $t$ to $g(s)$ with $g'(0)\ne0$. The holomorphic inverse-function
theorem makes $\mathbb C\{t\}\to\mathbb C\{s\}$, $t\mapsto g(s)$,
an isomorphism. Every element of $I$ maps to zero, so $I=0$.
This proves smoothness and reducedness even though $h^2(T_X)=1$.
The classifying map is consequently a local isomorphism of germs;
completeness of the Kuranishi deformation then gives the claimed
miniversality and exhaustion of sufficiently small deformations.
\end{proof}

\subsection{Re-derivation of Proposition~\ref{fam:src:B.4.1}}\label{app:rederive:fam-src-B.4.1}
\begin{proof}
\ref{fam:src:A.1.1} and \ref{fam:src:A.2.1} give base identity, $n=0$ and fibre-linear part $I$. Let $\ell_1,\ell_2$ be the integral defects of a lift of its translating section to the universal cover of $B^\circ$. Finite extension gives $N_j\ell_j=0$. The saturated image calculations in \ref{fam:src:A.1.2} give the integral equality $\ker N_j=(I-A_j)\Lambda$. The cusp valuation argument (B2.3) gives $\ell_0\in\Lambda_{\rm tor}$. Write
\[
 \ell_1=ae_1+bf_1,\qquad \ell_2=ce_2+df_2.
\]
The $\widehat u$ coordinate of $\ell_1+A_1\ell_2$ is $a-c-2d$, so the cusp condition is exactly $a=c+2d$. Changing the global lift by a period $\lambda$ changes this pair by the simultaneous coboundary $((I-A_1)\lambda,(I-A_2)\lambda)$. In the four displayed coefficient coordinates its matrix is
\[
 C=\begin{pmatrix}-6&1&-1&0\\2&-1&0&0\\0&1&1&0\\-3&0&-1&0\end{pmatrix}.
\]
The lattice of all allowed pairs has basis matrix
$K=\left(\begin{smallmatrix}0&1&2\\1&0&0\\0&1&0\\0&0&1\end{smallmatrix}\right)$. The first three columns of $C$ expressed in this basis have determinant one. Hence every allowed pair is an integral coboundary; the computation is exact, including its saturation, and is independently checked by Hermite form in \texttt{ab\_exact.py}. After changing the lift, all its defects vanish.

The corrected lift extends as an equivariant holomorphic vector on the finite discs. A holomorphic local lift exists on the whole root disc by lifting the translating section through its covering vector bundle. Its defect is in $\ker N_j=(I-A_j)\Lambda$; subtracting an integral period kills that defect while preserving holomorphy. The corrected global lift and this local defect-zero lift differ by an invariant period, which is holomorphic through the centre. At the cusp, the local toric logarithm in (B2.3) is $sq+\eta(t_c)$. The chosen global lift differs from it by $\Pi_c(s)\lambda=sB_0\bar\lambda+C(t_c)\bar\lambda+\lambda_{\rm tor}$. Defect zero gives exactly $q+B_0\bar\lambda=0$, so the logarithmic terms cancel and the remaining vector is holomorphic. This proves the extension at the cusp as well. We therefore obtain a global section of $\mathcal V$. Its exact sequence in \ref{fam:src:A.3.1} gives $H^0(B,\mathcal V)=\mathbb Ce$. Conversely each such section defines the displayed translations on the finite covers and torus multiplication at the cusp, hence a global automorphism. Its kernel consists of global invariant periods. The common fixed lattice is $\mathbb Zd$, and $\Pi_cd=e$, giving kernel exactly $\mathbb Z$. This proves the full group assertion and its connectedness.
\end{proof}


\section{Tangent cohomology and deformations}

We retain the standard filling data $(\ell _0,\ell _1,\ell _2)=(0,1,-1)$
and an arbitrary admissible $c$. All sheaves in this section are analytic,
and all cohomology groups have coefficients in $\C$. The integral twists
below mean the actual line bundles
\[
 M_n=\OO_X\bigl(n(S_1-S_2)\bigr),\qquad n\in\Z.
\]
We set $\mathfrak g=g_\gamma$. Thus, in the positively normalized
exponential coordinate of \ref{inv:C3.2},
$M_n=\exp(n\mathfrak g)$. The computations use the root base to retain the
nonreduced fibres. Passing to its coarse space is a separate, exact
invariant operation; it is never identified with taking the dual of an
invariant lattice.

\begin{proposition}[The root base and its explicit cohomology]\label{inv:D1.1}
\StatementRecord{inv:D1.1}{proved}{Historical events: GPT-6 Astra (OpenAI), ChatGPT Pro mode, effort Pro, S6\_09\_PAPERS, 17 November 2026: targeted proof check. Also checked in this writing. Pre-continuation writing events: GPT-6 Astra (OpenAI), ChatGPT Pro mode, effort Pro, S6\_10\_PAPERS, 18 November 2026: targeted proof check. Continuation metadata events: GPT-6 Astra (OpenAI), ChatGPT Pro mode, effort Pro, S6\_10\_PAPERS continuation, 19 November 2026: used and rechecked in this run; status-vocabulary reconciliation; metadata reconciliation only; no new mathematical check. \emph{Pre-continuation recorded proof scope.} Read full root-base, logarithmic tangent, direct-image and twisted-Dolbeault proofs. Re-derived coarse degree function, local normalization quotient, resonance equation, Fourier contraction comparison, joint order-three coefficient 123ik/(64pi), stated local modules, and jump-locus finiteness. Exact Koszul and twist-program computations retain their separately recorded replay scopes.}{the hypotheses, notation and definitions stated here; the explicit constructions and calculations in the proof}

Let $\rho:\mathcal B\to B$ be the analytic root curve with orders $3$ and
$4$ at $p_1,p_2$. Write $R_1,R_2$ for its tautological root line bundles.
Its Picard group is infinite cyclic, with generator
$R=R_1\otimes R_2^{-1}$, and
\[
 R_1=R^4,\quad R_2=R^3,\quad \rho^*\OO_B(1)=R^{12},
 \quad K_{\mathcal B}=R^{-7}.
\]
We abbreviate $R^n$ to $\OO_{\mathcal B}(n)$. Choose sections
$x\in H^0(R^3)$ and $y\in H^0(R^4)$ vanishing at $p_2$ and $p_1$,
respectively. For every integer $n$, put
\[
 r_1(n)\equiv n\pmod3,\quad 0\le r_1(n)<3,
 \qquad r_2(n)\equiv3n\pmod4,\quad0\le r_2(n)<4,
\]
\[
 d(n)=\frac{n-4r_1(n)-3r_2(n)}{12}.
\]
Then
\[
 \rho_*\OO_{\mathcal B}(n)=\OO_B(d(n)),\qquad
 H^i(\mathcal B,\OO(n))=H^i(B,\OO(d(n))).                 \tag{D.1}
\]
Equivalently, $h^0(\OO(n))$ is the number of pairs
$(a,b)\in\Z_{\ge0}^2$ with $3a+4b=n$, and
$h^1(\OO(n))=h^0(\OO(-n-7))$.
\end{proposition}



The root-curve Picard group and graded cohomology are
\citet[Lemma~6.1]{DuEtAl2026v1}. We re-derive them and retain the explicit
coarse invariant-descent function $d(n)$ in the same calculation.

\begin{proof}
\begin{enumerate}[label=\arabic*.,leftmargin=*]
\item On the root charts the group acts on a disc coordinate by a primitive
$m_j$th root of unity. The stabilizer character of any line bundle is
removed uniquely by tensoring with $R_1^{-r_1}R_2^{-r_2}$, with the
exponents in the displayed ranges. A line bundle with trivial stabilizer
characters descends: an equivariant nonzero local frame can be averaged,
and its invariant frame gives the descent transition function. It is
therefore the pullback of a unique line bundle on $B$. Consequently
\[
 \operatorname{Pic}(\mathcal B)=
 \langle H,R_1,R_2\mid 3R_1=H,\;4R_2=H\rangle=\Z R.
\]
Here and below additive notation for this presentation means tensor
product. The determinant of the coordinate change $t=s^m$ gives
$K_{\mathcal B}=\rho^*K_B\otimes R_1^2R_2^3=R^{-24+8+9}$.

\item An invariant section of $R_j^{r_j}$ on its chart is $s^{r_j}$ times a
holomorphic function of $t=s^{m_j}$. It follows, by the local invariant
comparison on overlaps, that
$\rho_*(\rho^*\OO(k)R_1^{r_1}R_2^{r_2})=\OO(k)$.
In characteristic zero taking finite-group invariants is exact: the
averaging idempotent splits the invariant submodule. Thus there are no
higher direct images for $\rho$, and (D.1) follows. This proves the
cohomology comparison in the direction used here, not merely equality
of Euler characteristics. Finally every monomial of degree $n$ is
$x^{r_2}y^{r_1}$ times a homogeneous polynomial of degree $d(n)$ in
$x^4,y^3$. This gives the stated count of sections. The negative-degree
Cech complex on $x\ne0$, $y\ne0$ has cokernel represented by
$x^{-a}y^{-b}$ with $a,b\ge1$ and $3a+4b=-n$; replacing $(a,b)$ by
$(a-1,b-1)$ proves the $h^1$ formula without an additional duality
assumption. The two calculations agree exactly for every integer $n$;
the program \texttt{code/d\_exact.py} also compares them for
$-240\le n\le240$.
\end{enumerate}
\end{proof}

\begin{proposition}[Vertical and logarithmic tangent bundles]\label{inv:D1.2}
\StatementRecord{inv:D1.2}{proved}{Historical events: GPT-6 Astra (OpenAI), ChatGPT Pro mode, effort Pro, S6\_09\_PAPERS, 17 November 2026: targeted proof check. Also checked in this writing. Pre-continuation writing events: GPT-6 Astra (OpenAI), ChatGPT Pro mode, effort Pro, S6\_10\_PAPERS, 18 November 2026: targeted proof check; used and rechecked in this run. Continuation metadata events: GPT-6 Astra (OpenAI), ChatGPT Pro mode, effort Pro, S6\_10\_PAPERS continuation, 19 November 2026: used and rechecked in this run; status-vocabulary reconciliation; metadata reconciliation only; no new mathematical check. \emph{Pre-continuation recorded proof scope.} Read full root-base, logarithmic tangent, direct-image and twisted-Dolbeault proofs. Re-derived coarse degree function, local normalization quotient, resonance equation, Fourier contraction comparison, joint order-three coefficient 123ik/(64pi), stated local modules, and jump-locus finiteness. Exact Koszul and twist-program computations retain their separately recorded replay scopes. Matched the previously successful unchanged d\_exact.py and d\_koszul.py replays to this full tangent/Koszul proof:481 weighted cohomology counts, root-character invariants, polynomial module relations and the joint coefficient123i/(64pi). The earlier event saying replay pending is superseded at exactly this finite-algebra scope; its complete analytic comparison review remains separate.}{the hypotheses, notation and definitions stated here; \cref{inv:D1.1}}

There is a canonical representable proper flat holomorphic map
$\bar f:X\to\mathcal B$ with $f=\rho\bar f$. If
\[
 \mathcal T=T_X(-\log W),\qquad
 \mathcal V=\ker\bigl(\mathcal T\to
                   \bar f^*T_{\mathcal B}(-p_0)\bigr),
\]
then, for rank-two vector bundles $E,G$ on $\mathcal B$,
\[
 \mathcal V=\bar f^*E,\quad E=\OO\oplus\OO(-1),\qquad
 G=R^1\bar f_*\OO_X=\OO\oplus\OO(-1),                 \tag{D.2}
\]
\[
 R^0\bar f_*\OO_X=\OO,\qquad R^2\bar f_*\OO_X=\OO(-1),
\]
and the following is an exact sequence of locally free sheaves on $X$:
\[
 0\longrightarrow\bar f^*E\longrightarrow\mathcal T
 \longrightarrow\bar f^*\OO(-5)\longrightarrow0.       \tag{D.3}
\]
The trivial summands in $E,G$ are generated by the vertical field $\xi$
and the constant integral cocycle $\gamma$, respectively.
\end{proposition}



The root lift, vertical splitting, logarithmic quotient and
structure direct images are re-derived after reading
\citet[Lemmas~6.2 and 6.7; Propositions~6.6 and 6.10]{DuEtAl2026v1}.
The complete root and ordinary-sheaf arguments are retained below.

\begin{proof}
\begin{enumerate}[label=\arabic*.,leftmargin=*]
\item The divisors $S_j$ and their canonical sections give root data
$\OO_X(S_j)^{m_j}=f^*\OO_B(p_j)$, hence the indicated map to the root
curve. Over its order-$m_j$ chart the fibre product is the original
smooth torus family $J'_j$, since its free affine cyclic quotient is
exactly $f^{-1}(D_j)$. The comparison $J'_j\to X$ is finite etale;
in particular pullback is faithful and the stated local comparisons
descend. Away from the two root points the map is the original $f$.
Properness and flatness may be checked after this atlas: properness is
the properness of the original fillings, and flatness holds because a
nonconstant base parameter is a non-zero-divisor in every smooth
total-space local ring. Fibres have dimension two, including the cusp.

\item On a smooth torus chart a vertical field is a holomorphic linear
combination of the translation fields. Its coefficients are constant
on each compact fibre and therefore depend only on the base. The
vertical derivative of an affine transition is its linear part $R_g$;
the translating part, including the logarithmic regluing, has vertical
derivative the identity. At the cusp the two fields
$x_1\partial_{x_1},x_2\partial_{x_2}$ extend as logarithmic vector
fields in every toric chart. The deck maps multiply each $x_i$ by a
base-dependent unit and a power of the base parameter. Their vertical
logarithmic derivative fixes both coefficient directions. The
exponential overlap sends $\partial_{\zeta_i}$ to
$2\pi i x_i\partial_{x_i}$. These transitions depend only on the base
and define $E$; they prove the actual bundle identity
$\mathcal V=\bar f^*E$.

\item We check the kernel and surjectivity at every type of cusp point. At
$t=xy$, the kernel of $df$ is freely generated by
$x\partial_x-y\partial_y$ and $\partial_z$. At $t=xyz$ it is generated
by $x\partial_x-y\partial_y$ and $x\partial_x-z\partial_z$.
Indeed the equation $yz a+xz b+xy c=0$ in the regular local ring forces
$x\mid a,y\mid b,z\mid c$; after division the only relation is that
the three coefficients sum to zero. The logarithmic map sends
$\sum a_i z_i\partial_{z_i}$ to $(\sum a_i)t\partial_t$ and is
surjective. At a finite root point it is the ordinary relative tangent
sequence of the smooth map to $s$, not the singular differential of
$s^{m_j}$. These comparisons prove (D.3), since
$T_{\mathcal B}(-p_0)=\OO(7-12)$.

\item The invariant vector $e_2$ defines a subbundle $\OO\xi\subset E$.
It remains nonzero as a vector of the logarithmic vertical bundle at
the cusp, even where its image in the ordinary tangent fibre is zero.
The quotient line has automorphy factors $-1/\tau$ and $1/\tau$ under
the order-three and order-four generators. The modular discriminant
has weight twelve and no zero in the upper half-plane. Therefore its
reciprocal is a meromorphic section of the twelfth tensor power of
this quotient line. On the cusp chart
$\tau=(\log q)/(2\pi i)+h(q)$ with $h$ holomorphic, so the product
expansion of the discriminant is $q$ times a holomorphic unit.
There are no other zeros or poles, including on the finite root
charts. The quotient line thus has twelfth power
$\OO(-p_0)=\OO(-12)$, and the torsion-free Picard group in \ref{inv:D1.1} makes
it $\OO(-1)$. Its extension by $\OO$ splits because
$H^1(\mathcal B,\OO(1))=0$.

\item For completeness the structure-sheaf computation at the cusp uses the
normalization $P=\mathrm dP_6$, the three double curves and the two
triple points. Their exact normalization complex has global section
complex
\[
 \C\xrightarrow{0}\C^3
 \xrightarrow{(a,b,c)\mapsto(a-b+c,a-b+c)}\C^2.
\]
Exactness at the two local models is the elementary gluing of functions
on the coordinate branches; all higher cohomology of the rational
normalization, the rational double curves and the points vanishes.
The complex therefore gives $(h^0,h^1,h^2)=(1,2,1)$. Its product is
the Cech product of the two-triangle torus and gives a nondegenerate
alternating pairing $H^1\wedge H^1\to H^2$. This also follows directly
from the two opposite oriented triangles: the two coordinate cocycles
have determinant one on their oriented fundamental cycle. These are
the local comparison and product in Claude and Alpoge (2026,
Lemma 9.12), used here on the same normalization and with both triple
points retained.

\item On all other geometric fibres of $\bar f$ the cohomology is that of a
smooth two-torus. Constancy of these three dimensions and proper flat
analytic base change give locally free direct images on the atlas.
Their canonical comparison maps commute with the finite group actions
and with overlap restriction; hence they descend to $\mathcal B$.
The determinant of (D.3), using $W=\bar f^*p_0$, gives
\[
 \omega_{X/\mathcal B}=\bar f^*(\det E)^{-1}.
\]
Relative Serre duality, applied on that same atlas and descended by
its trace pairing, gives $R^2\bar f_*\OO_X=\det E=\OO(-1)$.
Cup product is an isomorphism $\det G\to R^2\bar f_*\OO_X$, since it
is so on every geometric fibre by the preceding calculation.

\item The class $\gamma$ is nowhere zero in $G$. On a smooth marked fibre,
$H^1(\OO)$ is the quotient of constant lattice cocycles by the two
period rows. If $\gamma$ were a combination of those rows, evaluation
on $\widehat w,\widehat\delta$ would make both coefficients zero,
contradicting $\gamma(\widehat\gamma)=1$. The same argument applies
on each finite root cover, and affine translations act trivially on
these classes. At the cusp $\gamma$ vanishes on the angular sublattice
and descends to a nonzero character of its rank-two deck lattice.
In the normalization complex it is the corresponding nonzero
one-cocycle on the two-triangle torus. Its cup product with a complementary
integral character is nonzero, so it remains nonzero there. Thus
$\OO\gamma\subset G$ is a subbundle; its quotient is $\OO(-1)$ by
the determinant just computed. Once again $H^1(\OO(1))=0$ splits the
extension. This proves every assertion in (D.2).
\end{enumerate}
\end{proof}

\begin{proof}[Koszul and direct-image proof of Theorem~\ref{fam:st:A.5.2}]
\begin{enumerate}[label=\arabic*.,leftmargin=*]
\item Put $L=\OO_{\mathcal B}(-5)$. The first connecting map of (D.3) is
the holomorphic bundle map
$\kappa:L\to E\otimes G$. In the splittings fixed in \ref{inv:D1.2} the
coefficient in $\OO\xi\otimes\OO\gamma$ lies in $H^0(\OO(5))=0$.
The remaining coefficients lie in
$H^0(\OO(4))=\C y$, twice, and $H^0(\OO(3))=\C x$.
Therefore, as a map from $E^*$ to $G$ with its $L$ factor understood,
\[
 \kappa=\begin{pmatrix}0&a y\\ b y&c_0x\end{pmatrix}. \tag{D.5}
\]
We use $c_0$ for this coefficient to distinguish it from the deformation
parameter. The logarithmic period derivative at the cusp has leading
term the invertible matrix $B_0/(2\pi i)$. Hence
$\det\kappa=-ab y^2$ is nonzero there, so $ab\ne0$.
At the order-three point $\tau-\rho$ has order one on the root chart:
$j=1728t$ and the elliptic ramification order of $j$ at $\rho$ is
three. The marked torus period derivative is consequently nonzero
there: the two constant identity columns in the normalized period
matrix exclude an infinitesimal linear-frame change that could
otherwise conceal a nonzero derivative. In (D.5), $y=0$ and
$x\ne0$ at that point, so $c_0\ne0$.
These arguments establish the three nonzero coefficients; they do not
infer them from a global cohomology dimension. Rescaling the two
quotient frames and the source frame makes all three coefficients one.

\item The next connecting map is cup product with $\kappa$. In suitable
bases it is
\[
 L\otimes G=\OO(-5)\oplus\OO(-6)
 \xrightarrow{\ \delta_1\ }
 E\otimes\det G=\OO(-1)\oplus\OO(-2),\qquad
 \delta_1=\begin{pmatrix}-y&0\\-x&y\end{pmatrix}.       \tag{D.6}
\]
The sign follows by wedging first with $\gamma$ and then with the
second generator of $G$; reversing both sign conventions leaves the
cokernels unchanged. As its determinant is $-y^2$, it is injective as
a morphism of sheaves. The first connecting map is
$\delta_0=(0,y,y,x)^t$ and is a subbundle inclusion, since $x,y$ have
no common zero. Its cokernel is
\[
 \OO\oplus\OO(-1)\oplus\OO(2).                       \tag{D.7}
\]
Indeed one constant linear combination of the two $\OO(-1)$ summands
is unused, and the remaining quotient is given by the exact sequence
$0\to\OO(-5)\xrightarrow{(y,x)}\OO(-1)\oplus\OO(-2)
\xrightarrow{(x,-y)}\OO(2)\to0$.

\item The cokernel of (D.6) is supported at the order-three root point. There
$x$ is a unit, and elementary row and column operations identify it
with
\[
 Q=\OO(-1)\big|_{2p_1},                               \tag{D.8}
\]
where $2p_1$ means the length-two divisor on the root chart, not the
double coarse point. The generator is the first target component,
and its annihilator is $y^2$. The long direct-image sequence now gives
\[
 R^0\bar f_*\mathcal T=E,\quad
 R^1\bar f_*\mathcal T=\OO\oplus\OO(-1)\oplus\OO(2),
\]
\[
 0\longrightarrow Q\longrightarrow R^2\bar f_*\mathcal T
 \longrightarrow\OO(-6)\longrightarrow0.              \tag{D.9}
\]
This last extension splits, since $\operatorname{Ext}^1(\OO(-6),Q)
=H^1(Q(6))=0$; its splitting is not asserted canonical.

\item Let $\nu:P=\mathrm dP_6\to W\hookrightarrow X$ be the normalization
immersion. There is an exact sequence
\[
 0\longrightarrow\mathcal T\longrightarrow T_X
 \longrightarrow\nu_*K_P\longrightarrow0.             \tag{D.10}
\]
To verify the quotient, in a chart $W=(z_1\cdots z_r=0)$ the map takes
a vector field to its normal coefficient on each of the $r$ separate
branches. The coefficients lift independently; its kernel is precisely
the vector fields tangent to every branch. Thus the quotient is the
pushforward of the normal line of $\nu$, not $\OO_W(W)$.
The form $d\log x_1\wedge d\log x_2\wedge dq$ becomes a unit times
$dz_1\wedge dz_2\wedge dz_3$ in every unimodular chart and is invariant
under every cusp deck transformation. It trivializes $\nu^*K_X$.
Adjunction for the immersion gives its normal line $K_P$, proving
(D.10). Since $P$ is rational, its $K_P$ cohomology is $(0,0,1)$ by
Serre duality. Finite pushforward preserves that cohomology.

\item It remains to identify the extension at $p_0$ in degree two. We give
a fibre-sensitive calculation. The kernel of
\[
 \Omega_X^1|_W\longrightarrow\nu_*\Omega_P^1
\]
is $\nu_*\OO_P\otimes N_W^*$. At $xy=0$ its sections are
$a(y,z)y\,dx+b(x,z)x\,dy$; at $xyz=0$ they are
$a(y,z)yz\,dx+b(x,z)xz\,dy+c(x,y)xy\,dz$.
These identifications glue by the common conormal $dq$.
Since $N_W$ and $K_X|_W$ are trivial and $H^0(P,\Omega_P^1)=0$, the
Serre-dual space of $H^2(W,T_X|_W)$ is one-dimensional.
Top-degree base change identifies this group with the residue fibre
of $(R^2\bar f_*T_X)_{p_0}$. Its generic rank is one by (D.9).
A finite module over the discrete valuation ring $\C\{q\}$ of rank
one and residue dimension one is free of rank one. Thus the extension
by $\C_{p_0}$ increases the degree of its line part by $p_0$, without
adding torsion. We obtain the complete root direct images
\[
 R^0\bar f_*T_X=\OO\oplus\OO(-1),\quad
 R^1\bar f_*T_X=\OO\oplus\OO(-1)\oplus\OO(2),
\]
\[
 R^2\bar f_*T_X=\OO(6)\oplus Q.                       \tag{D.11}
\]
On the order-three chart, an invariant coefficient of the generator
of $Q$ must have exponent $2$ modulo $3$, whereas $Q$ only retains
exponents $0,1$. Therefore $\rho_*Q=0$. Applying the exact $\rho_*$
to (D.11) and using
$d(0)=0,d(-1)=d(2)=-1,d(6)=0$ proves (D.4).
Finally Leray over a curve has only columns zero and one. The usual
cohomology of these displayed line bundles gives the four asserted
dimensions. This is a direct-image computation, not an inference from
the tangent Euler characteristic.
\end{enumerate}
\end{proof}

\begin{proof}[Actual-gluing Kodaira--Spencer calculation for Theorem~\ref{fam:st:A.5.2}]
\begin{enumerate}[label=\arabic*.,leftmargin=*]
\item All line bundles other than the trivial summands in
$R^q\bar f_*\mathcal V=E\otimes R^q\bar f_*\OO_X$ are acyclic:
they have weights $-1,-2$. Hence $H^1(X,\mathcal V)=
\C(\xi\otimes\gamma)$ and its other positive-degree groups vanish.
The three direct images of $\bar f^*\OO(-5)$ have weights
$-5$, $(-5,-6)$, and $-6$, all acyclic by \ref{inv:D1.1}. Consequently (D.3)
induces an isomorphism in every cohomology degree. Sequence (D.10),
whose quotient has no cohomology in degrees zero and one, then gives
the claimed isomorphism $H^1(\mathcal V)\to H^1(T_X)$.

\item We use the Cech convention in which the derivative of a transition
arrow is composed with the inverse of the central arrow. A lattice
transition is translation by $\Pi_c\lambda$, and its derivative in
$c$ is $\gamma(\lambda)\xi$. The linear deck matrices are independent
of $c$. On the finite charts the affine translations contribute
$\gamma(v_j)\xi/m_j$; their $m_j$th iterates give
$\gamma(v_j)\xi$, precisely the lattice cocycle. On the overlap with
the smooth part the derivative of the logarithmic translating section
is $\gamma(v_j)(\log s_j)\xi/(2\pi i)$. Its Cech coboundary accounts
for the difference of these representatives. On the cusp the deck
multiplier differentiates to $\gamma(\lambda)\xi$ in logarithmic
coordinates, and the exponential overlap identifies it with the same
cocycle. Thus the global Kodaira--Spencer class is
$[\xi\otimes\gamma]$, with a minus sign if the tangent-lift convention
$[\bar\partial h]$ is used instead. The sign is recorded; nonvanishing
does not depend on it.

\item Nonvanishing can also be checked on a regular torus. In the quotient
description of $H^1(\OO)$ used above, $[\gamma]\ne0$; the first
coefficient of (D.5) is zero, so $\xi\gamma$ cannot be a scalar
multiple of the base Kodaira--Spencer class. Its image in
$H^1(T_X|_F)$ therefore remains nonzero. Since $H^1(X,T_X)$ is
one-dimensional by \ref{inv:D1.3}, this proves the asserted isomorphism.
\end{enumerate}
\end{proof}

\begin{proof}[Kuranishi ideal calculation for Theorem~\ref{fam:st:A.5.2}]
\begin{enumerate}[label=\arabic*.,leftmargin=*]
\item Let $\alpha$ be the generator in \ref{inv:D2.1}. The constructed holomorphic
family gives an actual convergent deformation tangent to $\alpha$.
Differentiating its integrability equation to second order gives
$\bar\partial\alpha_2=-[\alpha,\alpha]/2$, so the cohomology class
of $[\alpha,\alpha]$ is zero. Bilinearity proves that the bracket
$H^1(T_X)\times H^1(T_X)\to H^2(T_X)$ vanishes on every pair.
This argument alone would not exclude higher obstructions, so we
make the analytic argument as well. A Kuranishi germ is a closed
analytic subgerm of its tangent space $H^1(T_X)=\C$. Its classifying
map from the actual parameter disc has nonzero derivative by \ref{inv:D2.1}.
Composing with this embedding gives a holomorphic map of discs with
nonzero derivative; the inverse function theorem makes its image
contain a neighbourhood of zero in $\C$. Every equation of the
Kuranishi germ vanishes on that neighbourhood and is therefore zero.
The germ, including its analytic structure, is the smooth line.
The same inverse-function argument makes the classifying map an
isomorphism of germs. Pulling back the semiuniversal deformation
therefore realizes every sufficiently small deformation in the
$c$-family. This is semiuniversality; the positive-dimensional
automorphism group is not being discarded or interpreted as a
fine-moduli assertion.
\end{enumerate}
\end{proof}




\section{The integral parameter comparison}

\subsection{The blind control: local charts and coherent direct images}

The proof in this section was written before opening either version of
arXiv:2610.00250.  Only the construction data of AC26 are used for the
threefold.  Write $B=\mathbb P^1$, let $S_1,S_2$ be the reduced fibres of
multiplicities $3,4$, let $W$ be the cusp fibre, and let $D$ be its reduced
double locus.  Thus $D$ consists of three rational curves joining the two
triple points.  Set
\[
 L=f^*\mathcal O_B(-1)\otimes\mathcal O_X(2S_1+S_2),\quad
 F=\mathcal O_X(S_1+S_2),\quad
 E=\ker(df),\quad J=\mathcal I_D\otimes F.
\]
Let $\widehat B$ denote the analytic root orbifold with orders $3,4$ at
$p_1,p_2$, let $\pi:\widehat B\to B$ be its coarse map, and let
$h:X\to\widehat B$ be defined by the canonical sections of $S_1,S_2$.
In particular $f=\pi h$.  This notation means descent on the ordinary
root discs by the specified finite cyclic actions.  Taking invariants is
exact over $\mathbb C$, by averaging.  Denote by $\widehat L,\widehat F$
the root line bundles whose pullbacks are $L,F$.

\begin{theorem}\label{night:P5.1}\label{stmt:P5.1}
\StatementRecord{night:P5.1}{proved}{Historical events: GPT-6 Astra (OpenAI), ChatGPT Pro mode, effort Pro, S6\_06\_BLIND, 16--18 November 2026: targeted proof check. Also checked in this writing. Pre-continuation writing events: GPT-6 Astra (OpenAI), ChatGPT Pro mode, effort Pro, S6\_10\_PAPERS, 18 November 2026: targeted proof check. Continuation metadata events: GPT-6 Astra (OpenAI), ChatGPT Pro mode, effort Pro, S6\_10\_PAPERS continuation, 19 November 2026: used and rechecked in this run; status-vocabulary reconciliation; metadata reconciliation only; no new mathematical check. \emph{Pre-continuation recorded proof scope.} Read the full cp25 re-derivation and re-derived the local syzygy/root-cohomology, KS period equation, integral translation-sheaf, or order-two obstruction calculation stated in the respective proof. These are credited DGKX comparisons; the event does not assert an independent historical discovery.}{the definitions and exact calculations in this statement}

The sheaf $E$ is locally free of rank two and is the pullback of the
translation Lie bundle $V$ on $\widehat B$.  There is an exact sequence
\begin{equation}\label{eq:P5tangent}
 0\longrightarrow E\longrightarrow T_X\longrightarrow J\longrightarrow0.
\end{equation}
\end{theorem}
\begin{proof}
\begin{enumerate}[label=\arabic*.,leftmargin=*]
\item On a finite root chart $t=s^m$, the image of the differential is
$(s^{m-1})f^*T_B$.  The root chart maps étale to $X$ because its affine
cyclic action is free.  Thus the vanishing orders along $S_1,S_2$ are
$2,3$.

\item At a cusp double point use $q=xy$ in coordinates $(x,y,z)$.  The image of
$dq$ is $(x,y)$, and its kernel is freely generated by
$x\partial_x-y\partial_y$ and $\partial_z$: reducing $ya+xb=0$ modulo
$x$ and $y$ gives $a=xA,b=-yA$.  At a triple point $q=xyz$, reduction of
$yz a+xz b+xy c=0$ modulo $x,y,z$ gives
$a=xA,b=yB,c=zC$, and cancellation gives $A+B+C=0$.  Hence its kernel is
freely generated by $(x,-y,0),(x,0,-z)$ and its image is
$(xy,xz,yz)$, the reduced ideal of the three axes.  At a smooth point of
$W$ the differential is surjective.  Consequently its global image is
\[
 \mathcal I_D\otimes f^*T_B\otimes\mathcal O_X(-2S_1-3S_2)
 =\mathcal I_D\otimes\mathcal O_X(S_1+S_2),
\]
using $3S_1=f^*(p_1)$ and $4S_2=f^*(p_2)$.  This proves the exact sequence
including the scheme structure at a triple point.  On the toric charts
the two vertical logarithmic torus fields give exactly this kernel.
Period translations and collar translations have identity vertical
derivative; finite generators act by their vertical linear matrices.
The evaluations of the translation Lie bundle therefore glue to
$h^*V\simeq E$.
\end{enumerate}
\end{proof}

\begin{theorem}\label{night:P5.2}\label{stmt:P5.2}
\StatementRecord{night:P5.2}{proved}{Historical events: GPT-6 Astra (OpenAI), ChatGPT Pro mode, effort Pro, S6\_06\_BLIND, 16--18 November 2026: targeted proof check. Also checked in this writing. Pre-continuation writing events: GPT-6 Astra (OpenAI), ChatGPT Pro mode, effort Pro, S6\_10\_PAPERS, 18 November 2026: targeted proof check. Continuation metadata events: GPT-6 Astra (OpenAI), ChatGPT Pro mode, effort Pro, S6\_10\_PAPERS continuation, 19 November 2026: used and rechecked in this run; status-vocabulary reconciliation; metadata reconciliation only; no new mathematical check. \emph{Pre-continuation recorded proof scope.} Read the full cp25 re-derivation and re-derived the local syzygy/root-cohomology, KS period equation, integral translation-sheaf, or order-two obstruction calculation stated in the respective proof. These are credited DGKX comparisons; the event does not assert an independent historical discovery.}{P5.1}

On the root orbifold,
\[
 V\simeq\mathcal O\oplus\widehat L,
 \qquad R^qh_*\mathcal O_X\simeq
 \bigwedge^q(\mathcal O\oplus\widehat L)\quad(0\leq q\leq2).
\]
On $B$, the direct images used below are
\[
\begin{array}{c|ccc}
 & f_* & R^1f_* &R^2f_*\\ \hline
 \mathcal O_X&\mathcal O&\mathcal O\oplus\mathcal O(-1)&\mathcal O(-1)\\
 L&\mathcal O(-1)&\mathcal O(-1)^{\oplus2}&\mathcal O(-1)\\
 F&\mathcal O&\mathcal O^{\oplus2}&\mathcal O\\
 E&\mathcal O\oplus\mathcal O(-1)&\mathcal O\oplus\mathcal O(-1)^{\oplus3}
       &\mathcal O(-1)^{\oplus2}\\
 J&\mathcal O(-1)&\mathcal O(-1)^{\oplus2}&\mathcal O .
\end{array}
\]
\end{theorem}
\begin{proof}
\begin{enumerate}[label=\arabic*.,leftmargin=*]
\item Applying the homology matrices of the construction to the last two
period columns gives the vertical linear factors
\[
 R_1=\begin{pmatrix}-1/\tau&0\\(1-\mu)/\tau&1\end{pmatrix},\qquad
 R_2=\begin{pmatrix}1/\tau&0\\-\mu/\tau&1\end{pmatrix}.
\]
The second coordinate line is fixed.  The other eigenvalues at the root
points are $\zeta_1^2,\zeta_2$, where
$\zeta_j=\exp(-2\pi i/m_j)$.  At the cusp both logarithmic torus fields
extend in the abstract bundle $E$ by Statement P5.1.

\item The modular discriminant
$\Delta(\tau)=q_\tau\prod_{n\geq1}(1-q_\tau^n)^{24}$ has weight $12$,
has no zero on the upper half-plane, and has a simple zero at the cusp.
The determinant factor of $R_g$ is the reciprocal modular factor, up to
a sign which disappears in its twelfth power.  Thus $\Delta^{-1}$ is a
meromorphic section of $(\det V)^{12}$ with precisely a simple pole at
$p_0$.  Its two root characters determine
\begin{equation}\label{eq:P5L}
 \det V=\widehat L
 =\pi^*\mathcal O_B(-1)\otimes
   \mathcal O_{\widehat B}(2p_1/3+p_2/4).
\end{equation}
Indeed the right side has the prescribed root characters and its
twelfth power is $\pi^*\mathcal O_B(-1)$; the quotient descends to
$\mathbb P^1$ and has trivial twelfth power, so is trivial.

\item For every root line bundle the invariant direct image is
\begin{equation}\label{eq:P5floor}
 \pi_*\bigl(\pi^*\mathcal O_B(d)\otimes
 \mathcal O_{\widehat B}(ap_1/3+bp_2/4)\bigr)
 =\mathcal O_B\bigl(d+\lfloor a/3\rfloor+\lfloor b/4\rfloor\bigr).
\end{equation}
This follows on a root disc by writing its invariant meromorphic
coefficients in integral powers of $t=s^m$; the largest allowed pole
order is $\lfloor a/m\rfloor$.  It also proves the formula for negative
integers.  The invariant line in $V$ gives
$0\to\mathcal O\to V\to\widehat L\to0$.  Its extension group is
$H^1(\widehat B,\widehat L^{-1})=H^1(B,\mathcal O(-1))=0$, so it splits.

\item On a smooth torus, compute $H^1(\mathcal O)$ as the quotient of flat
cohomology by the two period rows.  In that quotient $\gamma,u$ are a
basis, with
\[
 w=-6\mu\gamma-\tau u,\qquad \delta=-\beta\gamma-\mu u.
\]
The transitions are triangular with diagonal $(1,\det R_g)$; for example
$T_1u=-u-w$ gives second column
$(6\mu(g_1z),\tau(g_1z)-1)^{\mathsf t}$.  The image of $\gamma$ is
nowhere zero: comparison of the $w,\delta$ coefficients in the period
rows shows that $\gamma$ never lies in their span.

\item We check the extension and base change at the cusp.  Its normalization
$S=dP_6$ and its double curves give the exact branch complex
\begin{equation}\label{eq:P5normalization}
0\to\mathcal O_W\to\nu_*\mathcal O_S\to
 \bigoplus_{i=1}^3\mathcal O_{D_i}\to
 \mathbb C_P\oplus\mathbb C_Q\to0.
\end{equation}
At $xy=0$ this is the fibre-product identity; at $xyz=0$ it is
inclusion-exclusion of the three convergent power-series rings.  The
restriction differences have the single alternating compatibility
$g_{12}(0)-g_{13}(0)+g_{23}(0)=0$, which suffices to reconstruct the
three branch functions.  All the strata have vanishing higher coherent
cohomology.  Their global complex is
\[
 \mathbb C\xrightarrow{0}\mathbb C^3
 \xrightarrow{\left(\begin{smallmatrix}1&-1&1\\1&-1&1\end{smallmatrix}\right)}
 \mathbb C^2.
\]
It gives $(h^0,h^1,h^2)(\mathcal O_W)=(1,2,1)$.  It is the cellular
cochain complex of the two-triangle torus.  Its product in degrees
$(1,1)$ is the nonzero exterior product to the oriented top class.

\item The constant branch complex identifies
$H^1(W,\mathbb C)\to H^1(W,\mathcal O_W)$ as an isomorphism.  A loop
recording a component translation in the toroidal cover has its class
in $\Lambda/\Lambda_{\rm tor}$; $\gamma,u$ evaluate as its two
coordinates.  Locally constant Cech cocycles on this cover, viewed as
holomorphic cocycles, therefore extend the two classes through the
cusp, and their central values are the basis just calculated.

\item On each ordinary or root disc, the local ring of the smooth total
space is torsion free over the base DVR, hence flat.  Properness is part
of the construction and every fibre has dimension two.  The structure
sheaf is flat and its fibre cohomology dimensions are the constant
values $(1,2,1)$.  Grauert's cohomology and base-change theorem, \S7
Satz 5, applies.  We checked the 1963 corrigendum, whose correction in
\S7 concerns Satz 2.  Thus the direct images are locally free and commute
with base change.  The extension
$0\to\mathcal O\to R^1h_*\mathcal O_X\to\widehat L\to0$ exists on all
charts and splits by the preceding extension-group calculation.  The
relative cup product identifies $R^2h_*\mathcal O_X$ with $\widehat L$,
being nonzero both on smooth tori and on $W$.  Connected fibres and the
maximum principle give $h_*\mathcal O_X=\mathcal O$.

\item For a line bundle $A$ on $\widehat B$, projection and descent give
\[
 R^qf_*h^*A=\pi_*\bigl(A\otimes
              \bigwedge^q(\mathcal O\oplus\widehat L)\bigr).
\]
On a root disc this means tensoring followed by the averaging
projector, so it includes the nonreduced coarse fibres.  Formula
\eqref{eq:P5floor} gives the first four rows in the statement.

\item For the last row one must compute restriction to $D$.  The conductor
square gives
\[
0\to\mathcal O_W\to\nu_*\mathcal O_S\oplus\mathcal O_D
\to\nu_*\mathcal O_{\partial S}\to0,
\]
where $\partial S$ is the hexagon.  Here $h^1(\mathcal O_{\partial S})=1$
and $h^1(\mathcal O_D)=2$.  The equality conditions at the two triple
points give incidence matrix
\[
 \begin{pmatrix}1&-1&0\\0&1&-1\\1&-1&0\\0&1&-1\end{pmatrix}.
\]
Orient the three theta-graph edges from $P$ to $Q$.  The hexagon
traversal maps to $e_1-e_2+e_3-e_1+e_2-e_3=0$.  The constant-to-coherent
branch comparison is an isomorphism in degree one, so the induced map
$H^1(D,\mathcal O_D)\to H^1(\partial S,\mathcal O_{\partial S})$ is
zero.  The conductor sequence then makes
$H^1(W,\mathcal O_W)\to H^1(D,\mathcal O_D)$ an isomorphism.

\item The canonical section of $F$ is nonzero near $W$, so $F|_D=\mathcal O_D$.
In $0\to J\to F\to\mathcal O_D\to0$, the maps of direct images
$\mathcal O\to\mathbb C_{p_0}$ and
$\mathcal O^{\oplus2}\to\mathbb C_{p_0}^{\oplus2}$ are evaluation and
evaluation followed by this isomorphism.  They are surjective, with
kernels $\mathcal O(-p_0)$ in the indicated ranks.  Finally
$H^2(D,\mathcal O_D)=0$ and $R^2f_*F=\mathcal O$.  This gives the last row.
\end{enumerate}
\end{proof}

\begin{theorem}\label{night:P5.3}\label{stmt:P5.3}
\StatementRecord{night:P5.3}{proved}{Historical events: GPT-6 Astra (OpenAI), ChatGPT Pro mode, effort Pro, S6\_06\_BLIND, 16--18 November 2026: targeted proof check. Also checked in this writing. Pre-continuation writing events: GPT-6 Astra (OpenAI), ChatGPT Pro mode, effort Pro, S6\_10\_PAPERS, 18 November 2026: targeted proof check. Continuation metadata events: GPT-6 Astra (OpenAI), ChatGPT Pro mode, effort Pro, S6\_10\_PAPERS continuation, 19 November 2026: used and rechecked in this run; status-vocabulary reconciliation; metadata reconciliation only; no new mathematical check. \emph{Pre-continuation recorded proof scope.} Read the full cp25 re-derivation and re-derived the local syzygy/root-cohomology, KS period equation, integral translation-sheaf, or order-two obstruction calculation stated in the respective proof. These are credited DGKX comparisons; the event does not assert an independent historical discovery.}{P5.2}

For every admissible $c$,
\[
 \bigl(h^0(T_{X_c}),h^1(T_{X_c}),h^2(T_{X_c}),h^3(T_{X_c})\bigr)
 =(1,1,1,0).
\]
\end{theorem}
\begin{proof}
\begin{enumerate}[label=\arabic*.,leftmargin=*]
\item The Laurent--Cech calculation on the two standard charts of
$\mathbb P^1$ gives $h^0(\mathcal O(d))=\max(d+1,0)$ and
$h^1(\mathcal O(d))=\max(-d-1,0)$.  Leray has only columns zero and one,
so has no higher differential.  Statement P5.2 gives
\[
 H^q(X,E)=(\mathbb C,\mathbb C,0,0),\qquad
 H^q(X,J)=(0,0,\mathbb C,0).
\]
The long exact sequence of \eqref{eq:P5tangent} proves the assertion;
every adjacent group which could change these dimensions is zero.
The bundle and incidence computations are independently recomputed in
\texttt{code/part5/bundle\_cohomology.py}.
\end{enumerate}
\end{proof}

\begin{theorem}\label{night:P5.4}\label{stmt:P5.4}
\StatementRecord{night:P5.4}{proved}{Historical events: GPT-6 Astra (OpenAI), ChatGPT Pro mode, effort Pro, S6\_06\_BLIND, 16--18 November 2026: targeted proof check. Also checked in this writing. Pre-continuation writing events: GPT-6 Astra (OpenAI), ChatGPT Pro mode, effort Pro, S6\_10\_PAPERS, 18 November 2026: targeted proof check. Continuation metadata events: GPT-6 Astra (OpenAI), ChatGPT Pro mode, effort Pro, S6\_10\_PAPERS continuation, 19 November 2026: used and rechecked in this run; status-vocabulary reconciliation; metadata reconciliation only; no new mathematical check. \emph{Pre-continuation recorded proof scope.} Read the full cp25 re-derivation and re-derived the local syzygy/root-cohomology, KS period equation, integral translation-sheaf, or order-two obstruction calculation stated in the respective proof. These are credited DGKX comparisons; the event does not assert an independent historical discovery.}{P5.3}

The germ of the $c$-parameter family at each admissible member is a
miniversal deformation germ.  The Kuranishi germ is a smooth disc.
\end{theorem}
\begin{proof}
\begin{enumerate}[label=\arabic*.,leftmargin=*]
\item The period and gluing formulas depend holomorphically on $c$, and the
strict admissibility condition is open, so a small $c$-disc gives a
proper holomorphic family.  A first-order trivialization of its
deformation in the direction $dc=1$ would carry $f$ to
$f+\epsilon v$, with $v\in H^0(X,f^*T_B)$.  Since $f_*\mathcal O_X$
is $\mathcal O_B$, this is a base vector field.  On a simply connected
regular base chart the integral marking is constant under a first-order
trivialization and is the identity at the reference point.  Translation
does not change periods.  Its infinitesimal linear coordinate change
$H(t)$ must therefore satisfy
\[
 \partial_c\Pi+v(t)\partial_t\Pi=H(t)\Pi.
\]
The last two period columns are the constant identity columns, so
$H=0$.  The upper entry of the second column gives $v\tau'=0$; the
constructed nonconstant $\tau$ has nonzero derivative on the regular
chart, and thus $v=0$ everywhere.  The lower entry of the first column
now says $1=0$.  This proves that the Kodaira--Spencer map is nonzero,
and by P5.3 it is an isomorphism.

\item Douady's \S VII.8 principal theorem gives a finite-dimensional versal
analytic germ, and \S VIII.1 Theorem 1 and its corollary give minimality
and uniqueness of its tangent map.  First-order changes of holomorphic
transition maps are tangent-valued Cech cocycles and changes of
trivialization are coboundaries, so its tangent space here is $H^1(T)$.
Embed the minimal germ as the zero space of an ideal
$I\subset\mathbb C\{z\}$.  Versality gives a map from our disc with
invertible derivative.  Substitution by a convergent series with
nonzero derivative is injective and annihilates $I$, forcing $I=0$.
The inverse-function theorem then identifies the germs.  This also
excludes a nonreduced embedded germ.
\end{enumerate}
\end{proof}

\subsection{All automorphisms and the integral parameter period}

\begin{theorem}\label{night:P5.5}\label{stmt:P5.5}
\StatementRecord{night:P5.5}{proved}{Historical events: GPT-6 Astra (OpenAI), ChatGPT Pro mode, effort Pro, S6\_06\_BLIND, 16--18 November 2026: targeted proof check. Also checked in this writing. Pre-continuation writing events: GPT-6 Astra (OpenAI), ChatGPT Pro mode, effort Pro, S6\_10\_PAPERS, 18 November 2026: targeted proof check. Continuation metadata events: GPT-6 Astra (OpenAI), ChatGPT Pro mode, effort Pro, S6\_10\_PAPERS continuation, 19 November 2026: used and rechecked in this run; status-vocabulary reconciliation; metadata reconciliation only; no new mathematical check. \emph{Pre-continuation recorded proof scope.} Read the full cp25 re-derivation and re-derived the local syzygy/root-cohomology, KS period equation, integral translation-sheaf, or order-two obstruction calculation stated in the respective proof. These are credited DGKX comparisons; the event does not assert an independent historical discovery.}{the definitions and exact calculations in this statement}

For every admissible twist of the $(3,4)$ construction,
\[
 \operatorname{Aut}(X_c)=\mathbb C/\mathbb Z\simeq\mathbb C^*.
\]
The parameter $a\bmod\mathbb Z$ acts by translation by $a\widehat\delta$.
\end{theorem}
\begin{proof}
\begin{enumerate}[label=\arabic*.,leftmargin=*]
\item Taking determinants in \eqref{eq:P5tangent}, whose quotient ideal has
codimension-two support, gives
\[
 K_X^{-1}=L\otimes F
 =f^*\mathcal O(-1)\otimes\mathcal O(3S_1+2S_2)
 =\mathcal O(2S_2),\qquad K_X^{-2}=f^*\mathcal O_B(1).
\]
The complete linear system of $K_X^{-2}$ is exactly $f$, since its
sections are the two sections pulled back from $B$.  Every automorphism
therefore preserves $f$ up to a base Möbius transformation.  The two
multiplicities and the cusp distinguish its three special fibres, so
that transformation is the identity.

\item A holomorphic torus automorphism is affine: its lifted derivative is
lattice-periodic, descends to a holomorphic matrix-valued function on a
compact torus, and hence is constant.  Integrating gives an affine lift.
The locally constant integral linear part must commute with
\[
 A_1=\begin{pmatrix}1&0&0&0\\6&0&1&0\\-6&-1&-1&0\\-2&1&0&1\end{pmatrix},\quad
 A_2=\begin{pmatrix}1&0&0&0\\0&0&-1&0\\-6&1&0&0\\3&0&1&1\end{pmatrix}.
\]
Solving $PA_j=A_jP$ over $\mathbb Q$ gives
$P=aI+bE_{41}$; the complete rational linear elimination is in
\texttt{code/part5/centralizer.py}.  Integrality gives $a,b\in\mathbb Z$,
and invertibility gives $a=\pm1$.  The equation $H\Pi=\Pi P$ first gives
$H=aI$ from the identity columns and then $b=0$ from the first column.

\item The linear part $-I$ is excluded by the first multiple fibre.  Normalize
the root base change $t=s^3$.  The automorphism lifts, fixing $s$; on the
central covering torus it is $z\mapsto-z+b$.  Commutation with the affine
deck transformation gives
\[
 (I-A_1)b=2v_1/3\pmod\Lambda.
\]
The torus character $\gamma:\Lambda_{\mathbb R}/\Lambda\to\mathbb R/\mathbb Z$
annihilates $I-A_1$, whereas the right side has character $2l_1/3$.
Freeness implies $3\nmid l_1$, a contradiction.  Thus every automorphism
is fibrewise translation.

\item We compute all components of the translation group.  Let $j$ be the
inclusion of the cusp-punctured orbifold and set
$\mathcal H=j_*\Lambda$.  Its cusp stalk is
$\ker(M_0-I)=\mathbb Z\widehat w\oplus\mathbb Z\widehat\delta$.
Let $\mathcal A$ be the sheaf of fibrewise translations extending to
the local compactification.  There is an exact sequence
\begin{equation}\label{eq:P5exponential}
0\longrightarrow\mathcal H\longrightarrow V\longrightarrow\mathcal A
\longrightarrow0.
\end{equation}
For the local check at the cusp, use the toroidal covering whose
component lattice is $\Gamma=\Lambda/\Lambda_{\rm tor}$.  A translation
acts identically on the nearby fibre lattice and on $\Gamma$, and
therefore lifts.  Its meridian remains bounding a disc through a smooth
point of $W$, which verifies the same assertion for the full local
cover.  On the punctured torus the lift has form $x\mapsto a(q)x$.
It sends the component $E_0$ to $E_v$; the index-one map
$B_0:\Gamma\to\mathbb Z^2$ is an isomorphism, so composing with the
unique corresponding deck shear makes it preserve $E_0$.  It then
preserves the torus $E_0\setminus\partial E_0$.  Evaluation at $(1,1)$
in the chart $(\mathbb C^*)^2\times\Delta_q$ shows that $a_i(q)$ are
holomorphic units at zero, and hence have logarithms.  Conversely their
multiplication extends on every toric chart and commutes with the deck
shears.  The logarithm kernel consists precisely of those periods with
no logarithmic term, namely $\ker(M_0-I)$.  On the finite root charts
ordinary logarithms prove the same exact sequence with their finite
equivariance retained.

\item Here is the integral lattice calculation, including torsion.  Put
\[
 M_0=(A_1A_2)^{-1}=\begin{pmatrix}1&0&0&0\\0&1&0&0\\0&1&1&0\\-1&0&0&1\end{pmatrix}.
\]
A cocycle on the punctured orbifold is a pair $(x_1,x_2)$ with
$N_1x_1=N_2x_2=0$, where
\[
 N_1=\begin{pmatrix}3&0&0&0\\6&0&0&0\\-12&0&0&0\\0&2&1&3\end{pmatrix},\qquad
 N_2=\begin{pmatrix}4&0&0&0\\12&0&0&0\\-12&0&0&0\\0&2&2&4\end{pmatrix}.
\]
Extension over the cusp says that $x_1+A_1x_2$ vanishes in
$\operatorname{coker}(M_0-I)$, so its first two coordinates vanish.
The complete integral solution is
\[
 x_1=(0,-b,2b-3c,c),\qquad x_2=(0,-2a-b,b,a),\quad a,b,c\in\mathbb Z.
\]
The principal cocycle of $(g,u,w,d)$ has parameters
\[
 \begin{pmatrix}a\\b\\c\end{pmatrix}
 =\begin{pmatrix}3&0&1\\-6&1&-1\\-2&1&0\end{pmatrix}
  \begin{pmatrix}g\\u\\w\end{pmatrix}.
\]
This matrix has determinant $-1$; its inverse is explicitly
$g=c-a-b$, $u=3c-2a-2b$, $w=4a+3b-3c$.  Hence
$H^1(\widehat B,\mathcal H)=0$ integrally.  The independent total-complex
calculation in \texttt{code/part5/orbifold\_lattice\_complex.py} prints
every differential and checks the same vanishing.

\item Sequence \eqref{eq:P5exponential} now makes every global translation a
global logarithm in $H^0(\widehat B,V)=\mathbb C$.  Its kernel is
$\Lambda^{A_1,A_2}=\mathbb Z\widehat\delta$.  The flow extends by
commutation with the affine finite actions and by torus multiplication
on the cusp charts.  This proves existence, faithfulness, and
completeness of the asserted group.
\end{enumerate}
\end{proof}

\begin{theorem}\label{night:P5.6}\label{stmt:P5.6}
\StatementRecord{night:P5.6}{proved}{Historical events: GPT-6 Astra (OpenAI), ChatGPT Pro mode, effort Pro, S6\_06\_BLIND, 16--18 November 2026: targeted proof check. Also checked in this writing. Pre-continuation writing events: GPT-6 Astra (OpenAI), ChatGPT Pro mode, effort Pro, S6\_10\_PAPERS, 18 November 2026: targeted proof check. Continuation metadata events: GPT-6 Astra (OpenAI), ChatGPT Pro mode, effort Pro, S6\_10\_PAPERS continuation, 19 November 2026: used and rechecked in this run; status-vocabulary reconciliation; metadata reconciliation only; no new mathematical check. \emph{Pre-continuation recorded proof scope.} Read the full cp25 re-derivation and re-derived the local syzygy/root-cohomology, KS period equation, integral translation-sheaf, or order-two obstruction calculation stated in the respective proof. These are credited DGKX comparisons; the event does not assert an independent historical discovery.}{the definitions and exact calculations in this statement}

For fixed admissible twist data, $X_c\simeq X_{c'}$ precisely when
$c-c'\in2\mathbb Z$.  This fixes the integral period of $c$; the
construction still permits an additive choice of $\beta^{\rm part}$.
\end{theorem}
\begin{proof}
\begin{enumerate}[label=\arabic*.,leftmargin=*]
\item The intrinsic pencil argument and the centralizer computation give
$P=aI+bE_{41}$ with $a=\pm1$ and $b=a(c-c')$.  The same root-fibre
character excludes $a=-1$.  Thus $b=c-c'$ is an integer.

\item Compute $H^2(\widehat B,\mathcal H)$ with a pair of pants, two root discs,
and a cusp disc.  The pair of pants has cochains $\Lambda\to\Lambda^2$
with differential $((A_1-I),(A_2-I))$.  A root disc has the periodic
cyclic resolution with alternating differentials $A_j-I,N_j$; the
cusp disc has $\ker(M_0-I)$ in degree zero.  An annulus has
$\Lambda\xrightarrow{A_j-I}\Lambda$.  The degree-one restriction from
the pair of pants is
\[
 (x_1,x_2)\longmapsto(x_1,x_2,-M_0(x_1+A_1x_2)).
\]
The total-complex ranks in degrees zero through three are $14,28,20,8$.
The script \texttt{orbifold\_lattice\_complex.py} prints its three
integer differentials, checks their zero products, and computes
$H^0=\mathbb Z$, $H^1=0$, $H^2=\mathbb Z\oplus\mathbb Z/2$.

\item The degree-two presentation also reduces by hand.  The invariant
lattices have bases $(\widehat\delta,\varepsilon_j)$, where
$\varepsilon_1=(1,2,-4,0)$ and $\varepsilon_2=(1,3,-3,0)$.
The five generators $(\delta_1,\varepsilon_1,\delta_2,\varepsilon_2,h_c)$
have relations
\begin{equation}\label{eq:P5H2}
 \delta_1=0,\qquad 2\delta_2=0,\qquad
 3\varepsilon_1=h_c=4\varepsilon_2.
\end{equation}
For the constant sheaf $\mathbb Z$, the same computation has generators
$z_1,z_2,h_c$ and relations $3z_1=h_c=4z_2$; its primitive generator is
$z_1-z_2$.  Inclusion $\mathbb Z\widehat\delta\to\mathcal H$ sends it
to the nonzero order-two class $\delta_2$.  The character $\gamma$
sends $\varepsilon_j$ to $z_j$ and kills $\delta_j$.

\item Both $H^1(\widehat B,V)$ and $H^2(\widehat B,V)$ vanish, by
$\pi_*V=\mathcal O\oplus\mathcal O(-1)$ and the dimension of the coarse
curve.  Consequently \eqref{eq:P5exponential} identifies compactified
translation torsors $H^1(\widehat B,\mathcal A)$ with
$H^2(\widehat B,\mathcal H)$.  This identification takes logarithms of
local translation transitions and their lattice-valued Cech
coboundaries; the two vanishings also reverse that construction.

\item The cyclic norm of the translation $v_j/m_j$ is $v_j$, while the cusp
collar has the opposite boundary orientation.  Thus the torsor
character is
\[
 l_1z_1+l_2z_2-l_0h_c
 =(4l_1+3l_2-12l_0)(z_1-z_2)=-p(z_1-z_2).
\]
The shear $S_b-I=b\widehat\delta\otimes\gamma$ changes its class by
$-bp\delta_2$.  Freeness forces $l_2$ odd and thus $p$ odd.  The class
is preserved precisely for even $b$.

\item For an integral shear the Jacobian charts themselves identify: the
period lattices differ by $S_b$, and the bounded cusp matrix changes by
an integer in its $\beta$ entry.  Every toric unit
$\exp(2\pi iC\lambda)$ is therefore unchanged.  The fan, lattice, and
toric gluings are the same, and on root discs the basis change
identifies the smooth Jacobian and its cyclic action.  Equality of the
torsor classes now gives a holomorphic isomorphism for the fixed sheaf
$\mathcal A$, which consists precisely of translations extending over
the cusp.  It extends on all three special charts.  Conversely every
isomorphism preserves the class, so forces $b$ even.
\end{enumerate}
\end{proof}
\section{Periods, compact subgroups, and the transcendence question}

We use a Weierstrass equation over $\overline\Q$ and its invariant
differential. Its period lattice is $\Lambda=\Z\omega_1+\Z\omega_2$,
with $\Im(\omega_2/\omega_1)>0$. The quasi-periods are defined by
$\zeta(z+\omega_j)=\zeta(z)+\eta_j$. We extend
$\eta:\Lambda\to\C$ real-linearly to $\C$. The discriminant in
this section is the discriminant of this equation, equivalently of
the lattice in the specified differential. The quantity $N$ is
independent of changing this algebraic equation and differential.

\begin{proposition}[sigma expression and the third-kind period]\label{ari:C.7.1}
\StatementRecord{ari:C.7.1}{proved}{Historical events: GPT-6 Astra (OpenAI), ChatGPT Pro mode, effort Pro, S6\_09\_PAPERS, 17 November 2026: targeted proof check. Also checked in this writing. Pre-continuation writing events: GPT-6 Astra (OpenAI), ChatGPT Pro mode, effort Pro, S6\_10\_PAPERS, 18 November 2026: targeted proof check. Continuation metadata events: GPT-6 Astra (OpenAI), ChatGPT Pro mode, effort Pro, S6\_10\_PAPERS continuation, 19 November 2026: used and rechecked in this run; status-vocabulary reconciliation; metadata reconciliation only; no new mathematical check. \emph{Pre-continuation recorded proof scope.} Re-derived the sigma norm, algebraic third-kind differential, exact semiabelian exponential, periods and logarithm identities, including the minus-point divisor convention.}{the hypotheses, notation and definitions stated here; \citet[Definitions 2.3--2.4 and
Section 6]{bertolin2022}; \cref{ari:eq:cp-periods}}

Let $P\in E(\overline\Q)$, with $P\ne O$ and $2P\ne O$, and
let $u$ be an elliptic logarithm of $P$. Then
\begin{equation}\label{ari:eq:cp-sigma}
 2\log N(P)=\frac16\log|\Delta|
       +\Re\{\Lambda_1+u(\zeta(u)-\eta(u))\},
 \qquad \Lambda_1=2\log\sigma(u)-u\zeta(u).
\end{equation}
The choices of logarithms change the expression inside the real
part by purely imaginary integral multiples of $2\pi$.
There is an extension $G_P$ of $E$ by $\mathbb G_m$, defined
over $\overline\Q$ and corresponding, under the principal
polarization, to $P$, for which $\Lambda_1$, modulo a logarithm
of an algebraic number, is the negative of a relative period
of $[\Z\to G_P]$. Changing an algebraic differential of the
third kind adds an algebraic multiple of $u$ and algebraic
endpoint terms.
\end{proposition}


\begin{proof}
\begin{enumerate}[label=\arabic*.,leftmargin=*]
\item The sigma--theta product, or its divisor and automorphy law
together with its derivative at zero, gives
\begin{equation}\label{ari:eq:cp-klein}
 N(P)=|\Delta|^{1/12}|\sigma(u)|
                   \exp\!\left(-\frac12\Re(u\eta(u))\right).
\end{equation}
To check the constant, for $\Lambda=\Z+\tau\Z$ one uses
$\sigma(z)=e^{\eta_1z^2/2}\vartheta_1(z,\tau)/(2\pi\eta(\tau)^3)$,
$\Delta=(2\pi)^{12}\eta(\tau)^{24}$, and
$\eta(z)=\eta_1z-2\pi i\Im z/T$. Substitution gives exactly
the defining Gaussian in $N$. Homogeneity gives the general
lattice formula. Taking logarithms proves \eqref{ari:eq:cp-sigma}.

\item We make the period identification explicit. Put
\[
 f_u(z)=\frac{\sigma(z+u)}{\sigma(z)\sigma(u)}e^{-\zeta(u)z}.
\]
Its logarithmic differential is
\[
 \frac{df_u}{f_u}
 =\frac12\frac{\wp'(z)-\wp'(u)}{\wp(z)-\wp(u)}\,dz.
\]
This is the pullback of an algebraic differential with residues
$+1$ at $-P$ and $-1$ at $O$. The apparent $0/0$ at $z=u$
has a removable singularity when $P$ is not $2$-torsion.
The rational cocycle $f_u(z+z')/(f_u(z)f_u(z'))$ defines
the algebraic extension; in its rational trivialization the
exponential map has the form
\[
 \exp_{G_P}(z,v)=(\exp_E z,e^v f_u(z)).
\]
Its periods include
\begin{equation}\label{ari:eq:cp-periods}
 (0,2\pi i),\qquad
 (\omega_j,\zeta(u)\omega_j-\eta_j u),\quad j=1,2.
\end{equation}
These signs follow directly from the sigma automorphy law.
The identity
$\sigma(2u)=-\wp'(u)\sigma(u)^4$ implies
\[
 \log f_u(u)=\Lambda_1+\log(-\wp'(u)).
\]
For an algebraic lift $R=(P,r)$, $r\in\overline\Q^{\times}$,
the second coordinate of $\log_{G_P}R$ is
\begin{equation}\label{ari:eq:cp-v}
 v=\log r-\Lambda_1-\log(-\wp'(u)).
\end{equation}
It is a period of the relative homology defining the indicated
$1$-motive. This also avoids treating an integral with an endpoint
at a residue pole as an unregularized ordinary integral.
The description agrees with the generalized logarithm and
Serre function in \citet[Definitions 2.3--2.4 and
Section 6]{bertolin2022}; the signs in
\eqref{ari:eq:cp-periods} are for the exponential rather than the
quasi-periods of $f_u$ itself.
\end{enumerate}
\end{proof}

\begin{proposition}[the compact-subgroup criterion and the real formula]\label{ari:C.7.2}
\StatementRecord{ari:C.7.2}{proved}{Historical events: GPT-6 Astra (OpenAI), ChatGPT Pro mode, effort Pro, S6\_09\_PAPERS, 17 November 2026: targeted proof check. Also checked in this writing. Pre-continuation writing events: GPT-6 Astra (OpenAI), ChatGPT Pro mode, effort Pro, S6\_10\_PAPERS, 18 November 2026: targeted proof check. Continuation mathematical and source checks: GPT-6 Astra (OpenAI), ChatGPT Pro mode, effort Pro, S6\_10\_PAPERS continuation, 19 November 2026: targeted proof check. Continuation metadata events: GPT-6 Astra (OpenAI), ChatGPT Pro mode, effort Pro, S6\_10\_PAPERS continuation, 19 November 2026: used and rechecked in this run; status-vocabulary reconciliation; metadata reconciliation only; no new mathematical check. \emph{Pre-continuation recorded proof scope.} Re-derived the real compactness character, rank-three compact kernel and exact algebraicity equivalence for the actual compact lift.}{the hypotheses, notation and definitions stated here; \cref{ari:C.7.1}; \cref{ari:eq:cp-periods}; \cref{ari:eq:cp-v}; \cref{ari:eq:cp-chi}; \citet[slides 5--8]{bertrand2017}}

For $P$ as in Statement \ref{ari:C.7.1}, $N(P)$ is algebraic if and only
if the fibre of $G_P$ above $P$ contains an algebraic point of
the maximal compact subgroup $G_P(\C)^c$.
For the family in the stated conventions, if $t>1$ is real, choose
$\tau=iT$, $T>1$, and the zero $\mu=1/2-iy$, $0<y<T/2$.
With $x=2y/T$, one has
\[
 N(P)^4=3^{3/4}(t-1)^{-1/4}
       \left|\frac{\vartheta_1(x,i/T)}{\eta(i/T)}\right|.
\]
\end{proposition}


\begin{proof}
\begin{enumerate}[label=\arabic*.,leftmargin=*]
\item On the analytic universal cover of $G_P$ define the real
homomorphism
\[
 \chi(z,v)=\Re\{v-\zeta(u)z+u\eta(z)\}.
\]
It vanishes on the periods \eqref{ari:eq:cp-periods}, hence descends
to $G_P(\C)$. Its restriction to the torus is $\log|\cdot|$.
The kernel is compact: in the universal cover it is a real
three-dimensional vector space modulo the rank-three period
lattice. Every compact subgroup has zero image in $\R$, so
this kernel is exactly $G_P(\C)^c$.
The identity
$\Re(u\eta(z))=\Re(z\eta(u))$ follows from the Legendre relation;
equivalently $\eta(z)=sz+\pi\bar z/\operatorname{area}(\Lambda)$
for a complex constant $s$. Substituting
\eqref{ari:eq:cp-v} at $z=u$ gives
\begin{equation}\label{ari:eq:cp-chi}
 \chi(R)=\log|r|-\log|\wp'(u)|
            -2\log N(P)+\frac16\log|\Delta|.
\end{equation}
If $r$ is algebraic and $\chi(R)=0$, this equation makes
$N(P)^2$ algebraic. Conversely, if $N(P)$ is algebraic,
choose the positive real algebraic number
$r=|\wp'(u)|\,|\Delta|^{-1/6}N(P)^2$.
Then \eqref{ari:eq:cp-chi} vanishes. This proves both directions
with the chosen rational trivialization of the fibre.
The compact interpretation is the one discussed by
\citet[slides 5--8]{bertrand2017}; here the calculation fixes
its normalization on the diagonal $P=P$.

\item For the real formula, the rectangular elliptic curve has its
two zeros of $\wp$ on $1/2+i\R$; along that edge $\wp$ runs
between the positive root $e_1$ and the negative root $e_2$.
Thus the indicated representative exists. Translate $2\mu$
by $-1$ and apply the modular substitution $\tau\mapsto-1/\tau$.
The resulting elliptic argument is $-2y/T$, which is real.
The Gaussian is therefore one. Evenness of $N$ and the
duplication identity prove the displayed formula. This removes
the Gaussian from this particular expression but makes no
algebraicity assertion about its theta value.
\end{enumerate}
\end{proof}

\subsection{Motivic groups from the point and extension data}

\citet[Construction~3.4, p.~12, Lemma~3.5, p.~13, and
Proposition~4.12, pp.~19--20]{bertolin2022} compute the motivic
group from its abelian point and extension data and the central
Weil-pairing bracket. For the $1$-motives used here, let $B$ be the
smallest abelian subvariety generated by these data, and let $Z(1)$
be the central torus obtained from that bracket and the remaining
algebraic lift. Their dimension formula is
\begin{equation}\label{ari:eq:cp-MTdim}
 \dim G_{\mathrm{mot}}(M)
 =\dim G_{\mathrm{mot}}(A)+2\dim B+\dim Z(1).
\end{equation}
The bracket is described in their Section~6.1, pp.~26--27, and
their corrected Table~2, pp.~28--29, gives unipotent dimension
three in the dependent nondeficient case. The following calculation
determines these groups through the displayed polarization maps;
the arithmetic question for the metric norm is recorded in the
proved implication with its explicitly open hypothesis in \cref{ari:C.7.4}.

\begin{proposition}[the motivic groups needed on the diagonal]\label{ari:C.7.3}
\StatementRecord{ari:C.7.3}{proved}{Historical events: GPT-6 Astra (OpenAI), ChatGPT Pro mode, effort Pro, S6\_09\_PAPERS, 17 November 2026: targeted proof check. Also checked in this writing. Pre-continuation writing events: GPT-6 Astra (OpenAI), ChatGPT Pro mode, effort Pro, S6\_10\_PAPERS, 18 November 2026: targeted proof check. Continuation mathematical and source checks: GPT-6 Astra (OpenAI), ChatGPT Pro mode, effort Pro, S6\_10\_PAPERS continuation, 19 November 2026: computation or formal certificate; independent re-proof; used and rechecked in this run. Continuation metadata events: GPT-6 Astra (OpenAI), ChatGPT Pro mode, effort Pro, S6\_10\_PAPERS continuation, 19 November 2026: used and rechecked in this run; status-vocabulary reconciliation; metadata reconciliation only; no new mathematical check. \emph{Pre-continuation recorded proof scope.} Re-derived the pure and unipotent dimensions from the point-extension subvariety and alternating bracket, including the conjugate dependent center and quotient-lift parameter. Matched full attached Bertolin passages to primary arXiv2010.15170v2: Definitions2.3-2.4p8, Construction3.4p12, Lemma3.5p13, Proposition4.12pp19-20, Remark4.13p20, bracket Section6.1pp26-27, errataTable2pp28-29. Conjectures are not assumed.}{the hypotheses, notation and definitions stated here; \cref{ari:eq:cp-MTdim}}

Let $P$ be nontorsion and let $R$ be any algebraic lift above
$P$ in $G_P$. Put $M_R=[\Z\to G_P]$, $1\mapsto R$.
Its connected motivic group is, after a choice of Levi subgroup,
the semidirect product of $G_{\mathrm{mot}}(E)$ and the
three-dimensional Heisenberg group attached to $H_1(E)$ and
its principal polarization. Its dimension is $5$ for CM and
$7$ without CM.

For a non-real embedding put
$\mathcal M=M_R\oplus\overline{M_R}$.
If $E$ and $\bar E$ are not isogenous, its dimension is $13$.
If they are isogenous, choose an isogeny
$\phi:\bar E\to E$ of degree $d$ and write
$Q=\phi(\bar P)$. Let $r$ be the dimension over
$F=\End(E)\otimes\Q$ of the span of $P,Q$, modulo torsion.
Then
\[
\begin{array}{c|c|c|c}
\text{case}&\dim B&\dim Z&\dim G_{\mathrm{mot}}(\mathcal M)\\ \hline
r=2,\ E\text{ CM}&2&2&8\\
r=2,\ E\text{ not CM}&2&2&10\\
r=1,\ E\text{ CM}&1&1+e&5+e\\
r=1,\ E\text{ not CM}&1&1+e&7+e.
\end{array}
\]
Here $e$ is either zero or one: it is zero precisely when the
algebraic lift in the one-dimensional quotient of the two
central tori by the bracket torus is torsion. This criterion is
independent of the choices made modulo isogeny. In the real
case we use the single motive $M_R$; no extra quotient torus
is necessary.
\end{proposition}


\begin{proof}
\begin{enumerate}[label=\arabic*.,leftmargin=*]
\item All curves, polarizations, points, and lifts are algebraic over
$\overline\Q$, as required by Construction 3.4 and Lemma 3.5.
For $M_R$ the smallest abelian subvariety containing $(P,P)$ is
the diagonal $B=E$, since $P$ is nontorsion. If $\psi$ is the
polarization pairing, the bracket on this diagonal is
\[
 [(x,x),(y,y)]=\psi(x,y)-\psi(y,x)=2\psi(x,y).
\]
It is nonzero and spans $\Q(1)$. Hence $Z'=Z=\mathbb G_m$;
there is no deficient or Ribet exception, including for CM.
The pure group is $\mathrm{GL}_2$ without CM and
$\operatorname{Res}_{F/\Q}\mathbb G_m$ for the imaginary
quadratic CM field. Their dimensions are four and two.
This proves the single-motive assertion, including its Lie
bracket and Levi action: the action on the centre is the
polarization multiplier.

\item For the direct sum only the two diagonal central coordinates
occur. If the two elliptic curves are not isogenous, they
cannot be CM conjugates. The pure group is
$\{(g_1,g_2)\in\mathrm{GL}_2^2:\det g_1=\det g_2\}$,
of dimension seven. The point data generate $E\times\bar E$,
and the two brackets are independent, giving $7+4+2=13$.

\item If the curves are isogenous, transport the second polarized
factor through $\phi$. Its polarization pairing becomes
$\psi/d$. The smallest abelian subvariety has dimension $r$.
For $r=2$, independent variation in its two factors makes the
two diagonal brackets span both central coordinates; thus
$Z'=Z=\mathbb G_m^2$. The unipotent group is the product
of two Heisenberg groups, with the same pure group acting
on both standard factors.
For $r=1$, after multiplying the points by a common integer
we may write $Q=aP$, $a\in F^\times$. The bracket image is
the one-dimensional line with proportional central coordinates
\begin{equation}\label{ari:eq:cp-positive-centre}
 (1,c),\qquad c=\frac{a a^\dagger}{d}>0,
\end{equation}
where $\dagger$ is the Rosati involution. On an elliptic curve
$aa^\dagger$ is a positive rational number: a square in the
non-CM case and the positive field norm in the CM case.
Construction 3.4(3) adds a further central direction exactly
when the image of the algebraic lift in the quotient torus
is nontorsion. This is the integer $e$ above.
The unipotent Lie algebra is the Heisenberg algebra, with one
additional central Tate line when $e=1$. Formula
\eqref{ari:eq:cp-MTdim} proves every entry of the table. This
description computes the group even when the lift contributes
a Kummer factor; specifying only the span of $P,Q$ would
not determine that last factor.
\end{enumerate}
\end{proof}





% Compendium-only ari:C.7.5: apparatus: verification/diagnostic record retained with source computation




\begin{proposition}[The period-conjecture implication for the local norm]\label{ari:C.7.4}
\StatementRecord{ari:C.7.4}{proved under hypotheses H}{Historical events: Claude Fable 5.2 (Anthropic), effort high, 17 November 2026: carried as background without fault found. GPT-6 Astra (OpenAI), ChatGPT Pro mode, effort Pro, S6\_09\_PAPERS, 17 November 2026: targeted proof check. Also checked in this writing. Pre-continuation writing events: GPT-6 Astra (OpenAI), ChatGPT Pro mode, effort Pro, S6\_10\_PAPERS, 18 November 2026: targeted proof check. Continuation mathematical and source checks: GPT-6 Astra (OpenAI), ChatGPT Pro mode, effort Pro, S6\_10\_PAPERS continuation, 19 November 2026: computation or formal certificate; independent re-proof; used and rechecked in this run. Continuation metadata events: GPT-6 Astra (OpenAI), ChatGPT Pro mode, effort Pro, S6\_10\_PAPERS continuation, 19 November 2026: used and rechecked in this run; status-vocabulary reconciliation; metadata reconciliation only; no new mathematical check. \emph{Continuation recorded proof scope.} Complete written derivation before reading the target proof, compared in the continuation. The compact-character equation, actual CM conjugate isogeny, Heisenberg central quotient and dimension contradiction were checked; five exact algebraic identities passed in the fresh program \path{code/continuation_root/ari_period_exact.py}. The explicit period hypothesis H remains open.}{\cref{ari:C.7.2}, \cref{ari:C.7.3}, \cref{ari:eq:cp-v}, \cref{ari:eq:cp-positive-centre}; \cite[Construction~3.4, p.~12, Lemma~3.5, p.~13, Conjecture~4.7, p.~18, Proposition~4.12 and Remark~4.13, pp.~19--20]{bertolin2022}}
Let $E/\overline\Q$ be an elliptic curve and let
$P\in E(\overline\Q)$ be nontorsion. Use the local norm $N(P)$,
the algebraic extension $G_P$, and its algebraic fibre coordinate
from \cref{ari:C.7.1}, \cref{ari:C.7.2}.
The hypothesis $H$ is the Relative Semi-abelian Conjecture
\cite[Conjecture~4.7, p.~18]{bertolin2022}, restricted to
$G=G_P\times G_{\bar P}$ and the algebraic lifts used below.
Explicitly, for those lifts and their associated $B_Q,Z(1)$, it gives
\[
 \operatorname{trdeg}_{\overline\Q(\Omega_G)}
 \overline\Q\bigl(\Omega_G,R,\bar R,
                  \widetilde{\log}_G(R,\bar R)\bigr)
 \geq 2\dim B_Q+\dim Z(1),
\]
where $\Omega_G$ is the full de Rham--Betti period matrix,
including the second-kind periods. Under $H$, $N(P)$ is
transcendental. This applies to every algebraic parameter
$t\ne0,1$ with nontorsion $P_t$, including the conductor-$361$
CM point identified below. Hypothesis $H$ remains open for
these points. Grothendieck's period conjecture for the algebraic
$1$-motives in this proof implies $H$
\cite[Proposition~4.12, pp.~19--20]{bertolin2022}.
\end{proposition}

\begin{proof}
\begin{enumerate}[label=\arabic*.,leftmargin=*]
\item Suppose that $N(P)$ is algebraic. The compact-character
identity of \cref{ari:C.7.2} supplies an algebraic point
$R=(P,r)\in G_P(\C)^c$: we take the positive algebraic number
$r=|\wp'(u)|\,|\Delta|^{-1/6}N(P)^2$ in the specified
algebraic fibre coordinate. The conjugate point gives the
corresponding algebraic lift $\bar R$ over $\bar P$.
Set $M_R=[\Z\to G_P]$, $1\mapsto R$, and
$\mathcal M=M_R\oplus\overline{M_R}$.

\item Put $K_0=\overline\Q(\Omega_G)$. Write
$P=[u]$ with periods $\omega_1,\omega_2$, quasi-periods
$\eta_1,\eta_2$, and the logarithmic coordinate $v$ of
\eqref{ari:eq:cp-v}. The full matrix contains
\[
 \kappa_j=\zeta(u)\omega_j-\eta_j u,\qquad j=1,2.
\]
The coefficient matrix for $(\zeta(u),u)$ is
$\left(\begin{smallmatrix}\omega_1&-\eta_1\\
\omega_2&-\eta_2\end{smallmatrix}\right)$; its determinant
is the nonzero Legendre determinant. Hence $u,\zeta(u)$
belong to $K_0$. The same calculation in the conjugate factor
places their conjugates there. This is the explicit elliptic
instance of \cite[Remark~4.13, p.~20]{bertolin2022}.

\item Write $u=a\omega_1+b\omega_2$ with $a,b\in\R$.
Solving this system together with its conjugate gives
\begin{equation}\label{ari:eq:cp-real-linear}
 \eta(u)=
 \frac{(u\bar\omega_2-\bar u\omega_2)\eta_1
       +(\bar u\omega_1-u\bar\omega_1)\eta_2}
      {\omega_1\bar\omega_2-\bar\omega_1\omega_2}.
\end{equation}
The denominator is nonzero by real-linear independence of the
periods. Thus $u\eta(u)$ and its conjugate belong to $K_0$.
The compact-character equation is
$\Re(v-u\zeta(u)+u\eta(u))=0$, so
\begin{equation}\label{ari:eq:cp-relation}
 v+\bar v=u\zeta(u)+\bar u\overline{\zeta(u)}
                -u\eta(u)-\bar u\overline{\eta(u)}\in K_0.
\end{equation}
Each conjugate is a coordinate of the conjugate algebraic
$1$-motive, whose full periods were included in $K_0$.

\item The period torsor of $\mathcal M$ maps to the period
torsor of $[0\to G]$. Projection of the abelian point-extension
variety $B$ to the extension-coordinate variety is an isogeny:
the point coordinates are the same algebraic data as the
extension coordinates under the specified polarizations.
Equivalently, in the product $E\times\bar E$, the projected
point $(P,\bar P)$ is the sum of the two extension parameters
$(P,0)$ and $(0,\bar P)$. The formula for $B_Q$ in
\cite[equation~(4.4), p.~19]{bertolin2022} gives $\dim B_Q=0$.
There is therefore no weight-$-1$ kernel. The full kernel on
the connected motivic groups is the weight-$-2$ vector group
$\operatorname{Lie}Z$, as in
\cite[Construction~3.4 and Lemma~3.5, pp.~12--13]{bertolin2022}.
After fixing the base periods its fibre is an affine torsor
over $K_0$ of dimension $\dim Z$.

\item In the positive Tate basis of each algebraic factor,
the central action changes $(v,\bar v)$ by
$2\pi i(c_1,c_2)$. Replacing a conjugated negative Betti loop
by the positive loop changes its period column, while leaving
the analytic coordinate $\bar v$ fixed. The polarization
transport through $\phi:\bar E\to E$ of degree $d$ is
$\phi^tJ_E\phi=dJ_{\bar E}$. Thus the bracket line computed
in \cref{ari:C.7.3} has coordinates
$(1,aa^\dagger/d)$ with $aa^\dagger/d>0$; in the independent
case the bracket fills both central coordinates. In both cases
$(c_1,c_2)\mapsto c_1+c_2$ is nonzero on the actual central
group. Any extra quotient Tate direction preserves this
nonzero restriction. Consequently \eqref{ari:eq:cp-relation}
cuts a proper affine hyperplane in that torsor and forces the
relative transcendence degree to be at most $\dim Z-1$.
The same doubled construction applies to real data; its
rank-one bracket still has positive coefficient.

\item All noncentral generalized logarithm coordinates are in
$K_0$ by steps 2--3, and $R,\bar R$ are algebraic. Hypothesis
$H$ gives the lower bound $\dim Z$, contradicting step 5.
This proves the stated implication. The source's
Proposition~4.12 supplies it from the period conjecture for
these algebraic $1$-motives. No assertion of $H$ is used
elsewhere to remove its stated open status.
\end{enumerate}
\end{proof}

\subsection{CM theta normalization and further unconditional deductions}

The holomorphic theta normalization proposed for the CM interpolation approach differs from the
norm occurring in the stated problem. We make that distinction explicit,
including the finite-character ambiguity in the multiplication law.

\begin{proposition}[the exact CM multiplication law]\label{ari:C.7.6}
\StatementRecord{ari:C.7.6}{proved}{Historical events: GPT-6 Astra (OpenAI), ChatGPT Pro mode, effort Pro, S6\_09\_PAPERS, 17 November 2026: targeted proof check. Also checked in this writing. Pre-continuation writing events: GPT-6 Astra (OpenAI), ChatGPT Pro mode, effort Pro, S6\_10\_PAPERS, 18 November 2026: targeted proof check. Continuation metadata events: GPT-6 Astra (OpenAI), ChatGPT Pro mode, effort Pro, S6\_10\_PAPERS continuation, 19 November 2026: used and rechecked in this run; status-vocabulary reconciliation; metadata reconciliation only; no new mathematical check. \emph{Pre-continuation recorded proof scope.} Re-derived the CM theta square law by the normalized divisor and leading coefficient, algebraicity of the coefficient s2 and the square-lattice obstruction for the unsquared divisor.}{the hypotheses, notation and definitions stated here; the explicit constructions and calculations in the proof}

Let $E/\overline{\Q}$ have complex multiplication by an order
$\mathcal O$, and use an algebraic invariant differential. Write
$A=\operatorname{area}(\Lambda)$ and
\[
 \eta(z)=s_2z+\frac{\pi}{A}\bar z,\qquad
 \theta(z)=\sigma(z)e^{-s_2z^2/2}.
\]
Then $s_2\in\overline{\Q}$. For every nonzero
$\alpha\in\mathcal O$, with $d=\alpha\bar\alpha$, there is
$F_\alpha\in\overline{\Q}(E)$ such that
\begin{equation}\label{ari:eq:ct-CMmult}
 \theta(\alpha z)^2=F_\alpha(\exp_E z)\theta(z)^{2d}.
\end{equation}
Its divisor and leading Laurent coefficient are
\[
 \operatorname{div}(F_\alpha)=2[\alpha]^*(O)-2d(O),\qquad
 F_\alpha(\exp_Ez)=\alpha^2z^{2-2d}(1+O(z^2)).
\]
Thus the unsquared quotient is an algebraic function on $E$, but
it need not be a rational function on $E$. For a nontorsion
$P=\exp_Eu\in E(\overline{\Q})$, algebraicity of $\theta(u)$
would imply algebraicity of $\theta(\alpha u)$ for all
$\alpha\in\mathcal O$. The additive group $\mathcal O u$ has
rank two, not more than four.
\end{proposition}


\begin{proof}
\begin{enumerate}[label=\arabic*.,leftmargin=*]
\item The Legendre relation gives the displayed real-linear decomposition
of $\eta$. The sigma automorphy formula gives
\[
 \theta(z+\omega)=\epsilon(\omega)
  \exp\left(\frac{\pi\bar\omega}{A}
                         (z+\omega/2)\right)\theta(z),
 \qquad \epsilon(\omega)\in\{1,-1\}.
\]
For $\alpha\Lambda\subset\Lambda$, the exponential multipliers
of $\theta(\alpha z)/\theta(z)^d$ cancel because
$\alpha\bar\alpha=d$. Its remaining multiplier is a sign.
Squaring therefore gives a meromorphic function on $E$, with the
stated divisor. This divisor is defined over $\overline{\Q}$.
The existence of a rational function with this divisor over $\C$
implies its existence over $\overline{\Q}$: the relevant
one-dimensional Riemann--Roch space is obtained by extension of
scalars. Normalize its leading Laurent coefficient, in the formal
parameter of the algebraic differential, to $\alpha^2$. This
identifies it with the squared analytic quotient. The action of
$\alpha$ on this differential is algebraic, so $\alpha^2$ is an
allowed normalization.

\item Choose a nonreal $\alpha\in\mathcal O$. Since the normalized
sigma series is $\sigma(z)=z+O(z^5)$ with algebraic Laurent
coefficients, the next coefficient of the squared quotient is
\[
 \alpha^2z^{2-2d}
       \{1-s_2(\alpha^2-d)z^2+O(z^4)\}.
\]
The coefficients of the rational function in this algebraic formal
parameter are algebraic. As $\alpha^2-d\ne0$, this proves
$s_2\in\overline{\Q}$ without an unnormalized appeal to a CM
weight-two form. It also proves that the functions in
\eqref{ari:eq:ct-CMmult} are algebraic over the chosen field.
At nontorsion $u$, none of the arguments $\alpha u$ is a lattice
point; evaluating the identity proves the algebraicity implication.
The injective map $\mathcal O\to\C$, $\alpha\mapsto\alpha u$,
proves the rank assertion.

\item The qualification about rationality is necessary. For a square
lattice, $s_2=0$ and $\alpha=1+i$ give
$F_{1+i}=2i\wp$. The unsquared quotient has a simple zero at
the nonzero point in $\ker(1+i)$ and a simple pole at $O$,
which is not a principal divisor on $E$; its square is the
rational function just displayed.
\end{enumerate}
\end{proof}

\begin{cited}[with derivation: two unconditional normalized values]\label{ari:C.7.7}
\StatementRecord{ari:C.7.7}{proved}{Historical events: GPT-6 Astra (OpenAI), ChatGPT Pro mode, effort Pro, S6\_09\_PAPERS, 17 November 2026: targeted proof check. This writing records metadata reconciliation only; no new mathematical check. Continuation metadata events: GPT-6 Astra (OpenAI), ChatGPT Pro mode, effort Pro, S6\_10\_PAPERS continuation, 19 November 2026: used and rechecked in this run; status-vocabulary reconciliation; metadata reconciliation only; no new mathematical check. \emph{Pre-continuation recorded proof scope.} Current-writing verification remains in progress. The earlier checks named above retain their recorded scope.}{the hypotheses, notation and definitions stated here; \cref{ari:C.7.6}; \citet[Theorem~1 and Corollary~1]{waldschmidt1979sigma}; \citet[Theorem~6.1]{huber2022}; \CompendiumRef{I}{app:rederive:ari-C.7.7}; \citet[final slide, comment (c)]{bertrand2017}; \citet[p.~15-03]{waldschmidt1988theta}; \cref{ari:C.7.1}; \cref{ari:eq:ct-Gaussian}; \citet[Corollary~2]{reyssat1982weierstrass}; \citet[p.~15-05]{waldschmidt1988theta}; \citet[Theorem 3.1.2]{bost2009}; \cref{ari:C.7.2}; \citet[Theorem 1.4 and Corollary 1.6]{nguyen2023}; \citet[Theorems 1.3 and 2.3]{xie2024algebraicity}; \cref{ari:C.7.9}}

Let $E/\overline{\Q}$ have an algebraic invariant differential,
and let $u$ be an elliptic logarithm of an algebraic nontorsion
point $P$. Then
\[
 \sigma(u)\exp\left(-\frac12u\zeta(u)\right)
\]
is transcendental. If $E$ has CM and $s_2$ is normalized as in
\ref{ari:C.7.6}, then $\zeta(u)-s_2u$ is also transcendental. The first
assertion concerns the displayed holomorphic normalization;
it does not assert transcendence of either $N(P)$ or
$\theta(u)=\sigma(u)e^{-s_2u^2/2}$.
\end{cited}
The expanded source derivation is retained in the proof compendium, Paper~I, label \nolinkurl{ari:C.7.7}.




The stated normalized values follow from \citet[Theorem~1 and Corollary~1]{waldschmidt1979sigma} and the analytic subgroup theorem in \citet[Theorem~6.1]{huber2022}. A complete re-derivation is given in \CompendiumRef{I}{app:rederive:ari-C.7.7}.


\paragraph{A theta criterion requiring its original normalization.}
An earlier draft applied the pointwise criterion printed in
\citet[final slide, comment (c)]{bertrand2017} to
$g(u,v)-s_2uv$, where
$g(u,v)=\log(\sigma(u+v)/(\sigma(u)\sigma(v)))$.
It inferred transcendence of $\theta(u)$ by taking $v=u$.
We withdraw that inference pending verification of the original
lift conventions in \citet[Corollary 3]{bertrand1984}.
Indeed replacing $v$ by $v+\omega$, $\omega\in\Lambda$,
changes $g(u,v)-s_2uv$, modulo $2\pi i\Z$, by
\[
 (\eta(\omega)-s_2\omega)u=\frac{\pi\bar\omega u}{A},
\]
whereas the condition $N\exp_Ev=f(\exp_Eu)$ with $f$
Rosati-antisymmetric depends only on the elliptic points.
Thus the two-sided statement cannot be used with unrestricted
lifts without checking its normalization. The author's slide
and the original paper's first two pages were accessible;
the full original Corollary 3 was not. We use neither direction
of that criterion here. The withdrawn text and the lift audit
are retained in the records.

A subsequent source check resolves the normalization of the
one-variable holomorphic assertion itself:
\citet[p.~15-03]{waldschmidt1988theta} explicitly states that
$\sigma(u)e^{-s_2u^2/2}$ is transcendental when the invariants
are algebraic, $E$ has CM, and $\exp_Eu$ is an algebraic
nontorsion point.  That source defines $s_2$ by the regularized
lattice sum and attributes the result to Bertrand's 1982
paper \emph{Valeurs de fonctions th\^eta et hauteurs p-adiques}.
We have checked this complete normalized statement in
Waldschmidt's published exposition; the full original 1982
proof has not been retrieved.  This source clarification
does not validate the unrestricted-lift inference withdrawn
above, and the theorem concerns the holomorphic value rather
than the metric norm.  \ref{ari:C.7.7} retains the two assertions whose
proofs have been fully checked here.

\paragraph{The archimedean term still requiring control.}
The exact comparison is
\begin{equation}\label{ari:eq:ct-Gaussian}
 N(P)=|\Delta|^{1/12}|\theta(u)|
                         \exp\left(-\frac{\pi|u|^2}{2A}\right).
\end{equation}
It follows from \ref{ari:C.7.1} and the real-linear decomposition of
$\eta$. The normalization proved transcendental in \ref{ari:C.7.7} instead
gives the exact identity
\[
 N(P)=|\Delta|^{1/12}
 \left|\sigma(u)e^{-u\zeta(u)/2}\right|
 \exp\left(\frac12\Re\{u(\zeta(u)-\eta(u))\}\right).
\]
Neither formula excludes algebraicity of $N(P)$: transcendence
of one complex factor does not control its absolute value or
its product with the remaining exponential.
The generic finite-order interpolation approach
supplies only the rank-two set $\mathcal O u$; it does not
meet the proposed rank-greater-than-four criterion. The
derivative approach would need the algebraicity of
$\theta'(u)/\theta(u)=\zeta(u)-s_2u$, which \ref{ari:C.7.7} excludes.
Replacing the norm by the holomorphic theta value would
therefore change the problem precisely by the factor in
\eqref{ari:eq:ct-Gaussian}.

The related independence theorem of
\citet[Corollary~2]{reyssat1982weierstrass} says that two of
$e^{i\pi u/\omega_1}$, $s(u)$ and $s'(u)$ are algebraically
independent, where $s(z)=\sigma(z)e^{-\eta_1z^2/(2\omega_1)}$.
For real $u$, hypothetical algebraicity of $N(P)$ makes
$s(u)$ algebraic, so the conclusion instead forces independence
of $e^{i\pi u/\omega_1}$ and
$\zeta(u)-\eta_1u/\omega_1$; it yields no contradiction.
The periodic-functional-equation theorem reported in
\citet[p.~15-05]{waldschmidt1988theta} assumes algebraic
$q=e^{i\pi\tau}$, whereas this $q$ is transcendental for
imaginary quadratic $\tau$ by Gelfond--Schneider.  These
specific checked applications therefore do not decide the
metric norm; their source scopes and the larger CM
independence tuples are recorded separately.

We also tested an additive-extension route.
\citet[Theorem 3.1.2]{bost2009} proves that an algebraic point
of the universal vector extension of an abelian variety which
lies in its maximal compact subgroup at one complex embedding
is torsion. The theorem's hypotheses require an algebraic
point of the vector extension, equivalently an algebraic line
bundle with an algebraic connection having unitary monodromy.
\ref{ari:C.7.2} supplies only an algebraic point in one fibre of a
multiplicative extension. It does not supply that algebraic
connection. Differentiating a metric value at one algebraic
point cannot establish algebraicity of the connection there
or globally. This is the exact missing implication in that
attempt.

Another route uses the Latt\`es map $f_E$ defined by
$x(2Q)=f_E(x(Q))$. Its canonical height is
$\widehat h_{f_E}(x(P))=2\widehat h_O(P)$. For rational
parameters C.6 therefore makes algebraicity of $N(P)$
equivalent to algebraicity of
$\exp(\widehat h_{f_E}(x(P)))$.
\citet[Theorem 1.4 and Corollary 1.6]{nguyen2023} settle the
corresponding multiplicative-height problem for nonintegrable
polynomials using B\"ottcher coordinates in the basin of
infinity. They do not apply to $f_E$. Indeed an elliptic
duplication Latt\`es map has no exceptional point: the inverse
images of any point contain the images of all successive
division points of an elliptic point, an infinite set. A
polynomial has the exceptional point $\infty$, so no conjugate
of $f_E$ is polynomial. The Julia set is the whole sphere
and the required basin of infinity is absent. This checks
the failed hypothesis, rather than treating a theorem for
polynomials as a theorem for all rational maps.

We also checked the arithmetic algebraicity criterion of
\citet[Theorems 1.3 and 2.3]{xie2024algebraicity}.
Besides sufficiently rapid height-normalized approximation by
points in a fixed number field, it requires the analytic germ
to be formally defined over that field. This hypothesis fails
for the compact leaf in \ref{ari:C.7.9}: its tangent has the proved
transcendental slope
$c=\zeta(u)-\eta_1u/\omega_1$.
A smooth formal germ defined over $\overline{\Q}$ has an
algebraic tangent; algebraic translations, coordinate changes
and isogenies preserve this obstruction. Assuming $N(P)$
algebraic supplies an algebraic point on the leaf and its
integer multiples, but not an algebraic formal germ. Attempting
to force convergence by successive division also loses the
fixed-number-field hypothesis: if $[2]^nQ_n=Q$ with all
$Q_n\in E(K)$ for one number field, then
$\widehat h_O(Q_n)=4^{-n}\widehat h_O(Q)>0$.
The $Q_n$ are distinct with bounded ample height, contrary to
Northcott finiteness. The full hypothesis audit and the opened
preprint are retained in the records.

\paragraph{A current source whose theorem text remains unchecked.}
The publisher's abstract for \citet{pham2025sigma} announces
complex and $p$-adic linear-independence results for semiabelian
varieties and applications to Weierstrass sigma functions.
We opened its metadata, abstract and reference list, but both
the full-HTML link and the publisher's citation-PDF link returned
the abstract/access page. No accessible author manuscript was
located. Consequently we could not check any theorem statement
or its hypotheses. This is a potentially relevant unverified
source: we make no claim either that it settles the Gaussian
norm question or that its results are inapplicable to it.



\subsection{The precise nonarchimedean comparison}
\begin{cited}[with derivation: a nonzero ordinary height]\label{ari:C.7.8}
\StatementRecord{ari:C.7.8}{proved}{Historical events: GPT-6 Astra (OpenAI), ChatGPT Pro mode, effort Pro, S6\_09\_PAPERS, 17 November 2026: targeted proof check. This writing records metadata reconciliation only; no new mathematical check. Continuation metadata events: GPT-6 Astra (OpenAI), ChatGPT Pro mode, effort Pro, S6\_10\_PAPERS continuation, 19 November 2026: used and rechecked in this run; status-vocabulary reconciliation; metadata reconciliation only; no new mathematical check. \emph{Pre-continuation recorded proof scope.} Current-writing verification remains in progress. The earlier checks named above retain their recorded scope.}{the hypotheses, notation and definitions stated here; \citet[Appendix A]{agboola2006}; \citet{agboola2006}; \CompendiumRef{I}{app:rederive:ari-C.7.8}; \cref{ari:C.7.6}; \citet[Theorem 2.1]{bertrand1995}; \cref{ari:C.7.3}; \cref{ari:C.7.4}; \citet[Theorem 3.1]{bertrand1995}; \citet{bost2009}; \cref{ari:C.7.2}}

Let $E/\Q$ have complex multiplication by the full ring of integers
of an imaginary quadratic field $K$ of class number one. Assume
that its functional-equation sign is $-1$, and let $p>3$ be a
prime of good ordinary reduction, $p\mathcal O_K=\mathfrak p
\mathfrak p^*$. For a nontorsion $P\in E(\Q)$ the
$\mathfrak p$-adic height and the cyclotomic $p$-adic height
of $P$ over $K$, in the normalization of
\citet[Appendix A]{agboola2006}, are nonzero.
In particular this applies to
\[
 E:y^2+y=x^3-38x+90,\qquad P=(0,9),\qquad p=5.
\]
\end{cited}
The expanded source derivation is retained in the proof compendium, Paper~I, label \nolinkurl{ari:C.7.8}.




The nonvanishing theorem is reproduced in Rubin's Appendix~A, Theorem~A.1 and Lemma~A.2 of \citet{agboola2006}; the exact hypotheses and application are retained. A complete re-derivation is given in \CompendiumRef{I}{app:rederive:ari-C.7.8}.


The comparison explains the missing step in an attempted
archimedean adaptation. The ordinary $p$-adic sigma construction
uses a holomorphic quadratic correction. In the CM normalization
the corresponding complex holomorphic function is $\theta$
of \ref{ari:C.7.6}, whereas the real metric also contains
$\exp(-\pi|u|^2/(2A))$. Theorem A.1 supplies no identity
between the ordinary $p$-adic height and that real factor.
Transferring the conclusion would require precisely such a
new argument; replacing the real factor by the holomorphic
CM correction does not preserve the norm being asked about.

There is also a definite limitation on the maximal-compact
approach. \citet[Theorem 2.1]{bertrand1995} constructs
nontrivial multiplicative extensions of polarized abelian
varieties with nontorsion algebraic points in their maximal
real compact subgroup. Its hypotheses include an endomorphism
not fixed by the Rosati involution and a nontorsion rational
point over a real number field. The construction uses a base
point related to the extension class by a Rosati-antisymmetric
endomorphism. It cannot be replaced by a theorem saying that
every algebraic compact point of every multiplicative
extension is torsion. For our diagonal class, an equality
$NP=f(P)$ with $N>0$ and $f+f^\dagger=0$ would force
$P$ torsion: $N-f$ has no zero eigenvalue, since
$(N-f)^\dagger(N-f)=N^2-f^2$ is positive. In the elliptic
case this is simply a nonzero isogeny. The diagonal positivity
used in \ref{ari:C.7.3}--\ref{ari:C.7.4} is therefore essential, but these
counterexamples neither prove nor refute the desired
unconditional diagonal statement. The additive result in
\citet[Theorem 3.1]{bertrand1995}, and its higher-dimensional
extension in \citet{bost2009}, requires a point in a vector
extension, a hypothesis not furnished by \ref{ari:C.7.2}.


\subsection{The fixed semiabelian extension and the diagonal subgroup attempt}

\begin{proposition}[the endomorphism and slope obstructions]\label{ari:C.7.9}
\StatementRecord{ari:C.7.9}{proved}{Historical events: Claude Fable 5.2 (Anthropic), effort high, 15 November 2026: carried as background without fault found. GPT-6 Astra (OpenAI), ChatGPT Pro mode, effort Pro, S6\_09\_PAPERS, 17 November 2026: targeted proof check. Also checked in this writing. Pre-continuation writing events: GPT-6 Astra (OpenAI), ChatGPT Pro mode, effort Pro, S6\_10\_PAPERS, 18 November 2026: targeted proof check. Continuation metadata events: GPT-6 Astra (OpenAI), ChatGPT Pro mode, effort Pro, S6\_10\_PAPERS continuation, 19 November 2026: used and rechecked in this run; status-vocabulary reconciliation; metadata reconciliation only; no new mathematical check. \emph{Pre-continuation recorded proof scope.} Read and re-derived the full extension-endomorphism calculation, doubling identity, real compact lift and transcendental slope proof. The exact analytic-subgroup primary locator remains to be matched; receiving text is being stated for the actual compact point. Matched the exact analytic subgroup input to Huber-Wuestholz Theorem6.1 printedp43 and the divided-period reduction pp43-44. The actual compact point and algebraicity equivalence are now stated without an unresolved algebraicity premise.}{the hypotheses, notation and definitions stated here; \cref{ari:C.7.1}; \cref{ari:eq:cd-compatibility}; \citep[Theorem 6.1]{huber2022}; \cref{ari:C.6.R1}; \cref{ari:C.7.2}}

Let $E/\overline\Q$ be an elliptic curve, let
$q\in E^\vee(\overline\Q)$ be nontorsion, and let
\[
 1\longrightarrow\mathbb G_m\longrightarrow G_q
       \longrightarrow E\longrightarrow0
\]
be the extension of class $q$. Then
$\operatorname{End}_{\overline\Q}(G_q)=\Z$, including when
$E$ has complex multiplication. In particular a noninteger
CM endomorphism of $E$ does not lift to an endomorphism of this
fixed extension.

For
\[
 E:y^2+y=x^3-38x+90,\qquad P=(0,9),\qquad Q=2P=(4,-2),
\]
one has
\[
 N(Q)=19^{1/4}N(P)^4.
\]
Use the differential $dx/(2y+1)$, choose the positive real
period $\omega_1$, and choose the real logarithm
$u\in(0,\omega_1/2)$ of $Q$. Put
\begin{equation}\label{ari:eq:cd-slope}
 c=\zeta(u)-\frac{\eta_1}{\omega_1}u,
 \qquad \eta_1=\eta(\omega_1).
\end{equation}
Then $c$ is a real transcendental number. The actual point
\[
 R_N=(Q,3\,19^{-1/2}N(Q)^2)\in G_Q(\C)
\]
is nontorsion, belongs to the maximal compact subgroup, and has
the real logarithm $(u,cu)$. Its algebraicity is equivalent to
the algebraicity of $N(P)$. In the algebraic logarithm coordinates
of \ref{ari:C.7.1}, the line $\ell_c=\C(1,c)$ has transcendental
slope. The proof establishes this slope statement from a period
of $G_Q$ without deciding the algebraicity of $R_N$.
\end{proposition}


\begin{proof}
\begin{enumerate}[label=\arabic*.,leftmargin=*]
\item The torus is the maximal connected affine subgroup of $G_q$.
Consequently every endomorphism of $G_q$ induces multiplication
by some $n\in\Z$ on $\mathbb G_m$ and an endomorphism
$\alpha$ on $E$. Functoriality of extension classes gives
\begin{equation}\label{ari:eq:cd-compatibility}
 \alpha^*q=nq\quad\text{in }E^\vee.
\end{equation}
Here $\alpha^*$ is the dual endomorphism, not $\alpha$
acting on an arbitrarily identified copy of $E$. To make the
convention explicit, choose a principal polarization
$\lambda:E\to E^\vee$, put $Q_0=\lambda^{-1}q$, and
write $\alpha^\dagger=\lambda^{-1}\alpha^*\lambda$.
Equation \eqref{ari:eq:cd-compatibility} says
$(\alpha^\dagger-[n])Q_0=O$. A nonzero endomorphism of an
elliptic curve is an isogeny and has finite kernel. Since
$Q_0$ is nontorsion, $\alpha^\dagger=[n]$, hence
$\alpha=[n]$ as well.

\item The difference between the given endomorphism and $[n]$ has
image in the torus and is trivial on the torus. It therefore
factors through an algebraic-group homomorphism
$E\to\mathbb G_m$. This homomorphism is zero, since a
morphism from the proper connected curve $E$ to the affine
group $\mathbb G_m$ is constant and a group homomorphism
preserves the identity. Thus there is no additional ambiguity
from $\operatorname{Hom}(E,\mathbb G_m)$, and the
endomorphism is exactly $[n]$. Conversely all integer
multiplications exist. This proves the first assertion.

\item For the concrete curve, the tangent slope at $P$ is $-2$,
giving $2P=(4,-2)$. Its discriminant is $-19^3$ and
$\psi_2(P)=2\cdot9+1=19$. The division identity in C.4
therefore gives the asserted formula for $N(Q)$, so replacing
$P$ by $Q$ preserves algebraicity of the norm. This doubling
also places the base point in the identity component of the
real elliptic group without any component assumption. In the
chosen differential the Weierstrass coordinates are
$(\wp,\wp')=(x,2y+1)$, with $g_2=152$, $g_3=-361$.
The negative discriminant means that the real cubic has one
real root. Along $(0,\omega_1/2)$, $\wp$ decreases from
$+\infty$ to that root and $\wp'<0$. Since
$\wp'(u)=2(-2)+1=-3$, the indicated real logarithm of $Q$
exists.

\item The period vectors of $G_Q$ in \ref{ari:C.7.1} include
\[
 (\omega_1,\zeta(u)\omega_1-\eta_1u)
       =\omega_1(1,c).
\]
These are algebraic logarithm coordinates: $dz$ and the
invariant differential $d\log r-df_u/f_u$ are algebraic,
because $df_u/f_u$ is the algebraic third-kind differential
displayed in \ref{ari:C.7.1}. Real invariants and real $u,\omega_1$
make $c$ real.

\item Suppose $c$ were algebraic. Then $\ell_c$ would be an
algebraic Lie subspace. Dividing the nonzero period
$\omega_1(1,c)$ by a sufficiently large positive integer
gives a vector in $\ell_c$ whose exponential is a nonzero
torsion point of $G_Q$, hence an algebraic point. The
analytic subgroup theorem
\citep[Theorem 6.1]{huber2022} supplies a positive-dimensional
connected algebraic subgroup with Lie algebra contained in
$\ell_c$. Its dimension is one, and its projection to $E$
is nontrivial. Such a subgroup must be an elliptic curve
$E'$ projecting isogenously to $E$: a connected
one-dimensional affine algebraic group cannot map
nontrivially to an elliptic curve. Its inclusion splits the
pullback of $G_Q$ to $E'$. The dual of that isogeny would
therefore kill the extension class $\lambda(Q)$, forcing
$Q$ torsion. This contradicts \ref{ari:C.6.R1}. We conclude that $c$
is transcendental. In particular the line is not defined over
$\overline\Q$. This argument uses the period vector alone;
it does not assume algebraicity of $N(P)$.

\item Define the positive real number
\[
 r=3\,19^{-1/2}N(Q)^2.
\]
Proposition~\ref{ari:C.7.2} places the actual point $R_N=(Q,r)$
in the maximal compact subgroup of $G_Q$. Since $Q$ is algebraic,
$R_N$ is algebraic exactly when $r$ is algebraic, equivalently
when $N(Q)$, and hence $N(P)$, is algebraic. The chosen real
logarithm has
\[
 f_u(u)=-\wp'(u)\sigma(u)^2e^{-u\zeta(u)}
        =3\sigma(u)^2e^{-u\zeta(u)}>0.
\]
Consequently $v=\log(r/f_u(u))$ is real and
$(u,v)$ is a logarithm of $R_N$, with no imaginary logarithm
ambiguity. Substituting this real logarithm into the exact
compactness equation $\chi=0$ in \ref{ari:C.7.2} gives
\[
 v=\zeta(u)u-u\eta(u)
   =u\left(\zeta(u)-\frac{\eta_1}{\omega_1}u\right)=cu.
\]
Its projection $Q$ is nontorsion, so $R_N$ is nontorsion.
The input missing from this proposed use of the analytic
subgroup theorem is a proper Lie subspace defined over
$\overline\Q$; the actual compact slope line is proved
not to have that property. The CM decomposition makes the
same obstruction visible as
\[
 c=\zeta(u)-s_2u-\frac{\pi}{A}u.
\]
The last term is the derivative of the real Gaussian
correction. It cannot be discarded while retaining the
compactness equation.
\end{enumerate}
\end{proof}

\paragraph{What functional transcendence gives in this construction.}
The functional theorem can be applied fully, but its conclusion
does not contradict this hypothetical algebraic compact point.
Indeed consider
\[
 f:\C\longrightarrow\operatorname{Lie}(G_Q),\qquad
 f(z)=(z,cz).
\]
No proper connected algebraic subgroup of $G_Q$ projects
onto $E$: as in the proof above, such a subgroup would split
the nontorsion extension after an isogeny. Every proper
connected subgroup is therefore contained in the torus,
whereas $\exp_{G_Q}f$ has nonconstant elliptic projection.
Thus $f$ is $G_Q$-nondegenerate in Definition 1.2 of
\citet{noguchi2024}. Their Theorem 1.3 gives
\[
 \dim_{\C}\overline{\{(\exp_{G_Q}f(z),f(z)):z\in\C\}}^{\rm Zar}
       \ge \dim G_Q+1=3.
\]
The graph is contained in $G_Q\times\ell_c$, of dimension
three, so equality holds. This is a statement over $\C$
about a family of meromorphic functions. It remains valid
whether or not the single value $\exp_{G_Q}f(u)$ is
algebraic. Under the norm-algebraicity hypothesis the integer
values $\exp_{G_Q}f(nu)=[n]R$ are algebraic, but they do
not give an algebraic relation holding on the entire graph.
There is no contradiction with its full dimension.
One would need an additional arithmetic specialization
argument, or a way to obtain a further functional relation
from that single metric value. Neither is supplied by the
checked functional theorem. The endomorphism calculation
above also prevents obtaining extra CM multiplication
directions inside this fixed $G_Q$ by lifting the
noninteger endomorphisms of $E$.


\subsection{A CM-stable enlargement and its Gaussian eigencoordinate}

The fixed-extension calculation in \ref{ari:C.7.9} does not prevent
enlarging the torus. We test that alternative explicitly.
Retain its curve, its point $Q=2P$, and its real logarithm
$u$. Write $\beta=\sqrt{-19}=i\sqrt{19}$, $d=19$,
$D=\zeta(u)-s_2u$, and $k=\pi/A$. All endomorphisms and
extensions in what follows are defined over $\overline\Q$.

\begin{proposition}[the enlarged CM action does not isolate the Gaussian]\label{ari:C.7.10}
\StatementRecord{ari:C.7.10}{proved}{Historical events: GPT-6 Astra (OpenAI), ChatGPT Pro mode, effort Pro, S6\_09\_PAPERS, 17 November 2026: targeted proof check. Also checked in this writing. Pre-continuation writing events: GPT-6 Astra (OpenAI), ChatGPT Pro mode, effort Pro, S6\_10\_PAPERS, 18 November 2026: computation or formal certificate; targeted proof check. Continuation metadata events: GPT-6 Astra (OpenAI), ChatGPT Pro mode, effort Pro, S6\_10\_PAPERS continuation, 19 November 2026: used and rechecked in this run; status-vocabulary reconciliation; metadata reconciliation only; no new mathematical check. \emph{Pre-continuation recorded proof scope.} Read and re-derived the full enlarged CM extension, exact differential matrix and projector gap, plus the complete conductor-361 polynomial witness after the old proof boundary. Exact polynomial replay is recorded separately. Read and freshly replayed unchanged C7\_cm\_explicit.gp. Both polynomial residuals are zero, all displayed gcds are one, exact s2/a/r2 values agree, and analytic comparisons differ by less than 6e-114. Exact identities prove the algebraic witness; decimal comparisons are diagnostic. Matched the toric Lie-line application to Huber-Wuestholz Theorem6.1p43. Receiving revision preserves the full polynomial witness inside the numbered proof and states the actual paired compact point.}{the hypotheses, notation and definitions stated here; \cref{ari:C.7.9}; \cref{ari:C.7.6}; \cref{ari:eq:ce-zeta}}

Let
\[
 \mathcal G=G_Q\mathbin{\times_E}G_{\beta Q},\qquad
 1\longrightarrow\mathbb G_m^2\longrightarrow\mathcal G
       \longrightarrow E\longrightarrow0.
\]
There is an endomorphism $B$ of $\mathcal G$ inducing
$\beta$ on $E$ and
$(t_1,t_2)\mapsto(t_2^{-1},t_1^d)$ on the torus.
It satisfies $B^2=[-d]$. Moreover
\[
 \operatorname{Hom}(\mathcal G,\mathbb G_m)=0,
 \qquad
 \operatorname{End}(\mathcal G)\otimes\Q=\Q(\beta).
\]
Thus $\mathcal G$ has no algebraic $\mathbb G_m$ quotient,
including after isogeny.

There is an algebraic compact point in the second extension
above $Q$ whose logarithm is
\[
 (u,v_2),\qquad v_2=-\beta uD+au,
 \qquad
 a=\frac1{2\beta}
       \frac{d}{du}\log F_\beta(Q)\in\overline\Q.
\]
Here the derivative means differentiation of the rational
function $F_\beta(\exp_E z)$ with respect to the algebraic
invariant differential, evaluated at $z=u$.
Pair this point with the actual compact point $R_N$ of
\ref{ari:C.7.9}. The resulting compact point of $\mathcal G(\C)$
has logarithm
\[
 U=(u,v_1,v_2),\qquad v_1=uD-ku^2.
\]
Its algebraicity is equivalent to the algebraicity of $N(P)$.
In these algebraic logarithm coordinates the differential of
$B$ and its $-\beta$ eigenprojection are
\[
 dB=\begin{pmatrix}
 \beta&0&0\\ a&0&-1\\ \beta a&d&0
 \end{pmatrix},\qquad
 \Pi_-:=\frac12\left(I-\frac{dB}{\beta}\right),
\]
\[
 \Pi_-U=-\frac{\pi u^2}{2A}(0,1,\beta).
\]
The first torus coordinate of $\exp(\Pi_-U)$ is exactly
the Gaussian $\exp(-\pi u^2/(2A))$. The projector is not
the differential of an algebraic quasi-endomorphism. No
sequence of rational polynomials in $B$, even allowing
inverses of isogenies, can converge to it as complex linear
operators. At least one of the two torus coordinates of
$\exp(\Pi_-U)$ is transcendental; the proof does not decide the
algebraicity of its first coordinate alone.
\end{proposition}


\begin{proof}
\begin{enumerate}[label=\arabic*.,leftmargin=*]
\item Use the principal polarization to identify the two extension
classes with $Q$ and $\beta Q$. They are linearly independent
over $\Q$ modulo torsion: any nonzero $m+n\beta$ is an
isogeny, so $(m+n\beta)Q$ cannot be torsion. A character
of $\mathbb G_m^2$ given by $(m,n)\in\Z^2$ extends to
$\mathcal G$ exactly when the pushout extension has zero
class, namely when $mQ+n\beta Q=O$. Hence the restriction
of every homomorphism $\mathcal G\to\mathbb G_m$ to
the torus is zero. Such a homomorphism then factors through
$E$ and is zero. This proves the asserted absence of
characters. An isogenous semiabelian variety cannot acquire
a nonzero character, because composing it with the isogeny
would give a nonzero character of $\mathcal G$.

\item The Rosati adjoint of $\beta$ is $-\beta$. Pullback by
$\beta$ sends the ordered extension classes to
$(-\beta Q,dQ)$, exactly as does pushout by the displayed
torus matrix. Equality of these extension classes provides
the endomorphism $B$. It is unique because
$\operatorname{Hom}(E,\mathbb G_m^2)=0$. Its square
induces $[-d]$ on the torus and on $E$, so the same
uniqueness gives $B^2=[-d]$. The map
\[
 \operatorname{End}(\mathcal G)\otimes\Q
     \longrightarrow\operatorname{End}(E)\otimes\Q
          =\Q(\beta)
\]
is injective: an endomorphism inducing zero on $E$ has image
in the torus, and the absence of characters makes it zero.
Its image contains $\beta$, hence equals $\Q(\beta)$.

\item Next construct the compact point without assuming anything
about $N(P)$. The identities in \ref{ari:C.7.6}, applied at $\beta$
and $1+\beta$, give the algebraic nonzero number
\begin{equation}\label{ari:eq:ce-ribet-r}
 r_2=\frac{\theta((1+\beta)u)}{\theta(u)\theta(\beta u)},
 \qquad r_2^2=\frac{F_{1+\beta}(Q)}{F_\beta(Q)}.
\end{equation}
Indeed $(1+\beta)(1+\bar\beta)=1+d$, so the powers
of $\theta(u)$ cancel. None of the arguments is a lattice
point since $Q$ is nontorsion. Written in sigma functions,
\[
 r_2=\frac{\sigma((1+\beta)u)}{\sigma(u)\sigma(\beta u)}
                         e^{-s_2\beta u^2}.
\]
The point $(Q,r_2)$ in the rational trivialization of
$G_{\beta Q}$ therefore has logarithm
\[
 (u,v_2),\qquad
 v_2=u\{\zeta(\beta u)-s_2\beta u\}.
\]
Differentiating \ref{ari:C.7.6} shows that
\begin{equation}\label{ari:eq:ce-zeta}
 \zeta(\beta u)-s_2\beta u=-\beta D+a,
 \qquad a\in\overline\Q,
\end{equation}
with the value of $a$ in the statement. The compactness
character of this extension vanishes, since
\[
 \Re\{v_2-\zeta(\beta u)u+\beta u\eta(u)\}
    =\Re\left(\frac{\beta\pi|u|^2}{A}\right)=0.
\]
This also gives a direct, fully normalized construction of
the antisymmetric compact lift; no unrestricted theta
criterion is being used.

\item For completeness, the matrix $dB$ can be recovered without
choosing a transcendental change of coordinates. Its torus
block is the stated integer matrix and its first entry is
$\beta$, so initially write its first column as
$(\beta,a_1,a_2)^t$, with $a_1,a_2$ algebraic.
The identity $B^2=[-d]$ gives $a_2=\beta a_1$.
For $\omega\in\Lambda$, set
\[
 \kappa_1(\omega)=\zeta(u)\omega-\eta(\omega)u,
 \quad
 \kappa_2(\omega)=\zeta(\beta u)\omega-\eta(\omega)\beta u.
\]
Equation \eqref{ari:eq:ce-zeta} and the real-linear formula for
$\eta$ imply
\[
 \kappa_1(\beta\omega)=a\omega-\kappa_2(\omega).
\]
Comparing the second coordinates of the corresponding period
vectors under $B$ gives
$(a_1-a)\omega_j\in2\pi i\Z$ for $j=1,2$.
Since $\omega_2/\omega_1$ is nonreal, both integers must
be zero, so $a_1=a$. This proves the matrix formula.

\item The actual first compact logarithm in \ref{ari:C.7.9} is
$v_1=uD-ku^2$. Pairing the two compact lifts over the same base
point gives $U$ with no algebraicity assumption. The second lift
is algebraic, so the paired point is algebraic exactly when its
first lift is algebraic.
Substitution into the displayed matrix yields
\[
 v_2+\beta v_1-au=-\beta ku^2,
 \qquad
 \Pi_-U=-\frac{ku^2}{2}(0,1,\beta),
\]
which proves the exact Gaussian identity.

\item The two eigenspaces of $dB$ have eigenvalues $\beta$ and
$-\beta$. Every algebraic quasi-endomorphism has differential
$a_0I+b_0dB$ for rational $a_0,b_0$, with eigenvalues
$a_0+b_0\beta$ and $a_0-b_0\beta$. The projector
$\Pi_-$ would require these values to be $0$ and $1$,
which is impossible. There is also a uniform obstruction to
approximation. For arbitrary real $a_0,b_0$, put
$x=a_0+b_0\beta$, so that its other eigenvalue is
$\bar x$. Then
\begin{equation}\label{ari:eq:ce-gap}
 \max\{|x|,|\bar x-1|\}\ge\frac12,
\end{equation}
because $1\le|\bar x|+|\bar x-1|=|x|+|\bar x-1|$.
These are exactly the two eigenvalue errors in approximating
$\Pi_-$. The spectral radius is bounded by every induced
operator norm, so no such sequence can converge to the
projector. Compositions and inverses of isogenies still belong
to the same endomorphism algebra and satisfy this bound.

\item Finally the image of $\Pi_-$ is the toric Lie line
$\C(0,1,\beta)$. It is not the Lie algebra of an algebraic
subtorus: an algebraic one-dimensional subtorus of
$\mathbb G_m^2$ has rational slope. Since $u\ne0$,
the projected vector is nonzero. The analytic subgroup
theorem in fact shows that its exponential cannot be an
algebraic point of both torus coordinates. Indeed its nonzero logarithm lies in the algebraic line
$\C(1,\beta)$, whereas a one-dimensional algebraic subtorus
has a rational slope. The analytic subgroup theorem would put
a positive-dimensional subtorus in that line for an algebraic
exponential, contradicting these two slope computations. The construction isolates the Gaussian in
a Lie eigencoordinate, without producing an algebraic
group map that evaluates it.

\item The spectral comparison with the ordinary prime $5$ is
explicit. The CM number $\alpha=(1+\beta)/2$ satisfies
$\alpha^2-\alpha+5=0$. Its two archimedean conjugates
both have absolute value $\sqrt5$. At the two places over
$5$, their normalized valuations are $0$ and $1$, because
the polynomial reduces to $X(X-1)$ and has product $5$.
The separation available in that nonarchimedean linear
algebra is absent in \eqref{ari:eq:ce-gap}. This comparison
does not turn a limiting operation on complex Lie spaces
into an algebraic-group endomorphism or transfer a
$p$-adic height conclusion to the real Gaussian.

\item \paragraph{An exact compact witness on the conductor-$361$ curve.}
In the preceding normalization the constants and the second
fibre coordinate are explicitly
\begin{equation}\label{ari:eq:ce-explicit-witness}
 s_2=-2,\qquad a=-\frac{21\beta}{703},\qquad
 r_2=-\frac{165}{74}+\frac{1839\beta}{1406}.
\end{equation}
We give polynomial certificates for these values, so they
do not depend on recognition from numerical approximations.
Put $f=4x^3-152x+361$ and
\[
 \begin{aligned}
 h={}&x^9-38x^8+437x^7-1444x^6-7942x^5\\
     &+82308x^4-274360x^3+390963x^2-130321x-130321,\\
 N={}&(-19x+76)h^2+f(h''h-(h')^2)+(6x^2-76)h'h,\\
 K={}&N'h-2Nh'.
 \end{aligned}
\]
Direct expansion gives
\[
 -K^2f=4N^3-152\cdot19^2Nh^4+361\cdot19^3h^6,
 \qquad \gcd(N,h)=1.
\]
Hence, in coordinates $Y=2y+1$, the map
\[
 (x,Y)\longmapsto
 \left(\frac{N}{19h^2},\frac{KY}{19\beta h^3}\right)
\]
is a degree-$19$ endomorphism fixing $O$, with pullback of
$dx/Y$ equal to $\beta\,dx/Y$. It is therefore the CM
endomorphism $\beta$. Divisors and the leading coefficient
at $O$ give $F_\beta=\beta^2h^2$. Its Laurent coefficient
of relative degree two is $-76$. Comparison with \ref{ari:C.7.6}
gives $38s_2=-76$, proving $s_2=-2$.

\item For the other kernel put $h_+=(H_0+\beta H_1)/10$, where
\[
 \begin{aligned}
 H_0={}&10x^8-380x^7+4864x^6-28880x^5+72561x^4\\
       &+54872x^3-795644x^2+1824494x-1433531,\\
 H_1={}&20x^7-304x^6+760x^5+10469x^4-83752x^3\\
       &+246924x^2-315514x+130321.
 \end{aligned}
\]
Writing $\bar h_+$ for replacement of $\beta$ by $-\beta$,
the finite polynomial identities and coprimality checks are
\[
 \begin{gathered}
 19xh^2-N=5f h_+\bar h_+,\qquad
 h_+\mid K+19\beta h^3,\\
 \gcd(h_+,fh)=\gcd(h_+,h_+')=1.
 \end{gathered}
\]
At a root of $h_+$ the displayed CM map sends $(x,Y)$ to
$(x,-Y)$, so these are precisely the eight pairs of
non-$2$-torsion points in $\ker(1+\beta)$. The other three
nonzero kernel points are $E[2]\setminus\{O\}$, since
$1+\beta=2\alpha$ in the CM order. Its norm is $20$,
which verifies the count. Thus
\[
 F_{1+\beta}=\frac{(1+\beta)^2}{4}f h_+^2.
\]
Taking the square ratio in \eqref{ari:eq:ce-ribet-r} and fixing
the sign by its expansion at $u=0$ gives the exact rational
function
\[
 r_2=-\frac{1+\beta}{2\beta}\frac{\wp'(u)h_+(\wp(u))}
                                             {h(\wp(u))}.
\]
At $Q$, the needed evaluations are
\[
 h(4)=-37,\qquad h'(4)=7,\qquad
 h_+(4)=\frac{829-279\beta}{10},\qquad\wp'(u)=-3.
\]
They give exactly \eqref{ari:eq:ce-explicit-witness}. In
particular $(Q,r_2)$ is an explicit nontorsion algebraic
point in the maximal compact subgroup of $G_{\beta Q}$.
Its torus-coordinate norm is
$r_2\bar r_2=974799/26011=4-\wp(\beta u)$.

\item As an independent numerical check, evaluation at $110$
decimal digits gives
\[
 \begin{aligned}
 c&=1.14793408195385447403276278660663356664777041472753\ldots,\\
 e^{-\pi u^2/(2A)}
  &=0.68133647525194864004469591167703191986051023921164\ldots.
 \end{aligned}
\]
The sigma quotient and the zeta expression for $r_2$ and $a$
agree with their exact values to more than $110$ digits.
The exact polynomial identities, evaluations and numerical
outputs are retained in \path{code/C7_cm_explicit.gp} and
its output. The polynomial construction proves the values;
the numerical comparison only checks the analytic branch.
\end{enumerate}
\end{proof}

\subsection{Complexifying the Gaussian: values, jets, and the rank criterion}

The antiholomorphic Gaussian can be made holomorphic by introducing a
second variable. This gives more algebraic values under the hypothesis
in C.7(c), but the derivative and growth hypotheses of the resulting
transcendence arguments must still be checked. We do this on the
conductor-$361$ model, where all constants and conjugations are fixed.

\begin{proposition}[the complexified route and its exact obstructions]\label{ari:C.7.11}
\StatementRecord{ari:C.7.11}{proved}{Historical events: GPT-6 Astra (OpenAI), ChatGPT Pro mode, effort Pro, S6\_09\_PAPERS, 17 November 2026: targeted proof check. Also checked in this writing. Pre-continuation writing events: GPT-6 Astra (OpenAI), ChatGPT Pro mode, effort Pro, S6\_10\_PAPERS, 18 November 2026: targeted proof check; used and rechecked in this run. Continuation metadata events: GPT-6 Astra (OpenAI), ChatGPT Pro mode, effort Pro, S6\_10\_PAPERS continuation, 19 November 2026: used and rechecked in this run; status-vocabulary reconciliation; metadata reconciliation only; no new mathematical check. \emph{Pre-continuation recorded proof scope.} Read and re-derived the full complexified norm, rank-four module, fixed-field multiplier identities, first and second derivative formulas, algebraic-independence argument and exact finite-order inequality audit. Matched the universal-vector-extension slope application to Huber-Wuestholz Theorem6.1p43 and its period reduction. Receiving revision now states the exact field Qbar(k) of the normalized second derivative, with no algebraicity premise. Closed the external-criterion source gates by reading Waldschmidt2009 Theorems3.6–3.7 pp21–22 and Waldschmidt1975 Theorem1 p24/full proof25–27, with visual strict-inequality check. Matched the derivative-ring obstruction and the exact real-plane packing/order inequality; no algebraicity conclusion for the diagonal norm is added.}{the hypotheses, notation and definitions stated here; \cref{ari:C.7.9}; \cref{ari:C.7.1}; \cref{ari:C.7.6}; \cref{ari:eq:cc-fixedfield}; \citep[Theorem 6.1]{huber2022}; \cref{ari:C.7.7}}

Use the curve, point and real logarithm of \ref{ari:C.7.9}:
\[
 E:y^2+y=x^3-38x+90,\qquad Q=(4,-2)=\exp_Eu,
 \qquad 0<u<\omega_1/2.
\]
Put $\beta=i\sqrt{19}$, $k=\pi/A$, $s_2=-2$, and
\[
 \theta(z)=\sigma(z)e^{-s_2z^2/2},\qquad
 D(z)=\zeta(z)-s_2z,\qquad
 \mathcal F(z,w)=\sqrt{19}\,\theta(z)\theta(w)e^{-kzw},
 \qquad H=\mathcal F^2.
\]
Then $\mathcal F(z,\bar z)=N(z)^2$, and $\mathcal F$ is an
entire function of two complex variables of order at most two.
It is periodic under $(z,w)\mapsto(z+\omega,w+\bar\omega)$,
$\omega\in\Lambda$. The subgroup
\[
 \Gamma=\{(\alpha u+\omega,\bar\alpha u+\bar\omega):
                   \alpha\in\Z[\beta],\ \omega\in\Lambda\}
 \subset\C^2
\]
has rank four over $\Z$ and lies in the real plane
$\{(z,\bar z):z\in\C\}$. There is one number field $K$ such
that, with $H_0=H(u,u)=N(Q)^4$, every value of $H$ on
$\Gamma$ belongs to $K[H_0]$. More precisely, the nonzero values
are algebraic multiples of the powers $H_0^{\alpha\bar\alpha}$,
as specified in \eqref{ari:eq:cc-fixedfield}.

For every nonzero $\alpha\in\Z[\beta]$ there is an algebraic
number $a_\alpha$ such that
\[
 D(\alpha u)=\bar\alpha D(u)+a_\alpha.
\]
At $p_\alpha=(\alpha u,\bar\alpha u)$ the algebraic constant
direction $V_\alpha=\alpha\partial_z-\bar\alpha\partial_w$
satisfies
\begin{align*}
 V_\alpha\log\mathcal F(p_\alpha)
    &=\alpha a_\alpha-\bar\alpha\bar a_\alpha
                         \in\overline{\Q},\\
 V_\alpha^2\log\mathcal F(p_\alpha)
    &=-\alpha^2\{\wp(\alpha u)+s_2\}
      -\bar\alpha^2\{\wp(\bar\alpha u)+s_2\}
      +2\alpha\bar\alpha k.
\end{align*}
The number $k$ is transcendental. For every nonzero $\alpha$,
the exact normalized second derivative satisfies
\[
 \overline\Q\!\left(\frac{V_\alpha^2H(p_\alpha)}{H(p_\alpha)}\right)
      =\overline\Q(k).
\] In particular,
\[
 (\partial_z-\partial_w)H(u,u)=0,\qquad
 \frac{(\partial_z-\partial_w)^2H(u,u)}{H(u,u)}=4k-8.
\]
Thus the normalized first derivative is algebraic and the
normalized second derivative generates $\overline\Q(k)$. Moreover the rank-four set does not satisfy the
explicit derivative-free finite-order criterion audited below.
These are obstructions to the stated applications of those
criteria; they do not decide the algebraicity of $N(Q)$.
\end{proposition}



\begin{proof}
\begin{enumerate}[label=\arabic*.,leftmargin=*]
\item The real model and its algebraic differential give real sigma
coefficients and a conjugation-stable lattice. In particular
$\overline{\theta(z)}=\theta(\bar z)$. The discriminant is
$-19^3$, so \ref{ari:C.7.1} gives the asserted diagonal identity.
The sigma function and the exponential of the quadratic
polynomial $-kzw$ have order at most two. From \ref{ari:C.7.6},
\[
 \theta(z+\omega)=\epsilon(\omega)
         e^{k\bar\omega(z+\omega/2)}\theta(z),
 \qquad \epsilon(\bar\omega)=\epsilon(\omega)\in\{1,-1\}.
\]
Multiplying this formula and its conjugate version, the two
signs multiply to one and the exponent cancels the change of
$-kzw$. This proves periodicity on all of $\C^2$, not only
on the real diagonal.

\item If $\alpha u+\omega=0$, then $[\alpha]Q=O$.
Nontorsion of $Q$ and the fact that a nonzero elliptic
endomorphism is an isogeny imply $\alpha=0$, and then
$\omega=0$. Thus the rank-two module $\Z[\beta]u$ and the
rank-two period lattice are independent over $\Z$.
This proves the rank assertion, with no extra algebraic
independence assumption on $u$ and the periods.

\item For $d=\alpha\bar\alpha$, the squared CM law \ref{ari:C.7.6} gives
the exact fixed-field identity
\begin{equation}\label{ari:eq:cc-fixedfield}
 H(\alpha u,\bar\alpha u)
   =19^{1-d}F_\alpha(Q)\overline{F_\alpha(Q)}H(u,u)^d.
\end{equation}
All the endomorphisms and their normalized rational functions
$F_\alpha$ are defined over a single number field: choose
a field containing the coefficients of the finitely many
generators of the CM order, and close it under conjugation.
The uniqueness of the Laurent normalization in \ref{ari:C.7.6} gives
the same field of definition for every $F_\alpha$.
Writing that fixed number field as $K$, formula
\eqref{ari:eq:cc-fixedfield} places all values in $K[H_0]$,
with no algebraicity assumption on $H_0$.
The value for $\alpha=0$ is zero, and period shifts do not
change it. Squaring was included because algebraicity of the
individual positive square roots alone would not establish
their membership in one fixed number field.

\item Logarithmically differentiating the squared CM law at $u$
gives
\[
 a_\alpha=\frac{1}{2\alpha}
     \frac{(d/dz)F_\alpha(\exp_Ez)|_{z=u}}{F_\alpha(Q)}
                  \in\overline{\Q}.
\]
The denominator is nonzero since $Q$ is nontorsion. The
derivation $d/dz$ is the algebraic invariant derivation on
$E$, so the displayed derivative has algebraic value.
Since $D(u)$ is real, conjugation gives the corresponding
formula at $\bar\alpha u$. On a neighborhood where
$\mathcal F\ne0$, its logarithmic gradient and Hessian are
\[
 (\partial_z\log\mathcal F,\partial_w\log\mathcal F)
       =(D(z)-kw,D(w)-kz),\qquad
 \operatorname{Hess}\log\mathcal F=
 \begin{pmatrix}-\wp(z)-s_2&-k\\-k&-\wp(w)-s_2\end{pmatrix}.
\]
Evaluation on $(\alpha,-\bar\alpha)$ proves both identities
in the statement. In particular the Gaussian cancels in the
first derivative and reappears in the second derivative.

\item To prove $k$ transcendental, recall the real period formula
$\eta_1=(s_2+k)\omega_1$. If $k$ were algebraic, the
nonzero period vector $(\omega_1,\eta_1)$ of the algebraic
universal vector extension of $E$ would lie on the algebraic
line of slope $s_2+k$. Dividing this period by a sufficiently
large integer supplies a nonzero algebraic torsion exponential.
The analytic subgroup theorem
\citep[Theorem 6.1]{huber2022} would give a one-dimensional
algebraic subgroup with that Lie line. Its projection to $E$
is nonzero, so it is an elliptic curve isogenous to $E$ and
splits the vector extension after that isogeny. This is
impossible: its nonzero class in $H^1(E,\mathcal O_E)$ stays
nonzero under an isogeny, whose pullback has a left inverse
after dividing the trace by the degree. These are the same
algebraic logarithm coordinates and extension as in \ref{ari:C.7.7};
a fixed algebraic linear change of coordinates leaves the
argument unchanged.

\item Write the two logarithmic derivatives as
$V_\alpha\log\mathcal F=B_\alpha$ and
$V_\alpha^2\log\mathcal F=A_\alpha+2dk$, where
$A_\alpha,B_\alpha$ are algebraic. Since $H=\mathcal F^2$,
\[
 \frac{V_\alpha^2H(p_\alpha)}{H(p_\alpha)}
       =4B_\alpha^2+2A_\alpha+4dk.
\]
Here $d>0$, so the displayed affine expression both belongs to
$\overline\Q(k)$ and recovers $k$ by subtraction of its algebraic
constant and division by $4d$. This proves the asserted field
equality; $H(p_\alpha)$ is nonzero because $Q$ is nontorsion.
At $\alpha=1$, $a_1=0$, $\wp(u)=4$ and $s_2=-2$,
which give the last displayed formula in the statement.

\item \paragraph{The precise several-variable criteria that were checked.}
The differential criterion of
\citet[Theorems 3.6 and 3.7]{waldschmidt2009auxiliary}
requires a finite algebraically independent family whose
partial derivatives belong to the polynomial ring generated
by that family over one number field. Algebraic values of
the full family then give algebraic values of all iterated
partial derivatives. The calculation above gives the exact value-field obstruction:
the field containing $H(p_\alpha)$ and its second directional
derivative contains $k$ as soon as the nonzero value is inverted.
In particular its normalized second jet is transcendental. The issue persists
although the first cancelling direction is algebraic; the
direction also varies with $\alpha$.

\item For a criterion without differential equations we opened
\citet[Theorem 1, p.~24]{waldschmidt1975dimension}, including
its proof and a visual check of the displayed inequalities.
The source uses $d$ meromorphic functions on $\C^n$, growth and
arithmetic-size exponents $\rho_1,\ldots,\rho_d$, and sampling
sets of cardinality $N^{\ell-o(1)}$ in a ball of radius $O(N)$,
with separation at least $N^{-\lambda-o(1)}$. Its strict numerical
requirement is
\begin{equation}\label{ari:eq:cc-rankcriterion}
 (d-1)\ell>\sum_{j=1}^d\rho_j
                     +2d(n-1)(\lambda+1).
\end{equation}
The theorem also requires all function values to belong to
one number field and bounds on denominators in the entire
representations of the meromorphic functions. We do not
silently discard these hypotheses.

\item The natural independent triple here is
$\wp(z),\wp(w),H(z,w)$, all of order at most two.
The first two are algebraically independent. If $H$ satisfied
a polynomial over $\C(\wp(z),\wp(w))$, fix a generic pair
$(z,w)$ and translate only $z$ by $m\omega$, with
$0\ne\omega\in\Lambda$. The coefficients of the polynomial
remain constant, whereas the absolute values of its proposed
roots $H(z+m\omega,w)$ grow as
$\exp(k|\omega|^2m^2+O(|m|))$. This contradicts the finite
set of roots of a nonzero polynomial. Thus the triple is
independent.

\item The growth-order lower envelope is $\rho_j=2$. Substitution of
this envelope and $n=2$, $d=3$, $\ell=4$ into
\eqref{ari:eq:cc-rankcriterion} gives the required inequality
\[
 8>12+6\lambda,
\]
which fails even for $\lambda=0$. The actual geometry is
less favorable: the $\asymp N^4$ points obtained by bounding
the four integer coefficients lie in a ball of radius $O(N)$
in a real plane. Packing separated discs in that plane gives
$N^4\ll N^{2+2\lambda+o(1)}$, hence $\lambda\ge1$.
Removing the rank-two set of poles does not change this count.
Larger growth exponents only worsen the inequality. This
checks the particular published criterion rather than asserting
that every possible multivariable transcendence method fails.
The missing unconditional step in C.7(c) remains the exclusion
of an algebraic value on the diagonal compact leaf.
\end{enumerate}
\end{proof}

\subsection{Division values and a variable-degree interpolation test}


\begin{proposition}[division ratios, scalar models and translated functions]\label{ari:C.7.12}
\StatementRecord{ari:C.7.12}{proved}{Historical events: Claude Fable 5.2 (Anthropic), effort high, 15 November 2026: carried as background without fault found. GPT-6 Astra (OpenAI), ChatGPT Pro mode, effort Pro, S6\_09\_PAPERS, 17 November 2026: targeted proof check. Also checked in this writing. Pre-continuation writing events: GPT-6 Astra (OpenAI), ChatGPT Pro mode, effort Pro, S6\_10\_PAPERS, 18 November 2026: targeted proof check; used and rechecked in this run. Continuation metadata events: GPT-6 Astra (OpenAI), ChatGPT Pro mode, effort Pro, S6\_10\_PAPERS continuation, 19 November 2026: used and rechecked in this run; status-vocabulary reconciliation; metadata reconciliation only; no new mathematical check. \emph{Pre-continuation recorded proof scope.} Read the full division-value, height, translate-field and criterion audit, including all continuation paragraphs. The receiving statement uses unconditional algebraic ratios and scalar models; no unresolved algebraicity hypothesis is adopted. Re-derived the full unconditional replacement: exact CM weight and algebraic ratio, scalar-parameter transformation, n4/2n6 degree and bounded-height proofs, algebraic normalized-translate scalars, transcendence-degree-three proof and n5 group-division bound. Actual analytic values are recovered at lambda=N(Q), whose algebraicity remains open. Separate agent independently re-derived the four delicate steps: uniqueness/additivity of kappa, sampling projection nu, cancellation fixing elliptic-ratio scalars at lambda=1, and n3 times n2 joint degree. Incorporated its bounded numerator-window correction and its algebraic group-law treatment of meromorphic-section exceptional points. This event has exactly that bounded mathematical scope. Closed the listed criterion-source gates: GMW Lemma1.0 p98/Theorem2.1 pp101–105 visually; Herblot Definitions2.4–2.5 p11/Theorem6.1 pp31–32; Gasbarri Definition3.7 pp14–15/Theorem4.16 pp25–26; Adamczewski–Faverjon system1.2/Theorem1.1 p56; Einsiedler–Everest–Ward Theorem5.1 pp9–10; Waldschmidt1988 Theorem4.1/Corollary4.3 pp385–386 visually; Roy–Waldschmidt quotient definitions/main theorem pp359–364. Scope is the exact inequalities and assumptions used in the receiving audit.}{\cref{ari:C.4.2}; \cref{ari:C.7.1}; \cref{ari:C.7.6};
\cref{ari:C.7.9}; \cref{ari:C.7.11}; the exact divisor and
height computations below; the cited criteria at their stated scopes}

Let $E:y^2+y=x^3-38x+90$, $Q=(4,-2)$ and $N_0=N(Q)$.
Use the positive real logarithm $u$ and period $\omega_1$ of
\ref{ari:C.7.9}. Put
\[
 \delta=19^{1/4},\qquad b=\eta_1/\omega_1,\qquad
 g(z)=\delta\sigma(z)e^{-bz^2/2}.
\]
Thus $|g(z)|=N(\exp_Ez)$ on the real axis. Let $\mathcal D$
be the CM division hull of $Q$. For $R\in\mathcal D$, the
equations $[n]R=[\alpha]Q$, $n>0$, determine a unique
$\kappa_R=\alpha/n\in\operatorname{End}(E)\otimes\Q$.
Set $w_R=\kappa_R\bar\kappa_R$. For $R\ne O$,
\[
 A_R:=\frac{N(R)}{N_0^{w_R}}\in\overline\Q_{>0}.
\]
In particular, for $[n]R=[m]Q$ with $m\ne0$,
\begin{equation}\label{ari:eq:division-ratio-model}
 A_R^{n^2}=\delta^{n^2-m^2}
                  \frac{|\psi_m(Q)|}{|\psi_n(R)|}.
\end{equation}
For each positive algebraic number $\lambda$, define
\[
 N_\lambda(R)=A_R\lambda^{w_R}
       =N(R)\left(\frac{\lambda}{N_0}\right)^{w_R}.
\]
All these values are algebraic. For $[n]R=[m]Q$, $|m|\le n$,
$R\ne O$, there are constants $D,H$ depending on $\lambda,E,Q$
such that
\[
 h(N_\lambda(R))\le H,\qquad
 [\Q(R,N_\lambda(R)):\Q]\le
 \begin{cases}Dn^4,&R\in E(\R),\\2Dn^6,&R\in E(\C).
 \end{cases}
\]
The transformation is exact for every positive real parameter:
$N_\lambda(R)/N_\mu(R)=(\lambda/\mu)^{w_R}$ and
$N_{N_0}(R)=N(R)$.

For any finite collection of real division-hull logarithms $a$,
the functions $g(z)$ and $g(z+a)/g(a)$, with $\exp_Ea\ne O$,
have function-field transcendence degree at most three over
$\overline\Q$ after adjoining the elliptic coordinates.
The three functions $g(z),g(z+v)/g(v),g(z-v)/g(v)$, $v=u/q$,
attain three. The normalized holomorphic extensions
$F_a(z)/N(a)^2$ of $N(\exp_E(z+a))^2$ for fixed complex
division-hull shifts belong to the same type of field.
The unnormalized translates have the same upper bound, attained by
the displayed triple, over the algebraic closure $L$ of $\overline\Q(N_0)$ in $\C$.

For the scalar models on the real division grids
$z=(mu+k\omega_1)/n$, $|m|\le Cn$, $k$ modulo $n$, with
fixed $C$, every fixed finite family of these values has a common
degree envelope $O(n^5)$ and bounded heights. For the actual values the same construction
gives a relative degree envelope over a fixed finite extension
of $\Q(N_0)$. The algebraicity of $N_0$ remains the open question
of \ref{ari:C.7.4}. The explicit interpolation and tangent
calculations below record what these bounds establish.
\end{proposition}

\begin{proof}
\begin{enumerate}[label=\arabic*.,leftmargin=*]
\item \emph{The algebraic ratios and their exact change of parameter.}
Two equations $[n]R=[\alpha]Q$ and $[n']R=[\alpha']Q$
give $[n'\alpha-n\alpha']Q=O$. A nonzero elliptic endomorphism
has finite kernel, whereas $Q$ is nontorsion. Thus
$\alpha/n=\alpha'/n'$, proving uniqueness. The resulting
map $R\mapsto\kappa_R$ is additive, and its kernel is the
torsion subgroup. Every such $R$ is algebraic, since it lies
in a finite isogeny fibre above an algebraic point.

\item The division identity and the squared CM law give, for
$d=\alpha\bar\alpha$ and $\alpha\ne0$,
\[
 N(R)^{n^2}=\frac{N(nR)\delta^{n^2-1}}{|\psi_n(R)|},
 \qquad
 N(\alpha Q)^2=\delta^{2-2d}|F_\alpha(Q)|N_0^{2d}.
\]
Consequently
\[
 A_R^{2n^2}=\frac{\delta^{2n^2-2d}|F_\alpha(Q)|}
                         {|\psi_n(R)|^2}\in\overline\Q_{>0}.
\]
The positive root proves $A_R\in\overline\Q$. For integral
$\alpha=m$, the usual multiplication identity is
$N(mQ)=N_0^{m^2}|\psi_m(Q)|/\delta^{m^2-1}$;
substitution proves \eqref{ari:eq:division-ratio-model}.
For nonzero $n$-torsion $R$, take the limit at the simple zero
of $\psi_n$ in the uniformizing coordinate. Since
$N(z)=\delta|z|+o(|z|)$ at zero,
\[
 A_R^{n^2}=N(R)^{n^2}
       =\frac{n\delta^{n^2}}{|\partial_z\psi_n(R)|}.
\]
The invariant derivation is algebraic, so this also is algebraic.
All exponents $w_R$ are rational. The asserted scalar models,
their parameter-change identity and the specialization at $N_0$
follow directly. In particular their definition is independent
of every presentation of $R$ by a division equation.

\item Put $K_\lambda=\Q(\delta,\lambda)$ and
$D=[K_\lambda:\Q]$. For integral $m\ne0$,
\[
 N_\lambda(mQ)=\lambda^{m^2}
                     |\psi_m(Q)|/\delta^{m^2-1}\in K_\lambda,
\]
because $Q$ is rational. The fibre of $[n]$ has $n^2$ points,
so adjoining $R$ costs at most $n^2$. For real $R$ the absolute
value of $\psi_n(R)$ differs from it by a sign. The division
equation for $N_\lambda(R)$ costs at most another $n^2$.
For complex $R$, adjoin $R$ and $\bar R$, at cost at most
$n^4$, and use
\[
 N_\lambda(R)^{2n^2}
 =\frac{N_\lambda(mQ)^2\delta^{2(n^2-1)}}
              {\psi_n(R)\psi_n(\bar R)}.
\]
This gives $2Dn^6$. At nonzero torsion points the derivative
formula proves the same bounds.

\item \emph{Uniform heights.}
We give the coefficient estimate needed here, so the height
assertion does not rely on an unspecified division-polynomial
bound.  On the fixed integral model, $\psi_n$ is a polynomial
in $x$ for odd $n$, and $(2y+1)$ times a polynomial in $x$
for even $n$.  These polynomials have degree $O(n^2)$.
Choose a fixed circle $|x|=A>1$ avoiding the finite branch values.
On its compact inverse image in $E(\C)$, $N$ has a positive
minimum $m_A$, and $|2y+1|$ has a positive minimum.  Let $M$
be the maximum of $N$ on the whole compact elliptic curve.
The rearranged division identity gives
\[
 |\psi_n(R)|\leq M\delta^{n^2-1}m_A^{-n^2}
       \qquad(|x(R)|=A).
\]
Cauchy's coefficient estimate, and division by $2y+1$ for
even $n$, show that the logarithm of the coefficient
$\ell^1$ norm is $O_E(n^2)$.  The same assertion holds for
$\partial_z\psi_n$, because differentiation and the fixed
Weierstrass equation add only polynomial factors in $n$.

\item The elementary height estimate for a fixed degree-four
duplication map, followed by telescoping, gives
$h(x(R))=2\widehat h(R)+O_E(1)$ uniformly over
$E(\overline\Q)$.  Since
$\widehat h(R)=(m/n)^2\widehat h(Q)$, the coordinate heights
are uniformly bounded when $|m|\leq n$.
At each nonarchimedean place the ultrametric inequality bounds a
polynomial value by its largest monomial; at each archimedean
place use the sum of the monomials. Summing these logarithmic
bounds by the product formula gives the polynomial evaluation
height inequality, and hence
$h(\psi_n(R))=O_{E,Q}(n^2)$, and the same estimate for its
invariant derivative at torsion points.  Also
\[
 h(N_\lambda(mQ))\leq m^2h(\lambda)+O_{E,Q}(m^2).
\]
Apply the same division identity to $N_\lambda$ and divide its height by $n^2$.
For a complex value one uses
$h(|\alpha|)\leq h(\alpha)$, which follows by squaring and
using $|\alpha|^2=\alpha\bar\alpha$.
This proves a uniform bound for $h(N_\lambda(R))$.

\item \emph{Normalized translates and their algebraic scalars.}
For a real logarithm $a$ of a nonzero division point put
\[
 h_a(z)=\frac{\sigma(z+a)}{\sigma(z)\sigma(a)}e^{-baz}.
\]
The exact identities are
\[
 g(z+a)=\delta^{-1}g(z)g(a)h_a(z),\qquad
 h_a(z+\omega_1)=h_a(z),\quad
 h_a(z+\omega_2)=e^{-2\pi ia/\omega_1}h_a(z).
\]
They follow by substitution from sigma automorphy and Legendre's
relation. A real point in the CM division hull lies in the
ordinary division hull: conjugating $[n]R=[\alpha]Q$ gives
$[\alpha-\bar\alpha]Q=O$, hence $\alpha=\bar\alpha\in\Z$.
The real group is connected because the discriminant is negative.
Its division logarithms are therefore exactly
$\widetilde{\mathcal D}=\Q u+\Q\omega_1$.
The coefficient map
\[
 \nu:\widetilde{\mathcal D}\longrightarrow\Q,
 \qquad \nu(ru+s\omega_1)=r
\]
is well defined and additive: a rational relation between $u$
and $\omega_1$ would make $Q$ torsion.

\item For any positive real $\lambda$ define, on this additive group,
\[
 \widetilde g_\lambda(z)
   =g(z)(\lambda/N_0)^{\nu(z)^2},\qquad
 h_{a,\lambda}(z)
   =h_a(z)(\lambda/N_0)^{2\nu(z)\nu(a)}.
\]
At nonzero division values, the absolute value of
$\widetilde g_\lambda$ is $N_\lambda$ and its sign is the real
sign of $g$. Thus these values are algebraic for algebraic
$\lambda$. The identity
\[
 h_{a,\lambda}(z)=
 \frac{\delta\widetilde g_\lambda(z+a)}
      {\widetilde g_\lambda(z)\widetilde g_\lambda(a)}
\]
is proved by expanding $(\nu(z)+\nu(a))^2$.
It specifies the exact transformation on the sampling group.

\item Choose a common denominator $q$ for finitely many shifts,
write $a_j=(m_ju+k_j\omega_1)/q$, and put $v=u/q$.
The quotient $h_{a_j}/h_v^{m_j}$ is elliptic for
$\Lambda_q=\Z\omega_1+\Z q\omega_2$. This finite isogeny
of $E$ has a model over $\overline\Q$ with the stated
differential. The quotient has an algebraic divisor, hence
is a rational function over $\overline\Q$ times a scalar.
The scalar is algebraic. Indeed the parameter factors cancel:
\[
 \frac{h_{a_j}(z)}{h_v(z)^{m_j}}
  =\frac{h_{a_j,1}(z)}{h_{v,1}(z)^{m_j}}
       \qquad(z\in\widetilde{\mathcal D}).
\]
Choose a real division point avoiding the finite zeros and poles.
The right side and the value of the chosen rational function
are algebraic and nonzero, so their quotient is algebraic.
This proves that $g(z)$ and all normalized real translates are
algebraic over $\overline\Q(E)(g,h_v)$, after the fixed finite
isogeny. The upper transcendence-degree bound is three.

\item For the lower bound, work even over $\C$. The multiplier of
$h_v$ under $q\omega_2$ is $e^{-2\pi iu/\omega_1}$,
which has infinite order. A polynomial relation over the elliptic
field, evaluated at a generic $z$ and translated by $mq\omega_2$,
would have infinitely many distinct roots obtained by these
powers. Thus $h_v$ is transcendental over that field. Moreover
\[
 |g(z+mq\omega_2)|=\exp(c m^2+O(|m|)),
 \qquad c=\pi q^2\Im(\omega_2/\omega_1)>0.
\]
In a proposed polynomial relation for $g$ over the elliptic field
and $h_v$, clear denominators. The nonzero leading coefficient
is a polynomial evaluated along an irrational rotation; it is
bounded away from zero on a subsequence, and all other
coefficients remain bounded there. Quadratic growth makes the
highest power of $g$ dominate, a contradiction. Finally
\[
 \wp(z)=\wp(v)-
 \frac{\delta^2g(z+v)g(z-v)}{g(z)^2g(v)^2}
\]
recovers the elliptic coordinate from the three stated translates,
and $h_v=\delta g(z+v)/(g(z)g(v))$. They therefore attain three.
The constants $g(a)$ lie in $L$, by the already proved ratios.
Multiplication by these constants gives the exact map from
normalized to raw translates, proving the assertion over $L$.

\item For a complex CM division shift $a$, the holomorphic extension
from the real axis is
\[
 F_a(z)=\frac{N(a)^2}{\delta^2}g(z)^2h_a(z)h_{\bar a}(z).
\]
Its multiplier depends on $a+\bar a\in\Q u+\Q\omega_1$.
Choose $q,m$ with $a+\bar a=(mu+k\omega_1)/q$.
Then $h_ah_{\bar a}/h_v^m$ is elliptic on $\C/\Lambda_q$
with algebraic divisor. Its scalar is algebraic by the same
evaluation argument, using the exact identity on real division
points
\[
 h_a(z)h_{\bar a}(z)(\lambda/N_0)^{2\nu(z)m/q}
 =\frac{\delta^2N_\lambda(\exp_E(z+a))^2}
        {N_\lambda(\exp_Ea)^2\widetilde g_\lambda(z)^2}.
\]
Here $w_{\exp_E(z+a)}=|\kappa_{\exp_Ea}+\nu(z)|^2$;
expanding the square proves the identity. Set $\lambda=1$
and divide by $h_{v,1}^m$. This proves the normalized complex
translate assertion with all scalar factors accounted for.

\item \emph{Joint degrees and heights on the sampling group.}
In the algebraic meromorphic section coordinate of
\ref{ari:C.7.1},
\[
 z\longmapsto(\exp_Ez,h_v(z))
   =\exp_{G_V}\bigl(z,(\zeta(v)-bv)z\bigr),\qquad
 V=\exp_Ev.
\]
For the scalar model, multiply its torus coordinate by
$(\lambda/N_0)^{2\nu(z)\nu(v)}$. This factor is a character
of the additive sampling group, so the resulting map
$\Phi_\lambda:\widetilde{\mathcal D}\to G_V(\C)$ is an
abstract group homomorphism. Its real period is $\omega_1$,
because $\nu(\omega_1)=0$ and the displayed analytic map
has that period. For algebraic $\lambda$, its values in the displayed section
chart are algebraic by the formula for $h_{v,\lambda}$ above.
For an exceptional sampling point at a zero or pole of that
section, write $z=z_1+z_2$ with both $z_j$ avoiding the finite
exceptional set modulo the elliptic lattice. Such a choice exists
in the infinite division group. Then
$\Phi_\lambda(z)=\Phi_\lambda(z_1)\Phi_\lambda(z_2)$ is
algebraic by the algebraic group law. This also covers the identity.

\item For $z=(mu+k\omega_1)/n$,
$[n]\Phi_\lambda(z)=[m]\Phi_\lambda(u)$.
Choose a fixed number field containing the coefficients of
$G_V$, the point $\Phi_\lambda(u)$ and the finitely many
fixed constants. Multiplication by $n$ on $G_V$ has degree
$n^3$, since its kernel is an extension of $E[n]$ by
$\mathbb G_m[n]$. Adjoining $\widetilde g_\lambda(z)$ costs
at most $n^2$ by the norm division equation. Passing to the
fixed isogeny and adjoining fixed constants changes only the
constant in $O(n^5)$. The same argument for any positive real
$\lambda$ works over a fixed finite extension of $\Q(\lambda)$,
because all the fixed division values are algebraic over that
field. At $\lambda=N_0$ it gives the claimed relative envelope.

\item The height proof above applies with $|m|\le Cn$ for any fixed
$C$, changing only its constant. It bounds the model values at
$z$ and every fixed real translate $z+a$. Their quotients bound
$h_{v,\lambda}(z)$. Coordinates on the fixed isogenous elliptic
curve have bounded height by the fixed isogeny height comparison.
The normalized complex translates are fixed rational functions
of those coordinates, $\widetilde g_\lambda^2$ and a fixed power
of $h_{v,\lambda}$, with algebraic coefficients just proved.
The polynomial evaluation inequality bounds their heights.
Any removable zero or pole in these expressions lies in a fixed
finite set and is handled by its actual model value. Taking the
positive square root for a fixed complex-shifted norm divides
height by two and costs at most a factor two in degree.

\item \emph{The exact numerical and tangent checks.}
The sufficient inequality in Gramain--Mignotte--Waldschmidt
\citep[Theorem2.1,pp.101--102]{gramain1986values} has the form
\[
 d(N)\prod_{j=1}^k\max\{d(N)\psi_j(N),\theta_j(R(N))\}
 \le c_0\mu(N)^{k-1}\varphi(N)^k,
\]
\[
 \varphi(N)=\log\frac{R(N)^2+r(N)r(N+1)}
                         {R(N)(r(N)+r(N+1))}.
\]
On fixed nested discs the growth and geometric factors are
constant. The scalar-model height envelopes are constant too.
The numerical comparison of these envelopes is therefore
$d^{k+1}\ll\mu^{k-1}$. At denominator $n$ the real group
has $O(n^2)$ sample points, and through denominator $N$ it
has $O(N^3)$. Substitution of the proved envelopes gives
$N^{12}$ against $N^3$ for two functions and $N^{20}$ against
$N^6$ for three. These upper-envelope calculations provide
no such inequality and assert no lower bound for the actual
degrees. For actual analytic samples the base field is the
displayed extension of $\Q(N_0)$; its number-field status is
exactly the unresolved algebraicity issue. Extra normalized
translates remain in the function field of transcendence degree
three proved above.

\item The arithmetic-germ criteria
\citep[Theorem6.1]{herblot2012meromorphic}
\citep[Theorem4.16]{gasbarri2010analytic}
can be compared with the actual logarithmic tangent. At a real
nontorsion division point $z$, the exact formula is
\[
 \frac1{g(z)}\frac{dg}{dx}
       =\frac{\zeta(z)-bz}{\wp'(z)}.
\]
Its numerator is transcendental by the period-line proof of
\ref{ari:C.7.9}, applied to this algebraic nontorsion point;
its denominator is nonzero algebraic. A formal graph defined
over $\overline\Q$ would have an algebraic nonzero constant
value and algebraic first derivative, contradicting their
displayed ratio. The scalar change on the sampling group has
not changed this exact tangent of the analytic graph.

\item For the Mahler system in
\citep[Theorem1.1 and system(1.2)]{adamczewski2017mahler},
the needed algebraic Taylor coefficients and algebraic regular
argument can be tested directly: the coefficient of $z^3$ in
$g$ is $-\delta b/2$, and $b=s_2+k$ is transcendental by
\ref{ari:C.7.11}. The logarithm $u$ is transcendental too:
an algebraic $u\ne0$ would give an algebraic point $(u,Q)$
on the diagonal Lie line of $\mathbb G_a\times E$; the analytic
subgroup theorem would produce an algebraic one-dimensional
subgroup projecting nontrivially to both factors, which is
impossible since an elliptic curve has no nonconstant map to
$\mathbb G_a$ and $\mathbb G_a$ has no nonzero group map to $E$.
The duplication identity acts by the explicit elliptic
multiplication map. The entropy identity in
\citep[Theorem5.1]{einsiedler2001entropy} uses the sequence
built from the denominators of $x(2^nQ)$ with scale $4^n$.
These precise inputs supply no evaluation of $N_0$ here.

\item Finally the period-refined rank comparison
\citep[Theorem4.1 and Corollary4.3]{waldschmidt1988periodrank}
\citep[main theorem,p.363]{roy1993periodrank}
has coefficient $\alpha(G)=2\dim G-2d_0(G)-d_1(G)$,
where the last terms count additive and multiplicative
quotients. Both vanish for $G_Q$. An additive homomorphism
vanishes on its torus and then on the proper quotient $E$;
a nonzero torus character extending to $G_Q$ would kill a
nonzero integer multiple of its nontorsion extension class.
Thus $\alpha(G_Q)=4$.

\item The line $\ell_c$ has period rank one. A general period is
$(\omega,\zeta(u)\omega-\eta(\omega)u+2\pi i j)$.
Writing $\omega=a\omega_1+d\omega_2$, membership in
$\ell_c$ gives $du/\omega_1\in\Z$ by Legendre's relation.
Since $u/\omega_1$ is irrational, $d=j=0$. The numerical
rank bound in the cited theorem is
$(4-1)\dim\ell_c=3$. The period and the actual logarithm
$(u,cu)$ are independent over $\Q$ and span rank two;
the algebraicity of the latter's exponential remains exactly
the point-algebraicity question in \ref{ari:C.7.9}. These
explicit ranks leave that question open.
\end{enumerate}
\end{proof}


\section{Limits and open points: the coefficient atlas}

\paragraph{what is and is not explicit}
\label{wb:rem:bcpr-literal-gauge}
Choose global $\bar\partial$-closed forms
\begin{equation}
 G\in\mathcal A^{0,1}(X),\qquad
 Q\in\mathcal A^{1,1}(X)
\end{equation}
representing $\mathfrak g$ and $\mathfrak q$.
\Cref{inv:A3.2} says that
$i_b^*[Q]_{\bar\partial}$ is the Dolbeault class of the explicit
translation-invariant
form displayed in \CompendiumRef{I}{wb:eq:bcpr-q-fibre-expanded}; it does not say that an arbitrary
representative $Q$ restricts literally to that form.  A literal global
coefficient formula would require additional $C^\infty$ gauge choices and
horizontal correction forms on the disc--pair-of-pants cover, compatible
with the two affine cyclic conjugacies and the toric cusp.  Those choices
are not part of the canonical data.  No such literal gauge is needed for
the class and product calculations below.


\Uses{the hypotheses, notation and definitions stated here; \cref{inv:A3.2}; \CompendiumRef{I}{wb:eq:bcpr-q-fibre-expanded}}


\paragraph{literal global gauge boundary}
\label{wb:rem:bcpr-global-gauge-boundary}
The Bott--Chern algebra \cref{inv:C3.3} and the fibre-coordinate
restrictions displayed in
\CompendiumRef{I}{wb:eq:bcpr-gamma-form}, \CompendiumRef{I}{wb:eq:bcpr-u-form}, \CompendiumRef{I}{wb:eq:bcpr-q-fibre}, \CompendiumRef{I}{wb:eq:bcpr-q-fibre-expanded}, \CompendiumRef{I}{wb:eq:bcpr-q-gamma-product}, \CompendiumRef{I}{wb:eq:bcpr-q-square}
are complete.  The canonical data do not select a preferred literal global
coefficient atlas for the smooth forms $G$ and $Q$.  Such an atlas would
require explicit smooth local trivializations and horizontal correction
functions through both finite affine fillings and the cusp.  After fixing
$\mathfrak g$ and one coherent-Leray lift $\mathfrak q$, changing those
smooth gauges changes $G$ and $Q$ by Dolbeault coboundaries but changes
none of the cohomology classes, nonvanishing arguments, or products proved
above.  This is a limitation only on a preferred global coefficient
presentation, not on the representatives as cohomology classes or on the
Bott--Chern algebra.


\Uses{the hypotheses, notation and definitions stated here; \cref{inv:C3.3}; \CompendiumRef{I}{wb:eq:bcpr-gamma-form}, \CompendiumRef{I}{wb:eq:bcpr-u-form}, \CompendiumRef{I}{wb:eq:bcpr-q-fibre}, \CompendiumRef{I}{wb:eq:bcpr-q-fibre-expanded}, \CompendiumRef{I}{wb:eq:bcpr-q-gamma-product}, \CompendiumRef{I}{wb:eq:bcpr-q-square}}


\begin{question}[Open: A literal global coefficient atlas for the generators]\label{inv:C3.8}
\StatementRecord{inv:C3.8}{open}{Historical events: Claude Fable 5.2 (Anthropic), effort high, 12 November 2026: targeted proof check. Claude Fable 5.2 (Anthropic), effort high, 16 November 2026: carried as background without fault found. Claude Fable 5.2 (Anthropic), effort high, 17 November 2026: carried as background without fault found. GPT-6 Astra (OpenAI), ChatGPT Pro mode, effort Pro, S6\_09\_PAPERS, 17 November 2026: targeted proof check. Also checked in this writing. Pre-continuation writing events: GPT-6 Astra (OpenAI), ChatGPT Pro mode, effort Pro, S6\_10\_PAPERS, 18 November 2026: targeted proof check. Continuation metadata events: GPT-6 Astra (OpenAI), ChatGPT Pro mode, effort Pro, S6\_10\_PAPERS continuation, 19 November 2026: used and rechecked in this run; status-vocabulary reconciliation; metadata reconciliation only; no new mathematical check. \emph{Pre-continuation recorded proof scope.} Re-derived the harmonic Riesz representative, positivity of its denominator, exact restriction coefficient and invariance under a unitary change of basis; added a receiving proof for these settled assertions. The literal common-atlas coefficient evaluation remains open; no elliptic inverse was numerically evaluated.}{\cref{inv:A3.2}, \cref{inv:C3.1}, \cref{inv:C3.2}, \cref{inv:C3.7}}

The invariant computations and all the metric types named above are
settled for $X$.  The real restriction map
\[
r_b:H_A^{1,1}(X;\R)\longrightarrow H_A^{1,1}(F_b;\R)
\]
has image $\R\eta$ by \ref{inv:A3.2} and \ref{inv:C3.7}.  Its source has real dimension
two and its target has real dimension four.  The map is induced by pullback of $\partial\bar\partial$-closed forms followed by the quotient by $\partial$- and $\bar\partial$-exact forms on $F_b$. Thus it specifies the class of each restricted representative through that exact quotient map.

The formula \eqref{inv:eq:C-G-Green} supplies a global $G$ after a metric
choice.  For the standard primitive marking, $G_\gamma=G_0/12$
represents $g_\gamma$ by \ref{inv:C3.2}.  There is a precise operator
construction of $Q$ as well.  Fix a Hermitian metric on $X$ and a
regular fibre $F$, identify $\eta$ with its real invariant form on
$F$, and choose an orthonormal basis $e_1,e_2$ of the Dolbeault-harmonic
$(1,1)$ forms for $\Delta_{\bar\partial}$.  Define
\[
\ell(e)=-\frac1{12}\int_F i_F^*e\wedge\eta,\qquad
Q=\frac{\overline{\ell(e_1)}e_1+
            \overline{\ell(e_2)}e_2}
          {|\ell(e_1)|^2+|\ell(e_2)|^2}.
\]
The denominator is positive by \ref{inv:A3.2} and
$\int_F\eta^2=-12$.  The expression is independent of the
orthonormal basis, is $\bar\partial$-closed, and has fibre
Dolbeault class $\eta$.  Indeed the image of fibre restriction on
$H^{1,1}_{\bar\partial}(X)$ is $\C\eta$, with kernel
$\C d_1g$; the displayed functional measures exactly the coefficient
of $\eta$, and the formula makes that coefficient one.

These two formulas fix global smooth representatives by explicit
elliptic operators and a finite normalization.  We have not evaluated
their coefficient functions on a common atlas through the three
special fibres.  The remaining task is therefore precise: compute
the harmonic projector and the two displayed restriction integrals
for one fully specified global Hermitian metric, or construct and
evaluate a Cech representative and partition-of-unity correction
with the same restrictions.  No numerical computation of those
infinite-dimensional elliptic inverses was run.  The coefficient functions on that specified common atlas remain open.
\end{question}

\begin{proof}
\begin{enumerate}[label=\arabic*.,leftmargin=*]
\item We verify the settled operator assertions in this open question.
The Green formula and $\exp([G_0])=f^*\OO(1)$ were proved in
\ref{inv:C3.1}; the primitive marking in \ref{inv:C3.2} then gives
$[G_0]/12=g_\gamma$.

\item Restriction of a Dolbeault-harmonic representative is a closed
Dolbeault form on the regular fibre. By \ref{inv:A3.2}, its class
is $c\eta$. Integration against the closed $(1,1)$ form $\eta$
annihilates Dolbeault-exact changes by Stokes' theorem. Since
$\int_F\eta^2=-12$, the displayed functional gives $\ell(e)=c$.
The image is $\C\eta$, hence $\ell$ is nonzero on the
 two-dimensional harmonic space and its squared norm
$|\ell(e_1)|^2+|\ell(e_2)|^2$ is positive.

\item For another orthonormal basis $e'_j=\sum_k U_{jk}e_k$ with
$U$ unitary, put $l_k=\ell(e_k)$. Then $l'=Ul$ and
\[
 \sum_j\overline{l'_j}e'_j
 =\sum_{j,k,r}\overline{U_{jk}}\bar l_k U_{jr}e_r
 =\sum_k\bar l_ke_k,
 \qquad \sum_j|l'_j|^2=\sum_k|l_k|^2.
\]
Thus $Q$ is independent of the orthonormal basis. It is a linear
combination of harmonic forms and $\ell(Q)=1$, so its fibre class
is precisely $\eta$. These computations establish the operator
representatives and their restrictions. They leave the stated
evaluation of their coefficient functions on a common global
atlas unevaluated.
\end{enumerate}
\end{proof}

\Uses{the hypotheses, notation and definitions stated here; \cref{inv:A3.2}; \cref{inv:C3.7}; \cref{inv:eq:C-G-Green}; \cref{inv:C3.2}}

\section{Open arithmetic questions for the normalized norm}\label{ari:open:norm}



Without hypothesis $H$, the arithmetic nature of the metric norm remains open at the
nontorsion points of \cref{ari:C.7.4}. An explicit small instance is
\[
 E:y^2=x^3-140x+784,\qquad P=(0,28),\qquad j=-3375,
 \qquad N(P)=2e^{-\widehat h(P)}.
\]
Here $P$ is a primitive generator modulo the rational
$2$-torsion and
\[
 \widehat h(P)=
 0.46736567718147128000752020126284678363562034370331\ldots.
\]
The exact unresolved assertion is that $2e^{-\widehat h(P)}$
is transcendental. The corresponding example of conductor
$361$, $y^2+y=x^3-38x+90$, $P=(0,9)$, has
$N=e^{-\widehat h(P)}$ and no finite local contribution.
The relative-period proof isolates its obstruction: an
unconditional theorem excluding
\eqref{ari:eq:cp-relation} for the positive diagonal central
direction. A theorem about linear $1$-period relations
does not supply that exclusion. The attempted approaches and the exact CM theta comparison are printed in the preceding period and CM-normalization sections;
their source checks and computational records are retained in the audit files.
These arguments do not resolve the stated exclusion.

\clearpage


\clearpage
\begin{thebibliography}{129}
\providecommand{\natexlab}[1]{#1}
\providecommand{\url}[1]{\texttt{#1}}
\expandafter\ifx\csname urlstyle\endcsname\relax
  \providecommand{\doi}[1]{doi: #1}\else
  \providecommand{\doi}{doi: \begingroup \urlstyle{rm}\Url}\fi

\bibitem[Abasheva and D{\'e}ev(2024)]{AbashevaDeev2024AB}
Anna Abasheva and Rodion D{\'e}ev.
\newblock Complex surfaces with many algebraic structures.
\newblock \emph{International Mathematics Research Notices}, 2024\penalty0
  (9):\penalty0 7379--7400, 2024.
\newblock \doi{10.1093/imrn/rnad190}.
\newblock URL \url{https://arxiv.org/pdf/2303.10764}.
\newblock Verified in this writing on 18 November 2026 at these loci only.
  arXiv2303.10764, Section2.2 pp4–6.. Scope: Read diagonal Hopf multipliers,
  the line-bundle quotient parameter, exponential map and explicit translation
  extension formulas.

\bibitem[Abramovich et~al.(2002)Abramovich, Karu, Matsuki, and
  W{\l}odarczyk]{AKMW02AB}
Dan Abramovich, Kalle Karu, Kenji Matsuki, and Jaros{\l}aw W{\l}odarczyk.
\newblock Torification and factorization of birational maps.
\newblock \emph{Journal of the American Mathematical Society}, 15\penalty0
  (3):\penalty0 531--572, 2002.
\newblock \doi{10.1090/S0894-0347-02-00396-X}.
\newblock URL \url{https://personal.math.ubc.ca/~karu/papers/semist.pdf}.
\newblock Verified in this writing on 18 November 2026 at these loci only.
  Theorem1.6.1p543,Theorem2.7.1pp551-552,Sections5.5-5.6pp564-566. Scope:
  Toroidal star factorization;\allowbreak{} finite cone-complex refinement and analytic
  monomial realization.

\bibitem[Adamczewski and Faverjon(2017)]{adamczewski2017mahler}
Boris Adamczewski and Colin Faverjon.
\newblock M\'ethode de {Mahler}: relations lin\'eaires, transcendance et
  applications aux nombres automatiques.
\newblock \emph{Proceedings of the London Mathematical Society}, 115:\penalty0
  55--90, 2017.
\newblock \doi{10.1112/plms.12038}.
\newblock URL
  \url{https://faverjon.perso.math.cnrs.fr/AdamczewskiFaverjon_mahler1.pdf}.
\newblock Verified in this writing on 18 November 2026 at these loci only.
  Published author PDF, system(1.2), regular-point definition and Theorem1.1
  p56 (PDFp2).. Scope: The source requires algebraic-coefficient power series
  and an algebraic regular argument. The actual g has transcendental cubic
  coefficient and its given elliptic logarithm u is transcendental, as
  separately proved in C7.12.

\bibitem[Agboola and Howard(2006)]{agboola2006}
Adebisi Agboola and Benjamin Howard.
\newblock Anticyclotomic {Iwasawa} theory of {CM} elliptic curves.
\newblock \emph{Annales de l'Institut Fourier}, 56\penalty0 (4):\penalty0
  1001--1048, 2006.
\newblock URL
  \url{https://web.math.ucsb.edu/~agboola/papers/bisi_aif_2006.pdf}.
\newblock Verified in this writing on 18 November 2026 at these loci only.
  Published PDF pp1001–1002;\allowbreak{} Rubin AppendixA setup p1042, TheoremA.1 and
  LemmaA.2 with proof p1043 (PDFpp2–3,43–44).. Scope: Recorded
  maximal-order CM, class number one, p>3 good ordinary setup and the appendix
  sign convention. TheoremA.1 gives the one-prime height;\allowbreak{} LemmaA.2 gives the
  relation among the three projected heights. No nondegeneracy theorem for an
  entire height pairing is attributed to this passage.

\bibitem[Angella(2018)]{Angella2018}
Daniele Angella.
\newblock Hodge numbers of a hypothetical complex structure on {$S^6$}.
\newblock \emph{Differential Geometry and its Applications}, 57:\penalty0
  105--120, 2018.
\newblock \doi{10.1016/j.difgeo.2017.10.012}.
\newblock URL \url{https://arxiv.org/abs/1705.10518}.
\newblock Verified in this writing on 18 November 2026 at these loci only.
  arXiv1705.10518v2,10November2017: §4 pp13–15 and Figure5 p16. Scope: Read
  each displayed boundary-entry morphism, the exact-sequence alternating sum
  and the relation hBC22=2hBC21\ensuremath{-}2h01+2, together with Figure5.
  Under h20=h02=0 and h01=1, the formula yields hBC22=2h. The exact-sequence
  inputs and specialization agree with the manuscript. The preprint version
  is2017;\allowbreak{} the cited journal publication is2018. This is a bounded formula
  comparison, not a claim that the article proves existence of a complex
  structure on S6.

\bibitem[{Anthropic}(2026)]{AnthropicEffort2026}
{Anthropic}.
\newblock Effort.
\newblock Claude Platform official documentation, accessed 17 November 2026,
  2026.
\newblock URL
  \url{https://platform.claude.com/docs/en/build-with-claude/effort}.
\newblock Verified in this writing on 18 November 2026 at these loci only. How
  effort works and Effort levels;\allowbreak{} lines138–155.. Scope: Response-token effort
  control, including thinking.

\bibitem[Atiyah and Singer(1968)]{AtiyahSinger1968IIIE}
Michael~F. Atiyah and Isadore~M. Singer.
\newblock The index of elliptic operators: Iii.
\newblock \emph{Annals of Mathematics}, 87\penalty0 (3):\penalty0 546--604,
  1968.
\newblock \doi{10.2307/1970717}.
\newblock URL
  \url{https://math.uchicago.edu/~shmuel/tom-readings/AS%20III.pdf}.
\newblock Verified in this writing on 18 November 2026 at these loci only.
  Holomorphic Lefschetz Theorem4.6p566;\allowbreak{}equations4.4-4.5p565. Scope: Original
  scan rendered and visually read. Compact complex manifold, holomorphic vector
  bundle, finite group action, fixed-locus normal eigenbundles. Application
  uses roots of unity then equality of Laurent characters. The index of
  elliptic operators III, original pp564–565 (PDF20–21),
  equations4.1–4.2, Riemann–Roch Theorem4.3 and its displayed proof..
  Scope: Original scan visually opened. The Dolbeault symbol and Todd-class
  calculation applies to the compact complex manifold and holomorphic bundle as
  stated. This adds Theorem4.3 to the separately recorded holomorphic Lefschetz
  loci, without claiming a full reading of the60-page article.

\bibitem[B{\u a}nic{\u a} and St{\u a}n{\u a}{\c s}il{\u a}(1976)]{BS76}
Constantin B{\u a}nic{\u a} and Octavian St{\u a}n{\u a}{\c s}il{\u a}.
\newblock \emph{Algebraic Methods in the Global Theory of Complex Spaces}.
\newblock Editura Academiei and John Wiley \& Sons, Bucharest and London, 1976.
\newblock URL \url{https://djvu.online/file/2hQiOnz8WuoQk}.
\newblock Verified in this writing on 18 November 2026 at these loci only. 1976
  English book, Chapter III: Theorem3.4 pp116–120;\allowbreak{} Corollary3.5 pp120–121;\allowbreak{}
  Theorem4.1 statement p124;\allowbreak{} Theorems4.10 and4.12 pp132–135. Scope: Opened
  the original-book text reproduction. Read the canonical fibre/tensor and
  arbitrary-base-change morphisms, the vanishing corollary proof, the finite
  free-complex statement, the global base-change criterion and
  continuity/Euler-characteristic proof. Properness and relative flatness are
  explicit;\allowbreak{} the converse constant-dimension criterion requires reduced base.
  Theorem4.1 p124 supplies the functorial bounded finite free complex under
  flatness;\allowbreak{} its full pp124–130 construction was not separately read
  end-to-end. Theorem3.4 proof includes completed local modules and exact
  tensor maps;\allowbreak{} no algebraicity or projectivity is imposed. BS76push is the
  same-book bibliographic alias, not a second source reading.

\bibitem[Belabas et~al.(2019)Belabas, Bernardi, and
  Perrin-Riou]{CgzBelabas2019}
Karim Belabas, Dominique Bernardi, and Bernadette Perrin-Riou.
\newblock La constante de {Manin} et le degr\'e modulaire d'une courbe
  elliptique.
\newblock \emph{Publications math\'ematiques de Besan\c con}, 2019/2:\penalty0
  81--103, 2019.
\newblock \doi{10.5802/pmb.37}.
\newblock URL \url{https://www.numdam.org/item/10.5802/pmb.37.pdf}.
\newblock Verified in this writing on 18 November 2026 at these loci only.
  AppendixB,p102,unnumbered26569a1/a2block. Scope: Manin1,modular degree59368
  visually read.

\bibitem[Bertolin et~al.(2022)Bertolin, Philippon, Saha, and
  Saha]{bertolin2022}
Cristiana Bertolin, Patrice Philippon, Biswajyoti Saha, and Ekata Saha.
\newblock Semi-abelian analogues of {Schanuel} conjecture and applications.
\newblock \emph{Journal of Algebra}, 596:\penalty0 250--288, 2022.
\newblock URL \url{https://arxiv.org/pdf/2010.15170v2}.
\newblock Verified in this writing on 18 November 2026 at these loci only.
  arXiv2010.15170v2, Definitions2.3-2.4p8;\allowbreak{} Construction3.4p12 and Lemma3.5p13;\allowbreak{}
  Proposition4.12pp19-20 and Remark4.13p20;\allowbreak{} Section6.1pp26-27 and
  errataTable2pp28-29. Scope: Semiabelian exponential and generalized period
  coordinates;\allowbreak{} exact unipotent dimension 2dimB+dimZ, alternating bracket and
  antisymmetric lift;\allowbreak{} C.7.4 now states Conjecture 4.7 as its explicit open
  hypothesis H;\allowbreak{} the group and matrix identities remain unconditional. Reopened
  19 November 2026: full arXiv PDF pp. 8, 12--13, 18--20. Relevant attached TeX
  passages were read in full and matched to the primary PDF locators.

\bibitem[Bertolini et~al.(2012)Bertolini, Darmon, and Prasanna]{CrBDP2012}
Massimo Bertolini, Henri Darmon, and Kartik Prasanna.
\newblock {$p$}-adic {Rankin} {L}-series and rational points on {CM} elliptic
  curves.
\newblock \emph{Pacific Journal of Mathematics}, 260\penalty0 (2):\penalty0
  261--303, 2012.
\newblock \doi{10.2140/pjm.2012.260.261}.
\newblock URL
  \url{https://www.math.mcgill.ca/darmon/pub/Articles/Research/52.BDP2/bdp-pjm.pdf}.
\newblock Verified in this writing on 18 November 2026 at these loci only.
  Published PDF: Introduction setup p261 and Theorem1 with all preceding
  prime/CM hypotheses p262 (PDFpp2–3).. Scope: The Katz p-adic
  value/formal-logarithm identity modulo K-star in analytic rank one at a split
  good prime for maximal-order CM. Corrected the bibliography title to the
  exact title in this DOI-bearing PDF.

\bibitem[Bertrand(1984)]{bertrand1984}
Daniel Bertrand.
\newblock Propri{\'e}t{\'e}s arithm{\'e}tiques de fonctions th{\^e}ta {\`a}
  plusieurs variables.
\newblock In \emph{Number Theory, Noordwijkerhout 1983}, volume 1068 of
  \emph{Lecture Notes in Mathematics}, pages 17--22. Springer, 1984.
\newblock \doi{10.1007/BFb0099438}.
\newblock URL \url{https://link.springer.com/chapter/10.1007/BFb0099438}.
\newblock Unverified in this writing at this assembly stage;\allowbreak{} no mathematical
  dependence is authorized by an earlier source-note alone.

\bibitem[Bertrand(1995)]{bertrand1995}
Daniel Bertrand.
\newblock Points rationnels sur les sous-groupes compacts des groupes
  alg{\'e}briques.
\newblock \emph{Experimental Mathematics}, 4\penalty0 (2):\penalty0 145--151,
  1995.
\newblock \doi{10.1080/10586458.1995.10504315}.
\newblock URL
  \url{https://www.maths.tcd.ie/EMIS/journals/EM/restricted/4/4.2/bertrand.ps}.
\newblock Verified in this writing on 18 November 2026 at these loci only.
  Original public PostScript, converted locally to PDF;\allowbreak{} all7published
  pages145–151, including Theorem2.1, Proposition2.2 and Theorem3.1 with
  proof.. Scope: Read the explicit antisymmetric endomorphism h-h-dual, the
  resulting Ribet section and the stated lifting/compactness relation.
  Conversion of the original PostScript supplies this reading;\allowbreak{} no modern source
  substitution.

\bibitem[Bertrand(2017)]{bertrand2017}
Daniel Bertrand.
\newblock Transcendence problems related to heights.
\newblock Fields Institute, Heights and Applications to Unlikely Intersections,
  February 13--17, 2017.
\newblock URL
  \url{https://webusers.imj-prg.fr/~daniel.bertrand/Recherche/rpdf/DB_Fieldspost.pdf}.
\newblock Verified in this writing on 18 November 2026 at these loci only.
  Author Fields Institute slides,13–17February2017, all15slides read;\allowbreak{}
  operative loci slides5–8 and final comment(c).. Scope: The explicit real
  compactness determinant and Green-function formula are matched. The pointwise
  conjecture, known CM antisymmetric case and functional results remain
  attributed with their own stated status.

\bibitem[Bierstone and Milman(1995)]{BM95AB}
Edward Bierstone and Pierre~D. Milman.
\newblock Canonical desingularization in characteristic zero by blowing up the
  maximum strata of a local invariant, 1995.
\newblock URL \url{https://arxiv.org/pdf/alg-geom/9508005}.
\newblock Verified in this writing on 18 November 2026 at these loci only.
  Section0.2(2)p3;\allowbreak{}Theorem1.6 and noncompact analytic extension
  p5;\allowbreak{}Remarks1.7(3)p6 and1.16p11;\allowbreak{}Theorem1.10p7. Scope: Analytic resolution
  isomorphic over regular locus, locally finite construction and
  normal-crossing principalization;\allowbreak{} now used to place the radial-chain
  triangulation in a smooth ambient manifold.

\bibitem[Bierstone and Milman(1997)]{BM97}
Edward Bierstone and Pierre~D. Milman.
\newblock Canonical desingularization in characteristic zero by blowing up the
  maximum strata of a local invariant.
\newblock \emph{Inventiones Mathematicae}, 128:\penalty0 207--302, 1997.
\newblock \doi{10.1007/s002220050141}.
\newblock URL \url{https://arxiv.org/pdf/alg-geom/9508005}.
\newblock Verified in this writing on 18 November 2026 at these loci only.
  ChapterI preprint alg-geom/9508005: Theorem1.6 and following analytic
  paragraph p5;\allowbreak{} Remark1.7 p6;\allowbreak{} Theorem1.10 p7;\allowbreak{} analytic-category setup (0.2) p3
  was opened in the earlier BM95AB record.. Scope: The downloaded edition
  contains30pages. Operative references to its absent Section13 were replaced
  with the inspected ChapterI statements. Smooth-centre analytic resolution and
  ideal principalization are used with their explicit local/proper-neighborhood
  hypotheses;\allowbreak{} no reading of Theorem13.3 is claimed. Theorem1.6 and following
  analytic-neighborhood paragraph, preprint p.5;\allowbreak{} ambient analytic scope (0.2),
  p.3. Scope: Opened embedded resolution by smooth centers, simultaneous NC
  final divisor, and the relatively compact or locally finite analytic
  formulation. Compactness of the central divisor gives the finite modification
  used here. Bibliographic identity: Exact title, authors and arXiv
  alg-geom/9508005;\allowbreak{} the theorem locators explicitly refer to the preprint, not
  journal pages. This is the same reading event, with no additional independent
  credit.

\bibitem[Biliotti and Ghigi(2021)]{BiliottiGhigi2021Root}
Leonardo Biliotti and Alessandro Ghigi.
\newblock Meromorphic limits of automorphisms.
\newblock \emph{Transformation Groups}, 26:\penalty0 1147--1168, 2021.
\newblock \doi{10.1007/s00031-020-09551-x}.
\newblock URL \url{https://arxiv.org/abs/1901.10724}.
\newblock Verified in this writing on 18 November 2026 at these loci only.
  arXiv1901.10724, Theorem1.2 and its surrounding graph-compactification setup
  p3.. Scope: The standing Fujiki class C assumption is retained in the
  receiving comparison;\allowbreak{} this event checks the exact statement and context.

\bibitem[Bost and K{\"u}nnemann(2009)]{bost2009}
Jean-Beno{\^i}t Bost and Klaus K{\"u}nnemann.
\newblock Hermitian vector bundles and extension groups on arithmetic schemes
  {II}. {The} arithmetic {Atiyah} extension.
\newblock \emph{Ast{\'e}risque}, 327:\penalty0 361--424, 2009.
\newblock URL \url{https://www.numdam.org/article/AST_2009__327__361_0.pdf}.
\newblock Verified in this writing on 18 November 2026 at these loci only.
  Published source title and abstract pp361–362;\allowbreak{} Theorems3.1.1–3.1.2 with
  surrounding equivalence discussion, p386 (PDFp27).. Scope: The compact-point
  criterion is stated for the universal vector extension of an abelian variety
  over a number field. The precise extension and its dual
  line-bundle/connection interpretation are retained;\allowbreak{} no identity with the
  diagonal metric norm is supplied by this theorem.

\bibitem[Bourdon et~al.(2017)Bourdon, Clark, and
  Stankewicz]{BourdonClarkStankewicz2017}
Abbey Bourdon, Pete~L. Clark, and James Stankewicz.
\newblock Torsion points on {CM} elliptic curves over real number fields.
\newblock \emph{Transactions of the American Mathematical Society},
  369\penalty0 (12):\penalty0 8457--8496, 2017.
\newblock \doi{10.1090/tran/6905}.
\newblock URL \url{https://arxiv.org/pdf/1411.2742}.
\newblock Verified in this writing on 18 November 2026 at these loci only.
  arXiv:1411.2742v3, Theorem4.15 and Remark4.16, p23.. Scope: Parish torsion
  restriction for F=Q(j(E)), and its separately stated CM-field extension. The
  source theorem statement was matched;\allowbreak{} no reading of Parish original proof is
  claimed.

\bibitem[Buchdahl(1999)]{Buchdahl1999}
Nicholas Buchdahl.
\newblock On compact kähler surfaces.
\newblock \emph{Annales de l'Institut Fourier}, 49\penalty0 (1):\penalty0
  287--302, 1999.
\newblock \doi{10.5802/aif.1674}.
\newblock URL \url{https://www.numdam.org/article/AIF_1999__49_1_287_0.pdf}.
\newblock Verified in this writing on 18 November 2026 at these loci only.
  Theorem11p297. Scope: Compact complex surface with even first Betti number
  admits a Kahler metric Theorem 11, printed p.297. Scope: Opened and matched
  the Kahler conclusion for a compact complex surface with even first Betti
  number.

\bibitem[Cai et~al.(2014)Cai, Shu, and Tian]{CgzCai2014}
Li~Cai, Jie Shu, and Ye~Tian.
\newblock Explicit {Gross--Zagier} and {Waldspurger} formulae.
\newblock \emph{Algebra \& Number Theory}, 8\penalty0 (10):\penalty0
  2523--2572, 2014.
\newblock \doi{10.2140/ant.2014.8.2523}.
\newblock URL \url{https://msp.org/ant/2014/8-10/ant-v8-n10-p05-p.pdf}.
\newblock Verified in this writing on 18 November 2026 at these loci only.
  Theorem1.1pp2524-2525. Scope: Primitive ring class character, coprime
  conductor, split Heegner primes;\allowbreak{} exact GZ pairing formula.

\bibitem[Call and Silverman(1993)]{callSilverman1993}
Gregory~S. Call and Joseph~H. Silverman.
\newblock Canonical heights on varieties with morphisms.
\newblock \emph{Compositio Mathematica}, 89\penalty0 (2):\penalty0 163--205,
  1993.
\newblock URL \url{https://www.numdam.org/item/CM_1993__89_2_163_0.pdf}.
\newblock Verified in this writing on 18 November 2026 at these loci only.
  Theorem4.1 and following Tate remark, journal pp183–184 (PDFpp22–23),
  both opened visually.. Scope: The height ratio limit uses a degree-one height
  on the base. The following elliptic-family remark gives the divisor-height
  formula with bounded error. Reading of the original Tate paper is not
  claimed.

\bibitem[Campana et~al.(2020)Campana, Demailly, and Peternell]{CDP2020}
Fr{\'e}d{\'e}ric Campana, Jean-Pierre Demailly, and Thomas Peternell.
\newblock The algebraic dimension of compact complex threefolds with vanishing
  second betti number.
\newblock \emph{Compositio Mathematica}, 156\penalty0 (4):\penalty0 679--696,
  2020.
\newblock \doi{10.1112/S0010437X19007802}.
\newblock Verified in this writing on 18 November 2026 at these loci only.
  Attached-source edition: Theorem 2.1 p.1;\allowbreak{} Theorem 2.2, Corollary 2.3,
  Proposition 2.4 p.2;\allowbreak{} also source lines272-350. Scope: Opened the actual
  printed hypotheses and conclusions and the full-fibre cotangent reduction.
  Bibliographic identity: Same Campana–Demailly–Peternell 2020 title and
  DOI10.1112/S0010437X19007802;\allowbreak{} attached-source pagination remains explicitly
  identified. This is the same reading event, with no additional independent
  credit. Attached-source edition: Notation4.5, Lemma4.6 and Proposition4.7 p7;\allowbreak{}
  Proposition4.11 and full proof pp8-9;\allowbreak{} Proposition7.3 statement p14 and full
  proof pp15-17, especially the reducible-curve inference p16. Scope: Read the
  exact conductor identities and the complete displayed proof of
  Proposition4.7, the full Proposition4.11 proof on pp8-9, and the full
  Proposition7.3 proof. Pages7 and16 were also rendered and visually read.
  Lemma4.6 cites a separate proof and is not treated as a full source-proof
  reading;\allowbreak{} the receiving finite-duality and discrepancy arguments are explicit.
  The reducible-curve sentence on p16 is the precisely rejected intermediate
  inference. These are attached-source edition pages, not unverified publisher
  pagination. Bibliographic identity: Same Campana–Demailly–Peternell 2020
  title and DOI10.1112/S0010437X19007802;\allowbreak{} attached-source pagination remains
  explicitly identified. This is the same reading event, with no additional
  independent credit.

\bibitem[Cartan(1952)]{Cartan1952B_AB}
Henri Cartan.
\newblock Faisceaux analytiques sur les vari{\'e}t{\'e}s de {Stein}:
  d{\'e}monstration des th{\'e}or{\`e}mes fondamentaux.
\newblock \emph{S{\'e}minaire Henri Cartan}, 4:\penalty0 1--15, 1952.
\newblock URL \url{https://www.numdam.org/article/SHC_1951-1952__4__A19_0.pdf}.
\newblock Verified in this writing on 18 November 2026 at these loci only.
  Expose19,Section5,pp19-14--19-15 (PDFpp15-16);\allowbreak{}Section4 preparatory lemmas
  pp19-11--19-13. Scope: TheoremB coherent cohomology vanishing on Stein
  manifolds, including degree one gluing.

\bibitem[Castella(2025)]{CrCastella2025}
Francesc Castella.
\newblock Tamagawa number conjecture for {CM} modular forms and
  {R}ankin--{S}elberg convolutions.
\newblock \emph{Proceedings of the London Mathematical Society}, 131\penalty0
  (4):\penalty0 e70089, 2025.
\newblock \doi{10.1112/plms.70089}.
\newblock URL \url{https://web.math.ucsb.edu/~castella/Katz.pdf}.
\newblock Verified in this writing on 18 November 2026 at these loci only.
  Current author PDF: TheoremA/Remark1.1.1 p2;\allowbreak{} period conventions and
  L(E/K)=L(E/Q)\textasciicircum{}2 explanation p3.. Scope: Recorded all
  full-CM, abelian torsion-extension, split-prime, p prime to6hK and conductor,
  and exact-different hypotheses. This is a bounded source check of the stated
  context, with no extension to the excluded -16 application or the inert
  prime3 in the last row.

\bibitem[Chen(2026)]{Chen2026HodgeRoot}
Kuan-Wen Chen.
\newblock On the {Hodge} structures of global smoothings of normal crossing
  varieties, 2026.
\newblock URL \url{https://arxiv.org/abs/2404.19229v2}.
\newblock Verified in this writing on 18 November 2026 at these loci only.
  arXiv2404.19229v2, Section2.1 p8.. Scope: Read the same-dimensional
  domination statement for the ddbar lemma and the Kahler case. Scope is the
  statement used in the receiving text.

\bibitem[{Claude} and {L. Alp{\"o}ge}(2026)]{AC2026}
{Claude} and {L. Alp{\"o}ge}.
\newblock The (3,4,infinity) modular family of 2-tori, completed at its three
  special points, is a complex structure on {$S^6$}.
\newblock 108-page manuscript, 2026.
\newblock URL \url{https://alpo.ge/s6.pdf}.
\newblock Verified in this writing on 18 November 2026 at these loci only.
  Original attached108-page PDF: Lemma9.12 pp72–74;\allowbreak{} Proposition9.13
  pp74–76;\allowbreak{} base-change clarification p92;\allowbreak{} Theorem10.5(b), proof Steps1 and4
  p88.. Scope: Complete inv:A1.0 known structure-sheaf input, Picard
  identification and injectivity on every full scheme fibre for inv:C3.1.
  Current reading does not credit the specialized BC/Green operator formulas to
  AC26. Exact attached 108-page edition SHA256
  283bba102dd1d5dc346af81b28145bdaaea6654398d5032e76e97bafb9a858f2: section10.5
  and Lemma10.7 p.90;\allowbreak{} rational component count and bielliptic product p.91;\allowbreak{}
  section10.6 starts p.94. Scope: Opened printed pp.90-93, including the actual
  coinvariant map, pair sequence, rank formula and stated hypotheses of the
  bielliptic example. The section10.6 heading was located on p.94;\allowbreak{} no complete
  section10.6 reading is claimed by this record. Original supplied manuscript:
  Remark1.2 p4 and Remark9.26 p83. Scope: Read the two exact remarks repeating
  the HKP Corollary1.3 bound and their stated scope. Only this comparison with
  the universal vector-field bound is recorded here;\allowbreak{} it is separate from the
  parent’s other AC26 source readings. Attached original108-page manuscript:
  Proposition9.24 with full proof pp82–83;\allowbreak{} full Remarks9.25–9.26 p83.
  Scope: Read the original PDF text and visually inspected both printed
  pages82–83. Opened the complete Proposition9.24 proof: actual etale
  quotient, reference-fan cocharacter, full fixed P1, and
  weights+1,\ensuremath{-}1. Read the complete finite-subgroup fixed-lattice
  criterion in Remark9.25 and the alternative two-point computation in
  Remark9.26. These establish the credited original reference-fan scope. The
  full arbitrary-basis fan classification, periodic flops and all-fan
  cap-freeness application are independently proved in the receiving F35. No
  entire-source or independent-authority credit is inferred. PDF SHA256
  283bba102dd1d5dc346af81b28145bdaaea6654398d5032e76e97bafb9a858f2 Same two
  authors, exact title, year2026, original URL and attached108-page edition;\allowbreak{}
  all aliases share only this bounded opening.

\bibitem[Cossec et~al.(2024)Cossec, Dolgachev, and Liedtke]{CDL24}
Fran\c{c}ois Cossec, Igor Dolgachev, and Christian Liedtke.
\newblock Enriques surfaces {I}, 2024.
\newblock URL
  \url{https://sites.lsa.umich.edu/idolga/wp-content/uploads/sites/1334/2024/08/EnriquesOne.pdf}.
\newblock Verified in this writing on 18 November 2026 at these loci only.
  April19,2024 manuscript:
  Proposition0.9.2p143;\allowbreak{}Section4.2p370;\allowbreak{}Lemma4.7.8,eq4.7.4,Theorem4.7.9p445.
  Scope: Exact relative-Picard/Brauer obstruction, Br(C)=0, H2(K,A)=0, and
  nontrivial elliptic-Jacobian local-to-global surjectivity.

\bibitem[{Creative Commons}(2026)]{CC01Universal}
{Creative Commons}.
\newblock Cc0 1.0 universal.
\newblock Official documentation;\allowbreak{} accessed 18 November 2026, 2026.
\newblock URL \url{https://creativecommons.org/publicdomain/zero/1.0/}.
\newblock Verified in this writing on 18 November 2026 at these loci only.
  Official deed, No Copyright, items1–2.. Scope: Public-domain dedication
  identified by the official deed.

\bibitem[Cremona and Siksek(2006)]{CrCremonaSiksek2006}
John~E. Cremona and Samir Siksek.
\newblock Computing a lower bound for the canonical height on elliptic curves
  over {$\mathbb Q$}.
\newblock In \emph{Algorithmic Number Theory}, volume 4076 of \emph{Lecture
  Notes in Computer Science}, pages 275--286. Springer, 2006.
\newblock URL \url{https://johncremona.github.io/papers/HeightLB.pdf}.
\newblock Verified in this writing on 18 November 2026 at these loci only.
  Author paper pp1–8: good-reduction subgroup setup, Proposition1 statement
  p4 with full proof p5, Propositions2–3 with complete proofs, and Theorem1
  p8.. Scope: Read the explicit good-reduction subgroup, lcm Tamagawa
  multiplier, local-height convention, finite denominator sum, real
  elliptic-logarithm isomorphism and complete interval-inverse-image criterion.
  These are the operative inputs to the freshly replayed saturation
  calculation. This record claims these exact pages, not the entire paper.

\bibitem[Darmon(2009)]{CrDarmon2009}
Henri Darmon.
\newblock Cycles on modular varieties and rational points on elliptic curves,
  2009.
\newblock URL
  \url{https://www.math.mcgill.ca/darmon/pub/Articles/Expository/13.Oberwolfach/paper.pdf}.
\newblock Verified in this writing on 18 November 2026 at these loci only.
  Section3pp4-5,Assumption3.1 and displayed canonical theta series. Scope:
  Class-number-one CM fields, conductorD\textasciicircum{}2, Hecke character
  and full coefficient formula.

\bibitem[de~Cataldo et~al.(2016)de~Cataldo, Migliorini, and
  Musta\c{t}\u{a}]{DMM16}
Mark Andrea~A. de~Cataldo, Luca Migliorini, and Mircea Musta\c{t}\u{a}.
\newblock Combinatorics and topology of proper toric maps, 2016.
\newblock URL \url{https://arxiv.org/pdf/1407.3497}.
\newblock Verified in this writing on 18 November 2026 at these loci only.
  arXiv1407.3497v2,Section2p6. Scope: Properness iff inverse image of fan
  support equals source support;\allowbreak{} explicit lattice-induced coordinate morphism.

\bibitem[Demailly(2012)]{Demailly2012E}
Jean-Pierre Demailly.
\newblock Complex analytic and differential geometry, 2012.
\newblock URL
  \url{https://www-fourier.univ-grenoble-alpes.fr/~demailly/manuscripts/agbook.pdf}.
\newblock Verified in this writing on 18 November 2026 at these loci only. 21
  June 2012 edition, ChapterII: Theorem4.19 pp.95-96;\allowbreak{} Theorem7.6 and
  Corollary7.7 p.113;\allowbreak{} Definition7.9 and Theorem7.12 pp.114-115;\allowbreak{} Theorem8.8
  p.118. Scope: Opened the finite local projection and cone bound giving its
  unique zero fibre, normal-curve regularity, finite proper normalization, and
  proper-image analyticity. The actual curve through both specified points is
  constructed in E.2.23.

\bibitem[DeMarco and Mavraki(2020)]{demarco2017}
Laura DeMarco and Niki~Myrto Mavraki.
\newblock Variation of canonical height and equidistribution.
\newblock \emph{American Journal of Mathematics}, 142\penalty0 (2):\penalty0
  443--473, 2020.
\newblock URL \url{https://arxiv.org/pdf/1701.07947}.
\newblock Verified in this writing on 18 November 2026 at these loci only.
  arXiv:1701.07947v2: Theorem1.1/Corollary1.2 p2;\allowbreak{} Section2.1 pp5–6;\allowbreak{}
  Theorem2.2 p7 and2.3 p8;\allowbreak{} Proposition3.1 p10;\allowbreak{} Section3.3 pp13–14;\allowbreak{} metric
  pp15–16;\allowbreak{} final proof Section4.3 p23.. Scope: Local height conventions and
  discriminant term, finite-place equality/continuous remainder, semipositive
  height metric and torsion-specialization equidistribution. The receiving
  calculation retains degree1/12 and divides curvature by that degree to obtain
  probability measure12m\_P;\allowbreak{} source prose about curvature is read with this
  required mass normalization.

\bibitem[Dloussky(2005)]{Dloussky2005AB}
Georges Dloussky.
\newblock On surfaces of class {VII}$_0^+$ with numerically anticanonical
  divisor, 2005.
\newblock URL \url{https://arxiv.org/pdf/math/0406387}.
\newblock Verified in this writing on 18 November 2026 at these loci only.
  Corollary2.25pp16-17;\allowbreak{}Lemma2.26p17;\allowbreak{}Proposition2.27 and full proof pp17-19.
  Scope: Positive-b2 minimal classVII with NAC divisor: cycle existence,
  coefficient-one cycle has no branch, Inoue/Inoue-Hirzebruch alternatives.

\bibitem[Dloussky and Teleman(2012)]{DlousskyTeleman2012AB}
Georges Dloussky and Andrei Teleman.
\newblock Infinite bubbling in non-{K}{\"a}hlerian geometry.
\newblock \emph{Mathematische Annalen}, 353\penalty0 (4):\penalty0 1283--1314,
  2012.
\newblock \doi{10.1007/s00208-011-0713-9}.
\newblock URL \url{https://arxiv.org/pdf/1012.5247}.
\newblock Verified in this writing on 18 November 2026 at these loci only.
  arXiv1012.5247v2,Proposition2.3 and full proof pp6-7. Scope: Relative Douady
  persistence of simply exceptional divisors under holomorphic surface
  submersions;\allowbreak{} normal sequence and vanishing stated explicitly.

\bibitem[Dloussky and Teleman(2020)]{DlousskyTeleman2020D}
Georges Dloussky and Andrei Teleman.
\newblock Smooth deformations of singular contractions of class {VII} surfaces.
\newblock \emph{Mathematische Zeitschrift}, 296:\penalty0 1521--1537, 2020.
\newblock \doi{10.1007/s00209-020-02481-0}.
\newblock URL \url{https://arxiv.org/pdf/1803.07631}.
\newblock Verified in this writing on 18 November 2026 at these loci only.
  arXiv1803.07631v2,Section3.2,Proposition3.3p11 and paragraph before
  Proposition3.4. Scope: Regularity of local-to-global deformation germ gives
  analytic smoothability/effectivity;\allowbreak{} no assertion that nonisolated section
  satisfies isolated theorem.

\bibitem[Doi(2009)]{Doi2009Root}
Mamoru Doi.
\newblock Gluing construction of compact complex surfaces with trivial
  canonical bundle.
\newblock \emph{Journal of the Mathematical Society of Japan}, 61\penalty0
  (3):\penalty0 853--884, 2009.
\newblock \doi{10.2969/jmsj/06130853}.
\newblock URL
  \url{https://www.jstage.jst.go.jp/article/jmath/61/3/61_3_853/_pdf/-char/en}.
\newblock Verified in this writing on 18 November 2026 at these loci only.
  Published pp880–883: Theorem5.2 statement, Example5.3 in full, and
  Section5.3 through Theorem5.5 and following discussion.. Scope: Matched the
  dual holomorphic normal bundles, anticanonical double divisors and opposite
  residues. Theorem5.5 asserts a smooth real six-dimensional total manifold
  with a smooth map and fibre complex structures continuous off the double
  locus. The preceding gluing map and convergence expression5.1 were read. No
  complete PDE proof from the earlier sections is claimed.

\bibitem[Du et~al.(2026{\natexlab{a}})Du, Guo, Kusakabe, and Xie]{DuEtAl2026v1}
Yun-Heng Du, Bin Guo, Yuta Kusakabe, and Song-Yan Xie.
\newblock Oka geometry and deformations of torus-fibred complex six-spheres.
\newblock arXiv:2610.00250v1, superseded version, 2026{\natexlab{a}}.
\newblock URL \url{https://arxiv.org/abs/2610.00250v1}.
\newblock Verified in this writing on 18 November 2026 at these loci only.
  arXiv2610.00250v1, Theorems1.3–1.4 p3;\allowbreak{} Lemma6.11 and full proof p81;\allowbreak{}
  Theorem6.12 and full proof pp81–82;\allowbreak{} Corollary6.13 and proof p82;\allowbreak{}
  Proposition6.16 and full proof p85;\allowbreak{} Proposition7.8 and full proof pp91–92;\allowbreak{}
  Theorem7.14 and full proof p95;\allowbreak{} PropositionsC.1–C.2 with full proofs p102;\allowbreak{}
  PropositionC.4 and full proof p103.. Scope: Matched the supplied TeX labels
  against the original107-page PDF. The branch-normal quotient, one-variable
  Kuranishi local-ring substitution, even-shift maps and quadratic composition
  correction were read at these loci. Earlier receiving-proof rederivation
  events remain separately scoped. Version1 Theorem1.4 is the exact2Z
  classification;\allowbreak{} the arXiv primary indexed record confirms23September2026. No
  complete reading of all107pages is claimed. Bibliographic identity: Identical
  four authors, version-specific title and arXiv2610.00250v1;\allowbreak{} original PDF
  matched the supplied TeX at the listed passages. Only those bounded readings
  are inherited. This is the same reading event, with no additional independent
  credit.

\bibitem[Du et~al.(2026{\natexlab{b}})Du, Guo, Kusakabe, and Xie]{DuEtAl2026v2}
Yun-Heng Du, Bin Guo, Yuta Kusakabe, and Song-Yan Xie.
\newblock Oka geometry of complex six-spheres.
\newblock arXiv:2610.00250v2, 2026{\natexlab{b}}.
\newblock URL \url{https://arxiv.org/abs/2610.00250v2}.
\newblock Verified in this writing on 18 November 2026 at these loci only.
  arXiv2610.00250v2, title/abstract/contents p1 and full introduction pp1–3;\allowbreak{}
  complete supplied source section headings and term searches.. Scope:
  The30-page PDF and supplied source omit the version1 tangent, Kuranishi,
  automorphism and exact-period sections. This comparison records their absence
  and makes no inference about a reason. The Oka conclusions on pp1–3 were
  read as source statements, without a whole-proof reading claim. Bibliographic
  identity: Identical four authors, version-specific title and
  arXiv2610.00250v2;\allowbreak{} original PDF matched the supplied TeX at the listed
  passages. Only those bounded readings are inherited. This is the same reading
  event, with no additional independent credit.

\bibitem[Du et~al.(2026{\natexlab{c}})Du, Guo, Kusakabe, and
  Xie]{DuGuoKusakabeXie2026v1}
Yun-Heng Du, Bin Guo, Yuta Kusakabe, and Song-Yan Xie.
\newblock Oka geometry and deformations of torus-fibred complex six-spheres,
  2026{\natexlab{c}}.
\newblock URL \url{https://arxiv.org/abs/2610.00250v1}.
\newblock Verified in this writing on 18 November 2026 at these loci only.
  arXiv2610.00250v1, Theorems1.3–1.4 p3;\allowbreak{} Lemma6.11 and full proof p81;\allowbreak{}
  Theorem6.12 and full proof pp81–82;\allowbreak{} Corollary6.13 and proof p82;\allowbreak{}
  Proposition6.16 and full proof p85;\allowbreak{} Proposition7.8 and full proof pp91–92;\allowbreak{}
  Theorem7.14 and full proof p95;\allowbreak{} PropositionsC.1–C.2 with full proofs p102;\allowbreak{}
  PropositionC.4 and full proof p103.. Scope: Matched the supplied TeX labels
  against the original107-page PDF. The branch-normal quotient, one-variable
  Kuranishi local-ring substitution, even-shift maps and quadratic composition
  correction were read at these loci. Earlier receiving-proof rederivation
  events remain separately scoped. Version1 Theorem1.4 is the exact2Z
  classification;\allowbreak{} the arXiv primary indexed record confirms23September2026. No
  complete reading of all107pages is claimed. Bibliographic identity: Identical
  four authors, version-specific title and arXiv2610.00250v1;\allowbreak{} original PDF
  matched the supplied TeX at the listed passages. Only those bounded readings
  are inherited. This is the same reading event, with no additional independent
  credit.

\bibitem[Duke and Imamo{\u g}lu(2008)]{DukeImamoglu2008}
William Duke and {\"O}zlem Imamo{\u g}lu.
\newblock The zeros of the {Weierstrass} {$\wp$}-function and hypergeometric
  series.
\newblock \emph{Mathematische Annalen}, 340\penalty0 (4):\penalty0 897--905,
  2008.
\newblock \doi{10.1007/s00208-007-0174-3}.
\newblock URL \url{https://www.math.ucla.edu/~wdduke/preprints/}.
\newblock Verified in this writing on 18 November 2026 at these loci only.
  Author manuscript pp1–3: equations(1)–(5), Theorem1 and
  Corollaries1–2.. Scope: The map x=1-1728/j and the hypergeometric quotient
  for the zero are matched with the receiving period convention. Bounded
  statement check;\allowbreak{} later proof sections were not newly read in this event.

\bibitem[Eichler and Zagier(1982)]{EichlerZagier1982}
Martin Eichler and Don Zagier.
\newblock On the zeros of the {Weierstrass} {$\wp$}-function.
\newblock \emph{Mathematische Annalen}, 258\penalty0 (4):\penalty0 399--407,
  1982.
\newblock \doi{10.1007/BF01453974}.
\newblock URL
  \url{https://people.mpim-bonn.mpg.de/zagier/files/doi/10.1007/BF01453974/fulltext.pdf}.
\newblock Verified in this writing on 18 November 2026 at these loci only.
  Published scan, main theorem p399 (PDFp1) opened visually;\allowbreak{} corollary p400
  read in extracted text.. Scope: Exact zero formula for the Weierstrass
  function, including logarithmic constant, factor144*pi*i*sqrt(6), vertical
  integral and the normalized weight6 and weight12 forms. The source formula
  and branch conventions are matched;\allowbreak{} no full-paper re-proof credit.

\bibitem[Einsiedler et~al.(2001)Einsiedler, Everest, and
  Ward]{einsiedler2001entropy}
Manfred Einsiedler, Graham Everest, and Thomas Ward.
\newblock Entropy and the canonical height.
\newblock \emph{Journal of Number Theory}, 91\penalty0 (2):\penalty0 256--273,
  2001.
\newblock \doi{10.1006/jnth.2001.2682}.
\newblock URL \url{https://people.math.ethz.ch/~einsiedl/height.pdf}.
\newblock Verified in this writing on 18 November 2026 at these loci only.
  Author PDF, Theorem5.1 and complete proof pp9–10;\allowbreak{} adjacent Theorem5.2
  pp10–11 read for its different normalization.. Scope: Theorem5.1 uses
  denominator multipliers bn from
  x(2\textasciicircum{}nQ)=an/bn\textasciicircum{}2 and time
  scale4\textasciicircum{}n, with entropy equal to canonical height. The cited
  identity gives no algebraicity evaluation of exp(-height).

\bibitem[Engel(2026)]{Engel2026}
Philip Engel.
\newblock Complex structures on {$S^6$}.
\newblock arXiv:2609.38442v1, 2026.
\newblock URL \url{https://arxiv.org/abs/2609.38442}.
\newblock Verified in this writing on 18 November 2026 at these loci only.
  Attached-source edition Lemma2.1pp3-4, Corollary2.2p4, Conjecture8.2p22.
  Scope: Exact T=t6P line Hom(T*L,L), section height1/6, local corrections4/3
  and1/2, and literal proposed pinching Attached original S6.tex, whole
  Remark8.5 at source lines1738–1760;\allowbreak{} matched whole remark on p24 of the
  attached-source PDF build.. Scope: Read the exact fixed elliptic surface,
  degree-zero generic-fibre line and translation data, Hom line pi*O(plus or
  minus E), uniform contracting/expanding scaling and Mumford/log-transform
  construction. This verifies the precise input scope of F50/F51. The PDF is
  explicitly an attached-source build, not an independently downloaded journal
  edition. Attached-source build: Lemma7.4 and complete proof pp18–19;\allowbreak{} cusp
  clause p19. Scope: Read the full finite-cap and cusp specialization lemma.
  The finite indices use the literal affine quotient presentation, intersection
  determinants and evenness;\allowbreak{} the cusp clause asserts its H2/H3 Mayer–Vietoris
  computation in one sentence. The current iv:rank\_two\_specialization proof
  supplies its complete integral matrices, actual unit-preserving family map
  and primitive cross-line construction. No omitted computation is represented
  as printed in Engel. PDF SHA256
  c9c17426430e44582dfe3fceca5b1486174d5181396a14be43a6a960f4f514a0 Same
  identified author/title/arXiv work or original edition;\allowbreak{} current verification
  applies only to the exact opened edition and locators stated here.
  Attached-source rendered edition: Remark8.3 pp22–23, including its explicit
  physical C* action and Aut0 assertion p23;\allowbreak{} full Remark8.5 p24. Scope: Read
  the actual printed passages. Remark8.3 describes the original scalar action,
  fixed edge and Seifert quotient;\allowbreak{} it does not print the present catalogue-wide
  finite-kernel/global-endpoint calculation. Remark8.5 specifies the full-fibre
  Hom linearization, contracting/expanding scalar, Mumford filling and
  logarithmic transforms;\allowbreak{} the current16-profile/26-pair sphere classification
  is not asserted there. Receiving proofs cite these known construction
  portions and retain their actual application maps. Same attached Philip Engel
  paper;\allowbreak{} the verified scope is exactly the two opened remarks and pages, not
  the entire work.

\bibitem[Engel et~al.(2025)Engel, de~Gaay~Fortman, and Schreieder]{EGFS25}
Philip Engel, Olivier de~Gaay~Fortman, and Stefan Schreieder.
\newblock Combinatorics and {H}odge theory of degenerations of abelian
  varieties: A survey of the {M}umford construction, 2025.
\newblock URL \url{https://arxiv.org/html/2507.15695v3}.
\newblock Verified in this writing on 18 November 2026 at these loci only.
  arXiv2507.15695v3, complete Constructions3.3 and3.8, Remark3.5 and
  Proposition3.7 with proof, printedpp21–25.. Scope: Read the
  positive-definite symmetric bilinear form, periodic rational tiling, quotient
  action, exact H1 extension and canonical polarization. The multivariable
  extension retains symmetric semidefinite forms with positive-definite sum and
  the explicit flat fan morphism. These hypotheses are retained in the
  receiving comparison.

\bibitem[Enokizono and Hashizume(2024)]{EH2024E}
Makoto Enokizono and Kenta Hashizume.
\newblock Semistable reduction for complex analytic spaces, 2024.
\newblock URL \url{https://arxiv.org/pdf/2404.05124}.
\newblock Verified in this writing on 18 November 2026 at these loci only.
  arXiv2404.05124v3, Definition2.13 and Lemma2.14 statement p6, full proof p7;\allowbreak{}
  adjacent Lemma2.15 and Definition2.16.. Scope: Strict analytic toroidal
  proper modifications correspond to fan subdivisions through the displayed
  Specan gluing. The projective hypotheses of the paper main theorem are not
  imported into this local analytic lemma.

\bibitem[Erban and Van~Gorder(2021)]{ErbanVanGorder2021}
Radek Erban and Robert~A. Van~Gorder.
\newblock On the sum of positive divisors functions.
\newblock \emph{Research in Number Theory}, 7\penalty0 (25), 2021.
\newblock \doi{10.1007/s40993-021-00240-6}.
\newblock URL \url{https://people.maths.ox.ac.uk/erban/papers/paperRNT.pdf}.
\newblock Primary PDF opened and page1 visually read on 19 November 2026.
  Equations(1.1)--(1.4), p.1, state the Eisenstein series normalizations and
  Ramanujan differential system. In the receiving notation x=exp(2pi i tau),
  P=E2, Q=E4 and R=E6;\allowbreak{} the explicit derivative conversion is printed in PaperI.
  This is a bounded source verification of those equations.

\bibitem[Fantechi and G{\"o}ttsche(2010)]{FG10}
Barbara Fantechi and Lothar G{\"o}ttsche.
\newblock Riemann--roch theorems and elliptic genus for virtually smooth
  schemes.
\newblock \emph{Geometry \& Topology}, 14\penalty0 (1):\penalty0 83--115, 2010.
\newblock \doi{10.2140/gt.2010.14.83}.
\newblock URL \url{https://msp.org/gt/2010/14-1/gt-v14-n1-p02-s.pdf}.
\newblock Verified in this writing on 18 November 2026 at these loci only.
  Publishedpp100–102, complete Corollary5.4, Proposition5.7, Theorem5.8 and
  Remark5.9 with their displayed proofs and local-complete-intersection setup..
  Scope: Read virtual Euler/top-Chern equality and lci deformation invariance
  at the stated proper scheme scope. The receiving analytic calculation is
  written through its own tangent/cotangent complex and characteristic classes;\allowbreak{}
  the source supplies the exact algebraic counterpart.

\bibitem[Friedman(2016)]{Friedman2016AB}
Robert Friedman.
\newblock On the geometry of anticanonical pairs, 2016.
\newblock URL \url{https://arxiv.org/pdf/1502.02560}.
\newblock Verified in this writing on 18 November 2026 at these loci only.
  Definition1.1pp4-5;\allowbreak{}Lemma1.2p5;\allowbreak{}Lemma2.2p11;\allowbreak{}Lemma2.7p14. Scope: Rational
  anticanonical-pair charge equals Euler of complement, interior/corner blowup
  changes, nonnegative charge and zero precisely for toric pairs.

\bibitem[Fujiki(2009)]{Fujiki2009D}
Akira Fujiki.
\newblock Automorphisms of parabolic {Inoue} surfaces, 2009.
\newblock URL \url{https://arxiv.org/pdf/0903.5374}.
\newblock Verified in this writing on 18 November 2026 at these loci only.
  Theorem1.1pp1-2;\allowbreak{}Section2 equations1-8p4;\allowbreak{}equations9-13p5;\allowbreak{}proof of
  Theorem1.1p8. Scope: Exact parabolic toric charts, normal line degree, full
  automorphism group and C-star scaling.

\bibitem[Fujiki and Pontecorvo(2010)]{FujikiPontecorvo2010AB}
Akira Fujiki and Massimiliano Pontecorvo.
\newblock Anti-self-dual bihermitian structures on {Inoue} surfaces.
\newblock \emph{Journal of Differential Geometry}, 85\penalty0 (1):\penalty0
  15--71, 2010.
\newblock \doi{10.4310/jdg/1284557925}.
\newblock URL
  \url{https://archive.ymsc.tsinghua.edu.cn/pacm_download/21/834-JDG-85-1-A2-fujiki.pdf}.
\newblock Verified in this writing on 18 November 2026 at these loci only.
  Section3.2pp20-21;\allowbreak{}Section3.3 equations8-10 and Lemmas3.2-3.3pp21-22;\allowbreak{}proof of
  Lemma3.2p68. Scope: Parabolic/hyperbolic/half Inoue curves and canonical
  divisors;\allowbreak{} half Inoue has twisted anticanonical divisor and empty untwisted
  system;\allowbreak{} disconnected anticanonical classification with stated fundamental
  group.

\bibitem[Fujimoto(1988)]{Fujimoto1988Root}
Yoshio Fujimoto.
\newblock Logarithmic transformations on elliptic fiber spaces.
\newblock \emph{Journal of Mathematics of Kyoto University}, 28\penalty0
  (1):\penalty0 91--110, 1988.
\newblock URL
  \url{https://scispace.com/pdf/logarithmic-transformations-on-elliptic-fiber-spaces-1uinrvp8jj.pdf}.
\newblock Verified in this writing on 18 November 2026 at these loci only.
  Reproduced original journal text: full standing definition pp.91-92;\allowbreak{}
  Theorem2.8 and Corollary2.10 pp.96-97;\allowbreak{} Example3.4(1) p.99. Scope: Opened and
  matched all standing hypotheses: a proper surjective holomorphic map between
  complex manifolds, connected fibres and smooth elliptic general fibre, with
  no projectivity or section hypothesis. The existing compact Hopf-threefold
  map satisfies those hypotheses and explicitly refutes the printed
  divisibility and fundamental-group assertions. The proposed cyclic cover
  normalizes to m disjoint global graphs. Example3.4(1) supplies prior
  occurrence of a Kodaira-dimension-one component;\allowbreak{} no use of the refuted
  assertions supports the proof. Bibliographic identity: Identical author,
  title, journal28(1),pp91–110,1988 and reproduced source URL. Only the
  recorded pp91–92,96–97,99 passages are inherited. This is the same
  reading event, with no additional independent credit.

\bibitem[Fujimoto(1990)]{Fujimoto1990AB}
Yoshio Fujimoto.
\newblock Generalized logarithmic transformations---elliptic bundle case.
\newblock In K.~Konno, M.-H. Saito, and S.~Usui, editors, \emph{Proceedings of
  the Meeting and the Workshop ``Algebraic Geometry and Hodge Theory'', Vol.
  I}, volume~16 of \emph{Hokkaido University Technical Report Series in
  Mathematics}, pages 81--102. Hokkaido University, 1990.
\newblock \doi{10.14943/5135}.
\newblock URL
  \url{https://eprints.lib.hokudai.ac.jp/repo/huscap/all/5450/16.pdf}.
\newblock Verified in this writing on 18 November 2026 at these loci only.
  Contribution in a263-page conference volume: Definition1.1 printedpp85–86
  (PDF89–90), preceding introduction printedpp81–84, and printedpp93–94
  (PDF97–98).. Scope: Read the Cartier logarithmic cocycle definition,
  codimension-at-least-three extension statement and Example1.15 cocycle1/(xy),
  including the stated1990 relative-compactification question. No general
  nonexistence claim is inferred.

\bibitem[Gasbarri(2010)]{gasbarri2010analytic}
Carlo Gasbarri.
\newblock Analytic subvarieties with many rational points.
\newblock \emph{Mathematische Annalen}, 346\penalty0 (1):\penalty0 199--243,
  2010.
\newblock \doi{10.1007/s00208-009-0394-9}.
\newblock URL \url{https://preprints.ihes.fr/storage/newscheiderlang.pdf}.
\newblock Verified in this writing on 18 November 2026 at these loci only.
  Definition3.1 p12;\allowbreak{} Definition3.7 pp14–15;\allowbreak{} Theorem4.16 and its proof
  pp25–26.. Scope: LG-germs require a formal subscheme over a number field,
  local analyticity and integral/factorial coefficient control. The present
  exact tangent obstruction is checked before invoking those bounds.

\bibitem[Gauduchon(1985)]{Gauduchon1985}
Paul Gauduchon.
\newblock Les m\'etriques standard d'une surface complexe compacte \`a premier
  nombre de {Betti} pair.
\newblock \emph{Ast\'erisque}, 126:\penalty0 129--135, 1985.
\newblock URL \url{https://numdam.org/item/AST_1985__126__129_0/}.
\newblock Verified in this writing on 18 November 2026 at these loci only.
  Original Astérisque126, pp129–133: operator conventions pp130–131;\allowbreak{}
  maximum-principle consequences p132;\allowbreak{} Theorem2 and complete proof pp132–133.
  Scope: Read the actual conformal transformation of L and L*, the
  one-dimensional adjoint kernel, the Fredholm-index argument, sign argument
  and strict positivity lemma. In complex dimension3 the conformal factor is
  the positive kernel generator raised to1/(m\ensuremath{-}1), yielding a
  Gauduchon metric unique up to positive scaling inside the fixed conformal
  class. This is the general almost-complex metric theorem, not only the
  even-b1 surface application in Theorem1.

\bibitem[Graf and Kov{\'a}cs(2014)]{GrafKovacs2014Root}
Patrick Graf and S{\'a}ndor~J. Kov{\'a}cs.
\newblock Potentially {Du Bois} spaces.
\newblock \emph{Journal of Singularities}, 8:\penalty0 117--134, 2014.
\newblock \doi{10.5427/jsing.2014.8i}.
\newblock URL
  \url{https://sites.math.washington.edu/~kovacs/current/papers/Graf_Kovacs__Potentially_Du_Bois_spaces.pdf}.
\newblock Verified in this writing on 18 November 2026 at these loci only.
  Published pp117–118 (PDF1–2): abstract, equation1.1, Definition1.3 and
  Theorem1.4.. Scope: Read the proper/projective complex context of the
  cohomological surjectivity statement and the Du Bois pair theorem. The
  receiving analytic comparison has its own displayed proof.

\bibitem[Gramain et~al.(1986)Gramain, Mignotte, and
  Waldschmidt]{gramain1986values}
Fran\c{c}ois Gramain, Maurice Mignotte, and Michel Waldschmidt.
\newblock Valeurs alg\'ebriques de fonctions analytiques.
\newblock \emph{Acta Arithmetica}, 47\penalty0 (2):\penalty0 97--121, 1986.
\newblock \doi{10.4064/aa-47-2-97-121}.
\newblock URL \url{https://matwbn.icm.edu.pl/ksiazki/aa/aa47/aa4721.pdf}.
\newblock Verified in this writing on 18 November 2026 at these loci only.
  Lemma1.0 p98;\allowbreak{} Theorem2.1 hypotheses and statement pp101–102;\allowbreak{} full proof
  pp102–105. PDF spreads2–5 opened visually.. Scope: Exact joint-field
  degree, height and growth inequality and the Liouville/Siegel factors used in
  the C7.12 envelope audit. Only the displayed model envelopes are compared;\allowbreak{} no
  lower bound for actual degree and no transcendence of N0 is inferred.

\bibitem[Grauert(1960{\natexlab{a}})]{Grauert1960D}
Hans Grauert.
\newblock Ein theorem der analytischen garbentheorie und die modulräume
  komplexer strukturen.
\newblock \emph{Publications Mathématiques de l'IHÉS}, 5:\penalty0 5--64,
  1960{\natexlab{a}}.
\newblock \doi{10.1007/BF02684746}.
\newblock URL \url{https://www.numdam.org/article/PMIHES_1960__5__5_0.pdf}.
\newblock Verified in this writing on 18 November 2026 at these loci only.
  Section6 HauptsatzeI,II,IIap59;\allowbreak{}Section7 Satz3p62,Satz5p64. Scope: Coherence,
  analytic formal functions, upper semicontinuity and constant-dimension base
  change, with flatness hypothesis retained Section 7, Satz 5, pp. 63–64.
  Scope: Opened full Satz 5. Also read corrigendum 1963 pp. 35–36
  (https://www.numdam.org/item/10.1007/BF02684291.pdf);\allowbreak{} it does not alter this
  statement. Author/title and edition matched to the opened item. Verification
  applies only to its recorded locator and edition;\allowbreak{} aliasing does not transfer
  preprint pagination to a published edition.

\bibitem[Grauert(1960{\natexlab{b}})]{Grauert60}
Hans Grauert.
\newblock Ein theorem der analytischen garbentheorie und die modulräume
  komplexer strukturen.
\newblock \emph{Publications Mathématiques de l'IHÉS}, 5:\penalty0 5--64,
  1960{\natexlab{b}}.
\newblock \doi{10.1007/BF02684746}.
\newblock URL \url{https://www.numdam.org/article/PMIHES_1960__5__5_0.pdf}.
\newblock Verified in this writing on 18 November 2026 at these loci only.
  Section 6, Hauptsatze I and IIa, printed p.59. Scope: Opened coherence and
  the analytic formal-functions comparison. The proof uses both degrees zero
  and one and explicitly checks cofinality of the exceptional and maximal-ideal
  filtrations. Section6 HauptsatzeI,II,IIap59;\allowbreak{}Section7 Satz3p62,Satz5p64.
  Scope: Coherence, analytic formal functions, upper semicontinuity and
  constant-dimension base change, with flatness hypothesis retained
  Bibliographic identity: Same 1960 Grauert direct-image article and printed
  Section6/7 locators;\allowbreak{} existing sign-source record names this same item. This
  is the same reading event, with no additional independent credit. Section 7,
  Satz 5, pp. 63–64. Scope: Opened full Satz 5. Also read corrigendum 1963
  pp. 35–36 (https://www.numdam.org/item/10.1007/BF02684291.pdf);\allowbreak{} it does not
  alter this statement. Author/title and edition matched to the opened item.
  Verification applies only to its recorded locator and edition;\allowbreak{} aliasing does
  not transfer preprint pagination to a published edition. Bibliographic
  identity: Same 1960 Grauert direct-image article and printed Section6/7
  locators;\allowbreak{} existing sign-source record names this same item. This is the same
  reading event, with no additional independent credit.

\bibitem[Grauert(1963)]{Grauert1963ErratumD}
Hans Grauert.
\newblock Berichtigung: Ein theorem der analytischen garbentheorie und die
  modulräume komplexer strukturen.
\newblock \emph{Publications Mathématiques de l'IHÉS}, 16:\penalty0 35--36,
  1963.
\newblock URL \url{https://www.numdam.org/item/PMIHES_1963__16__35_0.pdf}.
\newblock Verified in this writing on 18 November 2026 at these loci only.
  pp35-36,corrections1-2. Scope: Proof-radius and induction corrections;\allowbreak{}
  Section7 Satz2 exceptional set need not be closed. The receiving proof uses
  direct-image coherence and flat constant-dimension base change, not the
  corrected exceptional-set assertion.

\bibitem[Gross(1991)]{CrGross1991}
Benedict~H. Gross.
\newblock Kolyvagin's work on modular elliptic curves.
\newblock In \emph{L-functions and Arithmetic}, volume 153 of \emph{London
  Mathematical Society Lecture Note Series}, pages 235--256. Cambridge
  University Press, 1991.
\newblock URL
  \url{https://wstein.org/papers/bib/gross-kolyvagins_work_on_modular_elliptic_curves.pdf}.
\newblock Verified in this writing on 18 November 2026 at these loci only.
  Original scan: Section1 setup/formula1.1 p235;\allowbreak{} Theorem1.3 p236;\allowbreak{} opening of
  Section2 p237. PDFspreads1–2 opened visually.. Scope: Rank one and finite
  Tate–Shafarevich group from a nontorsion Heegner point. The non-CM
  restriction starts in the subsequent illustrative proof, not in Theorem1.3.
  No full re-proof of Kolyvagin is claimed.

\bibitem[Gross and Zagier(1986)]{CrGrossZagier1986}
Benedict~H. Gross and Don~B. Zagier.
\newblock Heegner points and derivatives of {L}-series.
\newblock \emph{Inventiones Mathematicae}, 84\penalty0 (2):\penalty0 225--320,
  1986.
\newblock \doi{10.1007/BF01388809}.
\newblock URL
  \url{https://archive.mpim-bonn.mpg.de/id/eprint/3578/1/preprint_1985_30.pdf}.
\newblock Verified in this writing on 18 November 2026 at these loci only.
  Original MPIM85-30 preprint: I.7.3 pp11–12 (PDFspread8) and its final proof
  V.2 p129 (PDFspread67), both opened visually;\allowbreak{} setup pp125–126 and
  TheoremV.2.1 p127 read in extracted text.. Scope: The rational nonzero factor
  connecting the derivative, a rational point height and a rational
  differential period. The exact trace and period-product step on p129 is
  matched;\allowbreak{} the surrounding BSD conjectures are not assumptions in the receiving
  proof.

\bibitem[Herblot(2012)]{herblot2012meromorphic}
Mathilde Herblot.
\newblock Algebraic points on meromorphic curves, 2012.
\newblock URL \url{https://webusers.imj-prg.fr/~mathilde.herblot/publi/SL.pdf}.
\newblock Verified in this writing on 18 November 2026 at these loci only.
  Definitions2.4–2.5 p11;\allowbreak{} Definitions5.4–5.5 and Theorem6.1 pp31–32..
  Scope: Arithmetic formal subschemes over residue number fields and their
  evaluation maps;\allowbreak{} compared with the exact transcendental logarithmic tangent
  of the proposed analytic graph. This records the theorem statement and
  defining hypotheses, not a re-proof of Herblot.

\bibitem[Hochschild and Serre(1953)]{HS53}
Gerhard Hochschild and Jean-Pierre Serre.
\newblock Cohomology of group extensions.
\newblock \emph{Transactions of the American Mathematical Society}, 74\penalty0
  (1):\penalty0 110--134, 1953.
\newblock \doi{10.1090/S0002-9947-1953-0052438-8}.
\newblock URL
  \url{https://www.sas.rochester.edu/mth/sites/doug-ravenel/otherpapers/hochschild-serre.pdf}.
\newblock Verified in this writing on 18 November 2026 at these loci only.
  ChapterIII Section4,Theorem2 and m=1 remark,pp129-130. Scope:
  Inflation-restriction-transgression exact five-term sequence, including final
  H2 map.

\bibitem[Honda and Viaclovsky(2025)]{HondaViaclovsky2025D}
Nobuhiro Honda and Jeff Viaclovsky.
\newblock Some remarks on fibrations in complex geometry.
\newblock \emph{Journal of Mathematical Study}, 58\penalty0 (4):\penalty0
  552--574, 2025.
\newblock \doi{10.4208/jms.v58n4.25.11}.
\newblock URL \url{https://arxiv.org/abs/2504.14420}.
\newblock Verified in this writing on 18 November 2026 at these loci only.
  arXiv2504.14420v1,Theorem1.2p1;\allowbreak{}Theorem3.2pp8-9;\allowbreak{}equation3.9p9. Scope: Opened
  exact printed converse from identity monodromy to equal b1. Receiving B4.17
  supplies explicit compact Hopf counterexample;\allowbreak{} preserves the inequality and
  equality-to-monodromy direction. Version 1: Proposition 3.1 p.7;\allowbreak{} section3.3
  Corollary3.7 and proof p.11. Scope: Opened the analytic deformation
  retraction and full-multiplicity fibre class, and the H1 invariant-cycle
  equality for an irreducible support. These page numbers refer to arXiv v1,
  not the differently paginated journal edition.

\bibitem[Huber and W{\"u}stholz(2023)]{huber2022}
Annette Huber and Gisbert W{\"u}stholz.
\newblock \emph{Transcendence and linear relations of 1-periods}, volume 284 of
  \emph{Memoirs of the American Mathematical Society}.
\newblock American Mathematical Society, 2023.
\newblock URL \url{https://arxiv.org/abs/1805.10104}.
\newblock Verified in this writing on 18 November 2026 at these loci only.
  arXiv1805.10104v8, Theorems6.1 and6.2 printedp43, period-division reduction
  and proofp44 (PDF59-60). Scope: Connected commutative algebraic groups over
  Qbar;\allowbreak{} nonzero algebraic exponential in a Qbar Lie subspace supplies an
  algebraic subgroup. The divided nonzero period case is explicitly justified
  in the source. Matched to the torus, compact-extension and
  universal-vector-extension slope arguments.

\bibitem[Illman(1972)]{Illman1972}
S{\"o}ren Illman.
\newblock \emph{Equivariant Algebraic Topology}.
\newblock PhD thesis, Princeton University, 1972.
\newblock URL
  \url{https://ncatlab.org/nlab/files/Illman_EquivariantAlgebraicTopology1972.pdf}.
\newblock Verified in this writing on 18 November 2026 at these loci only.
  Original1972 thesis scan: Theorem3.1 printedp48 (PDF53) opened visually;\allowbreak{}
  following proof and Lemma3.2 through printedp51 read by OCR.. Scope: The
  equivariant triangulation statement for differentiable G-manifolds is
  matched. The application is a compact smooth manifold with a finite
  involution, yielding finite-dimensional equivariant cells. This event does
  not claim a reading of the entire triangulation proof.

\bibitem[Illusie et~al.(2005)Illusie, Kato, and Nakayama]{IKN2005E}
Luc Illusie, Kazuya Kato, and Chikara Nakayama.
\newblock Quasi-unipotent logarithmic {Riemann--Hilbert} correspondences.
\newblock \emph{Journal of Mathematical Sciences, the University of Tokyo},
  12\penalty0 (1):\penalty0 1--66, 2005.
\newblock URL \url{https://www.ms.u-tokyo.ac.jp/journal/pdf/jms120101.pdf}.
\newblock Verified in this writing on 18 November 2026 at these loci only.
  Introduction pp1–4;\allowbreak{} Proposition6.1 and proof pp30–32;\allowbreak{} Theorem6.2 and full
  proof pp32–33;\allowbreak{} Theorem6.3 and full proof p34;\allowbreak{} Remarks6.3.1–3 pp34–35;\allowbreak{}
  Theorem6.4 p35.. Scope: Proper separated log-smooth fs analytic hypotheses,
  the log-smooth-base or exact-morphism alternative, torsion-free
  characteristic cokernels and unipotent coefficients are retained. Ordinary
  log-de-Rham and nilpotent Gauss–Manin conclusions were matched at this
  precise scope.

\bibitem[Illusie et~al.(2007)Illusie, Kato, and Nakayama]{IKN2007ErratumE}
Luc Illusie, Kazuya Kato, and Chikara Nakayama.
\newblock Erratum to ``quasi-unipotent logarithmic {Riemann--Hilbert}
  correspondences''.
\newblock \emph{Journal of Mathematical Sciences, the University of Tokyo},
  14\penalty0 (1):\penalty0 113--116, 2007.
\newblock URL \url{https://www.ms.u-tokyo.ac.jp/journal/pdf/jms140105.pdf}.
\newblock Verified in this writing on 18 November 2026 at these loci only.
  Entire mathematical note, journal pp113–115 (PDFpp1–3);\allowbreak{} addresses on
  PDFp4 were not read.. Scope: Theorem7.1 correction and full proof read,
  including the use of Theorem6.2(2), saturation, etale localization and
  flatness.

\bibitem[Istrati et~al.(2022)Istrati, Otiman, Pontecorvo, and
  Ruggiero]{IstratiEtAl2022Root}
Nicolina Istrati, Alexandra Otiman, Massimiliano Pontecorvo, and Matteo
  Ruggiero.
\newblock Toric {Kato} manifolds.
\newblock \emph{Journal de l'{\'E}cole polytechnique---Math{\'e}matiques},
  9:\penalty0 1347--1395, 2022.
\newblock \doi{10.5802/jep.208}.
\newblock URL \url{https://www.numdam.org/item/JEP_2022__9__1347_0.pdf}.
\newblock Verified in this writing on 18 November 2026 at these loci only.
  Section2
  pp1351-1353;\allowbreak{}Theorem5.2p1366;\allowbreak{}Proposition5.6pp1369-1370;\allowbreak{}Theorem7.1p1373;\allowbreak{}Proposition8.4p1377;\allowbreak{}Theorem8.5p1378;\allowbreak{}Theorem8.8p1380;\allowbreak{}Example9.10p1386.
  Scope: Kato collar/deck construction, exact monomial quotient, negative
  normal multipliers, proper Nakamura family and H0 forms/H1O in the stated
  toric hyperbolic scope.

\bibitem[Jetchev et~al.(2008)Jetchev, Lauter, and Stein]{CrJetchev2008}
Dimitar Jetchev, Kristin Lauter, and William Stein.
\newblock Explicit {H}eegner points: {K}olyvagin's conjecture and non-trivial
  elements in the {S}hafarevich--{T}ate group.
\newblock \emph{Journal of Number Theory}, 128:\penalty0 3394--3410, 2008.
\newblock URL \url{https://arxiv.org/pdf/0707.0032}.
\newblock Verified in this writing on 18 November 2026 at these loci only.
  arXiv:0707.0032v1, Introduction Theorem1.1 p2.. Scope: Primary research
  statement of equality of analytic and Mordell–Weil ranks for the verified
  simple-zero rows. Finite Sha is supported separately by Gross Theorem1.3.

\bibitem[Johansson(2016)]{Johansson2016}
Fredrik Johansson.
\newblock Arb: Efficient arbitrary-precision midpoint-radius interval
  arithmetic, 2016.
\newblock URL \url{https://arxiv.org/abs/1611.02831}.
\newblock Verified in this writing on 18 November 2026 at these loci only.
  arXiv1611.02831v1, Section2 pp2–4 and complete Section3 pp4–6, with
  Section2.2 exact-accuracy tests and Section3.2 radius arithmetic.. Scope:
  Matched midpoint-radius representation, arbitrary-precision dyadic midpoints,
  automatically upward-rounded magnitude bounds and the explicit
  compiler/elementary-operation assumptions. The received interval proof uses
  enclosure comparisons;\allowbreak{} this source reading adds no independent
  code-correctness or replay claim.

\bibitem[Kajiwara et~al.(2008{\natexlab{a}})Kajiwara, Kato, and
  Nakayama]{KKN2008PicardE}
Takeshi Kajiwara, Kazuya Kato, and Chikara Nakayama.
\newblock Analytic log {Picard} varieties.
\newblock \emph{Nagoya Mathematical Journal}, 191:\penalty0 149--180,
  2008{\natexlab{a}}.
\newblock \doi{10.1017/S0027763000025940}.
\newblock URL
  \url{https://www.cambridge.org/core/services/aop-cambridge-core/content/view/E8B336E3252DE1F448583CC5865DB1DA/S0027763000025940a.pdf/analytic_log_picard_varieties.pdf}.
\newblock Verified in this writing on 18 November 2026 at these loci only.
  Sections 7–9 in full, pp.164–174;\allowbreak{} Definition 7.1 pp.164–165;\allowbreak{} Section
  7.2 p.165;\allowbreak{} Section 8 hypotheses pp.165–166;\allowbreak{} Assumption 8.1 p.166;\allowbreak{} Theorem
  8.2 pp.166–167 and proof pp.168–172;\allowbreak{} Corollary 8.3 p.167 and proof
  pp.172–173;\allowbreak{} Section 9.3 pp.173–174. Scope: Goodness includes comparison
  isomorphisms, strict Hodge filtration conditions, and a log Hodge structure.
  Section8 retains properness, separatedness, log smoothness, exactness,
  torsion-free characteristic cokernels and Assumption8.1. Its maps give
  H1(Gm)\textasciicircum{}0 into the log Picard sheaf into
  H1(Gm,log)\textasciicircum{}0;\allowbreak{} H2 goodness enters the quotient and
  Corollary8.3. Section9.3 supplies goodness over a log-smooth base for the
  stated projective or log-Kahler maps. Section7.2 explicitly leaves
  equivalence of its alternate comparison definition unproved. The receiving
  application must retain its actual comparison maps and verified hypotheses.

\bibitem[Kajiwara et~al.(2008{\natexlab{b}})Kajiwara, Kato, and
  Nakayama]{KKN2008Root}
Takeshi Kajiwara, Kazuya Kato, and Chikara Nakayama.
\newblock Logarithmic abelian varieties, {Part I}: Complex analytic theory.
\newblock \emph{Journal of Mathematical Sciences, the University of Tokyo},
  15\penalty0 (1):\penalty0 69--193, 2008{\natexlab{b}}.
\newblock URL \url{https://www.ms.u-tokyo.ac.jp/journal/pdf/jms150104.pdf}.
\newblock Verified in this writing on 18 November 2026 at these loci only.
  Definition 1.2.3 pp.77–78;\allowbreak{} Definition 1.2.6 p.79;\allowbreak{} Lemma 1.2.9 and full
  proof pp.80–81;\allowbreak{} Sections 1.3.1–1.3.4 pp.81–82;\allowbreak{} Sections 5.1.1–5.1.3
  p.170;\allowbreak{} Theorem 5.2.1 statement p.171;\allowbreak{} construction begins pp.171–172;\allowbreak{}
  Theorem 5.3.2 statement p.180 and deduction from Proposition 5.3.5 p.181.
  Scope: Admissibility tests every face and nonnegative functional on its
  quotient lattices;\allowbreak{} nondegeneracy also tests the residual unit pairing. Lemma
  1.2.9 proves base-change preservation by identifying the saturated face
  kernels. The quotient uses the bounded-character subsheaf. The proper-model
  map is explicitly A(Sigma)=V(Sigma)/Y into the specified A. Theorem 5.2.1
  supplies invariant finite-orbit complete fans;\allowbreak{} 5.3.2 supplies their
  represented log-smooth proper models. Polarization enters the additional
  projective-model conclusion in 5.1.3. The receiving B4.24 ghost matrix and
  residual tests match these definitions.

\bibitem[Kamenova(n.d.)]{KamenovaToriRoot}
Ljudmila Kamenova.
\newblock Degenerations of 2-dimensional tori.
\newblock Author-hosted manuscript, n.d.
\newblock URL
  \url{https://www.math.stonybrook.edu/~kamenova/homepage_files/degen.pdf}.
\newblock Verified in this writing on 18 November 2026 at these loci only.
  Author four-page manuscript, all4pages;\allowbreak{} Theorem3.1 and proof pp2–3.. Scope:
  Proper family with Kahler total space and trivial canonical class hypotheses
  were read with the complete proof. These hypotheses are recorded in the
  receiving application.

\bibitem[Katzarkov et~al.(2025)Katzarkov, Lee, Lupercio, and
  Meersseman]{KatzarkovEtAl2025Root}
Ludmil Katzarkov, Kyoung-Seog Lee, Ernesto Lupercio, and Laurent Meersseman.
\newblock On {Hodge} polynomials for nonalgebraic complex manifolds.
\newblock \emph{Proceedings of the National Academy of Sciences}, 122\penalty0
  (15):\penalty0 e2415722122, 2025.
\newblock \doi{10.1073/pnas.2415722122}.
\newblock URL \url{https://www.pnas.org/doi/10.1073/pnas.2415722122}.
\newblock Verified in this writing on 18 November 2026 at these loci only.
  Publisher full text reproduced at PMC12012539: Example1 and Theorem4,
  lines352–358.. Scope: The explicit Hopf-surface and elliptic-product
  compactifications share an analytic open set. Theorem4 separately requires
  definable analytic isomorphism of the open subsets;\allowbreak{} that hypothesis is not
  supplied for the present family.

\bibitem[Knudson(2004)]{Knudson2004}
Kevin~P. Knudson.
\newblock Invariant chains and the homology of quotient spaces, 2004.
\newblock URL \url{https://arxiv.org/pdf/math/0409133}.
\newblock Verified in this writing on 18 November 2026 at these loci only.
  arXiv math/0409133, Theorem4.1 and full proof p10;\allowbreak{} Theorem4.2 and full proof
  p11.. Scope: Prime-order Smith fixed-point conclusions for finite-dimensional
  G-CW complexes, and the finite-complex acyclic quotient statement. The
  dimensions and coefficient fields are retained in the actual application.

\bibitem[Levy(2008)]{Levy08}
Roni~N. Levy.
\newblock Riemann--roch theorem for higher bivariant {K}-functors.
\newblock \emph{Annales de l'Institut Fourier}, 58\penalty0 (2):\penalty0
  571--601, 2008.
\newblock \doi{10.5802/aif.2361}.
\newblock URL \url{https://www.numdam.org/article/AIF_2008__58_2_571_0.pdf}.
\newblock Verified in this writing on 18 November 2026 at these loci only.
  Lemmas3.1-3.2,Section3.1,pp588-590. Scope: Exact comparison functor and
  global continuous bundle complex via exact-summand stabilization.

\bibitem[\L{}ojasiewicz(1964)]{Loj64}
Stanis{\l}aw \L{}ojasiewicz.
\newblock Triangulation of semi-analytic sets.
\newblock \emph{Annali della Scuola Normale Superiore di Pisa, Classe di
  Scienze}, 18\penalty0 (4):\penalty0 449--474, 1964.
\newblock URL \url{https://www.numdam.org/item/ASNSP_1964_3_18_4_449_0.pdf}.
\newblock Verified in this writing on 18 November 2026 at these loci only.
  Section3Theorem1p463,Theorem2p464. Scope: Compatible triangulation of locally
  finite semianalytic subsets of real analytic manifold;\allowbreak{} singular-atlas
  application separately tracked.

\bibitem[Manetti(1991)]{Manetti1991D}
Marco Manetti.
\newblock Normal degenerations of the complex projective plane.
\newblock \emph{Journal f{\"u}r die reine und angewandte Mathematik},
  419:\penalty0 89--118, 1991.
\newblock \doi{10.1515/crll.1991.419.89}.
\newblock URL \url{https://gdz.sub.uni-goettingen.de/id/PPN243919689_0419}.
\newblock Verified in this writing on 18 November 2026 at these loci only.
  Lemma1,p93;\allowbreak{} introductory scope on same page. Scope: Original journal scan
  visually opened: reduced compact complex space with isolated singularities
  and H2(Theta)=0 gives smooth analytic local-to-global deformation germ. Full
  spectral-sequence proof read.

\bibitem[Matsuki(2000)]{Matsuki00ErrAB}
Kenji Matsuki.
\newblock Erratum to the paper ``a note on the factorization theorem of toric
  birational maps after {Morelli} and its toroidal extension'', 2000.
\newblock URL \url{https://arxiv.org/abs/math/0006038}.
\newblock Verified in this writing on 18 November 2026 at these loci only.
  Section2,p2,Weak Factorization Theorem for Toroidal Birational Morphisms.
  Scope: Corrected smooth-stratum blowup/blowdown factorization over the
  target;\allowbreak{} no use of the refuted strong-factorization algorithm.

\bibitem[McHugh(2018)]{McHugh2018}
Andrew McHugh.
\newblock Bott--chern--aeppli, dolbeault and fr{\"o}licher on compact complex
  3-folds.
\newblock arXiv:1708.03251v4, 2018.
\newblock URL \url{https://arxiv.org/abs/1708.03251v4}.
\newblock Verified in this writing on 18 November 2026 at these loci only.
  arXiv1708.03251v4,15November2018: Theorem6.3 p20;\allowbreak{} complete displayed
  derivation pp20–23;\allowbreak{} Proposition6.4 and equation(7) p23;\allowbreak{} table p24 and §7.1
  pp24–25. Scope: Read the real pluriharmonic exact-sequence dimension
  identity and the displayed Frolicher/BC-Aeppli derivation through
  equation(7). Substitution of k\textasciicircum{}(p,q)=0 in intermediate
  degrees, b1=0, H\textasciicircum{}(1,1)(R)=0 and E2\textasciicircum{}(3,1)=0
  gives hBC11=2,hBC12=h,hBC22=\ensuremath{-}1+(h+1)+h=2h. The actual table
  boundary entries match. This checks the cited formulas in their manuscript
  specialization, not every assertion in the source paper.

\bibitem[McMullen(1996)]{McMullen1996Complex}
Curtis~T. McMullen.
\newblock Complex manifolds: Math 241 course outline.
\newblock University of California, Berkeley, course notes, 1996.
\newblock URL
  \url{https://people.math.harvard.edu/~ctm/home/text/class/berkeley/241/96/course/course.pdf}.
\newblock Verified in this writing on 18 November 2026 at these loci only.
  Author course notes, numbered item67, p15, together with surrounding
  items63–66.. Scope: Infinitesimal deformation space H1(Theta) and the curve
  duality with quadratic differentials. This is a course-note statement check,
  not a new proof of general Kuranishi theory.

\bibitem[Meersseman(2011)]{Meersseman2011Root}
Laurent Meersseman.
\newblock Foliated structure of the {Kuranishi} space and isomorphisms of
  deformation families of compact complex manifolds.
\newblock \emph{Annales scientifiques de l'{\'E}cole Normale Sup{\'e}rieure},
  44\penalty0 (3):\penalty0 495--525, 2011.
\newblock \doi{10.24033/asens.2148}.
\newblock URL \url{https://www.numdam.org/item/10.24033/asens.2148.pdf}.
\newblock Verified in this writing on 18 November 2026 at these loci only.
  Published p511 (PDF18): Theorem3, Corollaries3–4 and proof opening;\allowbreak{}
  preceding examples on PDF17.. Scope: Reduced base and locally constant h0 are
  retained in the local-isomorphism statement. The later proof pages were not
  read in this event.

\bibitem[Merkulov(1995)]{Merkulov1995D}
Sergey~A. Merkulov.
\newblock Moduli spaces of compact complex submanifolds of complex fibered
  manifolds.
\newblock \emph{Mathematical Proceedings of the Cambridge Philosophical
  Society}, 118\penalty0 (1):\penalty0 71--91, 1995.
\newblock \doi{10.1017/S0305004100073473}.
\newblock URL \url{https://doi.org/10.1017/S0305004100073473}.
\newblock Verified in this writing on 18 November 2026 at these loci only.
  Publisher extract for Moduli spaces of compact complex submanifolds of
  complex fibered manifolds, Mathematical Proceedings118(1),1995,pp71–91,
  DOI10.1017/S0305004100073473;\allowbreak{} abstract/extract paragraph on Kodaira
  completeness.. Scope: The opened extract states the compact-submanifold
  H1(normal)=0 completeness criterion, smooth parameter-space dimension
  H0(normal), and tangent-space identification. This is an extract-only check;\allowbreak{}
  the full article was not read.

\bibitem[Miller(2011)]{miller2011}
Robert~L. Miller.
\newblock Proving the {Birch} and {Swinnerton-Dyer} conjecture for specific
  elliptic curves of analytic rank zero and one.
\newblock \emph{LMS Journal of Computation and Mathematics}, 14:\penalty0
  327--350, 2011.
\newblock \doi{10.1112/S1461157011000180}.
\newblock URL \url{https://arxiv.org/pdf/1010.2431v3}.
\newblock Verified in this writing on 18 November 2026 at these loci only.
  arXiv:1010.2431v3: Definition1.1/Theorem1.2 p2;\allowbreak{} isogeny-invariance
  explanation p2;\allowbreak{} Section7 opening/Theorems7.1–7.2 p16;\allowbreak{} complete exception
  Table9 p23 and its setup p22.. Scope: The three analytic-rank-one curves have
  conductors3136,3136,361, below5000 and outside every exception conductor.
  This proves full BSD for those curves;\allowbreak{} the separate exact local data and
  certified height/derivative intervals isolate the positive integer Sha order
  as1.

\bibitem[Movasati(2024)]{Movasati24}
Hossein Movasati.
\newblock A course in complex geometry and holomorphic foliations.
\newblock Author version dated 13 May 2024, 2024.
\newblock URL \url{https://w3.impa.br/~hossein/myarticles/ComplexGeometry.pdf}.
\newblock Verified in this writing on 18 November 2026 at these loci only.
  Author version dated13May2024: Theorem8.5 and its complete displayed proof,
  printed pp80–81 (PDFpp94–95);\allowbreak{} browser PDF lines4700–4758.. Scope: Stein
  factorization of a proper holomorphic map, including connected fibres, finite
  second map, direct-image identity and uniqueness. This is the exact
  factorization used for the proper divisor-to-curve map in P2.6.

\bibitem[Nakayama and Ogus(2010)]{NO10}
Chikara Nakayama and Arthur Ogus.
\newblock Relative rounding in toric and logarithmic geometry.
\newblock \emph{Geometry \& Topology}, 14\penalty0 (4):\penalty0 2189--2241,
  2010.
\newblock \doi{10.2140/gt.2010.14.2189}.
\newblock URL \url{https://msp.org/gt/2010/14-4/gt-v14-n4-p09-s.pdf}.
\newblock Verified in this writing on 18 November 2026 at these loci only.
  Theorem5.1,Corollary5.2,Proposition5.3p2226. Scope: Proper separated exact
  relatively smooth log analytic morphism with fine base and relatively
  coherent source gives oriented topological fibre bundle;\allowbreak{} local triangulation
  is also stated.

\bibitem[Nakayama(2002)]{Nakayama2002Root}
Noboru Nakayama.
\newblock Global structure of an elliptic fibration.
\newblock \emph{Publications of the Research Institute for Mathematical
  Sciences}, 38\penalty0 (3):\penalty0 451--649, 2002.
\newblock \doi{10.2977/PRIMS/1145476270}.
\newblock URL \url{https://ems.press/content/serial-article-files/40808?nt=1}.
\newblock Verified in this writing on 18 November 2026 at these loci only.
  Published Theorem7.1.3 and its setup, p604, plus proof Steps1–4 pp604–606
  (PDF154–156);\allowbreak{} all three pages opened visually.. Scope: Read the locally
  bimeromorphically projective elliptic-fibration hypotheses,
  codimension-at-least-two omitted locus, equivalent extension and BP
  conditions, and explicit flattening morphism in Step3. The nondegenerate
  triangular pullback map in Step4 is also read. No complete seven-step proof
  review is claimed.

\bibitem[{National Institute of Standards and Technology}(2026)]{DLMFTheta2026}
{National Institute of Standards and Technology}.
\newblock {NIST Digital Library of Mathematical Functions}, chapters 20 and 23,
  2026.
\newblock URL \url{https://dlmf.nist.gov/}.
\newblock Verified in this writing on 18 November 2026 at these loci only. NIST
  DLMF version1.2.8: equations20.2.1-2,20.2.6,20.2.12;\allowbreak{} 20.4.3-5
  and20.4.8;\allowbreak{}20.5.1;\allowbreak{}20.7.5;\allowbreak{}20.13.1;\allowbreak{}23.3.9;\allowbreak{}23.6.1-2 and23.6.8-9;\allowbreak{}23.17.8.
  Additional continuation reading on 19 November 2026:
  equations20.7.26,20.7.30,23.5.2--3,23.18.4--6 and20.13.1 were opened and
  matched to the exact period, theta and eta transformations printed in PaperI.
  Scope: Opened every listed exact equation on the primary NIST pages,
  including its displayed nome and period conventions. The receiving conversion
  is vartheta\_j(z,tau)=theta\_j(pi*z|tau) and 2omega1=1,2omega3=tau. Under
  this conversion the heat coefficient is 1/(4pi*i) and the half-period value
  is
  pi\textasciicircum{}2(theta3\textasciicircum{}4+theta4\textasciicircum{}4)/3.
  Differentiating the locally normally convergent products yields eta'/eta and
  the Laurent coefficient of the logarithmic theta derivative;\allowbreak{} the proof
  explicitly gives this conversion and the elliptic principal-part comparison.
  This bounded formula-source reading does not claim a reading of the books
  cited by DLMF.

\bibitem[Nguyen(2023)]{nguyen2023}
Khoa~D. Nguyen.
\newblock Transcendence of polynomial canonical heights.
\newblock \emph{Mathematische Annalen}, 387:\penalty0 1--15, 2023.
\newblock URL \url{https://arxiv.org/pdf/2112.14937}.
\newblock Verified in this writing on 18 November 2026 at these loci only.
  arXiv2112.14937: Definition1.3 and Bottcher-domain setup p2;\allowbreak{} Theorem1.4 and
  Corollary1.6, p3.. Scope: Polynomial hypotheses, functional basin, common
  degree and non-integrability are recorded exactly. The receiving discussion
  gives the explicit elliptic duplication quotient and its ramification instead
  of applying this polynomial result to it.

\bibitem[Nishiguchi(1988)]{Nishiguchi1988IIRoot}
Kenji Nishiguchi.
\newblock Degeneration of {K3} surfaces, {II}.
\newblock Max-Planck-Institut f{\"u}r Mathematik preprint MPI/88-48, 1988.
\newblock URL
  \url{https://archive.mpim-bonn.mpg.de/id/eprint/3390/1/preprint_1988_48.pdf}.
\newblock Verified in this writing on 18 November 2026 at these loci only.
  MPI/88-48 original scan, introduction printedpp1–2 (PDF3–4);\allowbreak{} Theorem1.4
  and its source explanation printedp7 (PDF9);\allowbreak{} Example2.1 in full printedp13
  (PDF15), visually checked.. Scope: Read the K3 semistable-modification
  theorem with Hopf and CB components and the explicit diagonal Hopf quotient
  glued to an elliptic K3 fibre. Receiving citations to Sections2–3 are
  narrowed to these actual passages. This is a source statement/example
  reading, not an independent proof of the cited predecessor theorem.

\bibitem[Noguchi(2024)]{noguchi2024}
Junjiro Noguchi.
\newblock Analytic {Ax--Schanuel} for semi-abelian varieties and {Nevanlinna}
  theory.
\newblock \emph{Journal of the Mathematical Society of Japan}, 76\penalty0
  (1):\penalty0 1--22, 2024.
\newblock \doi{10.2969/jmsj/89588958}.
\newblock URL \url{https://www.ms.u-tokyo.ac.jp/~noguchi/An-axsch.pdf}.
\newblock Verified in this writing on 18 November 2026 at these loci only.
  Author PDF: subgroup condition1.2 and Theorem1.3 p2;\allowbreak{} Section3 proof pp7–9;\allowbreak{}
  jet description4.3–4.5 p10.. Scope: The entire-map semi-abelian
  Ax–Schanuel statement and its exact subgroup nondegeneracy hypothesis match
  the fixed-extension functional application. This event is a source reading,
  not an independent check of every external input to that proof.

\bibitem[Oda(1978)]{Oda1978D}
Tadao Oda.
\newblock \emph{Lectures on Torus Embeddings and Applications}.
\newblock Number~58 in Tata Institute of Fundamental Research Lectures on
  Mathematics and Physics. Springer-Verlag, Berlin, Heidelberg, New York, 1978.
\newblock URL
  \url{https://mathweb.tifr.res.in/Documents/Publications/Lectures/tifr58.pdf}.
\newblock Verified in this writing on 18 November 2026 at these loci only.
  Section13pp113-118;\allowbreak{}Theorem15.1p130 and following paragraphp131. Scope:
  Explicit toric Hopf degeneration and hyperbolic Inoue complement as R times
  the torus mapping-torus bundle.

\bibitem[{OpenAI}(2026{\natexlab{a}})]{OpenAIChatGPTPro2026}
{OpenAI}.
\newblock Gpt-6 and other models in chatgpt.
\newblock Official documentation;\allowbreak{} accessed 18 November 2026,
  2026{\natexlab{a}}.
\newblock URL
  \url{https://help.openai.com/en/articles/20001354-gpt-6-and-other-models-in-chatgpt}.
\newblock Verified in this writing on 18 November 2026 at these loci only.
  Choose a thinking level;\allowbreak{} lines25–36 in the opened page.. Scope: The
  documented Pro selection and GPT-6 Astra model mapping.

\bibitem[{OpenAI}(2026{\natexlab{b}})]{OpenAIReasoning2026}
{OpenAI}.
\newblock Reasoning models.
\newblock Official API documentation, accessed 17 November 2026,
  2026{\natexlab{b}}.
\newblock URL \url{https://developers.openai.com/api/docs/guides/reasoning}.
\newblock Verified in this writing on 18 November 2026 at these loci only.
  Reasoning effort and Reasoning mode;\allowbreak{} lines1274–1303 in the opened page..
  Scope: Reasoning allocation controls and the named configuration fields.

\bibitem[Parusi{\'n}ski(1990)]{Parusinski1990E}
Adam Parusi{\'n}ski.
\newblock Subanalytic functions.
\newblock Technical Report MPI/90-58, Max-Planck-Institut f{\"u}r Mathematik,
  1990.
\newblock URL
  \url{https://archive.mpim-bonn.mpg.de/id/eprint/1622/1/preprint_1990_58.pdf}.
\newblock Verified in this writing on 18 November 2026 at these loci only. MPIM
  preprint, introduction printedp1 (PDF3);\allowbreak{} Theorem1.1 printedp2;\allowbreak{} Corollary1.2
  with complete proof printedpp3–6 (PDF5–8).. Scope: The compact sets and
  local blowup neighborhoods, finite coefficient ideal, decreasing fibre excess
  and smooth-centre principalization were read. Receiving references to pp3–5
  are corrected to pp3–6.

\bibitem[Pham(2025)]{pham2025sigma}
Duc~Hiep Pham.
\newblock Semiabelian varieties and transcendence on {Weierstrass} sigma
  functions.
\newblock \emph{Glasgow Mathematical Journal}, 67\penalty0 (1):\penalty0
  107--113, 2025.
\newblock \doi{10.1017/S0017089524000272}.
\newblock URL
  \url{https://www.cambridge.org/core/journals/glasgow-mathematical-journal/article/abs/semiabelian-varieties-and-transcendence-on-weierstrass-sigma-functions/D291E48168188D71FF20D939CAFDD651}.
\newblock Unverified in this writing at this assembly stage;\allowbreak{} no mathematical
  dependence is authorized by an earlier source-note alone.

\bibitem[Pinkham(1974)]{Pinkham1974Root}
Henry~C. Pinkham.
\newblock \emph{Deformations of algebraic varieties with {$G_m$} action}.
\newblock Number~20 in Ast\'erisque. Soci\'et\'e math\'ematique de France,
  1974.
\newblock URL \url{https://www.numdam.org/item/AST_1974__20__1_0/}.
\newblock Verified in this writing on 18 November 2026 at these loci only.
  Remark7.6(iii)p46;\allowbreak{}Section9.1-9.4pp63-65. Scope: Projectively normal
  hyperplane cone sweeping and degree-eight elliptic cone smoothing by the
  Veronese quadric;\allowbreak{} exact degree-eight second component.

\bibitem[{python-flint developers}(2026)]{PythonFlint2026}
{python-flint developers}.
\newblock Python-flint documentation: Arb real ball arithmetic, 2026.
\newblock URL \url{https://python-flint.readthedocs.io/en/latest/arb.html}.
\newblock Verified in this writing on 18 November 2026 at these loci only.
  Official python-flint0.9.0 arb documentation, constructor;\allowbreak{}
  abs\_lower,abs\_upper;\allowbreak{} lower,upper;\allowbreak{} man\_exp;\allowbreak{} mid\_rad\_10exp;\allowbreak{} str sections,
  accessed18November2026.. Scope: Matched exact binary tuple/rational-string
  construction, upward radius rounding and outward endpoint rounding. Decimal
  display is an enclosing output representation;\allowbreak{} exact finite endpoint
  extraction uses man\_exp. These semantics were read against the replayed
  interval certificate and period script. Documentation reading grants no new
  numerical replay credit.

\bibitem[Reyssat(1982)]{reyssat1982weierstrass}
Eric Reyssat.
\newblock Propri\'et\'es d'ind\'ependance alg\'ebrique de nombres li\'es aux
  fonctions de weierstrass.
\newblock \emph{Acta Arithmetica}, 41\penalty0 (3):\penalty0 291--310, 1982.
\newblock \doi{10.4064/aa-41-3-291-310}.
\newblock URL \url{https://matwbn.icm.edu.pl/ksiazki/aa/aa41/aa4137.pdf}.
\newblock Verified in this writing on 18 November 2026 at these loci only.
  Primary indexed PDF text, introduction and Corollary2 printedp292;\allowbreak{} Section5.1
  proof printedp301 was inspected as indexed text.. Scope: Matched the source
  definitions of g\_u,omega and s\_omega and the three quantities in the
  transcendence-degree-at-least-two statement. Local PDF download failed;\allowbreak{} no
  visual scan check or independent verification of the OCR half-argument
  identity is claimed.

\bibitem[Roy and Waldschmidt(1993)]{roy1993periodrank}
Damien Roy and Michel Waldschmidt.
\newblock Autour du th\'eor\`eme du sous-groupe alg\'ebrique.
\newblock \emph{Canadian Mathematical Bulletin}, 36\penalty0 (3):\penalty0
  358--367, 1993.
\newblock \doi{10.4153/CMB-1993-049-8}.
\newblock URL
  \url{https://www.cambridge.org/core/journals/canadian-mathematical-bulletin/article/autour-du-theoreme-du-sousgroupe-algebrique/A1E35F3798690D084664ACD83D1557E8}.
\newblock Verified in this writing on 18 November 2026 at these loci only.
  Canad.Math.Bull.36(3)(1993): quotient definitions p359;\allowbreak{} S\_K and alpha\_K
  definitions p360;\allowbreak{} CM lemma pp361–362;\allowbreak{} main theorem p363 and proof
  pp363–364.. Scope: The coefficient alpha(G)=2dimG-2d0(G)-d1(G) uses
  additive and multiplicative quotient dimensions. It equals4 for the specified
  nontrivial semiabelian extension;\allowbreak{} the current period rank is1, giving the
  displayed bound3.

\bibitem[Rubin(1991)]{CrRubin1991}
Karl Rubin.
\newblock The ``main conjectures'' of {I}wasawa theory for imaginary quadratic
  fields.
\newblock \emph{Inventiones Mathematicae}, 103:\penalty0 25--68, 1991.
\newblock \doi{10.1007/BF01239508}.
\newblock URL \url{https://doi.org/10.1007/BF01239508}.
\newblock Unverified in this writing at this assembly stage;\allowbreak{} no mathematical
  dependence is authorized by an earlier source-note alone.

\bibitem[Rubin(1999)]{CrRubin1999}
Karl Rubin.
\newblock Elliptic curves with complex multiplication and the conjecture of
  {B}irch and {S}winnerton-{D}yer.
\newblock In \emph{Arithmetic Theory of Elliptic Curves}, volume 1716 of
  \emph{Lecture Notes in Mathematics}, pages 167--234. Springer, 1999.
\newblock URL \url{https://swc-math.github.io/aws/1999/99RubinCM.pdf}.
\newblock Verified in this writing on 18 November 2026 at these loci only.
  Definition5.1/Proposition5.3 and proof p15;\allowbreak{} maximal-order convention p16;\allowbreak{}
  Theorem10.1 and complete Section10 argument pp47–50.. Scope: Coates–Wiles
  central-value/rank implication, with the explicit isogeny to maximal order.
  Applied by contrapositive to the nontorsion points already verified;\allowbreak{} it
  supplies no rank-one leading-factor formula.

\bibitem[Schnell(2019)]{Schnell19}
Christian Schnell.
\newblock Lecture 4: The {Appel--Humbert} theorem.
\newblock MAT 615, Stony Brook University, lecture of 6 February, 2019.
\newblock URL
  \url{https://www.math.stonybrook.edu/~cschnell/mat615/lectures/lecture4.pdf}.
\newblock Verified in this writing on 18 November 2026 at these loci only.
  Theorem4.1p1;\allowbreak{} equation4.3p2;\allowbreak{} Theorem4.4p3. Scope: Appell-Humbert exact
  quotient multipliers, effective quotient and projectivity;\allowbreak{} PDF has February6
  heading and no printed year.

\bibitem[Schr{\"o}er and Siebert(2002{\natexlab{a}})]{SS00}
Stefan Schr{\"o}er and Bernd Siebert.
\newblock Irreducible degenerations of primary {K}odaira surfaces.
\newblock In \emph{Complex geometry (G{\"o}ttingen, 2000)}, pages 193--222.
  Springer, Berlin, 2002{\natexlab{a}}.
\newblock URL
  \url{https://www.math.uni-duesseldorf.de/~schroeer/publications_pdf/kodaira.pdf}.
\newblock Verified in this writing on 18 November 2026 at these loci only.
  arXivmath/0001077v2, Definition2.1 p4;\allowbreak{} Propositions3.1–3.5 and full proofs
  pp6–10;\allowbreak{} Theorem3.7 and full proof pp10–11;\allowbreak{} Theorem4.4 and full proof p14;\allowbreak{}
  Theorem5.4 and full proof pp18–20;\allowbreak{} Proposition6.3 and full proof
  pp22–23.. Scope: Read the Hopf conductor maps, normal-line characters and
  toric quotient smoothings. The operative receiving comparisons use Section3.
  The printed final generator in Theorem4.4 is recorded separately: (Phi
  Psi)\textasciicircum{}w acts trivially on
  U/<Phi\textasciicircum{}w,Psi\textasciicircum{}w>, while Phi Psi has orderw;\allowbreak{}
  no use is made of the printed generator. Admissibility retains all
  Definition2.1 hypotheses. Author bibliography entry13 identifies the
  published chapter;\allowbreak{} the listed pages refer expressly to the25-page preprint.

\bibitem[Schr{\"o}er and Siebert(2002{\natexlab{b}})]{SchroeerSiebert2002Root}
Stefan Schr{\"o}er and Bernd Siebert.
\newblock Irreducible degenerations of primary {Kodaira} surfaces.
\newblock In \emph{Complex geometry ({G}{\"o}ttingen, 2000)}, pages 193--222.
  Springer, Berlin, 2002{\natexlab{b}}.
\newblock URL
  \url{https://www.math.uni-duesseldorf.de/~schroeer/publications_pdf/kodaira.pdf}.
\newblock Verified in this writing on 18 November 2026 at these loci only.
  arXivmath/0001077v2, Definition2.1 p4;\allowbreak{} Propositions3.1–3.5 and full proofs
  pp6–10;\allowbreak{} Theorem3.7 and full proof pp10–11;\allowbreak{} Theorem4.4 and full proof p14;\allowbreak{}
  Theorem5.4 and full proof pp18–20;\allowbreak{} Proposition6.3 and full proof
  pp22–23.. Scope: Read the Hopf conductor maps, normal-line characters and
  toric quotient smoothings. The operative receiving comparisons use Section3.
  The printed final generator in Theorem4.4 is recorded separately: (Phi
  Psi)\textasciicircum{}w acts trivially on
  U/<Phi\textasciicircum{}w,Psi\textasciicircum{}w>, while Phi Psi has orderw;\allowbreak{}
  no use is made of the printed generator. Admissibility retains all
  Definition2.1 hypotheses. Author bibliography entry13 identifies the
  published chapter;\allowbreak{} the listed pages refer expressly to the25-page preprint.
  Bibliographic identity: Identical authors, chapter, published page
  range193–222 and author PDF URL;\allowbreak{} author bibliography entry13 identifies the
  chapter. Only explicitly preprint-page loci are inherited. This is the same
  reading event, with no additional independent credit.

\bibitem[Shin(2015)]{Shin2015}
Dong~Hwa Shin.
\newblock Generation of class fields by {Siegel--Ramachandra} invariants.
\newblock \emph{Journal of the Korean Mathematical Society}, 52\penalty0
  (5):\penalty0 907--928, 2015.
\newblock URL \url{https://arxiv.org/pdf/1009.2253}.
\newblock Verified in this writing on 18 November 2026 at these loci only.
  arXiv:1009.2253v6: equation(1), p1;\allowbreak{} modular-function setup p2;\allowbreak{} Remark2.2 and
  complete Proposition2.3 with its proof pointers, pp3–4.. Scope: The product
  convention, q-order, lattice-translation factor and modular-unit exponent
  match the receiving Siegel-function identities. This checks the explicit
  source statement;\allowbreak{} the cited earlier book proofs were not newly opened.

\bibitem[Shioda(1990)]{Shioda90}
Tetsuji Shioda.
\newblock On the {M}ordell--{W}eil lattices.
\newblock \emph{Commentarii Mathematici Universitatis Sancti Pauli},
  39\penalty0 (2):\penalty0 211--240, 1990.
\newblock \doi{10.14992/00009994}.
\newblock URL
  \url{https://www.researchgate.net/profile/Tetsuji-Shioda/publication/247043151_On_the_Mordell-Weil_Lattices/links/5704b7fe08ae13eb88b69184/On-the-Mordell-Weil-Lattices.pdf}.
\newblock Verified in this writing on 18 November 2026 at these loci only.
  Author-uploaded 37-page edition: Theorem1.3p5;\allowbreak{} Theorem4.1p10;\allowbreak{} Theorem8.4p20;\allowbreak{}
  Theorem8.6p21;\allowbreak{} table8.16p22;\allowbreak{} Proposition8.12/eq8.25p24. Scope: Exact NS/T
  map, Picard pullback, positive height pairing, intersection/Cartan inverse
  formula, I\_n correction and base-change degree. Original journal repository
  located but fetch429;\allowbreak{} pagination explicitly belongs to author-uploaded
  edition. Original journal edition39(2)(1990), Theorem8.6/eq8.12 p.228;\allowbreak{}
  table8.16 p.229;\allowbreak{} Proposition8.12/eq8.25 p.231. Scope: Downloaded the original
  Rikkyo scan and visually read printed pp.228–231. Read the complete
  height-formula proof and the complete degree-of-field-extension proof.
  Checked local component conventions and the positive correction sign on
  solving for section intersection. These locators are journal pages, not the
  separate37-page author-uploaded pagination. PDF SHA256
  2d21b4174b910f3d0b0122c478cbc6de52ad82f04014c7ba44374aff244ed75f Author/title
  and opened edition validated. Verification applies only to these exact
  locators and edition;\allowbreak{} no pagination or proof scope transfers to another
  edition.

\bibitem[Silverman(1983)]{silverman1983specialization}
Joseph~H. Silverman.
\newblock Heights and the specialization map for families of abelian varieties.
\newblock \emph{Journal f{\"u}r die reine und angewandte Mathematik},
  342:\penalty0 197--211, 1983.
\newblock \doi{10.1515/crll.1983.342.197}.
\newblock URL \url{https://eudml.org/doc/152542}.
\newblock Unverified in this writing at this assembly stage;\allowbreak{} no mathematical
  dependence is authorized by an earlier source-note alone.

\bibitem[Silverman(2006)]{silverman2006heights}
Joseph~H. Silverman.
\newblock An introduction to height functions.
\newblock MSRI workshop lecture slides, January 2006, 2006.
\newblock URL
  \url{https://slmath-production-assets.sfo3.digitaloceanspaces.com/attachments/workshops/301/HtSurveyMSRIJan06.pdf}.
\newblock Verified in this writing on 18 November 2026 at these loci only.
  Author MSRI lecture, printed slides21–22 (PDF26–27): Variation of the
  Canonical Height in Families I–II.. Scope: Read the uniform
  canonical-versus-Weil-height bound with base-height term and the section
  specialization limit. The operative bounded-remainder proof in the paper
  remains the separately opened DeMarco–Mavraki theorems;\allowbreak{} the inaccessible
  original Tate and1983Silverman articles retain their explicit background-only
  status.

\bibitem[Stein(2007)]{CgzStein2007}
William Stein.
\newblock \emph{Modular Forms: A Computational Approach}, volume~79 of
  \emph{Graduate Studies in Mathematics}.
\newblock American Mathematical Society, 2007.
\newblock URL \url{https://wstein.org/books/modform/stein-modform.pdf}.
\newblock Verified in this writing on 18 November 2026 at these loci only.
  Author-hosted book: Theorem9.18 and complete proof pp171–173;\allowbreak{}
  Corollaries9.19–9.20 pp173–174 (PDFpp185–188).. Scope: Sturm
  coefficient bounds, including fixed character. The receiving
  characteristic-zero identity follows also from its displayed
  coset-product/valence proof;\allowbreak{} all coefficient comparisons reach the required
  bounds.

\bibitem[Stelzig(2021)]{Stelzig2021}
Jonas Stelzig.
\newblock On the structure of double complexes.
\newblock \emph{Journal of the London Mathematical Society}, 104\penalty0
  (2):\penalty0 956--988, 2021.
\newblock \doi{10.1112/jlms.12453}.
\newblock Verified in this writing on 18 November 2026 at these loci only.
  arXiv1812.00865v3,8March2021: Theorem3 p7 and complete displayed proof/Lemma4
  pp8–11;\allowbreak{} cohomology definitions pp13–15;\allowbreak{} Proposition6 and Corollary7 p16;\allowbreak{}
  Lemma8 p17. Scope: Read the bounded double-complex decomposition, including
  arbitrary cardinal multiplicities and the explicit splitting/quiver argument,
  then the exact maps on row/column/total/BC/Aeppli cohomology. Even zigzag
  multiplicity is rank of the appropriate Frölicher differential;\allowbreak{} odd zigzags
  supply total cohomology;\allowbreak{} BC and Aeppli count respectively incoming and
  outgoing corners. The actual maps are induced by the chosen noncanonical
  complex isomorphism;\allowbreak{} cardinal multiplicities are canonical. This applies to
  the bounded bidegree complex of smooth forms without requiring every cochain
  space to be finite-dimensional.

\bibitem[Sutherland(2023)]{Sutherland2023CM}
Andrew~V. Sutherland.
\newblock The {Hilbert} class polynomial.
\newblock MIT 18.783 Elliptic Curves, Lecture 20, November 2023.
\newblock URL
  \url{https://math.mit.edu/classes/18.783/2023/LectureNotes20.pdf}.
\newblock Verified in this writing on 18 November 2026 at these loci only.
  MIT18.783 Fall2023 Lecture20, Theorem20.12 with its full displayed proof and
  Corollary20.13, p6;\allowbreak{} Galois/class-group discussion p7.. Scope: CM j-invariants
  are algebraic integers. The proof uses a cyclic prime-degree self-isogeny and
  the monic diagonal modular polynomial;\allowbreak{} the conjugation map on endomorphism
  rings is written explicitly.

\bibitem[Sz\'ekelyhidi et~al.(2017)Sz\'ekelyhidi, Tosatti, and
  Weinkove]{STW2017}
G\'abor Sz\'ekelyhidi, Valentino Tosatti, and Ben Weinkove.
\newblock Gauduchon metrics with prescribed volume form, 2017.
\newblock URL \url{https://arxiv.org/abs/1503.04491}.
\newblock Verified in this writing on 18 November 2026 at these loci only.
  arXiv1503.04491v2,8December2017: Ricci and Bott–Chern conventions p2;\allowbreak{}
  Theorem1.3 p3;\allowbreak{} exact ansatz(1.6) and reduction to Theorem1.4 p4. Scope: Read
  the theorem with prescribed real closed(1,1) representative of the
  Chern–Ricci class, fixed Aeppli class of
  omega\textasciicircum{}(n\ensuremath{-}1), and the explicit
  partial-differential-equation reduction. The paper uses [Ric] as its
  Bott–Chern first Chern class, whereas the manuscript uses iF/(2pi);\allowbreak{}
  matching requires [Ric]=2pi c1\_BC. Thus the target Ric=pi f*omega\_FS
  corresponds to the normalized input c1\_BC=(1/2)[f*omega\_FS];\allowbreak{} the computed
  Chern class remains the manuscript input and is not re-derived by this source
  event. This bounded source match does not claim an independent reproof of
  the29-page PDE estimates.

\bibitem[Tate(1983)]{tate1983variation}
John Tate.
\newblock Variation of the canonical height of a point depending on a
  parameter.
\newblock \emph{American Journal of Mathematics}, 105\penalty0 (1):\penalty0
  287--294, 1983.
\newblock \doi{10.2307/2374389}.
\newblock URL \url{https://www.jstor.org/stable/2374389}.
\newblock Unverified in this writing at this assembly stage;\allowbreak{} no mathematical
  dependence is authorized by an earlier source-note alone.

\bibitem[Teleman(2010)]{Teleman2010AB}
Andrei Teleman.
\newblock Instantons and curves on class {VII} surfaces.
\newblock \emph{Annals of Mathematics}, 172\penalty0 (3):\penalty0 1749--1804,
  2010.
\newblock \doi{10.4007/annals.2010.172.1749}.
\newblock URL
  \url{https://annals.math.princeton.edu/wp-content/uploads/annals-v172-n3-p07-p.pdf}.
\newblock Verified in this writing on 18 November 2026 at these loci only.
  Published Introduction pp1750–1751, including the explicit Kato
  modification and spherical-shell gluing description and Theorem0.1.. Scope:
  Matched the concrete modification plus boundary identification construction
  used in the receiving Kato-family comparison. The general gauge-theoretic
  proof elsewhere in the54-page article was not read in this event.

\bibitem[Viaclovsky(2026{\natexlab{a}})]{ViaclovskyLog2026}
Jeff Viaclovsky.
\newblock Logarithmic transformations on the six-sphere.
\newblock arXiv:2610.39102v1, 2026{\natexlab{a}}.
\newblock URL \url{https://arxiv.org/abs/2610.39102}.
\newblock Verified in this writing on 18 November 2026 at these loci only.
  Attached October28,2026 source build: Section2.1 equations(2.12)-(2.17),
  pp.4-6;\allowbreak{} full Proposition2.3 proof pp.5-6;\allowbreak{} Theorem3.2 statement and
  equation(3.9), p.7. Scope: Read the entire indicated Section2.1 source locus
  and rendered pages5-7. The actual affine origin change conjugates the
  punctured deck action to its linear form. The invariant period defines the
  logarithmic translation, its branch change is integral, and the displayed
  equivariance gives the quotient biholomorphism with inverse translation by
  its negative. The invariant covector proves freeness for the exact coprime
  data, giving the stated multiplicity. Read the exact simply connected
  selection in Theorem3.2;\allowbreak{} its entire subsequent proof is not credited by this
  source event. All listed keys name the same attached author/title/arXiv
  source, and only the recorded edition and locators are transferred. PDF
  SHA256 8a067e37512fb46d02bcad9f8d6650ef43b77fc59e13c4599371da6111f393c4
  Author/title and attached or reproduced edition validated. Current
  verification applies only to these exact locators;\allowbreak{} no full-proof or
  different-edition scope is inferred. Section2.1 (2.11)–(2.12),
  Proposition2.3;\allowbreak{} Propositions4.1,4.3 and finite-filling normal-torus formulas.
  Scope: The exact free affine action, logarithmic clutching and integral
  normal-torus maps for D2.5. Bibliographic identity: Same attached
  arXiv2610.39102v1 author source. Only the recorded local logarithmic-filling
  statements are inherited. This is the same reading event, with no additional
  independent credit. Lemma 2.1, label mono, lines 288–315;\allowbreak{} invariant
  covectors/forms and the finite-filling setup, lines 316–390. Scope: Full
  determinant-one change of basis and its three displayed monodromy matrices,
  primitive invariant circle/covector and exact filling coordinates. This
  extends the prior equations (2.11)–(2.12) and logarithmic-clutch reading.
  Bibliographic identity: Same attached arXiv2610.39102v1 author source. Only
  the recorded local logarithmic-filling statements are inherited. This is the
  same reading event, with no additional independent credit. Attached
  October28,2026 source build: Section2.1 equations(2.12)-(2.17), pp.4-6;\allowbreak{} full
  Proposition2.3 proof pp.5-6;\allowbreak{} Theorem3.2 statement and equation(3.9), p.7.
  Scope: Read the entire indicated Section2.1 source locus and rendered
  pages5-7. The actual affine origin change conjugates the punctured deck
  action to its linear form. The invariant period defines the logarithmic
  translation, its branch change is integral, and the displayed equivariance
  gives the quotient biholomorphism with inverse translation by its negative.
  The invariant covector proves freeness for the exact coprime data, giving the
  stated multiplicity. Read the exact simply connected selection in Theorem3.2;\allowbreak{}
  its entire subsequent proof is not credited by this source event. All listed
  keys name the same attached author/title/arXiv source, and only the recorded
  edition and locators are transferred. PDF SHA256
  8a067e37512fb46d02bcad9f8d6650ef43b77fc59e13c4599371da6111f393c4 Author/title
  and attached or reproduced edition validated. Current verification applies
  only to these exact locators;\allowbreak{} no full-proof or different-edition scope is
  inferred. Bibliographic identity: Same attached arXiv2610.39102v1 author
  source. Only the recorded local logarithmic-filling statements are inherited.
  This is the same reading event, with no additional independent credit.

\bibitem[Viaclovsky(2026{\natexlab{b}})]{ViaclovskyTwo2026}
Jeff Viaclovsky.
\newblock A two-parameter family of complex structures on {$S^6$}.
\newblock arXiv:2609.33785v1, 2026{\natexlab{b}}.
\newblock URL \url{https://arxiv.org/abs/2609.33785}.
\newblock Verified in this writing on 18 November 2026 at these loci only.
  Attached-source edition: Proposition 3.1 p.6;\allowbreak{} Proposition 6.4 p.19. Scope:
  Opened the rank-one/rank-two Mumford models, their normalized surfaces and
  primitive monodromy, and the simultaneous parameter-bidisc construction.
  Bibliographic identity: Same attached arXiv2609.33785v1 author source;\allowbreak{}
  source-edition locators remain unchanged. This is the same reading event,
  with no additional independent credit. Attached-source build: Lemma2.1
  pp3–4;\allowbreak{} Lemma3.2 and complete proof p7;\allowbreak{} Section3.3 equations3.11–3.14 p8.
  Scope: Read the full source-fibre, section, height and doubling proof, the
  exact fourth-root coordinates and nowhere-zero lifted map, its O-section
  frame correction, and the complete primitive ell/4 affine quotient and
  logarithmic inverse. Visually opened p3. These exact source and cap maps
  support the F18 degree-four page calculation. The attached PDF title is
  Complex structures on S6;\allowbreak{} the supplied canonical bibliography title is its
  two-parameter-family variant, identified by the same author and arXiv
  identifier. PDF SHA256
  99545c911a6a68d919aec11322b898aaf145f94ae7f4747e6ad1e66104b77846 Same
  identified author/title/arXiv work or original edition;\allowbreak{} current verification
  applies only to the exact opened edition and locators stated here.
  Bibliographic identity: Same attached arXiv2609.33785v1 author source;\allowbreak{}
  source-edition locators remain unchanged. This is the same reading event,
  with no additional independent credit.

\bibitem[Waldschmidt(1975)]{waldschmidt1975dimension}
Michel Waldschmidt.
\newblock Dimension alg{\'e}brique de sous-groupes analytiques de
  vari{\'e}t{\'e}s de groupe.
\newblock \emph{Annales de l'Institut Fourier}, 25\penalty0 (1):\penalty0
  23--33, 1975.
\newblock \doi{10.5802/aif.543}.
\newblock URL
  \url{https://webusers.imj-prg.fr/~michel.waldschmidt/articles/pdf/Fourier.pdf}.
\newblock Verified in this writing on 18 November 2026 at these loci only.
  Theorem1 p24, complete proof pp25–27;\allowbreak{} strict inequality on p24 also opened
  visually.. Scope: Fixed number field, growth/size exponents, point count and
  separation exponent. For the independent triple on C2 the necessary displayed
  numerical inequality fails already at order2;\allowbreak{} packing in the actual real
  plane strengthens the failure.

\bibitem[Waldschmidt(1979)]{waldschmidt1979sigma}
Michel Waldschmidt.
\newblock Nombres transcendants et fonctions sigma de {Weierstrass}.
\newblock \emph{Comptes Rendus Mathematiques de l'Academie des Sciences,
  Canada}, 1\penalty0 (2):\penalty0 111--114, 1979.
\newblock URL
  \url{https://webusers.imj-prg.fr/~michel.waldschmidt/articles/pdf/GDL/SigmaNoteCanada1979-eng.pdf}.
\newblock Verified in this writing on 18 November 2026 at these loci only.
  Author-hosted June2015 English translation by K.Senthil Kumar and S.Kannappan
  of the1979 note, all3pages;\allowbreak{} Theorem1 and Corollaries1–2 p3.. Scope: The
  sigma quotient with its zeta exponential and the
  nontorsion/algebraic-invariant hypotheses are matched. The translation date
  is distinguished from the1979 article date.

\bibitem[Waldschmidt(1988{\natexlab{a}})]{waldschmidt1988periodrank}
Michel Waldschmidt.
\newblock On the transcendence methods of {Gelfond} and {Schneider} in several
  variables.
\newblock In Alan Baker, editor, \emph{New Advances in Transcendence Theory},
  chapter~24, pages 375--398. Cambridge University Press, 1988{\natexlab{a}}.
\newblock URL
  \url{https://webusers.imj-prg.fr/~michel.waldschmidt/articles/pdf/Durham1986TMGSSV.pdf}.
\newblock Verified in this writing on 18 November 2026 at these loci only.
  Original scanned chapter, Theorem4.1 pp385–386 and Corollary4.3 p386,
  opened visually;\allowbreak{} introductory group factors p375 and period-rank argument
  pp392–395 also read visually.. Scope: Exact period correction kappa and
  dimension/rank coefficients used in the compact-line audit. The
  multiplicative factors in a product decomposition and multiplicative
  quotients are tracked through the specified group.

\bibitem[Waldschmidt(1988{\natexlab{b}})]{waldschmidt1988theta}
Michel Waldschmidt.
\newblock Fonctions th\^eta et nombres transcendants.
\newblock S\'eminaire de Th\'eorie des Nombres de Bordeaux, 1987--1988,
  expos\'e 15, pp. 1--7, 1988{\natexlab{b}}.
\newblock URL
  \url{https://webusers.imj-prg.fr/~michel.waldschmidt/articles/pdf/PPN320141322_0016__0_0.pdf}.
\newblock Verified in this writing on 18 November 2026 at these loci only.
  Original scanned lecture, printed pp15-03 and15-05 (PDFpp4 and6), both opened
  visually.. Scope: The CM holomorphic sigma normalization and the algebraic-q
  hypothesis of the periodic-functional-equation theorem were checked directly.
  The sigma factor and the Hermitian metric factor remain related by the
  explicit Gaussian in the receiving text.

\bibitem[Waldschmidt(2009)]{waldschmidt2009auxiliary}
Michel Waldschmidt.
\newblock Auxiliary functions in transcendental number theory.
\newblock \emph{The Ramanujan Journal}, 20\penalty0 (3):\penalty0 341--373,
  2009.
\newblock \doi{10.1007/s11139-009-9204-y}.
\newblock URL
  \url{https://webusers.imj-prg.fr/~michel.waldschmidt/articles/pdf/SASTRAaftp.pdf}.
\newblock Verified in this writing on 18 November 2026 at these loci only.
  arXiv:0908.4024v1, Theorems3.6–3.7 pp21–22.. Scope: Several-variable
  Schneider–Lang and Bombieri inputs require an algebraically independent
  finite-order family closed under all partial derivatives over one fixed
  number field. The present normalized second jet recovers the proved
  transcendental k.

\bibitem[Xiao(1991)]{Xiao91}
Gang Xiao.
\newblock Elliptic surfaces with fixed jacobian.
\newblock Technical Report MPI/91-30, Max-Planck-Institut f{\"u}r Mathematik,
  1991.
\newblock URL
  \url{https://archive.mpim-bonn.mpg.de/id/eprint/2953/1/preprint_1991_30.pdf}.
\newblock Verified in this writing on 18 November 2026 at these loci only.
  MPI91-30,Theorem3.4p14;\allowbreak{}Section4 equation4p21;\allowbreak{}Theorem4.6p25. Scope:
  Finite-order translating twists, surjective localization for nontrivial
  Jacobian, Sha=(Q/Z)\textasciicircum{}(b2-rho) for singular nonisotrivial
  Jacobian.

\bibitem[Xie(2024)]{xie2024algebraicity}
Junyi Xie.
\newblock Algebraicity criteria, invariant subvarieties and transcendence
  problems from arithmetic dynamics.
\newblock \emph{Peking Mathematical Journal}, 7:\penalty0 345--398, 2024.
\newblock \doi{10.1007/s42543-022-00059-9}.
\newblock URL \url{https://arxiv.org/pdf/2202.09032}.
\newblock Verified in this writing on 18 November 2026 at these loci only.
  arXiv2202.09032: Theorem1.3 and its setup p3;\allowbreak{} Theorem2.3 and full proof
  pp12–14;\allowbreak{} Theorem2.5 and full proof pp14–15.. Scope: Recorded the
  K-defined formal germ, K-rational approximating points and positive local
  approximation/height ratio. These hypotheses are retained explicitly in the
  discussion of the attempted application.

\bibitem[Yang(2026)]{Yang2026}
Xiufan Yang.
\newblock Compact complex threefolds fibred by two-dimensional complex tori.
\newblock arXiv:2609.29666v1, 2026.
\newblock URL \url{https://arxiv.org/abs/2609.29666}.
\newblock Verified in this writing on 18 November 2026 at these loci only.
  Attached repaired v4.5 source build: Lemma7.1 and Theorem7.2 with complete
  proof pp40–41. Scope: Read the actual invariant closure, all finite/cusp
  attaching relations and the entire exact presentation proof, including the
  Smith-minor calculation. Visually opened p40. This source already supplies
  the canonical Z/2 and Z/3 groups and the stronger full all-twist
  presentation. The current F19 quotient map is an explicit rederivation. No
  pagination from this attached repaired edition is assigned to arXiv v1. PDF
  SHA256 7601378a2250cd9e7a92f36afbd97c542b251a57711538735f9517db443e0533 Same
  identified author/title/arXiv work or original edition;\allowbreak{} current verification
  applies only to the exact opened edition and locators stated here. Theorem
  8.7, pp.57–59, and Corollary 8.8, p.59, in the locally rendered attached
  REPAIRED\_v4\_5 source;\allowbreak{} section 8.4. Scope: Read the full theorem and proof
  in the attached source and the rendered PDF pages57–59;\allowbreak{} visually opened
  p.57. Verified the exact period action exp(2pi i
  D(t,s)nu)t\textasciicircum{}nu, uniform logarithmic properness bounds, the
  finite fundamental polydiscs, the canonical form descent, the lattice index
  and maximal-simplex orbit Euler count. For rank3 primitive standard grid this
  gives6. Only font maps were added to render the attached body;\allowbreak{} the pagination
  is explicitly this attached build, not an unverified arXiv edition. Source
  SHA256 5061ce4e6cda8ed7ba3bb2a7a597d50e018bfd53fb9e49e7b38317f628c332a5;\allowbreak{} PDF
  SHA256 7601378a2250cd9e7a92f36afbd97c542b251a57711538735f9517db443e0533
  Author/title and edition matched to the opened item. Verification applies
  only to its recorded locator and edition;\allowbreak{} aliasing does not transfer preprint
  pagination to a published edition.

\bibitem[Yuan et~al.(2011)Yuan, Zhang, and Zhang]{CgzYZZ2011}
Xinyi Yuan, Shou-Wu Zhang, and Wei Zhang.
\newblock Gross--zagier formula on shimura curves.
\newblock Verified in this writing on 18 November 2026 at these loci only.
  Section7.1.1pp218-219,Theorems7.1.1-7.1.2,Proposition7.1.3. Scope:
  Relative-to-base-field height, bilinear pairing without1/2, Poincare
  polarization compatibility, 2011.
\newblock URL
  \url{https://www.researchgate.net/profile/Xinyi-Yuan-11/publication/267551160_Gross-Zagier_Formula_On_Shimura_Curves/links/548e842c0cf225bf66a5ff13/Gross-Zagier-Formula-On-Shimura-Curves.pdf?origin=publication_detail}.

\end{thebibliography}

\end{document}
